\begin{center}
  {\Huge \textbf{חומש בראשית}}
\end{center}
\vspace{1.5cm}

\addparasha{חומש בראשית | פרשת בראשית}
% ---------- פרק 1 | פסוק 1 ----------
 \textbf{א} בְּרֵאשִׁ֖ית בָּרָ֣א אֱלֹהִ֑ים אֵ֥ת הַשָּׁמַ֖יִם וְאֵ֥ת הָאָֽרֶץ׃ \Rashi{\textbf{בראשית.}אמר רבי יצחק לא היה צריך להתחיל את התורה אלא מהחדש הזה לכם, שהיא מצוה ראשונה שנצטוו בה ישראל, ומה טעם פתח בבראשית? משום כח מעשיו הגיד לעמו לתת להם נחלת גוים (תהילים קי"א), שאם יאמרו אמות העולם לישראל לסטים אתם, שכבשתם ארצות שבעה גוים, הם אומרים להם כל הארץ של הקב"ה היא, הוא בראה ונתנה לאשר ישר בעיניו, ברצונו נתנה להם, וברצונו נטלה מהם ונתנה לנו: בראשית ברא אין המקרא הזה אומר אלא דרשני, כמו שדרשוהו רבותינו בשביל התורה שנקראת ראשית דרכו (משלי ח'), ובשביל ישראל שנקראו ראשית תבואתו (ירמיה ב'); ואם באת לפרשו כפשוטו, כך פרשהו בראשית בריאת שמים וארץ, והארץ היתה תהו ובהו וחשך ויאמר אלהים יהי אור ולא בא המקרא להורות סדר הבריאה לומר שאלו קדמו, שאם בא להורות כך, היה לו לכתב בראשונה ברא את השמים וגו' שאין לך ראשית במקרא שאינו דבוק לתבה שלאחריו, כמו בראשית ממלכת יהויקים (שם כ"ז), ראשית ממלכתו (בראשית י'), ראשית דגנך (דברים י"ח, ד'), – אף כאן אתה אומר בראשית ברא אלהים וגו', כמו בראשית ברא; ודומה לו תחלת דבר ה' בהושע (הושע א'), כלומר תחלת דבורו של הקב"ה בהושע, ויאמר ה' אל הושע וגו'. וא"ת להורות בא שאלו תחלה נבראו, ופרושו בראשית הכל ברא אלו – ויש לך מקראות שמקצרים לשונם וממעטים תבה אחת, כמו כי לא סגר דלתי בטני (איוב ג'), ולא פרש מי הסוגר וכמו ישא את חיל דמשק (ישעיהו ח'), ולא פרש מי ישאנו; וכמו אם יחרוש בבקרים (עמוס ו'), ולא פרש אם יחרוש אדם בבקרים, וכמו מגיד מראשית אחרית (ישעיהו מ"ו), ולא פרש מגיד מראשית דבר אחרית דבר – א"כ, תמה על עצמך, שהרי המים קדמו, שהרי כתיב ורוח אלהים מרחפת על פני המים, ועדיין לא גלה המקרא, בריאת המים מתי היתה, הא למדת, שקדמו המים לארץ, ועוד, שהשמים מאש ומים נבראו, על כרחך לא למד המקרא סדר המקדמים והמאחרים כלום: ברא אלהים ולא נאמר ברא ה', שבתחלה עלה במחשבה לבראתו במדת הדין, ראה שאין העולם מתקים, הקדים מדת רחמים ושתפה למה"ד, והיינו דכתיב ביום עשות ה' אלהים ארץ ושמים:}%
% ---------- פרק 1 | פסוק 2 ----------
 \textbf{ב} וְהָאָ֗רֶץ הָיְתָ֥ה תֹ֙הוּ֙ וָבֹ֔הוּ וְחֹ֖שֶׁךְ עַל־פְּנֵ֣י תְה֑וֹם וְר֣וּחַ אֱלֹהִ֔ים מְרַחֶ֖פֶת עַל־פְּנֵ֥י הַמָּֽיִם׃ \Rashi{\textbf{תהו ובהו.}תהו לשון תמה ושממון, שאדם תוהא ומשתומם על בהו שבה: תהו. אשטורדי"שון בלע"ז: בהו. לשון רקות וצדו: על פני תהום. על פני המים שעל הארץ: ורוח אלהים מרחפת. כסא הכבוד עומד באויר ומרחף על פני המים ברוח פיו של הקב"ה ובמאמרו, כיונה המרחפת על הקן, אקוב"טיר בלע"ז:}%
% ---------- פרק 1 | פסוק 3 ----------
 \textbf{ג} וַיֹּ֥אמֶר אֱלֹהִ֖ים יְהִ֣י א֑וֹר וַֽיְהִי־אֽוֹר׃ %
% ---------- פרק 1 | פסוק 4 ----------
 \textbf{ד} וַיַּ֧רְא אֱלֹהִ֛ים אֶת־הָא֖וֹר כִּי־ט֑וֹב וַיַּבְדֵּ֣ל אֱלֹהִ֔ים בֵּ֥ין הָא֖וֹר וּבֵ֥ין הַחֹֽשֶׁךְ׃ \Rashi{\textbf{וירא אלהים את האור כי טוב ויבדל.}אף בזה אנו צריכין לדברי אגדה, ראהו שאינו כדאי להשתמש בו רשעים, והבדילו לצדיקים לעתיד לבא. ולפי פשוטו כך פרשהו, ראהו כי טוב, ואין נאה לו ולחשך שיהיו משתמשים בערבוביא, וקבע לזה תחומו ביום, ולזה תחומו בלילה:}%
% ---------- פרק 1 | פסוק 5 ----------
 \textbf{ה} וַיִּקְרָ֨א אֱלֹהִ֤ים ׀ לָאוֹר֙ י֔וֹם וְלַחֹ֖שֶׁךְ קָ֣רָא לָ֑יְלָה וַֽיְהִי־עֶ֥רֶב וַֽיְהִי־בֹ֖קֶר י֥וֹם אֶחָֽד׃ \{פ\} \Rashi{\textbf{יום אחד.}לפי סדר לשון הפרשה היה לו לכתב יום ראשון, כמו שכתוב בשאר הימים שני, שלישי, רביעי, למה כתב אחד? לפי שהיה הקב"ה יחיד בעולמו, שלא נבראו המלאכים עד יום שני, כך מפרש בב"ר:}%
% ---------- פרק 1 | פסוק 6 ----------
 \textbf{ו} וַיֹּ֣אמֶר אֱלֹהִ֔ים יְהִ֥י רָקִ֖יעַ בְּת֣וֹךְ הַמָּ֑יִם וִיהִ֣י מַבְדִּ֔יל בֵּ֥ין מַ֖יִם לָמָֽיִם׃ \Rashi{\textbf{יהי רקיע.}יחזק הרקיע, שאף על פי שנבראו שמים ביום א' עדין לחים היו וקרשו בשני מגערת הקב"ה באמרו יהי רקיע, וזהו שכתב (איוב כ"ו) עמודי שמים ירופפו – כל יום ראשון, ובשני יתמהו מגערתו, כאדם שמשתומם ועומד מגערת המאים עליו: בתוך המים. באמצע המים, שיש הפרש בין מים העליונים לרקיע כמו בין הרקיע למים שעל הארץ, הא למדת, שהם תלויים במאמרו של מלך:}%
% ---------- פרק 1 | פסוק 7 ----------
 \textbf{ז} וַיַּ֣עַשׂ אֱלֹהִים֮ אֶת־הָרָקִ֒יעַ֒ וַיַּבְדֵּ֗ל בֵּ֤ין הַמַּ֙יִם֙ אֲשֶׁר֙ מִתַּ֣חַת לָרָקִ֔יעַ וּבֵ֣ין הַמַּ֔יִם אֲשֶׁ֖ר מֵעַ֣ל לָרָקִ֑יעַ וַֽיְהִי־כֵֽן׃ \Rashi{\textbf{ויעש אלהים את הרקיע.}תקנו על עמדו, והיא עשיתו, כמו ועשתה את צפרניה (דברים כ"א): מעל לרקיע. על הרקיע לא נאמר אלא מעל לרקיע, לפי שהן תלויין באויר. ומפני מה לא נאמר כי טוב ביום שני? לפי שלא היה נגמר מלאכת המים עד יום שלישי, והרי התחיל בה בשני, ודבר שלא נגמר אינו במלואו וטובו; ובשלישי שנגמר מלאכת המים והתחיל וגמר מלאכה אחרת, כפל בו כי טוב שתי פעמים, אחת לגמר מלאכת השני ואחת לגמר מלאכת היום:}%
% ---------- פרק 1 | פסוק 8 ----------
 \textbf{ח} וַיִּקְרָ֧א אֱלֹהִ֛ים לָֽרָקִ֖יעַ שָׁמָ֑יִם וַֽיְהִי־עֶ֥רֶב וַֽיְהִי־בֹ֖קֶר י֥וֹם שֵׁנִֽי׃ \{פ\} \Rashi{\textbf{ויקרא אלהים לרקיע שמים.}שא מים, שם מים, אש ומים, שערבן זה בזה ועשה מהם שמים:}%
% ---------- פרק 1 | פסוק 9 ----------
 \textbf{ט} וַיֹּ֣אמֶר אֱלֹהִ֗ים יִקָּו֨וּ הַמַּ֜יִם מִתַּ֤חַת הַשָּׁמַ֙יִם֙ אֶל־מָק֣וֹם אֶחָ֔ד וְתֵרָאֶ֖ה הַיַּבָּשָׁ֑ה וַֽיְהִי־כֵֽן׃ \Rashi{\textbf{יקוו המים.}שהיו שטוחין על פני כל הארץ והקום באוקינוס הוא הים הגדול שבכל הימים:}%
% ---------- פרק 1 | פסוק 10 ----------
 \textbf{י} וַיִּקְרָ֨א אֱלֹהִ֤ים ׀ לַיַּבָּשָׁה֙ אֶ֔רֶץ וּלְמִקְוֵ֥ה הַמַּ֖יִם קָרָ֣א יַמִּ֑ים וַיַּ֥רְא אֱלֹהִ֖ים כִּי־טֽוֹב׃ \Rashi{\textbf{קרא ימים.}והלא ים אחד הוא? אלא אינו דומה טעם דג העולה מן הים בעכו לטעם דג העולה מן הים באספמיא:}%
% ---------- פרק 1 | פסוק 11 ----------
 \textbf{יא} וַיֹּ֣אמֶר אֱלֹהִ֗ים תַּֽדְשֵׁ֤א הָאָ֙רֶץ֙ דֶּ֗שֶׁא עֵ֚שֶׂב מַזְרִ֣יעַ זֶ֔רַע עֵ֣ץ פְּרִ֞י עֹ֤שֶׂה פְּרִי֙ לְמִינ֔וֹ אֲשֶׁ֥ר זַרְעוֹ־ב֖וֹ עַל־הָאָ֑רֶץ וַֽיְהִי־כֵֽן׃ \Rashi{\textbf{תדשא הארץ דשא עשב.}לא דשא לשון עשב ולא עשב לשון דשא, ולא היה לשון המקרא לומר תעשיב הארץ, שמיני דשאים מחלקין, כל אחד לעצמו נקרא עשב פלוני, ואין לשון למדבר לומר דשא פלוני, שלשון דשא הוא לבישת הארץ כשהיא מתמלאת בדשאים: תדשא הארץ. תתמלא ותתכסה לבוש עשבים. בלשון לעז נקרא דשא ארבריץ, כלן בערבוביא, וכל שרש לעצמו נקרא עשב: מזריע זרע. שיגדל בו זרעו לזרע ממנו במקום אחר: עץ פרי. שיהא טעם העץ כטעם הפרי, והיא לא עשתה כן, אלא ותוצא הארץ עץ עושה פרי ולא העץ פרי, לפיכך כשנתקלל אדם על עונו נפקדה גם היא על עונה ונתקללה: אשר זרעו בו. הן גרעיני כל פרי שמהן האילן צומח כשנוטעין אותו:}%
% ---------- פרק 1 | פסוק 12 ----------
 \textbf{יב} וַתּוֹצֵ֨א הָאָ֜רֶץ דֶּ֠שֶׁא עֵ֣שֶׂב מַזְרִ֤יעַ זֶ֙רַע֙ לְמִינֵ֔הוּ וְעֵ֧ץ עֹֽשֶׂה־פְּרִ֛י אֲשֶׁ֥ר זַרְעוֹ־ב֖וֹ לְמִינֵ֑הוּ וַיַּ֥רְא אֱלֹהִ֖ים כִּי־טֽוֹב׃ \Rashi{\textbf{ותוצא הארץ וגו'.}אע"פ שלא נאמר למינהו בדשאין בצווייהן, שמעו שנצטוו האילנות על כך, נשאו ק"ו בעצמן, כמפרש באגדה בשחיטת חלין:}%
% ---------- פרק 1 | פסוק 13 ----------
 \textbf{יג} וַֽיְהִי־עֶ֥רֶב וַֽיְהִי־בֹ֖קֶר י֥וֹם שְׁלִישִֽׁי׃ \{פ\} %
% ---------- פרק 1 | פסוק 14 ----------
 \textbf{יד} וַיֹּ֣אמֶר אֱלֹהִ֗ים יְהִ֤י מְאֹרֹת֙ בִּרְקִ֣יעַ הַשָּׁמַ֔יִם לְהַבְדִּ֕יל בֵּ֥ין הַיּ֖וֹם וּבֵ֣ין הַלָּ֑יְלָה וְהָי֤וּ לְאֹתֹת֙ וּלְמ֣וֹעֲדִ֔ים וּלְיָמִ֖ים וְשָׁנִֽים׃ \Rashi{\textbf{יהי מארת וגו'.}מיום ראשון נבראו וברביעי צוה עליהם להתלות ברקיע, וכן כל תולדות שמים וארץ נבראו מיום ראשון, וכל אחד ואחד נקבע ביום שנגזר עליו, הוא שכתוב את השמים, לרבות תולדותיהם, ואת הארץ, לרבות תולדותיה: יהי מארת. חסר וי"ו כתיב, על שהוא יום מארה לפל אסכרה בתינוקות, הוא ששנינו בד' היו מתענים על אסכרה שלא תפל בתינוקות: להבדיל בין היום ובין הלילה. משנגנז האור הראשון, אבל בשבעת ימי בראשית שמשו האור והחשך הראשונים יחד בין ביום ובין בלילה: והיו לאותות. כשהמאורות לוקין, סימן רע הוא לעולם, שנאמר מאתות השמים אל תחתו (ירמיה י'), בעשותכם רצון הקב"ה אין אתם צריכין לדאג מן הפרענות: ולמועדים. על שם העתיד, שעתידים ישראל להצטוות על המועדות והם נמנים למולד הלבנה: ולימים. שמוש החמה חצי יום ושמוש הלבנה חציו, הרי יום שלם: ושנים. לסוף שס"ה ימים (ס"א ורביע יום), יגמרו מהלכתם בי"ב מזלות המשרתים אותם, והיא שנה, וחוזרים ומתחילים פעם שניה לסבב בגלגל כמהלכן הראשון:}%
% ---------- פרק 1 | פסוק 15 ----------
 \textbf{טו} וְהָי֤וּ לִמְאוֹרֹת֙ בִּרְקִ֣יעַ הַשָּׁמַ֔יִם לְהָאִ֖יר עַל־הָאָ֑רֶץ וַֽיְהִי־כֵֽן׃ \Rashi{\textbf{והיו למאורת.}עוד זאת ישמשו, שיאירו לעולם:}%
% ---------- פרק 1 | פסוק 16 ----------
 \textbf{טז} וַיַּ֣עַשׂ אֱלֹהִ֔ים אֶת־שְׁנֵ֥י הַמְּאֹרֹ֖ת הַגְּדֹלִ֑ים אֶת־הַמָּא֤וֹר הַגָּדֹל֙ לְמֶמְשֶׁ֣לֶת הַיּ֔וֹם וְאֶת־הַמָּא֤וֹר הַקָּטֹן֙ לְמֶמְשֶׁ֣לֶת הַלַּ֔יְלָה וְאֵ֖ת הַכּוֹכָבִֽים׃ \Rashi{\textbf{המארת הגדולים.}שוים נבראו ונתמעטה הלבנה על שקטרגה ואמרה א"א לשני מלכים שישתמשו בכתר אחד: ואת הכוכבים. על ידי שמעט את הלבנה הרבה צבאיה להפיס דעתה:}%
% ---------- פרק 1 | פסוק 17 ----------
 \textbf{יז} וַיִּתֵּ֥ן אֹתָ֛ם אֱלֹהִ֖ים בִּרְקִ֣יעַ הַשָּׁמָ֑יִם לְהָאִ֖יר עַל־הָאָֽרֶץ׃ %
% ---------- פרק 1 | פסוק 18 ----------
 \textbf{יח} וְלִמְשֹׁל֙ בַּיּ֣וֹם וּבַלַּ֔יְלָה וּֽלְהַבְדִּ֔יל בֵּ֥ין הָא֖וֹר וּבֵ֣ין הַחֹ֑שֶׁךְ וַיַּ֥רְא אֱלֹהִ֖ים כִּי־טֽוֹב׃ %
% ---------- פרק 1 | פסוק 19 ----------
 \textbf{יט} וַֽיְהִי־עֶ֥רֶב וַֽיְהִי־בֹ֖קֶר י֥וֹם רְבִיעִֽי׃ \{פ\} %
% ---------- פרק 1 | פסוק 20 ----------
 \textbf{כ} וַיֹּ֣אמֶר אֱלֹהִ֔ים יִשְׁרְצ֣וּ הַמַּ֔יִם שֶׁ֖רֶץ נֶ֣פֶשׁ חַיָּ֑ה וְעוֹף֙ יְעוֹפֵ֣ף עַל־הָאָ֔רֶץ עַל־פְּנֵ֖י רְקִ֥יעַ הַשָּׁמָֽיִם׃ \Rashi{\textbf{נפש חיה.}שיהא בה חיות: שרץ. כל דבר חי שאינו גבוה מן הארץ קרוי שרץ: בעוף, כגון זבובים, בשקצים, כגון נמלים וחפושים ותולעים, בבריות, כגון חלד ועכבר וחמט וכיוצא בהם, וכל הדגים:}%
% ---------- פרק 1 | פסוק 21 ----------
 \textbf{כא} וַיִּבְרָ֣א אֱלֹהִ֔ים אֶת־הַתַּנִּינִ֖ם הַגְּדֹלִ֑ים וְאֵ֣ת כׇּל־נֶ֣פֶשׁ הַֽחַיָּ֣ה ׀ הָֽרֹמֶ֡שֶׂת אֲשֶׁר֩ שָׁרְצ֨וּ הַמַּ֜יִם לְמִֽינֵהֶ֗ם וְאֵ֨ת כׇּל־ע֤וֹף כָּנָף֙ לְמִינֵ֔הוּ וַיַּ֥רְא אֱלֹהִ֖ים כִּי־טֽוֹב׃ \Rashi{\textbf{התנינם.}דגים גדולים שבים, ובדברי אגדה הוא לויתן ובן זוגו שבראם זכר ונקבה והרג את הנקבה ומלחה לצדיקים לעתיד לבא, שאם יפרו וירבו לא יתקים העולם בפניהם: נפש החיה. שיש בה חיות:}%
% ---------- פרק 1 | פסוק 22 ----------
 \textbf{כב} וַיְבָ֧רֶךְ אֹתָ֛ם אֱלֹהִ֖ים לֵאמֹ֑ר פְּר֣וּ וּרְב֗וּ וּמִלְא֤וּ אֶת־הַמַּ֙יִם֙ בַּיַּמִּ֔ים וְהָע֖וֹף יִ֥רֶב בָּאָֽרֶץ׃ \Rashi{\textbf{ויברך אותם.}לפי שמחסרים אותם וצדין מהם ואוכלין אותם, הצרכו לברכה, ואף החיות הצרכו לברכה, אלא מפני הנחש שעתיד לקללה, לכך לא ברכן, שלא יהא הוא בכלל: פרו. ל' פרי כלו' עשו פרות: ורבו. אם לא אמר אלא פרו, היה אחד מוליד אחד ולא יותר, ובא ורבו שאחד מוליד הרבה:}%
% ---------- פרק 1 | פסוק 23 ----------
 \textbf{כג} וַֽיְהִי־עֶ֥רֶב וַֽיְהִי־בֹ֖קֶר י֥וֹם חֲמִישִֽׁי׃ \{פ\} %
% ---------- פרק 1 | פסוק 24 ----------
 \textbf{כד} וַיֹּ֣אמֶר אֱלֹהִ֗ים תּוֹצֵ֨א הָאָ֜רֶץ נֶ֤פֶשׁ חַיָּה֙ לְמִינָ֔הּ בְּהֵמָ֥ה וָרֶ֛מֶשׂ וְחַֽיְתוֹ־אֶ֖רֶץ לְמִינָ֑הּ וַֽיְהִי־כֵֽן׃ \Rashi{\textbf{תוצא הארץ.}הוא שפרשתי שהכל נברא מיום ראשון, ולא הצרכו אלא להוציאם: נפש חיה. שיש בה חיות: ורמש. הם שרצים, שהם נמוכים ורומשים על הארץ ונראים כאלו נגררים שאין הלוכן נכר. כל לשון רמש ושרץ בלשוננו קונמו"בריש בלע"ז:}%
% ---------- פרק 1 | פסוק 25 ----------
 \textbf{כה} וַיַּ֣עַשׂ אֱלֹהִים֩ אֶת־חַיַּ֨ת הָאָ֜רֶץ לְמִינָ֗הּ וְאֶת־הַבְּהֵמָה֙ לְמִינָ֔הּ וְאֵ֛ת כׇּל־רֶ֥מֶשׂ הָֽאֲדָמָ֖ה לְמִינֵ֑הוּ וַיַּ֥רְא אֱלֹהִ֖ים כִּי־טֽוֹב׃ \Rashi{\textbf{ויעש.}תקנם בצביונם ובקומתן:}%
% ---------- פרק 1 | פסוק 26 ----------
 \textbf{כו} וַיֹּ֣אמֶר אֱלֹהִ֔ים נַֽעֲשֶׂ֥ה אָדָ֛ם בְּצַלְמֵ֖נוּ כִּדְמוּתֵ֑נוּ וְיִרְדּוּ֩ בִדְגַ֨ת הַיָּ֜ם וּבְע֣וֹף הַשָּׁמַ֗יִם וּבַבְּהֵמָה֙ וּבְכׇל־הָאָ֔רֶץ וּבְכׇל־הָרֶ֖מֶשׂ הָֽרֹמֵ֥שׂ עַל־הָאָֽרֶץ׃ \Rashi{\textbf{נעשה אדם.}ענותנותו של הקב"ה למדנו מכאן, לפי שאדם הוא בדמות המלאכים ויתקנאו בו, לפיכך נמלך בהם, וכשהוא דן את המלכים הוא נמלך בפמליא שלו, שכן מצינו באחאב, שאמר לו מיכה ראיתי את ה' ישב על כסאו וכל צבא השמים עומדים עליו מימינו ומשמאלו (מלכים א כ"ב), וכי יש ימין ושמאל לפניו? אלא אלו מימינים לזכות ואלו משמאילים לחובה, וכן בגזרת עירין פתגמא ובמאמר קדישין שאלתא (דניאל ד'), אף כאן בפמליא שלו נטל רשות, אמר להם יש בעליונים כדמותי; אם אין כדמותי בתחתונים הרי יש קנאה במעשה בראשית: נעשה אדם. אע"פ שלא סיעוהו ביצירתו ויש מקום למינים לרדות, לא נמנע הכתוב מללמד דרך ארץ ומדת ענוה שיהא הגדול נמלך ונוטל רשות מן הקטן; ואם כתב אעשה אדם, לא למדנו שיהא מדבר עם בית דינו, אלא עם עצמו, ותשובת המינים כתב בצדו, ויברא את האדם, ולא כתב ויבראו: בצלמנו. בדפוס שלנו: כדמותנו. להבין ולהשכיל: וירדו בדגת הים. יש בלשון הזה לשון רדוי ולשון ירידה; זכה, רודה בחיות ובבהמות, לא זכה, נעשה ירוד לפניהם והחיה מושלת בו:}%
% ---------- פרק 1 | פסוק 27 ----------
 \textbf{כז} וַיִּבְרָ֨א אֱלֹהִ֤ים ׀ אֶת־הָֽאָדָם֙ בְּצַלְמ֔וֹ בְּצֶ֥לֶם אֱלֹהִ֖ים בָּרָ֣א אֹת֑וֹ זָכָ֥ר וּנְקֵבָ֖ה בָּרָ֥א אֹתָֽם׃ \Rashi{\textbf{ויברא אלהים את האדם בצלמו.}בדפוס העשוי לו, שהכל נברא במאמר והוא נברא בידים, שנאמר ותשת עלי כפכה (תהילים קל"ט); נעשה בחותם כמטבע העשויה על ידי רשם שקורין קוי"ן בלע"ז וכן הוא אומר תתהפך כחמר חותם (איוב ל"ח): בצלם אלהים ברא אותו. פרש לך שאותו צלם המתקן לו, צלם דיוקן יוצרו הוא: זכר ונקבה ברא אותם. ולהלן הוא אומר ויקח אחת מצלעתיו וגו' (בראשית ב')? מדרש אגדה שבראו שני פרצופים בבריאה ראשונה, ואחר כך חלקו. ופשוטו של מקרא, כאן הודיעך שנבראו שניהם בששי ולא פי' לך כיצד בריתן, ופרש לך במקום אחר:}%
% ---------- פרק 1 | פסוק 28 ----------
 \textbf{כח} וַיְבָ֣רֶךְ אֹתָם֮ אֱלֹהִים֒ וַיֹּ֨אמֶר לָהֶ֜ם אֱלֹהִ֗ים פְּר֥וּ וּרְב֛וּ וּמִלְא֥וּ אֶת־הָאָ֖רֶץ וְכִבְשֻׁ֑הָ וּרְד֞וּ בִּדְגַ֤ת הַיָּם֙ וּבְע֣וֹף הַשָּׁמַ֔יִם וּבְכׇל־חַיָּ֖ה הָֽרֹמֶ֥שֶׂת עַל־הָאָֽרֶץ׃ \Rashi{\textbf{וכבשה.}חסר וי"ו ללמדך שהזכר כובש את הנקבה, שלא תהא יצאנית, ועוד ללמדך שהאיש שדרכו לכבש מצוה על פריה ורביה ולא האשה:}%
% ---------- פרק 1 | פסוק 29 ----------
 \textbf{כט} וַיֹּ֣אמֶר אֱלֹהִ֗ים הִנֵּה֩ נָתַ֨תִּי לָכֶ֜ם אֶת־כׇּל־עֵ֣שֶׂב ׀ זֹרֵ֣עַ זֶ֗רַע אֲשֶׁר֙ עַל־פְּנֵ֣י כׇל־הָאָ֔רֶץ וְאֶת־כׇּל־הָעֵ֛ץ אֲשֶׁר־בּ֥וֹ פְרִי־עֵ֖ץ זֹרֵ֣עַ זָ֑רַע לָכֶ֥ם יִֽהְיֶ֖ה לְאׇכְלָֽה׃ \Rashi{\textbf{לכם יהיה לאכלה ולכל חית הארץ.}השוה להם בהמות וחיות למאכל, ולא הרשה לאדם ולאשתו להמית בריה ולאכל בשר, אך כל ירק עשב יאכלו יחד כלם, וכשבאו בני נח התיר להם בשר, שנאמר כל רמש אשר הוא חי וגו' כירק עשב שהתרתי לאדם הראשון נתתי לכם את כל:}%
% ---------- פרק 1 | פסוק 30 ----------
 \textbf{ל} וּֽלְכׇל־חַיַּ֣ת הָ֠אָ֠רֶץ וּלְכׇל־ע֨וֹף הַשָּׁמַ֜יִם וּלְכֹ֣ל ׀ רוֹמֵ֣שׂ עַל־הָאָ֗רֶץ אֲשֶׁר־בּוֹ֙ נֶ֣פֶשׁ חַיָּ֔ה אֶת־כׇּל־יֶ֥רֶק עֵ֖שֶׂב לְאׇכְלָ֑ה וַֽיְהִי־כֵֽן׃ %
% ---------- פרק 1 | פסוק 31 ----------
 \textbf{לא} וַיַּ֤רְא אֱלֹהִים֙ אֶת־כׇּל־אֲשֶׁ֣ר עָשָׂ֔ה וְהִנֵּה־ט֖וֹב מְאֹ֑ד וַֽיְהִי־עֶ֥רֶב וַֽיְהִי־בֹ֖קֶר י֥וֹם הַשִּׁשִּֽׁי׃ \{פ\} \Rashi{\textbf{יום הששי.}הוסיף ה' בששי בגמר מעשה בראשית, לומר שהתנה עמהם, על מנת שיקבלו עליהם ישראל חמשה חמשי תורה. דבר אחר יום הששי, כלם תלוים ועומדים עד יום הששי, הוא ו' בסיון המוכן למתן תורה:}%
% ---------- פרק 2 | פסוק 1 ----------
 \textbf{א} וַיְכֻלּ֛וּ הַשָּׁמַ֥יִם וְהָאָ֖רֶץ וְכׇל־צְבָאָֽם׃ %
% ---------- פרק 2 | פסוק 2 ----------
 \textbf{ב} וַיְכַ֤ל אֱלֹהִים֙ בַּיּ֣וֹם הַשְּׁבִיעִ֔י מְלַאכְתּ֖וֹ אֲשֶׁ֣ר עָשָׂ֑ה וַיִּשְׁבֹּת֙ בַּיּ֣וֹם הַשְּׁבִיעִ֔י מִכׇּל־מְלַאכְתּ֖וֹ אֲשֶׁ֥ר עָשָֽׂה׃ \Rashi{\textbf{ויכל אלהים ביום השביעי.}ר' שמעון אומר בשר ודם, שאינו יודע עתיו ורגעיו צריך להוסיף מחל על הקדש, הקב"ה שיודע עתיו ורגעיו, נכנס בו כחוט השערה, ונראה כאלו כלה בו ביום. ד"א מה היה העולם חסר? מנוחה, באת שבת באת מנוחה, כלתה ונגמרה המלאכה:}%
% ---------- פרק 2 | פסוק 3 ----------
 \textbf{ג} וַיְבָ֤רֶךְ אֱלֹהִים֙ אֶת־י֣וֹם הַשְּׁבִיעִ֔י וַיְקַדֵּ֖שׁ אֹת֑וֹ כִּ֣י ב֤וֹ שָׁבַת֙ מִכׇּל־מְלַאכְתּ֔וֹ אֲשֶׁר־בָּרָ֥א אֱלֹהִ֖ים לַעֲשֽׂוֹת׃ \{פ\} \Rashi{\textbf{ויברך, ויקדש.}ברכו במן, שכל ימות השבוע יורד להם עמר לגלגלת, ובששי לחם משנה, וקדשו במן, שלא ירד כלל בשבת, והמקרא כתוב על העתיד: אשר ברא אלהים לעשות. המלאכה שהיתה ראויה לעשות בשבת, כפל ועשאה בששי כמו שמפרש בב"ר:}%
% ---------- פרק 2 | פסוק 4 ----------
 \textbf{ד} אֵ֣לֶּה תוֹלְד֧וֹת הַשָּׁמַ֛יִם וְהָאָ֖רֶץ בְּהִבָּֽרְאָ֑ם בְּי֗וֹם עֲשׂ֛וֹת יְהֹוָ֥ה אֱלֹהִ֖ים אֶ֥רֶץ וְשָׁמָֽיִם׃ \Rashi{\textbf{אלה.}האמורים למעלה: תולדות השמים והארץ בהבראם ביום עשות ה'. למדך שכלם נבראו בראשון. דבר אחר בהבראם, בה' בראם, שנ' ביה ה' צור עולמים (ישעיהו כ"ו), בב' אותיות הללו של השם יצר שני עולמים, ולמדך כאן שהעולם הזה נברא בה"א, רמז שירדו למטה לראות שחת כה"א זאת שסתומה מכל צדדים ופתוחה למטה לרדת דרך שם:}%
% ---------- פרק 2 | פסוק 5 ----------
 \textbf{ה} וְכֹ֣ל ׀ שִׂ֣יחַ הַשָּׂדֶ֗ה טֶ֚רֶם יִֽהְיֶ֣ה בָאָ֔רֶץ וְכׇל־עֵ֥שֶׂב הַשָּׂדֶ֖ה טֶ֣רֶם יִצְמָ֑ח כִּי֩ לֹ֨א הִמְטִ֜יר יְהֹוָ֤ה אֱלֹהִים֙ עַל־הָאָ֔רֶץ וְאָדָ֣ם אַ֔יִן לַֽעֲבֹ֖ד אֶת־הָֽאֲדָמָֽה׃ \Rashi{\textbf{טרם יהיה בארץ.}כל טרם שבמקרא לשון עד לא הוא, ואינו לשון קדם, ואינו נפעל לומר הטרים, כאשר יאמר הקדים וזה מוכיח, ועוד אחר כי טרם תיראון (שמות ט) עדין לא תיראון, ואף זה תפרש עדין לא היה בארץ, כשנגמרה בריאת העולם בששי קדם שנברא אדם, וכל עשב השדה טרם יצמח, עדין לא צמח, ובג' שכתוב תוצא הארץ, לא יצאו, אלא על פתח קרקע עמדו עד יום ששי: כי לא המטיר. ומ"ט לא המטיר? לפי שאדם אין לעבד את האדמה, ואין מכיר בטובתם של גשמים, וכשבא אדם וידע שהם צרך לעולם, התפלל עליהם וירדו וצמחו האילנות והדשאים: ה' אלהים. ה' הוא שמו, אלהים שהוא שליט ושופט על כל, וכן פרוש זה בכל מקום לפי פשוטו, ה' שהוא אלהים:}%
% ---------- פרק 2 | פסוק 6 ----------
 \textbf{ו} וְאֵ֖ד יַֽעֲלֶ֣ה מִן־הָאָ֑רֶץ וְהִשְׁקָ֖ה אֶֽת־כׇּל־פְּנֵ֥י הָֽאֲדָמָֽה׃ \Rashi{\textbf{ואד יעלה.}לענין בריאתו של אדם, העלה התהום והשקה עננים לשרות העפר ונברא אדם, כגבל זה, שנותן מים ואח"כ לש את העסה, אף כאן, והשקה, ואח"כ וייצר:}%
% ---------- פרק 2 | פסוק 7 ----------
 \textbf{ז} וַיִּ֩יצֶר֩ יְהֹוָ֨ה אֱלֹהִ֜ים אֶת־הָֽאָדָ֗ם עָפָר֙ מִן־הָ֣אֲדָמָ֔ה וַיִּפַּ֥ח בְּאַפָּ֖יו נִשְׁמַ֣ת חַיִּ֑ים וַיְהִ֥י הָֽאָדָ֖ם לְנֶ֥פֶשׁ חַיָּֽה׃ \Rashi{\textbf{וייצר.}שתי יצירות, יצירה לעולם הזה ויצירה לתחית המתים, אבל בבהמה שאינה עומדת לדין לא נכתב ביצירתה שני יודי"ן: עפר מן האדמה. צבר עפרו מכל האדמה מארבע רוחות, שכל מקום שימות שם תהא קולטתו לקבורה. ד"א נטל עפרו ממקום שנאמר בו מזבח אדמה תעשה לי (שמות כ'), הלואי תהיה לו כפרה ויוכל לעמד: ויפח באפיו. עשאו מן התחתונים ומן העליונים, גוף מן התחתונים ונשמה מן העליונים, לפי שביום ראשון נבראו שמים וארץ, בשני ברא רקיע לעליונים, בשלישי תראה היבשה לתחתונים, ברביעי ברא מאורות לעליונים, בחמישי ישרצו המים לתחתונים, הזקק בששי לבראות בו בעליונים ובתחתונים, ואם לאו יש קנאה במעשה בראשית, שיהיו אלו רבים על אלו בבריאת יום אחד: לנפש חיה. אף בהמה וחיה נקראו נפש חיה, אך זו של אדם חיה שבכלן, שנתוסף בו דעה ודבור:}%
% ---------- פרק 2 | פסוק 8 ----------
 \textbf{ח} וַיִּטַּ֞ע יְהֹוָ֧ה אֱלֹהִ֛ים גַּן־בְּעֵ֖דֶן מִקֶּ֑דֶם וַיָּ֣שֶׂם שָׁ֔ם אֶת־הָֽאָדָ֖ם אֲשֶׁ֥ר יָצָֽר׃ \Rashi{\textbf{מקדם.}במזרחו של עדן נטע את הגן. ואם תאמר, הרי כבר כתב ויברא את האדם וגו', ראיתי בבריתא של ר"א בנו של ר' יוסי הגלילי מל"ב מדות שהתורה נדרשת, וזו אחת מהן, כלל שלאחריו מעשה הוא פרטו של ראשון, ויברא את האדם זהו כלל, סתם בריאתו מהיכן וסתם מעשיו, חזר ופרש וייצר ה' אלהים וגו', ויצמח לו גן עדן ויניחהו בגן עדן ויפל עליו תרדמה, השומע סבור שהוא מעשה אחר, ואינו אלא פרטו של ראשון; וכן אצל הבהמה חזר וכתב ויצר ה' וגו' מן האדמה כל חית השדה, כדי לפרש ויבא אל האדם לקרות שם וללמד על העופות שנבראו מן הרקק:}%
% ---------- פרק 2 | פסוק 9 ----------
 \textbf{ט} וַיַּצְמַ֞ח יְהֹוָ֤ה אֱלֹהִים֙ מִן־הָ֣אֲדָמָ֔ה כׇּל־עֵ֛ץ נֶחְמָ֥ד לְמַרְאֶ֖ה וְט֣וֹב לְמַאֲכָ֑ל וְעֵ֤ץ הַֽחַיִּים֙ בְּת֣וֹךְ הַגָּ֔ן וְעֵ֕ץ הַדַּ֖עַת ט֥וֹב וָרָֽע׃ \Rashi{\textbf{ויצמח.}לענין הגן הכתוב מדבר: בתוך הגן. באמצע:}%
% ---------- פרק 2 | פסוק 10 ----------
 \textbf{י} וְנָהָר֙ יֹצֵ֣א מֵעֵ֔דֶן לְהַשְׁק֖וֹת אֶת־הַגָּ֑ן וּמִשָּׁם֙ יִפָּרֵ֔ד וְהָיָ֖ה לְאַרְבָּעָ֥ה רָאשִֽׁים׃ %
% ---------- פרק 2 | פסוק 11 ----------
 \textbf{יא} שֵׁ֥ם הָֽאֶחָ֖ד פִּישׁ֑וֹן ה֣וּא הַסֹּבֵ֗ב אֵ֚ת כׇּל־אֶ֣רֶץ הַֽחֲוִילָ֔ה אֲשֶׁר־שָׁ֖ם הַזָּהָֽב׃ \Rashi{\textbf{פישון.}הוא נילוס נהר מצרים, ועל שם שמימיו מתברכין ועולין ומשקין את הארץ נקרא פישון כמו ופשו פרשיו (חבקוק א') ד"א פישון, שהוא מגדל פשתן, שנאמר על מצרים (ישעיהו י"ט) ובשו עבדי פשתים:}%
% ---------- פרק 2 | פסוק 12 ----------
 \textbf{יב} וּֽזְהַ֛ב הָאָ֥רֶץ הַהִ֖וא ט֑וֹב שָׁ֥ם הַבְּדֹ֖לַח וְאֶ֥בֶן הַשֹּֽׁהַם׃ %
% ---------- פרק 2 | פסוק 13 ----------
 \textbf{יג} וְשֵֽׁם־הַנָּהָ֥ר הַשֵּׁנִ֖י גִּיח֑וֹן ה֣וּא הַסּוֹבֵ֔ב אֵ֖ת כׇּל־אֶ֥רֶץ כּֽוּשׁ׃ \Rashi{\textbf{גיחון.}שהיה הולך והומה, והמיתו גדולה מאד, כמו כי יגח (שמות כ"א), שמנגח והולך והומה:}%
% ---------- פרק 2 | פסוק 14 ----------
 \textbf{יד} וְשֵׁ֨ם הַנָּהָ֤ר הַשְּׁלִישִׁי֙ חִדֶּ֔קֶל ה֥וּא הַֽהֹלֵ֖ךְ קִדְמַ֣ת אַשּׁ֑וּר וְהַנָּהָ֥ר הָֽרְבִיעִ֖י ה֥וּא פְרָֽת׃ \Rashi{\textbf{חדקל.}שמימיו חדין וקלין: פרת. שמימיו פרין ורבין ומברין את האדם: כוש ואשור. עדין לא היו, וכתב המקרא על שם העתיד: קדמת אשור. למזרחה של אשור: הוא פרת. החשוב על כלם, הנזכר על שם א"י:}%
% ---------- פרק 2 | פסוק 15 ----------
 \textbf{טו} וַיִּקַּ֛ח יְהֹוָ֥ה אֱלֹהִ֖ים אֶת־הָֽאָדָ֑ם וַיַּנִּחֵ֣הוּ בְגַן־עֵ֔דֶן לְעׇבְדָ֖הּ וּלְשׇׁמְרָֽהּ׃ \Rashi{\textbf{ויקח.}לקחו בדברים נאים ופתהו לכנס (בראשית רבה):}%
% ---------- פרק 2 | פסוק 16 ----------
 \textbf{טז} וַיְצַו֙ יְהֹוָ֣ה אֱלֹהִ֔ים עַל־הָֽאָדָ֖ם לֵאמֹ֑ר מִכֹּ֥ל עֵֽץ־הַגָּ֖ן אָכֹ֥ל תֹּאכֵֽל׃ %
% ---------- פרק 2 | פסוק 17 ----------
 \textbf{יז} וּמֵעֵ֗ץ הַדַּ֙עַת֙ ט֣וֹב וָרָ֔ע לֹ֥א תֹאכַ֖ל מִמֶּ֑נּוּ כִּ֗י בְּי֛וֹם אֲכׇלְךָ֥ מִמֶּ֖נּוּ מ֥וֹת תָּמֽוּת׃ %
% ---------- פרק 2 | פסוק 18 ----------
 \textbf{יח} וַיֹּ֙אמֶר֙ יְהֹוָ֣ה אֱלֹהִ֔ים לֹא־ט֛וֹב הֱי֥וֹת הָֽאָדָ֖ם לְבַדּ֑וֹ אֶֽעֱשֶׂה־לּ֥וֹ עֵ֖זֶר כְּנֶגְדּֽוֹ׃ \Rashi{\textbf{לא טוב היות וגו'.}שלא יאמרו שתי רשויות הן, הקב"ה יחיד בעליונים ואין לו זוג, וזה יחיד בתחתונים ואין לו זוג (בראשית רבה): עזר כנגדו. זכה – עזר; לא זכה – כנגדו להלחם:}%
% ---------- פרק 2 | פסוק 19 ----------
 \textbf{יט} וַיִּ֩צֶר֩ יְהֹוָ֨ה אֱלֹהִ֜ים מִן־הָֽאֲדָמָ֗ה כׇּל־חַיַּ֤ת הַשָּׂדֶה֙ וְאֵת֙ כׇּל־ע֣וֹף הַשָּׁמַ֔יִם וַיָּבֵא֙ אֶל־הָ֣אָדָ֔ם לִרְא֖וֹת מַה־יִּקְרָא־ל֑וֹ וְכֹל֩ אֲשֶׁ֨ר יִקְרָא־ל֧וֹ הָֽאָדָ֛ם נֶ֥פֶשׁ חַיָּ֖ה ה֥וּא שְׁמֽוֹ׃ \Rashi{\textbf{וייצר מן האדמה.}היא יצירה היא עשיה האמורה למעלה ויעש אלהים את חית הארץ וגו', אלא בא ופרש שהעופות נבראו מן הרקק, לפי שאמר למעלה מן המים נבראו וכאן אמר מן הארץ נבראו (חולין כ"ז), ועוד למדך כאן שבשעת יצירתן מיד בו ביום הביאם אל האדם לקרות להם שם. ובדברי אגדה יצירה זו לשון רדוי וכבוש, כמו כי תצור אל עיר (דברים כ"ד), שכבשן תחת ידו של אדם (בראשית רבה): וכל אשר יקרא לו האדם נפש חיה וגו'. סרסהו ופרשהו, כל נפש חיה אשר יקרא לו האדם שם, הוא שמו לעולם:}%
% ---------- פרק 2 | פסוק 20 ----------
 \textbf{כ} וַיִּקְרָ֨א הָֽאָדָ֜ם שֵׁמ֗וֹת לְכׇל־הַבְּהֵמָה֙ וּלְע֣וֹף הַשָּׁמַ֔יִם וּלְכֹ֖ל חַיַּ֣ת הַשָּׂדֶ֑ה וּלְאָדָ֕ם לֹֽא־מָצָ֥א עֵ֖זֶר כְּנֶגְדּֽוֹ׃ \Rashi{\textbf{ולאדם לא מצא עזר.}ויפל ה' אלהים תרדמה. כשהביאן הביאן לפניו כל מין ומין זכר ונקבה, אמר לכלם יש בן זוג ולי אין בן זוג, מיד ויפל (בראשית רבה):}%
% ---------- פרק 2 | פסוק 21 ----------
 \textbf{כא} וַיַּפֵּל֩ יְהֹוָ֨ה אֱלֹהִ֧ים ׀ תַּרְדֵּמָ֛ה עַל־הָאָדָ֖ם וַיִּישָׁ֑ן וַיִּקַּ֗ח אַחַת֙ מִצַּלְעֹתָ֔יו וַיִּסְגֹּ֥ר בָּשָׂ֖ר תַּחְתֶּֽנָּה׃ \Rashi{\textbf{מצלעותיו.}מסטריו, כמו ולצלע המשכן (שמות כ"ו), זהו שאמרו שני פרצופים נבראו (עירובין י"ז): ויסגור. מקום החתך: ויישן ויקח. שלא יראה חתיכת הבשר שממנו נבראת ותתבזה עליו:}%
% ---------- פרק 2 | פסוק 22 ----------
 \textbf{כב} וַיִּ֩בֶן֩ יְהֹוָ֨ה אֱלֹהִ֧ים ׀ אֶֽת־הַצֵּלָ֛ע אֲשֶׁר־לָקַ֥ח מִן־הָֽאָדָ֖ם לְאִשָּׁ֑ה וַיְבִאֶ֖הָ אֶל־הָֽאָדָֽם׃ \Rashi{\textbf{ויבן.}כבנין, רחבה מלמטה וקצרה מלמעלה לקבל הולד כאוצר של חטים, שהוא רחב מלמטה וקצר מלמעלה שלא יכביד משאו על קירותיו: ויבן את הצלע לאשה. להיות אשה, כמו ויעש אותו גדעון לאפוד (שופטים ח'), להיות אפוד:}%
% ---------- פרק 2 | פסוק 23 ----------
 \textbf{כג} וַיֹּ֘אמֶר֮ הָֽאָדָם֒ זֹ֣את הַפַּ֗עַם עֶ֚צֶם מֵֽעֲצָמַ֔י וּבָשָׂ֖ר מִבְּשָׂרִ֑י לְזֹאת֙ יִקָּרֵ֣א אִשָּׁ֔ה כִּ֥י מֵאִ֖ישׁ לֻֽקְחָה־זֹּֽאת׃ \Rashi{\textbf{זאת הפעם.}מלמד שבא אדם על כל בהמה וחיה, ולא נתקררה דעתו בהם (יבמות ס"ג): לזאת יקרא אשה כי מאיש וגו'. לשון נופל על לשון מכאן שנברא העולם בלשון הקדש (בראשית רבה):}%
% ---------- פרק 2 | פסוק 24 ----------
 \textbf{כד} עַל־כֵּן֙ יַֽעֲזׇב־אִ֔ישׁ אֶת־אָבִ֖יו וְאֶת־אִמּ֑וֹ וְדָבַ֣ק בְּאִשְׁתּ֔וֹ וְהָי֖וּ לְבָשָׂ֥ר אֶחָֽד׃ \Rashi{\textbf{על כן יעזב איש.}רוח הקדש אומרה כן לאסר על בני נח העריות (סנהדרין נ"ז): לבשר אחד. הולד נוצר על ידי שניהם ושם נעשה בשרם אחד (שם נ"ח):}%
% ---------- פרק 2 | פסוק 25 ----------
 \textbf{כה} וַיִּֽהְי֤וּ שְׁנֵיהֶם֙ עֲרוּמִּ֔ים הָֽאָדָ֖ם וְאִשְׁתּ֑וֹ וְלֹ֖א יִתְבֹּשָֽׁשׁוּ׃ \Rashi{\textbf{ולא יתבוששו.}שלא היו יודעים דרך צניעות להבחין בין טוב לרע, ואע"פ שנתנה בו דעה לקרות לו שמות, לא נתן בו יצר הרע עד אכלו מן העץ ונכנס בו יצר הרע וידע מה בין טוב לרע (בראשית רבה):}%
% ---------- פרק 3 | פסוק 1 ----------
 \textbf{א} וְהַנָּחָשׁ֙ הָיָ֣ה עָר֔וּם מִכֹּל֙ חַיַּ֣ת הַשָּׂדֶ֔ה אֲשֶׁ֥ר עָשָׂ֖ה יְהֹוָ֣ה אֱלֹהִ֑ים וַיֹּ֙אמֶר֙ אֶל־הָ֣אִשָּׁ֔ה אַ֚ף כִּֽי־אָמַ֣ר אֱלֹהִ֔ים לֹ֣א תֹֽאכְל֔וּ מִכֹּ֖ל עֵ֥ץ הַגָּֽן׃ \Rashi{\textbf{והנחש היה ערום.}מה ענין זה לכאן? היה לו לסמך ויעש לאדם ולאשתו כתנות עור וילבשם, אלא למדך מאיזו עצה קפץ הנחש עליהם, ראה אותם ערמים ועוסקים בתשמיש לעין כל ונתאוה לה: ערום מכל. לפי ערמתו וגדלתו היתה מפלתו, ערום מכל, ארור מכל (בראשית רבה): אף כי אמר וגו'. שמא אמר לכם לא תאכלו מכל וגו', ואע"פ שראה אותם אוכלים משאר פרות, הרבה עליה דברים כדי שתשיבנו ויבא לדבר באותו העץ:}%
% ---------- פרק 3 | פסוק 2 ----------
 \textbf{ב} וַתֹּ֥אמֶר הָֽאִשָּׁ֖ה אֶל־הַנָּחָ֑שׁ מִפְּרִ֥י עֵֽץ־הַגָּ֖ן נֹאכֵֽל׃ %
% ---------- פרק 3 | פסוק 3 ----------
 \textbf{ג} וּמִפְּרִ֣י הָעֵץ֮ אֲשֶׁ֣ר בְּתוֹךְ־הַגָּן֒ אָמַ֣ר אֱלֹהִ֗ים לֹ֤א תֹֽאכְלוּ֙ מִמֶּ֔נּוּ וְלֹ֥א תִגְּע֖וּ בּ֑וֹ פֶּן־תְּמֻתֽוּן׃ \Rashi{\textbf{ולא תגעו בו.}הוסיפה על הצווי, לפיכך באה לידי גרעון, הוא שנאמר אל תוסף על דבריו (משלי ל'):}%
% ---------- פרק 3 | פסוק 4 ----------
 \textbf{ד} וַיֹּ֥אמֶר הַנָּחָ֖שׁ אֶל־הָֽאִשָּׁ֑ה לֹֽא־מ֖וֹת תְּמֻתֽוּן׃ \Rashi{\textbf{לא מות תמותון.}דחפה עד שנגעה בו, אמר לה כשם שאין מיתה בנגיעה, כך אין מיתה באכילה (בראשית רבה):}%
% ---------- פרק 3 | פסוק 5 ----------
 \textbf{ה} כִּ֚י יֹדֵ֣עַ אֱלֹהִ֔ים כִּ֗י בְּיוֹם֙ אֲכׇלְכֶ֣ם מִמֶּ֔נּוּ וְנִפְקְח֖וּ עֵֽינֵיכֶ֑ם וִהְיִיתֶם֙ כֵּֽאלֹהִ֔ים יֹדְעֵ֖י ט֥וֹב וָרָֽע׃ \Rashi{\textbf{כי ידע.}כל אמן שונא את בני אמנותו, מן העץ אכל וברא את העולם (בראשית רבה): והייתם כאלהים. יוצרי עולמות (בראשית רבה):}%
% ---------- פרק 3 | פסוק 6 ----------
 \textbf{ו} וַתֵּ֣רֶא הָֽאִשָּׁ֡ה כִּ֣י טוֹב֩ הָעֵ֨ץ לְמַאֲכָ֜ל וְכִ֧י תַֽאֲוָה־ה֣וּא לָעֵינַ֗יִם וְנֶחְמָ֤ד הָעֵץ֙ לְהַשְׂכִּ֔יל וַתִּקַּ֥ח מִפִּרְי֖וֹ וַתֹּאכַ֑ל וַתִּתֵּ֧ן גַּם־לְאִישָׁ֛הּ עִמָּ֖הּ וַיֹּאכַֽל׃ \Rashi{\textbf{ותרא האשה.}ראתה דבריו של נחש והנאו לה והאמינתו (בראשית רבה): כי טוב העץ. להיות כאלהים: וכי תאוה הוא לעינים. כמו שאמר לה ונפקחו עיניכם: ונחמד להשכיל. כמו שאמר לה יודעי טוב ורע: ותתן גם לאישה. שלא תמות היא ויחיה הוא, וישא אחרת: גם. לרבות בהמה וחיה:}%
% ---------- פרק 3 | פסוק 7 ----------
 \textbf{ז} וַתִּפָּקַ֙חְנָה֙ עֵינֵ֣י שְׁנֵיהֶ֔ם וַיֵּ֣דְע֔וּ כִּ֥י עֵֽירֻמִּ֖ם הֵ֑ם וַֽיִּתְפְּרוּ֙ עֲלֵ֣ה תְאֵנָ֔ה וַיַּעֲשׂ֥וּ לָהֶ֖ם חֲגֹרֹֽת׃ \Rashi{\textbf{ותפקחנה וגו'.}לענין החכמה דבר הכתוב ולא לענין ראיה ממש, וסוף המקרא מוכיח: וידעו כי עירמים הם. אף הסומא יודע כשהוא ערום, אלא מהו וידעו כי עירמם הם? מצוה אחת היתה בידם ונתערטלו הימנה (בראשית רבה): עלה תאנה. הוא העץ שאכלו ממנו, בדבר שנתקלקלו בו נתקנו (סנהדרין ע'), אבל שאר העצים מנעום מלטל עליהם. ומפני מה לא נתפרסם העץ? שאין הקב"ה חפץ להונות בריה, שלא יכלימוה ויאמרו זהו שלקה העולם על ידו, מדרש רבי תנחומא:}%
% ---------- פרק 3 | פסוק 8 ----------
 \textbf{ח} וַֽיִּשְׁמְע֞וּ אֶת־ק֨וֹל יְהֹוָ֧ה אֱלֹהִ֛ים מִתְהַלֵּ֥ךְ בַּגָּ֖ן לְר֣וּחַ הַיּ֑וֹם וַיִּתְחַבֵּ֨א הָֽאָדָ֜ם וְאִשְׁתּ֗וֹ מִפְּנֵי֙ יְהֹוָ֣ה אֱלֹהִ֔ים בְּת֖וֹךְ עֵ֥ץ הַגָּֽן׃ \Rashi{\textbf{וישמעו.}יש מדרשי אגדה רבים וכבר סדרום רבותינו על מכונם בב"ר ובשאר מדרשות; ואני לא באתי אלא לפשוטו של מקרא ולאגדה המישבת דברי המקרא דבר דבור על אפניו: וישמעו. מה שמעו? שמעו את קול הקב"ה שהיה מתהלך בגן: לרוח היום. לאותו רוח שהשמש באה לשם, וזו היא מערבית; שלפנות ערב חמה במערב, והם סרחו בעשירית (סנהדרין ל"ח):}%
% ---------- פרק 3 | פסוק 9 ----------
 \textbf{ט} וַיִּקְרָ֛א יְהֹוָ֥ה אֱלֹהִ֖ים אֶל־הָֽאָדָ֑ם וַיֹּ֥אמֶר ל֖וֹ אַיֶּֽכָּה׃ \Rashi{\textbf{איכה.}יודע היה היכן הוא אלא לכנס עמו בדברים, שלא יהא נבהל להשיב אם יענישהו פתאום (בראשית רבה), וכן בקין אמר לו אי הבל אחיך (בראשית ד), וכן בבלעם מי האנשים האלה עמך (במדבר כ"ב), לכנס עמהם בדברים, וכן בחזקיה בשלוחי מראדך בלאדן:}%
% ---------- פרק 3 | פסוק 10 ----------
 \textbf{י} וַיֹּ֕אמֶר אֶת־קֹלְךָ֥ שָׁמַ֖עְתִּי בַּגָּ֑ן וָאִירָ֛א כִּֽי־עֵירֹ֥ם אָנֹ֖כִי וָאֵחָבֵֽא׃ %
% ---------- פרק 3 | פסוק 11 ----------
 \textbf{יא} וַיֹּ֕אמֶר מִ֚י הִגִּ֣יד לְךָ֔ כִּ֥י עֵירֹ֖ם אָ֑תָּה הֲמִן־הָעֵ֗ץ אֲשֶׁ֧ר צִוִּיתִ֛יךָ לְבִלְתִּ֥י אֲכׇל־מִמֶּ֖נּוּ אָכָֽלְתָּ׃ \Rashi{\textbf{מי הגיד לך.}מאין לך לדעת מה בשת יש בעומד ערם? המן העץ. בתמיה:}%
% ---------- פרק 3 | פסוק 12 ----------
 \textbf{יב} וַיֹּ֖אמֶר הָֽאָדָ֑ם הָֽאִשָּׁה֙ אֲשֶׁ֣ר נָתַ֣תָּה עִמָּדִ֔י הִ֛וא נָֽתְנָה־לִּ֥י מִן־הָעֵ֖ץ וָאֹכֵֽל׃ \Rashi{\textbf{אשר נתתה עמדי.}כאן כפר בטובה (עבודה זרה ה):}%
% ---------- פרק 3 | פסוק 13 ----------
 \textbf{יג} וַיֹּ֨אמֶר יְהֹוָ֧ה אֱלֹהִ֛ים לָאִשָּׁ֖ה מַה־זֹּ֣את עָשִׂ֑ית וַתֹּ֙אמֶר֙ הָֽאִשָּׁ֔ה הַנָּחָ֥שׁ הִשִּׁיאַ֖נִי וָאֹכֵֽל׃ \Rashi{\textbf{השיאני.}הטעני, כמו אל ישיא אתכם חזקיהו (דברי הימים ב' ל"ב):}%
% ---------- פרק 3 | פסוק 14 ----------
 \textbf{יד} וַיֹּ֩אמֶר֩ יְהֹוָ֨ה אֱלֹהִ֥ים ׀ אֶֽל־הַנָּחָשׁ֮ כִּ֣י עָשִׂ֣יתָ זֹּאת֒ אָר֤וּר אַתָּה֙ מִכׇּל־הַבְּהֵמָ֔ה וּמִכֹּ֖ל חַיַּ֣ת הַשָּׂדֶ֑ה עַל־גְּחֹנְךָ֣ תֵלֵ֔ךְ וְעָפָ֥ר תֹּאכַ֖ל כׇּל־יְמֵ֥י חַיֶּֽיךָ׃ \Rashi{\textbf{כי עשית זאת.}מכאן שאין מהפכים בזכותו של מסית, שאלו שאלו למה עשית זאת? היה לו להשיב דברי הרב ודברי התלמיד דברי מי שומעין? (סנהדרין כ"ט): מכל הבהמה ומכל חית השדה. אם מבהמה נתקלל מחיה לא כל שכן, העמידו רבותינו מדרש זה במסכת בכורות (דף ח') ללמד שימי עבורו של נחש שבע שנים: על גחונך תלך. רגלים היו לו ונקצצו:}%
% ---------- פרק 3 | פסוק 15 ----------
 \textbf{טו} וְאֵיבָ֣ה ׀ אָשִׁ֗ית בֵּֽינְךָ֙ וּבֵ֣ין הָֽאִשָּׁ֔ה וּבֵ֥ין זַרְעֲךָ֖ וּבֵ֣ין זַרְעָ֑הּ ה֚וּא יְשׁוּפְךָ֣ רֹ֔אשׁ וְאַתָּ֖ה תְּשׁוּפֶ֥נּוּ עָקֵֽב׃ \{ס\} \Rashi{\textbf{ואיבה אשית.}אתה לא נתכונת אלא שימות אדם כשיאכל הוא תחלה ותשא את חוה, ולא באת לדבר אל חוה תחלה אלא לפי שהנשים קלות להתפתות ויודעות לפתות את בעליהן, לפיכך ואיבה אשית: ישופך. יכתתך (סוטה ט') כמו ואכת אתו (דברים ט), ותרגומו ושפית יתי': ואתה תשופנו עקב. לא יהא לך קומה ותשכנו בעקבו, ואף משם תמיתנו, ולשון תשופנו כמו נשף בהם (ישעיהו מ') כשהנחש בא לנשך הוא נושף כמין שריקה ולפי שהלשון נופל על הלשון כתיב לשון נשיפה בשניהם:}%
% ---------- פרק 3 | פסוק 16 ----------
 \textbf{טז} אֶֽל־הָאִשָּׁ֣ה אָמַ֗ר הַרְבָּ֤ה אַרְבֶּה֙ עִצְּבוֹנֵ֣ךְ וְהֵֽרֹנֵ֔ךְ בְּעֶ֖צֶב תֵּֽלְדִ֣י בָנִ֑ים וְאֶל־אִישֵׁךְ֙ תְּשׁ֣וּקָתֵ֔ךְ וְה֖וּא יִמְשׇׁל־בָּֽךְ׃ \{ס\} \Rashi{\textbf{עצבונך.}זה צער גדול בנים (עירובין ק'): והרנך. זה צער העבור: בעצב תלדי בנים. זה צער הלדה: ואל אישך תשוקתך. לתשמיש, ואעפ"כ אין לך מצח לתבעו בפה, אלא הוא ימשל בך, הכל ממנו ולא ממך: תשוקתך. תאותך, כמו נפש שוקקה (ישעיהו כ"ט):}%
% ---------- פרק 3 | פסוק 17 ----------
 \textbf{יז} וּלְאָדָ֣ם אָמַ֗ר כִּֽי־שָׁמַ֘עְתָּ֮ לְק֣וֹל אִשְׁתֶּ֒ךָ֒ וַתֹּ֙אכַל֙ מִן־הָעֵ֔ץ אֲשֶׁ֤ר צִוִּיתִ֙יךָ֙ לֵאמֹ֔ר לֹ֥א תֹאכַ֖ל מִמֶּ֑נּוּ אֲרוּרָ֤ה הָֽאֲדָמָה֙ בַּֽעֲבוּרֶ֔ךָ בְּעִצָּבוֹן֙ תֹּֽאכְלֶ֔נָּה כֹּ֖ל יְמֵ֥י חַיֶּֽיךָ׃ \Rashi{\textbf{ארורה האדמה בעבורך.}מעלה לך דברים ארורים, כגון זבובים, פרעושים ונמלים, משל ליוצא לתרבות רעה והבריות מקללות שדים שינק מהם:}%
% ---------- פרק 3 | פסוק 18 ----------
 \textbf{יח} וְק֥וֹץ וְדַרְדַּ֖ר תַּצְמִ֣יחַֽ לָ֑ךְ וְאָכַלְתָּ֖ אֶת־עֵ֥שֶׂב הַשָּׂדֶֽה׃ \Rashi{\textbf{וקוץ ודרדר תצמיח לך.}הארץ כשתזרענה מיני זרעים, תצמיח קוץ ודרדר קנדס ועכביות, והן נאכלין על ידי תקון (ביצה ל"ד): ואכלת את עשב השדה. ומה קללה היא זו? והלא בברכה נאמר לו הנה נתתי לכם את כל עשב זרע זרע וגו? אלא מה אמור כאן בראש הענין ארורה האדמה בעבורך בעצבון תאכלנה, ואחר העצבון וקוץ ודרדר תצמיח לך, כשתזרענה קטנית או ירקות גנה היא תצמיח לך קוצים ודרדרים ושאר עשבי השדה ועל כרחך תאכלם:}%
% ---------- פרק 3 | פסוק 19 ----------
 \textbf{יט} בְּזֵעַ֤ת אַפֶּ֙יךָ֙ תֹּ֣אכַל לֶ֔חֶם עַ֤ד שֽׁוּבְךָ֙ אֶל־הָ֣אֲדָמָ֔ה כִּ֥י מִמֶּ֖נָּה לֻקָּ֑חְתָּ כִּֽי־עָפָ֣ר אַ֔תָּה וְאֶל־עָפָ֖ר תָּשֽׁוּב׃ \Rashi{\textbf{בזעת אפיך.}לאחר שתטרח בו הרבה:}%
% ---------- פרק 3 | פסוק 20 ----------
 \textbf{כ} וַיִּקְרָ֧א הָֽאָדָ֛ם שֵׁ֥ם אִשְׁתּ֖וֹ חַוָּ֑ה כִּ֛י הִ֥וא הָֽיְתָ֖ה אֵ֥ם כׇּל־חָֽי׃ \Rashi{\textbf{ויקרא האדם.}חזר הכתוב לענינו הראשון ויקרא האדם שמות, ולא הפסיק אלא ללמדך שעל ידי קריאת שמות נזדוגה לו חוה, כמו שכתוב ולאדם לא מצא עזר כנגדו, לפיכך, ויפל תרדמה, ועל ידי שכתב ויהיו ערומים, סמך לו פרשת הנחש, להודיעך שמתוך שראה ערותה וראה אותם עסוקים בתשמיש נתאוה לה ובא עליהם במחשבה ובמרמה: חוה. נופל על לשון חיה, שמחיה את ולדותיה, כאשר תאמר מה הוה לאדם (קהלת ב'), בלשון היה:}%
% ---------- פרק 3 | פסוק 21 ----------
 \textbf{כא} וַיַּ֩עַשׂ֩ יְהֹוָ֨ה אֱלֹהִ֜ים לְאָדָ֧ם וּלְאִשְׁתּ֛וֹ כׇּתְנ֥וֹת ע֖וֹר וַיַּלְבִּשֵֽׁם׃ \{פ\} \Rashi{\textbf{כתנות עור.}יש דברי אגדה אומרים חלקים כצפרן היו מדבקים על עורן וי"א דבר הבא מן העור, כגון צמר הארנבים שהוא רך וחם ועשה להם כתנות ממנו:}%
% ---------- פרק 3 | פסוק 22 ----------
 \textbf{כב} וַיֹּ֣אמֶר ׀ יְהֹוָ֣ה אֱלֹהִ֗ים הֵ֤ן הָֽאָדָם֙ הָיָה֙ כְּאַחַ֣ד מִמֶּ֔נּוּ לָדַ֖עַת ט֣וֹב וָרָ֑ע וְעַתָּ֣ה ׀ פֶּן־יִשְׁלַ֣ח יָד֗וֹ וְלָקַח֙ גַּ֚ם מֵעֵ֣ץ הַֽחַיִּ֔ים וְאָכַ֖ל וָחַ֥י לְעֹלָֽם׃ \Rashi{\textbf{היה כאחד ממנו.}הרי הוא יחיד בתחתונים כמו שאני יחיד בעליונים, ומה היא יחידותו? לדעת טוב ורע, מה שאין כן בבהמה וחיה ועתה פן ישלח ידו. ומשיחיה לעולם הרי הוא קרוב להטעות הבריות אחריו, ולומר אף הוא אלוה. ויש מדרשי אגדה, אבל אין מישבין על פשוטו:}%
% ---------- פרק 3 | פסוק 23 ----------
 \textbf{כג} וַֽיְשַׁלְּחֵ֛הוּ יְהֹוָ֥ה אֱלֹהִ֖ים מִגַּן־עֵ֑דֶן לַֽעֲבֹד֙ אֶת־הָ֣אֲדָמָ֔ה אֲשֶׁ֥ר לֻקַּ֖ח מִשָּֽׁם׃ %
% ---------- פרק 3 | פסוק 24 ----------
 \textbf{כד} וַיְגָ֖רֶשׁ אֶת־הָֽאָדָ֑ם וַיַּשְׁכֵּן֩ מִקֶּ֨דֶם לְגַן־עֵ֜דֶן אֶת־הַכְּרֻבִ֗ים וְאֵ֨ת לַ֤הַט הַחֶ֙רֶב֙ הַמִּתְהַפֶּ֔כֶת לִשְׁמֹ֕ר אֶת־דֶּ֖רֶךְ עֵ֥ץ הַֽחַיִּֽים׃ \{ס\} \Rashi{\textbf{מקדם לגן עדן.}במזרחו של גן עדן חוץ לגן: את הכרובים. מלאכי חבלה: החרב המתהפכת. ולה להט לאים עליו מלכנס עוד לגן תרגום להט שנן, כמו שלף שננא ובלשון לעז לא"מא ומדרשי אגדה יש, ואני איני בא אלא לפשוטו:}%
% ---------- פרק 4 | פסוק 1 ----------
 \textbf{א} וְהָ֣אָדָ֔ם יָדַ֖ע אֶת־חַוָּ֣ה אִשְׁתּ֑וֹ וַתַּ֙הַר֙ וַתֵּ֣לֶד אֶת־קַ֔יִן וַתֹּ֕אמֶר קָנִ֥יתִי אִ֖ישׁ אֶת־יְהֹוָֽה׃ \Rashi{\textbf{והאדם ידע.}כבר קדם הענין של מעלה, קדם שחטא ונטרד מגן עדן, וכן ההריון והלדה, שאם כתב וידע אדם, נשמע שלאחר שנטרד היו לו בנים: קין. על שם קניתי: את ה'. כמו עם ה', כשברא אותי ואת אישי, הוא לבדו בראנו אבל בזה שתפים אנו עמו (נדה ל"א): את קין את אחיו את הבל. ג' אתים רבויים הם, מלמד, שתאומה נולדה עם קין, ועם הבל נולדו שתים, לכך נאמר ותסף (בראשית רבה):}%
% ---------- פרק 4 | פסוק 2 ----------
 \textbf{ב} וַתֹּ֣סֶף לָלֶ֔דֶת אֶת־אָחִ֖יו אֶת־הָ֑בֶל וַֽיְהִי־הֶ֙בֶל֙ רֹ֣עֵה צֹ֔אן וְקַ֕יִן הָיָ֖ה עֹבֵ֥ד אֲדָמָֽה׃ \Rashi{\textbf{רעה צאן.}לפי שנתקללה האדמה פרש לו מעבודתה:}%
% ---------- פרק 4 | פסוק 3 ----------
 \textbf{ג} וַֽיְהִ֖י מִקֵּ֣ץ יָמִ֑ים וַיָּבֵ֨א קַ֜יִן מִפְּרִ֧י הָֽאֲדָמָ֛ה מִנְחָ֖ה לַֽיהֹוָֽה׃ \Rashi{\textbf{מפרי האדמה.}מן הגרוע, ויש אגדה שאומרת זרע פשתן היה:}%
% ---------- פרק 4 | פסוק 4 ----------
 \textbf{ד} וְהֶ֨בֶל הֵבִ֥יא גַם־ה֛וּא מִבְּכֹר֥וֹת צֹאנ֖וֹ וּמֵֽחֶלְבֵהֶ֑ן וַיִּ֣שַׁע יְהֹוָ֔ה אֶל־הֶ֖בֶל וְאֶל־מִנְחָתֽוֹ׃ \Rashi{\textbf{וישע.}ויפן, וכן ואל מנחתו לא שעה – לא פנה, וכן ולא ישעה (ישעיהו י"ז) – לא יפנה, וכן שעה מעליו (איוב י"ד) פנה מעליו: וישע. ירדה אש ולחכה מנחתו:}%
% ---------- פרק 4 | פסוק 5 ----------
 \textbf{ה} וְאֶל־קַ֥יִן וְאֶל־מִנְחָת֖וֹ לֹ֣א שָׁעָ֑ה וַיִּ֤חַר לְקַ֙יִן֙ מְאֹ֔ד וַֽיִּפְּל֖וּ פָּנָֽיו׃ %
% ---------- פרק 4 | פסוק 6 ----------
 \textbf{ו} וַיֹּ֥אמֶר יְהֹוָ֖ה אֶל־קָ֑יִן לָ֚מָּה חָ֣רָה לָ֔ךְ וְלָ֖מָּה נָפְל֥וּ פָנֶֽיךָ׃ %
% ---------- פרק 4 | פסוק 7 ----------
 \textbf{ז} הֲל֤וֹא אִם־תֵּיטִיב֙ שְׂאֵ֔ת וְאִם֙ לֹ֣א תֵיטִ֔יב לַפֶּ֖תַח חַטָּ֣את רֹבֵ֑ץ וְאֵלֶ֙יךָ֙ תְּשׁ֣וּקָת֔וֹ וְאַתָּ֖ה תִּמְשׇׁל־בּֽוֹ׃ \Rashi{\textbf{הלא אם תיטיב.}כתרגומו פרושו: לפתח חטאת רובץ. לפתח קברך חטאך שמור: ואליך תשוקתו. של חטאת הוא יצר הרע, תמיד שוקק ומתאוה להכשילך: ואתה תמשל בו. אם תרצה תתגבר עליו:}%
% ---------- פרק 4 | פסוק 8 ----------
 \textbf{ח} וַיֹּ֥אמֶר קַ֖יִן אֶל־הֶ֣בֶל אָחִ֑יו וַֽיְהִי֙ בִּהְיוֹתָ֣ם בַּשָּׂדֶ֔ה וַיָּ֥קׇם קַ֛יִן אֶל־הֶ֥בֶל אָחִ֖יו וַיַּהַרְגֵֽהוּ׃ \Rashi{\textbf{ויאמר קין אל הבל.}נכנס עמו בדברי ריב ומצה להתעולל עליו להרגו. ויש בזה מדרשי אגדה, אך זה ישובו של מקרא:}%
% ---------- פרק 4 | פסוק 9 ----------
 \textbf{ט} וַיֹּ֤אמֶר יְהֹוָה֙ אֶל־קַ֔יִן אֵ֖י הֶ֣בֶל אָחִ֑יךָ וַיֹּ֙אמֶר֙ לֹ֣א יָדַ֔עְתִּי הֲשֹׁמֵ֥ר אָחִ֖י אָנֹֽכִי׃ \Rashi{\textbf{אי הבל אחיך.}לכנס עמו בדברי נחת, אולי ישוב ויאמר: אני הרגתיו וחטאתי לך: לא ידעתי. נעשה כגונב דעת העליונה: השומר אחי. לשון תמה הוא, וכן כל ה"א הנקודה בחטף פתח:}%
% ---------- פרק 4 | פסוק 10 ----------
 \textbf{י} וַיֹּ֖אמֶר מֶ֣ה עָשִׂ֑יתָ ק֚וֹל דְּמֵ֣י אָחִ֔יךָ צֹעֲקִ֥ים אֵלַ֖י מִן־הָֽאֲדָמָֽה׃ \Rashi{\textbf{דמי אחיך.}דמו ודם זרעיותיו. דבר אחר, שעשה בו פצעים הרבה, שלא היה יודע מהיכן נפשו יוצאה (סנהדרין ל"ז):}%
% ---------- פרק 4 | פסוק 11 ----------
 \textbf{יא} וְעַתָּ֖ה אָר֣וּר אָ֑תָּה מִן־הָֽאֲדָמָה֙ אֲשֶׁ֣ר פָּצְתָ֣ה אֶת־פִּ֔יהָ לָקַ֛חַת אֶת־דְּמֵ֥י אָחִ֖יךָ מִיָּדֶֽךָ׃ \Rashi{\textbf{ארור אתה מן האדמה.}יותר ממה שנתקללה היא כבר בעונה, וגם בזו הוסיפה לחטא, אשר פצתה את פיה לקחת את דמי אחיך וגו'. והנני מוסיף לה קללה אצלך, לא תסף תת כחה:}%
% ---------- פרק 4 | פסוק 12 ----------
 \textbf{יב} כִּ֤י תַֽעֲבֹד֙ אֶת־הָ֣אֲדָמָ֔ה לֹֽא־תֹסֵ֥ף תֵּת־כֹּחָ֖הּ לָ֑ךְ נָ֥ע וָנָ֖ד תִּֽהְיֶ֥ה בָאָֽרֶץ׃ \Rashi{\textbf{נע ונד.}אין לך רשות לדור במקום אחד:}%
% ---------- פרק 4 | פסוק 13 ----------
 \textbf{יג} וַיֹּ֥אמֶר קַ֖יִן אֶל־יְהֹוָ֑ה גָּד֥וֹל עֲוֺנִ֖י מִנְּשֹֽׂא*(בספרי ספרד ואשכנז מִנְּשֽׂוֹא)׃ \Rashi{\textbf{עוני מנשוא.}בתמיה, אתה טוען עליונים ותחתונים, ועוני אי אפשר לטען? (בראשית רבה)}%
% ---------- פרק 4 | פסוק 14 ----------
 \textbf{יד} הֵן֩ גֵּרַ֨שְׁתָּ אֹתִ֜י הַיּ֗וֹם מֵעַל֙ פְּנֵ֣י הָֽאֲדָמָ֔ה וּמִפָּנֶ֖יךָ אֶסָּתֵ֑ר וְהָיִ֜יתִי נָ֤ע וָנָד֙ בָּאָ֔רֶץ וְהָיָ֥ה כׇל־מֹצְאִ֖י יַֽהַרְגֵֽנִי׃ %
% ---------- פרק 4 | פסוק 15 ----------
 \textbf{טו} וַיֹּ֧אמֶר ל֣וֹ יְהֹוָ֗ה לָכֵן֙ כׇּל־הֹרֵ֣ג קַ֔יִן שִׁבְעָתַ֖יִם יֻקָּ֑ם וַיָּ֨שֶׂם יְהֹוָ֤ה לְקַ֙יִן֙ א֔וֹת לְבִלְתִּ֥י הַכּוֹת־אֹת֖וֹ כׇּל־מֹצְאֽוֹ׃ \Rashi{\textbf{לכן כל הורג קין.}זה אחד מן המקראות שקצרו דבריהם ורמזו ולא פרשו, לכן כל הרג קין לשון גערה, כה יעשה לו, כך וכך ענשו ולא פרש ענשו: שבעתים יקם. איני רוצה להנקם מקין עכשו, לסוף שבעה דורות אני נוקם נקמתי ממנו, שיעמד למך מבני בניו ויהרגהו, וסוף המקרא שאמר שבעתים יקם, והיא נקמת הבל מקין, למדנו שתחלת מקרא לשון גערה היא, שלא תהא בריה מזיקתו, וכיוצא בו ויאמר דוד כל מכה יבסי ויגע בצנור (שמואל ב ה'), ולא פרש מה יעשה לו, אבל דבר הכתוב ברמז כל מכה יבסי ויגע בצנור ויקרב אל השער ויכבשנו ואת העורים וגו', וגם אותם יכה על אשר אמרו העור והפסח לא יבא דוד אל תוך הבית – המכה את אלו אני אעשנו ראש ושר, כאן קצר דבריו, ובדברי הימים פרש יהיה לראש ולשר: וישם ה' לקין אות. חקק לו אות משמו במצחו. ס"א דבר אחר, כל מוצאי יהרגני, הבהמות והחיות, אבל בני אדם עדין לא היו שיירא מהם, רק אביו ואמו, ומהם לא היה ירא שיהרגוהו, אלא אמר עד עכשו היה פחדתי על כל החיות כמ"ש ומוראכם וגו', ועכשו בשביל עון זה לא ייראו ממני החיות ויהרגוני, מיד וישם ה' לקין אות – החזיר את מוראו על החיות:}%
% ---------- פרק 4 | פסוק 16 ----------
 \textbf{טז} וַיֵּ֥צֵא קַ֖יִן מִלִּפְנֵ֣י יְהֹוָ֑ה וַיֵּ֥שֶׁב בְּאֶֽרֶץ־נ֖וֹד קִדְמַת־עֵֽדֶן׃ \Rashi{\textbf{ויצא קין.}יצא בהכנעה כגונב דעת העליונה: בארץ נוד. בארץ שכל הגולים נדים שם: קדמת עדן. שם גלה אביו כשגרש מגן עדן, שנאמר וישכן מקדם לגן עדן את שמירת דרך מבוא הגן, שיש ללמד שהיה אדם שם, ומצינו רוח מזרחית קולטת בכל מקום את הרוצחים שנאמר אז יבדיל משה וגומר, מזרחה שמש (דברים ד'), דבר אחר בארץ נוד, כל מקום שהולך היתה הארץ מזדעזעה תחתיו, והבריות אומרים סורו מעליו, זהו שהרג את אחיו:}%
% ---------- פרק 4 | פסוק 17 ----------
 \textbf{יז} וַיֵּ֤דַע קַ֙יִן֙ אֶת־אִשְׁתּ֔וֹ וַתַּ֖הַר וַתֵּ֣לֶד אֶת־חֲנ֑וֹךְ וַֽיְהִי֙ בֹּ֣נֶה עִ֔יר וַיִּקְרָא֙ שֵׁ֣ם הָעִ֔יר כְּשֵׁ֖ם בְּנ֥וֹ חֲנֽוֹךְ׃ \Rashi{ויהי קין בונה עיר ויקרא שם העיר לזכר בנו חנוך:}%
% ---------- פרק 4 | פסוק 18 ----------
 \textbf{יח} וַיִּוָּלֵ֤ד לַֽחֲנוֹךְ֙ אֶת־עִירָ֔ד וְעִירָ֕ד יָלַ֖ד אֶת־מְחֽוּיָאֵ֑ל וּמְחִיָּיאֵ֗ל יָלַד֙ אֶת־מְת֣וּשָׁאֵ֔ל וּמְתוּשָׁאֵ֖ל יָלַ֥ד אֶת־לָֽמֶךְ׃ \Rashi{\textbf{ועירד ילד.}יש מקום שהוא אומר בזכר הוליד ויש מקום שהוא אומר ילד, שהלדה משמשת שתי לשונות, לדת האשה ניש"טרא בלעז וזריעת תולדות האיש אינזי"ראר בלעז. כשהוא אומר הוליד בלשון הפעיל, מדבר בלדת האשה – פלוני הוליד את אשתו בן או בת, כשהוא אומר ילד, מדבר בזריעת האיש:}%
% ---------- פרק 4 | פסוק 19 ----------
 \textbf{יט} וַיִּֽקַּֽח־ל֥וֹ לֶ֖מֶךְ שְׁתֵּ֣י נָשִׁ֑ים שֵׁ֤ם הָֽאַחַת֙ עָדָ֔ה וְשֵׁ֥ם הַשֵּׁנִ֖ית צִלָּֽה׃ \Rashi{\textbf{ויקח לו למך.}לא היה לו לפרש כל זה, אלא ללמדנו מסוף הענין, שקים הקב"ה הבטחתו שאמר שבעתים יקם קין, עמד למך לאחר שהוליד בנים ועשה דור שביעי והרג את קין, זהו שאמר כי איש הרגתי לפצעי וגומר: שתי נשים. כך היה דרכן של דור המבול אחת לפריה ורביה ואחת לתשמיש, זו שהיא לתשמיש משקה כוס של עקרין כדי שתעקר ומקשטת ככלה ומאכילה מעדנים, וחברתה נזופה ואבלה כאלמנה, וזהו שפרש איוב רעה עקרה לא תלד ואלמנה לא ייטיב (איוב כ"ד), כמו שמפרש באגדת חלק (ובב"ר): עדה. היא של פריה ורביה על שם שמגנה עליו ומוסרת מאצלו, עדה תרגום של סורה: צלה. היא של תשמיש, על שם שיושבת תמיד בצלו, דברי אגדה הם בב"ר:}%
% ---------- פרק 4 | פסוק 20 ----------
 \textbf{כ} וַתֵּ֥לֶד עָדָ֖ה אֶת־יָבָ֑ל ה֣וּא הָיָ֔ה אֲבִ֕י יֹשֵׁ֥ב אֹ֖הֶל וּמִקְנֶֽה׃ \Rashi{\textbf{אבי ישב אהל ומקנה.}הוא היה הראשון לרועי בהמות במדברות, יושב אהלים חדש כאן וחדש כאן בשביל מרעה צאנו, וכשכלה המרעה במקום זה הולך ותוקע אהלו במקום אחר. ומדרש אגדה בונה בתים לע"ז, כמה דאת אמר סמל הקנאה המקנה (יחזקאל ח'), וכן אחיו תופש כנור ועוגב לזמר לע"ז:}%
% ---------- פרק 4 | פסוק 21 ----------
 \textbf{כא} וְשֵׁ֥ם אָחִ֖יו יוּבָ֑ל ה֣וּא הָיָ֔ה אֲבִ֕י כׇּל־תֹּפֵ֥שׂ כִּנּ֖וֹר וְעוּגָֽב׃ %
% ---------- פרק 4 | פסוק 22 ----------
 \textbf{כב} וְצִלָּ֣ה גַם־הִ֗וא יָֽלְדָה֙ אֶת־תּ֣וּבַל קַ֔יִן לֹטֵ֕שׁ כׇּל־חֹרֵ֥שׁ נְחֹ֖שֶׁת וּבַרְזֶ֑ל וַֽאֲח֥וֹת תּֽוּבַל־קַ֖יִן נַֽעֲמָֽה׃ \Rashi{\textbf{תובל קין.}תובל אמנתו של קין, תובל לשון תבלין, תבל והתקין אמנותו של קין לעשות כלי זין לרוצחים: לטש כל חרש נחשת וברזל. מחדד אמנות נחשת וברזל כמו ילטש עיניו לי (איוב ט"ז). חורש אינו לשון פעל אלא לשון פועל, שהרי נקוד קמץ קטן וטעמו למטה, כלומר, מחדד ומצחצח כל כלי אמנות נחשת וברזל: נעמה. היא אשתו של נח:}%
% ---------- פרק 4 | פסוק 23 ----------
 \textbf{כג} וַיֹּ֨אמֶר לֶ֜מֶךְ לְנָשָׁ֗יו עָדָ֤ה וְצִלָּה֙ שְׁמַ֣עַן קוֹלִ֔י נְשֵׁ֣י לֶ֔מֶךְ הַאְזֵ֖נָּה אִמְרָתִ֑י כִּ֣י אִ֤ישׁ הָרַ֙גְתִּי֙ לְפִצְעִ֔י וְיֶ֖לֶד לְחַבֻּרָתִֽי׃ \Rashi{\textbf{שמען קולי.}שהיו פורשות ממנו מתשמיש, לפי שהרג את קין ואת תובל קין בנו, שהיה למך סומא ותובל קין מושכו, וראה את קין ונדמה לו כחיה ואמר לאביו למשך בקשת והרגו, וכיון שידע שהוא קין זקנו הכה כף אל כף וספק את בנו ביניהם והרגו, והיו נשיו פורשות ממנו והוא מפיסן: שמען קולי. להשמע לי לתשמיש, וכי איש אשר הרגתי לפצעי הוא נהרג, וכי אני פצעתיו מזיד שיהא הפצע קרוי על שמי? וילד אשר הרגתי, לחברתי נהרג? כלומר, על ידי חברתי? בתמיה, והלא שוגג אני ולא מזיד, לא זהו פצעי ולא זהו חברתי: פצע. מכת חרב או חץ מקא"דורה בלעז:}%
% ---------- פרק 4 | פסוק 24 ----------
 \textbf{כד} כִּ֥י שִׁבְעָתַ֖יִם יֻקַּם־קָ֑יִן וְלֶ֖מֶךְ שִׁבְעִ֥ים וְשִׁבְעָֽה׃ \Rashi{\textbf{כי שבעתים יקם קין.}קין שהרג מזיד נתלה לו עד שבעה דורות, אני שהרגתי שוגג לא כל שכן שיתלה לי שביעיות הרבה: שבעים ושבעה. לשון רבוי שביעיות אחז לו, כך דרש רבי תנחומא. ומדרש בראשית רבה לא הרג למך כלום, ונשיו פורשות ממנו משקימו פריה ורביה, לפי שנגזרה גזרה לכלות זרעו של קין לאחר שבעה דורות, אמרו מה אנו יולדות לבהלה? למחר המבול בא ושוטף את הכל והוא אומר להן וכי איש הרגתי לפצעי? וכי אני הרגתי את הבל, שהיה איש בקומה וילד בשנים, שיהא זרעי כלה באותו עון? ומה קין שהרג נתלה לו שבעה דורות, אני שלא הרגתי, לא כ"ש שיתלו לי שביעיות הרבה? וזהו קל וחמר של שטות, אם כן אין הקב"ה גובה את חובו ומקים את דברו:}%
% ---------- פרק 4 | פסוק 25 ----------
 \textbf{כה} וַיֵּ֨דַע אָדָ֥ם עוֹד֙ אֶת־אִשְׁתּ֔וֹ וַתֵּ֣לֶד בֵּ֔ן וַתִּקְרָ֥א אֶת־שְׁמ֖וֹ שֵׁ֑ת כִּ֣י שָֽׁת־לִ֤י אֱלֹהִים֙ זֶ֣רַע אַחֵ֔ר תַּ֣חַת הֶ֔בֶל כִּ֥י הֲרָג֖וֹ קָֽיִן׃ \Rashi{\textbf{וידע אדם וגו'.}בא לו למך אצל אדם הראשון וקבל על נשיו, אמר להם: וכי עליכם לדקדק על גזרתו של מקום? אתם עשו מצותכם והוא יעשה את שלו. אמרו לו: קשט עצמך תחלה, והלא פרשת מאשתך זה מאה ושלשים שנה משנקנסה מיתה על ידך: מיד וידע אדם וגומר. ומהו עוד, ללמד שנתוספה לו תאוה על תאותו (בראשית רבה):}%
% ---------- פרק 4 | פסוק 26 ----------
 \textbf{כו} וּלְשֵׁ֤ת גַּם־הוּא֙ יֻלַּד־בֵּ֔ן וַיִּקְרָ֥א אֶת־שְׁמ֖וֹ אֱנ֑וֹשׁ אָ֣ז הוּחַ֔ל לִקְרֹ֖א בְּשֵׁ֥ם יְהֹוָֽה׃ \{ס\} \Rashi{\textbf{אז הוחל.}לשון חלין, לקרא את שמות האדם ואת שמות העצבים בשמו של הקב"ה לעשותן ע"ז ולקרותן אלהות:}%
% ---------- פרק 5 | פסוק 1 ----------
 \textbf{א} זֶ֣ה סֵ֔פֶר*(בספרי תימן סֵ֔פֶר בסמ״ך גדולה) תּוֹלְדֹ֖ת אָדָ֑ם בְּי֗וֹם בְּרֹ֤א אֱלֹהִים֙ אָדָ֔ם בִּדְמ֥וּת אֱלֹהִ֖ים עָשָׂ֥ה אֹתֽוֹ׃ \Rashi{\textbf{זה ספר תולדות אדם.}זו היא ספירת תולדות אדם, ומדרשי אגדה יש רבים: ביום ברא וגו'.  מגיד שביום שנברא הוליד:}%
% ---------- פרק 5 | פסוק 2 ----------
 \textbf{ב} זָכָ֥ר וּנְקֵבָ֖ה בְּרָאָ֑ם וַיְבָ֣רֶךְ אֹתָ֗ם וַיִּקְרָ֤א אֶת־שְׁמָם֙ אָדָ֔ם בְּי֖וֹם הִבָּֽרְאָֽם׃ %
% ---------- פרק 5 | פסוק 3 ----------
 \textbf{ג} וַיְחִ֣י אָדָ֗ם שְׁלֹשִׁ֤ים וּמְאַת֙ שָׁנָ֔ה וַיּ֥וֹלֶד בִּדְמוּת֖וֹ כְּצַלְמ֑וֹ וַיִּקְרָ֥א אֶת־שְׁמ֖וֹ שֵֽׁת׃ \Rashi{\textbf{שלשים ומאת שנה.}עד כאן פרש מן האשה (בראשית רבה):}%
% ---------- פרק 5 | פסוק 4 ----------
 \textbf{ד} וַיִּֽהְי֣וּ יְמֵי־אָדָ֗ם אַֽחֲרֵי֙ הוֹלִיד֣וֹ אֶת־שֵׁ֔ת שְׁמֹנֶ֥ה מֵאֹ֖ת שָׁנָ֑ה וַיּ֥וֹלֶד בָּנִ֖ים וּבָנֽוֹת׃ %
% ---------- פרק 5 | פסוק 5 ----------
 \textbf{ה} וַיִּֽהְי֞וּ כׇּל־יְמֵ֤י אָדָם֙ אֲשֶׁר־חַ֔י תְּשַׁ֤ע מֵאוֹת֙ שָׁנָ֔ה וּשְׁלֹשִׁ֖ים שָׁנָ֑ה וַיָּמֹֽת׃ \{ס\} %
% ---------- פרק 5 | פסוק 6 ----------
 \textbf{ו} וַֽיְחִי־שֵׁ֕ת חָמֵ֥שׁ שָׁנִ֖ים וּמְאַ֣ת שָׁנָ֑ה וַיּ֖וֹלֶד אֶת־אֱנֽוֹשׁ׃ %
% ---------- פרק 5 | פסוק 7 ----------
 \textbf{ז} וַֽיְחִי־שֵׁ֗ת אַֽחֲרֵי֙ הוֹלִיד֣וֹ אֶת־אֱנ֔וֹשׁ שֶׁ֣בַע שָׁנִ֔ים וּשְׁמֹנֶ֥ה מֵא֖וֹת שָׁנָ֑ה וַיּ֥וֹלֶד בָּנִ֖ים וּבָנֽוֹת׃ %
% ---------- פרק 5 | פסוק 8 ----------
 \textbf{ח} וַיִּֽהְיוּ֙ כׇּל־יְמֵי־שֵׁ֔ת שְׁתֵּ֤ים עֶשְׂרֵה֙ שָׁנָ֔ה וּתְשַׁ֥ע מֵא֖וֹת שָׁנָ֑ה וַיָּמֹֽת׃ \{ס\} %
% ---------- פרק 5 | פסוק 9 ----------
 \textbf{ט} וַיְחִ֥י אֱנ֖וֹשׁ תִּשְׁעִ֣ים שָׁנָ֑ה וַיּ֖וֹלֶד אֶת־קֵינָֽן׃ %
% ---------- פרק 5 | פסוק 10 ----------
 \textbf{י} וַיְחִ֣י אֱנ֗וֹשׁ אַֽחֲרֵי֙ הוֹלִיד֣וֹ אֶת־קֵינָ֔ן חֲמֵ֤שׁ עֶשְׂרֵה֙ שָׁנָ֔ה וּשְׁמֹנֶ֥ה מֵא֖וֹת שָׁנָ֑ה וַיּ֥וֹלֶד בָּנִ֖ים וּבָנֽוֹת׃ %
% ---------- פרק 5 | פסוק 11 ----------
 \textbf{יא} וַיִּֽהְיוּ֙ כׇּל־יְמֵ֣י אֱנ֔וֹשׁ חָמֵ֣שׁ שָׁנִ֔ים וּתְשַׁ֥ע מֵא֖וֹת שָׁנָ֑ה וַיָּמֹֽת׃ \{ס\} %
% ---------- פרק 5 | פסוק 12 ----------
 \textbf{יב} וַיְחִ֥י קֵינָ֖ן שִׁבְעִ֣ים שָׁנָ֑ה וַיּ֖וֹלֶד אֶת־מַֽהֲלַלְאֵֽל׃ %
% ---------- פרק 5 | פסוק 13 ----------
 \textbf{יג} וַיְחִ֣י קֵינָ֗ן אַחֲרֵי֙ הוֹלִיד֣וֹ אֶת־מַֽהֲלַלְאֵ֔ל אַרְבָּעִ֣ים שָׁנָ֔ה וּשְׁמֹנֶ֥ה מֵא֖וֹת שָׁנָ֑ה וַיּ֥וֹלֶד בָּנִ֖ים וּבָנֽוֹת׃ %
% ---------- פרק 5 | פסוק 14 ----------
 \textbf{יד} וַיִּֽהְיוּ֙ כׇּל־יְמֵ֣י קֵינָ֔ן עֶ֣שֶׂר שָׁנִ֔ים וּתְשַׁ֥ע מֵא֖וֹת שָׁנָ֑ה וַיָּמֹֽת׃ \{ס\} %
% ---------- פרק 5 | פסוק 15 ----------
 \textbf{טו} וַיְחִ֣י מַֽהֲלַלְאֵ֔ל חָמֵ֥שׁ שָׁנִ֖ים וְשִׁשִּׁ֣ים שָׁנָ֑ה וַיּ֖וֹלֶד אֶת־יָֽרֶד׃ %
% ---------- פרק 5 | פסוק 16 ----------
 \textbf{טז} וַיְחִ֣י מַֽהֲלַלְאֵ֗ל אַֽחֲרֵי֙ הוֹלִיד֣וֹ אֶת־יֶ֔רֶד שְׁלֹשִׁ֣ים שָׁנָ֔ה וּשְׁמֹנֶ֥ה מֵא֖וֹת שָׁנָ֑ה וַיּ֥וֹלֶד בָּנִ֖ים וּבָנֽוֹת׃ %
% ---------- פרק 5 | פסוק 17 ----------
 \textbf{יז} וַיִּהְיוּ֙ כׇּל־יְמֵ֣י מַהֲלַלְאֵ֔ל חָמֵ֤שׁ וְתִשְׁעִים֙ שָׁנָ֔ה וּשְׁמֹנֶ֥ה מֵא֖וֹת שָׁנָ֑ה וַיָּמֹֽת׃ \{ס\} %
% ---------- פרק 5 | פסוק 18 ----------
 \textbf{יח} וַֽיְחִי־יֶ֕רֶד שְׁתַּ֧יִם וְשִׁשִּׁ֛ים שָׁנָ֖ה וּמְאַ֣ת שָׁנָ֑ה וַיּ֖וֹלֶד אֶת־חֲנֽוֹךְ׃ %
% ---------- פרק 5 | פסוק 19 ----------
 \textbf{יט} וַֽיְחִי־יֶ֗רֶד אַֽחֲרֵי֙ הוֹלִיד֣וֹ אֶת־חֲנ֔וֹךְ שְׁמֹנֶ֥ה מֵא֖וֹת שָׁנָ֑ה וַיּ֥וֹלֶד בָּנִ֖ים וּבָנֽוֹת׃ %
% ---------- פרק 5 | פסוק 20 ----------
 \textbf{כ} וַיִּֽהְיוּ֙ כׇּל־יְמֵי־יֶ֔רֶד שְׁתַּ֤יִם וְשִׁשִּׁים֙ שָׁנָ֔ה וּתְשַׁ֥ע מֵא֖וֹת שָׁנָ֑ה וַיָּמֹֽת׃ \{ס\} %
% ---------- פרק 5 | פסוק 21 ----------
 \textbf{כא} וַיְחִ֣י חֲנ֔וֹךְ חָמֵ֥שׁ וְשִׁשִּׁ֖ים שָׁנָ֑ה וַיּ֖וֹלֶד אֶת־מְתוּשָֽׁלַח׃ %
% ---------- פרק 5 | פסוק 22 ----------
 \textbf{כב} וַיִּתְהַלֵּ֨ךְ חֲנ֜וֹךְ אֶת־הָֽאֱלֹהִ֗ים אַֽחֲרֵי֙ הוֹלִיד֣וֹ אֶת־מְתוּשֶׁ֔לַח שְׁלֹ֥שׁ מֵא֖וֹת שָׁנָ֑ה וַיּ֥וֹלֶד בָּנִ֖ים וּבָנֽוֹת׃ %
% ---------- פרק 5 | פסוק 23 ----------
 \textbf{כג} וַיְהִ֖י כׇּל־יְמֵ֣י חֲנ֑וֹךְ חָמֵ֤שׁ וְשִׁשִּׁים֙ שָׁנָ֔ה וּשְׁלֹ֥שׁ מֵא֖וֹת שָׁנָֽה׃ %
% ---------- פרק 5 | פסוק 24 ----------
 \textbf{כד} וַיִּתְהַלֵּ֥ךְ חֲנ֖וֹךְ אֶת־הָֽאֱלֹהִ֑ים וְאֵינֶ֕נּוּ כִּֽי־לָקַ֥ח אֹת֖וֹ אֱלֹהִֽים׃ \{ס\} \Rashi{\textbf{ויתהלך חנוך.}צדיק היה וקל בדעתו לשוב להרשיע, לפיכך מהר הקב"ה וסלקו והמיתו קדם זמנו, וזהו ששנה הכתוב במיתתו לכתב ואיננו בעולם – למלאות שנותיו (בראשית רבה): כי לקח אותו. לפני זמנו, כמו הנני לקח ממך את מחמד עיניך (יחזקאל כ"ד):}%
% ---------- פרק 5 | פסוק 25 ----------
 \textbf{כה} וַיְחִ֣י מְתוּשֶׁ֔לַח שֶׁ֧בַע וּשְׁמֹנִ֛ים שָׁנָ֖ה וּמְאַ֣ת שָׁנָ֑ה וַיּ֖וֹלֶד אֶת־לָֽמֶךְ׃ %
% ---------- פרק 5 | פסוק 26 ----------
 \textbf{כו} וַיְחִ֣י מְתוּשֶׁ֗לַח אַֽחֲרֵי֙ הוֹלִיד֣וֹ אֶת־לֶ֔מֶךְ שְׁתַּ֤יִם וּשְׁמוֹנִים֙ שָׁנָ֔ה וּשְׁבַ֥ע מֵא֖וֹת שָׁנָ֑ה וַיּ֥וֹלֶד בָּנִ֖ים וּבָנֽוֹת׃ %
% ---------- פרק 5 | פסוק 27 ----------
 \textbf{כז} וַיִּהְיוּ֙ כׇּל־יְמֵ֣י מְתוּשֶׁ֔לַח תֵּ֤שַׁע וְשִׁשִּׁים֙ שָׁנָ֔ה וּתְשַׁ֥ע מֵא֖וֹת שָׁנָ֑ה וַיָּמֹֽת׃ \{ס\} %
% ---------- פרק 5 | פסוק 28 ----------
 \textbf{כח} וַֽיְחִי־לֶ֕מֶךְ שְׁתַּ֧יִם וּשְׁמֹנִ֛ים שָׁנָ֖ה וּמְאַ֣ת שָׁנָ֑ה וַיּ֖וֹלֶד בֵּֽן׃ \Rashi{\textbf{ויולד בן.}שממנו נבנה העולם (בראשית רבה):}%
% ---------- פרק 5 | פסוק 29 ----------
 \textbf{כט} וַיִּקְרָ֧א אֶת־שְׁמ֛וֹ נֹ֖חַ לֵאמֹ֑ר זֶ֞֠ה יְנַחֲמֵ֤נוּ מִֽמַּעֲשֵׂ֙נוּ֙ וּמֵעִצְּב֣וֹן יָדֵ֔ינוּ מִן־הָ֣אֲדָמָ֔ה אֲשֶׁ֥ר אֵֽרְרָ֖הּ יְהֹוָֽה׃ \Rashi{\textbf{זה ינחמנו.}ינח ממנו את עצבון ידינו. עד שלא בא נח לא היה להם כלי מחרשה והוא הכין להם והיתה הארץ מוציאה קוצים ודרדרים כשזורעים חטים, מקללתו של אדם הראשון, ובימי נח נחה, וזהו ינחמנו, ואם לא תפרשהו כך אין טעם הלשון נופל על השם ואתה צריך לקרות שמו מנחם:}%
% ---------- פרק 5 | פסוק 30 ----------
 \textbf{ל} וַֽיְחִי־לֶ֗מֶךְ אַֽחֲרֵי֙ הוֹלִיד֣וֹ אֶת־נֹ֔חַ חָמֵ֤שׁ וְתִשְׁעִים֙ שָׁנָ֔ה וַחֲמֵ֥שׁ מֵאֹ֖ת שָׁנָ֑ה וַיּ֥וֹלֶד בָּנִ֖ים וּבָנֽוֹת׃ %
% ---------- פרק 5 | פסוק 31 ----------
 \textbf{לא} וַֽיְהִי֙ כׇּל־יְמֵי־לֶ֔מֶךְ שֶׁ֤בַע וְשִׁבְעִים֙ שָׁנָ֔ה וּשְׁבַ֥ע מֵא֖וֹת שָׁנָ֑ה וַיָּמֹֽת׃ \{ס\} %
% ---------- פרק 5 | פסוק 32 ----------
 \textbf{לב} וַֽיְהִי־נֹ֕חַ בֶּן־חֲמֵ֥שׁ מֵא֖וֹת שָׁנָ֑ה וַיּ֣וֹלֶד נֹ֔חַ אֶת־שֵׁ֖ם אֶת־חָ֥ם וְאֶת־יָֽפֶת׃ \Rashi{\textbf{בן חמש מאות שנה.}אמר רבי יודן מה טעם כל הדורות הולידו לק' שנה וזה לחמש מאות? אמר הקב"ה אם רשעים הם יאבדו במים ורע לצדיק זה, ואם צדיקים הם אטריח עליו לעשות תבות הרבה, כבש את מעינו ולא הוליד עד חמש מאות שנה, כדי שלא יהא יפת הגדול שבבניו ראוי לענשין לפני המבול (בראשית רבה), דכתיב כי הנער בן מאה שנה ימות (ישעיהו ס"ה) ראוי לענש לעתיד, וכן לפני מתן תורה: את שם את חם ואת יפת. והלא יפת הגדול הוא? אלא בתחלה אתה דורש את שהוא צדיק ונולד כשהוא מהול ושאברהם יצא ממנו וכו' בבראשית רבה:}%
% ---------- פרק 6 | פסוק 1 ----------
 \textbf{א} וַֽיְהִי֙ כִּֽי־הֵחֵ֣ל הָֽאָדָ֔ם לָרֹ֖ב עַל־פְּנֵ֣י הָֽאֲדָמָ֑ה וּבָנ֖וֹת יֻלְּד֥וּ לָהֶֽם׃ %
% ---------- פרק 6 | פסוק 2 ----------
 \textbf{ב} וַיִּרְא֤וּ בְנֵי־הָֽאֱלֹהִים֙ אֶת־בְּנ֣וֹת הָֽאָדָ֔ם כִּ֥י טֹבֹ֖ת הֵ֑נָּה וַיִּקְח֤וּ לָהֶם֙ נָשִׁ֔ים מִכֹּ֖ל אֲשֶׁ֥ר בָּחָֽרוּ׃ \Rashi{\textbf{בני האלהים.}בני השרים והשופטים. דבר אחר בני האלהים, הם השרים ההולכים בשליחותו של מקום, אף הם היו מתערבים בהם; כל אלהים שבמקרא לשון מרות, וזה יוכיח ואתה תהיה לו לאלהים (שמות ד'), ראה נתתיך אלהים (שם ז'): כי טבת הנה. אמר רבי יודן טבת כתיב, כשהיו מטיבין אותה מקשטת לכנס לחפה, היה גדול נכנס ובועלה תחלה (בראשית רבה): מכל אשר בחרו. אף בעולת בעל, אף הזכר והבהמה:}%
% ---------- פרק 6 | פסוק 3 ----------
 \textbf{ג} וַיֹּ֣אמֶר יְהֹוָ֗ה לֹֽא־יָד֨וֹן רוּחִ֤י בָֽאָדָם֙ לְעֹלָ֔ם בְּשַׁגַּ֖ם ה֣וּא בָשָׂ֑ר וְהָי֣וּ יָמָ֔יו מֵאָ֥ה וְעֶשְׂרִ֖ים שָׁנָֽה׃ \Rashi{\textbf{לא ידון רוחי באדם.}לא יתרעם ויריב רוחי עלי בשביל האדם לעלם. לארך ימים; הנה רוחי נדון בקרבי אם להשחית ואם לרחם, לא יהיה מדון זה ברוחי לעולם, כלומר לארך ימים: בשגם הוא בשר. כמו בשגם, כלומר, בשביל שגם זאת בו, שהוא בשר ואף על פי כן אינו נכנע לפני, ומה אם יהיה אש או דבר קשה? כיוצא בו עד שקמתי דבורה (שופטים ה') כמו שקמתי וכן שאתה מדבר עמי (שם ו') כמו שאתה, אף בשגם כמו בשגם; והיו ימיו וגו'. עד ק"ך שנה אאריך להם אפי, ואם לא ישובו אביא עליהם מבול. ואם תאמר משנולד יפת עד המבול אינו אלא מאה שנה, אין מקדם ומאחר בתורה (פסחים ו') – כבר היתה הגזרה גזורה עשרים שנה קדם שהוליד נח תולדות, וכן מצינו בסדר עולם. יש מדרשי אגדה רבים בלא ידון, אבל זה הוא צחצוח פשוטו:}%
% ---------- פרק 6 | פסוק 4 ----------
 \textbf{ד} הַנְּפִלִ֞ים הָי֣וּ בָאָ֘רֶץ֮ בַּיָּמִ֣ים הָהֵם֒ וְגַ֣ם אַֽחֲרֵי־כֵ֗ן אֲשֶׁ֨ר יָבֹ֜אוּ בְּנֵ֤י הָֽאֱלֹהִים֙ אֶל־בְּנ֣וֹת הָֽאָדָ֔ם וְיָלְד֖וּ לָהֶ֑ם הֵ֧מָּה הַגִּבֹּרִ֛ים אֲשֶׁ֥ר מֵעוֹלָ֖ם אַנְשֵׁ֥י הַשֵּֽׁם׃ \{פ\} \Rashi{\textbf{הנפלים.}על שם שנפלו והפילו את העולם, ובלשון עברית לשון ענקים הוא: בימים ההם. בימי דור אנוש ובני קין (בראשית רבה): וגם אחרי כן. אף על פי שראו באבדן של דור אנוש, שעלה אוקינוס והציף שליש העולם לא נכנע דור המבול ללמד מהם: אשר יבאו. היו יולדות ענקים כמותם: הגבורים. למרד במקום: אנשי השם. אותן שנקבו בשמות עירד מחויאל מתושאל שנקבו על שם אבדן, שנמחו והתשו. דבר אחר אנשי שממון, ששממו את העולם:}%
% ---------- פרק 6 | פסוק 5 ----------
 \textbf{ה} וַיַּ֣רְא יְהֹוָ֔ה כִּ֥י רַבָּ֛ה רָעַ֥ת הָאָדָ֖ם בָּאָ֑רֶץ וְכׇל־יֵ֙צֶר֙ מַחְשְׁבֹ֣ת לִבּ֔וֹ רַ֥ק רַ֖ע כׇּל־הַיּֽוֹם׃ %
% ---------- פרק 6 | פסוק 6 ----------
 \textbf{ו} וַיִּנָּ֣חֶם יְהֹוָ֔ה כִּֽי־עָשָׂ֥ה אֶת־הָֽאָדָ֖ם בָּאָ֑רֶץ וַיִּתְעַצֵּ֖ב אֶל־לִבּֽוֹ׃ \Rashi{\textbf{וינחם ה' כי עשה.}נחמה היתה לפניו שבראו בתחתונים, שאלו היה מן העליונים היה ממרידן (בבראשית רבה): ויתעצב. האדם אל לבו של מקום, עלה במחשבתו של מקום להעציבו, זהו תרגום אנקלוס. דבר אחר וינחם – נהפכה מחשבתו של מקום ממדת רחמים למדת הדין, עלה במחשבה לפניו מה לעשות באדם שעשה בארץ, וכן כל לשון נחום שבמקרא לשון נמלך מה לעשות, ובן אדם ויתנחם (במדבר כ"ג) ועל עבדיו יתנחם (דברים ל"ב) וינחם ה' על הרעה (שמות ל"ב) נחמתי כי המלכתי (שמואל א' ט"ו) כלם לשון מחשבה אחרת הם: ויתעצב אל לבו. נתאבל על אבדן מעשה ידיו, כמו נעצב המלך על בנו (שם ב' י"ט), וזו כתבתי לתשובת המינים גוי אחד ששאל את רבי יהושע בן קרחה, אמר לו אין אתם מודים שהקב"ה רואה את הנולד? אמר לו הן, אמר לו והא כתיב ויתעצב אל לבו? אמר לו נולד לך בן זכר מימיך? אמר לו הן, אמר לו ומה עשית? אמר לו שמחתי ושמחתי את הכל, אמר לו ולא היית יודע שסופו למות? אמר לו בשעת חדותא חדותא בשעת אבלא אבלא, אמר לו כך מעשה הקב"ה, אף על פי שגלוי לפניו שסופן לחטוא ולאבדן לא נמנע מלבראן בשביל הצדיקים העתידים לעמד מהם:}%
% ---------- פרק 6 | פסוק 7 ----------
 \textbf{ז} וַיֹּ֣אמֶר יְהֹוָ֗ה אֶמְחֶ֨ה אֶת־הָאָדָ֤ם אֲשֶׁר־בָּרָ֙אתִי֙ מֵעַל֙ פְּנֵ֣י הָֽאֲדָמָ֔ה מֵֽאָדָם֙ עַד־בְּהֵמָ֔ה עַד־רֶ֖מֶשׂ וְעַד־ע֣וֹף הַשָּׁמָ֑יִם כִּ֥י נִחַ֖מְתִּי כִּ֥י עֲשִׂיתִֽם׃ \Rashi{\textbf{ויאמר ה' אמחה את האדם.}הוא עפר ואביא עליו מים ואמחה אותו, לכך נאמר לשון מחוי: מאדם עד בהמה. אף הם השחיתו דרכם (בראשית רבה). דבר אחר הכל נברא בשביל אדם וכיון שהוא כלה מה צרך באלו: כי נחמתי כי עשיתם. חשבתי מה לעשות על אשר עשיתים:}%
% ---------- פרק 6 | פסוק 8 ----------
 \textbf{ח} וְנֹ֕חַ מָ֥צָא חֵ֖ן בְּעֵינֵ֥י יְהֹוָֽה׃ \{פ\} %
\addparasha{חומש בראשית | פרשת נח}
% ---------- פרק 6 | פסוק 9 ----------
 \textbf{ט} אֵ֚לֶּה תּוֹלְדֹ֣ת נֹ֔חַ נֹ֗חַ אִ֥ישׁ צַדִּ֛יק תָּמִ֥ים הָיָ֖ה בְּדֹֽרֹתָ֑יו אֶת־הָֽאֱלֹהִ֖ים הִֽתְהַלֶּךְ־נֹֽחַ׃ \Rashi{\textbf{אלה תולדת נח נח איש צדיק.}הואיל והזכירו ספר בשבחו, שנאמר זכר צדיק לברכה (משלי י'). דבר אחר למדך שעקר תולדותיהם של צדיקים מעשים טובים: בדרותיו. יש מרבותינו דורשים אותו לשבח, כל שכן אלו היה בדור צדיקים היה צדיק יותר; ויש שדורשים אותו לגנאי, לפי דורו היה צדיק ואלו היה בדורו של אברהם לא היה נחשב לכלום (סנה' ק"ח): את האלהים התהלך נח. ובאברהם הוא אומר אשר התהלכתי לפניו? (ברא' כ"ד), נח היה צריך סעד לתמכו, אבל אברהם היה מתחזק ומהלך בצדקו מאליו: התהלך. לשון עבר, וזהו שמושו של לשון, בלשון כבד משמשת להבא ולשעבר בלשון א', קום התהלך (שם י"ג) להבא, התהלך נח לשעבר; התפלל בעד עבדיך (שמואל א' י"ב) להבא, ובא והתפלל אל הבית הזה (מ"א ח') לשון עבר, אלא שהוי"ו שבראשו הפכו להבא:}%
% ---------- פרק 6 | פסוק 10 ----------
 \textbf{י} וַיּ֥וֹלֶד נֹ֖חַ שְׁלֹשָׁ֣ה בָנִ֑ים אֶת־שֵׁ֖ם אֶת־חָ֥ם וְאֶת־יָֽפֶת׃ %
% ---------- פרק 6 | פסוק 11 ----------
 \textbf{יא} וַתִּשָּׁחֵ֥ת הָאָ֖רֶץ לִפְנֵ֣י הָֽאֱלֹהִ֑ים וַתִּמָּלֵ֥א הָאָ֖רֶץ חָמָֽס׃ \Rashi{\textbf{ותשחת.}לשון ערוה וע"ז (סנה' נ"ז) כמו פן תשחתון (דב' ד'), כי השחית כל בשר וגו': ותמלא הארץ חמס. גזל:}%
% ---------- פרק 6 | פסוק 12 ----------
 \textbf{יב} וַיַּ֧רְא אֱלֹהִ֛ים אֶת־הָאָ֖רֶץ וְהִנֵּ֣ה נִשְׁחָ֑תָה כִּֽי־הִשְׁחִ֧ית כׇּל־בָּשָׂ֛ר אֶת־דַּרְכּ֖וֹ עַל־הָאָֽרֶץ׃ \{ס\} \Rashi{\textbf{כי השחית כל בשר.}אפלו בהמה חיה ועוף נזקקין לשאינן מינן:}%
% ---------- פרק 6 | פסוק 13 ----------
 \textbf{יג} וַיֹּ֨אמֶר אֱלֹהִ֜ים לְנֹ֗חַ קֵ֤ץ כׇּל־בָּשָׂר֙ בָּ֣א לְפָנַ֔י כִּֽי־מָלְאָ֥ה הָאָ֛רֶץ חָמָ֖ס מִפְּנֵיהֶ֑ם וְהִנְנִ֥י מַשְׁחִיתָ֖ם אֶת־הָאָֽרֶץ׃ \Rashi{\textbf{קץ כל בשר.}כל מקום שאתה מוצא זנות וע"ז אנדרולומוסיא באה לעולם והורגת טובים ורעים. כי מלאה הארץ חמס. לא נחתם גזר דינם אלא על הגזל (סנה' ק"ח): את הארץ. כמו מן הארץ, ודומה לו כצאתי את העיר (שמ' ט'), מן העיר; חלה את רגליו (מ"א ט"ו), מן רגליו. ד"א את הארץ – עם הארץ, שאף ג' טפחים של עמק המחרשה נמחו ונטשטשו:}%
% ---------- פרק 6 | פסוק 14 ----------
 \textbf{יד} עֲשֵׂ֤ה לְךָ֙ תֵּבַ֣ת עֲצֵי־גֹ֔פֶר קִנִּ֖ים תַּֽעֲשֶׂ֣ה אֶת־הַתֵּבָ֑ה וְכָֽפַרְתָּ֥ אֹתָ֛הּ מִבַּ֥יִת וּמִח֖וּץ בַּכֹּֽפֶר׃ \Rashi{\textbf{עשה לך תבת.}הרבה רוח והצלה לפניו, ולמה הטריחו בבנין זה? כדי שיראוהו אנשי דור המבול עוסק בה ק"כ שנה, ושואלין אותו מה זאת לך, והוא אומר להם עתיד הקב"ה להביא מבול לעולם, אולי ישובו: עצי גופר. כך שמו. ולמה ממין זה? על שם גפרית שנגזר עליהם למחות בו: קנים. מדורים מדורים לכל בהמה וחיה: בכפר. זפת בלשון ארמי, ומצינו בתלמוד כופרא בתבתו של משה, על ידי שהיו המים תשים דיה בחמר מבפנים וזפת מבחוץ; ועוד, כדי שלא יריח אותו צדיק ריח רע של זפת, אבל כאן מפני חזק המים זפתה מבית ומחוץ:}%
% ---------- פרק 6 | פסוק 15 ----------
 \textbf{טו} וְזֶ֕ה אֲשֶׁ֥ר תַּֽעֲשֶׂ֖ה אֹתָ֑הּ שְׁלֹ֧שׁ מֵא֣וֹת אַמָּ֗ה אֹ֚רֶךְ הַתֵּבָ֔ה חֲמִשִּׁ֤ים אַמָּה֙ רׇחְבָּ֔הּ וּשְׁלֹשִׁ֥ים אַמָּ֖ה קוֹמָתָֽהּ׃ %
% ---------- פרק 6 | פסוק 16 ----------
 \textbf{טז} צֹ֣הַר ׀ תַּעֲשֶׂ֣ה לַתֵּבָ֗ה וְאֶל־אַמָּה֙ תְּכַלֶּ֣נָּה מִלְמַ֔עְלָה וּפֶ֥תַח הַתֵּבָ֖ה בְּצִדָּ֣הּ תָּשִׂ֑ים תַּחְתִּיִּ֛ם שְׁנִיִּ֥ם וּשְׁלִשִׁ֖ים תַּֽעֲשֶֽׂהָ׃ \Rashi{\textbf{צהר.}י"א חלון וי"א אבן טובה המאירה להם (ב"ר שם): ואל אמה תכלנה מלמעלה. כסויה משפע ועולה עד שהוא קצר מלמעלה ועומד על אמה כדי שיזובו המים למטה מכאן ומכאן: בצדה תשים. שלא יפלו הגשמים בה: תחתים שנים ושלשים. ג' עליות זו על גב זו, עליונים לאדם, אמצעים למדור, תחתיים לזבל (שם):}%
% ---------- פרק 6 | פסוק 17 ----------
 \textbf{יז} וַאֲנִ֗י הִנְנִי֩ מֵבִ֨יא אֶת־הַמַּבּ֥וּל מַ֙יִם֙ עַל־הָאָ֔רֶץ לְשַׁחֵ֣ת כׇּל־בָּשָׂ֗ר אֲשֶׁר־בּוֹ֙ ר֣וּחַ חַיִּ֔ים מִתַּ֖חַת הַשָּׁמָ֑יִם כֹּ֥ל אֲשֶׁר־בָּאָ֖רֶץ יִגְוָֽע׃ \Rashi{\textbf{ואני הנני מביא.}הנני מוכן להסכים עם אותם שזרזוני ואמרו לפני כבר מה אנוש כי תזכרנו (תה' ח'): מבול. שבלה את הכל, שבלבל את הכל, שהוביל את הכל מן הגבוה לנמוך; וזהו לשון אנקלוס שתרגם טופנא, שהציף את הכל, והביאה לבבל שהיא עמקה, לכך נקראת שנער, שננערו שם כל מתי מבול:}%
% ---------- פרק 6 | פסוק 18 ----------
 \textbf{יח} וַהֲקִמֹתִ֥י אֶת־בְּרִיתִ֖י אִתָּ֑ךְ וּבָאתָ֙ אֶל־הַתֵּבָ֔ה אַתָּ֕ה וּבָנֶ֛יךָ וְאִשְׁתְּךָ֥ וּנְשֵֽׁי־בָנֶ֖יךָ אִתָּֽךְ׃ \Rashi{\textbf{והקמתי את בריתי.}ברית היה צריך על הפרות שלא ירקבו ויעפשו ושלא יהרגוהו רשעים שבדור (ב"ר): אתה ובניך ואשתך. האנשים לבד והנשים לבד; מכאן שנאסרו בתשמיש המטה: (שם):}%
% ---------- פרק 6 | פסוק 19 ----------
 \textbf{יט} וּמִכׇּל־הָ֠חַ֠י מִֽכׇּל־בָּשָׂ֞ר שְׁנַ֧יִם מִכֹּ֛ל תָּבִ֥יא אֶל־הַתֵּבָ֖ה לְהַחֲיֹ֣ת אִתָּ֑ךְ זָכָ֥ר וּנְקֵבָ֖ה יִֽהְיֽוּ׃ \Rashi{\textbf{ומכל החי.}אפלו שדים (ב"ר): שנים מכל. מן הפחות שבהם לא פחות משנים, אחד זכר ואחד נקבה:}%
% ---------- פרק 6 | פסוק 20 ----------
 \textbf{כ} מֵהָע֣וֹף לְמִינֵ֗הוּ וּמִן־הַבְּהֵמָה֙ לְמִינָ֔הּ מִכֹּ֛ל רֶ֥מֶשׂ הָֽאֲדָמָ֖ה לְמִינֵ֑הוּ שְׁנַ֧יִם מִכֹּ֛ל יָבֹ֥אוּ אֵלֶ֖יךָ לְהַֽחֲיֽוֹת׃ \Rashi{\textbf{מהעוף למינהו.}אותן שדבקו במיניהם ולא השחיתו דרכם, ומאליהם באו, וכל שהתבה קולטתו הכניס (ב"ר):}%
% ---------- פרק 6 | פסוק 21 ----------
 \textbf{כא} וְאַתָּ֣ה קַח־לְךָ֗ מִכׇּל־מַֽאֲכָל֙ אֲשֶׁ֣ר יֵֽאָכֵ֔ל וְאָסַפְתָּ֖ אֵלֶ֑יךָ וְהָיָ֥ה לְךָ֛ וְלָהֶ֖ם לְאׇכְלָֽה׃ %
% ---------- פרק 6 | פסוק 22 ----------
 \textbf{כב} וַיַּ֖עַשׂ נֹ֑חַ כְּ֠כֹ֠ל אֲשֶׁ֨ר צִוָּ֥ה אֹת֛וֹ אֱלֹהִ֖ים כֵּ֥ן עָשָֽׂה׃ \Rashi{\textbf{ויעש נח.}זה בנין התבה:}%
% ---------- פרק 7 | פסוק 1 ----------
 \textbf{א} וַיֹּ֤אמֶר יְהֹוָה֙ לְנֹ֔חַ בֹּֽא־אַתָּ֥ה וְכׇל־בֵּיתְךָ֖ אֶל־הַתֵּבָ֑ה כִּֽי־אֹתְךָ֥ רָאִ֛יתִי צַדִּ֥יק לְפָנַ֖י בַּדּ֥וֹר הַזֶּֽה׃ \Rashi{\textbf{ראיתי צדיק.}ולא נאמר צדיק תמים, מכאן שאומרים מקצת שבחו של אדם בפניו וכלו שלא בפניו (ב"ר):}%
% ---------- פרק 7 | פסוק 2 ----------
 \textbf{ב} מִכֹּ֣ל ׀ הַבְּהֵמָ֣ה הַטְּהוֹרָ֗ה תִּֽקַּח־לְךָ֛ שִׁבְעָ֥ה שִׁבְעָ֖ה אִ֣ישׁ וְאִשְׁתּ֑וֹ וּמִן־הַבְּהֵמָ֡ה אֲ֠שֶׁ֠ר לֹ֣א טְהֹרָ֥ה הִ֛וא שְׁנַ֖יִם אִ֥ישׁ וְאִשְׁתּֽוֹ׃ \Rashi{\textbf{הטהורה.}העתידה להיות טהורה לישראל; למדנו שלמד נח תורה (שם פ' כ"ו): שבעה שבעה. כדי שיקריב מהם קרבן בצאתו:}%
% ---------- פרק 7 | פסוק 3 ----------
 \textbf{ג} גַּ֣ם מֵע֧וֹף הַשָּׁמַ֛יִם שִׁבְעָ֥ה שִׁבְעָ֖ה זָכָ֣ר וּנְקֵבָ֑ה לְחַיּ֥וֹת זֶ֖רַע עַל־פְּנֵ֥י כׇל־הָאָֽרֶץ׃ \Rashi{\textbf{גם מעוף השמים וגו'.}בטהורים הכתוב מדבר, ולמד סתום מן המפרש:}%
% ---------- פרק 7 | פסוק 4 ----------
 \textbf{ד} כִּי֩ לְיָמִ֨ים ע֜וֹד שִׁבְעָ֗ה אָֽנֹכִי֙ מַמְטִ֣יר עַל־הָאָ֔רֶץ אַרְבָּעִ֣ים י֔וֹם וְאַרְבָּעִ֖ים לָ֑יְלָה וּמָחִ֗יתִי אֶֽת־כׇּל־הַיְקוּם֙ אֲשֶׁ֣ר עָשִׂ֔יתִי מֵעַ֖ל פְּנֵ֥י הָֽאֲדָמָֽה׃ \Rashi{\textbf{כי לימים עוד שבעה.}אלו ז' ימי אבלו של מתושלח הצדיק שחס הקב"ה על כבודו ועכב את הפרענות. צא וחשב שנותיו של מתושלח ותמצא שהם כלים בשנת ת"ר שנה לחיי נח (סנה' ק"ח): כי לימים עוד. מהו עוד? זמן אחר זמן זה נוסף על ק"כ שנה: ארבעים יום. כנגד יצירת הולד, שקלקלו להטריח ליוצרם לצור צורת ממזרים:}%
% ---------- פרק 7 | פסוק 5 ----------
 \textbf{ה} וַיַּ֖עַשׂ נֹ֑חַ כְּכֹ֥ל אֲשֶׁר־צִוָּ֖הוּ יְהֹוָֽה׃ \Rashi{\textbf{ויעש נח.}זה ביאתו לתבה:}%
% ---------- פרק 7 | פסוק 6 ----------
 \textbf{ו} וְנֹ֕חַ בֶּן־שֵׁ֥שׁ מֵא֖וֹת שָׁנָ֑ה וְהַמַּבּ֣וּל הָיָ֔ה מַ֖יִם עַל־הָאָֽרֶץ׃ %
% ---------- פרק 7 | פסוק 7 ----------
 \textbf{ז} וַיָּ֣בֹא נֹ֗חַ וּ֠בָנָ֠יו וְאִשְׁתּ֧וֹ וּנְשֵֽׁי־בָנָ֛יו אִתּ֖וֹ אֶל־הַתֵּבָ֑ה מִפְּנֵ֖י מֵ֥י הַמַּבּֽוּל׃ \Rashi{\textbf{נח ובניו.}האנשים לבד והנשים לבד, לפי שנאסרו בתשמיש המטה מפני שהעולם שרוי בצער: מפני מי המבול. אף נח מקטני אמנה היה, מאמין ואינו מאמין שיבא המבול, ולא נכנס לתבה עד שדחקוהו המים:}%
% ---------- פרק 7 | פסוק 8 ----------
 \textbf{ח} מִן־הַבְּהֵמָה֙ הַטְּהוֹרָ֔ה וּמִ֨ן־הַבְּהֵמָ֔ה אֲשֶׁ֥ר אֵינֶ֖נָּה טְהֹרָ֑ה וּמִ֨ן־הָע֔וֹף וְכֹ֥ל אֲשֶׁר־רֹמֵ֖שׂ עַל־הָֽאֲדָמָֽה׃ %
% ---------- פרק 7 | פסוק 9 ----------
 \textbf{ט} שְׁנַ֨יִם שְׁנַ֜יִם בָּ֧אוּ אֶל־נֹ֛חַ אֶל־הַתֵּבָ֖ה זָכָ֣ר וּנְקֵבָ֑ה כַּֽאֲשֶׁ֛ר צִוָּ֥ה אֱלֹהִ֖ים אֶת־נֹֽחַ׃ \Rashi{\textbf{שנים שנים.}כלם השוו במנין זה – מן הפחות היו שנים: באו אל נח. מאליהן:}%
% ---------- פרק 7 | פסוק 10 ----------
 \textbf{י} וַיְהִ֖י לְשִׁבְעַ֣ת הַיָּמִ֑ים וּמֵ֣י הַמַּבּ֔וּל הָי֖וּ עַל־הָאָֽרֶץ׃ %
% ---------- פרק 7 | פסוק 11 ----------
 \textbf{יא} בִּשְׁנַ֨ת שֵׁשׁ־מֵא֤וֹת שָׁנָה֙ לְחַיֵּי־נֹ֔חַ בַּחֹ֙דֶשׁ֙ הַשֵּׁנִ֔י בְּשִׁבְעָֽה־עָשָׂ֥ר י֖וֹם לַחֹ֑דֶשׁ בַּיּ֣וֹם הַזֶּ֗ה נִבְקְעוּ֙ כׇּֽל־מַעְיְנֹת֙*(בספרי ספרד ואשכנז מַעְיְנוֹת֙) תְּה֣וֹם רַבָּ֔ה וַאֲרֻבֹּ֥ת הַשָּׁמַ֖יִם נִפְתָּֽחוּ׃ \Rashi{\textbf{בחודש השני.}רבי אליעזר אומר זה מרחשון; רבי יהושע אומר זה איר (ר"ה י"א): נבקעו. להוציא מימיהן: תהום רבה. מדה כנגד מדה, הם חטאו ברבה רעת האדם ולקו בתהום רבה:}%
% ---------- פרק 7 | פסוק 12 ----------
 \textbf{יב} וַיְהִ֥י הַגֶּ֖שֶׁם עַל־הָאָ֑רֶץ אַרְבָּעִ֣ים י֔וֹם וְאַרְבָּעִ֖ים לָֽיְלָה׃ \Rashi{\textbf{ויהי הגשם על הארץ.}ולהלן הוא אומר ויהי המבול? אלא כשהורידן הורידן ברחמים, שאם יחזרו יהיו גשמי ברכה; כשלא חזרו היו למבול: ארבעים יום וגו'. אין יום ראשון מן המנין, לפי שאין לילו עמו שהרי כתיב ביום הזה נבקעו כל מעינת, נמצאו ארבעים יום כלים בכ"ח בכסלו לרבי אליעזר, שהחדשים נמנין כסדרן אחד מלא ואחד חסר, הרי י"ב ממרחשון וכ"ח מכסלו:}%
% ---------- פרק 7 | פסוק 13 ----------
 \textbf{יג} בְּעֶ֨צֶם הַיּ֤וֹם הַזֶּה֙ בָּ֣א נֹ֔חַ וְשֵׁם־וְחָ֥ם וָיֶ֖פֶת בְּנֵי־נֹ֑חַ וְאֵ֣שֶׁת נֹ֗חַ וּשְׁלֹ֧שֶׁת נְשֵֽׁי־בָנָ֛יו אִתָּ֖ם אֶל־הַתֵּבָֽה׃ \Rashi{\textbf{בעצם היום הזה.}למדך הכתוב שהיו בני דורו אומרים אלו אנו רואים אותו נכנס לתבה, אנו שוברין אותה והורגין אותו, אמר הקב"ה אני מכניסו לעיני כלם ונראה דבר מי יקום:}%
% ---------- פרק 7 | פסוק 14 ----------
 \textbf{יד} הֵ֜מָּה וְכׇל־הַֽחַיָּ֣ה לְמִינָ֗הּ וְכׇל־הַבְּהֵמָה֙ לְמִינָ֔הּ וְכׇל־הָרֶ֛מֶשׂ הָרֹמֵ֥שׂ עַל־הָאָ֖רֶץ לְמִינֵ֑הוּ וְכׇל־הָע֣וֹף לְמִינֵ֔הוּ כֹּ֖ל צִפּ֥וֹר כׇּל־כָּנָֽף׃ \Rashi{\textbf{צפור כל כנף.}דבוק הוא, צפור של כל מין כנף, לרבות חגבים (כנף לשון נוצה, כמו ושסע אתו בכנפיו, שאפלו נוצתה עולה, אף כאן צפור כל מין מראית נוצה):}%
% ---------- פרק 7 | פסוק 15 ----------
 \textbf{טו} וַיָּבֹ֥אוּ אֶל־נֹ֖חַ אֶל־הַתֵּבָ֑ה שְׁנַ֤יִם שְׁנַ֙יִם֙ מִכׇּל־הַבָּשָׂ֔ר אֲשֶׁר־בּ֖וֹ ר֥וּחַ חַיִּֽים׃ %
% ---------- פרק 7 | פסוק 16 ----------
 \textbf{טז} וְהַבָּאִ֗ים זָכָ֨ר וּנְקֵבָ֤ה מִכׇּל־בָּשָׂר֙ בָּ֔אוּ כַּֽאֲשֶׁ֛ר צִוָּ֥ה אֹת֖וֹ אֱלֹהִ֑ים וַיִּסְגֹּ֥ר יְהֹוָ֖ה בַּֽעֲדֽוֹ׃ \Rashi{\textbf{ויסגור ה' בעדו.}הגן עליו שלא ישברוה, הקיף התבה דבים ואריות והיו הורגים בהם; ופשוטו של מקרא סגר כנגדו מן המים, וכן כל בעד שבמקרא לשון כנגד הוא, בעד כל רחם (בראשית כ׳:י״ח), בעדך ובעד בניך (מ"ב ד'), עור בעד עור (איוב ב'), מגן בעדי (תהלים ג'), התפלל בעד עבדיך (שמואל א' י"ב) – כנגד עבדיך:}%
% ---------- פרק 7 | פסוק 17 ----------
 \textbf{יז} וַֽיְהִ֧י הַמַּבּ֛וּל אַרְבָּעִ֥ים י֖וֹם עַל־הָאָ֑רֶץ וַיִּרְבּ֣וּ הַמַּ֗יִם וַיִּשְׂאוּ֙ אֶת־הַתֵּבָ֔ה וַתָּ֖רׇם מֵעַ֥ל הָאָֽרֶץ׃ \Rashi{\textbf{ותרם מעל הארץ.}משקעת היתה במים י"א אמה כספינה טעונה המשקעת מקצתה במים; ומקראות שלפנינו יוכיחו:}%
% ---------- פרק 7 | פסוק 18 ----------
 \textbf{יח} וַיִּגְבְּר֥וּ הַמַּ֛יִם וַיִּרְבּ֥וּ מְאֹ֖ד עַל־הָאָ֑רֶץ וַתֵּ֥לֶךְ הַתֵּבָ֖ה עַל־פְּנֵ֥י הַמָּֽיִם׃ \Rashi{\textbf{ויגברו.}מאליהן:}%
% ---------- פרק 7 | פסוק 19 ----------
 \textbf{יט} וְהַמַּ֗יִם גָּ֥בְר֛וּ מְאֹ֥ד מְאֹ֖ד עַל־הָאָ֑רֶץ וַיְכֻסּ֗וּ כׇּל־הֶֽהָרִים֙ הַגְּבֹהִ֔ים אֲשֶׁר־תַּ֖חַת כׇּל־הַשָּׁמָֽיִם׃ %
% ---------- פרק 7 | פסוק 20 ----------
 \textbf{כ} חֲמֵ֨שׁ עֶשְׂרֵ֤ה אַמָּה֙ מִלְמַ֔עְלָה גָּבְר֖וּ הַמָּ֑יִם וַיְכֻסּ֖וּ הֶהָרִֽים׃ \Rashi{\textbf{חמש עשרה אמה מלמעלה.}למעלה של גבה כל ההרים, לאחר שהשוו המים לראשי ההרים:}%
% ---------- פרק 7 | פסוק 21 ----------
 \textbf{כא} וַיִּגְוַ֞ע כׇּל־בָּשָׂ֣ר ׀ הָרֹמֵ֣שׂ עַל־הָאָ֗רֶץ בָּע֤וֹף וּבַבְּהֵמָה֙ וּבַ֣חַיָּ֔ה וּבְכׇל־הַשֶּׁ֖רֶץ הַשֹּׁרֵ֣ץ עַל־הָאָ֑רֶץ וְכֹ֖ל הָאָדָֽם׃ %
% ---------- פרק 7 | פסוק 22 ----------
 \textbf{כב} כֹּ֡ל אֲשֶׁר֩ נִשְׁמַת־ר֨וּחַ חַיִּ֜ים בְּאַפָּ֗יו מִכֹּ֛ל אֲשֶׁ֥ר בֶּחָֽרָבָ֖ה מֵֽתוּ׃ \Rashi{\textbf{נשמת רוח חיים.}נשמה של רוח חיים: אשר בחרבה. ולא דגים שבים (סנה' ק"ח):}%
% ---------- פרק 7 | פסוק 23 ----------
 \textbf{כג} וַיִּ֜מַח אֶֽת־כׇּל־הַיְק֣וּם ׀ אֲשֶׁ֣ר ׀ עַל־פְּנֵ֣י הָֽאֲדָמָ֗ה מֵאָדָ֤ם עַד־בְּהֵמָה֙ עַד־רֶ֙מֶשׂ֙ וְעַד־ע֣וֹף הַשָּׁמַ֔יִם וַיִּמָּח֖וּ מִן־הָאָ֑רֶץ וַיִּשָּׁ֧אֶר אַךְ־נֹ֛חַ וַֽאֲשֶׁ֥ר אִתּ֖וֹ בַּתֵּבָֽה׃ \Rashi{\textbf{וימח.}לשון ויפעל הוא ואינו לשון ויפעל, והוא מגזרת ויפן ויבן כל תבה שסופה ה"א, כגון בנה, מחה, קנה, כשהוא נותן וי"ו יו"ד בראשה, נקוד בחירק תחת היו"ד: אך נח. לבד נח, זהו פשוטו; ומ"א גונח וכוהה דם מטרח הבהמות והחיות. וי"א שאחר מזונות לארי והכישו, ועליו נאמר הן צדיק בארץ ישלם (משלי יא):}%
% ---------- פרק 7 | פסוק 24 ----------
 \textbf{כד} וַיִּגְבְּר֥וּ הַמַּ֖יִם עַל־הָאָ֑רֶץ חֲמִשִּׁ֥ים וּמְאַ֖ת יֽוֹם׃ %
% ---------- פרק 8 | פסוק 1 ----------
 \textbf{א} וַיִּזְכֹּ֤ר אֱלֹהִים֙ אֶת־נֹ֔חַ וְאֵ֤ת כׇּל־הַֽחַיָּה֙ וְאֶת־כׇּל־הַבְּהֵמָ֔ה אֲשֶׁ֥ר אִתּ֖וֹ בַּתֵּבָ֑ה וַיַּעֲבֵ֨ר אֱלֹהִ֥ים ר֙וּחַ֙ עַל־הָאָ֔רֶץ וַיָּשֹׁ֖כּוּ הַמָּֽיִם׃ \Rashi{\textbf{ויזכור אלהים.}זה השם מדת הדין הוא, ונהפכה למדת רחמים על ידי תפלת הצדיקים; ורשעתן של רשעים הופכת מדת רחמים למדת הדין, שנאמר וירא ה' כי רבה רעת האדם, ויאמר ה' אמחה, (ברא' ו') והוא שם מדת רחמים: ויזכור אלהים את נח וגו'. מה זכר להם לבהמות? זכות, שלא השחיתו דרכם קדם לכן ושלא שמשו בתבה: ויעבר אלהים רוח. רוח תנחומין והנחה עברה לפניו: על הארץ. על עסקי הארץ: וישכו. כמו כשך חמת המלך (אסתר ב') לשון הנחת חמה:}%
% ---------- פרק 8 | פסוק 2 ----------
 \textbf{ב} וַיִּסָּֽכְרוּ֙ מַעְיְנֹ֣ת תְּה֔וֹם וַֽאֲרֻבֹּ֖ת הַשָּׁמָ֑יִם וַיִּכָּלֵ֥א הַגֶּ֖שֶׁם מִן־הַשָּׁמָֽיִם׃ \Rashi{\textbf{ויסכרו מעינות.}כשנפתחו כתיב כל מעינת; וכאן אין כתיב כל, לפי שנשתירו מהם אותן שיש בהם צרך לעולם, כגון חמי טבריא וכיוצא בהן: ויכלא. וימנע כמו לא תכלא רחמיך (תהלים מ'), לא יכלה ממך (ברא' כג):}%
% ---------- פרק 8 | פסוק 3 ----------
 \textbf{ג} וַיָּשֻׁ֧בוּ הַמַּ֛יִם מֵעַ֥ל הָאָ֖רֶץ הָל֣וֹךְ וָשׁ֑וֹב וַיַּחְסְר֣וּ הַמַּ֔יִם מִקְצֵ֕ה חֲמִשִּׁ֥ים וּמְאַ֖ת יֽוֹם׃ \Rashi{\textbf{מקצה חמשים ומאת יום.}התחילו לחסר והוא אחד בסיון. כיצד? בכ"ז בכסלו פסקו הגשמים, הרי ג' מכסלו וכ"ט מטבת הרי ל"ב, ושבט ואדר וניסן ואיר קי"ח, הרי ק"נ:}%
% ---------- פרק 8 | פסוק 4 ----------
 \textbf{ד} וַתָּ֤נַח הַתֵּבָה֙ בַּחֹ֣דֶשׁ הַשְּׁבִיעִ֔י בְּשִׁבְעָה־עָשָׂ֥ר י֖וֹם לַחֹ֑דֶשׁ עַ֖ל הָרֵ֥י אֲרָרָֽט׃ \Rashi{\textbf{בחדש השביעי.}סיון, והוא שביעי לכסלו שבו פסקו הגשמים: בשבעה עשר יום. מכאן אתה למד שהיתה התבה משקעת במים י"א אמה, שהרי כתיב בעשירי באחד לחדש נראו ראשי ההרים, זה אב, שהיא עשירי למרחשון לירידת גשמים, והם היו גבוהים על ההרים חמש עשרה אמה, וחסרו מיום אחד בסיון עד אחד באב חמש עשרה אמה לששים יום, הרי אמה לד' ימים, נמצא, שבי"ו בסיון לא חסרו אלא ד' אמות ונחה התבה ליום המחרת; למדת, שהיתה משקעת י"א אמה במים שעל ראשי ההרים:}%
% ---------- פרק 8 | פסוק 5 ----------
 \textbf{ה} וְהַמַּ֗יִם הָיוּ֙ הָל֣וֹךְ וְחָס֔וֹר עַ֖ד הַחֹ֣דֶשׁ הָֽעֲשִׂירִ֑י בָּֽעֲשִׂירִי֙ בְּאֶחָ֣ד לַחֹ֔דֶשׁ נִרְא֖וּ רָאשֵׁ֥י הֶֽהָרִֽים׃ \Rashi{\textbf{בעשירי נראו ראשי ההרים.}זה אב, שהוא עשירי למרחשון שהתחיל הגשם. ואם תאמר הוא אלול, ועשירי לכסלו שפסק הגשם, כשם שאתה אומר בחדש השביעי סיון, והוא שביעי להפסקה – אי אפשר לומר כן, על כרחך שביעי אי אתה מונה אלא להפסקה, שהרי לא כלו ארבעים יום של ירידת גשמים, ומאה וחמשים של תגברת המים עד אחד בסיון; ואם אתה אומר שביעי לירידה, אין זה סיון, והעשירי אי אפשר למנות אלא לירידה, שאם אתה אומר להפסקה והוא אלול, אי אתה מוצא בראשון באחד לחדש חרבו המים מעל הארץ; שהרי מקץ ארבעים יום משנראו ראשי ההרים שלח את העורב וכ"א יום הוחיל בשליחות היונה, הרי ששים יום משנראו ראשי ההרים עד שחרבו פני האדמה. ואם תאמר באלול נראו, נמצא שחרבו במרחשון והוא קורא אותו ראשון, ואין זה אלא תשרי, שהוא ראשון לבריאת עולם, ולרבי יהושע הוא ניסן:}%
% ---------- פרק 8 | פסוק 6 ----------
 \textbf{ו} וַיְהִ֕י מִקֵּ֖ץ אַרְבָּעִ֣ים י֑וֹם וַיִּפְתַּ֣ח נֹ֔חַ אֶת־חַלּ֥וֹן הַתֵּבָ֖ה אֲשֶׁ֥ר עָשָֽׂה׃ \Rashi{\textbf{מקץ ארבעים יום.}משנראו ראשי ההרים: את חלון התבה אשר עשה. לצהר, ולא זה פתח התבה העשוי לביאה ויציאה:}%
% ---------- פרק 8 | פסוק 7 ----------
 \textbf{ז} וַיְשַׁלַּ֖ח אֶת־הָֽעֹרֵ֑ב וַיֵּצֵ֤א יָצוֹא֙ וָשׁ֔וֹב עַד־יְבֹ֥שֶׁת הַמַּ֖יִם מֵעַ֥ל הָאָֽרֶץ׃ \Rashi{\textbf{יצוא ושוב.}הולך ומקיף סביבות התבה, ולא הלך בשליחותו, שהיה חושדו על בת זוגו, כמו ששנינו באגדת חלק (סנה' ק"ח): עד יבשת המים. פשוטו כמשמעו. אבל מדרש אגדה מוכן היה העורב לשליחות אחרת בעצירת גשמים בימי אליהו, שנאמר והערבים מבאים לו לחם ובשר (מ"א י"ז):}%
% ---------- פרק 8 | פסוק 8 ----------
 \textbf{ח} וַיְשַׁלַּ֥ח אֶת־הַיּוֹנָ֖ה מֵאִתּ֑וֹ לִרְאוֹת֙ הֲקַ֣לּוּ הַמַּ֔יִם מֵעַ֖ל פְּנֵ֥י הָֽאֲדָמָֽה׃ \Rashi{\textbf{וישלח את היונה.}לסוף ז' ימים שהרי כתיב ויחל עוד ז' ימים אחרים, מכלל זה אתה למד שאף בראשונה הוחיל ז' ימים: וישלח. אין זה לשון שליחות אלא לשון שלוח, שלחה ללכת לדרכה, ובזו יראה אם קלו המים, שאם תמצא מנוח לא תשוב אליו:}%
% ---------- פרק 8 | פסוק 9 ----------
 \textbf{ט} וְלֹֽא־מָצְאָה֩ הַיּוֹנָ֨ה מָנ֜וֹחַ לְכַף־רַגְלָ֗הּ וַתָּ֤שׇׁב אֵלָיו֙ אֶל־הַתֵּבָ֔ה כִּי־מַ֖יִם עַל־פְּנֵ֣י כׇל־הָאָ֑רֶץ וַיִּשְׁלַ֤ח יָדוֹ֙ וַיִּקָּחֶ֔הָ וַיָּבֵ֥א אֹתָ֛הּ אֵלָ֖יו אֶל־הַתֵּבָֽה׃ %
% ---------- פרק 8 | פסוק 10 ----------
 \textbf{י} וַיָּ֣חֶל ע֔וֹד שִׁבְעַ֥ת יָמִ֖ים אֲחֵרִ֑ים וַיֹּ֛סֶף שַׁלַּ֥ח אֶת־הַיּוֹנָ֖ה מִן־הַתֵּבָֽה׃ \Rashi{\textbf{ויחל.}לשון המתנה, וכן לי שמעו ויחלו (איוב כ"ט), והרבה יש במקרא:}%
% ---------- פרק 8 | פסוק 11 ----------
 \textbf{יא} וַתָּבֹ֨א אֵלָ֤יו הַיּוֹנָה֙ לְעֵ֣ת עֶ֔רֶב וְהִנֵּ֥ה עֲלֵה־זַ֖יִת טָרָ֣ף בְּפִ֑יהָ וַיֵּ֣דַע נֹ֔חַ כִּי־קַ֥לּוּ הַמַּ֖יִם מֵעַ֥ל הָאָֽרֶץ׃ \Rashi{\textbf{טרף בפיה.}אומר אני שזכר היה, לכן קוראו פעמים לשון זכר ופעמים לשון נקבה, לפי שכל יונה שבמקרא לשון נקבה, כמו כיונים על אפיקי מים רחצות (שיר ה'), כיוני הגאיות כלם המות (יחז' ז') וכמו כיונה פותה (הושע ז'): טרף. חטף. ומדרש אגדה (עיר' י"ח) לשון מזון; ודרשו בפיה לשון מאמר, אמרה יהיו מזונותי מרורין כזית בידו של הקב"ה ולא מתוקין כדבש בידי בשר ודם:}%
% ---------- פרק 8 | פסוק 12 ----------
 \textbf{יב} וַיִּיָּ֣חֶל ע֔וֹד שִׁבְעַ֥ת יָמִ֖ים אֲחֵרִ֑ים וַיְשַׁלַּח֙ אֶת־הַיּוֹנָ֔ה וְלֹֽא־יָסְפָ֥ה שׁוּב־אֵלָ֖יו עֽוֹד׃ \Rashi{\textbf{וייחל.}הוא לשון ויחל, אלא שזה לשון ויפעל וזה לשון ויתפעל, ויחל – וימתן; וייחל – ויתמתן:}%
% ---------- פרק 8 | פסוק 13 ----------
 \textbf{יג} וַ֠יְהִ֠י בְּאַחַ֨ת וְשֵׁשׁ־מֵא֜וֹת שָׁנָ֗ה בָּֽרִאשׁוֹן֙ בְּאֶחָ֣ד לַחֹ֔דֶשׁ חָֽרְב֥וּ הַמַּ֖יִם מֵעַ֣ל הָאָ֑רֶץ וַיָּ֤סַר נֹ֙חַ֙ אֶת־מִכְסֵ֣ה הַתֵּבָ֔ה וַיַּ֕רְא וְהִנֵּ֥ה חָֽרְב֖וּ פְּנֵ֥י הָֽאֲדָמָֽה׃ \Rashi{\textbf{בראשון.}לרבי אליעזר הוא תשרי ולרבי יהושע הוא ניסן (ר"ה י"א): חרבו. נעשה כמין טיט, שקרמו פניה של מעלה:}%
% ---------- פרק 8 | פסוק 14 ----------
 \textbf{יד} וּבַחֹ֙דֶשׁ֙ הַשֵּׁנִ֔י בְּשִׁבְעָ֧ה וְעֶשְׂרִ֛ים י֖וֹם לַחֹ֑דֶשׁ יָבְשָׁ֖ה הָאָֽרֶץ׃ \{ס\} \Rashi{\textbf{בשבעה ועשרים.}ירידתן בחדש השני בי"ז בחדש, אלו י"א ימים שהחמה יתרה על הלבנה, שמשפט דור המבול שנה תמימה היה: יבשה. נעשה גריד כהלכתה:}%
% ---------- פרק 8 | פסוק 15 ----------
 \textbf{טו} וַיְדַבֵּ֥ר אֱלֹהִ֖ים אֶל־נֹ֥חַ לֵאמֹֽר׃ %
% ---------- פרק 8 | פסוק 16 ----------
 \textbf{טז} צֵ֖א מִן־הַתֵּבָ֑ה אַתָּ֕ה וְאִשְׁתְּךָ֛ וּבָנֶ֥יךָ וּנְשֵֽׁי־בָנֶ֖יךָ אִתָּֽךְ׃ \Rashi{\textbf{אתה ואשתך וגו'.}איש ואשתו, כאן התיר להם תשמיש המטה:}%
% ---------- פרק 8 | פסוק 17 ----------
 \textbf{יז} כׇּל־הַחַיָּ֨ה אֲשֶֽׁר־אִתְּךָ֜ מִכׇּל־בָּשָׂ֗ר בָּע֧וֹף וּבַבְּהֵמָ֛ה וּבְכׇל־הָרֶ֛מֶשׂ הָרֹמֵ֥שׂ עַל־הָאָ֖רֶץ (הוצא) [הַיְצֵ֣א] אִתָּ֑ךְ וְשָֽׁרְצ֣וּ בָאָ֔רֶץ וּפָר֥וּ וְרָב֖וּ עַל־הָאָֽרֶץ׃ \Rashi{\textbf{הוצא.}הוצא כתיב, היצא קרי, היצא – אמר להם שיצאו; הוצא – אם אינם רוצים לצאת הוציאם אתה: ושרצו בארץ. ולא בתבה, מגיד שאף הבהמה והעוף נאסרו בתשמיש:}%
% ---------- פרק 8 | פסוק 18 ----------
 \textbf{יח} וַיֵּ֖צֵא־נֹ֑חַ וּבָנָ֛יו וְאִשְׁתּ֥וֹ וּנְשֵֽׁי־בָנָ֖יו אִתּֽוֹ׃ %
% ---------- פרק 8 | פסוק 19 ----------
 \textbf{יט} כׇּל־הַֽחַיָּ֗ה כׇּל־הָרֶ֙מֶשׂ֙ וְכׇל־הָע֔וֹף כֹּ֖ל רוֹמֵ֣שׂ עַל־הָאָ֑רֶץ לְמִשְׁפְּחֹ֣תֵיהֶ֔ם יָצְא֖וּ מִן־הַתֵּבָֽה׃ \Rashi{\textbf{למשפחותיהם.}קבלו עליהם על מנת לדבק במינן:}%
% ---------- פרק 8 | פסוק 20 ----------
 \textbf{כ} וַיִּ֥בֶן נֹ֛חַ מִזְבֵּ֖חַ לַֽיהֹוָ֑ה וַיִּקַּ֞ח מִכֹּ֣ל ׀ הַבְּהֵמָ֣ה הַטְּהֹרָ֗ה וּמִכֹּל֙ הָע֣וֹף הַטָּה֔וֹר וַיַּ֥עַל עֹלֹ֖ת בַּמִּזְבֵּֽחַ׃ \Rashi{\textbf{מכל הבהמה הטהורה.}אמר לא צוה לי הקב"ה להכניס מאלו ז' ז' אלא כדי להקריב קרבן מהם:}%
% ---------- פרק 8 | פסוק 21 ----------
 \textbf{כא} וַיָּ֣רַח יְהֹוָה֮ אֶת־רֵ֣יחַ הַנִּיחֹ֒חַ֒ וַיֹּ֨אמֶר יְהֹוָ֜ה אֶל־לִבּ֗וֹ לֹֽא־אֹ֠סִ֠ף לְקַלֵּ֨ל ע֤וֹד אֶת־הָֽאֲדָמָה֙ בַּעֲב֣וּר הָֽאָדָ֔ם כִּ֠י יֵ֣צֶר לֵ֧ב הָאָדָ֛ם רַ֖ע מִנְּעֻרָ֑יו וְלֹֽא־אֹסִ֥ף ע֛וֹד לְהַכּ֥וֹת אֶת־כׇּל־חַ֖י כַּֽאֲשֶׁ֥ר עָשִֽׂיתִי׃ \Rashi{\textbf{מנעריו.}מנעריו כתיב, משננער לצאת ממעי אמו נתן בו יצר הרע: לא אסף ולא אסף. כפל הדבר לשבועה; הוא שכתוב אשר נשבעתי מעבר מי נח, ולא מצינו בה שבועה, אלא זו שכפל דבריו והיא שבועה וכן דרשו חכמים במס' שבועות:}%
% ---------- פרק 8 | פסוק 22 ----------
 \textbf{כב} עֹ֖ד כׇּל־יְמֵ֣י הָאָ֑רֶץ זֶ֡רַע וְ֠קָצִ֠יר וְקֹ֨ר וָחֹ֜ם וְקַ֧יִץ וָחֹ֛רֶף וְי֥וֹם וָלַ֖יְלָה לֹ֥א יִשְׁבֹּֽתוּ׃ \Rashi{\textbf{עוד כל ימי הארץ וגו' לא ישבתו.}ו' עתים הללו שני חדשים לכל אחד ואחד, כמו ששנינו חצי תשרי ומרחשון וחצי כסלו זרע, חצי כסלו וטבת וחצי שבט קר (חרף) (ב"מ ק"ו): קור. קשה מחרף: חרף. עת זרע שעורים וקטנית החריפין להתבשל מהר, והוא חצי שבט ואדר וחצי ניסן: קציר. חצי ניסן ואיר וחצי סיון: קיץ. חצי סיון תמוז וחצי אב, הוא זמן לקיטת תאנים וזמן שמיבשים אותן בשדות, ושמו קיץ, כמו הלחם והקיץ לאכול הנערים (ש"ב ט"ז): חום. הוא סוף ימות החמה, חצי אב ואלול וחצי תשרי, שהעולם חם ביותר; כמו ששנינו במסכת יומא שלהי קייטא קשי מקייטא: ויום ולילה לא ישבתו. מכלל ששבתו כל ימות המבול, שלא שמשו המזלות, ולא נכר בין יום ובין לילה (ב"ר פכ"ה ופ' ל"ד): לא ישבתו. לא יפסקו כל אלה מלהתנהג כסדרן:}%
% ---------- פרק 9 | פסוק 1 ----------
 \textbf{א} וַיְבָ֣רֶךְ אֱלֹהִ֔ים אֶת־נֹ֖חַ וְאֶת־בָּנָ֑יו וַיֹּ֧אמֶר לָהֶ֛ם פְּר֥וּ וּרְב֖וּ וּמִלְא֥וּ אֶת־הָאָֽרֶץ׃ %
% ---------- פרק 9 | פסוק 2 ----------
 \textbf{ב} וּמוֹרַאֲכֶ֤ם וְחִתְּכֶם֙ יִֽהְיֶ֔ה עַ֚ל כׇּל־חַיַּ֣ת הָאָ֔רֶץ וְעַ֖ל כׇּל־ע֣וֹף הַשָּׁמָ֑יִם בְּכֹל֩ אֲשֶׁ֨ר תִּרְמֹ֧שׂ הָֽאֲדָמָ֛ה וּֽבְכׇל־דְּגֵ֥י הַיָּ֖ם בְּיֶדְכֶ֥ם נִתָּֽנוּ׃ \Rashi{\textbf{וחתכם.}ואימתכם, כמו תראו חתת (איוב ו') ואגדה לשון חיות שכל זמן שתינוק בן יומו חי אין אתה צריך לשמרו מן העכברים; עוג מלך הבשן מת, צריך לשמרו מן העכברים, שנאמר ומוראכם וחתכם יהיה, אימתי יהיה מוראכם על החיות? כל זמן שאתם חיים (שבת קנ"א):}%
% ---------- פרק 9 | פסוק 3 ----------
 \textbf{ג} כׇּל־רֶ֙מֶשׂ֙ אֲשֶׁ֣ר הוּא־חַ֔י לָכֶ֥ם יִהְיֶ֖ה לְאׇכְלָ֑ה כְּיֶ֣רֶק עֵ֔שֶׂב נָתַ֥תִּי לָכֶ֖ם אֶת־כֹּֽל׃ \Rashi{\textbf{לכם יהיה לאכלה.}שלא הרשיתי לאדם הראשון בשר אלא ירק עשב; ולכם – כירק עשב שהפקרתי לאדם הראשון, נתתי לכם את כל (סנה' נ"ט):}%
% ---------- פרק 9 | פסוק 4 ----------
 \textbf{ד} אַךְ־בָּשָׂ֕ר בְּנַפְשׁ֥וֹ דָמ֖וֹ לֹ֥א תֹאכֵֽלוּ׃ \Rashi{\textbf{בשר בנפשו.}אסר להם אבר מן החי, כלומר כל זמן שנפשו בו, לא תאכלו הבשר: בנפשו דמו. בעוד נפשו בו בשר בנפשו לא תאכלו. הרי אבר מן החי, ואף בנפשו דמו לא תאכלו, הרי דם מן החי:}%
% ---------- פרק 9 | פסוק 5 ----------
 \textbf{ה} וְאַ֨ךְ אֶת־דִּמְכֶ֤ם לְנַפְשֹֽׁתֵיכֶם֙ אֶדְרֹ֔שׁ מִיַּ֥ד כׇּל־חַיָּ֖ה אֶדְרְשֶׁ֑נּוּ וּמִיַּ֣ד הָֽאָדָ֗ם מִיַּד֙ אִ֣ישׁ אָחִ֔יו אֶדְרֹ֖שׁ אֶת־נֶ֥פֶשׁ הָֽאָדָֽם׃ \Rashi{\textbf{ואך את דמכם.}אע"פ שהתרתי לכם נטילת נשמה בבהמה, את דמכם אדרש מהשופך דם עצמו: לנפשתיכם. אף החונק עצמו, אף על פי שלא יצא ממנו דם: מיד כל חיה. לפי שחטאו דור המבול והפקרו למאכל חיות רעות לשלט בהן, שנאמר נמשל כבהמות נדמו (תהלים מ"ט), לפיכך הצרך להזהיר עליהן את החיות: ומיד האדם. מיד ההורג במזיד ואין עדים, אני אדרש: מיד איש אחיו. שהוא אוהב לו כאח והרג שוגג, אני אדרש אם לא יגלה ויבקש על עונו למחל, שאף השוגג צריך כפרה, ואם אין עדים לחיבו גלות והוא אינו נכנע, הקב"ה דורש ממנו, כמו שדרשו רבותינו "והאלהים אנה לידו" – במסכת מכות – הקב"ה מזמנן לפנדק אחד וכו':}%
% ---------- פרק 9 | פסוק 6 ----------
 \textbf{ו} שֹׁפֵךְ֙ דַּ֣ם הָֽאָדָ֔ם בָּֽאָדָ֖ם דָּמ֣וֹ יִשָּׁפֵ֑ךְ כִּ֚י בְּצֶ֣לֶם אֱלֹהִ֔ים עָשָׂ֖ה אֶת־הָאָדָֽם׃ \Rashi{\textbf{באדם דמו ישפך.}אם יש עדים המיתוהו אתם, למה? כי בצלם אלהים וגו': עשה את האדם. זה מקרא חסר, וצריך להיות עשה העושה את האדם וכן הרבה במקרא:}%
% ---------- פרק 9 | פסוק 7 ----------
 \textbf{ז} וְאַתֶּ֖ם פְּר֣וּ וּרְב֑וּ שִׁרְצ֥וּ בָאָ֖רֶץ וּרְבוּ־בָֽהּ׃ \{ס\} \Rashi{\textbf{ואתם פרו ורבו.}לפי פשוטו הראשונה לברכה וכאן לצווי. ולפי מדרשו להקיש מי שאינו עוסק בפריה ורביה לשופך דמים (יבמות ס"ג):}%
% ---------- פרק 9 | פסוק 8 ----------
 \textbf{ח} וַיֹּ֤אמֶר אֱלֹהִים֙ אֶל־נֹ֔חַ וְאֶל־בָּנָ֥יו אִתּ֖וֹ לֵאמֹֽר׃ %
% ---------- פרק 9 | פסוק 9 ----------
 \textbf{ט} וַאֲנִ֕י הִנְנִ֥י מֵקִ֛ים אֶת־בְּרִיתִ֖י אִתְּכֶ֑ם וְאֶֽת־זַרְעֲכֶ֖ם אַֽחֲרֵיכֶֽם׃ \Rashi{\textbf{ואני הנני.}מסכים אני עמך, שהיה נח דואג לעסק בפריה ורביה עד שהבטיחו הקב"ה שלא לשחת העולם עוד, וכן עשה ובאחרונה אמר לו הנני מסכים לעשות קיום וחזוק ברית להבטחתי ואתן לך אות:}%
% ---------- פרק 9 | פסוק 10 ----------
 \textbf{י} וְאֵ֨ת כׇּל־נֶ֤פֶשׁ הַֽחַיָּה֙ אֲשֶׁ֣ר אִתְּכֶ֔ם בָּע֧וֹף בַּבְּהֵמָ֛ה וּֽבְכׇל־חַיַּ֥ת הָאָ֖רֶץ אִתְּכֶ֑ם מִכֹּל֙ יֹצְאֵ֣י הַתֵּבָ֔ה לְכֹ֖ל חַיַּ֥ת הָאָֽרֶץ׃ \Rashi{\textbf{חית הארץ אתכם.}הם המתהלכים עם הבריות: מכל יוצאי התבה. להביא שקצים ורמשים: חית הארץ. להביא המזיקין, שאינן בכלל החיה אשר אתכם, שאין הלוכן עם הבריות:}%
% ---------- פרק 9 | פסוק 11 ----------
 \textbf{יא} וַהֲקִמֹתִ֤י אֶת־בְּרִיתִי֙ אִתְּכֶ֔ם וְלֹֽא־יִכָּרֵ֧ת כׇּל־בָּשָׂ֛ר ע֖וֹד מִמֵּ֣י הַמַּבּ֑וּל וְלֹֽא־יִהְיֶ֥ה ע֛וֹד מַבּ֖וּל לְשַׁחֵ֥ת הָאָֽרֶץ׃ \Rashi{\textbf{והקמותי.}אעשה קיום לבריתי. ומהו קיומו? אות הקשת כמו שמסים והולך:}%
% ---------- פרק 9 | פסוק 12 ----------
 \textbf{יב} וַיֹּ֣אמֶר אֱלֹהִ֗ים זֹ֤את אֽוֹת־הַבְּרִית֙ אֲשֶׁר־אֲנִ֣י נֹתֵ֗ן בֵּינִי֙ וּבֵ֣ינֵיכֶ֔ם וּבֵ֛ין כׇּל־נֶ֥פֶשׁ חַיָּ֖ה אֲשֶׁ֣ר אִתְּכֶ֑ם לְדֹרֹ֖ת עוֹלָֽם׃ \Rashi{\textbf{לדרת עולם.}נכתב חסר, שיש דורות שלא הצרכו לאות לפי שצדיקים גמורים היו, כמו דורו של חזקיהו מלך יהודה ודורו של רבי שמעון בן יוחאי (ב"ר):}%
% ---------- פרק 9 | פסוק 13 ----------
 \textbf{יג} אֶת־קַשְׁתִּ֕י נָתַ֖תִּי בֶּֽעָנָ֑ן וְהָֽיְתָה֙ לְא֣וֹת בְּרִ֔ית בֵּינִ֖י וּבֵ֥ין הָאָֽרֶץ׃ %
% ---------- פרק 9 | פסוק 14 ----------
 \textbf{יד} וְהָיָ֕ה בְּעַֽנְנִ֥י עָנָ֖ן עַל־הָאָ֑רֶץ וְנִרְאֲתָ֥ה הַקֶּ֖שֶׁת בֶּעָנָֽן׃ \Rashi{\textbf{בענני ענן.}כשתעלה במחשבה לפני להביא חשך ואבדון לעולם:}%
% ---------- פרק 9 | פסוק 15 ----------
 \textbf{טו} וְזָכַרְתִּ֣י אֶת־בְּרִיתִ֗י אֲשֶׁ֤ר בֵּינִי֙ וּבֵ֣ינֵיכֶ֔ם וּבֵ֛ין כׇּל־נֶ֥פֶשׁ חַיָּ֖ה בְּכׇל־בָּשָׂ֑ר וְלֹֽא־יִֽהְיֶ֨ה ע֤וֹד הַמַּ֙יִם֙ לְמַבּ֔וּל לְשַׁחֵ֖ת כׇּל־בָּשָֽׂר׃ %
% ---------- פרק 9 | פסוק 16 ----------
 \textbf{טז} וְהָיְתָ֥ה הַקֶּ֖שֶׁת בֶּֽעָנָ֑ן וּרְאִיתִ֗יהָ לִזְכֹּר֙ בְּרִ֣ית עוֹלָ֔ם בֵּ֣ין אֱלֹהִ֔ים וּבֵין֙ כׇּל־נֶ֣פֶשׁ חַיָּ֔ה בְּכׇל־בָּשָׂ֖ר אֲשֶׁ֥ר עַל־הָאָֽרֶץ׃ \Rashi{\textbf{בין אלהים ובין כל נפש חיה.}בין מדת הדין של מעלה וביניכם, שהיה לו לכתב ביני ובין כל נפש חיה, אלא זהו מדרשו כשתבא מדת הדין לקטרג עליכם לחיב אתכם, אני רואה את האות ונזכר:}%
% ---------- פרק 9 | פסוק 17 ----------
 \textbf{יז} וַיֹּ֥אמֶר אֱלֹהִ֖ים אֶל־נֹ֑חַ זֹ֤את אֽוֹת־הַבְּרִית֙ אֲשֶׁ֣ר הֲקִמֹ֔תִי בֵּינִ֕י וּבֵ֥ין כׇּל־בָּשָׂ֖ר אֲשֶׁ֥ר עַל־הָאָֽרֶץ׃ \{פ\} \Rashi{\textbf{זאת אות הברית.}הראהו הקשת ואמר לו הרי האות שאמרתי:}%
% ---------- פרק 9 | פסוק 18 ----------
 \textbf{יח} וַיִּֽהְי֣וּ בְנֵי־נֹ֗חַ הַיֹּֽצְאִים֙ מִן־הַתֵּבָ֔ה שֵׁ֖ם וְחָ֣ם וָיָ֑פֶת וְחָ֕ם ה֖וּא אֲבִ֥י כְנָֽעַן׃ \Rashi{\textbf{וחם הוא אבי כנען.}למה הצרך לומר כאן? לפי שהפרשה עסוקה ובאה בשכרותו של נח שקלקל בה חם ועל ידו נתקלל כנען, ועדין לא כתב תולדות חם, ולא ידענו שכנען בנו – לפיכך הצרך לומר כאן וחם הוא אבי כנען:}%
% ---------- פרק 9 | פסוק 19 ----------
 \textbf{יט} שְׁלֹשָׁ֥ה אֵ֖לֶּה בְּנֵי־נֹ֑חַ וּמֵאֵ֖לֶּה נָֽפְצָ֥ה כׇל־הָאָֽרֶץ׃ %
% ---------- פרק 9 | פסוק 20 ----------
 \textbf{כ} וַיָּ֥חֶל נֹ֖חַ אִ֣ישׁ הָֽאֲדָמָ֑ה וַיִּטַּ֖ע כָּֽרֶם׃ %
% ---------- פרק 9 | פסוק 21 ----------
 \textbf{כא} וַיֵּ֥שְׁתְּ מִן־הַיַּ֖יִן וַיִּשְׁכָּ֑ר וַיִּתְגַּ֖ל בְּת֥וֹךְ אׇהֳלֹֽה׃ \Rashi{\textbf{ויתגל.}לשון ויתפעל: אהלה. אהלה כתיב, רמז לי' שבטים שנקראו על שם שומרון, שנקראת אהלה שגלו על עסקי היין שנאמר השתים במזרקי יין (עמוס ו'):}%
% ---------- פרק 9 | פסוק 22 ----------
 \textbf{כב} וַיַּ֗רְא חָ֚ם אֲבִ֣י כְנַ֔עַן אֵ֖ת עֶרְוַ֣ת אָבִ֑יו וַיַּגֵּ֥ד לִשְׁנֵֽי־אֶחָ֖יו בַּחֽוּץ׃ \Rashi{\textbf{וירא חם אבי כנען.}יש מרבותינו אומרים כנען ראה והגיד לאביו, לכך הזכר על הדבר ונתקלל (ב"ר): וירא את ערות אביו. יש אומרים סרסו ויש אומרים רבעו (סנה' ע'):}%
% ---------- פרק 9 | פסוק 23 ----------
 \textbf{כג} וַיִּקַּח֩ שֵׁ֨ם וָיֶ֜פֶת אֶת־הַשִּׂמְלָ֗ה וַיָּשִׂ֙ימוּ֙ עַל־שְׁכֶ֣ם שְׁנֵיהֶ֔ם וַיֵּֽלְכוּ֙ אֲחֹ֣רַנִּ֔ית וַיְכַסּ֕וּ אֵ֖ת עֶרְוַ֣ת אֲבִיהֶ֑ם וּפְנֵיהֶם֙ אֲחֹ֣רַנִּ֔ית וְעֶרְוַ֥ת אֲבִיהֶ֖ם לֹ֥א רָאֽוּ׃ %
% ---------- פרק 9 | פסוק 24 ----------
 \textbf{כד} וַיִּ֥יקֶץ נֹ֖חַ מִיֵּינ֑וֹ וַיֵּ֕דַע אֵ֛ת אֲשֶׁר־עָ֥שָׂה ל֖וֹ בְּנ֥וֹ הַקָּטָֽן׃ \Rashi{\textbf{בנו הקטן.}הפסול והבזוי, כמו הנה קטן נתתיך בגוים בזוי באדם (ירמיה מ"ט):}%
% ---------- פרק 9 | פסוק 25 ----------
 \textbf{כה} וַיֹּ֖אמֶר אָר֣וּר כְּנָ֑עַן עֶ֥בֶד עֲבָדִ֖ים יִֽהְיֶ֥ה לְאֶחָֽיו׃ \Rashi{\textbf{ארור כנען.}אתה גרמת לי שלא אוליד בן רביעי אחר לשמשני, ארור בנך רביעי להיות משמש את זרעם של אלו הגדולים, שהוטל עליהם טרח עבודתי מעתה. ומה ראה חם שסרסו? אמר להם לאחיו אדם הראשון שני בנים היו לו, והרג זה את זה בשביל ירשת העולם, ואבינו יש לו ג' בנים ועודנו מבקש בן רביעי:}%
% ---------- פרק 9 | פסוק 26 ----------
 \textbf{כו} וַיֹּ֕אמֶר בָּר֥וּךְ יְהֹוָ֖ה אֱלֹ֣הֵי שֵׁ֑ם וִיהִ֥י כְנַ֖עַן עֶ֥בֶד לָֽמוֹ׃ \Rashi{\textbf{ברוך ה' אלהי שם.}שעתיד לשמר הבטחתו לזרעו, לתת להם את ארץ כנען: ויהי. להם כנען למס עובד:}%
% ---------- פרק 9 | פסוק 27 ----------
 \textbf{כז} יַ֤פְתְּ אֱלֹהִים֙ לְיֶ֔פֶת וְיִשְׁכֹּ֖ן בְּאׇֽהֳלֵי־שֵׁ֑ם וִיהִ֥י כְנַ֖עַן עֶ֥בֶד לָֽמוֹ׃ \Rashi{\textbf{יפת אלהים ליפת.}מתרגם יפתי, ירחיב: וישכן באהלי שם. ישרה שכינתו בישראל. ומדרש חכמים אע"פ שיפת אלהים ליפת – שבנה כרש, שהיה מבני יפת, בית שני – לא שרתה בו שכינה; והיכן שרתה? במקדש ראשון שבנה שלמה שהיה מבני שם: ויהי כנען עבד למו. אף משיגלו בני שם ימכרו להם עבדים מבני כנען:}%
% ---------- פרק 9 | פסוק 28 ----------
 \textbf{כח} וַֽיְחִי־נֹ֖חַ אַחַ֣ר הַמַּבּ֑וּל שְׁלֹ֤שׁ מֵאוֹת֙ שָׁנָ֔ה וַֽחֲמִשִּׁ֖ים שָׁנָֽה׃ %
% ---------- פרק 9 | פסוק 29 ----------
 \textbf{כט} וַיִּֽהְיוּ֙*(בספרי ספרד ואשכנז וַֽיְהִי֙) כׇּל־יְמֵי־נֹ֔חַ תְּשַׁ֤ע מֵאוֹת֙ שָׁנָ֔ה וַחֲמִשִּׁ֖ים שָׁנָ֑ה וַיָּמֹֽת׃ \{פ\} %
% ---------- פרק 11 | פסוק 1 ----------
 \textbf{א} וַיְהִ֥י כׇל־הָאָ֖רֶץ שָׂפָ֣ה אֶחָ֑ת וּדְבָרִ֖ים אֲחָדִֽים׃ \Rashi{\textbf{שפה אחת.}לשון הקדש: ודברים אחדים. באו בעצה אחת ואמרו לא כל הימנו שיבר לו את העליונים, נעלה לרקיע ונעשה עמו מלחמה. דבר אחר על יחידו של עולם. דבר אחר ודברים אחדים – אמרו אחת לאלף ותרנ"ו שנים הרקיע מתמוטט כשם שעשה בימי המבול, באו ונעשה לו סמיכות (ב"ר):}%
% ---------- פרק 11 | פסוק 2 ----------
 \textbf{ב} וַיְהִ֖י בְּנׇסְעָ֣ם מִקֶּ֑דֶם וַֽיִּמְצְא֥וּ בִקְעָ֛ה בְּאֶ֥רֶץ שִׁנְעָ֖ר וַיֵּ֥שְׁבוּ שָֽׁם׃ \Rashi{\textbf{בנסעם מקדם.}שהיו יושבים שם, כדכתיב למעלה ויהי מושבם וגו' הר הקדם ונסעו משם לתור להם מקום להחזיק את כלם ולא מצאו אלא שנער:}%
% ---------- פרק 11 | פסוק 3 ----------
 \textbf{ג} וַיֹּאמְר֞וּ אִ֣ישׁ אֶל־רֵעֵ֗הוּ הָ֚בָה נִלְבְּנָ֣ה לְבֵנִ֔ים וְנִשְׂרְפָ֖ה לִשְׂרֵפָ֑ה וַתְּהִ֨י לָהֶ֤ם הַלְּבֵנָה֙ לְאָ֔בֶן וְהַ֣חֵמָ֔ר הָיָ֥ה לָהֶ֖ם לַחֹֽמֶר׃ \Rashi{\textbf{איש אל רעהו.}אמה לאמה, מצרים לכוש, וכוש לפוט ופוט לכנען: הבה. הזמינו עצמכם. כל הבה לשון הזמנה הוא, שמכינים עצמן ומתחברים למלאכה או לעצה או למשא. הבה הזמינו, אפר"לייר בלע"ז: לבנים. שאין אבנים בבבל, שהיא בקעה: ונשרפה לשרפה. כך עושין הלבנים שקורין טוויל"ש בל', שורפים אותן בכבשן: לחמר. לטוח הקיר:}%
% ---------- פרק 11 | פסוק 4 ----------
 \textbf{ד} וַיֹּאמְר֞וּ הָ֣בָה ׀ נִבְנֶה־לָּ֣נוּ עִ֗יר וּמִגְדָּל֙ וְרֹאשׁ֣וֹ בַשָּׁמַ֔יִם וְנַֽעֲשֶׂה־לָּ֖נוּ שֵׁ֑ם פֶּן־נָפ֖וּץ עַל־פְּנֵ֥י כׇל־הָאָֽרֶץ׃ \Rashi{\textbf{פן נפוץ.}שלא יביא עלינו שום מכה להפיצנו מכאן:}%
% ---------- פרק 11 | פסוק 5 ----------
 \textbf{ה} וַיֵּ֣רֶד יְהֹוָ֔ה לִרְאֹ֥ת אֶת־הָעִ֖יר וְאֶת־הַמִּגְדָּ֑ל אֲשֶׁ֥ר בָּנ֖וּ בְּנֵ֥י הָאָדָֽם׃ \Rashi{\textbf{וירד ה' לראות.}לא הצרך לכך, אלא בא ללמד לדינים שלא ירשיעו הנדון עד שיראו ויבינו, מדרש רבי תנחומא: בני האדם. אלא בני מי? שמא בני חמורים וגמלים? אלא בני אדם הראשון שכפה את הטובה ואמר האשה אשר נתת עמדי, אף אלו כפו בטובה למרד במי שהשפיעם טובה, ומלטם מן המבול:}%
% ---------- פרק 11 | פסוק 6 ----------
 \textbf{ו} וַיֹּ֣אמֶר יְהֹוָ֗ה הֵ֣ן עַ֤ם אֶחָד֙ וְשָׂפָ֤ה אַחַת֙ לְכֻלָּ֔ם וְזֶ֖ה הַחִלָּ֣ם לַעֲשׂ֑וֹת וְעַתָּה֙ לֹֽא־יִבָּצֵ֣ר מֵהֶ֔ם כֹּ֛ל אֲשֶׁ֥ר יָזְמ֖וּ לַֽעֲשֽׂוֹת׃ \Rashi{\textbf{הן עם אחד.}כל טובה זו יש עמהן, שעם אחד הם, ושפה אחת לכלן ודבר זה החלו לעשות: החלם. כמו אמרם, עשותם – להתחיל הם לעשות: לא יבצר מהם וגו' לעשות. בתמיה. יבצר לשון מניעה, כתרגומו; ודומה לו יבצר רוח נגידים (תה' ע"ו):}%
% ---------- פרק 11 | פסוק 7 ----------
 \textbf{ז} הָ֚בָה נֵֽרְדָ֔ה וְנָבְלָ֥ה שָׁ֖ם שְׂפָתָ֑ם אֲשֶׁר֙ לֹ֣א יִשְׁמְע֔וּ אִ֖ישׁ שְׂפַ֥ת רֵעֵֽהוּ׃ \Rashi{\textbf{הבה נרדה.}בבית דינו נמלך מענותנותו יתרה: הבה. מדה כנגד מדה – הם אמרו הבה נבנה, והוא כנגדם מדד ואמר הבה נרדה: ונבלה. ונבלבל; נו"ן משמש בלשון רבים, וה"א אחרונה יתרה, כה"א של נרדה: לא ישמעו. זה שואל לבנה וזה מביא טיט, וזה עומד עליו ופוצע את מחו:}%
% ---------- פרק 11 | פסוק 8 ----------
 \textbf{ח} וַיָּ֨פֶץ יְהֹוָ֥ה אֹתָ֛ם מִשָּׁ֖ם עַל־פְּנֵ֣י כׇל־הָאָ֑רֶץ וַֽיַּחְדְּל֖וּ לִבְנֹ֥ת הָעִֽיר׃ \Rashi{\textbf{ויפץ ה' אתם משם.}בעולם הזה; מה שאמרו פן נפוץ, נתקים עליהם, הוא שאמר שלמה מגורת רשע היא תבואנו (משלי י'):}%
% ---------- פרק 11 | פסוק 9 ----------
 \textbf{ט} עַל־כֵּ֞ן קָרָ֤א שְׁמָהּ֙ בָּבֶ֔ל כִּי־שָׁ֛ם בָּלַ֥ל יְהֹוָ֖ה שְׂפַ֣ת כׇּל־הָאָ֑רֶץ וּמִשָּׁם֙ הֱפִיצָ֣ם יְהֹוָ֔ה עַל־פְּנֵ֖י כׇּל־הָאָֽרֶץ׃ \{פ\} \Rashi{\textbf{ומשם הפיצם.}למד שאין להם חלק לעולם הבא; וכי איזו קשה, של דור המבול או של דור הפלגה? אלו לא פשטו יד בעקר ואלו פשטו יד בעקר להלחם בו, ואלו נשטפו ואלו לא נאבדו מן העולם? אלא שדור המבול היו גזלנים והיתה מריבה ביניהם, לכך נאבדו; ואלו היו נוהגים אהבה ורעות ביניהם, שנ' שפה אחת ודברים אחדים, למדת ששנאוי המחלקת וגדול השלום:}%
% ---------- פרק 11 | פסוק 10 ----------
 \textbf{י} אֵ֚לֶּה תּוֹלְדֹ֣ת שֵׁ֔ם שֵׁ֚ם בֶּן־מְאַ֣ת שָׁנָ֔ה וַיּ֖וֹלֶד אֶת־אַרְפַּכְשָׁ֑ד שְׁנָתַ֖יִם אַחַ֥ר הַמַּבּֽוּל׃ \Rashi{\textbf{שם בן מאת שנה.}כשהוליד את ארפכשד שנתים אחר המבול:}%
% ---------- פרק 11 | פסוק 11 ----------
 \textbf{יא} וַֽיְחִי־שֵׁ֗ם אַֽחֲרֵי֙ הוֹלִיד֣וֹ אֶת־אַרְפַּכְשָׁ֔ד חֲמֵ֥שׁ מֵא֖וֹת שָׁנָ֑ה וַיּ֥וֹלֶד בָּנִ֖ים וּבָנֽוֹת׃ \{ס\} %
% ---------- פרק 11 | פסוק 12 ----------
 \textbf{יב} וְאַרְפַּכְשַׁ֣ד חַ֔י חָמֵ֥שׁ וּשְׁלֹשִׁ֖ים שָׁנָ֑ה וַיּ֖וֹלֶד אֶת־שָֽׁלַח׃ %
% ---------- פרק 11 | פסוק 13 ----------
 \textbf{יג} וַיְחִ֣י אַרְפַּכְשַׁ֗ד אַֽחֲרֵי֙ הוֹלִיד֣וֹ אֶת־שֶׁ֔לַח שָׁלֹ֣שׁ שָׁנִ֔ים וְאַרְבַּ֥ע מֵא֖וֹת שָׁנָ֑ה וַיּ֥וֹלֶד בָּנִ֖ים וּבָנֽוֹת׃ \{ס\} %
% ---------- פרק 11 | פסוק 14 ----------
 \textbf{יד} וְשֶׁ֥לַח חַ֖י שְׁלֹשִׁ֣ים שָׁנָ֑ה וַיּ֖וֹלֶד אֶת־עֵֽבֶר׃ %
% ---------- פרק 11 | פסוק 15 ----------
 \textbf{טו} וַֽיְחִי־שֶׁ֗לַח אַחֲרֵי֙ הוֹלִיד֣וֹ אֶת־עֵ֔בֶר שָׁלֹ֣שׁ שָׁנִ֔ים וְאַרְבַּ֥ע מֵא֖וֹת שָׁנָ֑ה וַיּ֥וֹלֶד בָּנִ֖ים וּבָנֽוֹת׃ \{ס\} %
% ---------- פרק 11 | פסוק 16 ----------
 \textbf{טז} וַֽיְחִי־עֵ֕בֶר אַרְבַּ֥ע וּשְׁלֹשִׁ֖ים שָׁנָ֑ה וַיּ֖וֹלֶד אֶת־פָּֽלֶג׃ %
% ---------- פרק 11 | פסוק 17 ----------
 \textbf{יז} וַֽיְחִי־עֵ֗בֶר אַחֲרֵי֙ הוֹלִיד֣וֹ אֶת־פֶּ֔לֶג שְׁלֹשִׁ֣ים שָׁנָ֔ה וְאַרְבַּ֥ע מֵא֖וֹת שָׁנָ֑ה וַיּ֥וֹלֶד בָּנִ֖ים וּבָנֽוֹת׃ \{ס\} %
% ---------- פרק 11 | פסוק 18 ----------
 \textbf{יח} וַֽיְחִי־פֶ֖לֶג שְׁלֹשִׁ֣ים שָׁנָ֑ה וַיּ֖וֹלֶד אֶת־רְעֽוּ׃ %
% ---------- פרק 11 | פסוק 19 ----------
 \textbf{יט} וַֽיְחִי־פֶ֗לֶג אַחֲרֵי֙ הוֹלִיד֣וֹ אֶת־רְע֔וּ תֵּ֥שַׁע שָׁנִ֖ים וּמָאתַ֣יִם שָׁנָ֑ה וַיּ֥וֹלֶד בָּנִ֖ים וּבָנֽוֹת׃ \{ס\} %
% ---------- פרק 11 | פסוק 20 ----------
 \textbf{כ} וַיְחִ֣י רְע֔וּ שְׁתַּ֥יִם וּשְׁלֹשִׁ֖ים שָׁנָ֑ה וַיּ֖וֹלֶד אֶת־שְׂרֽוּג׃ %
% ---------- פרק 11 | פסוק 21 ----------
 \textbf{כא} וַיְחִ֣י רְע֗וּ אַחֲרֵי֙ הוֹלִיד֣וֹ אֶת־שְׂר֔וּג שֶׁ֥בַע שָׁנִ֖ים וּמָאתַ֣יִם שָׁנָ֑ה וַיּ֥וֹלֶד בָּנִ֖ים וּבָנֽוֹת׃ \{ס\} %
% ---------- פרק 11 | פסוק 22 ----------
 \textbf{כב} וַיְחִ֥י שְׂר֖וּג שְׁלֹשִׁ֣ים שָׁנָ֑ה וַיּ֖וֹלֶד אֶת־נָחֽוֹר׃ %
% ---------- פרק 11 | פסוק 23 ----------
 \textbf{כג} וַיְחִ֣י שְׂר֗וּג אַחֲרֵ֛י הוֹלִיד֥וֹ אֶת־נָח֖וֹר מָאתַ֣יִם שָׁנָ֑ה וַיּ֥וֹלֶד בָּנִ֖ים וּבָנֽוֹת׃ \{ס\} %
% ---------- פרק 11 | פסוק 24 ----------
 \textbf{כד} וַיְחִ֣י נָח֔וֹר תֵּ֥שַׁע וְעֶשְׂרִ֖ים שָׁנָ֑ה וַיּ֖וֹלֶד אֶת־תָּֽרַח׃ %
% ---------- פרק 11 | פסוק 25 ----------
 \textbf{כה} וַיְחִ֣י נָח֗וֹר אַחֲרֵי֙ הוֹלִיד֣וֹ אֶת־תֶּ֔רַח תְּשַֽׁע־עֶשְׂרֵ֥ה שָׁנָ֖ה וּמְאַ֣ת שָׁנָ֑ה וַיּ֥וֹלֶד בָּנִ֖ים וּבָנֽוֹת׃ \{ס\} %
% ---------- פרק 11 | פסוק 26 ----------
 \textbf{כו} וַֽיְחִי־תֶ֖רַח שִׁבְעִ֣ים שָׁנָ֑ה וַיּ֙וֹלֶד֙ אֶת־אַבְרָ֔ם אֶת־נָח֖וֹר וְאֶת־הָרָֽן׃ %
% ---------- פרק 11 | פסוק 27 ----------
 \textbf{כז} וְאֵ֙לֶּה֙ תּוֹלְדֹ֣ת תֶּ֔רַח תֶּ֚רַח הוֹלִ֣יד אֶת־אַבְרָ֔ם אֶת־נָח֖וֹר וְאֶת־הָרָ֑ן וְהָרָ֖ן הוֹלִ֥יד אֶת־לֽוֹט׃ %
% ---------- פרק 11 | פסוק 28 ----------
 \textbf{כח} וַיָּ֣מׇת הָרָ֔ן עַל־פְּנֵ֖י תֶּ֣רַח אָבִ֑יו בְּאֶ֥רֶץ מוֹלַדְתּ֖וֹ בְּא֥וּר כַּשְׂדִּֽים׃ \Rashi{\textbf{על פני תרח אביו.}בחיי אביו. ומדרש אגדה אומר שעל ידי אביו מת; שקבל תרח על אברם בנו לפני נמרוד על שכתת את צלמיו, והשליכו לכבשן האש, והרן יושב ואומר בלבו אם אברם נוצח אני משלו, ואם נמרוד נוצח אני משלו; וכשנצל אברם, אמרו לו להרן משל מי אתה? אמר להם הרן משל אברם אני, השליכוהו לכבשן האש ונשרף, וזהו אור כשדים (ב"ר). ומנחם פרש אור – בקעה; וכן בארים כבדו ה' (ישעיה כ"ד), וכן מאורת צפעוני (שם י"א) כל חר ובקע עמק קרוי אור:}%
% ---------- פרק 11 | פסוק 29 ----------
 \textbf{כט} וַיִּקַּ֨ח אַבְרָ֧ם וְנָח֛וֹר לָהֶ֖ם נָשִׁ֑ים שֵׁ֤ם אֵֽשֶׁת־אַבְרָם֙ שָׂרָ֔י וְשֵׁ֤ם אֵֽשֶׁת־נָחוֹר֙ מִלְכָּ֔ה בַּת־הָרָ֥ן אֲבִֽי־מִלְכָּ֖ה וַֽאֲבִ֥י יִסְכָּֽה׃ \Rashi{\textbf{יסכה.}זו שרה, על שם שסוכה ברוח הקדש ושהכל סוכין ביפיה; ועוד יסכה לשון נסיכות, כמו שרה לשון שררה:}%
% ---------- פרק 11 | פסוק 30 ----------
 \textbf{ל} וַתְּהִ֥י שָׂרַ֖י עֲקָרָ֑ה אֵ֥ין לָ֖הּ וָלָֽד׃ %
% ---------- פרק 11 | פסוק 31 ----------
 \textbf{לא} וַיִּקַּ֨ח תֶּ֜רַח אֶת־אַבְרָ֣ם בְּנ֗וֹ וְאֶת־ל֤וֹט בֶּן־הָרָן֙ בֶּן־בְּנ֔וֹ וְאֵת֙ שָׂרַ֣י כַּלָּת֔וֹ אֵ֖שֶׁת אַבְרָ֣ם בְּנ֑וֹ וַיֵּצְא֨וּ אִתָּ֜ם מֵא֣וּר כַּשְׂדִּ֗ים לָלֶ֙כֶת֙ אַ֣רְצָה כְּנַ֔עַן וַיָּבֹ֥אוּ עַד־חָרָ֖ן וַיֵּ֥שְׁבוּ שָֽׁם׃ \Rashi{\textbf{ויצאו אתם.}ויצאו תרח ואברם עם לוט ושרי:}%
% ---------- פרק 11 | פסוק 32 ----------
 \textbf{לב} וַיִּהְי֣וּ יְמֵי־תֶ֔רַח חָמֵ֥שׁ שָׁנִ֖ים וּמָאתַ֣יִם שָׁנָ֑ה וַיָּ֥מׇת תֶּ֖רַח בְּחָרָֽן׃ \{פ\} \Rashi{\textbf{וימת תרח בחרן.}לאחר שיצא אברם מחרן ובא לארץ כנען, והיה שם יותר מששים שנה; שהרי כתוב ואברם בן חמש שנים ושבעים שנה בצאתו מחרן ותרח בן שבעים שנה כשנולד אברם, הרי קמ"ה לתרח; כשיצא אברם מחרן עדין נשארו משנותיו הרבה, ולמה הקדים הכתוב מיתתו של תרח ליציאתו של אברם? שלא יהא הדבר מפרסם לכל ויאמרו לא קים אברם את כבוד אביו, שהניחו זקן והלך לו, לפיכך קראו הכתוב מת, שהרשעים אף בחייהם קרויים מתים, והצדיקים אף במיתתן קרויים חיים, שנאמר ובניהו בן יהוידע בן איש חי (שמואל ב' כ"ג): (בחרן הנו"ן הפוכה לומר לך, עד אברם חרון אף של מקום):}%
\addparasha{חומש בראשית | פרשת לך לך}
% ---------- פרק 12 | פסוק 1 ----------
 \textbf{א} וַיֹּ֤אמֶר יְהֹוָה֙ אֶל־אַבְרָ֔ם לֶךְ־לְךָ֛ מֵאַרְצְךָ֥ וּמִמּֽוֹלַדְתְּךָ֖ וּמִבֵּ֣ית אָבִ֑יךָ אֶל־הָאָ֖רֶץ אֲשֶׁ֥ר אַרְאֶֽךָּ׃ \Rashi{\textbf{לך לך.}להנאתך ולטובתך, שם אעשך לגוי גדול, כאן אי אתה זוכה לבנים, ועוד שאודיע טבעך בעולם:}%
% ---------- פרק 12 | פסוק 2 ----------
 \textbf{ב} וְאֶֽעֶשְׂךָ֙ לְג֣וֹי גָּד֔וֹל וַאֲבָ֣רֶכְךָ֔ וַאֲגַדְּלָ֖ה שְׁמֶ֑ךָ וֶהְיֵ֖ה בְּרָכָֽה׃ \Rashi{\textbf{ואעשך לגוי גדול.}לפי שהדרך גורמת לשלשה דברים, ממעטת פריה ורביה וממעטת את הממון וממעטת את השם, לכך הזקק לשלשה ברכות הללו, שהבטיחו על הבנים ועל הממון ועל השם: ואברכך. בממון (בראשית רבה): והיה ברכה. הברכות נתונות בידך; עד עכשו היו בידי, ברכתי לאדם ונח, ומעכשו אתה תברך את אשר תחפץ (בראשית רבה). דבר אחר ואעשך לגוי גדול זה שאומרים אלהי אברהם, ואברכך זה שאומרים אלהי יצחק, ואגדלה שמך זה שאומרים אלהי יעקב. יכול יהיו חותמין בכלן, תלמוד לומר והיה ברכה, בך חותמין ולא בהם: מארצך. והלא כבר יצא משם עם אביו ובא עד חרן? אלא כך אמר לו התרחק עוד משם וצא מבית אביך: אשר אראך. לא גלה לו הארץ מיד כדי לחבבה בעיניו, ולתת לו שכר על כל דבור ודבור; כיוצא בו, את בנך את יחידך אשר אהבת את יצחק (ברא' כ"ב), כיוצא בו על אחד ההרים אשר אמר אליך, וכיו"ב וקרא אליה את הקריאה אשר אנכי דבר אליך (יונה ג'):}%
% ---------- פרק 12 | פסוק 3 ----------
 \textbf{ג} וַאֲבָֽרְכָה֙ מְבָ֣רְכֶ֔יךָ וּמְקַלֶּלְךָ֖ אָאֹ֑ר וְנִבְרְכ֣וּ בְךָ֔ כֹּ֖ל מִשְׁפְּחֹ֥ת הָאֲדָמָֽה׃ \Rashi{\textbf{ונברכו בך.}יש אגדות רבות, וזהו פשוטו, אדם אומר לבנו תהא כאברהם, וכן כל ונברכו בך שבמקרא, וזה מוכיח בך יברך ישראל לאמר ישמך אלהים כאפרים וכמנשה (בר' מ"ח):}%
% ---------- פרק 12 | פסוק 4 ----------
 \textbf{ד} וַיֵּ֣לֶךְ אַבְרָ֗ם כַּאֲשֶׁ֨ר דִּבֶּ֤ר אֵלָיו֙ יְהֹוָ֔ה וַיֵּ֥לֶךְ אִתּ֖וֹ ל֑וֹט וְאַבְרָ֗ם בֶּן־חָמֵ֤שׁ שָׁנִים֙ וְשִׁבְעִ֣ים שָׁנָ֔ה בְּצֵאת֖וֹ מֵחָרָֽן׃ %
% ---------- פרק 12 | פסוק 5 ----------
 \textbf{ה} וַיִּקַּ֣ח אַבְרָם֩ אֶת־שָׂרַ֨י אִשְׁתּ֜וֹ וְאֶת־ל֣וֹט בֶּן־אָחִ֗יו וְאֶת־כׇּל־רְכוּשָׁם֙ אֲשֶׁ֣ר רָכָ֔שׁוּ וְאֶת־הַנֶּ֖פֶשׁ אֲשֶׁר־עָשׂ֣וּ בְחָרָ֑ן וַיֵּצְא֗וּ לָלֶ֙כֶת֙ אַ֣רְצָה כְּנַ֔עַן וַיָּבֹ֖אוּ אַ֥רְצָה כְּנָֽעַן׃ \Rashi{\textbf{אשר עשו בחרן.}שהכניסן תחת כנפי השכינה; אברהם מגיר את האנשים ושרה מגירת הנשים, ומעלה עליהם הכתוב כאלו עשאום; ופשוטו של מקרא עבדים ושפחות שקנו להם, כמו עשה את כל הכבד הזה (שם ל"א), וישראל עשה חיל (במדבר כד יח), לשון קונה וכונס:}%
% ---------- פרק 12 | פסוק 6 ----------
 \textbf{ו} וַיַּעֲבֹ֤ר אַבְרָם֙ בָּאָ֔רֶץ עַ֚ד מְק֣וֹם שְׁכֶ֔ם עַ֖ד אֵל֣וֹן מוֹרֶ֑ה וְהַֽכְּנַעֲנִ֖י אָ֥ז בָּאָֽרֶץ׃ \Rashi{\textbf{ויעבר אברם בארץ.}נכנס לתוכה: עד מקום שכם. להתפלל על בני יעקב, כשיבאו להלחם בשכם: אלון מורה. הוא שכם. הראהו הר גריזים והר עיבל, ששם קבלו ישראל שבועת התורה: והכנעני אז בארץ. היה הולך וכובש את ארץ ישראל מזרעו של שם, שבחלקו של שם נפלה כשחלק נח את הארץ לבניו, שנאמר ומלכי צדק מלך שלם (ברא' י"ב), לפיכך ויאמר ה' אל אברם לזרעך אתן את הארץ הזאת, עתיד אני להחזירה לבניך שהם מזרעו של שם:}%
% ---------- פרק 12 | פסוק 7 ----------
 \textbf{ז} וַיֵּרָ֤א יְהֹוָה֙ אֶל־אַבְרָ֔ם וַיֹּ֕אמֶר לְזַ֨רְעֲךָ֔ אֶתֵּ֖ן אֶת־הָאָ֣רֶץ הַזֹּ֑את וַיִּ֤בֶן שָׁם֙ מִזְבֵּ֔חַ לַיהֹוָ֖ה הַנִּרְאֶ֥ה אֵלָֽיו׃ \Rashi{\textbf{ויבן שם מזבח.}על בשורת הזרע ועל בשורת ארץ ישראל:}%
% ---------- פרק 12 | פסוק 8 ----------
 \textbf{ח} וַיַּעְתֵּ֨ק מִשָּׁ֜ם הָהָ֗רָה מִקֶּ֛דֶם לְבֵֽית־אֵ֖ל וַיֵּ֣ט אׇהֳלֹ֑ה בֵּֽית־אֵ֤ל מִיָּם֙ וְהָעַ֣י מִקֶּ֔דֶם וַיִּֽבֶן־שָׁ֤ם מִזְבֵּ֙חַ֙ לַֽיהֹוָ֔ה וַיִּקְרָ֖א בְּשֵׁ֥ם יְהֹוָֽה׃ \Rashi{\textbf{ויעתק משם.}אהלו: מקדם לבית אל. במזרחה של בית אל נמצאת בית אל במערבו, הוא שנאמר בית אל מים: אהלו. אהלה כתיב, בתחלה נטה את אהל אשתו ואחר כך את שלו (בראשית רבה): ויבן שם מזבח. נתנבא שעתידין בניו להכשל שם על עון עכן, והתפלל שם עליהם:}%
% ---------- פרק 12 | פסוק 9 ----------
 \textbf{ט} וַיִּסַּ֣ע אַבְרָ֔ם הָל֥וֹךְ וְנָס֖וֹעַ הַנֶּֽגְבָּה׃ \{פ\} \Rashi{\textbf{הלוך ונסוע.}לפרקים יושב כאן חדש או יותר, ונוסע משם ונוטה אהלו במקום אחר, וכל מסעיו הנגבה, ללכת לדרומה של ארץ ישראל והיא לצד ירושלים, שהוא בחלקו של יהודה, שנטלו בדרומה של ארץ ישראל להר המוריה שהיא נחלתו (בראשית רבה):}%
% ---------- פרק 12 | פסוק 10 ----------
 \textbf{י} וַיְהִ֥י רָעָ֖ב בָּאָ֑רֶץ וַיֵּ֨רֶד אַבְרָ֤ם מִצְרַ֙יְמָה֙ לָג֣וּר שָׁ֔ם כִּֽי־כָבֵ֥ד הָרָעָ֖ב בָּאָֽרֶץ׃ \Rashi{\textbf{רעב בארץ.}באותה הארץ לבדה, לנסותו, אם יהרהר אחר דבריו של הקב"ה שאמר ללכת אל ארץ כנען, ועכשו משיאו לצאת ממנה:}%
% ---------- פרק 12 | פסוק 11 ----------
 \textbf{יא} וַיְהִ֕י כַּאֲשֶׁ֥ר הִקְרִ֖יב לָב֣וֹא מִצְרָ֑יְמָה וַיֹּ֙אמֶר֙ אֶל־שָׂרַ֣י אִשְׁתּ֔וֹ הִנֵּה־נָ֣א יָדַ֔עְתִּי כִּ֛י אִשָּׁ֥ה יְפַת־מַרְאֶ֖ה אָֽתְּ׃ \Rashi{\textbf{הנה נא ידעתי.}מדרש אגדה עד עכשו לא הכיר בה מתוך צניעות שבשניהם, ועכשו הכיר בה על ידי מעשה. דבר אחר, מנהג העולם שעל ידי טרח הדרך אדם מתבזה, וזאת עמדה ביפיה. ופשוטו של מקרא הנה נא הגיע השעה שיש לדאג על יפיך, ידעתי זה ימים רבים כי יפת מראה את, ועכשו אנו באים בין אנשים שחרים ומכערים, אחיהם של כושים, ולא הרגלו באשה יפה; ודומה לו הנה נא אדני סורו נא (בר' י"ט):}%
% ---------- פרק 12 | פסוק 12 ----------
 \textbf{יב} וְהָיָ֗ה כִּֽי־יִרְא֤וּ אֹתָךְ֙ הַמִּצְרִ֔ים וְאָמְר֖וּ אִשְׁתּ֣וֹ זֹ֑את וְהָרְג֥וּ אֹתִ֖י וְאֹתָ֥ךְ יְחַיּֽוּ׃ %
% ---------- פרק 12 | פסוק 13 ----------
 \textbf{יג} אִמְרִי־נָ֖א אֲחֹ֣תִי אָ֑תְּ לְמַ֙עַן֙ יִֽיטַב־לִ֣י בַעֲבוּרֵ֔ךְ וְחָיְתָ֥ה נַפְשִׁ֖י בִּגְלָלֵֽךְ׃ \Rashi{\textbf{למען ייטב לי בעבורך.}יתנו לי מתנות:}%
% ---------- פרק 12 | פסוק 14 ----------
 \textbf{יד} וַיְהִ֕י כְּב֥וֹא אַבְרָ֖ם מִצְרָ֑יְמָה וַיִּרְא֤וּ הַמִּצְרִים֙ אֶת־הָ֣אִשָּׁ֔ה כִּֽי־יָפָ֥ה הִ֖וא מְאֹֽד׃ \Rashi{\textbf{ויהי כבוא אברם מצרימה.}היה לו לומר כבואם מצרימה? אלא למד שהטמין אותה בתבה, ועל ידי שתבעו את המכס פתחו וראו אותה:}%
% ---------- פרק 12 | פסוק 15 ----------
 \textbf{טו} וַיִּרְא֤וּ אֹתָהּ֙ שָׂרֵ֣י פַרְעֹ֔ה וַיְהַֽלְל֥וּ אֹתָ֖הּ אֶל־פַּרְעֹ֑ה וַתֻּקַּ֥ח הָאִשָּׁ֖ה בֵּ֥ית פַּרְעֹֽה׃ \Rashi{\textbf{ויהללו אותה אל פרעה.}הללוה ביניהם לומר הגונה זו למלך:}%
% ---------- פרק 12 | פסוק 16 ----------
 \textbf{טז} וּלְאַבְרָ֥ם הֵיטִ֖יב בַּעֲבוּרָ֑הּ וַֽיְהִי־ל֤וֹ צֹאן־וּבָקָר֙ וַחֲמֹרִ֔ים וַעֲבָדִים֙ וּשְׁפָחֹ֔ת וַאֲתֹנֹ֖ת וּגְמַלִּֽים׃ \Rashi{\textbf{ולאברם היטיב.}פרעה בעבורה:}%
% ---------- פרק 12 | פסוק 17 ----------
 \textbf{יז} וַיְנַגַּ֨ע יְהֹוָ֧ה ׀ אֶת־פַּרְעֹ֛ה נְגָעִ֥ים גְּדֹלִ֖ים וְאֶת־בֵּית֑וֹ עַל־דְּבַ֥ר שָׂרַ֖י אֵ֥שֶׁת אַבְרָֽם׃ \Rashi{\textbf{וינגע ה' וגו'.}במכת ראתן לקה, שהתשמיש קשה לו (בראשית רבה): (ואת ביתו כתרגומו ועל אנש ביתה ומדרשו לרבות כתליו ועמודיו וכליו ברש"י ישן): על דבר שרי. על פי דבורה; אומרת למלאך הך, והוא מכה;}%
% ---------- פרק 12 | פסוק 18 ----------
 \textbf{יח} וַיִּקְרָ֤א פַרְעֹה֙ לְאַבְרָ֔ם וַיֹּ֕אמֶר מַה־זֹּ֖את עָשִׂ֣יתָ לִּ֑י לָ֚מָּה לֹא־הִגַּ֣דְתָּ לִּ֔י כִּ֥י אִשְׁתְּךָ֖ הִֽוא׃ %
% ---------- פרק 12 | פסוק 19 ----------
 \textbf{יט} לָמָ֤ה אָמַ֙רְתָּ֙ אֲחֹ֣תִי הִ֔וא וָאֶקַּ֥ח אֹתָ֛הּ לִ֖י לְאִשָּׁ֑ה וְעַתָּ֕ה הִנֵּ֥ה אִשְׁתְּךָ֖ קַ֥ח וָלֵֽךְ׃ \Rashi{\textbf{קח ולך.}לא כאבימלך שאמר לו הנה ארצי לפניך, אלא אמר לו לך ואל תעמד, שהמצרים שטופי זמה הם, שנאמר וזרמת סוסים זרמתם: (יחזקאל כ"ג):}%
% ---------- פרק 12 | פסוק 20 ----------
 \textbf{כ} וַיְצַ֥ו עָלָ֛יו פַּרְעֹ֖ה אֲנָשִׁ֑ים וַֽיְשַׁלְּח֥וּ אֹת֛וֹ וְאֶת־אִשְׁתּ֖וֹ וְאֶת־כׇּל־אֲשֶׁר־לֽוֹ׃ \Rashi{\textbf{ויצו עליו.}על אודותיו, לשלחו ולשמרו: וישלחו. כתרגומו ואלויאו:}%
% ---------- פרק 13 | פסוק 1 ----------
 \textbf{א} וַיַּ֩עַל֩ אַבְרָ֨ם מִמִּצְרַ֜יִם ה֠וּא וְאִשְׁתּ֧וֹ וְכׇל־אֲשֶׁר־ל֛וֹ וְל֥וֹט עִמּ֖וֹ הַנֶּֽגְבָּה׃ \Rashi{\textbf{ויעל אברם וגו' הנגבה.}לבא לדרומה של ארץ ישראל, כמו שאמר למעלה הלוך ונסוע הנגבה, להר המוריה; ומכל מקום כשהוא הולך ממצרים לארץ כנען, מדרום לצפון הוא מהלך, שארץ מצרים בדרומה של ארץ ישראל כמו שמוכיח במסעות ובגבולי הארץ:}%
% ---------- פרק 13 | פסוק 2 ----------
 \textbf{ב} וְאַבְרָ֖ם כָּבֵ֣ד מְאֹ֑ד בַּמִּקְנֶ֕ה בַּכֶּ֖סֶף וּבַזָּהָֽב׃ \Rashi{\textbf{כבד מאד.}טעון משאות:}%
% ---------- פרק 13 | פסוק 3 ----------
 \textbf{ג} וַיֵּ֙לֶךְ֙ לְמַסָּעָ֔יו מִנֶּ֖גֶב וְעַד־בֵּֽית־אֵ֑ל עַד־הַמָּק֗וֹם אֲשֶׁר־הָ֨יָה שָׁ֤ם אׇֽהֳלֹה֙ בַּתְּחִלָּ֔ה בֵּ֥ין בֵּֽית־אֵ֖ל וּבֵ֥ין הָעָֽי׃ \Rashi{\textbf{וילך למסעיו.}כשחזר ממצרים לארץ כנען היה הולך ולן באכסניות שלן בהם בהליכתו למצרים, למדך דרך ארץ, שלא ישנה אדם מאכסניא שלו. דבר אחר בחזרתו פרע הקפותיו (בראשית רבה): מנגב. ארץ מצרים בדרומה של ארץ כנען:}%
% ---------- פרק 13 | פסוק 4 ----------
 \textbf{ד} אֶל־מְקוֹם֙ הַמִּזְבֵּ֔חַ אֲשֶׁר־עָ֥שָׂה שָׁ֖ם בָּרִאשֹׁנָ֑ה וַיִּקְרָ֥א שָׁ֛ם אַבְרָ֖ם בְּשֵׁ֥ם יְהֹוָֽה׃ \Rashi{\textbf{אשר עשה שם בראשנה.}ואשר קרא שם אברם בשם ה' וגם יש לומר ויקרא שם עכשו בשם ה':}%
% ---------- פרק 13 | פסוק 5 ----------
 \textbf{ה} וְגַ֨ם־לְל֔וֹט הַהֹלֵ֖ךְ אֶת־אַבְרָ֑ם הָיָ֥ה צֹאן־וּבָקָ֖ר וְאֹהָלִֽים׃ \Rashi{\textbf{ההלך את אברם.}מי גרם שהיה לו זאת? הליכתו עם אברם:}%
% ---------- פרק 13 | פסוק 6 ----------
 \textbf{ו} וְלֹא־נָשָׂ֥א אֹתָ֛ם הָאָ֖רֶץ לָשֶׁ֣בֶת יַחְדָּ֑ו כִּֽי־הָיָ֤ה רְכוּשָׁם֙ רָ֔ב וְלֹ֥א יָֽכְל֖וּ לָשֶׁ֥בֶת יַחְדָּֽו׃ \Rashi{\textbf{ולא נשא אתם.}לא היתה יכולה להספיק מרעה למקניהם, ולשון קצר הוא, וצריך להוסיף עליו, כמו ולא נשא אותם מרעה הארץ, לפיכך כתב ולא נשא בלשון זכר:}%
% ---------- פרק 13 | פסוק 7 ----------
 \textbf{ז} וַֽיְהִי־רִ֗יב בֵּ֚ין רֹעֵ֣י מִקְנֵֽה־אַבְרָ֔ם וּבֵ֖ין רֹעֵ֣י מִקְנֵה־ל֑וֹט וְהַֽכְּנַעֲנִי֙ וְהַפְּרִזִּ֔י אָ֖ז יֹשֵׁ֥ב בָּאָֽרֶץ׃ \Rashi{\textbf{ויהי ריב.}לפי שהיו רועים של לוט רשעים ומרעים בהמתם בשדות אחרים, ורועי אברם מוכיחים אותם על הגזל, והם אומרים נתנה הארץ לאברם, ולו אין יורש, ולוט יורשו, ואין זה גזל, והכתוב אומר והכנעני והפרזי אז ישב בארץ, ולא זכה בה אברם עדין (ב"):}%
% ---------- פרק 13 | פסוק 8 ----------
 \textbf{ח} וַיֹּ֨אמֶר אַבְרָ֜ם אֶל־ל֗וֹט אַל־נָ֨א תְהִ֤י מְרִיבָה֙ בֵּינִ֣י וּבֵינֶ֔ךָ וּבֵ֥ין רֹעַ֖י וּבֵ֣ין רֹעֶ֑יךָ כִּֽי־אֲנָשִׁ֥ים אַחִ֖ים אֲנָֽחְנוּ׃ \Rashi{\textbf{אנשים אחים.}קרובים, ומדרש אגדה דומין בקלסתר פנים:}%
% ---------- פרק 13 | פסוק 9 ----------
 \textbf{ט} הֲלֹ֤א כׇל־הָאָ֙רֶץ֙ לְפָנֶ֔יךָ הִפָּ֥רֶד נָ֖א מֵעָלָ֑י אִם־הַשְּׂמֹ֣אל וְאֵימִ֔נָה וְאִם־הַיָּמִ֖ין וְאַשְׂמְאִֽילָה׃ \Rashi{\textbf{אם השמאל ואימנה.}בכל אשר תשב, לא אתרחק ממך ואעמד לך למגן ולעזר, וסוף דבר הצרך לו, שנאמר וישמע אברם כי נשבה אחיו וגו': ואימנה. אימין את עצמי, כמו ואשמאילה אשמאיל את עצמי. ואם תאמר היה לו לנקד ואימינה, כך מצינו במקום אחר אם יש להמין (שמואל ב י"ד) ואין נקוד להימין:}%
% ---------- פרק 13 | פסוק 10 ----------
 \textbf{י} וַיִּשָּׂא־ל֣וֹט אֶת־עֵינָ֗יו וַיַּרְא֙ אֶת־כׇּל־כִּכַּ֣ר הַיַּרְדֵּ֔ן כִּ֥י כֻלָּ֖הּ מַשְׁקֶ֑ה לִפְנֵ֣י ׀ שַׁחֵ֣ת יְהֹוָ֗ה אֶת־סְדֹם֙ וְאֶת־עֲמֹרָ֔ה כְּגַן־יְהֹוָה֙ כְּאֶ֣רֶץ מִצְרַ֔יִם בֹּאֲכָ֖ה צֹֽעַר׃ \Rashi{\textbf{כי כלה משקה.}ארץ נחלי מים: לפני שחת ה' את סדום ואת עמרה. היה אותו מישור: כגן ה'. לאילנות: כארץ מצרים. לזרעים (בראשית רבה): באכה צער. עד צער. ומדרש אגדה דורשה לגנאי, על שהיו שטופי זמה בחר לו לוט בשכונתם:}%
% ---------- פרק 13 | פסוק 11 ----------
 \textbf{יא} וַיִּבְחַר־ל֣וֹ ל֗וֹט אֵ֚ת כׇּל־כִּכַּ֣ר הַיַּרְדֵּ֔ן וַיִּסַּ֥ע ל֖וֹט מִקֶּ֑דֶם וַיִּפָּ֣רְד֔וּ אִ֖ישׁ מֵעַ֥ל אָחִֽיו׃ \Rashi{\textbf{ככר.}מישור, כתרגומו: מקדם. נסע מאצל אברם והלך לו למערבו של אברם; נמצא נוסע ממזרח למערב. ומדרש אגדה הסיע עצמו מקדמונו של עולם; אמר אי אפשי לא באברם ולא באלהיו:}%
% ---------- פרק 13 | פסוק 12 ----------
 \textbf{יב} אַבְרָ֖ם יָשַׁ֣ב בְּאֶֽרֶץ־כְּנָ֑עַן וְל֗וֹט יָשַׁב֙ בְּעָרֵ֣י הַכִּכָּ֔ר וַיֶּאֱהַ֖ל עַד־סְדֹֽם׃ \Rashi{\textbf{ויאהל.}נטה אהלים לרועיו ולמקנהו עד סדום:}%
% ---------- פרק 13 | פסוק 13 ----------
 \textbf{יג} וְאַנְשֵׁ֣י סְדֹ֔ם רָעִ֖ים וְחַטָּאִ֑ים לַיהֹוָ֖ה מְאֹֽד׃ \Rashi{\textbf{ואנשי סדום רעים.}ואף על פי כן לא נמנע לוט מלשכן עמהם. ורבותינו למדו מכאן (יומא ל"ח) ושם רשעים ירקב (משלי י'): רעים. בגופם: וחטאים. בממונם (סנה' ק"ט): לה' מאד. יודעים רבונם, ומתכונים למרד בו:}%
% ---------- פרק 13 | פסוק 14 ----------
 \textbf{יד} וַֽיהֹוָ֞ה אָמַ֣ר אֶל־אַבְרָ֗ם אַחֲרֵי֙ הִפָּֽרֶד־ל֣וֹט מֵֽעִמּ֔וֹ שָׂ֣א נָ֤א עֵינֶ֙יךָ֙ וּרְאֵ֔ה מִן־הַמָּק֖וֹם אֲשֶׁר־אַתָּ֣ה שָׁ֑ם צָפֹ֥נָה וָנֶ֖גְבָּה וָקֵ֥דְמָה וָיָֽמָּה׃ \Rashi{\textbf{אחרי הפרד לוט.}כל זמן שהרשע עמו היה הדבור פורש ממנו:}%
% ---------- פרק 13 | פסוק 15 ----------
 \textbf{טו} כִּ֧י אֶת־כׇּל־הָאָ֛רֶץ אֲשֶׁר־אַתָּ֥ה רֹאֶ֖ה לְךָ֣ אֶתְּנֶ֑נָּה וּֽלְזַרְעֲךָ֖ עַד־עוֹלָֽם׃ %
% ---------- פרק 13 | פסוק 16 ----------
 \textbf{טז} וְשַׂמְתִּ֥י אֶֽת־זַרְעֲךָ֖ כַּעֲפַ֣ר הָאָ֑רֶץ אֲשֶׁ֣ר ׀ אִם־יוּכַ֣ל אִ֗ישׁ לִמְנוֹת֙ אֶת־עֲפַ֣ר הָאָ֔רֶץ גַּֽם־זַרְעֲךָ֖ יִמָּנֶֽה׃ \Rashi{\textbf{אשר אם יוכל איש.}כשם שאי אפשר לעפר להמנות כך זרעך לא ימנה:}%
% ---------- פרק 13 | פסוק 17 ----------
 \textbf{יז} ק֚וּם הִתְהַלֵּ֣ךְ בָּאָ֔רֶץ לְאׇרְכָּ֖הּ וּלְרׇחְבָּ֑הּ כִּ֥י לְךָ֖ אֶתְּנֶֽנָּה׃ %
% ---------- פרק 13 | פסוק 18 ----------
 \textbf{יח} וַיֶּאֱהַ֣ל אַבְרָ֗ם וַיָּבֹ֛א וַיֵּ֛שֶׁב בְּאֵלֹנֵ֥י מַמְרֵ֖א אֲשֶׁ֣ר בְּחֶבְר֑וֹן וַיִּֽבֶן־שָׁ֥ם מִזְבֵּ֖חַ לַֽיהֹוָֽה׃ \{פ\} \Rashi{\textbf{ממרא.}שם אדם:}%
% ---------- פרק 14 | פסוק 1 ----------
 \textbf{א} וַיְהִ֗י בִּימֵי֙ אַמְרָפֶ֣ל מֶֽלֶךְ־שִׁנְעָ֔ר אַרְי֖וֹךְ מֶ֣לֶךְ אֶלָּסָ֑ר כְּדׇרְלָעֹ֙מֶר֙ מֶ֣לֶךְ עֵילָ֔ם וְתִדְעָ֖ל מֶ֥לֶךְ גּוֹיִֽם׃ \Rashi{\textbf{אמרפל.}הוא נמרוד, שאמר לאברהם פל לתוך כבשן האש (בראשית רבה): מלך גוים. מקום יש ששמו גוים, על שם שנתקבצו שמה מכמה אמות ומקומות והמליכו איש עליהם ושמו תדעל:}%
% ---------- פרק 14 | פסוק 2 ----------
 \textbf{ב} עָשׂ֣וּ מִלְחָמָ֗ה אֶת־בֶּ֙רַע֙ מֶ֣לֶךְ סְדֹ֔ם וְאֶת־בִּרְשַׁ֖ע מֶ֣לֶךְ עֲמֹרָ֑ה שִׁנְאָ֣ב ׀ מֶ֣לֶךְ אַדְמָ֗ה וְשֶׁמְאֵ֙בֶר֙ מֶ֣לֶךְ (צביים) [צְבוֹיִ֔ם] וּמֶ֥לֶךְ בֶּ֖לַע הִיא־צֹֽעַר׃ \Rashi{\textbf{ברע.}רע לשמים ורע לבריות: ברשע. שנתעלה ברשע: שנאב. שונא אביו שבשמים: שמאבר. שם אבר לעוף ולקפץ ולמרד בהקב"ה: בלע. שם העיר:}%
% ---------- פרק 14 | פסוק 3 ----------
 \textbf{ג} כׇּל־אֵ֙לֶּה֙ חָֽבְר֔וּ אֶל־עֵ֖מֶק הַשִּׂדִּ֑ים ה֖וּא יָ֥ם הַמֶּֽלַח׃ \Rashi{\textbf{עמק השדים.}כך שמו, על שם שהיו בו שדות הרבה: הוא ים המלח. לאחר זמן נמשך הים לתוכו ונעשה ים המלח. ומדרש אגדה אומר שנתבקעו הצורים סביבותיו ונמשכו יאורים לתוכו:}%
% ---------- פרק 14 | פסוק 4 ----------
 \textbf{ד} שְׁתֵּ֤ים עֶשְׂרֵה֙ שָׁנָ֔ה עָבְד֖וּ אֶת־כְּדׇרְלָעֹ֑מֶר וּשְׁלֹשׁ־עֶשְׂרֵ֥ה שָׁנָ֖ה מָרָֽדוּ׃ \Rashi{\textbf{שתים עשרה שנה עבדו.}חמשה מלכים הללו את כדרלעמר:}%
% ---------- פרק 14 | פסוק 5 ----------
 \textbf{ה} וּבְאַרְבַּע֩ עֶשְׂרֵ֨ה שָׁנָ֜ה בָּ֣א כְדׇרְלָעֹ֗מֶר וְהַמְּלָכִים֙ אֲשֶׁ֣ר אִתּ֔וֹ וַיַּכּ֤וּ אֶת־רְפָאִים֙ בְּעַשְׁתְּרֹ֣ת קַרְנַ֔יִם וְאֶת־הַזּוּזִ֖ים בְּהָ֑ם וְאֵת֙ הָֽאֵימִ֔ים בְּשָׁוֵ֖ה קִרְיָתָֽיִם׃ \Rashi{\textbf{ובארבע עשרה שנה.}למרדן: בא כדרלעמר. לפי שהוא היה בעל המעשה נכנס בעבי הקורה: והמלכים. אלה שלשה מלכים: הזוזים. הם זמזמים:}%
% ---------- פרק 14 | פסוק 6 ----------
 \textbf{ו} וְאֶת־הַחֹרִ֖י בְּהַרְרָ֣ם שֵׂעִ֑יר עַ֚ד אֵ֣יל פָּארָ֔ן אֲשֶׁ֖ר עַל־הַמִּדְבָּֽר׃ \Rashi{\textbf{בהררם.}בהר שלהם: איל פארן. כתרגומו מישור. ואומר אני שאין איל לשון מישור, אלא מישור של פארן איל שמו, ושל ממרא אלוני שמו, ושל ירדן ככר שמו, ושל שטים אבל שמו, אבל השטים, וכן בעל גד בעל שמו; וכלם מתרגמין מישור, וכל א' שמו עליו: על המדבר. אצל המדבר, כמו ועליו מטה מנשה (במדבר ב'):}%
% ---------- פרק 14 | פסוק 7 ----------
 \textbf{ז} וַ֠יָּשֻׁ֠בוּ וַיָּבֹ֜אוּ אֶל־עֵ֤ין מִשְׁפָּט֙ הִ֣וא קָדֵ֔שׁ וַיַּכּ֕וּ אֶֽת־כׇּל־שְׂדֵ֖ה הָעֲמָלֵקִ֑י וְגַם֙ אֶת־הָ֣אֱמֹרִ֔י הַיֹּשֵׁ֖ב בְּחַֽצְצֹ֥ן תָּמָֽר׃ \Rashi{\textbf{עין משפט היא קדש.}על שם העתיד, שעתידין משה ואהרן להשפט שם על עסקי אותו העין והם מי מריבה. ואנקלוס תרגמו כפשוטו מקום שהיו בני המדינה מתקבצים שם לכל משפט: שדה העמלקי. עדין לא נולד עמלק, ונקרא על שם העתיד: בחצצון תמר. הוא עין גדי, מקרא מלא בדברי הימים (ב' כ') ביהושפט:}%
% ---------- פרק 14 | פסוק 8 ----------
 \textbf{ח} וַיֵּצֵ֨א מֶֽלֶךְ־סְדֹ֜ם וּמֶ֣לֶךְ עֲמֹרָ֗ה וּמֶ֤לֶךְ אַדְמָה֙ וּמֶ֣לֶךְ (צביים) [צְבוֹיִ֔ם] וּמֶ֥לֶךְ בֶּ֖לַע הִוא־צֹ֑עַר וַיַּֽעַרְכ֤וּ אִתָּם֙ מִלְחָמָ֔ה בְּעֵ֖מֶק הַשִּׂדִּֽים׃ %
% ---------- פרק 14 | פסוק 9 ----------
 \textbf{ט} אֵ֣ת כְּדׇרְלָעֹ֜מֶר מֶ֣לֶךְ עֵילָ֗ם וְתִדְעָל֙ מֶ֣לֶךְ גּוֹיִ֔ם וְאַמְרָפֶל֙ מֶ֣לֶךְ שִׁנְעָ֔ר וְאַרְי֖וֹךְ מֶ֣לֶךְ אֶלָּסָ֑ר אַרְבָּעָ֥ה מְלָכִ֖ים אֶת־הַחֲמִשָּֽׁה׃ \Rashi{\textbf{ארבעה מלכים וגו'.}ואף על פי כן נצחו המעטים, להודיעך שגבורים היו, ואף על פי כן לא נמנע אברהם מלרדף אחריהם:}%
% ---------- פרק 14 | פסוק 10 ----------
 \textbf{י} וְעֵ֣מֶק הַשִּׂדִּ֗ים בֶּֽאֱרֹ֤ת בֶּאֱרֹת֙ חֵמָ֔ר וַיָּנֻ֛סוּ מֶֽלֶךְ־סְדֹ֥ם וַעֲמֹרָ֖ה וַיִּפְּלוּ־שָׁ֑מָּה וְהַנִּשְׁאָרִ֖ים הֶ֥רָה נָּֽסוּ׃ \Rashi{\textbf{בארת בארת חמר.}בארות הרבה היו שם, שנוטלין משם אדמה לטיט של בנין. ומדרש אגדה שהיה הטיט מגבל בהם, ונעשה נס למלך סדום שיצא משם, לפי שהיו באמות מקצתן שלא היו מאמינין שנצל אברהם מאור כשדים מכבשן האש, וכיון שיצא זה מן החמר, האמינו באברם למפרע: הרה נסו. להר נסו, הרה, כמו להר; כל תבה שצריכה למ"ד בתחלתה, הטיל לה ה"א בסופה. ויש חלוק בין הרה לההרה, שה"א שבסוף התבה עומדת במקום למ"ד שבראשה, אבל אינה עומדת במקום למ"ד ונקודה פתח תחתיה, והרי הרה כמו להר או כמו אל הר, ואינו מפרש לאיזה הר, אלא שכל אחד נס באשר מצא הר תחלה, וכשהוא נותן ה"א בראשה לכתב ההרה או המדברה, פתרונו כמו אל ההר או כמו לההר, ומשמע לאותו הר הידוע ומפרש בפרשה:}%
% ---------- פרק 14 | פסוק 11 ----------
 \textbf{יא} וַ֠יִּקְח֠וּ אֶת־כׇּל־רְכֻ֨שׁ סְדֹ֧ם וַעֲמֹרָ֛ה וְאֶת־כׇּל־אׇכְלָ֖ם וַיֵּלֵֽכוּ׃ %
% ---------- פרק 14 | פסוק 12 ----------
 \textbf{יב} וַיִּקְח֨וּ אֶת־ל֧וֹט וְאֶת־רְכֻשׁ֛וֹ בֶּן־אֲחִ֥י אַבְרָ֖ם וַיֵּלֵ֑כוּ וְה֥וּא יֹשֵׁ֖ב בִּסְדֹֽם׃ \Rashi{\textbf{והוא ישב בסדום.}מי גרם לו זאת? ישיבתו בסדום:}%
% ---------- פרק 14 | פסוק 13 ----------
 \textbf{יג} וַיָּבֹא֙ הַפָּלִ֔יט וַיַּגֵּ֖ד לְאַבְרָ֣ם הָעִבְרִ֑י וְהוּא֩ שֹׁכֵ֨ן בְּאֵֽלֹנֵ֜י מַמְרֵ֣א הָאֱמֹרִ֗י אֲחִ֤י אֶשְׁכֹּל֙ וַאֲחִ֣י עָנֵ֔ר וְהֵ֖ם בַּעֲלֵ֥י בְרִית־אַבְרָֽם׃ \Rashi{\textbf{ויבא הפליט.}לפי פשוטו זה עוג, שפלט מן המלחמה, והוא שכתוב כי רק עוג נשאר מיתר הרפאים (דבר' ג') וזהו נשאר, שלא הרגוהו אמרפל וחבריו כשהכו הרפאים בעשתרות קרנים, תנחומא ומדרש ב"ר זה עוג, שפלט מדור המבול; וזהו מיתר הרפאים, שנאמר הנפלים היו בארץ וגו' (לעיל ו') ומתכון שיהרג אברם וישא את שרה: העברי. שבא מעבר הנהר (בראשית רבה): בעלי ברית אברם. שכרתו עמו ברית:}%
% ---------- פרק 14 | פסוק 14 ----------
 \textbf{יד} וַיִּשְׁמַ֣ע אַבְרָ֔ם כִּ֥י נִשְׁבָּ֖ה אָחִ֑יו וַיָּ֨רֶק אֶת־חֲנִיכָ֜יו יְלִידֵ֣י בֵית֗וֹ שְׁמֹנָ֤ה עָשָׂר֙ וּשְׁלֹ֣שׁ מֵא֔וֹת וַיִּרְדֹּ֖ף עַד־דָּֽן׃ \Rashi{\textbf{וירק.}כתרגומו וזריז, וכן והריקתי אחריכם חרב (ויקרא כ"ו) – אזדין בחרבי עליכם, וכן אריק חרבי (שמות ט"ו), וכן והרק חנית וסגר (תה' ל"ה): חניכיו. חנכו כתיב (ס"א קרי), זה אליעזר שחנכו למצות והוא לשון התחלת כניסת האדם או כלי לאמנות שהוא עתיד לעמד בה, וכן חנך לנער (משלי כ"ב), חנכת המזבח (במד' ז'), חנכת הבית (תה' ל') ובלע"ז קורין לו אינצ"נייר: שמנה עשר וגו'. רבותינו אמרו אליעזר לבדו היה (נדרים ל"ב), והוא מנין גימטריא של שמו: עד דן. שם תשש כחו, שראה שעתידין בניו להעמיד שם עגל (סנה' צ"ו):}%
% ---------- פרק 14 | פסוק 15 ----------
 \textbf{טו} וַיֵּחָלֵ֨ק עֲלֵיהֶ֧ם ׀ לַ֛יְלָה ה֥וּא וַעֲבָדָ֖יו וַיַּכֵּ֑ם וַֽיִּרְדְּפֵם֙ עַד־חוֹבָ֔ה אֲשֶׁ֥ר מִשְּׂמֹ֖אל לְדַמָּֽשֶׂק׃ \Rashi{\textbf{ויחלק עליהם.}לפי פשוטו סרס המקרא "ויחלק הוא ועבדיו עליהם לילה" כדרך הרודפים שמתפלגים אחר הנרדפים, כשבורחין זה לכאן וזה לכאן: לילה. כלומר אחר שחשכה לא נמנע מלרדפם. ומדרש אגדה שנחלק הלילה, ובחציו הראשון נעשה לו נס, וחציו השני נשמר ובא לו לחצות לילה של מצרים: עד חובה. אין מקום ששמו חובה, ודן קורא חובה, על שם עבודה זרה שעתידה להיות שם:}%
% ---------- פרק 14 | פסוק 16 ----------
 \textbf{טז} וַיָּ֕שֶׁב אֵ֖ת כׇּל־הָרְכֻ֑שׁ וְגַם֩ אֶת־ל֨וֹט אָחִ֤יו וּרְכֻשׁוֹ֙ הֵשִׁ֔יב וְגַ֥ם אֶת־הַנָּשִׁ֖ים וְאֶת־הָעָֽם׃ %
% ---------- פרק 14 | פסוק 17 ----------
 \textbf{יז} וַיֵּצֵ֣א מֶֽלֶךְ־סְדֹם֮ לִקְרָאתוֹ֒ אַחֲרֵ֣י שׁוּב֗וֹ מֵֽהַכּוֹת֙ אֶת־כְּדׇרְלָעֹ֔מֶר וְאֶת־הַמְּלָכִ֖ים אֲשֶׁ֣ר אִתּ֑וֹ אֶל־עֵ֣מֶק שָׁוֵ֔ה ה֖וּא עֵ֥מֶק הַמֶּֽלֶךְ׃ \Rashi{\textbf{עמק שוה.}כך שמו, כתרגומו למישר מפנא – פנוי מאילנות ומכל מכשול: עמק המלך. בית ריסא דמלכא בית ריס א' שהוא שלשים קנים, שהיה מיחד למלך לצחק שם. ומדרש אגדה עמק, שהשוו שם כל האמות והמליכו את אברהם עליהם לנשיא אלהים ולקצין:}%
% ---------- פרק 14 | פסוק 18 ----------
 \textbf{יח} וּמַלְכִּי־צֶ֙דֶק֙ מֶ֣לֶךְ שָׁלֵ֔ם הוֹצִ֖יא לֶ֣חֶם וָיָ֑יִן וְה֥וּא כֹהֵ֖ן לְאֵ֥ל עֶלְיֽוֹן׃ \Rashi{\textbf{ומלכי צדק.}מ"א הוא שם בן נח (נדרים ל"ב): לחם ויין. כך עושים ליגיעי מלחמה, והראה לו שאין בלבו עליו על שהרג את בניו. ומ"א רמז לו על המנחות ועל הנסכים שיקריבו שם בניו:}%
% ---------- פרק 14 | פסוק 19 ----------
 \textbf{יט} וַֽיְבָרְכֵ֖הוּ וַיֹּאמַ֑ר בָּר֤וּךְ אַבְרָם֙ לְאֵ֣ל עֶלְי֔וֹן קֹנֵ֖ה שָׁמַ֥יִם וָאָֽרֶץ׃ \Rashi{\textbf{קנה שמים וארץ.}כמו עשה שמים וארץ (תה' קל"ד); על ידי עשיתן קנאן להיות שלו:}%
% ---------- פרק 14 | פסוק 20 ----------
 \textbf{כ} וּבָרוּךְ֙ אֵ֣ל עֶלְי֔וֹן אֲשֶׁר־מִגֵּ֥ן צָרֶ֖יךָ בְּיָדֶ֑ךָ וַיִּתֶּן־ל֥וֹ מַעֲשֵׂ֖ר מִכֹּֽל׃ \Rashi{\textbf{אשר מגן.}אשר הסגיר, וכן אמגנך ישראל (הושע י"א): ויתן לו. אברהם מעשר מכל אשר לו, לפי שהיה כהן:}%
% ---------- פרק 14 | פסוק 21 ----------
 \textbf{כא} וַיֹּ֥אמֶר מֶֽלֶךְ־סְדֹ֖ם אֶל־אַבְרָ֑ם תֶּן־לִ֣י הַנֶּ֔פֶשׁ וְהָרְכֻ֖שׁ קַֽח־לָֽךְ׃ \Rashi{\textbf{תן לי הנפש.}מן השבי שלי שהצלת, החזר לי הגופים לבדם:}%
% ---------- פרק 14 | פסוק 22 ----------
 \textbf{כב} וַיֹּ֥אמֶר אַבְרָ֖ם אֶל־מֶ֣לֶךְ סְדֹ֑ם הֲרִמֹ֨תִי יָדִ֤י אֶל־יְהֹוָה֙ אֵ֣ל עֶלְי֔וֹן קֹנֵ֖ה שָׁמַ֥יִם וָאָֽרֶץ׃ \Rashi{\textbf{הרימתי ידי.}לשון שבועה, מרים אני את ידי לאל עליון, וכן בי נשבעתי (בר' כ"ב), נשבע אני, וכן נתתי כסף השדה קח ממני (שם כ"ג), נותן אני לך כסף השדה, קחהו ממני:}%
% ---------- פרק 14 | פסוק 23 ----------
 \textbf{כג} אִם־מִחוּט֙ וְעַ֣ד שְׂרֽוֹךְ־נַ֔עַל וְאִם־אֶקַּ֖ח מִכׇּל־אֲשֶׁר־לָ֑ךְ וְלֹ֣א תֹאמַ֔ר אֲנִ֖י הֶעֱשַׁ֥רְתִּי אֶת־אַבְרָֽם׃ \Rashi{\textbf{אם מחוט ועד שרוך נעל.}אעכב לעצמי מן השבי: ואם אקח מכל אשר לך. ואם תאמר לתת לי שכר מבית גנזיך, לא אקח: ולא תאמר וגו'. שהקדוש ברוך הוא הבטיחני לעשרני, שנאמר ואברכך וגו':}%
% ---------- פרק 14 | פסוק 24 ----------
 \textbf{כד} בִּלְעָדַ֗י רַ֚ק אֲשֶׁ֣ר אָֽכְל֣וּ הַנְּעָרִ֔ים וְחֵ֙לֶק֙ הָֽאֲנָשִׁ֔ים אֲשֶׁ֥ר הָלְכ֖וּ אִתִּ֑י עָנֵר֙ אֶשְׁכֹּ֣ל וּמַמְרֵ֔א הֵ֖ם יִקְח֥וּ חֶלְקָֽם׃ \{ס\} \Rashi{\textbf{הנערים.}עבדי אשר הלכו אתי ועוד ענר אשכל וממרא וגו'; אף על פי שעבדי נכנסו למלחמה, שנאמר הוא ועבדיו ויכם, וענר וחבריו ישבו על הכלים לשמר, אפלו הכי הם יקחו חלקם. וממנו למד דוד, שאמר כחלק הירד במלחמה וכחלק הישב על הכלים יחדו יחלקו (שמואל א ל'), ולכך נאמר ויהי מהיום ההוא ומעלה וישמה לחק ולמשפט (שם), ולא נאמר והלאה, לפי שכבר נתן החק בימי אברם:}%
% ---------- פרק 15 | פסוק 1 ----------
 \textbf{א} אַחַ֣ר ׀ הַדְּבָרִ֣ים הָאֵ֗לֶּה הָיָ֤ה דְבַר־יְהֹוָה֙ אֶל־אַבְרָ֔ם בַּֽמַּחֲזֶ֖ה לֵאמֹ֑ר אַל־תִּירָ֣א אַבְרָ֗ם אָנֹכִי֙ מָגֵ֣ן לָ֔ךְ שְׂכָרְךָ֖ הַרְבֵּ֥ה מְאֹֽד׃ \Rashi{\textbf{אחר הדברים האלה.}כל מקום שנאמר אחר – סמוך, אחרי – מפלג (בראשית רבה). אחר הדברים האלה אחר שנעשה לו נס זה, שהרג את המלכים והיה דואג ואומר, שמא קבלתי שכר על כל צדקותי, לכך אמר לו המקום אל תירא אברם אנכי מגן לך, מן הענש, שלא תענש על כל אותן נפשות שהרגת; ומה שאתה דואג על קבול שכרך, שכרך הרבה מאד (בראשית רבה):}%
% ---------- פרק 15 | פסוק 2 ----------
 \textbf{ב} וַיֹּ֣אמֶר אַבְרָ֗ם אֲדֹנָ֤י יֱהֹוִה֙ מַה־תִּתֶּן־לִ֔י וְאָנֹכִ֖י הוֹלֵ֣ךְ עֲרִירִ֑י וּבֶן־מֶ֣שֶׁק בֵּיתִ֔י ה֖וּא דַּמֶּ֥שֶׂק אֱלִיעֶֽזֶר׃ \Rashi{\textbf{הולך ערירי.}מנחם בן סרוק פרשו לשון יורש, וחבר לו ער וענה (מלאכי ב'), ערירי, בלא יורש, כאשר תאמר ובכל תבואתי תשרש (איוב ל"א), תעקר שרשיה, כך לשון ערירי חסר בנים, ובלע"ז דישאנפנטיש. ולי נראה ער וענה מגזרת ולבי ער (שיר השירים ה'), וערירי לשון חרבן, וכן ערו ערו (תה' קל"ז), וכן ערות יסוד (חבקוק ג'), וכן ערער תתערער (ירמ' נ"א), וכן כי ארזה ערה (צפניה ב'): ובן משק ביתי. כתרגומו, שכל ביתי נזון על פיו, כמו ועל פיך ישק (בר' מלכים א), אפוטרופוס שלי, ואלו היה לי בן, היה בני ממנה על שלי: דמשק. לפי התרגום מדמשק היה, ולפי מ"א שרדף המלכים עד דמשק. ובתלמוד דרשו נוטריקון דולה ומשקה מתורת רבו לאחרים:}%
% ---------- פרק 15 | פסוק 3 ----------
 \textbf{ג} וַיֹּ֣אמֶר אַבְרָ֔ם הֵ֣ן לִ֔י לֹ֥א נָתַ֖תָּה זָ֑רַע וְהִנֵּ֥ה בֶן־בֵּיתִ֖י יוֹרֵ֥שׁ אֹתִֽי׃ \Rashi{\textbf{הן לי לא נתת זרע.}ומה תועלת בכל אשר תתן לי?}%
% ---------- פרק 15 | פסוק 4 ----------
 \textbf{ד} וְהִנֵּ֨ה דְבַר־יְהֹוָ֤ה אֵלָיו֙ לֵאמֹ֔ר לֹ֥א יִֽירָשְׁךָ֖ זֶ֑ה כִּי־אִם֙ אֲשֶׁ֣ר יֵצֵ֣א מִמֵּעֶ֔יךָ ה֖וּא יִֽירָשֶֽׁךָ׃ %
% ---------- פרק 15 | פסוק 5 ----------
 \textbf{ה} וַיּוֹצֵ֨א אֹת֜וֹ הַח֗וּצָה וַיֹּ֙אמֶר֙ הַבֶּט־נָ֣א הַשָּׁמַ֗יְמָה וּסְפֹר֙ הַכּ֣וֹכָבִ֔ים אִם־תּוּכַ֖ל לִסְפֹּ֣ר אֹתָ֑ם וַיֹּ֣אמֶר ל֔וֹ כֹּ֥ה יִהְיֶ֖ה זַרְעֶֽךָ׃ \Rashi{\textbf{ויוצא אתו החוצה.}לפי פשוטו הוציאו מאהלו לחוץ לראות הכוכבים, ולפי מדרשו אמר לו צא מאצטגנינות שלך שראית במזלות שאינך עתיד להעמיד בן, אברם אין לו בן, אבל אברהם יש לו בן, שרי לא תלד, אבל שרה תלד; אני קורא לכם שם אחר וישתנה המזל. ד"א הוציאו מחללו של עולם והגביהו למעלה מן הכוכבים, וזהו לשון הבטה מלמעלה למטה:}%
% ---------- פרק 15 | פסוק 6 ----------
 \textbf{ו} וְהֶאֱמִ֖ן בַּֽיהֹוָ֑ה וַיַּחְשְׁבֶ֥הָ לּ֖וֹ צְדָקָֽה׃ \Rashi{\textbf{והאמין בה'.}לא שאל לו אות על זאת; אבל על ירשת הארץ שאל לו אות ואמר לו במה אדע. ויחשבה לו צדקה. הקב"ה חשבה לאברם לזכות ולצדקה על האמנה שהאמין בו. דבר אחר במה אדע, לא שאל לו אות אלא אמר לפניו, הודיעני באיזה זכות יתקימו בה, אמר לו הקב"ה, בזכות הקרבנות:}%
% ---------- פרק 15 | פסוק 7 ----------
 \textbf{ז} וַיֹּ֖אמֶר אֵלָ֑יו אֲנִ֣י יְהֹוָ֗ה אֲשֶׁ֤ר הוֹצֵאתִ֙יךָ֙ מֵא֣וּר כַּשְׂדִּ֔ים לָ֧תֶת לְךָ֛ אֶת־הָאָ֥רֶץ הַזֹּ֖את לְרִשְׁתָּֽהּ׃ %
% ---------- פרק 15 | פסוק 8 ----------
 \textbf{ח} וַיֹּאמַ֑ר אֲדֹנָ֣י יֱהֹוִ֔ה בַּמָּ֥ה אֵדַ֖ע כִּ֥י אִֽירָשֶֽׁנָּה׃ %
% ---------- פרק 15 | פסוק 9 ----------
 \textbf{ט} וַיֹּ֣אמֶר אֵלָ֗יו קְחָ֥ה לִי֙ עֶגְלָ֣ה מְשֻׁלֶּ֔שֶׁת וְעֵ֥ז מְשֻׁלֶּ֖שֶׁת וְאַ֣יִל מְשֻׁלָּ֑שׁ וְתֹ֖ר וְגוֹזָֽל׃ \Rashi{\textbf{עגלה משלשת.}שלשה עגלים; רמז לשלשה פרים, פר יום הכפורים ופר העלם דבר של צבור ועגלה ערופה: ועז משלשת. רמז לשעיר הנעשה בפנים ושעירי מוספין של מועד ושעיר חטאת יחיד: ואיל משלש. אשם ודאי ואשם תלוי וכבשה של חטאת יחיד: ותור וגוזל. תור ובן יונה (בראשית רבה):}%
% ---------- פרק 15 | פסוק 10 ----------
 \textbf{י} וַיִּֽקַּֽח־ל֣וֹ אֶת־כׇּל־אֵ֗לֶּה וַיְבַתֵּ֤ר אֹתָם֙ בַּתָּ֔וֶךְ וַיִּתֵּ֥ן אִישׁ־בִּתְר֖וֹ לִקְרַ֣את רֵעֵ֑הוּ וְאֶת־הַצִּפֹּ֖ר לֹ֥א בָתָֽר׃ \Rashi{\textbf{ויבתר אתם.}חלק כל אחד לשני חלקים; ואין המקרא יוצא מידי פשוטו, לפי שהיה כורת ברית עמו לשמר הבטחתו להוריש לבניו את הארץ כדכתיב ביום ההוא כרת ה' את אברם ברית לאמר וגומר, ודרך כורתי ברית לחלק בהמה ולעבר בין בתריה, כמה שנאמר להלן העברים בין בתרי העגל (ירמיה ל"ד), אף כאן תנור עשן ולפיד אש אשר עבר בין הגזרים הוא שלוחו של שכינה, שהוא אש. ואת הצפור לא בתר. לפי שהאמות נמשלו לפרים ואילים ושעירים, שנאמר סבבוני פרים רבים וגו' (תהל' כ"ב), ואומר האיל אשר ראית בעל הקרנים מלכי מדי ופרס (דניאל ח'), ואומר והצפיר השעיר מלך יון (שם), וישראל נמשלו לבני יונה, שנאמר יונתי בחגוי הסלע (שיר השירים ב'), לפיכך בתר הבהמות, רמז שיהיו האמות כלים והולכים, ואת הצפור לא בתר, רמז שיהיו ישראל קימין לעולם:}%
% ---------- פרק 15 | פסוק 11 ----------
 \textbf{יא} וַיֵּ֥רֶד הָעַ֖יִט עַל־הַפְּגָרִ֑ים וַיַּשֵּׁ֥ב אֹתָ֖ם אַבְרָֽם׃ \Rashi{\textbf{העיט.}הוא עוף, על שם שהוא עט ושואף אל הנבלות לטוש עלי אכל, כמו ותעט אל השלל (שמואל א ט"ו): על הפגרים. על הבתרים. הפגרים מתרגמינן פגליא, אלא מתוך שהרגלו לתרגם "איש בתרו" ויהב פלגיא, נתחלף להם תבת פגליא לפלגיא, ותרגמו הפגרים פלגיא; וכל המתרגם כן טועה, לפי שאין להקיש בתרים לפגרים, שבתרים תרגומו פלגיא, ופגרים תרגומו פגליא, לשון פגול, כמו פגול הוא (ויקרא ז') ל' פגר: וישב. לשון נשיבה והפרחה, כמו ישב רוחו (תה' קמ"ז), רמז שיבא דוד בן ישי לכלותם, ואין מניחין אותו מן השמים עד שיבא מלך המשיח:}%
% ---------- פרק 15 | פסוק 12 ----------
 \textbf{יב} וַיְהִ֤י הַשֶּׁ֙מֶשׁ֙ לָב֔וֹא וְתַרְדֵּמָ֖ה נָפְלָ֣ה עַל־אַבְרָ֑ם וְהִנֵּ֥ה אֵימָ֛ה חֲשֵׁכָ֥ה גְדֹלָ֖ה נֹפֶ֥לֶת עָלָֽיו׃ \Rashi{\textbf{והנה אימה וגו'.}רמז לצרות וחשך של גליות:}%
% ---------- פרק 15 | פסוק 13 ----------
 \textbf{יג} וַיֹּ֣אמֶר לְאַבְרָ֗ם יָדֹ֨עַ תֵּדַ֜ע כִּי־גֵ֣ר ׀ יִהְיֶ֣ה זַרְעֲךָ֗ בְּאֶ֙רֶץ֙ לֹ֣א לָהֶ֔ם וַעֲבָד֖וּם וְעִנּ֣וּ אֹתָ֑ם אַרְבַּ֥ע מֵא֖וֹת שָׁנָֽה׃ \Rashi{\textbf{כי גר יהיה זרעך.}משנולד יצחק עד שיצאו ישראל ממצרים ארבע מאות שנה כיצד? יצחק בן ששים שנה כשנולד יעקב, ויעקב כשירד למצרים אמר ימי שני מגורי שלשים ומאת שנה, הרי ק"צ, ובמצרים היו מאתים ועשר כמנין רד"ו, הרי ארבע מאות שנה. ואם תאמר במצרים היו ארבע מאות, הרי קהת מיורדי מצרים היה, צא וחשב שנותיו של קהת ושל עמרם ושמונים של משה שהיה כשיצאו ישראל ממצרים, אין אתה מוצא אלא שלש מאות וחמשים, ואתה צריך להוציא מהן כל השנים שחי קהת אחר לדת עמרם ושחי עמרם אחר לדת משה: בארץ לא להם. לא נאמר בארץ מצרים, אלא לא להם – ומשנולד יצחק ויגר אברהם וגו' (בר' כ"א), וביצחק גור בארץ הזאת (שם כ"ו), ויעקב גר בארץ חם (תה' ק"ה), לגור בארץ באנו (בר' מ"ז):}%
% ---------- פרק 15 | פסוק 14 ----------
 \textbf{יד} וְגַ֧ם אֶת־הַגּ֛וֹי אֲשֶׁ֥ר יַעֲבֹ֖דוּ דָּ֣ן אָנֹ֑כִי וְאַחֲרֵי־כֵ֥ן יֵצְא֖וּ בִּרְכֻ֥שׁ גָּדֽוֹל׃ \Rashi{\textbf{וגם את הגוי.}וגם לרבות ד' מלכיות, שאף הן כלים על ששעבדו את ישראל: דן אנכי. בעשר מכות: ברכש גדול. בממון גדול, כמו שנאמר וינצלו את מצרים (שמות י"ב):}%
% ---------- פרק 15 | פסוק 15 ----------
 \textbf{טו} וְאַתָּ֛ה תָּב֥וֹא אֶל־אֲבֹתֶ֖יךָ בְּשָׁל֑וֹם תִּקָּבֵ֖ר בְּשֵׂיבָ֥ה טוֹבָֽה׃ \Rashi{\textbf{ואתה תבוא.}ולא תראה כל אלה: אל אבותיך. אביו עובד ע"ז והוא מבשרו שיבא אליו? למדך שעשה תרח תשובה: תקבר בשיבה טובה. בשרו שיעשה ישמעאל תשובה בימיו, ולא יצא עשו לתרבות רעה בימיו; ולפיכך מת ה' שנים קדם זמנו, ובו ביום מרד עשו:}%
% ---------- פרק 15 | פסוק 16 ----------
 \textbf{טז} וְד֥וֹר רְבִיעִ֖י יָשׁ֣וּבוּ הֵ֑נָּה כִּ֧י לֹא־שָׁלֵ֛ם עֲוֺ֥ן הָאֱמֹרִ֖י עַד־הֵֽנָּה׃ \Rashi{\textbf{ודור רביעי.}לאחר שיגלו למצרים יהיו שם ג' דורות, והרביעי ישובו לארץ הזאת, לפי שבארץ כנען היה מדבר עמו, וכרת ברית זו, כדכתיב: לתת לך את הארץ הזאת לרשתה. וכן היה: יעקב ירד למצרים, צא וחשב דורותיו: יהודה, פרץ, וחצרון, וכלב בן חצרון מבאי הארץ היה: כי לא שלם עון האמרי. להיות משתלח מארצו עד אותו זמן, שאין הקב"ה נפרע מאמה עד שתתמלא סאתה, שנאמר בסאסאה בשלחה תריבנה (יש' כ"ז):}%
% ---------- פרק 15 | פסוק 17 ----------
 \textbf{יז} וַיְהִ֤י הַשֶּׁ֙מֶשׁ֙ בָּ֔אָה וַעֲלָטָ֖ה הָיָ֑ה וְהִנֵּ֨ה תַנּ֤וּר עָשָׁן֙ וְלַפִּ֣יד אֵ֔שׁ אֲשֶׁ֣ר עָבַ֔ר בֵּ֖ין הַגְּזָרִ֥ים הָאֵֽלֶּה׃ \Rashi{\textbf{ויהי השמש באה.}כמו ויהי הם מריקים שקיהם (בראשית מ"ב), ויהי הם קברים איש (מלכים ב י"ג), כלומר, ויהי דבר זה: השמש באה. שקעה: ועלטה היה. חשך היום: והנה תנור עשן וגו'. רמז לו שיפלו המלכיות בגיהנם: באה. טעמו למעלה, לכך הוא מבאר שבא כבר. ואם היה טעמו למטה, באל"ף, היה מבאר כשהיא שוקעת; וא"א לומר כן, שהרי כבר כתיב ויהי השמש לבוא, והעברת תנור עשן לאחר מכאן היתה, נמצא שכבר שקעה; וזה חלוק בכל תבה לשון נקבה שיסודה שתי אותיות, כמו בא, קם, שב: כשהטעם למעלה, לשון עבר הוא, כגון זה, וכגון ורחל באה, קמה אלמתי, הנה שבה יבמתך, וכשהטעם למטה הוא לשון הווה, דבר שנעשה עכשו והולך, כמו באה עם הצאן, בערב היא באה ובבקר היא שבה:}%
% ---------- פרק 15 | פסוק 18 ----------
 \textbf{יח} בַּיּ֣וֹם הַה֗וּא כָּרַ֧ת יְהֹוָ֛ה אֶת־אַבְרָ֖ם בְּרִ֣ית לֵאמֹ֑ר לְזַרְעֲךָ֗ נָתַ֙תִּי֙ אֶת־הָאָ֣רֶץ הַזֹּ֔את מִנְּהַ֣ר מִצְרַ֔יִם עַד־הַנָּהָ֥ר הַגָּדֹ֖ל נְהַר־פְּרָֽת׃ \Rashi{\textbf{לזרעך נתתי.}אמירתו של הקב"ה כאלו היא עשויה: הנהר הגדול. לפי שהוא דבוק לארץ ישראל קוראהו גדול, אע"פ שהוא מאחר בארבעה נהרות היוצאים מעדן, שנאמר והנהר הרביעי הוא פרת. משל הדיוט: עבד מלך מלך, הדבק לשחור וישתחוו לך (בראשית רבה):}%
% ---------- פרק 15 | פסוק 19 ----------
 \textbf{יט} אֶת־הַקֵּינִי֙ וְאֶת־הַקְּנִזִּ֔י וְאֵ֖ת הַקַּדְמֹנִֽי׃ \Rashi{\textbf{את הקני.}עשר אמות יש כאן ולא נתן להם אלא שבעה גוים, והשלשה אדום ומואב ועמון, והם קיני, קנזי, וקדמוני עתידים להיות ירשה לעתיד (בראשית רבה), שנאמר: אדום ומואב משלוח ידם ובני עמון משמעתם (ישעיה י"א):}%
% ---------- פרק 15 | פסוק 20 ----------
 \textbf{כ} וְאֶת־הַחִתִּ֥י וְאֶת־הַפְּרִזִּ֖י וְאֶת־הָרְפָאִֽים׃ \Rashi{\textbf{ואת הרפאים.}ארץ עוג, שנאמר בה: ההוא יקרא ארץ רפאים (דברים ג'):}%
% ---------- פרק 15 | פסוק 21 ----------
 \textbf{כא} וְאֶת־הָֽאֱמֹרִי֙ וְאֶת־הַֽכְּנַעֲנִ֔י וְאֶת־הַגִּרְגָּשִׁ֖י וְאֶת־הַיְבוּסִֽי׃ \{ס\} %
% ---------- פרק 16 | פסוק 1 ----------
 \textbf{א} וְשָׂרַי֙ אֵ֣שֶׁת אַבְרָ֔ם לֹ֥א יָלְדָ֖ה ל֑וֹ וְלָ֛הּ שִׁפְחָ֥ה מִצְרִ֖ית וּשְׁמָ֥הּ הָגָֽר׃ \Rashi{\textbf{שפחה מצרית.}בת פרעה היתה, כשראה נסים שנעשו לשרה אמר: מוטב שתהא בתי שפחה בבית זה ולא גבירה בבית אחר:}%
% ---------- פרק 16 | פסוק 2 ----------
 \textbf{ב} וַתֹּ֨אמֶר שָׂרַ֜י אֶל־אַבְרָ֗ם הִנֵּה־נָ֞א עֲצָרַ֤נִי יְהֹוָה֙ מִלֶּ֔דֶת בֹּא־נָא֙ אֶל־שִׁפְחָתִ֔י אוּלַ֥י אִבָּנֶ֖ה מִמֶּ֑נָּה וַיִּשְׁמַ֥ע אַבְרָ֖ם לְק֥וֹל שָׂרָֽי׃ \Rashi{\textbf{אולי אבנה ממנה.}למד על מי שאין לו בנים שאינו בנוי אלא הרוס: אבנה ממנה. בזכות שאכניס צרתי לתוך ביתי: לקול שרי. לרוח הקדש שבה (בראשית רבה):}%
% ---------- פרק 16 | פסוק 3 ----------
 \textbf{ג} וַתִּקַּ֞ח שָׂרַ֣י אֵֽשֶׁת־אַבְרָ֗ם אֶת־הָגָ֤ר הַמִּצְרִית֙ שִׁפְחָתָ֔הּ מִקֵּץ֙ עֶ֣שֶׂר שָׁנִ֔ים לְשֶׁ֥בֶת אַבְרָ֖ם בְּאֶ֣רֶץ כְּנָ֑עַן וַתִּתֵּ֥ן אֹתָ֛הּ לְאַבְרָ֥ם אִישָׁ֖הּ ל֥וֹ לְאִשָּֽׁה׃ \Rashi{\textbf{ותקח שרי.}לקחתה בדברים, אשריך שזכית לדבק בגוף קדוש כזה (בראשית רבה): מקץ עשר שנים. מועד הקבוע לאשה ששהתה עשר שנים ולא ילדה לבעלה, חיב לשא אחרת: לשבת אברם וגו'. מגיד שאין ישיבת חוצה לארץ עולה לו מן המנין (יב' ס"ד), לפי שלא נאמר לו ואעשך לגוי גדול, עד שיבא לארץ ישראל:}%
% ---------- פרק 16 | פסוק 4 ----------
 \textbf{ד} וַיָּבֹ֥א אֶל־הָגָ֖ר וַתַּ֑הַר וַתֵּ֙רֶא֙ כִּ֣י הָרָ֔תָה וַתֵּקַ֥ל גְּבִרְתָּ֖הּ בְּעֵינֶֽיהָ׃ \Rashi{\textbf{ויבא אל הגר ותהר.}מביאה ראשונה (בראשית רבה): ותקל גברתה בעיניה. אמרה שרי זו אין סתרה כגלויה מראה עצמה כאלו היא צדקת ואינה צדקת, שלא זכתה להריון כל השנים הללו, ואני נתעברתי מביאה ראשונה (בראשית רבה):}%
% ---------- פרק 16 | פסוק 5 ----------
 \textbf{ה} וַתֹּ֨אמֶר שָׂרַ֣י אֶל־אַבְרָם֮ חֲמָסִ֣י עָלֶ֒יךָ֒ אָנֹכִ֗י נָתַ֤תִּי שִׁפְחָתִי֙ בְּחֵיקֶ֔ךָ וַתֵּ֙רֶא֙ כִּ֣י הָרָ֔תָה וָאֵקַ֖ל בְּעֵינֶ֑יהָ יִשְׁפֹּ֥ט יְהֹוָ֖ה בֵּינִ֥י וּבֵינֶֽיׄךָ׃ \Rashi{\textbf{חמסי עליך.}חמס העשוי לי, עליך אני מטיל הענש; כשהתפללת להקב"ה מה תתן לי ואנכי הולך ערירי, לא התפללת אלא עליך, והיה לך להתפלל על שנינו, והייתי אני נפקדת עמך, ועוד דבריך אתה חומס ממני, שאתה שומע בזיוני ושותק (בראשית רבה): אנכי נתתי שפחתי וגו' ביני וביניך. כל ביניך שבמקרא חסר, וזה מלא, קרי ביה ובניך, שהכניסה עין הרע בעבורה של הגר והפילה עברה; הוא שהמלאך אומר להגר הנך הרה, והלא כבר הרתה, והוא מבשר לה שתהר? אלא מלמד שהפילה הריון הראשון:}%
% ---------- פרק 16 | פסוק 6 ----------
 \textbf{ו} וַיֹּ֨אמֶר אַבְרָ֜ם אֶל־שָׂרַ֗י הִנֵּ֤ה שִׁפְחָתֵךְ֙ בְּיָדֵ֔ךְ עֲשִׂי־לָ֖הּ הַטּ֣וֹב בְּעֵינָ֑יִךְ וַתְּעַנֶּ֣הָ שָׂרַ֔י וַתִּבְרַ֖ח מִפָּנֶֽיהָ׃ \Rashi{\textbf{ותענה שרי.}היתה משעבדת בה בקשי (בראשית רבה):}%
% ---------- פרק 16 | פסוק 7 ----------
 \textbf{ז} וַֽיִּמְצָאָ֞הּ מַלְאַ֧ךְ יְהֹוָ֛ה עַל־עֵ֥ין הַמַּ֖יִם בַּמִּדְבָּ֑ר עַל־הָעַ֖יִן בְּדֶ֥רֶךְ שֽׁוּר׃ %
% ---------- פרק 16 | פסוק 8 ----------
 \textbf{ח} וַיֹּאמַ֗ר הָגָ֞ר שִׁפְחַ֥ת שָׂרַ֛י אֵֽי־מִזֶּ֥ה בָ֖את וְאָ֣נָה תֵלֵ֑כִי וַתֹּ֕אמֶר מִפְּנֵי֙ שָׂרַ֣י גְּבִרְתִּ֔י אָנֹכִ֖י בֹּרַֽחַת׃ \Rashi{\textbf{מזה באת.}מהיכן באת? יודע היה, אלא לתן לה פתח לכנס עמה בדברים. ולשון אי מזה: איה המקום שתאמר עליו מזה אני באה?}%
% ---------- פרק 16 | פסוק 9 ----------
 \textbf{ט} וַיֹּ֤אמֶר לָהּ֙ מַלְאַ֣ךְ יְהֹוָ֔ה שׁ֖וּבִי אֶל־גְּבִרְתֵּ֑ךְ וְהִתְעַנִּ֖י תַּ֥חַת יָדֶֽיהָ׃ \Rashi{\textbf{ויאמר לה מלאך וגו'.}על כל אמירה היה שלוח לה מלאך אחר, לכך נאמר מלאך בכל אמירה ואמירה:}%
% ---------- פרק 16 | פסוק 10 ----------
 \textbf{י} וַיֹּ֤אמֶר לָהּ֙ מַלְאַ֣ךְ יְהֹוָ֔ה הַרְבָּ֥ה אַרְבֶּ֖ה אֶת־זַרְעֵ֑ךְ וְלֹ֥א יִסָּפֵ֖ר מֵרֹֽב׃ %
% ---------- פרק 16 | פסוק 11 ----------
 \textbf{יא} וַיֹּ֤אמֶר לָהּ֙ מַלְאַ֣ךְ יְהֹוָ֔ה הִנָּ֥ךְ הָרָ֖ה וְיֹלַ֣דְתְּ בֵּ֑ן וְקָרָ֤את שְׁמוֹ֙ יִשְׁמָעֵ֔אל כִּֽי־שָׁמַ֥ע יְהֹוָ֖ה אֶל־עׇנְיֵֽךְ׃ \Rashi{\textbf{הנך הרה.}כשתשובי תהרי, כמו הנך הרה דאשת מנוח: וילדת בן. כמו ויולדת, ודומה לו ישבת בלבנון מקננתי בארזים (ירמיה כ"ב): וקראת שמו. צווי הוא, כמו שאומר לזכר וקראת את שמו יצחק (בר' י"ז):}%
% ---------- פרק 16 | פסוק 12 ----------
 \textbf{יב} וְה֤וּא יִהְיֶה֙ פֶּ֣רֶא אָדָ֔ם יָד֣וֹ בַכֹּ֔ל וְיַ֥ד כֹּ֖ל בּ֑וֹ וְעַל־פְּנֵ֥י כׇל־אֶחָ֖יו יִשְׁכֹּֽן׃ \Rashi{\textbf{פרא אדם.}אוהב מדברות לצוד חיות, כמו שכתוב וישב במדבר ויהי רבה קשת (בראשית כ"א): ידו בכל. לסטים: ויד כל בו. הכל שונאין אותו ומתגרין בו: ועל פני כל אחיו ישכון. שיהיה זרעו גדול (בראשית רבה):}%
% ---------- פרק 16 | פסוק 13 ----------
 \textbf{יג} וַתִּקְרָ֤א שֵׁם־יְהֹוָה֙ הַדֹּבֵ֣ר אֵלֶ֔יהָ אַתָּ֖ה אֵ֣ל רֳאִ֑י כִּ֣י אָֽמְרָ֗ה הֲגַ֥ם הֲלֹ֛ם רָאִ֖יתִי אַחֲרֵ֥י רֹאִֽי׃ \Rashi{\textbf{אתה אל ראי.}נקוד חטף קמץ, מפני שהוא שם דבר, אלוה הראיה, שרואה בעלבון של עלובין: הגם הלם. לשון תמה; וכי סבורה הייתי שאף הלום במדברות ראיתי שלוחו של מקום אחרי ראי אותם בביתו של אברהם ששם הייתי רגילה לראות מלאכים? ותדע שהיתה רגילה לראותם, שהרי מנוח ראה את המלאך פעם אחת ואמר מות נמות (שופ' י"ג), וזו ראתה ד' זה אחר זה ולא חרדה:}%
% ---------- פרק 16 | פסוק 14 ----------
 \textbf{יד} עַל־כֵּן֙ קָרָ֣א לַבְּאֵ֔ר בְּאֵ֥ר לַחַ֖י רֹאִ֑י הִנֵּ֥ה בֵין־קָדֵ֖שׁ וּבֵ֥ין בָּֽרֶד׃ \Rashi{\textbf{באר לחי.}כתרגומו:}%
% ---------- פרק 16 | פסוק 15 ----------
 \textbf{טו} וַתֵּ֧לֶד הָגָ֛ר לְאַבְרָ֖ם בֵּ֑ן וַיִּקְרָ֨א אַבְרָ֧ם שֶׁם־בְּנ֛וֹ אֲשֶׁר־יָלְדָ֥ה הָגָ֖ר יִשְׁמָעֵֽאל׃ \Rashi{\textbf{ויקרא אברם שם וגו'.}אע"פ שלא שמע אברם דברי המלאך שאמר וקראת שמו ישמעאל, שרתה רוח הקדש עליו וקראו ישמעאל:}%
% ---------- פרק 16 | פסוק 16 ----------
 \textbf{טז} וְאַבְרָ֕ם בֶּן־שְׁמֹנִ֥ים שָׁנָ֖ה וְשֵׁ֣שׁ שָׁנִ֑ים בְּלֶֽדֶת־הָגָ֥ר אֶת־יִשְׁמָעֵ֖אל לְאַבְרָֽם׃ \{ס\} \Rashi{\textbf{ואברם בן שמונים וגו'.}לשבחו של ישמעאל נכתב, להודיע שהיה בן י"ג שנה כשנמול ולא עכב:}%
% ---------- פרק 17 | פסוק 1 ----------
 \textbf{א} וַיְהִ֣י אַבְרָ֔ם בֶּן־תִּשְׁעִ֥ים שָׁנָ֖ה וְתֵ֣שַׁע שָׁנִ֑ים וַיֵּרָ֨א יְהֹוָ֜ה אֶל־אַבְרָ֗ם וַיֹּ֤אמֶר אֵלָיו֙ אֲנִי־אֵ֣ל שַׁדַּ֔י הִתְהַלֵּ֥ךְ לְפָנַ֖י וֶהְיֵ֥ה תָמִֽים׃ \Rashi{\textbf{שדי.}אני הוא שיש די באלהותי לכל בריה, לפיכך התהלך לפני ואהיה לך לאלוה ולפטרון; וכן כל מקום שהוא במקרא פרושו כך: די שלו, והכל לפי הענין: התהלך לפני. כתרגומו פלח קדמי, הדבק בעבודתי: והיה תמים. אף זה צווי אחר צווי, היה שלם בכל נסיונותי; ולפי מדרשו התהלך לפני במצות מילה, ובדבר הזה תהיה תמים, שכל זמן שהערלה בך אתה בעל מום לפני. דבר אחר והיה תמים, (עכשו אתה חסר ה' אברים, ב' עינים, ב' אזנים, וראש הגויה) שאוסיף לך אות על שמך, ויהיו מנין אותיותיך רמ"ח כמנין אבריך:}%
% ---------- פרק 17 | פסוק 2 ----------
 \textbf{ב} וְאֶתְּנָ֥ה בְרִיתִ֖י בֵּינִ֣י וּבֵינֶ֑ךָ וְאַרְבֶּ֥ה אוֹתְךָ֖ בִּמְאֹ֥ד מְאֹֽד׃ \Rashi{\textbf{ואתנה בריתי.}ברית של אהבה וברית הארץ להורישה לך על ידי מצוה זו:}%
% ---------- פרק 17 | פסוק 3 ----------
 \textbf{ג} וַיִּפֹּ֥ל אַבְרָ֖ם עַל־פָּנָ֑יו וַיְדַבֵּ֥ר אִתּ֛וֹ אֱלֹהִ֖ים לֵאמֹֽר׃ \Rashi{\textbf{ויפל אברם על פניו.}ממורא השכינה; שעד שלא מל לא היה בו כח לעמד ורוח הקדש נצבת עליו, וזהו שנאמר בבלעם נפל וגלוי עינים, בבריתא דרבי אליעזר:}%
% ---------- פרק 17 | פסוק 4 ----------
 \textbf{ד} אֲנִ֕י הִנֵּ֥ה בְרִיתִ֖י אִתָּ֑ךְ וְהָיִ֕יתָ לְאַ֖ב הֲמ֥וֹן גּוֹיִֽם׃ %
% ---------- פרק 17 | פסוק 5 ----------
 \textbf{ה} וְלֹא־יִקָּרֵ֥א ע֛וֹד אֶת־שִׁמְךָ֖ אַבְרָ֑ם וְהָיָ֤ה שִׁמְךָ֙ אַבְרָהָ֔ם כִּ֛י אַב־הֲמ֥וֹן גּוֹיִ֖ם נְתַתִּֽיךָ׃ \Rashi{\textbf{כי אב המון גוים.}לשון נוטריקון של שמו; ורי"ש שהיתה בו בתחלה, שלא היה אב אלא לארם שהוא מקומו ועכשו אב לכל העולם, לא זזה ממקומה, שאף יו"ד של שרי נתרעמה על השכינה עד שנתוספה ליהושע, שנאמר ויקרא משה להושע בן נון יהושע (במדבר י"ג):}%
% ---------- פרק 17 | פסוק 6 ----------
 \textbf{ו} וְהִפְרֵתִ֤י אֹֽתְךָ֙ בִּמְאֹ֣ד מְאֹ֔ד וּנְתַתִּ֖יךָ לְגוֹיִ֑ם וּמְלָכִ֖ים מִמְּךָ֥ יֵצֵֽאוּ׃ \Rashi{\textbf{ונתתיך לגוים.}ישראל ואדום, שהרי ישמעאל כבר היה לו, ולא היה מבשרו עליו:}%
% ---------- פרק 17 | פסוק 7 ----------
 \textbf{ז} וַהֲקִמֹתִ֨י אֶת־בְּרִיתִ֜י בֵּינִ֣י וּבֵינֶ֗ךָ וּבֵ֨ין זַרְעֲךָ֧ אַחֲרֶ֛יךָ לְדֹרֹתָ֖ם לִבְרִ֣ית עוֹלָ֑ם לִהְי֤וֹת לְךָ֙ לֵֽאלֹהִ֔ים וּֽלְזַרְעֲךָ֖ אַחֲרֶֽיךָ׃ \Rashi{\textbf{והקמתי את בריתי.}ומה היא הברית? להיות לך לאלהים:}%
% ---------- פרק 17 | פסוק 8 ----------
 \textbf{ח} וְנָתַתִּ֣י לְ֠ךָ֠ וּלְזַרְעֲךָ֨ אַחֲרֶ֜יךָ אֵ֣ת ׀ אֶ֣רֶץ מְגֻרֶ֗יךָ אֵ֚ת כׇּל־אֶ֣רֶץ כְּנַ֔עַן לַאֲחֻזַּ֖ת עוֹלָ֑ם וְהָיִ֥יתִי לָהֶ֖ם לֵאלֹהִֽים׃ \Rashi{\textbf{לאחזת עולם.}ושם אהיה לכם לאלהים, אבל הדר בחוצה לארץ דומה כמי שאין לו אלוה:}%
% ---------- פרק 17 | פסוק 9 ----------
 \textbf{ט} וַיֹּ֤אמֶר אֱלֹהִים֙ אֶל־אַבְרָהָ֔ם וְאַתָּ֖ה אֶת־בְּרִיתִ֣י תִשְׁמֹ֑ר אַתָּ֛ה וְזַרְעֲךָ֥ אַֽחֲרֶ֖יךָ לְדֹרֹתָֽם׃ \Rashi{\textbf{ואתה.}וי"ו זו מוסיף על ענין ראשון, אני הנה בריתי אתך, ואתה הוי זהיר לשמרו, ומה היא שמירתו? זאת בריתי אשר תשמרו המול לכם וגו':}%
% ---------- פרק 17 | פסוק 10 ----------
 \textbf{י} זֹ֣את בְּרִיתִ֞י אֲשֶׁ֣ר תִּשְׁמְר֗וּ בֵּינִי֙ וּבֵ֣ינֵיכֶ֔ם וּבֵ֥ין זַרְעֲךָ֖ אַחֲרֶ֑יךָ הִמּ֥וֹל לָכֶ֖ם כׇּל־זָכָֽר׃ \Rashi{\textbf{ביני וביניכם וגו'.}אותם של עכשו: ובין זרעך אחריך. העתידין להולד: המול. כמו להמול, כמו שאתה אומר עשות, כמו לעשות:}%
% ---------- פרק 17 | פסוק 11 ----------
 \textbf{יא} וּנְמַלְתֶּ֕ם אֵ֖ת בְּשַׂ֣ר עׇרְלַתְכֶ֑ם וְהָיָה֙ לְא֣וֹת בְּרִ֔ית בֵּינִ֖י וּבֵינֵיכֶֽם׃ %
% ---------- פרק 17 | פסוק 12 ----------
 \textbf{יב} וּבֶן־שְׁמֹנַ֣ת יָמִ֗ים יִמּ֥וֹל לָכֶ֛ם כׇּל־זָכָ֖ר לְדֹרֹתֵיכֶ֑ם יְלִ֣יד בָּ֔יִת וּמִקְנַת־כֶּ֙סֶף֙ מִכֹּ֣ל בֶּן־נֵכָ֔ר אֲשֶׁ֛ר לֹ֥א מִֽזַּרְעֲךָ֖ הֽוּא׃ \Rashi{\textbf{יליד בית.}שילדתו השפחה בבית: מקנת כסף. שקנאו משנולד:}%
% ---------- פרק 17 | פסוק 13 ----------
 \textbf{יג} הִמּ֧וֹל ׀ יִמּ֛וֹל יְלִ֥יד בֵּֽיתְךָ֖ וּמִקְנַ֣ת כַּסְפֶּ֑ךָ וְהָיְתָ֧ה בְרִיתִ֛י בִּבְשַׂרְכֶ֖ם לִבְרִ֥ית עוֹלָֽם׃ \Rashi{\textbf{המול ימול יליד ביתך.}כאן כפל עליו ולא אמר לשמונה ימים ללמדך שיש יליד בית נמול לאחר שמונה ימים (ס"א לאחד) כמו שמפרש במסכת שבת (דף קל"ה ב'):}%
% ---------- פרק 17 | פסוק 14 ----------
 \textbf{יד} וְעָרֵ֣ל ׀ זָכָ֗ר אֲשֶׁ֤ר לֹֽא־יִמּוֹל֙ אֶת־בְּשַׂ֣ר עׇרְלָת֔וֹ וְנִכְרְתָ֛ה הַנֶּ֥פֶשׁ הַהִ֖וא מֵעַמֶּ֑יהָ אֶת־בְּרִיתִ֖י הֵפַֽר׃ \{ס\} \Rashi{\textbf{וערל זכר.}כאן למד שהמילה באותו מקום שהוא נכר בין זכר לנקבה (בראשית רבה): אשר לא ימול. משיגיע לכלל ענשין ונכרתה; אבל אביו אין ענוש עליו כרת, אבל עובר בעשה: ונכרתה הנפש. הולך ערירי ומת קדם זמנו (שבת כ"ד):}%
% ---------- פרק 17 | פסוק 15 ----------
 \textbf{טו} וַיֹּ֤אמֶר אֱלֹהִים֙ אֶל־אַבְרָהָ֔ם שָׂרַ֣י אִשְׁתְּךָ֔ לֹא־תִקְרָ֥א אֶת־שְׁמָ֖הּ שָׂרָ֑י כִּ֥י שָׂרָ֖ה שְׁמָֽהּ׃ \Rashi{\textbf{לא תקרא את שמה שרי.}דמשמע שרי, לי ולא לאחרים, כי שרה סתם שמה, שתהא שרה על כל (ברכות י"ג):}%
% ---------- פרק 17 | פסוק 16 ----------
 \textbf{טז} וּבֵרַכְתִּ֣י אֹתָ֔הּ וְגַ֨ם נָתַ֧תִּי מִמֶּ֛נָּה לְךָ֖ בֵּ֑ן וּבֵֽרַכְתִּ֙יהָ֙ וְהָֽיְתָ֣ה לְגוֹיִ֔ם מַלְכֵ֥י עַמִּ֖ים מִמֶּ֥נָּה יִהְיֽוּ׃ \Rashi{\textbf{וברכתי אותה.}ומה היא הברכה? שחזרה לנערותה, שנאמר היתה לי עדנה (בר' י"ח): וברכתיה. בהנקת שדים כשנצרכה לכך ביום משתה של יצחק, שהיו מרננין עליהם שהביאו אסופי מן השוק ואומרים בננו הוא, והביאה כל אחת בנה עמה ומנקתה לא הביאה, והיא הניקה את כלם; הוא שנאמר היניקה בנים שרה, ב"ר רמזו במקצת:}%
% ---------- פרק 17 | פסוק 17 ----------
 \textbf{יז} וַיִּפֹּ֧ל אַבְרָהָ֛ם עַל־פָּנָ֖יו וַיִּצְחָ֑ק וַיֹּ֣אמֶר בְּלִבּ֗וֹ הַלְּבֶ֤ן מֵאָֽה־שָׁנָה֙ יִוָּלֵ֔ד וְאִ֨ם־שָׂרָ֔ה הֲבַת־תִּשְׁעִ֥ים שָׁנָ֖ה תֵּלֵֽד׃ \Rashi{\textbf{ויפל אברהם על פניו ויצחק.}זה ת"א, לשון שמחה, וחדי, ושל שרה לשון מחוך; למדת שאברהם האמין ושמח ושרה לא האמינה ולגלגה, וזהו שהקפיד הקב"ה על שרה ולא הקפיד על אברהם: הלבן. יש תמיהות שהן קימות, כמו הנגלה נגליתי? (שמואל א ב'), הראה אתה (יח' ח'), אף זו היא קימת; וכך אמר בלבו, הנעשה חסד זה לאחר, מה שהקב"ה עושה לי? ואם שרה הבת תשעים שנה. היתה כדאי לילד; ואף על פי שדורות הראשונים היו מולידין בני ת"ק שנה, בימי אברהם נתמעטו השנים כבר ובא תשות כח לעולם, צא ולמד מעשרה דורות שמנח ועד אברהם, שמהרו תולדותיהן בני ס' ובני ע':}%
% ---------- פרק 17 | פסוק 18 ----------
 \textbf{יח} וַיֹּ֥אמֶר אַבְרָהָ֖ם אֶל־הָֽאֱלֹהִ֑ים ל֥וּ יִשְׁמָעֵ֖אל יִחְיֶ֥ה לְפָנֶֽיךָ׃ \Rashi{\textbf{לו ישמעאל יחיה.}הלואי שיחיה ישמעאל; איני כדאי לקבל מתן שכר כזה: יחיה לפניך. יחיה ביראתך, כמו התהלך לפני – פלח קדמי:}%
% ---------- פרק 17 | פסוק 19 ----------
 \textbf{יט} וַיֹּ֣אמֶר אֱלֹהִ֗ים אֲבָל֙ שָׂרָ֣ה אִשְׁתְּךָ֗ יֹלֶ֤דֶת לְךָ֙ בֵּ֔ן וְקָרָ֥אתָ אֶת־שְׁמ֖וֹ יִצְחָ֑ק וַהֲקִמֹתִ֨י אֶת־בְּרִיתִ֥י אִתּ֛וֹ לִבְרִ֥ית עוֹלָ֖ם לְזַרְע֥וֹ אַחֲרָֽיו׃ \Rashi{\textbf{אבל.}לשון אמתת דברים, וכן אבל אשמים אנחנו (בראשית מ"ב), אבל בן אין לה (מלכים ב ד'): וקראת את שמו יצחק. על שם הצחוק, וי"א על שם העשרה נסיונות וצ' שנה של שרה וח' ימים שנמול, וק' שנה של אברהם: והקימתי את בריתי אתו. למה נאמר, כבר כתיב ואתה את בריתי תשמר אתה וזרעך וגו'? אלא לפי שאומר והקמותי וגו', יכול בני ישמעאל ובני קטורה בכלל הקיום, ת"ל והקמתי את בריתי אתו – ולא עם אחרים: ואת בריתי אקים את יצחק. למה נאמר? אלא למד שהיה קדוש מבטן. ד"א אמר ר' אבא, מכאן למד ק"ו בן הגבירה מבן האמה; כתיב הנה ברכתי אתו והפריתי אתו והרביתי אתו, זה ישמעאל, וק"ו ואת בריתי אקים את יצחק: ואת בריתי. ברית המילה תהא מסורה לזרעו של יצחק:}%
% ---------- פרק 17 | פסוק 20 ----------
 \textbf{כ} וּֽלְיִשְׁמָעֵאל֮ שְׁמַעְתִּ֒יךָ֒ הִנֵּ֣ה ׀ בֵּרַ֣כְתִּי אֹת֗וֹ וְהִפְרֵיתִ֥י אֹת֛וֹ וְהִרְבֵּיתִ֥י אֹת֖וֹ בִּמְאֹ֣ד מְאֹ֑ד שְׁנֵים־עָשָׂ֤ר נְשִׂיאִם֙ יוֹלִ֔יד וּנְתַתִּ֖יו לְג֥וֹי גָּדֽוֹל׃ \Rashi{\textbf{שנים עשר נשיאים.}כעננים יכלו; כמו נשיאים ורוח (משלי כ"ה):}%
% ---------- פרק 17 | פסוק 21 ----------
 \textbf{כא} וְאֶת־בְּרִיתִ֖י אָקִ֣ים אֶת־יִצְחָ֑ק אֲשֶׁר֩ תֵּלֵ֨ד לְךָ֤ שָׂרָה֙ לַמּוֹעֵ֣ד הַזֶּ֔ה בַּשָּׁנָ֖ה הָאַחֶֽרֶת׃ %
% ---------- פרק 17 | פסוק 22 ----------
 \textbf{כב} וַיְכַ֖ל לְדַבֵּ֣ר אִתּ֑וֹ וַיַּ֣עַל אֱלֹהִ֔ים מֵעַ֖ל אַבְרָהָֽם׃ \Rashi{\textbf{מעל אברהם.}לשון נקיה הוא כלפי שכינה, ולמדנו שהצדיקים מרכבתו של מקום (בראשית רבה):}%
% ---------- פרק 17 | פסוק 23 ----------
 \textbf{כג} וַיִּקַּ֨ח אַבְרָהָ֜ם אֶת־יִשְׁמָעֵ֣אל בְּנ֗וֹ וְאֵ֨ת כׇּל־יְלִידֵ֤י בֵיתוֹ֙ וְאֵת֙ כׇּל־מִקְנַ֣ת כַּסְפּ֔וֹ כׇּל־זָכָ֕ר בְּאַנְשֵׁ֖י בֵּ֣ית אַבְרָהָ֑ם וַיָּ֜מׇל אֶת־בְּשַׂ֣ר עׇרְלָתָ֗ם בְּעֶ֙צֶם֙ הַיּ֣וֹם הַזֶּ֔ה כַּאֲשֶׁ֛ר דִּבֶּ֥ר אִתּ֖וֹ אֱלֹהִֽים׃ \Rashi{\textbf{בעצם היום.}בו ביום שנצטוה, ביום ולא בלילה; לא נתירא לא מן הגוים ולא מן הלצנים, ושלא יהיו אויביו ובני דורו אומרים אלו ראינוהו לא הנחנוהו למול ולקים מצותו של מקום: וימל. לשון ויפעל:}%
% ---------- פרק 17 | פסוק 24 ----------
 \textbf{כד} וְאַ֨בְרָהָ֔ם בֶּן־תִּשְׁעִ֥ים וָתֵ֖שַׁע שָׁנָ֑ה בְּהִמֹּל֖וֹ בְּשַׂ֥ר עׇרְלָתֽוֹ׃ \Rashi{\textbf{בהמולו.}בהפעלו, כמו בהבראם (בראשית ב'):}%
% ---------- פרק 17 | פסוק 25 ----------
 \textbf{כה} וְיִשְׁמָעֵ֣אל בְּנ֔וֹ בֶּן־שְׁלֹ֥שׁ עֶשְׂרֵ֖ה שָׁנָ֑ה בְּהִ֨מֹּל֔וֹ אֵ֖ת בְּשַׂ֥ר עׇרְלָתֽוֹ׃ \Rashi{\textbf{בהמולו את בשר ערלתו.}באברהם לא נאמר "את", לפי שלא היה חסר אלא חתוך בשר, שכבר נתמעך על ידי תשמיש; אבל ישמעאל שהיה ילד, הזקק לחתך ערלה ולפרע המילה, לכך נאמר בו "את" (בראשית רבה):}%
% ---------- פרק 17 | פסוק 26 ----------
 \textbf{כו} בְּעֶ֙צֶם֙ הַיּ֣וֹם הַזֶּ֔ה נִמּ֖וֹל אַבְרָהָ֑ם וְיִשְׁמָעֵ֖אל בְּנֽוֹ׃ \Rashi{\textbf{בעצם היום הזה.}שמלאו לאברהם צ"ט שנה ולישמעאל י"ג שנים נמול אברהם וישמעאל בנו:}%
% ---------- פרק 17 | פסוק 27 ----------
 \textbf{כז} וְכׇל־אַנְשֵׁ֤י בֵיתוֹ֙ יְלִ֣יד בָּ֔יִת וּמִקְנַת־כֶּ֖סֶף מֵאֵ֣ת בֶּן־נֵכָ֑ר נִמֹּ֖לוּ אִתּֽוֹ׃ \{פ\} %
\addparasha{חומש בראשית | פרשת ויירא}
% ---------- פרק 18 | פסוק 1 ----------
 \textbf{א} וַיֵּרָ֤א אֵלָיו֙ יְהֹוָ֔ה בְּאֵלֹנֵ֖י מַמְרֵ֑א וְה֛וּא יֹשֵׁ֥ב פֶּֽתַח־הָאֹ֖הֶל כְּחֹ֥ם הַיּֽוֹם׃ \Rashi{\textbf{וירא אליו.}לבקר את החולה. אמר רבי חמא בר חנינא, יום שלישי למילתו היה, ובא הקב"ה ושאל בשלומו (בבא מציעא פ"ו): באלוני ממרא. הוא שנתן לו עצה על המילה, לפיכך נגלה עליו בחלקו (בראשית רבה): יושב. ישב כתיב, בקש לעמד, אמר לו הקב"ה שב ואני אעמד, ואתה סימן לבניך שעתיד אני להתיצב בעדת הדינים והן יושבין, שנאמר אלהים נצב בעדת אל (תהילים פ"ב) (בראשית רבה): פתח האהל. לראות אם יש עובר ושב ויכניסם בביתו: כחום היום. הוציא הקב"ה חמה מנרתיקה, שלא להטריחו באורחים, ולפי שראהו מצטער שלא היו אורחים באים, הביא המלאכים עליו בדמות אנשים (בבא מציעא שם):}%
% ---------- פרק 18 | פסוק 2 ----------
 \textbf{ב} וַיִּשָּׂ֤א עֵינָיו֙ וַיַּ֔רְא וְהִנֵּה֙ שְׁלֹשָׁ֣ה אֲנָשִׁ֔ים נִצָּבִ֖ים עָלָ֑יו וַיַּ֗רְא וַיָּ֤רׇץ לִקְרָאתָם֙ מִפֶּ֣תַח הָאֹ֔הֶל וַיִּשְׁתַּ֖חוּ אָֽרְצָה׃ \Rashi{\textbf{והנה שלשה אנשים.}אחד לבשר את שרה ואחד להפך את סדום ואחד לרפאות את אברהם, שאין מלאך אחד עושה שתי שליחיות (בראשית רבה). תדע לך, שכן כל הפרשה הוא מזכירן בלשון רבים ויאכלו, ויאמרו אליו, ובבשורה נאמר ויאמר שוב אשוב אליך, ובהפיכת סדום הוא אומר כי לא אוכל לעשות דבר לבלתי הפכי, ורפאל שרפא את אברהם הלך משם להציל את לוט; הוא שנאמר ויהי כהוציאם אתם החוצה ויאמר המלט על נפשך, למדת שהאחד היה מציל (בראשית רבה): נצבים עליו. לפניו, אבל לשון נקיה הוא כלפי המלאכים: וירא. מהו וירא וירא שתי פעמים? הראשון כמשמעו והשני לשון הבנה, נסתכל שהיו נצבים במקום אחד והבין שלא היו רוצים להטריחו, ואף על פי שיודעים היו שיצא לקראתם עמדו במקומם לכבודו, להראותו שלא רצו להטריחו, וקדם הוא ורץ לקראתם, בבבא מציעא, כתיב נצבים עליו, וכתיב וירץ לקראתם? כד חזיוהו דהוה שרי ואסר פרשו הימנו, מיד וירץ לקראתם:}%
% ---------- פרק 18 | פסוק 3 ----------
 \textbf{ג} וַיֹּאמַ֑ר אֲדֹנָ֗י אִם־נָ֨א מָצָ֤אתִי חֵן֙ בְּעֵינֶ֔יךָ אַל־נָ֥א תַעֲבֹ֖ר מֵעַ֥ל עַבְדֶּֽךָ׃ \Rashi{\textbf{ויאמר אדני אם נא וגו'.}לגדול שבהם אמר וקראם כלם אדונים ולגדול אמר אל נא תעבר, וכיון שלא יעבר הוא יעמדו חבריו עמו, ובלשון זה הוא חל (שבועות ל"ה). ד"א קדש הוא, והיה אומר להקב"ה להמתין לו עד שירוץ ויכניס את האורחים ואע"פ שכתוב אחר וירץ לקראתם, האמירה קדם לכן היתה, ודרך המקראות לדבר כן, כמו שפרשתי אצל לא ידון רוחי באדם, שנכתב אחר ויולד נח, ואי אפשר לומר אלא שהיתה הגזרה קדם ללדה כ' שנה, ושתי הלשונות בב"ר:}%
% ---------- פרק 18 | פסוק 4 ----------
 \textbf{ד} יֻקַּֽח־נָ֣א מְעַט־מַ֔יִם וְרַחֲצ֖וּ רַגְלֵיכֶ֑ם וְהִֽשָּׁעֲנ֖וּ תַּ֥חַת הָעֵֽץ׃ \Rashi{\textbf{יקח נא.}על ידי שליח, והקב"ה שלם לבניו על ידי שליח, שנאמר וירם משה את ידו ויך את הסלע (במדבר כ'): ורחצו רגליכם. כסבור שהם ערביים שמשתחוים לאבק רגליהם והקפיד שלא להכניס ע"ז לביתו; אבל לוט שלא הקפיד, הקדים לינה לרחיצה, שנאמר ולינו ורחצו רגליכם: תחת העץ. תחת האילן:}%
% ---------- פרק 18 | פסוק 5 ----------
 \textbf{ה} וְאֶקְחָ֨ה פַת־לֶ֜חֶם וְסַעֲד֤וּ לִבְּכֶם֙ אַחַ֣ר תַּעֲבֹ֔רוּ כִּֽי־עַל־כֵּ֥ן עֲבַרְתֶּ֖ם עַֽל־עַבְדְּכֶ֑ם וַיֹּ֣אמְר֔וּ כֵּ֥ן תַּעֲשֶׂ֖ה כַּאֲשֶׁ֥ר דִּבַּֽרְתָּ׃ \Rashi{\textbf{וסעדו לבכם.}בתורה, בנביאים ובכתובים מצינו דפתא סעדתא דלבא. בתורה וסעדו לבכם, בנביאים סעד לבך פת לחם (שופטים י"ט), בכתובים ולחם לבב אנוש יסעד (תהל' ק"ד). אמר רבי חמא, לבבכם אין כתיב כאן אלא לבכם, מגיד שאין יצר הרע שולט במלאכים (בראשית רבה): אחר תעבורו. אחר כן תלכו: כי על כן עברתם. כי הדבר הזה אני מבקש מכם מאחר שעברתם עלי לכבודי: כי על כן. כמו על אשר, וכן כל כי על כן שבמקרא, כי על כן באו בצל קרתי (בראשית י"ט), כי על כן ראיתי פניך (בראשית ל״ג:י׳), כי על כן לא נתתיה (שם ל"ח), כי על כן ידעת חנתנו (במדבר י'):}%
% ---------- פרק 18 | פסוק 6 ----------
 \textbf{ו} וַיְמַהֵ֧ר אַבְרָהָ֛ם הָאֹ֖הֱלָה אֶל־שָׂרָ֑ה וַיֹּ֗אמֶר מַהֲרִ֞י שְׁלֹ֤שׁ סְאִים֙ קֶ֣מַח סֹ֔לֶת ל֖וּשִׁי וַעֲשִׂ֥י עֻגֽוֹת׃ \Rashi{\textbf{קמח סלת.}סלת לעגות, קמח לעמילן של טבחים, לכסות את הקדרה לשאב את הזהמא:}%
% ---------- פרק 18 | פסוק 7 ----------
 \textbf{ז} וְאֶל־הַבָּקָ֖ר רָ֣ץ אַבְרָהָ֑ם וַיִּקַּ֨ח בֶּן־בָּקָ֜ר רַ֤ךְ וָטוֹב֙ וַיִּתֵּ֣ן אֶל־הַנַּ֔עַר וַיְמַהֵ֖ר לַעֲשׂ֥וֹת אֹתֽוֹ׃ \Rashi{\textbf{בן בקר רך וטוב.}ג' פרים היו, כדי להאכילן ג' לשונות בחרדל (בבא מציעא פ"ז): אל הנער. זה ישמעאל, לחנכו במצות (בראשית רבה):}%
% ---------- פרק 18 | פסוק 8 ----------
 \textbf{ח} וַיִּקַּ֨ח חֶמְאָ֜ה וְחָלָ֗ב וּבֶן־הַבָּקָר֙ אֲשֶׁ֣ר עָשָׂ֔ה וַיִּתֵּ֖ן לִפְנֵיהֶ֑ם וְהֽוּא־עֹמֵ֧ד עֲלֵיהֶ֛ם תַּ֥חַת הָעֵ֖ץ וַיֹּאכֵֽלוּ׃ \Rashi{\textbf{ויקח חמאה וגו'.}ולחם לא הביא, לפי שפרסה שרה נדה, שחזר לה ארח כנשים אותו היום, ונטמאת העסה (בבא מציעא שם): חמאה. שמן החלב שקולטין מעל פניו: ובן הבקר אשר עשה. אשר תקן; קמא קמא שתקן אמטי ואייתי קמיהו: ויאכלו. נראו כמי שאכלו, מכאן שלא ישנה אדם מן המנהג (בבא מציעא שם):}%
% ---------- פרק 18 | פסוק 9 ----------
 \textbf{ט} וַיֹּאמְר֣וּ אֵׄלָ֔יׄוׄ אַיֵּ֖ה שָׂרָ֣ה אִשְׁתֶּ֑ךָ וַיֹּ֖אמֶר הִנֵּ֥ה בָאֹֽהֶל׃ \Rashi{\textbf{ויאמרו אליו.}נקוד על אי"ו שבאליו, ותניא ר"ש בן אלעזר אומר כל מקום שכתב רבה על הנקדה, אתה דורש הכתב; וכאן הנקדה רבה על הכתב אתה דורש הנקדה, שאף לשרה שאלו איו אברהם, למדנו שישאל אדם באכסניא שלו לאיש על האשה ולאשה על האיש. בבבא מציעא אומרים יודעים היו מלאכי השרת שרה אמנו היכן היתה, אלא להודיע שצנועה היתה, כדי לחבבה על בעלה. אמר רבי יוסי בר חנינא, כדי לשגר לה כוס של ברכה: הנה באהל. צנועה היא:}%
% ---------- פרק 18 | פסוק 10 ----------
 \textbf{י} וַיֹּ֗אמֶר שׁ֣וֹב אָשׁ֤וּב אֵלֶ֙יךָ֙ כָּעֵ֣ת חַיָּ֔ה וְהִנֵּה־בֵ֖ן לְשָׂרָ֣ה אִשְׁתֶּ֑ךָ וְשָׂרָ֥ה שֹׁמַ֛עַת פֶּ֥תַח הָאֹ֖הֶל וְה֥וּא אַחֲרָֽיו׃ \Rashi{\textbf{שוב אשוב.}לא בשרו המלאך שישוב אליו, אלא בשליחותו של מקום אמר לו, כמו ויאמר לה מלאך ה' הרבה ארבה, והוא אין בידו להרבות, אלא בשליחותו של מקום, אף כאן בשליחותו של מקום אמר לו כן (ב"ר). אלישע אמר לשונמית למועד הזה כעת חיה את חבקת בן ותאמר אל אדני איש האלהים אל תכזב בשפחתך, אותן המלאכים שבשרו את שרה אמרו למועד אשוב, אמר לה אלישע אותם המלאכים שהם חיים וקימים לעולם אמרו למועד אשוב, אבל אני בשר ודם שהיום חי ומחר מת – בין חי ובין מת למועד הזה וגו': כעת חיה. כעת הזאת לשנה הבאה, ופסח היה, ולפסח הבא נולד יצחק, מדלא קרינן כעת אלא כעת: כעת חיה. כעת הזאת שתהא חיה לכם, שתהיו כלכם שלמים וקימים: והוא אחריו. הפתח היה אחר המלאך:}%
% ---------- פרק 18 | פסוק 11 ----------
 \textbf{יא} וְאַבְרָהָ֤ם וְשָׂרָה֙ זְקֵנִ֔ים בָּאִ֖ים בַּיָּמִ֑ים חָדַל֙ לִהְי֣וֹת לְשָׂרָ֔ה אֹ֖רַח כַּנָּשִֽׁים׃ \Rashi{\textbf{חדל להיות.}פסק ממנה: אורח כנשים. ארח נדות:}%
% ---------- פרק 18 | פסוק 12 ----------
 \textbf{יב} וַתִּצְחַ֥ק שָׂרָ֖ה בְּקִרְבָּ֣הּ לֵאמֹ֑ר אַחֲרֵ֤י בְלֹתִי֙ הָֽיְתָה־לִּ֣י עֶדְנָ֔ה וַֽאדֹנִ֖י זָקֵֽן׃ \Rashi{\textbf{בקרבה.}מסתכלת במעיה ואמרה אפשר הקרבים הללו טעונין ולד? השדים הללו שצמקו מושכין חלב? תנחומא: עדנה. צחצוח בשר, ולשון משנה משיר את השער ומעדן את הבשר. דבר אחר לשון עדן; זמן וסת נדות:}%
% ---------- פרק 18 | פסוק 13 ----------
 \textbf{יג} וַיֹּ֥אמֶר יְהֹוָ֖ה אֶל־אַבְרָהָ֑ם לָ֣מָּה זֶּה֩ צָחֲקָ֨ה שָׂרָ֜ה לֵאמֹ֗ר הַאַ֥ף אֻמְנָ֛ם אֵלֵ֖ד וַאֲנִ֥י זָקַֽנְתִּי׃ \Rashi{\textbf{האף אמנם.}הגם אמת אלד?  ואני זקנתי שנה הכתוב מפני השלום, שהרי היא אמרה ואדני זקן:}%
% ---------- פרק 18 | פסוק 14 ----------
 \textbf{יד} הֲיִפָּלֵ֥א מֵיְהֹוָ֖ה דָּבָ֑ר לַמּוֹעֵ֞ד אָשׁ֥וּב אֵלֶ֛יךָ כָּעֵ֥ת חַיָּ֖ה וּלְשָׂרָ֥ה בֵֽן׃ \Rashi{\textbf{היפלא.}כתרגומו היתכסי? וכי שום דבר מפלא ומפרד ומכסה ממני מלעשות כרצוני: למועד. לאותו מועד המיחד שקבעתי לך אתמול, למועד הזה בשנה האחרת:}%
% ---------- פרק 18 | פסוק 15 ----------
 \textbf{טו} וַתְּכַחֵ֨שׁ שָׂרָ֧ה ׀ לֵאמֹ֛ר לֹ֥א צָחַ֖קְתִּי כִּ֣י ׀ יָרֵ֑אָה וַיֹּ֥אמֶֽר ׀ לֹ֖א כִּ֥י צָחָֽקְתְּ׃ \Rashi{\textbf{כי יראה וגו' כי צחקת.}הראשון משמש לשון דהא, שנותן טעם לדבר ותכחש שרה – לפי שיראה; והשני משמש בלשון אלא, ויאמר לא כדברך הוא אלא צחקת, שאמרו רבותינו כי משמש בארבע לשונות, אי, דלמא, אלא, דהא:}%
% ---------- פרק 18 | פסוק 16 ----------
 \textbf{טז} וַיָּקֻ֤מוּ מִשָּׁם֙ הָֽאֲנָשִׁ֔ים וַיַּשְׁקִ֖פוּ עַל־פְּנֵ֣י סְדֹ֑ם וְאַ֨בְרָהָ֔ם הֹלֵ֥ךְ עִמָּ֖ם לְשַׁלְּחָֽם׃ \Rashi{\textbf{וישקיפו.}כל השקפה שבמקרא לרעה חוץ מהשקיפה ממעון קדשך (דברים כ"ו), שגדול כח מתנות עניים שהופך מדת הרגז לרחמים: לשלחם. ללוותם; כסבור אורחים הם:}%
% ---------- פרק 18 | פסוק 17 ----------
 \textbf{יז} וַֽיהֹוָ֖ה אָמָ֑ר הַֽמְכַסֶּ֤ה אֲנִי֙ מֵֽאַבְרָהָ֔ם אֲשֶׁ֖ר אֲנִ֥י עֹשֶֽׂה׃ \Rashi{\textbf{המכסה אני.}בתמיה: אשר אני עושה בסדום. לא יפה לי לעשות דבר זה שלא מדעתו. אני נתתי לו את הארץ הזאת, וחמשה כרכין הללו שלו הן, שנאמר: גבול הכנעני מצידן באכה סדמה ועמרה וגו' (בראשית י'). קראתי אותו אברהם – אב המון גוים, ואשמיד את הבנים ולא אודיע לאב, שהוא אוהבי?}%
% ---------- פרק 18 | פסוק 18 ----------
 \textbf{יח} וְאַ֨בְרָהָ֔ם הָי֧וֹ יִֽהְיֶ֛ה לְג֥וֹי גָּד֖וֹל וְעָצ֑וּם וְנִ֨בְרְכוּ־ב֔וֹ כֹּ֖ל גּוֹיֵ֥י הָאָֽרֶץ׃ \Rashi{\textbf{ואברהם היו יהיה.}מ"א: זכר צדיק לברכה, הואיל והזכירו – ברכו. ופשוטו: וכי ממנו אני מעלים, והרי הוא חביב לפני להיות לגוי גדול ולהתברך בו כל גויי הארץ.}%
% ---------- פרק 18 | פסוק 19 ----------
 \textbf{יט} כִּ֣י יְדַעְתִּ֗יו לְמַ֩עַן֩ אֲשֶׁ֨ר יְצַוֶּ֜ה אֶת־בָּנָ֤יו וְאֶת־בֵּיתוֹ֙ אַחֲרָ֔יו וְשָֽׁמְרוּ֙ דֶּ֣רֶךְ יְהֹוָ֔ה לַעֲשׂ֥וֹת צְדָקָ֖ה וּמִשְׁפָּ֑ט לְמַ֗עַן הָבִ֤יא יְהֹוָה֙ עַל־אַבְרָהָ֔ם אֵ֥ת אֲשֶׁר־דִּבֶּ֖ר עָלָֽיו׃ \Rashi{\textbf{כי ידעתיו.}לשון חבה, כמו מודע לאישה (רות ב'), הלא בעז מדעתנו (שם ג'), ואדעך בשם (שמות ל"ג). ואמנם עקר לשון כלם אינו אלא לשון ידיעה, שהמחבב את האדם מקרבו אצלו ויודעו ומכירו, ולמה ידעתיו? למען אשר יצוה. לפי שהוא מצוה את בניו עלי לשמר דרכי. ואם תפרשהו כתרגומו, יודע אני בו שיצוה את בניו וגו', אין למען נופל על הלשון: יצוה. לשון הווה, כמו ככה יעשה איוב (איוב א'): למען הביא. כך הוא מצוה לבניו, שמרו דרך ה' כדי שיביא ה' על אברהם וגו'; על בית אברהם לא נאמר אלא על אברהם, למדנו, כל המעמיד בן צדיק כאלו אינו מת:}%
% ---------- פרק 18 | פסוק 20 ----------
 \textbf{כ} וַיֹּ֣אמֶר יְהֹוָ֔ה זַעֲקַ֛ת סְדֹ֥ם וַעֲמֹרָ֖ה כִּי־רָ֑בָּה וְחַ֨טָּאתָ֔ם כִּ֥י כָבְדָ֖ה מְאֹֽד׃ \Rashi{\textbf{ויאמר ה'.}אל אברהם שעשה כאשר אמר, שלא יכסה ממנו: כי רבה. כל רבה שבמקרא הטעם למטה בבי"ת, לפי שהן מתרגמין גדולה, או גדלה והולכת, אבל זה טעמו למעלה ברי"ש, לפי שמתרגם גדלה כבר כמו שפרשתי ויהי השמש באה, הנה שבה יבמתך:}%
% ---------- פרק 18 | פסוק 21 ----------
 \textbf{כא} אֵֽרְדָה־נָּ֣א וְאֶרְאֶ֔ה הַכְּצַעֲקָתָ֛הּ הַבָּ֥אָה אֵלַ֖י עָשׂ֣וּ ׀ כָּלָ֑ה וְאִם־לֹ֖א אֵדָֽעָה׃ \Rashi{\textbf{ארדה נא.}למד לדינים שלא יפסקו דיני נפשות אלא בראיה, הכל כמו שפרשתי בפרשת הפלגה. דבר אחר, ארדה נא לסוף מעשיהם: הכצעקתה. של מדינה: הבאה אלי עשו. וכן עומדים במרדם, כלה אני עושה בהם, ואם לא יעמדו במרדן אדעה מה אעשה להפרע מהן ביסורין ולא אכלה אותן. וכיוצא בו מצינו במקום אחר, ועתה הורד עדיך מעליך ואדעה מה אעשה לך (שמות ל"ג); ולפיכך יש הפסק נקדת פסיק בין עשו לכלה, כדי להפריד תבה מחברתה. ורבותינו דרשו הכצעקתה – צעקת ריבה אחת שהרגו מיתה משנה על שנתנה מזון לעני, כמפרש בחלק:}%
% ---------- פרק 18 | פסוק 22 ----------
 \textbf{כב} וַיִּפְנ֤וּ מִשָּׁם֙ הָֽאֲנָשִׁ֔ים וַיֵּלְכ֖וּ סְדֹ֑מָה וְאַ֨בְרָהָ֔ם עוֹדֶ֥נּוּ עֹמֵ֖ד לִפְנֵ֥י יְהֹוָֽה׃ \Rashi{\textbf{ויפנו משם.}ממקום שאברהם לום שם: ואברהם עודנו עומד לפני ה'. והלא לא הלך לעמד לפניו, אלא הקב"ה בא אצלו ואמר לו, זעקת סדם ועמרה כי רבה, והיה לו לכתב "וה' עודנו עומד על אברהם"? אלא תקון סופרים הוא זה (אשר הפכוהו רז"ל לכתב כן) (בראשית רבה):}%
% ---------- פרק 18 | פסוק 23 ----------
 \textbf{כג} וַיִּגַּ֥שׁ אַבְרָהָ֖ם וַיֹּאמַ֑ר הַאַ֣ף תִּסְפֶּ֔ה צַדִּ֖יק עִם־רָשָֽׁע׃ \Rashi{\textbf{ויגש אברהם.}מצינו הגשה למלחמה, ויגש יואב וגו' (שמואל ב י'), והגשה לפיוס ויגש אליו יהודה, והגשה לתפלה, ויגש אליהו הנביא (מלכים א י"ח), ולכל אלה נכנס אברהם לדבר קשות ולפיוס ולתפלה: האף תספה. הגם תספה. ולתרגום של אנקלוס, שתרגמו לשון רגז, כך פרושו, האף ישיאך שתספה צדיק עם רשע:}%
% ---------- פרק 18 | פסוק 24 ----------
 \textbf{כד} אוּלַ֥י יֵ֛שׁ חֲמִשִּׁ֥ים צַדִּיקִ֖ם בְּת֣וֹךְ הָעִ֑יר הַאַ֤ף תִּסְפֶּה֙ וְלֹא־תִשָּׂ֣א לַמָּק֔וֹם לְמַ֛עַן חֲמִשִּׁ֥ים הַצַּדִּיקִ֖ם אֲשֶׁ֥ר בְּקִרְבָּֽהּ׃ \Rashi{\textbf{אולי יש חמשים צדיקים.}עשרה צדיקים לכל כרך וכרך, כי ה' מקומות יש. ואם תאמר: לא יצילו הצדיקים את הרשעים, למה תמית הצדיקים?}%
% ---------- פרק 18 | פסוק 25 ----------
 \textbf{כה} חָלִ֨לָה לְּךָ֜ מֵעֲשֹׂ֣ת ׀ כַּדָּבָ֣ר הַזֶּ֗ה לְהָמִ֤ית צַדִּיק֙ עִם־רָשָׁ֔ע וְהָיָ֥ה כַצַּדִּ֖יק כָּרָשָׁ֑ע חָלִ֣לָה לָּ֔ךְ הֲשֹׁפֵט֙ כׇּל־הָאָ֔רֶץ לֹ֥א יַעֲשֶׂ֖ה מִשְׁפָּֽט׃ \Rashi{\textbf{חלילה לך.}חלין הוא לך, יאמרו כך הוא אמנותו, שוטף הכל, צדיקים ורשעים, כך עשית לדור המבול ולדור הפלגה: כדבר הזה. לא הוא ולא כיוצא בו: חלילה לך. לעולם הבא: השופט כל הארץ. נקוד בחטף פתח ה"א של השופט, לשון תמיה, וכי מי שהוא שופט לא יעשה משפט אמת:}%
% ---------- פרק 18 | פסוק 26 ----------
 \textbf{כו} וַיֹּ֣אמֶר יְהֹוָ֔ה אִם־אֶמְצָ֥א בִסְדֹ֛ם חֲמִשִּׁ֥ים צַדִּיקִ֖ם בְּת֣וֹךְ הָעִ֑יר וְנָשָׂ֥אתִי לְכׇל־הַמָּק֖וֹם בַּעֲבוּרָֽם׃ \Rashi{\textbf{אם אמצא בסדם וגו' לכל המקום.}לכל הכרכים; ולפי שסדום היתה מטרופולין וחשובה מכלם, תלה בה הכתוב:}%
% ---------- פרק 18 | פסוק 27 ----------
 \textbf{כז} וַיַּ֥עַן אַבְרָהָ֖ם וַיֹּאמַ֑ר הִנֵּה־נָ֤א הוֹאַ֙לְתִּי֙ לְדַבֵּ֣ר אֶל־אֲדֹנָ֔י וְאָנֹכִ֖י עָפָ֥ר וָאֵֽפֶר׃ \Rashi{\textbf{הואלתי.}רציתי כמו ויואל משה (שמות ב'): ואנכי עפר ואפר. וכבר הייתי ראוי להיות עפר על ידי המלכים, ואפר על ידי נמרוד, לולי רחמיך אשר עמדו לי:}%
% ---------- פרק 18 | פסוק 28 ----------
 \textbf{כח} א֠וּלַ֠י יַחְסְר֞וּן חֲמִשִּׁ֤ים הַצַּדִּיקִם֙ חֲמִשָּׁ֔ה הֲתַשְׁחִ֥ית בַּחֲמִשָּׁ֖ה אֶת־כׇּל־הָעִ֑יר וַיֹּ֙אמֶר֙ לֹ֣א אַשְׁחִ֔ית אִם־אֶמְצָ֣א שָׁ֔ם אַרְבָּעִ֖ים וַחֲמִשָּֽׁה׃ \Rashi{\textbf{התשחית בחמשה.}והלא הן ט' לכל כרך, ואתה צדיקו של עולם תצטרף עמהם:}%
% ---------- פרק 18 | פסוק 29 ----------
 \textbf{כט} וַיֹּ֨סֶף ע֜וֹד לְדַבֵּ֤ר אֵלָיו֙ וַיֹּאמַ֔ר אוּלַ֛י יִמָּצְא֥וּן שָׁ֖ם אַרְבָּעִ֑ים וַיֹּ֙אמֶר֙ לֹ֣א אֶֽעֱשֶׂ֔ה בַּעֲב֖וּר הָאַרְבָּעִֽים׃ \Rashi{\textbf{אולי ימצאון שם ארבעים.}וימלטו ד' הכרכים וכן ל' יצילו ג' מהם או כ' יצילו ב' מהם או י' יצילו אחד מהם:}%
% ---------- פרק 18 | פסוק 30 ----------
 \textbf{ל} וַ֠יֹּ֠אמֶר אַל־נָ֞א יִ֤חַר לַֽאדֹנָי֙ וַאֲדַבֵּ֔רָה אוּלַ֛י יִמָּצְא֥וּן שָׁ֖ם שְׁלֹשִׁ֑ים וַיֹּ֙אמֶר֙ לֹ֣א אֶֽעֱשֶׂ֔ה אִם־אֶמְצָ֥א שָׁ֖ם שְׁלֹשִֽׁים׃ %
% ---------- פרק 18 | פסוק 31 ----------
 \textbf{לא} וַיֹּ֗אמֶר הִנֵּֽה־נָ֤א הוֹאַ֙לְתִּי֙ לְדַבֵּ֣ר אֶל־אֲדֹנָ֔י אוּלַ֛י יִמָּצְא֥וּן שָׁ֖ם עֶשְׂרִ֑ים וַיֹּ֙אמֶר֙ לֹ֣א אַשְׁחִ֔ית בַּעֲב֖וּר הָֽעֶשְׂרִֽים׃ \Rashi{\textbf{הואלתי.}רציתי, כמו ויואל משה (שמות ב'):}%
% ---------- פרק 18 | פסוק 32 ----------
 \textbf{לב} וַ֠יֹּ֠אמֶר אַל־נָ֞א יִ֤חַר לַֽאדֹנָי֙ וַאֲדַבְּרָ֣ה אַךְ־הַפַּ֔עַם אוּלַ֛י יִמָּצְא֥וּן שָׁ֖ם עֲשָׂרָ֑ה וַיֹּ֙אמֶר֙ לֹ֣א אַשְׁחִ֔ית בַּעֲב֖וּר הָעֲשָׂרָֽה׃ \Rashi{\textbf{אולי ימצאון שם עשרה.}על הפחות לא בקש, אמר דור המבול היו ח', נח ובניו ונשיהם, ולא הצילו על דורם; ועל ט' על ידי צרוף כבר בקש ולא מצא:}%
% ---------- פרק 18 | פסוק 33 ----------
 \textbf{לג} וַיֵּ֣לֶךְ יְהֹוָ֔ה כַּאֲשֶׁ֣ר כִּלָּ֔ה לְדַבֵּ֖ר אֶל־אַבְרָהָ֑ם וְאַבְרָהָ֖ם שָׁ֥ב לִמְקֹמֽוֹ׃ \Rashi{\textbf{וילך ה' וגו'.}כיון שנשתתק הסנגור הלך לו הדין: ואברהם שב למקומו. נסתלק הדין נסתלק הסנגור, והקטיגור מקטרג ולפיכך ויבאו שני המלאכים סדמה, להשחית:}%
% ---------- פרק 19 | פסוק 1 ----------
 \textbf{א} וַ֠יָּבֹ֠אוּ שְׁנֵ֨י הַמַּלְאָכִ֤ים סְדֹ֙מָה֙ בָּעֶ֔רֶב וְל֖וֹט יֹשֵׁ֣ב בְּשַֽׁעַר־סְדֹ֑ם וַיַּרְא־לוֹט֙ וַיָּ֣קׇם לִקְרָאתָ֔ם וַיִּשְׁתַּ֥חוּ אַפַּ֖יִם אָֽרְצָה׃ \Rashi{\textbf{שני.}א' להשחית את סדום וא' להציל את לוט, והוא אותו שבא לרפאות אברהם, והשלישי שבא לבשר את שרה, כיון שעשה שליחותו נסתלק לו: המלאכים. ולהלן קראם אנשים? כשהיתה שכינה עמהם קראם אנשים; דבר אחר אצל אברהם שכחו גדול והיו המלאכים תדירין אצלו כאנשים קראם אנשים, ואצל לוט קראם מלאכים (בראשית רבה): בערב. וכי כל כך שהו המלאכים מחברון לסדום? אלא מלאכי רחמים היו וממתינים שמא יוכל אברהם ללמד עליהם סנגוריא: ולוט ישב בשער סדום. ישב כתיב, אותו היום מנוהו שופט עליהם (בראשית רבה): וירא לוט וגו'. מבית אברהם למד לחזר על האורחים:}%
% ---------- פרק 19 | פסוק 2 ----------
 \textbf{ב} וַיֹּ֜אמֶר הִנֶּ֣ה נָּא־אֲדֹנַ֗י ס֣וּרוּ נָ֠א אֶל־בֵּ֨ית עַבְדְּכֶ֤ם וְלִ֙ינוּ֙ וְרַחֲצ֣וּ רַגְלֵיכֶ֔ם וְהִשְׁכַּמְתֶּ֖ם וַהֲלַכְתֶּ֣ם לְדַרְכְּכֶ֑ם וַיֹּאמְר֣וּ לֹּ֔א כִּ֥י בָרְח֖וֹב נָלִֽין׃ \Rashi{\textbf{הנה נא אדני.}הנה נא אתם אדונים לי אחר שעברתם עלי. דבר אחר הנה נא צריכים אתם לתת לב על הרשעים הללו שלא יכירו בכם וזו היא עצה נכונה: סורו נא. עקמו את הדרך לביתי דרך עקלתון שלא יכירו שאתם נכנסין שם, לכך נאמר סורו (בראשית רבה): ולינו ורחצו רגליכם. וכי דרכן של בני אדם ללון תחלה ואחר כך רוחץ? ועוד, שהרי אברהם אמר להם תחלה רחצו רגליכם? אלא כך אמר לוט, אם כשיבאו אנשי סדום ויראו שכבר רחצו רגליהם יעלילו עלי ויאמרו כבר עברו שני ימים או שלשה שבאו לביתך ולא הודעתנו, לפיכך אמר, מוטב שיתעכבו כאן באבק רגליהם שיהיו נראין כמו שבאו עכשו; לפיכך אמר לינו תחלה ואחר כך רחצו: ויאמרו לא. ולאברהם אמרו כן תעשה, מכאן שמסרבין לקטן ואין מסרבין לגדול (בראשית רבה): כי ברחוב נלין. הרי כי משמש בלשון אלא, שאמרו לא נסור אל ביתך אלא ברחובה של עיר נלין:}%
% ---------- פרק 19 | פסוק 3 ----------
 \textbf{ג} וַיִּפְצַר־בָּ֣ם מְאֹ֔ד וַיָּסֻ֣רוּ אֵלָ֔יו וַיָּבֹ֖אוּ אֶל־בֵּית֑וֹ וַיַּ֤עַשׂ לָהֶם֙ מִשְׁתֶּ֔ה וּמַצּ֥וֹת אָפָ֖ה וַיֹּאכֵֽלוּ׃ \Rashi{\textbf{ויסרו אליו.}עקמו את הדרך לצד ביתו: ומצות אפה. פסח היה:}%
% ---------- פרק 19 | פסוק 4 ----------
 \textbf{ד} טֶ֘רֶם֮ יִשְׁכָּ֒בוּ֒ וְאַנְשֵׁ֨י הָעִ֜יר אַנְשֵׁ֤י סְדֹם֙ נָסַ֣בּוּ עַל־הַבַּ֔יִת מִנַּ֖עַר וְעַד־זָקֵ֑ן כׇּל־הָעָ֖ם מִקָּצֶֽה׃ \Rashi{\textbf{טרם ישכבו ואנשי העיר אנשי סדום.}כך נדרש בב"ר טרם ישכבו ואנשי העיר היו בפיהם של מלאכים, שהיו שואלים ללוט מה טיבם ומעשיהם והוא אומר להם רבם רשעים; עודם מדברים בהם ואנשי סדום וגו'. ופשוטו של מקרא ואנשי העיר אנשי רשע, נסבו על הבית; על שהיו רשעים נקראים אנשי "סדום", כמו שאמר הכתוב ואנשי סדם רעים וחטאים: כל העם מקצה. מקצה העיר עד הקצה, שאין אחד מהם מוחה בידם, שאפלו צדיק אחד אין בהם:}%
% ---------- פרק 19 | פסוק 5 ----------
 \textbf{ה} וַיִּקְרְא֤וּ אֶל־לוֹט֙ וַיֹּ֣אמְרוּ ל֔וֹ אַיֵּ֧ה הָאֲנָשִׁ֛ים אֲשֶׁר־בָּ֥אוּ אֵלֶ֖יךָ הַלָּ֑יְלָה הוֹצִיאֵ֣ם אֵלֵ֔ינוּ וְנֵדְעָ֖ה אֹתָֽם׃ \Rashi{\textbf{ונדעה אותם.}במשכב זכר כמו אשר לא ידעו איש (בראשית רבה):}%
% ---------- פרק 19 | פסוק 6 ----------
 \textbf{ו} וַיֵּצֵ֧א אֲלֵהֶ֛ם ל֖וֹט הַפֶּ֑תְחָה וְהַדֶּ֖לֶת סָגַ֥ר אַחֲרָֽיו׃ %
% ---------- פרק 19 | פסוק 7 ----------
 \textbf{ז} וַיֹּאמַ֑ר אַל־נָ֥א אַחַ֖י תָּרֵֽעוּ׃ %
% ---------- פרק 19 | פסוק 8 ----------
 \textbf{ח} הִנֵּה־נָ֨א לִ֜י שְׁתֵּ֣י בָנ֗וֹת אֲשֶׁ֤ר לֹֽא־יָדְעוּ֙ אִ֔ישׁ אוֹצִֽיאָה־נָּ֤א אֶתְהֶן֙ אֲלֵיכֶ֔ם וַעֲשׂ֣וּ לָהֶ֔ן כַּטּ֖וֹב בְּעֵינֵיכֶ֑ם רַ֠ק לָֽאֲנָשִׁ֤ים הָאֵל֙ אַל־תַּעֲשׂ֣וּ דָבָ֔ר כִּֽי־עַל־כֵּ֥ן בָּ֖אוּ בְּצֵ֥ל קֹרָתִֽי׃ \Rashi{\textbf{האל.}כמו האלה: כי על כן באו. כי הטובה הזאת תעשו לכבודי על אשר באו בצל קרתי; תרגום בטלל שרותי, תרגום של קורה – שרותא:}%
% ---------- פרק 19 | פסוק 9 ----------
 \textbf{ט} וַיֹּאמְר֣וּ ׀ גֶּשׁ־הָ֗לְאָה וַיֹּֽאמְרוּ֙ הָאֶחָ֤ד בָּֽא־לָגוּר֙ וַיִּשְׁפֹּ֣ט שָׁפ֔וֹט עַתָּ֕ה נָרַ֥ע לְךָ֖ מֵהֶ֑ם וַיִּפְצְר֨וּ בָאִ֤ישׁ בְּלוֹט֙ מְאֹ֔ד וַֽיִּגְּשׁ֖וּ לִשְׁבֹּ֥ר הַדָּֽלֶת׃ \Rashi{\textbf{ויאמרו גש הלאה.}קרב להלאה, כלומר, התקרב לצדדין והתרחק ממנו, וכן כל הלאה שבמקרא לשון רחוק, כמו זרה הלאה (במדבר י"ז), הנה החצי ממך והלאה (שמואל א' כ'), גש הלאה, (ישעיהו מ"ט) – המשך להלן, בלשון לעז טריטי"דנוש, ודבר נזיפה הוא, לומר – אין אנו חוששין לך. ודומה לו קרב אליך אל תגש בי (ישעיהו ס"ה), וכן גשה לי ואשבה (שם מ"ט) – המשך לצדדין בעבורי ואשב אצלך; אתה מליץ על האורחין, איך מלאך לבך? על שאמר להם על הבנות, אמרו לו גש הלאה, לשון נחת; ועל שהיה מליץ על האורחים – האחד בא לגור, אדם נכרי יחידי אתה בינינו, שבאת לגור, וישפוט שפוט, ונעשית מוכיח אותנו? הדלת. היא הסובבת לנעל ולפתח:}%
% ---------- פרק 19 | פסוק 10 ----------
 \textbf{י} וַיִּשְׁלְח֤וּ הָֽאֲנָשִׁים֙ אֶת־יָדָ֔ם וַיָּבִ֧יאוּ אֶת־ל֛וֹט אֲלֵיהֶ֖ם הַבָּ֑יְתָה וְאֶת־הַדֶּ֖לֶת סָגָֽרוּ׃ %
% ---------- פרק 19 | פסוק 11 ----------
 \textbf{יא} וְֽאֶת־הָאֲנָשִׁ֞ים אֲשֶׁר־פֶּ֣תַח הַבַּ֗יִת הִכּוּ֙ בַּסַּנְוֵרִ֔ים מִקָּטֹ֖ן וְעַד־גָּד֑וֹל וַיִּלְא֖וּ לִמְצֹ֥א הַפָּֽתַח׃ \Rashi{\textbf{פתח.}הוא החלל שבו נכנסין ויוצאין: בסנורים. מכת עורון: מקטן ועד גדול. הקטנים התחילו בעברה תחלה, שנאמר מנער ועד זקן, לפיכך התחילה הפרענות מהם (בראשית רבה):}%
% ---------- פרק 19 | פסוק 12 ----------
 \textbf{יב} וַיֹּאמְר֨וּ הָאֲנָשִׁ֜ים אֶל־ל֗וֹט עֹ֚ד מִֽי־לְךָ֣ פֹ֔ה חָתָן֙ וּבָנֶ֣יךָ וּבְנֹתֶ֔יךָ וְכֹ֥ל אֲשֶׁר־לְךָ֖ בָּעִ֑יר הוֹצֵ֖א מִן־הַמָּקֽוֹם׃ \Rashi{\textbf{עוד מי לך פה.}פשוטו של מקרא מי יש לך עוד בעיר הזאת חוץ מאשתך ובנותיך שבבית: חתן ובניך ובנתיך. ואם יש לך חתן או בנים ובנות הוצא מן המקום: ובניך. בני בנותיך הנשואות. ומ"א עד – מאחר שעושין נבלה כזאת, מי לך פתחון פה ללמד סנגוריא עליהם? שכל הלילה היה מליץ עליהם טובות, קרי ביה מי לך פה:}%
% ---------- פרק 19 | פסוק 13 ----------
 \textbf{יג} כִּֽי־מַשְׁחִתִ֣ים אֲנַ֔חְנוּ אֶת־הַמָּק֖וֹם הַזֶּ֑ה כִּֽי־גָדְלָ֤ה צַעֲקָתָם֙ אֶת־פְּנֵ֣י יְהֹוָ֔ה וַיְשַׁלְּחֵ֥נוּ יְהֹוָ֖ה לְשַׁחֲתָֽהּ׃ %
% ---------- פרק 19 | פסוק 14 ----------
 \textbf{יד} וַיֵּצֵ֨א ל֜וֹט וַיְדַבֵּ֣ר ׀ אֶל־חֲתָנָ֣יו ׀ לֹקְחֵ֣י בְנֹתָ֗יו וַיֹּ֙אמֶר֙ ק֤וּמוּ צְּאוּ֙ מִן־הַמָּק֣וֹם הַזֶּ֔ה כִּֽי־מַשְׁחִ֥ית יְהֹוָ֖ה אֶת־הָעִ֑יר וַיְהִ֥י כִמְצַחֵ֖ק בְּעֵינֵ֥י חֲתָנָֽיו׃ \Rashi{\textbf{חתניו.}שתי בנות נשואות היו לו בעיר: לוקחי בנותיו. שאותן שבבית ארוסות להם:}%
% ---------- פרק 19 | פסוק 15 ----------
 \textbf{טו} וּכְמוֹ֙ הַשַּׁ֣חַר עָלָ֔ה וַיָּאִ֥יצוּ הַמַּלְאָכִ֖ים בְּל֣וֹט לֵאמֹ֑ר קוּם֩ קַ֨ח אֶֽת־אִשְׁתְּךָ֜ וְאֶת־שְׁתֵּ֤י בְנֹתֶ֙יךָ֙ הַנִּמְצָאֹ֔ת פֶּן־תִּסָּפֶ֖ה בַּעֲוֺ֥ן הָעִֽיר׃ \Rashi{\textbf{ויאיצו.}כתרגומו ודחיקו – מהרוהו: הנמצאות. המזמנות לך בבית להצילם. ומ"א יש וזה ישובו של מקרא: תספה. תהיה כלה, עד תם כל הדור (דברים ב') מתרגם עד דסף כל דרא:}%
% ---------- פרק 19 | פסוק 16 ----------
 \textbf{טז} וַֽיִּתְמַהְמָ֓הּ ׀ וַיַּחֲזִ֨יקוּ הָאֲנָשִׁ֜ים בְּיָד֣וֹ וּבְיַד־אִשְׁתּ֗וֹ וּבְיַד֙ שְׁתֵּ֣י בְנֹתָ֔יו בְּחֶמְלַ֥ת יְהֹוָ֖ה עָלָ֑יו וַיֹּצִאֻ֥הוּ וַיַּנִּחֻ֖הוּ מִח֥וּץ לָעִֽיר׃ \Rashi{\textbf{ויתמהמה.}כדי להציל את ממונו: ויחזיקו. אחד מהם היה שליח להצילו וחברו להפך סדום, לכך נאמר "ויאמר" המלט, ולא נאמר ויאמרו (בראשית רבה):}%
% ---------- פרק 19 | פסוק 17 ----------
 \textbf{יז} וַיְהִי֩ כְהוֹצִיאָ֨ם אֹתָ֜ם הַח֗וּצָה וַיֹּ֙אמֶר֙ הִמָּלֵ֣ט עַל־נַפְשֶׁ֔ךָ אַל־תַּבִּ֣יט אַחֲרֶ֔יךָ וְאַֽל־תַּעֲמֹ֖ד בְּכׇל־הַכִּכָּ֑ר הָהָ֥רָה הִמָּלֵ֖ט פֶּן־תִּסָּפֶֽה׃ \Rashi{\textbf{המלט על נפשך.}דיך להציל נפשות, אל תחוס על הממון: אל תביט אחריך. אתה הרשעת עמהם ובזכות אברהם אתה נצל; אינך כדאי לראות בפרענותם ואתה נצל: בכל הככר. ככר הירדן: ההרה המלט. אצל אברהם ברח, שהוא יושב בהר, שנאמר ויעתק משם ההרה (בראשית י"ב), ואף עכשו היה יושב שם, שנאמר אל המקום אשר היה שם אהלו בתחלה, ואע"פ שכתוב ויאהל אברהם וגו' (שם י"ג), אהלים הרבה היו לו ונמשכו עד חברון: המלט. לשון השמטה, וכן כל המלטה שבמקרא אשמוצי"ר בלע"ז וכן והמליטה זכר (ישעיהו ס"ו), שנשמט העבר מן הרחם, כצפור נמלטה (תהילים קכ"ד), לא יכלו מלט משא (ישעיהו מ"ו), להשמיט משא הרעי שבנקביהם:}%
% ---------- פרק 19 | פסוק 18 ----------
 \textbf{יח} וַיֹּ֥אמֶר ל֖וֹט אֲלֵהֶ֑ם אַל־נָ֖א אֲדֹנָֽי׃ \Rashi{\textbf{אל נא אדני.}רבותינו אמרו (שבועות ל"ב) ששם זה קדש, שנאמר בו להחיות את נפשי – מי שיש בידו להמית ולהחיות; ותרגומו בבעו כען ה': אל נא. אל תאמרו אלי להמלט ההרה: נא. לשון בקשה:}%
% ---------- פרק 19 | פסוק 19 ----------
 \textbf{יט} הִנֵּה־נָ֠א מָצָ֨א עַבְדְּךָ֣ חֵן֮ בְּעֵינֶ֒יךָ֒ וַתַּגְדֵּ֣ל חַסְדְּךָ֗ אֲשֶׁ֤ר עָשִׂ֙יתָ֙ עִמָּדִ֔י לְהַחֲי֖וֹת אֶת־נַפְשִׁ֑י וְאָנֹכִ֗י לֹ֤א אוּכַל֙ לְהִמָּלֵ֣ט הָהָ֔רָה פֶּן־תִּדְבָּקַ֥נִי הָרָעָ֖ה וָמַֽתִּי׃ \Rashi{\textbf{פן תדבקני הרעה.}כשהייתי אצל אנשי סדום היה הקב"ה רואה מעשי ומעשה בני העיר, והייתי נראה צדיק וכדאי להנצל, וכשאבא אצל צדיק אני כרשע, וכן אמרה הצרפית לאליהו באת אלי להזכיר את עוני (מלכים א י"ז), עד שלא באת אצלי היה הקב"ה רואה מעשי ומעשה עמי, ואני צדקת ביניהם, ומשבאת אצלי, לפי מעשיך אני רשעה:}%
% ---------- פרק 19 | פסוק 20 ----------
 \textbf{כ} הִנֵּה־נָ֠א הָעִ֨יר הַזֹּ֧את קְרֹבָ֛ה לָנ֥וּס שָׁ֖מָּה וְהִ֣וא מִצְעָ֑ר אִמָּלְטָ֨ה נָּ֜א שָׁ֗מָּה הֲלֹ֥א מִצְעָ֛ר הִ֖וא וּתְחִ֥י נַפְשִֽׁי׃ \Rashi{\textbf{העיר הזאת קרובה.}קרובה ישיבתה, נתישבה מקרוב, לפיכך לא נתמלאה סאתה עדין (שבת י'). ומה היא קריבתה? מדור הפלגה שנתפלגו האנשים והתחילו להתישב איש איש במקומו, והיא היתה בשנת מות פלג, ומשם ועד כאן נ"ב שנה, שפלג מת בשנת מ"ח לאברהם. כיצד? פלג חי אחרי הולידו את רעו ר"ט שנה, צא מהם ל"ב כשנולד שרוג, ומשרוג עד שנולד נחור ל', הרי ס"ב, ומנחור עד שנולד תרח כ"ט, הרי צ"א, ומשם עד שנולד אברהם ע', הרי קס"א, תן להם מ"ח הרי ר"ט, ואותה שנה היתה שנת הפלגה, וכשנחרבה סדום היה אברהם בן צ"ט שנה, הרי מדור הפלגה עד כאן נ"ב שנה; וצער אחרה ישיבתה אחרי ישיבת סדום וחברותיה שנה אחת, הוא שנאמר אמלטה נא, נא בגימטריא נ"א: הלא מצער היא. והלא עונותיה מועטין ויכול אתה להניחה:  ותחי נפשי בה, זהו מדרשו. ופשוטו של מקרא הלא עיר קטנה היא ואנשים בה מעט, אין לך להקפיד אם תניחנה, ותחי נפשי בה:}%
% ---------- פרק 19 | פסוק 21 ----------
 \textbf{כא} וַיֹּ֣אמֶר אֵלָ֔יו הִנֵּה֙ נָשָׂ֣אתִי פָנֶ֔יךָ גַּ֖ם לַדָּבָ֣ר הַזֶּ֑ה לְבִלְתִּ֛י הׇפְכִּ֥י אֶת־הָעִ֖יר אֲשֶׁ֥ר דִּבַּֽרְתָּ׃ \Rashi{\textbf{גם לדבר הזה.}לא דיך שאתה נצל, אלא אף כל העיר אציל בגללך: הפכי. הופך אני, כמו עד באי (בראשית מ"ח), אחרי ראי (בראשית ט״ז:י״ג), מדי דברי בו (ירמיהו ל"א):}%
% ---------- פרק 19 | פסוק 22 ----------
 \textbf{כב} מַהֵר֙ הִמָּלֵ֣ט שָׁ֔מָּה כִּ֣י לֹ֤א אוּכַל֙ לַעֲשׂ֣וֹת דָּבָ֔ר עַד־בֹּאֲךָ֖ שָׁ֑מָּה עַל־כֵּ֛ן קָרָ֥א שֵׁם־הָעִ֖יר צֽוֹעַר׃ \Rashi{\textbf{כי לא אוכל לעשות.}זה ענשן של מלאכים על שאמרו כי משחתים אנחנו, ותלו הדבר בעצמן, לפיכך לא זזו משם עד שהזקקו לומר שאין הדבר ברשותן: כי לא אוכל. לשון יחיד; מכאן אתה למד שהאחד הופך והאחד מציל, שאין ב' מלאכים נשלחים לדבר א': על כן קרא שם העיר צוער. על שם והוא מצער:}%
% ---------- פרק 19 | פסוק 23 ----------
 \textbf{כג} הַשֶּׁ֖מֶשׁ יָצָ֣א עַל־הָאָ֑רֶץ וְל֖וֹט בָּ֥א צֹֽעֲרָה׃ %
% ---------- פרק 19 | פסוק 24 ----------
 \textbf{כד} וַֽיהֹוָ֗ה הִמְטִ֧יר עַל־סְדֹ֛ם וְעַל־עֲמֹרָ֖ה גׇּפְרִ֣ית וָאֵ֑שׁ מֵאֵ֥ת יְהֹוָ֖ה מִן־הַשָּׁמָֽיִם׃ \Rashi{\textbf{וה' המטיר.}כל מקום שנאמר וה' – הוא ובית דינו: המטיר על סדום. בעלות השחר, כמו שנאמר וכמו השחר עלה, שעה שהלבנה עומדת ברקיע עם החמה, לפי שהיו מהם עובדין לחמה ומהם ללבנה, אמר הקב"ה אם אפרע מהם ביום, יהיו עובדי לבנה אומרים אלו היה בלילה, כשהלבנה מושלת, לא היינו חרבין, ואם אפרע מהם בלילה, יהיו עובדי החמה אומרים אלו היה ביום, כשהחמה מושלת, לא היינו חרבין; לכך כתיב וכמו השחר עלה, ונפרע מהם בשעה שהחמה והלבנה מושלות: המטיר גפרית ואש. בתחלה מטר, ונעשה גפרית ואש: מאת ה'. דרך המקראות לדבר כן, כמו נשי למך, ולא אמר נשי; וכן אמר דוד קחו עמכם את עבדי אדניכם (מלכים א א'), ולא אמר מעבדי; וכן אחשורוש אמר בשם המלך, ולא אמר בשמי; אף כאן אמר, מאת ה' ולא אמר מאתי: מן השמים. הוא שאמר הכתוב כי בם ידין עמים וגו' (איוב ל"ו), כשבא ליסר הבריות מביא עליהם אש מן השמים, כמו שעשה לסדום, וכשבא להוריד המן מן השמים הנני ממטיר לכם לחם מן השמים (שמות ט"ז):}%
% ---------- פרק 19 | פסוק 25 ----------
 \textbf{כה} וַֽיַּהֲפֹךְ֙ אֶת־הֶעָרִ֣ים הָאֵ֔ל וְאֵ֖ת כׇּל־הַכִּכָּ֑ר וְאֵת֙ כׇּל־יֹשְׁבֵ֣י הֶעָרִ֔ים וְצֶ֖מַח הָאֲדָמָֽה׃ \Rashi{\textbf{ויהפך את הערים וגו'.}ארבעתם יושבות בסלע אחד והפכן מלמעלה למטה, שנאמר בחלמיש שלח ידו וגו' (איוב כ'):}%
% ---------- פרק 19 | פסוק 26 ----------
 \textbf{כו} וַתַּבֵּ֥ט אִשְׁתּ֖וֹ מֵאַחֲרָ֑יו וַתְּהִ֖י נְצִ֥יב מֶֽלַח׃ \Rashi{\textbf{ותבט אשתו מאחריו.}של לוט: ותהי נציב מלח. במלח חטאה ובמלח לקתה; אמר לה תני מעט מלח לאורחים הללו, אמרה לו אף המנהג הרע הזה אתה בא להנהיג במקום הזה (בראשית רבה):}%
% ---------- פרק 19 | פסוק 27 ----------
 \textbf{כז} וַיַּשְׁכֵּ֥ם אַבְרָהָ֖ם בַּבֹּ֑קֶר אֶ֨ל־הַמָּק֔וֹם אֲשֶׁר־עָ֥מַד שָׁ֖ם אֶת־פְּנֵ֥י יְהֹוָֽה׃ %
% ---------- פרק 19 | פסוק 28 ----------
 \textbf{כח} וַיַּשְׁקֵ֗ף עַל־פְּנֵ֤י סְדֹם֙ וַעֲמֹרָ֔ה וְעַֽל־כׇּל־פְּנֵ֖י אֶ֣רֶץ הַכִּכָּ֑ר וַיַּ֗רְא וְהִנֵּ֤ה עָלָה֙ קִיטֹ֣ר הָאָ֔רֶץ כְּקִיטֹ֖ר הַכִּבְשָֽׁן׃ \Rashi{\textbf{קיטור.}תמור של עשן, טורק"א בלע"ז: הכבשן. חפירה ששורפין בה את האבנים לסיד, וכן כל כבשן שבתורה:}%
% ---------- פרק 19 | פסוק 29 ----------
 \textbf{כט} וַיְהִ֗י בְּשַׁחֵ֤ת אֱלֹהִים֙ אֶת־עָרֵ֣י הַכִּכָּ֔ר וַיִּזְכֹּ֥ר אֱלֹהִ֖ים אֶת־אַבְרָהָ֑ם וַיְשַׁלַּ֤ח אֶת־לוֹט֙ מִתּ֣וֹךְ הַהֲפֵכָ֔ה בַּהֲפֹךְ֙ אֶת־הֶ֣עָרִ֔ים אֲשֶׁר־יָשַׁ֥ב בָּהֵ֖ן לֽוֹט׃ \Rashi{\textbf{ויזכור אלהים את אברהם.}מהו זכירתו של אברהם על לוט? נזכר שהיה לוט יודע ששרה אשתו של אברהם ושמע שאמר אברהם במצרים על שרה אחותי היא, ולא גלה הדבר, שהיה חס עליו, לפיכך חס הקב"ה עליו (בראשית רבה):}%
% ---------- פרק 19 | פסוק 30 ----------
 \textbf{ל} וַיַּ֩עַל֩ ל֨וֹט מִצּ֜וֹעַר וַיֵּ֣שֶׁב בָּהָ֗ר וּשְׁתֵּ֤י בְנֹתָיו֙ עִמּ֔וֹ כִּ֥י יָרֵ֖א לָשֶׁ֣בֶת בְּצ֑וֹעַר וַיֵּ֙שֶׁב֙ בַּמְּעָרָ֔ה ה֖וּא וּשְׁתֵּ֥י בְנֹתָֽיו׃ \Rashi{\textbf{כי ירא לשבת בצוער.}לפי שהיתה קרובה לסדום:}%
% ---------- פרק 19 | פסוק 31 ----------
 \textbf{לא} וַתֹּ֧אמֶר הַבְּכִירָ֛ה אֶל־הַצְּעִירָ֖ה אָבִ֣ינוּ זָקֵ֑ן וְאִ֨ישׁ אֵ֤ין בָּאָ֙רֶץ֙ לָב֣וֹא עָלֵ֔ינוּ כְּדֶ֖רֶךְ כׇּל־הָאָֽרֶץ׃ \Rashi{\textbf{אבינו זקן.}ואם לא עכשו אימתי? שמא ימות או יפסק מהוליד: ואיש אין בארץ. סבורות היו שכל העולם נחרב כמו בדור המבול: (בראשית רבה):}%
% ---------- פרק 19 | פסוק 32 ----------
 \textbf{לב} לְכָ֨ה נַשְׁקֶ֧ה אֶת־אָבִ֛ינוּ יַ֖יִן וְנִשְׁכְּבָ֣ה עִמּ֑וֹ וּנְחַיֶּ֥ה מֵאָבִ֖ינוּ זָֽרַע׃ %
% ---------- פרק 19 | פסוק 33 ----------
 \textbf{לג} וַתַּשְׁקֶ֧יןָ אֶת־אֲבִיהֶ֛ן יַ֖יִן בַּלַּ֣יְלָה ה֑וּא וַתָּבֹ֤א הַבְּכִירָה֙ וַתִּשְׁכַּ֣ב אֶת־אָבִ֔יהָ וְלֹֽא־יָדַ֥ע בְּשִׁכְבָ֖הּ וּבְקוּׄמָֽהּ׃ \Rashi{\textbf{ותשקין וגו'.}יין נזדמן להן במערה להוציא מהן שני אמות: ותשכב את אביה. ובצעירה כתיב ותשכב עמו? צעירה, לפי שלא פתחה בזנות אלא אחותה למדתה חסך עליה הכתוב ולא פרש גנותה; אבל בכירה שפתחה בזנות פרסמה הכתוב במפרש. ובקומה. של בכירה נקוד, לומר, שבקומה ידע, ואעפ"כ לא נשמר ליל שני מלשתות. אמר רבי לוי כל מי שהוא להוט אחר בולמוס של עריות לסוף מאכילים אותו מבשרו:}%
% ---------- פרק 19 | פסוק 34 ----------
 \textbf{לד} וַֽיְהִי֙ מִֽמׇּחֳרָ֔ת וַתֹּ֤אמֶר הַבְּכִירָה֙ אֶל־הַצְּעִירָ֔ה הֵן־שָׁכַ֥בְתִּי אֶ֖מֶשׁ אֶת־אָבִ֑י נַשְׁקֶ֨נּוּ יַ֜יִן גַּם־הַלַּ֗יְלָה וּבֹ֙אִי֙ שִׁכְבִ֣י עִמּ֔וֹ וּנְחַיֶּ֥ה מֵאָבִ֖ינוּ זָֽרַע׃ %
% ---------- פרק 19 | פסוק 35 ----------
 \textbf{לה} וַתַּשְׁקֶ֜יןָ גַּ֣ם בַּלַּ֧יְלָה הַה֛וּא אֶת־אֲבִיהֶ֖ן יָ֑יִן וַתָּ֤קׇם הַצְּעִירָה֙ וַתִּשְׁכַּ֣ב עִמּ֔וֹ וְלֹֽא־יָדַ֥ע בְּשִׁכְבָ֖הּ וּבְקֻמָֽהּ׃ %
% ---------- פרק 19 | פסוק 36 ----------
 \textbf{לו} וַֽתַּהֲרֶ֛יןָ שְׁתֵּ֥י בְנֽוֹת־ל֖וֹט מֵאֲבִיהֶֽן׃ \Rashi{\textbf{ותהרין וגו'.}אע"פ שאין האשה מתעברת מביאה ראשונה, אלו שלטו בעצמן והוציאו ערותן לחוץ ונתעברו מביאה ראשונה:}%
% ---------- פרק 19 | פסוק 37 ----------
 \textbf{לז} וַתֵּ֤לֶד הַבְּכִירָה֙ בֵּ֔ן וַתִּקְרָ֥א שְׁמ֖וֹ מוֹאָ֑ב ה֥וּא אֲבִֽי־מוֹאָ֖ב עַד־הַיּֽוֹם׃ \Rashi{\textbf{מואב.}זו שלא היתה צנועה פרשה שמאביה הוא, אבל צעירה קראתו בלשון נקיה, וקבלה שכר בימי משה, שנאמר בבני עמון אל תתגר בם (דברים ב') – כלל, ובמואב לא הזהיר אלא שלא ילחם בם, אבל לצערן התיר לו:}%
% ---------- פרק 19 | פסוק 38 ----------
 \textbf{לח} וְהַצְּעִירָ֤ה גַם־הִוא֙ יָ֣לְדָה בֵּ֔ן וַתִּקְרָ֥א שְׁמ֖וֹ בֶּן־עַמִּ֑י ה֛וּא אֲבִ֥י בְנֵֽי־עַמּ֖וֹן עַד־הַיּֽוֹם׃ \{ס\} %
% ---------- פרק 20 | פסוק 1 ----------
 \textbf{א} וַיִּסַּ֨ע מִשָּׁ֤ם אַבְרָהָם֙ אַ֣רְצָה הַנֶּ֔גֶב וַיֵּ֥שֶׁב בֵּין־קָדֵ֖שׁ וּבֵ֣ין שׁ֑וּר וַיָּ֖גׇר בִּגְרָֽר׃ \Rashi{\textbf{ויסע משם אברהם.}כשראה שחרבו הכרכים ופסקו העוברים והשבים, נסע לו משם. ד"א להתרחק מלוט, שיצא עליו שם רע, שבא על בנותיו:}%
% ---------- פרק 20 | פסוק 2 ----------
 \textbf{ב} וַיֹּ֧אמֶר אַבְרָהָ֛ם אֶל־שָׂרָ֥ה אִשְׁתּ֖וֹ אֲחֹ֣תִי הִ֑וא וַיִּשְׁלַ֗ח אֲבִימֶ֙לֶךְ֙ מֶ֣לֶךְ גְּרָ֔ר וַיִּקַּ֖ח אֶת־שָׂרָֽה׃ \Rashi{\textbf{ויאמר אברהם.}כאן לא נטל רשות אלא על כרחה שלא בטובתה, לפי שכבר לקחה לבית פרעה על ידי כן: אל שרה אשתו. על שרה אשתו, כיוצא בו אל הלקח ארון וגו' ואל מות חמיה (שמואל א ד'), שניהם בלשון על:}%
% ---------- פרק 20 | פסוק 3 ----------
 \textbf{ג} וַיָּבֹ֧א אֱלֹהִ֛ים אֶל־אֲבִימֶ֖לֶךְ בַּחֲל֣וֹם הַלָּ֑יְלָה וַיֹּ֣אמֶר ל֗וֹ הִנְּךָ֥ מֵת֙ עַל־הָאִשָּׁ֣ה אֲשֶׁר־לָקַ֔חְתָּ וְהִ֖וא בְּעֻ֥לַת בָּֽעַל׃ %
% ---------- פרק 20 | פסוק 4 ----------
 \textbf{ד} וַאֲבִימֶ֕לֶךְ לֹ֥א קָרַ֖ב אֵלֶ֑יהָ וַיֹּאמַ֕ר אֲדֹנָ֕י הֲג֥וֹי גַּם־צַדִּ֖יק תַּהֲרֹֽג׃ \Rashi{\textbf{לא קרב אליה.}המלאך מנעו, כמו שנאמר לא נתתיך לנגע אליה: הגוי גם צדיק תהרג. אף אם הוא צדיק תהרגנו, שמא כך דרכך לאבד האמות חנם? כך עשית לדור המבול ולדור הפלגה, אף אני אומר שהרגתם על לא דבר, כמו שאתה אומר להרגני:}%
% ---------- פרק 20 | פסוק 5 ----------
 \textbf{ה} הֲלֹ֨א ה֤וּא אָֽמַר־לִי֙ אֲחֹ֣תִי הִ֔וא וְהִֽיא־גַם־הִ֥וא אָֽמְרָ֖ה אָחִ֣י ה֑וּא בְּתׇם־לְבָבִ֛י וּבְנִקְיֹ֥ן כַּפַּ֖י עָשִׂ֥יתִי זֹֽאת׃ \Rashi{\textbf{גם הוא.}לרבות עבדים וגמלים וחמרים שלה, את כלם שאלתי ואמרו לי אחיה הוא: בתם לבבי. שלא דמיתי לחטא: ובנקיון כפי. נקי אני מן החטא, שלא נגעתי בה:}%
% ---------- פרק 20 | פסוק 6 ----------
 \textbf{ו} וַיֹּ֩אמֶר֩ אֵלָ֨יו הָֽאֱלֹהִ֜ים בַּחֲלֹ֗ם גַּ֣ם אָנֹכִ֤י יָדַ֙עְתִּי֙ כִּ֤י בְתׇם־לְבָבְךָ֙ עָשִׂ֣יתָ זֹּ֔את וָאֶחְשֹׂ֧ךְ גַּם־אָנֹכִ֛י אֽוֹתְךָ֖ מֵחֲטוֹ־לִ֑י עַל־כֵּ֥ן לֹא־נְתַתִּ֖יךָ לִנְגֹּ֥עַ אֵלֶֽיהָ׃ \Rashi{\textbf{ידעתי כי בתם לבבך וגו'.}אמת שלא דמית מתחלה לחטא אבל נקיון כפים אין כאן: לא נתתיך. לא ממך היה שלא נגעת בה, אלא חשכתי אני אותך מחטא ולא נתתי לך כח, וכן ולא נתנו אלהים (בראשית ל״א:ז׳), וכן ולא נתנו אביה לבוא (שופטים ט"ו):}%
% ---------- פרק 20 | פסוק 7 ----------
 \textbf{ז} וְעַתָּ֗ה הָשֵׁ֤ב אֵֽשֶׁת־הָאִישׁ֙ כִּֽי־נָבִ֣יא ה֔וּא וְיִתְפַּלֵּ֥ל בַּֽעַדְךָ֖ וֶֽחְיֵ֑ה וְאִם־אֵֽינְךָ֣ מֵשִׁ֔יב דַּ֚ע כִּי־מ֣וֹת תָּמ֔וּת אַתָּ֖ה וְכׇל־אֲשֶׁר־לָֽךְ׃ \Rashi{\textbf{השב אשת האיש.}ואל תהא סבור שמא תתגנה בעיניו ולא יקבלנה או ישנאך ולא יתפלל עליך: כי נביא הוא. ויודע שלא נגעת בה, לפיכך ויתפלל בעדך (בבא קמא צ"ב):}%
% ---------- פרק 20 | פסוק 8 ----------
 \textbf{ח} וַיַּשְׁכֵּ֨ם אֲבִימֶ֜לֶךְ בַּבֹּ֗קֶר וַיִּקְרָא֙ לְכׇל־עֲבָדָ֔יו וַיְדַבֵּ֛ר אֶת־כׇּל־הַדְּבָרִ֥ים הָאֵ֖לֶּה בְּאׇזְנֵיהֶ֑ם וַיִּֽירְא֥וּ הָאֲנָשִׁ֖ים מְאֹֽד׃ %
% ---------- פרק 20 | פסוק 9 ----------
 \textbf{ט} וַיִּקְרָ֨א אֲבִימֶ֜לֶךְ לְאַבְרָהָ֗ם וַיֹּ֨אמֶר ל֜וֹ מֶֽה־עָשִׂ֤יתָ לָּ֙נוּ֙ וּמֶֽה־חָטָ֣אתִי לָ֔ךְ כִּֽי־הֵבֵ֧אתָ עָלַ֛י וְעַל־מַמְלַכְתִּ֖י חֲטָאָ֣ה גְדֹלָ֑ה מַעֲשִׂים֙ אֲשֶׁ֣ר לֹא־יֵֽעָשׂ֔וּ עָשִׂ֖יתָ עִמָּדִֽי׃ \Rashi{\textbf{מעשים אשר לא יעשו.}מכה אשר לא הרגלה לבא על בריה באה לנו על ידך, עצירת כל נקבים של זרע ושל קטנים ורעי ואזנים וחטם:}%
% ---------- פרק 20 | פסוק 10 ----------
 \textbf{י} וַיֹּ֥אמֶר אֲבִימֶ֖לֶךְ אֶל־אַבְרָהָ֑ם מָ֣ה רָאִ֔יתָ כִּ֥י עָשִׂ֖יתָ אֶת־הַדָּבָ֥ר הַזֶּֽה׃ %
% ---------- פרק 20 | פסוק 11 ----------
 \textbf{יא} וַיֹּ֙אמֶר֙ אַבְרָהָ֔ם כִּ֣י אָמַ֗רְתִּי רַ֚ק אֵין־יִרְאַ֣ת אֱלֹהִ֔ים בַּמָּק֖וֹם הַזֶּ֑ה וַהֲרָג֖וּנִי עַל־דְּבַ֥ר אִשְׁתִּֽי׃ \Rashi{\textbf{רק אין יראת אלהים.}אכסנאי שבא לעיר, על עסקי אכילה ושתיה שואלין אותו או על עסקי אשתו שואלין אותו? אשתך היא או אחותך היא? (בבא קמא שם):}%
% ---------- פרק 20 | פסוק 12 ----------
 \textbf{יב} וְגַם־אׇמְנָ֗ה אֲחֹתִ֤י בַת־אָבִי֙ הִ֔וא אַ֖ךְ לֹ֣א בַת־אִמִּ֑י וַתְּהִי־לִ֖י לְאִשָּֽׁה׃ \Rashi{\textbf{אחתי בת אבי הוא.}ובת אב מתרת לבן נח, שאין אבהות לגוי, וכדי לאמת דבריו השיבם כן. ואם תאמר והלא בת אחיו היתה? בני בנים הרי הן כבנים, והרי היא בתו של תרח וכך הוא אומר ללוט כי אנשים אחים אנחנו (בראשית י״ג:ח׳): אך לא בת אמי. הרן מאם אחרת היה:}%
% ---------- פרק 20 | פסוק 13 ----------
 \textbf{יג} וַיְהִ֞י כַּאֲשֶׁ֧ר הִתְע֣וּ אֹתִ֗י אֱלֹהִים֮ מִבֵּ֣ית אָבִי֒ וָאֹמַ֣ר לָ֔הּ זֶ֣ה חַסְדֵּ֔ךְ אֲשֶׁ֥ר תַּעֲשִׂ֖י עִמָּדִ֑י אֶ֤ל כׇּל־הַמָּקוֹם֙ אֲשֶׁ֣ר נָב֣וֹא שָׁ֔מָּה אִמְרִי־לִ֖י אָחִ֥י הֽוּא׃ \Rashi{\textbf{ויהי כאשר התעו וגו'.}אנקלוס תרגם מה שתרגם; ויש לישבו עוד דבר דבור על אפניו, כשהוציאני הקב"ה מבית אבי להיות משוטט ונד ממקום למקום, ידעתי שאעבר במקום רשעים ואמר לה זה חסדך: כאשר התעו. לשון רבים. ואל תתמה, כי הרבה מקומות לשון אלהות ולשון מרות קרוי בלשון רבים, אשר הלכו אלהים (שמואל ב ז'), אלהים חיים (דברים ה'), אלהים קדשים (יהושע כ"ד), וכל לשון אדנות לשון רבים; וכן ויקח אדני יוסף (בראשית ל"ט), ואדני האדנים (דברים י'), אדני הארץ (בראשית מ"ב), וכן בעליו עמו (שמות כ"ב), והועד בבעליו (שם כ"א). ואם תאמר מהו לשון התעו? כל הגולה ממקומו ואינו מישב קרוי תועה, כמו ותלך ותתע (בראשית כ"א), תעיתי כשה אבד (תהילים קי"ט), יתעו לבלי אכל (איוב ל"ח) – יצאו ויתעו לבקש אכלם: אמרי לי. עלי, וכן וישאלו אנשי המקום לאשתו (בראשית כ"ו), על אשתו, וכן ואמר פרעה לבני ישראל (שמות י"ד), כמו על בני ישראל; פן יאמרו לי אשה הרגתהו (שופטים ט'):}%
% ---------- פרק 20 | פסוק 14 ----------
 \textbf{יד} וַיִּקַּ֨ח אֲבִימֶ֜לֶךְ צֹ֣אן וּבָקָ֗ר וַעֲבָדִים֙ וּשְׁפָחֹ֔ת וַיִּתֵּ֖ן לְאַבְרָהָ֑ם וַיָּ֣שֶׁב ל֔וֹ אֵ֖ת שָׂרָ֥ה אִשְׁתּֽוֹ׃ \Rashi{\textbf{ויתן לאברהם.}כדי שיתפיס ויתפלל עליו:}%
% ---------- פרק 20 | פסוק 15 ----------
 \textbf{טו} וַיֹּ֣אמֶר אֲבִימֶ֔לֶךְ הִנֵּ֥ה אַרְצִ֖י לְפָנֶ֑יךָ בַּטּ֥וֹב בְּעֵינֶ֖יךָ שֵֽׁב׃ \Rashi{\textbf{הנה ארצי לפניך.}אבל פרעה אמר לו הנה אשתך קח ולך, לפי שנתירא, שהמצריים שטופי זמה:}%
% ---------- פרק 20 | פסוק 16 ----------
 \textbf{טז} וּלְשָׂרָ֣ה אָמַ֗ר הִנֵּ֨ה נָתַ֜תִּי אֶ֤לֶף כֶּ֙סֶף֙ לְאָחִ֔יךְ הִנֵּ֤ה הוּא־לָךְ֙ כְּס֣וּת עֵינַ֔יִם לְכֹ֖ל אֲשֶׁ֣ר אִתָּ֑ךְ וְאֵ֥ת כֹּ֖ל וְנֹכָֽחַת׃ \Rashi{\textbf{ולשרה אמר.}אבימלך לכבודה כדי לפיסה, הנה עשיתי לך כבוד זה, נתתי ממון לאחיך שאמרת עליו אחי הוא, הנה הממון והכבוד הזה לך כסות עינים: לכל אשר אתך. יכסו עיניהם שלא יקלוך, שאלו השיבותיך ריקנית, יש להם לומר לאחר שנתעלל בה החזירה; עכשו שהצרכתי לבזבז ממון ולפיסך, יודעים יהיו שעל כרחי השיבותיך ועל ידי נס: ואת כל. ועם כל באי עולם: ונכחת. יהי לך פתחון פה להתוכח ולהראות דברים נכרים הללו; ולשון הוכחה בכל מקום ברור דברים ובלע"ז אשפרו"ביר, ואנקלוס תרגם בפנים אחרים, ולשון המקרא כך הוא נופל על התרגום, הנה לך כסות של כבוד על העינים שלי ששלטו בך ובכל אשר אתך, ועל כן תרגמו וחזית יתיך וית כל דעמך. ויש מ"א אבל ישוב לשון המקרא פרשתי:}%
% ---------- פרק 20 | פסוק 17 ----------
 \textbf{יז} וַיִּתְפַּלֵּ֥ל אַבְרָהָ֖ם אֶל־הָאֱלֹהִ֑ים וַיִּרְפָּ֨א אֱלֹהִ֜ים אֶת־אֲבִימֶ֧לֶךְ וְאֶת־אִשְׁתּ֛וֹ וְאַמְהֹתָ֖יו וַיֵּלֵֽדוּ׃ \Rashi{\textbf{וילדו.}כתרגומו ואתרוחו, נפתחו נקביהם והוציאו, והיא לדה שלהם:}%
% ---------- פרק 20 | פסוק 18 ----------
 \textbf{יח} כִּֽי־עָצֹ֤ר עָצַר֙ יְהֹוָ֔ה בְּעַ֥ד כׇּל־רֶ֖חֶם לְבֵ֣ית אֲבִימֶ֑לֶךְ עַל־דְּבַ֥ר שָׂרָ֖ה אֵ֥שֶׁת אַבְרָהָֽם׃ \{ס\} \Rashi{\textbf{בעד כל רחם.}כנגד כל פתח: על דבר שרה. על פי דבורה של שרה:}%
% ---------- פרק 21 | פסוק 1 ----------
 \textbf{א} וַֽיהֹוָ֛ה פָּקַ֥ד אֶת־שָׂרָ֖ה כַּאֲשֶׁ֣ר אָמָ֑ר וַיַּ֧עַשׂ יְהֹוָ֛ה לְשָׂרָ֖ה כַּאֲשֶׁ֥ר דִּבֵּֽר׃ \Rashi{\textbf{וה' פקד את שרה וגו'.}סמך פרשה זו ללמדך, שכל המבקש רחמים על חברו והוא צריך לאותו דבר, הוא נענה תחלה (בבא קמא צ"ב), שנאמר ויתפלל וגו', וסמיך ליה וה' פקד את שרה, שפקדה כבר קדם שרפא את אבימלך: פקד את שרה כאשר אמר. בהריון: כאשר דבר. בלדה. והיכן היא אמירה, והיכן הוא דבור? אמירה ויאמר אלהים אבל שרה אשתך וגו '(בראשית י"ז), דבור היה דבר ה' אל אברם בברית בין הבתרים, ושם נאמר לא יירשך זה וגו' (שם ט"ו), והביא היורש משרה: ויעש ה' לשרה כאשר דבר. לאברהם:}%
% ---------- פרק 21 | פסוק 2 ----------
 \textbf{ב} וַתַּ֩הַר֩ וַתֵּ֨לֶד שָׂרָ֧ה לְאַבְרָהָ֛ם בֵּ֖ן לִזְקֻנָ֑יו לַמּוֹעֵ֕ד אֲשֶׁר־דִּבֶּ֥ר אֹת֖וֹ אֱלֹהִֽים׃ \Rashi{\textbf{למועד אשר דבר אתו.}די מליל יתה, את המועד דבר וקבע כשאמר לו למועד אשוב אליך (שם י"ח), שרט לו שריטה בכתל, אמר לו כשתגיע חמה לשריטה זו בשנה האחרת, תלד: לזקניו. היה זיו איקונין שלו דומה לו:}%
% ---------- פרק 21 | פסוק 3 ----------
 \textbf{ג} וַיִּקְרָ֨א אַבְרָהָ֜ם אֶֽת־שֶׁם־בְּנ֧וֹ הַנּֽוֹלַד־ל֛וֹ אֲשֶׁר־יָלְדָה־לּ֥וֹ שָׂרָ֖ה יִצְחָֽק׃ %
% ---------- פרק 21 | פסוק 4 ----------
 \textbf{ד} וַיָּ֤מׇל אַבְרָהָם֙ אֶת־יִצְחָ֣ק בְּנ֔וֹ בֶּן־שְׁמֹנַ֖ת יָמִ֑ים כַּאֲשֶׁ֛ר צִוָּ֥ה אֹת֖וֹ אֱלֹהִֽים׃ %
% ---------- פרק 21 | פסוק 5 ----------
 \textbf{ה} וְאַבְרָהָ֖ם בֶּן־מְאַ֣ת שָׁנָ֑ה בְּהִוָּ֣לֶד ל֔וֹ אֵ֖ת יִצְחָ֥ק בְּנֽוֹ׃ %
% ---------- פרק 21 | פסוק 6 ----------
 \textbf{ו} וַתֹּ֣אמֶר שָׂרָ֔ה צְחֹ֕ק עָ֥שָׂה לִ֖י אֱלֹהִ֑ים כׇּל־הַשֹּׁמֵ֖עַ יִֽצְחַק־לִֽי׃ \Rashi{\textbf{יצחק לי.}ישמח עלי. ומ"א הרבה עקרות נפקדו עמה, הרבה חולים נתרפאו בו ביום, הרבה תפלות נענו עמה, ורב שחוק היה בעולם:}%
% ---------- פרק 21 | פסוק 7 ----------
 \textbf{ז} וַתֹּ֗אמֶר מִ֤י מִלֵּל֙ לְאַבְרָהָ֔ם הֵינִ֥יקָה בָנִ֖ים שָׂרָ֑ה כִּֽי־יָלַ֥דְתִּי בֵ֖ן לִזְקֻנָֽיו׃ \Rashi{\textbf{מי מלל לאברהם.}לשון שבח וחשיבות, ראו, מי הוא וכמה הוא שומר הבטחתו הקב"ה מבטיח ועושה. מלל. שנה הכתוב ולא אמר דבר – גימטריא שלו ק', כלומר לסוף מאה לאברהם: היניקה בנים שרה. ומהו בנים לשון רבים? ביום המשתה הביאו השרות בניהן עמהן והניקה אותם, שהיו אומרות לא ילדה שרה, אלא אסופי הביאה מן השוק:}%
% ---------- פרק 21 | פסוק 8 ----------
 \textbf{ח} וַיִּגְדַּ֥ל הַיֶּ֖לֶד וַיִּגָּמַ֑ל וַיַּ֤עַשׂ אַבְרָהָם֙ מִשְׁתֶּ֣ה גָד֔וֹל בְּי֖וֹם הִגָּמֵ֥ל אֶת־יִצְחָֽק׃ \Rashi{\textbf{ויגמל.}לסוף כ"ד חדש (גיטין ע"ה): משתה גדול. שהיו שם גדולי הדור שם ועבר ואבימלך (בראשית רבה):}%
% ---------- פרק 21 | פסוק 9 ----------
 \textbf{ט} וַתֵּ֨רֶא שָׂרָ֜ה אֶֽת־בֶּן־הָגָ֧ר הַמִּצְרִ֛ית אֲשֶׁר־יָלְדָ֥ה לְאַבְרָהָ֖ם מְצַחֵֽק׃ \Rashi{\textbf{מצחק.}לשון ע"ז, כמו שנאמר ויקמו לצחק (שמות ל"ב), ד"א לשון גלוי עריות, כמה דתימא לצחק בי (בראשית ל"ט), ד"א לשון רציחה, כמו יקומו נא הנערים וישחקו לפנינו וגו' (שמואל ב ב'), שהיה מריב עם יצחק על הירשה ואומר אני בכור ונוטל פי שנים, ויוצאים בשדה ונוטל קשתו ויורה בו חצים, כמה דתימא כמתלהלה הירה זקים וגו' ואמר הלא משחק אני (משלי כ"ו):}%
% ---------- פרק 21 | פסוק 10 ----------
 \textbf{י} וַתֹּ֙אמֶר֙ לְאַבְרָהָ֔ם גָּרֵ֛שׁ הָאָמָ֥ה הַזֹּ֖את וְאֶת־בְּנָ֑הּ כִּ֣י לֹ֤א יִירַשׁ֙ בֶּן־הָאָמָ֣ה הַזֹּ֔את עִם־בְּנִ֖י עִם־יִצְחָֽק׃ \Rashi{\textbf{עם בני עם יצחק.}מכיון שהוא בני אפלו אם אינו הגון כיצחק, או הגון כיצחק אפלו אינו בני, אין זה כדאי לירש עמו, ק"ו עם בני עם יצחק, ששתיהן בו (בראשית רבה):}%
% ---------- פרק 21 | פסוק 11 ----------
 \textbf{יא} וַיֵּ֧רַע הַדָּבָ֛ר מְאֹ֖ד בְּעֵינֵ֣י אַבְרָהָ֑ם עַ֖ל אוֹדֹ֥ת בְּנֽוֹ׃ \Rashi{\textbf{על אודות בנו.}ששמע שיצא לתרבות רעה (בראשית רבה), ופשוטו על שאמרה לו לשלחו:}%
% ---------- פרק 21 | פסוק 12 ----------
 \textbf{יב} וַיֹּ֨אמֶר אֱלֹהִ֜ים אֶל־אַבְרָהָ֗ם אַל־יֵרַ֤ע בְּעֵינֶ֙יךָ֙ עַל־הַנַּ֣עַר וְעַל־אֲמָתֶ֔ךָ כֹּל֩ אֲשֶׁ֨ר תֹּאמַ֥ר אֵלֶ֛יךָ שָׂרָ֖ה שְׁמַ֣ע בְּקֹלָ֑הּ כִּ֣י בְיִצְחָ֔ק יִקָּרֵ֥א לְךָ֖ זָֽרַע׃ \Rashi{\textbf{שמע בקולה.}למדנו שהיה אברהם טפל לשרה בנביאות:}%
% ---------- פרק 21 | פסוק 13 ----------
 \textbf{יג} וְגַ֥ם אֶת־בֶּן־הָאָמָ֖ה לְג֣וֹי אֲשִׂימֶ֑נּוּ כִּ֥י זַרְעֲךָ֖ הֽוּא׃ %
% ---------- פרק 21 | פסוק 14 ----------
 \textbf{יד} וַיַּשְׁכֵּ֣ם אַבְרָהָ֣ם ׀ בַּבֹּ֡קֶר וַיִּֽקַּֽח־לֶ֩חֶם֩ וְחֵ֨מַת מַ֜יִם וַיִּתֵּ֣ן אֶל־הָ֠גָ֠ר שָׂ֧ם עַל־שִׁכְמָ֛הּ וְאֶת־הַיֶּ֖לֶד וַֽיְשַׁלְּחֶ֑הָ וַתֵּ֣לֶךְ וַתֵּ֔תַע בְּמִדְבַּ֖ר בְּאֵ֥ר שָֽׁבַע׃ \Rashi{\textbf{לחם וחמת מים.}ולא כסף וזהב, לפי שהיה שונאו על שיצא לתרבות רעה: ואת הילד. אף הילד שם על שכמה, שהכניסה בו שרה עין רעה ואחזתו חמה ולא יכול לילך ברגליו (בראשית רבה): ותלך ותתע. חזרה לגלולי בית אביה (פרקי דרבי אליעזר פ' ל'):}%
% ---------- פרק 21 | פסוק 15 ----------
 \textbf{טו} וַיִּכְל֥וּ הַמַּ֖יִם מִן־הַחֵ֑מֶת וַתַּשְׁלֵ֣ךְ אֶת־הַיֶּ֔לֶד תַּ֖חַת אַחַ֥ד הַשִּׂיחִֽם׃ \Rashi{\textbf{ויכלו המים.}לפי שדרך חולים לשתות הרבה:}%
% ---------- פרק 21 | פסוק 16 ----------
 \textbf{טז} וַתֵּ֩לֶךְ֩ וַתֵּ֨שֶׁב לָ֜הּ מִנֶּ֗גֶד הַרְחֵק֙ כִּמְטַחֲוֵ֣י קֶ֔שֶׁת כִּ֣י אָֽמְרָ֔ה אַל־אֶרְאֶ֖ה בְּמ֣וֹת הַיָּ֑לֶד וַתֵּ֣שֶׁב מִנֶּ֔גֶד וַתִּשָּׂ֥א אֶת־קֹלָ֖הּ וַתֵּֽבְךְּ׃ \Rashi{\textbf{מנגד.}מרחוק: כמטחוי קשת. כשתי טיחות (בראשית רבה), והוא לשון ירית חץ, בלשון משנה שהטיח באשתו, על שם שהזרע יורה כחץ; ואם תאמר היה לו לכתב כמטחי קשת, משפט הוי"ו לכנס לכאן, כמו בחגוי הסלע (שיר השירים ב'), מגזרת והיתה אדמת יהודה למצרים לחגא (ישעיהו י"ט), ומגזרת יחוגו וינועו כשכור (תהילים ק"ז), וכן קצוי ארץ (תהילים ס"ה), מגזרת קצה: ותשב מנגד. כיון שקרב למות הוסיפה להתרחק:}%
% ---------- פרק 21 | פסוק 17 ----------
 \textbf{יז} וַיִּשְׁמַ֣ע אֱלֹהִים֮ אֶת־ק֣וֹל הַנַּ֒עַר֒ וַיִּקְרָא֩ מַלְאַ֨ךְ אֱלֹהִ֤ים ׀ אֶל־הָגָר֙ מִן־הַשָּׁמַ֔יִם וַיֹּ֥אמֶר לָ֖הּ מַה־לָּ֣ךְ הָגָ֑ר אַל־תִּ֣ירְאִ֔י כִּֽי־שָׁמַ֧ע אֱלֹהִ֛ים אֶל־ק֥וֹל הַנַּ֖עַר בַּאֲשֶׁ֥ר הוּא־שָֽׁם׃ \Rashi{\textbf{את קול הנער.}מכאן שיפה תפלת החולה מתפלת אחרים עליו והיא קודמת להתקבל (בראשית רבה): באשר הוא שם. לפי מעשים שהוא עושה עכשו הוא נדון, ולא לפי מה שהוא עתיד לעשות (ראש השנה ט"ז), לפי שהיו מלאכי השרת מקטרגים ואומרים רבונו של עולם, מי שעתיד זרעו להמית בניך בצמא אתה מעלה לו באר, והוא משיבם עכשו מה הוא, צדיק או רשע? אמרו לו צדיק, אמר להם לפי מעשיו של עכשו אני דנו, וזהו באשר הוא שם. והיכן המית את ישראל בצמא? כשהגלם נבוכדנצר, שנאמר משא בערב וגו' לקראת צמא התיו מים וגו' (ישעיהו כ"א), כשהיו מוליכין אותם אצל ערביים, היו ישראל אומרים לשוביהם בבקשה מכם הוליכונו אצל בני דודנו ישמעאל וירחמו עלינו, שנאמר ארחות דדנים (שם) – אל תקרי דדנים אלא דודים – ואלו יוצאים לקראתם ומביאין להם בשר ודג מלוח ונודות נפוחים, כסבורים ישראל שהם מלאים מים, וכשמכניסו לתוך פיו ופותחו, הרוח נכנסת בגופו ומת:}%
% ---------- פרק 21 | פסוק 18 ----------
 \textbf{יח} ק֚וּמִי שְׂאִ֣י אֶת־הַנַּ֔עַר וְהַחֲזִ֥יקִי אֶת־יָדֵ֖ךְ בּ֑וֹ כִּֽי־לְג֥וֹי גָּד֖וֹל אֲשִׂימֶֽנּוּ׃ %
% ---------- פרק 21 | פסוק 19 ----------
 \textbf{יט} וַיִּפְקַ֤ח אֱלֹהִים֙ אֶת־עֵינֶ֔יהָ וַתֵּ֖רֶא בְּאֵ֣ר מָ֑יִם וַתֵּ֜לֶךְ וַתְּמַלֵּ֤א אֶת־הַחֵ֙מֶת֙ מַ֔יִם וַתַּ֖שְׁקְ אֶת־הַנָּֽעַר׃ %
% ---------- פרק 21 | פסוק 20 ----------
 \textbf{כ} וַיְהִ֧י אֱלֹהִ֛ים אֶת־הַנַּ֖עַר וַיִּגְדָּ֑ל וַיֵּ֙שֶׁב֙ בַּמִּדְבָּ֔ר וַיְהִ֖י רֹבֶ֥ה קַשָּֽׁת׃ \Rashi{\textbf{רבה קשת.}יורה חצים בקשת: קשת. על שם האמנות, כמו חמר, גמל, ציד, לפיכך השי"ן מדגשת. היה יושב במדבר ומלסטם את העוברים, הוא שנאמר ידו בכל וגומר:}%
% ---------- פרק 21 | פסוק 21 ----------
 \textbf{כא} וַיֵּ֖שֶׁב בְּמִדְבַּ֣ר פָּארָ֑ן וַתִּֽקַּֽח־ל֥וֹ אִמּ֛וֹ אִשָּׁ֖ה מֵאֶ֥רֶץ מִצְרָֽיִם׃ \{פ\} \Rashi{\textbf{מארץ מצרים.}ממקום גדוליה, שנאמר ולה שפחה מצרית וגו' (בראשית ט״ז:א׳), הינו דאמרי אינשי זרק חוטרא לאוירא, אעיקריה קאי:}%
% ---------- פרק 21 | פסוק 22 ----------
 \textbf{כב} וַֽיְהִי֙ בָּעֵ֣ת הַהִ֔וא וַיֹּ֣אמֶר אֲבִימֶ֗לֶךְ וּפִיכֹל֙ שַׂר־צְבָא֔וֹ אֶל־אַבְרָהָ֖ם לֵאמֹ֑ר אֱלֹהִ֣ים עִמְּךָ֔ בְּכֹ֥ל אֲשֶׁר־אַתָּ֖ה עֹשֶֽׂה׃ \Rashi{\textbf{אלהים עמך.}לפי שראו שיצא משכונת סדום לשלום, ועם המלכים נלחם ונפלו בידו, ונפקדה אשתו לזקוניו:}%
% ---------- פרק 21 | פסוק 23 ----------
 \textbf{כג} וְעַתָּ֗ה הִשָּׁ֨בְעָה לִּ֤י בֵֽאלֹהִים֙ הֵ֔נָּה אִם־תִּשְׁקֹ֣ר לִ֔י וּלְנִינִ֖י וּלְנֶכְדִּ֑י כַּחֶ֜סֶד אֲשֶׁר־עָשִׂ֤יתִי עִמְּךָ֙ תַּעֲשֶׂ֣ה עִמָּדִ֔י וְעִם־הָאָ֖רֶץ אֲשֶׁר־גַּ֥רְתָּה בָּֽהּ׃ \Rashi{\textbf{ולניני ולנכדי.}עד כאן רחמי האב על הבן: כחסד אשר עשיתי עמך תעשה עמדי. שאמרתי לך הנה ארצי לפניך (בראשית רבה):}%
% ---------- פרק 21 | פסוק 24 ----------
 \textbf{כד} וַיֹּ֙אמֶר֙ אַבְרָהָ֔ם אָנֹכִ֖י אִשָּׁבֵֽעַ׃ %
% ---------- פרק 21 | פסוק 25 ----------
 \textbf{כה} וְהוֹכִ֥חַ אַבְרָהָ֖ם אֶת־אֲבִימֶ֑לֶךְ עַל־אֹדוֹת֙ בְּאֵ֣ר הַמַּ֔יִם אֲשֶׁ֥ר גָּזְל֖וּ עַבְדֵ֥י אֲבִימֶֽלֶךְ׃ \Rashi{\textbf{והוכח.}נתוכח עמו על כך:}%
% ---------- פרק 21 | פסוק 26 ----------
 \textbf{כו} וַיֹּ֣אמֶר אֲבִימֶ֔לֶךְ לֹ֣א יָדַ֔עְתִּי מִ֥י עָשָׂ֖ה אֶת־הַדָּבָ֣ר הַזֶּ֑ה וְגַם־אַתָּ֞ה לֹא־הִגַּ֣דְתָּ לִּ֗י וְגַ֧ם אָנֹכִ֛י לֹ֥א שָׁמַ֖עְתִּי בִּלְתִּ֥י הַיּֽוֹם׃ %
% ---------- פרק 21 | פסוק 27 ----------
 \textbf{כז} וַיִּקַּ֤ח אַבְרָהָם֙ צֹ֣אן וּבָקָ֔ר וַיִּתֵּ֖ן לַאֲבִימֶ֑לֶךְ וַיִּכְרְת֥וּ שְׁנֵיהֶ֖ם בְּרִֽית׃ %
% ---------- פרק 21 | פסוק 28 ----------
 \textbf{כח} וַיַּצֵּ֣ב אַבְרָהָ֗ם אֶת־שֶׁ֛בַע כִּבְשֹׂ֥ת הַצֹּ֖אן לְבַדְּהֶֽן׃ %
% ---------- פרק 21 | פסוק 29 ----------
 \textbf{כט} וַיֹּ֥אמֶר אֲבִימֶ֖לֶךְ אֶל־אַבְרָהָ֑ם מָ֣ה הֵ֗נָּה שֶׁ֤בַע כְּבָשֹׂת֙ הָאֵ֔לֶּה אֲשֶׁ֥ר הִצַּ֖בְתָּ לְבַדָּֽנָה׃ %
% ---------- פרק 21 | פסוק 30 ----------
 \textbf{ל} וַיֹּ֕אמֶר כִּ֚י אֶת־שֶׁ֣בַע כְּבָשֹׂ֔ת תִּקַּ֖ח מִיָּדִ֑י בַּעֲבוּר֙ תִּֽהְיֶה־לִּ֣י לְעֵדָ֔ה כִּ֥י חָפַ֖רְתִּי אֶת־הַבְּאֵ֥ר הַזֹּֽאת׃ \Rashi{\textbf{בעבור תהיה לי.}זאת: לעדה. לשון עדות של נקבה, כמו ועדה המצבה (בראשית ל"א): כי חפרתי את הבאר. מריבים היו עליה רועי אבימלך ואומרים אנחנו חפרנוה, אמרו ביניהם כל מי שיתראה על הבאר ויעלו המים לקראתו, שלו היא, ועלו לקראת אברהם:}%
% ---------- פרק 21 | פסוק 31 ----------
 \textbf{לא} עַל־כֵּ֗ן קָרָ֛א לַמָּק֥וֹם הַה֖וּא בְּאֵ֣ר שָׁ֑בַע כִּ֛י שָׁ֥ם נִשְׁבְּע֖וּ שְׁנֵיהֶֽם׃ %
% ---------- פרק 21 | פסוק 32 ----------
 \textbf{לב} וַיִּכְרְת֥וּ בְרִ֖ית בִּבְאֵ֣ר שָׁ֑בַע וַיָּ֣קׇם אֲבִימֶ֗לֶךְ וּפִיכֹל֙ שַׂר־צְבָא֔וֹ וַיָּשֻׁ֖בוּ אֶל־אֶ֥רֶץ פְּלִשְׁתִּֽים׃ %
% ---------- פרק 21 | פסוק 33 ----------
 \textbf{לג} וַיִּטַּ֥ע אֶ֖שֶׁל בִּבְאֵ֣ר שָׁ֑בַע וַיִּ֨קְרָא־שָׁ֔ם בְּשֵׁ֥ם יְהֹוָ֖ה אֵ֥ל עוֹלָֽם׃ \Rashi{\textbf{אשל.}רב ושמואל, חד אמר פרדס להביא ממנו פרות לאורחים בסעודה, וחד אמר פנדק לאכסניא ובו כל מיני פרות. ומצינו לשון נטיעה באהלים, שנאמר ויטע אהלי אפדנו (דניאל י"א): ויקרא שם וגו'. על ידי אותו אשל נקרא שמו של הקב"ה אלוה לכל העולם, לאחר שאוכלים ושותים אומר להם ברכו למי שאכלתם משלו, סבורים אתם שמשלי אכלתם? משל מי שאמר והיה העולם אכלתם (סוטה י'):}%
% ---------- פרק 21 | פסוק 34 ----------
 \textbf{לד} וַיָּ֧גׇר אַבְרָהָ֛ם בְּאֶ֥רֶץ פְּלִשְׁתִּ֖ים יָמִ֥ים רַבִּֽים׃ \{פ\} \Rashi{\textbf{ימים רבים.}מרבים על של חברון. בחברון עשה כ"ה שנה וכאן כ"ו, שהרי בן ע"ה שנה היה בצאתו מחרן, אותה שנה ויבא וישב באלני ממרא, שלא מצינו קדם לכן שנתישב אלא שם, שבכל מקומותיו היה כאורח חונה ונוסע והולך, שנאמר ויעבר אברם, ויעתק משם, ויהי רעב בארץ וירד אברם מצרימה (בראשית י"ב), ובמצרים לא עשה אלא שלשה חדשים שהרי שלחו פרעה מיד, וילך למסעיו עד ויבא וישב באלני ממרא אשר בחברון (שם י"ג), שם ישב עד שנהפכה סדום, מיד ויסע משם אברהם מפני בושה של לוט ובא לארץ פלשתים, ובן צ"ט שנה היה, שהרי בשלישי למילתו באו אצלו המלאכים, הרי כ"ה שנה, וכאן כתיב ימים רבים, מרבים על הראשונים, ולא בא הכתוב לסתם אלא לפרש, ואם היו מרבים עליהם שתי שנים או יותר היה מפרשם, ועל כרחך אינם יתרים יותר משנה, הרי כ"ו שנה; מיד יצא משם וחזר לחברון, ואותה שנה קדמה לפני עקדתו של יצחק י"ב שנים, בסדר עולם:}%
% ---------- פרק 22 | פסוק 1 ----------
 \textbf{א} וַיְהִ֗י אַחַר֙ הַדְּבָרִ֣ים הָאֵ֔לֶּה וְהָ֣אֱלֹהִ֔ים נִסָּ֖ה אֶת־אַבְרָהָ֑ם וַיֹּ֣אמֶר אֵלָ֔יו אַבְרָהָ֖ם וַיֹּ֥אמֶר הִנֵּֽנִי׃ \Rashi{\textbf{אחר הדברים האלה.}יש מרבותינו אומרים (סנהדרין פ"ט) אחר דבריו של שטן, שהיה מקטרג ואומר מכל סעודה שעשה אברהם לא הקריב לפניך פר אחד או איל אחד; אמר לו כלום עשה אלא בשביל בנו, אלו הייתי אומר לו זבח אותו לפני לא היה מעכב; וי"א אחר דבריו של ישמעאל, שהיה מתפאר על יצחק שמל בן י"ג שנה ולא מחה, אמר לו יצחק באבר א' אתה מיראני? אלו אמר לי הקב"ה זבח עצמך לפני, לא הייתי מעכב. הנני. כך היא עניתם של חסידים, לשון ענוה הוא ולשון זמון:}%
% ---------- פרק 22 | פסוק 2 ----------
 \textbf{ב} וַיֹּ֡אמֶר קַח־נָ֠א אֶת־בִּנְךָ֨ אֶת־יְחִֽידְךָ֤ אֲשֶׁר־אָהַ֙בְתָּ֙ אֶת־יִצְחָ֔ק וְלֶ֨ךְ־לְךָ֔ אֶל־אֶ֖רֶץ הַמֹּרִיָּ֑ה וְהַעֲלֵ֤הוּ שָׁם֙ לְעֹלָ֔ה עַ֚ל אַחַ֣ד הֶֽהָרִ֔ים אֲשֶׁ֖ר אֹמַ֥ר אֵלֶֽיךָ׃ \Rashi{\textbf{קח נא.}אין נא אלא לשון בקשה, אמר לו בבקשה ממך, עמד לי בזה הנסיון, שלא יאמרו הראשונות לא היה בהן ממש: את בנך. אמר לו שני בנים יש לי, אמר לו את יחידך; אמר לו זה יחיד לאמו וזה יחיד לאמו, אמר לו אשר אהבת; אמר לו שניהם אני אוהב, אמר לו את יצחק (סנהדרין פ"ט). ולמה לא גלה לו מתחלה? שלא לערבבו פתאום ותזוח דעתו עליו ותטרף, וכדי לחבב עליו את המצוה ולתן לו שכר על כל דבור ודבור: ארץ המוריה. ירושלים, וכן בדברי הימים לבנות את בית ה' בירושלם בהר המוריה (דברי הימים ב ג'). ורבותינו פרשו על שם שמשם הוראה יוצאה לישראל. ואנקלוס תרגמו על שם עבודת הקטרת שיש בו מור נרד ושאר בשמים: והעלהו. לא אמר לו שחטהו, לפי שלא היה חפץ הקב"ה לשחטו אלא להעלהו להר לעשותו עולה, ומשהעלהו אמר לו הורידהו: אחד ההרים. הקב"ה מתהא הצדיקים ואחר כך מגלה להם; וכל זה כדי להרבות שכרן, וכן אל הארץ אשר אראך (בראשית י"ב), וכן ביונה וקרא עליה את הקריאה (יונה ג ב):}%
% ---------- פרק 22 | פסוק 3 ----------
 \textbf{ג} וַיַּשְׁכֵּ֨ם אַבְרָהָ֜ם בַּבֹּ֗קֶר וַֽיַּחֲבֹשׁ֙ אֶת־חֲמֹר֔וֹ וַיִּקַּ֞ח אֶת־שְׁנֵ֤י נְעָרָיו֙ אִתּ֔וֹ וְאֵ֖ת יִצְחָ֣ק בְּנ֑וֹ וַיְבַקַּע֙ עֲצֵ֣י עֹלָ֔ה וַיָּ֣קׇם וַיֵּ֔לֶךְ אֶל־הַמָּק֖וֹם אֲשֶׁר־אָֽמַר־ל֥וֹ הָאֱלֹהִֽים׃ \Rashi{\textbf{וישכם.}נזדרז למצוה (פסחים ד'): ויחבש. הוא בעצמו, ולא צוה לאחד מעבדיו, שהאהבה מקלקלת השורה: את שני נעריו. ישמעאל ואליעזר, שאין אדם חשוב רשאי לצאת לדרך בלא ב' אנשים (בראשית רבה), שאם יצטרך האחד לנקביו ויתרחק, יהיה השני עמו: ויבקע. תרגום וצלח, כמו וצלחו הירדן (שמואל ב' י"ט), לשון בקוע, פינדר"א בלע"ז:}%
% ---------- פרק 22 | פסוק 4 ----------
 \textbf{ד} בַּיּ֣וֹם הַשְּׁלִישִׁ֗י וַיִּשָּׂ֨א אַבְרָהָ֧ם אֶת־עֵינָ֛יו וַיַּ֥רְא אֶת־הַמָּק֖וֹם מֵרָחֹֽק׃ \Rashi{\textbf{ביום השלישי.}למה אחר מלהראותו מיד? כדי שלא יאמרו הממו וערבבו פתאום וטרף דעתו, ואלו היה לו שהות להמלך אל לבו לא היה עושה: וירא את המקום. ראה ענן קשור על ההר:}%
% ---------- פרק 22 | פסוק 5 ----------
 \textbf{ה} וַיֹּ֨אמֶר אַבְרָהָ֜ם אֶל־נְעָרָ֗יו שְׁבוּ־לָכֶ֥ם פֹּה֙ עִֽם־הַחֲמ֔וֹר וַאֲנִ֣י וְהַנַּ֔עַר נֵלְכָ֖ה עַד־כֹּ֑ה וְנִֽשְׁתַּחֲוֶ֖ה וְנָשׁ֥וּבָה אֲלֵיכֶֽם׃ \Rashi{\textbf{עד כה.}כלומר דרך מעט למקום אשר לפנינו. ומ"א אראה היכן הוא מה שאמר לי המקום כה יהיה זרעך: ונשובה. נתנבא שישובו שניהם:}%
% ---------- פרק 22 | פסוק 6 ----------
 \textbf{ו} וַיִּקַּ֨ח אַבְרָהָ֜ם אֶת־עֲצֵ֣י הָעֹלָ֗ה וַיָּ֙שֶׂם֙ עַל־יִצְחָ֣ק בְּנ֔וֹ וַיִּקַּ֣ח בְּיָד֔וֹ אֶת־הָאֵ֖שׁ וְאֶת־הַֽמַּאֲכֶ֑לֶת וַיֵּלְכ֥וּ שְׁנֵיהֶ֖ם יַחְדָּֽו׃ \Rashi{\textbf{המאכלת.}סכין, על שם שאוכלת את הבשר, כמה דתימא וחרבי תאכל בשר (דברים ל"ב), ושמכשרת בשר לאכילה. ד"א זאת נקראת מאכלת על שם שישראל אוכלים מתן שכרה: וילכו שניהם יחדו. אברהם שהיה יודע שהולך לשחט את בנו היה הולך ברצון ושמחה כיצחק, שלא היה מרגיש בדבר:}%
% ---------- פרק 22 | פסוק 7 ----------
 \textbf{ז} וַיֹּ֨אמֶר יִצְחָ֜ק אֶל־אַבְרָהָ֤ם אָבִיו֙ וַיֹּ֣אמֶר אָבִ֔י וַיֹּ֖אמֶר הִנֶּ֣נִּֽי בְנִ֑י וַיֹּ֗אמֶר הִנֵּ֤ה הָאֵשׁ֙ וְהָ֣עֵצִ֔ים וְאַיֵּ֥ה הַשֶּׂ֖ה לְעֹלָֽה׃ %
% ---------- פרק 22 | פסוק 8 ----------
 \textbf{ח} וַיֹּ֙אמֶר֙ אַבְרָהָ֔ם אֱלֹהִ֞ים יִרְאֶה־לּ֥וֹ הַשֶּׂ֛ה לְעֹלָ֖ה בְּנִ֑י וַיֵּלְכ֥וּ שְׁנֵיהֶ֖ם יַחְדָּֽו׃ \Rashi{\textbf{יראה לו השה.}כלומר יראה ויבחר לו השה, ואם אין שה, לעולה בני. ואף על פי שהבין יצחק שהוא הולך להשחט, וילכו שניהם יחדו. בלב שוה:}%
% ---------- פרק 22 | פסוק 9 ----------
 \textbf{ט} וַיָּבֹ֗אוּ אֶֽל־הַמָּקוֹם֮ אֲשֶׁ֣ר אָֽמַר־ל֣וֹ הָאֱלֹהִים֒ וַיִּ֨בֶן שָׁ֤ם אַבְרָהָם֙ אֶת־הַמִּזְבֵּ֔חַ וַֽיַּעֲרֹ֖ךְ אֶת־הָעֵצִ֑ים וַֽיַּעֲקֹד֙ אֶת־יִצְחָ֣ק בְּנ֔וֹ וַיָּ֤שֶׂם אֹתוֹ֙ עַל־הַמִּזְבֵּ֔חַ מִמַּ֖עַל לָעֵצִֽים׃ \Rashi{\textbf{ויעקד.}ידיו ורגליו מאחוריו; הידים ורגלים ביחד היא עקידה (שבת נ"ד), והוא לשון עקדים, שהיו קרסוליהם לבנים, מקום שעוקדין אותן בו היה נכר:}%
% ---------- פרק 22 | פסוק 10 ----------
 \textbf{י} וַיִּשְׁלַ֤ח אַבְרָהָם֙ אֶת־יָד֔וֹ וַיִּקַּ֖ח אֶת־הַֽמַּאֲכֶ֑לֶת לִשְׁחֹ֖ט אֶת־בְּנֽוֹ׃ %
% ---------- פרק 22 | פסוק 11 ----------
 \textbf{יא} וַיִּקְרָ֨א אֵלָ֜יו מַלְאַ֤ךְ יְהֹוָה֙ מִן־הַשָּׁמַ֔יִם וַיֹּ֖אמֶר אַבְרָהָ֣ם ׀ אַבְרָהָ֑ם וַיֹּ֖אמֶר הִנֵּֽנִי׃ \Rashi{\textbf{אברהם אברהם.}לשון חבה הוא, שכופל את שמו:}%
% ---------- פרק 22 | פסוק 12 ----------
 \textbf{יב} וַיֹּ֗אמֶר אַל־תִּשְׁלַ֤ח יָֽדְךָ֙ אֶל־הַנַּ֔עַר וְאַל־תַּ֥עַשׂ ל֖וֹ מְא֑וּמָה כִּ֣י ׀ עַתָּ֣ה יָדַ֗עְתִּי כִּֽי־יְרֵ֤א אֱלֹהִים֙ אַ֔תָּה וְלֹ֥א חָשַׂ֛כְתָּ אֶת־בִּנְךָ֥ אֶת־יְחִידְךָ֖ מִמֶּֽנִּי׃ \Rashi{\textbf{אל תשלח.}לשחט; אמר לו אם כן לחנם באתי לכאן, אעשה בו חבלה ואוציא ממנו מעט דם, אמר לו אל תעש לו מאומה – אל תעש בו מום: כי עתה ידעתי. אמר רבי אבא, אמר לו אברהם אפרש לפניך את שיחתי, אתמול אמרת לי כי ביצחק יקרא לך זרע, וחזרת ואמרת קח נא את בנך, עכשו אתה אומר אל תשלח ידך אל הנער? אמר לו הקב"ה, לא אחלל בריתי ומוצא שפתי לא אשנה (תהילים פ"ט); כשאמרתי לך "קח", מוצא שפתי לא אשנה – לא אמרתי לך שחטהו אלא העלהו, אסקתיה, אחתיה. כי עתה ידעתי. מעתה יש לי מה להשיב לשטן ולאמות התמהים מה היא חבתי אצלך; יש לי פתחון פה עכשו, שרואים כי ירא אלהים אתה:}%
% ---------- פרק 22 | פסוק 13 ----------
 \textbf{יג} וַיִּשָּׂ֨א אַבְרָהָ֜ם אֶת־עֵינָ֗יו וַיַּרְא֙ וְהִנֵּה־אַ֔יִל אַחַ֕ר נֶאֱחַ֥ז בַּסְּבַ֖ךְ בְּקַרְנָ֑יו וַיֵּ֤לֶךְ אַבְרָהָם֙ וַיִּקַּ֣ח אֶת־הָאַ֔יִל וַיַּעֲלֵ֥הוּ לְעֹלָ֖ה תַּ֥חַת בְּנֽוֹ׃ \Rashi{\textbf{והנה איל.}מוכן היה לכך מששת ימי בראשית: אחר. אחרי שאמר לו המלאך אל תשלח ידך, ראהו כשהוא נאחז, והוא שמתרגמינן וזקף אברהם ית עינוהי בתר אלין: בסבך. אילן: בקרניו. שהיה רץ אצל אברהם, והשטן סובכו ומערבבו באילנות: תחת בנו. מאחר שכתוב ויעלהו לעלה, לא חסר המקרא כלום, מהו תחת בנו? על כל עבודה שעשה ממנו היה מתפלל ואומר יהי רצון שתהא זו כאלו היא עשויה בבני, כאלו בני שחוט, כאלו דמו זרוק, כאלו הוא מפשט, כאלו הוא נקטר ונעשה דשן:}%
% ---------- פרק 22 | פסוק 14 ----------
 \textbf{יד} וַיִּקְרָ֧א אַבְרָהָ֛ם שֵֽׁם־הַמָּק֥וֹם הַה֖וּא יְהֹוָ֣ה ׀ יִרְאֶ֑ה אֲשֶׁר֙ יֵאָמֵ֣ר הַיּ֔וֹם בְּהַ֥ר יְהֹוָ֖ה יֵרָאֶֽה׃ \Rashi{\textbf{ה' יראה.}פשוטו כתרגומו, ה' יבחר ויראה לו את המקום הזה להשרות בו שכינתו ולהקריב כאן קרבנות: אשר יאמר היום. שיאמרו לימי הדורות עליו בהר זה יראה הקב"ה לעמו: היום. הימים העתידין, כמו עד היום הזה שבכל המקרא, שכל הדורות הבאים הקוראים את המקרא הזה אומרים עד היום הזה על היום שעומדים בו. ומ"א ה' יראה עקדה זו לסלח לישראל בכל שנה ולהצילם מן הפרענות, כדי שיאמר היום הזה – בכל הדורות הבאים – בהר ה' יראה אפרו של יצחק צבור ועומד לכפרה:}%
% ---------- פרק 22 | פסוק 15 ----------
 \textbf{טו} וַיִּקְרָ֛א מַלְאַ֥ךְ יְהֹוָ֖ה אֶל־אַבְרָהָ֑ם שֵׁנִ֖ית מִן־הַשָּׁמָֽיִם׃ %
% ---------- פרק 22 | פסוק 16 ----------
 \textbf{טז} וַיֹּ֕אמֶר בִּ֥י נִשְׁבַּ֖עְתִּי נְאֻם־יְהֹוָ֑ה כִּ֗י יַ֚עַן אֲשֶׁ֤ר עָשִׂ֙יתָ֙ אֶת־הַדָּבָ֣ר הַזֶּ֔ה וְלֹ֥א חָשַׂ֖כְתָּ אֶת־בִּנְךָ֥ אֶת־יְחִידֶֽךָ׃ %
% ---------- פרק 22 | פסוק 17 ----------
 \textbf{יז} כִּֽי־בָרֵ֣ךְ אֲבָרֶכְךָ֗ וְהַרְבָּ֨ה אַרְבֶּ֤ה אֶֽת־זַרְעֲךָ֙ כְּכוֹכְבֵ֣י הַשָּׁמַ֔יִם וְכַח֕וֹל אֲשֶׁ֖ר עַל־שְׂפַ֣ת הַיָּ֑ם וְיִרַ֣שׁ זַרְעֲךָ֔ אֵ֖ת שַׁ֥עַר אֹיְבָֽיו׃ \Rashi{\textbf{ברך אברכך.}אחת לאב ואחת לבן: והרבה ארבה. אחת לאב ואחת לבן:}%
% ---------- פרק 22 | פסוק 18 ----------
 \textbf{יח} וְהִתְבָּרְכ֣וּ בְזַרְעֲךָ֔ כֹּ֖ל גּוֹיֵ֣י הָאָ֑רֶץ עֵ֕קֶב אֲשֶׁ֥ר שָׁמַ֖עְתָּ בְּקֹלִֽי׃ %
% ---------- פרק 22 | פסוק 19 ----------
 \textbf{יט} וַיָּ֤שׇׁב אַבְרָהָם֙ אֶל־נְעָרָ֔יו וַיָּקֻ֛מוּ וַיֵּלְכ֥וּ יַחְדָּ֖ו אֶל־בְּאֵ֣ר שָׁ֑בַע וַיֵּ֥שֶׁב אַבְרָהָ֖ם בִּבְאֵ֥ר שָֽׁבַע׃ \{פ\} \Rashi{\textbf{וישב אברהם בבאר שבע.}לא ישיבה ממש, שהרי בחברון היה יושב, י"ב שנים לפני עקדתו של יצחק יצא מבאר שבע והלך לו לחברון, כמו שנאמר ויגר אברהם בארץ פלשתים ימים רבים, מרבים משל חברון הראשונים, והם כ"ו שנה, כמו שפרשנו למעלה:}%
% ---------- פרק 22 | פסוק 20 ----------
 \textbf{כ} וַיְהִ֗י אַחֲרֵי֙ הַדְּבָרִ֣ים הָאֵ֔לֶּה וַיֻּגַּ֥ד לְאַבְרָהָ֖ם לֵאמֹ֑ר הִ֠נֵּ֠ה יָלְדָ֨ה מִלְכָּ֥ה גַם־הִ֛וא בָּנִ֖ים לְנָח֥וֹר אָחִֽיךָ׃ \Rashi{\textbf{אחרי הדברים האלה ויגד וגו'.}בשובו מהר המוריה היה אברהם מהרהר ואומר אלו היה בני שחוט כבר, היה הולך בלא בנים, היה לי להשיאו אשה מבנות ענר אשכול וממרא, בשרו הקב"ה שנולדה רבקה בת זוגו, וזהו הדברים האלה, הרהורי דברים שהיו על ידי עקדה: גם היא. אף היא השות משפחותיה למשפחות אברהם י"ב, מה אברהם י"ב שבטים שיצאו מיעקב, ח' בני הגבירות וד' בני שפחות, אף אלו ח' בני גבירות וד' בני פילגש:}%
% ---------- פרק 22 | פסוק 21 ----------
 \textbf{כא} אֶת־ע֥וּץ בְּכֹר֖וֹ וְאֶת־בּ֣וּז אָחִ֑יו וְאֶת־קְמוּאֵ֖ל אֲבִ֥י אֲרָֽם׃ %
% ---------- פרק 22 | פסוק 22 ----------
 \textbf{כב} וְאֶת־כֶּ֣שֶׂד וְאֶת־חֲז֔וֹ וְאֶת־פִּלְדָּ֖שׁ וְאֶת־יִדְלָ֑ף וְאֵ֖ת בְּתוּאֵֽל׃ %
% ---------- פרק 22 | פסוק 23 ----------
 \textbf{כג} וּבְתוּאֵ֖ל יָלַ֣ד אֶת־רִבְקָ֑ה שְׁמֹנָ֥ה אֵ֙לֶּה֙ יָלְדָ֣ה מִלְכָּ֔ה לְנָח֖וֹר אֲחִ֥י אַבְרָהָֽם׃ \Rashi{\textbf{ובתואל ילד את רבקה.}כל היחוסין הללו לא נכתבו אלא בשביל פסוק זה:}%
% ---------- פרק 22 | פסוק 24 ----------
 \textbf{כד} וּפִֽילַגְשׁ֖וֹ וּשְׁמָ֣הּ רְאוּמָ֑ה וַתֵּ֤לֶד גַּם־הִוא֙ אֶת־טֶ֣בַח וְאֶת־גַּ֔חַם וְאֶת־תַּ֖חַשׁ וְאֶֽת־מַעֲכָֽה׃ \{פ\} %
\addparasha{חומש בראשית | פרשת חיי שרה}
% ---------- פרק 23 | פסוק 1 ----------
 \textbf{א} וַיִּהְיוּ֙ חַיֵּ֣י שָׂרָ֔ה מֵאָ֥ה שָׁנָ֛ה וְעֶשְׂרִ֥ים שָׁנָ֖ה וְשֶׁ֣בַע שָׁנִ֑ים שְׁנֵ֖י חַיֵּ֥י שָׂרָֽה׃ \Rashi{\textbf{ויהיו חיי שרה מאה שנה ועשרים שנה ושבע שנים.}לכך נכתב שנה בכל כלל וכלל, לומר לך שכל אחד נדרש לעצמו, בת ק' כבת כ' לחטא, מה בת כ' לא חטאה, שהרי אינה בת ענשין, אף בת ק' בלא חטא, ובת כ' כבת ז' ליפי: שני חיי שרה. כלן שוין לטובה:}%
% ---------- פרק 23 | פסוק 2 ----------
 \textbf{ב} וַתָּ֣מׇת שָׂרָ֗ה בְּקִרְיַ֥ת אַרְבַּ֛ע הִ֥וא חֶבְר֖וֹן בְּאֶ֣רֶץ כְּנָ֑עַן וַיָּבֹא֙ אַבְרָהָ֔ם לִסְפֹּ֥ד לְשָׂרָ֖ה וְלִבְכֹּתָֽהּ׃ \Rashi{\textbf{בקרית ארבע.}על שם ארבע ענקים שהיו שם, אחימן ששי ותלמי ואביהם. ד"א על שם ארבע זוגות שנקברו שם איש ואשתו, אדם וחוה, אברהם ושרה, יצחק ורבקה, יעקב ולאה (בראשית רבה): ויבא אברהם. מבאר שבע: לספוד לשרה ולבכתה. ונסמכה מיתת שרה לעקדת יצחק לפי שעל ידי בשורת העקדה, שנזדמן בנה לשחיטה וכמעט שלא נשחט, פרחה נשמתה ממנה ומתה:}%
% ---------- פרק 23 | פסוק 3 ----------
 \textbf{ג} וַיָּ֙קׇם֙ אַבְרָהָ֔ם מֵעַ֖ל פְּנֵ֣י מֵת֑וֹ וַיְדַבֵּ֥ר אֶל־בְּנֵי־חֵ֖ת לֵאמֹֽר׃ %
% ---------- פרק 23 | פסוק 4 ----------
 \textbf{ד} גֵּר־וְתוֹשָׁ֥ב אָנֹכִ֖י עִמָּכֶ֑ם תְּנ֨וּ לִ֤י אֲחֻזַּת־קֶ֙בֶר֙ עִמָּכֶ֔ם וְאֶקְבְּרָ֥ה מֵתִ֖י מִלְּפָנָֽי׃ \Rashi{\textbf{גר ותושב אנכי עמכם.}גר מארץ אחרת ונתישבתי עמכם. ומדרש אגדה אם תרצו הריני גר, ואם לאו אהיה תושב ואטלנה מן הדין, שאמר לי הקב"ה לזרעך אתן את הארץ הזאת: אחוזת קבר. אחזת קרקע לבית הקברות:}%
% ---------- פרק 23 | פסוק 5 ----------
 \textbf{ה} וַיַּעֲנ֧וּ בְנֵי־חֵ֛ת אֶת־אַבְרָהָ֖ם לֵאמֹ֥ר לֽוֹ׃ %
% ---------- פרק 23 | פסוק 6 ----------
 \textbf{ו} שְׁמָעֵ֣נוּ ׀ אֲדֹנִ֗י נְשִׂ֨יא אֱלֹהִ֤ים אַתָּה֙ בְּתוֹכֵ֔נוּ בְּמִבְחַ֣ר קְבָרֵ֔ינוּ קְבֹ֖ר אֶת־מֵתֶ֑ךָ אִ֣ישׁ מִמֶּ֔נּוּ אֶת־קִבְר֛וֹ לֹֽא־יִכְלֶ֥ה מִמְּךָ֖ מִקְּבֹ֥ר מֵתֶֽךָ׃ \Rashi{\textbf{לא יכלה.}לא ימנע כמו לא תכלא רחמיך (תהילים מ') וכמו ויכלא הגשם (בראשית ח'):}%
% ---------- פרק 23 | פסוק 7 ----------
 \textbf{ז} וַיָּ֧קׇם אַבְרָהָ֛ם וַיִּשְׁתַּ֥חוּ לְעַם־הָאָ֖רֶץ לִבְנֵי־חֵֽת׃ %
% ---------- פרק 23 | פסוק 8 ----------
 \textbf{ח} וַיְדַבֵּ֥ר אִתָּ֖ם לֵאמֹ֑ר אִם־יֵ֣שׁ אֶֽת־נַפְשְׁכֶ֗ם לִקְבֹּ֤ר אֶת־מֵתִי֙ מִלְּפָנַ֔י שְׁמָע֕וּנִי וּפִגְעוּ־לִ֖י בְּעֶפְר֥וֹן בֶּן־צֹֽחַר׃ \Rashi{\textbf{נפשכם.}רצונכם: ופגעו לי. לשון בקשה כמו: אל תפגעי בי (רות א'):}%
% ---------- פרק 23 | פסוק 9 ----------
 \textbf{ט} וְיִתֶּן־לִ֗י אֶת־מְעָרַ֤ת הַמַּכְפֵּלָה֙ אֲשֶׁר־ל֔וֹ אֲשֶׁ֖ר בִּקְצֵ֣ה שָׂדֵ֑הוּ בְּכֶ֨סֶף מָלֵ֜א יִתְּנֶ֥נָּה לִּ֛י בְּתוֹכְכֶ֖ם לַאֲחֻזַּת־קָֽבֶר׃ \Rashi{\textbf{המכפלה.}בית ועליה על גביו. ד"א שכפולה בזוגות (עירובין נ"ג): בכסף מלא. שלם, כל שויה; וכן דוד אמר לארונה בכסף מלא (דברי הימים א כ"א):}%
% ---------- פרק 23 | פסוק 10 ----------
 \textbf{י} וְעֶפְר֥וֹן יֹשֵׁ֖ב בְּת֣וֹךְ בְּנֵי־חֵ֑ת וַיַּ֩עַן֩ עֶפְר֨וֹן הַחִתִּ֤י אֶת־אַבְרָהָם֙ בְּאׇזְנֵ֣י בְנֵי־חֵ֔ת לְכֹ֛ל בָּאֵ֥י שַֽׁעַר־עִיר֖וֹ לֵאמֹֽר׃ \Rashi{\textbf{ועפרון ישב.}כתיב חסר, אותו היום מנוהו שוטר עליהם; מפני חשיבותו של אברהם שהיה צריך לו, עלה לגדלה: לכל באי שער עירו. שכלן בטלו ממלאכתן ובאו לגמל חסד לשרה:}%
% ---------- פרק 23 | פסוק 11 ----------
 \textbf{יא} לֹֽא־אֲדֹנִ֣י שְׁמָעֵ֔נִי הַשָּׂדֶה֙ נָתַ֣תִּי לָ֔ךְ וְהַמְּעָרָ֥ה אֲשֶׁר־בּ֖וֹ לְךָ֣ נְתַתִּ֑יהָ לְעֵינֵ֧י בְנֵי־עַמִּ֛י נְתַתִּ֥יהָ לָּ֖ךְ קְבֹ֥ר מֵתֶֽךָ׃ \Rashi{\textbf{לא אדני.}לא תקנה אותה בדמים: נתתי לך. הרי הוא כמו שנתתיה לך:}%
% ---------- פרק 23 | פסוק 12 ----------
 \textbf{יב} וַיִּשְׁתַּ֙חוּ֙ אַבְרָהָ֔ם לִפְנֵ֖י עַ֥ם הָאָֽרֶץ׃ %
% ---------- פרק 23 | פסוק 13 ----------
 \textbf{יג} וַיְדַבֵּ֨ר אֶל־עֶפְר֜וֹן בְּאׇזְנֵ֤י עַם־הָאָ֙רֶץ֙ לֵאמֹ֔ר אַ֛ךְ אִם־אַתָּ֥ה ל֖וּ שְׁמָעֵ֑נִי נָתַ֜תִּי כֶּ֤סֶף הַשָּׂדֶה֙ קַ֣ח מִמֶּ֔נִּי וְאֶקְבְּרָ֥ה אֶת־מֵתִ֖י שָֽׁמָּה׃ \Rashi{\textbf{אך אם אתה לו שמעני.}אתה אומר לי לשמע לך ולקח בחנם, אני אי אפשי בכך: אך אם אתה לו שמעני. הלואי ותשמעני: נתתי. דני"ש בלע"ז, מוכן הוא אצלי, והלואי נתתי לך כבר:}%
% ---------- פרק 23 | פסוק 14 ----------
 \textbf{יד} וַיַּ֧עַן עֶפְר֛וֹן אֶת־אַבְרָהָ֖ם לֵאמֹ֥ר לֽוֹ׃ %
% ---------- פרק 23 | פסוק 15 ----------
 \textbf{טו} אֲדֹנִ֣י שְׁמָעֵ֔נִי אֶ֩רֶץ֩ אַרְבַּ֨ע מֵאֹ֧ת שֶֽׁקֶל־כֶּ֛סֶף בֵּינִ֥י וּבֵֽינְךָ֖ מַה־הִ֑וא וְאֶת־מֵתְךָ֖ קְבֹֽר׃ \Rashi{\textbf{ביני ובינך.}בין שני אוהבים כמונו מה היא חשובה לכלום, אלא הנח את המכר ואת מתך קבר (בראשית רבה):}%
% ---------- פרק 23 | פסוק 16 ----------
 \textbf{טז} וַיִּשְׁמַ֣ע אַבְרָהָם֮ אֶל־עֶפְרוֹן֒ וַיִּשְׁקֹ֤ל אַבְרָהָם֙ לְעֶפְרֹ֔ן אֶת־הַכֶּ֕סֶף אֲשֶׁ֥ר דִּבֶּ֖ר בְּאׇזְנֵ֣י בְנֵי־חֵ֑ת אַרְבַּ֤ע מֵאוֹת֙ שֶׁ֣קֶל כֶּ֔סֶף עֹבֵ֖ר לַסֹּחֵֽר׃ \Rashi{\textbf{וישקל אברהם לעפרן.}חסר וי"ו; לפי שאמר הרבה, ואפלו מעט לא עשה (בבא מציעא פ"ז), שנטל ממנו שקלים גדולים שהן קנטרין, שנאמר עבר לסחר, שמתקבלים בשקל בכל מקום (שם ובכורות נ'), ויש מקום ששקליהן גדולים, שהן קנטרין צנטינא"רש בלע"ז:}%
% ---------- פרק 23 | פסוק 17 ----------
 \textbf{יז} וַיָּ֣קׇם ׀ שְׂדֵ֣ה עֶפְר֗וֹן אֲשֶׁר֙ בַּמַּכְפֵּלָ֔ה אֲשֶׁ֖ר לִפְנֵ֣י מַמְרֵ֑א הַשָּׂדֶה֙ וְהַמְּעָרָ֣ה אֲשֶׁר־בּ֔וֹ וְכׇל־הָעֵץ֙ אֲשֶׁ֣ר בַּשָּׂדֶ֔ה אֲשֶׁ֥ר בְּכׇל־גְּבֻל֖וֹ סָבִֽיב׃ \Rashi{\textbf{ויקם שדה עפרון.}תקומה היתה לו, שיצא מיד הדיוט ליד מלך; ופשוטו של מקרא ויקם השדה והמערה אשר בו וכל העץ לאברהם למקנה וגו':}%
% ---------- פרק 23 | פסוק 18 ----------
 \textbf{יח} לְאַבְרָהָ֥ם לְמִקְנָ֖ה לְעֵינֵ֣י בְנֵי־חֵ֑ת בְּכֹ֖ל בָּאֵ֥י שַֽׁעַר־עִירֽוֹ׃ \Rashi{\textbf{בכל באי שער עירו.}בקרב כלם ובמעמד כלם הקנהו לו:}%
% ---------- פרק 23 | פסוק 19 ----------
 \textbf{יט} וְאַחֲרֵי־כֵן֩ קָבַ֨ר אַבְרָהָ֜ם אֶת־שָׂרָ֣ה אִשְׁתּ֗וֹ אֶל־מְעָרַ֞ת שְׂדֵ֧ה הַמַּכְפֵּלָ֛ה עַל־פְּנֵ֥י מַמְרֵ֖א הִ֣וא חֶבְר֑וֹן בְּאֶ֖רֶץ כְּנָֽעַן׃ %
% ---------- פרק 23 | פסוק 20 ----------
 \textbf{כ} וַיָּ֨קׇם הַשָּׂדֶ֜ה וְהַמְּעָרָ֧ה אֲשֶׁר־בּ֛וֹ לְאַבְרָהָ֖ם לַאֲחֻזַּת־קָ֑בֶר מֵאֵ֖ת בְּנֵי־חֵֽת׃ \{ס\} %
% ---------- פרק 24 | פסוק 1 ----------
 \textbf{א} וְאַבְרָהָ֣ם זָקֵ֔ן בָּ֖א בַּיָּמִ֑ים וַֽיהֹוָ֛ה בֵּרַ֥ךְ אֶת־אַבְרָהָ֖ם בַּכֹּֽל׃ \Rashi{\textbf{ברך את אברהם בכל.}בכל עולה בגימטריא בן, ומאחר שהיה לו בן היה צריך להשיאו אשה:}%
% ---------- פרק 24 | פסוק 2 ----------
 \textbf{ב} וַיֹּ֣אמֶר אַבְרָהָ֗ם אֶל־עַבְדּוֹ֙ זְקַ֣ן בֵּית֔וֹ הַמֹּשֵׁ֖ל בְּכׇל־אֲשֶׁר־ל֑וֹ שִֽׂים־נָ֥א יָדְךָ֖ תַּ֥חַת יְרֵכִֽי׃ \Rashi{\textbf{זקן ביתו.}לפי שהוא דבוק נקוד זקן: תחת ירכי. לפי שהנשבע צריך שיטל בידו חפץ של מצוה, כגון ספר תורה או תפלין (שבועות ל"ח), והמילה היתה מצוה ראשונה לו ובאה לו על ידי צער והיתה חביבה עליו ונטלה:}%
% ---------- פרק 24 | פסוק 3 ----------
 \textbf{ג} וְאַשְׁבִּ֣יעֲךָ֔ בַּֽיהֹוָה֙ אֱלֹהֵ֣י הַשָּׁמַ֔יִם וֵֽאלֹהֵ֖י הָאָ֑רֶץ אֲשֶׁ֨ר לֹֽא־תִקַּ֤ח אִשָּׁה֙ לִבְנִ֔י מִבְּנוֹת֙ הַֽכְּנַעֲנִ֔י אֲשֶׁ֥ר אָנֹכִ֖י יוֹשֵׁ֥ב בְּקִרְבּֽוֹ׃ %
% ---------- פרק 24 | פסוק 4 ----------
 \textbf{ד} כִּ֧י אֶל־אַרְצִ֛י וְאֶל־מוֹלַדְתִּ֖י תֵּלֵ֑ךְ וְלָקַחְתָּ֥ אִשָּׁ֖ה לִבְנִ֥י לְיִצְחָֽק׃ %
% ---------- פרק 24 | פסוק 5 ----------
 \textbf{ה} וַיֹּ֤אמֶר אֵלָיו֙ הָעֶ֔בֶד אוּלַי֙ לֹא־תֹאבֶ֣ה הָֽאִשָּׁ֔ה לָלֶ֥כֶת אַחֲרַ֖י אֶל־הָאָ֣רֶץ הַזֹּ֑את הֶֽהָשֵׁ֤ב אָשִׁיב֙ אֶת־בִּנְךָ֔ אֶל־הָאָ֖רֶץ אֲשֶׁר־יָצָ֥אתָ מִשָּֽׁם׃ %
% ---------- פרק 24 | פסוק 6 ----------
 \textbf{ו} וַיֹּ֥אמֶר אֵלָ֖יו אַבְרָהָ֑ם הִשָּׁ֣מֶר לְךָ֔ פֶּן־תָּשִׁ֥יב אֶת־בְּנִ֖י שָֽׁמָּה׃ %
% ---------- פרק 24 | פסוק 7 ----------
 \textbf{ז} יְהֹוָ֣ה ׀ אֱלֹהֵ֣י הַשָּׁמַ֗יִם אֲשֶׁ֨ר לְקָחַ֜נִי מִבֵּ֣ית אָבִי֮ וּמֵאֶ֣רֶץ מֽוֹלַדְתִּי֒ וַאֲשֶׁ֨ר דִּבֶּר־לִ֜י וַאֲשֶׁ֤ר נִֽשְׁבַּֽע־לִי֙ לֵאמֹ֔ר לְזַ֨רְעֲךָ֔ אֶתֵּ֖ן אֶת־הָאָ֣רֶץ הַזֹּ֑את ה֗וּא יִשְׁלַ֤ח מַלְאָכוֹ֙ לְפָנֶ֔יךָ וְלָקַחְתָּ֥ אִשָּׁ֛ה לִבְנִ֖י מִשָּֽׁם׃ \Rashi{\textbf{ה' אלהי השמים אשר לקחני מבית אבי.}ולא אמר "ואלהי הארץ", ולמעלה אמר ואשביעך וגו'? אמר לו עכשו הוא אלהי השמים ואלהי הארץ, שהרגלתיו בפי הבריות, אבל כשלקחני מבית אבי היה אלהי השמים ולא אלהי הארץ, שלא היו באי עולם מכירים בו ושמו לא היה רגיל בארץ: מבית אבי. מחרן: ומארץ מולדתי. מאור כשדים: ואשר דבר לי. לצרכי; כמו אשר דבר עלי (מלכים א' ב'), וכן כל לי ולו ולהם הסמוכים אצל דבור מפרשים בלשון על, ותרגום שלהם עלי, עלוהי, עליהון, שאין נופל אצל דבור לשון לי ולו ולהם, אלא אלי, אליו, אליהם, ותרגום שלהם עמי, עמיה, עמהון, אבל אצל אמירה נופל לשון לי ולו ולהם: ואשר נשבע לי. בין הבתרים:}%
% ---------- פרק 24 | פסוק 8 ----------
 \textbf{ח} וְאִם־לֹ֨א תֹאבֶ֤ה הָֽאִשָּׁה֙ לָלֶ֣כֶת אַחֲרֶ֔יךָ וְנִקִּ֕יתָ מִשְּׁבֻעָתִ֖י זֹ֑את רַ֣ק אֶת־בְּנִ֔י לֹ֥א תָשֵׁ֖ב שָֽׁמָּה׃ \Rashi{\textbf{ונקית משבועתי וגו'.}וקח לו אשה מבנות ענר אשכול וממרא: רק את בני וגו'. רק מעוט הוא; בני אינו חוזר, אבל יעקב בן בני סופו לחזר:}%
% ---------- פרק 24 | פסוק 9 ----------
 \textbf{ט} וַיָּ֤שֶׂם הָעֶ֙בֶד֙ אֶת־יָד֔וֹ תַּ֛חַת יֶ֥רֶךְ אַבְרָהָ֖ם אֲדֹנָ֑יו וַיִּשָּׁ֣בַֽע ל֔וֹ עַל־הַדָּבָ֖ר הַזֶּֽה׃ %
% ---------- פרק 24 | פסוק 10 ----------
 \textbf{י} וַיִּקַּ֣ח הָ֠עֶ֠בֶד עֲשָׂרָ֨ה גְמַלִּ֜ים מִגְּמַלֵּ֤י אֲדֹנָיו֙ וַיֵּ֔לֶךְ וְכׇל־ט֥וּב אֲדֹנָ֖יו בְּיָד֑וֹ וַיָּ֗קׇם וַיֵּ֛לֶךְ אֶל־אֲרַ֥ם נַֽהֲרַ֖יִם אֶל־עִ֥יר נָחֽוֹר׃ \Rashi{\textbf{מגמלי אדוניו.}נכרין היו משאר גמלים, שהיו יוצאין זמומין מפני הגזל, שלא ירעו בשדות אחרים (בראשית רבה): וכל טוב אדוניו בידו. שטר מתנה כתב ליצחק על כל אשר לו, כדי שיקפצו לשלח לו בתם (בראשית רבה): ארם נהרים. בין שתי נהרות יושבת:}%
% ---------- פרק 24 | פסוק 11 ----------
 \textbf{יא} וַיַּבְרֵ֧ךְ הַגְּמַלִּ֛ים מִח֥וּץ לָעִ֖יר אֶל־בְּאֵ֣ר הַמָּ֑יִם לְעֵ֣ת עֶ֔רֶב לְעֵ֖ת צֵ֥את הַשֹּׁאֲבֹֽת׃ \Rashi{\textbf{ויברך הגמלים.}הרביצם:}%
% ---------- פרק 24 | פסוק 12 ----------
 \textbf{יב} וַיֹּאמַ֓ר ׀ יְהֹוָ֗ה אֱלֹהֵי֙ אֲדֹנִ֣י אַבְרָהָ֔ם הַקְרֵה־נָ֥א לְפָנַ֖י הַיּ֑וֹם וַעֲשֵׂה־חֶ֕סֶד עִ֖ם אֲדֹנִ֥י אַבְרָהָֽם׃ %
% ---------- פרק 24 | פסוק 13 ----------
 \textbf{יג} הִנֵּ֛ה אָנֹכִ֥י נִצָּ֖ב עַל־עֵ֣ין הַמָּ֑יִם וּבְנוֹת֙ אַנְשֵׁ֣י הָעִ֔יר יֹצְאֹ֖ת לִשְׁאֹ֥ב מָֽיִם׃ %
% ---------- פרק 24 | פסוק 14 ----------
 \textbf{יד} וְהָיָ֣ה הַֽנַּעֲרָ֗ אֲשֶׁ֨ר אֹמַ֤ר אֵלֶ֙יהָ֙ הַטִּי־נָ֤א כַדֵּךְ֙ וְאֶשְׁתֶּ֔ה וְאָמְרָ֣ה שְׁתֵ֔ה וְגַם־גְּמַלֶּ֖יךָ אַשְׁקֶ֑ה אֹתָ֤הּ הֹכַ֙חְתָּ֙ לְעַבְדְּךָ֣ לְיִצְחָ֔ק וּבָ֣הּ אֵדַ֔ע כִּי־עָשִׂ֥יתָ חֶ֖סֶד עִם־אֲדֹנִֽי׃ \Rashi{\textbf{אתה הכחת.}ראויה היא לו שתהא גומלת חסדים, וכדאי לכנס בביתו של אברהם; ולשון הכחת בררת, אפרו"בישט בלע"ז: ובה אדע. לשון תחנה, הודע לי בה: כי עשית חסד. אם תהיה ממשפחתו והוגנת לו אדע כי עשית חסד:}%
% ---------- פרק 24 | פסוק 15 ----------
 \textbf{טו} וַֽיְהִי־ה֗וּא טֶ֘רֶם֮ כִּלָּ֣ה לְדַבֵּר֒ וְהִנֵּ֧ה רִבְקָ֣ה יֹצֵ֗את אֲשֶׁ֤ר יֻלְּדָה֙ לִבְתוּאֵ֣ל בֶּן־מִלְכָּ֔ה אֵ֥שֶׁת נָח֖וֹר אֲחִ֣י אַבְרָהָ֑ם וְכַדָּ֖הּ עַל־שִׁכְמָֽהּ׃ %
% ---------- פרק 24 | פסוק 16 ----------
 \textbf{טז} וְהַֽנַּעֲרָ֗ טֹבַ֤ת מַרְאֶה֙ מְאֹ֔ד בְּתוּלָ֕ה וְאִ֖ישׁ לֹ֣א יְדָעָ֑הּ וַתֵּ֣רֶד הָעַ֔יְנָה וַתְּמַלֵּ֥א כַדָּ֖הּ וַתָּֽעַל׃ \Rashi{\textbf{בתולה.}ממקום בתולים (בראשית רבה): ואיש לא ידעה. שלא כדרכה, לפי שבנות הגוים היו משמרות מקום בתוליהן ומפקירות עצמן ממקום אחר, העיד על זו שנקיה מכל:}%
% ---------- פרק 24 | פסוק 17 ----------
 \textbf{יז} וַיָּ֥רׇץ הָעֶ֖בֶד לִקְרָאתָ֑הּ וַיֹּ֕אמֶר הַגְמִיאִ֥ינִי נָ֛א מְעַט־מַ֖יִם מִכַּדֵּֽךְ׃ \Rashi{\textbf{וירץ העבד לקראתה.}לפי שראה שעלו המים לקראתה (בראשית רבה): הגמיאיני נא. לשון גמיעה, הומיי"ר בלע"ז:}%
% ---------- פרק 24 | פסוק 18 ----------
 \textbf{יח} וַתֹּ֖אמֶר שְׁתֵ֣ה אֲדֹנִ֑י וַתְּמַהֵ֗ר וַתֹּ֧רֶד כַּדָּ֛הּ עַל־יָדָ֖הּ וַתַּשְׁקֵֽהוּ׃ \Rashi{\textbf{ותרד כדה.}מעל שכמה:}%
% ---------- פרק 24 | פסוק 19 ----------
 \textbf{יט} וַתְּכַ֖ל לְהַשְׁקֹת֑וֹ וַתֹּ֗אמֶר גַּ֤ם לִגְמַלֶּ֙יךָ֙ אֶשְׁאָ֔ב עַ֥ד אִם־כִּלּ֖וּ לִשְׁתֹּֽת׃ \Rashi{\textbf{עד אם כלו.}הרי אם משמש בלשון אשר, אם כלו די ספוקון, שזו היא גמר שתיתן, כששתו די ספוקן:}%
% ---------- פרק 24 | פסוק 20 ----------
 \textbf{כ} וַתְּמַהֵ֗ר וַתְּעַ֤ר כַּדָּהּ֙ אֶל־הַשֹּׁ֔קֶת וַתָּ֥רׇץ ע֛וֹד אֶֽל־הַבְּאֵ֖ר לִשְׁאֹ֑ב וַתִּשְׁאַ֖ב לְכׇל־גְּמַלָּֽיו׃ \Rashi{\textbf{ותער.}לשון נפיצה, והרבה יש בלשון משנה המערה מכלי אל כלי, ובמקרא יש לו דומה אל תער נפשי (תהילים קמא), אשר הערה למות נפשו (ישעיהו נ"ג): השקת. אבן חלולה ששותים בה הגמלים:}%
% ---------- פרק 24 | פסוק 21 ----------
 \textbf{כא} וְהָאִ֥ישׁ מִשְׁתָּאֵ֖ה לָ֑הּ מַחֲרִ֕ישׁ לָדַ֗עַת הַֽהִצְלִ֧יחַ יְהֹוָ֛ה דַּרְכּ֖וֹ אִם־לֹֽא׃ \Rashi{\textbf{משתאה.}לשון שאיה, כמו שאו ערים, תשאה שממה (שם ו'): משתאה. משתומם ומתבהל על שראה דברו קרוב להצליח, אבל אינו יודע אם ממשפחת אברהם היא אם לאו. ואל תתמה בת' של משתאה, שאין לך תבה שתחלת יסודה שי"ן ומדברת בלשון מתפעל שאין תי"ו מפרידה בין שתי אותיות של עקר היסוד כגון משתאה, משתולל מגזרת שולל, וישתומם מגזרת שממה, וישתמר חקות עמרי (מיכה ו') מגזרת וישמר, אף כאן משתאה מגזרת תשאה, וכשם שאתה מוצא לשון משומם באדם נבהל ונאלם ובעל מחשבות, כמו על יומו נשמו אחרנים (איוב י"ח), שמו שמים (ירמיהו ב'), אשתומם כשעה חדא (דניאל ד'), כך תפרש לשון שאיה באדם בהול ובעל מחשבות (ואנקלוס תרגם לשון שהיה, וגברא שהי, שוהא ועומד במקום אחד לראות ההצליח ה' דרכו; ואין לתרגם שתי, שהרי אינו לשון שתיה, שאין אלף נופלת בלשון שתיה): משתאה לה. משתומם עליה, כמו אמרי לי אחי הוא (בראשית כ׳:י״ג) וכמו וישאלו אנשי המקום לאשתו (שם כ"ו):}%
% ---------- פרק 24 | פסוק 22 ----------
 \textbf{כב} וַיְהִ֗י כַּאֲשֶׁ֨ר כִּלּ֤וּ הַגְּמַלִּים֙ לִשְׁתּ֔וֹת וַיִּקַּ֤ח הָאִישׁ֙ נֶ֣זֶם זָהָ֔ב בֶּ֖קַע מִשְׁקָל֑וֹ וּשְׁנֵ֤י צְמִידִים֙ עַל־יָדֶ֔יהָ עֲשָׂרָ֥ה זָהָ֖ב מִשְׁקָלָֽם׃ \Rashi{\textbf{בקע.}רמז לשקלי ישראל בקע לגלגלת: ושני צמידים. רמז לשני לוחות מצמדות: עשרה זהב משקלם. רמז לעשרת הדברות שבהן:}%
% ---------- פרק 24 | פסוק 23 ----------
 \textbf{כג} וַיֹּ֙אמֶר֙ בַּת־מִ֣י אַ֔תְּ הַגִּ֥ידִי נָ֖א לִ֑י הֲיֵ֧שׁ בֵּית־אָבִ֛יךְ מָק֥וֹם לָ֖נוּ לָלִֽין׃ \Rashi{\textbf{ויאמר בת מי את.}לאחר שנתן לה שאלה, לפי שהיה בטוח בזכותו של אברהם שהצליח הקב"ה דרכו: ללין. לינה אחת, לין שם דבר; והיא אמרה ללון – כמה לינות:}%
% ---------- פרק 24 | פסוק 24 ----------
 \textbf{כד} וַתֹּ֣אמֶר אֵלָ֔יו בַּת־בְּתוּאֵ֖ל אָנֹ֑כִי בֶּן־מִלְכָּ֕ה אֲשֶׁ֥ר יָלְדָ֖ה לְנָחֽוֹר׃ \Rashi{\textbf{בת בתואל.}השיבתו על ראשון ראשון ועל אחרון אחרון:}%
% ---------- פרק 24 | פסוק 25 ----------
 \textbf{כה} וַתֹּ֣אמֶר אֵלָ֔יו גַּם־תֶּ֥בֶן גַּם־מִסְפּ֖וֹא רַ֣ב עִמָּ֑נוּ גַּם־מָק֖וֹם לָלֽוּן׃ \Rashi{\textbf{מספוא.}כל מאכל הגמלים קרוי מספוא, כגון תבן ושעורים:}%
% ---------- פרק 24 | פסוק 26 ----------
 \textbf{כו} וַיִּקֹּ֣ד הָאִ֔ישׁ וַיִּשְׁתַּ֖חוּ לַֽיהֹוָֽה׃ %
% ---------- פרק 24 | פסוק 27 ----------
 \textbf{כז} וַיֹּ֗אמֶר בָּר֤וּךְ יְהֹוָה֙ אֱלֹהֵי֙ אֲדֹנִ֣י אַבְרָהָ֔ם אֲ֠שֶׁ֠ר לֹֽא־עָזַ֥ב חַסְדּ֛וֹ וַאֲמִתּ֖וֹ מֵעִ֣ם אֲדֹנִ֑י אָנֹכִ֗י בַּדֶּ֙רֶךְ֙ נָחַ֣נִי יְהֹוָ֔ה בֵּ֖ית אֲחֵ֥י אֲדֹנִֽי׃ \Rashi{\textbf{בדרך.}דרך המזמן, דרך הישר, באותו דרך שהייתי צריך, וכן כל בי"ת ולמ"ד וה"א המשמשים בראש התבה ונקודים בפתח, מדברים בדבר הפשוט שנזכר כבר במקום אחר, או שהוא מברר ונכר באיזו הוא מדבר:}%
% ---------- פרק 24 | פסוק 28 ----------
 \textbf{כח} וַתָּ֙רׇץ֙ הַֽנַּעֲרָ֔ וַתַּגֵּ֖ד לְבֵ֣ית אִמָּ֑הּ כַּדְּבָרִ֖ים הָאֵֽלֶּה׃ \Rashi{\textbf{לבית אמה.}דרך הנשים היתה להיות להן בית לישב בו למלאכתן, ואין הבת מגדת אלא לאמה:}%
% ---------- פרק 24 | פסוק 29 ----------
 \textbf{כט} וּלְרִבְקָ֥ה אָ֖ח וּשְׁמ֣וֹ לָבָ֑ן וַיָּ֨רׇץ לָבָ֧ן אֶל־הָאִ֛ישׁ הַח֖וּצָה אֶל־הָעָֽיִן׃ \Rashi{\textbf{וירץ.}למה רץ ועל מה רץ? ויהי כראות את הנזם, אמר עשיר הוא זה, ונתן עיניו בממון:}%
% ---------- פרק 24 | פסוק 30 ----------
 \textbf{ל} וַיְהִ֣י ׀ כִּרְאֹ֣ת אֶת־הַנֶּ֗זֶם וְֽאֶת־הַצְּמִדִים֮ עַל־יְדֵ֣י אֲחֹתוֹ֒ וּכְשׇׁמְע֗וֹ אֶת־דִּבְרֵ֞י רִבְקָ֤ה אֲחֹתוֹ֙ לֵאמֹ֔ר כֹּֽה־דִבֶּ֥ר אֵלַ֖י הָאִ֑ישׁ וַיָּבֹא֙ אֶל־הָאִ֔ישׁ וְהִנֵּ֛ה עֹמֵ֥ד עַל־הַגְּמַלִּ֖ים עַל־הָעָֽיִן׃ \Rashi{\textbf{על הגמלים.}לשמרן, כמו והוא עומד עליהם (בראשית י״ח:ח׳), לשמשם:}%
% ---------- פרק 24 | פסוק 31 ----------
 \textbf{לא} וַיֹּ֕אמֶר בּ֖וֹא בְּר֣וּךְ יְהֹוָ֑ה לָ֤מָּה תַעֲמֹד֙ בַּח֔וּץ וְאָנֹכִי֙ פִּנִּ֣יתִי הַבַּ֔יִת וּמָק֖וֹם לַגְּמַלִּֽים׃ \Rashi{\textbf{פניתי הבית.}מעבודת אלילים (בראשית רבה):}%
% ---------- פרק 24 | פסוק 32 ----------
 \textbf{לב} וַיָּבֹ֤א הָאִישׁ֙ הַבַּ֔יְתָה וַיְפַתַּ֖ח הַגְּמַלִּ֑ים וַיִּתֵּ֨ן תֶּ֤בֶן וּמִסְפּוֹא֙ לַגְּמַלִּ֔ים וּמַ֙יִם֙ לִרְחֹ֣ץ רַגְלָ֔יו וְרַגְלֵ֥י הָאֲנָשִׁ֖ים אֲשֶׁ֥ר אִתּֽוֹ׃ \Rashi{\textbf{ויפתח.}התיר זמם שלהם, שהיה סותם את פיהם שלא ירעו בדרך בשדות אחרים:}%
% ---------- פרק 24 | פסוק 33 ----------
 \textbf{לג} (ויישם) [וַיּוּשַׂ֤ם] לְפָנָיו֙ לֶאֱכֹ֔ל וַיֹּ֙אמֶר֙ לֹ֣א אֹכַ֔ל עַ֥ד אִם־דִּבַּ֖רְתִּי דְּבָרָ֑י וַיֹּ֖אמֶר דַּבֵּֽר׃ \Rashi{\textbf{עד אם דברתי.}הרי אם משמש בלשון אשר ובלשון כי, כמו: עד כי יבא שילה (בראשית מ"ט) וזהו שאמרו חז"ל: כי משמש בד' לשונות, והאחד אי, והוא אם:}%
% ---------- פרק 24 | פסוק 34 ----------
 \textbf{לד} וַיֹּאמַ֑ר עֶ֥בֶד אַבְרָהָ֖ם אָנֹֽכִי׃ %
% ---------- פרק 24 | פסוק 35 ----------
 \textbf{לה} וַיהֹוָ֞ה בֵּרַ֧ךְ אֶת־אֲדֹנִ֛י מְאֹ֖ד וַיִּגְדָּ֑ל וַיִּתֶּן־ל֞וֹ צֹ֤אן וּבָקָר֙ וְכֶ֣סֶף וְזָהָ֔ב וַעֲבָדִם֙ וּשְׁפָחֹ֔ת וּגְמַלִּ֖ים וַחֲמֹרִֽים׃ %
% ---------- פרק 24 | פסוק 36 ----------
 \textbf{לו} וַתֵּ֡לֶד שָׂרָה֩ אֵ֨שֶׁת אֲדֹנִ֥י בֵן֙ לַֽאדֹנִ֔י אַחֲרֵ֖י זִקְנָתָ֑הּ וַיִּתֶּן־ל֖וֹ אֶת־כׇּל־אֲשֶׁר־לֽוֹ׃ \Rashi{\textbf{ויתן לו את כל אשר לו.}שטר המתנה הראה להם:}%
% ---------- פרק 24 | פסוק 37 ----------
 \textbf{לז} וַיַּשְׁבִּעֵ֥נִי אֲדֹנִ֖י לֵאמֹ֑ר לֹא־תִקַּ֤ח אִשָּׁה֙ לִבְנִ֔י מִבְּנוֹת֙ הַֽכְּנַעֲנִ֔י אֲשֶׁ֥ר אָנֹכִ֖י יֹשֵׁ֥ב בְּאַרְצֽוֹ׃ \Rashi{\textbf{לא תקח אשה לבני מבנות הכנעני.}אם לא תלך תחלה אל בית אבי ולא תאבה ללכת אחריך:}%
% ---------- פרק 24 | פסוק 38 ----------
 \textbf{לח} אִם־לֹ֧א אֶל־בֵּית־אָבִ֛י תֵּלֵ֖ךְ וְאֶל־מִשְׁפַּחְתִּ֑י וְלָקַחְתָּ֥ אִשָּׁ֖ה לִבְנִֽי׃ %
% ---------- פרק 24 | פסוק 39 ----------
 \textbf{לט} וָאֹמַ֖ר אֶל־אֲדֹנִ֑י אֻלַ֛י לֹא־תֵלֵ֥ךְ הָאִשָּׁ֖ה אַחֲרָֽי׃ \Rashi{\textbf{אלי לא תלך האשה.}אלי כתיב, בת היתה לו לאליעזר, והיה מחזר למצא עלה שיאמר לו אברהם לפנות אליו להשיאו בתו, אמר לו אברהם בני ברוך ואתה ארור ואין ארור מדבק בברוך:}%
% ---------- פרק 24 | פסוק 40 ----------
 \textbf{מ} וַיֹּ֖אמֶר אֵלָ֑י יְהֹוָ֞ה אֲשֶׁר־הִתְהַלַּ֣כְתִּי לְפָנָ֗יו יִשְׁלַ֨ח מַלְאָכ֤וֹ אִתָּךְ֙ וְהִצְלִ֣יחַ דַּרְכֶּ֔ךָ וְלָקַחְתָּ֤ אִשָּׁה֙ לִבְנִ֔י מִמִּשְׁפַּחְתִּ֖י וּמִבֵּ֥ית אָבִֽי׃ %
% ---------- פרק 24 | פסוק 41 ----------
 \textbf{מא} אָ֤ז תִּנָּקֶה֙ מֵאָ֣לָתִ֔י כִּ֥י תָב֖וֹא אֶל־מִשְׁפַּחְתִּ֑י וְאִם־לֹ֤א יִתְּנוּ֙ לָ֔ךְ וְהָיִ֥יתָ נָקִ֖י מֵאָלָתִֽי׃ %
% ---------- פרק 24 | פסוק 42 ----------
 \textbf{מב} וָאָבֹ֥א הַיּ֖וֹם אֶל־הָעָ֑יִן וָאֹמַ֗ר יְהֹוָה֙ אֱלֹהֵי֙ אֲדֹנִ֣י אַבְרָהָ֔ם אִם־יֶשְׁךָ־נָּא֙ מַצְלִ֣יחַ דַּרְכִּ֔י אֲשֶׁ֥ר אָנֹכִ֖י הֹלֵ֥ךְ עָלֶֽיהָ׃ \Rashi{\textbf{ואבא היום.}היום יצאתי והיום באתי, מכאן שקפצה לו הארץ. אמר רבי אחא יפה שיחתן של עבדי אבות לפני המקום מתורתן של בנים, שהרי פרשה של אליעזר כפולה בתורה והרבה גופי תורה לא נתנו אלא ברמיזה (בראשית רבה):}%
% ---------- פרק 24 | פסוק 43 ----------
 \textbf{מג} הִנֵּ֛ה אָנֹכִ֥י נִצָּ֖ב עַל־עֵ֣ין הַמָּ֑יִם וְהָיָ֤ה הָֽעַלְמָה֙ הַיֹּצֵ֣את לִשְׁאֹ֔ב וְאָמַרְתִּ֣י אֵלֶ֔יהָ הַשְׁקִֽינִי־נָ֥א מְעַט־מַ֖יִם מִכַּדֵּֽךְ׃ %
% ---------- פרק 24 | פסוק 44 ----------
 \textbf{מד} וְאָמְרָ֤ה אֵלַי֙ גַּם־אַתָּ֣ה שְׁתֵ֔ה וְגַ֥ם לִגְמַלֶּ֖יךָ אֶשְׁאָ֑ב הִ֣וא הָֽאִשָּׁ֔ה אֲשֶׁר־הֹכִ֥יחַ יְהֹוָ֖ה לְבֶן־אֲדֹנִֽי׃ \Rashi{\textbf{גם אתה.}גם לרבות אנשים שעמו: הוכיח. ברר והודיע, וכן כל הוכחה שבמקרא ברור דבר:}%
% ---------- פרק 24 | פסוק 45 ----------
 \textbf{מה} אֲנִי֩ טֶ֨רֶם אֲכַלֶּ֜ה לְדַבֵּ֣ר אֶל־לִבִּ֗י וְהִנֵּ֨ה רִבְקָ֤ה יֹצֵאת֙ וְכַדָּ֣הּ עַל־שִׁכְמָ֔הּ וַתֵּ֥רֶד הָעַ֖יְנָה וַתִּשְׁאָ֑ב וָאֹמַ֥ר אֵלֶ֖יהָ הַשְׁקִ֥ינִי נָֽא׃ \Rashi{\textbf{טרם אכלה.}טרם שאני מכלה, וכן כל לשון הווה, פעמים שהוא מדבר בלשון עבר ויכול לכתב טרם כליתי, ופעמים שמדבר בלשון עתיד, כמו כי אמר איוב, הרי לשון עבר, ככה יעשה איוב (איוב א'), הרי לשון עתיד, ופרוש שניהם לשון הווה, כי אומר היה איוב אולי חטאו בני וגו', והיה עושה כך:}%
% ---------- פרק 24 | פסוק 46 ----------
 \textbf{מו} וַתְּמַהֵ֗ר וַתּ֤וֹרֶד כַּדָּהּ֙ מֵֽעָלֶ֔יהָ וַתֹּ֣אמֶר שְׁתֵ֔ה וְגַם־גְּמַלֶּ֖יךָ אַשְׁקֶ֑ה וָאֵ֕שְׁתְּ וְגַ֥ם הַגְּמַלִּ֖ים הִשְׁקָֽתָה׃ %
% ---------- פרק 24 | פסוק 47 ----------
 \textbf{מז} וָאֶשְׁאַ֣ל אֹתָ֗הּ וָאֹמַר֮ בַּת־מִ֣י אַתְּ֒ וַתֹּ֗אמֶר בַּת־בְּתוּאֵל֙ בֶּן־נָח֔וֹר אֲשֶׁ֥ר יָֽלְדָה־לּ֖וֹ מִלְכָּ֑ה וָאָשִׂ֤ם הַנֶּ֙זֶם֙ עַל־אַפָּ֔הּ וְהַצְּמִידִ֖ים עַל־יָדֶֽיהָ׃ \Rashi{\textbf{ואשאל.}ואשם. שנה הסדר, שהרי הוא תחלה נתן ואחר כך שאל, אלא שלא יתפשוהו בדבריו ויאמרו האיך נתת לה ועדין אינך יודע מי היא:}%
% ---------- פרק 24 | פסוק 48 ----------
 \textbf{מח} וָאֶקֹּ֥ד וָֽאֶשְׁתַּחֲוֶ֖ה לַיהֹוָ֑ה וָאֲבָרֵ֗ךְ אֶת־יְהֹוָה֙ אֱלֹהֵי֙ אֲדֹנִ֣י אַבְרָהָ֔ם אֲשֶׁ֤ר הִנְחַ֙נִי֙ בְּדֶ֣רֶךְ אֱמֶ֔ת לָקַ֛חַת אֶת־בַּת־אֲחִ֥י אֲדֹנִ֖י לִבְנֽוֹ׃ %
% ---------- פרק 24 | פסוק 49 ----------
 \textbf{מט} וְ֠עַתָּ֠ה אִם־יֶשְׁכֶ֨ם עֹשִׂ֜ים חֶ֧סֶד וֶֽאֱמֶ֛ת אֶת־אֲדֹנִ֖י הַגִּ֣ידוּ לִ֑י וְאִם־לֹ֕א הַגִּ֣ידוּ לִ֔י וְאֶפְנֶ֥ה עַל־יָמִ֖ין א֥וֹ עַל־שְׂמֹֽאל׃ \Rashi{\textbf{על ימין.}מבנות ישמעאל: על שמאל. מבנות לוט, שהיה יושב לשמאלו של אברהם:}%
% ---------- פרק 24 | פסוק 50 ----------
 \textbf{נ} וַיַּ֨עַן לָבָ֤ן וּבְתוּאֵל֙ וַיֹּ֣אמְר֔וּ מֵיְהֹוָ֖ה יָצָ֣א הַדָּבָ֑ר לֹ֥א נוּכַ֛ל דַּבֵּ֥ר אֵלֶ֖יךָ רַ֥ע אוֹ־טֽוֹב׃ \Rashi{\textbf{ויען לבן ובתואל.}רשע היה, וקפץ להשיב לפני אביו: לא נוכל דבר אליך. למאן בדבר הזה לא על ידי תשובת דבר רע ולא על ידי תשובת דבר הגון, ונכר שמה' יצא הדבר, לפי דבריך שזמנה לך (בראשית רבה):}%
% ---------- פרק 24 | פסוק 51 ----------
 \textbf{נא} הִנֵּֽה־רִבְקָ֥ה לְפָנֶ֖יךָ קַ֣ח וָלֵ֑ךְ וּתְהִ֤י אִשָּׁה֙ לְבֶן־אֲדֹנֶ֔יךָ כַּאֲשֶׁ֖ר דִּבֶּ֥ר יְהֹוָֽה׃ %
% ---------- פרק 24 | פסוק 52 ----------
 \textbf{נב} וַיְהִ֕י כַּאֲשֶׁ֥ר שָׁמַ֛ע עֶ֥בֶד אַבְרָהָ֖ם אֶת־דִּבְרֵיהֶ֑ם וַיִּשְׁתַּ֥חוּ אַ֖רְצָה לַֽיהֹוָֽה׃ \Rashi{\textbf{וישתחו ארצה.}מכאן שמודים על בשורה טובה:}%
% ---------- פרק 24 | פסוק 53 ----------
 \textbf{נג} וַיּוֹצֵ֨א הָעֶ֜בֶד כְּלֵי־כֶ֨סֶף וּכְלֵ֤י זָהָב֙ וּבְגָדִ֔ים וַיִּתֵּ֖ן לְרִבְקָ֑ה וּמִ֨גְדָּנֹ֔ת נָתַ֥ן לְאָחִ֖יהָ וּלְאִמָּֽהּ׃ \Rashi{\textbf{ומגדנות.}לשון מגדים, שהביא עמו מיני פרות של ארץ ישראל:}%
% ---------- פרק 24 | פסוק 54 ----------
 \textbf{נד} וַיֹּאכְל֣וּ וַיִּשְׁתּ֗וּ ה֛וּא וְהָאֲנָשִׁ֥ים אֲשֶׁר־עִמּ֖וֹ וַיָּלִ֑ינוּ וַיָּק֣וּמוּ בַבֹּ֔קֶר וַיֹּ֖אמֶר שַׁלְּחֻ֥נִי לַֽאדֹנִֽי׃ \Rashi{\textbf{וילינו.}כל לינה שבמקרא לינת לילה אחד:}%
% ---------- פרק 24 | פסוק 55 ----------
 \textbf{נה} וַיֹּ֤אמֶר אָחִ֙יהָ֙ וְאִמָּ֔הּ תֵּשֵׁ֨ב הַנַּעֲרָ֥ אִתָּ֛נוּ יָמִ֖ים א֣וֹ עָשׂ֑וֹר אַחַ֖ר תֵּלֵֽךְ׃ \Rashi{\textbf{ויאמר אחיה ואמה.}ובתואל היכן היה? היה רוצה לעכב ובא מלאך והמיתו: ימים. שנה, כמו ימים תהיה גאלתו (ויקרא כ"ה), שכך נותנין לבתולה זמן י"ב חדש לפרנס עצמה בתכשיטים (כתובות נ"ז): או עשור. י' חדשים; ואם תאמר ימים ממש, אין דרך המבקשים לבקש דבר מעט, ואם לא תרצה תן לנו מרבה מזה:}%
% ---------- פרק 24 | פסוק 56 ----------
 \textbf{נו} וַיֹּ֤אמֶר אֲלֵהֶם֙ אַל־תְּאַחֲר֣וּ אֹתִ֔י וַֽיהֹוָ֖ה הִצְלִ֣יחַ דַּרְכִּ֑י שַׁלְּח֕וּנִי וְאֵלְכָ֖ה לַֽאדֹנִֽי׃ %
% ---------- פרק 24 | פסוק 57 ----------
 \textbf{נז} וַיֹּאמְר֖וּ נִקְרָ֣א לַֽנַּעֲרָ֑ וְנִשְׁאֲלָ֖ה אֶת־פִּֽיהָ׃ \Rashi{\textbf{ונשאלה את פיה.}מכאן שאין משיאין את האשה אלא מדעתה (בראשית רבה):}%
% ---------- פרק 24 | פסוק 58 ----------
 \textbf{נח} וַיִּקְרְא֤וּ לְרִבְקָה֙ וַיֹּאמְר֣וּ אֵלֶ֔יהָ הֲתֵלְכִ֖י עִם־הָאִ֣ישׁ הַזֶּ֑ה וַתֹּ֖אמֶר אֵלֵֽךְ׃ \Rashi{\textbf{ותאמר אלך.}מעצמי, ואף אם אינכם רוצים:}%
% ---------- פרק 24 | פסוק 59 ----------
 \textbf{נט} וַֽיְשַׁלְּח֛וּ אֶת־רִבְקָ֥ה אֲחֹתָ֖ם וְאֶת־מֵנִקְתָּ֑הּ וְאֶת־עֶ֥בֶד אַבְרָהָ֖ם וְאֶת־אֲנָשָֽׁיו׃ %
% ---------- פרק 24 | פסוק 60 ----------
 \textbf{ס} וַיְבָרְכ֤וּ אֶת־רִבְקָה֙ וַיֹּ֣אמְרוּ לָ֔הּ אֲחֹתֵ֕נוּ אַ֥תְּ הֲיִ֖י לְאַלְפֵ֣י רְבָבָ֑ה וְיִירַ֣שׁ זַרְעֵ֔ךְ אֵ֖ת שַׁ֥עַר שֹׂנְאָֽיו׃ \Rashi{\textbf{את היי לאלפי רבבה.}את וזרעך תקבלו אותה ברכה שנאמר לאברהם בהר המוריה והרבה ארבה את זרעך וגו', יהי רצון שיהא אותו הזרע ממך ולא מאשה אחרת:}%
% ---------- פרק 24 | פסוק 61 ----------
 \textbf{סא} וַתָּ֨קׇם רִבְקָ֜ה וְנַעֲרֹתֶ֗יהָ וַתִּרְכַּ֙בְנָה֙ עַל־הַגְּמַלִּ֔ים וַתֵּלַ֖כְנָה אַחֲרֵ֣י הָאִ֑ישׁ וַיִּקַּ֥ח הָעֶ֛בֶד אֶת־רִבְקָ֖ה וַיֵּלַֽךְ׃ %
% ---------- פרק 24 | פסוק 62 ----------
 \textbf{סב} וְיִצְחָק֙ בָּ֣א מִבּ֔וֹא בְּאֵ֥ר לַחַ֖י רֹאִ֑י וְה֥וּא יוֹשֵׁ֖ב בְּאֶ֥רֶץ הַנֶּֽגֶב׃ \Rashi{\textbf{מבוא באר לחי רואי.}שהלך להביא הגר לאברהם אביו שישאנה (בראשית רבה): יושב בארץ הנגב. קרוב לאותו באר, שנאמר ויסע משם אברהם ארצה הנגב וישב בין קדש ובין שור (בראשית כ') ושם היה הבאר, שנאמר הנה בין קדש ובין ברד (שם ט"ז):}%
% ---------- פרק 24 | פסוק 63 ----------
 \textbf{סג} וַיֵּצֵ֥א יִצְחָ֛ק לָשׂ֥וּחַ בַּשָּׂדֶ֖ה לִפְנ֣וֹת עָ֑רֶב וַיִּשָּׂ֤א עֵינָיו֙ וַיַּ֔רְא וְהִנֵּ֥ה גְמַלִּ֖ים בָּאִֽים׃ \Rashi{\textbf{לשוח.}לשון תפלה (בראשית רבה), כמו ישפך שיחו (תהילים ק"ב):}%
% ---------- פרק 24 | פסוק 64 ----------
 \textbf{סד} וַתִּשָּׂ֤א רִבְקָה֙ אֶת־עֵינֶ֔יהָ וַתֵּ֖רֶא אֶת־יִצְחָ֑ק וַתִּפֹּ֖ל מֵעַ֥ל הַגָּמָֽל׃ \Rashi{\textbf{ותרא את יצחק.}ראתה אותו הדור ותוהא מפניו (בראשית רבה): ותפל. השמיטה עצמה לארץ, כתרגומו ואתרכינת, הטתה עצמה לארץ ולא הגיעה עד הקרקע, כמו הטי נא כדך, ארכיני, ויט שמים (תהילים י"ח), וארכין, לשון מטה לארץ, ודומה לו כי יפל לא יוטל (תהילים ל"ז), כלומר אם יטה לארץ לא יגיע עד הקרקע:}%
% ---------- פרק 24 | פסוק 65 ----------
 \textbf{סה} וַתֹּ֣אמֶר אֶל־הָעֶ֗בֶד מִֽי־הָאִ֤ישׁ הַלָּזֶה֙ הַהֹלֵ֤ךְ בַּשָּׂדֶה֙ לִקְרָאתֵ֔נוּ וַיֹּ֥אמֶר הָעֶ֖בֶד ה֣וּא אֲדֹנִ֑י וַתִּקַּ֥ח הַצָּעִ֖יף וַתִּתְכָּֽס׃ \Rashi{\textbf{ותתכס.}לשון ותתפעל, כמו ותקבר ותשבר:}%
% ---------- פרק 24 | פסוק 66 ----------
 \textbf{סו} וַיְסַפֵּ֥ר הָעֶ֖בֶד לְיִצְחָ֑ק אֵ֥ת כׇּל־הַדְּבָרִ֖ים אֲשֶׁ֥ר עָשָֽׂה׃ \Rashi{\textbf{ויספר העבד.}גלה לו נסים שנעשו לו, שקפצה לו הארץ ושנזדמנה לו רבקה בתפלתו (בראשית רבה):}%
% ---------- פרק 24 | פסוק 67 ----------
 \textbf{סז} וַיְבִאֶ֣הָ יִצְחָ֗ק הָאֹ֙הֱלָה֙ שָׂרָ֣ה אִמּ֔וֹ וַיִּקַּ֧ח אֶת־רִבְקָ֛ה וַתְּהִי־ל֥וֹ לְאִשָּׁ֖ה וַיֶּאֱהָבֶ֑הָ וַיִּנָּחֵ֥ם יִצְחָ֖ק אַחֲרֵ֥י אִמּֽוֹ׃ \{פ\} \Rashi{\textbf{האהלה שרה אמו.}ויביאה האהלה ונעשית דגמת שרה אמו, כלומר והרי היא שרה אמו, שכל זמן ששרה קימת היה נר דלוק מערב שבת לערב שבת וברכה מצויה בעסה וענן קשור על האהל, ומשמתה פסקו, וכשבאת רבקה חזרו (בראשית רבה): אחרי אמו. דרך ארץ כל זמן שאמו של אדם קימת, כרוך הוא אצלה; ומשמתה, הוא מתנחם באשתו:}%
% ---------- פרק 25 | פסוק 1 ----------
 \textbf{א} וַיֹּ֧סֶף אַבְרָהָ֛ם וַיִּקַּ֥ח אִשָּׁ֖ה וּשְׁמָ֥הּ קְטוּרָֽה׃ \Rashi{\textbf{קטורה.}זו הגר, ונקראת קטורה על שנאים מעשיה כקטרת (בראשית רבה), ושקשרה פתחה, שלא נזדוגה לאדם מיום שפרשה מאברהם:}%
% ---------- פרק 25 | פסוק 2 ----------
 \textbf{ב} וַתֵּ֣לֶד ל֗וֹ אֶת־זִמְרָן֙ וְאֶת־יׇקְשָׁ֔ן וְאֶת־מְדָ֖ן וְאֶת־מִדְיָ֑ן וְאֶת־יִשְׁבָּ֖ק וְאֶת־שֽׁוּחַ׃ %
% ---------- פרק 25 | פסוק 3 ----------
 \textbf{ג} וְיׇקְשָׁ֣ן יָלַ֔ד אֶת־שְׁבָ֖א וְאֶת־דְּדָ֑ן וּבְנֵ֣י דְדָ֔ן הָי֛וּ אַשּׁוּרִ֥ם וּלְטוּשִׁ֖ם וּלְאֻמִּֽים׃ \Rashi{\textbf{אשורים ולטושים.}שם ראשי אמות (בראשית רבה), ותרגום של אנקלוס אין לי לישבו על לשון המקרא שפרש למשירין, לשון מחנה, ואם תאמר שאינו כן מפני האל"ף שאינה יסודית, הרי לנו תבות שאין בראשם אל"ף ונתוספה אל"ף בראשם, כמו חומת אנך (עמוס ז) שהוא מן נכה רגלים, וכמו אסוך שמן (מ"ב ד), שהוא מן ורחצת וסכת: ולטושם. הם בעלי אהלים המתפזרים אנה ואנה ונוסעים איש באהלי אפדנו, וכן הוא אומר והנם נטשים על פני כל הארץ (שמואל א ל') שכן למ"ד ונו"ן מתחלפות זו בזו):}%
% ---------- פרק 25 | פסוק 4 ----------
 \textbf{ד} וּבְנֵ֣י מִדְיָ֗ן עֵיפָ֤ה וָעֵ֙פֶר֙ וַחֲנֹ֔ךְ וַאֲבִידָ֖ע וְאֶלְדָּעָ֑ה כׇּל־אֵ֖לֶּה בְּנֵ֥י קְטוּרָֽה׃ %
% ---------- פרק 25 | פסוק 5 ----------
 \textbf{ה} וַיִּתֵּ֧ן אַבְרָהָ֛ם אֶת־כׇּל־אֲשֶׁר־ל֖וֹ לְיִצְחָֽק׃ \Rashi{\textbf{ויתן אברהם וגו'.}אמר ר' נחמיה ברכה דיאתיקי, שאמר לו הקב"ה לאברהם והיה ברכה, הברכות מסורות בידך לברך את מי שתרצה, ואברהם מסרן ליצחק (בראשית רבה):}%
% ---------- פרק 25 | פסוק 6 ----------
 \textbf{ו} וְלִבְנֵ֤י הַפִּֽילַגְשִׁים֙ אֲשֶׁ֣ר לְאַבְרָהָ֔ם נָתַ֥ן אַבְרָהָ֖ם מַתָּנֹ֑ת וַֽיְשַׁלְּחֵ֞ם מֵעַ֨ל יִצְחָ֤ק בְּנוֹ֙ בְּעוֹדֶ֣נּוּ חַ֔י קֵ֖דְמָה אֶל־אֶ֥רֶץ קֶֽדֶם׃ \Rashi{\textbf{הפילגשים.}חסר כתיב, שלא היתה אלא פלגש אחת, היא הגר, היא קטורה (בראשית רבה). נשים בכתבה, פילגשים בלא כתבה, כדאמרינן בסנהדרין, בנשים ופילגשים דדוד: נתן אברהם מתנות. פרשו רבותינו שם טמאה מסר להם (סנהדרין צ"א). ד"א מה שנתן לו על אודות שרה ושאר מתנות שנתנו לו, הכל נתן להם, שלא רצה להנות מהם:}%
% ---------- פרק 25 | פסוק 7 ----------
 \textbf{ז} וְאֵ֗לֶּה יְמֵ֛י שְׁנֵֽי־חַיֵּ֥י אַבְרָהָ֖ם אֲשֶׁר־חָ֑י מְאַ֥ת שָׁנָ֛ה וְשִׁבְעִ֥ים שָׁנָ֖ה וְחָמֵ֥שׁ שָׁנִֽים׃ \Rashi{\textbf{מאת שנה ושבעים שנה וחמש שנים.}בן ק' כבן ע' ובן ע' כבן ה' בלא חטא:}%
% ---------- פרק 25 | פסוק 8 ----------
 \textbf{ח} וַיִּגְוַ֨ע וַיָּ֧מׇת אַבְרָהָ֛ם בְּשֵׂיבָ֥ה טוֹבָ֖ה זָקֵ֣ן וְשָׂבֵ֑עַ וַיֵּאָ֖סֶף אֶל־עַמָּֽיו׃ %
% ---------- פרק 25 | פסוק 9 ----------
 \textbf{ט} וַיִּקְבְּר֨וּ אֹת֜וֹ יִצְחָ֤ק וְיִשְׁמָעֵאל֙ בָּנָ֔יו אֶל־מְעָרַ֖ת הַמַּכְפֵּלָ֑ה אֶל־שְׂדֵ֞ה עֶפְרֹ֤ן בֶּן־צֹ֙חַר֙ הַֽחִתִּ֔י אֲשֶׁ֖ר עַל־פְּנֵ֥י מַמְרֵֽא׃ \Rashi{\textbf{יצחק וישמעאל.}מכאן שעשה ישמעאל תשובה והוליך את יצחק לפניו, והיא שיבה טובה שנאמר באברהם (בראשית רבה):}%
% ---------- פרק 25 | פסוק 10 ----------
 \textbf{י} הַשָּׂדֶ֛ה אֲשֶׁר־קָנָ֥ה אַבְרָהָ֖ם מֵאֵ֣ת בְּנֵי־חֵ֑ת שָׁ֛מָּה קֻבַּ֥ר אַבְרָהָ֖ם וְשָׂרָ֥ה אִשְׁתּֽוֹ׃ %
% ---------- פרק 25 | פסוק 11 ----------
 \textbf{יא} וַיְהִ֗י אַחֲרֵי֙ מ֣וֹת אַבְרָהָ֔ם וַיְבָ֥רֶךְ אֱלֹהִ֖ים אֶת־יִצְחָ֣ק בְּנ֑וֹ וַיֵּ֣שֶׁב יִצְחָ֔ק עִם־בְּאֵ֥ר לַחַ֖י רֹאִֽי׃ \{פ\} \Rashi{\textbf{ויהי אחרי מות אברהם ויברך וגו'.}נחמו תנחומי אבלים. ד"א אע"פ שמסר הקב"ה את הברכות לאברהם, נתירא לברך את יצחק מפני שצפה את עשו יוצא ממנו, אמר יבא בעל הברכות ויברך את אשר ייטב בעיניו, ובא הקב"ה וברכו (סוטה י"ד):}%
% ---------- פרק 25 | פסוק 12 ----------
 \textbf{יב} וְאֵ֛לֶּה תֹּלְדֹ֥ת יִשְׁמָעֵ֖אל בֶּן־אַבְרָהָ֑ם אֲשֶׁ֨ר יָלְדָ֜ה הָגָ֧ר הַמִּצְרִ֛ית שִׁפְחַ֥ת שָׂרָ֖ה לְאַבְרָהָֽם׃ %
% ---------- פרק 25 | פסוק 13 ----------
 \textbf{יג} וְאֵ֗לֶּה שְׁמוֹת֙ בְּנֵ֣י יִשְׁמָעֵ֔אל בִּשְׁמֹתָ֖ם לְתוֹלְדֹתָ֑ם בְּכֹ֤ר יִשְׁמָעֵאל֙ נְבָיֹ֔ת וְקֵדָ֥ר וְאַדְבְּאֵ֖ל וּמִבְשָֽׂם׃ \Rashi{\textbf{בשמותם לתולדותם.}סדר לידתן זה אחר זה:}%
% ---------- פרק 25 | פסוק 14 ----------
 \textbf{יד} וּמִשְׁמָ֥ע וְדוּמָ֖ה וּמַשָּֽׂא׃ %
% ---------- פרק 25 | פסוק 15 ----------
 \textbf{טו} חֲדַ֣ד וְתֵימָ֔א יְט֥וּר נָפִ֖ישׁ וָקֵֽדְמָה׃ %
% ---------- פרק 25 | פסוק 16 ----------
 \textbf{טז} אֵ֣לֶּה הֵ֞ם בְּנֵ֤י יִשְׁמָעֵאל֙ וְאֵ֣לֶּה שְׁמֹתָ֔ם בְּחַצְרֵיהֶ֖ם וּבְטִֽירֹתָ֑ם שְׁנֵים־עָשָׂ֥ר נְשִׂיאִ֖ם לְאֻמֹּתָֽם׃ \Rashi{\textbf{בחצריהם.}כרכים שאין להם חומה, ותרגומו בפצחיהון שהם מפצחים, לשון פתיחה, כמו פצחו ורננו (תהילים צ"ח):}%
% ---------- פרק 25 | פסוק 17 ----------
 \textbf{יז} וְאֵ֗לֶּה שְׁנֵי֙ חַיֵּ֣י יִשְׁמָעֵ֔אל מְאַ֥ת שָׁנָ֛ה וּשְׁלֹשִׁ֥ים שָׁנָ֖ה וְשֶׁ֣בַע שָׁנִ֑ים וַיִּגְוַ֣ע וַיָּ֔מׇת וַיֵּאָ֖סֶף אֶל־עַמָּֽיו׃ \Rashi{\textbf{ואלה שני חיי ישמעאל וגו'.}אמר רבי חיא בר אבא למה נמנו שנותיו של ישמעאל, כדי ליחס בהם שנותיו של יעקב משנותיו של ישמעאל, למדנו ששמש יעקב בבית עבר י"ד שנה כשפרש מאביו קדם שבא אצל לבן, שהרי כשפרש יעקב מאביו מת ישמעאל, שנאמר וילך עשו אל ישמעאל וגו' (בראשית רבה כ"ה), כמו שמפרש בסוף מגלה נקראת: ויגוע. לא נאמרה גויעה אלא בצדיקים:}%
% ---------- פרק 25 | פסוק 18 ----------
 \textbf{יח} וַיִּשְׁכְּנ֨וּ מֵֽחֲוִילָ֜ה עַד־שׁ֗וּר אֲשֶׁר֙ עַל־פְּנֵ֣י מִצְרַ֔יִם בֹּאֲכָ֖ה אַשּׁ֑וּרָה עַל־פְּנֵ֥י כׇל־אֶחָ֖יו נָפָֽל׃ \{פ\} \Rashi{\textbf{נפל.}שכן, כמו ומדין ועמלק ובני קדם נפלים בעמק (שופטים ז'). כאן הוא אומר לשון נפילה, ולהלן אומר על פני כל אחיו ישכן (בראשית ט״ז:י״ב)? עד שלא מת אברהם ישכן, משמת אברהם נפל:}%
\addparasha{חומש בראשית | פרשת תולדות}
% ---------- פרק 25 | פסוק 19 ----------
 \textbf{יט} וְאֵ֛לֶּה תּוֹלְדֹ֥ת יִצְחָ֖ק בֶּן־אַבְרָהָ֑ם אַבְרָהָ֖ם הוֹלִ֥יד אֶת־יִצְחָֽק׃ \Rashi{\textbf{ואלה תולדות יצחק.}יעקב ועשו האמורים בפרשה: אברהם הוליד את יצחק. על ידי שכתב הכתוב יצחק בן אברהם הזקק לומר אברהם הוליד את יצחק; לפי שהיו ליצני הדור אומרים מאבימלך נתעברה שרה, שהרי כמה שנים שהתה עם אברהם ולא נתעברה הימנו; מה עשה הקב"ה? צר קלסתר פניו של יצחק דומה לאברהם, והעידו הכל אברהם הוליד את יצחק, וזהו שכתוב כאן יצחק בן אברהם, שהרי עדות יש שאברהם הוליד את יצחק:}%
% ---------- פרק 25 | פסוק 20 ----------
 \textbf{כ} וַיְהִ֤י יִצְחָק֙ בֶּן־אַרְבָּעִ֣ים שָׁנָ֔ה בְּקַחְתּ֣וֹ אֶת־רִבְקָ֗ה בַּת־בְּתוּאֵל֙ הָֽאֲרַמִּ֔י מִפַּדַּ֖ן אֲרָ֑ם אֲח֛וֹת לָבָ֥ן הָאֲרַמִּ֖י ל֥וֹ לְאִשָּֽׁה׃ \Rashi{\textbf{בן ארבעים שנה.}שהרי כשבא אברהם מהר המוריה נתבשר שנולדה רבקה, ויצחק היה בן ל"ז שנה, שהרי בו בפרק מתה שרה, ומשנולד יצחק עד העקדה שמתה שרה, ל"ז שנה היו – כי בת צ' היתה כשנולד יצחק, ובת קכ"ז כשמתה – שנאמר ויהיו חיי שרה וגו' הרי ליצחק ל"ז שנים, ובו בפרק נולדה רבקה; המתין לה עד שתהא ראויה לביאה ג' שנים ונשאה: בת בתואל מפדן ארם אחות לבן. וכי עדין לא נכתב שהיא בת בתואל ואחות לבן ומפדן ארם? אלא להגיד שבחה, שהיתה בת רשע ואחות רשע ומקומה אנשי רשע ולא למדה ממעשיהם: מפדן ארם. על שם ששני ארם היו, ארם נהרים וארם צובה, קורא אותו פדן, לשון צמד בקר, תרגום פדן תורין; ויש פותרין פדן ארם כמו שדה ארם, שבלשון ישמעאל קורין לשדה פדן:}%
% ---------- פרק 25 | פסוק 21 ----------
 \textbf{כא} וַיֶּעְתַּ֨ר יִצְחָ֤ק לַֽיהֹוָה֙ לְנֹ֣כַח אִשְׁתּ֔וֹ כִּ֥י עֲקָרָ֖ה הִ֑וא וַיֵּעָ֤תֶר לוֹ֙ יְהֹוָ֔ה וַתַּ֖הַר רִבְקָ֥ה אִשְׁתּֽוֹ׃ \Rashi{\textbf{ויעתר.}הרבה והפציר בתפלה: ויעתר לו. נתפצר ונתפיס ונתפתה לו. ואומר אני, כל לשון עתר לשון הפצרה ורבוי הוא, וכן ועתר ענן הקטרת (יחז' ח') – מרבית עלית העשן, וכן והעתרתם עלי דבריכם (שם ל"ה), וכן ונעתרות נשיקות שונא (משלי כ"ז) – דומות למרבות והנם למשא; אנקרי"שרא בלע"ז: לנכח אשתו. זה עומד בזוית זו ומתפלל וזו עומדת בזוית זו ומתפללת: ויעתר לו. לו ולא לה, שאין דומה תפלת צדיק בן צדיק לתפלת צדיק בן רשע, לפיכך לו ולא לה (יבמות ס"ד):}%
% ---------- פרק 25 | פסוק 22 ----------
 \textbf{כב} וַיִּתְרֹֽצְצ֤וּ הַבָּנִים֙ בְּקִרְבָּ֔הּ וַתֹּ֣אמֶר אִם־כֵּ֔ן לָ֥מָּה זֶּ֖ה אָנֹ֑כִי וַתֵּ֖לֶךְ לִדְרֹ֥שׁ אֶת־יְהֹוָֽה׃ \Rashi{\textbf{ויתרוצצו.}ע"כ המקרא הזה אומר דרשני, שסתם מה היא רציצה זו וכתב אם כן למה זה אנכי? רבותינו דרשוהו לשון ריצה; כשהיתה עוברת על פתחי תורה של שם ועבר יעקב רץ ומפרכס לצאת, עוברת על פתח עבודת אלילים, עשו מפרכס לצאת. דבר אחר מתרוצצים זה עם זה ומריבים בנחלת שני עולמות. ותאמר אם כן. גדול צער העבור. למה זה אנכי. מתאוה ומתפללת על הריון: ותלך לדרוש. לבית מדרשו של שם: לדרוש את ה'. שיגיד לה מה תהא בסופה:}%
% ---------- פרק 25 | פסוק 23 ----------
 \textbf{כג} וַיֹּ֨אמֶר יְהֹוָ֜ה לָ֗הּ שְׁנֵ֤י (גיים) [גוֹיִם֙] בְּבִטְנֵ֔ךְ וּשְׁנֵ֣י לְאֻמִּ֔ים מִמֵּעַ֖יִךְ יִפָּרֵ֑דוּ וּלְאֹם֙ מִלְאֹ֣ם יֶֽאֱמָ֔ץ וְרַ֖ב יַעֲבֹ֥ד צָעִֽיר׃ \Rashi{\textbf{ויאמר ה' לה.}על ידי שליח; לשם נאמר ברוח הקדש והוא אמר לה (בראשית רבה): שני גוים בבטנך. גיים כתיב; אלו אנטונינוס ורבי, שלא פסקו מעל שלחנם לא צנון ולא חזרת לא בימות החמה ולא בימות הגשמים (עבודה זרה י"א): ושני לאמים. אין לאם אלא מלכות: ממעיך יפרדו. מן המעים הם נפרדים זה לרשעו וזה לתמו: מלאם יאמץ. לא ישוו בגדלה, כשזה קם זה נופל, וכן הוא אומר אמלאה החרבה (יחזקאל כ"ו), לא נתמלאה צור אלא מחרבנה של ירושלים:}%
% ---------- פרק 25 | פסוק 24 ----------
 \textbf{כד} וַיִּמְלְא֥וּ יָמֶ֖יהָ לָלֶ֑דֶת וְהִנֵּ֥ה תוֹמִ֖ם בְּבִטְנָֽהּ׃ \Rashi{\textbf{וימלאו ימיה.}אבל בתמר כתיב ויהי בעת לדתה (בראשית ל״ח:כ״ז), שלא מלאו ימיה, כי לז' חדשים ילדתם: והנה תומם. חסר, ובתמר תאומים מלא, לפי ששניהם צדיקים; אבל כאן אחד צדיק ואחד רשע:}%
% ---------- פרק 25 | פסוק 25 ----------
 \textbf{כה} וַיֵּצֵ֤א הָרִאשׁוֹן֙ אַדְמוֹנִ֔י כֻּלּ֖וֹ כְּאַדֶּ֣רֶת שֵׂעָ֑ר וַיִּקְרְא֥וּ שְׁמ֖וֹ עֵשָֽׂו׃ \Rashi{\textbf{אדמוני.}סימן הוא שיהא שופך דמים (בראשית רבה): כלו כאדרת שער. מלא שער כטלית של צמר המלאה שער, פלוקי"דא בלע"ז: ויקראו שמו עשו. הכל קראו לו כן, לפי שהיה נעשה ונגמר בשערו כבן שנים הרבה:}%
% ---------- פרק 25 | פסוק 26 ----------
 \textbf{כו} וְאַֽחֲרֵי־כֵ֞ן יָצָ֣א אָחִ֗יו וְיָד֤וֹ אֹחֶ֙זֶת֙ בַּעֲקֵ֣ב עֵשָׂ֔ו וַיִּקְרָ֥א שְׁמ֖וֹ יַעֲקֹ֑ב וְיִצְחָ֛ק בֶּן־שִׁשִּׁ֥ים שָׁנָ֖ה בְּלֶ֥דֶת אֹתָֽם׃ \Rashi{\textbf{ואחרי כן יצא אחיו וגו'.}שמעתי מ"א הדורשו לפי פשוטו, בדין היה אוחז בו לעכבו, יעקב נוצר מטפה ראשונה ועשו מן השניה; צא ולמד משפופרת שפיה קצרה, תן לה שתי אבנים זו אחר זו – הנכנסת ראשונה תצא אחרונה והנכנסת אחרונה תצא ראשונה; נמצא עשו הנוצר באחרונה יצא ראשון ויעקב שנוצר ראשונה יצא אחרון, ויעקב בא לעכבו, שיהא ראשון ללדה כראשון ליצירה, ויפטר את רחמה ויטל את הבכורה מן הדין: בעקב עשו. סימן שאין זה מספיק לגמר מלכותו עד שזה עומד ונוטלה הימנו: ויקרא שמו יעקב. הקב"ה; ד"א אביו קרא לו יעקב על שם אחיזת העקב: בן ששים שנה. י' שנים משנשאה עד שנעשית בת י"ג שנה וראויה להריון, וי' שנים הללו צפה והמתין לה, כמו שעשה אביו לשרה; כיון שלא נתעברה, ידע שהיא עקרה, והתפלל עליה (יבמות ס"ד), ושפחה לא רצה לשא, לפי שנתקדש בהר המוריה להיות עולה תמימה:}%
% ---------- פרק 25 | פסוק 27 ----------
 \textbf{כז} וַֽיִּגְדְּלוּ֙ הַנְּעָרִ֔ים וַיְהִ֣י עֵשָׂ֗ו אִ֛ישׁ יֹדֵ֥עַ צַ֖יִד אִ֣ישׁ שָׂדֶ֑ה וְיַעֲקֹב֙ אִ֣ישׁ תָּ֔ם יֹשֵׁ֖ב אֹהָלִֽים׃ \Rashi{\textbf{ויגדלו … ויהי עשו.}כל זמן שהיו קטנים, לא היו נכרים במעשיהם, ואין אדם מדקדק בהם מה טיבם; כיון שנעשו בני שלש עשרה שנה, זה פרש לבתי מדרשות וזה פרש לע"ז: ידע ציד. לצוד ולרמות את אביו בפיו ושואלו אבא, האיך מעשרין את המלח ואת התבן? כסבור אביו שהוא מדקדק במצות (תנחומא): איש שדה. כמשמעו, אדם בטל וצודה בקשתו חיות ועופות: תם. אינו בקי בכל אלה, כלבו כן פיו, מי שאינו חריף לרמות קרוי תם: ישב אהלים. אהלו של שם ואהלו של עבר:}%
% ---------- פרק 25 | פסוק 28 ----------
 \textbf{כח} וַיֶּאֱהַ֥ב יִצְחָ֛ק אֶת־עֵשָׂ֖ו כִּי־צַ֣יִד בְּפִ֑יו וְרִבְקָ֖ה אֹהֶ֥בֶת אֶֽת־יַעֲקֹֽב׃ \Rashi{\textbf{בפיו.}כתרגומו בפיו של יצחק. ומדרשו בפיו של עשו, שהיה צד אותו ומרמהו בדבריו:}%
% ---------- פרק 25 | פסוק 29 ----------
 \textbf{כט} וַיָּ֥זֶד יַעֲקֹ֖ב נָזִ֑יד וַיָּבֹ֥א עֵשָׂ֛ו מִן־הַשָּׂדֶ֖ה וְה֥וּא עָיֵֽף׃ \Rashi{\textbf{ויזד.}לשון בשול, כתרגומו: והוא עיף. ברציחה, כמה דתימא כי עיפה נפשי להרגים (ירמיהו ד') (בראשית רבה):}%
% ---------- פרק 25 | פסוק 30 ----------
 \textbf{ל} וַיֹּ֨אמֶר עֵשָׂ֜ו אֶֽל־יַעֲקֹ֗ב הַלְעִיטֵ֤נִי נָא֙ מִן־הָאָדֹ֤ם הָאָדֹם֙ הַזֶּ֔ה כִּ֥י עָיֵ֖ף אָנֹ֑כִי עַל־כֵּ֥ן קָרָֽא־שְׁמ֖וֹ אֱדֽוֹם׃ \Rashi{\textbf{הלעיטני.}אפתח פי ושפך הרבה לתוכה, כמו ששנינו אין אובסין את הגמל אבל מלעיטין אותו: מן האדם האדם. עדשים אדמות, ואותו היום מת אברהם, שלא יראה את עשו בן בנו יוצא לתרבות רעה, ואין זו שיבה טובה שהבטיחו הקב"ה; לפיכך קצר הקב"ה ה' שנים משנותיו, שיצחק חי ק"פ שנה וזה קע"ה, ובשל יעקב עדשים להברות את האבל. ולמה עדשים? שדומות לגלגל, שהאבלות גלגל החוזר בעולם (ועוד מה עדשים אין להם פה, כך האבל אין לו פה, שאסור לדבר, ולפיכך המנהג להברות האבל בתחלת מאכלו ביצים, שהם עגלים ואין להם פה, כך אבל אין לו פה, כדאמרינן במועד קטן (מועד קטן כא:) אבל כל שלשה ימים הראשונים אינו משיב שלום לכל אדם וכ"ש שאינו שואל בתחלה, מג' ועד ז' משיב ואינו שואל וכו'; ברש"י ישן):}%
% ---------- פרק 25 | פסוק 31 ----------
 \textbf{לא} וַיֹּ֖אמֶר יַעֲקֹ֑ב מִכְרָ֥ה כַיּ֛וֹם אֶת־בְּכֹרָתְךָ֖ לִֽי׃ \Rashi{\textbf{מכרה כיום.}כתרגומו כיום דלהן, כיום שהוא ברור, כך מכר לי מכירה ברורה: בכרתך. לפי שהעבודה בבכורות, אמר יעקב אין רשע זה כדאי שיקריב להקב"ה:}%
% ---------- פרק 25 | פסוק 32 ----------
 \textbf{לב} וַיֹּ֣אמֶר עֵשָׂ֔ו הִנֵּ֛ה אָנֹכִ֥י הוֹלֵ֖ךְ לָמ֑וּת וְלָמָּה־זֶּ֥ה לִ֖י בְּכֹרָֽה׃ \Rashi{\textbf{הנה אנכי הולך למות.}(מתנודדת והולכת היא הבכורה שלא תהא כל עת העבודה בבכורות, כי שבט לוי יטל אותה, ועוד) אמר עשו: מה טיבה של עבודה זו? אמר לו: כמה אזהרות וענשין ומיתות תלויין בה כאותה ששנינו, אלו הן שבמיתה: שתויי יין ופרועי ראש; אמר: אני הולך למות על ידה, אם כן מה חפץ לי בה?}%
% ---------- פרק 25 | פסוק 33 ----------
 \textbf{לג} וַיֹּ֣אמֶר יַעֲקֹ֗ב הִשָּׁ֤בְעָה לִּי֙ כַּיּ֔וֹם וַיִּשָּׁבַ֖ע ל֑וֹ וַיִּמְכֹּ֥ר אֶת־בְּכֹרָת֖וֹ לְיַעֲקֹֽב׃ %
% ---------- פרק 25 | פסוק 34 ----------
 \textbf{לד} וְיַעֲקֹ֞ב נָתַ֣ן לְעֵשָׂ֗ו לֶ֚חֶם וּנְזִ֣יד עֲדָשִׁ֔ים וַיֹּ֣אכַל וַיֵּ֔שְׁתְּ וַיָּ֖קׇם וַיֵּלַ֑ךְ וַיִּ֥בֶז עֵשָׂ֖ו אֶת־הַבְּכֹרָֽה׃ \{פ\} \Rashi{\textbf{ויבז עשו.}העיד הכתוב על רשעו שבזה עבודתו של מקום:}%
% ---------- פרק 26 | פסוק 1 ----------
 \textbf{א} וַיְהִ֤י רָעָב֙ בָּאָ֔רֶץ מִלְּבַד֙ הָרָעָ֣ב הָרִאשׁ֔וֹן אֲשֶׁ֥ר הָיָ֖ה בִּימֵ֣י אַבְרָהָ֑ם וַיֵּ֧לֶךְ יִצְחָ֛ק אֶל־אֲבִימֶ֥לֶךְ מֶֽלֶךְ־פְּלִשְׁתִּ֖ים גְּרָֽרָה׃ %
% ---------- פרק 26 | פסוק 2 ----------
 \textbf{ב} וַיֵּרָ֤א אֵלָיו֙ יְהֹוָ֔ה וַיֹּ֖אמֶר אַל־תֵּרֵ֣ד מִצְרָ֑יְמָה שְׁכֹ֣ן בָּאָ֔רֶץ אֲשֶׁ֖ר אֹמַ֥ר אֵלֶֽיךָ׃ \Rashi{\textbf{אל תרד מצרימה.}שהיה דעתו לרדת מצרימה, כמו שירד אביו בימי הרעב; אמר לו: אל תרד מצרימה, שאתה עולה תמימה, ואין חוצה לארץ כדאי לך:}%
% ---------- פרק 26 | פסוק 3 ----------
 \textbf{ג} גּ֚וּר בָּאָ֣רֶץ הַזֹּ֔את וְאֶֽהְיֶ֥ה עִמְּךָ֖ וַאֲבָרְכֶ֑ךָּ כִּֽי־לְךָ֣ וּֽלְזַרְעֲךָ֗ אֶתֵּן֙ אֶת־כׇּל־הָֽאֲרָצֹ֣ת הָאֵ֔ל וַהֲקִֽמֹתִי֙ אֶת־הַשְּׁבֻעָ֔ה אֲשֶׁ֥ר נִשְׁבַּ֖עְתִּי לְאַבְרָהָ֥ם אָבִֽיךָ׃ \Rashi{\textbf{האל.}כמו האלה:}%
% ---------- פרק 26 | פסוק 4 ----------
 \textbf{ד} וְהִרְבֵּיתִ֤י אֶֽת־זַרְעֲךָ֙ כְּכוֹכְבֵ֣י הַשָּׁמַ֔יִם וְנָתַתִּ֣י לְזַרְעֲךָ֔ אֵ֥ת כׇּל־הָאֲרָצֹ֖ת הָאֵ֑ל וְהִתְבָּרְכ֣וּ בְזַרְעֲךָ֔ כֹּ֖ל גּוֹיֵ֥י הָאָֽרֶץ׃ \Rashi{\textbf{והתברכו בזרעך.}אדם אומר לבנו יהא זרעך כזרעו של יצחק, וכן בכל המקרא, וזה אב לכלן בך יברך ישראל לאמר ישמך וגו' (בראשית מ״ח:כ׳); ואף לענין הקללה מצינו כן והיתה האשה לאלה (במדבר ה'), שהמקלל שונאו אומר תהא כפלונית, וכן והנחתם שמכם לשבועה לבחירי (ישעיהו ס"ה), שהנשבע אומר אהא כפלוני אם עשיתי כך וכך:}%
% ---------- פרק 26 | פסוק 5 ----------
 \textbf{ה} עֵ֕קֶב אֲשֶׁר־שָׁמַ֥ע אַבְרָהָ֖ם בְּקֹלִ֑י וַיִּשְׁמֹר֙ מִשְׁמַרְתִּ֔י מִצְוֺתַ֖י חֻקּוֹתַ֥י וְתוֹרֹתָֽי׃ \Rashi{\textbf{שמע אברהם בקולי.}כשנסיתי אותו: וישמר משמרתי. גזרות להרחקה על האזהרות שבתורה, כגון שניות לעריות ושבות לשבת (יבמות כ"א): מצותי. דברים שאלו לא נכתבו ראוין הם להצטוות, כגון גזל ושפיכות דמים: חקותי. דברים שיצר הרע ואמות העולם משיבין עליהם, כגון אכילת חזיר ולבישת שעטנז, שאין טעם בדבר, אלא גזרת המלך וחקותיו על עבדיו: ותורתי. להביא תורה שבעל פה הלכה למשה מסיני (בראשית רבה):}%
% ---------- פרק 26 | פסוק 6 ----------
 \textbf{ו} וַיֵּ֥שֶׁב יִצְחָ֖ק בִּגְרָֽר׃ %
% ---------- פרק 26 | פסוק 7 ----------
 \textbf{ז} וַֽיִּשְׁאֲל֞וּ אַנְשֵׁ֤י הַמָּקוֹם֙ לְאִשְׁתּ֔וֹ וַיֹּ֖אמֶר אֲחֹ֣תִי הִ֑וא כִּ֤י יָרֵא֙ לֵאמֹ֣ר אִשְׁתִּ֔י פֶּן־יַֽהַרְגֻ֜נִי אַנְשֵׁ֤י הַמָּקוֹם֙ עַל־רִבְקָ֔ה כִּֽי־טוֹבַ֥ת מַרְאֶ֖ה הִֽוא׃ \Rashi{\textbf{לאשתו.}על אשתו, כמו אמרי לי אחי הוא (בראשית כ׳:י״ג):}%
% ---------- פרק 26 | פסוק 8 ----------
 \textbf{ח} וַיְהִ֗י כִּ֣י אָֽרְכוּ־ל֥וֹ שָׁם֙ הַיָּמִ֔ים וַיַּשְׁקֵ֗ף אֲבִימֶ֙לֶךְ֙ מֶ֣לֶךְ פְּלִשְׁתִּ֔ים בְּעַ֖ד הַֽחַלּ֑וֹן וַיַּ֗רְא וְהִנֵּ֤ה יִצְחָק֙ מְצַחֵ֔ק אֵ֖ת רִבְקָ֥ה אִשְׁתּֽוֹ׃ \Rashi{\textbf{כי ארכו.}אמר, מעתה אין לי לדאג, מאחר שלא אנסוה עד עכשו, ולא נזהר להיות נשמר: וישקף אבימלך וגומר. ראהו משמש מטתו:}%
% ---------- פרק 26 | פסוק 9 ----------
 \textbf{ט} וַיִּקְרָ֨א אֲבִימֶ֜לֶךְ לְיִצְחָ֗ק וַיֹּ֙אמֶר֙ אַ֣ךְ הִנֵּ֤ה אִשְׁתְּךָ֙ הִ֔וא וְאֵ֥יךְ אָמַ֖רְתָּ אֲחֹ֣תִי הִ֑וא וַיֹּ֤אמֶר אֵלָיו֙ יִצְחָ֔ק כִּ֣י אָמַ֔רְתִּי פֶּן־אָמ֖וּת עָלֶֽיהָ׃ %
% ---------- פרק 26 | פסוק 10 ----------
 \textbf{י} וַיֹּ֣אמֶר אֲבִימֶ֔לֶךְ מַה־זֹּ֖את עָשִׂ֣יתָ לָּ֑נוּ כִּ֠מְעַ֠ט שָׁכַ֞ב אַחַ֤ד הָעָם֙ אֶת־אִשְׁתֶּ֔ךָ וְהֵבֵאתָ֥ עָלֵ֖ינוּ אָשָֽׁם׃ \Rashi{\textbf{אחד העם.}המיחד בעם, זה המלך: והבאת עלינו אשם. אם שכב כבר, הבאת עלינו אשם:}%
% ---------- פרק 26 | פסוק 11 ----------
 \textbf{יא} וַיְצַ֣ו אֲבִימֶ֔לֶךְ אֶת־כׇּל־הָעָ֖ם לֵאמֹ֑ר הַנֹּגֵ֜עַ בָּאִ֥ישׁ הַזֶּ֛ה וּבְאִשְׁתּ֖וֹ מ֥וֹת יוּמָֽת׃ %
% ---------- פרק 26 | פסוק 12 ----------
 \textbf{יב} וַיִּזְרַ֤ע יִצְחָק֙ בָּאָ֣רֶץ הַהִ֔וא וַיִּמְצָ֛א בַּשָּׁנָ֥ה הַהִ֖וא מֵאָ֣ה שְׁעָרִ֑ים וַֽיְבָרְכֵ֖הוּ יְהֹוָֽה׃ \Rashi{\textbf{בארץ ההוא.}אע"פ שאינו חשובה כארץ ישראל עצמה, כארץ שבעה גוים (בראשית רבה): בשנה ההיא. אע"פ שאינה כתקנה, שהיתה שנת רעבון (בראשית רבה): בארץ ההיא בשנה ההיא. שניהם למה? לומר, שהארץ קשה והשנה קשה: מאה שערים. שאמדוה כמה ראויה לעשות, ועשתה על אחת שאמדוה מאה. ורבותינו אמרו אמד זה למעשרות היה:}%
% ---------- פרק 26 | פסוק 13 ----------
 \textbf{יג} וַיִּגְדַּ֖ל הָאִ֑ישׁ וַיֵּ֤לֶךְ הָלוֹךְ֙ וְגָדֵ֔ל עַ֥ד כִּֽי־גָדַ֖ל מְאֹֽד׃ \Rashi{\textbf{כי גדל מאוד.}שהיו אומרים זבל פרדותיו של יצחק ולא כספו וזהבו של אבימלך (בראשית רבה):}%
% ---------- פרק 26 | פסוק 14 ----------
 \textbf{יד} וַֽיְהִי־ל֤וֹ מִקְנֵה־צֹאן֙ וּמִקְנֵ֣ה בָקָ֔ר וַעֲבֻדָּ֖ה רַבָּ֑ה וַיְקַנְא֥וּ אֹת֖וֹ פְּלִשְׁתִּֽים׃ \Rashi{\textbf{ועבדה רבה.}פעלה רבה בלשון לע"ז אוברי"נא, עבודה משמע עבודה אחת; עבדה משמע פעלה רבה:}%
% ---------- פרק 26 | פסוק 15 ----------
 \textbf{טו} וְכׇל־הַבְּאֵרֹ֗ת אֲשֶׁ֤ר חָֽפְרוּ֙ עַבְדֵ֣י אָבִ֔יו בִּימֵ֖י אַבְרָהָ֣ם אָבִ֑יו סִתְּמ֣וּם פְּלִשְׁתִּ֔ים וַיְמַלְא֖וּם עָפָֽר׃ \Rashi{\textbf{סתמום פלשתים.}מפני שאמרו תקלה הם לנו מפני הגיסות הבאות עלינו, טמונון פלשתאי לשון סתימה, ובלשון התלמוד מטמטם את הלב:}%
% ---------- פרק 26 | פסוק 16 ----------
 \textbf{טז} וַיֹּ֥אמֶר אֲבִימֶ֖לֶךְ אֶל־יִצְחָ֑ק לֵ֚ךְ מֵֽעִמָּ֔נוּ כִּֽי־עָצַ֥מְתָּ מִמֶּ֖נּוּ מְאֹֽד׃ %
% ---------- פרק 26 | פסוק 17 ----------
 \textbf{יז} וַיֵּ֥לֶךְ מִשָּׁ֖ם יִצְחָ֑ק וַיִּ֥חַן בְּנַֽחַל־גְּרָ֖ר וַיֵּ֥שֶׁב שָֽׁם׃ \Rashi{\textbf{בנחל גרר.}רחוק מן העיר:}%
% ---------- פרק 26 | פסוק 18 ----------
 \textbf{יח} וַיָּ֨שׇׁב יִצְחָ֜ק וַיַּחְפֹּ֣ר ׀ אֶת־בְּאֵרֹ֣ת הַמַּ֗יִם אֲשֶׁ֤ר חָֽפְרוּ֙ בִּימֵי֙ אַבְרָהָ֣ם אָבִ֔יו וַיְסַתְּמ֣וּם פְּלִשְׁתִּ֔ים אַחֲרֵ֖י מ֣וֹת אַבְרָהָ֑ם וַיִּקְרָ֤א לָהֶן֙ שֵׁמ֔וֹת כַּשֵּׁמֹ֕ת אֲשֶׁר־קָרָ֥א לָהֶ֖ן אָבִֽיו׃ \Rashi{\textbf{וישב ויחפר.}הבארות אשר חפרו בימי אברהם אביו ופלשתים סתמום, מקדם שנסע יצחק מגרר, חזר וחפרן:}%
% ---------- פרק 26 | פסוק 19 ----------
 \textbf{יט} וַיַּחְפְּר֥וּ עַבְדֵֽי־יִצְחָ֖ק בַּנָּ֑חַל וַיִּ֨מְצְאוּ־שָׁ֔ם בְּאֵ֖ר מַ֥יִם חַיִּֽים׃ %
% ---------- פרק 26 | פסוק 20 ----------
 \textbf{כ} וַיָּרִ֜יבוּ רֹעֵ֣י גְרָ֗ר עִם־רֹעֵ֥י יִצְחָ֛ק לֵאמֹ֖ר לָ֣נוּ הַמָּ֑יִם וַיִּקְרָ֤א שֵֽׁם־הַבְּאֵר֙ עֵ֔שֶׂק כִּ֥י הִֽתְעַשְּׂק֖וּ עִמּֽוֹ׃ \Rashi{\textbf{עשק.}ערעור: כי התעשקו עמו. נתעשקו עמו עליה במריבה וערעור:}%
% ---------- פרק 26 | פסוק 21 ----------
 \textbf{כא} וַֽיַּחְפְּרוּ֙ בְּאֵ֣ר אַחֶ֔רֶת וַיָּרִ֖יבוּ גַּם־עָלֶ֑יהָ וַיִּקְרָ֥א שְׁמָ֖הּ שִׂטְנָֽה׃ \Rashi{\textbf{שטנה.}נושימנ"ט:}%
% ---------- פרק 26 | פסוק 22 ----------
 \textbf{כב} וַיַּעְתֵּ֣ק מִשָּׁ֗ם וַיַּחְפֹּר֙ בְּאֵ֣ר אַחֶ֔רֶת וְלֹ֥א רָב֖וּ עָלֶ֑יהָ וַיִּקְרָ֤א שְׁמָהּ֙ רְחֹב֔וֹת וַיֹּ֗אמֶר כִּֽי־עַתָּ֞ה הִרְחִ֧יב יְהֹוָ֛ה לָ֖נוּ וּפָרִ֥ינוּ בָאָֽרֶץ׃ \Rashi{\textbf{ופרינו בארץ.}כתרגומו וניפוש בארעא:}%
% ---------- פרק 26 | פסוק 23 ----------
 \textbf{כג} וַיַּ֥עַל מִשָּׁ֖ם בְּאֵ֥ר שָֽׁבַע׃ %
% ---------- פרק 26 | פסוק 24 ----------
 \textbf{כד} וַיֵּרָ֨א אֵלָ֤יו יְהֹוָה֙ בַּלַּ֣יְלָה הַה֔וּא וַיֹּ֕אמֶר אָנֹכִ֕י אֱלֹהֵ֖י אַבְרָהָ֣ם אָבִ֑יךָ אַל־תִּירָא֙ כִּֽי־אִתְּךָ֣ אָנֹ֔כִי וּבֵֽרַכְתִּ֙יךָ֙ וְהִרְבֵּיתִ֣י אֶֽת־זַרְעֲךָ֔ בַּעֲב֖וּר אַבְרָהָ֥ם עַבְדִּֽי׃ %
% ---------- פרק 26 | פסוק 25 ----------
 \textbf{כה} וַיִּ֧בֶן שָׁ֣ם מִזְבֵּ֗חַ וַיִּקְרָא֙ בְּשֵׁ֣ם יְהֹוָ֔ה וַיֶּט־שָׁ֖ם אׇהֳל֑וֹ וַיִּכְרוּ־שָׁ֥ם עַבְדֵי־יִצְחָ֖ק בְּאֵֽר׃ %
% ---------- פרק 26 | פסוק 26 ----------
 \textbf{כו} וַאֲבִימֶ֕לֶךְ הָלַ֥ךְ אֵלָ֖יו מִגְּרָ֑ר וַאֲחֻזַּת֙ מֵֽרֵעֵ֔הוּ וּפִיכֹ֖ל שַׂר־צְבָאֽוֹ׃ \Rashi{\textbf{ואחזת מרעהו.}כתרגומו וסיעת מרחמוהי, סיעת מאוהביו. ויש פותרין מרעהו מ' מיסוד התבה, כמו שלשים מרעים דשמשון, כדי שתהיה תבת ואחזת דבוקה; אבל אין דרך ארץ לדבר על המלכות כן, סיעת אוהביו, שאם כן, כל סיעת אוהביו הוליך עמו, ולא היה לו אלא סיעה אחת של אוהבים – לכן יש לפתרו בלשון הראשון; ואל תתמה על תי"ו של אחזת ואע"פ שאין תבה סמוכה, יש דגמתה במקרא עזרת מצר (תהילים ס'), ושכרת ולא מיין (ישעיהו נ"א): אחזת. לשון קביצה ואגדה שנאחזין יחד:}%
% ---------- פרק 26 | פסוק 27 ----------
 \textbf{כז} וַיֹּ֤אמֶר אֲלֵהֶם֙ יִצְחָ֔ק מַדּ֖וּעַ בָּאתֶ֣ם אֵלָ֑י וְאַתֶּם֙ שְׂנֵאתֶ֣ם אֹתִ֔י וַתְּשַׁלְּח֖וּנִי מֵאִתְּכֶֽם׃ %
% ---------- פרק 26 | פסוק 28 ----------
 \textbf{כח} וַיֹּאמְר֗וּ רָא֣וֹ רָאִ֘ינוּ֮ כִּֽי־הָיָ֣ה יְהֹוָ֣ה ׀ עִמָּךְ֒ וַנֹּ֗אמֶר תְּהִ֨י נָ֥א אָלָ֛ה בֵּינוֹתֵ֖ינוּ בֵּינֵ֣ינוּ וּבֵינֶ֑ךָ וְנִכְרְתָ֥ה בְרִ֖ית עִמָּֽךְ׃ \Rashi{\textbf{ראו ראינו.}ראינו באביך, ראינו בך: תהי נא אלה בינותינו וגו'. האלה אשר בינותינו מימי אביך תהי גם עתה בינינו וביניך:}%
% ---------- פרק 26 | פסוק 29 ----------
 \textbf{כט} אִם־תַּעֲשֵׂ֨ה עִמָּ֜נוּ רָעָ֗ה כַּאֲשֶׁר֙ לֹ֣א נְגַֽעֲנ֔וּךָ וְכַאֲשֶׁ֨ר עָשִׂ֤ינוּ עִמְּךָ֙ רַק־ט֔וֹב וַנְּשַׁלֵּֽחֲךָ֖ בְּשָׁל֑וֹם אַתָּ֥ה עַתָּ֖ה בְּר֥וּךְ יְהֹוָֽה׃ \Rashi{\textbf{לא נגענוך.}כשאמרנו לך לך מעמנו: אתה. גם אתה עשה לנו כמו כן:}%
% ---------- פרק 26 | פסוק 30 ----------
 \textbf{ל} וַיַּ֤עַשׂ לָהֶם֙ מִשְׁתֶּ֔ה וַיֹּאכְל֖וּ וַיִּשְׁתּֽוּ׃ %
% ---------- פרק 26 | פסוק 31 ----------
 \textbf{לא} וַיַּשְׁכִּ֣ימוּ בַבֹּ֔קֶר וַיִּשָּׁבְע֖וּ אִ֣ישׁ לְאָחִ֑יו וַיְשַׁלְּחֵ֣ם יִצְחָ֔ק וַיֵּלְכ֥וּ מֵאִתּ֖וֹ בְּשָׁלֽוֹם׃ %
% ---------- פרק 26 | פסוק 32 ----------
 \textbf{לב} וַיְהִ֣י ׀ בַּיּ֣וֹם הַה֗וּא וַיָּבֹ֙אוּ֙ עַבְדֵ֣י יִצְחָ֔ק וַיַּגִּ֣דוּ ל֔וֹ עַל־אֹד֥וֹת הַבְּאֵ֖ר אֲשֶׁ֣ר חָפָ֑רוּ וַיֹּ֥אמְרוּ ל֖וֹ מָצָ֥אנוּ מָֽיִם׃ %
% ---------- פרק 26 | פסוק 33 ----------
 \textbf{לג} וַיִּקְרָ֥א אֹתָ֖הּ שִׁבְעָ֑ה עַל־כֵּ֤ן שֵׁם־הָעִיר֙ בְּאֵ֣ר שֶׁ֔בַע עַ֖ד הַיּ֥וֹם הַזֶּֽה׃ \{ס\} \Rashi{\textbf{שבעה.}על שם הברית:}%
% ---------- פרק 26 | פסוק 34 ----------
 \textbf{לד} וַיְהִ֤י עֵשָׂו֙ בֶּן־אַרְבָּעִ֣ים שָׁנָ֔ה וַיִּקַּ֤ח אִשָּׁה֙ אֶת־יְהוּדִ֔ית בַּת־בְּאֵרִ֖י הַֽחִתִּ֑י וְאֶת־בָּ֣שְׂמַ֔ת בַּת־אֵילֹ֖ן הַֽחִתִּֽי׃ \Rashi{\textbf{בן ארבעים שנה.}עשו היה נמשל לחזיר, שנאמר יכרסמנה חזיר מיער (תהילים פ'), החזיר הזה כשהוא שוכב, פושט טלפיו, לומר ראו שאני טהור; כך אלו גוזלים וחומסים ומראים עצמם כשרים; כל מ' שנה היה עשו צד נשים מתחת בעליהן ומענה אותם, כשהיה בן מ' אמר אבא בן מ' שנה נשא אשה, אף אני כן:}%
% ---------- פרק 26 | פסוק 35 ----------
 \textbf{לה} וַתִּהְיֶ֖יןָ מֹ֣רַת ר֑וּחַ לְיִצְחָ֖ק וּלְרִבְקָֽה׃ \{ס\} \Rashi{\textbf{מרת רוח.}לשון המראות רוח, כמו ממרים הייתם, כל מעשיהן היו להכעיס ולעצבון: ליצחק ולרבקה. שהיו עובדות ע"ז (בראשית רבה):}%
% ---------- פרק 27 | פסוק 1 ----------
 \textbf{א} וַֽיְהִי֙ כִּֽי־זָקֵ֣ן יִצְחָ֔ק וַתִּכְהֶ֥יןָ עֵינָ֖יו מֵרְאֹ֑ת וַיִּקְרָ֞א אֶת־עֵשָׂ֣ו ׀ בְּנ֣וֹ הַגָּדֹ֗ל וַיֹּ֤אמֶר אֵלָיו֙ בְּנִ֔י וַיֹּ֥אמֶר אֵלָ֖יו הִנֵּֽנִי׃ \Rashi{\textbf{ותכהין.}בעשנן של אלו. ד"א כשנעקד ע"ג המזבח והיה אביו רוצה לשחטו, באותה שעה נפתחו השמים וראו מלאכי השרת והיו בוכים, וירדו דמעותיהם ונפלו על עיניו, לפיכך כהו עיניו. ד"א כדי שיטל יעקב את הברכות:}%
% ---------- פרק 27 | פסוק 2 ----------
 \textbf{ב} וַיֹּ֕אמֶר הִנֵּה־נָ֖א זָקַ֑נְתִּי לֹ֥א יָדַ֖עְתִּי י֥וֹם מוֹתִֽי׃ \Rashi{\textbf{לא ידעתי יום מותי.}אמר רבי יהושע בן קרחה אם מגיע אדם לפרק אבותיו ידאג חמש שנים לפניהן וחמש לאחר כן; ויצחק היה בן קכ"ג, אמר שמא לפרק אמי אני מגיע והיא מתה בת קכ"ז והריני בן ה' שנים סמוך לפרקה; לפיכך לא ידעתי יום מותי, שמא לפרק אמי, שמא לפרק אבא:}%
% ---------- פרק 27 | פסוק 3 ----------
 \textbf{ג} וְעַתָּה֙ שָׂא־נָ֣א כֵלֶ֔יךָ תֶּלְיְךָ֖ וְקַשְׁתֶּ֑ךָ וְצֵא֙ הַשָּׂדֶ֔ה וְצ֥וּדָה לִּ֖י (צידה) [צָֽיִד]׃ \Rashi{\textbf{שא נא.}לשון השחזה, כאותה ששנינו אין משחיזין את הסכין, אבל משיאה על גבי חברתה (ביצה כ"ח), חדד סכינך ושחט יפה, שלא תאכילני נבלה (בראשית רבה): תליך. חרבך שדרך לתלותה: וצודה לי. מן ההפקר ולא מן הגזל:}%
% ---------- פרק 27 | פסוק 4 ----------
 \textbf{ד} וַעֲשֵׂה־לִ֨י מַטְעַמִּ֜ים כַּאֲשֶׁ֥ר אָהַ֛בְתִּי וְהָבִ֥יאָה לִּ֖י וְאֹכֵ֑לָה בַּעֲב֛וּר תְּבָרֶכְךָ֥ נַפְשִׁ֖י בְּטֶ֥רֶם אָמֽוּת׃ %
% ---------- פרק 27 | פסוק 5 ----------
 \textbf{ה} וְרִבְקָ֣ה שֹׁמַ֔עַת בְּדַבֵּ֣ר יִצְחָ֔ק אֶל־עֵשָׂ֖ו בְּנ֑וֹ וַיֵּ֤לֶךְ עֵשָׂו֙ הַשָּׂדֶ֔ה לָצ֥וּד צַ֖יִד לְהָבִֽיא׃ \Rashi{\textbf{לצוד ציד להביא.}מהו להביא? אם לא ימצא ציד יביא מן הגזל:}%
% ---------- פרק 27 | פסוק 6 ----------
 \textbf{ו} וְרִבְקָה֙ אָֽמְרָ֔ה אֶל־יַעֲקֹ֥ב בְּנָ֖הּ לֵאמֹ֑ר הִנֵּ֤ה שָׁמַ֙עְתִּי֙ אֶת־אָבִ֔יךָ מְדַבֵּ֛ר אֶל־עֵשָׂ֥ו אָחִ֖יךָ לֵאמֹֽר׃ %
% ---------- פרק 27 | פסוק 7 ----------
 \textbf{ז} הָבִ֨יאָה לִּ֥י צַ֛יִד וַעֲשֵׂה־לִ֥י מַטְעַמִּ֖ים וְאֹכֵ֑לָה וַאֲבָרֶכְכָ֛ה לִפְנֵ֥י יְהֹוָ֖ה לִפְנֵ֥י מוֹתִֽי׃ \Rashi{\textbf{לפני ה'.}ברשותו, שיסכים על ידי:}%
% ---------- פרק 27 | פסוק 8 ----------
 \textbf{ח} וְעַתָּ֥ה בְנִ֖י שְׁמַ֣ע בְּקֹלִ֑י לַאֲשֶׁ֥ר אֲנִ֖י מְצַוָּ֥ה אֹתָֽךְ׃ %
% ---------- פרק 27 | פסוק 9 ----------
 \textbf{ט} לֶךְ־נָא֙ אֶל־הַצֹּ֔אן וְקַֽח־לִ֣י מִשָּׁ֗ם שְׁנֵ֛י גְּדָיֵ֥י עִזִּ֖ים טֹבִ֑ים וְאֶֽעֱשֶׂ֨ה אֹתָ֧ם מַטְעַמִּ֛ים לְאָבִ֖יךָ כַּאֲשֶׁ֥ר אָהֵֽב׃ \Rashi{\textbf{וקח לי.}משלי הם ואינם גזל, שכך כתב לה יצחק בכתבתה לטל שני גדיי עזים בכל יום (בראשית רבה): שני גדיי עזים. וכי שני גדיי עזים היה מאכל של יצחק? אלא האחד הקריב לפסחו והאחד עשה מטעמים; בפרקי רבי אליעזר: כאשר אהב. כי טעם הגדי כטעם הצבי:}%
% ---------- פרק 27 | פסוק 10 ----------
 \textbf{י} וְהֵבֵאתָ֥ לְאָבִ֖יךָ וְאָכָ֑ל בַּעֲבֻ֛ר אֲשֶׁ֥ר יְבָרֶכְךָ֖ לִפְנֵ֥י מוֹתֽוֹ׃ %
% ---------- פרק 27 | פסוק 11 ----------
 \textbf{יא} וַיֹּ֣אמֶר יַעֲקֹ֔ב אֶל־רִבְקָ֖ה אִמּ֑וֹ הֵ֣ן עֵשָׂ֤ו אָחִי֙ אִ֣ישׁ שָׂעִ֔ר וְאָנֹכִ֖י אִ֥ישׁ חָלָֽק׃ \Rashi{\textbf{איש שעיר.}בעל שער:}%
% ---------- פרק 27 | פסוק 12 ----------
 \textbf{יב} אוּלַ֤י יְמֻשֵּׁ֙נִי֙ אָבִ֔י וְהָיִ֥יתִי בְעֵינָ֖יו כִּמְתַעְתֵּ֑עַ וְהֵבֵאתִ֥י עָלַ֛י קְלָלָ֖ה וְלֹ֥א בְרָכָֽה׃ \Rashi{\textbf{ימשני.}כמו ממשש בצהרים:}%
% ---------- פרק 27 | פסוק 13 ----------
 \textbf{יג} וַתֹּ֤אמֶר לוֹ֙ אִמּ֔וֹ עָלַ֥י קִלְלָתְךָ֖ בְּנִ֑י אַ֛ךְ שְׁמַ֥ע בְּקֹלִ֖י וְלֵ֥ךְ קַֽח־לִֽי׃ %
% ---------- פרק 27 | פסוק 14 ----------
 \textbf{יד} וַיֵּ֙לֶךְ֙ וַיִּקַּ֔ח וַיָּבֵ֖א לְאִמּ֑וֹ וַתַּ֤עַשׂ אִמּוֹ֙ מַטְעַמִּ֔ים כַּאֲשֶׁ֖ר אָהֵ֥ב אָבִֽיו׃ %
% ---------- פרק 27 | פסוק 15 ----------
 \textbf{טו} וַתִּקַּ֣ח רִ֠בְקָ֠ה אֶת־בִּגְדֵ֨י עֵשָׂ֜ו בְּנָ֤הּ הַגָּדֹל֙ הַחֲמֻדֹ֔ת אֲשֶׁ֥ר אִתָּ֖הּ בַּבָּ֑יִת וַתַּלְבֵּ֥שׁ אֶֽת־יַעֲקֹ֖ב בְּנָ֥הּ הַקָּטָֽן׃ \Rashi{\textbf{החמדת.}הנקיות, כתרגומו דכיתא, ד"א שחמד אותן מן נמרוד: אשר אתה בבית. והלא כמה נשים היו לו, והוא מפקיד אצל אמו? אלא שהיה בקי במעשיהן וחושדן:}%
% ---------- פרק 27 | פסוק 16 ----------
 \textbf{טז} וְאֵ֗ת עֹרֹת֙ גְּדָיֵ֣י הָֽעִזִּ֔ים הִלְבִּ֖ישָׁה עַל־יָדָ֑יו וְעַ֖ל חֶלְקַ֥ת צַוָּארָֽיו׃ %
% ---------- פרק 27 | פסוק 17 ----------
 \textbf{יז} וַתִּתֵּ֧ן אֶת־הַמַּטְעַמִּ֛ים וְאֶת־הַלֶּ֖חֶם אֲשֶׁ֣ר עָשָׂ֑תָה בְּיַ֖ד יַעֲקֹ֥ב בְּנָֽהּ׃ %
% ---------- פרק 27 | פסוק 18 ----------
 \textbf{יח} וַיָּבֹ֥א אֶל־אָבִ֖יו וַיֹּ֣אמֶר אָבִ֑י וַיֹּ֣אמֶר הִנֶּ֔נִּי מִ֥י אַתָּ֖ה בְּנִֽי׃ %
% ---------- פרק 27 | פסוק 19 ----------
 \textbf{יט} וַיֹּ֨אמֶר יַעֲקֹ֜ב אֶל־אָבִ֗יו אָנֹכִי֙ עֵשָׂ֣ו בְּכֹרֶ֔ךָ עָשִׂ֕יתִי כַּאֲשֶׁ֥ר דִּבַּ֖רְתָּ אֵלָ֑י קֽוּם־נָ֣א שְׁבָ֗ה וְאׇכְלָה֙ מִצֵּידִ֔י בַּעֲב֖וּר תְּבָרְכַ֥נִּי נַפְשֶֽׁךָ׃ \Rashi{\textbf{אנכי עשו בכורך.}אנכי המביא לך ועשו הוא בכורך: עשיתי. כמה דברים כאשר דברת אלי: שבה. לשון מסב על השלחן, לכך מתרגם אסתחר:}%
% ---------- פרק 27 | פסוק 20 ----------
 \textbf{כ} וַיֹּ֤אמֶר יִצְחָק֙ אֶל־בְּנ֔וֹ מַה־זֶּ֛ה מִהַ֥רְתָּ לִמְצֹ֖א בְּנִ֑י וַיֹּ֕אמֶר כִּ֥י הִקְרָ֛ה יְהֹוָ֥ה אֱלֹהֶ֖יךָ לְפָנָֽי׃ %
% ---------- פרק 27 | פסוק 21 ----------
 \textbf{כא} וַיֹּ֤אמֶר יִצְחָק֙ אֶֽל־יַעֲקֹ֔ב גְּשָׁה־נָּ֥א וַאֲמֻֽשְׁךָ֖ בְּנִ֑י הַֽאַתָּ֥ה זֶ֛ה בְּנִ֥י עֵשָׂ֖ו אִם־לֹֽא׃ \Rashi{\textbf{גשה נא ואמשך.}אמר יצחק בלבו אין דרך עשו להיות שם שמים שגור בפיו, וזה אמר כי הקרה ה' אלהיך:}%
% ---------- פרק 27 | פסוק 22 ----------
 \textbf{כב} וַיִּגַּ֧שׁ יַעֲקֹ֛ב אֶל־יִצְחָ֥ק אָבִ֖יו וַיְמֻשֵּׁ֑הוּ וַיֹּ֗אמֶר הַקֹּל֙ ק֣וֹל יַעֲקֹ֔ב וְהַיָּדַ֖יִם יְדֵ֥י עֵשָֽׂו׃ \Rashi{\textbf{קול יעקב.}שמדבר בלשון תחנונים קום נא, אבל עשו קנטוריא דבר יקם אבי:}%
% ---------- פרק 27 | פסוק 23 ----------
 \textbf{כג} וְלֹ֣א הִכִּיר֔וֹ כִּֽי־הָי֣וּ יָדָ֗יו כִּידֵ֛י עֵשָׂ֥ו אָחִ֖יו שְׂעִרֹ֑ת וַֽיְבָרְכֵֽהוּ׃ %
% ---------- פרק 27 | פסוק 24 ----------
 \textbf{כד} וַיֹּ֕אמֶר אַתָּ֥ה זֶ֖ה בְּנִ֣י עֵשָׂ֑ו וַיֹּ֖אמֶר אָֽנִי׃ \Rashi{\textbf{ויאמר אני.}לא אמר אני עשו אלא אני:}%
% ---------- פרק 27 | פסוק 25 ----------
 \textbf{כה} וַיֹּ֗אמֶר הַגִּ֤שָׁה לִּי֙ וְאֹֽכְלָה֙ מִצֵּ֣יד בְּנִ֔י לְמַ֥עַן תְּבָֽרֶכְךָ֖ נַפְשִׁ֑י וַיַּגֶּשׁ־לוֹ֙ וַיֹּאכַ֔ל וַיָּ֧בֵא ל֦וֹ יַ֖יִן וַיֵּֽשְׁתְּ׃ %
% ---------- פרק 27 | פסוק 26 ----------
 \textbf{כו} וַיֹּ֥אמֶר אֵלָ֖יו יִצְחָ֣ק אָבִ֑יו גְּשָׁה־נָּ֥א וּשְׁקָה־לִּ֖י בְּנִֽי׃ %
% ---------- פרק 27 | פסוק 27 ----------
 \textbf{כז} וַיִּגַּשׁ֙ וַיִּשַּׁק־ל֔וֹ וַיָּ֛רַח אֶת־רֵ֥יחַ בְּגָדָ֖יו וַֽיְבָרְכֵ֑הוּ וַיֹּ֗אמֶר רְאֵה֙ רֵ֣יחַ בְּנִ֔י כְּרֵ֣יחַ שָׂדֶ֔ה אֲשֶׁ֥ר בֵּרְכ֖וֹ יְהֹוָֽה׃ \Rashi{\textbf{וירח וגו'.}והלא אין ריח רע יותר משטף העזים? אלא מלמד שנכנסה עמו ריח גן עדן; כריח שדה אשר ברכו ה'. שנתן בו ריח טוב וזהו שדה תפוחים, כך דרשו רז"ל:}%
% ---------- פרק 27 | פסוק 28 ----------
 \textbf{כח} וְיִֽתֶּן־לְךָ֙ הָאֱלֹהִ֔ים מִטַּל֙ הַשָּׁמַ֔יִם וּמִשְׁמַנֵּ֖י הָאָ֑רֶץ וְרֹ֥ב דָּגָ֖ן וְתִירֹֽשׁ׃ \Rashi{\textbf{ויתן לך.}יתן ויחזר ויתן (בראשית רבה). ולפי פשוטו מוסב לענין הראשון: ראה ריח בני שנתן לו הקב"ה כריח שדה וגו' ועוד יתן לך מטל השמים וגו': מטל השמים. כמשמעו; ומדרשי אגדה יש להרבה פנים: האלהים. מהו האלהים? בדין; אם ראוי לך יתן לך ואם לאו לא יתן לך, אבל לעשו אמר משמני הארץ יהיה מושבך, בין צדיק בין רשע יתן לך; וממנו למד שלמה, כשעשה הבית סדר תפלתו. ישראל שהוא בעל אמונה ומצדיק עליו את הדין לא יקרא עליך תגר, לפיכך ונתת לאיש בכל דרכיו אשר תדע את לבבו (מלכים א, ח'), אבל נכרי מחסר אמנה, לפיכך אמר אתה תשמע השמים, ועשית ככל אשר יקרא אליך הנכרי (שם), בין ראוי בין שאינו ראוי תן לו, כדי שלא יקרא עליך תגר:}%
% ---------- פרק 27 | פסוק 29 ----------
 \textbf{כט} יַֽעַבְד֣וּךָ עַמִּ֗ים (וישתחו) [וְיִֽשְׁתַּחֲו֤וּ] לְךָ֙ לְאֻמִּ֔ים הֱוֵ֤ה גְבִיר֙ לְאַחֶ֔יךָ וְיִשְׁתַּחֲו֥וּ לְךָ֖ בְּנֵ֣י אִמֶּ֑ךָ אֹרְרֶ֣יךָ אָר֔וּר וּֽמְבָרְכֶ֖יךָ בָּרֽוּךְ׃ \Rashi{\textbf{בני אמך.}ויעקב אמר ליהודה בני אביך, לפי שהיו לו בנים מכמה אמהות, וכאן, שלא נשא אלא אשה אחת, אומר בני אמך (בראשית רבה): ארריך ארור ומברכיך ברוך. ובבלעם הוא אומר מברכיך ברוך וארריך ארור? הצדיקים תחלתם יסורים וסופן שלוה ואורריהם ומצעריהם קודמים למברכיהם, לפיכך יצחק הקדים קללת אוררים לברכת מברכים; הרשעים תחלתן שלוה וסופן יסורין, לפיכך בלעם הקדים ברכה לקללה (בראשית רבה):}%
% ---------- פרק 27 | פסוק 30 ----------
 \textbf{ל} וַיְהִ֗י כַּאֲשֶׁ֨ר כִּלָּ֣ה יִצְחָק֮ לְבָרֵ֣ךְ אֶֽת־יַעֲקֹב֒ וַיְהִ֗י אַ֣ךְ יָצֹ֤א יָצָא֙ יַעֲקֹ֔ב מֵאֵ֥ת פְּנֵ֖י יִצְחָ֣ק אָבִ֑יו וְעֵשָׂ֣ו אָחִ֔יו בָּ֖א מִצֵּידֽוֹ׃ \Rashi{\textbf{יצא יצא.}זה יוצא וזה בא:}%
% ---------- פרק 27 | פסוק 31 ----------
 \textbf{לא} וַיַּ֤עַשׂ גַּם־הוּא֙ מַטְעַמִּ֔ים וַיָּבֵ֖א לְאָבִ֑יו וַיֹּ֣אמֶר לְאָבִ֗יו יָקֻ֤ם אָבִי֙ וְיֹאכַל֙ מִצֵּ֣יד בְּנ֔וֹ בַּעֲבֻ֖ר תְּבָרְכַ֥נִּי נַפְשֶֽׁךָ׃ %
% ---------- פרק 27 | פסוק 32 ----------
 \textbf{לב} וַיֹּ֥אמֶר ל֛וֹ יִצְחָ֥ק אָבִ֖יו מִי־אָ֑תָּה וַיֹּ֕אמֶר אֲנִ֛י בִּנְךָ֥ בְכֹֽרְךָ֖ עֵשָֽׂו׃ %
% ---------- פרק 27 | פסוק 33 ----------
 \textbf{לג} וַיֶּחֱרַ֨ד יִצְחָ֣ק חֲרָדָה֮ גְּדֹלָ֣ה עַד־מְאֹד֒ וַיֹּ֡אמֶר מִֽי־אֵפ֡וֹא ה֣וּא הַצָּֽד־צַ֩יִד֩ וַיָּ֨בֵא לִ֜י וָאֹכַ֥ל מִכֹּ֛ל בְּטֶ֥רֶם תָּב֖וֹא וָאֲבָרְכֵ֑הוּ גַּם־בָּר֖וּךְ יִהְיֶֽה׃ \Rashi{\textbf{ויחרד.}כתרגומו ותוה, לשון תמיה. ומדרשו ראה גיהנם פתוחה מתחתיו: מי אפוא. לשון לעצמו משמש עם כמה דברים; איפה – איה פה, מי הוא ואיפה הוא הצד ציד? ואכל מכל. מכל טעמים שבקשתי לטעם, טעמתי בו (בראשית רבה): גם ברוך יהיה. שלא תאמר אלולי שרמה יעקב לאביו לא נטל את הברכות, לכך הסכים וברכו מדעתו (בראשית רבה):}%
% ---------- פרק 27 | פסוק 34 ----------
 \textbf{לד} כִּשְׁמֹ֤עַ עֵשָׂו֙ אֶת־דִּבְרֵ֣י אָבִ֔יו וַיִּצְעַ֣ק צְעָקָ֔ה גְּדֹלָ֥ה וּמָרָ֖ה עַד־מְאֹ֑ד וַיֹּ֣אמֶר לְאָבִ֔יו בָּרְכֵ֥נִי גַם־אָ֖נִי אָבִֽי׃ %
% ---------- פרק 27 | פסוק 35 ----------
 \textbf{לה} וַיֹּ֕אמֶר בָּ֥א אָחִ֖יךָ בְּמִרְמָ֑ה וַיִּקַּ֖ח בִּרְכָתֶֽךָ׃ \Rashi{\textbf{במרמה.}בחכמה:}%
% ---------- פרק 27 | פסוק 36 ----------
 \textbf{לו} וַיֹּ֡אמֶר הֲכִי֩ קָרָ֨א שְׁמ֜וֹ יַעֲקֹ֗ב וַֽיַּעְקְבֵ֙נִי֙ זֶ֣ה פַעֲמַ֔יִם אֶת־בְּכֹרָתִ֣י לָקָ֔ח וְהִנֵּ֥ה עַתָּ֖ה לָקַ֣ח בִּרְכָתִ֑י וַיֹּאמַ֕ר הֲלֹא־אָצַ֥לְתָּ לִּ֖י בְּרָכָֽה׃ \Rashi{\textbf{הכי קרא שמו.}לשון תמה הוא, כמו הכי אחי אתה, שמא לכך נקרא שמו יעקב על שם סופו שהוא עתיד לעקבני? תנחומא. למה חרד יצחק? אמר שמא עון יש בי שברכתי קטן לפני גדול ושניתי סדר היחס, התחיל עשו מצעק ויעקבני זה פעמים, אמר לו אביו מה עשה לך? אמר לו את בכרתי לקח, אמר בכך הייתי מצר וחרד שמא עברתי על שורת הדין, עכשו לבכור ברכתי, גם ברוך יהיה: ויעקבני. כתרגומו וכמני – ארבני; וארב – וכמן. ויש מתרגמין וחכמני, נתחכם לי: אצלת. לשון הפרשה כמו ויאצל:}%
% ---------- פרק 27 | פסוק 37 ----------
 \textbf{לז} וַיַּ֨עַן יִצְחָ֜ק וַיֹּ֣אמֶר לְעֵשָׂ֗ו הֵ֣ן גְּבִ֞יר שַׂמְתִּ֥יו לָךְ֙ וְאֶת־כׇּל־אֶחָ֗יו נָתַ֤תִּי לוֹ֙ לַעֲבָדִ֔ים וְדָגָ֥ן וְתִירֹ֖שׁ סְמַכְתִּ֑יו וּלְכָ֣ה אֵפ֔וֹא מָ֥ה אֶֽעֱשֶׂ֖ה בְּנִֽי׃ \Rashi{\textbf{הן גביר.}ברכה זו שביעית היא, והוא עושה אותה ראשונה? אלא אמר לו מה תועלת לך בברכה? אם תקנה נכסים, שלו הם, שהרי גביר שמתיו לך, ומה שקנה עבד קנה רבו: ולכה אפא מה אעשה. איה איפא אבקש מה לעשות לך:}%
% ---------- פרק 27 | פסוק 38 ----------
 \textbf{לח} וַיֹּ֨אמֶר עֵשָׂ֜ו אֶל־אָבִ֗יו הַֽבְרָכָ֨ה אַחַ֤ת הִֽוא־לְךָ֙ אָבִ֔י בָּרְכֵ֥נִי גַם־אָ֖נִי אָבִ֑י וַיִּשָּׂ֥א עֵשָׂ֛ו קֹל֖וֹ וַיֵּֽבְךְּ׃ \Rashi{\textbf{הברכה אחת.}ה"א זו משמשת לשון תמיה, כמו הבמחנים? השמנה היא? הכמות נבל?}%
% ---------- פרק 27 | פסוק 39 ----------
 \textbf{לט} וַיַּ֛עַן יִצְחָ֥ק אָבִ֖יו וַיֹּ֣אמֶר אֵלָ֑יו הִנֵּ֞ה מִשְׁמַנֵּ֤י הָאָ֙רֶץ֙ יִהְיֶ֣ה מֽוֹשָׁבֶ֔ךָ וּמִטַּ֥ל הַשָּׁמַ֖יִם מֵעָֽל׃ \Rashi{\textbf{משמני ארץ וגו'.}זו איטליאה של יון:}%
% ---------- פרק 27 | פסוק 40 ----------
 \textbf{מ} וְעַל־חַרְבְּךָ֣ תִֽחְיֶ֔ה וְאֶת־אָחִ֖יךָ תַּעֲבֹ֑ד וְהָיָה֙ כַּאֲשֶׁ֣ר תָּרִ֔יד וּפָרַקְתָּ֥ עֻלּ֖וֹ מֵעַ֥ל צַוָּארֶֽךָ׃ \Rashi{\textbf{ועל חרבך.}כמו בחרבך, יש על במקום ב', כמו עמדתם על חרבכם – בחרבכם; על צבאתם – בצבאתם: תריד. לשון צער כמו אריד בשיחי, כלומר, כשיעברו ישראל על התורה ויהיה לך פתחון פה להצטער על הברכות שנטל, ופרקת עלו וגומר:}%
% ---------- פרק 27 | פסוק 41 ----------
 \textbf{מא} וַיִּשְׂטֹ֤ם עֵשָׂו֙ אֶֽת־יַעֲקֹ֔ב עַ֨ל־הַבְּרָכָ֔ה אֲשֶׁ֥ר בֵּרְכ֖וֹ אָבִ֑יו וַיֹּ֨אמֶר עֵשָׂ֜ו בְּלִבּ֗וֹ יִקְרְבוּ֙ יְמֵי֙ אֵ֣בֶל אָבִ֔י וְאַֽהַרְגָ֖ה אֶת־יַעֲקֹ֥ב אָחִֽי׃ \Rashi{\textbf{יקרבו ימי אבל אבי.}כמשמעו, שלא אצער את אבא, ומדרשי אגדה לכמה פנים יש:}%
% ---------- פרק 27 | פסוק 42 ----------
 \textbf{מב} וַיֻּגַּ֣ד לְרִבְקָ֔ה אֶת־דִּבְרֵ֥י עֵשָׂ֖ו בְּנָ֣הּ הַגָּדֹ֑ל וַתִּשְׁלַ֞ח וַתִּקְרָ֤א לְיַעֲקֹב֙ בְּנָ֣הּ הַקָּטָ֔ן וַתֹּ֣אמֶר אֵלָ֔יו הִנֵּה֙ עֵשָׂ֣ו אָחִ֔יךָ מִתְנַחֵ֥ם לְךָ֖ לְהׇרְגֶֽךָ׃ \Rashi{\textbf{ויגד לרבקה.}ברוח הקדש הגד לה מה שעשו מהרהר בלבו: מתנחם לך. נחם על האחוה לחשב מחשבה אחרת להתנכר לך ולהרגך. ומדרש אגדה כבר אתה מת בעיניו ושתה עליך כוס של תנחומים. ולפי פשוטו לשון תנחומים, מתנחם הוא על הברכות בהריגתך:}%
% ---------- פרק 27 | פסוק 43 ----------
 \textbf{מג} וְעַתָּ֥ה בְנִ֖י שְׁמַ֣ע בְּקֹלִ֑י וְק֧וּם בְּרַח־לְךָ֛ אֶל־לָבָ֥ן אָחִ֖י חָרָֽנָה׃ %
% ---------- פרק 27 | פסוק 44 ----------
 \textbf{מד} וְיָשַׁבְתָּ֥ עִמּ֖וֹ יָמִ֣ים אֲחָדִ֑ים עַ֥ד אֲשֶׁר־תָּשׁ֖וּב חֲמַ֥ת אָחִֽיךָ׃ \Rashi{\textbf{אחדים.}מעטים:}%
% ---------- פרק 27 | פסוק 45 ----------
 \textbf{מה} עַד־שׁ֨וּב אַף־אָחִ֜יךָ מִמְּךָ֗ וְשָׁכַח֙ אֵ֣ת אֲשֶׁר־עָשִׂ֣יתָ לּ֔וֹ וְשָׁלַחְתִּ֖י וּלְקַחְתִּ֣יךָ מִשָּׁ֑ם לָמָ֥ה אֶשְׁכַּ֛ל גַּם־שְׁנֵיכֶ֖ם י֥וֹם אֶחָֽד׃ \Rashi{\textbf{למה אשכל.}אהיה שכולה משניכם. הקובר את בניו קרוי שכול, וכן ביעקב אמר כאשר שכלתי שכלתי: גם שניכם. אם יקום עליך ואתה תהרגנו יעמדו בניו ויהרגוך; ורוח הקדש נזרקה בה ונתנבאה שביום א' ימותו, כמו שמפרש בפרק המקנא לאשתו:}%
% ---------- פרק 27 | פסוק 46 ----------
 \textbf{מו} וַתֹּ֤אמֶר רִבְקָה֙ אֶל־יִצְחָ֔ק קַ֣צְתִּי בְחַיַּ֔י מִפְּנֵ֖י בְּנ֣וֹת חֵ֑ת אִם־לֹקֵ֣חַ יַ֠עֲקֹ֠ב אִשָּׁ֨ה מִבְּנֽוֹת־חֵ֤ת כָּאֵ֙לֶּה֙ מִבְּנ֣וֹת הָאָ֔רֶץ לָ֥מָּה לִּ֖י חַיִּֽים׃ \Rashi{\textbf{קצתי בחיי.}מאסתי בחיי:}%
% ---------- פרק 28 | פסוק 1 ----------
 \textbf{א} וַיִּקְרָ֥א יִצְחָ֛ק אֶֽל־יַעֲקֹ֖ב וַיְבָ֣רֶךְ אֹת֑וֹ וַיְצַוֵּ֙הוּ֙ וַיֹּ֣אמֶר ל֔וֹ לֹֽא־תִקַּ֥ח אִשָּׁ֖ה מִבְּנ֥וֹת כְּנָֽעַן׃ %
% ---------- פרק 28 | פסוק 2 ----------
 \textbf{ב} ק֥וּם לֵךְ֙ פַּדֶּ֣נָֽה אֲרָ֔ם בֵּ֥יתָה בְתוּאֵ֖ל אֲבִ֣י אִמֶּ֑ךָ וְקַח־לְךָ֤ מִשָּׁם֙ אִשָּׁ֔ה מִבְּנ֥וֹת לָבָ֖ן אֲחִ֥י אִמֶּֽךָ׃ \Rashi{\textbf{פדנה.}כמו לפדן: ביתה בתואל. לבית בתואל, כל תבה שצריכה למ"ד בתחלתה הטיל לה ה"א בסופה (יבמות י"ג):}%
% ---------- פרק 28 | פסוק 3 ----------
 \textbf{ג} וְאֵ֤ל שַׁדַּי֙ יְבָרֵ֣ךְ אֹֽתְךָ֔ וְיַפְרְךָ֖ וְיַרְבֶּ֑ךָ וְהָיִ֖יתָ לִקְהַ֥ל עַמִּֽים׃ \Rashi{\textbf{ואל שדי.}מי שדי בברכותיו למתברכין מפיו, יברך אותך:}%
% ---------- פרק 28 | פסוק 4 ----------
 \textbf{ד} וְיִֽתֶּן־לְךָ֙ אֶת־בִּרְכַּ֣ת אַבְרָהָ֔ם לְךָ֖ וּלְזַרְעֲךָ֣ אִתָּ֑ךְ לְרִשְׁתְּךָ֙ אֶת־אֶ֣רֶץ מְגֻרֶ֔יךָ אֲשֶׁר־נָתַ֥ן אֱלֹהִ֖ים לְאַבְרָהָֽם׃ \Rashi{\textbf{את ברכת אברהם.}שאמר לו ואעשך לגוי גדול והתברכו בזרעך יהיו אותן ברכות אמורות בשבילך – ממך יצא אותו הגוי ואותו הזרע המברך:}%
% ---------- פרק 28 | פסוק 5 ----------
 \textbf{ה} וַיִּשְׁלַ֤ח יִצְחָק֙ אֶֽת־יַעֲקֹ֔ב וַיֵּ֖לֶךְ פַּדֶּ֣נָֽה אֲרָ֑ם אֶל־לָבָ֤ן בֶּן־בְּתוּאֵל֙ הָֽאֲרַמִּ֔י אֲחִ֣י רִבְקָ֔ה אֵ֥ם יַעֲקֹ֖ב וְעֵשָֽׂו׃ \Rashi{\textbf{אם יעקב ועשו.}איני יודע מה מלמדנו:}%
% ---------- פרק 28 | פסוק 6 ----------
 \textbf{ו} וַיַּ֣רְא עֵשָׂ֗ו כִּֽי־בֵרַ֣ךְ יִצְחָק֮ אֶֽת־יַעֲקֹב֒ וְשִׁלַּ֤ח אֹתוֹ֙ פַּדֶּ֣נָֽה אֲרָ֔ם לָקַֽחַת־ל֥וֹ מִשָּׁ֖ם אִשָּׁ֑ה בְּבָרְכ֣וֹ אֹת֔וֹ וַיְצַ֤ו עָלָיו֙ לֵאמֹ֔ר לֹֽא־תִקַּ֥ח אִשָּׁ֖ה מִבְּנ֥וֹת כְּנָֽעַן׃ %
% ---------- פרק 28 | פסוק 7 ----------
 \textbf{ז} וַיִּשְׁמַ֣ע יַעֲקֹ֔ב אֶל־אָבִ֖יו וְאֶל־אִמּ֑וֹ וַיֵּ֖לֶךְ פַּדֶּ֥נָֽה אֲרָֽם׃ \Rashi{\textbf{וישמע יעקב.}מחבר לענין של מעלה, וירא עשו כי ברך יצחק וגו' וכי שלח אותו פדנה ארם, וכי שמע יעקב אל אביו והלך פדנה ארם, וכי רעות בנות כנען, והלך גם הוא אל ישמעאל:}%
% ---------- פרק 28 | פסוק 8 ----------
 \textbf{ח} וַיַּ֣רְא עֵשָׂ֔ו כִּ֥י רָע֖וֹת בְּנ֣וֹת כְּנָ֑עַן בְּעֵינֵ֖י יִצְחָ֥ק אָבִֽיו׃ %
% ---------- פרק 28 | פסוק 9 ----------
 \textbf{ט} וַיֵּ֥לֶךְ עֵשָׂ֖ו אֶל־יִשְׁמָעֵ֑אל וַיִּקַּ֡ח אֶֽת־מָחֲלַ֣ת ׀ בַּת־יִשְׁמָעֵ֨אל בֶּן־אַבְרָהָ֜ם אֲח֧וֹת נְבָי֛וֹת עַל־נָשָׁ֖יו ל֥וֹ לְאִשָּֽׁה׃ \{ס\} \Rashi{\textbf{אחות נביות.}ממשמע שנאמר בת ישמעאל איני יודע שהיא אחות נביות? אלא למדנו שמת ישמעאל משיעדה לעשו קדם נשואיה והשיאה נביות אחיה, ולמדנו שהיה יעקב באותו הפרק בן ס"ג שנים, שהרי ישמעאל בן ע"ד שנים היה כשנולד יעקב, י"ד שנה היה גדול ישמעאל מיצחק, ויצחק בן ס' שנה בלדת אותם הרי ע"ד, ושנותיו היו קל"ז, שנאמר ואלה שני חיי ישמעאל וגו', נמצא יעקב כשמת ישמעאל בן ס"ג שנים היה, ולמדנו מכאן שנטמן בבית עבר י"ד שנה ואחר כך הלך לחרן, שהרי לא שהה בבית לבן לפני לדתו של יוסף אלא י"ד שנה, שנאמר עבדתיך ארבע עשרה שנה בשתי בנתיך ושש שנים בצאנך, ושכר הצאן משנולד יוסף היה, שנאמר ויהי כאשר ילדה רחל את יוסף וגו'; ויוסף בן ל' שנה היה כשמלך, ומשם עד שירד יעקב למצרים ט' שנים, ז' של שבע וב' של רעב, ויעקב אמר לפרעה ימי שני מגורי שלשים ומאת שנה – צא וחשב י"ד שלפני לדת יוסף ושלשים של יוסף ותשע משמלך עד שבא יעקב, הרי נ"ג, וכשפרש מאביו היה בן ס"ג, הרי קי"ו, והוא אומר שלשים ומאת שנה, הרי חסרים י"ד שנים, הא למדת שאחר שקבל הברכות נטמן בבית עבר י"ד שנים, אבל לא נענש בזכות התורה; שהרי לא פרש יוסף מאביו אלא כ"ב שנה, דהינו מי"ז עד ל"ט, כנגד כ"ב שפרש יעקב מאביו ולא כבדו, והם כ' שנים בבית לבן ושתי שנים ששהה בדרך, כדכתיב ויבן לו בית ולמקנהו עשה סכת, ופר' רז"ל מזה הפסוק ששהה י"ח חדשים בדרך – דבית הוה בימות הגשמים וסכת הוו בימות החמה – ולחשבון הפסוקים שחשבנו לעיל משפרש מאביו עד שירד למצרים שהיה בן ק"ל שנים, ששם אנו מוצאים עוד י"ד שנים – אלא ודאי נטמן בבית עבר בהליכתו לבית לבן ללמד תורה ממנו, ובשביל זכות התורה לא נענש עליהם ולא פרש יוסף ממנו אלא כ"ב שנה – מדה כנגד מדה – ע"כ ברש"י ישן: על נשיו. הוסיף רשעה על רשעתו, שלא גרש את הראשונות:}%
\addparasha{חומש בראשית | פרשת ויצא}
% ---------- פרק 28 | פסוק 10 ----------
 \textbf{י} וַיֵּצֵ֥א יַעֲקֹ֖ב מִבְּאֵ֣ר שָׁ֑בַע וַיֵּ֖לֶךְ חָרָֽנָה׃ \Rashi{\textbf{ויצא יעקב.}על ידי שבשביל שרעות בנות כנען בעיני יצחק אביו הלך עשו אל ישמעאל, הפסיק הענין בפרשתו של יעקב, וכתיב וירא עשו כי ברך וגו', ומשגמר חזר לענין הראשון: ויצא. לא היה צריך לכתב אלא וילך יעקב חרנה, ולמה הזכיר יציאתו? אלא מגיד שיציאת צדיק מן המקום עושה רשם, שבזמן שהצדיק בעיר, הוא הודה הוא זיוה הוא הדרה; יצא משם, פנה הודה פנה זיוה פנה הדרה. וכן ותצא מן המקום האמור בנעמי ורות (רות א'): וילך חרנה. יצא ללכת לחרן:}%
% ---------- פרק 28 | פסוק 11 ----------
 \textbf{יא} וַיִּפְגַּ֨ע בַּמָּק֜וֹם וַיָּ֤לֶן שָׁם֙ כִּי־בָ֣א הַשֶּׁ֔מֶשׁ וַיִּקַּח֙ מֵאַבְנֵ֣י הַמָּק֔וֹם וַיָּ֖שֶׂם מְרַֽאֲשֹׁתָ֑יו וַיִּשְׁכַּ֖ב בַּמָּק֥וֹם הַהֽוּא׃ \Rashi{\textbf{ויפגע במקום.}לא הזכיר הכתוב באיזה מקום אלא במקום – הנזכר במקום אחר, הוא הר המוריה, שנאמר בו וירא את המקום מרחק: ויפגע. כמו ופגע ביריחו ופגע בדבשת (יהושע ט"ז וי"ט) ורבותינו פרשו לשון תפלה (ברכות כ"ו), כמו ואל תפגע בי (ירמיהו ז'), ולמדנו שתקן תפלת ערבית. ושנה הכתוב ולא כתב ויתפלל, ללמדך שקפצה לו הארץ, כמו שמפרש בפרק גיד הנשה (חולין צ"א): כי בא השמש. היה לו לכתב ויבא השמש וילן שם, כי בא השמש משמע ששקעה לו חמה פתאום, שלא בעונתה, כדי שילין שם: וישם מראשותיו. עשאן כמין מרזב סביב לראשו, שירא מפני חיות רעות; התחילו מריבות זו את זו, זאת אומרת עלי יניח צדיק את ראשו וזאת אומרת עלי יניח; מיד עשאן הקב"ה אבן אחת, וזהו שנאמר ויקח את האבן אשר שם מראשתיו: וישכב במקום ההוא. לשון מעוט; באותו מקום שכב אבל י"ד שנים ששמש בבית עבר לא שכב בלילה, שהיה עוסק בתורה:}%
% ---------- פרק 28 | פסוק 12 ----------
 \textbf{יב} וַֽיַּחֲלֹ֗ם וְהִנֵּ֤ה סֻלָּם֙ מֻצָּ֣ב אַ֔רְצָה וְרֹאשׁ֖וֹ מַגִּ֣יעַ הַשָּׁמָ֑יְמָה וְהִנֵּה֙ מַלְאֲכֵ֣י אֱלֹהִ֔ים עֹלִ֥ים וְיֹרְדִ֖ים בּֽוֹ׃ \Rashi{\textbf{עולים וירדים.}עולים תחלה ואחר כך יורדים? מלאכים שלווהו בארץ אין יוצאים חוצה לארץ, ועלו לרקיע וירדו מלאכי חוצה לארץ ללוותו:}%
% ---------- פרק 28 | פסוק 13 ----------
 \textbf{יג} וְהִנֵּ֨ה יְהֹוָ֜ה נִצָּ֣ב עָלָיו֮ וַיֹּאמַר֒ אֲנִ֣י יְהֹוָ֗ה אֱלֹהֵי֙ אַבְרָהָ֣ם אָבִ֔יךָ וֵאלֹהֵ֖י יִצְחָ֑ק הָאָ֗רֶץ אֲשֶׁ֤ר אַתָּה֙ שֹׁכֵ֣ב עָלֶ֔יהָ לְךָ֥ אֶתְּנֶ֖נָּה וּלְזַרְעֶֽךָ׃ \Rashi{\textbf{נצב עליו.}לשמרו: ואלהי יצחק. אע"פ שלא מצינו במקרא שיחד הקב"ה שמו על הצדיקים בחייהם לכתב אלהי פלוני, משום שנאמר הן בקדשו לא יאמין (איוב ט"ו), כאן יחד שמו על יצחק, לפי שכהו עיניו וכלוא בבית, והרי הוא כמת ויצר הרע פסק ממנו, תנחומא: שכב עליה. קפל הקב"ה כל ארץ ישראל תחתיו, רמז לו שתהא נוחה לכבש לבניו כד' אמות שזה מקומו של אדם (חולין שם):}%
% ---------- פרק 28 | פסוק 14 ----------
 \textbf{יד} וְהָיָ֤ה זַרְעֲךָ֙ כַּעֲפַ֣ר הָאָ֔רֶץ וּפָרַצְתָּ֛ יָ֥מָּה וָקֵ֖דְמָה וְצָפֹ֣נָה וָנֶ֑גְבָּה וְנִבְרְכ֥וּ בְךָ֛ כׇּל־מִשְׁפְּחֹ֥ת הָאֲדָמָ֖ה וּבְזַרְעֶֽךָ׃ \Rashi{\textbf{ופרצת.}וחזקת, כמו וכן יפרץ:}%
% ---------- פרק 28 | פסוק 15 ----------
 \textbf{טו} וְהִנֵּ֨ה אָנֹכִ֜י עִמָּ֗ךְ וּשְׁמַרְתִּ֙יךָ֙ בְּכֹ֣ל אֲשֶׁר־תֵּלֵ֔ךְ וַהֲשִׁ֣בֹתִ֔יךָ אֶל־הָאֲדָמָ֖ה הַזֹּ֑את כִּ֚י לֹ֣א אֶֽעֱזׇבְךָ֔ עַ֚ד אֲשֶׁ֣ר אִם־עָשִׂ֔יתִי אֵ֥ת אֲשֶׁר־דִּבַּ֖רְתִּי לָֽךְ׃ \Rashi{\textbf{אנכי עמך.}לפי שהיה ירא מעשו ומלבן: עד אשר אם עשיתי. אם משמש בלשון כי (גיטין צ'): דברתי לך. לצרכך ועליך; מה שהבטחתי לאברהם על זרעו, לך הבטחתיו ולא לעשו, שלא אמרתי לו כי יצחק יקרא לך זרע אלא כי ביצחק, ולא כל יצחק, וכן כל לי ולך ולו ולהם הסמוכים אצל דבור משמשים לשון על, וזה יוכיח, שהרי עם יעקב לא דבר קדם לכן:}%
% ---------- פרק 28 | פסוק 16 ----------
 \textbf{טז} וַיִּיקַ֣ץ יַעֲקֹב֮ מִשְּׁנָתוֹ֒ וַיֹּ֕אמֶר אָכֵן֙ יֵ֣שׁ יְהֹוָ֔ה בַּמָּק֖וֹם הַזֶּ֑ה וְאָנֹכִ֖י לֹ֥א יָדָֽעְתִּי׃ \Rashi{\textbf{ואנכי לא ידעתי.}שאם ידעתי, לא ישנתי במקום קדוש כזה:}%
% ---------- פרק 28 | פסוק 17 ----------
 \textbf{יז} וַיִּירָא֙ וַיֹּאמַ֔ר מַה־נּוֹרָ֖א הַמָּק֣וֹם הַזֶּ֑ה אֵ֣ין זֶ֗ה כִּ֚י אִם־בֵּ֣ית אֱלֹהִ֔ים וְזֶ֖ה שַׁ֥עַר הַשָּׁמָֽיִם׃ \Rashi{\textbf{כי אם בית אלהים.}א"ר אלעזר בשם רבי יוסי בן זמרא, הסלם הזה עומד בבאר שבע ואמצע שפועו מגיע כנגד בית המקדש, שבאר שבע עומד בדרומה של יהודה, וירושלים בצפונה, בגבול שבין יהודה ובנימין, ובית אל היה בצפון של נחלת בנימין, בגבול שבין בנימין ובין בני יוסף; נמצא סלם שרגליו בבאר שבע וראשו בבית אל מגיע אמצע שפועו נגד ירושלים; וכלפי שאמרו רבותינו שאמר הקב"ה צדיק זה בא לבית מלוני ויפטר בלא לינה, ואמרו יעקב קראה לירושלים בית אל וזו לוז היא ולא ירושלים, ומהיכן למדו לומר כן? אני אומר שנעקר הר המוריה ובא לכאן, וזו היא קפיצת הארץ האמורה בשחיטת חלין, שבא בית המקדש לקראתו עד בית אל, וזהו ויפגע במקום. ואם תאמר כשעבר יעקב על בית המקדש, מדוע לא עכבו שם? איהו לא יהב לביה להתפלל במקום שהתפללו אבותיו, ומן השמים יעכבוהו? איהו עד חרן אזל, כדאמרינן בפרק גיד הנשה; וקרא מוכיח וילך חרנה; כי מטא לחרן אמר, אפשר שעברתי על מקום שהתפללו אבותי ולא התפללתי בו? יהב דעתיה למהדר וחזר עד בית אל וקפצה לו הארץ. האי בית אל לא הסמוך לעי אלא לירושלים, ועל שם יהיה בית אלהים קראו בית אל, והוא הר המוריה שהתפלל בו אברהם, והוא שדה שהתפלל בו יצחק, כדכתיב לשוח בשדה; דהכי אמרינן בפסחים "אל הר ד' ואל בית אלהי יעקב" – מאי שנא יעקב? אלא לא כאברהם שקראו הר, דכתיב בהר ד' יראה, ולא כיצחק שקראו שדה, דכתיב לשוח בשדה, אלא כיעקב שקראו בית (ע"כ פרש"י מדיק): מה נורא. תרגום מה דחילו אתרא הדין, דחילו שם דבר הוא, כמו סוכלתנו וכסו למלבש: וזה שער השמים. מקום תפלה לעלות תפלתם השמימה. ומדרשו, שבית המקדש של מעלה מכון כנגד בית המקדש של מטה:}%
% ---------- פרק 28 | פסוק 18 ----------
 \textbf{יח} וַיַּשְׁכֵּ֨ם יַעֲקֹ֜ב בַּבֹּ֗קֶר וַיִּקַּ֤ח אֶת־הָאֶ֙בֶן֙ אֲשֶׁר־שָׂ֣ם מְרַֽאֲשֹׁתָ֔יו וַיָּ֥שֶׂם אֹתָ֖הּ מַצֵּבָ֑ה וַיִּצֹ֥ק שֶׁ֖מֶן עַל־רֹאשָֽׁהּ׃ %
% ---------- פרק 28 | פסוק 19 ----------
 \textbf{יט} וַיִּקְרָ֛א אֶת־שֵֽׁם־הַמָּק֥וֹם הַה֖וּא בֵּֽית־אֵ֑ל וְאוּלָ֛ם ל֥וּז שֵׁם־הָעִ֖יר לָרִאשֹׁנָֽה׃ %
% ---------- פרק 28 | פסוק 20 ----------
 \textbf{כ} וַיִּדַּ֥ר יַעֲקֹ֖ב נֶ֣דֶר לֵאמֹ֑ר אִם־יִהְיֶ֨ה אֱלֹהִ֜ים עִמָּדִ֗י וּשְׁמָרַ֙נִי֙ בַּדֶּ֤רֶךְ הַזֶּה֙ אֲשֶׁ֣ר אָנֹכִ֣י הוֹלֵ֔ךְ וְנָֽתַן־לִ֥י לֶ֛חֶם לֶאֱכֹ֖ל וּבֶ֥גֶד לִלְבֹּֽשׁ׃ \Rashi{\textbf{אם יהיה אלהים עמדי.}אם ישמר לי הבטחות הללו שהבטיחני להיות עמדי, כמו שאמר לי והנה אנכי עמך: ושמרני. כמו שאמר לי ושמרתיך בכל אשר תלך: ונתן לי לחם לאכל. כמו שאמר כי לא אעזבך, והמבקש לחם הוא קרוי נעזב, שנאמר ולא ראיתי צדיק נעזב וזרעו מבקש לחם (תהילים ל"ז):}%
% ---------- פרק 28 | פסוק 21 ----------
 \textbf{כא} וְשַׁבְתִּ֥י בְשָׁל֖וֹם אֶל־בֵּ֣ית אָבִ֑י וְהָיָ֧ה יְהֹוָ֛ה לִ֖י לֵאלֹהִֽים׃ \Rashi{\textbf{ושבתי.}כמו שאמר לי והשבתיך אל האדמה: בשלום. שלם מן החטא, שלא אלמד מדרכי לבן: והיה ה' לי לאלהים. שיחול שמו עלי מתחלה ועד סוף, שלא ימצא פסול בזרעי, כמו שנאמר אשר דברתי לך; והבטחה זו הבטיח לאברהם, שנאמר להיות לך לאלהים ולזרעך אחריך (בראשית י״ז:ז׳):}%
% ---------- פרק 28 | פסוק 22 ----------
 \textbf{כב} וְהָאֶ֣בֶן הַזֹּ֗את אֲשֶׁר־שַׂ֙מְתִּי֙ מַצֵּבָ֔ה יִהְיֶ֖ה בֵּ֣ית אֱלֹהִ֑ים וְכֹל֙ אֲשֶׁ֣ר תִּתֶּן־לִ֔י עַשֵּׂ֖ר אֲעַשְּׂרֶ֥נּוּ לָֽךְ׃ \Rashi{\textbf{והאבן הזאת.}כך תפרש וי"ו זו של והאבן: אם תעשה לי את אלה, אני אעשה זאת: והאבן הזאת אשר שמתי מצבה וגו'. כתרגומו אהי פלח עלה קדם ה', וכן עשה בשובו מפדן ארם; כשאמר לו קום עלה בית אל (שם ל"ה י"ד) מה נאמר שם? ויצב יעקב מצבה ויסך עליה נסך:}%
% ---------- פרק 29 | פסוק 1 ----------
 \textbf{א} וַיִּשָּׂ֥א יַעֲקֹ֖ב רַגְלָ֑יו וַיֵּ֖לֶךְ אַ֥רְצָה בְנֵי־קֶֽדֶם׃ \Rashi{\textbf{וישא יעקב רגליו.}משנתבשר בשורה טובה שהבטח בשמירה נשא לבו את רגליו ונעשה קל ללכת, כך מפרש בב"ר:}%
% ---------- פרק 29 | פסוק 2 ----------
 \textbf{ב} וַיַּ֞רְא וְהִנֵּ֧ה בְאֵ֣ר בַּשָּׂדֶ֗ה וְהִנֵּה־שָׁ֞ם שְׁלֹשָׁ֤ה עֶדְרֵי־צֹאן֙ רֹבְצִ֣ים עָלֶ֔יהָ כִּ֚י מִן־הַבְּאֵ֣ר הַהִ֔וא יַשְׁק֖וּ הָעֲדָרִ֑ים וְהָאֶ֥בֶן גְּדֹלָ֖ה עַל־פִּ֥י הַבְּאֵֽר׃ \Rashi{\textbf{ישקו העדרים.}משקים הרועים את העדרים, והמקרא דבר בלשון קצרה:}%
% ---------- פרק 29 | פסוק 3 ----------
 \textbf{ג} וְנֶאֶסְפוּ־שָׁ֣מָּה כׇל־הָעֲדָרִ֗ים וְגָלְל֤וּ אֶת־הָאֶ֙בֶן֙ מֵעַל֙ פִּ֣י הַבְּאֵ֔ר וְהִשְׁק֖וּ אֶת־הַצֹּ֑אן וְהֵשִׁ֧יבוּ אֶת־הָאֶ֛בֶן עַל־פִּ֥י הַבְּאֵ֖ר לִמְקֹמָֽהּ׃ \Rashi{\textbf{ונאספו.}רגילים היו להאסף, לפי שהיתה האבן גדולה: וגללו. וגוללים, ותרגומו ומגנדרין, כל לשון הווה משתנה לדבר בלשון עתיד ובלשון עבר, לפי שכל דבר ההווה תמיד כבר היה ועתיד להיות: והשיבו. תרגומו ומתיבין:}%
% ---------- פרק 29 | פסוק 4 ----------
 \textbf{ד} וַיֹּ֤אמֶר לָהֶם֙ יַעֲקֹ֔ב אַחַ֖י מֵאַ֣יִן אַתֶּ֑ם וַיֹּ֣אמְר֔וּ מֵחָרָ֖ן אֲנָֽחְנוּ׃ %
% ---------- פרק 29 | פסוק 5 ----------
 \textbf{ה} וַיֹּ֣אמֶר לָהֶ֔ם הַיְדַעְתֶּ֖ם אֶת־לָבָ֣ן בֶּן־נָח֑וֹר וַיֹּאמְר֖וּ יָדָֽעְנוּ׃ %
% ---------- פרק 29 | פסוק 6 ----------
 \textbf{ו} וַיֹּ֥אמֶר לָהֶ֖ם הֲשָׁל֣וֹם ל֑וֹ וַיֹּאמְר֣וּ שָׁל֔וֹם וְהִנֵּה֙ רָחֵ֣ל בִּתּ֔וֹ בָּאָ֖ה עִם־הַצֹּֽאן׃ \Rashi{\textbf{באה עם הצאן.}הטעם באל"ף ותרגומו אתיא; ורחל באה – הטעם למעלה בבי"ת ותרגומו אתת; הראשון לשון עושה והשני לשון עשתה:}%
% ---------- פרק 29 | פסוק 7 ----------
 \textbf{ז} וַיֹּ֗אמֶר הֵ֥ן עוֹד֙ הַיּ֣וֹם גָּד֔וֹל לֹא־עֵ֖ת הֵאָסֵ֣ף הַמִּקְנֶ֑ה הַשְׁק֥וּ הַצֹּ֖אן וּלְכ֥וּ רְעֽוּ׃ \Rashi{\textbf{הן עוד היום גדול.}לפי שראה אותם רובצים, כסבור שרוצים לאסף המקנה הביתה ולא ירעו עוד, אמר להן הן עוד היום גדול, כלומר אם שכירים אתם לא שלמתם פעלת היום, ואם הבהמות שלכם אע"פ כן לא עת האסף המקנה וגו' (בראשית רבה):}%
% ---------- פרק 29 | פסוק 8 ----------
 \textbf{ח} וַיֹּאמְרוּ֮ לֹ֣א נוּכַל֒ עַ֣ד אֲשֶׁ֤ר יֵאָֽסְפוּ֙ כׇּל־הָ֣עֲדָרִ֔ים וְגָֽלְלוּ֙ אֶת־הָאֶ֔בֶן מֵעַ֖ל פִּ֣י הַבְּאֵ֑ר וְהִשְׁקִ֖ינוּ הַצֹּֽאן׃ \Rashi{\textbf{לא נוכל.}להשקות, לפי שהאבן גדולה: וגללו. זה מתרגם ויגנדרון, לפי שהוא לשון עתיד:}%
% ---------- פרק 29 | פסוק 9 ----------
 \textbf{ט} עוֹדֶ֖נּוּ מְדַבֵּ֣ר עִמָּ֑ם וְרָחֵ֣ל ׀ בָּ֗אָה עִם־הַצֹּאן֙ אֲשֶׁ֣ר לְאָבִ֔יהָ כִּ֥י רֹעָ֖ה הִֽוא׃ %
% ---------- פרק 29 | פסוק 10 ----------
 \textbf{י} וַיְהִ֡י כַּאֲשֶׁר֩ רָאָ֨ה יַעֲקֹ֜ב אֶת־רָחֵ֗ל בַּת־לָבָן֙ אֲחִ֣י אִמּ֔וֹ וְאֶת־צֹ֥אן לָבָ֖ן אֲחִ֣י אִמּ֑וֹ וַיִּגַּ֣שׁ יַעֲקֹ֗ב וַיָּ֤גֶל אֶת־הָאֶ֙בֶן֙ מֵעַל֙ פִּ֣י הַבְּאֵ֔ר וַיַּ֕שְׁקְ אֶת־צֹ֥אן לָבָ֖ן אֲחִ֥י אִמּֽוֹ׃ \Rashi{\textbf{ויגש יעקב ויגל.}כמי שמעביר את הפקק מעל פי צלוחית; להודיעך שכחו גדול (בראשית רבה):}%
% ---------- פרק 29 | פסוק 11 ----------
 \textbf{יא} וַיִּשַּׁ֥ק יַעֲקֹ֖ב לְרָחֵ֑ל וַיִּשָּׂ֥א אֶת־קֹל֖וֹ וַיֵּֽבְךְּ׃ \Rashi{\textbf{ויבך.}לפי שצפה ברוח הקדש שאינה נכנסת עמו לקבורה. ד"א לפי שבא בידים רקניות; אמר, אליעזר עבד אבי אבא היו בידיו נזמים וצמידים ומגדנות ואני אין בידי כלום; לפי שרדף אליפז בן עשו במצות אביו אחריו להרגו והשיגו, ולפי שגדל אליפז בחיקו של יצחק, משך ידיו. אמר לו מה אעשה לצווי של אבא? אמר לו יעקב טל מה שבידי, והעני חשוב כמת:}%
% ---------- פרק 29 | פסוק 12 ----------
 \textbf{יב} וַיַּגֵּ֨ד יַעֲקֹ֜ב לְרָחֵ֗ל כִּ֣י אֲחִ֤י אָבִ֙יהָ֙ ה֔וּא וְכִ֥י בֶן־רִבְקָ֖ה ה֑וּא וַתָּ֖רׇץ וַתַּגֵּ֥ד לְאָבִֽיהָ׃ \Rashi{\textbf{כי אחי אביה הוא.}קרוב לאביה, כמו אנשים אחים אנחנו. ומדרשו, אם לרמאות הוא בא, גם אני אחיו ברמאות, ואם אדם כשר הוא, גם אני בן רבקה אחותו הכשרה: ותגד לאביה. לפי שאמה מתה, ולא היה לה להגיד אלא לו (בראשית רבה):}%
% ---------- פרק 29 | פסוק 13 ----------
 \textbf{יג} וַיְהִי֩ כִשְׁמֹ֨עַ לָבָ֜ן אֶת־שֵׁ֣מַע ׀ יַעֲקֹ֣ב בֶּן־אֲחֹת֗וֹ וַיָּ֤רׇץ לִקְרָאתוֹ֙ וַיְחַבֶּק־לוֹ֙ וַיְנַשֶּׁק־ל֔וֹ וַיְבִיאֵ֖הוּ אֶל־בֵּית֑וֹ וַיְסַפֵּ֣ר לְלָבָ֔ן אֵ֥ת כׇּל־הַדְּבָרִ֖ים הָאֵֽלֶּה׃ \Rashi{\textbf{וירץ לקראתו.}כסבור ממון הוא טעון, שהרי עבד הבית בא לכאן בעשרה גמלים טעונים: ויחבק. כשלא ראה עמו כלום, אמר שמא זהובים הביא והנם בחיקו: וינשק לו. אמר שמא מרגליות הביא והם בפיו, ב"ר: ויספר ללבן. שלא בא אלא מתוך אנס אחיו, ושנטלו ממונו ממנו:}%
% ---------- פרק 29 | פסוק 14 ----------
 \textbf{יד} וַיֹּ֤אמֶר לוֹ֙ לָבָ֔ן אַ֛ךְ עַצְמִ֥י וּבְשָׂרִ֖י אָ֑תָּה וַיֵּ֥שֶׁב עִמּ֖וֹ חֹ֥דֶשׁ יָמִֽים׃ \Rashi{\textbf{אך עצמי ובשרי.}מעתה אין לי לאספך הביתה, הואיל ואין בידך כלום, אלא מפני קורבה אטפל בך חדש ימים, וכן עשה, ואף זו לא לחנם, שהיה רועה צאנו:}%
% ---------- פרק 29 | פסוק 15 ----------
 \textbf{טו} וַיֹּ֤אמֶר לָבָן֙ לְיַעֲקֹ֔ב הֲכִי־אָחִ֣י אַ֔תָּה וַעֲבַדְתַּ֖נִי חִנָּ֑ם הַגִּ֥ידָה לִּ֖י מַה־מַּשְׂכֻּרְתֶּֽךָ׃ \Rashi{\textbf{הכי אחי אתה.}לשון תמה, וכי בשביל שאחי אתה תעבדני חנם: ועבדתני. כמו ותעבדני; וכן כל תבה שהיא לשון עבר הוסיף וי"ו בראשה, והיא הופכת התבה להבא:}%
% ---------- פרק 29 | פסוק 16 ----------
 \textbf{טז} וּלְלָבָ֖ן שְׁתֵּ֣י בָנ֑וֹת שֵׁ֤ם הַגְּדֹלָה֙ לֵאָ֔ה וְשֵׁ֥ם הַקְּטַנָּ֖ה רָחֵֽל׃ %
% ---------- פרק 29 | פסוק 17 ----------
 \textbf{יז} וְעֵינֵ֥י לֵאָ֖ה רַכּ֑וֹת וְרָחֵל֙ הָֽיְתָ֔ה יְפַת־תֹּ֖אַר וִיפַ֥ת מַרְאֶֽה׃ \Rashi{\textbf{רכות.}שהיתה סבורה לעלות בגורלו של עשו ובוכה, שהיו הכל אומרים שני בנים לרבקה ושתי בנות ללבן, הגדולה לגדול והקטנה לקטן (בבא בתרא קכ"ג): תאר. הוא צורת הפרצוף, לשון יתארהו בשרד: (ישעיהו מ"ד, י"ג), קונ"פאש בלע"ז: מראה. הוא זיו קלסתר:}%
% ---------- פרק 29 | פסוק 18 ----------
 \textbf{יח} וַיֶּאֱהַ֥ב יַעֲקֹ֖ב אֶת־רָחֵ֑ל וַיֹּ֗אמֶר אֶֽעֱבׇדְךָ֙ שֶׁ֣בַע שָׁנִ֔ים בְּרָחֵ֥ל בִּתְּךָ֖ הַקְּטַנָּֽה׃ \Rashi{\textbf{אעבדך שבע שנים.}הם ימים אחדים שאמרה לו אמו וישבת עמו ימים אחדים (בראשית רבה); ותדע שכן הוא, שהרי כתיב ויהיו בעיניו כימים אחדים (בראשית כ"ט): ברחל בתך הקטנה. כל הסימנים הללו למה? לפי שהוא יודע בו שהוא רמאי, אמר לו, אעבדך ברחל, ושמא תאמר רחל אחרת מן השוק, תלמוד לומר בתך, ושמא תאמר אחליף ללאה שמה ואקרא שמה רחל, תלמוד לומר הקטנה; ואף על פי כן לא הועיל, שהרי רמהו:}%
% ---------- פרק 29 | פסוק 19 ----------
 \textbf{יט} וַיֹּ֣אמֶר לָבָ֗ן ט֚וֹב תִּתִּ֣י אֹתָ֣הּ לָ֔ךְ מִתִּתִּ֥י אֹתָ֖הּ לְאִ֣ישׁ אַחֵ֑ר שְׁבָ֖ה עִמָּדִֽי׃ %
% ---------- פרק 29 | פסוק 20 ----------
 \textbf{כ} וַיַּעֲבֹ֧ד יַעֲקֹ֛ב בְּרָחֵ֖ל שֶׁ֣בַע שָׁנִ֑ים וַיִּהְי֤וּ בְעֵינָיו֙ כְּיָמִ֣ים אֲחָדִ֔ים בְּאַהֲבָת֖וֹ אֹתָֽהּ׃ %
% ---------- פרק 29 | פסוק 21 ----------
 \textbf{כא} וַיֹּ֨אמֶר יַעֲקֹ֤ב אֶל־לָבָן֙ הָבָ֣ה אֶת־אִשְׁתִּ֔י כִּ֥י מָלְא֖וּ יָמָ֑י וְאָב֖וֹאָה אֵלֶֽיהָ׃ \Rashi{\textbf{מלאו ימי.}שאמרה לי אמי, ועוד מלאו ימי, שהרי אני בן פ"ד שנה ואימתי אעמיד י"ב שבטים? וזהו שאמר ואבואה אליה, והלא קל שבקלים אינו אומר כן? אלא להוליד תולדות אמר כן:}%
% ---------- פרק 29 | פסוק 22 ----------
 \textbf{כב} וַיֶּאֱסֹ֥ף לָבָ֛ן אֶת־כׇּל־אַנְשֵׁ֥י הַמָּק֖וֹם וַיַּ֥עַשׂ מִשְׁתֶּֽה׃ %
% ---------- פרק 29 | פסוק 23 ----------
 \textbf{כג} וַיְהִ֣י בָעֶ֔רֶב וַיִּקַּח֙ אֶת־לֵאָ֣ה בִתּ֔וֹ וַיָּבֵ֥א אֹתָ֖הּ אֵלָ֑יו וַיָּבֹ֖א אֵלֶֽיהָ׃ %
% ---------- פרק 29 | פסוק 24 ----------
 \textbf{כד} וַיִּתֵּ֤ן לָבָן֙ לָ֔הּ אֶת־זִלְפָּ֖ה שִׁפְחָת֑וֹ לְלֵאָ֥ה בִתּ֖וֹ שִׁפְחָֽה׃ %
% ---------- פרק 29 | פסוק 25 ----------
 \textbf{כה} וַיְהִ֣י בַבֹּ֔קֶר וְהִנֵּה־הִ֖וא לֵאָ֑ה וַיֹּ֣אמֶר אֶל־לָבָ֗ן מַה־זֹּאת֙ עָשִׂ֣יתָ לִּ֔י הֲלֹ֤א בְרָחֵל֙ עָבַ֣דְתִּי עִמָּ֔ךְ וְלָ֖מָּה רִמִּיתָֽנִי׃ \Rashi{\textbf{ויהי בבקר והנה היא לאה.}אבל בלילה לא היתה לאה, לפי שמסר יעקב סימנים לרחל, וכשראתה רחל שמכניסין לו לאה אמרה: עכשו תכלם אחותי, עמדה ומסרה לה אותן סימנים (מגילה י"ג):}%
% ---------- פרק 29 | פסוק 26 ----------
 \textbf{כו} וַיֹּ֣אמֶר לָבָ֔ן לֹא־יֵעָשֶׂ֥ה כֵ֖ן בִּמְקוֹמֵ֑נוּ לָתֵ֥ת הַצְּעִירָ֖ה לִפְנֵ֥י הַבְּכִירָֽה׃ %
% ---------- פרק 29 | פסוק 27 ----------
 \textbf{כז} מַלֵּ֖א שְׁבֻ֣עַ זֹ֑את וְנִתְּנָ֨ה לְךָ֜ גַּם־אֶת־זֹ֗את בַּעֲבֹדָה֙ אֲשֶׁ֣ר תַּעֲבֹ֣ד עִמָּדִ֔י ע֖וֹד שֶֽׁבַע־שָׁנִ֥ים אֲחֵרֽוֹת׃ \Rashi{\textbf{מלא שבע זאת.}דבוק הוא, שהרי נקוד בחטף, שבוע של זאת, והן שבעת ימי המשתה, בתלמוד ירוש' במועד קטן (וא"א לומר שבוע ממש, שאם כן היה צריך להנקד בפתח השי"ן; ועוד ששבוע לשון זכר, כדכתיב שבעה שבעת תספר לך, לפיכך אין משמע שבוע אלא שבעה, שטיי"נא בלעז): ונתנה לך. לשון רבים, כמו נרדה ונבלה ונשרפה (בראשית י"א), אף זה לשון ונתן: גם את זאת. מיד לאחר שבעת ימי המשתה, ותעבד לאחר נשואיה:}%
% ---------- פרק 29 | פסוק 28 ----------
 \textbf{כח} וַיַּ֤עַשׂ יַעֲקֹב֙ כֵּ֔ן וַיְמַלֵּ֖א שְׁבֻ֣עַ זֹ֑את וַיִּתֶּן־ל֛וֹ אֶת־רָחֵ֥ל בִּתּ֖וֹ ל֥וֹ לְאִשָּֽׁה׃ %
% ---------- פרק 29 | פסוק 29 ----------
 \textbf{כט} וַיִּתֵּ֤ן לָבָן֙ לְרָחֵ֣ל בִּתּ֔וֹ אֶת־בִּלְהָ֖ה שִׁפְחָת֑וֹ לָ֖הּ לְשִׁפְחָֽה׃ %
% ---------- פרק 29 | פסוק 30 ----------
 \textbf{ל} וַיָּבֹא֙ גַּ֣ם אֶל־רָחֵ֔ל וַיֶּאֱהַ֥ב גַּֽם־אֶת־רָחֵ֖ל מִלֵּאָ֑ה וַיַּעֲבֹ֣ד עִמּ֔וֹ ע֖וֹד שֶֽׁבַע־שָׁנִ֥ים אֲחֵרֽוֹת׃ \Rashi{\textbf{עוד שבע שנים אחרות.}אחרות הקישן לראשונות, מה ראשונות באמונה אף האחרונות באמונה, ואע"פ שברמאות בא עליו (בראשית רבה):}%
% ---------- פרק 29 | פסוק 31 ----------
 \textbf{לא} וַיַּ֤רְא יְהֹוָה֙ כִּֽי־שְׂנוּאָ֣ה לֵאָ֔ה וַיִּפְתַּ֖ח אֶת־רַחְמָ֑הּ וְרָחֵ֖ל עֲקָרָֽה׃ %
% ---------- פרק 29 | פסוק 32 ----------
 \textbf{לב} וַתַּ֤הַר לֵאָה֙ וַתֵּ֣לֶד בֵּ֔ן וַתִּקְרָ֥א שְׁמ֖וֹ רְאוּבֵ֑ן כִּ֣י אָֽמְרָ֗ה כִּֽי־רָאָ֤ה יְהֹוָה֙ בְּעׇנְיִ֔י כִּ֥י עַתָּ֖ה יֶאֱהָבַ֥נִי אִישִֽׁי׃ \Rashi{(ותקרא שמו ראובן רבותינו פרשו, אמרה ראו מה בין בני לבן חמי שמכר הבכורה ליעקב, וזה לא מכרה ליוסף ולא ערער עליו, ולא עוד שלא ערער עליו אלא שבקש להוציאו מן הבור) (ברכות ז'):}%
% ---------- פרק 29 | פסוק 33 ----------
 \textbf{לג} וַתַּ֣הַר עוֹד֮ וַתֵּ֣לֶד בֵּן֒ וַתֹּ֗אמֶר כִּֽי־שָׁמַ֤ע יְהֹוָה֙ כִּֽי־שְׂנוּאָ֣ה אָנֹ֔כִי וַיִּתֶּן־לִ֖י גַּם־אֶת־זֶ֑ה וַתִּקְרָ֥א שְׁמ֖וֹ שִׁמְעֽוֹן׃ %
% ---------- פרק 29 | פסוק 34 ----------
 \textbf{לד} וַתַּ֣הַר עוֹד֮ וַתֵּ֣לֶד בֵּן֒ וַתֹּ֗אמֶר עַתָּ֤ה הַפַּ֙עַם֙ יִלָּוֶ֤ה אִישִׁי֙ אֵלַ֔י כִּֽי־יָלַ֥דְתִּי ל֖וֹ שְׁלֹשָׁ֣ה בָנִ֑ים עַל־כֵּ֥ן קָרָֽא־שְׁמ֖וֹ לֵוִֽי׃ \Rashi{\textbf{הפעם ילוה אישי.}לפי שהאמהות נביאות היו, ויודעות שי"ב שבטים יוצאים מיעקב וד' נשים ישא, אמרה, מעתה אין לו פתחון פה עלי, שהרי נטלתי כל חלקי בבנים (ברבות ס'): על כן. כל מי שנאמר בו על כן מרבה באכלוסין, חוץ מלוי שהארון היה מכלה בהם: קרא שמו לוי. בכלם כתיב ותקרא, וזה כתב בו קרא, ויש מ"א באלה הדברים רבה ששלח הקב"ה גבריאל והביאו לפניו וקרא לו שם זה, ונתן לו כ"ד מתנות כהנה, ועל שם שלוהו במתנות קראו לוי:}%
% ---------- פרק 29 | פסוק 35 ----------
 \textbf{לה} וַתַּ֨הַר ע֜וֹד וַתֵּ֣לֶד בֵּ֗ן וַתֹּ֙אמֶר֙ הַפַּ֙עַם֙ אוֹדֶ֣ה אֶת־יְהֹוָ֔ה עַל־כֵּ֛ן קָרְאָ֥ה שְׁמ֖וֹ יְהוּדָ֑ה וַֽתַּעֲמֹ֖ד מִלֶּֽדֶת׃ \Rashi{\textbf{הפעם אודה.}שנטלתי יותר מחלקי, מעתה יש לי להודות:}%
% ---------- פרק 30 | פסוק 1 ----------
 \textbf{א} וַתֵּ֣רֶא רָחֵ֗ל כִּ֣י לֹ֤א יָֽלְדָה֙ לְיַעֲקֹ֔ב וַתְּקַנֵּ֥א רָחֵ֖ל בַּאֲחֹתָ֑הּ וַתֹּ֤אמֶר אֶֽל־יַעֲקֹב֙ הָֽבָה־לִּ֣י בָנִ֔ים וְאִם־אַ֖יִן מֵתָ֥ה אָנֹֽכִי׃ \Rashi{\textbf{ותקנא רחל באחותה.}קנאה במעשיה הטובים, אמרה אלולי שצדקה ממני לא זכתה לבנים (בראשית רבה): הבה לי. וכי כך עשה אביך לאמך? והלא התפלל עליה: מתה אנכי. מכאן למי שאין לו בנים שחשוב כמת (בראשית רבה):}%
% ---------- פרק 30 | פסוק 2 ----------
 \textbf{ב} וַיִּֽחַר־אַ֥ף יַעֲקֹ֖ב בְּרָחֵ֑ל וַיֹּ֗אמֶר הֲתַ֤חַת אֱלֹהִים֙ אָנֹ֔כִי אֲשֶׁר־מָנַ֥ע מִמֵּ֖ךְ פְּרִי־בָֽטֶן׃ \Rashi{\textbf{התחת.}וכי במקומו אני? אשר מנע ממך. את אומרת שאעשה כאבא, אני איני כאבא, אבא לא היו לו בנים, אני יש לי בנים; ממך מנע ולא ממני:}%
% ---------- פרק 30 | פסוק 3 ----------
 \textbf{ג} וַתֹּ֕אמֶר הִנֵּ֛ה אֲמָתִ֥י בִלְהָ֖ה בֹּ֣א אֵלֶ֑יהָ וְתֵלֵד֙ עַל־בִּרְכַּ֔י וְאִבָּנֶ֥ה גַם־אָנֹכִ֖י מִמֶּֽנָּה׃ \Rashi{\textbf{על ברכי.}כתרגומו ואנא ארבי: ואבנה גם אנכי. מהו גם? אמרה לו זקנך אברהם היו לו בנים מהגר וחגר מתניו כנגד שרה, אמר לה זקנתי הכניסה צרתה לביתה, אמרה לו אם הדבר הזה מעכב, הנה אמתי: ואבנה גם אנכי. כשרה:}%
% ---------- פרק 30 | פסוק 4 ----------
 \textbf{ד} וַתִּתֶּן־ל֛וֹ אֶת־בִּלְהָ֥ה שִׁפְחָתָ֖הּ לְאִשָּׁ֑ה וַיָּבֹ֥א אֵלֶ֖יהָ יַעֲקֹֽב׃ %
% ---------- פרק 30 | פסוק 5 ----------
 \textbf{ה} וַתַּ֣הַר בִּלְהָ֔ה וַתֵּ֥לֶד לְיַעֲקֹ֖ב בֵּֽן׃ %
% ---------- פרק 30 | פסוק 6 ----------
 \textbf{ו} וַתֹּ֤אמֶר רָחֵל֙ דָּנַ֣נִּי אֱלֹהִ֔ים וְגַם֙ שָׁמַ֣ע בְּקֹלִ֔י וַיִּתֶּן־לִ֖י בֵּ֑ן עַל־כֵּ֛ן קָרְאָ֥ה שְׁמ֖וֹ דָּֽן׃ \Rashi{\textbf{דנני אלהים.}דנני וחיבני וזכני (בראשית רבה):}%
% ---------- פרק 30 | פסוק 7 ----------
 \textbf{ז} וַתַּ֣הַר ע֔וֹד וַתֵּ֕לֶד בִּלְהָ֖ה שִׁפְחַ֣ת רָחֵ֑ל בֵּ֥ן שֵׁנִ֖י לְיַעֲקֹֽב׃ %
% ---------- פרק 30 | פסוק 8 ----------
 \textbf{ח} וַתֹּ֣אמֶר רָחֵ֗ל נַפְתּוּלֵ֨י אֱלֹהִ֧ים ׀ נִפְתַּ֛לְתִּי עִם־אֲחֹתִ֖י גַּם־יָכֹ֑לְתִּי וַתִּקְרָ֥א שְׁמ֖וֹ נַפְתָּלִֽי׃ \Rashi{\textbf{נפתולי אלהים.}מנחם בן סרוק פרשו במחברת צמיד פתיל, חבורים מאת המקום נתחברתי עם אחותי לזכות לבנים. ואני מפרשו לשון עקש ופתלתל (דברים ל"ב) – נתעקשתי והפצרתי פצירות ונפתולים הרבה למקום, להיות שוה לאחותי: גם יכלתי. הסכים על ידי; ואנקלוס תרגם לשון תפלה, כמו נפתולי אלהים נתפלתי בקשות החביבות לפניו; נתקבלתי, ונעתרתי כאחותי: נפתלתי. נתקבלה תפלתי, ומ"א יש רבים בלשון נוטריקון:}%
% ---------- פרק 30 | פסוק 9 ----------
 \textbf{ט} וַתֵּ֣רֶא לֵאָ֔ה כִּ֥י עָמְדָ֖ה מִלֶּ֑דֶת וַתִּקַּח֙ אֶת־זִלְפָּ֣ה שִׁפְחָתָ֔הּ וַתִּתֵּ֥ן אֹתָ֛הּ לְיַעֲקֹ֖ב לְאִשָּֽׁה׃ %
% ---------- פרק 30 | פסוק 10 ----------
 \textbf{י} וַתֵּ֗לֶד זִלְפָּ֛ה שִׁפְחַ֥ת לֵאָ֖ה לְיַעֲקֹ֥ב בֵּֽן׃ \Rashi{\textbf{ותלד זלפה.}בכלן נאמר הריון חוץ מזלפה, לפי שהיתה בחורה מכלן ותינקת בשנים ואין הריון נכר בה; וכדי לרמות ליעקב, נתנה לבן ללאה, שלא יבין שמכניסין לו את לאה, שכך מנהג לתן שפחה הגדולה לגדולה וקטנה לקטנה:}%
% ---------- פרק 30 | פסוק 11 ----------
 \textbf{יא} וַתֹּ֥אמֶר לֵאָ֖ה (בגד) [בָּ֣א גָ֑ד] וַתִּקְרָ֥א אֶת־שְׁמ֖וֹ גָּֽד׃ \Rashi{\textbf{בא גד.}בא מזל טוב, כמו גד גדי וסנוק לא (שבת ס"ז), ודומה לו הערכים לגד שלחן (ישעיהו ס"ה), ומ"א שנולד מהול, כמו גדו אילנא (דניאל ד'), ולא ידעתי על מה נכתבה תבה אחת; ד"א למה נקראת תבה אחת בגד, כמו בגדת בי, כשבאת אל שפחתי, כאיש שבגד באשת נעורים:}%
% ---------- פרק 30 | פסוק 12 ----------
 \textbf{יב} וַתֵּ֗לֶד זִלְפָּה֙ שִׁפְחַ֣ת לֵאָ֔ה בֵּ֥ן שֵׁנִ֖י לְיַעֲקֹֽב׃ %
% ---------- פרק 30 | פסוק 13 ----------
 \textbf{יג} וַתֹּ֣אמֶר לֵאָ֔ה בְּאׇשְׁרִ֕י כִּ֥י אִשְּׁר֖וּנִי בָּנ֑וֹת וַתִּקְרָ֥א אֶת־שְׁמ֖וֹ אָשֵֽׁר׃ %
% ---------- פרק 30 | פסוק 14 ----------
 \textbf{יד} וַיֵּ֨לֶךְ רְאוּבֵ֜ן בִּימֵ֣י קְצִיר־חִטִּ֗ים וַיִּמְצָ֤א דֽוּדָאִים֙ בַּשָּׂדֶ֔ה וַיָּבֵ֣א אֹתָ֔ם אֶל־לֵאָ֖ה אִמּ֑וֹ וַתֹּ֤אמֶר רָחֵל֙ אֶל־לֵאָ֔ה תְּנִי־נָ֣א לִ֔י מִדּוּדָאֵ֖י בְּנֵֽךְ׃ \Rashi{\textbf{בימי קציר חטים.}להגיד שבחן של שבטים; שעת הקציר היה ולא פשט ידו בגזל להביא חטים ושעורים, אלא דבר ההפקר, שאין אדם מקפיד בו: דודאים. סיגלי, עשב הוא (סנהדרין צ"ט), ובלשון ישמעאל יסמי"ן:}%
% ---------- פרק 30 | פסוק 15 ----------
 \textbf{טו} וַתֹּ֣אמֶר לָ֗הּ הַמְעַט֙ קַחְתֵּ֣ךְ אֶת־אִישִׁ֔י וְלָקַ֕חַת גַּ֥ם אֶת־דּוּדָאֵ֖י בְּנִ֑י וַתֹּ֣אמֶר רָחֵ֗ל לָכֵן֙ יִשְׁכַּ֤ב עִמָּךְ֙ הַלַּ֔יְלָה תַּ֖חַת דּוּדָאֵ֥י בְנֵֽךְ׃ \Rashi{\textbf{ולקחת גם את דודאי בני.}בתמיה, ולעשות עוד זאת לקח גם את דודאי בני; ותרגומו ולמיסב: לכן ישכב עמך הלילה. שלי היתה שכיבת לילה זו, ואני נותנה לך תחת דודאי בנך, ולפי שזלזלה במשכב הצדיק לא זכתה להקבר עמו (נדה ל"א):}%
% ---------- פרק 30 | פסוק 16 ----------
 \textbf{טז} וַיָּבֹ֨א יַעֲקֹ֣ב מִן־הַשָּׂדֶה֮ בָּעֶ֒רֶב֒ וַתֵּצֵ֨א לֵאָ֜ה לִקְרָאת֗וֹ וַתֹּ֙אמֶר֙ אֵלַ֣י תָּב֔וֹא כִּ֚י שָׂכֹ֣ר שְׂכַרְתִּ֔יךָ בְּדוּדָאֵ֖י בְּנִ֑י וַיִּשְׁכַּ֥ב עִמָּ֖הּ בַּלַּ֥יְלָה הֽוּא׃ \Rashi{\textbf{שכר שכרתיך.}נתתי לרחל שכרה: בלילה הוא. הקב"ה סיעו, שיצא משם יששכר (בראשית רבה):}%
% ---------- פרק 30 | פסוק 17 ----------
 \textbf{יז} וַיִּשְׁמַ֥ע אֱלֹהִ֖ים אֶל־לֵאָ֑ה וַתַּ֛הַר וַתֵּ֥לֶד לְיַעֲקֹ֖ב בֵּ֥ן חֲמִישִֽׁי׃ \Rashi{\textbf{וישמע אלהים אל לאה.}שהיתה מתאוה ומחזרת להרבות שבטים:}%
% ---------- פרק 30 | פסוק 18 ----------
 \textbf{יח} וַתֹּ֣אמֶר לֵאָ֗ה נָתַ֤ן אֱלֹהִים֙ שְׂכָרִ֔י אֲשֶׁר־נָתַ֥תִּי שִׁפְחָתִ֖י לְאִישִׁ֑י וַתִּקְרָ֥א שְׁמ֖וֹ יִשָּׂשכָֽר׃ %
% ---------- פרק 30 | פסוק 19 ----------
 \textbf{יט} וַתַּ֤הַר עוֹד֙ לֵאָ֔ה וַתֵּ֥לֶד בֵּן־שִׁשִּׁ֖י לְיַעֲקֹֽב׃ %
% ---------- פרק 30 | פסוק 20 ----------
 \textbf{כ} וַתֹּ֣אמֶר לֵאָ֗ה זְבָדַ֨נִי אֱלֹהִ֥ים ׀ אֹתִי֮ זֵ֣בֶד טוֹב֒ הַפַּ֙עַם֙ יִזְבְּלֵ֣נִי אִישִׁ֔י כִּֽי־יָלַ֥דְתִּי ל֖וֹ שִׁשָּׁ֣ה בָנִ֑ים וַתִּקְרָ֥א אֶת־שְׁמ֖וֹ זְבֻלֽוּן׃ \Rashi{\textbf{זבד טוב.}כתרגומו: יזבלני. ל' בית זבל (מלכים א ח' י"ג), הירבריי"א בלעז, בית מדור, מעתה לא תהא עקר דירתו אלא עמי שיש לי בנים כנגד כל נשיו:}%
% ---------- פרק 30 | פסוק 21 ----------
 \textbf{כא} וְאַחַ֖ר יָ֣לְדָה בַּ֑ת וַתִּקְרָ֥א אֶת־שְׁמָ֖הּ דִּינָֽה׃ \Rashi{\textbf{דינה.}פרשו רבותינו שדנה לאה דין בעצמה, אם זה זכר לא תהא רחל אחותי כאחת השפחות, והתפללה עליו ונהפך לנקבה (ברכות ס'):}%
% ---------- פרק 30 | פסוק 22 ----------
 \textbf{כב} וַיִּזְכֹּ֥ר אֱלֹהִ֖ים אֶת־רָחֵ֑ל וַיִּשְׁמַ֤ע אֵלֶ֙יהָ֙ אֱלֹהִ֔ים וַיִּפְתַּ֖ח אֶת־רַחְמָֽהּ׃ \Rashi{\textbf{ויזכור אלהים את רחל.}זכר לה שמסרה סימניה לאחותה ושהיתה מצרה שמא תעלה בגורלו של עשו, שמא יגרשנה יעקב לפי שאין לה בנים, ואף עשו הרשע כך עלה בלבו כששמע שאין לה בנים; הוא שיסד הפיט האדמון כבט שלא חלה, צבה לקחתה לו ונתבהלה:}%
% ---------- פרק 30 | פסוק 23 ----------
 \textbf{כג} וַתַּ֖הַר וַתֵּ֣לֶד בֵּ֑ן וַתֹּ֕אמֶר אָסַ֥ף אֱלֹהִ֖ים אֶת־חֶרְפָּתִֽי׃ \Rashi{\textbf{אסף.}הכניסה במקום שלא תראה, וכן אסף חרפתנו (ישעיהו ד'), ולא יאסף הביתה (שמות ט'), אספו נגהם (יואל ד'), וירחך לא יאסף (ישעיהו ס') – לא יטמן: חרפתי. שהייתי לחרפה שאני עקרה והיו אומרים עלי שאעלה לחלקו של עשו הרשע. ואגדה: כל זמן שאין לאשה בן, אין לה במי לתלות סרחונה, משיש לה בן תולה בו: מי שבר כלי זה? בנך. מי אכל תאנים אלו? בנך:}%
% ---------- פרק 30 | פסוק 24 ----------
 \textbf{כד} וַתִּקְרָ֧א אֶת־שְׁמ֛וֹ יוֹסֵ֖ף לֵאמֹ֑ר יֹסֵ֧ף יְהֹוָ֛ה לִ֖י בֵּ֥ן אַחֵֽר׃ \Rashi{\textbf{יסף ה' לי בן אחר.}יודעת היתה בנבואה שאין יעקב עתיד להעמיד אלא שנים עשר שבטים, אמרה יהי רצון שאותו שהוא עתיד להעמיד יהא ממני, לכך לא נתפללה אלא על בן אחר:}%
% ---------- פרק 30 | פסוק 25 ----------
 \textbf{כה} וַיְהִ֕י כַּאֲשֶׁ֛ר יָלְדָ֥ה רָחֵ֖ל אֶת־יוֹסֵ֑ף וַיֹּ֤אמֶר יַעֲקֹב֙ אֶל־לָבָ֔ן שַׁלְּחֵ֙נִי֙ וְאֵ֣לְכָ֔ה אֶל־מְקוֹמִ֖י וּלְאַרְצִֽי׃ \Rashi{\textbf{כאשר ילדה רחל את יוסף.}משנולד שטנו של עשו, שנאמר והיה בית יעקב אש ובית יוסף להבה ובית עשו לקש (עובדיה א'), אש בלא להבה אינו שולט למרחוק, משנולד יוסף בטח יעקב בהקב"ה ורצה לשוב:}%
% ---------- פרק 30 | פסוק 26 ----------
 \textbf{כו} תְּנָ֞ה אֶת־נָשַׁ֣י וְאֶת־יְלָדַ֗י אֲשֶׁ֨ר עָבַ֧דְתִּי אֹֽתְךָ֛ בָּהֵ֖ן וְאֵלֵ֑כָה כִּ֚י אַתָּ֣ה יָדַ֔עְתָּ אֶת־עֲבֹדָתִ֖י אֲשֶׁ֥ר עֲבַדְתִּֽיךָ׃ \Rashi{\textbf{תנה את נשי וגו'.}איני רוצה לצאת כי אם ברשות:}%
% ---------- פרק 30 | פסוק 27 ----------
 \textbf{כז} וַיֹּ֤אמֶר אֵלָיו֙ לָבָ֔ן אִם־נָ֛א מָצָ֥אתִי חֵ֖ן בְּעֵינֶ֑יךָ נִחַ֕שְׁתִּי וַיְבָרְכֵ֥נִי יְהֹוָ֖ה בִּגְלָלֶֽךָ׃ \Rashi{\textbf{נחשתי.}מנחש היה, נסיתי בנחוש שלי שעל ידך באה לי ברכה; כשבאת לכאן לא היו לי בנים, שנאמר והנה רחל בתו באה עם הצאן (בראשית כ״ט:ו׳), אפשר יש לו בנים והוא שולח בתו אצל הרועים? עכשיו היו לו בנים, שנאמר וישמע את דברי בני לבן (שם ל"א):}%
% ---------- פרק 30 | פסוק 28 ----------
 \textbf{כח} וַיֹּאמַ֑ר נׇקְבָ֧ה שְׂכָרְךָ֛ עָלַ֖י וְאֶתֵּֽנָה׃ \Rashi{\textbf{נקבה.}כתרגומו פרש אגרך:}%
% ---------- פרק 30 | פסוק 29 ----------
 \textbf{כט} וַיֹּ֣אמֶר אֵלָ֔יו אַתָּ֣ה יָדַ֔עְתָּ אֵ֖ת אֲשֶׁ֣ר עֲבַדְתִּ֑יךָ וְאֵ֛ת אֲשֶׁר־הָיָ֥ה מִקְנְךָ֖ אִתִּֽי׃ \Rashi{\textbf{ואת אשר היה מקנך אתי.}את חשבון מעוט מקנך שבא לידי מתחלה, כמה היו:}%
% ---------- פרק 30 | פסוק 30 ----------
 \textbf{ל} כִּ֡י מְעַט֩ אֲשֶׁר־הָיָ֨ה לְךָ֤ לְפָנַי֙ וַיִּפְרֹ֣ץ לָרֹ֔ב וַיְבָ֧רֶךְ יְהֹוָ֛ה אֹתְךָ֖ לְרַגְלִ֑י וְעַתָּ֗ה מָתַ֛י אֶֽעֱשֶׂ֥ה גַם־אָנֹכִ֖י לְבֵיתִֽי׃ \Rashi{\textbf{לרגלי.}עם רגלי, בשביל ביאת רגלי באת אצלך הברכה, כמו העם אשר ברגליך (שמות י"א), לעם אשר ברגלי (שופטים ח') – הבאים עמי: גם אנכי לביתי. לצרך ביתי; עכשו אין עושין לצרכי אלא בני וצריך אני להיות עושה גם אני עמהם לסמכן, וזהו גם:}%
% ---------- פרק 30 | פסוק 31 ----------
 \textbf{לא} וַיֹּ֖אמֶר מָ֣ה אֶתֶּן־לָ֑ךְ וַיֹּ֤אמֶר יַעֲקֹב֙ לֹא־תִתֶּן־לִ֣י מְא֔וּמָה אִם־תַּֽעֲשֶׂה־לִּי֙ הַדָּבָ֣ר הַזֶּ֔ה אָשׁ֛וּבָה אֶרְעֶ֥ה צֹֽאנְךָ֖ אֶשְׁמֹֽר׃ %
% ---------- פרק 30 | פסוק 32 ----------
 \textbf{לב} אֶֽעֱבֹ֨ר בְּכׇל־צֹֽאנְךָ֜ הַיּ֗וֹם הָסֵ֨ר מִשָּׁ֜ם כׇּל־שֶׂ֣ה ׀ נָקֹ֣ד וְטָל֗וּא וְכׇל־שֶׂה־חוּם֙ בַּכְּשָׂבִ֔ים וְטָל֥וּא וְנָקֹ֖ד בָּעִזִּ֑ים וְהָיָ֖ה שְׂכָרִֽי׃ \Rashi{\textbf{נקד.}מנמר בחברבורות דקות כמו נקדות, פיונטור"א בלעז: וטלוא. לשון טלאים – חברבורות רחבות: חום. שחום, דומה לאדם, רו"ש בלעז; לשון משנה שחמתית ונמצאת לבנה, לענין התבואה: והיה שכרי. אותן שיולדו מכאן ולהבא נקדים וטלואים בעזים ושחומים בכשבים, יהיו שלי, ואותן שישנן עכשו, הפרש מהם והפקידם ביד בניך, שלא תאמר לי על הנולדים מעתה אלו היו שם מתחלה, ועוד, שלא תאמר לי, על ידי הזכרים שהן נקדים וטלואים תלדנה הנקבות דגמתן מכאן ואילך:}%
% ---------- פרק 30 | פסוק 33 ----------
 \textbf{לג} וְעָֽנְתָה־בִּ֤י צִדְקָתִי֙ בְּי֣וֹם מָחָ֔ר כִּֽי־תָב֥וֹא עַל־שְׂכָרִ֖י לְפָנֶ֑יךָ כֹּ֣ל אֲשֶׁר־אֵינֶ֩נּוּ֩ נָקֹ֨ד וְטָל֜וּא בָּֽעִזִּ֗ים וְחוּם֙ בַּכְּשָׂבִ֔ים גָּנ֥וּב ה֖וּא אִתִּֽי׃ \Rashi{\textbf{וענתה בי וגו'.}אם תחשדני שאני נוטל משלך כלום, תענה בי צדקתי כי תבא צדקתי ותעיד על שכרי לפניך, שלא תמצא בעדרי כי אם נקדים וטלואים, וכל שתמצא בהן שאינו נקד או טלוא או חום, בידוע שגנבתיו לך ובגנבה הוא שרוי אצלי:}%
% ---------- פרק 30 | פסוק 34 ----------
 \textbf{לד} וַיֹּ֥אמֶר לָבָ֖ן הֵ֑ן ל֖וּ יְהִ֥י כִדְבָרֶֽךָ׃ \Rashi{\textbf{הן.}ל' קבלת דברים: לו יהי כדברך. הלואי שתחפץ בכך:}%
% ---------- פרק 30 | פסוק 35 ----------
 \textbf{לה} וַיָּ֣סַר בַּיּוֹם֩ הַה֨וּא אֶת־הַתְּיָשִׁ֜ים הָֽעֲקֻדִּ֣ים וְהַטְּלֻאִ֗ים וְאֵ֤ת כׇּל־הָֽעִזִּים֙ הַנְּקֻדּ֣וֹת וְהַטְּלֻאֹ֔ת כֹּ֤ל אֲשֶׁר־לָבָן֙ בּ֔וֹ וְכׇל־ח֖וּם בַּכְּשָׂבִ֑ים וַיִּתֵּ֖ן בְּיַד־בָּנָֽיו׃ \Rashi{\textbf{ויסר.}לבן ביום ההוא: את התישים. עזים זכרים: כל אשר לבן בו. כל אשר היתה בו חברבורית לבנה: ויתן. לבן ביד בניו:}%
% ---------- פרק 30 | פסוק 36 ----------
 \textbf{לו} וַיָּ֗שֶׂם דֶּ֚רֶךְ שְׁלֹ֣שֶׁת יָמִ֔ים בֵּינ֖וֹ וּבֵ֣ין יַעֲקֹ֑ב וְיַעֲקֹ֗ב רֹעֶ֛ה אֶת־צֹ֥אן לָבָ֖ן הַנּוֹתָרֹֽת׃ \Rashi{\textbf{הנותרות.}הרעועות שבהן, החולות והעקרות שאינן אלא שירים, אותן מסר לו:}%
% ---------- פרק 30 | פסוק 37 ----------
 \textbf{לז} וַיִּֽקַּֽח־ל֣וֹ יַעֲקֹ֗ב מַקַּ֥ל לִבְנֶ֛ה לַ֖ח וְל֣וּז וְעַרְמ֑וֹן וַיְפַצֵּ֤ל בָּהֵן֙ פְּצָל֣וֹת לְבָנ֔וֹת מַחְשֹׂף֙ הַלָּבָ֔ן אֲשֶׁ֖ר עַל־הַמַּקְלֽוֹת׃ \Rashi{\textbf{מקל לבנה.}עץ הוא, ושמו לבנה, כמא דתימא תחת אלון ולבנה (הושע ד'), ואומר אני הוא שקורין טרינבל"א בלע"ז, שהוא לבן: לח. כשהוא רטב: לוז. ועוד לקח מקל לוז, עץ שגדלין בו אגוזים דקים, קולדו"י בלעז: וערמון. קשטני"יר בלעז: פצלות. קלופים קלופים, שהיה עושהו מנמר: מחשוף הלבן. גלוי לבן של מקל, כשהיה קולפו היה נראה ונגלה לבן שלו במקום הקלוף:}%
% ---------- פרק 30 | פסוק 38 ----------
 \textbf{לח} וַיַּצֵּ֗ג אֶת־הַמַּקְלוֹת֙ אֲשֶׁ֣ר פִּצֵּ֔ל בָּרְהָטִ֖ים בְּשִֽׁקְת֣וֹת הַמָּ֑יִם אֲשֶׁר֩ תָּבֹ֨אןָ הַצֹּ֤אן לִשְׁתּוֹת֙ לְנֹ֣כַח הַצֹּ֔אן וַיֵּחַ֖מְנָה בְּבֹאָ֥ן לִשְׁתּֽוֹת׃ \Rashi{\textbf{ויצג.}תרגומו ודעיץ, לשון תחיבה ונעיצה הוא בלשון ארמי, והרבה יש בתלמוד דצה ושלפה, דץ ביה מידי, דצה כמו דעצה, אלא שמקצר את לשונו: ברהטים. במרוצות המים, בברכות העשויות בארץ להשקות שם הצאן: אשר תבאן וגו'. ברהטים אשר תבאנה הצאן לשתות שם הציג המקלות לנכח הצאן: ויחמנה. הבהמה רואה את המקלות והיא נרתעת לאחוריה והזכר רובעה ויולדת כיוצא בו. רבי הושעיה אומר, המים נעשין זרע במעיהן ולא היו צריכות לזכר, וזהו ויחמנה וגו' (בראשית רבה):}%
% ---------- פרק 30 | פסוק 39 ----------
 \textbf{לט} וַיֶּחֱמ֥וּ הַצֹּ֖אן אֶל־הַמַּקְל֑וֹת וַתֵּלַ֣דְןָ הַצֹּ֔אן עֲקֻדִּ֥ים נְקֻדִּ֖ים וּטְלֻאִֽים׃ \Rashi{\textbf{אל המקלות.}אל מראות המקלות: עקדים. משנים במקום עקידתם, הם קרסלי ידיהם ורגליהם:}%
% ---------- פרק 30 | פסוק 40 ----------
 \textbf{מ} וְהַכְּשָׂבִים֮ הִפְרִ֣יד יַעֲקֹב֒ וַ֠יִּתֵּ֠ן פְּנֵ֨י הַצֹּ֧אן אֶל־עָקֹ֛ד וְכׇל־ח֖וּם בְּצֹ֣אן לָבָ֑ן וַיָּֽשֶׁת־ל֤וֹ עֲדָרִים֙ לְבַדּ֔וֹ וְלֹ֥א שָׁתָ֖ם עַל־צֹ֥אן לָבָֽן׃ \Rashi{\textbf{והכשבים הפריד יעקב.}הנולדים עקודים נקדים הבדיל והפריש לעצמן ועשה אותן עדר עדר לבדו, והוליך אותו העדר העקוד לפני הצאן, ופני הצאן ההולכים אחריהם צופות אליהם; וזהו שאמר ויתן פני הצאן אל עקד, שהיו פני הצאן אל העקודים, ואל כל חום שמצא בצאן לבן: וישת לו עדרים. כמו שפרשתי:}%
% ---------- פרק 30 | פסוק 41 ----------
 \textbf{מא} וְהָיָ֗ה בְּכׇל־יַחֵם֮ הַצֹּ֣אן הַמְקֻשָּׁרוֹת֒ וְשָׂ֨ם יַעֲקֹ֧ב אֶת־הַמַּקְל֛וֹת לְעֵינֵ֥י הַצֹּ֖אן בָּרְהָטִ֑ים לְיַחְמֵ֖נָּה בַּמַּקְלֽוֹת׃ \Rashi{\textbf{המקשרות.}כתרגומו הבכירות, ואין לי עד במקרא; ומנחם חברו עם אחיתפל בקשרים (שמואל ב ט"ו), ויהי הקשר אמץ (שם), אותן המתקשרות יחד למהר עבורן:}%
% ---------- פרק 30 | פסוק 42 ----------
 \textbf{מב} וּבְהַעֲטִ֥יף הַצֹּ֖אן לֹ֣א יָשִׂ֑ים וְהָיָ֤ה הָעֲטֻפִים֙ לְלָבָ֔ן וְהַקְּשֻׁרִ֖ים לְיַעֲקֹֽב׃ \Rashi{\textbf{ובהעטיף.}לשון אחור, כתרגומו ובלקישות. ומנחם חברו עם המחלצות והמעטפות (ישעיהו ג'), לשון עטיפת כסות, כלומר מתעטפות בעורן וצמרן ואינן מתאוות להתיחם ע"י הזכרים:}%
% ---------- פרק 30 | פסוק 43 ----------
 \textbf{מג} וַיִּפְרֹ֥ץ הָאִ֖ישׁ מְאֹ֣ד מְאֹ֑ד וַֽיְהִי־לוֹ֙ צֹ֣אן רַבּ֔וֹת וּשְׁפָחוֹת֙ וַעֲבָדִ֔ים וּגְמַלִּ֖ים וַחֲמֹרִֽים׃ \Rashi{\textbf{צאן רבות.}פרות ורבות משאר צאן: ושפחות ועבדים. מוכר צאנו בדמים יקרים ולוקח לו כל אלה:}%
% ---------- פרק 31 | פסוק 1 ----------
 \textbf{א} וַיִּשְׁמַ֗ע אֶת־דִּבְרֵ֤י בְנֵֽי־לָבָן֙ לֵאמֹ֔ר לָקַ֣ח יַעֲקֹ֔ב אֵ֖ת כׇּל־אֲשֶׁ֣ר לְאָבִ֑ינוּ וּמֵאֲשֶׁ֣ר לְאָבִ֔ינוּ עָשָׂ֕ה אֵ֥ת כׇּל־הַכָּבֹ֖ד הַזֶּֽה׃ \Rashi{\textbf{עשה.}כנס, כמו ויעש חיל ויך את עמלק (שמואל א י"ד):}%
% ---------- פרק 31 | פסוק 2 ----------
 \textbf{ב} וַיַּ֥רְא יַעֲקֹ֖ב אֶת־פְּנֵ֣י לָבָ֑ן וְהִנֵּ֥ה אֵינֶ֛נּוּ עִמּ֖וֹ כִּתְמ֥וֹל שִׁלְשֽׁוֹם׃ %
% ---------- פרק 31 | פסוק 3 ----------
 \textbf{ג} וַיֹּ֤אמֶר יְהֹוָה֙ אֶֽל־יַעֲקֹ֔ב שׁ֛וּב אֶל־אֶ֥רֶץ אֲבוֹתֶ֖יךָ וּלְמוֹלַדְתֶּ֑ךָ וְאֶֽהְיֶ֖ה עִמָּֽךְ׃ \Rashi{\textbf{שוב אל ארץ אבותיך.}ושם אהיה עמך, אבל בעודך מחבר לטמא אי אפשר להשרות שכינתי עליך (בראשית רבה):}%
% ---------- פרק 31 | פסוק 4 ----------
 \textbf{ד} וַיִּשְׁלַ֣ח יַעֲקֹ֔ב וַיִּקְרָ֖א לְרָחֵ֣ל וּלְלֵאָ֑ה הַשָּׂדֶ֖ה אֶל־צֹאנֽוֹ׃ \Rashi{\textbf{ויקרא לרחל וללאה.}לרחל תחלה ואחר כך ללאה, שהיא היתה עקרת הבית שבשבילה נזדוג יעקב עם לבן, ואף בניה של לאה מודים בדבר, שהרי בעז ובית דינו משבט יהודה אומרים כרחל וכלאה אשר בנו שתיהם וגו' – הקדימו רחל ללאה:}%
% ---------- פרק 31 | פסוק 5 ----------
 \textbf{ה} וַיֹּ֣אמֶר לָהֶ֗ן רֹאֶ֤ה אָנֹכִי֙ אֶת־פְּנֵ֣י אֲבִיכֶ֔ן כִּֽי־אֵינֶ֥נּוּ אֵלַ֖י כִּתְמֹ֣ל שִׁלְשֹׁ֑ם וֵֽאלֹהֵ֣י אָבִ֔י הָיָ֖ה עִמָּדִֽי׃ %
% ---------- פרק 31 | פסוק 6 ----------
 \textbf{ו} וְאַתֵּ֖נָה יְדַעְתֶּ֑ן כִּ֚י בְּכׇל־כֹּחִ֔י עָבַ֖דְתִּי אֶת־אֲבִיכֶֽן׃ %
% ---------- פרק 31 | פסוק 7 ----------
 \textbf{ז} וַאֲבִיכֶן֙ הֵ֣תֶל בִּ֔י וְהֶחֱלִ֥ף אֶת־מַשְׂכֻּרְתִּ֖י עֲשֶׂ֣רֶת מֹנִ֑ים וְלֹֽא־נְתָנ֣וֹ אֱלֹהִ֔ים לְהָרַ֖ע עִמָּדִֽי׃ \Rashi{\textbf{עשרת מנים.}אין מונים פחות מעשרה: מנים. לשון סכום, כלל החשבון והן עשיריות; למדנו שהחליף תנאו ק' פעמים:}%
% ---------- פרק 31 | פסוק 8 ----------
 \textbf{ח} אִם־כֹּ֣ה יֹאמַ֗ר נְקֻדִּים֙ יִהְיֶ֣ה שְׂכָרֶ֔ךָ וְיָלְד֥וּ כׇל־הַצֹּ֖אן נְקֻדִּ֑ים וְאִם־כֹּ֣ה יֹאמַ֗ר עֲקֻדִּים֙ יִהְיֶ֣ה שְׂכָרֶ֔ךָ וְיָלְד֥וּ כׇל־הַצֹּ֖אן עֲקֻדִּֽים׃ %
% ---------- פרק 31 | פסוק 9 ----------
 \textbf{ט} וַיַּצֵּ֧ל אֱלֹהִ֛ים אֶת־מִקְנֵ֥ה אֲבִיכֶ֖ם וַיִּתֶּן־לִֽי׃ %
% ---------- פרק 31 | פסוק 10 ----------
 \textbf{י} וַיְהִ֗י בְּעֵת֙ יַחֵ֣ם הַצֹּ֔אן וָאֶשָּׂ֥א עֵינַ֛י וָאֵ֖רֶא בַּחֲל֑וֹם וְהִנֵּ֤ה הָֽעַתֻּדִים֙ הָעֹלִ֣ים עַל־הַצֹּ֔אן עֲקֻדִּ֥ים נְקֻדִּ֖ים וּבְרֻדִּֽים׃ \Rashi{\textbf{והנה העתדים.}אע"פ שהבדילם לבן כלם שלא יתעברו הצאן דגמתן, היו המלאכים מביאין אותן מעדר המסור ביד בני לבן לעדר שביד יעקב (בראשית רבה): וברדים. כתרגומו ופציחין, פייש"יר בלע"ז, חוט של לבן מקיף את גופו סביב, חברבורות שלו פתוחה ומפלשת מזו אל זו; ואין להביא עד מן המקרא:}%
% ---------- פרק 31 | פסוק 11 ----------
 \textbf{יא} וַיֹּ֨אמֶר אֵלַ֜י מַלְאַ֧ךְ הָאֱלֹהִ֛ים בַּחֲל֖וֹם יַֽעֲקֹ֑ב וָאֹמַ֖ר הִנֵּֽנִי׃ %
% ---------- פרק 31 | פסוק 12 ----------
 \textbf{יב} וַיֹּ֗אמֶר שָׂא־נָ֨א עֵינֶ֤יךָ וּרְאֵה֙ כׇּל־הָֽעַתֻּדִים֙ הָעֹלִ֣ים עַל־הַצֹּ֔אן עֲקֻדִּ֥ים נְקֻדִּ֖ים וּבְרֻדִּ֑ים כִּ֣י רָאִ֔יתִי אֵ֛ת כׇּל־אֲשֶׁ֥ר לָבָ֖ן עֹ֥שֶׂה לָּֽךְ׃ %
% ---------- פרק 31 | פסוק 13 ----------
 \textbf{יג} אָנֹכִ֤י הָאֵל֙ בֵּֽית־אֵ֔ל אֲשֶׁ֨ר מָשַׁ֤חְתָּ שָּׁם֙ מַצֵּבָ֔ה אֲשֶׁ֨ר נָדַ֥רְתָּ לִּ֛י שָׁ֖ם נֶ֑דֶר עַתָּ֗ה ק֥וּם צֵא֙ מִן־הָאָ֣רֶץ הַזֹּ֔את וְשׁ֖וּב אֶל־אֶ֥רֶץ מוֹלַדְתֶּֽךָ׃ \Rashi{\textbf{האל בית אל.}כמו אל בית אל, הה"א יתרה, ודרך מקראות לדבר כן, כמו כי אתם באים אל הארץ כנען (במדבר ל"ד): משחת שם. ל' רבוי וגדלה כשנמשח למלכות, כך ויצק שמן על ראשה – להיות משוחה למזבח: אשר נדרת לי. וצריך אתה לשלמו, שאמרת יהיה בית אלהים, שתקריב שם קרבנות:}%
% ---------- פרק 31 | פסוק 14 ----------
 \textbf{יד} וַתַּ֤עַן רָחֵל֙ וְלֵאָ֔ה וַתֹּאמַ֖רְנָה ל֑וֹ הַע֥וֹד לָ֛נוּ חֵ֥לֶק וְנַחֲלָ֖ה בְּבֵ֥ית אָבִֽינוּ׃ \Rashi{\textbf{העוד לנו.}למה נעכב על ידך מלשוב? כלום אנו מיחלות לירש מנכסי אבינו כלום בין הזכרים:}%
% ---------- פרק 31 | פסוק 15 ----------
 \textbf{טו} הֲל֧וֹא נׇכְרִיּ֛וֹת נֶחְשַׁ֥בְנוּ ל֖וֹ כִּ֣י מְכָרָ֑נוּ וַיֹּ֥אכַל גַּם־אָכ֖וֹל אֶת־כַּסְפֵּֽנוּ׃ \Rashi{\textbf{הלא נכריות נחשבנו לו.}אפלו בשעה שדרך בני אדם לתת נדוניה לבנותיו – בשעת נשואין – נהג עמנו כנכריות, כי מכרנו לך בשכר הפעלה: את כספנו. שעכב דמי שכר פעלתך:}%
% ---------- פרק 31 | פסוק 16 ----------
 \textbf{טז} כִּ֣י כׇל־הָעֹ֗שֶׁר אֲשֶׁ֨ר הִצִּ֤יל אֱלֹהִים֙ מֵֽאָבִ֔ינוּ לָ֥נוּ ה֖וּא וּלְבָנֵ֑ינוּ וְעַתָּ֗ה כֹּל֩ אֲשֶׁ֨ר אָמַ֧ר אֱלֹהִ֛ים אֵלֶ֖יךָ עֲשֵֽׂה׃ \Rashi{\textbf{כי כל העשר.}כי זה משמש בלשון אלא; כלומר, משל אבינו אין לנו כלום, אלא מה שהציל הקב"ה מאבינו שלנו הוא: הציל. לשון הפריש, וכן כל לשון הצלה שבמקרא לשון הפרשה, שמפרישו מן הרעה ומן האויב:}%
% ---------- פרק 31 | פסוק 17 ----------
 \textbf{יז} וַיָּ֖קׇם יַעֲקֹ֑ב וַיִּשָּׂ֛א אֶת־בָּנָ֥יו וְאֶת־נָשָׁ֖יו עַל־הַגְּמַלִּֽים׃ \Rashi{\textbf{את בניו ואת נשיו.}הקדים זכרים לנקבות, ועשו הקדים נקבות לזכרים, שנאמר ויקח עשו את נשיו ואת בניו וגומר:}%
% ---------- פרק 31 | פסוק 18 ----------
 \textbf{יח} וַיִּנְהַ֣ג אֶת־כׇּל־מִקְנֵ֗הוּ וְאֶת־כׇּל־רְכֻשׁוֹ֙ אֲשֶׁ֣ר רָכָ֔שׁ מִקְנֵה֙ קִנְיָנ֔וֹ אֲשֶׁ֥ר רָכַ֖שׁ בְּפַדַּ֣ן אֲרָ֑ם לָב֛וֹא אֶל־יִצְחָ֥ק אָבִ֖יו אַ֥רְצָה כְּנָֽעַן׃ \Rashi{\textbf{מקנה קנינו.}מה שקנה מצאנו, עבדים ושפחות וגמלים וחמורים:}%
% ---------- פרק 31 | פסוק 19 ----------
 \textbf{יט} וְלָבָ֣ן הָלַ֔ךְ לִגְזֹ֖ז אֶת־צֹאנ֑וֹ וַתִּגְנֹ֣ב רָחֵ֔ל אֶת־הַתְּרָפִ֖ים אֲשֶׁ֥ר לְאָבִֽיהָ׃ \Rashi{\textbf{לגזוז את צאנו.}שנתן ביד בניו דרך שלשת ימים בינו ובין יעקב: ותגנב רחל את התרפים. להפריש את אביה מע"ז נתכונה (בראשית רבה):}%
% ---------- פרק 31 | פסוק 20 ----------
 \textbf{כ} וַיִּגְנֹ֣ב יַעֲקֹ֔ב אֶת־לֵ֥ב לָבָ֖ן הָאֲרַמִּ֑י עַל־בְּלִי֙ הִגִּ֣יד ל֔וֹ כִּ֥י בֹרֵ֖חַ הֽוּא׃ %
% ---------- פרק 31 | פסוק 21 ----------
 \textbf{כא} וַיִּבְרַ֥ח הוּא֙ וְכׇל־אֲשֶׁר־ל֔וֹ וַיָּ֖קׇם וַיַּעֲבֹ֣ר אֶת־הַנָּהָ֑ר וַיָּ֥שֶׂם אֶת־פָּנָ֖יו הַ֥ר הַגִּלְעָֽד׃ %
% ---------- פרק 31 | פסוק 22 ----------
 \textbf{כב} וַיֻּגַּ֥ד לְלָבָ֖ן בַּיּ֣וֹם הַשְּׁלִישִׁ֑י כִּ֥י בָרַ֖ח יַעֲקֹֽב׃ \Rashi{\textbf{ביום השלישי.}שהרי דרך שלשת ימים היה ביניהם:}%
% ---------- פרק 31 | פסוק 23 ----------
 \textbf{כג} וַיִּקַּ֤ח אֶת־אֶחָיו֙ עִמּ֔וֹ וַיִּרְדֹּ֣ף אַחֲרָ֔יו דֶּ֖רֶךְ שִׁבְעַ֣ת יָמִ֑ים וַיַּדְבֵּ֥ק אֹת֖וֹ בְּהַ֥ר הַגִּלְעָֽד׃ \Rashi{\textbf{את אחיו.}קרוביו: דרך שבעת ימים. כל אותן ג' ימים שהלך המגיד להגיד ללבן הלך יעקב לדרכו, נמצא, יעקב רחוק מלבן ששה ימים, ובשביעי השיגו לבן. למדנו שכל מה שהלך יעקב בשבעה ימים הלך לבן ביום אחד (שנאמר וירדף אחריו דרך שבעת ימים, ולא נאמר וירדף אחריו שבעת ימים):}%
% ---------- פרק 31 | פסוק 24 ----------
 \textbf{כד} וַיָּבֹ֧א אֱלֹהִ֛ים אֶל־לָבָ֥ן הָאֲרַמִּ֖י בַּחֲלֹ֣ם הַלָּ֑יְלָה וַיֹּ֣אמֶר ל֗וֹ הִשָּׁ֧מֶר לְךָ֛ פֶּן־תְּדַבֵּ֥ר עִֽם־יַעֲקֹ֖ב מִטּ֥וֹב עַד־רָֽע׃ \Rashi{\textbf{מטוב עד רע.}כל טובתן של רשעים רעה היא אצל הצדיקים (יבמות ק"ג):}%
% ---------- פרק 31 | פסוק 25 ----------
 \textbf{כה} וַיַּשֵּׂ֥ג לָבָ֖ן אֶֽת־יַעֲקֹ֑ב וְיַעֲקֹ֗ב תָּקַ֤ע אֶֽת־אׇהֳלוֹ֙ בָּהָ֔ר וְלָבָ֛ן תָּקַ֥ע אֶת־אֶחָ֖יו בְּהַ֥ר הַגִּלְעָֽד׃ %
% ---------- פרק 31 | פסוק 26 ----------
 \textbf{כו} וַיֹּ֤אמֶר לָבָן֙ לְיַעֲקֹ֔ב מֶ֣ה עָשִׂ֔יתָ וַתִּגְנֹ֖ב אֶת־לְבָבִ֑י וַתְּנַהֵג֙ אֶת־בְּנֹתַ֔י כִּשְׁבֻי֖וֹת חָֽרֶב׃ \Rashi{\textbf{כשביות חרב.}כל חיל הבא למלחמה קרוי חרב:}%
% ---------- פרק 31 | פסוק 27 ----------
 \textbf{כז} לָ֤מָּה נַחְבֵּ֙אתָ֙ לִבְרֹ֔חַ וַתִּגְנֹ֖ב אֹתִ֑י וְלֹא־הִגַּ֣דְתָּ לִּ֔י וָֽאֲשַׁלֵּחֲךָ֛ בְּשִׂמְחָ֥ה וּבְשִׁרִ֖ים בְּתֹ֥ף וּבְכִנּֽוֹר׃ \Rashi{\textbf{ותגנב אתי.}גנבת את דעתי:}%
% ---------- פרק 31 | פסוק 28 ----------
 \textbf{כח} וְלֹ֣א נְטַשְׁתַּ֔נִי לְנַשֵּׁ֥ק לְבָנַ֖י וְלִבְנֹתָ֑י עַתָּ֖ה הִסְכַּ֥לְתָּֽ עֲשֽׂוֹ׃ %
% ---------- פרק 31 | פסוק 29 ----------
 \textbf{כט} יֶשׁ־לְאֵ֣ל יָדִ֔י לַעֲשׂ֥וֹת עִמָּכֶ֖ם רָ֑ע וֵֽאלֹהֵ֨י אֲבִיכֶ֜ם אֶ֣מֶשׁ ׀ אָמַ֧ר אֵלַ֣י לֵאמֹ֗ר הִשָּׁ֧מֶר לְךָ֛ מִדַּבֵּ֥ר עִֽם־יַעֲקֹ֖ב מִטּ֥וֹב עַד־רָֽע׃ \Rashi{\textbf{יש לאל ידי.}יש כח וחיל בידי לעשות עמכם רע, וכל אל שהוא ל' קדש על שם עזוז ורב אונים הוא:}%
% ---------- פרק 31 | פסוק 30 ----------
 \textbf{ל} וְעַתָּה֙ הָלֹ֣ךְ הָלַ֔כְתָּ כִּֽי־נִכְסֹ֥ף נִכְסַ֖פְתָּה לְבֵ֣ית אָבִ֑יךָ לָ֥מָּה גָנַ֖בְתָּ אֶת־אֱלֹהָֽי׃ \Rashi{\textbf{נכספתה.}חמדת, והרבה יש במקרא, נכספה וגם כלתה נפשי (תהילים פ"ד), למעשה ידיך תכסף (איוב י"ד):}%
% ---------- פרק 31 | פסוק 31 ----------
 \textbf{לא} וַיַּ֥עַן יַעֲקֹ֖ב וַיֹּ֣אמֶר לְלָבָ֑ן כִּ֣י יָרֵ֔אתִי כִּ֣י אָמַ֔רְתִּי פֶּן־תִּגְזֹ֥ל אֶת־בְּנוֹתֶ֖יךָ מֵעִמִּֽי׃ \Rashi{\textbf{כי יראתי וגו'.}השיבו על ראשון ראשון, שאמר לו, ותנהג את בנתי וגו':}%
% ---------- פרק 31 | פסוק 32 ----------
 \textbf{לב} עִ֠ם אֲשֶׁ֨ר תִּמְצָ֣א אֶת־אֱלֹהֶ֘יךָ֮ לֹ֣א יִֽחְיֶה֒ נֶ֣גֶד אַחֵ֧ינוּ הַֽכֶּר־לְךָ֛ מָ֥ה עִמָּדִ֖י וְקַֽח־לָ֑ךְ וְלֹֽא־יָדַ֣ע יַעֲקֹ֔ב כִּ֥י רָחֵ֖ל גְּנָבָֽתַם׃ \Rashi{\textbf{לא יחיה.}ומאותה קללה מתה רחל בדרך (בראשית רבה) מה עמדי. משלך:}%
% ---------- פרק 31 | פסוק 33 ----------
 \textbf{לג} וַיָּבֹ֨א לָבָ֜ן בְּאֹ֥הֶל יַעֲקֹ֣ב ׀ וּבְאֹ֣הֶל לֵאָ֗ה וּבְאֹ֛הֶל שְׁתֵּ֥י הָאֲמָהֹ֖ת וְלֹ֣א מָצָ֑א וַיֵּצֵא֙ מֵאֹ֣הֶל לֵאָ֔ה וַיָּבֹ֖א בְּאֹ֥הֶל רָחֵֽל׃ \Rashi{\textbf{באהל יעקב.}הוא אהל רחל, שהיה יעקב תדיר אצלה; וכן הוא אומר בני רחל אשת יעקב ובכלן לא נאמר אשת יעקב: ויבא באהל רחל. כשיצא מאהל לאה חזר לו לאהל רחל קדם שיחפש באהל האמהות, וכל כך למה? שהיה מכיר בה שהיא משמשנית:}%
% ---------- פרק 31 | פסוק 34 ----------
 \textbf{לד} וְרָחֵ֞ל לָקְחָ֣ה אֶת־הַתְּרָפִ֗ים וַתְּשִׂמֵ֛ם בְּכַ֥ר הַגָּמָ֖ל וַתֵּ֣שֶׁב עֲלֵיהֶ֑ם וַיְמַשֵּׁ֥שׁ לָבָ֛ן אֶת־כׇּל־הָאֹ֖הֶל וְלֹ֥א מָצָֽא׃ \Rashi{\textbf{בכר הגמל.}לשון כרים וכסתות, כתרגומו בעביטא דגמלא, והיא מרדעת העשויה כמין כר; ובערובין (דף ט"ז) שנינו הקיפוה בעביטין, והן עביטי גמלים, בשט"ו בלעז:}%
% ---------- פרק 31 | פסוק 35 ----------
 \textbf{לה} וַתֹּ֣אמֶר אֶל־אָבִ֗יהָ אַל־יִ֙חַר֙ בְּעֵינֵ֣י אֲדֹנִ֔י כִּ֣י ל֤וֹא אוּכַל֙ לָק֣וּם מִפָּנֶ֔יךָ כִּי־דֶ֥רֶךְ נָשִׁ֖ים לִ֑י וַיְחַפֵּ֕שׂ וְלֹ֥א מָצָ֖א אֶת־הַתְּרָפִֽים׃ %
% ---------- פרק 31 | פסוק 36 ----------
 \textbf{לו} וַיִּ֥חַר לְיַעֲקֹ֖ב וַיָּ֣רֶב בְּלָבָ֑ן וַיַּ֤עַן יַעֲקֹב֙ וַיֹּ֣אמֶר לְלָבָ֔ן מַה־פִּשְׁעִי֙ מַ֣ה חַטָּאתִ֔י כִּ֥י דָלַ֖קְתָּ אַחֲרָֽי׃ \Rashi{\textbf{דלקת.}רדפת, כמו על ההרים דלקנו (איכה ד'), וכמו מדלק אחרי פלשתים (שמואל א י"ז):}%
% ---------- פרק 31 | פסוק 37 ----------
 \textbf{לז} כִּֽי־מִשַּׁ֣שְׁתָּ אֶת־כׇּל־כֵּלַ֗י מַה־מָּצָ֙אתָ֙ מִכֹּ֣ל כְּלֵי־בֵיתֶ֔ךָ שִׂ֣ים כֹּ֔ה נֶ֥גֶד אַחַ֖י וְאַחֶ֑יךָ וְיוֹכִ֖יחוּ בֵּ֥ין שְׁנֵֽינוּ׃ \Rashi{\textbf{ויוכיחו.}ויבררו עם מי הדין, אפרובי"ר בלעז:}%
% ---------- פרק 31 | פסוק 38 ----------
 \textbf{לח} זֶה֩ עֶשְׂרִ֨ים שָׁנָ֤ה אָנֹכִי֙ עִמָּ֔ךְ רְחֵלֶ֥יךָ וְעִזֶּ֖יךָ לֹ֣א שִׁכֵּ֑לוּ וְאֵילֵ֥י צֹאנְךָ֖ לֹ֥א אָכָֽלְתִּי׃ \Rashi{\textbf{לא שכלו.}לא הפילו עבורם, כמו רחם משכיל (הושע ט'), תפלט פרתו ולא תשכל (איוב כ"א): ואילי צאנך. מכאן אמרו איל בן יומו קרוי איל, שאם לא כן, מה שבחו? אלים לא אכל, אבל כבשים אכל, אם כן גזלן הוא (בבא קמא ס"ה):}%
% ---------- פרק 31 | פסוק 39 ----------
 \textbf{לט} טְרֵפָה֙ לֹא־הֵבֵ֣אתִי אֵלֶ֔יךָ אָנֹכִ֣י אֲחַטֶּ֔נָּה מִיָּדִ֖י תְּבַקְשֶׁ֑נָּה גְּנֻֽבְתִ֣י י֔וֹם וּגְנֻֽבְתִ֖י לָֽיְלָה׃ \Rashi{\textbf{טרפה.}על ידי ארי וזאב: אנכי אחטנה. לשון קלע באבן אל השערה ולא יחטא (שופטים כ'), אני ובני שלמה חטאים (מלכים א א') – חסרים, אנכי אחסרנה, אם חסרה חסרה לי, שמידי תבקשנה: (אנכי אחטנה תרגומו דהות שגיא ממנינא, שהיתה נפקדת ומחסרת, כמו ולא נפקד ממנו איש (במדבר ל"א), תרגומו לא שגא): גנבתי יום וגנבתי לילה. גנובת יום או גנובת לילה, הכל שלמתי: גנבתי. כמו רבתי בגוים שרתי במדינות (איכה א'), מלאתי משפט (ישעיהו א') אהבתי לדוש (הושע י'):}%
% ---------- פרק 31 | פסוק 40 ----------
 \textbf{מ} הָיִ֧יתִי בַיּ֛וֹם אֲכָלַ֥נִי חֹ֖רֶב וְקֶ֣רַח בַּלָּ֑יְלָה וַתִּדַּ֥ד שְׁנָתִ֖י מֵֽעֵינָֽי׃ \Rashi{\textbf{אכלני חרב.}לשון אש אכלה: וקרח. כמו משליך קרחו (תהילים קמ"ז), תרגומו גלידא: שנתי. לשון שנה:}%
% ---------- פרק 31 | פסוק 41 ----------
 \textbf{מא} זֶה־לִּ֞י עֶשְׂרִ֣ים שָׁנָה֮ בְּבֵיתֶ֒ךָ֒ עֲבַדְתִּ֜יךָ אַרְבַּֽע־עֶשְׂרֵ֤ה שָׁנָה֙ בִּשְׁתֵּ֣י בְנֹתֶ֔יךָ וְשֵׁ֥שׁ שָׁנִ֖ים בְּצֹאנֶ֑ךָ וַתַּחֲלֵ֥ף אֶת־מַשְׂכֻּרְתִּ֖י עֲשֶׂ֥רֶת מֹנִֽים׃ \Rashi{\textbf{ותחלף את משכרתי.}היית משנה תנאי שבינינו, מנקד לטלוא, ומעקודים לברדים:}%
% ---------- פרק 31 | פסוק 42 ----------
 \textbf{מב} לוּלֵ֡י אֱלֹהֵ֣י אָבִי֩ אֱלֹהֵ֨י אַבְרָהָ֜ם וּפַ֤חַד יִצְחָק֙ הָ֣יָה לִ֔י כִּ֥י עַתָּ֖ה רֵיקָ֣ם שִׁלַּחְתָּ֑נִי אֶת־עׇנְיִ֞י וְאֶת־יְגִ֧יעַ כַּפַּ֛י רָאָ֥ה אֱלֹהִ֖ים וַיּ֥וֹכַח אָֽמֶשׁ׃ \Rashi{\textbf{ופחד יצחק.}לא רצה לומר אלהי יצחק, שאין הקב"ה מיחד שמו על הצדיקים בחייהם, ואע"פ שאמר לו בצאתו מבאר שבע אני ה' אלהי אברהם אביך ואלהי יצחק, בשביל שכהו עיניו, והרי הוא כמת; יעקב נתירא לומר אלהי, ואמר ופחד: ויוכח. לשון תוכחה הוא ולא לשון הוכחה:}%
% ---------- פרק 31 | פסוק 43 ----------
 \textbf{מג} וַיַּ֨עַן לָבָ֜ן וַיֹּ֣אמֶר אֶֽל־יַעֲקֹ֗ב הַבָּנ֨וֹת בְּנֹתַ֜י וְהַבָּנִ֤ים בָּנַי֙ וְהַצֹּ֣אן צֹאנִ֔י וְכֹ֛ל אֲשֶׁר־אַתָּ֥ה רֹאֶ֖ה לִי־ה֑וּא וְלִבְנֹתַ֞י מָֽה־אֶעֱשֶׂ֤ה לָאֵ֙לֶּה֙ הַיּ֔וֹם א֥וֹ לִבְנֵיהֶ֖ן אֲשֶׁ֥ר יָלָֽדוּ׃ \Rashi{\textbf{מה אעשה לאלה.}איך תעלה על לבי להרע להן:}%
% ---------- פרק 31 | פסוק 44 ----------
 \textbf{מד} וְעַתָּ֗ה לְכָ֛ה נִכְרְתָ֥ה בְרִ֖ית אֲנִ֣י וָאָ֑תָּה וְהָיָ֥ה לְעֵ֖ד בֵּינִ֥י וּבֵינֶֽךָ׃ \Rashi{\textbf{והיה לעד.}הקב"ה:}%
% ---------- פרק 31 | פסוק 45 ----------
 \textbf{מה} וַיִּקַּ֥ח יַעֲקֹ֖ב אָ֑בֶן וַיְרִימֶ֖הָ מַצֵּבָֽה׃ %
% ---------- פרק 31 | פסוק 46 ----------
 \textbf{מו} וַיֹּ֨אמֶר יַעֲקֹ֤ב לְאֶחָיו֙ לִקְט֣וּ אֲבָנִ֔ים וַיִּקְח֥וּ אֲבָנִ֖ים וַיַּֽעֲשׂוּ־גָ֑ל וַיֹּ֥אכְלוּ שָׁ֖ם עַל־הַגָּֽל׃ \Rashi{\textbf{לאחיו.}הם בניו שהיו לו אחים, נגשים אליו לצרה ולמלחמה (בראשית רבה):}%
% ---------- פרק 31 | פסוק 47 ----------
 \textbf{מז} וַיִּקְרָא־ל֣וֹ לָבָ֔ן יְגַ֖ר שָׂהֲדוּתָ֑א וְיַֽעֲקֹ֔ב קָ֥רָא ל֖וֹ גַּלְעֵֽד׃ \Rashi{\textbf{יגר שהדותא.}תרגומו של גלעד: גלעד. גל עד:}%
% ---------- פרק 31 | פסוק 48 ----------
 \textbf{מח} וַיֹּ֣אמֶר לָבָ֔ן הַגַּ֨ל הַזֶּ֥ה עֵ֛ד בֵּינִ֥י וּבֵינְךָ֖ הַיּ֑וֹם עַל־כֵּ֥ן קָרָֽא־שְׁמ֖וֹ גַּלְעֵֽד׃ %
% ---------- פרק 31 | פסוק 49 ----------
 \textbf{מט} וְהַמִּצְפָּה֙ אֲשֶׁ֣ר אָמַ֔ר יִ֥צֶף יְהֹוָ֖ה בֵּינִ֣י וּבֵינֶ֑ךָ כִּ֥י נִסָּתֵ֖ר אִ֥ישׁ מֵרֵעֵֽהוּ׃ \Rashi{\textbf{והמצפה אשר אמר וגו'.}והמצפה אשר בהר הגלעד, כמו שכתוב ויעבר את מצפה גלעד (שופטים י"א), ולמה נקראת שמה מצפה? לפי שאמר אחד מהם לחברו, יצף ה' ביני ובינך אם תעבר את הברית: כי נסתר. ולא נראה איש את רעהו:}%
% ---------- פרק 31 | פסוק 50 ----------
 \textbf{נ} אִם־תְּעַנֶּ֣ה אֶת־בְּנֹתַ֗י וְאִם־תִּקַּ֤ח נָשִׁים֙ עַל־בְּנֹתַ֔י אֵ֥ין אִ֖ישׁ עִמָּ֑נוּ רְאֵ֕ה אֱלֹהִ֥ים עֵ֖ד בֵּינִ֥י וּבֵינֶֽךָ׃ \Rashi{\textbf{בנתי, בנתי.}שתי פעמים, אף בלהה וזלפה בנותיו היו מפילגש: אם תענה את בנתי. למנע מהן עונת תשמיש (יומא ע"ז):}%
% ---------- פרק 31 | פסוק 51 ----------
 \textbf{נא} וַיֹּ֥אמֶר לָבָ֖ן לְיַעֲקֹ֑ב הִנֵּ֣ה ׀ הַגַּ֣ל הַזֶּ֗ה וְהִנֵּה֙ הַמַּצֵּבָ֔ה אֲשֶׁ֥ר יָרִ֖יתִי בֵּינִ֥י וּבֵינֶֽךָ׃ \Rashi{\textbf{יריתי.}כמו ירה בים, כזה שהוא יורה החץ:}%
% ---------- פרק 31 | פסוק 52 ----------
 \textbf{נב} עֵ֚ד הַגַּ֣ל הַזֶּ֔ה וְעֵדָ֖ה הַמַּצֵּבָ֑ה אִם־אָ֗נִי לֹֽא־אֶעֱבֹ֤ר אֵלֶ֙יךָ֙ אֶת־הַגַּ֣ל הַזֶּ֔ה וְאִם־אַ֠תָּ֠ה לֹא־תַעֲבֹ֨ר אֵלַ֜י אֶת־הַגַּ֥ל הַזֶּ֛ה וְאֶת־הַמַּצֵּבָ֥ה הַזֹּ֖את לְרָעָֽה׃ \Rashi{\textbf{אם אני.}הרי אם משמש בלשון אשר, כמו עד אם דברתי דברי: לרעה. לרעה אי אתה עובר, אבל אתה עובר לפרקמטיה:}%
% ---------- פרק 31 | פסוק 53 ----------
 \textbf{נג} אֱלֹהֵ֨י אַבְרָהָ֜ם וֵֽאלֹהֵ֤י נָחוֹר֙ יִשְׁפְּט֣וּ בֵינֵ֔ינוּ אֱלֹהֵ֖י אֲבִיהֶ֑ם וַיִּשָּׁבַ֣ע יַעֲקֹ֔ב בְּפַ֖חַד אָבִ֥יו יִצְחָֽק׃ \Rashi{\textbf{אלהי אברהם.}קדש (בראשית רבה): ואלהי נחור. חל: אלהי אביהם. חל:}%
% ---------- פרק 31 | פסוק 54 ----------
 \textbf{נד} וַיִּזְבַּ֨ח יַעֲקֹ֥ב זֶ֙בַח֙ בָּהָ֔ר וַיִּקְרָ֥א לְאֶחָ֖יו לֶאֱכׇל־לָ֑חֶם וַיֹּ֣אכְלוּ לֶ֔חֶם וַיָּלִ֖ינוּ בָּהָֽר׃ \Rashi{\textbf{ויזבח יעקב זבח.}שחט בהמות למשתה: לאחיו. לאוהביו שעם לבן: לאכל לחם. כל דבר מאכל קרוי לחם, כמו עבד לחם רב (דניאל ה'), נשחיתה עץ בלחמו (ירמיהו י"א):}%
% ---------- פרק 32 | פסוק 1 ----------
 \textbf{א} וַיַּשְׁכֵּ֨ם לָבָ֜ן בַּבֹּ֗קֶר וַיְנַשֵּׁ֧ק לְבָנָ֛יו וְלִבְנוֹתָ֖יו וַיְבָ֣רֶךְ אֶתְהֶ֑ם וַיֵּ֛לֶךְ וַיָּ֥שׇׁב לָבָ֖ן לִמְקֹמֽוֹ׃ %
% ---------- פרק 32 | פסוק 2 ----------
 \textbf{ב} וְיַעֲקֹ֖ב הָלַ֣ךְ לְדַרְכּ֑וֹ וַיִּפְגְּעוּ־ב֖וֹ מַלְאֲכֵ֥י אֱלֹהִֽים׃ \Rashi{\textbf{ויפגעו בו מלאכי אלהים.}מלאכים של ארץ ישראל באו לקראתו ללוותו לארץ:}%
% ---------- פרק 32 | פסוק 3 ----------
 \textbf{ג} וַיֹּ֤אמֶר יַעֲקֹב֙ כַּאֲשֶׁ֣ר רָאָ֔ם מַחֲנֵ֥ה אֱלֹהִ֖ים זֶ֑ה וַיִּקְרָ֛א שֵֽׁם־הַמָּק֥וֹם הַה֖וּא מַֽחֲנָֽיִם׃ \{פ\} \Rashi{\textbf{מחנים.}שתי מחנות, של חוצה לארץ שבאו עמו עד כאן ושל ארץ ישראל שבאו לקראתו (ע' תנחומא פ' וישלח):}%
\addparasha{חומש בראשית | פרשת וישלח}
% ---------- פרק 32 | פסוק 4 ----------
 \textbf{ד} וַיִּשְׁלַ֨ח יַעֲקֹ֤ב מַלְאָכִים֙ לְפָנָ֔יו אֶל־עֵשָׂ֖ו אָחִ֑יו אַ֥רְצָה שֵׂעִ֖יר שְׂדֵ֥ה אֱדֽוֹם׃ \Rashi{\textbf{וישלח יעקב מלאכים.}מלאכים ממש (בראשית רבה): ארצה שעיר. לארץ שעיר כל תבה שצריכה למ"ד בתחלתה הטיל לה הכתוב ה"א בסופה:}%
% ---------- פרק 32 | פסוק 5 ----------
 \textbf{ה} וַיְצַ֤ו אֹתָם֙ לֵאמֹ֔ר כֹּ֣ה תֹאמְר֔וּן לַֽאדֹנִ֖י לְעֵשָׂ֑ו כֹּ֤ה אָמַר֙ עַבְדְּךָ֣ יַעֲקֹ֔ב עִם־לָבָ֣ן גַּ֔רְתִּי וָאֵחַ֖ר עַד־עָֽתָּה׃ \Rashi{\textbf{גרתי.}לא נעשיתי שר וחשוב אלא גר, אינך כדאי לשנא אותי על ברכת אביך שברכני הוה גביר לאחיך, שהרי לא נתקימה בי. ד"א גרתי בגימטריא תרי"ג, כלומר, עם לבן גרתי ותרי"ג מצות שמרתי ולא למדתי ממעשיו הרעים:}%
% ---------- פרק 32 | פסוק 6 ----------
 \textbf{ו} וַֽיְהִי־לִי֙ שׁ֣וֹר וַחֲמ֔וֹר צֹ֖אן וְעֶ֣בֶד וְשִׁפְחָ֑ה וָֽאֶשְׁלְחָה֙ לְהַגִּ֣יד לַֽאדֹנִ֔י לִמְצֹא־חֵ֖ן בְּעֵינֶֽיךָ׃ \Rashi{\textbf{ויהי לי שור וחמור.}אבא אמר לי מטל השמים ומשמני הארץ, זו אינה לא מן השמים ולא מן הארץ: שור וחמור. דרך ארץ לומר על שורים הרבה שור – אדם אומר לחברו בלילה קרא התרנגול ואינו אומר קראו התרנגולים: ואשלחה להגיד לאדני. להודיע שאני בא אליך: למצא חן בעיניך. שאני שלם עמך ומבקש אהבתך:}%
% ---------- פרק 32 | פסוק 7 ----------
 \textbf{ז} וַיָּשֻׁ֙בוּ֙ הַמַּלְאָכִ֔ים אֶֽל־יַעֲקֹ֖ב לֵאמֹ֑ר בָּ֤אנוּ אֶל־אָחִ֙יךָ֙ אֶל־עֵשָׂ֔ו וְגַם֙ הֹלֵ֣ךְ לִקְרָֽאתְךָ֔ וְאַרְבַּע־מֵא֥וֹת אִ֖ישׁ עִמּֽוֹ׃ \Rashi{\textbf{באנו אל אחיך אל עשו.}שהיית אומר אחי הוא, אבל הוא נוהג עמך כעשו הרשע, עודנו בשנאתו (בראשית רבה פ'):}%
% ---------- פרק 32 | פסוק 8 ----------
 \textbf{ח} וַיִּירָ֧א יַעֲקֹ֛ב מְאֹ֖ד וַיֵּ֣צֶר ל֑וֹ וַיַּ֜חַץ אֶת־הָעָ֣ם אֲשֶׁר־אִתּ֗וֹ וְאֶת־הַצֹּ֧אן וְאֶת־הַבָּקָ֛ר וְהַגְּמַלִּ֖ים לִשְׁנֵ֥י מַחֲנֽוֹת׃ \Rashi{\textbf{ויירא ויצר.}ויירא שמא יהרג, ויצר לו אם יהרג הוא את אחרים (בראשית רבה ותנחומא):}%
% ---------- פרק 32 | פסוק 9 ----------
 \textbf{ט} וַיֹּ֕אמֶר אִם־יָב֥וֹא עֵשָׂ֛ו אֶל־הַמַּחֲנֶ֥ה הָאַחַ֖ת וְהִכָּ֑הוּ וְהָיָ֛ה הַמַּחֲנֶ֥ה הַנִּשְׁאָ֖ר לִפְלֵיטָֽה׃ \Rashi{\textbf{המחנה האחת והכהו.}מחנה משמש לשון זכר ולשון נקבה, אם תחנה עלי מחנה (תהלים כ"ז), הרי נקבה, המחנה הזה (בראשית ל"ג), זכר, וכן יש שאר דברים משמשים לשון זכר ולשון נקבה, השמש יצא על הארץ (בראשית י״ט:כ״ג), מקצה השמים מוצאו (תהלים י"ט), הרי לשון זכר; והשמש זרחה על המים (מלכים ב ג'), הרי לשון נקבה. וכן רוח, והנה רוח גדולה באה, הרי לשון נקבה, ויגע בארבע פנות הבית (איוב א'), הרי לשון זכר, ורוח גדולה וחזק מפרק הרים (מלכים א' י"ט), הרי לשון זכר ולשון נקבה. וכן אש, ואש יצאה מאת ה' (במדבר ט"ז), לשון נקבה, אש להט (תהילים ק"ד), לשון זכר: והיה המחנה הנשאר לפליטה. על כרחו, כי אלחם עמו. התקין עצמו לשלשה דברים, לדורון, לתפלה ולמלחמה. לדורון, ותעבר המנחה על פניו; לתפלה, אלהי אבי אברהם; למלחמה, והיה המחנה הנשאר לפליטה:}%
% ---------- פרק 32 | פסוק 10 ----------
 \textbf{י} וַיֹּ֘אמֶר֮ יַעֲקֹב֒ אֱלֹהֵי֙ אָבִ֣י אַבְרָהָ֔ם וֵאלֹהֵ֖י אָבִ֣י יִצְחָ֑ק יְהֹוָ֞ה הָאֹמֵ֣ר אֵלַ֗י שׁ֧וּב לְאַרְצְךָ֛ וּלְמוֹלַדְתְּךָ֖ וְאֵיטִ֥יבָה עִמָּֽךְ׃ \Rashi{\textbf{ואלהי אבי יצחק.}ולהלן הוא אומר ופחד יצחק? ועוד ומהו שחזר והזכיר שם המיחד? היה לו לכתב האומר אלי שוב לארצך וגו'. אלא כך אמר יעקב לפני הקב"ה, שתי הבטחות הבטחתני, אחת בצאתי מבית אבי מבאר שבע, שאמרת לי אני ה' אלהי אברהם אביך ואלהי יצחק, ושם אמרת לי ושמרתיך בכל אשר תלך, ובבית לבן אמרת לי (בראשית ל"א) שוב אל ארץ אבותיך ולמולדתך ואהיה עמך, ושם נגלית אלי בשם המיחד לבדו, שנאמר ויאמר ה' אל יעקב שוב אל ארץ אבותיך וגו', בשתי הבטחות אלו אני בא לפניך:}%
% ---------- פרק 32 | פסוק 11 ----------
 \textbf{יא} קָטֹ֜נְתִּי מִכֹּ֤ל הַחֲסָדִים֙ וּמִכׇּל־הָ֣אֱמֶ֔ת אֲשֶׁ֥ר עָשִׂ֖יתָ אֶת־עַבְדֶּ֑ךָ כִּ֣י בְמַקְלִ֗י עָבַ֙רְתִּי֙ אֶת־הַיַּרְדֵּ֣ן הַזֶּ֔ה וְעַתָּ֥ה הָיִ֖יתִי לִשְׁנֵ֥י מַחֲנֽוֹת׃ \Rashi{\textbf{קטנתי מכל החסדים.}נתמעטו זכיותי על ידי החסדים והאמת שעשית עמי, לכך אני ירא, שמא משהבטחתני נתקלקלתי בחטא ויגרם לי להמסר ביד עשו (שבת ל"ב): ומכל האמת. אמתת דבריך, ששמרת לי כל הבטחות שהבטחתני: כי במקלי. לא היה עמי לא כסף ולא זהב ולא מקנה אלא מקלי לבדו. ומ"א נתן מקלו בירדן ונבקע הירדן:}%
% ---------- פרק 32 | פסוק 12 ----------
 \textbf{יב} הַצִּילֵ֥נִי נָ֛א מִיַּ֥ד אָחִ֖י מִיַּ֣ד עֵשָׂ֑ו כִּֽי־יָרֵ֤א אָנֹכִי֙ אֹת֔וֹ פֶּן־יָב֣וֹא וְהִכַּ֔נִי אֵ֖ם עַל־בָּנִֽים׃ \Rashi{\textbf{מיד אחי מיד עשו.}מיד אחי שאין נוהג עמי כאח אלא כעשו הרשע:}%
% ---------- פרק 32 | פסוק 13 ----------
 \textbf{יג} וְאַתָּ֣ה אָמַ֔רְתָּ הֵיטֵ֥ב אֵיטִ֖יב עִמָּ֑ךְ וְשַׂמְתִּ֤י אֶֽת־זַרְעֲךָ֙ כְּח֣וֹל הַיָּ֔ם אֲשֶׁ֥ר לֹא־יִסָּפֵ֖ר מֵרֹֽב׃ \Rashi{\textbf{היטב איטיב.}היטב בזכותך, איטיב בזכות אבותיך (בראשית רבה): ושמתי את זרעך כחול הים. והיכן א"ל כן, והלא לא א"ל אלא והיה זרעך כעפר הארץ (שם כ"ח)? אלא שא"ל כי לא אעזבך עד אשר אם עשיתי את אשר דברתי לך (שם), ולאברהם אמר והרבה ארבה את זרעך ככוכבי השמים וכחול אשר על שפת הים:}%
% ---------- פרק 32 | פסוק 14 ----------
 \textbf{יד} וַיָּ֥לֶן שָׁ֖ם בַּלַּ֣יְלָה הַה֑וּא וַיִּקַּ֞ח מִן־הַבָּ֧א בְיָד֛וֹ מִנְחָ֖ה לְעֵשָׂ֥ו אָחִֽיו׃ \Rashi{\textbf{הבא בידו.}ברשותו, וכן ויקח את כל ארצו מידו (במדבר כ"א). ומ"א מן הבא בידו – אבנים טובות ומרגליות, שאדם צר בצרור ונושאם בידו (ד"א מן הבא בידו, מן החלין, שנטל מעשר, כמה דאת אמר עשר אעשרנו לך, והדר לקח מנחה):}%
% ---------- פרק 32 | פסוק 15 ----------
 \textbf{טו} עִזִּ֣ים מָאתַ֔יִם וּתְיָשִׁ֖ים עֶשְׂרִ֑ים רְחֵלִ֥ים מָאתַ֖יִם וְאֵילִ֥ים עֶשְׂרִֽים׃ \Rashi{\textbf{עזים מאתים ותישים עשרים.}מאתים עזים צריכות עשרים תישים, וכן כלם, הזכרים כדי צרך הנקבות; ובב"ר דורש מכאן לעונה האמורה בתורה, הטילים בכל יום, הפועלים שתים בשבת, החמרים אחת בשבת, הגמלים אחת לשלשים יום, הספנים אחת לששה חדשים; ואיני יודע לכון המדרש הזה בכוון, אך נראה בעיני שלמדנו מכאן שאין העונה שוה בכל אדם אלא לפי טרח המטל עליו, שמצינו כאן שמסר לכל תיש עשר עזים וכן לכל איל; לפי שהם פנויים ממלאכה דרכם להרבות תשמיש ולעבר עשר נקבות, ובהמה משנתעברה אינה מקבלת זכר, ופרים שעוסקין במלאכה, לא מסר לזכר אלא ארבע נקבות, ולחמור שהולך בדרך רחוקה שתי נקבות לזכר, ולגמלים שהולכים דרך יותר רחוקה נקבה אחת לזכר:}%
% ---------- פרק 32 | פסוק 16 ----------
 \textbf{טז} גְּמַלִּ֧ים מֵינִיק֛וֹת וּבְנֵיהֶ֖ם שְׁלֹשִׁ֑ים פָּר֤וֹת אַרְבָּעִים֙ וּפָרִ֣ים עֲשָׂרָ֔ה אֲתֹנֹ֣ת עֶשְׂרִ֔ים וַעְיָרִ֖ם עֲשָׂרָֽה׃ \Rashi{\textbf{גמלים מיניקות שלשים, ובניהם.}עמהם: ומ"א: ובניהם – בנאיהם, זכר כנגד נקבה, ולפי שצנוע בתשמיש לא פרסמו הכתוב (בראשית רבה): ועירם. חמורים זכרים:}%
% ---------- פרק 32 | פסוק 17 ----------
 \textbf{יז} וַיִּתֵּן֙ בְּיַד־עֲבָדָ֔יו עֵ֥דֶר עֵ֖דֶר לְבַדּ֑וֹ וַיֹּ֤אמֶר אֶל־עֲבָדָיו֙ עִבְר֣וּ לְפָנַ֔י וְרֶ֣וַח תָּשִׂ֔ימוּ בֵּ֥ין עֵ֖דֶר וּבֵ֥ין עֵֽדֶר׃ \Rashi{\textbf{עדר עדר לבדו.}כל מין ומין לעצמו: עברו לפני. דרך יום או פחות ואני אבא אחריכם: ורוח תשימו. עדר לפני חברו מלא עין, כדי להשביע עינו של אותו רשע ולתוהו על רבוי הדורון:}%
% ---------- פרק 32 | פסוק 18 ----------
 \textbf{יח} וַיְצַ֥ו אֶת־הָרִאשׁ֖וֹן לֵאמֹ֑ר כִּ֣י יִֽפְגׇשְׁךָ֞ עֵשָׂ֣ו אָחִ֗י וּשְׁאֵֽלְךָ֙ לֵאמֹ֔ר לְמִי־אַ֙תָּה֙ וְאָ֣נָה תֵלֵ֔ךְ וּלְמִ֖י אֵ֥לֶּה לְפָנֶֽיךָ׃ \Rashi{\textbf{למי אתה.}של מי אתה, מי שולחך, ותרגום דמאן את? ולמי אלה לפניך. ואלה שלפניך של מי הם; למי המנחה הזאת שלוחה? למ"ד משמשת בראש התבה במקום של, כמו: וכל אשר אתה ראה לי הוא (בראשית ל״א:מ״ג) – שלי הוא; לה' הארץ ומלואה (תהילים כ"ד) – של ה':}%
% ---------- פרק 32 | פסוק 19 ----------
 \textbf{יט} וְאָֽמַרְתָּ֙ לְעַבְדְּךָ֣ לְיַעֲקֹ֔ב מִנְחָ֥ה הִוא֙ שְׁלוּחָ֔ה לַֽאדֹנִ֖י לְעֵשָׂ֑ו וְהִנֵּ֥ה גַם־ה֖וּא אַחֲרֵֽינוּ׃ \Rashi{\textbf{ואמרת לעבדך ליעקב.}על ראשון ראשון ועל אחרון אחרון; ששאלת למי אתה? לעבדך ליעקב אני – ותרגומו דעבדך דיעקב – וששאלת ולמי אלה לפניך, מנחה היא שלוחה וגו': והנה גם הוא. יעקב:}%
% ---------- פרק 32 | פסוק 20 ----------
 \textbf{כ} וַיְצַ֞ו גַּ֣ם אֶת־הַשֵּׁנִ֗י גַּ֚ם אֶת־הַשְּׁלִישִׁ֔י גַּ֚ם אֶת־כׇּל־הַהֹ֣לְכִ֔ים אַחֲרֵ֥י הָעֲדָרִ֖ים לֵאמֹ֑ר כַּדָּבָ֤ר הַזֶּה֙ תְּדַבְּר֣וּן אֶל־עֵשָׂ֔ו בְּמֹצַאֲכֶ֖ם אֹתֽוֹ׃ %
% ---------- פרק 32 | פסוק 21 ----------
 \textbf{כא} וַאֲמַרְתֶּ֕ם גַּ֗ם הִנֵּ֛ה עַבְדְּךָ֥ יַעֲקֹ֖ב אַחֲרֵ֑ינוּ כִּֽי־אָמַ֞ר אֲכַפְּרָ֣ה פָנָ֗יו בַּמִּנְחָה֙ הַהֹלֶ֣כֶת לְפָנָ֔י וְאַחֲרֵי־כֵן֙ אֶרְאֶ֣ה פָנָ֔יו אוּלַ֖י יִשָּׂ֥א פָנָֽי׃ \Rashi{\textbf{אכפרה פניו.}אבטל רגזו; וכן וכפר בריתכם את מות (ישעיהו כ"ח), לא תוכלי כפרה (שם מ"ז). ונראה בעיני שכל כפרה שאצל עון וחטא ואצל פנים כלן לשון קנוח והעברה הן, ולשון ארמי הוא, והרבה בתלמוד וכפר ידיה, בעי לכפורי ידי בההוא גברא, וגם בלשון המקרא נקראים המזרקים של קדש כפורי זהב (עזרא א'), על שם שהכהן מקנח ידיו בהן בשפת המזרק:}%
% ---------- פרק 32 | פסוק 22 ----------
 \textbf{כב} וַתַּעֲבֹ֥ר הַמִּנְחָ֖ה עַל־פָּנָ֑יו וְה֛וּא לָ֥ן בַּלַּֽיְלָה־הַה֖וּא בַּֽמַּחֲנֶֽה׃ \Rashi{\textbf{על פניו.}כמו לפניו, וכן חמס ושד ישמע בה על פני תמיד (ירמיהו ו'), וכן המכעסים אתי על פני (ישעיהו ס"ה); ומ"א על פניו, אף הוא שרוי בכעס, שהיה צריך לכל זה (בראשית רבה):}%
% ---------- פרק 32 | פסוק 23 ----------
 \textbf{כג} וַיָּ֣קׇם ׀ בַּלַּ֣יְלָה ה֗וּא וַיִּקַּ֞ח אֶת־שְׁתֵּ֤י נָשָׁיו֙ וְאֶת־שְׁתֵּ֣י שִׁפְחֹתָ֔יו וְאֶת־אַחַ֥ד עָשָׂ֖ר יְלָדָ֑יו וַֽיַּעֲבֹ֔ר אֵ֖ת מַעֲבַ֥ר יַבֹּֽק׃ \Rashi{\textbf{ואת אחד עשר ילדיו.}ודינה היכן היתה? נתנה בתבה ונעל בפניה, שלא יתן בה עשו עיניו, ולכך נענש יעקב שמנעה מאחיו, שמא תחזירנו למוטב, ונפלה ביד שכם (בראשית רבה): יבק. שם הנהר:}%
% ---------- פרק 32 | פסוק 24 ----------
 \textbf{כד} וַיִּ֨קָּחֵ֔ם וַיַּֽעֲבִרֵ֖ם אֶת־הַנָּ֑חַל וַֽיַּעֲבֵ֖ר אֶת־אֲשֶׁר־לֽוֹ׃ \Rashi{\textbf{את אשר לו.}הבהמה והמטלטלים; עשה עצמו כגשר נוטל מכאן ומניח כאן (בראשית רבה):}%
% ---------- פרק 32 | פסוק 25 ----------
 \textbf{כה} וַיִּוָּתֵ֥ר יַעֲקֹ֖ב לְבַדּ֑וֹ וַיֵּאָבֵ֥ק אִישׁ֙ עִמּ֔וֹ עַ֖ד עֲל֥וֹת הַשָּֽׁחַר׃ \Rashi{\textbf{ויותר יעקב.}שכח פכים קטנים וחזר עליהם (חולין צ"א): ויאבק איש. מנחם פי' ויתעפר איש, לשון אבק, שהיו מעלים עפר ברגליהם ע"י נענועם. ולי נראה שהוא לשון ויתקשר, ולשון ארמי הוא, בתר דאביקו ביה, ואביק ליה מיבק – לשון עניבה, שכן דרך שנים שמתעצמים להפיל איש את רעהו, שחובקו ואובקו בזרועותיו. ופרשו רז"ל שהוא שרו של עשו (בראשית רבה):}%
% ---------- פרק 32 | פסוק 26 ----------
 \textbf{כו} וַיַּ֗רְא כִּ֣י לֹ֤א יָכֹל֙ ל֔וֹ וַיִּגַּ֖ע בְּכַף־יְרֵכ֑וֹ וַתֵּ֙קַע֙ כַּף־יֶ֣רֶךְ יַעֲקֹ֔ב בְּהֵאָֽבְק֖וֹ עִמּֽוֹ׃ \Rashi{\textbf{ויגע בכף ירכו.}קולית הירך התקוע בקלבוסית קרוי כף, ע"ש שהבשר שעליה כמין כף של קדרה: ותקע. נתקעקעה ממקום מחברתה, ודומה לו פן תקע נפשי ממך (ירמיהו ו'), לשון הסרה. ובמשנה לקעקע ביצתן – לשרש שרשיהן:}%
% ---------- פרק 32 | פסוק 27 ----------
 \textbf{כז} וַיֹּ֣אמֶר שַׁלְּחֵ֔נִי כִּ֥י עָלָ֖ה הַשָּׁ֑חַר וַיֹּ֙אמֶר֙ לֹ֣א אֲשַֽׁלֵּחֲךָ֔ כִּ֖י אִם־בֵּרַכְתָּֽנִי׃ \Rashi{\textbf{כי עלה השחר.}וצריך אני לומר שירה ביום (חולין צ"א): ברכתני. הודה לי על הברכות שברכני אבי, שעשו מערער עליהן:}%
% ---------- פרק 32 | פסוק 28 ----------
 \textbf{כח} וַיֹּ֥אמֶר אֵלָ֖יו מַה־שְּׁמֶ֑ךָ וַיֹּ֖אמֶר יַעֲקֹֽב׃ %
% ---------- פרק 32 | פסוק 29 ----------
 \textbf{כט} וַיֹּ֗אמֶר לֹ֤א יַעֲקֹב֙ יֵאָמֵ֥ר עוֹד֙ שִׁמְךָ֔ כִּ֖י אִם־יִשְׂרָאֵ֑ל כִּֽי־שָׂרִ֧יתָ עִם־אֱלֹהִ֛ים וְעִם־אֲנָשִׁ֖ים וַתּוּכָֽל׃ \Rashi{\textbf{לא יעקב.}לא יאמר עוד שהברכות באו לך בעקבה ורמיה כי אם בשררה וגלוי פנים, וסופך שהקב"ה נגלה אליך בבית אל ומחליף שמך ושם הוא מברכך, ואני שם אהיה ואודה לך עליהן, וזה שכתוב וישר אל מלאך ויכל בכה ויתחנן לו (הושע י"ב) – בכה המלאך ויתחנן לו, ומה נתחנן לו? בית אל ימצאנו ושם ידבר עמנו, המתן לי עד שידבר עמנו שם; ולא רצה יעקב, ועל כרחו הודה לו עליהן, וזהו ויברך אתו שם, שהיה מתחנן להמתין לו ולא רצה: ועם אנשים. עשו ולבן: ותוכל. להם:}%
% ---------- פרק 32 | פסוק 30 ----------
 \textbf{ל} וַיִּשְׁאַ֣ל יַעֲקֹ֗ב וַיֹּ֙אמֶר֙ הַגִּֽידָה־נָּ֣א שְׁמֶ֔ךָ וַיֹּ֕אמֶר לָ֥מָּה זֶּ֖ה תִּשְׁאַ֣ל לִשְׁמִ֑י וַיְבָ֥רֶךְ אֹת֖וֹ שָֽׁם׃ \Rashi{\textbf{למה זה תשאל.}אין לנו שם קבוע, משתנים שמותנו, הכל לפי מצות עבודת השליחות שאנו משתלחים (בראשית רבה):}%
% ---------- פרק 32 | פסוק 31 ----------
 \textbf{לא} וַיִּקְרָ֧א יַעֲקֹ֛ב שֵׁ֥ם הַמָּק֖וֹם פְּנִיאֵ֑ל כִּֽי־רָאִ֤יתִי אֱלֹהִים֙ פָּנִ֣ים אֶל־פָּנִ֔ים וַתִּנָּצֵ֖ל נַפְשִֽׁי׃ %
% ---------- פרק 32 | פסוק 32 ----------
 \textbf{לב} וַיִּֽזְרַֽח־ל֣וֹ הַשֶּׁ֔מֶשׁ כַּאֲשֶׁ֥ר עָבַ֖ר אֶת־פְּנוּאֵ֑ל וְה֥וּא צֹלֵ֖עַ עַל־יְרֵכֽוֹ׃ \Rashi{\textbf{ויזרח לו השמש.}לשון בני אדם הוא כשהגענו למקום פלוני האיר לנו השחר, זהו פשוטו; ומ"א ויזרח לו לצרכו – לרפאות את צלעתו, כמה דתימא שמש צדקה ומרפא בכנפיה (מלאכי ג'); ואותן שעות שמהרה לשקע בשבילו כשיצא מבאר שבע מהרה לזרח בשבילו: והוא צלע. היה צולע כשזרחה השמש:}%
% ---------- פרק 32 | פסוק 33 ----------
 \textbf{לג} עַל־כֵּ֡ן לֹֽא־יֹאכְל֨וּ בְנֵֽי־יִשְׂרָאֵ֜ל אֶת־גִּ֣יד הַנָּשֶׁ֗ה אֲשֶׁר֙ עַל־כַּ֣ף הַיָּרֵ֔ךְ עַ֖ד הַיּ֣וֹם הַזֶּ֑ה כִּ֤י נָגַע֙ בְּכַף־יֶ֣רֶךְ יַעֲקֹ֔ב בְּגִ֖יד הַנָּשֶֽׁה׃ \Rashi{\textbf{גיד הנשה.}למה נקרא שמו גיד הנשה? לפי שנשה ממקומו ועלה, והוא לשון קפיצה, וכן נשתה גבורתם (ירמיהו נ"א), וכן כי נשני אלהים את כל עמלי (בראשית מ"א):}%
% ---------- פרק 33 | פסוק 1 ----------
 \textbf{א} וַיִּשָּׂ֨א יַעֲקֹ֜ב עֵינָ֗יו וַיַּרְא֙ וְהִנֵּ֣ה עֵשָׂ֣ו בָּ֔א וְעִמּ֕וֹ אַרְבַּ֥ע מֵא֖וֹת אִ֑ישׁ וַיַּ֣חַץ אֶת־הַיְלָדִ֗ים עַל־לֵאָה֙ וְעַל־רָחֵ֔ל וְעַ֖ל שְׁתֵּ֥י הַשְּׁפָחֽוֹת׃ %
% ---------- פרק 33 | פסוק 2 ----------
 \textbf{ב} וַיָּ֧שֶׂם אֶת־הַשְּׁפָח֛וֹת וְאֶת־יַלְדֵיהֶ֖ן רִֽאשֹׁנָ֑ה וְאֶת־לֵאָ֤ה וִֽילָדֶ֙יהָ֙ אַחֲרֹנִ֔ים וְאֶת־רָחֵ֥ל וְאֶת־יוֹסֵ֖ף אַחֲרֹנִֽים׃ \Rashi{\textbf{ואת לאה וילדיה אחרנים.}אחרון אחרון חביב:}%
% ---------- פרק 33 | פסוק 3 ----------
 \textbf{ג} וְה֖וּא עָבַ֣ר לִפְנֵיהֶ֑ם וַיִּשְׁתַּ֤חוּ אַ֙רְצָה֙ שֶׁ֣בַע פְּעָמִ֔ים עַד־גִּשְׁתּ֖וֹ עַד־אָחִֽיו׃ \Rashi{\textbf{עבר לפניהם.}אמר אם יבא אותו רשע להלחם ילחם בי תחלה:}%
% ---------- פרק 33 | פסוק 4 ----------
 \textbf{ד} וַיָּ֨רׇץ עֵשָׂ֤ו לִקְרָאתוֹ֙ וַֽיְחַבְּקֵ֔הוּ וַיִּפֹּ֥ל עַל־צַוָּארָ֖ו וַׄיִּׄשָּׁׄקֵ֑ׄהׄוּׄ וַיִּבְכּֽוּ׃ \Rashi{\textbf{ויחבקהו.}נתגלגלו רחמיו, כשראהו משתחוה כל השתחואות הללו: וישקהו. נקוד עליו; ויש חולקין בדבר הזה בבריתא דספרי, יש שדרשו נקדה זו שלא נשקו בכל לבו, אמר רבי שמעון בן יוחאי, הלכה היא בידוע שעשו שונא ליעקב, אלא שנכמרו רחמיו באותה שעה ונשקו בכל לבו (ספרי במדבר):}%
% ---------- פרק 33 | פסוק 5 ----------
 \textbf{ה} וַיִּשָּׂ֣א אֶת־עֵינָ֗יו וַיַּ֤רְא אֶת־הַנָּשִׁים֙ וְאֶת־הַיְלָדִ֔ים וַיֹּ֖אמֶר מִי־אֵ֣לֶּה לָּ֑ךְ וַיֹּאמַ֕ר הַיְלָדִ֕ים אֲשֶׁר־חָנַ֥ן אֱלֹהִ֖ים אֶת־עַבְדֶּֽךָ׃ \Rashi{\textbf{מי אלה לך.}מי אלה להיות שלך:}%
% ---------- פרק 33 | פסוק 6 ----------
 \textbf{ו} וַתִּגַּ֧שְׁןָ הַשְּׁפָח֛וֹת הֵ֥נָּה וְיַלְדֵיהֶ֖ן וַתִּֽשְׁתַּחֲוֶֽיןָ׃ %
% ---------- פרק 33 | פסוק 7 ----------
 \textbf{ז} וַתִּגַּ֧שׁ גַּם־לֵאָ֛ה וִילָדֶ֖יהָ וַיִּֽשְׁתַּחֲו֑וּ וְאַחַ֗ר נִגַּ֥שׁ יוֹסֵ֛ף וְרָחֵ֖ל וַיִּֽשְׁתַּחֲוֽוּ׃ \Rashi{\textbf{נגש יוסף ורחל.}בכלן האמהות נגשות לפני הבנים, אבל ברחל יוסף נגש לפניה, אמר אמי יפת תאר, שמא יתלה בה עיניו אותו רשע, אעמד כנגדה ואעכבנו מלהסתכל בה; מכאן זכה יוסף לברכת עלי עין:}%
% ---------- פרק 33 | פסוק 8 ----------
 \textbf{ח} וַיֹּ֕אמֶר מִ֥י לְךָ֛ כׇּל־הַמַּחֲנֶ֥ה הַזֶּ֖ה אֲשֶׁ֣ר פָּגָ֑שְׁתִּי וַיֹּ֕אמֶר לִמְצֹא־חֵ֖ן בְּעֵינֵ֥י אֲדֹנִֽי׃ \Rashi{\textbf{מי לך כל המחנה.}מי כל המחנה אשר פגשתי, שהוא שלך? כלומר למה הוא לך? פשוטו של מקרא על מוליכי המנחה, ומדרשו כתות של מלאכים פגע, שהיו דוחפין אותו ואת אנשיו ואומרים להם של מי אתם? והם אומרים להם של עשו, והם אומרים הכו, הכו, ואלו אומרים הניחו, בנו של יצחק הוא, ולא היו משגיחים עליו; בן בנו של אברהם הוא, ולא היו משגיחין; אחיו של יעקב הוא, אומרים להם אם כן, משלנו אתם:}%
% ---------- פרק 33 | פסוק 9 ----------
 \textbf{ט} וַיֹּ֥אמֶר עֵשָׂ֖ו יֶשׁ־לִ֣י רָ֑ב אָחִ֕י יְהִ֥י לְךָ֖ אֲשֶׁר־לָֽךְ׃ \Rashi{\textbf{יהי לך אשר לך.}כאן הודה לו על הברכות (בראשית רבה):}%
% ---------- פרק 33 | פסוק 10 ----------
 \textbf{י} וַיֹּ֣אמֶר יַעֲקֹ֗ב אַל־נָא֙ אִם־נָ֨א מָצָ֤אתִי חֵן֙ בְּעֵינֶ֔יךָ וְלָקַחְתָּ֥ מִנְחָתִ֖י מִיָּדִ֑י כִּ֣י עַל־כֵּ֞ן רָאִ֣יתִי פָנֶ֗יךָ כִּרְאֹ֛ת פְּנֵ֥י אֱלֹהִ֖ים וַתִּרְצֵֽנִי׃ \Rashi{\textbf{אל נא.}אל נא תאמר לי כן: אם נא מצאתי חן בעיניך ולקחת מנחתי מידי כי על כן ראיתי פניך וגו'. כי כדאי והגון לך שתקבל מנחתי על אשר ראיתי פניך, והן חשובין לי כראית פני המלאך, שראיתי שר שלך, ועוד על שנתרצית למחל על סרחני. ולמה הזכיר לו ראית המלאך? כדי שיתירא הימנו ויאמר ראה מלאכים ונצל, איני יכול לו מעתה: ותרצני. נתפיסת לי, וכן כל רצון שבמקרא לשון פיוס, אפייצימנ"טו בלעז, כי לא לרצון יהיה לכם (ויקרא כ"ב) – הקרבנות באות לפיס ולרצות, וכן שפתי צדיק ידעון רצון (משלי י') – יודעים לפיס ולרצות:}%
% ---------- פרק 33 | פסוק 11 ----------
 \textbf{יא} קַח־נָ֤א אֶת־בִּרְכָתִי֙ אֲשֶׁ֣ר הֻבָ֣את לָ֔ךְ כִּֽי־חַנַּ֥נִי אֱלֹהִ֖ים וְכִ֣י יֶשׁ־לִי־כֹ֑ל וַיִּפְצַר־בּ֖וֹ וַיִּקָּֽח׃ \Rashi{\textbf{ברכתי.}מנחתי, מנחה זו הבאה על ראית פנים ולפרקים אינה באה אלא לשאלת שלום, וכל ברכה שהיא לראית פנים – כגון ויברך יעקב את פרעה (בראשית מ״ז:ז׳), עשו אתי ברכה דסנחריב (מלכים ב י"ח), וכן לשאל לו לשלום ולברכו דתעי מלך חמת (שמואל ב ח') – כלם לשון ברכת שלום הן, שקורין בלעז שלו"איר, אף זו ברכתי מו"ן שלו"ד: אשר הבאת לך. לא טרחת בה ואני יגעתי להגיעה עד שבאה לידך (בראשית רבה): חנני. נו"ן ראשונה מדגשת, לפי שהיא משמשת במקום שתי נוני"ן, שהיה לו לומר חננני, שאין חנן בלא שתי נוני"ן, והשלישית לשמוש, כמו עשני, זבדני: יש לי כל. כל ספוקי; ועשו דבר בלשון גאוה יש לי רב, יותר ויותר מכדי צרכי:}%
% ---------- פרק 33 | פסוק 12 ----------
 \textbf{יב} וַיֹּ֖אמֶר נִסְעָ֣ה וְנֵלֵ֑כָה וְאֵלְכָ֖ה לְנֶגְדֶּֽךָ׃ \Rashi{\textbf{נסעה.}כמו שמעה, סלחה, שהוא כמו שמע, סלח, אף כאן נסעה כמו נסע, והנו"ן יסוד בתבה, ותרגום של אנקלוס טול ונהך, עשו אמר ליעקב נסע מכאן ונלך: ואלכה לנגדך. בשוה לך; טובה זו אעשה לך, שאאריך ימי מהלכתי, ללכת לאט כאשר אתה צריך, וזהו לנגדך – בשוה לך:}%
% ---------- פרק 33 | פסוק 13 ----------
 \textbf{יג} וַיֹּ֣אמֶר אֵלָ֗יו אֲדֹנִ֤י יֹדֵ֙עַ֙ כִּֽי־הַיְלָדִ֣ים רַכִּ֔ים וְהַצֹּ֥אן וְהַבָּקָ֖ר עָל֣וֹת עָלָ֑י וּדְפָקוּם֙ י֣וֹם אֶחָ֔ד וָמֵ֖תוּ כׇּל־הַצֹּֽאן׃ \Rashi{\textbf{עלות עלי.}הצאן והבקר, שהן עלות, מטלות עלי לנהלן לאט: עלות. מגדלות עולליהן, לשון עולל ויונק (איכה ב'), עול ימים (ישעיהו ס"ה), ושתי פרות עלות (שמואל א ו') ובלעז אנפנטיי"ש: ודפקום יום אחד. ליגעם בדרך במרוצה ומתו כל הצאן: ודפקום. כמו קול דודי דופק (שיר השירים), נוקש בדלת:}%
% ---------- פרק 33 | פסוק 14 ----------
 \textbf{יד} יַעֲבׇר־נָ֥א אֲדֹנִ֖י לִפְנֵ֣י עַבְדּ֑וֹ וַאֲנִ֞י אֶֽתְנָהֲלָ֣ה לְאִטִּ֗י לְרֶ֨גֶל הַמְּלָאכָ֤ה אֲשֶׁר־לְפָנַי֙ וּלְרֶ֣גֶל הַיְלָדִ֔ים עַ֛ד אֲשֶׁר־אָבֹ֥א אֶל־אֲדֹנִ֖י שֵׂעִֽירָה׃ \Rashi{\textbf{יעבר נא אדני.}אל נא תאריך ימי הליכתך, עבר כפי דרכך ואף אם תתרחק: אתנהלה. אתנהל, ה"א יתירה, כמו ארדה, אשמעה: לאטי. לאט שלי, לשון נחת, ההולכים לאט, לאט לי לנער (שמואל ב י"ח). לאטי הלמ"ד מן היסוד, ואינה משמשת – אתנהל נחת שלי: לרגל המלאכה. לפי צרך הליכת רגלי המלאכה המטלת עלי להוליך: ולרגל הילדים. לפי רגליהם שהם יכולים לילך: עד אשר אבא אל אדני שעירה. הרחיב לו הדרך, שלא היה דעתו ללכת אלא עד סכות; אמר אם דעתו לעשות לי רעה, ימתין עד בואי אצלו, והוא לא הלך; ואימתי ילך? בימי המשיח (בראשית רבה), שנאמר ועלו מושעים בהר ציון לשפט את הר עשו (עובדיה א'); ומ"א יש לפרשה זו רבים:}%
% ---------- פרק 33 | פסוק 15 ----------
 \textbf{טו} וַיֹּ֣אמֶר עֵשָׂ֔ו אַצִּֽיגָה־נָּ֣א עִמְּךָ֔ מִן־הָעָ֖ם אֲשֶׁ֣ר אִתִּ֑י וַיֹּ֙אמֶר֙ לָ֣מָּה זֶּ֔ה אֶמְצָא־חֵ֖ן בְּעֵינֵ֥י אֲדֹנִֽי׃ \Rashi{\textbf{ויאמר למה זה.}תעשה לי טובה זו שאיני צריך לה: אמצא חן בעיני אדני. ולא תשלם לי עתה שום גמול:}%
% ---------- פרק 33 | פסוק 16 ----------
 \textbf{טז} וַיָּ֩שׇׁב֩ בַּיּ֨וֹם הַה֥וּא עֵשָׂ֛ו לְדַרְכּ֖וֹ שֵׂעִֽירָה׃ \Rashi{\textbf{וישב ביום ההוא עשו לדרכו.}עשו לבדו, וד' מאות איש שהלכו עמו נשמטו מאצלו אחד אחד, והיכן פרע להם הקב"ה? בימי דוד, שנאמר כי אם ארבע מאות איש נער אשר רכבו על הגמלים (שמואל א ל') (בראשית רבה):}%
% ---------- פרק 33 | פסוק 17 ----------
 \textbf{יז} וְיַעֲקֹב֙ נָסַ֣ע סֻכֹּ֔תָה וַיִּ֥בֶן ל֖וֹ בָּ֑יִת וּלְמִקְנֵ֙הוּ֙ עָשָׂ֣ה סֻכֹּ֔ת עַל־כֵּ֛ן קָרָ֥א שֵׁם־הַמָּק֖וֹם סֻכּֽוֹת׃ \{ס\} \Rashi{\textbf{ויבן לו בית.}שהה שם י"ח חדש, קיץ וחרף וקיץ (מגילה י"ז), סכות – קיץ; בית – חרף, סכות – קיץ:}%
% ---------- פרק 33 | פסוק 18 ----------
 \textbf{יח} וַיָּבֹא֩ יַעֲקֹ֨ב שָׁלֵ֜ם עִ֣יר שְׁכֶ֗ם אֲשֶׁר֙ בְּאֶ֣רֶץ כְּנַ֔עַן בְּבֹא֖וֹ מִפַּדַּ֣ן אֲרָ֑ם וַיִּ֖חַן אֶת־פְּנֵ֥י הָעִֽיר׃ \Rashi{\textbf{שלם.}שלם בגופו, שנתרפא מצלעתו; שלם בממונו, שלא חסר כלום מכל אותו דורון; שלם בתורתו, שלא שכח תלמודו בבית לבן (שבת ל"ג): עיר שכם. כמו לעיר, וכמוהו עד בואנה בית לחם (רות א'): בבאו מפדן ארם. כאדם האומר לחברו יצא פלוני מבין שני אריות ובא שלם, אף כאן ויבא שלם מפדן ארם, מלבן ומעשו שנזדוגו לו בדרך:}%
% ---------- פרק 33 | פסוק 19 ----------
 \textbf{יט} וַיִּ֜קֶן אֶת־חֶלְקַ֣ת הַשָּׂדֶ֗ה אֲשֶׁ֤ר נָֽטָה־שָׁם֙ אׇהֳל֔וֹ מִיַּ֥ד בְּנֵֽי־חֲמ֖וֹר אֲבִ֣י שְׁכֶ֑ם בְּמֵאָ֖ה קְשִׂיטָֽה׃ \Rashi{\textbf{קשיטה.}מעה. אר"ע כשהלכתי לכרכי הים היו קורין למעה קשיטה (ותרגומו חורפן – טובים, חריפים בכל מקום, כגון עובר לסוחר):}%
% ---------- פרק 33 | פסוק 20 ----------
 \textbf{כ} וַיַּצֶּב־שָׁ֖ם מִזְבֵּ֑חַ וַיִּ֨קְרָא־ל֔וֹ אֵ֖ל אֱלֹהֵ֥י יִשְׂרָאֵֽל׃ \{ס\} \Rashi{\textbf{ויקרא לו אל אלהי ישראל.}לא שהמזבח קרוי אלהי ישראל אלא על שם שהיה הקב"ה עמו והצילו קרא שם המזבח על שם הנס, להיות שבחו של מקום נזכר בקריאת השם, כלומר מי שהוא אל, הוא הקב"ה, הוא לאלהים לי, ששמי ישראל. וכן מצינו במשה ויקרא שמו ה' נסי (שמות י"ז), לא שהמזבח קרוי ה', אלא על שם הנס קרא שם המזבח להזכיר שבחו של הקדוש ברוך הוא, ה' הוא נסי. ורבותנו דרשו שהקב"ה קראו ליעקב אל (מגילה י"ח); ודברי תורה כפטיש יפצץ סלע (ע' שבת פ"ח) – מתחלקים לכמה טעמים, ואני לישב פשוטו של מקרא באתי:}%
% ---------- פרק 34 | פסוק 1 ----------
 \textbf{א} וַתֵּצֵ֤א דִינָה֙ בַּת־לֵאָ֔ה אֲשֶׁ֥ר יָלְדָ֖ה לְיַעֲקֹ֑ב לִרְא֖וֹת בִּבְנ֥וֹת הָאָֽרֶץ׃ \Rashi{\textbf{בת לאה.}ולא בת יעקב? אלא על שם יציאתה נקראת בת לאה, שאף היא יצאנית היתה (בראשית רבה), שנאמר ותצא לאה לקראתו (ועליה משלו המשל כאמה כבתה):}%
% ---------- פרק 34 | פסוק 2 ----------
 \textbf{ב} וַיַּ֨רְא אֹתָ֜הּ שְׁכֶ֧ם בֶּן־חֲמ֛וֹר הַֽחִוִּ֖י נְשִׂ֣יא הָאָ֑רֶץ וַיִּקַּ֥ח אֹתָ֛הּ וַיִּשְׁכַּ֥ב אֹתָ֖הּ וַיְעַנֶּֽהָ׃ \Rashi{\textbf{וישכב אתה.}כדרכה: ויענה. שלא כדרכה (בראשית רבה):}%
% ---------- פרק 34 | פסוק 3 ----------
 \textbf{ג} וַתִּדְבַּ֣ק נַפְשׁ֔וֹ בְּדִינָ֖ה בַּֽת־יַעֲקֹ֑ב וַיֶּֽאֱהַב֙ אֶת־הַֽנַּעֲרָ֔ וַיְדַבֵּ֖ר עַל־לֵ֥ב הַֽנַּעֲרָֽ׃ \Rashi{\textbf{על לב הנערה.}דברים המתישבים על הלב; ראי, אביך בחלקת שדה קטנה כמה ממון בזבז, אני אשיאך ותקנה העיר וכל שדותיה:}%
% ---------- פרק 34 | פסוק 4 ----------
 \textbf{ד} וַיֹּ֣אמֶר שְׁכֶ֔ם אֶל־חֲמ֥וֹר אָבִ֖יו לֵאמֹ֑ר קַֽח־לִ֛י אֶת־הַיַּלְדָּ֥ה הַזֹּ֖את לְאִשָּֽׁה׃ %
% ---------- פרק 34 | פסוק 5 ----------
 \textbf{ה} וְיַעֲקֹ֣ב שָׁמַ֗ע כִּ֤י טִמֵּא֙ אֶת־דִּינָ֣ה בִתּ֔וֹ וּבָנָ֛יו הָי֥וּ אֶת־מִקְנֵ֖הוּ בַּשָּׂדֶ֑ה וְהֶחֱרִ֥שׁ יַעֲקֹ֖ב עַד־בֹּאָֽם׃ %
% ---------- פרק 34 | פסוק 6 ----------
 \textbf{ו} וַיֵּצֵ֛א חֲמ֥וֹר אֲבִֽי־שְׁכֶ֖ם אֶֽל־יַעֲקֹ֑ב לְדַבֵּ֖ר אִתּֽוֹ׃ %
% ---------- פרק 34 | פסוק 7 ----------
 \textbf{ז} וּבְנֵ֨י יַעֲקֹ֜ב בָּ֤אוּ מִן־הַשָּׂדֶה֙ כְּשׇׁמְעָ֔ם וַיִּֽתְעַצְּבוּ֙ הָֽאֲנָשִׁ֔ים וַיִּ֥חַר לָהֶ֖ם מְאֹ֑ד כִּֽי־נְבָלָ֞ה עָשָׂ֣ה בְיִשְׂרָאֵ֗ל לִשְׁכַּב֙ אֶת־בַּֽת־יַעֲקֹ֔ב וְכֵ֖ן לֹ֥א יֵעָשֶֽׂה׃ \Rashi{\textbf{וכן לא יעשה.}לענות את הבתולות, שהאמות גדרו עצמן מן העריות על ידי המבול (בראשית רבה):}%
% ---------- פרק 34 | פסוק 8 ----------
 \textbf{ח} וַיְדַבֵּ֥ר חֲמ֖וֹר אִתָּ֣ם לֵאמֹ֑ר שְׁכֶ֣ם בְּנִ֗י חָֽשְׁקָ֤ה נַפְשׁוֹ֙ בְּבִתְּכֶ֔ם תְּנ֨וּ נָ֥א אֹתָ֛הּ ל֖וֹ לְאִשָּֽׁה׃ \Rashi{\textbf{חשקה.}חפצה:}%
% ---------- פרק 34 | פסוק 9 ----------
 \textbf{ט} וְהִֽתְחַתְּנ֖וּ אֹתָ֑נוּ בְּנֹֽתֵיכֶם֙ תִּתְּנוּ־לָ֔נוּ וְאֶת־בְּנֹתֵ֖ינוּ תִּקְח֥וּ לָכֶֽם׃ %
% ---------- פרק 34 | פסוק 10 ----------
 \textbf{י} וְאִתָּ֖נוּ תֵּשֵׁ֑בוּ וְהָאָ֙רֶץ֙ תִּהְיֶ֣ה לִפְנֵיכֶ֔ם שְׁבוּ֙ וּסְחָר֔וּהָ וְהֵֽאָחֲז֖וּ בָּֽהּ׃ %
% ---------- פרק 34 | פסוק 11 ----------
 \textbf{יא} וַיֹּ֤אמֶר שְׁכֶם֙ אֶל־אָבִ֣יהָ וְאֶל־אַחֶ֔יהָ אֶמְצָא־חֵ֖ן בְּעֵינֵיכֶ֑ם וַאֲשֶׁ֥ר תֹּאמְר֛וּ אֵלַ֖י אֶתֵּֽן׃ %
% ---------- פרק 34 | פסוק 12 ----------
 \textbf{יב} הַרְבּ֨וּ עָלַ֤י מְאֹד֙ מֹ֣הַר וּמַתָּ֔ן וְאֶ֨תְּנָ֔ה כַּאֲשֶׁ֥ר תֹּאמְר֖וּ אֵלָ֑י וּתְנוּ־לִ֥י אֶת־הַֽנַּעֲרָ֖ לְאִשָּֽׁה׃ \Rashi{\textbf{מהר.}כתבה:}%
% ---------- פרק 34 | פסוק 13 ----------
 \textbf{יג} וַיַּעֲנ֨וּ בְנֵֽי־יַעֲקֹ֜ב אֶת־שְׁכֶ֨ם וְאֶת־חֲמ֥וֹר אָבִ֛יו בְּמִרְמָ֖ה וַיְדַבֵּ֑רוּ אֲשֶׁ֣ר טִמֵּ֔א אֵ֖ת דִּינָ֥ה אֲחֹתָֽם׃ \Rashi{\textbf{במרמה.}בחכמה: אשר טמא. הכתוב אומר שלא היתה רמיה, שהרי טמא את דינה אחתם (בראשית רבה):}%
% ---------- פרק 34 | פסוק 14 ----------
 \textbf{יד} וַיֹּאמְר֣וּ אֲלֵיהֶ֗ם לֹ֤א נוּכַל֙ לַעֲשׂוֹת֙ הַדָּבָ֣ר הַזֶּ֔ה לָתֵת֙ אֶת־אֲחֹתֵ֔נוּ לְאִ֖ישׁ אֲשֶׁר־ל֣וֹ עׇרְלָ֑ה כִּֽי־חֶרְפָּ֥ה הִ֖וא לָֽנוּ׃ \Rashi{\textbf{חרפה הוא.}שמץ פסול הוא אצלנו; הבא לחרף חברו הוא אומר לו ערל אתה או בן ערל. חרפה בכל מקום גדוף:}%
% ---------- פרק 34 | פסוק 15 ----------
 \textbf{טו} אַךְ־בְּזֹ֖את נֵא֣וֹת לָכֶ֑ם אִ֚ם תִּהְי֣וּ כָמֹ֔נוּ לְהִמֹּ֥ל לָכֶ֖ם כׇּל־זָכָֽר׃ \Rashi{\textbf{נאות לכם.}נתרצה לכם, לשון ויאתו (מלכים ב י"ב): להמול. להיות נמול. אינו לשון לפעל אלא לשון להפעל:}%
% ---------- פרק 34 | פסוק 16 ----------
 \textbf{טז} וְנָתַ֤נּוּ אֶת־בְּנֹתֵ֙ינוּ֙ לָכֶ֔ם וְאֶת־בְּנֹתֵיכֶ֖ם נִֽקַּֽח־לָ֑נוּ וְיָשַׁ֣בְנוּ אִתְּכֶ֔ם וְהָיִ֖ינוּ לְעַ֥ם אֶחָֽד׃ \Rashi{\textbf{ונתנו.}נו"ן שניה מדגשת לפי שהיא משמשת במקום שתי נוני"ן, ונתננו: ואת בנתיכם נקח לנו. אתה מוצא בתנאי שאמר חמור ליעקב ובתשובת בני יעקב לחמור, שתלו החשיבות בבני יעקב לקח בנות שכם את שיבחרו להם, ובנותיהם יתנו להם לפי דעתם, דכתיב ונתנו את בנתינו – לפי דעתנו, ואת בנתיכם נקח לנו – ככל אשר נחפץ; וכשדברו חמור ושכם בנו אל יושבי עירם הפכו הדברים, את בנתם נקח לנו לנשים ואת בנתינו נתן להם, כדי לרצותם שיאותו להמול:}%
% ---------- פרק 34 | פסוק 17 ----------
 \textbf{יז} וְאִם־לֹ֧א תִשְׁמְע֛וּ אֵלֵ֖ינוּ לְהִמּ֑וֹל וְלָקַ֥חְנוּ אֶת־בִּתֵּ֖נוּ וְהָלָֽכְנוּ׃ %
% ---------- פרק 34 | פסוק 18 ----------
 \textbf{יח} וַיִּֽיטְב֥וּ דִבְרֵיהֶ֖ם בְּעֵינֵ֣י חֲמ֑וֹר וּבְעֵינֵ֖י שְׁכֶ֥ם בֶּן־חֲמֽוֹר׃ %
% ---------- פרק 34 | פסוק 19 ----------
 \textbf{יט} וְלֹֽא־אֵחַ֤ר הַנַּ֙עַר֙ לַעֲשׂ֣וֹת הַדָּבָ֔ר כִּ֥י חָפֵ֖ץ בְּבַֽת־יַעֲקֹ֑ב וְה֣וּא נִכְבָּ֔ד מִכֹּ֖ל בֵּ֥ית אָבִֽיו׃ %
% ---------- פרק 34 | פסוק 20 ----------
 \textbf{כ} וַיָּבֹ֥א חֲמ֛וֹר וּשְׁכֶ֥ם בְּנ֖וֹ אֶל־שַׁ֣עַר עִירָ֑ם וַֽיְדַבְּר֛וּ אֶל־אַנְשֵׁ֥י עִירָ֖ם לֵאמֹֽר׃ %
% ---------- פרק 34 | פסוק 21 ----------
 \textbf{כא} הָאֲנָשִׁ֨ים הָאֵ֜לֶּה שְֽׁלֵמִ֧ים הֵ֣ם אִתָּ֗נוּ וְיֵשְׁב֤וּ בָאָ֙רֶץ֙ וְיִסְחֲר֣וּ אֹתָ֔הּ וְהָאָ֛רֶץ הִנֵּ֥ה רַֽחֲבַת־יָדַ֖יִם לִפְנֵיהֶ֑ם אֶת־בְּנֹתָם֙ נִקַּֽח־לָ֣נוּ לְנָשִׁ֔ים וְאֶת־בְּנֹתֵ֖ינוּ נִתֵּ֥ן לָהֶֽם׃ \Rashi{\textbf{שלמים.}בשלום ובלב שלם: והארץ הנה רחבת ידים. כאדם שידו רחבה וותרנית; כלומר, אל תפסידו כלום – פרקמטיא הרבה באה לכאן ואין לה קונים:}%
% ---------- פרק 34 | פסוק 22 ----------
 \textbf{כב} אַךְ־בְּ֠זֹ֠את יֵאֹ֨תוּ לָ֤נוּ הָאֲנָשִׁים֙ לָשֶׁ֣בֶת אִתָּ֔נוּ לִהְי֖וֹת לְעַ֣ם אֶחָ֑ד בְּהִמּ֥וֹל לָ֙נוּ֙ כׇּל־זָכָ֔ר כַּאֲשֶׁ֖ר הֵ֥ם נִמֹּלִֽים׃ \Rashi{\textbf{בהמול.}בהיות נמול:}%
% ---------- פרק 34 | פסוק 23 ----------
 \textbf{כג} מִקְנֵהֶ֤ם וְקִנְיָנָם֙ וְכׇל־בְּהֶמְתָּ֔ם הֲל֥וֹא לָ֖נוּ הֵ֑ם אַ֚ךְ נֵא֣וֹתָה לָהֶ֔ם וְיֵשְׁב֖וּ אִתָּֽנוּ׃ \Rashi{\textbf{אך נאותה להם.}לדבר זה, ועל ידי כן ישבו אתנו:}%
% ---------- פרק 34 | פסוק 24 ----------
 \textbf{כד} וַיִּשְׁמְע֤וּ אֶל־חֲמוֹר֙ וְאֶל־שְׁכֶ֣ם בְּנ֔וֹ כׇּל־יֹצְאֵ֖י שַׁ֣עַר עִיר֑וֹ וַיִּמֹּ֙לוּ֙ כׇּל־זָכָ֔ר כׇּל־יֹצְאֵ֖י שַׁ֥עַר עִירֽוֹ׃ %
% ---------- פרק 34 | פסוק 25 ----------
 \textbf{כה} וַיְהִי֩ בַיּ֨וֹם הַשְּׁלִישִׁ֜י בִּֽהְיוֹתָ֣ם כֹּֽאֲבִ֗ים וַיִּקְח֣וּ שְׁנֵֽי־בְנֵי־יַ֠עֲקֹ֠ב שִׁמְע֨וֹן וְלֵוִ֜י אֲחֵ֤י דִינָה֙ אִ֣ישׁ חַרְבּ֔וֹ וַיָּבֹ֥אוּ עַל־הָעִ֖יר בֶּ֑טַח וַיַּֽהַרְג֖וּ כׇּל־זָכָֽר׃ \Rashi{\textbf{שני בני יעקב.}בניו היו, ואע"פ כן נהגו עצמן שמעון ולוי כשאר אנשים שאינם בניו, שלא נטלו עצה הימנו (בראשית רבה): אחי דינה. לפי שמסרו עצמן עליה נקראו אחיה: בטח. שהיו כואבים ומ"א בטוחים היו על כחו של זקן (בראשית רבה):}%
% ---------- פרק 34 | פסוק 26 ----------
 \textbf{כו} וְאֶת־חֲמוֹר֙ וְאֶת־שְׁכֶ֣ם בְּנ֔וֹ הָרְג֖וּ לְפִי־חָ֑רֶב וַיִּקְח֧וּ אֶת־דִּינָ֛ה מִבֵּ֥ית שְׁכֶ֖ם וַיֵּצֵֽאוּ׃ %
% ---------- פרק 34 | פסוק 27 ----------
 \textbf{כז} בְּנֵ֣י יַעֲקֹ֗ב בָּ֚אוּ עַל־הַ֣חֲלָלִ֔ים וַיָּבֹ֖זּוּ הָעִ֑יר אֲשֶׁ֥ר טִמְּא֖וּ אֲחוֹתָֽם׃ \Rashi{\textbf{על החללים.}לפשט החללים:}%
% ---------- פרק 34 | פסוק 28 ----------
 \textbf{כח} אֶת־צֹאנָ֥ם וְאֶת־בְּקָרָ֖ם וְאֶת־חֲמֹרֵיהֶ֑ם וְאֵ֧ת אֲשֶׁר־בָּעִ֛יר וְאֶת־אֲשֶׁ֥ר בַּשָּׂדֶ֖ה לָקָֽחוּ׃ %
% ---------- פרק 34 | פסוק 29 ----------
 \textbf{כט} וְאֶת־כׇּל־חֵילָ֤ם וְאֶת־כׇּל־טַפָּם֙ וְאֶת־נְשֵׁיהֶ֔ם שָׁב֖וּ וַיָּבֹ֑זּוּ וְאֵ֖ת כׇּל־אֲשֶׁ֥ר בַּבָּֽיִת׃ \Rashi{\textbf{חילם.}ממונם וכן עשה לי את החיל הזה (דברים ח'), וישראל עשה חיל (במדבר כ"ד), ועזבו לאחרים חילם (תהילים מ"ט): שבו. לשון שביה, לפיכך טעמו מלרע:}%
% ---------- פרק 34 | פסוק 30 ----------
 \textbf{ל} וַיֹּ֨אמֶר יַעֲקֹ֜ב אֶל־שִׁמְע֣וֹן וְאֶל־לֵוִי֮ עֲכַרְתֶּ֣ם אֹתִי֒ לְהַבְאִישֵׁ֙נִי֙ בְּיֹשֵׁ֣ב הָאָ֔רֶץ בַּֽכְּנַעֲנִ֖י וּבַפְּרִזִּ֑י וַאֲנִי֙ מְתֵ֣י מִסְפָּ֔ר וְנֶאֶסְפ֤וּ עָלַי֙ וְהִכּ֔וּנִי וְנִשְׁמַדְתִּ֖י אֲנִ֥י וּבֵיתִֽי׃ \Rashi{\textbf{עכרתם.}לשון מים עכורים, אין דעתי צלולה עכשו (ברכות כ"ה). ואגדה, צלולה היתה החבית ועכרתם אותה (בראשית רבה) – מסרת היתה ביד כנענים שיפלו ביד בני יעקב, אלא שהיו אומרים עד אשר תפרה ונחלת את הארץ (שמות כ"ג), לפיכך היו שותקין: מתי מספר. אנשים מועטים:}%
% ---------- פרק 34 | פסוק 31 ----------
 \textbf{לא} וַיֹּאמְר֑וּ הַכְזוֹנָ֕ה יַעֲשֶׂ֖ה אֶת־אֲחוֹתֵֽנוּ׃ \{פ\} \Rashi{\textbf{הכזונה.}הפקר: את אחותנו. ית אחתנא:}%
% ---------- פרק 35 | פסוק 1 ----------
 \textbf{א} וַיֹּ֤אמֶר אֱלֹהִים֙ אֶֽל־יַעֲקֹ֔ב ק֛וּם עֲלֵ֥ה בֵֽית־אֵ֖ל וְשֶׁב־שָׁ֑ם וַעֲשֵׂה־שָׁ֣ם מִזְבֵּ֔חַ לָאֵל֙ הַנִּרְאֶ֣ה אֵלֶ֔יךָ בְּבׇ֨רְחֲךָ֔ מִפְּנֵ֖י עֵשָׂ֥ו אָחִֽיךָ׃ \Rashi{\textbf{קום עלה.}לפי שאחרת בדרך נענשת ובא לך זאת מבתך (תנחומא):}%
% ---------- פרק 35 | פסוק 2 ----------
 \textbf{ב} וַיֹּ֤אמֶר יַעֲקֹב֙ אֶל־בֵּית֔וֹ וְאֶ֖ל כׇּל־אֲשֶׁ֣ר עִמּ֑וֹ הָסִ֜רוּ אֶת־אֱלֹהֵ֤י הַנֵּכָר֙ אֲשֶׁ֣ר בְּתֹכְכֶ֔ם וְהִֽטַּהֲר֔וּ וְהַחֲלִ֖יפוּ שִׂמְלֹתֵיכֶֽם׃ \Rashi{\textbf{הנכר.}שיש בידכם משלל של שכם: והטהרו. מע"ז: והחליפו שמלתיכם. שמא יש בידכם כסות של ע"ז:}%
% ---------- פרק 35 | פסוק 3 ----------
 \textbf{ג} וְנָק֥וּמָה וְנַעֲלֶ֖ה בֵּֽית־אֵ֑ל וְאֶֽעֱשֶׂה־שָּׁ֣ם מִזְבֵּ֗חַ לָאֵ֞ל הָעֹנֶ֤ה אֹתִי֙ בְּי֣וֹם צָֽרָתִ֔י וַֽיְהִי֙ עִמָּדִ֔י בַּדֶּ֖רֶךְ אֲשֶׁ֥ר הָלָֽכְתִּי׃ %
% ---------- פרק 35 | פסוק 4 ----------
 \textbf{ד} וַיִּתְּנ֣וּ אֶֽל־יַעֲקֹ֗ב אֵ֣ת כׇּל־אֱלֹהֵ֤י הַנֵּכָר֙ אֲשֶׁ֣ר בְּיָדָ֔ם וְאֶת־הַנְּזָמִ֖ים אֲשֶׁ֣ר בְּאׇזְנֵיהֶ֑ם וַיִּטְמֹ֤ן אֹתָם֙ יַעֲקֹ֔ב תַּ֥חַת הָאֵלָ֖ה אֲשֶׁ֥ר עִם־שְׁכֶֽם׃ \Rashi{\textbf{האלה.}מין אילן סרק: עם שכם. אצל שכם:}%
% ---------- פרק 35 | פסוק 5 ----------
 \textbf{ה} וַיִּסָּ֑עוּ וַיְהִ֣י ׀ חִתַּ֣ת אֱלֹהִ֗ים עַל־הֶֽעָרִים֙ אֲשֶׁר֙ סְבִיב֣וֹתֵיהֶ֔ם וְלֹ֣א רָֽדְפ֔וּ אַחֲרֵ֖י בְּנֵ֥י יַעֲקֹֽב׃ \Rashi{\textbf{חתת.}פחד:}%
% ---------- פרק 35 | פסוק 6 ----------
 \textbf{ו} וַיָּבֹ֨א יַעֲקֹ֜ב ל֗וּזָה אֲשֶׁר֙ בְּאֶ֣רֶץ כְּנַ֔עַן הִ֖וא בֵּֽית־אֵ֑ל ה֖וּא וְכׇל־הָעָ֥ם אֲשֶׁר־עִמּֽוֹ׃ %
% ---------- פרק 35 | פסוק 7 ----------
 \textbf{ז} וַיִּ֤בֶן שָׁם֙ מִזְבֵּ֔חַ וַיִּקְרָא֙ לַמָּק֔וֹם אֵ֖ל בֵּֽית־אֵ֑ל כִּ֣י שָׁ֗ם נִגְל֤וּ אֵלָיו֙ הָֽאֱלֹהִ֔ים בְּבׇרְח֖וֹ מִפְּנֵ֥י אָחִֽיו׃ \Rashi{\textbf{אל בית אל.}הקב"ה בבית אל, גלוי שכינתו בבית אל. יש תבה חסרה בי"ת המשמשת בראשה, כמו הנה הוא בית מכיר בן עמיאל (שמואל ב ט') – כמו בבית מכיר; בית אביך (בראשית ל״ח:י״א) – כמו בבית אביך: נגלו אליו האלהים. במקומות הרבה יש שם אלהות ואדנות בלשון רבים, כמו אדני יוסף (שם ל"ט), אם בעליו עמו (שמות כ'), ולא נאמר בעלו, וכן אלהות שהוא לשון שופט ומרות נזכר בלשון רבים, אבל אחד מכל שאר השמות לא תמצא בלשון רבים:}%
% ---------- פרק 35 | פסוק 8 ----------
 \textbf{ח} וַתָּ֤מׇת דְּבֹרָה֙ מֵינֶ֣קֶת רִבְקָ֔ה וַתִּקָּבֵ֛ר מִתַּ֥חַת לְבֵֽית־אֵ֖ל תַּ֣חַת הָֽאַלּ֑וֹן וַיִּקְרָ֥א שְׁמ֖וֹ אַלּ֥וֹן בָּכֽוּת׃ \{פ\} \Rashi{\textbf{ותמת דבורה.}מה ענין דבורה בבית יעקב? אלא לפי שאמרה רבקה ליעקב ושלחתי ולקחתיך משם, שלחה דבורה אצלו לפדן ארם לצאת משם, ומתה בדרך; מדברי רבי משה הדרשן למדתיה: מתחת לבית אל. העיר יושבת בהר ונקברה ברגלי ההר: תחת האלון. בשפולי מישרא, שהיה מישור מלמעלה בשפוע ההר והקבורה מלמטה, ומישור של בית אל היו קורין לו אלון. ואגדה, נתבשר שם באבל שני, שהגד לו על אמו שמתה (בראשית רבה), ואלון בלשון יוני אחר, ולפי שהעלימו את יום מותה שלא יקללו הבריות כרס שיצא ממנו עשו, אף הכתוב לא פרסמה:}%
% ---------- פרק 35 | פסוק 9 ----------
 \textbf{ט} וַיֵּרָ֨א אֱלֹהִ֤ים אֶֽל־יַעֲקֹב֙ ע֔וֹד בְּבֹא֖וֹ מִפַּדַּ֣ן אֲרָ֑ם וַיְבָ֖רֶךְ אֹתֽוֹ׃ \Rashi{\textbf{עוד.}פעם שני במקום הזה, אחד בלכתו ואחד בשובו: ויברך אתו. ברכת אבלים (בראשית רבה):}%
% ---------- פרק 35 | פסוק 10 ----------
 \textbf{י} וַיֹּֽאמֶר־ל֥וֹ אֱלֹהִ֖ים שִׁמְךָ֣ יַעֲקֹ֑ב לֹֽא־יִקָּרֵא֩ שִׁמְךָ֨ ע֜וֹד יַעֲקֹ֗ב כִּ֤י אִם־יִשְׂרָאֵל֙ יִהְיֶ֣ה שְׁמֶ֔ךָ וַיִּקְרָ֥א אֶת־שְׁמ֖וֹ יִשְׂרָאֵֽל׃ \Rashi{\textbf{לא יקרא שמך עוד יעקב.}לשון אדם הבא במארב ועקבה, אלא לשון שר ונגיד:}%
% ---------- פרק 35 | פסוק 11 ----------
 \textbf{יא} וַיֹּ֩אמֶר֩ ל֨וֹ אֱלֹהִ֜ים אֲנִ֨י אֵ֤ל שַׁדַּי֙ פְּרֵ֣ה וּרְבֵ֔ה גּ֛וֹי וּקְהַ֥ל גּוֹיִ֖ם יִהְיֶ֣ה מִמֶּ֑ךָּ וּמְלָכִ֖ים מֵחֲלָצֶ֥יךָ יֵצֵֽאוּ׃ \Rashi{\textbf{אני אל שדי.}שאני כדי לברך, שהברכות שלי: פרה ורבה. על שם שעדין לא נולד בנימין, ואף על פי שכבר נתעברה ממנו: גוי. בנימין: גוים. מנשה ואפרים, שעתידים לצאת מיוסף, והם במנין השבטים (בראשית רבה): ומלכים. שאול ואיש בשת, שהיו משבט בנימין (בראשית רבה), שעדין לא נולדו (ופסוק זה דרשו אבנר כשהמליך איש בשת, ואף השבטים דרשוהו וקרבו בנימין, דכתיב איש ממנו לא יתן את בתו לבנימן לאשה (שופטים כ"א), וחזרו ואמרו אלמלא היה עולה מן השבטים לא היה הקב"ה אומר ליעקב ומלכים מחלציך יצאו: גוי וקהל גוים. שגוים עתידים בניו להעשות כמנין הגוים, שהם ע' אמות, וכן כל הסנהדרין שבעים; דבר אחר שעתידים בניו להקריב בשעת אסור במות כגוים בימי אליהו):}%
% ---------- פרק 35 | פסוק 12 ----------
 \textbf{יב} וְאֶת־הָאָ֗רֶץ אֲשֶׁ֥ר נָתַ֛תִּי לְאַבְרָהָ֥ם וּלְיִצְחָ֖ק לְךָ֣ אֶתְּנֶ֑נָּה וּֽלְזַרְעֲךָ֥ אַחֲרֶ֖יךָ אֶתֵּ֥ן אֶת־הָאָֽרֶץ׃ %
% ---------- פרק 35 | פסוק 13 ----------
 \textbf{יג} וַיַּ֥עַל מֵעָלָ֖יו אֱלֹהִ֑ים בַּמָּק֖וֹם אֲשֶׁר־דִּבֶּ֥ר אִתּֽוֹ׃ \Rashi{\textbf{במקום אשר דבר אתו.}איני יודע מה מלמדנו:}%
% ---------- פרק 35 | פסוק 14 ----------
 \textbf{יד} וַיַּצֵּ֨ב יַעֲקֹ֜ב מַצֵּבָ֗ה בַּמָּק֛וֹם אֲשֶׁר־דִּבֶּ֥ר אִתּ֖וֹ מַצֶּ֣בֶת אָ֑בֶן וַיַּסֵּ֤ךְ עָלֶ֙יהָ֙ נֶ֔סֶךְ וַיִּצֹ֥ק עָלֶ֖יהָ שָֽׁמֶן׃ %
% ---------- פרק 35 | פסוק 15 ----------
 \textbf{טו} וַיִּקְרָ֨א יַעֲקֹ֜ב אֶת־שֵׁ֣ם הַמָּק֗וֹם אֲשֶׁר֩ דִּבֶּ֨ר אִתּ֥וֹ שָׁ֛ם אֱלֹהִ֖ים בֵּֽית־אֵֽל׃ %
% ---------- פרק 35 | פסוק 16 ----------
 \textbf{טז} וַיִּסְעוּ֙ מִבֵּ֣ית אֵ֔ל וַֽיְהִי־ע֥וֹד כִּבְרַת־הָאָ֖רֶץ לָב֣וֹא אֶפְרָ֑תָה וַתֵּ֥לֶד רָחֵ֖ל וַתְּקַ֥שׁ בְּלִדְתָּֽהּ׃ \Rashi{\textbf{כברת הארץ.}מנחם פרש לשון כביר, רבוי, מהלך רב. ואגדה, בזמן שהארץ חלולה ומנקבת ככברה (בראשית רבה), שהניר מצוי, הסתו עבר והשרב עדין לא בא (שם ז), ואין זה פשוטו של מקרא, שהרי בנעמן מצינו וילך מאתו כברת ארץ (מלכים ב ה'). ואומר אני שהוא שם מדת קרקע, כמו מהלך פרסה או יותר, כמו שאתה אומר צמד כרם (ישעיהו ה'), חלקת השדה (בראשית ל"ג), כך במהלך אדם נותן שם מדה כברת ארץ:}%
% ---------- פרק 35 | פסוק 17 ----------
 \textbf{יז} וַיְהִ֥י בְהַקְשֹׁתָ֖הּ בְּלִדְתָּ֑הּ וַתֹּ֨אמֶר לָ֤הּ הַמְיַלֶּ֙דֶת֙ אַל־תִּ֣ירְאִ֔י כִּֽי־גַם־זֶ֥ה לָ֖ךְ בֵּֽן׃ \Rashi{\textbf{כי גם זה.}נוסף לך על יוסף; ורבותינו דרשו: עם כל שבט נולדה תאומה, ועם בנימין נולדה תאומה יתרה:}%
% ---------- פרק 35 | פסוק 18 ----------
 \textbf{יח} וַיְהִ֞י בְּצֵ֤את נַפְשָׁהּ֙ כִּ֣י מֵ֔תָה וַתִּקְרָ֥א שְׁמ֖וֹ בֶּן־אוֹנִ֑י וְאָבִ֖יו קָֽרָא־ל֥וֹ בִנְיָמִֽין׃ \Rashi{\textbf{בן אוני.}בן צערי: בנימין. נראה בעיני, לפי שהוא לבדו נולד בארץ כנען שהוא בנגב כשאדם בא מארם נהרים, כמו שנאמר בנגב בארץ כנען (במדבר ל"ג), הלוך ונסוע הנגבה (בראשית י"ב):  בנימין. בן ימין, לשון צפון וימין אתה בראתם (תהילים פ"ט), לפיכך הוא מלא:}%
% ---------- פרק 35 | פסוק 19 ----------
 \textbf{יט} וַתָּ֖מׇת רָחֵ֑ל וַתִּקָּבֵר֙ בְּדֶ֣רֶךְ אֶפְרָ֔תָה הִ֖וא בֵּ֥ית לָֽחֶם׃ %
% ---------- פרק 35 | פסוק 20 ----------
 \textbf{כ} וַיַּצֵּ֧ב יַעֲקֹ֛ב מַצֵּבָ֖ה עַל־קְבֻרָתָ֑הּ הִ֛וא מַצֶּ֥בֶת קְבֻֽרַת־רָחֵ֖ל עַד־הַיּֽוֹם׃ %
% ---------- פרק 35 | פסוק 21 ----------
 \textbf{כא} וַיִּסַּ֖ע יִשְׂרָאֵ֑ל וַיֵּ֣ט אׇֽהֳלֹ֔ה מֵהָ֖לְאָה לְמִגְדַּל־עֵֽדֶר׃ %
% ---------- פרק 35 | פסוק 22 ----------
 \textbf{כב} וַיְהִ֗י בִּשְׁכֹּ֤ן יִשְׂרָאֵל֙ בָּאָ֣רֶץ הַהִ֔וא וַיֵּ֣לֶךְ רְאוּבֵ֗֔ן וַיִּשְׁכַּ֕ב֙ אֶת־בִּלְהָ֖ה֙ פִּילֶ֣גֶשׁ אָבִ֑֔יו וַיִּשְׁמַ֖ע יִשְׂרָאֵ֑͏ֽל׃ \{פ\}וַיִּֽהְי֥וּ בְנֵֽי־יַעֲקֹ֖ב שְׁנֵ֥ים עָשָֽׂר׃ \Rashi{\textbf{בשכן ישראל בארץ ההיא.}עד שלא בא לחברון אצל יצחק ארעוהו כל אלה: וישכב. מתוך שבלבל משכבו, מעלה עליו הכתוב כאלו שכבה; ולמה בלבל וחלל יצועיו? שכשמתה רחל נטל יעקב מטתו, שהיתה נתונה תדיר באהל רחל ולא בשאר אהלים, ונתנה באהל בלהה; בא ראובן ותבע עלבון אמו, אמר אם אחות אמי היתה צרה לאמי, שפחת אחות אמי תהא צרה לאמי? לכך בלבל (שבת נ"ה): ויהיו בני יעקב שנים עשר. מתחיל לענין ראשון, משנולד בנימין נשלמה המטה ומעתה ראויים להמנות, ומנאן; ורבותינו דרשו, ללמדנו בא, שכלן שוין וכלן צדיקים, שלא חטא ראובן:}%
% ---------- פרק 35 | פסוק 23 ----------
 \textbf{כג} בְּנֵ֣י לֵאָ֔ה בְּכ֥וֹר יַעֲקֹ֖ב רְאוּבֵ֑ן וְשִׁמְעוֹן֙ וְלֵוִ֣י וִֽיהוּדָ֔ה וְיִשָּׂשכָ֖ר וּזְבֻלֽוּן׃ \Rashi{\textbf{בכור יעקב.}אפלו בשעת הקלקלה קראו בכור: בכור יעקב. בכור לנחלה, בכור לעבודה, בכור למנין, ולא נתנה בכורה ליוסף אלא לענין השבטים – שנעשה לשני שבטים:}%
% ---------- פרק 35 | פסוק 24 ----------
 \textbf{כד} בְּנֵ֣י רָחֵ֔ל יוֹסֵ֖ף וּבִנְיָמִֽן׃ %
% ---------- פרק 35 | פסוק 25 ----------
 \textbf{כה} וּבְנֵ֤י בִלְהָה֙ שִׁפְחַ֣ת רָחֵ֔ל דָּ֖ן וְנַפְתָּלִֽי׃ %
% ---------- פרק 35 | פסוק 26 ----------
 \textbf{כו} וּבְנֵ֥י זִלְפָּ֛ה שִׁפְחַ֥ת לֵאָ֖ה גָּ֣ד וְאָשֵׁ֑ר אֵ֚לֶּה בְּנֵ֣י יַעֲקֹ֔ב אֲשֶׁ֥ר יֻלַּד־ל֖וֹ בְּפַדַּ֥ן אֲרָֽם׃ %
% ---------- פרק 35 | פסוק 27 ----------
 \textbf{כז} וַיָּבֹ֤א יַעֲקֹב֙ אֶל־יִצְחָ֣ק אָבִ֔יו מַמְרֵ֖א קִרְיַ֣ת הָֽאַרְבַּ֑ע הִ֣וא חֶבְר֔וֹן אֲשֶׁר־גָּֽר־שָׁ֥ם אַבְרָהָ֖ם וְיִצְחָֽק׃ \Rashi{\textbf{ממרא.}שם המישור: קרית הארבע. שם העיר: ממרא קרית הארבע. אל מישור של קרית ארבע. וא"ת היה לו לכתוב ממרא הקרית ארבע? כן דרך המקרא בכל דבר ששמו כפול, כגון זה, וכגון בית לחם, אבי עזר, בית אל, אם הצרך להטיל בו ה"א, נותנה בראש התבה השניה – בית הלחמי (שמואל א ט"ז), בעפרת אבי העזרי (שופטים ו'), בנה חיאל בית האלי (מלכים א ט"ז):}%
% ---------- פרק 35 | פסוק 28 ----------
 \textbf{כח} וַיִּֽהְי֖וּ יְמֵ֣י יִצְחָ֑ק מְאַ֥ת שָׁנָ֖ה וּשְׁמֹנִ֥ים שָׁנָֽה׃ %
% ---------- פרק 35 | פסוק 29 ----------
 \textbf{כט} וַיִּגְוַ֨ע יִצְחָ֤ק וַיָּ֙מׇת֙ וַיֵּאָ֣סֶף אֶל־עַמָּ֔יו זָקֵ֖ן וּשְׂבַ֣ע יָמִ֑ים וַיִּקְבְּר֣וּ אֹת֔וֹ עֵשָׂ֥ו וְיַעֲקֹ֖ב בָּנָֽיו׃ \{פ\} \Rashi{\textbf{ויגוע יצחק.}אין מקדם ומאחר בתורה; מכירתו של יוסף קדמה למיתתו של יצחק י"ב שנה, שהרי כשנולד יעקב היה יצחק בן ששים שנה, שנאמר ויצחק בן ששים שנה וגו', ויצחק מת בשנת ק"כ ליעקב – אם תוציא ששים ממאה ושמונים שנה נשארו ק"כ – ויוסף נמכר בן י"ז שנה, ואותה שנה שנת מאה ושמונה ליעקב, כיצד? בן ששים ושלוש נתברך, ארבע עשרה שנה נטמן בבית עבר, הרי שבעים ושבע, וארבע עשרה עבד באשה, ובסוף ארבע עשרה נולד יוסף, שנאמר ויהי כאשר ילדה רחל את יוסף וגו', הרי תשעים ואחת, וי"ז עד שלא נמכר יוסף, הרי מאה ושמונה (עוד מפרש מן המקרא משנמכר יוסף עד שבא יעקב מצרימה כ"ב שנה שנ', ויוסף בן שלשים שנה וגו', וז' שנים שבע ושנתים רעב, הרי כ"ב, וכתיב ימי שני מגורי שלשים ומאת שנה, נמצא יעקב במכירתו ק"ח):}%
% ---------- פרק 36 | פסוק 1 ----------
 \textbf{א} וְאֵ֛לֶּה תֹּלְד֥וֹת עֵשָׂ֖ו ה֥וּא אֱדֽוֹם׃ %
% ---------- פרק 36 | פסוק 2 ----------
 \textbf{ב} עֵשָׂ֛ו לָקַ֥ח אֶת־נָשָׁ֖יו מִבְּנ֣וֹת כְּנָ֑עַן אֶת־עָדָ֗ה בַּת־אֵילוֹן֙ הַֽחִתִּ֔י וְאֶת־אׇהֳלִֽיבָמָה֙ בַּת־עֲנָ֔ה בַּת־צִבְע֖וֹן הַֽחִוִּֽי׃ \Rashi{\textbf{עדה בת אילון.}היא בשמת בת אילון, ונקראת בשמת על שם שהיתה מקטרת בשמים לע"ז: אהליבמה. היא יהודית, והוא כנה שמה יהודית לומר שהיא כופרת בע"ז כדי להטעות את אביו: בת ענה בת צבעון. אם בת ענה לא בת צבעון? ענה בנו של צבעון, שנאמר ואלה בני צבעון ואיה וענה? מלמד שבא צבעון על כלתו אשת ענה, ויצאה אהליבמה מבין שניהם, והודיעך הכתוב שכלן בני ממזרות היו:}%
% ---------- פרק 36 | פסוק 3 ----------
 \textbf{ג} וְאֶת־בָּשְׂמַ֥ת בַּת־יִשְׁמָעֵ֖אל אֲח֥וֹת נְבָיֽוֹת׃ \Rashi{\textbf{בשמת בת ישמעאל.}ולהלן קורא לה מחלת? מצינו באגדת מדרש ספר שמואל ג' מוחלים להן עונותיהן, גר שנתגיר, והעולה לגדלה, והנושא אשה. ולמד הטעם מכאן – לכך נקראת מחלת, שנמחלו עונותיו: אחות נביות. על שם שהוא השיאה לו משמת ישמעאל נקראת על שמו:}%
% ---------- פרק 36 | פסוק 4 ----------
 \textbf{ד} וַתֵּ֧לֶד עָדָ֛ה לְעֵשָׂ֖ו אֶת־אֱלִיפָ֑ז וּבָ֣שְׂמַ֔ת יָלְדָ֖ה אֶת־רְעוּאֵֽל׃ %
% ---------- פרק 36 | פסוק 5 ----------
 \textbf{ה} וְאׇהֳלִֽיבָמָה֙ יָֽלְדָ֔ה אֶת־[יְע֥וּשׁ] (יעיש) וְאֶת־יַעְלָ֖ם וְאֶת־קֹ֑רַח אֵ֚לֶּה בְּנֵ֣י עֵשָׂ֔ו אֲשֶׁ֥ר יֻלְּדוּ־ל֖וֹ בְּאֶ֥רֶץ כְּנָֽעַן׃ \Rashi{\textbf{ואהליבמה ילדה וגו'.}קרח זה ממזר היה, ובן אליפז היה, שבא על אשת אביו, אל אהליבמה אשת עשו, שהרי הוא מנוי עם אלופי אליפז בסוף הענין (בראשית רבה):}%
% ---------- פרק 36 | פסוק 6 ----------
 \textbf{ו} וַיִּקַּ֣ח עֵשָׂ֡ו אֶת־נָ֠שָׁ֠יו וְאֶת־בָּנָ֣יו וְאֶת־בְּנֹתָיו֮ וְאֶת־כׇּל־נַפְשׁ֣וֹת בֵּיתוֹ֒ וְאֶת־מִקְנֵ֣הוּ וְאֶת־כׇּל־בְּהֶמְתּ֗וֹ וְאֵת֙ כׇּל־קִנְיָנ֔וֹ אֲשֶׁ֥ר רָכַ֖שׁ בְּאֶ֣רֶץ כְּנָ֑עַן וַיֵּ֣לֶךְ אֶל־אֶ֔רֶץ מִפְּנֵ֖י יַעֲקֹ֥ב אָחִֽיו׃ \Rashi{\textbf{וילך אל ארץ.}לגור באשר ימצא:}%
% ---------- פרק 36 | פסוק 7 ----------
 \textbf{ז} כִּֽי־הָיָ֧ה רְכוּשָׁ֛ם רָ֖ב מִשֶּׁ֣בֶת יַחְדָּ֑ו וְלֹ֨א יָֽכְלָ֜ה אֶ֤רֶץ מְגֽוּרֵיהֶם֙ לָשֵׂ֣את אֹתָ֔ם מִפְּנֵ֖י מִקְנֵיהֶֽם׃ \Rashi{\textbf{ולא יכלה ארץ מגוריהם.}להספיק מרעה לבהמות שלהם (ב"ר). ומ"א: מפני יעקב אחיו. מפני שטר חוב של גזרת "כי גר יהיה זרעך" המטל על זרעו של יצחק, אמר אלך לי מכאן, אין לי חלק לא במתנה, שנתנה לו הארץ הזאת, ולא בפרעון השטר, ומפני הבושה שמכר בכורתו:}%
% ---------- פרק 36 | פסוק 8 ----------
 \textbf{ח} וַיֵּ֤שֶׁב עֵשָׂו֙ בְּהַ֣ר שֵׂעִ֔יר עֵשָׂ֖ו ה֥וּא אֱדֽוֹם׃ %
% ---------- פרק 36 | פסוק 9 ----------
 \textbf{ט} וְאֵ֛לֶּה תֹּלְד֥וֹת עֵשָׂ֖ו אֲבִ֣י אֱד֑וֹם בְּהַ֖ר שֵׂעִֽיר׃ \Rashi{\textbf{ואלה.}התולדות שהולידו בניו עכשו משהלך לשעיר:}%
% ---------- פרק 36 | פסוק 10 ----------
 \textbf{י} אֵ֖לֶּה שְׁמ֣וֹת בְּנֵֽי־עֵשָׂ֑ו אֱלִיפַ֗ז בֶּן־עָדָה֙ אֵ֣שֶׁת עֵשָׂ֔ו רְעוּאֵ֕ל בֶּן־בָּשְׂמַ֖ת אֵ֥שֶׁת עֵשָֽׂו׃ %
% ---------- פרק 36 | פסוק 11 ----------
 \textbf{יא} וַיִּהְי֖וּ בְּנֵ֣י אֱלִיפָ֑ז תֵּימָ֣ן אוֹמָ֔ר צְפ֥וֹ וְגַעְתָּ֖ם וּקְנַֽז׃ %
% ---------- פרק 36 | פסוק 12 ----------
 \textbf{יב} וְתִמְנַ֣ע ׀ הָיְתָ֣ה פִילֶ֗גֶשׁ לֶֽאֱלִיפַז֙ בֶּן־עֵשָׂ֔ו וַתֵּ֥לֶד לֶאֱלִיפַ֖ז אֶת־עֲמָלֵ֑ק אֵ֕לֶּה בְּנֵ֥י עָדָ֖ה אֵ֥שֶׁת עֵשָֽׂו׃ \Rashi{\textbf{ותמנע היתה פילגש.}להודיע גדלתו של אברהם, כמה היו תאבים לדבק בזרעו. תמנע זו בת אלופים היתה, שנאמר ואחות לוטן תמנע, ולוטן מאלופי יושבי שעיר היה מן החורים שישבו בה לפנים, אמרה איני זוכה להנשא לך, הלואי ואהיה פילגש. ובדברי הימים מונה אותה בבניו של אליפז, מלמד שבא על אשתו של שעיר ויצאה תמנע מביניהם, וכשגדלה נעשית פילגשו, וזהו ואחות לוטן תמנע, ולא מנאה עם בני שעיר שהיתה אחותו מן האם ולא מן האב:}%
% ---------- פרק 36 | פסוק 13 ----------
 \textbf{יג} וְאֵ֙לֶּה֙ בְּנֵ֣י רְעוּאֵ֔ל נַ֥חַת וָזֶ֖רַח שַׁמָּ֣ה וּמִזָּ֑ה אֵ֣לֶּה הָי֔וּ בְּנֵ֥י בָשְׂמַ֖ת אֵ֥שֶׁת עֵשָֽׂו׃ %
% ---------- פרק 36 | פסוק 14 ----------
 \textbf{יד} וְאֵ֣לֶּה הָי֗וּ בְּנֵ֨י אׇהֳלִיבָמָ֧ה בַת־עֲנָ֛ה בַּת־צִבְע֖וֹן אֵ֣שֶׁת עֵשָׂ֑ו וַתֵּ֣לֶד לְעֵשָׂ֔ו אֶת־[יְע֥וּשׁ] (יעיש) וְאֶת־יַעְלָ֖ם וְאֶת־קֹֽרַח׃ %
% ---------- פרק 36 | פסוק 15 ----------
 \textbf{טו} אֵ֖לֶּה אַלּוּפֵ֣י בְנֵֽי־עֵשָׂ֑ו בְּנֵ֤י אֱלִיפַז֙ בְּכ֣וֹר עֵשָׂ֔ו אַלּ֤וּף תֵּימָן֙ אַלּ֣וּף אוֹמָ֔ר אַלּ֥וּף צְפ֖וֹ אַלּ֥וּף קְנַֽז׃ \Rashi{\textbf{אלה אלופי בני עשו.}ראשי משפחות:}%
% ---------- פרק 36 | פסוק 16 ----------
 \textbf{טז} אַלּֽוּף־קֹ֛רַח אַלּ֥וּף גַּעְתָּ֖ם אַלּ֣וּף עֲמָלֵ֑ק אֵ֣לֶּה אַלּוּפֵ֤י אֱלִיפַז֙ בְּאֶ֣רֶץ אֱד֔וֹם אֵ֖לֶּה בְּנֵ֥י עָדָֽה׃ %
% ---------- פרק 36 | פסוק 17 ----------
 \textbf{יז} וְאֵ֗לֶּה בְּנֵ֤י רְעוּאֵל֙ בֶּן־עֵשָׂ֔ו אַלּ֥וּף נַ֙חַת֙ אַלּ֣וּף זֶ֔רַח אַלּ֥וּף שַׁמָּ֖ה אַלּ֣וּף מִזָּ֑ה אֵ֣לֶּה אַלּוּפֵ֤י רְעוּאֵל֙ בְּאֶ֣רֶץ אֱד֔וֹם אֵ֕לֶּה בְּנֵ֥י בָשְׂמַ֖ת אֵ֥שֶׁת עֵשָֽׂו׃ %
% ---------- פרק 36 | פסוק 18 ----------
 \textbf{יח} וְאֵ֗לֶּה בְּנֵ֤י אׇהֳלִֽיבָמָה֙ אֵ֣שֶׁת עֵשָׂ֔ו אַלּ֥וּף יְע֛וּשׁ אַלּ֥וּף יַעְלָ֖ם אַלּ֣וּף קֹ֑רַח אֵ֣לֶּה אַלּוּפֵ֞י אׇהֳלִֽיבָמָ֛ה בַּת־עֲנָ֖ה אֵ֥שֶׁת עֵשָֽׂו׃ %
% ---------- פרק 36 | פסוק 19 ----------
 \textbf{יט} אֵ֧לֶּה בְנֵי־עֵשָׂ֛ו וְאֵ֥לֶּה אַלּוּפֵיהֶ֖ם ה֥וּא אֱדֽוֹם׃ \{ס\} %
% ---------- פרק 36 | פסוק 20 ----------
 \textbf{כ} אֵ֤לֶּה בְנֵֽי־שֵׂעִיר֙ הַחֹרִ֔י יֹשְׁבֵ֖י הָאָ֑רֶץ לוֹטָ֥ן וְשׁוֹבָ֖ל וְצִבְע֥וֹן וַעֲנָֽה׃ \Rashi{\textbf{ישבי הארץ.}שהיו יושביה קדם שבא עשו לשם. ורבותינו דרשו (שבת פ"ה) שהיו בקיאין בישובה של ארץ; מלא קנה זה לזיתים, מלא קנה זה לגפנים, שהיו טועמין העפר ויודעין אי זו נטיעה ראויה לו:}%
% ---------- פרק 36 | פסוק 21 ----------
 \textbf{כא} וְדִשׁ֥וֹן וְאֵ֖צֶר וְדִישָׁ֑ן אֵ֣לֶּה אַלּוּפֵ֧י הַחֹרִ֛י בְּנֵ֥י שֵׂעִ֖יר בְּאֶ֥רֶץ אֱדֽוֹם׃ %
% ---------- פרק 36 | פסוק 22 ----------
 \textbf{כב} וַיִּהְי֥וּ בְנֵי־לוֹטָ֖ן חֹרִ֣י וְהֵימָ֑ם וַאֲח֥וֹת לוֹטָ֖ן תִּמְנָֽע׃ %
% ---------- פרק 36 | פסוק 23 ----------
 \textbf{כג} וְאֵ֙לֶּה֙ בְּנֵ֣י שׁוֹבָ֔ל עַלְוָ֥ן וּמָנַ֖חַת וְעֵיבָ֑ל שְׁפ֖וֹ וְאוֹנָֽם׃ %
% ---------- פרק 36 | פסוק 24 ----------
 \textbf{כד} וְאֵ֥לֶּה בְנֵֽי־צִבְע֖וֹן וְאַיָּ֣ה וַעֲנָ֑ה ה֣וּא עֲנָ֗ה אֲשֶׁ֨ר מָצָ֤א אֶת־הַיֵּמִם֙ בַּמִּדְבָּ֔ר בִּרְעֹת֥וֹ אֶת־הַחֲמֹרִ֖ים לְצִבְע֥וֹן אָבִֽיו׃ \Rashi{\textbf{ואיה וענה.}וי"ו יתרה, והוא כמו איה וענה, והרבה יש במקרא, תת וקדש וצבא מרמס (דניאל ח'), נרדם ורכב וסוס (תהילים ע"ו): הוא ענה. האמור למעלה שהוא אחיו של צבעון וכאן הוא קורא אותו בנו? מלמד שבא צבעון על אמו והוליד את ענה: את הימם. פרדים, הרביע חמור על סוס נקבה וילדה פרד, והוא היה ממזר והביא פסולין לעולם; ולמה נקרא שמם ימים? שאימתן מטלת על הבריות, דאמר רבי חנינא מימי לא שאלני אדם על מכת פרדה לבנה וחיה (חולין ז') (והלא קא חזינן דחיה? אל תקרי וחיה אלא וחיתה, כי המכה לא תרפא לעולם, ברש"י ישן) ולא הזקק לכתב לנו משפחות החורי אלא מפני תמנע, ולהודיע גדלת אברהם, כמו שפרשתי למעלה:}%
% ---------- פרק 36 | פסוק 25 ----------
 \textbf{כה} וְאֵ֥לֶּה בְנֵֽי־עֲנָ֖ה דִּשֹׁ֑ן וְאׇהֳלִיבָמָ֖ה בַּת־עֲנָֽה׃ %
% ---------- פרק 36 | פסוק 26 ----------
 \textbf{כו} וְאֵ֖לֶּה בְּנֵ֣י דִישָׁ֑ן חֶמְדָּ֥ן וְאֶשְׁבָּ֖ן וְיִתְרָ֥ן וּכְרָֽן׃ %
% ---------- פרק 36 | פסוק 27 ----------
 \textbf{כז} אֵ֖לֶּה בְּנֵי־אֵ֑צֶר בִּלְהָ֥ן וְזַעֲוָ֖ן וַעֲקָֽן׃ %
% ---------- פרק 36 | פסוק 28 ----------
 \textbf{כח} אֵ֥לֶּה בְנֵֽי־דִישָׁ֖ן ע֥וּץ וַאֲרָֽן׃ %
% ---------- פרק 36 | פסוק 29 ----------
 \textbf{כט} אֵ֖לֶּה אַלּוּפֵ֣י הַחֹרִ֑י אַלּ֤וּף לוֹטָן֙ אַלּ֣וּף שׁוֹבָ֔ל אַלּ֥וּף צִבְע֖וֹן אַלּ֥וּף עֲנָֽה׃ %
% ---------- פרק 36 | פסוק 30 ----------
 \textbf{ל} אַלּ֥וּף דִּשֹׁ֛ן אַלּ֥וּף אֵ֖צֶר אַלּ֣וּף דִּישָׁ֑ן אֵ֣לֶּה אַלּוּפֵ֧י הַחֹרִ֛י לְאַלֻּפֵיהֶ֖ם בְּאֶ֥רֶץ שֵׂעִֽיר׃ \{פ\} %
% ---------- פרק 36 | פסוק 31 ----------
 \textbf{לא} וְאֵ֙לֶּה֙ הַמְּלָכִ֔ים אֲשֶׁ֥ר מָלְכ֖וּ בְּאֶ֣רֶץ אֱד֑וֹם לִפְנֵ֥י מְלׇךְ־מֶ֖לֶךְ לִבְנֵ֥י יִשְׂרָאֵֽל׃ \Rashi{\textbf{ואלה המלכים וגו'.}שמונה היו וכנגדן העמיד יעקב ובטל מלכות עשו בימיהם, אלו הן, שאול ואיש בשת, דוד ושלמה, רחבעם, אביה, אסא, יהושפט; ובימי יורם בנו כתוב בימיו פשע אדום מתחת יד יהודה וימלכו עליהם מלך (מלכים ב ח'), ובימי שאול כתוב ומלך אין באדום נצב מלך (מלכים א כ"ב):}%
% ---------- פרק 36 | פסוק 32 ----------
 \textbf{לב} וַיִּמְלֹ֣ךְ בֶּאֱד֔וֹם בֶּ֖לַע בֶּן־בְּע֑וֹר וְשֵׁ֥ם עִיר֖וֹ דִּנְהָֽבָה׃ %
% ---------- פרק 36 | פסוק 33 ----------
 \textbf{לג} וַיָּ֖מׇת בָּ֑לַע וַיִּמְלֹ֣ךְ תַּחְתָּ֔יו יוֹבָ֥ב בֶּן־זֶ֖רַח מִבׇּצְרָֽה׃ \Rashi{\textbf{יובב בן זרח מבצרה.}בצרה מערי מואב היא, שנאמר ועל קריות ועל בצרה וגו' (ירמיהו מ"ח); ולפי שהעמידה מלך לאדום, עתידה ללקות עמהם, שנאמר כי זבח לה' בבצרה (ישעיהו ל"ד):}%
% ---------- פרק 36 | פסוק 34 ----------
 \textbf{לד} וַיָּ֖מׇת יוֹבָ֑ב וַיִּמְלֹ֣ךְ תַּחְתָּ֔יו חֻשָׁ֖ם מֵאֶ֥רֶץ הַתֵּימָנִֽי׃ %
% ---------- פרק 36 | פסוק 35 ----------
 \textbf{לה} וַיָּ֖מׇת חֻשָׁ֑ם וַיִּמְלֹ֨ךְ תַּחְתָּ֜יו הֲדַ֣ד בֶּן־בְּדַ֗ד הַמַּכֶּ֤ה אֶת־מִדְיָן֙ בִּשְׂדֵ֣ה מוֹאָ֔ב וְשֵׁ֥ם עִיר֖וֹ עֲוִֽית׃ \Rashi{\textbf{המכה את מדין בשדה מואב.}שבאו מדין על מואב למלחמה, והלך מלך אדום לעזר את מואב; ומכאן אנו למדים, שהיו מדין ומואב מריבים זה עם זה, ובימי בלעם עשו שלום להתקשר על ישראל:}%
% ---------- פרק 36 | פסוק 36 ----------
 \textbf{לו} וַיָּ֖מׇת הֲדָ֑ד וַיִּמְלֹ֣ךְ תַּחְתָּ֔יו שַׂמְלָ֖ה מִמַּשְׂרֵקָֽה׃ %
% ---------- פרק 36 | פסוק 37 ----------
 \textbf{לז} וַיָּ֖מׇת שַׂמְלָ֑ה וַיִּמְלֹ֣ךְ תַּחְתָּ֔יו שָׁא֖וּל מֵרְחֹב֥וֹת הַנָּהָֽר׃ %
% ---------- פרק 36 | פסוק 38 ----------
 \textbf{לח} וַיָּ֖מׇת שָׁא֑וּל וַיִּמְלֹ֣ךְ תַּחְתָּ֔יו בַּ֥עַל חָנָ֖ן בֶּן־עַכְבּֽוֹר׃ %
% ---------- פרק 36 | פסוק 39 ----------
 \textbf{לט} וַיָּ֘מׇת֮ בַּ֣עַל חָנָ֣ן בֶּן־עַכְבּוֹר֒ וַיִּמְלֹ֤ךְ תַּחְתָּיו֙ הֲדַ֔ר וְשֵׁ֥ם עִיר֖וֹ פָּ֑עוּ וְשֵׁ֨ם אִשְׁתּ֤וֹ מְהֵֽיטַבְאֵל֙ בַּת־מַטְרֵ֔ד בַּ֖ת מֵ֥י זָהָֽב׃ \Rashi{\textbf{בת מי זהב.}מהו זהב? עשיר היה, ואין זהב חשוב בעיניו לכלום:}%
% ---------- פרק 36 | פסוק 40 ----------
 \textbf{מ} וְ֠אֵ֠לֶּה שְׁמ֞וֹת אַלּוּפֵ֤י עֵשָׂו֙ לְמִשְׁפְּחֹתָ֔ם לִמְקֹמֹתָ֖ם בִּשְׁמֹתָ֑ם אַלּ֥וּף תִּמְנָ֛ע אַלּ֥וּף עַֽלְוָ֖ה אַלּ֥וּף יְתֵֽת׃ \Rashi{\textbf{ואלה שמות אלופי עשו.}שנקראו על שם מדינותיהם לאחר שמת הדר ופסקה מהם מלכות, והראשונים הנזכרים למעלה הם שמות תולדותם; וכן מפרש בדברי הימים וימת הדד ויהיו אלופי אדום אלוף תמנע וגומר:}%
% ---------- פרק 36 | פסוק 41 ----------
 \textbf{מא} אַלּ֧וּף אׇהֳלִיבָמָ֛ה אַלּ֥וּף אֵלָ֖ה אַלּ֥וּף פִּינֹֽן׃ %
% ---------- פרק 36 | פסוק 42 ----------
 \textbf{מב} אַלּ֥וּף קְנַ֛ז אַלּ֥וּף תֵּימָ֖ן אַלּ֥וּף מִבְצָֽר׃ %
% ---------- פרק 36 | פסוק 43 ----------
 \textbf{מג} אַלּ֥וּף מַגְדִּיאֵ֖ל אַלּ֣וּף עִירָ֑ם אֵ֣לֶּה ׀ אַלּוּפֵ֣י אֱד֗וֹם לְמֹֽשְׁבֹתָם֙ בְּאֶ֣רֶץ אֲחֻזָּתָ֔ם ה֥וּא עֵשָׂ֖ו אֲבִ֥י אֱדֽוֹם׃ \{פ\} \Rashi{\textbf{מגדיאל.}היא רומי:}%
\addparasha{חומש בראשית | פרשת וישב}
% ---------- פרק 37 | פסוק 1 ----------
 \textbf{א} וַיֵּ֣שֶׁב יַעֲקֹ֔ב בְּאֶ֖רֶץ מְגוּרֵ֣י אָבִ֑יו בְּאֶ֖רֶץ כְּנָֽעַן׃ \Rashi{\textbf{וישב יעקב וגו'.}אחר שכתב לך ישובי עשו ותולדותיו בדרך קצרה, שלא היו ספונים וחשובים לפרש היאך נתישבו וסדר מלחמותיהם איך הורישו את החרי, פרש לך ישובי יעקב ותולדותיו בדרך ארכה כל גלגולי סבתם, לפי שהם חשובים לפני המקום להאריך בהם, וכן אתה מוצא בי' דורות שמאדם ועד נח פלוני הוליד פלוני, וכשבא לנח האריך בו, וכן בי' דורות שמנח ועד אברהם קצר בהם, ומשהגיע אצל אברהם האריך בו. משל למרגלית שנפלה בין החול, אדם ממשמש בחול וכוברו בכברה עד שמוצא את המרגלית, ומשמצאה הוא משליך את הצרורות מידו ונוטל המרגלית. ד"א וישב יעקב, הפשתני הזה נכנסו גמליו טעונים פשתן, הפחמי תמה אנה יכנס כל הפשתן הזה? היה פקח אחד משיב לו נצוץ אחד יוצא ממפוח שלך ששורף את כלו, כך יעקב ראה את כל האלופים הכתובים למעלה, תמה ואמר מי יכול לכבש את כלן? מה כתיב למטה, אלה תולדות יעקב יוסף, דכתיב והיה בית יעקב אש ובית יוסף להבה ובית עשו לקש (עובדיה א') – נצוץ יוצא מיוסף שמכלה ושורף את כלם:}%
% ---------- פרק 37 | פסוק 2 ----------
 \textbf{ב} אֵ֣לֶּה ׀ תֹּלְד֣וֹת יַעֲקֹ֗ב יוֹסֵ֞ף בֶּן־שְׁבַֽע־עֶשְׂרֵ֤ה שָׁנָה֙ הָיָ֨ה רֹעֶ֤ה אֶת־אֶחָיו֙ בַּצֹּ֔אן וְה֣וּא נַ֗עַר אֶת־בְּנֵ֥י בִלְהָ֛ה וְאֶת־בְּנֵ֥י זִלְפָּ֖ה נְשֵׁ֣י אָבִ֑יו וַיָּבֵ֥א יוֹסֵ֛ף אֶת־דִּבָּתָ֥ם רָעָ֖ה אֶל־אֲבִיהֶֽם׃ \Rashi{\textbf{אלה תולדות יעקב.}ואלה של תולדות יעקב, אלו ישוביהם וגלגוליהם עד שבאו לכלל ישוב, סבה ראשונה יוסף בן י"ז וגומר, על ידי זה נתגלגלו וירדו למצרים, זהו אחר ישוב פשוטו של מקרא להיות דבר דבור על אפניו. ומ"א דורש תלה הכתוב תולדות יעקב ביוסף מפני כמה דברים, אחת, שכל עצמו של יעקב לא עבד אצל לבן אלא ברחל, ושהיה זיו איקונין של יוסף דומה לו, וכל מה שארע ליעקב ארע ליוסף: זה נשטם וזה נשטם, זה אחיו מבקש להרגו וזה אחיו מבקשים להרגו, וכן הרבה בב"ר. ועוד נדרש בו וישב, בקש יעקב לישב בשלוה, קפץ עליו רגזו של יוסף – צדיקים מבקשים לישב בשלוה, אמר הקב"ה לא דין לצדיקים מה שמתקן להם לעולם הבא, אלא שמבקשים לישב בשלוה בעולם הזה: והוא נער. שהיה עושה מעשה נערות, מתקן בשערו, ממשמש בעיניו, כדי שיהיה נראה יפה: את בני בלהה. כלומר ורגיל אצל בני בלהה, לפי שהיו אחיו מבזין אותן והוא מקרבן: את דבתם רעה. כל רעה שהיה רואה באחיו בני לאה היה מגיד לאביו, שהיו אוכלין אבר מן החי, ומזלזלין בבני השפחות לקרותן עבדים, וחשודים על העריות, ובשלשתן לקה: וישחטו שעיר עזים במכירתו ולא אכלוהו חי, ועל דבה שספר עליהם שקורין לאחיהם עבדים – "לעבד נמכר יוסף", ועל עריות שספר עליהם, ותשא אשת אדניו וגו: דבתם. כל לשון דבה פרלדי"ץ בלעז, כל מה שהיה יכול לדבר בהם רעה, היה מספר: דבה. לשון דובב שפתי ישנים (שיר ז'):}%
% ---------- פרק 37 | פסוק 3 ----------
 \textbf{ג} וְיִשְׂרָאֵ֗ל אָהַ֤ב אֶת־יוֹסֵף֙ מִכׇּל־בָּנָ֔יו כִּֽי־בֶן־זְקֻנִ֥ים ה֖וּא ל֑וֹ וְעָ֥שָׂה ל֖וֹ כְּתֹ֥נֶת פַּסִּֽים׃ \Rashi{\textbf{בן זקנים.}שנולד לו לעת זקנתו. ואנקלוס תרגם בר חכים הוא לה – כל מה שלמד משם ועבר מסר לו. ד"א, שהיה זיו איקונין שלו דומה לו: פסים. לשון כלי מילת, כמו כרפס ותכלת, וכמו כתנת הפסים דתמר ואמנון, ומ"א ע"ש צרותיו, שנמכר לפוטיפר ולסוחרים ולישמעאלים ולמדינים:}%
% ---------- פרק 37 | פסוק 4 ----------
 \textbf{ד} וַיִּרְא֣וּ אֶחָ֗יו כִּֽי־אֹת֞וֹ אָהַ֤ב אֲבִיהֶם֙ מִכׇּל־אֶחָ֔יו וַֽיִּשְׂנְא֖וּ אֹת֑וֹ וְלֹ֥א יָכְל֖וּ דַּבְּר֥וֹ לְשָׁלֹֽם׃ \Rashi{\textbf{ולא יכלו דברו לשלום.}מתוך גנותם למדנו שבחם, שלא דברו אחת בפה ואחת בלב: דברו. לדבר עמו:}%
% ---------- פרק 37 | פסוק 5 ----------
 \textbf{ה} וַיַּחֲלֹ֤ם יוֹסֵף֙ חֲל֔וֹם וַיַּגֵּ֖ד לְאֶחָ֑יו וַיּוֹסִ֥פוּ ע֖וֹד שְׂנֹ֥א אֹתֽוֹ׃ %
% ---------- פרק 37 | פסוק 6 ----------
 \textbf{ו} וַיֹּ֖אמֶר אֲלֵיהֶ֑ם שִׁמְעוּ־נָ֕א הַחֲל֥וֹם הַזֶּ֖ה אֲשֶׁ֥ר חָלָֽמְתִּי׃ %
% ---------- פרק 37 | פסוק 7 ----------
 \textbf{ז} וְ֠הִנֵּ֠ה אֲנַ֜חְנוּ מְאַלְּמִ֤ים אֲלֻמִּים֙ בְּת֣וֹךְ הַשָּׂדֶ֔ה וְהִנֵּ֛ה קָ֥מָה אֲלֻמָּתִ֖י וְגַם־נִצָּ֑בָה וְהִנֵּ֤ה תְסֻבֶּ֙ינָה֙ אֲלֻמֹּ֣תֵיכֶ֔ם וַתִּֽשְׁתַּחֲוֶ֖יןָ לַאֲלֻמָּתִֽי׃ \Rashi{\textbf{מאלמים אלמים.}כתרגומו מאסרין אסרין – עמרין, וכן נשא אלמתיו (תהלים קכ"ו), וכמוהו בלשון משנה והאלמות נוטל ומכריז: קמה אלמתי. נזקפה: וגם נצבה. לעמד על עמדה בזקיפה:}%
% ---------- פרק 37 | פסוק 8 ----------
 \textbf{ח} וַיֹּ֤אמְרוּ לוֹ֙ אֶחָ֔יו הֲמָלֹ֤ךְ תִּמְלֹךְ֙ עָלֵ֔ינוּ אִם־מָשׁ֥וֹל תִּמְשֹׁ֖ל בָּ֑נוּ וַיּוֹסִ֤פוּ עוֹד֙ שְׂנֹ֣א אֹת֔וֹ עַל־חֲלֹמֹתָ֖יו וְעַל־דְּבָרָֽיו׃ \Rashi{\textbf{ועל דבריו.}על דבתם רעה שהיה מביא לאביהם:}%
% ---------- פרק 37 | פסוק 9 ----------
 \textbf{ט} וַיַּחֲלֹ֥ם עוֹד֙ חֲל֣וֹם אַחֵ֔ר וַיְסַפֵּ֥ר אֹת֖וֹ לְאֶחָ֑יו וַיֹּ֗אמֶר הִנֵּ֨ה חָלַ֤מְתִּֽי חֲלוֹם֙ ע֔וֹד וְהִנֵּ֧ה הַשֶּׁ֣מֶשׁ וְהַיָּרֵ֗חַ וְאַחַ֤ד עָשָׂר֙ כּֽוֹכָבִ֔ים מִֽשְׁתַּחֲוִ֖ים לִֽי׃ %
% ---------- פרק 37 | פסוק 10 ----------
 \textbf{י} וַיְסַפֵּ֣ר אֶל־אָבִיו֮ וְאֶל־אֶחָיו֒ וַיִּגְעַר־בּ֣וֹ אָבִ֔יו וַיֹּ֣אמֶר ל֔וֹ מָ֛ה הַחֲל֥וֹם הַזֶּ֖ה אֲשֶׁ֣ר חָלָ֑מְתָּ הֲב֣וֹא נָב֗וֹא אֲנִי֙ וְאִמְּךָ֣ וְאַחֶ֔יךָ לְהִשְׁתַּחֲוֺ֥ת לְךָ֖ אָֽרְצָה׃ \Rashi{\textbf{ויספר אל אביו ואל אחיו.}לאחר שספר אותו לאחיו חזר וספרו לאביו בפניהם: ויגער בו. לפי שהיה מטיל שנאה עליו: הבוא נבוא. והלא אמך כבר מתה, והוא לא היה יודע שהדברים מגיעין לבלהה שגדלתו כאמו. ורבותינו למדו מכאן שאין חלום בלא דברים בטלים, ויעקב נתכון להוציא הדבר מלב בניו שלא יקניאוהו, לכך אמר לו הבוא נבוא וגו' – כשם שאי אפשר באמך כך השאר הוא בטל:}%
% ---------- פרק 37 | פסוק 11 ----------
 \textbf{יא} וַיְקַנְאוּ־ב֖וֹ אֶחָ֑יו וְאָבִ֖יו שָׁמַ֥ר אֶת־הַדָּבָֽר׃ \Rashi{\textbf{שמר את הדבר.}היה ממתין ומצפה מתי יבא, וכן שמר אמנים (ישעיהו כ"ו), וכן לא תשמר על חטאתי (איוב י"ד) – לא תמתין:}%
% ---------- פרק 37 | פסוק 12 ----------
 \textbf{יב} וַיֵּלְכ֖וּ אֶחָ֑יו לִרְע֛וֹת אֶׄתׄ־צֹ֥אן אֲבִיהֶ֖ם בִּשְׁכֶֽם׃ \Rashi{\textbf{לרעות את צאן.}נקוד על את, שלא הלכו אלא לרעות את עצמן:}%
% ---------- פרק 37 | פסוק 13 ----------
 \textbf{יג} וַיֹּ֨אמֶר יִשְׂרָאֵ֜ל אֶל־יוֹסֵ֗ף הֲל֤וֹא אַחֶ֙יךָ֙ רֹעִ֣ים בִּשְׁכֶ֔ם לְכָ֖ה וְאֶשְׁלָחֲךָ֣ אֲלֵיהֶ֑ם וַיֹּ֥אמֶר ל֖וֹ הִנֵּֽנִי׃ \Rashi{\textbf{הנני.}לשון ענוה וזריזות, נזדרז למצות אביו, ואף על פי שהיה יודע באחיו ששונאין אותו:}%
% ---------- פרק 37 | פסוק 14 ----------
 \textbf{יד} וַיֹּ֣אמֶר ל֗וֹ לֶךְ־נָ֨א רְאֵ֜ה אֶת־שְׁל֤וֹם אַחֶ֙יךָ֙ וְאֶת־שְׁל֣וֹם הַצֹּ֔אן וַהֲשִׁבֵ֖נִי דָּבָ֑ר וַיִּשְׁלָחֵ֙הוּ֙ מֵעֵ֣מֶק חֶבְר֔וֹן וַיָּבֹ֖א שְׁכֶֽמָה׃ \Rashi{\textbf{מעמק חברון.}והלא חברון בהר, שנאמר ויעלו בנגב ויבא עד חברון (במדבר י"ג), אלא מעצה עמקה של אותו צדיק הקבור בחברון, לקים מה שנאמר לאברהם בין הבתרים כי גר יהיה זרעך (בראשית ט״ו:י״ג): ויבא שכמה. מקום מוכן לפרענות, שם קלקלו השבטים, שם ענו את דינה, שם נחלקה מלכות בית דוד, שנאמר וילך רחבעם שכמה (מלכים א י"ב):}%
% ---------- פרק 37 | פסוק 15 ----------
 \textbf{טו} וַיִּמְצָאֵ֣הוּ אִ֔ישׁ וְהִנֵּ֥ה תֹעֶ֖ה בַּשָּׂדֶ֑ה וַיִּשְׁאָלֵ֧הוּ הָאִ֛ישׁ לֵאמֹ֖ר מַה־תְּבַקֵּֽשׁ׃ \Rashi{\textbf{וימצאהו איש.}זה גבריאל, שנאמר והאיש גבריאל (דניאל ט כא):}%
% ---------- פרק 37 | פסוק 16 ----------
 \textbf{טז} וַיֹּ֕אמֶר אֶת־אַחַ֖י אָנֹכִ֣י מְבַקֵּ֑שׁ הַגִּֽידָה־נָּ֣א לִ֔י אֵיפֹ֖ה הֵ֥ם רֹעִֽים׃ %
% ---------- פרק 37 | פסוק 17 ----------
 \textbf{יז} וַיֹּ֤אמֶר הָאִישׁ֙ נָסְע֣וּ מִזֶּ֔ה כִּ֤י שָׁמַ֙עְתִּי֙ אֹֽמְרִ֔ים נֵלְכָ֖ה דֹּתָ֑יְנָה וַיֵּ֤לֶךְ יוֹסֵף֙ אַחַ֣ר אֶחָ֔יו וַיִּמְצָאֵ֖ם בְּדֹתָֽן׃ \Rashi{\textbf{נסעו מזה.}הסיעו עצמן מן האחוה: נלכה דתינה. לבקש לך נכלי דתות שימיתוך בהם. ולפי פשוטו שם מקום הוא, ואין מקרא יוצא מידי פשוטו:}%
% ---------- פרק 37 | פסוק 18 ----------
 \textbf{יח} וַיִּרְא֥וּ אֹת֖וֹ מֵרָחֹ֑ק וּבְטֶ֙רֶם֙ יִקְרַ֣ב אֲלֵיהֶ֔ם וַיִּֽתְנַכְּל֥וּ אֹת֖וֹ לַהֲמִיתֽוֹ׃ \Rashi{\textbf{ויתנכלו.}נתמלאו נכלים וערמומית: אתו. כמו אתו, עמו, כלומר, אליו:}%
% ---------- פרק 37 | פסוק 19 ----------
 \textbf{יט} וַיֹּאמְר֖וּ אִ֣ישׁ אֶל־אָחִ֑יו הִנֵּ֗ה בַּ֛עַל הַחֲלֹמ֥וֹת הַלָּזֶ֖ה בָּֽא׃ %
% ---------- פרק 37 | פסוק 20 ----------
 \textbf{כ} וְעַתָּ֣ה ׀ לְכ֣וּ וְנַֽהַרְגֵ֗הוּ וְנַשְׁלִכֵ֙הוּ֙ בְּאַחַ֣ד הַבֹּר֔וֹת וְאָמַ֕רְנוּ חַיָּ֥ה רָעָ֖ה אֲכָלָ֑תְהוּ וְנִרְאֶ֕ה מַה־יִּהְי֖וּ חֲלֹמֹתָֽיו׃ \Rashi{\textbf{ונראה מה יהיו חלמתיו.}אמר רבי יצחק מקרא זה אומר דרשני, רוח הקדש אומרת כן, הם אומרים נהרגהו והכתוב מסים ונראה מה יהיו חלמתיו – נראה דבר מי יקום אם שלכם או שלי. וא"א שיאמרו הם ונראה מה יהיו חלמתיו, שמכיון שיהרגוהו בטלו חלומותיו:}%
% ---------- פרק 37 | פסוק 21 ----------
 \textbf{כא} וַיִּשְׁמַ֣ע רְאוּבֵ֔ן וַיַּצִּלֵ֖הוּ מִיָּדָ֑ם וַיֹּ֕אמֶר לֹ֥א נַכֶּ֖נּוּ נָֽפֶשׁ׃ \Rashi{\textbf{לא נכנו נפש.}מכת נפש, זו היא מיתה:}%
% ---------- פרק 37 | פסוק 22 ----------
 \textbf{כב} וַיֹּ֨אמֶר אֲלֵהֶ֣ם ׀ רְאוּבֵן֮ אַל־תִּשְׁפְּכוּ־דָם֒ הַשְׁלִ֣יכוּ אֹת֗וֹ אֶל־הַבּ֤וֹר הַזֶּה֙ אֲשֶׁ֣ר בַּמִּדְבָּ֔ר וְיָ֖ד אַל־תִּשְׁלְחוּ־ב֑וֹ לְמַ֗עַן הַצִּ֤יל אֹתוֹ֙ מִיָּדָ֔ם לַהֲשִׁיב֖וֹ אֶל־אָבִֽיו׃ \Rashi{\textbf{למען הציל אתו.}רוח הקדש מעידה על ראובן שלא אמר זאת אלא להציל אותו שיבא הוא ויעלנו משם, אמר אני בכור וגדול שבכלן, לא יתלה הסרחון אלא בי:}%
% ---------- פרק 37 | פסוק 23 ----------
 \textbf{כג} וַיְהִ֕י כַּֽאֲשֶׁר־בָּ֥א יוֹסֵ֖ף אֶל־אֶחָ֑יו וַיַּפְשִׁ֤יטוּ אֶת־יוֹסֵף֙ אֶת־כֻּתׇּנְתּ֔וֹ אֶת־כְּתֹ֥נֶת הַפַּסִּ֖ים אֲשֶׁ֥ר עָלָֽיו׃ \Rashi{\textbf{את כתנתו.}זה חלוק: את כתנת הפסים. הוא שהוסיף לו אביו יותר על אחיו:}%
% ---------- פרק 37 | פסוק 24 ----------
 \textbf{כד} וַיִּ֨קָּחֻ֔הוּ וַיַּשְׁלִ֥כוּ אֹת֖וֹ הַבֹּ֑רָה וְהַבּ֣וֹר רֵ֔ק אֵ֥ין בּ֖וֹ מָֽיִם׃ \Rashi{\textbf{והבור רק אין בו מים.}ממשמע שנאמר והבור רק, איני יודע שאין בו מים, מה תלמוד לומר אין בו מים? מים אין בו, אבל נחשים ועקרבים יש בו (בראשית רבה, שבת כ"ב):}%
% ---------- פרק 37 | פסוק 25 ----------
 \textbf{כה} וַיֵּשְׁבוּ֮ לֶֽאֱכׇל־לֶ֒חֶם֒ וַיִּשְׂא֤וּ עֵֽינֵיהֶם֙ וַיִּרְא֔וּ וְהִנֵּה֙ אֹרְחַ֣ת יִשְׁמְעֵאלִ֔ים בָּאָ֖ה מִגִּלְעָ֑ד וּגְמַלֵּיהֶ֣ם נֹֽשְׂאִ֗ים נְכֹאת֙ וּצְרִ֣י וָלֹ֔ט הוֹלְכִ֖ים לְהוֹרִ֥יד מִצְרָֽיְמָה׃ \Rashi{\textbf{ארחת.}כתרגומו שירת, על שם הולכי ארח: וגמליהם נשאים וגו'. למה פרסם הכתוב את משאם? להודיע מתן שכרן של צדיקים, שאין דרכן של ערביים לשאת אלא נפט ועטרן, שריחן רע, ולזה נזדמנו בשמים, שלא יזוק מריח רע: נכאת. כל כנוסי בשמים הרבה קרוי נכאת, וכן ויראם את כל בית נכתה (מלכים ב כ') – מרקחת בשמיו. ואנקלוס תרגם לשון שעוה: צרי. שרף הנוטף מעצי הקטף, והוא נטף הנמנה עם סמני הקטרת: ולט. לוטיתא שמו בלשון משנה. ורבותינו פרשוהו שרש עשב, ושמו אשטרולוזיאה במסכת נדה:}%
% ---------- פרק 37 | פסוק 26 ----------
 \textbf{כו} וַיֹּ֥אמֶר יְהוּדָ֖ה אֶל־אֶחָ֑יו מַה־בֶּ֗צַע כִּ֤י נַהֲרֹג֙ אֶת־אָחִ֔ינוּ וְכִסִּ֖ינוּ אֶת־דָּמֽוֹ׃ \Rashi{\textbf{מה בצע.}מה ממון, כתרגומו: וכסינו את דמו. ונעלים את מיתתו:}%
% ---------- פרק 37 | פסוק 27 ----------
 \textbf{כז} לְכ֞וּ וְנִמְכְּרֶ֣נּוּ לַיִּשְׁמְעֵאלִ֗ים וְיָדֵ֙נוּ֙ אַל־תְּהִי־ב֔וֹ כִּֽי־אָחִ֥ינוּ בְשָׂרֵ֖נוּ ה֑וּא וַֽיִּשְׁמְע֖וּ אֶחָֽיו׃ \Rashi{\textbf{וישמעו.}וקבילו מנה, וכל שמיעה שהיא קבלת דברים כגון זה, וכגון וישמע יעקב אל אביו, נעשה ונשמע – מתרגם נקבל; וכל שהיא שמיעת האזן, כגון וישמעו את קול ה' אלהים מתהלך בגן, ורבקה שמעת, וישמע ישראל, שמעתי את תלונת – כלן מתרגם ושמעו, ושמעת, ושמע, שמיע קדמי:}%
% ---------- פרק 37 | פסוק 28 ----------
 \textbf{כח} וַיַּֽעַבְרוּ֩ אֲנָשִׁ֨ים מִדְיָנִ֜ים סֹֽחֲרִ֗ים וַֽיִּמְשְׁכוּ֙ וַיַּֽעֲל֤וּ אֶת־יוֹסֵף֙ מִן־הַבּ֔וֹר וַיִּמְכְּר֧וּ אֶת־יוֹסֵ֛ף לַיִּשְׁמְעֵאלִ֖ים בְּעֶשְׂרִ֣ים כָּ֑סֶף וַיָּבִ֥יאוּ אֶת־יוֹסֵ֖ף מִצְרָֽיְמָה׃ \Rashi{\textbf{ויעברו אנשים מדינים.}זו היא שירה אחרת, והודיעך הכתוב שנמכר פעמים הרבה: וימשכו. בני יעקב את יוסף מן הבור וימכרוהו לישמעאלים, וישמעאלים למדינים, והמדינים למצרים:}%
% ---------- פרק 37 | פסוק 29 ----------
 \textbf{כט} וַיָּ֤שׇׁב רְאוּבֵן֙ אֶל־הַבּ֔וֹר וְהִנֵּ֥ה אֵין־יוֹסֵ֖ף בַּבּ֑וֹר וַיִּקְרַ֖ע אֶת־בְּגָדָֽיו׃ \Rashi{\textbf{וישב ראובן.}במכירתו לא היה שם, שהגיע יומו לילך ולשמש את אביו. ד"א עסוק היה בשקו ובתעניתו על שבלבל יצועי אביו:}%
% ---------- פרק 37 | פסוק 30 ----------
 \textbf{ל} וַיָּ֥שׇׁב אֶל־אֶחָ֖יו וַיֹּאמַ֑ר הַיֶּ֣לֶד אֵינֶ֔נּוּ וַאֲנִ֖י אָ֥נָה אֲנִי־בָֽא׃ \Rashi{\textbf{אנה אני בא.}אנה אברח מצערו של אבא?}%
% ---------- פרק 37 | פסוק 31 ----------
 \textbf{לא} וַיִּקְח֖וּ אֶת־כְּתֹ֣נֶת יוֹסֵ֑ף וַֽיִּשְׁחֲטוּ֙ שְׂעִ֣יר עִזִּ֔ים וַיִּטְבְּל֥וּ אֶת־הַכֻּתֹּ֖נֶת בַּדָּֽם׃ \Rashi{\textbf{שעיר עזים.}דמו דומה לשל אדם: הכתנת. זה שמה, וכשהיא דבוקה לתבה אחרת, כגון כתנת יוסף, כתנת פסים, כתנת בד, נקוד כתנת:}%
% ---------- פרק 37 | פסוק 32 ----------
 \textbf{לב} וַֽיְשַׁלְּח֞וּ אֶת־כְּתֹ֣נֶת הַפַּסִּ֗ים וַיָּבִ֙יאוּ֙ אֶל־אֲבִיהֶ֔ם וַיֹּאמְר֖וּ זֹ֣את מָצָ֑אנוּ הַכֶּר־נָ֗א הַכְּתֹ֧נֶת בִּנְךָ֛ הִ֖וא אִם־לֹֽא׃ %
% ---------- פרק 37 | פסוק 33 ----------
 \textbf{לג} וַיַּכִּירָ֤הּ וַיֹּ֙אמֶר֙ כְּתֹ֣נֶת בְּנִ֔י חַיָּ֥ה רָעָ֖ה אֲכָלָ֑תְהוּ טָרֹ֥ף טֹרַ֖ף יוֹסֵֽף׃ \Rashi{\textbf{ויאמר כתנת בני.}היא זו: חיה רעה אכלתהו. נצנצה בו רוח הקדש, סופו שתתגרה בו אשת פוטיפר. ולמה לא גלה לו הקב"ה? לפי שהחרימו וקללו את כל מי שיגלה, ושתפו להקב"ה עמהם (תנחומא), אבל יצחק היה יודע שהוא חי, אמר היאך אגלה והקב"ה אינו רוצה לגלות לו:}%
% ---------- פרק 37 | פסוק 34 ----------
 \textbf{לד} וַיִּקְרַ֤ע יַעֲקֹב֙ שִׂמְלֹתָ֔יו וַיָּ֥שֶׂם שַׂ֖ק בְּמׇתְנָ֑יו וַיִּתְאַבֵּ֥ל עַל־בְּנ֖וֹ יָמִ֥ים רַבִּֽים׃ \Rashi{\textbf{ימים רבים.}כ"ב שנה, משפרש ממנו עד שירד יעקב למצרים, שנאמר יוסף בן שבע עשרה שנה וגו', ובן שלשים שנה היה בעמדו לפני פרעה, ושבע שני השבע, ושנתים הרעב כשבא יעקב למצרים, הרי כ"ב שנה, כנגד כ"ב שנה שלא קים יעקב כבוד אב ואם (מגילה ט"ז) – כ' שנה שהיה בבית לבן, וב' שנה בדרך בשובו מבית לבן, שנה וחצי בסכות וששה חדשים בבית אל – וזהו שאמר ללבן זה לי עשרים שנה בביתך, לי הן, עלי, וסופי ללקות כנגדן:}%
% ---------- פרק 37 | פסוק 35 ----------
 \textbf{לה} וַיָּקֻ֩מוּ֩ כׇל־בָּנָ֨יו וְכׇל־בְּנֹתָ֜יו לְנַחֲמ֗וֹ וַיְמָאֵן֙ לְהִתְנַחֵ֔ם וַיֹּ֕אמֶר כִּֽי־אֵרֵ֧ד אֶל־בְּנִ֛י אָבֵ֖ל שְׁאֹ֑לָה וַיֵּ֥בְךְּ אֹת֖וֹ אָבִֽיו׃ \Rashi{\textbf{וכל בנתיו.}רבי יהודה אומר אחיות תאומות נולדו עם כל שבט ושבט ונשאום, רבי נחמיה אומר כנעניות היו, אלא מהו וכל בנתיו? – כלותיו, שאין אדם נמנע מלקרא לחתנו בנו ולכלתו בתו: וימאן להתנחם. אין אדם מקבל תנחומין על החי וסבור שמת, שעל המת נגזרה גזרה שישתכח מן הלב ולא על החי (בראשית רבה): ארד אל בני. כמו על בני, והרבה אל משמשין בלשון על, אל שאול ואל בית הדמים (שמואל ב כ"א), אל הלקח ארון האלהים ואל מות חמיה ואישה (שמואל א ד'): אבל שאלה. כפשוטו לשון קבר הוא – באבלי אקבר, ולא אתנחם כל ימי. ומדרשו, גיהנם; סימן זה היה מסור בידי מפי הגבורה, אם לא ימות אחד מבני בחיי, מבטח אני שאיני רואה גיהנם: ויבך אתו אביו. יצחק, בוכה היה מפני צרתו של יעקב אבל לא היה מתאבל, שהיה יודע שהוא חי:}%
% ---------- פרק 37 | פסוק 36 ----------
 \textbf{לו} וְהַ֨מְּדָנִ֔ים מָכְר֥וּ אֹת֖וֹ אֶל־מִצְרָ֑יִם לְפֽוֹטִיפַר֙ סְרִ֣יס פַּרְעֹ֔ה שַׂ֖ר הַטַּבָּחִֽים׃ \{פ\} \Rashi{\textbf{הטבחים.}שוחטי בהמות המלך:}%
% ---------- פרק 38 | פסוק 1 ----------
 \textbf{א} וַֽיְהִי֙ בָּעֵ֣ת הַהִ֔וא וַיֵּ֥רֶד יְהוּדָ֖ה מֵאֵ֣ת אֶחָ֑יו וַיֵּ֛ט עַד־אִ֥ישׁ עֲדֻלָּמִ֖י וּשְׁמ֥וֹ חִירָֽה׃ \Rashi{\textbf{ויהי בעת ההוא.}למה נסמכה פרשה זו לכאן, והפסיק בפרשתו של יוסף? ללמד שהורידוהו אחיו מגדלתו כשראו בצרת אביהם, אמרו: אתה אמרת למכרו, אלו אמרת להשיבו היינו שומעים לך: ויט. מאת אחיו: עד איש עדלמי. נשתתף עמו:}%
% ---------- פרק 38 | פסוק 2 ----------
 \textbf{ב} וַיַּרְא־שָׁ֧ם יְהוּדָ֛ה בַּת־אִ֥ישׁ כְּנַעֲנִ֖י וּשְׁמ֣וֹ שׁ֑וּעַ וַיִּקָּחֶ֖הָ וַיָּבֹ֥א אֵלֶֽיהָ׃ \Rashi{\textbf{כנעני.}תגרא (בראשית רבה):}%
% ---------- פרק 38 | פסוק 3 ----------
 \textbf{ג} וַתַּ֖הַר וַתֵּ֣לֶד בֵּ֑ן וַיִּקְרָ֥א אֶת־שְׁמ֖וֹ עֵֽר׃ %
% ---------- פרק 38 | פסוק 4 ----------
 \textbf{ד} וַתַּ֥הַר ע֖וֹד וַתֵּ֣לֶד בֵּ֑ן וַתִּקְרָ֥א אֶת־שְׁמ֖וֹ אוֹנָֽן׃ %
% ---------- פרק 38 | פסוק 5 ----------
 \textbf{ה} וַתֹּ֤סֶף עוֹד֙ וַתֵּ֣לֶד בֵּ֔ן וַתִּקְרָ֥א אֶת־שְׁמ֖וֹ שֵׁלָ֑ה וְהָיָ֥ה בִכְזִ֖יב בְּלִדְתָּ֥הּ אֹתֽוֹ׃ \Rashi{\textbf{והיה בכזיב.}שם המקום. ואומר אני על שם שפסקה מלדת נקרא כזיב, לשון היו תהיה לי כמו אכזב (ירמיהו ט"ו), אשר לא יכזבו מימיו (ישעיהו נ"ח), ואם לא כן מה בא להודיענו? ובב"ר ראיתי, ותקרא שמו שלה וגומר פסקת:}%
% ---------- פרק 38 | פסוק 6 ----------
 \textbf{ו} וַיִּקַּ֧ח יְהוּדָ֛ה אִשָּׁ֖ה לְעֵ֣ר בְּכוֹר֑וֹ וּשְׁמָ֖הּ תָּמָֽר׃ %
% ---------- פרק 38 | פסוק 7 ----------
 \textbf{ז} וַיְהִ֗י עֵ֚ר בְּכ֣וֹר יְהוּדָ֔ה רַ֖ע בְּעֵינֵ֣י יְהֹוָ֑ה וַיְמִתֵ֖הוּ יְהֹוָֽה׃ \Rashi{\textbf{רע בעיני ה'.}כרעתו של אונן, משחית זרעו, שנאמר באונן וימת גם אתו, כמיתתו של ער מיתתו של אונן. ולמה היה ער משחית זרעו? כדי שלא תתעבר ויכחיש יפיה:}%
% ---------- פרק 38 | פסוק 8 ----------
 \textbf{ח} וַיֹּ֤אמֶר יְהוּדָה֙ לְאוֹנָ֔ן בֹּ֛א אֶל־אֵ֥שֶׁת אָחִ֖יךָ וְיַבֵּ֣ם אֹתָ֑הּ וְהָקֵ֥ם זֶ֖רַע לְאָחִֽיךָ׃ \Rashi{\textbf{והקם זרע.}הבן יקרא על שם המת:}%
% ---------- פרק 38 | פסוק 9 ----------
 \textbf{ט} וַיֵּ֣דַע אוֹנָ֔ן כִּ֛י לֹּ֥א ל֖וֹ יִהְיֶ֣ה הַזָּ֑רַע וְהָיָ֞ה אִם־בָּ֨א אֶל־אֵ֤שֶׁת אָחִיו֙ וְשִׁחֵ֣ת אַ֔רְצָה לְבִלְתִּ֥י נְתׇן־זֶ֖רַע לְאָחִֽיו׃ \Rashi{\textbf{ושחת ארצה.}דש מבפנים וזורה מבחוץ (בראשית רבה):}%
% ---------- פרק 38 | פסוק 10 ----------
 \textbf{י} וַיֵּ֛רַע בְּעֵינֵ֥י יְהֹוָ֖ה אֲשֶׁ֣ר עָשָׂ֑ה וַיָּ֖מֶת גַּם־אֹתֽוֹ׃ %
% ---------- פרק 38 | פסוק 11 ----------
 \textbf{יא} וַיֹּ֣אמֶר יְהוּדָה֩ לְתָמָ֨ר כַּלָּת֜וֹ שְׁבִ֧י אַלְמָנָ֣ה בֵית־אָבִ֗יךְ עַד־יִגְדַּל֙ שֵׁלָ֣ה בְנִ֔י כִּ֣י אָמַ֔ר פֶּן־יָמ֥וּת גַּם־ה֖וּא כְּאֶחָ֑יו וַתֵּ֣לֶךְ תָּמָ֔ר וַתֵּ֖שֶׁב בֵּ֥ית אָבִֽיהָ׃ \Rashi{\textbf{כי אמר וגו'.}כלומר, דוחה היה אותה בקש, שלא היה בדעתו להשיאה לו: כי אמר פן ימות. מחזקת היא זו שימותו אנשיה (כתובות מ"ג):}%
% ---------- פרק 38 | פסוק 12 ----------
 \textbf{יב} וַיִּרְבּוּ֙ הַיָּמִ֔ים וַתָּ֖מׇת בַּת־שׁ֣וּעַ אֵֽשֶׁת־יְהוּדָ֑ה וַיִּנָּ֣חֶם יְהוּדָ֗ה וַיַּ֜עַל עַל־גֹּֽזְזֵ֤י צֹאנוֹ֙ ה֗וּא וְחִירָ֛ה רֵעֵ֥הוּ הָעֲדֻלָּמִ֖י תִּמְנָֽתָה׃ \Rashi{\textbf{ויעל על גוזזי צאנו.}ויעל תמנתה לעמד על גוזזי צאנו:}%
% ---------- פרק 38 | פסוק 13 ----------
 \textbf{יג} וַיֻּגַּ֥ד לְתָמָ֖ר לֵאמֹ֑ר הִנֵּ֥ה חָמִ֛יךְ עֹלֶ֥ה תִמְנָ֖תָה לָגֹ֥ז צֹאנֽוֹ׃ \Rashi{\textbf{עלה תמנתה.}ובשמשון הוא אומר וירד שמשון תמנתה (שופטים י"ד)? בשפוע ההר היתה יושבת, עולין לה מכאן ויורדין לה מכאן:}%
% ---------- פרק 38 | פסוק 14 ----------
 \textbf{יד} וַתָּ֩סַר֩ בִּגְדֵ֨י אַלְמְנוּתָ֜הּ מֵֽעָלֶ֗יהָ וַתְּכַ֤ס בַּצָּעִיף֙ וַתִּתְעַלָּ֔ף וַתֵּ֙שֶׁב֙ בְּפֶ֣תַח עֵינַ֔יִם אֲשֶׁ֖ר עַל־דֶּ֣רֶךְ תִּמְנָ֑תָה כִּ֤י רָאֲתָה֙ כִּֽי־גָדַ֣ל שֵׁלָ֔ה וְהִ֕וא לֹֽא־נִתְּנָ֥ה ל֖וֹ לְאִשָּֽׁה׃ \Rashi{\textbf{ותתעלף.}כסתה פניה, שלא יכיר בה: ותשב בפתח עינים. בפתיחת עינים, בפרשת דרכים שעל דרך תמנתה. ורבותינו דרשו בפתחו של אברהם אבינו, שכל עינים מצפות לראותו (סוטה י'): כי ראתה כי גדל שלה וגו'. לפיכך הפקירה עצמה אצל יהודה, שהיתה מתאוה להעמיד ממנו בנים:}%
% ---------- פרק 38 | פסוק 15 ----------
 \textbf{טו} וַיִּרְאֶ֣הָ יְהוּדָ֔ה וַֽיַּחְשְׁבֶ֖הָ לְזוֹנָ֑ה כִּ֥י כִסְּתָ֖ה פָּנֶֽיהָ׃ \Rashi{\textbf{ויחשבה לזונה.}לפי שיושבת בפרשת דרכים: כי כסתה פניה. ולא יכול לראותה ולהכירה. ומדרש רבותינו כי כסתה פניה, כשהיתה בבית חמיה היתה צנועה, לפיכך לא חשדה:}%
% ---------- פרק 38 | פסוק 16 ----------
 \textbf{טז} וַיֵּ֨ט אֵלֶ֜יהָ אֶל־הַדֶּ֗רֶךְ וַיֹּ֙אמֶר֙ הָֽבָה־נָּא֙ אָב֣וֹא אֵלַ֔יִךְ כִּ֚י לֹ֣א יָדַ֔ע כִּ֥י כַלָּת֖וֹ הִ֑וא וַתֹּ֙אמֶר֙ מַה־תִּתֶּן־לִ֔י כִּ֥י תָב֖וֹא אֵלָֽי׃ \Rashi{\textbf{ויט אליה אל הדרך.}מדרך שהיה בה, נטה אל הדרך אשר היא בה, ובלשון לעז דשטורנ"יר: הבה נא. הכיני עצמך ודעתך לכך. כל לשון הבה לשון הזמנה הוא, חוץ ממקום שיש לתרגמו בלשון נתינה, ואף אותן של הזמנה קרובים ללשון נתינה הם:}%
% ---------- פרק 38 | פסוק 17 ----------
 \textbf{יז} וַיֹּ֕אמֶר אָנֹכִ֛י אֲשַׁלַּ֥ח גְּדִֽי־עִזִּ֖ים מִן־הַצֹּ֑אן וַתֹּ֕אמֶר אִם־תִּתֵּ֥ן עֵרָב֖וֹן עַ֥ד שׇׁלְחֶֽךָ׃ \Rashi{\textbf{ערבון.}משכון:}%
% ---------- פרק 38 | פסוק 18 ----------
 \textbf{יח} וַיֹּ֗אמֶר מָ֣ה הָעֵֽרָבוֹן֮ אֲשֶׁ֣ר אֶתֶּן־לָךְ֒ וַתֹּ֗אמֶר חֹתָֽמְךָ֙ וּפְתִילֶ֔ךָ וּמַטְּךָ֖ אֲשֶׁ֣ר בְּיָדֶ֑ךָ וַיִּתֶּן־לָ֛הּ וַיָּבֹ֥א אֵלֶ֖יהָ וַתַּ֥הַר לֽוֹ׃ \Rashi{\textbf{חתמך ופתילך.}עזקתך ושושיפך – טבעת שאתה חותם בה ושמלתך שאתה מתכסה בה: ותהר לו. גבורים כיוצא בו, צדיקים כיוצא בו.}%
% ---------- פרק 38 | פסוק 19 ----------
 \textbf{יט} וַתָּ֣קׇם וַתֵּ֔לֶךְ וַתָּ֥סַר צְעִיפָ֖הּ מֵעָלֶ֑יהָ וַתִּלְבַּ֖שׁ בִּגְדֵ֥י אַלְמְנוּתָֽהּ׃ %
% ---------- פרק 38 | פסוק 20 ----------
 \textbf{כ} וַיִּשְׁלַ֨ח יְהוּדָ֜ה אֶת־גְּדִ֣י הָֽעִזִּ֗ים בְּיַד֙ רֵעֵ֣הוּ הָֽעֲדֻלָּמִ֔י לָקַ֥חַת הָעֵרָב֖וֹן מִיַּ֣ד הָאִשָּׁ֑ה וְלֹ֖א מְצָאָֽהּ׃ %
% ---------- פרק 38 | פסוק 21 ----------
 \textbf{כא} וַיִּשְׁאַ֞ל אֶת־אַנְשֵׁ֤י מְקֹמָהּ֙ לֵאמֹ֔ר אַיֵּ֧ה הַקְּדֵשָׁ֛ה הִ֥וא בָעֵינַ֖יִם עַל־הַדָּ֑רֶךְ וַיֹּ֣אמְר֔וּ לֹא־הָיְתָ֥ה בָזֶ֖ה קְדֵשָֽׁה׃ \Rashi{\textbf{הקדשה.}מקדשת ומזמנת לזנות:}%
% ---------- פרק 38 | פסוק 22 ----------
 \textbf{כב} וַיָּ֙שׇׁב֙ אֶל־יְהוּדָ֔ה וַיֹּ֖אמֶר לֹ֣א מְצָאתִ֑יהָ וְגַ֨ם אַנְשֵׁ֤י הַמָּקוֹם֙ אָֽמְר֔וּ לֹא־הָיְתָ֥ה בָזֶ֖ה קְדֵשָֽׁה׃ %
% ---------- פרק 38 | פסוק 23 ----------
 \textbf{כג} וַיֹּ֤אמֶר יְהוּדָה֙ תִּֽקַּֽח־לָ֔הּ פֶּ֖ן נִהְיֶ֣ה לָב֑וּז הִנֵּ֤ה שָׁלַ֙חְתִּי֙ הַגְּדִ֣י הַזֶּ֔ה וְאַתָּ֖ה לֹ֥א מְצָאתָֽהּ׃ \Rashi{\textbf{תקח לה.}יהיה שלה מה שבידה: פן נהיה לבוז. אם תבקשנה עוד יתפרסם הדבר ויהיה גנאי, כי מה עלי לעשות עוד לאמת דברי: הנה שלחתי הגדי הזה. לפי שרמה יהודה את אביו בגדי עזים שהטביל כתנת יוסף בדמו, רמוהו גם הוא בגדי עזים (בראשית רבה):}%
% ---------- פרק 38 | פסוק 24 ----------
 \textbf{כד} וַיְהִ֣י ׀ כְּמִשְׁלֹ֣שׁ חֳדָשִׁ֗ים וַיֻּגַּ֨ד לִֽיהוּדָ֤ה לֵאמֹר֙ זָֽנְתָה֙ תָּמָ֣ר כַּלָּתֶ֔ךָ וְגַ֛ם הִנֵּ֥ה הָרָ֖ה לִזְנוּנִ֑ים וַיֹּ֣אמֶר יְהוּדָ֔ה הוֹצִיא֖וּהָ וְתִשָּׂרֵֽף׃ \Rashi{\textbf{כמשלש חדשים.}רבו של ראשון ורבו של אחרון ואמצעי שלם (נדה מ"נ). ולשון כמשלש חדשים כהשתלש החדשים, כמו ומשלח מנות (אסתר ט'), משלוח ידם (ישעיהו י"א), וכן תרגם אנקלוס כתלתות ירחיא: הרה לזנונים. שם דבר, מעברת כמו אשה הרה, וכמו ברה כחמה: ותשרף. אמר אפרים מקשאה משום רבי מאיר, בתו של שם היתה, שהוא כהן, לפיכך דנוה בשרפה:}%
% ---------- פרק 38 | פסוק 25 ----------
 \textbf{כה} הִ֣וא מוּצֵ֗את וְהִ֨יא שָׁלְחָ֤ה אֶל־חָמִ֙יהָ֙ לֵאמֹ֔ר לְאִישׁ֙ אֲשֶׁר־אֵ֣לֶּה לּ֔וֹ אָנֹכִ֖י הָרָ֑ה וַתֹּ֙אמֶר֙ הַכֶּר־נָ֔א לְמִ֞י הַחֹתֶ֧מֶת וְהַפְּתִילִ֛ים וְהַמַּטֶּ֖ה הָאֵֽלֶּה׃ \Rashi{\textbf{הוא מוצאת.}להשרף: והיא שלחה אל חמיה. לא רצתה להלבין פניו ולומר ממך אני מעברת, אלא לאיש אשר אלה לו, אמרה אם יודה יודה מעצמו, ואם לאו, ישרפוני, ואל אלבין פניו. מכאן אמרו נוח לו לאדם שיפילוהו לכבשן האש ואל ילבין פני חברו ברבים: הכר נא. אין נא אלא לשון בקשה – הכר נא בוראך, ואל תאבד שלש נפשות:}%
% ---------- פרק 38 | פסוק 26 ----------
 \textbf{כו} וַיַּכֵּ֣ר יְהוּדָ֗ה וַיֹּ֙אמֶר֙ צָֽדְקָ֣ה מִמֶּ֔נִּי כִּֽי־עַל־כֵּ֥ן לֹא־נְתַתִּ֖יהָ לְשֵׁלָ֣ה בְנִ֑י וְלֹֽא־יָסַ֥ף ע֖וֹד לְדַעְתָּֽהּ׃ \Rashi{\textbf{צדקה.}בדבריה: ממני. היא מעברת. ורבותינו ז"ל דרשו שיצאה בת קול ואמרה ממני ומאתי יצאו הדברים, לפי שהיתה צנועה בבית חמיה, גזרתי שיצאו ממנה מלכים, ומשבט יהודה גזרתי להעמיד מלכים בישראל: כי על כן לא נתתיה. כי בדין עשתה, על אשר לא נתתיה לשלה בני: ולא יסף עוד. יש אומרים לא הוסיף, ויש אומרים לא פסק. וחברו גבי אלדד ומידד, ולא יספו, ומתרגמינן ולא פסקו:}%
% ---------- פרק 38 | פסוק 27 ----------
 \textbf{כז} וַיְהִ֖י בְּעֵ֣ת לִדְתָּ֑הּ וְהִנֵּ֥ה תְאוֹמִ֖ים בְּבִטְנָֽהּ׃ \Rashi{\textbf{בעת לדתה.}וברבקה הוא אומר וימלאו ימיה ללדת? להלן למלאים וכאן לחסרים: והנה תאומים. מלא, ולהלן תומם חסר, לפי שהאחד רשע, אבל אלו שניהם צדיקים:}%
% ---------- פרק 38 | פסוק 28 ----------
 \textbf{כח} וַיְהִ֥י בְלִדְתָּ֖הּ וַיִּתֶּן־יָ֑ד וַתִּקַּ֣ח הַמְיַלֶּ֗דֶת וַתִּקְשֹׁ֨ר עַל־יָד֤וֹ שָׁנִי֙ לֵאמֹ֔ר זֶ֖ה יָצָ֥א רִאשֹׁנָֽה׃ \Rashi{\textbf{ויתן יד.}הוציא האחד ידו לחוץ, ולאחר שקשרה על ידו השני החזירה:}%
% ---------- פרק 38 | פסוק 29 ----------
 \textbf{כט} וַיְהִ֣י ׀ כְּמֵשִׁ֣יב יָד֗וֹ וְהִנֵּה֙ יָצָ֣א אָחִ֔יו וַתֹּ֕אמֶר מַה־פָּרַ֖צְתָּ עָלֶ֣יךָ פָּ֑רֶץ וַיִּקְרָ֥א שְׁמ֖וֹ פָּֽרֶץ׃ \Rashi{\textbf{פרצת.}חזקת עליך חזק:}%
% ---------- פרק 38 | פסוק 30 ----------
 \textbf{ל} וְאַחַר֙ יָצָ֣א אָחִ֔יו אֲשֶׁ֥ר עַל־יָד֖וֹ הַשָּׁנִ֑י וַיִּקְרָ֥א שְׁמ֖וֹ זָֽרַח׃ \{ס\} %
% ---------- פרק 39 | פסוק 1 ----------
 \textbf{א} וְיוֹסֵ֖ף הוּרַ֣ד מִצְרָ֑יְמָה וַיִּקְנֵ֡הוּ פּוֹטִיפַר֩ סְרִ֨יס פַּרְעֹ֜ה שַׂ֤ר הַטַּבָּחִים֙ אִ֣ישׁ מִצְרִ֔י מִיַּד֙ הַיִּשְׁמְעֵאלִ֔ים אֲשֶׁ֥ר הוֹרִדֻ֖הוּ שָֽׁמָּה׃ %
% ---------- פרק 39 | פסוק 2 ----------
 \textbf{ב} וַיְהִ֤י יְהֹוָה֙ אֶת־יוֹסֵ֔ף וַיְהִ֖י אִ֣ישׁ מַצְלִ֑יחַ וַיְהִ֕י בְּבֵ֖ית אֲדֹנָ֥יו הַמִּצְרִֽי׃ %
% ---------- פרק 39 | פסוק 3 ----------
 \textbf{ג} וַיַּ֣רְא אֲדֹנָ֔יו כִּ֥י יְהֹוָ֖ה אִתּ֑וֹ וְכֹל֙ אֲשֶׁר־ה֣וּא עֹשֶׂ֔ה יְהֹוָ֖ה מַצְלִ֥יחַ בְּיָדֽוֹ׃ \Rashi{\textbf{כי ה' אתו.}שם שמים שגור בפיו:}%
% ---------- פרק 39 | פסוק 4 ----------
 \textbf{ד} וַיִּמְצָ֨א יוֹסֵ֥ף חֵ֛ן בְּעֵינָ֖יו וַיְשָׁ֣רֶת אֹת֑וֹ וַיַּפְקִדֵ֙הוּ֙ עַל־בֵּית֔וֹ וְכׇל־יֶשׁ־ל֖וֹ נָתַ֥ן בְּיָדֽוֹ׃ \Rashi{\textbf{וכל יש לו.}הרי לשון קצר, חסר "אשר":}%
% ---------- פרק 39 | פסוק 5 ----------
 \textbf{ה} וַיְהִ֡י מֵאָז֩ הִפְקִ֨יד אֹת֜וֹ בְּבֵית֗וֹ וְעַל֙ כׇּל־אֲשֶׁ֣ר יֶשׁ־ל֔וֹ וַיְבָ֧רֶךְ יְהֹוָ֛ה אֶת־בֵּ֥ית הַמִּצְרִ֖י בִּגְלַ֣ל יוֹסֵ֑ף וַיְהִ֞י בִּרְכַּ֤ת יְהֹוָה֙ בְּכׇל־אֲשֶׁ֣ר יֶשׁ־ל֔וֹ בַּבַּ֖יִת וּבַשָּׂדֶֽה׃ %
% ---------- פרק 39 | פסוק 6 ----------
 \textbf{ו} וַיַּעֲזֹ֣ב כׇּל־אֲשֶׁר־לוֹ֮ בְּיַד־יוֹסֵף֒ וְלֹא־יָדַ֤ע אִתּוֹ֙ מְא֔וּמָה כִּ֥י אִם־הַלֶּ֖חֶם אֲשֶׁר־ה֣וּא אוֹכֵ֑ל וַיְהִ֣י יוֹסֵ֔ף יְפֵה־תֹ֖אַר וִיפֵ֥ה מַרְאֶֽה׃ %
% ---------- פרק 39 | פסוק 7 ----------
 \textbf{ז} וַיְהִ֗י אַחַר֙ הַדְּבָרִ֣ים הָאֵ֔לֶּה וַתִּשָּׂ֧א אֵֽשֶׁת־אֲדֹנָ֛יו אֶת־עֵינֶ֖יהָ אֶל־יוֹסֵ֑ף וַתֹּ֖אמֶר שִׁכְבָ֥ה עִמִּֽי׃ %
% ---------- פרק 39 | פסוק 8 ----------
 \textbf{ח} וַיְמָאֵ֓ן ׀ וַיֹּ֙אמֶר֙ אֶל־אֵ֣שֶׁת אֲדֹנָ֔יו הֵ֣ן אֲדֹנִ֔י לֹא־יָדַ֥ע אִתִּ֖י מַה־בַּבָּ֑יִת וְכֹ֥ל אֲשֶׁר־יֶשׁ־ל֖וֹ נָתַ֥ן בְּיָדִֽי׃ %
% ---------- פרק 39 | פסוק 9 ----------
 \textbf{ט} אֵינֶ֨נּוּ גָד֜וֹל בַּבַּ֣יִת הַזֶּה֮ מִמֶּ֒נִּי֒ וְלֹֽא־חָשַׂ֤ךְ מִמֶּ֙נִּי֙ מְא֔וּמָה כִּ֥י אִם־אוֹתָ֖ךְ בַּאֲשֶׁ֣ר אַתְּ־אִשְׁתּ֑וֹ וְאֵ֨יךְ אֶֽעֱשֶׂ֜ה הָרָעָ֤ה הַגְּדֹלָה֙ הַזֹּ֔את וְחָטָ֖אתִי לֵֽאלֹהִֽים׃ \Rashi{\textbf{וחטאתי לאלהים.}בני נח נצטוו על העריות:}%
% ---------- פרק 39 | פסוק 10 ----------
 \textbf{י} וַיְהִ֕י כְּדַבְּרָ֥הּ אֶל־יוֹסֵ֖ף י֣וֹם ׀ י֑וֹם וְלֹא־שָׁמַ֥ע אֵלֶ֛יהָ לִשְׁכַּ֥ב אֶצְלָ֖הּ לִהְי֥וֹת עִמָּֽהּ׃ %
% ---------- פרק 39 | פסוק 11 ----------
 \textbf{יא} וַֽיְהִי֙ כְּהַיּ֣וֹם הַזֶּ֔ה וַיָּבֹ֥א הַבַּ֖יְתָה לַעֲשׂ֣וֹת מְלַאכְתּ֑וֹ וְאֵ֨ין אִ֜ישׁ מֵאַנְשֵׁ֥י הַבַּ֛יִת שָׁ֖ם בַּבָּֽיִת׃ \Rashi{\textbf{ויהי כהיום הזה.}כלומר ויהי כאשר הגיע יום מיחד, יום צחוק, יום איד שלהם, שהלכו כלם לבית ע"ז, אמרה אין לי יום הגון להזקק ליוסף כהיום הזה, אמרה להם חולה אני ואיני יכולה לילך (תנחומא): לעשות מלאכתו. רב ושמואל, חד אמר מלאכתו ממש, וחד אמר לעשות צרכיו עמה, אלא שנראית לו דמות דיוקנו של אביו וכו' כדאיתא במסכת סוטה (דף ל"ז):}%
% ---------- פרק 39 | פסוק 12 ----------
 \textbf{יב} וַתִּתְפְּשֵׂ֧הוּ בְּבִגְד֛וֹ לֵאמֹ֖ר שִׁכְבָ֣ה עִמִּ֑י וַיַּעֲזֹ֤ב בִּגְדוֹ֙ בְּיָדָ֔הּ וַיָּ֖נׇס וַיֵּצֵ֥א הַחֽוּצָה׃ %
% ---------- פרק 39 | פסוק 13 ----------
 \textbf{יג} וַֽיְהִי֙ כִּרְאוֹתָ֔הּ כִּֽי־עָזַ֥ב בִּגְד֖וֹ בְּיָדָ֑הּ וַיָּ֖נׇס הַחֽוּצָה׃ %
% ---------- פרק 39 | פסוק 14 ----------
 \textbf{יד} וַתִּקְרָ֞א לְאַנְשֵׁ֣י בֵיתָ֗הּ וַתֹּ֤אמֶר לָהֶם֙ לֵאמֹ֔ר רְא֗וּ הֵ֥בִיא לָ֛נוּ אִ֥ישׁ עִבְרִ֖י לְצַ֣חֶק בָּ֑נוּ בָּ֤א אֵלַי֙ לִשְׁכַּ֣ב עִמִּ֔י וָאֶקְרָ֖א בְּק֥וֹל גָּדֽוֹל׃ \Rashi{\textbf{ראו הביא לנו.}הרי זה לשון קצרה – הביא לנו, ולא פרש מי הביאו, ועל בעלה אומרת כן: עברי. מעבר הנהר, מבני עבר (בראשית רבה):}%
% ---------- פרק 39 | פסוק 15 ----------
 \textbf{טו} וַיְהִ֣י כְשׇׁמְע֔וֹ כִּֽי־הֲרִימֹ֥תִי קוֹלִ֖י וָאֶקְרָ֑א וַיַּעֲזֹ֤ב בִּגְדוֹ֙ אֶצְלִ֔י וַיָּ֖נׇס וַיֵּצֵ֥א הַחֽוּצָה׃ %
% ---------- פרק 39 | פסוק 16 ----------
 \textbf{טז} וַתַּנַּ֥ח בִּגְד֖וֹ אֶצְלָ֑הּ עַד־בּ֥וֹא אֲדֹנָ֖יו אֶל־בֵּיתֽוֹ׃ \Rashi{\textbf{אדניו.}של יוסף:}%
% ---------- פרק 39 | פסוק 17 ----------
 \textbf{יז} וַתְּדַבֵּ֣ר אֵלָ֔יו כַּדְּבָרִ֥ים הָאֵ֖לֶּה לֵאמֹ֑ר בָּֽא־אֵלַ֞י הָעֶ֧בֶד הָֽעִבְרִ֛י אֲשֶׁר־הֵבֵ֥אתָ לָּ֖נוּ לְצַ֥חֶק בִּֽי׃ \Rashi{\textbf{בא אלי.}לצחק בי העבד העברי אשר הבאת לנו:}%
% ---------- פרק 39 | פסוק 18 ----------
 \textbf{יח} וַיְהִ֕י כַּהֲרִימִ֥י קוֹלִ֖י וָאֶקְרָ֑א וַיַּעֲזֹ֥ב בִּגְד֛וֹ אֶצְלִ֖י וַיָּ֥נׇס הַחֽוּצָה׃ %
% ---------- פרק 39 | פסוק 19 ----------
 \textbf{יט} וַיְהִי֩ כִשְׁמֹ֨עַ אֲדֹנָ֜יו אֶת־דִּבְרֵ֣י אִשְׁתּ֗וֹ אֲשֶׁ֨ר דִּבְּרָ֤ה אֵלָיו֙ לֵאמֹ֔ר כַּדְּבָרִ֣ים הָאֵ֔לֶּה עָ֥שָׂה לִ֖י עַבְדֶּ֑ךָ וַיִּ֖חַר אַפּֽוֹ׃ \Rashi{\textbf{ויהי כשמע אדניו וגו'.}בשעת תשמיש אמרה לו כן, וזהו שאמרה כדברים האלה עשה לי עבדך, עניני תשמיש כאלה:}%
% ---------- פרק 39 | פסוק 20 ----------
 \textbf{כ} וַיִּקַּח֩ אֲדֹנֵ֨י יוֹסֵ֜ף אֹת֗וֹ וַֽיִּתְּנֵ֙הוּ֙ אֶל־בֵּ֣ית הַסֹּ֔הַר מְק֕וֹם אֲשֶׁר־[אֲסִירֵ֥י] (אסורי) הַמֶּ֖לֶךְ אֲסוּרִ֑ים וַֽיְהִי־שָׁ֖ם בְּבֵ֥ית הַסֹּֽהַר׃ %
% ---------- פרק 39 | פסוק 21 ----------
 \textbf{כא} וַיְהִ֤י יְהֹוָה֙ אֶת־יוֹסֵ֔ף וַיֵּ֥ט אֵלָ֖יו חָ֑סֶד וַיִּתֵּ֣ן חִנּ֔וֹ בְּעֵינֵ֖י שַׂ֥ר בֵּית־הַסֹּֽהַר׃ \Rashi{\textbf{ויט אליו חסד.}שהיה מקבל לכל רואיו, לשון "כלה נאה וחסודה" שבמשנה:}%
% ---------- פרק 39 | פסוק 22 ----------
 \textbf{כב} וַיִּתֵּ֞ן שַׂ֤ר בֵּית־הַסֹּ֙הַר֙ בְּיַד־יוֹסֵ֔ף אֵ֚ת כׇּל־הָ֣אֲסִירִ֔ם אֲשֶׁ֖ר בְּבֵ֣ית הַסֹּ֑הַר וְאֵ֨ת כׇּל־אֲשֶׁ֤ר עֹשִׂים֙ שָׁ֔ם ה֖וּא הָיָ֥ה עֹשֶֽׂה׃ \Rashi{\textbf{הוא היה עושה.}כתרגומו במימרה הוה מתעביד:}%
% ---------- פרק 39 | פסוק 23 ----------
 \textbf{כג} אֵ֣ין ׀ שַׂ֣ר בֵּית־הַסֹּ֗הַר רֹאֶ֤ה אֶֽת־כׇּל־מְא֙וּמָה֙ בְּיָד֔וֹ בַּאֲשֶׁ֥ר יְהֹוָ֖ה אִתּ֑וֹ וַֽאֲשֶׁר־ה֥וּא עֹשֶׂ֖ה יְהֹוָ֥ה מַצְלִֽיחַ׃ \{פ\} \Rashi{\textbf{באשר ה' אתו.}בשביל שה' אתו:}%
% ---------- פרק 40 | פסוק 1 ----------
 \textbf{א} וַיְהִ֗י אַחַר֙ הַדְּבָרִ֣ים הָאֵ֔לֶּה חָ֥טְא֛וּ מַשְׁקֵ֥ה מֶֽלֶךְ־מִצְרַ֖יִם וְהָאֹפֶ֑ה לַאֲדֹנֵיהֶ֖ם לְמֶ֥לֶךְ מִצְרָֽיִם׃ \Rashi{\textbf{אחר הדברים האלה.}לפי שהרגילה אותה ארורה את הצדיק בפי כלם לדבר בו ולגנותו, הביא להם הקב"ה סרחנם של אלו, שיפנו אליהם ולא אליו, ועוד שתבא הרוחה לצדיק על ידיהם: חטאו. זה נמצא זבוב בפיילי פוטירין שלו, וזה נמצא צרור בגלוסקין שלו: והאפה. את פת המלך, ואין לשון אפיה אלא בפת, ובלעז פישטו"ר:}%
% ---------- פרק 40 | פסוק 2 ----------
 \textbf{ב} וַיִּקְצֹ֣ף פַּרְעֹ֔ה עַ֖ל שְׁנֵ֣י סָרִיסָ֑יו עַ֚ל שַׂ֣ר הַמַּשְׁקִ֔ים וְעַ֖ל שַׂ֥ר הָאוֹפִֽים׃ %
% ---------- פרק 40 | פסוק 3 ----------
 \textbf{ג} וַיִּתֵּ֨ן אֹתָ֜ם בְּמִשְׁמַ֗ר בֵּ֛ית שַׂ֥ר הַטַּבָּחִ֖ים אֶל־בֵּ֣ית הַסֹּ֑הַר מְק֕וֹם אֲשֶׁ֥ר יוֹסֵ֖ף אָס֥וּר שָֽׁם׃ %
% ---------- פרק 40 | פסוק 4 ----------
 \textbf{ד} וַ֠יִּפְקֹ֠ד שַׂ֣ר הַטַּבָּחִ֧ים אֶת־יוֹסֵ֛ף אִתָּ֖ם וַיְשָׁ֣רֶת אֹתָ֑ם וַיִּהְי֥וּ יָמִ֖ים בְּמִשְׁמָֽר׃ \Rashi{\textbf{ויפקד שר הטבחים את יוסף.}להיות אתם: ויהיו ימים במשמר. שנים עשר חדש:}%
% ---------- פרק 40 | פסוק 5 ----------
 \textbf{ה} וַיַּֽחַלְמוּ֩ חֲל֨וֹם שְׁנֵיהֶ֜ם אִ֤ישׁ חֲלֹמוֹ֙ בְּלַ֣יְלָה אֶחָ֔ד אִ֖ישׁ כְּפִתְר֣וֹן חֲלֹמ֑וֹ הַמַּשְׁקֶ֣ה וְהָאֹפֶ֗ה אֲשֶׁר֙ לְמֶ֣לֶךְ מִצְרַ֔יִם אֲשֶׁ֥ר אֲסוּרִ֖ים בְּבֵ֥ית הַסֹּֽהַר׃ \Rashi{\textbf{ויחלמו חלום שניהם.}ויחלמו שניהם חלום, זהו פשוטו, ומדרשו כל א' חלם חלום שניהם – שחלם את חלומו ופתרון חברו, וזהו שנאמר וירא שר האפים כי טוב פתר: איש כפתרון חלומו. כל אחד חלם חלום הדומה לפתרון העתיד לבא עליהם:}%
% ---------- פרק 40 | פסוק 6 ----------
 \textbf{ו} וַיָּבֹ֧א אֲלֵיהֶ֛ם יוֹסֵ֖ף בַּבֹּ֑קֶר וַיַּ֣רְא אֹתָ֔ם וְהִנָּ֖ם זֹעֲפִֽים׃ \Rashi{\textbf{זעפים.}עצבים כמו סר וזעף (מלכים א כ'), זעף ה' אשא (מיכה ז'):}%
% ---------- פרק 40 | פסוק 7 ----------
 \textbf{ז} וַיִּשְׁאַ֞ל אֶת־סְרִיסֵ֣י פַרְעֹ֗ה אֲשֶׁ֨ר אִתּ֧וֹ בְמִשְׁמַ֛ר בֵּ֥ית אֲדֹנָ֖יו לֵאמֹ֑ר מַדּ֛וּעַ פְּנֵיכֶ֥ם רָעִ֖ים הַיּֽוֹם׃ %
% ---------- פרק 40 | פסוק 8 ----------
 \textbf{ח} וַיֹּאמְר֣וּ אֵלָ֔יו חֲל֣וֹם חָלַ֔מְנוּ וּפֹתֵ֖ר אֵ֣ין אֹת֑וֹ וַיֹּ֨אמֶר אֲלֵהֶ֜ם יוֹסֵ֗ף הֲל֤וֹא לֵֽאלֹהִים֙ פִּתְרֹנִ֔ים סַפְּרוּ־נָ֖א לִֽי׃ %
% ---------- פרק 40 | פסוק 9 ----------
 \textbf{ט} וַיְסַפֵּ֧ר שַֽׂר־הַמַּשְׁקִ֛ים אֶת־חֲלֹמ֖וֹ לְיוֹסֵ֑ף וַיֹּ֣אמֶר ל֔וֹ בַּחֲלוֹמִ֕י וְהִנֵּה־גֶ֖פֶן לְפָנָֽי׃ %
% ---------- פרק 40 | פסוק 10 ----------
 \textbf{י} וּבַגֶּ֖פֶן שְׁלֹשָׁ֣ה שָׂרִיגִ֑ם וְהִ֤וא כְפֹרַ֙חַת֙ עָלְתָ֣ה נִצָּ֔הּ הִבְשִׁ֥ילוּ אַשְׁכְּלֹתֶ֖יהָ עֲנָבִֽים׃ \Rashi{\textbf{שריגם.}זמורות ארוכות, שקורין ווידי"ץ בלעז: והוא כפרחת. דומה לפורחת והיא כפורחת – נדמה לי בחלומי כאלו היא פורחת, ואחר הפרח עלתה נצה, ונעשו סמדר, אשפנ"יר בלעז, ואחר כך הבשילו, והיא כד אפרחת אפקת לבלבין ע"כ תרגום של פורחת נץ גדול מפרח כדכתיב ובסר גמל יהיה נצה (ישעיהו י"ח), וכתיב ויצא פרח, והדר ויצץ ציץ (במדבר י"ז):}%
% ---------- פרק 40 | פסוק 11 ----------
 \textbf{יא} וְכ֥וֹס פַּרְעֹ֖ה בְּיָדִ֑י וָאֶקַּ֣ח אֶת־הָֽעֲנָבִ֗ים וָֽאֶשְׂחַ֤ט אֹתָם֙ אֶל־כּ֣וֹס פַּרְעֹ֔ה וָאֶתֵּ֥ן אֶת־הַכּ֖וֹס עַל־כַּ֥ף פַּרְעֹֽה׃ \Rashi{\textbf{ואשחט.}כתרגומו ועצרית, והרבה יש בלשון משנה:}%
% ---------- פרק 40 | פסוק 12 ----------
 \textbf{יב} וַיֹּ֤אמֶר לוֹ֙ יוֹסֵ֔ף זֶ֖ה פִּתְרֹנ֑וֹ שְׁלֹ֙שֶׁת֙ הַשָּׂ֣רִגִ֔ים שְׁלֹ֥שֶׁת יָמִ֖ים הֵֽם׃ \Rashi{\textbf{שלשת ימים הם.}סימן הם לך לשלשת ימים, ויש מדרשי אגדה הרבה (חולין צ"ב):}%
% ---------- פרק 40 | פסוק 13 ----------
 \textbf{יג} בְּע֣וֹד ׀ שְׁלֹ֣שֶׁת יָמִ֗ים יִשָּׂ֤א פַרְעֹה֙ אֶת־רֹאשֶׁ֔ךָ וַהֲשִֽׁיבְךָ֖ עַל־כַּנֶּ֑ךָ וְנָתַתָּ֤ כוֹס־פַּרְעֹה֙ בְּיָד֔וֹ כַּמִּשְׁפָּט֙ הָֽרִאשׁ֔וֹן אֲשֶׁ֥ר הָיִ֖יתָ מַשְׁקֵֽהוּ׃ \Rashi{\textbf{ישא פרעה את ראשך.}ל' חשבון, כשיפקד שאר עבדיו לשרת לפניו בסעודה, ימנה אותך עמהם: כנך. בסיס שלך ומושבך:}%
% ---------- פרק 40 | פסוק 14 ----------
 \textbf{יד} כִּ֧י אִם־זְכַרְתַּ֣נִי אִתְּךָ֗ כַּאֲשֶׁר֙ יִ֣יטַב לָ֔ךְ וְעָשִֽׂיתָ־נָּ֥א עִמָּדִ֖י חָ֑סֶד וְהִזְכַּרְתַּ֙נִי֙ אֶל־פַּרְעֹ֔ה וְהוֹצֵאתַ֖נִי מִן־הַבַּ֥יִת הַזֶּֽה׃ \Rashi{\textbf{כי אם זכרתני אתך.}אשר אם זכרתני אתך, מאחר שייטב לך כפתרוני: ועשית נא עמדי חסד. אין נא אלא לשון בקשה:}%
% ---------- פרק 40 | פסוק 15 ----------
 \textbf{טו} כִּֽי־גֻנֹּ֣ב גֻּנַּ֔בְתִּי מֵאֶ֖רֶץ הָעִבְרִ֑ים וְגַם־פֹּה֙ לֹא־עָשִׂ֣יתִֽי מְא֔וּמָה כִּֽי־שָׂמ֥וּ אֹתִ֖י בַּבּֽוֹר׃ %
% ---------- פרק 40 | פסוק 16 ----------
 \textbf{טז} וַיַּ֥רְא שַׂר־הָאֹפִ֖ים כִּ֣י ט֣וֹב פָּתָ֑ר וַיֹּ֙אמֶר֙ אֶל־יוֹסֵ֔ף אַף־אֲנִי֙ בַּחֲלוֹמִ֔י וְהִנֵּ֗ה שְׁלֹשָׁ֛ה סַלֵּ֥י חֹרִ֖י עַל־רֹאשִֽׁי׃ \Rashi{\textbf{סלי חורי.}סלים של נצרים קלופים חורין חורין, ובמקומנו יש הרבה, ודרך מוכרי פת כסנין, שקורין אובל"יש בלעז, לתתם באותם סלים:}%
% ---------- פרק 40 | פסוק 17 ----------
 \textbf{יז} וּבַסַּ֣ל הָֽעֶלְי֗וֹן מִכֹּ֛ל מַאֲכַ֥ל פַּרְעֹ֖ה מַעֲשֵׂ֣ה אֹפֶ֑ה וְהָע֗וֹף אֹכֵ֥ל אֹתָ֛ם מִן־הַסַּ֖ל מֵעַ֥ל רֹאשִֽׁי׃ %
% ---------- פרק 40 | פסוק 18 ----------
 \textbf{יח} וַיַּ֤עַן יוֹסֵף֙ וַיֹּ֔אמֶר זֶ֖ה פִּתְרֹנ֑וֹ שְׁלֹ֙שֶׁת֙ הַסַּלִּ֔ים שְׁלֹ֥שֶׁת יָמִ֖ים הֵֽם׃ %
% ---------- פרק 40 | פסוק 19 ----------
 \textbf{יט} בְּע֣וֹד ׀ שְׁלֹ֣שֶׁת יָמִ֗ים יִשָּׂ֨א פַרְעֹ֤ה אֶת־רֹֽאשְׁךָ֙ מֵֽעָלֶ֔יךָ וְתָלָ֥ה אוֹתְךָ֖ עַל־עֵ֑ץ וְאָכַ֥ל הָע֛וֹף אֶת־בְּשָׂרְךָ֖ מֵעָלֶֽיךָ׃ %
% ---------- פרק 40 | פסוק 20 ----------
 \textbf{כ} וַיְהִ֣י ׀ בַּיּ֣וֹם הַשְּׁלִישִׁ֗י י֚וֹם הֻלֶּ֣דֶת אֶת־פַּרְעֹ֔ה וַיַּ֥עַשׂ מִשְׁתֶּ֖ה לְכׇל־עֲבָדָ֑יו וַיִּשָּׂ֞א אֶת־רֹ֣אשׁ ׀ שַׂ֣ר הַמַּשְׁקִ֗ים וְאֶת־רֹ֛אשׁ שַׂ֥ר הָאֹפִ֖ים בְּת֥וֹךְ עֲבָדָֽיו׃ \Rashi{\textbf{יום הלדת את פרעה.}יום לידתו, וקורין לו יום גינוסיא. ולשון הלדת, לפי שאין הולד נוצר אלא על ידי אחרים, שהחיה מילדת את האשה, ועל כן החיה נקראת מילדת, וכן ומולדותיך ביום הולדת אותך (יחזקאל ט"ז) וכן אחרי הכבס את הנגע (ויקרא י"ג), שכבוסו על ידי אחרים: וישא את ראש וגו'. מנאם עם שאר עבדיו, שהיה מונה המשרתים שישרתו לו בסעודתו וזכר את אלו בתוכם, כמו שאו את ראש (במדבר א'), לשון מנין:}%
% ---------- פרק 40 | פסוק 21 ----------
 \textbf{כא} וַיָּ֛שֶׁב אֶת־שַׂ֥ר הַמַּשְׁקִ֖ים עַל־מַשְׁקֵ֑הוּ וַיִּתֵּ֥ן הַכּ֖וֹס עַל־כַּ֥ף פַּרְעֹֽה׃ %
% ---------- פרק 40 | פסוק 22 ----------
 \textbf{כב} וְאֵ֛ת שַׂ֥ר הָאֹפִ֖ים תָּלָ֑ה כַּאֲשֶׁ֥ר פָּתַ֛ר לָהֶ֖ם יוֹסֵֽף׃ %
% ---------- פרק 40 | פסוק 23 ----------
 \textbf{כג} וְלֹֽא־זָכַ֧ר שַֽׂר־הַמַּשְׁקִ֛ים אֶת־יוֹסֵ֖ף וַיִּשְׁכָּחֵֽהוּ׃ \{פ\} \Rashi{\textbf{ולא זכר שר המשקים.}בו ביום: וישכחהו. לאחר מכאן. מפני שתלה בו יוסף לזכרו, הזקק להיות אסור שתי שנים, שנאמר אשרי הגבר אשר שם ה' מבטחו ולא פנה אל רהבים (תהילים מ') – ולא בטח על מצרים, הקרויים רהב:}%
\addparasha{חומש בראשית | פרשת מקץ}
% ---------- פרק 41 | פסוק 1 ----------
 \textbf{א} וַיְהִ֕י מִקֵּ֖ץ שְׁנָתַ֣יִם יָמִ֑ים וּפַרְעֹ֣ה חֹלֵ֔ם וְהִנֵּ֖ה עֹמֵ֥ד עַל־הַיְאֹֽר׃ \Rashi{\textbf{ויהי מקץ.}כתרגומו מסוף, וכל לשון קץ סוף הוא: על היאר. כל שאר נהרות אינם קרויין יאורים חוץ מנילוס, מפני שכל הארץ עשוים יאורים יאורים בידי אדם ונילוס עולה בתוכם ומשקה אותם, לפי שאין גשמים יורדין במצרים תדיר כשאר ארצות:}%
% ---------- פרק 41 | פסוק 2 ----------
 \textbf{ב} וְהִנֵּ֣ה מִן־הַיְאֹ֗ר עֹלֹת֙ שֶׁ֣בַע פָּר֔וֹת יְפ֥וֹת מַרְאֶ֖ה וּבְרִיאֹ֣ת בָּשָׂ֑ר וַתִּרְעֶ֖ינָה בָּאָֽחוּ׃ \Rashi{\textbf{יפות מראה.}סימן הוא לימי שבע, שהבריות נראות יפות זו לזו, שאין עין בריה צרה בחברתה: באחו. באגם, מריש"ק בלעז, כמו ישגא אחו (איוב ח'):}%
% ---------- פרק 41 | פסוק 3 ----------
 \textbf{ג} וְהִנֵּ֞ה שֶׁ֧בַע פָּר֣וֹת אֲחֵר֗וֹת עֹל֤וֹת אַחֲרֵיהֶן֙ מִן־הַיְאֹ֔ר רָע֥וֹת מַרְאֶ֖ה וְדַקּ֣וֹת בָּשָׂ֑ר וַֽתַּעֲמֹ֛דְנָה אֵ֥צֶל הַפָּר֖וֹת עַל־שְׂפַ֥ת הַיְאֹֽר׃ \Rashi{\textbf{ודקות בשר.}טינב"ש בלעז, לשון דק:}%
% ---------- פרק 41 | פסוק 4 ----------
 \textbf{ד} וַתֹּאכַ֣לְנָה הַפָּר֗וֹת רָע֤וֹת הַמַּרְאֶה֙ וְדַקֹּ֣ת הַבָּשָׂ֔ר אֵ֚ת שֶׁ֣בַע הַפָּר֔וֹת יְפֹ֥ת הַמַּרְאֶ֖ה וְהַבְּרִיאֹ֑ת וַיִּיקַ֖ץ פַּרְעֹֽה׃ \Rashi{\textbf{ותאכלנה.}סימן שתהא כל שמחת השבע נשכחת בימי הרעב:}%
% ---------- פרק 41 | פסוק 5 ----------
 \textbf{ה} וַיִּישָׁ֕ן וַֽיַּחֲלֹ֖ם שֵׁנִ֑ית וְהִנֵּ֣ה ׀ שֶׁ֣בַע שִׁבֳּלִ֗ים עֹל֛וֹת בְּקָנֶ֥ה אֶחָ֖ד בְּרִיא֥וֹת וְטֹבֽוֹת׃ \Rashi{\textbf{בקנה אחת.}טוד"ל בלעז: בריאות. שיינ"ש בלעז:}%
% ---------- פרק 41 | פסוק 6 ----------
 \textbf{ו} וְהִנֵּה֙ שֶׁ֣בַע שִׁבֳּלִ֔ים דַּקּ֖וֹת וּשְׁדוּפֹ֣ת קָדִ֑ים צֹמְח֖וֹת אַחֲרֵיהֶֽן׃ \Rashi{\textbf{ושדופת.}הלייא"ש בלעז, שקיפן קדום, חבוטות, לשון משקוף החבוט תמיד על ידי הדלת המכה עליו: קדים. רוח מזרחית, שקורין בי"סא בלעז:}%
% ---------- פרק 41 | פסוק 7 ----------
 \textbf{ז} וַתִּבְלַ֙עְנָה֙ הַשִּׁבֳּלִ֣ים הַדַּקּ֔וֹת אֵ֚ת שֶׁ֣בַע הַֽשִּׁבֳּלִ֔ים הַבְּרִיא֖וֹת וְהַמְּלֵא֑וֹת וַיִּיקַ֥ץ פַּרְעֹ֖ה וְהִנֵּ֥ה חֲלֽוֹם׃ \Rashi{\textbf{הבריאות.}שיינ"ש בלעז: והנה חלום. והנה נשלם חלום שלם לפניו והצרך לפותרים:}%
% ---------- פרק 41 | פסוק 8 ----------
 \textbf{ח} וַיְהִ֤י בַבֹּ֙קֶר֙ וַתִּפָּ֣עֶם רוּח֔וֹ וַיִּשְׁלַ֗ח וַיִּקְרָ֛א אֶת־כׇּל־חַרְטֻמֵּ֥י מִצְרַ֖יִם וְאֶת־כׇּל־חֲכָמֶ֑יהָ וַיְסַפֵּ֨ר פַּרְעֹ֤ה לָהֶם֙ אֶת־חֲלֹמ֔וֹ וְאֵין־פּוֹתֵ֥ר אוֹתָ֖ם לְפַרְעֹֽה׃ \Rashi{\textbf{ותפעם רוחו.}ומטרפא רוחיה, מקשקשת בתוכו כפעמון; ובנבוכדנצר אומר, ותתפעם רוחו (דניאל ב'), לפי שהיו שם שתי פעימות, שכחת החלום והעלמת פתרונו (בראשית רבה): חרטמי. הנחרים בטימי מתים, ששואלין בעצמות. טימי הן עצמות בלשון ארמי, ובמשנה בית שהוא מלא טימיא (אהלות פי"ז), מלא עצמות: ואין פותר אותם לפרעה. פותרים היו אותם, אבל לא לפרעה, שלא היה קולן נכנס באזניו, ולא היה לו קורת רוח בפתרונם, שהיו אומרים שבע בנות אתה מוליד ושבע בנות אתה קובר (בראשית רבה):}%
% ---------- פרק 41 | פסוק 9 ----------
 \textbf{ט} וַיְדַבֵּר֙ שַׂ֣ר הַמַּשְׁקִ֔ים אֶת־פַּרְעֹ֖ה לֵאמֹ֑ר אֶת־חֲטָאַ֕י אֲנִ֖י מַזְכִּ֥יר הַיּֽוֹם׃ %
% ---------- פרק 41 | פסוק 10 ----------
 \textbf{י} פַּרְעֹ֖ה קָצַ֣ף עַל־עֲבָדָ֑יו וַיִּתֵּ֨ן אֹתִ֜י בְּמִשְׁמַ֗ר בֵּ֚ית שַׂ֣ר הַטַּבָּחִ֔ים אֹתִ֕י וְאֵ֖ת שַׂ֥ר הָאֹפִֽים׃ %
% ---------- פרק 41 | פסוק 11 ----------
 \textbf{יא} וַנַּֽחַלְמָ֥ה חֲל֛וֹם בְּלַ֥יְלָה אֶחָ֖ד אֲנִ֣י וָה֑וּא אִ֛ישׁ כְּפִתְר֥וֹן חֲלֹמ֖וֹ חָלָֽמְנוּ׃ \Rashi{\textbf{איש כפתרון חלמו.}חלום הראוי לפתרון שנפתר לנו ודומה לו:}%
% ---------- פרק 41 | פסוק 12 ----------
 \textbf{יב} וְשָׁ֨ם אִתָּ֜נוּ נַ֣עַר עִבְרִ֗י עֶ֚בֶד לְשַׂ֣ר הַטַּבָּחִ֔ים וַ֨נְּסַפֶּר־ל֔וֹ וַיִּפְתׇּר־לָ֖נוּ אֶת־חֲלֹמֹתֵ֑ינוּ אִ֥ישׁ כַּחֲלֹמ֖וֹ פָּתָֽר׃ \Rashi{\textbf{נער עברי עבד.}ארורים הרשעים, שאין טובתם שלמה, מזכירו בלשון בזיון: נער. שוטה ואין ראוי לגדלה: עברי. אפלו לשוננו אינו מכיר: עבד. וכתוב בנמוסי מצרים שאין עבד מולך ולא לובש בגדי שרים: איש כחלמו. לפי החלום וקרוב לענינו:}%
% ---------- פרק 41 | פסוק 13 ----------
 \textbf{יג} וַיְהִ֛י כַּאֲשֶׁ֥ר פָּֽתַר־לָ֖נוּ כֵּ֣ן הָיָ֑ה אֹתִ֛י הֵשִׁ֥יב עַל־כַּנִּ֖י וְאֹת֥וֹ תָלָֽה׃ \Rashi{\textbf{השיב על כני.}פרעה, הנזכר למעלה, כמו שאמר פרעה קצף על עבדיו; הרי מקרא קצר לשון, ולא פרש מי השיב, לפי שאין צריך לפרש; מי השיב? מי שבידו להשיב, והוא פרעה; וכן דרך כל מקראות קצרים, על מי שעליו לעשות הם סותמים את הדבר:}%
% ---------- פרק 41 | פסוק 14 ----------
 \textbf{יד} וַיִּשְׁלַ֤ח פַּרְעֹה֙ וַיִּקְרָ֣א אֶת־יוֹסֵ֔ף וַיְרִיצֻ֖הוּ מִן־הַבּ֑וֹר וַיְגַלַּח֙ וַיְחַלֵּ֣ף שִׂמְלֹתָ֔יו וַיָּבֹ֖א אֶל־פַּרְעֹֽה׃ \Rashi{\textbf{מן הבור.}מן בית הסהר שהוא עשוי כמין גמא, וכן כל בור שבמקרא לשון גמא הוא, ואף אם אין בו מים קרוי בור, פוש"א בלעז: ויגלח. מפני כבוד המלכות (בראשית רבה):}%
% ---------- פרק 41 | פסוק 15 ----------
 \textbf{טו} וַיֹּ֤אמֶר פַּרְעֹה֙ אֶל־יוֹסֵ֔ף חֲל֣וֹם חָלַ֔מְתִּי וּפֹתֵ֖ר אֵ֣ין אֹת֑וֹ וַאֲנִ֗י שָׁמַ֤עְתִּי עָלֶ֙יךָ֙ לֵאמֹ֔ר תִּשְׁמַ֥ע חֲל֖וֹם לִפְתֹּ֥ר אֹתֽוֹ׃ \Rashi{\textbf{תשמע חלום לפתר אתו.}תאזין ותבין חלום לפתר אותו: תשמע. לשון הבנה והאזנה, כמו שומע יוסף, אשר לא תשמע לשנו, אנטינ"דרא בלעז:}%
% ---------- פרק 41 | פסוק 16 ----------
 \textbf{טז} וַיַּ֨עַן יוֹסֵ֧ף אֶת־פַּרְעֹ֛ה לֵאמֹ֖ר בִּלְעָדָ֑י אֱלֹהִ֕ים יַעֲנֶ֖ה אֶת־שְׁל֥וֹם פַּרְעֹֽה׃ \Rashi{\textbf{בלעדי.}אין החכמה משלי, אלא אלהים יענה – יתן עניה בפי – לשלום פרעה:}%
% ---------- פרק 41 | פסוק 17 ----------
 \textbf{יז} וַיְדַבֵּ֥ר פַּרְעֹ֖ה אֶל־יוֹסֵ֑ף בַּחֲלֹמִ֕י הִנְנִ֥י עֹמֵ֖ד עַל־שְׂפַ֥ת הַיְאֹֽר׃ %
% ---------- פרק 41 | פסוק 18 ----------
 \textbf{יח} וְהִנֵּ֣ה מִן־הַיְאֹ֗ר עֹלֹת֙ שֶׁ֣בַע פָּר֔וֹת בְּרִיא֥וֹת בָּשָׂ֖ר וִיפֹ֣ת תֹּ֑אַר וַתִּרְעֶ֖ינָה בָּאָֽחוּ׃ %
% ---------- פרק 41 | פסוק 19 ----------
 \textbf{יט} וְהִנֵּ֞ה שֶֽׁבַע־פָּר֤וֹת אֲחֵרוֹת֙ עֹל֣וֹת אַחֲרֵיהֶ֔ן דַּלּ֨וֹת וְרָע֥וֹת תֹּ֛אַר מְאֹ֖ד וְרַקּ֣וֹת בָּשָׂ֑ר לֹֽא־רָאִ֧יתִי כָהֵ֛נָּה בְּכׇל־אֶ֥רֶץ מִצְרַ֖יִם לָרֹֽעַ׃ \Rashi{\textbf{דלות.}כחושות, כמו מדוע אתה ככה דל, דאמנון (שמואל ב י"ג): ורקות בשר. כל לשון רקות שבמקרא חסרי בשר, ובלעז בלוא"ש:}%
% ---------- פרק 41 | פסוק 20 ----------
 \textbf{כ} וַתֹּאכַ֙לְנָה֙ הַפָּר֔וֹת הָרַקּ֖וֹת וְהָרָע֑וֹת אֵ֣ת שֶׁ֧בַע הַפָּר֛וֹת הָרִאשֹׁנ֖וֹת הַבְּרִיאֹֽת׃ %
% ---------- פרק 41 | פסוק 21 ----------
 \textbf{כא} וַתָּבֹ֣אנָה אֶל־קִרְבֶּ֗נָה וְלֹ֤א נוֹדַע֙ כִּי־בָ֣אוּ אֶל־קִרְבֶּ֔נָה וּמַרְאֵיהֶ֣ן רַ֔ע כַּאֲשֶׁ֖ר בַּתְּחִלָּ֑ה וָאִיקָֽץ׃ %
% ---------- פרק 41 | פסוק 22 ----------
 \textbf{כב} וָאֵ֖רֶא בַּחֲלֹמִ֑י וְהִנֵּ֣ה ׀ שֶׁ֣בַע שִׁבֳּלִ֗ים עֹלֹ֛ת בְּקָנֶ֥ה אֶחָ֖ד מְלֵאֹ֥ת וְטֹבֽוֹת׃ %
% ---------- פרק 41 | פסוק 23 ----------
 \textbf{כג} וְהִנֵּה֙ שֶׁ֣בַע שִׁבֳּלִ֔ים צְנֻמ֥וֹת דַּקּ֖וֹת שְׁדֻפ֣וֹת קָדִ֑ים צֹמְח֖וֹת אַחֲרֵיהֶֽם׃ \Rashi{\textbf{צנמות.}צונמא בלשון ארמי סלע, הרי הן כעץ בלי לחלוח וקשות כסלע, ותרגומו נצן לקין, נצן – אין בהם אלא הנץ, לפי שנתרוקנו מן הזרע:}%
% ---------- פרק 41 | פסוק 24 ----------
 \textbf{כד} וַתִּבְלַ֙עְןָ֙ הַשִּׁבֳּלִ֣ים הַדַּקֹּ֔ת אֵ֛ת שֶׁ֥בַע הַֽשִּׁבֳּלִ֖ים הַטֹּב֑וֹת וָֽאֹמַר֙ אֶל־הַֽחַרְטֻמִּ֔ים וְאֵ֥ין מַגִּ֖יד לִֽי׃ %
% ---------- פרק 41 | פסוק 25 ----------
 \textbf{כה} וַיֹּ֤אמֶר יוֹסֵף֙ אֶל־פַּרְעֹ֔ה חֲל֥וֹם פַּרְעֹ֖ה אֶחָ֣ד ה֑וּא אֵ֣ת אֲשֶׁ֧ר הָאֱלֹהִ֛ים עֹשֶׂ֖ה הִגִּ֥יד לְפַרְעֹֽה׃ %
% ---------- פרק 41 | פסוק 26 ----------
 \textbf{כו} שֶׁ֧בַע פָּרֹ֣ת הַטֹּבֹ֗ת שֶׁ֤בַע שָׁנִים֙ הֵ֔נָּה וְשֶׁ֤בַע הַֽשִּׁבֳּלִים֙ הַטֹּבֹ֔ת שֶׁ֥בַע שָׁנִ֖ים הֵ֑נָּה חֲל֖וֹם אֶחָ֥ד הֽוּא׃ \Rashi{\textbf{שבע שנים ושבע שנים.}כלן אינן אלא שבע, ואשר נשנה החלום פעמים, לפי שהדבר מזמן כמו שפרש לו בסוף ועל השנות החלום וגו' בשבע שנים הטובות נאמר הגיד לפרעה לפי שהיה סמוך ובשבע שני רעב נאמר הראה את פרעה, לפי שהיה הדבר מפלג ורחוק, נופל בו לשון מראה:}%
% ---------- פרק 41 | פסוק 27 ----------
 \textbf{כז} וְשֶׁ֣בַע הַ֠פָּר֠וֹת הָֽרַקּ֨וֹת וְהָרָעֹ֜ת הָעֹלֹ֣ת אַחֲרֵיהֶ֗ן שֶׁ֤בַע שָׁנִים֙ הֵ֔נָּה וְשֶׁ֤בַע הַֽשִּׁבֳּלִים֙ הָרֵק֔וֹת שְׁדֻפ֖וֹת הַקָּדִ֑ים יִהְי֕וּ שֶׁ֖בַע שְׁנֵ֥י רָעָֽב׃ %
% ---------- פרק 41 | פסוק 28 ----------
 \textbf{כח} ה֣וּא הַדָּבָ֔ר אֲשֶׁ֥ר דִּבַּ֖רְתִּי אֶל־פַּרְעֹ֑ה אֲשֶׁ֧ר הָאֱלֹהִ֛ים עֹשֶׂ֖ה הֶרְאָ֥ה אֶת־פַּרְעֹֽה׃ %
% ---------- פרק 41 | פסוק 29 ----------
 \textbf{כט} הִנֵּ֛ה שֶׁ֥בַע שָׁנִ֖ים בָּא֑וֹת שָׂבָ֥ע גָּד֖וֹל בְּכׇל־אֶ֥רֶץ מִצְרָֽיִם׃ %
% ---------- פרק 41 | פסוק 30 ----------
 \textbf{ל} וְ֠קָ֠מוּ שֶׁ֜בַע שְׁנֵ֤י רָעָב֙ אַחֲרֵיהֶ֔ן וְנִשְׁכַּ֥ח כׇּל־הַשָּׂבָ֖ע בְּאֶ֣רֶץ מִצְרָ֑יִם וְכִלָּ֥ה הָרָעָ֖ב אֶת־הָאָֽרֶץ׃ \Rashi{\textbf{ונשכח כל השבע.}הוא פתרון הבליעה:}%
% ---------- פרק 41 | פסוק 31 ----------
 \textbf{לא} וְלֹֽא־יִוָּדַ֤ע הַשָּׂבָע֙ בָּאָ֔רֶץ מִפְּנֵ֛י הָרָעָ֥ב הַה֖וּא אַחֲרֵי־כֵ֑ן כִּֽי־כָבֵ֥ד ה֖וּא מְאֹֽד׃ \Rashi{\textbf{ולא יודע השבע.}הוא פתרון "ולא נודע כי באו אל קרבנה":}%
% ---------- פרק 41 | פסוק 32 ----------
 \textbf{לב} וְעַ֨ל הִשָּׁנ֧וֹת הַחֲל֛וֹם אֶל־פַּרְעֹ֖ה פַּעֲמָ֑יִם כִּֽי־נָכ֤וֹן הַדָּבָר֙ מֵעִ֣ם הָאֱלֹהִ֔ים וּמְמַהֵ֥ר הָאֱלֹהִ֖ים לַעֲשֹׂתֽוֹ׃ \Rashi{\textbf{נכון.}מזמן:}%
% ---------- פרק 41 | פסוק 33 ----------
 \textbf{לג} וְעַתָּה֙ יֵרֶ֣א פַרְעֹ֔ה אִ֖ישׁ נָב֣וֹן וְחָכָ֑ם וִישִׁיתֵ֖הוּ עַל־אֶ֥רֶץ מִצְרָֽיִם׃ %
% ---------- פרק 41 | פסוק 34 ----------
 \textbf{לד} יַעֲשֶׂ֣ה פַרְעֹ֔ה וְיַפְקֵ֥ד פְּקִדִ֖ים עַל־הָאָ֑רֶץ וְחִמֵּשׁ֙ אֶת־אֶ֣רֶץ מִצְרַ֔יִם בְּשֶׁ֖בַע שְׁנֵ֥י הַשָּׂבָֽע׃ \Rashi{\textbf{וחמש.}כתרגומו ויזרזון; וכן וחמשים (שמות י"ג):}%
% ---------- פרק 41 | פסוק 35 ----------
 \textbf{לה} וְיִקְבְּצ֗וּ אֶת־כׇּל־אֹ֙כֶל֙ הַשָּׁנִ֣ים הַטֹּב֔וֹת הַבָּאֹ֖ת הָאֵ֑לֶּה וְיִצְבְּרוּ־בָ֞ר תַּ֧חַת יַד־פַּרְעֹ֛ה אֹ֥כֶל בֶּעָרִ֖ים וְשָׁמָֽרוּ׃ \Rashi{\textbf{את כל אכל.}שם דבר הוא, לפיכך טעמו באל"ף ונקוד בפתח קטן, ואוכל שהוא פעל, כגון כי כל אכל חלב, טעמו למטה בכ"ף, ונקוד קמץ קטן: תחת יד פרעה. ברשותו ובאוצרותיו:}%
% ---------- פרק 41 | פסוק 36 ----------
 \textbf{לו} וְהָיָ֨ה הָאֹ֤כֶל לְפִקָּדוֹן֙ לָאָ֔רֶץ לְשֶׁ֙בַע֙ שְׁנֵ֣י הָרָעָ֔ב אֲשֶׁ֥ר תִּהְיֶ֖יןָ בְּאֶ֣רֶץ מִצְרָ֑יִם וְלֹֽא־תִכָּרֵ֥ת הָאָ֖רֶץ בָּרָעָֽב׃ \Rashi{\textbf{והיה האכל.}הצבור כשאר פקדון הגנוז לקיום הארץ:}%
% ---------- פרק 41 | פסוק 37 ----------
 \textbf{לז} וַיִּיטַ֥ב הַדָּבָ֖ר בְּעֵינֵ֣י פַרְעֹ֑ה וּבְעֵינֵ֖י כׇּל־עֲבָדָֽיו׃ %
% ---------- פרק 41 | פסוק 38 ----------
 \textbf{לח} וַיֹּ֥אמֶר פַּרְעֹ֖ה אֶל־עֲבָדָ֑יו הֲנִמְצָ֣א כָזֶ֔ה אִ֕ישׁ אֲשֶׁ֛ר ר֥וּחַ אֱלֹהִ֖ים בּֽוֹ׃ \Rashi{\textbf{הנמצא כזה.}הנשכח כדין? אם נלך ונבקשנו, הנמצא כמוהו? הנמצא לשון תמיהה; וכן כל ה"א המשמשת בראש תבה ונקודה בחטף פתח:}%
% ---------- פרק 41 | פסוק 39 ----------
 \textbf{לט} וַיֹּ֤אמֶר פַּרְעֹה֙ אֶל־יוֹסֵ֔ף אַחֲרֵ֨י הוֹדִ֧יעַ אֱלֹהִ֛ים אוֹתְךָ֖ אֶת־כׇּל־זֹ֑את אֵין־נָב֥וֹן וְחָכָ֖ם כָּמֽוֹךָ׃ \Rashi{\textbf{אין נבון וחכם כמוך.}לבקש איש נבון וחכם שאמרת, לא נמצא כמוך:}%
% ---------- פרק 41 | פסוק 40 ----------
 \textbf{מ} אַתָּה֙ תִּהְיֶ֣ה עַל־בֵּיתִ֔י וְעַל־פִּ֖יךָ יִשַּׁ֣ק כׇּל־עַמִּ֑י רַ֥ק הַכִּסֵּ֖א אֶגְדַּ֥ל מִמֶּֽךָּ׃ \Rashi{\textbf{ישק.}יתזן, יתפרנס – כל צרכי עמי יהיו נעשים על ידך, כמו נשקו בר (תהילים ב'), גרני"שון בלעז: רק הכסא. שיהיו קורין לי מלך: כסא. לשון שם המלוכה, כמו ויגדל את כסאו מכסא אדני המלך (מלכים א א'):}%
% ---------- פרק 41 | פסוק 41 ----------
 \textbf{מא} וַיֹּ֥אמֶר פַּרְעֹ֖ה אֶל־יוֹסֵ֑ף רְאֵה֙ נָתַ֣תִּי אֹֽתְךָ֔ עַ֖ל כׇּל־אֶ֥רֶץ מִצְרָֽיִם׃ \Rashi{\textbf{נתתי אתך.}מניתי יתך, ואע"פ כן לשון נתינה הוא, כמו ולתתך עליון (דברים כ"ו); בין לגדלה בין לשפלות נופל לשון נתינה עליו, כמו נתתי אתכם נבזים ושפלים (מלאכי ב'):}%
% ---------- פרק 41 | פסוק 42 ----------
 \textbf{מב} וַיָּ֨סַר פַּרְעֹ֤ה אֶת־טַבַּעְתּוֹ֙ מֵעַ֣ל יָד֔וֹ וַיִּתֵּ֥ן אֹתָ֖הּ עַל־יַ֣ד יוֹסֵ֑ף וַיַּלְבֵּ֤שׁ אֹתוֹ֙ בִּגְדֵי־שֵׁ֔שׁ וַיָּ֛שֶׂם רְבִ֥ד הַזָּהָ֖ב עַל־צַוָּארֽוֹ׃ \Rashi{\textbf{ויסר פרעה את טבעתו.}נתינת טבעת המלך היא אות למי שנותנה לו להיות שני לו לגדלה: בגדי שש. דבר חשיבות הוא במצרים: רביד. ענק, ועל שהוא רצוף בטבעות קרוי רביד וכן מרבדים רבדתי ערשי (משלי ז') – רצפתי ערשי מרצפות, בלשון משנה מקף רובדין של אבן, על הרבד שבעזרה, והיא רצפה:}%
% ---------- פרק 41 | פסוק 43 ----------
 \textbf{מג} וַיַּרְכֵּ֣ב אֹת֗וֹ בְּמִרְכֶּ֤בֶת הַמִּשְׁנֶה֙ אֲשֶׁר־ל֔וֹ וַיִּקְרְא֥וּ לְפָנָ֖יו אַבְרֵ֑ךְ וְנָת֣וֹן אֹת֔וֹ עַ֖ל כׇּל־אֶ֥רֶץ מִצְרָֽיִם׃ \Rashi{\textbf{במרכבת המשנה.}השניה למרכבתו, המהלכת אצל שלו: אברך. כתרגומו "דין אבא למלכא", רך בלשון ארמי (נ"א רומי) מלך; בהשתפין לא ריכא ולא בר ריכא, ובדברי אגדה דרש רבי יהודה אברך זה יוסף שהוא אב בחכמה ורך בשנים, אמר לו ר' יוסי בן דרמסקית עד מתי אתה מעות עלינו את הכתובים? אין אברך אלא לשון ברכים, שהכל היו נכנסין ויוצאין תחת ידו, כענין שנאמר, ונתון אותו וגו':}%
% ---------- פרק 41 | פסוק 44 ----------
 \textbf{מד} וַיֹּ֧אמֶר פַּרְעֹ֛ה אֶל־יוֹסֵ֖ף אֲנִ֣י פַרְעֹ֑ה וּבִלְעָדֶ֗יךָ לֹֽא־יָרִ֨ים אִ֧ישׁ אֶת־יָד֛וֹ וְאֶת־רַגְל֖וֹ בְּכׇל־אֶ֥רֶץ מִצְרָֽיִם׃ \Rashi{\textbf{אני פרעה.}שיש יכלת בידי לגזר גזרה על מלכותי, ואני גוזר שלא ירים איש ידו בלעדיך, – שלא ברשותך. ד"א אני פרעה, אני אהיה מלך, ובלעדיך וגו' וזהו דגמת רק הכסא: את ידו ואת רגלו. כתרגומו:}%
% ---------- פרק 41 | פסוק 45 ----------
 \textbf{מה} וַיִּקְרָ֨א פַרְעֹ֣ה שֵׁם־יוֹסֵף֮ צָֽפְנַ֣ת פַּעְנֵ֒חַ֒ וַיִּתֶּן־ל֣וֹ אֶת־אָֽסְנַ֗ת בַּת־פּ֥וֹטִי פֶ֛רַע*(בספרי תימן פּֽוֹטִיפֶ֛רַע בתיבה אחת) כֹּהֵ֥ן אֹ֖ן לְאִשָּׁ֑ה וַיֵּצֵ֥א יוֹסֵ֖ף עַל־אֶ֥רֶץ מִצְרָֽיִם׃ \Rashi{\textbf{צפנת פענח.}מפרש הצפונות; ואין לפענח דמיון במקרא: פוטיפרע. הוא פוטיפר, ונקרא פוטיפרע על שנסתרס מאליו, לפי שחמד את יוסף למשכב זכר (סוטה י"ג):}%
% ---------- פרק 41 | פסוק 46 ----------
 \textbf{מו} וְיוֹסֵף֙ בֶּן־שְׁלֹשִׁ֣ים שָׁנָ֔ה בְּעׇמְד֕וֹ לִפְנֵ֖י פַּרְעֹ֣ה מֶֽלֶךְ־מִצְרָ֑יִם וַיֵּצֵ֤א יוֹסֵף֙ מִלִּפְנֵ֣י פַרְעֹ֔ה וַֽיַּעֲבֹ֖ר בְּכׇל־אֶ֥רֶץ מִצְרָֽיִם׃ %
% ---------- פרק 41 | פסוק 47 ----------
 \textbf{מז} וַתַּ֣עַשׂ הָאָ֔רֶץ בְּשֶׁ֖בַע שְׁנֵ֣י הַשָּׂבָ֑ע לִקְמָצִֽים׃ \Rashi{\textbf{ותעש הארץ.}כתרגומו, ואין הלשון נעקר מלשון עשיה: לקמצים. קמץ על קמץ, יד על יד, היו אוצרים:}%
% ---------- פרק 41 | פסוק 48 ----------
 \textbf{מח} וַיִּקְבֹּ֞ץ אֶת־כׇּל־אֹ֣כֶל ׀ שֶׁ֣בַע שָׁנִ֗ים אֲשֶׁ֤ר הָיוּ֙ בְּאֶ֣רֶץ מִצְרַ֔יִם וַיִּתֶּן־אֹ֖כֶל בֶּעָרִ֑ים אֹ֧כֶל שְׂדֵה־הָעִ֛יר אֲשֶׁ֥ר סְבִיבֹתֶ֖יהָ נָתַ֥ן בְּתוֹכָֽהּ׃ \Rashi{\textbf{אכל שדה העיר אשר סביבתיה נתן בתוכה.}שכל ארץ וארץ מעמדת פרותיה, ונותנין בתבואה מעפר המקום, ומעמיד את התבואה מלרקב:}%
% ---------- פרק 41 | פסוק 49 ----------
 \textbf{מט} וַיִּצְבֹּ֨ר יוֹסֵ֥ף בָּ֛ר כְּח֥וֹל הַיָּ֖ם הַרְבֵּ֣ה מְאֹ֑ד עַ֛ד כִּי־חָדַ֥ל לִסְפֹּ֖ר כִּי־אֵ֥ין מִסְפָּֽר׃ \Rashi{\textbf{עד כי חדל לספר.}עד כי חדל לו הסופר לספר, והרי זה מקרא קצר: כי אין מספר. לפי שאין מספר, והרי כי משמש בלשון דהא:}%
% ---------- פרק 41 | פסוק 50 ----------
 \textbf{נ} וּלְיוֹסֵ֤ף יֻלַּד֙ שְׁנֵ֣י בָנִ֔ים בְּטֶ֥רֶם תָּב֖וֹא שְׁנַ֣ת הָרָעָ֑ב אֲשֶׁ֤ר יָֽלְדָה־לּוֹ֙ אָֽסְנַ֔ת בַּת־פּ֥וֹטִי פֶ֖רַע*(בספרי תימן פּֽוֹטִיפֶ֖רַע בתיבה אחת) כֹּהֵ֥ן אֽוֹן׃ \Rashi{\textbf{בטרם תבא שנת הרעב.}מכאן שאדם אסור לשמש מטתו בשני רעבון (תענית י"א):}%
% ---------- פרק 41 | פסוק 51 ----------
 \textbf{נא} וַיִּקְרָ֥א יוֹסֵ֛ף אֶת־שֵׁ֥ם הַבְּכ֖וֹר מְנַשֶּׁ֑ה כִּֽי־נַשַּׁ֤נִי אֱלֹהִים֙ אֶת־כׇּל־עֲמָלִ֔י וְאֵ֖ת כׇּל־בֵּ֥ית אָבִֽי׃ %
% ---------- פרק 41 | פסוק 52 ----------
 \textbf{נב} וְאֵ֛ת שֵׁ֥ם הַשֵּׁנִ֖י קָרָ֣א אֶפְרָ֑יִם כִּֽי־הִפְרַ֥נִי אֱלֹהִ֖ים בְּאֶ֥רֶץ עׇנְיִֽי׃ %
% ---------- פרק 41 | פסוק 53 ----------
 \textbf{נג} וַתִּכְלֶ֕ינָה שֶׁ֖בַע שְׁנֵ֣י הַשָּׂבָ֑ע אֲשֶׁ֥ר הָיָ֖ה בְּאֶ֥רֶץ מִצְרָֽיִם׃ %
% ---------- פרק 41 | פסוק 54 ----------
 \textbf{נד} וַתְּחִלֶּ֜ינָה שֶׁ֣בַע שְׁנֵ֤י הָרָעָב֙ לָב֔וֹא כַּאֲשֶׁ֖ר אָמַ֣ר יוֹסֵ֑ף וַיְהִ֤י רָעָב֙ בְּכׇל־הָ֣אֲרָצ֔וֹת וּבְכׇל־אֶ֥רֶץ מִצְרַ֖יִם הָ֥יָה לָֽחֶם׃ %
% ---------- פרק 41 | פסוק 55 ----------
 \textbf{נה} וַתִּרְעַב֙ כׇּל־אֶ֣רֶץ מִצְרַ֔יִם וַיִּצְעַ֥ק הָעָ֛ם אֶל־פַּרְעֹ֖ה לַלָּ֑חֶם וַיֹּ֨אמֶר פַּרְעֹ֤ה לְכׇל־מִצְרַ֙יִם֙ לְכ֣וּ אֶל־יוֹסֵ֔ף אֲשֶׁר־יֹאמַ֥ר לָכֶ֖ם תַּעֲשֽׂוּ׃ \Rashi{\textbf{ותרעב כל ארץ מצרים.}שהרקיבה תבואתם שאצרו חוץ משל יוסף: אשר יאמר לכם תעשו. לפי שהיה יוסף אומר להם שימולו, וכשבאו אצל פרעה ואומרים: כך הוא אומר לנו, אמר להם: ולמה לא צברתם בר, והלא הכריז לכם ששני הרעב באים? אמרו לו: אספנו הרבה והרקיבה, אמר להם: אם כן, כל אשר יאמר לכם תעשו; הרי גזר על התבואה והרקיבה, מה אם יגזר עלינו ונמות!}%
% ---------- פרק 41 | פסוק 56 ----------
 \textbf{נו} וְהָרָעָ֣ב הָיָ֔ה עַ֖ל כׇּל־פְּנֵ֣י הָאָ֑רֶץ וַיִּפְתַּ֨ח יוֹסֵ֜ף אֶֽת־כׇּל־אֲשֶׁ֤ר בָּהֶם֙ וַיִּשְׁבֹּ֣ר לְמִצְרַ֔יִם וַיֶּחֱזַ֥ק הָֽרָעָ֖ב בְּאֶ֥רֶץ מִצְרָֽיִם׃ \Rashi{\textbf{על כל פני הארץ.}מי הם פני הארץ אלו העשירים: את כל אשר בהם. כתרגומו די בהון עיבורא: וישבר למצרים. שבר לשון מכר ולשון קנין הוא; כאן משמש לשון מכר; שברו לנו מעט אכל, לשון קנין. ואל תאמר אינו כי אם בתבואה, שאף ביין וחלב מצינו ולכו שברו בלוא כסף ובלוא מחיר יין וחלב (ישעיהו נ"ה):}%
% ---------- פרק 41 | פסוק 57 ----------
 \textbf{נז} וְכׇל־הָאָ֙רֶץ֙ בָּ֣אוּ מִצְרַ֔יְמָה לִשְׁבֹּ֖ר אֶל־יוֹסֵ֑ף כִּֽי־חָזַ֥ק הָרָעָ֖ב בְּכׇל־הָאָֽרֶץ׃ \Rashi{וכל הארץ באו מצרימה אל יוסף לשבר, ואם תדרשהו כסדרו, היה צריך לכתב: לשבר מן יוסף:}%
% ---------- פרק 42 | פסוק 1 ----------
 \textbf{א} וַיַּ֣רְא יַעֲקֹ֔ב כִּ֥י יֶשׁ־שֶׁ֖בֶר בְּמִצְרָ֑יִם וַיֹּ֤אמֶר יַעֲקֹב֙ לְבָנָ֔יו לָ֖מָּה תִּתְרָאֽוּ׃ \Rashi{\textbf{וירא יעקב כי יש שבר במצרים.}ומהיכן ראה? והלא לא ראה אלא שמע, שנאמר הנה שמעתי וגו' ומהו וירא? ראה באספקלריא של קדש שעדין יש לו שבר במצרים; ולא היתה נבואה ממש להודיעו בפרוש שזה יוסף: למה תתראו. למה תראו עצמכם בפני בני ישמעאל ובני עשו כאלו אתם שבעים? באותה שעה עדין היה להם תבואה (תענית י'). ולי נראה פשוטו למה תתראו, למה יהיו הכל מסתכלין בכם ומתמיהים בכם, שאין אתם מבקשים לכם אכל בטרם שיכלה מה שבידכם. ומפי אחרים שמעתי, שהוא לשון כחישה, למה תהיו כחושים ברעב, ודומה לו ומרוה גם הוא יורא (משלי י"א):}%
% ---------- פרק 42 | פסוק 2 ----------
 \textbf{ב} וַיֹּ֕אמֶר הִנֵּ֣ה שָׁמַ֔עְתִּי כִּ֥י יֶשׁ־שֶׁ֖בֶר בְּמִצְרָ֑יִם רְדוּ־שָׁ֙מָּה֙ וְשִׁבְרוּ־לָ֣נוּ מִשָּׁ֔ם וְנִחְיֶ֖ה וְלֹ֥א נָמֽוּת׃ \Rashi{\textbf{רדו שמה.}ולא אמר לכו, רמז למאתים ועשר שנים שנשתעבדו למצרים כמנין רד"ו:}%
% ---------- פרק 42 | פסוק 3 ----------
 \textbf{ג} וַיֵּרְד֥וּ אֲחֵֽי־יוֹסֵ֖ף עֲשָׂרָ֑ה לִשְׁבֹּ֥ר בָּ֖ר מִמִּצְרָֽיִם׃ \Rashi{\textbf{וירדו אחי יוסף.}ולא כתב בני יעקב, מלמד שהיו מתחרטים במכירתו ונתנו לבם להתנהג עמו באחוה ולפדותו בכל ממון שיפסקו עליהם: עשרה. מה תלמוד לומר, והלא כתיב ואת בנימין אחי יוסף לא שלח? אלא לענין האחוה היו חלוקין לי' שלא היתה אהבת כלם ושנאת כלם שוה לו; אבל לענין לשבר בר כלם לב אחד להם (בראשית רבה):}%
% ---------- פרק 42 | פסוק 4 ----------
 \textbf{ד} וְאֶת־בִּנְיָמִין֙ אֲחִ֣י יוֹסֵ֔ף לֹא־שָׁלַ֥ח יַעֲקֹ֖ב אֶת־אֶחָ֑יו כִּ֣י אָמַ֔ר פֶּן־יִקְרָאֶ֖נּוּ אָסֽוֹן׃ \Rashi{\textbf{פן יקראנו אסון.}ובבית לא יקראנו? אמר רבי אליעזר בן יעקב, מכאן שהשטן מקטרג בשעת הסכנה (בראשית רבה):}%
% ---------- פרק 42 | פסוק 5 ----------
 \textbf{ה} וַיָּבֹ֙אוּ֙ בְּנֵ֣י יִשְׂרָאֵ֔ל לִשְׁבֹּ֖ר בְּת֣וֹךְ הַבָּאִ֑ים כִּֽי־הָיָ֥ה הָרָעָ֖ב בְּאֶ֥רֶץ כְּנָֽעַן׃ \Rashi{\textbf{בתוך הבאים.}מטמינין עצמם שלא יכירום, לפי שצוה להם אביהם שלא יתראו כלם בפתח א', אלא שיכנס כל א' בפתחו, כדי שלא תשלט בהם עין הרע, שכלם נאים וכלם גבורים:}%
% ---------- פרק 42 | פסוק 6 ----------
 \textbf{ו} וְיוֹסֵ֗ף ה֚וּא הַשַּׁלִּ֣יט עַל־הָאָ֔רֶץ ה֥וּא הַמַּשְׁבִּ֖יר לְכׇל־עַ֣ם הָאָ֑רֶץ וַיָּבֹ֙אוּ֙ אֲחֵ֣י יוֹסֵ֔ף וַיִּשְׁתַּֽחֲווּ־ל֥וֹ אַפַּ֖יִם אָֽרְצָה׃ \Rashi{\textbf{וישתחוו לו אפים.}נשתטחו לו על פניהם, וכן כל השתחואה פשוט ידים ורגלים הוא:}%
% ---------- פרק 42 | פסוק 7 ----------
 \textbf{ז} וַיַּ֥רְא יוֹסֵ֛ף אֶת־אֶחָ֖יו וַיַּכִּרֵ֑ם וַיִּתְנַכֵּ֨ר אֲלֵיהֶ֜ם וַיְדַבֵּ֧ר אִתָּ֣ם קָשׁ֗וֹת וַיֹּ֤אמֶר אֲלֵהֶם֙ מֵאַ֣יִן בָּאתֶ֔ם וַיֹּ֣אמְר֔וּ מֵאֶ֥רֶץ כְּנַ֖עַן לִשְׁבׇּר־אֹֽכֶל׃ \Rashi{\textbf{ויתנכר אליהם.}נעשה להם כנכרי בדברים לדבר קשות:}%
% ---------- פרק 42 | פסוק 8 ----------
 \textbf{ח} וַיַּכֵּ֥ר יוֹסֵ֖ף אֶת־אֶחָ֑יו וְהֵ֖ם לֹ֥א הִכִּרֻֽהוּ׃ \Rashi{\textbf{ויכר יוסף וגו'.}לפי שהניחם חתומי זקן: והם לא הכרהו. שיצא מאצלם בלא חתימת זקן, ועכשו בא בחתימת זקן (כתובות כ"ז): ומ"א ויכר יוסף את אחיו כשנמסרו בידו הכיר שהם אחיו ורחם עליהם, והם לא הכירוהו כשנפל בידם לנהג בו אחוה:}%
% ---------- פרק 42 | פסוק 9 ----------
 \textbf{ט} וַיִּזְכֹּ֣ר יוֹסֵ֔ף אֵ֚ת הַחֲלֹמ֔וֹת אֲשֶׁ֥ר חָלַ֖ם לָהֶ֑ם וַיֹּ֤אמֶר אֲלֵהֶם֙ מְרַגְּלִ֣ים אַתֶּ֔ם לִרְא֛וֹת אֶת־עֶרְוַ֥ת הָאָ֖רֶץ בָּאתֶֽם׃ \Rashi{\textbf{אשר חלם להם.}עליהם; וידע שנתקימו שהרי השתחוו לו: ערות הארץ. גלוי הארץ – מהיכן היא נוחה לכבש, כמו את מקרה הערה (ויקרא כ'), וכמו ערם ועריה (יחזקאל ט"ז) וכן כל ערוה שבמקרא ל' גלוי, ות"א בדקא דארעא, כמו בדק הבית (מלכים ב י"ב), רעוע הבית, אבל לא דקדק לפרשו אחר לשון המקרא:}%
% ---------- פרק 42 | פסוק 10 ----------
 \textbf{י} וַיֹּאמְר֥וּ אֵלָ֖יו לֹ֣א אֲדֹנִ֑י וַעֲבָדֶ֥יךָ בָּ֖אוּ לִשְׁבׇּר־אֹֽכֶל׃ \Rashi{\textbf{לא אדני.}לא תאמר כן, שהרי עבדיך באו לשבר אכל:}%
% ---------- פרק 42 | פסוק 11 ----------
 \textbf{יא} כֻּלָּ֕נוּ בְּנֵ֥י אִישׁ־אֶחָ֖ד נָ֑חְנוּ כֵּנִ֣ים אֲנַ֔חְנוּ לֹא־הָי֥וּ עֲבָדֶ֖יךָ מְרַגְּלִֽים׃ \Rashi{\textbf{כלנו בני איש אחד נחנו.}נצנצה בהם רוח הקדש וכללוהו עמהם, שאף הוא בן אביהם: כנים. אמתיים, כמו כן דברת (שמות י'), כן בנות צלפחד דברת (במדבר כ"ז), ועברתו לא כן בדיו (ישעיהו ט"ז):}%
% ---------- פרק 42 | פסוק 12 ----------
 \textbf{יב} וַיֹּ֖אמֶר אֲלֵהֶ֑ם לֹ֕א כִּֽי־עֶרְוַ֥ת הָאָ֖רֶץ בָּאתֶ֥ם לִרְאֽוֹת׃ \Rashi{\textbf{כי ערות הארץ באתם לראות.}שהרי נכנסתם בי' שערי העיר, למה לא נכנסתם בשער א'? (בראשית רבה):}%
% ---------- פרק 42 | פסוק 13 ----------
 \textbf{יג} וַיֹּאמְר֗וּ שְׁנֵ֣ים עָשָׂר֩ עֲבָדֶ֨יךָ אַחִ֧ים ׀ אֲנַ֛חְנוּ בְּנֵ֥י אִישׁ־אֶחָ֖ד בְּאֶ֣רֶץ כְּנָ֑עַן וְהִנֵּ֨ה הַקָּטֹ֤ן אֶת־אָבִ֙ינוּ֙ הַיּ֔וֹם וְהָאֶחָ֖ד אֵינֶֽנּוּ׃ \Rashi{\textbf{ויאמרו שנים עשר עבדיך וגו'.}ובשביל אותו א' שאיננו נתפזרנו בעיר לבקשו:}%
% ---------- פרק 42 | פסוק 14 ----------
 \textbf{יד} וַיֹּ֥אמֶר אֲלֵהֶ֖ם יוֹסֵ֑ף ה֗וּא אֲשֶׁ֨ר דִּבַּ֧רְתִּי אֲלֵכֶ֛ם לֵאמֹ֖ר מְרַגְּלִ֥ים אַתֶּֽם׃ \Rashi{\textbf{הוא אשר דברתי.}הדבר אשר דברתי שאתם מרגלים, הוא האמת והנכון, זהו לפי פשוטו. ומדרשו אמר להם ואלו מצאתם אותו ויפסקו עליכם ממון הרבה, תפדוהו? אמרו לו הן, אמר להם ואם יאמרו לכם שלא יחזירוהו בשום ממון, מה תעשו? אמרו לו לכך באנו, להרג או להרג, אמר להם הוא אשר דברתי אליכם, להרג בני העיר באתם, מנחש אני בגביע שלי ששנים מכם החריבו כרך גדול של שכם:}%
% ---------- פרק 42 | פסוק 15 ----------
 \textbf{טו} בְּזֹ֖את תִּבָּחֵ֑נוּ חֵ֤י פַרְעֹה֙ אִם־תֵּצְא֣וּ מִזֶּ֔ה כִּ֧י אִם־בְּב֛וֹא אֲחִיכֶ֥ם הַקָּטֹ֖ן הֵֽנָּה׃ \Rashi{\textbf{חי פרעה.}אם יחיה פרעה. כשהיה נשבע לשקר, היה נשבע בחיי פרעה: אם תצאו מזה. מן המקום הזה:}%
% ---------- פרק 42 | פסוק 16 ----------
 \textbf{טז} שִׁלְח֨וּ מִכֶּ֣ם אֶחָד֮ וְיִקַּ֣ח אֶת־אֲחִיכֶם֒ וְאַתֶּם֙ הֵאָ֣סְר֔וּ וְיִבָּֽחֲנוּ֙ דִּבְרֵיכֶ֔ם הַֽאֱמֶ֖ת אִתְּכֶ֑ם וְאִם־לֹ֕א חֵ֣י פַרְעֹ֔ה כִּ֥י מְרַגְּלִ֖ים אַתֶּֽם׃ \Rashi{\textbf{האמת אתכם.}אם אמת אתכם, לפיכך ה"א נקוד פתח, שהוא כמו בלשון תמה ואם לא תביאוהו, חי פרעה, כי מרגלים אתם:}%
% ---------- פרק 42 | פסוק 17 ----------
 \textbf{יז} וַיֶּאֱסֹ֥ף אֹתָ֛ם אֶל־מִשְׁמָ֖ר שְׁלֹ֥שֶׁת יָמִֽים׃ \Rashi{\textbf{משמר.}בית האסורים:}%
% ---------- פרק 42 | פסוק 18 ----------
 \textbf{יח} וַיֹּ֨אמֶר אֲלֵהֶ֤ם יוֹסֵף֙ בַּיּ֣וֹם הַשְּׁלִישִׁ֔י זֹ֥את עֲשׂ֖וּ וִֽחְי֑וּ אֶת־הָאֱלֹהִ֖ים אֲנִ֥י יָרֵֽא׃ %
% ---------- פרק 42 | פסוק 19 ----------
 \textbf{יט} אִם־כֵּנִ֣ים אַתֶּ֔ם אֲחִיכֶ֣ם אֶחָ֔ד יֵאָסֵ֖ר בְּבֵ֣ית מִשְׁמַרְכֶ֑ם וְאַתֶּם֙ לְכ֣וּ הָבִ֔יאוּ שֶׁ֖בֶר רַעֲב֥וֹן בָּתֵּיכֶֽם׃ \Rashi{\textbf{בבית משמרכם.}שאתם אסורים בו עכשו: ואתם לכו הביאו. לבית אביכם: שבר רעבון בתיכם. מה שקניתם לרעבון אנשי בתיכם:}%
% ---------- פרק 42 | פסוק 20 ----------
 \textbf{כ} וְאֶת־אֲחִיכֶ֤ם הַקָּטֹן֙ תָּבִ֣יאוּ אֵלַ֔י וְיֵאָמְנ֥וּ דִבְרֵיכֶ֖ם וְלֹ֣א תָמ֑וּתוּ וַיַּעֲשׂוּ־כֵֽן׃ \Rashi{\textbf{ויאמנו דבריכם.}יתאמתו ויתקימו; כמו אמן אמן (במדבר ה'), וכמו, יאמן נא דברך (מלכים א ח'):}%
% ---------- פרק 42 | פסוק 21 ----------
 \textbf{כא} וַיֹּאמְר֞וּ אִ֣ישׁ אֶל־אָחִ֗יו אֲבָל֮ אֲשֵׁמִ֣ים ׀ אֲנַ֘חְנוּ֮ עַל־אָחִ֒ינוּ֒ אֲשֶׁ֨ר רָאִ֜ינוּ צָרַ֥ת נַפְשׁ֛וֹ בְּהִתְחַֽנְנ֥וֹ אֵלֵ֖ינוּ וְלֹ֣א שָׁמָ֑עְנוּ עַל־כֵּן֙ בָּ֣אָה אֵלֵ֔ינוּ הַצָּרָ֖ה הַזֹּֽאת׃ \Rashi{\textbf{אבל.}כתרגומו בקושטא, וראיתי בב"ר לישנא דרומאה הוא אבל – ברם: באה אלינו. טעמו בבי"ת, לפי שהוא בלשון עבר, שכבר באה, ותרגומו אתת לנא:}%
% ---------- פרק 42 | פסוק 22 ----------
 \textbf{כב} וַיַּ֩עַן֩ רְאוּבֵ֨ן אֹתָ֜ם לֵאמֹ֗ר הֲלוֹא֩ אָמַ֨רְתִּי אֲלֵיכֶ֧ם ׀ לֵאמֹ֛ר אַל־תֶּחֶטְא֥וּ בַיֶּ֖לֶד וְלֹ֣א שְׁמַעְתֶּ֑ם וְגַם־דָּמ֖וֹ הִנֵּ֥ה נִדְרָֽשׁ׃ \Rashi{\textbf{וגם דמו.}אתין וגמין רבויין – דמו וגם דם הזקן:}%
% ---------- פרק 42 | פסוק 23 ----------
 \textbf{כג} וְהֵם֙ לֹ֣א יָֽדְע֔וּ כִּ֥י שֹׁמֵ֖עַ יוֹסֵ֑ף כִּ֥י הַמֵּלִ֖יץ בֵּינֹתָֽם׃ \Rashi{\textbf{והם לא ידעו כי שמע יוסף.}מבין לשונם, ובפניו היו מדברים כן: כי המליץ בינתם. כי כשהיו מדברים עמו היה המליץ ביניהם היודע לשון עברי ולשון מצרי, והיה מליץ דבריהם ליוסף ודברי יוסף להם, לכך היו סבורים שאין יוסף מכיר בלשון עברי: המליץ. זה מנשה:}%
% ---------- פרק 42 | פסוק 24 ----------
 \textbf{כד} וַיִּסֹּ֥ב מֵֽעֲלֵיהֶ֖ם וַיֵּ֑בְךְּ וַיָּ֤שׇׁב אֲלֵהֶם֙ וַיְדַבֵּ֣ר אֲלֵהֶ֔ם וַיִּקַּ֤ח מֵֽאִתָּם֙ אֶת־שִׁמְע֔וֹן וַיֶּאֱסֹ֥ר אֹת֖וֹ לְעֵינֵיהֶֽם׃ \Rashi{\textbf{ויסב מעליהם.}נתרחק מעליהם שלא יראוהו בוכה: ויבך. לפי ששמע שהיו מתחרטין: את שמעון. הוא השליכו לבור, הוא שאמר ללוי הנה בעל החלמות הלזה בא (בראשית רבה). דבר אחר נתכון יוסף להפרידו מלוי, שמא יתיעצו שניהם להרג אותו: ויאסר אתו לעיניהם. לא אסרו אלא לעיניהם, וכיון שיצאו הוציאו והאכילו והשקהו (בראשית רבה):}%
% ---------- פרק 42 | פסוק 25 ----------
 \textbf{כה} וַיְצַ֣ו יוֹסֵ֗ף וַיְמַלְא֣וּ אֶת־כְּלֵיהֶם֮ בָּר֒ וּלְהָשִׁ֤יב כַּסְפֵּיהֶם֙ אִ֣ישׁ אֶל־שַׂקּ֔וֹ וְלָתֵ֥ת לָהֶ֛ם צֵדָ֖ה לַדָּ֑רֶךְ וַיַּ֥עַשׂ לָהֶ֖ם כֵּֽן׃ %
% ---------- פרק 42 | פסוק 26 ----------
 \textbf{כו} וַיִּשְׂא֥וּ אֶת־שִׁבְרָ֖ם עַל־חֲמֹרֵיהֶ֑ם וַיֵּלְכ֖וּ מִשָּֽׁם׃ %
% ---------- פרק 42 | פסוק 27 ----------
 \textbf{כז} וַיִּפְתַּ֨ח הָאֶחָ֜ד אֶת־שַׂקּ֗וֹ לָתֵ֥ת מִסְפּ֛וֹא לַחֲמֹר֖וֹ בַּמָּל֑וֹן וַיַּרְא֙ אֶת־כַּסְפּ֔וֹ וְהִנֵּה־ה֖וּא בְּפִ֥י אַמְתַּחְתּֽוֹ׃ \Rashi{\textbf{ויפתח האחד.}הוא לוי, שנשאר יחיד משמעון בן זוגו: במלון. במקום שלנו בלילה: אמתחתו. הוא שק:}%
% ---------- פרק 42 | פסוק 28 ----------
 \textbf{כח} וַיֹּ֤אמֶר אֶל־אֶחָיו֙ הוּשַׁ֣ב כַּסְפִּ֔י וְגַ֖ם הִנֵּ֣ה בְאַמְתַּחְתִּ֑י וַיֵּצֵ֣א לִבָּ֗ם וַיֶּֽחֶרְד֞וּ אִ֤ישׁ אֶל־אָחִיו֙ לֵאמֹ֔ר מַה־זֹּ֛את עָשָׂ֥ה אֱלֹהִ֖ים לָֽנוּ׃ \Rashi{\textbf{וגם הנה באמתחתי.}גם הכסף בו עם התבואה: מה זאת עשה אלהים לנו. להביאנו לידי עלילה זו, שלא הושב אלא להתעולל עלינו:}%
% ---------- פרק 42 | פסוק 29 ----------
 \textbf{כט} וַיָּבֹ֛אוּ אֶל־יַעֲקֹ֥ב אֲבִיהֶ֖ם אַ֣רְצָה כְּנָ֑עַן וַיַּגִּ֣ידוּ ל֔וֹ אֵ֛ת כׇּל־הַקֹּרֹ֥ת אֹתָ֖ם לֵאמֹֽר׃ %
% ---------- פרק 42 | פסוק 30 ----------
 \textbf{ל} דִּ֠בֶּ֠ר הָאִ֨ישׁ אֲדֹנֵ֥י הָאָ֛רֶץ אִתָּ֖נוּ קָשׁ֑וֹת וַיִּתֵּ֣ן אֹתָ֔נוּ כִּֽמְרַגְּלִ֖ים אֶת־הָאָֽרֶץ׃ %
% ---------- פרק 42 | פסוק 31 ----------
 \textbf{לא} וַנֹּ֥אמֶר אֵלָ֖יו כֵּנִ֣ים אֲנָ֑חְנוּ לֹ֥א הָיִ֖ינוּ מְרַגְּלִֽים׃ %
% ---------- פרק 42 | פסוק 32 ----------
 \textbf{לב} שְׁנֵים־עָשָׂ֥ר אֲנַ֛חְנוּ אַחִ֖ים בְּנֵ֣י אָבִ֑ינוּ הָאֶחָ֣ד אֵינֶ֔נּוּ וְהַקָּטֹ֥ן הַיּ֛וֹם אֶת־אָבִ֖ינוּ בְּאֶ֥רֶץ כְּנָֽעַן׃ %
% ---------- פרק 42 | פסוק 33 ----------
 \textbf{לג} וַיֹּ֣אמֶר אֵלֵ֗ינוּ הָאִישׁ֙ אֲדֹנֵ֣י הָאָ֔רֶץ בְּזֹ֣את אֵדַ֔ע כִּ֥י כֵנִ֖ים אַתֶּ֑ם אֲחִיכֶ֤ם הָֽאֶחָד֙ הַנִּ֣יחוּ אִתִּ֔י וְאֶת־רַעֲב֥וֹן בָּתֵּיכֶ֖ם קְח֥וּ וָלֵֽכוּ׃ %
% ---------- פרק 42 | פסוק 34 ----------
 \textbf{לד} וְ֠הָבִ֠יאוּ אֶת־אֲחִיכֶ֣ם הַקָּטֹן֮ אֵלַי֒ וְאֵֽדְעָ֗ה כִּ֣י לֹ֤א מְרַגְּלִים֙ אַתֶּ֔ם כִּ֥י כֵנִ֖ים אַתֶּ֑ם אֶת־אֲחִיכֶם֙ אֶתֵּ֣ן לָכֶ֔ם וְאֶת־הָאָ֖רֶץ תִּסְחָֽרוּ׃ \Rashi{\textbf{ואת הארץ תסחרו.}תסובבו; וכל לשון סוחרים וסחורה על שם שמחזרים וסובבים אחר הפרקמטיא:}%
% ---------- פרק 42 | פסוק 35 ----------
 \textbf{לה} וַיְהִ֗י הֵ֚ם מְרִיקִ֣ים שַׂקֵּיהֶ֔ם וְהִנֵּה־אִ֥ישׁ צְרוֹר־כַּסְפּ֖וֹ בְּשַׂקּ֑וֹ וַיִּרְא֞וּ אֶת־צְרֹר֧וֹת כַּסְפֵּיהֶ֛ם הֵ֥מָּה וַאֲבִיהֶ֖ם וַיִּירָֽאוּ׃ \Rashi{\textbf{צרור כספו.}קשר כספו:}%
% ---------- פרק 42 | פסוק 36 ----------
 \textbf{לו} וַיֹּ֤אמֶר אֲלֵהֶם֙ יַעֲקֹ֣ב אֲבִיהֶ֔ם אֹתִ֖י שִׁכַּלְתֶּ֑ם יוֹסֵ֤ף אֵינֶ֙נּוּ֙ וְשִׁמְע֣וֹן אֵינֶ֔נּוּ וְאֶת־בִּנְיָמִ֣ן תִּקָּ֔חוּ עָלַ֖י הָי֥וּ כֻלָּֽנָה׃ \Rashi{\textbf{אתי שכלתם.}מלמד שחשדן שמא הרגוהו או מכרוהו כיוסף: שכלתם. כל מי שבניו אבודים קרוי שכול:}%
% ---------- פרק 42 | פסוק 37 ----------
 \textbf{לז} וַיֹּ֤אמֶר רְאוּבֵן֙ אֶל־אָבִ֣יו לֵאמֹ֔ר אֶת־שְׁנֵ֤י בָנַי֙ תָּמִ֔ית אִם־לֹ֥א אֲבִיאֶ֖נּוּ אֵלֶ֑יךָ תְּנָ֤ה אֹתוֹ֙ עַל־יָדִ֔י וַאֲנִ֖י אֲשִׁיבֶ֥נּוּ אֵלֶֽיךָ׃ %
% ---------- פרק 42 | פסוק 38 ----------
 \textbf{לח} וַיֹּ֕אמֶר לֹֽא־יֵרֵ֥ד בְּנִ֖י עִמָּכֶ֑ם כִּֽי־אָחִ֨יו מֵ֜ת וְה֧וּא לְבַדּ֣וֹ נִשְׁאָ֗ר וּקְרָאָ֤הוּ אָסוֹן֙ בַּדֶּ֙רֶךְ֙ אֲשֶׁ֣ר תֵּֽלְכוּ־בָ֔הּ וְהוֹרַדְתֶּ֧ם אֶת־שֵׂיבָתִ֛י בְּיָג֖וֹן שְׁאֽוֹלָה׃ \Rashi{\textbf{לא ירד בני עמכם.}לא קבל דבריו של ראובן, אמר בכור שוטה הוא זה, הוא אומר להמית בניו, וכי בניו הם ולא בני? (בראשית רבה):}%
% ---------- פרק 43 | פסוק 1 ----------
 \textbf{א} וְהָרָעָ֖ב כָּבֵ֥ד בָּאָֽרֶץ׃ %
% ---------- פרק 43 | פסוק 2 ----------
 \textbf{ב} וַיְהִ֗י כַּאֲשֶׁ֤ר כִּלּוּ֙ לֶאֱכֹ֣ל אֶת־הַשֶּׁ֔בֶר אֲשֶׁ֥ר הֵבִ֖יאוּ מִמִּצְרָ֑יִם וַיֹּ֤אמֶר אֲלֵיהֶם֙ אֲבִיהֶ֔ם שֻׁ֖בוּ שִׁבְרוּ־לָ֥נוּ מְעַט־אֹֽכֶל׃ \Rashi{\textbf{כאשר כלו לאכל.}יהודה אמר להם המתינו לזקן עד שתכלה פת מן הבית (בראשית רבה): כאשר כלו. כד שציאו, והמתרגם כד ספיקו טועה; כאשר כלו הגמלים לשתות מתרגם כד ספיקו – כששתו די ספוקם הוא גמר שתיתם; אבל זה כאשר כלו לאכל, כאשר תם האכל הוא, ומתרגמינן כד שציאו:}%
% ---------- פרק 43 | פסוק 3 ----------
 \textbf{ג} וַיֹּ֧אמֶר אֵלָ֛יו יְהוּדָ֖ה לֵאמֹ֑ר הָעֵ֣ד הֵעִד֩ בָּ֨נוּ הָאִ֤ישׁ לֵאמֹר֙ לֹֽא־תִרְא֣וּ פָנַ֔י בִּלְתִּ֖י אֲחִיכֶ֥ם אִתְּכֶֽם׃ \Rashi{\textbf{העד העד.}לשון התראה שסתם התראה מתרה בו בפני עדים, וכן העדתי באבותיכם (ירמיהו י"א); רד העד בעם (שמות י"ט): לא תראו פני בלתי אחיכם אתכם. לא תראוני בלא אחיכם אתכם; ואנקלוס תרגם אלהן כד אחוכון עמכון, ישב הדבר על אפנו ולא דקדק לתרגם אחר לשון המקרא:}%
% ---------- פרק 43 | פסוק 4 ----------
 \textbf{ד} אִם־יֶשְׁךָ֛ מְשַׁלֵּ֥חַ אֶת־אָחִ֖ינוּ אִתָּ֑נוּ נֵרְדָ֕ה וְנִשְׁבְּרָ֥ה לְךָ֖ אֹֽכֶל׃ %
% ---------- פרק 43 | פסוק 5 ----------
 \textbf{ה} וְאִם־אֵינְךָ֥ מְשַׁלֵּ֖חַ לֹ֣א נֵרֵ֑ד כִּֽי־הָאִ֞ישׁ אָמַ֤ר אֵלֵ֙ינוּ֙ לֹֽא־תִרְא֣וּ פָנַ֔י בִּלְתִּ֖י אֲחִיכֶ֥ם אִתְּכֶֽם׃ %
% ---------- פרק 43 | פסוק 6 ----------
 \textbf{ו} וַיֹּ֙אמֶר֙ יִשְׂרָאֵ֔ל לָמָ֥ה הֲרֵעֹתֶ֖ם לִ֑י לְהַגִּ֣יד לָאִ֔ישׁ הַע֥וֹד לָכֶ֖ם אָֽח׃ %
% ---------- פרק 43 | פסוק 7 ----------
 \textbf{ז} וַיֹּאמְר֡וּ שָׁא֣וֹל שָֽׁאַל־הָ֠אִ֠ישׁ לָ֣נוּ וּלְמֽוֹלַדְתֵּ֜נוּ לֵאמֹ֗ר הַע֨וֹד אֲבִיכֶ֥ם חַי֙ הֲיֵ֣שׁ לָכֶ֣ם אָ֔ח וַנַּ֨גֶּד־ל֔וֹ עַל־פִּ֖י הַדְּבָרִ֣ים הָאֵ֑לֶּה הֲיָד֣וֹעַ נֵדַ֔ע כִּ֣י יֹאמַ֔ר הוֹרִ֖ידוּ אֶת־אֲחִיכֶֽם׃ \Rashi{\textbf{לנו ולמולדתנו.}למשפחותנו. ומדרשו אפלו עצי עריסותנו גלה לנו: ונגד לו. שיש לנו אב ואח: על פי הדברים האלה. על פי שאלותיו אשר שאל, הזקקנו להגיד: כי יאמר. אשר יאמר, כי משמש בלשון אם ואם משמש בלשון אשר; הרי זה שמוש אחד מארבע לשונותיו שמשמש כי, והוא אם, שהרי כי זה כמו אם, כמו עד אם דברתי דברי (בראשית כ״ד:ל״ג):}%
% ---------- פרק 43 | פסוק 8 ----------
 \textbf{ח} וַיֹּ֨אמֶר יְהוּדָ֜ה אֶל־יִשְׂרָאֵ֣ל אָבִ֗יו שִׁלְחָ֥ה הַנַּ֛עַר אִתִּ֖י וְנָק֣וּמָה וְנֵלֵ֑כָה וְנִֽחְיֶה֙ וְלֹ֣א נָמ֔וּת גַּם־אֲנַ֥חְנוּ גַם־אַתָּ֖ה גַּם־טַפֵּֽנוּ׃ \Rashi{\textbf{ונחיה.}נצנצה בו רוח הקדש, על ידי הליכה זו תחי רוחך, שנאמר ותחי רוח יעקב אביהם (בראשית מ"ה): ולא נמות. ברעב; בנימין ספק יתפש ספק לא יתפש, ואנו כלנו מתים ברעב אם לא נלך, מוטב שתניח את הספק ותתפש את הודאי:}%
% ---------- פרק 43 | פסוק 9 ----------
 \textbf{ט} אָֽנֹכִי֙ אֶֽעֶרְבֶ֔נּוּ מִיָּדִ֖י תְּבַקְשֶׁ֑נּוּ אִם־לֹ֨א הֲבִיאֹתִ֤יו אֵלֶ֙יךָ֙ וְהִצַּגְתִּ֣יו לְפָנֶ֔יךָ וְחָטָ֥אתִֽי לְךָ֖ כׇּל־הַיָּמִֽים׃ \Rashi{\textbf{והצגתיו לפניך.}שלא אביאנו אליך מת, כי אם חי: וחטאתי לך כל הימים. לעולם הבא:}%
% ---------- פרק 43 | פסוק 10 ----------
 \textbf{י} כִּ֖י לוּלֵ֣א הִתְמַהְמָ֑הְנוּ כִּֽי־עַתָּ֥ה שַׁ֖בְנוּ זֶ֥ה פַעֲמָֽיִם׃ \Rashi{לולא התמהמנו על ידך, כּבר היינוּ שׁבים עם שׁמעוֹן ולא נצטערתּ כּל היּמים הלּלוּ:}%
% ---------- פרק 43 | פסוק 11 ----------
 \textbf{יא} וַיֹּ֨אמֶר אֲלֵהֶ֜ם יִשְׂרָאֵ֣ל אֲבִיהֶ֗ם אִם־כֵּ֣ן ׀ אֵפוֹא֮ זֹ֣את עֲשׂוּ֒ קְח֞וּ מִזִּמְרַ֤ת הָאָ֙רֶץ֙ בִּכְלֵיכֶ֔ם וְהוֹרִ֥ידוּ לָאִ֖ישׁ מִנְחָ֑ה מְעַ֤ט צֳרִי֙ וּמְעַ֣ט דְּבַ֔שׁ נְכֹ֣את וָלֹ֔ט בׇּטְנִ֖ים וּשְׁקֵדִֽים׃ \Rashi{\textbf{אפוא.}לשון יתר הוא, לתקן מלה בלשון עברי, אם כן אזדקק לעשות, שאשלחנו עמכם, צריך אני לחזר ולבקש איה פה תקנה ועצה להשיאכם, ואומר אני זאת עשו: מזמרת הארץ. מדמשבח בארעא – שהכל מזמרים עליו כשהוא בא לעולם: נכאת. שעוה: בטנים. לא ידעתי מה הם, ובפרושי א"ב של רבי מכיר ראיתי פשט"ציאס ודומה לי שהם אפרסקין:}%
% ---------- פרק 43 | פסוק 12 ----------
 \textbf{יב} וְכֶ֥סֶף מִשְׁנֶ֖ה קְח֣וּ בְיֶדְכֶ֑ם וְאֶת־הַכֶּ֜סֶף הַמּוּשָׁ֨ב בְּפִ֤י אַמְתְּחֹֽתֵיכֶם֙ תָּשִׁ֣יבוּ בְיֶדְכֶ֔ם אוּלַ֥י מִשְׁגֶּ֖ה הֽוּא׃ \Rashi{\textbf{וכסף משנה.}פי שנים כראשון: קחו בידכם. לשבר אכל, שמא הוקר השער: אולי משגה הוא. שמא הממנה על הבית שכחו שוגג:}%
% ---------- פרק 43 | פסוק 13 ----------
 \textbf{יג} וְאֶת־אֲחִיכֶ֖ם קָ֑חוּ וְק֖וּמוּ שׁ֥וּבוּ אֶל־הָאִֽישׁ׃ %
% ---------- פרק 43 | פסוק 14 ----------
 \textbf{יד} וְאֵ֣ל שַׁדַּ֗י יִתֵּ֨ן לָכֶ֤ם רַחֲמִים֙ לִפְנֵ֣י הָאִ֔ישׁ וְשִׁלַּ֥ח לָכֶ֛ם אֶת־אֲחִיכֶ֥ם אַחֵ֖ר וְאֶת־בִּנְיָמִ֑ין וַאֲנִ֕י כַּאֲשֶׁ֥ר שָׁכֹ֖לְתִּי שָׁכָֽלְתִּי׃ \Rashi{\textbf{ואל שדי.}מעתה אינכם חסרים כלום אלא תפלה, הריני מתפלל עליכם: אל שדי. שדי בנתינת רחמיו וכדי היכלת בידו לתן, יתן לכם רחמים, זהו פשוטו. ומדרשו מי שאמר לעולם די, יאמר די לצרותי, שלא שקטתי מנעורי, צרת לבן, צרת עשו, צרת רחל, צרת דינה, צרת יוסף, צרת שמעון, צרת בנימין: ושלח לכם. ויפטר לכון כתרגומו – יפטרנו מאסוריו, לשון לחפשי ישלחנו (שמות כ"א), ואינו נופל בתרגום לשון וישלח, שהרי לשם הם הולכים אצלו: את אחיכם. זה שמעון: אחר. רוח הקדש נזרקה בו, לרבות יוסף (בראשית רבה): ואני. עד שובכם אהיה שכול מספק: כאשר שכלתי. מיוסף ומשמעון: שכלתי. מבנימין:}%
% ---------- פרק 43 | פסוק 15 ----------
 \textbf{טו} וַיִּקְח֤וּ הָֽאֲנָשִׁים֙ אֶת־הַמִּנְחָ֣ה הַזֹּ֔את וּמִשְׁנֶה־כֶּ֛סֶף לָקְח֥וּ בְיָדָ֖ם וְאֶת־בִּנְיָמִ֑ן וַיָּקֻ֙מוּ֙ וַיֵּרְד֣וּ מִצְרַ֔יִם וַיַּֽעַמְד֖וּ לִפְנֵ֥י יוֹסֵֽף׃ \Rashi{\textbf{ואת בנימין.}מתרגמינן ודברו ית בנימין, לפי שאין לקיחת הכסף ולקיחת האדם שוה בלשון ארמי; בדבר הנקח ביד מתרגמינן ונסיב, ודבר הנקח בהנהגת דברים, מתרגמינן ודבר:}%
% ---------- פרק 43 | פסוק 16 ----------
 \textbf{טז} וַיַּ֨רְא יוֹסֵ֣ף אִתָּם֮ אֶת־בִּנְיָמִין֒ וַיֹּ֙אמֶר֙ לַֽאֲשֶׁ֣ר עַל־בֵּית֔וֹ הָבֵ֥א אֶת־הָאֲנָשִׁ֖ים הַבָּ֑יְתָה וּטְבֹ֤חַ טֶ֙בַח֙ וְהָכֵ֔ן כִּ֥י אִתִּ֛י יֹאכְל֥וּ הָאֲנָשִׁ֖ים בַּֽצׇּהֳרָֽיִם׃ \Rashi{\textbf{וטבח טבח והכן.}כמו ולטבח טבח ולהכין, ואין טבח לשון צווי, שהיה לו לומר וטבח: בצהרים. זה מתרגם בשירותא, שהוא לשון סעודה ראשונה בלשון ארמי ובלעז דיזנ"ר, ויש הרבה בתלמוד, שדא לכלבא שירותה; בצע אכולא שירותא, אבל כל תרגום של צהרים טיהרא:}%
% ---------- פרק 43 | פסוק 17 ----------
 \textbf{יז} וַיַּ֣עַשׂ הָאִ֔ישׁ כַּֽאֲשֶׁ֖ר אָמַ֣ר יוֹסֵ֑ף וַיָּבֵ֥א הָאִ֛ישׁ אֶת־הָאֲנָשִׁ֖ים בֵּ֥יתָה יוֹסֵֽף׃ %
% ---------- פרק 43 | פסוק 18 ----------
 \textbf{יח} וַיִּֽירְא֣וּ הָֽאֲנָשִׁ֗ים כִּ֣י הֽוּבְאוּ֮ בֵּ֣ית יוֹסֵף֒ וַיֹּאמְר֗וּ עַל־דְּבַ֤ר הַכֶּ֙סֶף֙ הַשָּׁ֤ב בְּאַמְתְּחֹתֵ֙ינוּ֙ בַּתְּחִלָּ֔ה אֲנַ֖חְנוּ מֽוּבָאִ֑ים לְהִתְגֹּלֵ֤ל עָלֵ֙ינוּ֙ וּלְהִתְנַפֵּ֣ל עָלֵ֔ינוּ וְלָקַ֧חַת אֹתָ֛נוּ לַעֲבָדִ֖ים וְאֶת־חֲמֹרֵֽינוּ׃ \Rashi{\textbf{וייראו האנשים.}כתוב הוא בשני יוד"ין, ותרגומו ודחילו: כי הובאו בית יוסף. ואין דרך שאר הבאים לשבר בר ללון בבית יוסף, כי אם בפנדקאות שבעיר: וייראו. שאין זה אלא לאספם אל משמר: אנחנו מובאים. אל תוך הבית הזה: להתגלל. להיות מתגלגלת עלינו עלילת הכסף ולהיותה נופלת עלינו; ואנקלוס שתרגם ולאסתקפא עלנא הוא לשון להתעולל, כדמתרגמינן עלילת דברים – תסקופי מלין, ולא תרגמו אחר לשון המקרא ולהתגולל שתרגם לאתרברבא, הוא לשון גלת הזהב (קהלת י"ב), והצב גלתה העלתה (נחום ב'), שהוא לשון מלכות:}%
% ---------- פרק 43 | פסוק 19 ----------
 \textbf{יט} וַֽיִּגְּשׁוּ֙ אֶל־הָאִ֔ישׁ אֲשֶׁ֖ר עַל־בֵּ֣ית יוֹסֵ֑ף וַיְדַבְּר֥וּ אֵלָ֖יו פֶּ֥תַח הַבָּֽיִת׃ %
% ---------- פרק 43 | פסוק 20 ----------
 \textbf{כ} וַיֹּאמְר֖וּ בִּ֣י אֲדֹנִ֑י יָרֹ֥ד יָרַ֛דְנוּ בַּתְּחִלָּ֖ה לִשְׁבׇּר־אֹֽכֶל׃ \Rashi{\textbf{בי אדני.}לשון בעיא ותחנונים הוא, בלשון ארמי ביא ביא (יבמות צ"ז:, סנהדרין ס"ד.): ירד ירדנו. ירידה היא לנו, רגילים היינו לפרנס אחרים, עכשו אנו צריכים לך (בראשית רבה):}%
% ---------- פרק 43 | פסוק 21 ----------
 \textbf{כא} וַיְהִ֞י כִּי־בָ֣אנוּ אֶל־הַמָּל֗וֹן וַֽנִּפְתְּחָה֙ אֶת־אַמְתְּחֹתֵ֔ינוּ וְהִנֵּ֤ה כֶֽסֶף־אִישׁ֙ בְּפִ֣י אַמְתַּחְתּ֔וֹ כַּסְפֵּ֖נוּ בְּמִשְׁקָל֑וֹ וַנָּ֥שֶׁב אֹת֖וֹ בְּיָדֵֽנוּ׃ %
% ---------- פרק 43 | פסוק 22 ----------
 \textbf{כב} וְכֶ֧סֶף אַחֵ֛ר הוֹרַ֥דְנוּ בְיָדֵ֖נוּ לִשְׁבׇּר־אֹ֑כֶל לֹ֣א יָדַ֔עְנוּ מִי־שָׂ֥ם כַּסְפֵּ֖נוּ בְּאַמְתְּחֹתֵֽינוּ׃ %
% ---------- פרק 43 | פסוק 23 ----------
 \textbf{כג} וַיֹּ֩אמֶר֩ שָׁל֨וֹם לָכֶ֜ם אַל־תִּירָ֗אוּ אֱלֹ֨הֵיכֶ֜ם וֵֽאלֹהֵ֤י אֲבִיכֶם֙ נָתַ֨ן לָכֶ֤ם מַטְמוֹן֙ בְּאַמְתְּחֹ֣תֵיכֶ֔ם כַּסְפְּכֶ֖ם בָּ֣א אֵלָ֑י וַיּוֹצֵ֥א אֲלֵהֶ֖ם אֶת־שִׁמְעֽוֹן׃ \Rashi{\textbf{אלהיכם.}בזכותכם, ואם אין זכותכם כדאי, ואלהי אביכם, בזכות אביכם נתן לכם מטמון:}%
% ---------- פרק 43 | פסוק 24 ----------
 \textbf{כד} וַיָּבֵ֥א הָאִ֛ישׁ אֶת־הָאֲנָשִׁ֖ים בֵּ֣יתָה יוֹסֵ֑ף וַיִּתֶּן־מַ֙יִם֙ וַיִּרְחֲצ֣וּ רַגְלֵיהֶ֔ם וַיִּתֵּ֥ן מִסְפּ֖וֹא לַחֲמֹֽרֵיהֶֽם׃ \Rashi{\textbf{ויבא האיש.}הבאה אחר הבאה, לפי שהיו דוחפים אותו לחוץ עד שדברו אליו פתח הבית; ומשאמר להם שלום לכם, נמשכו ובאו אחריו:}%
% ---------- פרק 43 | פסוק 25 ----------
 \textbf{כה} וַיָּכִ֙ינוּ֙ אֶת־הַמִּנְחָ֔ה עַד־בּ֥וֹא יוֹסֵ֖ף בַּֽצׇּהֳרָ֑יִם כִּ֣י שָֽׁמְע֔וּ כִּי־שָׁ֖ם יֹ֥אכְלוּ לָֽחֶם׃ \Rashi{\textbf{ויכינו.}הזמינו, עטרוה בכלים נאים:}%
% ---------- פרק 43 | פסוק 26 ----------
 \textbf{כו} וַיָּבֹ֤א יוֹסֵף֙ הַבַּ֔יְתָה וַיָּבִ֥יאּוּ ל֛וֹ אֶת־הַמִּנְחָ֥ה אֲשֶׁר־בְּיָדָ֖ם הַבָּ֑יְתָה וַיִּשְׁתַּחֲווּ־ל֖וֹ אָֽרְצָה׃ \Rashi{\textbf{הביתה.}מפרוזדור לטרקלין:}%
% ---------- פרק 43 | פסוק 27 ----------
 \textbf{כז} וַיִּשְׁאַ֤ל לָהֶם֙ לְשָׁל֔וֹם וַיֹּ֗אמֶר הֲשָׁל֛וֹם אֲבִיכֶ֥ם הַזָּקֵ֖ן אֲשֶׁ֣ר אֲמַרְתֶּ֑ם הַעוֹדֶ֖נּוּ חָֽי׃ %
% ---------- פרק 43 | פסוק 28 ----------
 \textbf{כח} וַיֹּאמְר֗וּ שָׁל֛וֹם לְעַבְדְּךָ֥ לְאָבִ֖ינוּ עוֹדֶ֣נּוּ חָ֑י וַֽיִּקְּד֖וּ (וישתחו) [וַיִּֽשְׁתַּחֲוֽוּ]׃ \Rashi{\textbf{ויקדו וישתחוו.}על שאלת שלום קידה – כפיפת קדקד; השתחואה – משתטח לארץ:}%
% ---------- פרק 43 | פסוק 29 ----------
 \textbf{כט} וַיִּשָּׂ֣א עֵינָ֗יו וַיַּ֞רְא אֶת־בִּנְיָמִ֣ין אָחִיו֮ בֶּן־אִמּוֹ֒ וַיֹּ֗אמֶר הֲזֶה֙ אֲחִיכֶ֣ם הַקָּטֹ֔ן אֲשֶׁ֥ר אֲמַרְתֶּ֖ם אֵלָ֑י וַיֹּאמַ֕ר אֱלֹהִ֥ים יׇחְנְךָ֖ בְּנִֽי׃ \Rashi{\textbf{אלהים יחנך בני.}בשאר שבטים שמענו חנינה אשר חנן אלהים את עבדך (בראשית ל"ג), ובנימין עדין לא נולד, לכך ברכו יוסף בחנינה:}%
% ---------- פרק 43 | פסוק 30 ----------
 \textbf{ל} וַיְמַהֵ֣ר יוֹסֵ֗ף כִּֽי־נִכְמְר֤וּ רַחֲמָיו֙ אֶל־אָחִ֔יו וַיְבַקֵּ֖שׁ לִבְכּ֑וֹת וַיָּבֹ֥א הַחַ֖דְרָה וַיֵּ֥בְךְּ שָֽׁמָּה׃ \Rashi{\textbf{כי נכמרו רחמיו.}שאלו יש לך אח מאם? אמר לו אח היה לי, ואיני יודע היכן הוא; יש לך בנים? אמר לו יש לי י', אמר לו מה שמם? אמר לו בלע ובכר וגו'; אמר מה טיבן של שמות הללו? אמר לו כלם על שם אחי והצרות אשר מצאוהו, בלע שנבלע בין האמות, בכר שהיה בכור לאמי, אשבאל ששבאו אל, גרא שגר באכסניא, ונעמן שהיה נעים ביותר, אחי וראש – אחי היה וראשי היה, מפים – מפי אבי למד, וחפים – שלא ראה חפתי ולא ראיתי אני חפתו, וארד – שירד לבין האמות, כדאיתא במסכת סוטה, מיד נכמרו רחמיו: נכמרו. נתחממו, ובלשון משנה על הכמר של זתים, ובלשון ארמי במכמר בשרא; ובמקרא עורנו כתנור נכמרו (איכה ה'), נתחממו ונקמטו קמטים קמטים מפני זלעפות רעב; כן דרך כל עור כשמחממין אותו נקמט ונכוץ:}%
% ---------- פרק 43 | פסוק 31 ----------
 \textbf{לא} וַיִּרְחַ֥ץ פָּנָ֖יו וַיֵּצֵ֑א וַיִּ֨תְאַפַּ֔ק וַיֹּ֖אמֶר שִׂ֥ימוּ לָֽחֶם׃ \Rashi{\textbf{ויתאפק.}נתאמץ; והוא לשון אפיקי מגנים (איוב מ"א), חזק, וכן ומזיח אפיקים רפה (שם י"ב):}%
% ---------- פרק 43 | פסוק 32 ----------
 \textbf{לב} וַיָּשִׂ֥ימוּ ל֛וֹ לְבַדּ֖וֹ וְלָהֶ֣ם לְבַדָּ֑ם וְלַמִּצְרִ֞ים הָאֹכְלִ֤ים אִתּוֹ֙ לְבַדָּ֔ם כִּי֩ לֹ֨א יוּכְל֜וּן הַמִּצְרִ֗ים לֶאֱכֹ֤ל אֶת־הָֽעִבְרִים֙ לֶ֔חֶם כִּי־תוֹעֵבָ֥ה הִ֖וא לְמִצְרָֽיִם׃ \Rashi{\textbf{כי תועבה הוא.}דבר שנאוי הוא למצרים לאכל את העברים, ואנקלוס נתן טעם לדבר:}%
% ---------- פרק 43 | פסוק 33 ----------
 \textbf{לג} וַיֵּשְׁב֣וּ לְפָנָ֔יו הַבְּכֹר֙ כִּבְכֹ֣רָת֔וֹ וְהַצָּעִ֖יר כִּצְעִרָת֑וֹ וַיִּתְמְה֥וּ הָאֲנָשִׁ֖ים אִ֥ישׁ אֶל־רֵעֵֽהוּ׃ \Rashi{\textbf{הבכור כבכרתו.}מכה בגביע וקורא ראובן, שמעון, לוי, יהודה, יששכר, וזבולן, בני אם אחת, הסבו כסדר הזה, שהיא סדר תולדותיכם, וכן כלם. כיון שהגיע לבנימין, אמר זה אין לו אם ואני אין לי אם, ישב אצלי:}%
% ---------- פרק 43 | פסוק 34 ----------
 \textbf{לד} וַיִּשָּׂ֨א מַשְׂאֹ֜ת מֵאֵ֣ת פָּנָיו֮ אֲלֵהֶם֒ וַתֵּ֜רֶב מַשְׂאַ֧ת בִּנְיָמִ֛ן מִמַּשְׂאֹ֥ת כֻּלָּ֖ם חָמֵ֣שׁ יָד֑וֹת וַיִּשְׁתּ֥וּ וַֽיִּשְׁכְּר֖וּ עִמּֽוֹ׃ \Rashi{\textbf{משאת.}מנות: חמש ידות. חלקו עם אחיו, ומשאת יוסף ואסנת ומנשה ואפרים: וישכרו עמו. ומיום שמכרוהו לא שתו יין ולא הוא שתה יין, ואותו היום שתו (בראשית רבה):}%
% ---------- פרק 44 | פסוק 1 ----------
 \textbf{א} וַיְצַ֞ו אֶת־אֲשֶׁ֣ר עַל־בֵּיתוֹ֮ לֵאמֹר֒ מַלֵּ֞א אֶת־אַמְתְּחֹ֤ת הָֽאֲנָשִׁים֙ אֹ֔כֶל כַּאֲשֶׁ֥ר יוּכְל֖וּן שְׂאֵ֑ת וְשִׂ֥ים כֶּֽסֶף־אִ֖ישׁ בְּפִ֥י אַמְתַּחְתּֽוֹ׃ %
% ---------- פרק 44 | פסוק 2 ----------
 \textbf{ב} וְאֶת־גְּבִיעִ֞י גְּבִ֣יעַ הַכֶּ֗סֶף תָּשִׂים֙ בְּפִי֙ אַמְתַּ֣חַת הַקָּטֹ֔ן וְאֵ֖ת כֶּ֣סֶף שִׁבְר֑וֹ וַיַּ֕עַשׂ כִּדְבַ֥ר יוֹסֵ֖ף אֲשֶׁ֥ר דִּבֵּֽר׃ \Rashi{\textbf{גביעי.}כוס ארך, וקורין לו מדיר"נש בלעז:}%
% ---------- פרק 44 | פסוק 3 ----------
 \textbf{ג} הַבֹּ֖קֶר א֑וֹר וְהָאֲנָשִׁ֣ים שֻׁלְּח֔וּ הֵ֖מָּה וַחֲמֹרֵיהֶֽם׃ %
% ---------- פרק 44 | פסוק 4 ----------
 \textbf{ד} הֵ֠ם יָֽצְא֣וּ אֶת־הָעִיר֮ לֹ֣א הִרְחִ֒יקוּ֒ וְיוֹסֵ֤ף אָמַר֙ לַֽאֲשֶׁ֣ר עַל־בֵּית֔וֹ ק֥וּם רְדֹ֖ף אַחֲרֵ֣י הָֽאֲנָשִׁ֑ים וְהִשַּׂגְתָּם֙ וְאָמַרְתָּ֣ אֲלֵהֶ֔ם לָ֛מָּה שִׁלַּמְתֶּ֥ם רָעָ֖ה תַּ֥חַת טוֹבָֽה׃ %
% ---------- פרק 44 | פסוק 5 ----------
 \textbf{ה} הֲל֣וֹא זֶ֗ה אֲשֶׁ֨ר יִשְׁתֶּ֤ה אֲדֹנִי֙ בּ֔וֹ וְה֕וּא נַחֵ֥שׁ יְנַחֵ֖שׁ בּ֑וֹ הֲרֵעֹתֶ֖ם אֲשֶׁ֥ר עֲשִׂיתֶֽם׃ %
% ---------- פרק 44 | פסוק 6 ----------
 \textbf{ו} וַֽיַּשִּׂגֵ֑ם וַיְדַבֵּ֣ר אֲלֵהֶ֔ם אֶת־הַדְּבָרִ֖ים הָאֵֽלֶּה׃ %
% ---------- פרק 44 | פסוק 7 ----------
 \textbf{ז} וַיֹּאמְר֣וּ אֵלָ֔יו לָ֚מָּה יְדַבֵּ֣ר אֲדֹנִ֔י כַּדְּבָרִ֖ים הָאֵ֑לֶּה חָלִ֙ילָה֙ לַעֲבָדֶ֔יךָ מֵעֲשׂ֖וֹת כַּדָּבָ֥ר הַזֶּֽה׃ \Rashi{\textbf{חלילה לעבדיך.}חלין הוא לנו, לשון גנאי, ותרגום חס לעבדיך – חס מאת הקדוש ברוך הוא יהי עלינו מעשות זאת, והרבה יש בתלמוד "חס ושלום":}%
% ---------- פרק 44 | פסוק 8 ----------
 \textbf{ח} הֵ֣ן כֶּ֗סֶף אֲשֶׁ֤ר מָצָ֙אנוּ֙ בְּפִ֣י אַמְתְּחֹתֵ֔ינוּ הֱשִׁיבֹ֥נוּ אֵלֶ֖יךָ מֵאֶ֣רֶץ כְּנָ֑עַן וְאֵ֗יךְ נִגְנֹב֙ מִבֵּ֣ית אֲדֹנֶ֔יךָ כֶּ֖סֶף א֥וֹ זָהָֽב׃ \Rashi{\textbf{הן כסף אשר מצאנו.}זה אחד מעשרה קל וחמר האמורים בתורה, וכלן מנויין בב"ר:}%
% ---------- פרק 44 | פסוק 9 ----------
 \textbf{ט} אֲשֶׁ֨ר יִמָּצֵ֥א אִתּ֛וֹ מֵעֲבָדֶ֖יךָ וָמֵ֑ת וְגַם־אֲנַ֕חְנוּ נִֽהְיֶ֥ה לַֽאדֹנִ֖י לַעֲבָדִֽים׃ %
% ---------- פרק 44 | פסוק 10 ----------
 \textbf{י} וַיֹּ֕אמֶר גַּם־עַתָּ֥ה כְדִבְרֵיכֶ֖ם כֶּן־ה֑וּא אֲשֶׁ֨ר יִמָּצֵ֤א אִתּוֹ֙ יִהְיֶה־לִּ֣י עָ֔בֶד וְאַתֶּ֖ם תִּהְי֥וּ נְקִיִּֽם׃ \Rashi{\textbf{גם עתה כדבריכם.}אף זו מן הדין אמת כדבריכם כן הוא, שכלכם חיבין בדבר – י' שנמצאת גנבה ביד א' מהם כלם נתפשים, אבל אני אעשה לכם לפנים משורת הדין, אשר ימצא אתו יהיה לי עבד:}%
% ---------- פרק 44 | פסוק 11 ----------
 \textbf{יא} וַֽיְמַהֲר֗וּ וַיּוֹרִ֛דוּ אִ֥ישׁ אֶת־אַמְתַּחְתּ֖וֹ אָ֑רְצָה וַֽיִּפְתְּח֖וּ אִ֥ישׁ אַמְתַּחְתּֽוֹ׃ %
% ---------- פרק 44 | פסוק 12 ----------
 \textbf{יב} וַיְחַפֵּ֕שׂ בַּגָּד֣וֹל הֵחֵ֔ל וּבַקָּטֹ֖ן כִּלָּ֑ה וַיִּמָּצֵא֙ הַגָּבִ֔יעַ בְּאַמְתַּ֖חַת בִּנְיָמִֽן׃ \Rashi{\textbf{בגדול החל.}שלא ירגישו שהיה יודע היכן הוא:}%
% ---------- פרק 44 | פסוק 13 ----------
 \textbf{יג} וַֽיִּקְרְע֖וּ שִׂמְלֹתָ֑ם וַֽיַּעֲמֹס֙ אִ֣ישׁ עַל־חֲמֹר֔וֹ וַיָּשֻׁ֖בוּ הָעִֽירָה׃ \Rashi{\textbf{ויעמס איש על חמרו.}בעלי זרוע היו, ולא הצרכו לסיע זה את זה לטען: וישובו העירה. מטרופולין היתה, והוא אומר העירה, העיר כל שהוא אלא שלא היתה חשובה בעיניהם אלא כעיר בינונית של י' בני אדם לענין המלחמה:}%
% ---------- פרק 44 | פסוק 14 ----------
 \textbf{יד} וַיָּבֹ֨א יְהוּדָ֤ה וְאֶחָיו֙ בֵּ֣יתָה יוֹסֵ֔ף וְה֖וּא עוֹדֶ֣נּוּ שָׁ֑ם וַיִּפְּל֥וּ לְפָנָ֖יו אָֽרְצָה׃ \Rashi{\textbf{עודנו שם.}שהיה ממתין להם:}%
% ---------- פרק 44 | פסוק 15 ----------
 \textbf{טו} וַיֹּ֤אמֶר לָהֶם֙ יוֹסֵ֔ף מָֽה־הַמַּעֲשֶׂ֥ה הַזֶּ֖ה אֲשֶׁ֣ר עֲשִׂיתֶ֑ם הֲל֣וֹא יְדַעְתֶּ֔ם כִּֽי־נַחֵ֧שׁ יְנַחֵ֛שׁ אִ֖ישׁ אֲשֶׁ֥ר כָּמֹֽנִי׃ \Rashi{\textbf{הלא ידעתם כי נחש ינחש וגו'.}הלא ידעתם כי איש חשוב כמוני יודע לנחש ולדעת מדעת ומסברא ובינה, כי אתם גנבתם הגביע:}%
% ---------- פרק 44 | פסוק 16 ----------
 \textbf{טז} וַיֹּ֣אמֶר יְהוּדָ֗ה מַה־נֹּאמַר֙ לַֽאדֹנִ֔י מַה־נְּדַבֵּ֖ר וּמַה־נִּצְטַדָּ֑ק הָאֱלֹהִ֗ים מָצָא֙ אֶת־עֲוֺ֣ן עֲבָדֶ֔יךָ הִנֶּ֤נּוּ עֲבָדִים֙ לַֽאדֹנִ֔י גַּם־אֲנַ֕חְנוּ גַּ֛ם אֲשֶׁר־נִמְצָ֥א הַגָּבִ֖יעַ בְּיָדֽוֹ׃ \Rashi{\textbf{האלהים מצא.}יודעים אנו שלא סרחנו, אבל מאת המקום נהיתה להביא לנו זאת, מצא בעל חוב מקום לגבות שטר חובו: מה נצטדק. לשון צדק, וכן כל תבה שתחלת יסודה צד"י והיא באה לדבר בלשון מתפעל או נתפעל, נותן טי"ת במקום תי"ו, ואינו נותנה לפני אות ראשונה של יסוד התבה אלא באמצע אותיות העקר, כגון נצטדק מגזרת צדק, ויצטבע מגזרת צבע, ויצטירו מגזרת וציר אמונים, הצטידנו מגזרת צדה לדרך. ותבה שתחלתה סמ"ך או שי"ן, כשהיא מתפעלת, התי"ו מפרדת את אותיות העקר, כגון ויסתבל החגב (קהלת י"ב), מגזרת סבל; משתכל הוית בקרניא (דניאל ז'), וישתמר חקות עמרי (מיכה ו') מגזרת שמר; וסר מרע משתולל (ישעיהו נ"ט), מגזרת מוליך יועצים שולל (איוב י"ב), מסתולל בעמי (שמות ט"ו), מגזרת דרך לא סלולה (ירמיהו י"ח):}%
% ---------- פרק 44 | פסוק 17 ----------
 \textbf{יז} וַיֹּ֕אמֶר חָלִ֣ילָה לִּ֔י מֵעֲשׂ֖וֹת זֹ֑את הָאִ֡ישׁ אֲשֶׁר֩ נִמְצָ֨א הַגָּבִ֜יעַ בְּיָד֗וֹ ה֚וּא יִהְיֶה־לִּ֣י עָ֔בֶד וְאַתֶּ֕ם עֲל֥וּ לְשָׁל֖וֹם אֶל־אֲבִיכֶֽם׃ \{ס\} %
\addparasha{חומש בראשית | פרשת ויגש}
% ---------- פרק 44 | פסוק 18 ----------
 \textbf{יח} וַיִּגַּ֨שׁ אֵלָ֜יו יְהוּדָ֗ה וַיֹּ֘אמֶר֮ בִּ֣י אֲדֹנִי֒ יְדַבֶּר־נָ֨א עַבְדְּךָ֤ דָבָר֙ בְּאׇזְנֵ֣י אֲדֹנִ֔י וְאַל־יִ֥חַר אַפְּךָ֖ בְּעַבְדֶּ֑ךָ כִּ֥י כָמ֖וֹךָ כְּפַרְעֹֽה׃ \Rashi{\textbf{ויגש אליו וגומר … דבר באזני אדני.}יכנסו דברי באזניך: ואל יחר אפך. מכאן אתה למד שדבר אליו קשות: כי כמוך כפרעה. חשוב אתה בעיני כמלך, זהו פשוטו. ומדרשו סופך ללקות עליו בצרעת כמו שלקה פרעה על ידי זקנתי שרה על לילה אחת שעכבה (בראשית רבה). דבר אחר מה פרעה גוזר ואינו מקים, מבטיח ואינו עושה, אף אתה כן; וכי זו היא שימת עין שאמרת לשום עינך עליו? דבר אחר, כי כמוך כפרעה, אם תקניטני אהרג אותך ואת אדונך (בראשית רבה):}%
% ---------- פרק 44 | פסוק 19 ----------
 \textbf{יט} אֲדֹנִ֣י שָׁאַ֔ל אֶת־עֲבָדָ֖יו לֵאמֹ֑ר הֲיֵשׁ־לָכֶ֥ם אָ֖ב אוֹ־אָֽח׃ \Rashi{\textbf{שאל את עבדיו.}מתחלה בעלילה באת עלינו, למה היה לך לשאל כל אלה? בתך היינו מבקשים? או אחותנו אתה מבקש? ואף על פי כן}%
% ---------- פרק 44 | פסוק 20 ----------
 \textbf{כ} וַנֹּ֙אמֶר֙ אֶל־אֲדֹנִ֔י יֶשׁ־לָ֙נוּ֙ אָ֣ב זָקֵ֔ן וְיֶ֥לֶד זְקֻנִ֖ים קָטָ֑ן וְאָחִ֣יו מֵ֔ת וַיִּוָּתֵ֨ר ה֧וּא לְבַדּ֛וֹ לְאִמּ֖וֹ וְאָבִ֥יו אֲהֵבֽוֹ׃ \Rashi{\textbf{ונאמר אל אדוני, לא כחדנו ממך דבר (בראשית רבה): ואחיו מת.}מפני היראה היה מוציא דבר שקר מפיו, אמר אם אומר לו שהוא קים, יאמר הביאוהו אצלי: לבדו לאמו. מאותו האם אין לו עוד אח:}%
% ---------- פרק 44 | פסוק 21 ----------
 \textbf{כא} וַתֹּ֙אמֶר֙ אֶל־עֲבָדֶ֔יךָ הוֹרִדֻ֖הוּ אֵלָ֑י וְאָשִׂ֥ימָה עֵינִ֖י עָלָֽיו׃ %
% ---------- פרק 44 | פסוק 22 ----------
 \textbf{כב} וַנֹּ֙אמֶר֙ אֶל־אֲדֹנִ֔י לֹא־יוּכַ֥ל הַנַּ֖עַר לַעֲזֹ֣ב אֶת־אָבִ֑יו וְעָזַ֥ב אֶת־אָבִ֖יו וָמֵֽת׃ \Rashi{\textbf{ועזב את אביו ומת.}אם יעזב את אביו, דואגים אנו שמא ימות בדרך, שהרי אמו בדרך מתה:}%
% ---------- פרק 44 | פסוק 23 ----------
 \textbf{כג} וַתֹּ֙אמֶר֙ אֶל־עֲבָדֶ֔יךָ אִם־לֹ֥א יֵרֵ֛ד אֲחִיכֶ֥ם הַקָּטֹ֖ן אִתְּכֶ֑ם לֹ֥א תֹסִפ֖וּן לִרְא֥וֹת פָּנָֽי׃ %
% ---------- פרק 44 | פסוק 24 ----------
 \textbf{כד} וַֽיְהִי֙ כִּ֣י עָלִ֔ינוּ אֶֽל־עַבְדְּךָ֖ אָבִ֑י וַנַּ֨גֶּד־ל֔וֹ אֵ֖ת דִּבְרֵ֥י אֲדֹנִֽי׃ %
% ---------- פרק 44 | פסוק 25 ----------
 \textbf{כה} וַיֹּ֖אמֶר אָבִ֑ינוּ שֻׁ֖בוּ שִׁבְרוּ־לָ֥נוּ מְעַט־אֹֽכֶל׃ %
% ---------- פרק 44 | פסוק 26 ----------
 \textbf{כו} וַנֹּ֕אמֶר לֹ֥א נוּכַ֖ל לָרֶ֑דֶת אִם־יֵשׁ֩ אָחִ֨ינוּ הַקָּטֹ֤ן אִתָּ֙נוּ֙ וְיָרַ֔דְנוּ כִּי־לֹ֣א נוּכַ֗ל לִרְאוֹת֙ פְּנֵ֣י הָאִ֔ישׁ וְאָחִ֥ינוּ הַקָּטֹ֖ן אֵינֶ֥נּוּ אִתָּֽנוּ׃ %
% ---------- פרק 44 | פסוק 27 ----------
 \textbf{כז} וַיֹּ֛אמֶר עַבְדְּךָ֥ אָבִ֖י אֵלֵ֑ינוּ אַתֶּ֣ם יְדַעְתֶּ֔ם כִּ֥י שְׁנַ֖יִם יָֽלְדָה־לִּ֥י אִשְׁתִּֽי׃ %
% ---------- פרק 44 | פסוק 28 ----------
 \textbf{כח} וַיֵּצֵ֤א הָֽאֶחָד֙ מֵֽאִתִּ֔י וָאֹמַ֕ר אַ֖ךְ טָרֹ֣ף טֹרָ֑ף וְלֹ֥א רְאִיתִ֖יו עַד־הֵֽנָּה׃ %
% ---------- פרק 44 | פסוק 29 ----------
 \textbf{כט} וּלְקַחְתֶּ֧ם גַּם־אֶת־זֶ֛ה מֵעִ֥ם פָּנַ֖י וְקָרָ֣הוּ אָס֑וֹן וְהֽוֹרַדְתֶּ֧ם אֶת־שֵׂיבָתִ֛י בְּרָעָ֖ה שְׁאֹֽלָה׃ \Rashi{\textbf{וקרהו אסון.}שהשטן מקטרג בשעת הסכנה (בראשית רבה): והורדתם את שיבתי וגו'. עכשו כשהוא אצלי, אני מתנחם בו על אמו ועל אחיו, ואם ימות זה, דומה עלי ששלשתן מתו ביום אחד:}%
% ---------- פרק 44 | פסוק 30 ----------
 \textbf{ל} וְעַתָּ֗ה כְּבֹאִי֙ אֶל־עַבְדְּךָ֣ אָבִ֔י וְהַנַּ֖עַר אֵינֶ֣נּוּ אִתָּ֑נוּ וְנַפְשׁ֖וֹ קְשׁוּרָ֥ה בְנַפְשֽׁוֹ׃ %
% ---------- פרק 44 | פסוק 31 ----------
 \textbf{לא} וְהָיָ֗ה כִּרְאוֹת֛וֹ כִּי־אֵ֥ין הַנַּ֖עַר וָמֵ֑ת וְהוֹרִ֨ידוּ עֲבָדֶ֜יךָ אֶת־שֵׂיבַ֨ת עַבְדְּךָ֥ אָבִ֛ינוּ בְּיָג֖וֹן שְׁאֹֽלָה׃ \Rashi{\textbf{והיה כראותו כי אין הנער ומת.}אביו מצרתו:}%
% ---------- פרק 44 | פסוק 32 ----------
 \textbf{לב} כִּ֤י עַבְדְּךָ֙ עָרַ֣ב אֶת־הַנַּ֔עַר מֵעִ֥ם אָבִ֖י לֵאמֹ֑ר אִם־לֹ֤א אֲבִיאֶ֙נּוּ֙ אֵלֶ֔יךָ וְחָטָ֥אתִי לְאָבִ֖י כׇּל־הַיָּמִֽים׃ \Rashi{\textbf{כי עבדך ערב את הנער.}ואם תאמר, למה אני נכנס לתגר יותר משאר אחי? הם כלם מבחוץ ואני נתקשרתי בקשר חזק להיות מנדה בב' עולמות:}%
% ---------- פרק 44 | פסוק 33 ----------
 \textbf{לג} וְעַתָּ֗ה יֵֽשֶׁב־נָ֤א עַבְדְּךָ֙ תַּ֣חַת הַנַּ֔עַר עֶ֖בֶד לַֽאדֹנִ֑י וְהַנַּ֖עַר יַ֥עַל עִם־אֶחָֽיו׃ \Rashi{\textbf{ישב נא עבדך וגו'.}לכל דבר אני מעלה ממנו, לגבורה ולמלחמה ולשמש (בראשית רבה):}%
% ---------- פרק 44 | פסוק 34 ----------
 \textbf{לד} כִּי־אֵיךְ֙ אֶֽעֱלֶ֣ה אֶל־אָבִ֔י וְהַנַּ֖עַר אֵינֶ֣נּוּ אִתִּ֑י פֶּ֚ן אֶרְאֶ֣ה בָרָ֔ע אֲשֶׁ֥ר יִמְצָ֖א אֶת־אָבִֽי׃ %
% ---------- פרק 45 | פסוק 1 ----------
 \textbf{א} וְלֹֽא־יָכֹ֨ל יוֹסֵ֜ף לְהִתְאַפֵּ֗ק לְכֹ֤ל הַנִּצָּבִים֙ עָלָ֔יו וַיִּקְרָ֕א הוֹצִ֥יאוּ כׇל־אִ֖ישׁ מֵעָלָ֑י וְלֹא־עָ֤מַד אִישׁ֙ אִתּ֔וֹ בְּהִתְוַדַּ֥ע יוֹסֵ֖ף אֶל־אֶחָֽיו׃ \Rashi{\textbf{ולא יכל יוסף להתאפק לכל הנצבים.}לא היה יכול לסבל שיהיו מצרים נצבים עליו ושומעין שאחיו מתבישין בהודעו להם:}%
% ---------- פרק 45 | פסוק 2 ----------
 \textbf{ב} וַיִּתֵּ֥ן אֶת־קֹל֖וֹ בִּבְכִ֑י וַיִּשְׁמְע֣וּ מִצְרַ֔יִם וַיִּשְׁמַ֖ע בֵּ֥ית פַּרְעֹֽה׃ \Rashi{\textbf{וישמע בית פרעה.}ביתו של פרעה, כלומר, עבדיו ובני ביתו, ואין זה לשון בית ממש, אלא כמו בית ישראל, בית יהודה, מישנ"דה בלעז:}%
% ---------- פרק 45 | פסוק 3 ----------
 \textbf{ג} וַיֹּ֨אמֶר יוֹסֵ֤ף אֶל־אֶחָיו֙ אֲנִ֣י יוֹסֵ֔ף הַע֥וֹד אָבִ֖י חָ֑י וְלֹֽא־יָכְל֤וּ אֶחָיו֙ לַעֲנ֣וֹת אֹת֔וֹ כִּ֥י נִבְהֲל֖וּ מִפָּנָֽיו׃ \Rashi{\textbf{נבהלו מפניו.}מפני הבושה:}%
% ---------- פרק 45 | פסוק 4 ----------
 \textbf{ד} וַיֹּ֨אמֶר יוֹסֵ֧ף אֶל־אֶחָ֛יו גְּשׁוּ־נָ֥א אֵלַ֖י וַיִּגָּ֑שׁוּ וַיֹּ֗אמֶר אֲנִי֙ יוֹסֵ֣ף אֲחִיכֶ֔ם אֲשֶׁר־מְכַרְתֶּ֥ם אֹתִ֖י מִצְרָֽיְמָה׃ \Rashi{\textbf{גשו נא אלי.}ראה אותם נסוגים לאחור, אמר עכשו אחי נכלמים, קרא להם בלשון רכה ותחנונים, והראה להם שהוא מהול (בראשית רבה):}%
% ---------- פרק 45 | פסוק 5 ----------
 \textbf{ה} וְעַתָּ֣ה ׀ אַל־תֵּעָ֣צְב֗וּ וְאַל־יִ֙חַר֙ בְּעֵ֣ינֵיכֶ֔ם כִּֽי־מְכַרְתֶּ֥ם אֹתִ֖י הֵ֑נָּה כִּ֣י לְמִֽחְיָ֔ה שְׁלָחַ֥נִי אֱלֹהִ֖ים לִפְנֵיכֶֽם׃ \Rashi{\textbf{למחיה.}להיות לכם למחיה:}%
% ---------- פרק 45 | פסוק 6 ----------
 \textbf{ו} כִּי־זֶ֛ה שְׁנָתַ֥יִם הָרָעָ֖ב בְּקֶ֣רֶב הָאָ֑רֶץ וְעוֹד֙ חָמֵ֣שׁ שָׁנִ֔ים אֲשֶׁ֥ר אֵין־חָרִ֖ישׁ וְקָצִֽיר׃ \Rashi{\textbf{כי זה שנתים הרעב.}עברו משני הרעב:}%
% ---------- פרק 45 | פסוק 7 ----------
 \textbf{ז} וַיִּשְׁלָחֵ֤נִי אֱלֹהִים֙ לִפְנֵיכֶ֔ם לָשׂ֥וּם לָכֶ֛ם שְׁאֵרִ֖ית בָּאָ֑רֶץ וּלְהַחֲי֣וֹת לָכֶ֔ם לִפְלֵיטָ֖ה גְּדֹלָֽה׃ %
% ---------- פרק 45 | פסוק 8 ----------
 \textbf{ח} וְעַתָּ֗ה לֹֽא־אַתֶּ֞ם שְׁלַחְתֶּ֤ם אֹתִי֙ הֵ֔נָּה כִּ֖י הָאֱלֹהִ֑ים וַיְשִׂימֵ֨נִֽי לְאָ֜ב לְפַרְעֹ֗ה וּלְאָדוֹן֙ לְכׇל־בֵּית֔וֹ וּמֹשֵׁ֖ל בְּכׇל־אֶ֥רֶץ מִצְרָֽיִם׃ \Rashi{\textbf{לאב.}לחבר ולפטרון:}%
% ---------- פרק 45 | פסוק 9 ----------
 \textbf{ט} מַהֲרוּ֮ וַעֲל֣וּ אֶל־אָבִי֒ וַאֲמַרְתֶּ֣ם אֵלָ֗יו כֹּ֤ה אָמַר֙ בִּנְךָ֣ יוֹסֵ֔ף שָׂמַ֧נִי אֱלֹהִ֛ים לְאָד֖וֹן לְכׇל־מִצְרָ֑יִם רְדָ֥ה אֵלַ֖י אַֽל־תַּעֲמֹֽד׃ \Rashi{\textbf{ועלו אל אבי.}ארץ ישראל גבוהה מכל הארצות:}%
% ---------- פרק 45 | פסוק 10 ----------
 \textbf{י} וְיָשַׁבְתָּ֣ בְאֶֽרֶץ־גֹּ֗שֶׁן וְהָיִ֤יתָ קָרוֹב֙ אֵלַ֔י אַתָּ֕ה וּבָנֶ֖יךָ וּבְנֵ֣י בָנֶ֑יךָ וְצֹאנְךָ֥ וּבְקָרְךָ֖ וְכׇל־אֲשֶׁר־לָֽךְ׃ %
% ---------- פרק 45 | פסוק 11 ----------
 \textbf{יא} וְכִלְכַּלְתִּ֤י אֹֽתְךָ֙ שָׁ֔ם כִּי־ע֛וֹד חָמֵ֥שׁ שָׁנִ֖ים רָעָ֑ב פֶּן־תִּוָּרֵ֛שׁ אַתָּ֥ה וּבֵֽיתְךָ֖ וְכׇל־אֲשֶׁר־לָֽךְ׃ \Rashi{\textbf{פן תורש.}דלמא תתמסכן, לשון מוריש ומעשיר:}%
% ---------- פרק 45 | פסוק 12 ----------
 \textbf{יב} וְהִנֵּ֤ה עֵֽינֵיכֶם֙ רֹא֔וֹת וְעֵינֵ֖י אָחִ֣י בִנְיָמִ֑ין כִּי־פִ֖י הַֽמְדַבֵּ֥ר אֲלֵיכֶֽם׃ \Rashi{\textbf{והנה עיניכם ראות.}בכבודי, ושאני אחיכם שאני מהול ככם, ועוד, כי פי המדבר אליכם בלשון הקדש (בראשית רבה): ועיני אחי בנימין. השוה את כלם יחד לומר, שכשם שאין לי שנאה על בנימין אחי, שהרי לא היה במכירתי, כך אין בלבי שנאה עליכם:}%
% ---------- פרק 45 | פסוק 13 ----------
 \textbf{יג} וְהִגַּדְתֶּ֣ם לְאָבִ֗י אֶת־כׇּל־כְּבוֹדִי֙ בְּמִצְרַ֔יִם וְאֵ֖ת כׇּל־אֲשֶׁ֣ר רְאִיתֶ֑ם וּמִֽהַרְתֶּ֛ם וְהוֹרַדְתֶּ֥ם אֶת־אָבִ֖י הֵֽנָּה׃ %
% ---------- פרק 45 | פסוק 14 ----------
 \textbf{יד} וַיִּפֹּ֛ל עַל־צַוְּארֵ֥י בִנְיָֽמִן־אָחִ֖יו וַיֵּ֑בְךְּ וּבִ֨נְיָמִ֔ן בָּכָ֖ה עַל־צַוָּארָֽיו׃ \Rashi{\textbf{ויפל על צוארי בנימין אחיו ויבך.}על שני מקדשות שעתידין להיות בחלקו של בנימין וסופן לחרב: ובנימין בכה על צואריו. על משכן שילה שעתיד להיות בחלקו של יוסף וסופו לחרב (שם):}%
% ---------- פרק 45 | פסוק 15 ----------
 \textbf{טו} וַיְנַשֵּׁ֥ק לְכׇל־אֶחָ֖יו וַיֵּ֣בְךְּ עֲלֵהֶ֑ם וְאַ֣חֲרֵי כֵ֔ן דִּבְּר֥וּ אֶחָ֖יו אִתּֽוֹ׃ \Rashi{\textbf{ואחרי כן.}מאחר שראוהו בוכה ולבו שלם עמהם דברו אחיו אתו שמתחלה היו בושים ממנו:}%
% ---------- פרק 45 | פסוק 16 ----------
 \textbf{טז} וְהַקֹּ֣ל נִשְׁמַ֗ע בֵּ֤ית פַּרְעֹה֙ לֵאמֹ֔ר בָּ֖אוּ אֲחֵ֣י יוֹסֵ֑ף וַיִּיטַב֙ בְּעֵינֵ֣י פַרְעֹ֔ה וּבְעֵינֵ֖י עֲבָדָֽיו׃ \Rashi{\textbf{והקל נשמע בית פרעה.}כמו בבית פרעה, וזהו לשון בית ממש:}%
% ---------- פרק 45 | פסוק 17 ----------
 \textbf{יז} וַיֹּ֤אמֶר פַּרְעֹה֙ אֶל־יוֹסֵ֔ף אֱמֹ֥ר אֶל־אַחֶ֖יךָ זֹ֣את עֲשׂ֑וּ טַֽעֲנוּ֙ אֶת־בְּעִ֣ירְכֶ֔ם וּלְכוּ־בֹ֖אוּ אַ֥רְצָה כְּנָֽעַן׃ \Rashi{\textbf{טענו את בעירכם.}תבואה:}%
% ---------- פרק 45 | פסוק 18 ----------
 \textbf{יח} וּקְח֧וּ אֶת־אֲבִיכֶ֛ם וְאֶת־בָּתֵּיכֶ֖ם וּבֹ֣אוּ אֵלָ֑י וְאֶתְּנָ֣ה לָכֶ֗ם אֶת־טוּב֙ אֶ֣רֶץ מִצְרַ֔יִם וְאִכְל֖וּ אֶת־חֵ֥לֶב הָאָֽרֶץ׃ \Rashi{\textbf{את טוב ארץ מצרים.}ארץ גשן, נבא ואינו יודע מה נבא, סופם לעשותה כמצולה שאין בה דגים: חלב הארץ. כל חלב לשון מיטב הוא:}%
% ---------- פרק 45 | פסוק 19 ----------
 \textbf{יט} וְאַתָּ֥ה צֻוֵּ֖יתָה זֹ֣את עֲשׂ֑וּ קְחוּ־לָכֶם֩ מֵאֶ֨רֶץ מִצְרַ֜יִם עֲגָל֗וֹת לְטַפְּכֶם֙ וְלִנְשֵׁיכֶ֔ם וּנְשָׂאתֶ֥ם אֶת־אֲבִיכֶ֖ם וּבָאתֶֽם׃ \Rashi{\textbf{ואתה צויתה.}מפי, לומר להם: זאת עשו. כך אמר להם, שברשותי הוא:}%
% ---------- פרק 45 | פסוק 20 ----------
 \textbf{כ} וְעֵ֣ינְכֶ֔ם אַל־תָּחֹ֖ס עַל־כְּלֵיכֶ֑ם כִּי־ט֛וּב כׇּל־אֶ֥רֶץ מִצְרַ֖יִם לָכֶ֥ם הֽוּא׃ %
% ---------- פרק 45 | פסוק 21 ----------
 \textbf{כא} וַיַּֽעֲשׂוּ־כֵן֙ בְּנֵ֣י יִשְׂרָאֵ֔ל וַיִּתֵּ֨ן לָהֶ֥ם יוֹסֵ֛ף עֲגָל֖וֹת עַל־פִּ֣י פַרְעֹ֑ה וַיִּתֵּ֥ן לָהֶ֛ם צֵדָ֖ה לַדָּֽרֶךְ׃ %
% ---------- פרק 45 | פסוק 22 ----------
 \textbf{כב} לְכֻלָּ֥ם נָתַ֛ן לָאִ֖ישׁ חֲלִפ֣וֹת שְׂמָלֹ֑ת וּלְבִנְיָמִ֤ן נָתַן֙ שְׁלֹ֣שׁ מֵא֣וֹת כֶּ֔סֶף וְחָמֵ֖שׁ חֲלִפֹ֥ת שְׂמָלֹֽת׃ %
% ---------- פרק 45 | פסוק 23 ----------
 \textbf{כג} וּלְאָבִ֞יו שָׁלַ֤ח כְּזֹאת֙ עֲשָׂרָ֣ה חֲמֹרִ֔ים נֹשְׂאִ֖ים מִטּ֣וּב מִצְרָ֑יִם וְעֶ֣שֶׂר אֲתֹנֹ֡ת נֹֽ֠שְׂאֹ֠ת בָּ֣ר וָלֶ֧חֶם וּמָז֛וֹן לְאָבִ֖יו לַדָּֽרֶךְ׃ \Rashi{\textbf{שלח כזאת.}כחשבון הזה, ומהו החשבון? עשרה חמרים וגו': מטוב מצרים. מצינו בתלמוד ששלח לו יין ישן שדעת זקנים נוחה הימנו, ומדרש אגדה גריסין של פול: בר ולחם. כתרגומו: ומזון. לפתן:}%
% ---------- פרק 45 | פסוק 24 ----------
 \textbf{כד} וַיְשַׁלַּ֥ח אֶת־אֶחָ֖יו וַיֵּלֵ֑כוּ וַיֹּ֣אמֶר אֲלֵהֶ֔ם אַֽל־תִּרְגְּז֖וּ בַּדָּֽרֶךְ׃ \Rashi{\textbf{אל תרגזו בדרך.}אל תתעסקו בדבר הלכה שלא תרגז עליכם הדרך, ד"א אל תפסיעו פסיעה גסה והכנסו בחמה לעיר; ולפי פשוטו של מקרא יש לומר, לפי שהיו נכלמים, היה דואג, שמא יריבו בדרך על דבר מכירתו, להתוכח זה עם זה ולומר על ידך נמכר, אתה ספרת לשון הרע עליו, וגרמת לנו לשנאתו:}%
% ---------- פרק 45 | פסוק 25 ----------
 \textbf{כה} וַֽיַּעֲל֖וּ מִמִּצְרָ֑יִם וַיָּבֹ֙אוּ֙ אֶ֣רֶץ כְּנַ֔עַן אֶֽל־יַעֲקֹ֖ב אֲבִיהֶֽם׃ %
% ---------- פרק 45 | פסוק 26 ----------
 \textbf{כו} וַיַּגִּ֨דוּ ל֜וֹ לֵאמֹ֗ר ע֚וֹד יוֹסֵ֣ף חַ֔י וְכִֽי־ה֥וּא מֹשֵׁ֖ל בְּכׇל־אֶ֣רֶץ מִצְרָ֑יִם וַיָּ֣פׇג לִבּ֔וֹ כִּ֥י לֹא־הֶאֱמִ֖ין לָהֶֽם׃ \Rashi{\textbf{וכי הוא משל.}ואשר הוא מושל: ויפג לבו. נחלף לבו והלך מלהאמין, לא היה לבו פונה אל הדברים, לשון מפיגין טעמן בלשון משנה וכמו מאין הפגות (איכה ג'), וריחו לא נמר (ירמיהו מ"ח), מתרגמינן וריחיה לא פג:}%
% ---------- פרק 45 | פסוק 27 ----------
 \textbf{כז} וַיְדַבְּר֣וּ אֵלָ֗יו אֵ֣ת כׇּל־דִּבְרֵ֤י יוֹסֵף֙ אֲשֶׁ֣ר דִּבֶּ֣ר אֲלֵהֶ֔ם וַיַּרְא֙ אֶת־הָ֣עֲגָל֔וֹת אֲשֶׁר־שָׁלַ֥ח יוֹסֵ֖ף לָשֵׂ֣את אֹת֑וֹ וַתְּחִ֕י ר֖וּחַ יַעֲקֹ֥ב אֲבִיהֶֽם׃ \Rashi{\textbf{את כל דברי יוסף.}סימן מסר להם במה היה עוסק כשפרש ממנו – בפרשת עגלה ערופה, זהו שנאמר וירא את העגלות אשר שלח יוסף, ולא נאמר אשר שלח פרעה: ותחי רוח יעקב. שרתה עליו שכינה, שפרשה ממנו:}%
% ---------- פרק 45 | פסוק 28 ----------
 \textbf{כח} וַיֹּ֙אמֶר֙ יִשְׂרָאֵ֔ל רַ֛ב עוֹד־יוֹסֵ֥ף בְּנִ֖י חָ֑י אֵֽלְכָ֥ה וְאֶרְאֶ֖נּוּ בְּטֶ֥רֶם אָמֽוּת׃ \Rashi{\textbf{רב.}רב לי עוד שמחה וחדוה הואיל ועוד יוסף בני חי:}%
% ---------- פרק 46 | פסוק 1 ----------
 \textbf{א} וַיִּסַּ֤ע יִשְׂרָאֵל֙ וְכׇל־אֲשֶׁר־ל֔וֹ וַיָּבֹ֖א בְּאֵ֣רָה שָּׁ֑בַע וַיִּזְבַּ֣ח זְבָחִ֔ים לֵאלֹהֵ֖י אָבִ֥יו יִצְחָֽק׃ \Rashi{\textbf{בארה שבע.}כמו לבאר שבע; ה"א בסוף תבה במקום למ"ד בתחלתה: לאלהי אביו יצחק. חיב אדם בכבוד אביו יותר מבכבוד זקנו, לפיכך תלה ביצחק ולא באברהם:}%
% ---------- פרק 46 | פסוק 2 ----------
 \textbf{ב} וַיֹּ֨אמֶר אֱלֹהִ֤ים ׀ לְיִשְׂרָאֵל֙ בְּמַרְאֹ֣ת הַלַּ֔יְלָה וַיֹּ֖אמֶר יַעֲקֹ֣ב ׀ יַעֲקֹ֑ב וַיֹּ֖אמֶר הִנֵּֽנִי׃ \Rashi{\textbf{יעקב יעקב.}לשון חבה:}%
% ---------- פרק 46 | פסוק 3 ----------
 \textbf{ג} וַיֹּ֕אמֶר אָנֹכִ֥י הָאֵ֖ל אֱלֹהֵ֣י אָבִ֑יךָ אַל־תִּירָא֙ מֵרְדָ֣ה מִצְרַ֔יְמָה כִּֽי־לְג֥וֹי גָּד֖וֹל אֲשִֽׂימְךָ֥ שָֽׁם׃ \Rashi{\textbf{אל תירא מרדה מצרימה.}לפי שהיה מצר על שנזקק לצאת לחוצה לארץ:}%
% ---------- פרק 46 | פסוק 4 ----------
 \textbf{ד} אָנֹכִ֗י אֵרֵ֤ד עִמְּךָ֙ מִצְרַ֔יְמָה וְאָנֹכִ֖י אַֽעַלְךָ֣ גַם־עָלֹ֑ה וְיוֹסֵ֕ף יָשִׁ֥ית יָד֖וֹ עַל־עֵינֶֽיךָ׃ \Rashi{\textbf{ואנכי אעלך.}הבטיחו להיות נקבר בארץ:}%
% ---------- פרק 46 | פסוק 5 ----------
 \textbf{ה} וַיָּ֥קׇם יַעֲקֹ֖ב מִבְּאֵ֣ר שָׁ֑בַע וַיִּשְׂא֨וּ בְנֵֽי־יִשְׂרָאֵ֜ל אֶת־יַעֲקֹ֣ב אֲבִיהֶ֗ם וְאֶת־טַפָּם֙ וְאֶת־נְשֵׁיהֶ֔ם בָּעֲגָל֕וֹת אֲשֶׁר־שָׁלַ֥ח פַּרְעֹ֖ה לָשֵׂ֥את אֹתֽוֹ׃ %
% ---------- פרק 46 | פסוק 6 ----------
 \textbf{ו} וַיִּקְח֣וּ אֶת־מִקְנֵיהֶ֗ם וְאֶת־רְכוּשָׁם֙ אֲשֶׁ֤ר רָֽכְשׁוּ֙ בְּאֶ֣רֶץ כְּנַ֔עַן וַיָּבֹ֖אוּ מִצְרָ֑יְמָה יַעֲקֹ֖ב וְכׇל־זַרְע֥וֹ אִתּֽוֹ׃ \Rashi{\textbf{אשר רכשו בארץ כנען.}אבל מה שרכש בפדן ארם נתן הכל לעשו בשביל חלקו במערת המכפלה; אמר נכסי חוצה לארץ אינן כדאי לי, וזהו אשר כריתי לי, העמיד לו צבורין של זהב וכסף כמין כרי ואמר לו טל את אלו:}%
% ---------- פרק 46 | פסוק 7 ----------
 \textbf{ז} בָּנָ֞יו וּבְנֵ֤י בָנָיו֙ אִתּ֔וֹ בְּנֹתָ֛יו וּבְנ֥וֹת בָּנָ֖יו וְכׇל־זַרְע֑וֹ הֵבִ֥יא אִתּ֖וֹ מִצְרָֽיְמָה׃ \{ס\} \Rashi{\textbf{ובנות בניו.}סרח בת אשר ויוכבד בת לוי:}%
% ---------- פרק 46 | פסוק 8 ----------
 \textbf{ח} וְאֵ֨לֶּה שְׁמ֧וֹת בְּנֵֽי־יִשְׂרָאֵ֛ל הַבָּאִ֥ים מִצְרַ֖יְמָה יַעֲקֹ֣ב וּבָנָ֑יו בְּכֹ֥ר יַעֲקֹ֖ב רְאוּבֵֽן׃ \Rashi{\textbf{הבאים מצרימה.}על שם השעה קורא להם הכתוב באים, ואין לתמה על אשר לא כתב אשר באו:}%
% ---------- פרק 46 | פסוק 9 ----------
 \textbf{ט} וּבְנֵ֖י רְאוּבֵ֑ן חֲנ֥וֹךְ וּפַלּ֖וּא וְחֶצְרֹ֥ן וְכַרְמִֽי׃ %
% ---------- פרק 46 | פסוק 10 ----------
 \textbf{י} וּבְנֵ֣י שִׁמְע֗וֹן יְמוּאֵ֧ל וְיָמִ֛ין וְאֹ֖הַד וְיָכִ֣ין וְצֹ֑חַר וְשָׁא֖וּל בֶּן־הַֽכְּנַעֲנִֽית׃ \Rashi{\textbf{בן הכנענית.}בן דינה שנבעלה לכנעני; כשהרגו את שכם לא היתה דינה רוצה לצאת עד שנשבע לה שמעון שישאנה (בראשית רבה):}%
% ---------- פרק 46 | פסוק 11 ----------
 \textbf{יא} וּבְנֵ֖י לֵוִ֑י גֵּרְשׁ֕וֹן קְהָ֖ת וּמְרָרִֽי׃ %
% ---------- פרק 46 | פסוק 12 ----------
 \textbf{יב} וּבְנֵ֣י יְהוּדָ֗ה עֵ֧ר וְאוֹנָ֛ן וְשֵׁלָ֖ה וָפֶ֣רֶץ וָזָ֑רַח וַיָּ֨מׇת עֵ֤ר וְאוֹנָן֙ בְּאֶ֣רֶץ כְּנַ֔עַן וַיִּהְי֥וּ בְנֵי־פֶ֖רֶץ חֶצְרֹ֥ן וְחָמֽוּל׃ %
% ---------- פרק 46 | פסוק 13 ----------
 \textbf{יג} וּבְנֵ֖י יִשָּׂשכָ֑ר תּוֹלָ֥ע וּפֻוָ֖ה וְי֥וֹב וְשִׁמְרֹֽן׃ %
% ---------- פרק 46 | פסוק 14 ----------
 \textbf{יד} וּבְנֵ֖י זְבֻל֑וּן סֶ֥רֶד וְאֵל֖וֹן וְיַחְלְאֵֽל׃ %
% ---------- פרק 46 | פסוק 15 ----------
 \textbf{טו} אֵ֣לֶּה ׀ בְּנֵ֣י לֵאָ֗ה אֲשֶׁ֨ר יָֽלְדָ֤ה לְיַעֲקֹב֙ בְּפַדַּ֣ן אֲרָ֔ם וְאֵ֖ת דִּינָ֣ה בִתּ֑וֹ כׇּל־נֶ֧פֶשׁ בָּנָ֛יו וּבְנוֹתָ֖יו שְׁלֹשִׁ֥ים וְשָׁלֹֽשׁ׃ \Rashi{\textbf{אלה בני לאה, ואת דינה בתו.}הזכרים תלה בלאה והנקבות תלה ביעקב, ללמדך, אשה מזרעת תחלה יולדת זכר, איש מזריע תחלה יולדת נקבה: שלשים ושלש. ובפרטן אי אתה מוצא אלא שלשים ושנים, אלא זו יוכבד שנולדה בין החומות בכניסתן לעיר, שנאמר אשר ילדה אתה ללוי במצרים – לדתה במצרים ואין הורתה במצרים:}%
% ---------- פרק 46 | פסוק 16 ----------
 \textbf{טז} וּבְנֵ֣י גָ֔ד צִפְי֥וֹן וְחַגִּ֖י שׁוּנִ֣י וְאֶצְבֹּ֑ן עֵרִ֥י וַֽאֲרוֹדִ֖י וְאַרְאֵלִֽי׃ %
% ---------- פרק 46 | פסוק 17 ----------
 \textbf{יז} וּבְנֵ֣י אָשֵׁ֗ר יִמְנָ֧ה וְיִשְׁוָ֛ה וְיִשְׁוִ֥י וּבְרִיעָ֖ה וְשֶׂ֣רַח אֲחֹתָ֑ם וּבְנֵ֣י בְרִיעָ֔ה חֶ֖בֶר וּמַלְכִּיאֵֽל׃ %
% ---------- פרק 46 | פסוק 18 ----------
 \textbf{יח} אֵ֚לֶּה בְּנֵ֣י זִלְפָּ֔ה אֲשֶׁר־נָתַ֥ן לָבָ֖ן לְלֵאָ֣ה בִתּ֑וֹ וַתֵּ֤לֶד אֶת־אֵ֙לֶּה֙ לְיַעֲקֹ֔ב שֵׁ֥שׁ עֶשְׂרֵ֖ה נָֽפֶשׁ׃ %
% ---------- פרק 46 | פסוק 19 ----------
 \textbf{יט} בְּנֵ֤י רָחֵל֙ אֵ֣שֶׁת יַעֲקֹ֔ב יוֹסֵ֖ף וּבִנְיָמִֽן׃ \Rashi{\textbf{בני רחל אשת יעקב.}ובכלן לא נאמר בהן אשת, אלא שהיתה עקרו של בית:}%
% ---------- פרק 46 | פסוק 20 ----------
 \textbf{כ} וַיִּוָּלֵ֣ד לְיוֹסֵף֮ בְּאֶ֣רֶץ מִצְרַ֒יִם֒ אֲשֶׁ֤ר יָֽלְדָה־לּוֹ֙ אָֽסְנַ֔ת בַּת־פּ֥וֹטִי פֶ֖רַע*(בספרי תימן פּֽוֹטִיפֶ֖רַע בתיבה אחת) כֹּהֵ֣ן אֹ֑ן אֶת־מְנַשֶּׁ֖ה וְאֶת־אֶפְרָֽיִם׃ %
% ---------- פרק 46 | פסוק 21 ----------
 \textbf{כא} וּבְנֵ֣י בִנְיָמִ֗ן בֶּ֤לַע וָבֶ֙כֶר֙ וְאַשְׁבֵּ֔ל גֵּרָ֥א וְנַעֲמָ֖ן אֵחִ֣י וָרֹ֑אשׁ מֻפִּ֥ים וְחֻפִּ֖ים וָאָֽרְדְּ׃ %
% ---------- פרק 46 | פסוק 22 ----------
 \textbf{כב} אֵ֚לֶּה בְּנֵ֣י רָחֵ֔ל אֲשֶׁ֥ר יֻלַּ֖ד לְיַעֲקֹ֑ב כׇּל־נֶ֖פֶשׁ אַרְבָּעָ֥ה עָשָֽׂר׃ %
% ---------- פרק 46 | פסוק 23 ----------
 \textbf{כג} וּבְנֵי־דָ֖ן חֻשִֽׁים׃ %
% ---------- פרק 46 | פסוק 24 ----------
 \textbf{כד} וּבְנֵ֖י נַפְתָּלִ֑י יַחְצְאֵ֥ל וְגוּנִ֖י וְיֵ֥צֶר וְשִׁלֵּֽם׃ %
% ---------- פרק 46 | פסוק 25 ----------
 \textbf{כה} אֵ֚לֶּה בְּנֵ֣י בִלְהָ֔ה אֲשֶׁר־נָתַ֥ן לָבָ֖ן לְרָחֵ֣ל בִּתּ֑וֹ וַתֵּ֧לֶד אֶת־אֵ֛לֶּה לְיַעֲקֹ֖ב כׇּל־נֶ֥פֶשׁ שִׁבְעָֽה׃ %
% ---------- פרק 46 | פסוק 26 ----------
 \textbf{כו} כׇּל־הַ֠נֶּ֠פֶשׁ הַבָּאָ֨ה לְיַעֲקֹ֤ב מִצְרַ֙יְמָה֙ יֹצְאֵ֣י יְרֵכ֔וֹ מִלְּבַ֖ד נְשֵׁ֣י בְנֵי־יַעֲקֹ֑ב כׇּל־נֶ֖פֶשׁ שִׁשִּׁ֥ים וָשֵֽׁשׁ׃ \Rashi{\textbf{כל הנפש הבאה ליעקב.}שיצאו מארץ כנען לבא למצרים; ואין "הבאה" זו לשון עבר אלא לשון הווה, כמו בערב היא באה (אסתר ב' ד'), וכמו והנה רחל בתו באה עם הצאן (בראשית כ"ט), לפיכך טעמו למטה באל"ף, לפי שכשיצאו לבא מארץ כנען לא היו אלא ששים ושש, והשני, כל הנפש לבית יעקב הבאה מצרימה שבעים – הוא לשון עבר, לפיכך טעמו למעלה בבי"ת, לפי שמשבאו שם היו שבעים, שמצאו שם יוסף ושני בניו, ונתוספה להם יוכבד בין החומות. ולדברי האומר תאומות נולדו עם השבטים, צריכים אנו לומר שמתו לפני ירידתן למצרים, שהרי לא נמנו כאן; מצאתי בויקרא רבה, עשו שש נפשות היו לו, והכתוב קורא אותן נפשות ביתו, לשון רבים, לפי שהיו עובדין לאלהות הרבה; יעקב שבעים היו לו, והכתוב קורא אותן נפש, לפי שהיה עובדים לאל אחד:}%
% ---------- פרק 46 | פסוק 27 ----------
 \textbf{כז} וּבְנֵ֥י יוֹסֵ֛ף אֲשֶׁר־יֻלַּד־ל֥וֹ בְמִצְרַ֖יִם נֶ֣פֶשׁ שְׁנָ֑יִם כׇּל־הַנֶּ֧פֶשׁ לְבֵֽית־יַעֲקֹ֛ב הַבָּ֥אָה מִצְרַ֖יְמָה שִׁבְעִֽים׃ \{ס\} %
% ---------- פרק 46 | פסוק 28 ----------
 \textbf{כח} וְאֶת־יְהוּדָ֞ה שָׁלַ֤ח לְפָנָיו֙ אֶל־יוֹסֵ֔ף לְהוֹרֹ֥ת לְפָנָ֖יו גֹּ֑שְׁנָה וַיָּבֹ֖אוּ אַ֥רְצָה גֹּֽשֶׁן׃ \Rashi{\textbf{להורת לפניו.}כתרגומו, לפנות לו מקום, ולהורות היאך יתישב בה: לפניו. קדם שיגיע לשם. ומדרש אגדה להורות לפניו – לתקן לו בית תלמוד שמשם תצא הוראה:}%
% ---------- פרק 46 | פסוק 29 ----------
 \textbf{כט} וַיֶּאְסֹ֤ר יוֹסֵף֙ מֶרְכַּבְתּ֔וֹ וַיַּ֛עַל לִקְרַֽאת־יִשְׂרָאֵ֥ל אָבִ֖יו גֹּ֑שְׁנָה וַיֵּרָ֣א אֵלָ֗יו וַיִּפֹּל֙ עַל־צַוָּארָ֔יו וַיֵּ֥בְךְּ עַל־צַוָּארָ֖יו עֽוֹד׃ \Rashi{\textbf{ויאסר יוסף מרכבתו.}הוא עצמו אסר את הסוסים למרכבה להזדרז לכבוד אביו: וירא אליו. יוסף נראה אל אביו: ויבך על צואריו עוד. לשון הרבות בכיה, וכן כי לא על איש ישים עוד (איוב ל"ד), לשון רבוי הוא – אינו שם עליו עלילות נוספות על חטאיו; אף כאן הרבה והוסיף בבכי יותר על הרגיל; אבל יעקב לא נפל על צוארי יוסף ולא נשקו, ואמרו רבותינו, שהיה קורא את שמע:}%
% ---------- פרק 46 | פסוק 30 ----------
 \textbf{ל} וַיֹּ֧אמֶר יִשְׂרָאֵ֛ל אֶל־יוֹסֵ֖ף אָמ֣וּתָה הַפָּ֑עַם אַחֲרֵי֙ רְאוֹתִ֣י אֶת־פָּנֶ֔יךָ כִּ֥י עוֹדְךָ֖ חָֽי׃ \Rashi{\textbf{אמותה הפעם.}פשוטו כתרגומו. ומדרשו סבור הייתי למות שתי מיתות, בעולם הזה ולעולם הבא, שנסתלקה ממני שכינה, והייתי אומר שיתבעני הקב"ה מיתתך, עכשו שעודך חי, לא אמות אלא פעם אחת:}%
% ---------- פרק 46 | פסוק 31 ----------
 \textbf{לא} וַיֹּ֨אמֶר יוֹסֵ֤ף אֶל־אֶחָיו֙ וְאֶל־בֵּ֣ית אָבִ֔יו אֶעֱלֶ֖ה וְאַגִּ֣ידָה לְפַרְעֹ֑ה וְאֹֽמְרָ֣ה אֵלָ֔יו אַחַ֧י וּבֵית־אָבִ֛י אֲשֶׁ֥ר בְּאֶֽרֶץ־כְּנַ֖עַן בָּ֥אוּ אֵלָֽי׃ \Rashi{\textbf{ואמרה אליו אחי וגו'.}ועוד אמר לו, והאנשים רועי צאן וגו':}%
% ---------- פרק 46 | פסוק 32 ----------
 \textbf{לב} וְהָאֲנָשִׁים֙ רֹ֣עֵי צֹ֔אן כִּֽי־אַנְשֵׁ֥י מִקְנֶ֖ה הָי֑וּ וְצֹאנָ֧ם וּבְקָרָ֛ם וְכׇל־אֲשֶׁ֥ר לָהֶ֖ם הֵבִֽיאוּ׃ %
% ---------- פרק 46 | פסוק 33 ----------
 \textbf{לג} וְהָיָ֕ה כִּֽי־יִקְרָ֥א לָכֶ֖ם פַּרְעֹ֑ה וְאָמַ֖ר מַה־מַּעֲשֵׂיכֶֽם׃ %
% ---------- פרק 46 | פסוק 34 ----------
 \textbf{לד} וַאֲמַרְתֶּ֗ם אַנְשֵׁ֨י מִקְנֶ֜ה הָי֤וּ עֲבָדֶ֙יךָ֙ מִנְּעוּרֵ֣ינוּ וְעַד־עַ֔תָּה גַּם־אֲנַ֖חְנוּ גַּם־אֲבֹתֵ֑ינוּ בַּעֲב֗וּר תֵּשְׁבוּ֙ בְּאֶ֣רֶץ גֹּ֔שֶׁן כִּֽי־תוֹעֲבַ֥ת מִצְרַ֖יִם כׇּל־רֹ֥עֵה צֹֽאן׃ \Rashi{\textbf{בעבור תשבו בארץ גשן.}והיא צריכה לכם שהיא ארץ מרעה, וכשתאמרו לו שאין אתם בקיאין במלאכה אחרת, ירחיקכם מעליו ויושיבכם שם: כי תועבת מצרים כל רועה צאן. לפי שהם להם אלהות:}%
% ---------- פרק 47 | פסוק 1 ----------
 \textbf{א} וַיָּבֹ֣א יוֹסֵף֮ וַיַּגֵּ֣ד לְפַרְעֹה֒ וַיֹּ֗אמֶר אָבִ֨י וְאַחַ֜י וְצֹאנָ֤ם וּבְקָרָם֙ וְכׇל־אֲשֶׁ֣ר לָהֶ֔ם בָּ֖אוּ מֵאֶ֣רֶץ כְּנָ֑עַן וְהִנָּ֖ם בְּאֶ֥רֶץ גֹּֽשֶׁן׃ %
% ---------- פרק 47 | פסוק 2 ----------
 \textbf{ב} וּמִקְצֵ֣ה אֶחָ֔יו לָקַ֖ח חֲמִשָּׁ֣ה אֲנָשִׁ֑ים וַיַּצִּגֵ֖ם לִפְנֵ֥י פַרְעֹֽה׃ \Rashi{\textbf{ומקצה אחיו.}מן הפחותים שבהם לגבורה, שאין נראים גבורים, שאם יראה אותם גבורים יעשה אותם אנשי מלחמתו, ואלה הם: ראובן, שמעון, לוי, יששכר ובנימין – אותן שלא כפל משה שמותם כשברכן; אבל שמות הגבורים כפל וזאת ליהודה שמע ה' קול יהודה; ולגד אמר ברוך מרחיב גד; ולנפתלי אמר נפתלי, ולדן אמר דן, וכן לזבולון, וכן לאשר; זהו לשון ב"ר, שהיא אגדת ארץ ישראל. אבל בתלמוד בבלית שלנו (בבא קמא צ"ב) מצינו, שאותן שכפל משה שמותן הם החלשים, ואותן הביא לפני פרעה, ויהודה שהכפל שמו לא הכפל משום חלשות, אלא טעם יש בדבר, כדאיתא בב"ק ובבריתא דספרי שנינו בה בוזאת הברכה, כמו תלמוד שלנו:}%
% ---------- פרק 47 | פסוק 3 ----------
 \textbf{ג} וַיֹּ֧אמֶר פַּרְעֹ֛ה אֶל־אֶחָ֖יו מַה־מַּעֲשֵׂיכֶ֑ם וַיֹּאמְר֣וּ אֶל־פַּרְעֹ֗ה רֹעֵ֥ה צֹאן֙ עֲבָדֶ֔יךָ גַּם־אֲנַ֖חְנוּ גַּם־אֲבוֹתֵֽינוּ׃ %
% ---------- פרק 47 | פסוק 4 ----------
 \textbf{ד} וַיֹּאמְר֣וּ אֶל־פַּרְעֹ֗ה לָג֣וּר בָּאָ֘רֶץ֮ בָּ֒אנוּ֒ כִּי־אֵ֣ין מִרְעֶ֗ה לַצֹּאן֙ אֲשֶׁ֣ר לַעֲבָדֶ֔יךָ כִּֽי־כָבֵ֥ד הָרָעָ֖ב בְּאֶ֣רֶץ כְּנָ֑עַן וְעַתָּ֛ה יֵֽשְׁבוּ־נָ֥א עֲבָדֶ֖יךָ בְּאֶ֥רֶץ גֹּֽשֶׁן׃ %
% ---------- פרק 47 | פסוק 5 ----------
 \textbf{ה} וַיֹּ֣אמֶר פַּרְעֹ֔ה אֶל־יוֹסֵ֖ף לֵאמֹ֑ר אָבִ֥יךָ וְאַחֶ֖יךָ בָּ֥אוּ אֵלֶֽיךָ׃ %
% ---------- פרק 47 | פסוק 6 ----------
 \textbf{ו} אֶ֤רֶץ מִצְרַ֙יִם֙ לְפָנֶ֣יךָ הִ֔וא בְּמֵיטַ֣ב הָאָ֔רֶץ הוֹשֵׁ֥ב אֶת־אָבִ֖יךָ וְאֶת־אַחֶ֑יךָ יֵשְׁבוּ֙ בְּאֶ֣רֶץ גֹּ֔שֶׁן וְאִם־יָדַ֗עְתָּ וְיֶשׁ־בָּם֙ אַנְשֵׁי־חַ֔יִל וְשַׂמְתָּ֛ם שָׂרֵ֥י מִקְנֶ֖ה עַל־אֲשֶׁר־לִֽי׃ \Rashi{\textbf{אנשי חיל.}בקיאין באמנותן לרעות צאן: על אשר לי. על צאן שלי:}%
% ---------- פרק 47 | פסוק 7 ----------
 \textbf{ז} וַיָּבֵ֤א יוֹסֵף֙ אֶת־יַֽעֲקֹ֣ב אָבִ֔יו וַיַּֽעֲמִדֵ֖הוּ לִפְנֵ֣י פַרְעֹ֑ה וַיְבָ֥רֶךְ יַעֲקֹ֖ב אֶת־פַּרְעֹֽה׃ \Rashi{\textbf{ויברך יעקב.}היא שאילת שלום כדרך כל הנראים לפני המלכים לפרקים, שלו"איר בלעז:}%
% ---------- פרק 47 | פסוק 8 ----------
 \textbf{ח} וַיֹּ֥אמֶר פַּרְעֹ֖ה אֶֽל־יַעֲקֹ֑ב כַּמָּ֕ה יְמֵ֖י שְׁנֵ֥י חַיֶּֽיךָ׃ %
% ---------- פרק 47 | פסוק 9 ----------
 \textbf{ט} וַיֹּ֤אמֶר יַעֲקֹב֙ אֶל־פַּרְעֹ֔ה יְמֵי֙ שְׁנֵ֣י מְגוּרַ֔י שְׁלֹשִׁ֥ים וּמְאַ֖ת שָׁנָ֑ה מְעַ֣ט וְרָעִ֗ים הָיוּ֙ יְמֵי֙ שְׁנֵ֣י חַיַּ֔י וְלֹ֣א הִשִּׂ֗יגוּ אֶת־יְמֵי֙ שְׁנֵי֙ חַיֵּ֣י אֲבֹתַ֔י בִּימֵ֖י מְגוּרֵיהֶֽם׃ \Rashi{\textbf{שני מגורי.}ימי גרותי; כל ימי הייתי גר בארץ: ולא השיגו. בטובה:}%
% ---------- פרק 47 | פסוק 10 ----------
 \textbf{י} וַיְבָ֥רֶךְ יַעֲקֹ֖ב אֶת־פַּרְעֹ֑ה וַיֵּצֵ֖א מִלִּפְנֵ֥י פַרְעֹֽה׃ \Rashi{\textbf{ויברך יעקב.}כדרך כל הנפטרים מלפני שרים מברכים אותם ונוטלים רשות, ומה ברכה ברכו? שיעלה נילוס לרגליו, לפי שאין מצרים שותה מי גשמים אלא נילוס עולה ומשקה, ומברכתו של יעקב ואילך היה פרעה בא אל נילוס והוא עולה לקראתו ומשקה את הארץ (תנחומא):}%
% ---------- פרק 47 | פסוק 11 ----------
 \textbf{יא} וַיּוֹשֵׁ֣ב יוֹסֵף֮ אֶת־אָבִ֣יו וְאֶת־אֶחָיו֒ וַיִּתֵּ֨ן לָהֶ֤ם אֲחֻזָּה֙ בְּאֶ֣רֶץ מִצְרַ֔יִם בְּמֵיטַ֥ב הָאָ֖רֶץ בְּאֶ֣רֶץ רַעְמְסֵ֑ס כַּאֲשֶׁ֖ר צִוָּ֥ה פַרְעֹֽה׃ \Rashi{\textbf{רעמסס.}מארץ גשן היא:}%
% ---------- פרק 47 | פסוק 12 ----------
 \textbf{יב} וַיְכַלְכֵּ֤ל יוֹסֵף֙ אֶת־אָבִ֣יו וְאֶת־אֶחָ֔יו וְאֵ֖ת כׇּל־בֵּ֣ית אָבִ֑יו לֶ֖חֶם לְפִ֥י הַטָּֽף׃ \Rashi{\textbf{לפי הטף.}לפי הצריך לכל בני ביתם:}%
% ---------- פרק 47 | פסוק 13 ----------
 \textbf{יג} וְלֶ֤חֶם אֵין֙ בְּכׇל־הָאָ֔רֶץ כִּֽי־כָבֵ֥ד הָרָעָ֖ב מְאֹ֑ד וַתֵּ֜לַהּ אֶ֤רֶץ מִצְרַ֙יִם֙ וְאֶ֣רֶץ כְּנַ֔עַן מִפְּנֵ֖י הָרָעָֽב׃ \Rashi{\textbf{ולחם אין בכל הארץ.}חוזר לענין הראשון – לתחלת שני הרעב: ותלה. כמו ותלאה, לשון עיפות, כתרגומו; ודומה לו כמתלהלה הירה זקים (משלי כ"ו):}%
% ---------- פרק 47 | פסוק 14 ----------
 \textbf{יד} וַיְלַקֵּ֣ט יוֹסֵ֗ף אֶת־כׇּל־הַכֶּ֙סֶף֙ הַנִּמְצָ֤א בְאֶֽרֶץ־מִצְרַ֙יִם֙ וּבְאֶ֣רֶץ כְּנַ֔עַן בַּשֶּׁ֖בֶר אֲשֶׁר־הֵ֣ם שֹׁבְרִ֑ים וַיָּבֵ֥א יוֹסֵ֛ף אֶת־הַכֶּ֖סֶף בֵּ֥יתָה פַרְעֹֽה׃ \Rashi{\textbf{בשבר אשר הם שברים.}נותנין לו את הכסף:}%
% ---------- פרק 47 | פסוק 15 ----------
 \textbf{טו} וַיִּתֹּ֣ם הַכֶּ֗סֶף מֵאֶ֣רֶץ מִצְרַ֘יִם֮ וּמֵאֶ֣רֶץ כְּנַ֒עַן֒ וַיָּבֹ֩אוּ֩ כׇל־מִצְרַ֨יִם אֶל־יוֹסֵ֤ף לֵאמֹר֙ הָֽבָה־לָּ֣נוּ לֶ֔חֶם וְלָ֥מָּה נָמ֖וּת נֶגְדֶּ֑ךָ כִּ֥י אָפֵ֖ס כָּֽסֶף׃ \Rashi{\textbf{אפס.}כתרגומו שלים:}%
% ---------- פרק 47 | פסוק 16 ----------
 \textbf{טז} וַיֹּ֤אמֶר יוֹסֵף֙ הָב֣וּ מִקְנֵיכֶ֔ם וְאֶתְּנָ֥ה לָכֶ֖ם בְּמִקְנֵיכֶ֑ם אִם־אָפֵ֖ס כָּֽסֶף׃ %
% ---------- פרק 47 | פסוק 17 ----------
 \textbf{יז} וַיָּבִ֣יאוּ אֶת־מִקְנֵיהֶם֮ אֶל־יוֹסֵף֒ וַיִּתֵּ֣ן לָהֶם֩ יוֹסֵ֨ף לֶ֜חֶם בַּסּוּסִ֗ים וּבְמִקְנֵ֥ה הַצֹּ֛אן וּבְמִקְנֵ֥ה הַבָּקָ֖ר וּבַחֲמֹרִ֑ים וַיְנַהֲלֵ֤ם בַּלֶּ֙חֶם֙ בְּכׇל־מִקְנֵהֶ֔ם בַּשָּׁנָ֖ה הַהִֽוא׃ \Rashi{\textbf{וינהלם.}כמו וינהגם, ודומה לו אין מנהל לה (ישעיהו נ"א), על מי מנחות ינהלני (תהילים כ"ג):}%
% ---------- פרק 47 | פסוק 18 ----------
 \textbf{יח} וַתִּתֹּם֮ הַשָּׁנָ֣ה הַהִוא֒ וַיָּבֹ֨אוּ אֵלָ֜יו בַּשָּׁנָ֣ה הַשֵּׁנִ֗ית וַיֹּ֤אמְרוּ לוֹ֙ לֹֽא־נְכַחֵ֣ד מֵֽאֲדֹנִ֔י כִּ֚י אִם־תַּ֣ם הַכֶּ֔סֶף וּמִקְנֵ֥ה הַבְּהֵמָ֖ה אֶל־אֲדֹנִ֑י לֹ֤א נִשְׁאַר֙ לִפְנֵ֣י אֲדֹנִ֔י בִּלְתִּ֥י אִם־גְּוִיָּתֵ֖נוּ וְאַדְמָתֵֽנוּ׃ \Rashi{\textbf{בשנה השנית.}לשני הרעב: כי אם תם הכסף וגו'. כי אשר תם הכסף והמקנה ובא הכל אל יד אדוני: בלתי אם גויתנו. כמו אם לא גויתנו:}%
% ---------- פרק 47 | פסוק 19 ----------
 \textbf{יט} לָ֧מָּה נָמ֣וּת לְעֵינֶ֗יךָ גַּם־אֲנַ֙חְנוּ֙ גַּ֣ם אַדְמָתֵ֔נוּ קְנֵֽה־אֹתָ֥נוּ וְאֶת־אַדְמָתֵ֖נוּ בַּלָּ֑חֶם וְנִֽהְיֶ֞ה אֲנַ֤חְנוּ וְאַדְמָתֵ֙נוּ֙ עֲבָדִ֣ים לְפַרְעֹ֔ה וְתֶן־זֶ֗רַע וְנִֽחְיֶה֙ וְלֹ֣א נָמ֔וּת וְהָאֲדָמָ֖ה לֹ֥א תֵשָֽׁם׃ \Rashi{\textbf{ותן זרע.}לזרע האדמה; ואף על פי שאמר יוסף ועוד חמש שנים אשר אין חריש וקציר, מכיון שבא יעקב למצרים באה ברכה לרגליו והתחיל לזרע, וכלה הרעב, וכן שנינו בתוספתא דסוטה: לא תשם. לא תהא שממה; לא תבור – לשון שדה בור, שאינו חרוש:}%
% ---------- פרק 47 | פסוק 20 ----------
 \textbf{כ} וַיִּ֨קֶן יוֹסֵ֜ף אֶת־כׇּל־אַדְמַ֤ת מִצְרַ֙יִם֙ לְפַרְעֹ֔ה כִּֽי־מָכְר֤וּ מִצְרַ֙יִם֙ אִ֣ישׁ שָׂדֵ֔הוּ כִּֽי־חָזַ֥ק עֲלֵהֶ֖ם הָרָעָ֑ב וַתְּהִ֥י הָאָ֖רֶץ לְפַרְעֹֽה׃ \Rashi{\textbf{ותהי הארץ לפרעה.}קנויה לו:}%
% ---------- פרק 47 | פסוק 21 ----------
 \textbf{כא} וְאֶ֨ת־הָעָ֔ם הֶעֱבִ֥יר אֹת֖וֹ לֶעָרִ֑ים מִקְצֵ֥ה גְבוּל־מִצְרַ֖יִם וְעַד־קָצֵֽהוּ׃ \Rashi{\textbf{ואת העם העביר.}יוסף מעיר לעיר לזכרון, שאין להם עוד חלק בארץ, והושיב של עיר זו בחברתה, ולא הצרך הכתוב לכתב זאת אלא להודיעך שבחו של יוסף שנתכון להסיר חרפה מעל אחיו, שלא יהיו קורין אותם גולים (חולין ס'): מקצה גבול מצרים וגו'. כן עשה לכל הערים אשר במלכות מצרים, מקצה גבולה ועד קצה גבולה:}%
% ---------- פרק 47 | פסוק 22 ----------
 \textbf{כב} רַ֛ק אַדְמַ֥ת הַכֹּהֲנִ֖ים לֹ֣א קָנָ֑ה כִּי֩ חֹ֨ק לַכֹּהֲנִ֜ים מֵאֵ֣ת פַּרְעֹ֗ה וְאָֽכְל֤וּ אֶת־חֻקָּם֙ אֲשֶׁ֨ר נָתַ֤ן לָהֶם֙ פַּרְעֹ֔ה עַל־כֵּ֕ן לֹ֥א מָכְר֖וּ אֶת־אַדְמָתָֽם׃ \Rashi{\textbf{הכהנים.}הכומרים, כל לשון כהן משרת לאלהות הוא, חוץ מאותן שהם לשון גדלה, כמו כהן מדין, כהן און: חק לכהנים. חק כך וכך לחם ליום:}%
% ---------- פרק 47 | פסוק 23 ----------
 \textbf{כג} וַיֹּ֤אמֶר יוֹסֵף֙ אֶל־הָעָ֔ם הֵן֩ קָנִ֨יתִי אֶתְכֶ֥ם הַיּ֛וֹם וְאֶת־אַדְמַתְכֶ֖ם לְפַרְעֹ֑ה הֵֽא־לָכֶ֣ם זֶ֔רַע וּזְרַעְתֶּ֖ם אֶת־הָאֲדָמָֽה׃ \Rashi{\textbf{הא.}כמו הנה, כמו וגם אני הא דרכך בראש נתתי (יחזקאל ט"ז):}%
% ---------- פרק 47 | פסוק 24 ----------
 \textbf{כד} וְהָיָה֙ בַּתְּבוּאֹ֔ת וּנְתַתֶּ֥ם חֲמִישִׁ֖ית לְפַרְעֹ֑ה וְאַרְבַּ֣ע הַיָּדֹ֡ת יִהְיֶ֣ה לָכֶם֩ לְזֶ֨רַע הַשָּׂדֶ֧ה וּֽלְאׇכְלְכֶ֛ם וְלַאֲשֶׁ֥ר בְּבָתֵּיכֶ֖ם וְלֶאֱכֹ֥ל לְטַפְּכֶֽם׃ \Rashi{\textbf{לזרע השדה.}שבכל שנה: ולאשר בבתיכם. ולאכל העבדים והשפחות אשר בבתיכם: לטפכם. בנים קטנים:}%
% ---------- פרק 47 | פסוק 25 ----------
 \textbf{כה} וַיֹּאמְר֖וּ הֶחֱיִתָ֑נוּ נִמְצָא־חֵן֙ בְּעֵינֵ֣י אֲדֹנִ֔י וְהָיִ֥ינוּ עֲבָדִ֖ים לְפַרְעֹֽה׃ \Rashi{\textbf{נמצא חן.}לעשות לנו זאת כמו שאמרת: והיינו עבדים לפרעה. להעלות לו המס הזה בכל שנה:}%
% ---------- פרק 47 | פסוק 26 ----------
 \textbf{כו} וַיָּ֣שֶׂם אֹתָ֣הּ יוֹסֵ֡ף לְחֹק֩ עַד־הַיּ֨וֹם הַזֶּ֜ה עַל־אַדְמַ֥ת מִצְרַ֛יִם לְפַרְעֹ֖ה לַחֹ֑מֶשׁ רַ֞ק אַדְמַ֤ת הַכֹּֽהֲנִים֙ לְבַדָּ֔ם לֹ֥א הָיְתָ֖ה לְפַרְעֹֽה׃ \Rashi{\textbf{לחק.}שלא יעבר:}%
% ---------- פרק 47 | פסוק 27 ----------
 \textbf{כז} וַיֵּ֧שֶׁב יִשְׂרָאֵ֛ל בְּאֶ֥רֶץ מִצְרַ֖יִם בְּאֶ֣רֶץ גֹּ֑שֶׁן וַיֵּאָחֲז֣וּ בָ֔הּ וַיִּפְר֥וּ וַיִּרְבּ֖וּ מְאֹֽד׃ \Rashi{\textbf{וישב ישראל בארץ מצרים.}והיכן? בארץ גשן שהיא מארץ מצרים: (ויאחזו בה. לשון אחזה):}%
\addparasha{חומש בראשית | פרשת ויחי}
% ---------- פרק 47 | פסוק 28 ----------
 \textbf{כח} וַיְחִ֤י יַעֲקֹב֙ בְּאֶ֣רֶץ מִצְרַ֔יִם שְׁבַ֥ע עֶשְׂרֵ֖ה שָׁנָ֑ה וַיְהִ֤י יְמֵֽי־יַעֲקֹב֙ שְׁנֵ֣י חַיָּ֔יו שֶׁ֣בַע שָׁנִ֔ים וְאַרְבָּעִ֥ים וּמְאַ֖ת שָׁנָֽה׃ \Rashi{\textbf{ויחי יעקב.}למה פרשה זו סתומה? לפי שכיון שנפטר יעקב אבינו נסתמו עיניהם ולבם של ישראל מצרת השעבוד, שהתחילו לשעבדם; דבר אחר: שבקש לגלות את הקץ לבניו, ונסתם ממנו. בב"ר:}%
% ---------- פרק 47 | פסוק 29 ----------
 \textbf{כט} וַיִּקְרְב֣וּ יְמֵֽי־יִשְׂרָאֵל֮ לָמוּת֒ וַיִּקְרָ֣א ׀ לִבְנ֣וֹ לְיוֹסֵ֗ף וַיֹּ֤אמֶר לוֹ֙ אִם־נָ֨א מָצָ֤אתִי חֵן֙ בְּעֵינֶ֔יךָ שִֽׂים־נָ֥א יָדְךָ֖ תַּ֣חַת יְרֵכִ֑י וְעָשִׂ֤יתָ עִמָּדִי֙ חֶ֣סֶד וֶאֱמֶ֔ת אַל־נָ֥א תִקְבְּרֵ֖נִי בְּמִצְרָֽיִם׃ \Rashi{\textbf{ויקרבו ימי ישראל למות.}כל מי שנאמרה בו קריבה למות, לא הגיע לימי אבותיו (יצחק חיה ק"פ, ויעקב קמ"ז, בדוד נאמר קריבה, אביו חיה פ' שנים והוא חיה ע'): ויקרא לבנו ליוסף. למי שהיה יכלת בידו לעשות: שים נא ידך. והשבע: חסד ואמת. חסד שעושין עם המתים הוא חסד של אמת, שאינו מצפה לתשלום גמול: אל נא תקברני במצרים. סופה להיות עפרה כנים, ומרחשין תחת גופי, ושאין מתי חוצה לארץ חיים אלא בצער גלגול מחלות, ושלא יעשוני מצרים עבודה זרה:}%
% ---------- פרק 47 | פסוק 30 ----------
 \textbf{ל} וְשָֽׁכַבְתִּי֙ עִם־אֲבֹתַ֔י וּנְשָׂאתַ֙נִי֙ מִמִּצְרַ֔יִם וּקְבַרְתַּ֖נִי בִּקְבֻרָתָ֑ם וַיֹּאמַ֕ר אָנֹכִ֖י אֶֽעֱשֶׂ֥ה כִדְבָרֶֽךָ׃ \Rashi{\textbf{ושכבתי עם אבותי.}וי"ו זו מחבר למעלה לתחלת המקרא – שים נא ידך תחת ירכי והשבע לי – ואני סופי לשכב עם אבותי, ואתה תשאני ממצרים. ואין לומר, ושכבתי עם אבותי – השכיבני עם אבותי במערה, שהרי כתיב אחריו, ונשאתני ממצרים וקברתני בקברתם; ועוד, מצינו בכל מקום לשון שכיבה עם אבותיו היא הגויעה ולא הקבורה, כמו: וישכב דוד עם אבתיו, ואחר כך ויקבר בעיר דוד (מלכים א ב'):}%
% ---------- פרק 47 | פסוק 31 ----------
 \textbf{לא} וַיֹּ֗אמֶר הִשָּֽׁבְעָה֙ לִ֔י וַיִּשָּׁבַ֖ע ל֑וֹ וַיִּשְׁתַּ֥חוּ יִשְׂרָאֵ֖ל עַל־רֹ֥אשׁ הַמִּטָּֽה׃ \{פ\} \Rashi{\textbf{וישתחו ישראל.}תעלא בעידניה סגיד ליה: על ראש המטה. הפך עצמו לצד השכינה. מכאן אמרו שהשכינה למעלה מראשותיו של חולה (שבת י"ב, נדרים מ'); ד"א על ראש המטה – על שהיתה מטתו שלמה ולא היה בה רשע, שהרי יוסף מלך הוא, ועוד שנשבה לבין הגוים, והרי הוא עומד בצדקו:}%
% ---------- פרק 48 | פסוק 1 ----------
 \textbf{א} וַיְהִ֗י אַחֲרֵי֙ הַדְּבָרִ֣ים הָאֵ֔לֶּה וַיֹּ֣אמֶר לְיוֹסֵ֔ף הִנֵּ֥ה אָבִ֖יךָ חֹלֶ֑ה וַיִּקַּ֞ח אֶת־שְׁנֵ֤י בָנָיו֙ עִמּ֔וֹ אֶת־מְנַשֶּׁ֖ה וְאֶת־אֶפְרָֽיִם׃ \Rashi{\textbf{ויאמר ליוסף.}אחד מן המגידים, והרי זה מקרא קצר; וי"א, אפרים היה רגיל לפני יעקב בתלמוד, וכשחלה יעקב בארץ גשן, הלך אפרים אצל אביו למצרים להגיד לו: ויקח את שני בניו עמו. כדי שיברכם יעקב לפני מותו:}%
% ---------- פרק 48 | פסוק 2 ----------
 \textbf{ב} וַיַּגֵּ֣ד לְיַעֲקֹ֔ב וַיֹּ֕אמֶר הִנֵּ֛ה בִּנְךָ֥ יוֹסֵ֖ף בָּ֣א אֵלֶ֑יךָ וַיִּתְחַזֵּק֙ יִשְׂרָאֵ֔ל וַיֵּ֖שֶׁב עַל־הַמִּטָּֽה׃ \Rashi{\textbf{ויגד.}המגיד ליעקב, ולא פרש מי; והרבה מקראות קצרי לשון: ויתחזק ישראל. אמר, אף על פי שהוא בני, מלך הוא, אחלק לו כבוד; מכאן שחולקין כבוד למלכות, וכן משה חלק כבוד למלכות – וירדו כל עבדיך אלה אלי (שמות י"א); וכן אליהו, וישנס מתניו וגו' (מלכים א י"ח):}%
% ---------- פרק 48 | פסוק 3 ----------
 \textbf{ג} וַיֹּ֤אמֶר יַעֲקֹב֙ אֶל־יוֹסֵ֔ף אֵ֥ל שַׁדַּ֛י נִרְאָֽה־אֵלַ֥י בְּל֖וּז בְּאֶ֣רֶץ כְּנָ֑עַן וַיְבָ֖רֶךְ אֹתִֽי׃ %
% ---------- פרק 48 | פסוק 4 ----------
 \textbf{ד} וַיֹּ֣אמֶר אֵלַ֗י הִנְנִ֤י מַפְרְךָ֙ וְהִרְבִּיתִ֔ךָ וּנְתַתִּ֖יךָ לִקְהַ֣ל עַמִּ֑ים וְנָ֨תַתִּ֜י אֶת־הָאָ֧רֶץ הַזֹּ֛את לְזַרְעֲךָ֥ אַחֲרֶ֖יךָ אֲחֻזַּ֥ת עוֹלָֽם׃ \Rashi{\textbf{ונתתיך לקהל עמים.}בשרני שעתידים לצאת ממני עוד קהל ועמים, ואף על פי שאמר לי גוי וקהל גוים, גוי אמר לי על בנימין; קהל גוים – הרי שנים לבד מבנימין, ושוב לא נולד לי בן, למדני שעתיד אחד משבטי להחלק, ועתה אותה מתנה אני נותן לך:}%
% ---------- פרק 48 | פסוק 5 ----------
 \textbf{ה} וְעַתָּ֡ה שְׁנֵֽי־בָנֶ֩יךָ֩ הַנּוֹלָדִ֨ים לְךָ֜ בְּאֶ֣רֶץ מִצְרַ֗יִם עַד־בֹּאִ֥י אֵלֶ֛יךָ מִצְרַ֖יְמָה לִי־הֵ֑ם אֶפְרַ֙יִם֙ וּמְנַשֶּׁ֔ה כִּרְאוּבֵ֥ן וְשִׁמְע֖וֹן יִֽהְיוּ־לִֽי׃ \Rashi{\textbf{הנולדים לך, עד באי אליך.}לפני בואי אליך, כלומר, שנולדו משפרשת ממני עד שבאתי אצלך: לי הם. בחשבון שאר בני הם, לטל חלק בארץ איש כנגדו:}%
% ---------- פרק 48 | פסוק 6 ----------
 \textbf{ו} וּמוֹלַדְתְּךָ֛ אֲשֶׁר־הוֹלַ֥דְתָּ אַחֲרֵיהֶ֖ם לְךָ֣ יִהְי֑וּ עַ֣ל שֵׁ֧ם אֲחֵיהֶ֛ם יִקָּרְא֖וּ בְּנַחֲלָתָֽם׃ \Rashi{\textbf{ומולדתך וגו'.}אם תוליד עוד, לא יהיו במנין בני, אלא בתוך שבטי אפרים ומנשה יהיו נכללים, ולא יהא להם שם בשבטים לענין הנחלה; ואף על פי שנחלקה הארץ למנין גלגלותם, כדכתיב לרב תרבה נחלתו (במדבר כ"ו), וכל איש ואיש נטל בשוה חוץ מן הבכורות, מכל מקום לא נקראו שבטים אלא אלו, (להטיל גורל הארץ למנין שמות השבטים ונשיא לכל שבט ושבט ודגלים לזה ולזה):}%
% ---------- פרק 48 | פסוק 7 ----------
 \textbf{ז} וַאֲנִ֣י ׀ בְּבֹאִ֣י מִפַּדָּ֗ן מֵ֩תָה֩ עָלַ֨י רָחֵ֜ל בְּאֶ֤רֶץ כְּנַ֙עַן֙ בַּדֶּ֔רֶךְ בְּע֥וֹד כִּבְרַת־אֶ֖רֶץ לָבֹ֣א אֶפְרָ֑תָה וָאֶקְבְּרֶ֤הָ שָּׁם֙ בְּדֶ֣רֶךְ אֶפְרָ֔ת הִ֖וא בֵּ֥ית לָֽחֶם׃ \Rashi{\textbf{ואני בבאי מפדן וגו'.}ואע"פ שאני מטריח עליך להוליכני להקבר בארץ כנען, ולא כך עשיתי לאמך, שהרי מתה סמוך לבית לחם: כברת ארץ. מדת ארץ, והם אלפים אמה כמדת תחום שבת, כדברי רבי משה הדרשן – ולא תאמר שעכבו עלי גשמים מלהוליכה ולקברה בחברון, עת הגריד היה, שהארץ חלולה ומנקבת ככברה: אקברה שם. ולא הולכתיה אפלו לבית לחם להכניסה לארץ, וידעתי שיש בלבך עלי; אבל דע לך שעל פי הדבור קברתיה שם, שתהא לעזרה לבניה כשיגלה אותם נבוזראדן, והיו עוברים דרך שם, יוצאת רחל על קברה ובוכה ומבקשת עליהם רחמים, שנאמר קול ברמה נשמע וגו' והקב"ה משיבה יש שכר לפעלתך נאם ה' ושבו בנים לגבולם (ירמיהו ל"א). ואנקלוס תרגם כרוב ארעא – כדי שעור חרישת יום; ואומר אני, שהיה להם קצב שהיו קורין אותו כדי מחרשה אחת, קורדיי"א בלעז, כדאמרינן כריב ותני; כמה דמסיק תעלא מבי כרבא:}%
% ---------- פרק 48 | פסוק 8 ----------
 \textbf{ח} וַיַּ֥רְא יִשְׂרָאֵ֖ל אֶת־בְּנֵ֣י יוֹסֵ֑ף וַיֹּ֖אמֶר מִי־אֵֽלֶּה׃ \Rashi{\textbf{וירא ישראל את בני יוסף.}בקש לברכם ונסתלקה שכינה ממנו, לפי שעתיד ירבעם ואחאב לצאת מאפרים ויהוא ובניו ממנשה: ויאמר מי אלה. מהיכן יצאו אלו, שאינן ראויין לברכה:}%
% ---------- פרק 48 | פסוק 9 ----------
 \textbf{ט} וַיֹּ֤אמֶר יוֹסֵף֙ אֶל־אָבִ֔יו בָּנַ֣י הֵ֔ם אֲשֶׁר־נָֽתַן־לִ֥י אֱלֹהִ֖ים בָּזֶ֑ה וַיֹּאמַ֕ר קָֽחֶם־נָ֥א אֵלַ֖י וַאֲבָרְכֵֽם׃ \Rashi{\textbf{בזה.}הראה לו שטר ארוסין ושטר כתבה, ובקש יוסף רחמים על הדבר, ונחה עליו רוח הקדש: ויאמר קחם נא אלי ואברכם. זהו שאמר הכתוב ואנכי תרגלתי לאפרים קחם על זרועתיו (הושע י"א) – תרגלתי רוחי ביעקב בשביל אפרים עד שלקחן על זרועותיו:}%
% ---------- פרק 48 | פסוק 10 ----------
 \textbf{י} וְעֵינֵ֤י יִשְׂרָאֵל֙ כָּבְד֣וּ מִזֹּ֔קֶן לֹ֥א יוּכַ֖ל לִרְא֑וֹת וַיַּגֵּ֤שׁ אֹתָם֙ אֵלָ֔יו וַיִּשַּׁ֥ק לָהֶ֖ם וַיְחַבֵּ֥ק לָהֶֽם׃ %
% ---------- פרק 48 | פסוק 11 ----------
 \textbf{יא} וַיֹּ֤אמֶר יִשְׂרָאֵל֙ אֶל־יוֹסֵ֔ף רְאֹ֥ה פָנֶ֖יךָ לֹ֣א פִלָּ֑לְתִּי וְהִנֵּ֨ה הֶרְאָ֥ה אֹתִ֛י אֱלֹהִ֖ים גַּ֥ם אֶת־זַרְעֶֽךָ׃ \Rashi{\textbf{לא פללתי.}לא מלאני לבי לחשב מחשבה שאראה פניך עוד. פללתי לשון מחשבה, כמו הביאי עצה עשי פלילה (ישעיהו ט"ז):}%
% ---------- פרק 48 | פסוק 12 ----------
 \textbf{יב} וַיּוֹצֵ֥א יוֹסֵ֛ף אֹתָ֖ם מֵעִ֣ם בִּרְכָּ֑יו וַיִּשְׁתַּ֥חוּ לְאַפָּ֖יו אָֽרְצָה׃ \Rashi{\textbf{ויוצא יוסף אתם.}לאחר שנשקם הוציאם יוסף מעם ברכיו, כדי לישבם זה לימין וזה לשמאל לסמך ידיו עליהם ולברכם: וישתחו לאפיו. כשחזר לאחוריו מלפני אביו:}%
% ---------- פרק 48 | פסוק 13 ----------
 \textbf{יג} וַיִּקַּ֣ח יוֹסֵף֮ אֶת־שְׁנֵיהֶם֒ אֶת־אֶפְרַ֤יִם בִּֽימִינוֹ֙ מִשְּׂמֹ֣אל יִשְׂרָאֵ֔ל וְאֶת־מְנַשֶּׁ֥ה בִשְׂמֹאל֖וֹ מִימִ֣ין יִשְׂרָאֵ֑ל וַיַּגֵּ֖שׁ אֵלָֽיו׃ \Rashi{\textbf{את אפרים בימינו משמאל ישראל.}הבא לקראת חברו, ימינו כנגד שמאל חברו, וכיון שהוא הבכור, מימן לברכה:}%
% ---------- פרק 48 | פסוק 14 ----------
 \textbf{יד} וַיִּשְׁלַח֩ יִשְׂרָאֵ֨ל אֶת־יְמִינ֜וֹ וַיָּ֨שֶׁת עַל־רֹ֤אשׁ אֶפְרַ֙יִם֙ וְה֣וּא הַצָּעִ֔יר וְאֶת־שְׂמֹאל֖וֹ עַל־רֹ֣אשׁ מְנַשֶּׁ֑ה שִׂכֵּל֙ אֶת־יָדָ֔יו כִּ֥י מְנַשֶּׁ֖ה הַבְּכֽוֹר׃ \Rashi{\textbf{שכל את ידיו.}כתרגומו, אחכימינון – בהשכל וחכמה השכיל את ידיו לכך ומדעת, כי יודע היה כי מנשה הבכור ואע"פ כן לא שת ימינו עליו:}%
% ---------- פרק 48 | פסוק 15 ----------
 \textbf{טו} וַיְבָ֥רֶךְ אֶת־יוֹסֵ֖ף וַיֹּאמַ֑ר הָֽאֱלֹהִ֡ים אֲשֶׁר֩ הִתְהַלְּכ֨וּ אֲבֹתַ֤י לְפָנָיו֙ אַבְרָהָ֣ם וְיִצְחָ֔ק הָֽאֱלֹהִים֙ הָרֹעֶ֣ה אֹתִ֔י מֵעוֹדִ֖י עַד־הַיּ֥וֹם הַזֶּֽה׃ %
% ---------- פרק 48 | פסוק 16 ----------
 \textbf{טז} הַמַּלְאָךְ֩ הַגֹּאֵ֨ל אֹתִ֜י מִכׇּל־רָ֗ע יְבָרֵךְ֮ אֶת־הַנְּעָרִים֒ וְיִקָּרֵ֤א בָהֶם֙ שְׁמִ֔י וְשֵׁ֥ם אֲבֹתַ֖י אַבְרָהָ֣ם וְיִצְחָ֑ק וְיִדְגּ֥וּ לָרֹ֖ב בְּקֶ֥רֶב הָאָֽרֶץ׃ \Rashi{\textbf{המלאך הגאל אתי.}מלאך הרגיל להשתלח אלי בצרתי; כענין שנאמר ויאמר אלי מלאך האלהים בחלום יעקב וגו' אנכי האל בית אל (בראשית ל"א): יברך את הנערים. מנשה ואפרים: וידגו. כדגים הללו שפרים ורבים ואין עין הרע שולטת בהם:}%
% ---------- פרק 48 | פסוק 17 ----------
 \textbf{יז} וַיַּ֣רְא יוֹסֵ֗ף כִּי־יָשִׁ֨ית אָבִ֧יו יַד־יְמִינ֛וֹ עַל־רֹ֥אשׁ אֶפְרַ֖יִם וַיֵּ֣רַע בְּעֵינָ֑יו וַיִּתְמֹ֣ךְ יַד־אָבִ֗יו לְהָסִ֥יר אֹתָ֛הּ מֵעַ֥ל רֹאשׁ־אֶפְרַ֖יִם עַל־רֹ֥אשׁ מְנַשֶּֽׁה׃ \Rashi{\textbf{ויתמך יד אביו.}הרימה מעל ראש בנו ותמכה בידו:}%
% ---------- פרק 48 | פסוק 18 ----------
 \textbf{יח} וַיֹּ֧אמֶר יוֹסֵ֛ף אֶל־אָבִ֖יו לֹא־כֵ֣ן אָבִ֑י כִּי־זֶ֣ה הַבְּכֹ֔ר שִׂ֥ים יְמִינְךָ֖ עַל־רֹאשֽׁוֹ׃ %
% ---------- פרק 48 | פסוק 19 ----------
 \textbf{יט} וַיְמָאֵ֣ן אָבִ֗יו וַיֹּ֙אמֶר֙ יָדַ֤עְתִּֽי בְנִי֙ יָדַ֔עְתִּי גַּם־ה֥וּא יִֽהְיֶה־לְּעָ֖ם וְגַם־ה֣וּא יִגְדָּ֑ל וְאוּלָ֗ם אָחִ֤יו הַקָּטֹן֙ יִגְדַּ֣ל מִמֶּ֔נּוּ וְזַרְע֖וֹ יִהְיֶ֥ה מְלֹֽא־הַגּוֹיִֽם׃ \Rashi{\textbf{ידעתי בני ידעתי.}שהוא הבכור: גם הוא יהיה לעם וגם הוא יגדל. שעתיד גדעון לצאת ממנו, שהקב"ה עושה נס על ידו: ואולם אחיו הקטן יגדל ממנו. שעתיד יהושע לצאת ממנו, שינחיל את הארץ וילמד תורה לישראל: וזרעו יהיה מלא הגוים. כל העולם יתמלא בצאת שמעו ושמו כשיעמיד חמה בגבעון וירח בעמק אילון:}%
% ---------- פרק 48 | פסוק 20 ----------
 \textbf{כ} וַיְבָ֨רְכֵ֜ם בַּיּ֣וֹם הַהוּא֮ לֵאמוֹר֒ בְּךָ֗ יְבָרֵ֤ךְ יִשְׂרָאֵל֙ לֵאמֹ֔ר יְשִֽׂמְךָ֣ אֱלֹהִ֔ים כְּאֶפְרַ֖יִם וְכִמְנַשֶּׁ֑ה וַיָּ֥שֶׂם אֶת־אֶפְרַ֖יִם לִפְנֵ֥י מְנַשֶּֽׁה׃ \Rashi{\textbf{בך יברך ישראל.}הבא לברך את בניו יברכם בברכתם ויאמר איש לבנו ישימך אלהים כאפרים וכמנשה: וישם את אפרים. בברכתו לפני מנשה, להקדימו בדגלים ובחנכת הנשיאים:}%
% ---------- פרק 48 | פסוק 21 ----------
 \textbf{כא} וַיֹּ֤אמֶר יִשְׂרָאֵל֙ אֶל־יוֹסֵ֔ף הִנֵּ֥ה אָנֹכִ֖י מֵ֑ת וְהָיָ֤ה אֱלֹהִים֙ עִמָּכֶ֔ם וְהֵשִׁ֣יב אֶתְכֶ֔ם אֶל־אֶ֖רֶץ אֲבֹתֵיכֶֽם׃ %
% ---------- פרק 48 | פסוק 22 ----------
 \textbf{כב} וַאֲנִ֞י נָתַ֧תִּֽי לְךָ֛ שְׁכֶ֥ם אַחַ֖ד עַל־אַחֶ֑יךָ אֲשֶׁ֤ר לָקַ֙חְתִּי֙ מִיַּ֣ד הָֽאֱמֹרִ֔י בְּחַרְבִּ֖י וּבְקַשְׁתִּֽי׃ \{פ\} \Rashi{\textbf{ואני נתתי לך.}לפי שאתה טורח להתעסק בקבורתי, וגם אני נתתי לך נחלה שתקבר בה, ואי זו? זו שכם, שנאמר, ואת עצמות יוסף אשר העלו ממצרים קברו בשכם (יהושע כ"ד): שכם אחד על אחיך. שכם ממש, היא תהיה לך חלק אחד יתרה על אחיך: בחרבי ובקשתי. כשהרגו שמעון ולוי את אנשי שכם נתכנסו כל סביבותיהם להזדוג להם, וחגר יעקב כלי מלחמה כנגדן; דבר אחר שכם אחד הוא הבכורה, שיטלו בניו שני חלקים, ושכם לשון חלק הוא, והרבה יש לו דומים במקרא כי תשיתמו שכם (תהילים כ"א) – תשית שונאי לפני לחלקים; אחלקה שכם (שם ס'), דרך ירצחו שכמה (הושע ו') – איש חלקו; לעבדו שכם אחד (צפניה ג'): אשר לקחתי מיד האמרי. מיד עשו, שעשה מעשה אמורי. דבר אחר שהיה צד אביו באמרי פיו: בחרבי ובקשתי. היא חכמתו ותפלתו:}%
% ---------- פרק 49 | פסוק 1 ----------
 \textbf{א} וַיִּקְרָ֥א יַעֲקֹ֖ב אֶל־בָּנָ֑יו וַיֹּ֗אמֶר הֵאָֽסְפוּ֙ וְאַגִּ֣ידָה לָכֶ֔ם אֵ֛ת אֲשֶׁר־יִקְרָ֥א אֶתְכֶ֖ם בְּאַחֲרִ֥ית הַיָּמִֽים׃ \Rashi{\textbf{ואגידה לכם.}בקש לגלות את הקץ ונסתלקה ממנו שכינה והתחיל אומר דברים אחרים:}%
% ---------- פרק 49 | פסוק 2 ----------
 \textbf{ב} הִקָּבְצ֥וּ וְשִׁמְע֖וּ בְּנֵ֣י יַעֲקֹ֑ב וְשִׁמְע֖וּ אֶל־יִשְׂרָאֵ֥ל אֲבִיכֶֽם׃ %
% ---------- פרק 49 | פסוק 3 ----------
 \textbf{ג} רְאוּבֵן֙ בְּכֹ֣רִי אַ֔תָּה כֹּחִ֖י וְרֵאשִׁ֣ית אוֹנִ֑י יֶ֥תֶר שְׂאֵ֖ת וְיֶ֥תֶר עָֽז׃ \Rashi{\textbf{וראשית אוני.}היא טפה ראשונה, שלא ראה קרי מימיו (בראשית רבה): אוני. כוחי, כמו מצאתי און לי (הושע י"ב), מרב אונים, ולאין אונים (ישעיהו מ'): יתר שאת. ראוי היית להיות יתר על אחיך בכהנה, לשון נשיאות כפים: ויתר עז. במלכות, כמו: ויתן עז למלכו (שמואל א ב'), ומי גרם לך להפסיד כל אלה?}%
% ---------- פרק 49 | פסוק 4 ----------
 \textbf{ד} פַּ֤חַז כַּמַּ֙יִם֙ אַל־תּוֹתַ֔ר כִּ֥י עָלִ֖יתָ מִשְׁכְּבֵ֣י אָבִ֑יךָ אָ֥ז חִלַּ֖לְתָּ יְצוּעִ֥י עָלָֽה׃ \{פ\} \Rashi{\textbf{פחז כמים.}הפחז והבהלה אשר מהרת להראות כעסך כמים הללו הממהרים למרוצתם, לכך אל תותר. אל תרבה לטל כל היתרות הללו שהיו ראויות לך, ומהו הפחז אשר פחזת? כי עלית משכבי אביך: אז חללת. אותו שעלה על יצועי – והיא שכינה שהיה דרכה להיות עולה על יצועי: פחז שם דבר הוא, לפיכך טעמו למעלה, וכלו נקוד פתח, ואלו היה לשון עבר, היה נקוד חציו קמץ וחציו פתח וטעמו למטה: יצועי. לשון משכב, על שם שמציעים אותו על ידי לבדין וסדינין; והרבה דומים לו אם אעלה על ערש יצועי (תהילים קל"ב); אם זכרתיך על יצועי (שם ס"ג):}%
% ---------- פרק 49 | פסוק 5 ----------
 \textbf{ה} שִׁמְע֥וֹן וְלֵוִ֖י אַחִ֑ים כְּלֵ֥י חָמָ֖ס מְכֵרֹתֵיהֶֽם׃ \Rashi{\textbf{שמעון ולוי אחים.}בעצה אחת על שכם ועל יוסף; ויאמרו איש אל אחיו, ועתה לכו ונהרגהו, מי הם? א"ת ראובן או יהודה, הרי לא הסכימו בהריגתו; א"ת בני השפחות, הרי לא היתה שנאתן שלמה, שנאמר והוא נער את בני בלהה ואת בני זלפה וגו', יששכר וזבולן לא היו מדברים בפני אחיהם הגדולים מהם; על כרחך שמעון ולוי הם שקראם אביהם אחים: כלי חמס. אמנות זו של רציחה, חמס הוא בידיהם – מברכת עשו היא זו, אמנות שלו היא – ואתם חמסתם אותה הימנו: מכרתיהם. לשון כלי זין, הסיף בלשון יוני מכי"ר, תנחומא. דבר אחר מכרתיהם – בארץ מגורתם נהגו עצמן בכלי חמס, כמו מכרתיך ומלדתיך (יחזקאל טז), וזהו תרגום של אנקלוס:}%
% ---------- פרק 49 | פסוק 6 ----------
 \textbf{ו} בְּסֹדָם֙ אַל־תָּבֹ֣א נַפְשִׁ֔י בִּקְהָלָ֖ם אַל־תֵּחַ֣ד כְּבֹדִ֑י כִּ֤י בְאַפָּם֙ הָ֣רְגוּ אִ֔ישׁ וּבִרְצֹנָ֖ם עִקְּרוּ־שֽׁוֹר׃ \Rashi{\textbf{בסדם אל תבא נפשי.}זה מעשה זמרי, כשנתקבצו שבטו של שמעון להביא את המדינית לפני משה ואמרו לו, זו אסורה או מתרת? אם תאמר אסורה, בת יתרו מי התירה לך? – אל יזכר שמי בדבר, זמרי בן סלוא נשיא בית אב לשמעוני ולא כתב בן יעקב: בקהלם. כשיקהיל קרח, שהוא משבטו של לוי, את כל העדה על משה ועל אהרן, אל תחד כבודי. שם, אל יתיחד עמהם שמי, שנאמר קרח בן יצהר בן קהת בן לוי, ולא נאמר בן יעקב; אבל בדברי הימים, כשנתיחסו בני קרח על הדוכן, נאמר בן קרח בן יצהר בן קהת בן לוי בן ישראל (דברי הימים א ו'): אל תחד כבודי. כבוד לשון זכר הוא, ועל כרחך אתה צריך לפרש כמדבר אל הכבוד ואומר אתה, כבודי אל תתיחד עמהם כמו לא תחד אתם בקבורה (ישעיה י"ד): כי באפם הרגו איש. אלו חמור ואנשי שכם, ואינן חשובין כלם אלא כאיש אחד; וכן הוא אומר בגדעון, והכית את מדין כאיש אחד (שופטים ו') וכן במצרים סוס ורכבו רמה בים, זהו מדרשו. ופשוטו, אנשים הרבה קורא איש – כל אחד לעצמו, באפם הרגו כל איש שכעסו עליו, וכן וילמד לטרף טרף אדם אכל (יחזקאל י"ט): וברצנם עקרו שור. רצו לעקר את יוסף שנקרא שור, שנאמר בכור שורו הדר לו (דברים ל"ג). עקרו אשרי"טיר בלעז, לשון את סוסיהם תעקר (יהושע י"א):}%
% ---------- פרק 49 | פסוק 7 ----------
 \textbf{ז} אָר֤וּר אַפָּם֙ כִּ֣י עָ֔ז וְעֶבְרָתָ֖ם כִּ֣י קָשָׁ֑תָה אֲחַלְּקֵ֣ם בְּיַעֲקֹ֔ב וַאֲפִיצֵ֖ם בְּיִשְׂרָאֵֽל׃ \{פ\} \Rashi{\textbf{ארור אפם כי עז.}אפלו בשעת תוכחה לא קלל אלא אפם, וזהו שאמר בלעם מה אקב לא קבה אל (במדבר כ"ג): אחלקם ביעקב. אפרידם זה מזה שלא יהא לוי במנין השבטים, והרי הם חלוקים. דבר אחר אין לך עניים וסופרים ומלמדי תינוקות אלא משמעון, כדי שיהיו נפוצים, ושבטו של לוי עשאו מחזר על הגרנות לתרומות ולמעשרות, נתן לו תפוצתו דרך כבוד:}%
% ---------- פרק 49 | פסוק 8 ----------
 \textbf{ח} יְהוּדָ֗ה אַתָּה֙ יוֹד֣וּךָ אַחֶ֔יךָ יָדְךָ֖ בְּעֹ֣רֶף אֹיְבֶ֑יךָ יִשְׁתַּחֲו֥וּ לְךָ֖ בְּנֵ֥י אָבִֽיךָ׃ \Rashi{\textbf{יהודה אתה יודוך אחיך.}לפי שהוכיח את הראשונים בקנטורים, התחיל יהודה נסוג לאחוריו, (שלא יוכיחנו על מעשה תמר) וקראו יעקב בדברי רצוי, יהודה לא אתה כמותם (בראשית רבה): ידך בערף איביך. בימי דוד – ואיבי תתה לי ערף (שמואל ב כ"ב): בני אביך. על שם שהיו מנשים הרבה, לא אמר בני אמך כדרך שאמר יצחק:}%
% ---------- פרק 49 | פסוק 9 ----------
 \textbf{ט} גּ֤וּר אַרְיֵה֙ יְהוּדָ֔ה מִטֶּ֖רֶף בְּנִ֣י עָלִ֑יתָ כָּרַ֨ע רָבַ֧ץ כְּאַרְיֵ֛ה וּכְלָבִ֖יא מִ֥י יְקִימֶֽנּוּ׃ \Rashi{\textbf{גור אריה.}על דוד נתנבא – בתחלה גור בהיות שאול מלך עלינו אתה היית המוציא והמביא את ישראל (שמואל ב ה') – ולבסוף אריה, כשהמליכוהו עליהם, וזהו שתרגם אנקלוס שלטון יהא בשרויא – בתחלתו: מטרף. ממה שחשדתיך בטרף טרף יוסף חיה רעה אכלתהו – וזהו יהודה שנמשל לאריה – בני עלית. סלקת את עצמך ואמרת מה בצע וגו', וכן בהריגת תמר שהודה, צדקה ממני, לפיכך כרע רבץ וגו', בימי שלמה איש תחת גפנו וגו' (מלכים א ה'):}%
% ---------- פרק 49 | פסוק 10 ----------
 \textbf{י} לֹֽא־יָס֥וּר שֵׁ֙בֶט֙ מִֽיהוּדָ֔ה וּמְחֹקֵ֖ק מִבֵּ֣ין רַגְלָ֑יו עַ֚ד כִּֽי־יָבֹ֣א שִׁילֹ֔ה וְל֖וֹ יִקְּהַ֥ת עַמִּֽים׃ \Rashi{\textbf{לא יסור שבט מיהודה.}מדוד ואילך, אלו ראשי גליות שבבבל, שרודים את העם בשבט שממנים על פי המלכות (בראשית רבה): ומחקק מבין רגליו. התלמידים, אלו נשיאי ארץ ישראל: עד כי יבא שילה. מלך המשיח שהמלוכה שלו, וכן תרגמו אנקלוס. ומדרש אגדה שילו – שי לו, שנאמר, יבילו שי למורא (תהילים ע"ו): ולו יקהת עמים. אספת העמים, שהיו"ד עקר היא ביסוד, כמו יפעתך (יחזקאל כ"ח), ופעמים שנופלת ממנו, וכמה אותיות משמשות בלשון זה, והם נקראים עקר נופל; כגון נו"ן של נוגף ושל נושך, ואל"ף שבואחותי באזניכם (איוב י"ג), ושבאבחת חרב (יחזקאל כ"א), ואסוך שמן (מלכים ב ד'), אף זה יקהת עמים אספת עמים, שנאמר אליו גוים ידרשו (ישעיהו י"א). ודומה לו עין תלעג לאב ותבז ליקהת אם (משלי ל'), לקבוץ קמטים שבפניה מפני זקנתה; ובתלמוד דיתבי ומקהו אקהתא בשוקי דנהרדעא במסכת יבמות; ויכול היה לומר, קהית עמים:}%
% ---------- פרק 49 | פסוק 11 ----------
 \textbf{יא} אֹסְרִ֤י לַגֶּ֙פֶן֙ עִירֹ֔ה וְלַשֹּׂרֵקָ֖ה בְּנִ֣י אֲתֹנ֑וֹ כִּבֵּ֤ס בַּיַּ֙יִן֙ לְבֻשׁ֔וֹ וּבְדַם־עֲנָבִ֖ים סוּתֹֽה׃ \Rashi{\textbf{אסרי לגפן עירה.}נתנבא על ארץ יהודה שתהא מושכת יין כמעין; איש יהודה יאסר לגפן עיר אחד ויטעננו מגפן אחת, ומשרק אחד בן אתון אחד: שרקה. זמורה ארכה, קוריירא בלעז: כבס ביין. כל זה לשון רבוי יין: סותה. לשון מין בגד הוא, ואין לו דמיון במקרא: אסרי. כמו אוסר, דגמת מקימי מעפר דל (תהילים קי"ג), הישבי בשמים (שם קכ"ג), וכן בני אתנו כענין זה. ואנקלוס תרגם במלך המשיח: גפן – הם ישראל; עירה – זו ירושלים; שרקה – ישראל, ואנכי נטעתיך שורק – (ירמיהו ב'): בני אתנו. יבנון היכלה, לשון שער האיתון בספר יחזקאל (יחזקאל מ'). ועוד תרגמו בפנים אחרים: גפן אלו צדיקים; בני אתנו – עבדי אוריתא באולפן, על שם רכבי אתנות צחרות: כבס ביין. יהא ארגון טב שצבועו דומה ליין, וצבעונין הוא לשון סותה, שהאשה לובשתן ומסיתה בהן את הזכר לתן עיניו בה; ואף רבותינו פרשו בתלמוד לשון הסתת שכרות במסכת כתבות (דף קי"א). ועל היין שמא תאמר אינו מרוה, תלמוד לומר סותה:}%
% ---------- פרק 49 | פסוק 12 ----------
 \textbf{יב} חַכְלִילִ֥י עֵינַ֖יִם מִיָּ֑יִן וּלְבֶן־שִׁנַּ֖יִם מֵחָלָֽב׃ \{פ\} \Rashi{\textbf{חכלילי.}לשון אדם, כתרגומו, וכן למי חכללות עינים (משלי כ"ג), שכן דרך שותי יין עיניהם מאדימין: מחלב. מרב חלב, שיהא בארצו מרעה טוב לעדרי צאן; וכן פרוש המקרא: אדם עינים יהא מרב יין ולבן שנים יהא מרב חלב. ולפי תרגומו עינים לשון הרים, שמשם צופים למרחוק, ועוד תרגמו בפנים אחרים, לשון מעינות וקלוח היקבים, נעווהי – יקבים שלו, ולשון ארמי הוא במסכת ע"ז (דף ע"ד), נעוא ארתחו יחורן בקעתה – תרגום שנים, לשון שני הסלעים:}%
% ---------- פרק 49 | פסוק 13 ----------
 \textbf{יג} זְבוּלֻ֕ן לְח֥וֹף יַמִּ֖ים יִשְׁכֹּ֑ן וְהוּא֙ לְח֣וֹף אֳנִיֹּ֔ת וְיַרְכָת֖וֹ עַל־צִידֹֽן׃ \{פ\} \Rashi{\textbf{לחוף ימים.}על חוף ימים תהיה ארצו (יבמות מ"ה); חוף, כתרגומו ספר, מרק"א בלעז, והוא יהיה מצוי תדיר על חוף אניות במקום הנמל, שאניות מביאות שם פרקמטיא, שהיה זבולן עוסק בפרקמטיא וממציא מזון לשבט יששכר, והם עוסקים בתורה. הוא שאמר משה שמח זבולן בצאתך ויששכר באהליך (דברים ל"ג), זבולן יוצא בפרקמטיא, ויששכר עוסק בתורה באהלים: וירכתו על צידן. סוף גבולו יהא סמוך לצידן: ירכתו. סופו, כמו ולירכתי המשכן:}%
% ---------- פרק 49 | פסוק 14 ----------
 \textbf{יד} יִשָּׂשכָ֖ר חֲמֹ֣ר גָּ֑רֶם רֹבֵ֖ץ בֵּ֥ין הַֽמִּשְׁפְּתָֽיִם׃ \Rashi{\textbf{יששכר חמר גרם.}חמור בעל עצמות, סובל על תורה, כחמור חזק שמטעינין אותו משאוי כבד: רובץ בין המשפתים. כחמור המהלך ביום ובלילה, ואין לו לינה בבית, וכשהוא רוצה לנוח, רובץ בין התחומין, בתחומי העירות שמוליך שם פרקמטיא:}%
% ---------- פרק 49 | פסוק 15 ----------
 \textbf{טו} וַיַּ֤רְא מְנֻחָה֙ כִּ֣י ט֔וֹב וְאֶת־הָאָ֖רֶץ כִּ֣י נָעֵ֑מָה וַיֵּ֤ט שִׁכְמוֹ֙ לִסְבֹּ֔ל וַיְהִ֖י לְמַס־עֹבֵֽד׃ \{ס\} \Rashi{\textbf{וירא מנחה כי טוב.}ראה לחלקו ארץ מברכת וטובה להוציא פרות: ויט שכמו לסבול על תורה: ויהי. לכל אחיו ישראל: למס עובד. לפסק להם הוראות של תורה וסדרי עבורין, שנאמר ומבני יששכר יודעי בינה לעתים לדעת מה יעשה ישראל ראשיהם מאתים (דברי הימים א י"ב), מאתים ראשי סנהדראות העמיד וכל אחיהם על פיהם: ויט שכמו. השפיל שכמו, כמו ויט שמים (שמואל ב כ"ב), הטו אזנכם (תהילים ע"ח). ואנקלוס תרגמו בפנים אחרים, ויט שכמו לסבול מלחמות ולכבש מחוזות, שהם יושבים על הספר, ויהי האויב כבוש תחתיו למס עובד:}%
% ---------- פרק 49 | פסוק 16 ----------
 \textbf{טז} דָּ֖ן יָדִ֣ין עַמּ֑וֹ כְּאַחַ֖ד שִׁבְטֵ֥י יִשְׂרָאֵֽל׃ \Rashi{\textbf{דן ידין עמו.}ינקם נקמת עמו מפלשתים כמו כי ידין ה' עמו (דברים ל"ב): כאחד שבטי ישראל. כל ישראל יהיו כאחד עמו, ואת כלם ידין, ועל שמשון נבא נבואה זו; ועוד יש לפרש כאחד שבטי ישראל, כמיחד שבשבטים, הוא דוד, שבא מיהודה:}%
% ---------- פרק 49 | פסוק 17 ----------
 \textbf{יז} יְהִי־דָן֙ נָחָ֣שׁ עֲלֵי־דֶ֔רֶךְ שְׁפִיפֹ֖ן עֲלֵי־אֹ֑רַח הַנֹּשֵׁךְ֙ עִקְּבֵי־ס֔וּס וַיִּפֹּ֥ל רֹכְב֖וֹ אָחֽוֹר׃ \Rashi{\textbf{שפיפן.}הוא נחש; ואומר אני שקרוי כן על שם שהוא נושף, כמו ואתה תשופנו עקב (בראשית ג'): הנשך עקבי סוס. כך דרכו של נחש, ודמהו לנחש הנושך עקבי סוס, ויפל רכבו אחור – שלא נגע בו. ודגמתו מצינו בשמשון: וילפת שמשון את שני עמודי התוך וגומר (שופטים ט"ז) – ושעל הגג מתו; ואנקלוס תרגם כחוי חורמן, שם מין נחש, שאין רפואה לנשיכתו, והוא צפעוני, וקרוי חורמן על שם שעושה הכל חרם; וכפיתנא – כמו פתן, יכמון – יארב:}%
% ---------- פרק 49 | פסוק 18 ----------
 \textbf{יח} לִישׁוּעָֽתְךָ֖ קִוִּ֥יתִי יְהֹוָֽה׃ \{ס\} \Rashi{\textbf{לישועתך קויתי ה'.}נתנבא שינקרו פלשתים את עיניו, וסופו לומר זכרני נא וחזקני נא אך הפעם וגו' (שופטים ט"ז):}%
% ---------- פרק 49 | פסוק 19 ----------
 \textbf{יט} גָּ֖ד גְּד֣וּד יְגוּדֶ֑נּוּ וְה֖וּא יָגֻ֥ד עָקֵֽב׃ \{ס\} \Rashi{\textbf{גד גדוד יגודנו.}כלם לשון גדוד הם, וכך חברו מנחם; ואם תאמר, אין גדוד בלא שני דלתי"ן, יש לומר גדוד שם דבר צריך שני דלתי"ן, שכן דרך תבה בת שתי אותיות לכפל בסופה, ואין יסודה אלא שתי אותיות, וכן אמר כצפור לנוד (משלי כ"ו), מגזרת שבעתי נדדים (איוב ז'), שם נפל שדוד (שופטים ה'), מגזרת ישוד צהרים (תהלים צ"א), אף יגד, יגודנו וגדוד מגזרה אחת הם; וכשהוא מדבר בלשון יפעל אינו כפול, כמו יגוד, ינוד, ירום, ישוד, ישוב, וכשהוא מתפעל או מפעיל אחרים הוא כפול, כמו יתגודד, יתרומם, יתבולל, יתעודד; ובלשון מפעיל, יתום ואלמנה יעודד (תהילים קמ"ו), לשובב יעקב אליו (ישעיה מ"ט), משבב נתיבות (שם נ"ח), יגודנו האמור כאן אינו לשון שיפעלוהו אחרים, אלא כמו יגוד הימנו, כמו בני יצאני (ירמיהו י') – יצאו ממני. גד גדוד יגודנו, גדודים יגודו הימנו שיעברו הירדן עם אחיהם למלחמה כל חלוץ עד שנכבשה הארץ: והוא יגוד עקב. כל גדודיו ישובו על עקבם לנחלתם שלקחו בעבר הירדן, ולא יפקד מהם איש: עקב. בדרכן ובמסלותם שהלכו, ישובו, כמו ועקבותיך לא נדעו (תהילים ע"ז), וכן בעקבי הצאן (שיר א'), בלשון לעז טרצי"אם:}%
% ---------- פרק 49 | פסוק 20 ----------
 \textbf{כ} מֵאָשֵׁ֖ר שְׁמֵנָ֣ה לַחְמ֑וֹ וְה֥וּא יִתֵּ֖ן מַֽעֲדַנֵּי־מֶֽלֶךְ׃ \{ס\} \Rashi{\textbf{מאשר שמנה לחמו.}מאכל הבא מחלקו של אשר יהא שמן, שיהיו זיתים מרבים בחלקו והוא מושך שמן כמעין; וכן ברכו משה וטבל בשמן רגלו, כמו ששנינו במנחות (דף פ"ה) פעם אחת הצרכו אנשי לודקיא לשמן וכו':}%
% ---------- פרק 49 | פסוק 21 ----------
 \textbf{כא} נַפְתָּלִ֖י אַיָּלָ֣ה שְׁלֻחָ֑ה הַנֹּתֵ֖ן אִמְרֵי־שָֽׁפֶר׃ \{ס\} \Rashi{\textbf{אילה שלחה.}זו בקעת גינוסר שהיא קלה לבשל פרותיה כאילה זו שהיא קלה לרוץ. אילה שלחה – אילה משלחת לרוץ: הנתן אמרי שפר. כתרגומו. ד"א על מלחמת סיסרא נתנבא – ולקחת עמך עשרת אלפים איש מבני נפתלי וגו' (שופטים ד'), והלכו שם בזריזות, וכן נאמר שם לשון שלוח בעמק שלח ברגליו: הנתן אמרי שפר. על ידם שרו דבורה וברק שירה. ורבותינו דרשוהו על יום קבורת יעקב כשערער עשו על המערה, במסכת סוטה (דף י"ג); ותרגומו יתרמי עדבה – יפל חבלו, והוא יודה על חלקו אמרים נאים ושבח:}%
% ---------- פרק 49 | פסוק 22 ----------
 \textbf{כב} בֵּ֤ן פֹּרָת֙ יוֹסֵ֔ף בֵּ֥ן פֹּרָ֖ת עֲלֵי־עָ֑יִן בָּנ֕וֹת צָעֲדָ֖ה עֲלֵי־שֽׁוּר׃ \Rashi{\textbf{בן פרת.}בן חן, והוא לשון ארמי, אפרין נמטיה לרבי שמעון, בסוף בבא מציעא (דף קי"ט): בן פרת עלי עין. חנו נטוי על העין הרואה אותו: בנת צעדה עלי שור. בנות מצרים היו צועדות (על החומה) להסתכל ביפיו, בנות הרבה, צעדה כל אחת ואחת, במקום שתוכל לראותו משם: עלי שור. על ראיתו, כמו אשורנו ולא קרוב (במדבר כ"ד), ומ"א יש רבים, וזה נוטה לישוב המקרא. פרת תי"ו שבו הוא תקון הלשון, כמו על דברת בני האדם (קהלת ג'): שור. כמו לשור: עלי שור. בשביל לשור, ותרגום של אנקלוס, בנות צעדה עלי שור תרין שבטין יפקון מבנוהי וכו' וכתב בנות על שם בנות מנשה, בנות צלפחד, שנטלו חלק בשני עברי הירדן, ברי דיסגי יוסף, פרת לשון פריה ורביה; ויש מ"א בו המתישבים על הלשון, בשעה שבא עשו לקראת יעקב, בכלן קדמו האמהות ללכת לפני בניהם להשתחוות, וברחל כתיב נגש יוסף ורחל וישתחוו, אמר יוסף רשע הזה עינו רמה, שמא יתן עיניו באמי, יצא לפניה ושרבב קומתו לכסותה, והוא שברכו אביו בן פרת – הגדלת עצמך יוסף עלי עין של עשו, לפיכך זכית לגדלה: בנות צעדה עלי שור. להסתכל בך בצאתך על מצרים, ועוד דרשוהו לענין שלא ישלט בזרעו עין הרע; ואף כשברך מנשה ואפרים ברכם כדגים, שאין עין הרע שולטת בהם:}%
% ---------- פרק 49 | פסוק 23 ----------
 \textbf{כג} וַֽיְמָרְרֻ֖הוּ וָרֹ֑בּוּ וַֽיִּשְׂטְמֻ֖הוּ בַּעֲלֵ֥י חִצִּֽים׃ \Rashi{\textbf{וימררהו ורבו.}וימררוהו אחיו, וימררוהו פוטיפר ואשתו, לאסרו; לשון וימררו את חייהם: ורבו. נעשו לו אחיו אנשי ריב. ואין הלשון הזה לשון פעלו, שאם כן היה לו לנקד "ורבו", כמו המה מי מריבה אשר רבו וגו' (במדבר כ'), ואף אם לשון רבית חצים הוא כך היה לו לנקד; אינו אלא לשון פועלו, כמו שמו שמים (ירמיהו ב'), שהוא לשון הושמו; וכן רומו מעט, שהוא לשון הורמו, אלא שלשון הורמו והושמו על ידי אחרים, ולשון שמו, רמו, רבו מאליהם הוא – משוממים את עצמם, נתרוממו מעצמם, נעשו אנשי ריב, וכן דמו ישבי אי, כמו נדמו, וכן תרגם אנקלוס, ונקמוהי: בעלי חצים. שלשונם כחץ, ותרגומו לשון ותהי המחצה, אותן שהיו ראויים לחלק עמו נחלה:}%
% ---------- פרק 49 | פסוק 24 ----------
 \textbf{כד} וַתֵּ֤שֶׁב בְּאֵיתָן֙ קַשְׁתּ֔וֹ וַיָּפֹ֖זּוּ זְרֹעֵ֣י יָדָ֑יו מִידֵי֙ אֲבִ֣יר יַעֲקֹ֔ב מִשָּׁ֥ם רֹעֶ֖ה אֶ֥בֶן יִשְׂרָאֵֽל׃ \Rashi{\textbf{ותשב באיתן קשתו.}נתישבה בחזק: קשתו. חזקו: ויפזו זרעי ידיו. זו היא נתינת טבעת על ידו, לשון זהב מופז, זאת היתה לו מידי הקב"ה שהוא אביר יעקב, ומשם עלה להיות רעה אבן ישראל – עקרן של ישראל, לשון האבן הראשה, לשון מלכות; ואנקלוס אף הוא כך תרגמו ותשב – ותבת בהון נביאותה, החלומות אשר חלם להם – על דקים אוריתא בסתרא, תוספת הוא, ולא מלשון עברי שבמקרא, – ושוי תוקפא רוחצנה תרגום של באיתן קשתו, וכך לשון התרגום על העברי: ותשב נבואתו בשביל שאיתנו של הקב"ה היתה לו לקשת ולמבטח, בכן יתרמא דהב על דרעוהי – לכך ויפזו זרעי ידיו, לשון פז: אבן ישראל. לשון נוטריקון אב ובן, אבהן ובנין – יעקב ובניו. ורבותינו דרשו ותשב באיתן קשתו על כבישת יצרו באשת אדניו, וקוראו קשת, על שם שהזרע יורה כחץ: ויפזו זרעי ידיו. כמו ויפוצו, שיצא הזרע מבין אצבעות ידיו: מידי אביר יעקב. שנראתה לו דמות דיוקנו של אביו וכו', כדאיתא בסוטה (דף ל"ו, ב'):}%
% ---------- פרק 49 | פסוק 25 ----------
 \textbf{כה} מֵאֵ֨ל אָבִ֜יךָ וְיַעְזְרֶ֗ךָּ וְאֵ֤ת שַׁדַּי֙ וִיבָ֣רְכֶ֔ךָּ בִּרְכֹ֤ת שָׁמַ֙יִם֙ מֵעָ֔ל בִּרְכֹ֥ת תְּה֖וֹם רֹבֶ֣צֶת תָּ֑חַת בִּרְכֹ֥ת שָׁדַ֖יִם וָרָֽחַם׃ \Rashi{\textbf{מאל אביך.}היתה לך זאת והוא יעזרך: ואת שדי. ועם הקב"ה היה לבך כשלא שמעת לדברי אדונתך, והוא יברכך: ברכת שדים ורחם. ברכתא דאבא ודאמא, כלומר, יתברכו המולידים והיולדות, שיהיו הזכרים מזריעין טפה הראויה להריון, והנקבות לא ישכלו את רחם שלהן להפיל עבריהן: שדים. ירה יירה מתרגמינן אשתדאה ישתדי, אף שדים כאן על שם שהזרע יורה כחץ:}%
% ---------- פרק 49 | פסוק 26 ----------
 \textbf{כו} בִּרְכֹ֣ת אָבִ֗יךָ גָּֽבְרוּ֙ עַל־בִּרְכֹ֣ת הוֹרַ֔י עַֽד־תַּאֲוַ֖ת גִּבְעֹ֣ת עוֹלָ֑ם תִּֽהְיֶ֙יןָ֙ לְרֹ֣אשׁ יוֹסֵ֔ף וּלְקׇדְקֹ֖ד נְזִ֥יר אֶחָֽיו׃ \{פ\} \Rashi{\textbf{ברכת אביך גברו וגו'.}הברכות שברכני הקב"ה גברו והלכו על הברכות שברך את הורי: עד תאות גבעת עולם. לפי שהברכות שלי גברו עד סוף גבולי גבעות עולם, שנתן לי ברכה פרוצה בלי מצרים, מגעת עד ארבע קצות העולם, שנאמר ופרצת ימה וקדמה וגו', מה שלא אמר לאברהם אבינו וליצחק; לאברהם אמר שא נא עיניך וראה צפונה וגו' כי את כל הארץ אשר אתה ראה לך אתננה, ולא הראהו אלא ארץ ישראל בלבד; ליצחק אמר לך ולזרעך אתן את כל הארצת האל והקמתי את השבעה וגו', וזהו שאמר ישעיה (ישעיהו נ"ח) והאכלתיך נחלת יעקב אביך, ולא אמר נחלת אברהם: הורי. לשון הריון, שהורוני במעי אמי, כמו הרה גבר (איוב ג'): עד תאות. עד קצות, כמו והתאויתם לכם לגבול קדמה (במדבר ל"ד), תתאו לבא חמת (שם): תאות. אשמו"לץ בלעז, כך חברו מנחם בן סרוק. ואנקלוס תרגם תאות עולם לשון תאוה וחמדה, וגבעות לשון מצקי ארץ שחמדתן אמו והזקיקתו לקבלם: תהיין. כלם לראש יוסף: נזיר אחיו. פרישא דאחוהי, שנבדל מאחיו, כמו וינזרו מקדשי בני ישראל (ויקרא כ"ב), נזרו אחור (ישעיהו א'):}%
% ---------- פרק 49 | פסוק 27 ----------
 \textbf{כז} בִּנְיָמִין֙ זְאֵ֣ב יִטְרָ֔ף בַּבֹּ֖קֶר יֹ֣אכַל עַ֑ד וְלָעֶ֖רֶב יְחַלֵּ֥ק שָׁלָֽל׃ \Rashi{\textbf{בנימין זאב יטרף.}זאב הוא אשר יטרף; נבא על שיהיו עתידין להיות חטפנין, וחטפתם לכם איש אשתו בפלגש בגבעה (שופטים כ"א), ונבא על שאול שהיה נוצח באויביו סביב, שנאמר ושאול לכד המלוכה וילחם במואב וגו' ובאדום וגו' ובכל אשר יפנה ירשיע (שמואל א י"ד): בבקר יאכל עד. לשון בזה ושלל, המתרגם עדאה, ועוד יש לו דומה בלשון עברית אז חלק עד שלל (ישעיהו ל"ג), ועל שאול הוא אומר שעמד בתחלת פריחתן וזריחתן של ישראל: ולערב יחלק שלל. אף משתשקע שמשן של ישראל על ידי נבוכדנצר, שיגלם לבבל, יחלק שלל, מרדכי ואסתר שהם מבנימין יחלקו את שלל המן, שנאמר הנה בית המן נתתי לאסתר. ואנקלוס תרגם על שלל הכהנים בקדשי המקדש:}%
% ---------- פרק 49 | פסוק 28 ----------
 \textbf{כח} כׇּל־אֵ֛לֶּה שִׁבְטֵ֥י יִשְׂרָאֵ֖ל שְׁנֵ֣ים עָשָׂ֑ר וְ֠זֹ֠את אֲשֶׁר־דִּבֶּ֨ר לָהֶ֤ם אֲבִיהֶם֙ וַיְבָ֣רֶךְ אוֹתָ֔ם אִ֛ישׁ אֲשֶׁ֥ר כְּבִרְכָת֖וֹ בֵּרַ֥ךְ אֹתָֽם׃ \Rashi{\textbf{וזאת אשר דבר להם אביהם ויברך אותם.}והלא יש מהם שלא ברכם אלא קנטרן? אלא כך פרושו: וזאת אשר דבר להם אביהם מה שנאמר בענין; יכול שלא ברך לראובן, שמעון ולוי, ת"ל ויברך אותם – כלם במשמע: איש אשר כברכתו. ברכה העתידה לבא על כל אחד ואחד: ברך אותם. לא היה לו לומר אלא איש אשר כברכתו ברך אותו, מה תלמוד לומר ברך אותם? לפי שנתן ליהודה גבורת ארי ולבנימין חטיפתו של זאב ולנפתלי קלותו של איל, יכול שלא כללן כלם בכל הברכות, תלמוד לומר ברך אותם:}%
% ---------- פרק 49 | פסוק 29 ----------
 \textbf{כט} וַיְצַ֣ו אוֹתָ֗ם וַיֹּ֤אמֶר אֲלֵהֶם֙ אֲנִי֙ נֶאֱסָ֣ף אֶל־עַמִּ֔י קִבְר֥וּ אֹתִ֖י אֶל־אֲבֹתָ֑י אֶ֨ל־הַמְּעָרָ֔ה אֲשֶׁ֥ר בִּשְׂדֵ֖ה עֶפְר֥וֹן הַֽחִתִּֽי׃ \Rashi{\textbf{נאסף אל עמי.}על שם שמכניסין הנפשות אל מקום גניזתן, שיש אסיפה בלשון עברי שהיא לשון הכנסה, כגון ואין איש מאסף אותי הביתה (שופטים י"ט), ואספתו אל תוך ביתך (דברים כ"ב), באספכם את תבואת הארץ (ויקרא כ"ג), הכנסתם לבית מפני הגשמים, באספך את מעשיך (שמות כ"ג), וכל אסיפה האמורה במיתה אף היא לשון הכנסה: אל אבתי. עם אבותי:}%
% ---------- פרק 49 | פסוק 30 ----------
 \textbf{ל} בַּמְּעָרָ֞ה אֲשֶׁ֨ר בִּשְׂדֵ֧ה הַמַּכְפֵּלָ֛ה אֲשֶׁ֥ר עַל־פְּנֵי־מַמְרֵ֖א בְּאֶ֣רֶץ כְּנָ֑עַן אֲשֶׁר֩ קָנָ֨ה אַבְרָהָ֜ם אֶת־הַשָּׂדֶ֗ה מֵאֵ֛ת עֶפְרֹ֥ן הַחִתִּ֖י לַאֲחֻזַּת־קָֽבֶר׃ %
% ---------- פרק 49 | פסוק 31 ----------
 \textbf{לא} שָׁ֣מָּה קָֽבְר֞וּ אֶת־אַבְרָהָ֗ם וְאֵת֙ שָׂרָ֣ה אִשְׁתּ֔וֹ שָׁ֚מָּה קָבְר֣וּ אֶת־יִצְחָ֔ק וְאֵ֖ת רִבְקָ֣ה אִשְׁתּ֑וֹ וְשָׁ֥מָּה קָבַ֖רְתִּי אֶת־לֵאָֽה׃ %
% ---------- פרק 49 | פסוק 32 ----------
 \textbf{לב} מִקְנֵ֧ה הַשָּׂדֶ֛ה וְהַמְּעָרָ֥ה אֲשֶׁר־בּ֖וֹ מֵאֵ֥ת בְּנֵי־חֵֽת׃ %
% ---------- פרק 49 | פסוק 33 ----------
 \textbf{לג} וַיְכַ֤ל יַעֲקֹב֙ לְצַוֺּ֣ת אֶת־בָּנָ֔יו וַיֶּאֱסֹ֥ף רַגְלָ֖יו אֶל־הַמִּטָּ֑ה וַיִּגְוַ֖ע וַיֵּאָ֥סֶף אֶל־עַמָּֽיו׃ \Rashi{\textbf{ויאסף רגליו.}הכניס רגליו: ויגוע ויאסף. ומיתה לא נאמרה בו, וא"ר יעקב אבינו לא מת:}%
% ---------- פרק 50 | פסוק 1 ----------
 \textbf{א} וַיִּפֹּ֥ל יוֹסֵ֖ף עַל־פְּנֵ֣י אָבִ֑יו וַיֵּ֥בְךְּ עָלָ֖יו וַיִּשַּׁק־לֽוֹ׃ %
% ---------- פרק 50 | פסוק 2 ----------
 \textbf{ב} וַיְצַ֨ו יוֹסֵ֤ף אֶת־עֲבָדָיו֙ אֶת־הָרֹ֣פְאִ֔ים לַחֲנֹ֖ט אֶת־אָבִ֑יו וַיַּחַנְט֥וּ הָרֹפְאִ֖ים אֶת־יִשְׂרָאֵֽל׃ \Rashi{\textbf{לחנט את אביו.}ענין מרקחת בשמים הוא:}%
% ---------- פרק 50 | פסוק 3 ----------
 \textbf{ג} וַיִּמְלְאוּ־לוֹ֙ אַרְבָּעִ֣ים י֔וֹם כִּ֛י כֵּ֥ן יִמְלְא֖וּ יְמֵ֣י הַחֲנֻטִ֑ים וַיִּבְכּ֥וּ אֹת֛וֹ מִצְרַ֖יִם שִׁבְעִ֥ים יֽוֹם׃ \Rashi{\textbf{וימלאו לו.}השלימו לו ימי חניטתו עד שמלאו לו מ' יום: ויבכו אתו מצרים שבעים יום. מ' לחניטה ול' לבכיה, לפי שבאה להם ברכה לרגלו, שכלה הרעב והיו מי נילוס מתברכין:}%
% ---------- פרק 50 | פסוק 4 ----------
 \textbf{ד} וַיַּֽעַבְרוּ֙ יְמֵ֣י בְכִית֔וֹ וַיְדַבֵּ֣ר יוֹסֵ֔ף אֶל־בֵּ֥ית פַּרְעֹ֖ה לֵאמֹ֑ר אִם־נָ֨א מָצָ֤אתִי חֵן֙ בְּעֵ֣ינֵיכֶ֔ם דַּבְּרוּ־נָ֕א בְּאׇזְנֵ֥י פַרְעֹ֖ה לֵאמֹֽר׃ %
% ---------- פרק 50 | פסוק 5 ----------
 \textbf{ה} אָבִ֞י הִשְׁבִּיעַ֣נִי לֵאמֹ֗ר הִנֵּ֣ה אָנֹכִי֮ מֵת֒ בְּקִבְרִ֗י אֲשֶׁ֨ר כָּרִ֤יתִי לִי֙ בְּאֶ֣רֶץ כְּנַ֔עַן שָׁ֖מָּה תִּקְבְּרֵ֑נִי וְעַתָּ֗ה אֶֽעֱלֶה־נָּ֛א וְאֶקְבְּרָ֥ה אֶת־אָבִ֖י וְאָשֽׁוּבָה׃ \Rashi{\textbf{אשר כריתי לי.}כפשוטו כמו כי יכרה איש. ומדרשו עוד מתישב על הלשון, כמו אשר קניתי, אר"ע כשהלכתי לכרכי הים היו קורין למכירה כירה, ועוד מדרשו לשון כרי דגור, שנטל יעקב כל כסף וזהב שהביא מבית לבן ועשה אותו כרי ואמר לעשו טל זה בשביל חלקך במערה:}%
% ---------- פרק 50 | פסוק 6 ----------
 \textbf{ו} וַיֹּ֖אמֶר פַּרְעֹ֑ה עֲלֵ֛ה וּקְבֹ֥ר אֶת־אָבִ֖יךָ כַּאֲשֶׁ֥ר הִשְׁבִּיעֶֽךָ׃ \Rashi{\textbf{כאשר השביעך.}ואם לא בשביל השבועה לא הייתי מניחך, אבל ירא לומר עבר על השבועה שלא יאמר אם כן אעבר על השבועה שנשבעתי לך, שלא אגלה על לשון הקדש שאני מכיר עודף על ע' לשון, ואתה אינך מכיר בו, כדאיתא במסכת סוטה:}%
% ---------- פרק 50 | פסוק 7 ----------
 \textbf{ז} וַיַּ֥עַל יוֹסֵ֖ף לִקְבֹּ֣ר אֶת־אָבִ֑יו וַיַּֽעֲל֨וּ אִתּ֜וֹ כׇּל־עַבְדֵ֤י פַרְעֹה֙ זִקְנֵ֣י בֵית֔וֹ וְכֹ֖ל זִקְנֵ֥י אֶֽרֶץ־מִצְרָֽיִם׃ %
% ---------- פרק 50 | פסוק 8 ----------
 \textbf{ח} וְכֹל֙ בֵּ֣ית יוֹסֵ֔ף וְאֶחָ֖יו וּבֵ֣ית אָבִ֑יו רַ֗ק טַפָּם֙ וְצֹאנָ֣ם וּבְקָרָ֔ם עָזְב֖וּ בְּאֶ֥רֶץ גֹּֽשֶׁן׃ %
% ---------- פרק 50 | פסוק 9 ----------
 \textbf{ט} וַיַּ֣עַל עִמּ֔וֹ גַּם־רֶ֖כֶב גַּם־פָּרָשִׁ֑ים וַיְהִ֥י הַֽמַּחֲנֶ֖ה כָּבֵ֥ד מְאֹֽד׃ %
% ---------- פרק 50 | פסוק 10 ----------
 \textbf{י} וַיָּבֹ֜אוּ עַד־גֹּ֣רֶן הָאָטָ֗ד אֲשֶׁר֙ בְּעֵ֣בֶר הַיַּרְדֵּ֔ן וַיִּ֨סְפְּדוּ־שָׁ֔ם מִסְפֵּ֛ד גָּד֥וֹל וְכָבֵ֖ד מְאֹ֑ד וַיַּ֧עַשׂ לְאָבִ֛יו אֵ֖בֶל שִׁבְעַ֥ת יָמִֽים׃ \Rashi{\textbf{גרן האטד.}מקף אטדין היה, ורבותינו דרשו (סוטה י"ג), על שם המארע, שבאו כל מלכי כנען ונשיאי ישמעאל למלחמה, וכיון שראו כתרו של יוסף תלוי בארונו של יעקב עמדו כלן ותלו בו כתריהם והקיפוהו כתרים כגרן המקף סיג של קוצים:}%
% ---------- פרק 50 | פסוק 11 ----------
 \textbf{יא} וַיַּ֡רְא יוֹשֵׁב֩ הָאָ֨רֶץ הַֽכְּנַעֲנִ֜י אֶת־הָאֵ֗בֶל בְּגֹ֙רֶן֙ הָֽאָטָ֔ד וַיֹּ֣אמְר֔וּ אֵֽבֶל־כָּבֵ֥ד זֶ֖ה לְמִצְרָ֑יִם עַל־כֵּ֞ן קָרָ֤א שְׁמָהּ֙ אָבֵ֣ל מִצְרַ֔יִם אֲשֶׁ֖ר בְּעֵ֥בֶר הַיַּרְדֵּֽן׃ %
% ---------- פרק 50 | פסוק 12 ----------
 \textbf{יב} וַיַּעֲשׂ֥וּ בָנָ֖יו ל֑וֹ כֵּ֖ן כַּאֲשֶׁ֥ר צִוָּֽם׃ \Rashi{\textbf{כאשר צום.}מהו אשר צום?}%
% ---------- פרק 50 | פסוק 13 ----------
 \textbf{יג} וַיִּשְׂא֨וּ אֹת֤וֹ בָנָיו֙ אַ֣רְצָה כְּנַ֔עַן וַיִּקְבְּר֣וּ אֹת֔וֹ בִּמְעָרַ֖ת שְׂדֵ֣ה הַמַּכְפֵּלָ֑ה אֲשֶׁ֣ר קָנָה֩ אַבְרָהָ֨ם אֶת־הַשָּׂדֶ֜ה לַאֲחֻזַּת־קֶ֗בֶר מֵאֵ֛ת עֶפְרֹ֥ן הַחִתִּ֖י עַל־פְּנֵ֥י מַמְרֵֽא׃ \Rashi{\textbf{וישאו אתו בניו.}ולא בני בניו, שכך צום, אל ישאו מטתי לא איש מצרי ולא אחד מבניכם, שהם מבנות כנען, אלא אתם, וקבע להם מקום ג' למזרח, וכן לד' רוחות, וכסדרן למסע מחנה של דגלים נקבעו כאן; לוי לא ישא, שהוא עתיד לשאת את הארון, ויוסף לא ישא, לפי שהוא מלך, מנשה ואפרים יהיו תחתיהם, וזהו איש על דגלו באתת – באות שמסר להם אביהם לשא מטתו:}%
% ---------- פרק 50 | פסוק 14 ----------
 \textbf{יד} וַיָּ֨שׇׁב יוֹסֵ֤ף מִצְרַ֙יְמָה֙ ה֣וּא וְאֶחָ֔יו וְכׇל־הָעֹלִ֥ים אִתּ֖וֹ לִקְבֹּ֣ר אֶת־אָבִ֑יו אַחֲרֵ֖י קׇבְר֥וֹ אֶת־אָבִֽיו׃ \Rashi{\textbf{הוא ואחיו וכל העלים אתו.}בחזרתן הקדים אחיו למצרים העולים אתו, ובהליכתן הקדים מצרים לאחיו, שנאמר ויעלו אתו כל עבדי פרעה וגו', ואחר כך כל בית יוסף ואחיו, אלא לפי שראו כבוד שעשו מלכי כנען שתלו כתריהם בארונו של יעקב נהגו בהם כבוד:}%
% ---------- פרק 50 | פסוק 15 ----------
 \textbf{טו} וַיִּרְא֤וּ אֲחֵֽי־יוֹסֵף֙ כִּי־מֵ֣ת אֲבִיהֶ֔ם וַיֹּ֣אמְר֔וּ ל֥וּ יִשְׂטְמֵ֖נוּ יוֹסֵ֑ף וְהָשֵׁ֤ב יָשִׁיב֙ לָ֔נוּ אֵ֚ת כׇּל־הָ֣רָעָ֔ה אֲשֶׁ֥ר גָּמַ֖לְנוּ אֹתֽוֹ׃ \Rashi{\textbf{ויראו אחי יוסף כי מת אביהם.}מהו ויראו? הכירו במיתתו אצל יוסף, שהיו רגילים לסעד על שלחנו של יוסף והיה מקרבן בשביל כבוד אביו, ומשמת יעקב לא קרבן (בראשית רבה): לו ישטמנו. שמא ישטמנו; לו מתחלק לענינים הרבה; יש לו משמש בלשון בקשה ולשון הלואי, כגון לו יהי כדברך, לו שמעני, ולו הואלנו, לו מתנו, ויש לו משמש בלשון אם ואולי, כגון לו חכמו (דברים ל"ב), לו הקשבת למצותי (ישעיהו מ"ח), ולו אנכי שקל על כפי (שמואל ב י"ח); ויש לו משמש בלשון שמא, לו ישטמנו, ואין לו עוד דומה במקרא, והוא לשון אולי, כמו אלי לא תלך האשה אחרי, לשון שמא הוא; ויש אולי לשון בקשה, כגון אולי יראה ה' בעיני (שמואל ב ט"ז), אולי ה' אותי (יהושע י"ד), הרי הוא כמו לו יהי כדברך, ויש אולי לשון אם, אולי יש נ' צדיקם:}%
% ---------- פרק 50 | פסוק 16 ----------
 \textbf{טז} וַיְצַוּ֕וּ אֶל־יוֹסֵ֖ף לֵאמֹ֑ר אָבִ֣יךָ צִוָּ֔ה לִפְנֵ֥י מוֹת֖וֹ לֵאמֹֽר׃ \Rashi{\textbf{ויצוו אל יוסף.}כמו ויצום אל בני ישראל – צוה למשה ולאהרן להיות שלוחים אל בני ישראל, אף זה ויצוו אל שלוחם, להיות שליח אל יוסף לומר לו כן. ואת מי צוו? את בני בלהה, שהיו רגילין אצלו, שנאמר והוא נער את בני בלהה: אביך צוה. שנו בדבר מפני השלום, כי לא צוה יעקב כן, שלא נחשד יוסף בעיניו (בראשית רבה, יבמות ס"ה):}%
% ---------- פרק 50 | פסוק 17 ----------
 \textbf{יז} כֹּֽה־תֹאמְר֣וּ לְיוֹסֵ֗ף אָ֣נָּ֡א שָׂ֣א נָ֠א פֶּ֣שַׁע אַחֶ֤יךָ וְחַטָּאתָם֙ כִּי־רָעָ֣ה גְמָל֔וּךָ וְעַתָּה֙ שָׂ֣א נָ֔א לְפֶ֥שַׁע עַבְדֵ֖י אֱלֹהֵ֣י אָבִ֑יךָ וַיֵּ֥בְךְּ יוֹסֵ֖ף בְּדַבְּרָ֥ם אֵלָֽיו׃ \Rashi{\textbf{שא נא לפשע עבדי אלהי אביך.}אם אביך מת, אלהיו קים, והם עבדיו:}%
% ---------- פרק 50 | פסוק 18 ----------
 \textbf{יח} וַיֵּלְכוּ֙ גַּם־אֶחָ֔יו וַֽיִּפְּל֖וּ לְפָנָ֑יו וַיֹּ֣אמְר֔וּ הִנֶּ֥נּֽוּ לְךָ֖ לַעֲבָדִֽים׃ \Rashi{\textbf{וילכו גם אחיו.}מוסף על השליחות:}%
% ---------- פרק 50 | פסוק 19 ----------
 \textbf{יט} וַיֹּ֧אמֶר אֲלֵהֶ֛ם יוֹסֵ֖ף אַל־תִּירָ֑אוּ כִּ֛י הֲתַ֥חַת אֱלֹהִ֖ים אָֽנִי׃ \Rashi{\textbf{כי התחת אלהים אני.}שמא במקומו אני? בתמיה, אם הייתי רוצה להרע לכם, כלום אני יכול? והלא אתם כלכם חשבתם עלי רעה, והקדוש ברוך הוא חשבה לטובה והיאך אני לבדי יכול להרע לכם?}%
% ---------- פרק 50 | פסוק 20 ----------
 \textbf{כ} וְאַתֶּ֕ם חֲשַׁבְתֶּ֥ם עָלַ֖י רָעָ֑ה אֱלֹהִים֙ חֲשָׁבָ֣הּ לְטֹבָ֔ה לְמַ֗עַן עֲשֹׂ֛ה כַּיּ֥וֹם הַזֶּ֖ה לְהַחֲיֹ֥ת עַם־רָֽב׃ %
% ---------- פרק 50 | פסוק 21 ----------
 \textbf{כא} וְעַתָּה֙ אַל־תִּירָ֔אוּ אָנֹכִ֛י אֲכַלְכֵּ֥ל אֶתְכֶ֖ם וְאֶֽת־טַפְּכֶ֑ם וַיְנַחֵ֣ם אוֹתָ֔ם וַיְדַבֵּ֖ר עַל־לִבָּֽם׃ \Rashi{\textbf{וידבר על לבם.}דברים המתקבלים על הלב – עד שלא ירדתם לכאן היו מרננים עלי שאני עבד, על ידיכם נודע שאני בן חרין, ואם אני הורג אתכם, מה הבריות אומרות? כת של בחורים ראה ונשתבח בהם ואמר אחי הם, ולבסוף הרג אותם; יש לך אח שהורג את אחיו? דבר אחר, עשרה נרות לא יכלו לכבות נר א' וכו':}%
% ---------- פרק 50 | פסוק 22 ----------
 \textbf{כב} וַיֵּ֤שֶׁב יוֹסֵף֙ בְּמִצְרַ֔יִם ה֖וּא וּבֵ֣ית אָבִ֑יו וַיְחִ֣י יוֹסֵ֔ף מֵאָ֥ה וָעֶ֖שֶׂר שָׁנִֽים׃ %
% ---------- פרק 50 | פסוק 23 ----------
 \textbf{כג} וַיַּ֤רְא יוֹסֵף֙ לְאֶפְרַ֔יִם בְּנֵ֖י שִׁלֵּשִׁ֑ים גַּ֗ם בְּנֵ֤י מָכִיר֙ בֶּן־מְנַשֶּׁ֔ה יֻלְּד֖וּ עַל־בִּרְכֵּ֥י יוֹסֵֽף׃ \Rashi{\textbf{על ברכי יוסף.}כתרגומו גדלן בין ברכיו:}%
% ---------- פרק 50 | פסוק 24 ----------
 \textbf{כד} וַיֹּ֤אמֶר יוֹסֵף֙ אֶל־אֶחָ֔יו אָנֹכִ֖י מֵ֑ת וֵֽאלֹהִ֞ים פָּקֹ֧ד יִפְקֹ֣ד אֶתְכֶ֗ם וְהֶעֱלָ֤ה אֶתְכֶם֙ מִן־הָאָ֣רֶץ הַזֹּ֔את אֶל־הָאָ֕רֶץ אֲשֶׁ֥ר נִשְׁבַּ֛ע לְאַבְרָהָ֥ם לְיִצְחָ֖ק וּֽלְיַעֲקֹֽב׃ %
% ---------- פרק 50 | פסוק 25 ----------
 \textbf{כה} וַיַּשְׁבַּ֣ע יוֹסֵ֔ף אֶת־בְּנֵ֥י יִשְׂרָאֵ֖ל לֵאמֹ֑ר פָּקֹ֨ד יִפְקֹ֤ד אֱלֹהִים֙ אֶתְכֶ֔ם וְהַעֲלִתֶ֥ם אֶת־עַצְמֹתַ֖י מִזֶּֽה׃ %
% ---------- פרק 50 | פסוק 26 ----------
 \textbf{כו} וַיָּ֣מׇת יוֹסֵ֔ף בֶּן־מֵאָ֥ה וָעֶ֖שֶׂר שָׁנִ֑ים וַיַּחַנְט֣וּ אֹת֔וֹ וַיִּ֥ישֶׂם בָּאָר֖וֹן בְּמִצְרָֽיִם׃ %


\newpage
\begin{center}
  {\Huge \textbf{חומש שמות}}
\end{center}
\vspace{1.5cm}

\addparasha{חומש שמות | פרשת שמות}
% ---------- פרק 1 | פסוק 1 ----------
 \textbf{א} וְאֵ֗לֶּה שְׁמוֹת֙ בְּנֵ֣י יִשְׂרָאֵ֔ל הַבָּאִ֖ים מִצְרָ֑יְמָה אֵ֣ת יַעֲקֹ֔ב אִ֥ישׁ וּבֵית֖וֹ בָּֽאוּ׃ \Rashi{\textbf{ואלה שמות בני ישראל.}אע"פ שמנאן בחייהם בשמותם, חזר ומנאם במיתתם, להודיע חבתם, שנמשלו לכוכבים, שמוציאם ומכניסם במספר ובשמותם (שמות רבה), שנ' "המוציא במספר צבאם לכלם בשם יקרא" (ישעיהו מ'):}%
% ---------- פרק 1 | פסוק 2 ----------
 \textbf{ב} רְאוּבֵ֣ן שִׁמְע֔וֹן לֵוִ֖י וִיהוּדָֽה׃ %
% ---------- פרק 1 | פסוק 3 ----------
 \textbf{ג} יִשָּׂשכָ֥ר זְבוּלֻ֖ן וּבִנְיָמִֽן׃ %
% ---------- פרק 1 | פסוק 4 ----------
 \textbf{ד} דָּ֥ן וְנַפְתָּלִ֖י גָּ֥ד וְאָשֵֽׁר׃ %
% ---------- פרק 1 | פסוק 5 ----------
 \textbf{ה} וַֽיְהִ֗י כׇּל־נֶ֛פֶשׁ יֹצְאֵ֥י יֶֽרֶךְ־יַעֲקֹ֖ב שִׁבְעִ֣ים נָ֑פֶשׁ וְיוֹסֵ֖ף הָיָ֥ה בְמִצְרָֽיִם׃ \Rashi{\textbf{ויוסף היה במצרים.}והלא הוא ובניו היו בכלל שבעים, ומה בא ללמדנו? וכי לא היינו יודעים שהוא היה במצרים? אלא להודיעך צדקתו של יוסף, הוא יוסף הרועה את צאן אביו, הוא יוסף שהיה במצרים, ונעשה מלך, ועומד בצדקו (שמות רבה):}%
% ---------- פרק 1 | פסוק 6 ----------
 \textbf{ו} וַיָּ֤מׇת יוֹסֵף֙ וְכׇל־אֶחָ֔יו וְכֹ֖ל הַדּ֥וֹר הַהֽוּא׃ %
% ---------- פרק 1 | פסוק 7 ----------
 \textbf{ז} וּבְנֵ֣י יִשְׂרָאֵ֗ל פָּר֧וּ וַֽיִּשְׁרְצ֛וּ וַיִּרְבּ֥וּ וַיַּֽעַצְמ֖וּ בִּמְאֹ֣ד מְאֹ֑ד וַתִּמָּלֵ֥א הָאָ֖רֶץ אֹתָֽם׃ \{פ\} \Rashi{\textbf{פרו.}שלא הפילו נשותיהם ולא מתו כשהם קטנים:  וישרצו. שהיו יולדות ששה בכרס אחד (שם):}%
% ---------- פרק 1 | פסוק 8 ----------
 \textbf{ח} וַיָּ֥קׇם מֶֽלֶךְ־חָדָ֖שׁ עַל־מִצְרָ֑יִם אֲשֶׁ֥ר לֹֽא־יָדַ֖ע אֶת־יוֹסֵֽף׃ \Rashi{\textbf{ויקם מלך חדש.}רב ושמואל, חד אמר חדש ממש, וחד אמר שנתחדשו גזרותיו: אשר לא ידע. עשה עצמו כאלו לא ידעו (סוטה י"א):}%
% ---------- פרק 1 | פסוק 9 ----------
 \textbf{ט} וַיֹּ֖אמֶר אֶל־עַמּ֑וֹ הִנֵּ֗ה עַ֚ם בְּנֵ֣י יִשְׂרָאֵ֔ל רַ֥ב וְעָצ֖וּם מִמֶּֽנּוּ׃ %
% ---------- פרק 1 | פסוק 10 ----------
 \textbf{י} הָ֥בָה נִֽתְחַכְּמָ֖ה ל֑וֹ פֶּן־יִרְבֶּ֗ה וְהָיָ֞ה כִּֽי־תִקְרֶ֤אנָה מִלְחָמָה֙ וְנוֹסַ֤ף גַּם־הוּא֙ עַל־שֹׂ֣נְאֵ֔ינוּ וְנִלְחַם־בָּ֖נוּ וְעָלָ֥ה מִן־הָאָֽרֶץ׃ \Rashi{\textbf{הבה נתחכמה.}כל הבה לשון הכנה והזמנה לדבר הוא, כלומר, הזמינו עצמכם לכך: נתחכמה לו. לעם; נתחכמה מה לעשות לו. ורבותינו דרשו נתחכם למושיען של ישראל לדונם במים, שכבר נשבע שלא יביא מבול לעולם (שם): ועלה מן הארץ. על כרחנו. ורבותינו דרשו כאדם שמקלל עצמו ותולה קללתו באחרים (שם), והרי הוא כאלו כתב ועלינו מן הארץ – והם יירשוה:}%
% ---------- פרק 1 | פסוק 11 ----------
 \textbf{יא} וַיָּשִׂ֤ימוּ עָלָיו֙ שָׂרֵ֣י מִסִּ֔ים לְמַ֥עַן עַנֹּת֖וֹ בְּסִבְלֹתָ֑ם וַיִּ֜בֶן עָרֵ֤י מִסְכְּנוֹת֙ לְפַרְעֹ֔ה אֶת־פִּתֹ֖ם וְאֶת־רַעַמְסֵֽס׃ \Rashi{\textbf{עליו.}על העם: מסים. לשון מס, שרים שגובין מהם המס, ומהו המס? שיבנו ערי מסכנות לפרעה: למען ענתו בסבלתם. של מצרים: ערי מסכנות. כתרגומו, וכן "לך בא אל הסוכן הזה" (ישעיהו כ"ב) – גזבר הממנה על האוצרות: את פתם ואת רעמסס. שלא היו ראויות מתחלה לכך, ועשאום חזקות ובצורות לאוצר:}%
% ---------- פרק 1 | פסוק 12 ----------
 \textbf{יב} וְכַאֲשֶׁר֙ יְעַנּ֣וּ אֹת֔וֹ כֵּ֥ן יִרְבֶּ֖ה וְכֵ֣ן יִפְרֹ֑ץ וַיָּקֻ֕צוּ מִפְּנֵ֖י בְּנֵ֥י יִשְׂרָאֵֽל׃ \Rashi{\textbf{וכאשר יענו אתו.}בכל מה שהם נותנין לב לענות, כן לב הקב"ה להרבות ולהפריץ: כן ירבה. כן רבה וכן פרץ. ומדרשו, רוח הקדש אומרת כן, אתם אומרים פן ירבה ואני אומר כן ירבה (סוטה י"א): ויקצו. קצו בחייהם. ורבותינו דרשו כקוצים היו בעיניהם (שם):}%
% ---------- פרק 1 | פסוק 13 ----------
 \textbf{יג} וַיַּעֲבִ֧דוּ מִצְרַ֛יִם אֶת־בְּנֵ֥י יִשְׂרָאֵ֖ל בְּפָֽרֶךְ׃ \Rashi{\textbf{בפרך.}בעבודה קשה, המפרכת את הגוף ומשברתו (שם):}%
% ---------- פרק 1 | פסוק 14 ----------
 \textbf{יד} וַיְמָרְר֨וּ אֶת־חַיֵּיהֶ֜ם בַּעֲבֹדָ֣ה קָשָׁ֗ה בְּחֹ֙מֶר֙ וּבִלְבֵנִ֔ים וּבְכׇל־עֲבֹדָ֖ה בַּשָּׂדֶ֑ה אֵ֚ת כׇּל־עֲבֹ֣דָתָ֔ם אֲשֶׁר־עָבְד֥וּ בָהֶ֖ם בְּפָֽרֶךְ׃ %
% ---------- פרק 1 | פסוק 15 ----------
 \textbf{טו} וַיֹּ֙אמֶר֙ מֶ֣לֶךְ מִצְרַ֔יִם לַֽמְיַלְּדֹ֖ת הָֽעִבְרִיֹּ֑ת אֲשֶׁ֨ר שֵׁ֤ם הָֽאַחַת֙ שִׁפְרָ֔ה וְשֵׁ֥ם הַשֵּׁנִ֖ית פּוּעָֽה׃ \Rashi{\textbf{למילדת.}הוא לשון מולידות, אלא שיש ל' קל ויש ל' כבד, כמו שובר ומשבר, דובר ומדבר, כך מוליד ומילד: שפרה. יוכבד, על שם שמשפרת את הולד (שם): פועה. זו מרים, שפועה ומדברת והוגה לולד (שם), כדרך הנשים המפיסות תינוק הבוכה. פועה. לשון צעקה, כמו "כיולדה אפעה" (ישעיה מ"ב):}%
% ---------- פרק 1 | פסוק 16 ----------
 \textbf{טז} וַיֹּ֗אמֶר בְּיַלֶּדְכֶן֙ אֶת־הָֽעִבְרִיּ֔וֹת וּרְאִיתֶ֖ן עַל־הָאׇבְנָ֑יִם אִם־בֵּ֥ן הוּא֙ וַהֲמִתֶּ֣ן אֹת֔וֹ וְאִם־בַּ֥ת הִ֖וא וָחָֽיָה׃ \Rashi{\textbf{בילדכן.}כמו בהולידכן: על האבנים. מושב האשה היולדת, במקום אחר (ישעיהו ל"ז), קוראו משבר; וכמוהו "עושה מלאכה על האבנים" (ירמיהו י"ח), – מושב כלי אמנות יוצר חרש: אם בן הוא וגו'. לא היה מקפיד אלא על הזכרים, שאמרו לו אצטגניניו שעתיד להולד בן המושיע אותם: וחיה. ותחיה:}%
% ---------- פרק 1 | פסוק 17 ----------
 \textbf{יז} וַתִּירֶ֤אןָ הַֽמְיַלְּדֹת֙ אֶת־הָ֣אֱלֹהִ֔ים וְלֹ֣א עָשׂ֔וּ כַּאֲשֶׁ֛ר דִּבֶּ֥ר אֲלֵיהֶ֖ן מֶ֣לֶךְ מִצְרָ֑יִם וַתְּחַיֶּ֖יןָ אֶת־הַיְלָדִֽים׃ \Rashi{\textbf{ותחיין את הילדים.}מספקות להם מים ומזון. תרגום הראשון וקימא, והשני וקימתן, לפי שלשון עבר לנקבות רבות תבה זו וכיוצא בה משמשת לשון פעלו ולשון פעלתם, כגון "ותאמרן איש מצרי" (שמות ב'), (לשון עבר) כמו ויאמרו לזכרים, "ותדברנה בפיכם" (ירמיהו מ"ד), לשון דברתן, כמו ותדברו לזכרים, וכן "ותחללנה אתי אל עמי" (יחזקאל י"ג), לשון עבר חללתן, כמו ותחללו לזכרים:}%
% ---------- פרק 1 | פסוק 18 ----------
 \textbf{יח} וַיִּקְרָ֤א מֶֽלֶךְ־מִצְרַ֙יִם֙ לַֽמְיַלְּדֹ֔ת וַיֹּ֣אמֶר לָהֶ֔ן מַדּ֥וּעַ עֲשִׂיתֶ֖ן הַדָּבָ֣ר הַזֶּ֑ה וַתְּחַיֶּ֖יןָ אֶת־הַיְלָדִֽים׃ %
% ---------- פרק 1 | פסוק 19 ----------
 \textbf{יט} וַתֹּאמַ֤רְןָ הַֽמְיַלְּדֹת֙ אֶל־פַּרְעֹ֔ה כִּ֣י לֹ֧א כַנָּשִׁ֛ים הַמִּצְרִיֹּ֖ת הָֽעִבְרִיֹּ֑ת כִּֽי־חָי֣וֹת הֵ֔נָּה בְּטֶ֨רֶם תָּב֧וֹא אֲלֵהֶ֛ן*(בכתר ארם צובה היה כתוב אֲלֵיהֶ֛ן) הַמְיַלֶּ֖דֶת וְיָלָֽדוּ׃ \Rashi{\textbf{כי חיות הנה.}בקיאות כמילדות, תרגום מילדות חיתא. ורבותינו דרשו הרי הן משולות כחיות השדה (סוטה י"א), שאינן צריכות מילדות. והיכן משולות לחיות? גור אריה, זאב יטרף (בראשית מ"ט), בכור שורו (דברים ל"ג), אילה שלחה (בראשית מ"ט), ומי שלא נכתב בו, הרי הכתוב כללן ויברך אותם (שם), ועוד כתיב מה אמך לביא (יחזקאל י"ט):}%
% ---------- פרק 1 | פסוק 20 ----------
 \textbf{כ} וַיֵּ֥יטֶב אֱלֹהִ֖ים לַֽמְיַלְּדֹ֑ת וַיִּ֧רֶב הָעָ֛ם וַיַּֽעַצְמ֖וּ מְאֹֽד׃ \Rashi{\textbf{וייטב.}הטיב להן. וזה חלוק בתבה שיסודה שתי אותיות ונתן לה וי"ו יו"ד בראשה, כשהיא באה לדבר לשון ויפעיל הוא נקוד היו"ד בצירי, שהוא קמץ קטן, כגון וייטב אלהים למילדת, "וירב בבת יהודה" (איכה ב'), – הרבה תאניה. וכן "ויגל השארית" דנבוזראדן (דברי הימים ב' ל"ו), – הגלה את השארית, "ויפן זנב אל זנב" (שופטים ט"ו) – הפנה הזנבות זו לזו, כל אלו לשון הפעיל את אחרים; וכשהוא מדבר בלשון ויפעל, הוא נקוד היו"ד בחירק, כגון "וייטב בעיניו" (ויקרא י'), לשון הוטב, וכן וירב העם, – נתרבה העם, "ויגל יהודה" (מלכים ב כ"ה) – הגלה יהודה, "ויפן כה וכה" (שמות ב׳:י״ב) – הפנה לכאן ולכאן. ואל תשיבני וילך, וישב, וירד, ויצא, לפי שאינן מגזרתן של אלו, שהרי היו"ד יסוד בהן, ילך, ישב, ירד, יצא – י' אות שלישית בו: וייטב אלהים למילדת. מהו הטובה?:}%
% ---------- פרק 1 | פסוק 21 ----------
 \textbf{כא} וַיְהִ֕י כִּֽי־יָרְא֥וּ הַֽמְיַלְּדֹ֖ת אֶת־הָאֱלֹהִ֑ים וַיַּ֥עַשׂ לָהֶ֖ם בָּתִּֽים׃ \Rashi{\textbf{ויעש להם בתים.}בתי כהנה ולויה ומלכות שקרויין בתים, כמו שכתוב: "לבנות את בית ה' ואת בית המלך" (מלכים א ט׳:א׳), כהנה ולויה מיוכבד ומלכות ממרים. כדאיתא במסכת סוטה:}%
% ---------- פרק 1 | פסוק 22 ----------
 \textbf{כב} וַיְצַ֣ו פַּרְעֹ֔ה לְכׇל־עַמּ֖וֹ לֵאמֹ֑ר כׇּל־הַבֵּ֣ן הַיִּלּ֗וֹד הַיְאֹ֙רָה֙ תַּשְׁלִיכֻ֔הוּ וְכׇל־הַבַּ֖ת תְּחַיּֽוּן׃ \{פ\} \Rashi{\textbf{לכל עמו.}אף עליהם גזר (סוטה י"ב), יום שנולד משה אמרו לו אצטגניניו, היום נולד מושיען, ואין אנו יודעים אם ממצרים אם מישראל, ורואין אנו שסופו ללקות במים, לפיכך גזר אותו היום אף על המצריים, שנאמר כל הבן הילוד, ולא נאמר הילוד לעברים; והם לא היו יודעים שסופו ללקות על מי מריבה:}%
% ---------- פרק 2 | פסוק 1 ----------
 \textbf{א} וַיֵּ֥לֶךְ אִ֖ישׁ מִבֵּ֣ית לֵוִ֑י וַיִּקַּ֖ח אֶת־בַּת־לֵוִֽי׃ \Rashi{\textbf{ויקח את בת לוי.}פרוש היה ממנה מפני גזרת פרעה, והחזירה ועשה בה לקוחין שניים, ואף היא נהפכה להיות נערה; ובת ק"ל שנה היתה, שנולדה בבואם למצרים בין החומות, ומאתים ועשר נשתהו שם, וכשיצאו היה משה בן שמונים שנה, אם כן כשנתעברה ממנו היתה בת מאה ושלושים וקורא אותה בת לוי (עי' סוטה י"ב, בבא בתרא קי"ט ושמות רבה):}%
% ---------- פרק 2 | פסוק 2 ----------
 \textbf{ב} וַתַּ֥הַר הָאִשָּׁ֖ה וַתֵּ֣לֶד בֵּ֑ן וַתֵּ֤רֶא אֹתוֹ֙ כִּי־ט֣וֹב ה֔וּא וַֽתִּצְפְּנֵ֖הוּ שְׁלֹשָׁ֥ה יְרָחִֽים׃ \Rashi{\textbf{כי טוב הוא.}כשנולד נתמלא הבית כלו אורה (סוטה י"ב):}%
% ---------- פרק 2 | פסוק 3 ----------
 \textbf{ג} וְלֹא־יָכְלָ֣ה עוֹד֮ הַצְּפִינוֹ֒ וַתִּֽקַּֽח־לוֹ֙ תֵּ֣בַת גֹּ֔מֶא וַתַּחְמְרָ֥הֿ בַחֵמָ֖ר וּבַזָּ֑פֶת וַתָּ֤שֶׂם בָּהּ֙ אֶת־הַיֶּ֔לֶד וַתָּ֥שֶׂם בַּסּ֖וּף עַל־שְׂפַ֥ת הַיְאֹֽר׃ \Rashi{\textbf{ולא יכלה עוד הצפינו.}שמנו לה המצריים מיום שהחזירה, והיא ילדתו לששה חדשים ויום אחד (שם), שהיולדת לשבעה יולדת למקטעין (נדה ל"ח), והם בדקו אחריה לסוף תשעה: גמא. גמי בלשון משנה (שבת ע"ח), ובלעז יונק"ו, ודבר רך הוא ועומד בפני רך ובפני קשה (סוטה י"ב): בחמר ובזפת. זפת מבחוץ וטיט מבפנים, כדי שלא יריח אותו צדיק ריח רע של זפת (שם): ותשם בסוף. הוא לשון אגם, רושיי"ל בלעז, ודומה לו "קנה וסוף קמלו" (ישעיהו י"ט):}%
% ---------- פרק 2 | פסוק 4 ----------
 \textbf{ד} וַתֵּתַצַּ֥ב אֲחֹת֖וֹ מֵרָחֹ֑ק לְדֵעָ֕ה מַה־יֵּעָשֶׂ֖ה לֽוֹ׃ %
% ---------- פרק 2 | פסוק 5 ----------
 \textbf{ה} וַתֵּ֤רֶד בַּת־פַּרְעֹה֙ לִרְחֹ֣ץ עַל־הַיְאֹ֔ר וְנַעֲרֹתֶ֥יהָ הֹלְכֹ֖ת עַל־יַ֣ד הַיְאֹ֑ר וַתֵּ֤רֶא אֶת־הַתֵּבָה֙ בְּת֣וֹךְ הַסּ֔וּף וַתִּשְׁלַ֥ח אֶת־אֲמָתָ֖הּ וַתִּקָּחֶֽהָ׃ \Rashi{\textbf{לרחץ על היאר.}סרס המקרא ופרשהו: ותרד בת פרעה על היאר לרחץ בו: על יד היאר. אצל היאור, כמו "ראו חלקת יואב אל ידי" (שמואל ב י"ד), והוא לשון יד ממש, שיד האדם סמוכה לו. ורבותינו אמרו, "הולכות" לשון מיתה (סוטה י"ב), כמו "הנה אנכי הולך למות" (בראשית כ"ה) – הולכות למות לפי שמחו בה; והכתוב מסיען, כי למה לנו לכתב ונערותיה הולכות: את אמתה. את שפחתה. ורבותינו דרשו (סוטה שם), לשון יד, אבל לפי דקדוק לשון הקדש היה לו להנקד אמתה, דגושה, והם דרשו את אמתה – את ידה, ונשתרבבה אמתה אמות הרבה:}%
% ---------- פרק 2 | פסוק 6 ----------
 \textbf{ו} וַתִּפְתַּח֙ וַתִּרְאֵ֣הוּ אֶת־הַיֶּ֔לֶד וְהִנֵּה־נַ֖עַר בֹּכֶ֑ה וַתַּחְמֹ֣ל עָלָ֔יו וַתֹּ֕אמֶר מִיַּלְדֵ֥י הָֽעִבְרִ֖ים זֶֽה׃ \Rashi{\textbf{ותפתח ותראהו.}את מי ראתה? את הילד, זהו פשוטו. ומדרשו (שם), שראתה עמו שכינה: והנה נער בכה. קולו כנער (שם):}%
% ---------- פרק 2 | פסוק 7 ----------
 \textbf{ז} וַתֹּ֣אמֶר אֲחֹתוֹ֮ אֶל־בַּת־פַּרְעֹה֒ הַאֵלֵ֗ךְ וְקָרָ֤אתִי לָךְ֙ אִשָּׁ֣ה מֵינֶ֔קֶת מִ֖ן הָעִבְרִיֹּ֑ת וְתֵינִ֥ק לָ֖ךְ אֶת־הַיָּֽלֶד׃ \Rashi{\textbf{מן העברית.}שהחזירתו על מצריות הרבה לינק ולא ינק, לפי שהיה עתיד לדבר עם השכינה (שם):}%
% ---------- פרק 2 | פסוק 8 ----------
 \textbf{ח} וַתֹּֽאמֶר־לָ֥הּ בַּת־פַּרְעֹ֖ה לֵ֑כִי וַתֵּ֙לֶךְ֙ הָֽעַלְמָ֔ה וַתִּקְרָ֖א אֶת־אֵ֥ם הַיָּֽלֶד׃ \Rashi{\textbf{ותלך העלמה.}הלכה בזריזות ועלמות כעלם (שם):}%
% ---------- פרק 2 | פסוק 9 ----------
 \textbf{ט} וַתֹּ֧אמֶר לָ֣הּ בַּת־פַּרְעֹ֗ה הֵילִ֜יכִי אֶת־הַיֶּ֤לֶד הַזֶּה֙ וְהֵינִקִ֣הוּ לִ֔י וַאֲנִ֖י אֶתֵּ֣ן אֶת־שְׂכָרֵ֑ךְ וַתִּקַּ֧ח הָאִשָּׁ֛ה הַיֶּ֖לֶד וַתְּנִיקֵֽהוּ׃ \Rashi{\textbf{היליכי.}נתנבאה ולא ידעה מה נתנבאה – הי שליכי:}%
% ---------- פרק 2 | פסוק 10 ----------
 \textbf{י} וַיִּגְדַּ֣ל הַיֶּ֗לֶד וַתְּבִאֵ֙הוּ֙ לְבַת־פַּרְעֹ֔ה וַֽיְהִי־לָ֖הּ לְבֵ֑ן וַתִּקְרָ֤א שְׁמוֹ֙ מֹשֶׁ֔ה וַתֹּ֕אמֶר כִּ֥י מִן־הַמַּ֖יִם מְשִׁיתִֽהוּ׃ \Rashi{\textbf{משיתהו.}שחליתה, הוא לשון הוצאה בלשון ארמי, "כמשחל בינתא מחלבא", ובלשון עברי משיתיהו לשון הסירותיו, כמו לא ימוש, לא משו, כך חברו מנחם. ואני אומר שאינו ממחברת מש וימוש אלא מגזרת משה, ולשון הוצאה הוא, וכן "ימשני ממים רבים", (שמואל ב כ"ב), שאלו היה ממחברת מש לא יתכן לומר משיתיהו אלא המישותיהו, כאשר יאמר מן קם הקימותי, ומן שב השיבותי, ומן בא הביאותי, או משתיהו, כמו "ומשתי את עון הארץ" (זכריה ג'), אבל משיתי אינו אלא מגזרת תבה שפעל שלה מיסד בה"א בסוף התבה, כגון משה, בנה, עשה, צוה, פנה, כשיבא לומר בהם פעלתי, תבא היו"ד במקום ה"א, כמו בניתי, עשיתי, צויתי:}%
% ---------- פרק 2 | פסוק 11 ----------
 \textbf{יא} וַיְהִ֣י ׀ בַּיָּמִ֣ים הָהֵ֗ם וַיִּגְדַּ֤ל מֹשֶׁה֙ וַיֵּצֵ֣א אֶל־אֶחָ֔יו וַיַּ֖רְא בְּסִבְלֹתָ֑ם וַיַּרְא֙ אִ֣ישׁ מִצְרִ֔י מַכֶּ֥ה אִישׁ־עִבְרִ֖י מֵאֶחָֽיו׃ \Rashi{\textbf{ויגדל משה.}והלא כבר כתב ויגדל הילד? אמר רבי יהודה ברבי אלעאי, הראשון לקומה והשני לגדלה, שמנהו פרעה על ביתו (ילקוט שמעוני): וירא בסבלתם. נתן עיניו ולבו להיות מצר עליהם (שמות רבה א'): איש מצרי. נוגש היה ממנה על שוטרי ישראל והיה מעמידם מקרות הגבר למלאכתם (שם): מכה איש עברי. מלקהו ורודהו. ובעלה של שלומית בת דברי היה ונתן עיניו בה, ובלילה העמידו והוציאו מביתו, והוא חזר ונכנס לבית ובא על אשתו, כסבורה שהוא בעלה, וחזר האיש לביתו והרגיש בדבר, וכשראה אותו מצרי שהרגיש בדבר, היה מכהו ורודהו כל היום (שם):}%
% ---------- פרק 2 | פסוק 12 ----------
 \textbf{יב} וַיִּ֤פֶן כֹּה֙ וָכֹ֔ה וַיַּ֖רְא כִּ֣י אֵ֣ין אִ֑ישׁ וַיַּךְ֙ אֶת־הַמִּצְרִ֔י וַֽיִּטְמְנֵ֖הוּ בַּחֽוֹל׃ \Rashi{\textbf{ויפן כה וכה.}ראה מה עשה לו בבית ומה עשה לו בשדה (שם). ולפי פשוטו כמשמעו: וירא כי אין איש. עתיד לצאת ממנו שיתגיר (ת"י):}%
% ---------- פרק 2 | פסוק 13 ----------
 \textbf{יג} וַיֵּצֵא֙ בַּיּ֣וֹם הַשֵּׁנִ֔י וְהִנֵּ֛ה שְׁנֵֽי־אֲנָשִׁ֥ים עִבְרִ֖ים נִצִּ֑ים וַיֹּ֙אמֶר֙ לָֽרָשָׁ֔ע לָ֥מָּה תַכֶּ֖ה רֵעֶֽךָ׃ \Rashi{\textbf{שני אנשים עברים.}דתן ואבירם, הם שהותירו מן המן (נדרים ס"ד): נצים. מריבים: למה תכה. אע"פ שלא הכהו נקרא רשע בהרמת יד (סנהדרין נ"ח): רעך. רשע כמותך (תנחומא):}%
% ---------- פרק 2 | פסוק 14 ----------
 \textbf{יד} וַ֠יֹּ֠אמֶר מִ֣י שָֽׂמְךָ֞ לְאִ֨ישׁ שַׂ֤ר וְשֹׁפֵט֙ עָלֵ֔ינוּ הַלְהׇרְגֵ֙נִי֙ אַתָּ֣ה אֹמֵ֔ר כַּאֲשֶׁ֥ר הָרַ֖גְתָּ אֶת־הַמִּצְרִ֑י וַיִּירָ֤א מֹשֶׁה֙ וַיֹּאמַ֔ר אָכֵ֖ן נוֹדַ֥ע הַדָּבָֽר׃ \Rashi{\textbf{מי שמך לאיש.}והנה עודך נער: הלהרגני אתה אמר. מכאן אנו למדים שהרגו בשם המפרש (שמות רבה א'): ויירא משה. כפשוטו. ומדרשו: דאג לו על שראה בישראל רשעים דלטורין, אמר, מעתה שמא אינם ראויין להגאל (שמות רבה א'): אכן נודע הדבר. כמשמעו. ומדרשו, נודע לי הדבר שהייתי תמה עליו, מה חטאו ישראל מכל שבעים אמות להיות נרדים בעבודת פרך, אבל רואה אני שהם ראויים לכך (שמות רבה א'):}%
% ---------- פרק 2 | פסוק 15 ----------
 \textbf{טו} וַיִּשְׁמַ֤ע פַּרְעֹה֙ אֶת־הַדָּבָ֣ר הַזֶּ֔ה וַיְבַקֵּ֖שׁ לַהֲרֹ֣ג אֶת־מֹשֶׁ֑ה וַיִּבְרַ֤ח מֹשֶׁה֙ מִפְּנֵ֣י פַרְעֹ֔ה וַיֵּ֥שֶׁב בְּאֶֽרֶץ־מִדְיָ֖ן וַיֵּ֥שֶׁב עַֽל־הַבְּאֵֽר׃ \Rashi{\textbf{וישמע פרעה.}הם הלשינו עליו (שמות רבה א'): ויבקש להרג את משה. מסרו לקוסטינר להרגו ולא שלטה בו החרב (שם), הוא שאמר משה "ויצלני מחרב פרעה" (שמות י״ח:ד׳): (וישב בארץ מדין. נתעכב שם; כמו: "וישב יעקב" (בראשית ל"ז):) וישב על הבאר. לשון ישיבה. למד מיעקב שנזדוג לו זווגו מן הבאר (שמות רבה א'):}%
% ---------- פרק 2 | פסוק 16 ----------
 \textbf{טז} וּלְכֹהֵ֥ן מִדְיָ֖ן שֶׁ֣בַע בָּנ֑וֹת וַתָּבֹ֣אנָה וַתִּדְלֶ֗נָה וַתְּמַלֶּ֙אנָה֙ אֶת־הָ֣רְהָטִ֔ים לְהַשְׁק֖וֹת צֹ֥אן אֲבִיהֶֽן׃ \Rashi{\textbf{ולכהן מדין.}רב שבהן; ופרש לו מע"ז ונדוהו מאצלם (שם): את הרהטים. את ברכות מרוצות המים העשויות בארץ:}%
% ---------- פרק 2 | פסוק 17 ----------
 \textbf{יז} וַיָּבֹ֥אוּ הָרֹעִ֖ים וַיְגָרְשׁ֑וּם וַיָּ֤קׇם מֹשֶׁה֙ וַיּ֣וֹשִׁעָ֔ן וַיַּ֖שְׁקְ אֶת־צֹאנָֽם׃ \Rashi{\textbf{ויגרשום.}מפני הנדוי (שם):}%
% ---------- פרק 2 | פסוק 18 ----------
 \textbf{יח} וַתָּבֹ֕אנָה אֶל־רְעוּאֵ֖ל אֲבִיהֶ֑ן וַיֹּ֕אמֶר מַדּ֛וּעַ מִהַרְתֶּ֥ן בֹּ֖א הַיּֽוֹם׃ %
% ---------- פרק 2 | פסוק 19 ----------
 \textbf{יט} וַתֹּאמַ֕רְןָ אִ֣ישׁ מִצְרִ֔י הִצִּילָ֖נוּ מִיַּ֣ד הָרֹעִ֑ים וְגַם־דָּלֹ֤ה דָלָה֙ לָ֔נוּ וַיַּ֖שְׁקְ אֶת־הַצֹּֽאן׃ %
% ---------- פרק 2 | פסוק 20 ----------
 \textbf{כ} וַיֹּ֥אמֶר אֶל־בְּנֹתָ֖יו וְאַיּ֑וֹ לָ֤מָּה זֶּה֙ עֲזַבְתֶּ֣ן אֶת־הָאִ֔ישׁ קִרְאֶ֥ן ל֖וֹ וְיֹ֥אכַל לָֽחֶם׃ \Rashi{\textbf{למה זה עזבתן.}הכיר בו שהוא מזרעו של יעקב, שהמים עולים לקראתו (שמות רבה א'): ויאכל לחם. שמא ישא אחת מכם; כמה דאת אמר: "כי אם הלחם אשר הוא אוכל" (בראשית ל"ט):}%
% ---------- פרק 2 | פסוק 21 ----------
 \textbf{כא} וַיּ֥וֹאֶל מֹשֶׁ֖ה לָשֶׁ֣בֶת אֶת־הָאִ֑ישׁ וַיִּתֵּ֛ן אֶת־צִפֹּרָ֥ה בִתּ֖וֹ לְמֹשֶֽׁה׃ \Rashi{\textbf{ויואל.}כתרגומו, ודומה לו "הואל נא ולין" (שופטים י"ט), "ולו הואלנו" (יהושע ז'), "הואלתי לדבר" (בראשית י"ח). ומדרשו: לשון אלה – נשבע לו שלא יזוז ממדין כי אם ברשותו (נדרים ס"ה):}%
% ---------- פרק 2 | פסוק 22 ----------
 \textbf{כב} וַתֵּ֣לֶד בֵּ֔ן וַיִּקְרָ֥א אֶת־שְׁמ֖וֹ גֵּרְשֹׁ֑ם כִּ֣י אָמַ֔ר גֵּ֣ר הָיִ֔יתִי בְּאֶ֖רֶץ נׇכְרִיָּֽה׃ \{פ\} %
% ---------- פרק 2 | פסוק 23 ----------
 \textbf{כג} וַיְהִי֩ בַיָּמִ֨ים הָֽרַבִּ֜ים הָהֵ֗ם וַיָּ֙מׇת֙ מֶ֣לֶךְ מִצְרַ֔יִם וַיֵּאָנְח֧וּ בְנֵֽי־יִשְׂרָאֵ֛ל מִן־הָעֲבֹדָ֖ה וַיִּזְעָ֑קוּ וַתַּ֧עַל שַׁוְעָתָ֛ם אֶל־הָאֱלֹהִ֖ים מִן־הָעֲבֹדָֽה׃ \Rashi{\textbf{ויהי בימים הרבים ההם.}שהיה משה גר במדין, וימת מלך מצרים, והצרכו ישראל לתשועה, ומשה היה רעה (שמות ג׳:א׳), ובאת תשועה על ידו, לכך נסמכו פרשיות הללו: וימת מלך מצרים. נצטרע והיה שוחט תינוקות ישראל ורוחץ בדמם (שמות רבה א'):}%
% ---------- פרק 2 | פסוק 24 ----------
 \textbf{כד} וַיִּשְׁמַ֥ע אֱלֹהִ֖ים אֶת־נַאֲקָתָ֑ם וַיִּזְכֹּ֤ר אֱלֹהִים֙ אֶת־בְּרִית֔וֹ אֶת־אַבְרָהָ֖ם אֶת־יִצְחָ֥ק וְאֶֽת־יַעֲקֹֽב׃ \Rashi{\textbf{נאקתם.}צעקתם, וכן "מעיר מתים ינאקו" (איוב כ"ד): את בריתו את אברהם. עם אברהם:}%
% ---------- פרק 2 | פסוק 25 ----------
 \textbf{כה} וַיַּ֥רְא אֱלֹהִ֖ים אֶת־בְּנֵ֣י יִשְׂרָאֵ֑ל וַיֵּ֖דַע אֱלֹהִֽים׃ \{ס\} \Rashi{\textbf{וידע אלהים.}נתן עליהם לב ולא העלים עיניו:}%
% ---------- פרק 3 | פסוק 1 ----------
 \textbf{א} וּמֹשֶׁ֗ה הָיָ֥ה רֹעֶ֛ה אֶת־צֹ֛אן יִתְר֥וֹ חֹתְנ֖וֹ כֹּהֵ֣ן מִדְיָ֑ן וַיִּנְהַ֤ג אֶת־הַצֹּאן֙ אַחַ֣ר הַמִּדְבָּ֔ר וַיָּבֹ֛א אֶל־הַ֥ר הָאֱלֹהִ֖ים חֹרֵֽבָה׃ \Rashi{\textbf{אחר המדבר.}להתרחק מן הגזל, שלא ירעו בשדות אחרים (שמות רבה ב'): אל הר האלהים. על שם העתיד:}%
% ---------- פרק 3 | פסוק 2 ----------
 \textbf{ב} וַ֠יֵּרָ֠א מַלְאַ֨ךְ יְהֹוָ֥ה אֵלָ֛יו בְּלַבַּת־אֵ֖שׁ מִתּ֣וֹךְ הַסְּנֶ֑ה וַיַּ֗רְא וְהִנֵּ֤ה הַסְּנֶה֙ בֹּעֵ֣ר בָּאֵ֔שׁ וְהַסְּנֶ֖ה אֵינֶ֥נּוּ אֻכָּֽל׃ \Rashi{\textbf{בלבת אש.}בשלהבת אש – לבו של אש, כמו "לב השמים" (דברים ד'), "בלב האלה" (שמואל ב י"ח); ואל תתמה על התי"ו, שיש לנו כיוצא בו, "מה אמלה לבתך" (יחזקאל ט"ז): מתוך הסנה. ולא אילן אחר, משום "עמו אנכי בצרה" (תהילים צא טו): אכל. נאכל, כמו "לא עבד בה" (דברים כ"א), "אשר לקח משם" (בראשית ג'):}%
% ---------- פרק 3 | פסוק 3 ----------
 \textbf{ג} וַיֹּ֣אמֶר מֹשֶׁ֔ה אָסֻֽרָה־נָּ֣א וְאֶרְאֶ֔ה אֶת־הַמַּרְאֶ֥ה הַגָּדֹ֖ל הַזֶּ֑ה מַדּ֖וּעַ לֹא־יִבְעַ֥ר הַסְּנֶֽה׃ \Rashi{\textbf{אסרה נא.}אסורה מכאן להתקרב שם:}%
% ---------- פרק 3 | פסוק 4 ----------
 \textbf{ד} וַיַּ֥רְא יְהֹוָ֖ה כִּ֣י סָ֣ר לִרְא֑וֹת וַיִּקְרָא֩ אֵלָ֨יו אֱלֹהִ֜ים מִתּ֣וֹךְ הַסְּנֶ֗ה וַיֹּ֛אמֶר מֹשֶׁ֥ה מֹשֶׁ֖ה וַיֹּ֥אמֶר הִנֵּֽנִי׃ %
% ---------- פרק 3 | פסוק 5 ----------
 \textbf{ה} וַיֹּ֖אמֶר אַל־תִּקְרַ֣ב הֲלֹ֑ם שַׁל־נְעָלֶ֙יךָ֙ מֵעַ֣ל רַגְלֶ֔יךָ כִּ֣י הַמָּק֗וֹם אֲשֶׁ֤ר אַתָּה֙ עוֹמֵ֣ד עָלָ֔יו אַדְמַת־קֹ֖דֶשׁ הֽוּא׃ \Rashi{\textbf{של.}שלף והוצא, כמו "ונשל הברזל" (דברים י"ט), "כי ישל זיתך" (שם כ"ח): אדמת קדש הוא. המקום:}%
% ---------- פרק 3 | פסוק 6 ----------
 \textbf{ו} וַיֹּ֗אמֶר אָנֹכִי֙ אֱלֹהֵ֣י אָבִ֔יךָ אֱלֹהֵ֧י אַבְרָהָ֛ם אֱלֹהֵ֥י יִצְחָ֖ק וֵאלֹהֵ֣י יַעֲקֹ֑ב וַיַּסְתֵּ֤ר מֹשֶׁה֙ פָּנָ֔יו כִּ֣י יָרֵ֔א מֵהַבִּ֖יט אֶל־הָאֱלֹהִֽים׃ %
% ---------- פרק 3 | פסוק 7 ----------
 \textbf{ז} וַיֹּ֣אמֶר יְהֹוָ֔ה רָאֹ֥ה רָאִ֛יתִי אֶת־עֳנִ֥י עַמִּ֖י אֲשֶׁ֣ר בְּמִצְרָ֑יִם וְאֶת־צַעֲקָתָ֤ם שָׁמַ֙עְתִּי֙ מִפְּנֵ֣י נֹֽגְשָׂ֔יו כִּ֥י יָדַ֖עְתִּי אֶת־מַכְאֹבָֽיו׃ \Rashi{\textbf{כי ידעתי את מכאביו.}כמו "וידע אלהים" (שמות ב'), כלומר כי שמתי לב להתבונן ולדעת את מכאוביו ולא העלמתי עיני לאטם אזני מצעקתם:}%
% ---------- פרק 3 | פסוק 8 ----------
 \textbf{ח} וָאֵרֵ֞ד לְהַצִּיל֣וֹ ׀ מִיַּ֣ד מִצְרַ֗יִם וּֽלְהַעֲלֹתוֹ֮ מִן־הָאָ֣רֶץ הַהִוא֒ אֶל־אֶ֤רֶץ טוֹבָה֙ וּרְחָבָ֔ה אֶל־אֶ֛רֶץ זָבַ֥ת חָלָ֖ב וּדְבָ֑שׁ אֶל־מְק֤וֹם הַֽכְּנַעֲנִי֙ וְהַ֣חִתִּ֔י וְהָֽאֱמֹרִי֙ וְהַפְּרִזִּ֔י וְהַחִוִּ֖י וְהַיְבוּסִֽי׃ %
% ---------- פרק 3 | פסוק 9 ----------
 \textbf{ט} וְעַתָּ֕ה הִנֵּ֛ה צַעֲקַ֥ת בְּנֵי־יִשְׂרָאֵ֖ל בָּ֣אָה אֵלָ֑י וְגַם־רָאִ֙יתִי֙ אֶת־הַלַּ֔חַץ אֲשֶׁ֥ר מִצְרַ֖יִם לֹחֲצִ֥ים אֹתָֽם׃ %
% ---------- פרק 3 | פסוק 10 ----------
 \textbf{י} וְעַתָּ֣ה לְכָ֔ה וְאֶֽשְׁלָחֲךָ֖ אֶל־פַּרְעֹ֑ה וְהוֹצֵ֛א אֶת־עַמִּ֥י בְנֵֽי־יִשְׂרָאֵ֖ל מִמִּצְרָֽיִם׃ \Rashi{\textbf{ועתה לכה ואשלחך אל פרעה.}ואם תאמר מה תועיל? והוצא את עמי. יועילו דבריך ותוציאם משם:}%
% ---------- פרק 3 | פסוק 11 ----------
 \textbf{יא} וַיֹּ֤אמֶר מֹשֶׁה֙ אֶל־הָ֣אֱלֹהִ֔ים מִ֣י אָנֹ֔כִי כִּ֥י אֵלֵ֖ךְ אֶל־פַּרְעֹ֑ה וְכִ֥י אוֹצִ֛יא אֶת־בְּנֵ֥י יִשְׂרָאֵ֖ל מִמִּצְרָֽיִם׃ \Rashi{\textbf{מי אנכי.}מה אני חשוב לדבר עם המלכים? וכי אוציא את בני ישראל. ואף אם חשוב אני, מה זכו ישראל שיעשה להם נס ואוציאם ממצרים?:}%
% ---------- פרק 3 | פסוק 12 ----------
 \textbf{יב} וַיֹּ֙אמֶר֙ כִּֽי־אֶֽהְיֶ֣ה עִמָּ֔ךְ וְזֶה־לְּךָ֣ הָא֔וֹת כִּ֥י אָנֹכִ֖י שְׁלַחְתִּ֑יךָ בְּהוֹצִֽיאֲךָ֤ אֶת־הָעָם֙ מִמִּצְרַ֔יִם תַּֽעַבְדוּן֙ אֶת־הָ֣אֱלֹהִ֔ים עַ֖ל הָהָ֥ר הַזֶּֽה׃ \Rashi{\textbf{ויאמר כי אהיה עמך.}השיבו על ראשון ראשון ועל אחרון אחרון; שאמרת מי אנכי כי אלך אל פרעה – לא שלך היא כי אם משלי, כי אהיה עמך – וזה המראה אשר ראית בסנה לך האות כי אנכי שלחתיך, וכדאי אני להציל, כאשר ראית הסנה עושה שליחותי ואיננו אכל, כך תלך בשליחותי ואינך נזוק; וששאלת מה זכות יש לישראל שיצאו ממצרים? דבר גדול יש לי על הוצאה זו, שהרי עתידים לקבל התורה על ההר הזה לסוף ג' חדשים שיצאו ממצרים. ד"א: כי אהיה עמך, וזה שתצליח בשליחותך לך האות על הבטחה אחרת, שאני מבטיחך שכשתוציאם ממצרים תעבדון אותי על ההר הזה, שתקבלו התורה עליו, והיא הזכות העומדת לישראל. ודגמת לשון זה מצינו: "וזה לך האות אכול השנה ספיח וגו'" (ישעיהו ל"ז), מפלת סנחריב תהיה לך אות על הבטחה אחרת, שארצכם חרבה מפרות ואני אברך הספיחים:}%
% ---------- פרק 3 | פסוק 13 ----------
 \textbf{יג} וַיֹּ֨אמֶר מֹשֶׁ֜ה אֶל־הָֽאֱלֹהִ֗ים הִנֵּ֨ה אָנֹכִ֣י בָא֮ אֶל־בְּנֵ֣י יִשְׂרָאֵל֒ וְאָמַרְתִּ֣י לָהֶ֔ם אֱלֹהֵ֥י אֲבוֹתֵיכֶ֖ם שְׁלָחַ֣נִי אֲלֵיכֶ֑ם וְאָֽמְרוּ־לִ֣י מַה־שְּׁמ֔וֹ מָ֥ה אֹמַ֖ר אֲלֵהֶֽם׃ %
% ---------- פרק 3 | פסוק 14 ----------
 \textbf{יד} וַיֹּ֤אמֶר אֱלֹהִים֙ אֶל־מֹשֶׁ֔ה אֶֽהְיֶ֖ה אֲשֶׁ֣ר אֶֽהְיֶ֑ה וַיֹּ֗אמֶר כֹּ֤ה תֹאמַר֙ לִבְנֵ֣י יִשְׂרָאֵ֔ל אֶֽהְיֶ֖ה שְׁלָחַ֥נִי אֲלֵיכֶֽם׃ \Rashi{\textbf{אהיה אשר אהיה.}אהיה עמם בצרה זו אשר אהיה עמם בשעבוד שאר מלכיות. אמר לפניו, רבונו של עולם, מה אני מזכיר להם צרה אחרת? דים בצרה זו, אמר לו יפה אמרת, כה תאמר וגו' (ברכות ט'):}%
% ---------- פרק 3 | פסוק 15 ----------
 \textbf{טו} וַיֹּ֩אמֶר֩ ע֨וֹד אֱלֹהִ֜ים אֶל־מֹשֶׁ֗ה כֹּֽה־תֹאמַר֮ אֶל־בְּנֵ֣י יִשְׂרָאֵל֒ יְהֹוָ֞ה אֱלֹהֵ֣י אֲבֹתֵיכֶ֗ם אֱלֹהֵ֨י אַבְרָהָ֜ם אֱלֹהֵ֥י יִצְחָ֛ק וֵאלֹהֵ֥י יַעֲקֹ֖ב שְׁלָחַ֣נִי אֲלֵיכֶ֑ם זֶה־שְּׁמִ֣י לְעֹלָ֔ם וְזֶ֥ה זִכְרִ֖י לְדֹ֥ר דֹּֽר׃ \Rashi{\textbf{זה שמי לעלם.}חסר וי"ו, לומר, העלימהו – שלא יקרא ככתבו (פסחים נ'): וזה זכרי. למדו היאך נקרא. וכן דוד הוא אומר, "ה' שמך לעולם ה' זכרך לדר ודר" (תהלים קל"ה):}%
% ---------- פרק 3 | פסוק 16 ----------
 \textbf{טז} לֵ֣ךְ וְאָֽסַפְתָּ֞ אֶת־זִקְנֵ֣י יִשְׂרָאֵ֗ל וְאָמַרְתָּ֤ אֲלֵהֶם֙ יְהֹוָ֞ה אֱלֹהֵ֤י אֲבֹֽתֵיכֶם֙ נִרְאָ֣ה אֵלַ֔י אֱלֹהֵ֧י אַבְרָהָ֛ם יִצְחָ֥ק וְיַעֲקֹ֖ב לֵאמֹ֑ר פָּקֹ֤ד פָּקַ֙דְתִּי֙ אֶתְכֶ֔ם וְאֶת־הֶעָשׂ֥וּי לָכֶ֖ם בְּמִצְרָֽיִם׃ \Rashi{\textbf{את זקני ישראל.}מיחדים לישיבה. ואם תאמר זקנים סתם, היאך אפשר לו לאסף זקנים של ששים רבוא? (יומא כ"ח)}%
% ---------- פרק 3 | פסוק 17 ----------
 \textbf{יז} וָאֹמַ֗ר אַעֲלֶ֣ה אֶתְכֶם֮ מֵעֳנִ֣י מִצְרַ֒יִם֒ אֶל־אֶ֤רֶץ הַֽכְּנַעֲנִי֙ וְהַ֣חִתִּ֔י וְהָֽאֱמֹרִי֙ וְהַפְּרִזִּ֔י וְהַחִוִּ֖י וְהַיְבוּסִ֑י אֶל־אֶ֛רֶץ זָבַ֥ת חָלָ֖ב וּדְבָֽשׁ׃ %
% ---------- פרק 3 | פסוק 18 ----------
 \textbf{יח} וְשָׁמְע֖וּ לְקֹלֶ֑ךָ וּבָאתָ֡ אַתָּה֩ וְזִקְנֵ֨י יִשְׂרָאֵ֜ל אֶל־מֶ֣לֶךְ מִצְרַ֗יִם וַאֲמַרְתֶּ֤ם אֵלָיו֙ יְהֹוָ֞ה אֱלֹהֵ֤י הָֽעִבְרִיִּים֙ נִקְרָ֣ה עָלֵ֔ינוּ וְעַתָּ֗ה נֵֽלְכָה־נָּ֞א דֶּ֣רֶךְ שְׁלֹ֤שֶׁת יָמִים֙ בַּמִּדְבָּ֔ר וְנִזְבְּחָ֖ה לַֽיהֹוָ֥ה אֱלֹהֵֽינוּ׃ \Rashi{\textbf{ושמעו לקלך.}מכיון שתאמר להם לשון זה, ישמעו לקולך, שכבר סימן זה מסור בידם מיעקב ומיוסף, שבלשון זה הם נגאלים, יעקב אמר להם "ואלהים פקד יפקד אתכם" (בראשית נ'), יוסף אמר להם "פקד יפקד אלהים אתכם" (שם): נקרה עלינו. לשון מקרה; וכן "ויקר אלהים", "ואנכי אקרה כה" (במדבר כ"ג) – אהא נקרה מאתו הלם: (אלהי העבריים. יו"ד יתירה, רמז לי' מכות. ברש"י ישן):}%
% ---------- פרק 3 | פסוק 19 ----------
 \textbf{יט} וַאֲנִ֣י יָדַ֔עְתִּי כִּ֠י לֹֽא־יִתֵּ֥ן אֶתְכֶ֛ם מֶ֥לֶךְ מִצְרַ֖יִם לַהֲלֹ֑ךְ וְלֹ֖א בְּיָ֥ד חֲזָקָֽה׃ \Rashi{\textbf{לא יתן אתכם מלך מצרים להלך.}אם אין אני מראה לו ידי החזקה; כלומר, כל עוד שאין אני מודיעו ידי החזקה, לא יתן אתכם להלך: לא יתן. לא ישבק כמו: "על כן לא נתתיך" (בראשית כ'), "ולא נתנו אלהים להרע עמדי" (שם ל"א), וכלן לשון נתינה הם. וי"מ ולא ביד חזקה — ולא בשביל שידו חזקה, "כי מאז אשלח את ידי והכיתי את מצרים וגו'", ומתרגמין אותו "ולא מן קדם דחילה תקיף". משמו של רבי יעקב ברבי מנחם נאמר לי:}%
% ---------- פרק 3 | פסוק 20 ----------
 \textbf{כ} וְשָׁלַחְתִּ֤י אֶת־יָדִי֙ וְהִכֵּיתִ֣י אֶת־מִצְרַ֔יִם בְּכֹל֙ נִפְלְאֹתַ֔י אֲשֶׁ֥ר אֶֽעֱשֶׂ֖ה בְּקִרְבּ֑וֹ וְאַחֲרֵי־כֵ֖ן יְשַׁלַּ֥ח אֶתְכֶֽם׃ %
% ---------- פרק 3 | פסוק 21 ----------
 \textbf{כא} וְנָתַתִּ֛י אֶת־חֵ֥ן הָֽעָם־הַזֶּ֖ה בְּעֵינֵ֣י מִצְרָ֑יִם וְהָיָה֙ כִּ֣י תֵֽלֵכ֔וּן לֹ֥א תֵלְכ֖וּ רֵיקָֽם׃ %
% ---------- פרק 3 | פסוק 22 ----------
 \textbf{כב} וְשָׁאֲלָ֨ה אִשָּׁ֤ה מִשְּׁכֶנְתָּהּ֙ וּמִגָּרַ֣ת בֵּיתָ֔הּ כְּלֵי־כֶ֛סֶף וּכְלֵ֥י זָהָ֖ב וּשְׂמָלֹ֑ת וְשַׂמְתֶּ֗ם עַל־בְּנֵיכֶם֙ וְעַל־בְּנֹ֣תֵיכֶ֔ם וְנִצַּלְתֶּ֖ם אֶת־מִצְרָֽיִם׃ \Rashi{\textbf{ומגרת ביתה.}מאותה שהיא גרה אתה בבית: ונצלתם. כתרגומו "ותרוקנון", וכן "וינצלו את מצרים" (שמות י"ב), "ויתנצלו בני ישראל את עדים" (שם ל"ג), והנו"ן בו יסוד; ומנחם חברו במחברת צד"י עם "ויצל אלהים את מקנה אביכם", "אשר הציל אלהים מאבינו" (בראשית ל"א), ולא יאמנו דבריו, כי אם לא היתה הנו"ן יסוד והיא נקודה בחירק, לא תהא משמשת בלשון ופעלתם אלא בלשון ונפעלתם, כמו "ונסחתם מעל האדמה" (דברים כ"ח), "ונתתם ביד אויב", ונגפתם לפני אויביכם (ויקרא כ"ו), "ונתכתם בתוכה" (יחזקאל כ"ב), "ואמרתם נצלנו" (ירמיהו ז'), לשון נפעלנו, וכל נו"ן שהיא באה בתבה לפרקים ונופלת ממנה – כנו"ן של נוגף, נושא, נותן, נושך – כשהיא מדברת לשון ופעלתם תנקד בחטף, כגון "ונשאתם את אביכם" (בראשית מ"ה), "ונתתם להם את ארץ הגלעד" (במדבר ל"ב), "ונמלתם את בשר ערלתכם" (בראשית י"ז), לכך אני אומר שזאת הנקודה בחירק מן היסוד היא, ויסוד שם דבר נצול, והוא מן הלשונות הכבדים, כמו דבור, כפור, למוד, כשידבר בלשון ופעלתם, ינקד בחירק כמו "ודברתם אל הסלע" (במדבר כ'), "וכפרתם את הבית" (יחזקאל מ"ה), "ולמדתם אותם את בניכם" (דברים י"א):}%
% ---------- פרק 4 | פסוק 1 ----------
 \textbf{א} וַיַּ֤עַן מֹשֶׁה֙ וַיֹּ֔אמֶר וְהֵן֙ לֹֽא־יַאֲמִ֣ינוּ לִ֔י וְלֹ֥א יִשְׁמְע֖וּ בְּקֹלִ֑י כִּ֣י יֹֽאמְר֔וּ לֹֽא־נִרְאָ֥ה אֵלֶ֖יךָ יְהֹוָֽה׃ %
% ---------- פרק 4 | פסוק 2 ----------
 \textbf{ב} וַיֹּ֧אמֶר אֵלָ֛יו יְהֹוָ֖ה (מזה) [מַה־זֶּ֣ה] בְיָדֶ֑ךָ וַיֹּ֖אמֶר מַטֶּֽה׃ \Rashi{\textbf{מזה בידך.}לכך נכתב תבה אחת, לדרש: מזה שבידך אתה חיב ללקות, שחשדת בכשרים; ופשוטו: כאדם שאומר לחברו מודה אתה שזו שלפניך אבן היא? אומר לו הן; אמר לו הריני עושה אותה עץ:}%
% ---------- פרק 4 | פסוק 3 ----------
 \textbf{ג} וַיֹּ֙אמֶר֙ הַשְׁלִיכֵ֣הוּ אַ֔רְצָה וַיַּשְׁלִכֵ֥הוּ אַ֖רְצָה וַיְהִ֣י לְנָחָ֑שׁ וַיָּ֥נׇס מֹשֶׁ֖ה מִפָּנָֽיו׃ \Rashi{\textbf{ויהי לנחש.}רמז לו שספר לשון הרע על ישראל ותפש אמנותו של נחש (שמות רבה):}%
% ---------- פרק 4 | פסוק 4 ----------
 \textbf{ד} וַיֹּ֤אמֶר יְהֹוָה֙ אֶל־מֹשֶׁ֔ה שְׁלַח֙ יָֽדְךָ֔ וֶאֱחֹ֖ז בִּזְנָב֑וֹ וַיִּשְׁלַ֤ח יָדוֹ֙ וַיַּ֣חֲזֶק בּ֔וֹ וַיְהִ֥י לְמַטֶּ֖ה בְּכַפּֽוֹ׃ \Rashi{\textbf{ויחזק בו.}לשון אחיזה הוא; והרבה יש במקרא, "ויחזיקו האנשים בידו" (בראשית י"ט), "והחזיקה במבשיו" (דברים כ"ה), "והחזקתי בזקנו" (שמואל א י"ז), (כל לשון חזוק הדבוק לבי"ת לשון אחיזה הוא):}%
% ---------- פרק 4 | פסוק 5 ----------
 \textbf{ה} לְמַ֣עַן יַאֲמִ֔ינוּ כִּֽי־נִרְאָ֥ה אֵלֶ֛יךָ יְהֹוָ֖ה אֱלֹהֵ֣י אֲבֹתָ֑ם אֱלֹהֵ֧י אַבְרָהָ֛ם אֱלֹהֵ֥י יִצְחָ֖ק וֵאלֹהֵ֥י יַעֲקֹֽב׃ %
% ---------- פרק 4 | פסוק 6 ----------
 \textbf{ו} וַיֹּ֩אמֶר֩ יְהֹוָ֨ה ל֜וֹ ע֗וֹד הָֽבֵא־נָ֤א יָֽדְךָ֙ בְּחֵיקֶ֔ךָ וַיָּבֵ֥א יָד֖וֹ בְּחֵיק֑וֹ וַיּ֣וֹצִאָ֔הּ וְהִנֵּ֥ה יָד֖וֹ מְצֹרַ֥עַת כַּשָּֽׁלֶג׃ \Rashi{\textbf{מצרעת כשלג.}דרך צרעת להיות לבנה – "ואם בהרת לבנה היא" (ויקרא י"ג); אף באות זה רמז לו שלשון הרע ספר באמרו לא יאמינו לי, לפיכך הלקהו בצרעת, כמו שלקתה מרים על לשון הרע (שבת צ"ז):}%
% ---------- פרק 4 | פסוק 7 ----------
 \textbf{ז} וַיֹּ֗אמֶר הָשֵׁ֤ב יָֽדְךָ֙ אֶל־חֵיקֶ֔ךָ וַיָּ֥שֶׁב יָד֖וֹ אֶל־חֵיק֑וֹ וַיּֽוֹצִאָהּ֙ מֵֽחֵיק֔וֹ וְהִנֵּה־שָׁ֖בָה כִּבְשָׂרֽוֹ׃ \Rashi{\textbf{ויוצאה מחיקו והנה שבה כבשרו.}מכאן שמדה טובה ממהרת לבא ממדת פרענות, שהרי בראשונה לא נאמר "מחיקו" (שם):}%
% ---------- פרק 4 | פסוק 8 ----------
 \textbf{ח} וְהָיָה֙ אִם־לֹ֣א יַאֲמִ֣ינוּ לָ֔ךְ וְלֹ֣א יִשְׁמְע֔וּ לְקֹ֖ל הָאֹ֣ת הָרִאשׁ֑וֹן וְהֶֽאֱמִ֔ינוּ לְקֹ֖ל הָאֹ֥ת הָאַחֲרֽוֹן׃ \Rashi{\textbf{והאמינו לקל האת האחרון.}משתאמר להם בשבילכם לקיתי, על שספרתי עליכם לשון הרע, יאמינו לך, שכבר למדו בכך שהמזדוגין להרע להם לוקים בנגעים, כגון פרעה ואבימלך בשביל שרה:}%
% ---------- פרק 4 | פסוק 9 ----------
 \textbf{ט} וְהָיָ֡ה אִם־לֹ֣א יַאֲמִ֡ינוּ גַּם֩ לִשְׁנֵ֨י הָאֹת֜וֹת הָאֵ֗לֶּה וְלֹ֤א יִשְׁמְעוּן֙ לְקֹלֶ֔ךָ וְלָקַחְתָּ֙ מִמֵּימֵ֣י הַיְאֹ֔ר וְשָׁפַכְתָּ֖ הַיַּבָּשָׁ֑ה וְהָי֤וּ הַמַּ֙יִם֙ אֲשֶׁ֣ר תִּקַּ֣ח מִן־הַיְאֹ֔ר וְהָי֥וּ לְדָ֖ם בַּיַּבָּֽשֶׁת׃ \Rashi{\textbf{ולקחת ממימי היאור.}רמז להם שבמכה ראשונה נפרע מאלהותם, (פירוש, כשהקב"ה נפרע מן האומות, נפרע מאלהותם תחלה, שהיו עובדים לנילוס המחיה אותם, והפכם לדם. ברש"י ישן): והיו המים וגו'. "והיו" "והיו", ב' פעמים נראה בעיני, אלו נאמר והיו המים אשר תקח מן היאור לדם ביבשת, שומע אני, שבידו הם נהפכים לדם – ואף כשירדו לארץ יהיו בהויתן, אבל עכשו מלמדנו שלא יהיו דם עד שיהיו ביבשת:}%
% ---------- פרק 4 | פסוק 10 ----------
 \textbf{י} וַיֹּ֨אמֶר מֹשֶׁ֣ה אֶל־יְהֹוָה֮ בִּ֣י אֲדֹנָי֒ לֹא֩ אִ֨ישׁ דְּבָרִ֜ים אָנֹ֗כִי גַּ֤ם מִתְּמוֹל֙ גַּ֣ם מִשִּׁלְשֹׁ֔ם גַּ֛ם מֵאָ֥ז דַּבֶּרְךָ֖ אֶל־עַבְדֶּ֑ךָ כִּ֧י כְבַד־פֶּ֛ה וּכְבַ֥ד לָשׁ֖וֹן אָנֹֽכִי׃ \Rashi{\textbf{גם מתמול וגו'.}למדנו שכל ז' ימים היה הקב"ה מפתה את משה בסנה לילך בשליחותו, "מתמול", "שלשם", "מאז דברך" הרי ג' ושלושה גמין רבויין הם, הרי ששה, והוא היה עומד ביום השביעי כשאמר לו זאת עוד שלח נא ביד תשלח, עד שחרה בו וקבל עליו (שמות רבה); וכל זה שלא היה רוצה לטל גדלה על אהרן אחיו שהיה גדול הימנו ונביא היה, שנאמר "הנגלה נגליתי אל בית אביך בהיותם במצרים" (שמואל א ב') – הוא אהרן, וכן "ואודע להם בארץ מצרים … ואמר אליהם איש שקוצי עיניו השליכו" (יחזקאל כ'), ואותה נבואה לאהרן נאמרה (תנחומא): כבד פה. בכבדות אני מדבר. ובלשון לעז בלב"ו:}%
% ---------- פרק 4 | פסוק 11 ----------
 \textbf{יא} וַיֹּ֨אמֶר יְהֹוָ֜ה אֵלָ֗יו מִ֣י שָׂ֣ם פֶּה֮ לָֽאָדָם֒ א֚וֹ מִֽי־יָשׂ֣וּם אִלֵּ֔ם א֣וֹ חֵרֵ֔שׁ א֥וֹ פִקֵּ֖חַ א֣וֹ עִוֵּ֑ר הֲלֹ֥א אָנֹכִ֖י יְהֹוָֽה׃ \Rashi{\textbf{מי שם פה וגו'.}מי למדך לדבר, כשהיית נדון לפני פרעה על המצרי? או מי ישום אלם. מי עשה פרעה אלם, שלא נתאמץ במצות הריגתך, ואת משרתיו חרשים שלא שמעו בצוותו עליך, ואת אספקלטורין ההורגים מי עשאם עורים שלא ראו כשברחת מן הבימה ונמלטת? הלא אנכי. ששמי ה' עשיתי כל זאת (שם):}%
% ---------- פרק 4 | פסוק 12 ----------
 \textbf{יב} וְעַתָּ֖ה לֵ֑ךְ וְאָנֹכִי֙ אֶֽהְיֶ֣ה עִם־פִּ֔יךָ וְהוֹרֵיתִ֖יךָ אֲשֶׁ֥ר תְּדַבֵּֽר׃ %
% ---------- פרק 4 | פסוק 13 ----------
 \textbf{יג} וַיֹּ֖אמֶר בִּ֣י אֲדֹנָ֑י שְֽׁלַֽח־נָ֖א בְּיַד־תִּשְׁלָֽח׃ \Rashi{\textbf{ביד תשלח.}ביד מי שאתה רגיל לשלח, והוא אהרן. ד"א: ביד אחר שתרצה לשלח, אין סופי להכניסם לארץ ולהיות גואלם לעתיד, יש לך שלוחים הרבה:}%
% ---------- פרק 4 | פסוק 14 ----------
 \textbf{יד} וַיִּֽחַר־אַ֨ף יְהֹוָ֜ה בְּמֹשֶׁ֗ה וַיֹּ֙אמֶר֙ הֲלֹ֨א אַהֲרֹ֤ן אָחִ֙יךָ֙ הַלֵּוִ֔י יָדַ֕עְתִּי כִּֽי־דַבֵּ֥ר יְדַבֵּ֖ר ה֑וּא וְגַ֤ם הִנֵּה־הוּא֙ יֹצֵ֣א לִקְרָאתֶ֔ךָ וְרָאֲךָ֖ וְשָׂמַ֥ח בְּלִבּֽוֹ׃ \Rashi{\textbf{ויחר אף.}רבי יהושע בן קרחה אומר כל חרון אף שבתורה עושה רשם וזה לא נאמר בו רשם ולא מצינו שבא ענש ע"י אותו חרון; אמר לו רבי יוסי אף בזו נאמר בו רשם – הלא אהרן אחיך הלוי – שהיה עתיד להיות לוי ולא כהן, והכהנה הייתי אומר לצאת ממך, מעתה לא יהיה כן אלא הוא יהיה כהן ואתה הלוי, שנאמר "ומשה איש האלהים בניו יקראו על שבט הלוי" (דהי"א כ"ג) (זבחים ק"ב): הנה הוא יצא לקראתך. כשתלך למצרים: וראך ושמח בלבו. לא כשאתה סבור, שיהא מקפיד עליך שאתה עולה לגדלה; ומשם זכה אהרן לעדי החשן הנתון על הלב (שבת קל"ט):}%
% ---------- פרק 4 | פסוק 15 ----------
 \textbf{טו} וְדִבַּרְתָּ֣ אֵלָ֔יו וְשַׂמְתָּ֥ אֶת־הַדְּבָרִ֖ים בְּפִ֑יו וְאָנֹכִ֗י אֶֽהְיֶ֤ה עִם־פִּ֙יךָ֙ וְעִם־פִּ֔יהוּ וְהוֹרֵיתִ֣י אֶתְכֶ֔ם אֵ֖ת אֲשֶׁ֥ר תַּעֲשֽׂוּן׃ %
% ---------- פרק 4 | פסוק 16 ----------
 \textbf{טז} וְדִבֶּר־ה֥וּא לְךָ֖ אֶל־הָעָ֑ם וְהָ֤יָה הוּא֙ יִֽהְיֶה־לְּךָ֣ לְפֶ֔ה וְאַתָּ֖ה תִּֽהְיֶה־לּ֥וֹ לֵֽאלֹהִֽים׃ \Rashi{\textbf{ודבר הוא לך.}בשבילך ידבר אל העם. וזה יוכיח על כל לך ולי ולו ולכם ולהם הסמוכים לדבור שכלם לשון על הם: יהיה לך לפה. למליץ, לפי שאתה כבד פה: לאלהים. לרב ולשר:}%
% ---------- פרק 4 | פסוק 17 ----------
 \textbf{יז} וְאֶת־הַמַּטֶּ֥ה הַזֶּ֖ה תִּקַּ֣ח בְּיָדֶ֑ךָ אֲשֶׁ֥ר תַּעֲשֶׂה־בּ֖וֹ אֶת־הָאֹתֹֽת׃ \{פ\} %
% ---------- פרק 4 | פסוק 18 ----------
 \textbf{יח} וַיֵּ֨לֶךְ מֹשֶׁ֜ה וַיָּ֣שׇׁב ׀ אֶל־יֶ֣תֶר חֹֽתְנ֗וֹ וַיֹּ֤אמֶר לוֹ֙ אֵ֣לְכָה נָּ֗א וְאָשׁ֙וּבָה֙ אֶל־אַחַ֣י אֲשֶׁר־בְּמִצְרַ֔יִם וְאֶרְאֶ֖ה הַעוֹדָ֣ם חַיִּ֑ים וַיֹּ֧אמֶר יִתְר֛וֹ לְמֹשֶׁ֖ה לֵ֥ךְ לְשָׁלֽוֹם׃ \Rashi{\textbf{וישב אל יתר חתנו.}לטל רשות, שהרי נשבע לו (נדרים ס"ה). וז' שמות היו לו, רעואל, יתר, יתרו, קיני, חובב, חבר, פוטיאל:}%
% ---------- פרק 4 | פסוק 19 ----------
 \textbf{יט} וַיֹּ֨אמֶר יְהֹוָ֤ה אֶל־מֹשֶׁה֙ בְּמִדְיָ֔ן לֵ֖ךְ שֻׁ֣ב מִצְרָ֑יִם כִּי־מֵ֙תוּ֙ כׇּל־הָ֣אֲנָשִׁ֔ים הַֽמְבַקְשִׁ֖ים אֶת־נַפְשֶֽׁךָ׃ \Rashi{\textbf{כי מתו כל האנשים.}מי הם? דתן ואבירם; חיים היו, אלא שירדו מנכסיהם והעני חשוב כמת:}%
% ---------- פרק 4 | פסוק 20 ----------
 \textbf{כ} וַיִּקַּ֨ח מֹשֶׁ֜ה אֶת־אִשְׁתּ֣וֹ וְאֶת־בָּנָ֗יו וַיַּרְכִּבֵם֙ עַֽל־הַחֲמֹ֔ר וַיָּ֖שׇׁב אַ֣רְצָה מִצְרָ֑יִם וַיִּקַּ֥ח מֹשֶׁ֛ה אֶת־מַטֵּ֥ה הָאֱלֹהִ֖ים בְּיָדֽוֹ׃ \Rashi{\textbf{על החמר.}חמור המיחד, הוא החמור שחבש אברהם לעקדת יצחק, והוא שעתיד מלך המשיח להגלות עליו, שנאמר "עני ורכב על חמור" (זכריה ט') (פרקי דרבי אליעזר ל"א): וישב ארצה מצרים, ויקח משה את מטה. אין מקדם ומאחר מדקדקים במקרא:}%
% ---------- פרק 4 | פסוק 21 ----------
 \textbf{כא} וַיֹּ֣אמֶר יְהֹוָה֮ אֶל־מֹשֶׁה֒ בְּלֶכְתְּךָ֙ לָשׁ֣וּב מִצְרַ֔יְמָה רְאֵ֗ה כׇּל־הַמֹּֽפְתִים֙ אֲשֶׁר־שַׂ֣מְתִּי בְיָדֶ֔ךָ וַעֲשִׂיתָ֖ם לִפְנֵ֣י פַרְעֹ֑ה וַאֲנִי֙ אֲחַזֵּ֣ק אֶת־לִבּ֔וֹ וְלֹ֥א יְשַׁלַּ֖ח אֶת־הָעָֽם׃ \Rashi{\textbf{בלכתך לשוב מצרימה וגו'.}דע שעל מנת כן תלך שתהא גבור בשליחותי לעשות כל מופתי לפני פרעה ולא תירא ממנו: אשר שמתי בידך. לא על שלוש אותות האמורות למעלה, שהרי לא לפני פרעה צוה לעשותם אלא לפני ישראל שיאמינו לו, ולא מצינו שעשאם לפניו, אלא מופתים שאני עתיד לשום בידך במצרים, כמו "כי ידבר אליכם פרעה וגו'" (שמות ז'); ואל תתמה על אשר כתיב "אשר שמתי", שכן משמעו, כשתדבר עמו כבר שמתים בידך:}%
% ---------- פרק 4 | פסוק 22 ----------
 \textbf{כב} וְאָמַרְתָּ֖ אֶל־פַּרְעֹ֑ה כֹּ֚ה אָמַ֣ר יְהֹוָ֔ה בְּנִ֥י בְכֹרִ֖י יִשְׂרָאֵֽל׃ \Rashi{\textbf{ואמרת אל פרעה.}כשתשמע שלבו חזק וימאן לשלח אמר לו כן: בני בכרי. לשון גדלה; כמו: "אף אני בכור אתנהו" (תהלים פ"ט), זהו פשוטו. ומדרשו: כאן חתם הקב"ה על מכירת הבכורה שלקח יעקב מעשו (בראשית רבה ס"ג):}%
% ---------- פרק 4 | פסוק 23 ----------
 \textbf{כג} וָאֹמַ֣ר אֵלֶ֗יךָ שַׁלַּ֤ח אֶת־בְּנִי֙ וְיַֽעַבְדֵ֔נִי וַתְּמָאֵ֖ן לְשַׁלְּח֑וֹ הִנֵּה֙ אָנֹכִ֣י הֹרֵ֔ג אֶת־בִּנְךָ֖ בְּכֹרֶֽךָ׃ \Rashi{\textbf{ואמר אליך.}בשליחותו של מקום שלח את בני: הנה אנכי הרג וגו'. היא מכה אחרונה, ובה התרהו תחלה מפני שהיא קשה; וזה הוא שאמר "הן אל ישגיב בכחו" – לפיכך "מי כמהו מורה" (איוב ל"ו); בשר ודם המבקש להנקם מחברו, מעלים את דבריו, שלא יבקש הצלה, אבל הקב"ה ישגיב בכחו ואין יכלת להמלט מידו כי אם בשובו אליו, לפיכך הוא מורהו ומתרה בו לשוב:}%
% ---------- פרק 4 | פסוק 24 ----------
 \textbf{כד} וַיְהִ֥י בַדֶּ֖רֶךְ בַּמָּל֑וֹן וַיִּפְגְּשֵׁ֣הוּ יְהֹוָ֔ה וַיְבַקֵּ֖שׁ הֲמִיתֽוֹ׃ \Rashi{\textbf{ויהי משה בדרך במלון: ויבקש המיתו.}(המלאך ל)משה: לפי שלא מל את אליעזר בנו, ועל שנתרשל נענש מיתה. תניא אמר רבי יוסי: חס ושלום, לא נתרשל אלא אמר, אמול ואצא לדרך, סכנה היא לתינוק עד שלושת ימים, אמול ואשהה ג' ימים, הקב"ה צוני לך שב מצרים ומפני מה נענש? לפי שנתעסק במלון תחלה (מכילתא). במסכת נדרים: והיה המלאך נעשה כמין נחש ובולעו מראשו ועד ירכיו וחוזר ובולעו מרגליו ועד אותו מקום, הבינה צפורה שבשביל המילה הוא (נדרים ל"ב):}%
% ---------- פרק 4 | פסוק 25 ----------
 \textbf{כה} וַתִּקַּ֨ח צִפֹּרָ֜ה צֹ֗ר וַתִּכְרֹת֙ אֶת־עׇרְלַ֣ת בְּנָ֔הּ וַתַּגַּ֖ע לְרַגְלָ֑יו וַתֹּ֕אמֶר כִּ֧י חֲתַן־דָּמִ֛ים אַתָּ֖ה לִֽי׃ \Rashi{\textbf{ותגע לרגליו.}השליכתו לפני רגליו של משה: ותאמר. על בנה: כי חתן דמים אתה לי. אתה הייתה גורם להיות החתן שלי נרצח עליך – הורג אישי אתה לי:}%
% ---------- פרק 4 | פסוק 26 ----------
 \textbf{כו} וַיִּ֖רֶף מִמֶּ֑נּוּ אָ֚ז אָֽמְרָ֔ה חֲתַ֥ן דָּמִ֖ים לַמּוּלֹֽת׃ \{פ\} \Rashi{\textbf{וירף.}המלאך ממנו, אז הבינה שעל המילה בא להרגו: אמרה חתן דמים למולת. חתני היה נרצח על דבר המילה: למולת. על דבר המולות. שם דבר הוא, והלמ"ד משמשת בלשון על, כמו "ואמר פרעה לבני ישראל" (שמות י"ד), ואנקלוס תרגם דמים על דם המילה:}%
% ---------- פרק 4 | פסוק 27 ----------
 \textbf{כז} וַיֹּ֤אמֶר יְהֹוָה֙ אֶֽל־אַהֲרֹ֔ן לֵ֛ךְ לִקְרַ֥את מֹשֶׁ֖ה הַמִּדְבָּ֑רָה וַיֵּ֗לֶךְ וַֽיִּפְגְּשֵׁ֛הוּ בְּהַ֥ר הָאֱלֹהִ֖ים וַיִּשַּׁק־לֽוֹ׃ %
% ---------- פרק 4 | פסוק 28 ----------
 \textbf{כח} וַיַּגֵּ֤ד מֹשֶׁה֙ לְאַֽהֲרֹ֔ן אֵ֛ת כׇּל־דִּבְרֵ֥י יְהֹוָ֖ה אֲשֶׁ֣ר שְׁלָח֑וֹ וְאֵ֥ת כׇּל־הָאֹתֹ֖ת אֲשֶׁ֥ר צִוָּֽהוּ׃ %
% ---------- פרק 4 | פסוק 29 ----------
 \textbf{כט} וַיֵּ֥לֶךְ מֹשֶׁ֖ה וְאַהֲרֹ֑ן וַיַּ֣אַסְפ֔וּ אֶת־כׇּל־זִקְנֵ֖י בְּנֵ֥י יִשְׂרָאֵֽל׃ %
% ---------- פרק 4 | פסוק 30 ----------
 \textbf{ל} וַיְדַבֵּ֣ר אַהֲרֹ֔ן אֵ֚ת כׇּל־הַדְּבָרִ֔ים אֲשֶׁר־דִּבֶּ֥ר יְהֹוָ֖ה אֶל־מֹשֶׁ֑ה וַיַּ֥עַשׂ הָאֹתֹ֖ת לְעֵינֵ֥י הָעָֽם׃ %
% ---------- פרק 4 | פסוק 31 ----------
 \textbf{לא} וַֽיַּאֲמֵ֖ן הָעָ֑ם וַֽיִּשְׁמְע֡וּ כִּֽי־פָקַ֨ד יְהֹוָ֜ה אֶת־בְּנֵ֣י יִשְׂרָאֵ֗ל וְכִ֤י רָאָה֙ אֶת־עׇנְיָ֔ם וַֽיִּקְּד֖וּ וַיִּֽשְׁתַּחֲוֽוּ׃ %
% ---------- פרק 5 | פסוק 1 ----------
 \textbf{א} וְאַחַ֗ר בָּ֚אוּ מֹשֶׁ֣ה וְאַהֲרֹ֔ן וַיֹּאמְר֖וּ אֶל־פַּרְעֹ֑ה כֹּֽה־אָמַ֤ר יְהֹוָה֙ אֱלֹהֵ֣י יִשְׂרָאֵ֔ל שַׁלַּח֙ אֶת־עַמִּ֔י וְיָחֹ֥גּוּ לִ֖י בַּמִּדְבָּֽר׃ \Rashi{\textbf{ואחר באו משה ואהרן וגו'.}אבל הזקנים נשמטו אחד אחד מאחר משה ואהרן, עד שנשמטו כלם קדם שהגיעו לפלטין, לפי שיראו ללכת; ובסיני נפרע להם, "ונגש משה לבדו אל ה' והם לא יגשו" (שמות כ"ד) – החזירם לאחוריהם (שמות רבה):}%
% ---------- פרק 5 | פסוק 2 ----------
 \textbf{ב} וַיֹּ֣אמֶר פַּרְעֹ֔ה מִ֤י יְהֹוָה֙ אֲשֶׁ֣ר אֶשְׁמַ֣ע בְּקֹל֔וֹ לְשַׁלַּ֖ח אֶת־יִשְׂרָאֵ֑ל לֹ֤א יָדַ֙עְתִּי֙ אֶת־יְהֹוָ֔ה וְגַ֥ם אֶת־יִשְׂרָאֵ֖ל לֹ֥א אֲשַׁלֵּֽחַ׃ %
% ---------- פרק 5 | פסוק 3 ----------
 \textbf{ג} וַיֹּ֣אמְר֔וּ אֱלֹהֵ֥י הָעִבְרִ֖ים נִקְרָ֣א עָלֵ֑ינוּ נֵ֣לְכָה נָּ֡א דֶּ֩רֶךְ֩ שְׁלֹ֨שֶׁת יָמִ֜ים בַּמִּדְבָּ֗ר וְנִזְבְּחָה֙ לַֽיהֹוָ֣ה אֱלֹהֵ֔ינוּ פֶּ֨ן־יִפְגָּעֵ֔נוּ בַּדֶּ֖בֶר א֥וֹ בֶחָֽרֶב׃ \Rashi{\textbf{פן יפגענו.}פן יפגעך היו צריכים לומר, אלא שחלקו כבוד למלכות. פגיעה זו לשון מקרה מות היא (שם):}%
% ---------- פרק 5 | פסוק 4 ----------
 \textbf{ד} וַיֹּ֤אמֶר אֲלֵהֶם֙ מֶ֣לֶךְ מִצְרַ֔יִם לָ֚מָּה מֹשֶׁ֣ה וְאַהֲרֹ֔ן תַּפְרִ֥יעוּ אֶת־הָעָ֖ם מִֽמַּעֲשָׂ֑יו לְכ֖וּ לְסִבְלֹתֵיכֶֽם׃ \Rashi{\textbf{תפריעו את העם ממעשיו.}תבדילו ותרחיקו אותם ממלאכתם, ששומעין לכם וסבורים לנוח מן המלאכה; וכן "פרעהו אל תעבר בו" (משלי ד') – רחקהו, וכן "ותפרעו כל עצתי" (שם א'), "כי פרע הוא" (שמות ל"ב) – נרחק ונתעב: לכו לסבלתיכם. לכו למלאכתכם שיש לכם לעשות בבתיכם, אבל מלאכת שעבוד מצרים לא היתה על שבטו של לוי; תדע לך, שהרי משה ואהרן יוצאים ובאים (שמות רבה):}%
% ---------- פרק 5 | פסוק 5 ----------
 \textbf{ה} וַיֹּ֣אמֶר פַּרְעֹ֔ה הֵן־רַבִּ֥ים עַתָּ֖ה עַ֣ם הָאָ֑רֶץ וְהִשְׁבַּתֶּ֥ם אֹתָ֖ם מִסִּבְלֹתָֽם׃ \Rashi{\textbf{הן רבים עתה עם הארץ.}שהעבודה מטלת עליהם, ואתם משביתים אותם מסבלותם, הפסד גדול הוא זה:}%
% ---------- פרק 5 | פסוק 6 ----------
 \textbf{ו} וַיְצַ֥ו פַּרְעֹ֖ה בַּיּ֣וֹם הַה֑וּא אֶת־הַנֹּגְשִׂ֣ים בָּעָ֔ם וְאֶת־שֹׁטְרָ֖יו לֵאמֹֽר׃ \Rashi{\textbf{הנגשים.}מצריים היו והשוטרים היו ישראלים; הנוגש ממנה על כמה שוטרים והשוטר ממנה לרדות בעושי המלאכה (שמות רבה):}%
% ---------- פרק 5 | פסוק 7 ----------
 \textbf{ז} לֹ֣א תֹאסִפ֞וּן לָתֵ֨ת תֶּ֧בֶן לָעָ֛ם לִלְבֹּ֥ן הַלְּבֵנִ֖ים כִּתְמ֣וֹל שִׁלְשֹׁ֑ם הֵ֚ם יֵֽלְכ֔וּ וְקֹשְׁשׁ֥וּ לָהֶ֖ם תֶּֽבֶן׃ \Rashi{\textbf{תבן.}אשטו"בלא, היו גובלין אותו עם הטיט: לבנים. טיוול"ש בלעז, שעושים מטיט ומיבשין אותן בחמה; ויש ששורפין אותן בכבשן: כתמול שלשום. כאשר הייתם עושים עד הנה; וקששו. ולקטו:}%
% ---------- פרק 5 | פסוק 8 ----------
 \textbf{ח} וְאֶת־מַתְכֹּ֨נֶת הַלְּבֵנִ֜ים אֲשֶׁ֣ר הֵם֩ עֹשִׂ֨ים תְּמ֤וֹל שִׁלְשֹׁם֙ תָּשִׂ֣ימוּ עֲלֵיהֶ֔ם לֹ֥א תִגְרְע֖וּ מִמֶּ֑נּוּ כִּֽי־נִרְפִּ֣ים הֵ֔ם עַל־כֵּ֗ן הֵ֤ם צֹֽעֲקִים֙ לֵאמֹ֔ר נֵלְכָ֖ה נִזְבְּחָ֥ה לֵאלֹהֵֽינוּ׃ \Rashi{\textbf{ואת מתכנת הלבנים.}סכום חשבון הלבנים שהיה כל אחד עושה ליום כשהיה התבן נתן להם, אותו סכום תשימו עליהם גם עתה, למען תכבד העבודה עליהם: כי נרפים. מן העבודה הם, לכך לבם פונה אל הבטלה וצועקים לאמר נלכה וגו': מתכנת. ותכן לבנים, "ולו נתכנו עלילות" (שמואל א ב'), "את הכסף המתכן" (מלכים ב י"ב), כלם לשון חשבון הם: נרפים. המלאכה רפויה בידם ועזובה מהם והם נרפים ממנה; רטריי"ט בלעז:}%
% ---------- פרק 5 | פסוק 9 ----------
 \textbf{ט} תִּכְבַּ֧ד הָעֲבֹדָ֛ה עַל־הָאֲנָשִׁ֖ים וְיַעֲשׂוּ־בָ֑הּ וְאַל־יִשְׁע֖וּ בְּדִבְרֵי־שָֽׁקֶר׃ \Rashi{\textbf{ואל ישעו בדברי שקר.}ואל יהגו וידברו תמיד בדברי רוח לאמר, נלכה נזבחה; ודומה לו "ואשעה בחקיך תמיד" (תהילים קי"ט); "למשל ולשנינה" (דברים כ"ח), מתרגמינן "ולשועי"; ויספר – "ואשתעי"; ואי אפשר לומר ישעו לשון "וישע ה' אל הבל", "ואל קין ואל מנחתו לא שעה" (בראשית ד'), ולפרש אל ישעו – אל יפנו, שאם כן היה לו לכתב ואל ישעו אל דברי שקר או לדברי שקר, כי כן גזרת כלם, "ישעה האדם על עושהו" (ישעיהו י"ז), "ולא ישעה אל המזבחות", (שם), "ולא שעו על קדוש ישראל" (שם ל"א), ולא מצאתי שמוש של בי"ת סמוכה לאחריהם, אבל אחר לשון דבור – כמתעסק לדבר בדבר – נופל לשון שמוש בי"ת, כגון: "הנדברים בך" (יחזקאל ל"ג), "ותדבר מרים ואהרן במשה" (במדבר י"ב), "המלאך הדובר בי" (זכריה ד'), "לדבר בם" (דברים י"א), "ואדברה בעדותיך" (תהילים קי"ט), אף כאן – אל ישעו בדברי שקר – אל יהיו נדברים בדברי שוא והבאי:}%
% ---------- פרק 5 | פסוק 10 ----------
 \textbf{י} וַיֵּ֨צְא֜וּ נֹגְשֵׂ֤י הָעָם֙ וְשֹׁ֣טְרָ֔יו וַיֹּאמְר֥וּ אֶל־הָעָ֖ם לֵאמֹ֑ר כֹּ֚ה אָמַ֣ר פַּרְעֹ֔ה אֵינֶ֛נִּי נֹתֵ֥ן לָכֶ֖ם תֶּֽבֶן׃ %
% ---------- פרק 5 | פסוק 11 ----------
 \textbf{יא} אַתֶּ֗ם לְכ֨וּ קְח֤וּ לָכֶם֙ תֶּ֔בֶן מֵאֲשֶׁ֖ר תִּמְצָ֑אוּ כִּ֣י אֵ֥ין נִגְרָ֛ע מֵעֲבֹדַתְכֶ֖ם דָּבָֽר׃ \Rashi{\textbf{אתם לכו קחו לכם תבן.}וצריכים אתם לילך בזריזות, כי אין נגרע מעבדתכם דבר, מכל סכום לבנים שהייתם עושים ליום בהיות התבן נתן לכם מזמן מבית המלך:}%
% ---------- פרק 5 | פסוק 12 ----------
 \textbf{יב} וַיָּ֥פֶץ הָעָ֖ם בְּכׇל־אֶ֣רֶץ מִצְרָ֑יִם לְקֹשֵׁ֥שׁ קַ֖שׁ לַתֶּֽבֶן׃ \Rashi{\textbf{לקשש קש לתבן.}לאסף אסיפה, ללקט לקט, לצרך תבן הטיט: קש. לשון לקוט, על שם שדבר המתפזר הוא וצריך לקוששו קרוי קש בשאר מקומות:}%
% ---------- פרק 5 | פסוק 13 ----------
 \textbf{יג} וְהַנֹּגְשִׂ֖ים אָצִ֣ים לֵאמֹ֑ר כַּלּ֤וּ מַעֲשֵׂיכֶם֙ דְּבַר־י֣וֹם בְּיוֹמ֔וֹ כַּאֲשֶׁ֖ר בִּהְי֥וֹת הַתֶּֽבֶן׃ \Rashi{\textbf{אצים.}דוחקים: דבר יום ביומו. חשבון של כל יום כלו ביומו, כאשר עשיתם בהיות התבן מוכן:}%
% ---------- פרק 5 | פסוק 14 ----------
 \textbf{יד} וַיֻּכּ֗וּ שֹֽׁטְרֵי֙ בְּנֵ֣י יִשְׂרָאֵ֔ל אֲשֶׁר־שָׂ֣מוּ עֲלֵהֶ֔ם נֹגְשֵׂ֥י פַרְעֹ֖ה לֵאמֹ֑ר מַדּ֡וּעַ לֹא֩ כִלִּיתֶ֨ם חׇקְכֶ֤ם לִלְבֹּן֙ כִּתְמ֣וֹל שִׁלְשֹׁ֔ם גַּם־תְּמ֖וֹל גַּם־הַיּֽוֹם׃ \Rashi{\textbf{ויכו שטרי בני ישראל.}השוטרים ישראלים היו וחסים על חבריהם מלדחקם, וכשהיו משלימים הלבנים לנוגשים שהם מצריים, והיה חסר מן הסכום, היו מלקין אותם על שלא דחקו את עושי המלאכה; לפיכך זכו אותם שוטרים להיות סנהדרין ונאצל מן הרוח אשר על משה והושם עליהם, שנאמר: "אספה לי שבעים איש מזקני ישראל וגו'" – מאותן שידעת הטובה שעשו במצרים – כי הם זקני העם ושוטריו (במדבר י"א), (שמות רבה): ויכו שטרי בני ישראל אשר שמו נגשי פרעה אותם לשוטרים עליהם, לאמר מדוע וגו'. למה ויכו? שהיו אומרים להם מדוע לא כליתם גם תמול גם היום חק הקצוב עליכם ללבן כתמול השלישי, שהוא יום שלפני אתמול; והוא היה בהיות התבן נתן להם: ויכו. לשון ויפעלו, הכו מיד אחרים – הנוגשים הכום:}%
% ---------- פרק 5 | פסוק 15 ----------
 \textbf{טו} וַיָּבֹ֗אוּ שֹֽׁטְרֵי֙ בְּנֵ֣י יִשְׂרָאֵ֔ל וַיִּצְעֲק֥וּ אֶל־פַּרְעֹ֖ה לֵאמֹ֑ר לָ֧מָּה תַעֲשֶׂ֦ה כֹ֖ה לַעֲבָדֶֽיךָ׃ %
% ---------- פרק 5 | פסוק 16 ----------
 \textbf{טז} תֶּ֗בֶן אֵ֤ין נִתָּן֙ לַעֲבָדֶ֔יךָ וּלְבֵנִ֛ים אֹמְרִ֥ים לָ֖נוּ עֲשׂ֑וּ וְהִנֵּ֧ה עֲבָדֶ֛יךָ מֻכִּ֖ים וְחָטָ֥את עַמֶּֽךָ׃ \Rashi{\textbf{ולבנים אומרים לנו.}הנוגשים עשו, כמנין הראשון: וחטאת עמך. אלו היה נקוד פתח, הייתי אומר שהוא דבוק, ודבר זה חטאת עמך הוא, עכשו שהוא קמץ, שם דבר הוא, וכך פרושו: ודבר זה מביא חטאת על עמך, כאלו כתוב וחטאת לעמך; כמו "כבאנה בית לחם" (רות א'), שהוא כמו לבית לחם, וכן הרבה:}%
% ---------- פרק 5 | פסוק 17 ----------
 \textbf{יז} וַיֹּ֛אמֶר נִרְפִּ֥ים אַתֶּ֖ם נִרְפִּ֑ים עַל־כֵּן֙ אַתֶּ֣ם אֹֽמְרִ֔ים נֵלְכָ֖ה נִזְבְּחָ֥ה לַֽיהֹוָֽה׃ %
% ---------- פרק 5 | פסוק 18 ----------
 \textbf{יח} וְעַתָּה֙ לְכ֣וּ עִבְד֔וּ וְתֶ֖בֶן לֹא־יִנָּתֵ֣ן לָכֶ֑ם וְתֹ֥כֶן לְבֵנִ֖ים תִּתֵּֽנוּ׃ \Rashi{\textbf{ותכן לבנים.}חשבון הלבנים; וכן "את הכסף המתכן" (מלכים ב י"ב) – המנוי, כמו שאמור בענין "ויצרו וימנו את הכסף":}%
% ---------- פרק 5 | פסוק 19 ----------
 \textbf{יט} וַיִּרְא֞וּ שֹֽׁטְרֵ֧י בְנֵֽי־יִשְׂרָאֵ֛ל אֹתָ֖ם בְּרָ֣ע לֵאמֹ֑ר לֹא־תִגְרְע֥וּ מִלִּבְנֵיכֶ֖ם דְּבַר־י֥וֹם בְּיוֹמֽוֹ׃ \Rashi{\textbf{ויראו שטרי בני ישראל.}את חבריהם הנרדים על ידם, ברע, ראו אותם ברעה וצרה המוצאת אותם, בהכבידם העבודה עליהם לאמר לא תגרעו וגו':}%
% ---------- פרק 5 | פסוק 20 ----------
 \textbf{כ} וַֽיִּפְגְּעוּ֙ אֶת־מֹשֶׁ֣ה וְאֶֽת־אַהֲרֹ֔ן נִצָּבִ֖ים לִקְרָאתָ֑ם בְּצֵאתָ֖ם מֵאֵ֥ת פַּרְעֹֽה׃ \Rashi{\textbf{ויפגעו.}אנשים מישראל את משה ואת אהרן וגו'. ורבותינו דרשו: כל "נצים" "ונצבים" דתן ואבירם היו, שנ' בהם "יצאו נצבים" (במדבר ט״ז:כ״ד), (נדרים ס"ד):}%
% ---------- פרק 5 | פסוק 21 ----------
 \textbf{כא} וַיֹּאמְר֣וּ אֲלֵהֶ֔ם יֵ֧רֶא יְהֹוָ֛ה עֲלֵיכֶ֖ם וְיִשְׁפֹּ֑ט אֲשֶׁ֧ר הִבְאַשְׁתֶּ֣ם אֶת־רֵיחֵ֗נוּ בְּעֵינֵ֤י פַרְעֹה֙ וּבְעֵינֵ֣י עֲבָדָ֔יו לָֽתֶת־חֶ֥רֶב בְּיָדָ֖ם לְהׇרְגֵֽנוּ׃ %
% ---------- פרק 5 | פסוק 22 ----------
 \textbf{כב} וַיָּ֧שׇׁב מֹשֶׁ֛ה אֶל־יְהֹוָ֖ה וַיֹּאמַ֑ר אֲדֹנָ֗י לָמָ֤ה הֲרֵעֹ֙תָה֙ לָעָ֣ם הַזֶּ֔ה לָ֥מָּה זֶּ֖ה שְׁלַחְתָּֽנִי׃ \Rashi{\textbf{למה הרעתה לעם הזה.}ואם תאמר מה אכפת לך? קובל אני על ששלחתני (שמות רבה):}%
% ---------- פרק 5 | פסוק 23 ----------
 \textbf{כג} וּמֵאָ֞ז בָּ֤אתִי אֶל־פַּרְעֹה֙ לְדַבֵּ֣ר בִּשְׁמֶ֔ךָ הֵרַ֖ע לָעָ֣ם הַזֶּ֑ה וְהַצֵּ֥ל לֹא־הִצַּ֖לְתָּ אֶת־עַמֶּֽךָ׃ \Rashi{\textbf{הרע.}לשון הפעיל הוא – הרבה רעה עליהם, ותרגומו "אבאיש":}%
% ---------- פרק 6 | פסוק 1 ----------
 \textbf{א} וַיֹּ֤אמֶר יְהֹוָה֙ אֶל־מֹשֶׁ֔ה עַתָּ֣ה תִרְאֶ֔ה אֲשֶׁ֥ר אֶֽעֱשֶׂ֖ה לְפַרְעֹ֑ה כִּ֣י בְיָ֤ד חֲזָקָה֙ יְשַׁלְּחֵ֔ם וּבְיָ֣ד חֲזָקָ֔ה יְגָרְשֵׁ֖ם מֵאַרְצֽוֹ׃ \{ס\} \Rashi{\textbf{עתה תראה וגו'.}הרהרת על מדותי, לא כאברהם שאמרתי לו "כי ביצחק יקרא לך זרע" (בראשית כ"א), ואחר כך אמרתי לו "העלהו לעולה" ולא הרהר אחרי; לפיכך עתה תראה. העשוי לפרעה תראה, ולא העשוי למלכי שבעה אמות כשאביאם לארץ (סנהדרין קי"א): כי ביד חזקה ישלחם. מפני ידי החזקה שתחזק עליו ישלחם: וביד חזקה יגרשם מארצו. על כרחם של ישראל יגרשם, לא יספיקו לעשות להם צדה; וכן הוא אומר "ותחזק מצרים על העם וגו'" (שמות י"ב):}%
\addparasha{חומש שמות | פרשת וארא}
% ---------- פרק 6 | פסוק 2 ----------
 \textbf{ב} וַיְדַבֵּ֥ר אֱלֹהִ֖ים אֶל־מֹשֶׁ֑ה וַיֹּ֥אמֶר אֵלָ֖יו אֲנִ֥י יְהֹוָֽה׃ \Rashi{\textbf{וידבר אלהים אל משה.}דבר אתו משפט, על שהקשה לדבר ולומר "למה הרעתה לעם הזה": ויאמר אליו אני ה'. נאמן לשלם שכר טוב למתהלכים לפני, ולא לחנם שלחתיך כי אם לקים דברי שדברתי לאבות הראשונים. ובלשון הזה מצינו שהוא נדרש בכמה מקומות אני ה' נאמן להפרע – כשהוא אמור אצל ענש – כגון "וחללת את שם אלהיך אני ה'" (ויקרא י״ט:י״ב), וכשהוא אומר אצל קיום מצוות – כגון "ושמרתם מצותי ועשיתם אותם אני ה'" (שם כ"ב) – נאמן לתן שכר:}%
% ---------- פרק 6 | פסוק 3 ----------
 \textbf{ג} וָאֵרָ֗א אֶל־אַבְרָהָ֛ם אֶל־יִצְחָ֥ק וְאֶֽל־יַעֲקֹ֖ב בְּאֵ֣ל שַׁדָּ֑י וּשְׁמִ֣י יְהֹוָ֔ה לֹ֥א נוֹדַ֖עְתִּי לָהֶֽם׃ \Rashi{\textbf{וארא.}אל האבות: באל שדי. הבטחתים הבטחות ובכלן אמרתי להם אני אל שדי: ושמי ה' לא נודעתי להם. לא הודעתי אין כתיב כאן אלא לא נודעתי, לא נכרתי להם במדת אמתות שלי, שעליה נקרא שמי ה', נאמן לאמת דברי, שהרי הבטחתים ולא קימתי:}%
% ---------- פרק 6 | פסוק 4 ----------
 \textbf{ד} וְגַ֨ם הֲקִמֹ֤תִי אֶת־בְּרִיתִי֙ אִתָּ֔ם לָתֵ֥ת לָהֶ֖ם אֶת־אֶ֣רֶץ כְּנָ֑עַן אֵ֛ת אֶ֥רֶץ מְגֻרֵיהֶ֖ם אֲשֶׁר־גָּ֥רוּ בָֽהּ׃ \Rashi{\textbf{וגם הקמתי את בריתי וגו'.}וגם כשנראיתי להם באל שדי הצבתי והעמדתי בריתי ביני וביניהם, לתת להם ארץ כנען. לאברהם בפרשת מילה נאמר "אני אל שדי וגו' ונתתי לך ולזרעך אחריך את ארץ מגריך" (בראשית י"ז), ליצחק "כי לך ולזרעך אתן את כל הארצת האל והקמתי את השבעה אשר נשבעתי לאברהם" (שם כ"ו), ואותה שבועה שנשבעתי לאברהם באל שדי אמרתי ליעקב – "אני אל שדי פרה ורבה וגו' ואת הארץ אשר וגו'" (שם ל"ה), הרי שנדרתי להם ולא קימתי:}\Peirush{לכם לא נאמר אלא להם מכאן לתחיית המתים מן התורה\Makor{סנהדרין צ:.}}%
% ---------- פרק 6 | פסוק 5 ----------
 \textbf{ה} וְגַ֣ם ׀ אֲנִ֣י שָׁמַ֗עְתִּי אֶֽת־נַאֲקַת֙ בְּנֵ֣י יִשְׂרָאֵ֔ל אֲשֶׁ֥ר מִצְרַ֖יִם מַעֲבִדִ֣ים אֹתָ֑ם וָאֶזְכֹּ֖ר אֶת־בְּרִיתִֽי׃ \Rashi{\textbf{וגם אני.}כמו שהצבתי והעמדתי הברית יש עלי לקים, לפיכך שמעתי את נאקת בני ישראל הנואקים, אשר מצרים מעבדים אתם ואזכר אותו הברית, כי בברית בין הבתרים אמרתי לו "וגם את הגוי אשר יעבדו דן אנכי" (בראשית ט"ו):}%
% ---------- פרק 6 | פסוק 6 ----------
 \textbf{ו} לָכֵ֞ן אֱמֹ֥ר לִבְנֵֽי־יִשְׂרָאֵל֮ אֲנִ֣י יְהֹוָה֒ וְהוֹצֵאתִ֣י אֶתְכֶ֗ם מִתַּ֙חַת֙ סִבְלֹ֣ת מִצְרַ֔יִם וְהִצַּלְתִּ֥י אֶתְכֶ֖ם מֵעֲבֹדָתָ֑ם וְגָאַלְתִּ֤י אֶתְכֶם֙ בִּזְר֣וֹעַ נְטוּיָ֔ה וּבִשְׁפָטִ֖ים גְּדֹלִֽים׃ \Rashi{\textbf{לכן.}על פי אותה השבועה: אמר לבני ישראל אני ה'. הנאמן בהבטחתי: והוצאתי אתכם. כי כן הבטחתי: "ואחרי כן יצאו ברכש גדול" (שם): סבלת. טרח משא מצרים:}\Peirush{נאמר כאן סבלות מצרים, ועל יוסף נאמר "הסירותי מסבל שכמו" (תהלים פא, ז), כמו שיוסף יצא מבית האסורים בראש השנה כך יצאנו ממצרים בראש השנה\Makor{ראש השנה יא:}}%
% ---------- פרק 6 | פסוק 7 ----------
 \textbf{ז} וְלָקַחְתִּ֨י אֶתְכֶ֥ם לִי֙ לְעָ֔ם וְהָיִ֥יתִי לָכֶ֖ם לֵֽאלֹהִ֑ים וִֽידַעְתֶּ֗ם כִּ֣י אֲנִ֤י יְהֹוָה֙ אֱלֹ֣הֵיכֶ֔ם הַמּוֹצִ֣יא אֶתְכֶ֔ם מִתַּ֖חַת סִבְל֥וֹת מִצְרָֽיִם׃ \Peirush{נאמר ולקחתי אתכם ונאמר "והבאתי אתכם" (ח), מקיש יציאת מצרים לכניסת הארץ מה כניסתם לארץ נכנסו שנים מששים ריבוא [יהושע וכלב] אף ביציאת מצרים יצאו שנים משישים ריבוא\Makor{סנהדרין קיא.}}\Peirush{ר' נחמיה סובר שמברכים מוציא לחם מן הארץ, משום שהמוציא זהו לשון הווה שנאמר המוציא אתכם\Makor{ברכות לח.}}%
% ---------- פרק 6 | פסוק 8 ----------
 \textbf{ח} וְהֵבֵאתִ֤י אֶתְכֶם֙ אֶל־הָאָ֔רֶץ אֲשֶׁ֤ר נָשָׂ֙אתִי֙ אֶת־יָדִ֔י לָתֵ֣ת אֹתָ֔הּ לְאַבְרָהָ֥ם לְיִצְחָ֖ק וּֽלְיַעֲקֹ֑ב וְנָתַתִּ֨י אֹתָ֥הּ לָכֶ֛ם מוֹרָשָׁ֖ה אֲנִ֥י יְהֹוָֽה׃ \Rashi{\textbf{נשאתי את ידי.}הרימותיה להשבע בכסאי:}\Peirush{ראה פס' ז'\Makor{סנהדרין קיא.}}\Peirush{נאמר לכם, מכאן שחלוקת הארץ היא באחוזים לפי יוצאי מצרים ולא לפי השבטים\Makor{בבא בתרא קיז: שם קיט:}}%
% ---------- פרק 6 | פסוק 9 ----------
 \textbf{ט} וַיְדַבֵּ֥ר מֹשֶׁ֛ה כֵּ֖ן אֶל־בְּנֵ֣י יִשְׂרָאֵ֑ל וְלֹ֤א שָֽׁמְעוּ֙ אֶל־מֹשֶׁ֔ה מִקֹּ֣צֶר ר֔וּחַ וּמֵעֲבֹדָ֖ה קָשָֽׁה׃ \{פ\} \Rashi{\textbf{ולא שמעו אל משה.}לא קבלו תנחומין: מקצר רוח. כל מי שהוא מצר, רוחו ונשימתו קצרה, ואינו יכול להאריך בנשימתו: קרוב לענין זה שמעתי בפרשה זו מרבי ברוך ברבי אליעזר, והביא לי ראיה ממקרא זה, "בפעם הזאת אודיעם את ידי ואת גבורתי וידעו כי שמי ה'" (ירמיהו ט"ז) – למדנו כשהקב"ה מאמן את דבריו אפלו לפרענות מודיע ששמו ה', וכל שכן האמנה לטובה. ורבותינו דרשוהו לענין של מעלה, שאמר משה למה הרעתה, אמר לו הקב"ה "חבל על דאבדין ולא משתכחין" – יש לי להתאונן על מיתת האבות, הרבה פעמים נגליתי אליהם באל שדי ולא אמרו לי מה שמך, ואתה אמרת מה שמו, מה אמר אליהם: וגם הקמתי וגו'. וכשבקש אברהם לקבר את שרה לא מצא קבר עד שקנה בדמים מרבים, וכן ביצחק ערערו עליו על הבארות אשר חפר, וכן יעקב ויקן את חלקת השדה, לנטות אהלו, ולא הרהרו אחר מדותי, ואתה אמרת למה הרעתה? ואין המדרש מתישב אחר המקרא מפני כמה דברים, אחת, שלא נאמר ושמי ה' לא שאלו לי, ואם תאמר לא הודיעם שכך שמו, הרי תחלה כשנגלה לאברהם בין הבתרים נאמר "אני ה' אשר הוצאתיך מאור כשדים" (בראשית ט"ו), ועוד היאך הסמיכה נמשכת בדברים שהוא סומך לכאן וגם אני שמעתי וגו' לכן אמר לבני ישראל? לכך אני אומר יתישב המקרא על פשוטו, דבור דבור על אפניו, והדרשה תדרש, שנאמר "הלוא כה דברי כאש נאם ה' וכפטיש יפוצץ סלע" (ירמיהו כ"ג) – מתחלק לכמה ניצוצות (סנהדרין ל"ד):}%
% ---------- פרק 6 | פסוק 10 ----------
 \textbf{י} וַיְדַבֵּ֥ר יְהֹוָ֖ה אֶל־מֹשֶׁ֥ה לֵּאמֹֽר׃ %
% ---------- פרק 6 | פסוק 11 ----------
 \textbf{יא} בֹּ֣א דַבֵּ֔ר אֶל־פַּרְעֹ֖ה מֶ֣לֶךְ מִצְרָ֑יִם וִֽישַׁלַּ֥ח אֶת־בְּנֵֽי־יִשְׂרָאֵ֖ל מֵאַרְצֽוֹ׃ %
% ---------- פרק 6 | פסוק 12 ----------
 \textbf{יב} וַיְדַבֵּ֣ר מֹשֶׁ֔ה לִפְנֵ֥י יְהֹוָ֖ה לֵאמֹ֑ר הֵ֤ן בְּנֵֽי־יִשְׂרָאֵל֙ לֹֽא־שָׁמְע֣וּ אֵלַ֔י וְאֵיךְ֙ יִשְׁמָעֵ֣נִי פַרְעֹ֔ה וַאֲנִ֖י עֲרַ֥ל שְׂפָתָֽיִם׃ \{פ\} \Rashi{\textbf{ואיך ישמעני פרעה.}זה אחד מעשרה קל וחמר שבתורה (בראשית רבה צ"ב): ערל שפתים. אטום שפתים; וכן כל לשון ערלה אני אומר שהוא אטום, "ערלה אזנם" (ירמיהו ו') – אטומה משמע, "ערלי לב" (שם ט') – אטומים מהבין, "שתה גם אתה והערל" (חבקוק ב') – והאטם משכרות כוס הקללה, "ערלת בשר" – שהגיד אטום ומכסה בה, "וערלתם ערלתו" (ויקרא י"ט) – עשו לו אטם וכסוי אסור שיבדיל בפני אכילתו, "שלש שנים יהיה לכם ערלים" (שם) – אטום ומכסה ומבדל מלאכלו:}%
% ---------- פרק 6 | פסוק 13 ----------
 \textbf{יג} וַיְדַבֵּ֣ר יְהֹוָה֮ אֶל־מֹשֶׁ֣ה וְאֶֽל־אַהֲרֹן֒ וַיְצַוֵּם֙ אֶל־בְּנֵ֣י יִשְׂרָאֵ֔ל וְאֶל־פַּרְעֹ֖ה מֶ֣לֶךְ מִצְרָ֑יִם לְהוֹצִ֥יא אֶת־בְּנֵֽי־יִשְׂרָאֵ֖ל מֵאֶ֥רֶץ מִצְרָֽיִם׃ \{ס\} \Rashi{\textbf{וידבר ה' אל משה ואל אהרן.}לפי שאמר משה אני ערל שפתים, צרף לו הקב"ה את אהרן להיות לו לפה ולמליץ: ויצום אל בני ישראל. צום עליהם להנהיגם בנחת ולסבל אותם (שמות רבה): ואל פרעה מלך מצרים. צום עליו לחלק לו כבוד בדבריהם, זה מדרשו (שם); ופשוטו, צום על דבר ישראל ועל שליחותו אל פרעה; ודבר הצווי מהו, מפרש בפרשה שניה לאחר סדר היחס, אלא מתוך שהזכיר משה ואהרן, הפסיק הענין באלה ראשי בית אבתם ללמדנו היאך נולדו משה ואהרן ובמי נתיחסו:}%
% ---------- פרק 6 | פסוק 14 ----------
 \textbf{יד} אֵ֖לֶּה רָאשֵׁ֣י בֵית־אֲבֹתָ֑ם בְּנֵ֨י רְאוּבֵ֜ן בְּכֹ֣ר יִשְׂרָאֵ֗ל חֲנ֤וֹךְ וּפַלּוּא֙ חֶצְרֹ֣ן וְכַרְמִ֔י אֵ֖לֶּה מִשְׁפְּחֹ֥ת רְאוּבֵֽן׃ \Rashi{\textbf{אלה ראשי בית אבתם.}מתוך שהזקק ליחס שבטו של לוי עד משה ואהרן בשביל משה ואהרן, התחיל ליחסם דרך תולדותם מראובן. ובפסיקתא רבתי ראיתי, לפי שקנטרם יעקב אבינו לשלושה שבטים הללו בשעת מותו, חזר הכתוב ויחסם כאן לבדם, לומר שחשובים הם:}%
% ---------- פרק 6 | פסוק 15 ----------
 \textbf{טו} וּבְנֵ֣י שִׁמְע֗וֹן יְמוּאֵ֨ל וְיָמִ֤ין וְאֹ֙הַד֙ וְיָכִ֣ין וְצֹ֔חַר וְשָׁא֖וּל בֶּן־הַֽכְּנַעֲנִ֑ית אֵ֖לֶּה מִשְׁפְּחֹ֥ת שִׁמְעֽוֹן׃ %
% ---------- פרק 6 | פסוק 16 ----------
 \textbf{טז} וְאֵ֨לֶּה שְׁמ֤וֹת בְּנֵֽי־לֵוִי֙ לְתֹ֣לְדֹתָ֔ם גֵּרְשׁ֕וֹן וּקְהָ֖ת וּמְרָרִ֑י וּשְׁנֵי֙ חַיֵּ֣י לֵוִ֔י שֶׁ֧בַע וּשְׁלֹשִׁ֛ים וּמְאַ֖ת שָׁנָֽה׃ \Rashi{\textbf{ושני חיי לוי וגו'.}למה נמנו שנותיו של לוי? להודיע כמה ימי השעבוד; שכל זמן שאחד מן השבטים קים לא היה שעבוד (שמות רבה א'), שנאמר "וימת יוסף וכל אחיו" (שמות א), ואחר כך "ויקם מלך חדש", ולוי האריך ימים על כלם:}%
% ---------- פרק 6 | פסוק 17 ----------
 \textbf{יז} בְּנֵ֥י גֵרְשׁ֛וֹן לִבְנִ֥י וְשִׁמְעִ֖י לְמִשְׁפְּחֹתָֽם׃ %
% ---------- פרק 6 | פסוק 18 ----------
 \textbf{יח} וּבְנֵ֣י קְהָ֔ת עַמְרָ֣ם וְיִצְהָ֔ר וְחֶבְר֖וֹן וְעֻזִּיאֵ֑ל וּשְׁנֵי֙ חַיֵּ֣י קְהָ֔ת שָׁלֹ֧שׁ וּשְׁלֹשִׁ֛ים וּמְאַ֖ת שָׁנָֽה׃ \Rashi{\textbf{ושני חיי קהת, ושני חיי עמרם וגו'.}מחשבון זה אנו למדים על מושב בני ישראל ארבע מאות שנה שאמר הכתוב (בראשית ט״ו:י״ג), שלא בארץ מצרים לבדה היו, אלא מיום שנולד יצחק; שהרי קהת מיורדי מצרים היה, חשב כל שנותיו ושנות עמרם ושמונים של משה, לא תמצאם ארבע מאות שנה, והרבה שנים נבלעים לבנים בשני האבות:}%
% ---------- פרק 6 | פסוק 19 ----------
 \textbf{יט} וּבְנֵ֥י מְרָרִ֖י מַחְלִ֣י וּמוּשִׁ֑י אֵ֛לֶּה מִשְׁפְּחֹ֥ת הַלֵּוִ֖י לְתֹלְדֹתָֽם׃ %
% ---------- פרק 6 | פסוק 20 ----------
 \textbf{כ} וַיִּקַּ֨ח עַמְרָ֜ם אֶת־יוֹכֶ֤בֶד דֹּֽדָתוֹ֙ ל֣וֹ לְאִשָּׁ֔ה וַתֵּ֣לֶד ל֔וֹ אֶֽת־אַהֲרֹ֖ן וְאֶת־מֹשֶׁ֑ה וּשְׁנֵי֙ חַיֵּ֣י עַמְרָ֔ם שֶׁ֧בַע וּשְׁלֹשִׁ֛ים וּמְאַ֖ת שָׁנָֽה׃ \Rashi{\textbf{יוכבד דודתו.}אחת אבוהי, בת לוי אחות קהת:}\Peirush{נאמר דודתו, אולי משמע דודתו מן האם, ומותר לשאת את אחות האם. לא, הכוונה לדודתו מן האב\Makor{סנהדרין נח:}}\Peirush{נאמר באחשוורוש שמלך על "שבע ועשרים ומאה מדינה" (אסתר א, א) ולומדים שבתחילה מלך על שבע ולבסוף מלך על עשרים ולבסוף מלך על מאה. ומקשים, נאמר בעמרם שבע ושלשים ומאת שנה, אך לא לומדים משם דרשה\Makor{מגילה יא.}}%
% ---------- פרק 6 | פסוק 21 ----------
 \textbf{כא} וּבְנֵ֖י יִצְהָ֑ר קֹ֥רַח וָנֶ֖פֶג וְזִכְרִֽי׃ %
% ---------- פרק 6 | פסוק 22 ----------
 \textbf{כב} וּבְנֵ֖י עֻזִּיאֵ֑ל מִֽישָׁאֵ֥ל וְאֶלְצָפָ֖ן וְסִתְרִֽי׃ %
% ---------- פרק 6 | פסוק 23 ----------
 \textbf{כג} וַיִּקַּ֨ח אַהֲרֹ֜ן אֶת־אֱלִישֶׁ֧בַע בַּת־עַמִּינָדָ֛ב אֲח֥וֹת נַחְשׁ֖וֹן ל֣וֹ לְאִשָּׁ֑ה וַתֵּ֣לֶד ל֗וֹ אֶת־נָדָב֙ וְאֶת־אֲבִיה֔וּא אֶת־אֶלְעָזָ֖ר וְאֶת־אִֽיתָמָֽר׃ \Rashi{\textbf{אחות נחשון.}מכאן למדו, הנושא אשה צריך לבדק באחיה (בבא בתרא ק"י):}\Peirush{הנושא אשה צריך שיבדוק באחיה שנאמר אחות נחשון\Makor{בבא בתרא קי.}}%
% ---------- פרק 6 | פסוק 24 ----------
 \textbf{כד} וּבְנֵ֣י קֹ֔רַח אַסִּ֥יר וְאֶלְקָנָ֖ה וַאֲבִיאָסָ֑ף אֵ֖לֶּה מִשְׁפְּחֹ֥ת הַקׇּרְחִֽי׃ %
% ---------- פרק 6 | פסוק 25 ----------
 \textbf{כה} וְאֶלְעָזָ֨ר בֶּֽן־אַהֲרֹ֜ן לָקַֽח־ל֨וֹ מִבְּנ֤וֹת פּֽוּטִיאֵל֙ ל֣וֹ לְאִשָּׁ֔ה וַתֵּ֥לֶד ל֖וֹ אֶת־פִּֽינְחָ֑ס אֵ֗לֶּה רָאשֵׁ֛י אֲב֥וֹת הַלְוִיִּ֖ם לְמִשְׁפְּחֹתָֽם׃ \Rashi{\textbf{מבנות פוטיאל.}מזרע יתרו שפטם עגלים לע"ז, ומזרע יוסף שפטפט ביצרו (סוטה מ"ג):}\Peirush{נאמר מבנות פוטיאל, מכאן שפנחס יצא מיתרו שפיטם עגלים לעבודת כוכבים, ומיוסף שפיטפט ביצרו [כבש אותו]\Makor{סוטה מג. בבא בתרא קט:}}%
% ---------- פרק 6 | פסוק 26 ----------
 \textbf{כו} ה֥וּא אַהֲרֹ֖ן וּמֹשֶׁ֑ה אֲשֶׁ֨ר אָמַ֤ר יְהֹוָה֙ לָהֶ֔ם הוֹצִ֜יאוּ אֶת־בְּנֵ֧י יִשְׂרָאֵ֛ל מֵאֶ֥רֶץ מִצְרַ֖יִם עַל־צִבְאֹתָֽם׃ \Rashi{\textbf{הוא אהרן ומשה.}אלו שהזכרו למעלה, שילדה יוכבד לעמרם. הוא אהרן ומשה. יש מקומות שמקדים אהרן למשה ויש מקומות שמקדים משה לאהרן, לומר לך ששקולין כאחד: על צבאתם. בצבאותם, כל צבאם לשבטיהם; יש על שאינו אלא במקום אות אחת, "ועל חרבך תחיה" (בראשית כ"ז), כמו בחרבך, "עמדתם על חרבכם" (יחזקאל ל"ג) – בחרבכם:}%
% ---------- פרק 6 | פסוק 27 ----------
 \textbf{כז} הֵ֗ם הַֽמְדַבְּרִים֙ אֶל־פַּרְעֹ֣ה מֶֽלֶךְ־מִצְרַ֔יִם לְהוֹצִ֥יא אֶת־בְּנֵֽי־יִשְׂרָאֵ֖ל מִמִּצְרָ֑יִם ה֥וּא מֹשֶׁ֖ה וְאַהֲרֹֽן׃ \Rashi{\textbf{הם המדברים וגו'.}הם שנצטוו, הם שקימו: הוא משה ואהרן. הם בשליחותם ובצדקתם מתחלה ועד סוף (מגילה י"א):}\Peirush{נאמר הוא אהרן ומשה, הן בצדקן מתחילתן ועד סופן\Makor{מגילה יא.}}%
% ---------- פרק 6 | פסוק 28 ----------
 \textbf{כח} וַיְהִ֗י בְּי֨וֹם דִּבֶּ֧ר יְהֹוָ֛ה אֶל־מֹשֶׁ֖ה בְּאֶ֥רֶץ מִצְרָֽיִם׃ \{ס\} \Rashi{\textbf{ויהי ביום דבר וגו'.}מחבר למקרא שלאחריו:}%
% ---------- פרק 6 | פסוק 29 ----------
 \textbf{כט} וַיְדַבֵּ֧ר יְהֹוָ֛ה אֶל־מֹשֶׁ֥ה לֵּאמֹ֖ר אֲנִ֣י יְהֹוָ֑ה דַּבֵּ֗ר אֶל־פַּרְעֹה֙ מֶ֣לֶךְ מִצְרַ֔יִם אֵ֛ת כׇּל־אֲשֶׁ֥ר אֲנִ֖י דֹּבֵ֥ר אֵלֶֽיךָ׃ \Rashi{\textbf{וידבר ה'.}הוא הדבור עצמו האמור למעלה בא דבר אל פרעה מלך מצרים, אלא מתוך שהפסיק הענין כדי ליחסם, חזר עליו להתחיל בו: אני ה'. כדאי אני לשלחך ולקים דברי שליחותי:}%
% ---------- פרק 6 | פסוק 30 ----------
 \textbf{ל} וַיֹּ֥אמֶר מֹשֶׁ֖ה לִפְנֵ֣י יְהֹוָ֑ה הֵ֤ן אֲנִי֙ עֲרַ֣ל שְׂפָתַ֔יִם וְאֵ֕יךְ יִשְׁמַ֥ע אֵלַ֖י פַּרְעֹֽה׃ \{פ\} \Rashi{\textbf{ויאמר משה לפני ה'.}היא האמירה שאמר למעלה הן בני ישראל לא שמעו אלי, ושנה הכתוב כאן, כיון שהפסיק הענין; וכן היא השיטה, כאדם האומר "נחזר על הראשונות":}%
% ---------- פרק 7 | פסוק 1 ----------
 \textbf{א} וַיֹּ֤אמֶר יְהֹוָה֙ אֶל־מֹשֶׁ֔ה רְאֵ֛ה נְתַתִּ֥יךָ אֱלֹהִ֖ים לְפַרְעֹ֑ה וְאַהֲרֹ֥ן אָחִ֖יךָ יִהְיֶ֥ה נְבִיאֶֽךָ׃ \Rashi{\textbf{נתתיך אלהים לפרעה.}שופט ורודה, לרדותו במכות ויסורין: יהיה נביאך. כתרגומו – "מתרגמנך". וכן כל לשון נבואה אדם המכריז ומשמיע לעם דברי תוכחות, והוא מגזרת "ניב שפתים" (ישעיהו נ"ז), "ינוב חכמה" (משלי י'), "ויכל מהתנבות" דשמואל (א' י'); ובלעז קוראין לו פרידיכ"ר:}%
% ---------- פרק 7 | פסוק 2 ----------
 \textbf{ב} אַתָּ֣ה תְדַבֵּ֔ר אֵ֖ת כׇּל־אֲשֶׁ֣ר אֲצַוֶּ֑ךָּ וְאַהֲרֹ֤ן אָחִ֙יךָ֙ יְדַבֵּ֣ר אֶל־פַּרְעֹ֔ה וְשִׁלַּ֥ח אֶת־בְּנֵֽי־יִשְׂרָאֵ֖ל מֵאַרְצֽוֹ׃ \Rashi{\textbf{אתה תדבר.}פעם אחת כל שליחות ושליחות כפי ששמעתו מפי, ואהרן אחיך ימליצנו ויטעימנו באזני פרעה:}%
% ---------- פרק 7 | פסוק 3 ----------
 \textbf{ג} וַאֲנִ֥י אַקְשֶׁ֖ה אֶת־לֵ֣ב פַּרְעֹ֑ה וְהִרְבֵּיתִ֧י אֶת־אֹתֹתַ֛י וְאֶת־מוֹפְתַ֖י בְּאֶ֥רֶץ מִצְרָֽיִם׃ \Rashi{\textbf{ואני אקשה.}מאחר שהרשיע והתריס כנגדי, וגלוי לפני שאין נחת רוח באמות לתת לב שלם לשוב, טוב שיתקשה לבו, למען הרבות בו אותותי, ותכירו אתם את גבורתי. וכן מדתו של הקב"ה, מביא פרענות על האמות כדי שישמעו ישראל וייראו, שנא' "הכרתי גוים נשמו פנותם … אמרתי אך תיראי אותי תקחי מוסר" (צפניה ג'); ואעפ"כ בחמש מכות הראשונות לא נאמר ויחזק ה' את לב פרעה, אלא ויחזק לב פרעה (תנחומא):}%
% ---------- פרק 7 | פסוק 4 ----------
 \textbf{ד} וְלֹֽא־יִשְׁמַ֤ע אֲלֵכֶם֙ פַּרְעֹ֔ה וְנָתַתִּ֥י אֶת־יָדִ֖י בְּמִצְרָ֑יִם וְהוֹצֵאתִ֨י אֶת־צִבְאֹתַ֜י אֶת־עַמִּ֤י בְנֵֽי־יִשְׂרָאֵל֙ מֵאֶ֣רֶץ מִצְרַ֔יִם בִּשְׁפָטִ֖ים גְּדֹלִֽים׃ \Rashi{\textbf{את ידי.}יד ממש, להכות בהם:}\Peirush{ניתן למחוק את השם 'צבאות' שלא נקרא אלא על שם ישראל שנאמר את צבאותי את עמי בני ישראל\Makor{שבועות לה:}}%
% ---------- פרק 7 | פסוק 5 ----------
 \textbf{ה} וְיָדְע֤וּ מִצְרַ֙יִם֙ כִּֽי־אֲנִ֣י יְהֹוָ֔ה בִּנְטֹתִ֥י אֶת־יָדִ֖י עַל־מִצְרָ֑יִם וְהוֹצֵאתִ֥י אֶת־בְּנֵֽי־יִשְׂרָאֵ֖ל מִתּוֹכָֽם׃ %
% ---------- פרק 7 | פסוק 6 ----------
 \textbf{ו} וַיַּ֥עַשׂ מֹשֶׁ֖ה וְאַהֲרֹ֑ן כַּאֲשֶׁ֨ר צִוָּ֧ה יְהֹוָ֛ה אֹתָ֖ם כֵּ֥ן עָשֽׂוּ׃ %
% ---------- פרק 7 | פסוק 7 ----------
 \textbf{ז} וּמֹשֶׁה֙ בֶּן־שְׁמֹנִ֣ים שָׁנָ֔ה וְאַֽהֲרֹ֔ן בֶּן־שָׁלֹ֥שׁ וּשְׁמֹנִ֖ים שָׁנָ֑ה בְּדַבְּרָ֖ם אֶל־פַּרְעֹֽה׃ \{פ\} %
% ---------- פרק 7 | פסוק 8 ----------
 \textbf{ח} וַיֹּ֣אמֶר יְהֹוָ֔ה אֶל־מֹשֶׁ֥ה וְאֶֽל־אַהֲרֹ֖ן לֵאמֹֽר׃ %
% ---------- פרק 7 | פסוק 9 ----------
 \textbf{ט} כִּי֩ יְדַבֵּ֨ר אֲלֵכֶ֤ם פַּרְעֹה֙ לֵאמֹ֔ר תְּנ֥וּ לָכֶ֖ם מוֹפֵ֑ת וְאָמַרְתָּ֣ אֶֽל־אַהֲרֹ֗ן קַ֧ח אֶֽת־מַטְּךָ֛ וְהַשְׁלֵ֥ךְ לִפְנֵֽי־פַרְעֹ֖ה יְהִ֥י לְתַנִּֽין׃ \Rashi{\textbf{מופת.}אות, להודיע שיש צרוך במי ששולח אתכם:}%
% ---------- פרק 7 | פסוק 10 ----------
 \textbf{י} וַיָּבֹ֨א מֹשֶׁ֤ה וְאַהֲרֹן֙ אֶל־פַּרְעֹ֔ה וַיַּ֣עֲשׂוּ כֵ֔ן כַּאֲשֶׁ֖ר צִוָּ֣ה יְהֹוָ֑ה וַיַּשְׁלֵ֨ךְ אַהֲרֹ֜ן אֶת־מַטֵּ֗הוּ לִפְנֵ֥י פַרְעֹ֛ה וְלִפְנֵ֥י עֲבָדָ֖יו וַיְהִ֥י לְתַנִּֽין׃ \Rashi{\textbf{לתנין.}נחש:}%
% ---------- פרק 7 | פסוק 11 ----------
 \textbf{יא} וַיִּקְרָא֙ גַּם־פַּרְעֹ֔ה לַֽחֲכָמִ֖ים וְלַֽמְכַשְּׁפִ֑ים וַיַּֽעֲשׂ֨וּ גַם־הֵ֜ם חַרְטֻמֵּ֥י מִצְרַ֛יִם בְּלַהֲטֵיהֶ֖ם כֵּֽן׃ \Rashi{\textbf{בלהטיהם.}בלחשיהון; ואין לו דמיון במקרא; ויש לדמות לו "להט החרב המתהפכת" (בראשית ג') – דומה שהיא מתהפכת על ידי לחש:}\Peirush{בלהטיהם אלו מעשה כשפים\Makor{סנהדרין סז:}}%
% ---------- פרק 7 | פסוק 12 ----------
 \textbf{יב} וַיַּשְׁלִ֙יכוּ֙ אִ֣ישׁ מַטֵּ֔הוּ וַיִּהְי֖וּ לְתַנִּינִ֑ם וַיִּבְלַ֥ע מַטֵּֽה־אַהֲרֹ֖ן אֶת־מַטֹּתָֽם׃ \Rashi{\textbf{ויבלע מטה אהרן.}מאחר שחזר ונעשה מטה, בלע את כלן (שבת צ"ז):}\Peirush{נס בתוך נס, לאחר שחזר ונעשה מטה בלע את המטות, ולא כשהוא תנין\Makor{שבת צז.}}%
% ---------- פרק 7 | פסוק 13 ----------
 \textbf{יג} וַיֶּֽחֱזַק֙ לֵ֣ב פַּרְעֹ֔ה וְלֹ֥א שָׁמַ֖ע אֲלֵהֶ֑ם כַּאֲשֶׁ֖ר דִּבֶּ֥ר יְהֹוָֽה׃ \{ס\} %
% ---------- פרק 7 | פסוק 14 ----------
 \textbf{יד} וַיֹּ֤אמֶר יְהֹוָה֙ אֶל־מֹשֶׁ֔ה כָּבֵ֖ד לֵ֣ב פַּרְעֹ֑ה מֵאֵ֖ן לְשַׁלַּ֥ח הָעָֽם׃ \Rashi{\textbf{כבד.}תרגומו יקיר ולא אתיקר, מפני שהוא שם דבר, כמו "כי כבד ממך הדבר" (שמות י"ח):}%
% ---------- פרק 7 | פסוק 15 ----------
 \textbf{טו} לֵ֣ךְ אֶל־פַּרְעֹ֞ה בַּבֹּ֗קֶר הִנֵּה֙ יֹצֵ֣א הַמַּ֔יְמָה וְנִצַּבְתָּ֥ לִקְרָאת֖וֹ עַל־שְׂפַ֣ת הַיְאֹ֑ר וְהַמַּטֶּ֛ה אֲשֶׁר־נֶהְפַּ֥ךְ לְנָחָ֖שׁ תִּקַּ֥ח בְּיָדֶֽךָ׃ \Rashi{\textbf{הנה יצא המימה.}לנקביו; שהיה עושה עצמו אלוה ואומר שאינו צריך לנקביו, ומשכים ויוצא לנילוס ועושה שם צרכיו (תנחומא):}\Peirush{אמר ר"ל מלך הוא והסבר לו פנים ורבי יוחנן אמר רשע הוא והעז פניך בו\Makor{זבחים קב.}}\Peirush{פרעה היה אמגושי\Makor{מועד קטן יח:}}%
% ---------- פרק 7 | פסוק 16 ----------
 \textbf{טז} וְאָמַרְתָּ֣ אֵלָ֗יו יְהֹוָ֞ה אֱלֹהֵ֤י הָעִבְרִים֙ שְׁלָחַ֤נִי אֵלֶ֙יךָ֙ לֵאמֹ֔ר שַׁלַּח֙ אֶת־עַמִּ֔י וְיַֽעַבְדֻ֖נִי בַּמִּדְבָּ֑ר וְהִנֵּ֥ה לֹא־שָׁמַ֖עְתָּ עַד־כֹּֽה׃ \Rashi{\textbf{עד כה.}עד הנה; ומדרשו: עד שתשמע ממני מכת בכורות שאפתח בה ב"כה" אמר ה' כחצות הלילה (שמות י״א:ד׳):}%
% ---------- פרק 7 | פסוק 17 ----------
 \textbf{יז} כֹּ֚ה אָמַ֣ר יְהֹוָ֔ה בְּזֹ֣את תֵּדַ֔ע כִּ֖י אֲנִ֣י יְהֹוָ֑ה הִנֵּ֨ה אָנֹכִ֜י מַכֶּ֣ה ׀ בַּמַּטֶּ֣ה אֲשֶׁר־בְּיָדִ֗י עַל־הַמַּ֛יִם אֲשֶׁ֥ר בַּיְאֹ֖ר וְנֶהֶפְכ֥וּ לְדָֽם׃ \Rashi{\textbf{ונהפכו לדם.}לפי שאין גשמים יורדים במצרים ונילוס עולה ומשקה את הארץ ומצרים עובדים לנילוס, לפיכך הלקה את יראתם ואח"כ הלקה אותם (תנחומא):}%
% ---------- פרק 7 | פסוק 18 ----------
 \textbf{יח} וְהַדָּגָ֧ה אֲשֶׁר־בַּיְאֹ֛ר תָּמ֖וּת וּבָאַ֣שׁ הַיְאֹ֑ר וְנִלְא֣וּ מִצְרַ֔יִם לִשְׁתּ֥וֹת מַ֖יִם מִן־הַיְאֹֽר׃ \{ס\} \Rashi{\textbf{ונלאו מצרים.}לבקש רפואה למי היאור שיהיו ראויין לשתות:}%
% ---------- פרק 7 | פסוק 19 ----------
 \textbf{יט} וַיֹּ֨אמֶר יְהֹוָ֜ה אֶל־מֹשֶׁ֗ה אֱמֹ֣ר אֶֽל־אַהֲרֹ֡ן קַ֣ח מַטְּךָ֣ וּנְטֵֽה־יָדְךָ֩ עַל־מֵימֵ֨י מִצְרַ֜יִם עַֽל־נַהֲרֹתָ֣ם ׀ עַל־יְאֹרֵיהֶ֣ם וְעַל־אַגְמֵיהֶ֗ם וְעַ֛ל כׇּל־מִקְוֵ֥ה מֵימֵיהֶ֖ם וְיִֽהְיוּ־דָ֑ם וְהָ֤יָה דָם֙ בְּכׇל־אֶ֣רֶץ מִצְרַ֔יִם וּבָעֵצִ֖ים וּבָאֲבָנִֽים׃ \Rashi{\textbf{אמר אל אהרן.}לפי שהגן היאור על משה כשנשלך לתוכו, לפיכך לא לקה על ידו לא בדם ולא בצפרדעים, ולקה על ידי אהרן (שמות רבה): נהרתם. הם נהרות המושכים, כעין נהרות שלנו: יאריהם. הם ברכות נגרים, העשויות בידי אדם משפת הנהר לשדות, ונילוס מימיו מתברכים ועולה דרך היאורים ומשקה השדות: אגמיהם. קבוצת מים שאינן נובעין ואין מושכין, אלא עומדין במקום אחד, וקורין לו אשטנ"ק: בכל ארץ מצרים. אף במרחצאות ובאמבטאות שבבתים: ובעצים ובאבנים. מים שבכלי עץ ובכלי אבן:}%
% ---------- פרק 7 | פסוק 20 ----------
 \textbf{כ} וַיַּֽעֲשׂוּ־כֵן֩ מֹשֶׁ֨ה וְאַהֲרֹ֜ן כַּאֲשֶׁ֣ר ׀ צִוָּ֣ה יְהֹוָ֗ה וַיָּ֤רֶם בַּמַּטֶּה֙ וַיַּ֤ךְ אֶת־הַמַּ֙יִם֙ אֲשֶׁ֣ר בַּיְאֹ֔ר לְעֵינֵ֣י פַרְעֹ֔ה וּלְעֵינֵ֖י עֲבָדָ֑יו וַיֵּהָ֥פְכ֛וּ כׇּל־הַמַּ֥יִם אֲשֶׁר־בַּיְאֹ֖ר לְדָֽם׃ %
% ---------- פרק 7 | פסוק 21 ----------
 \textbf{כא} וְהַדָּגָ֨ה אֲשֶׁר־בַּיְאֹ֥ר מֵ֙תָה֙ וַיִּבְאַ֣שׁ הַיְאֹ֔ר וְלֹא־יָכְל֣וּ מִצְרַ֔יִם לִשְׁתּ֥וֹת מַ֖יִם מִן־הַיְאֹ֑ר וַיְהִ֥י הַדָּ֖ם בְּכׇל־אֶ֥רֶץ מִצְרָֽיִם׃ \Peirush{הנודר לא לאכול דגה אסור בדגים קטנים וחייב בגדולים. אך קשה, הרי כתוב והדגה אשר ביאור מתה והכוונה גם לגדולים וגם לקטנים\Makor{נדרים נא:}}%
% ---------- פרק 7 | פסוק 22 ----------
 \textbf{כב} וַיַּֽעֲשׂוּ־כֵ֛ן חַרְטֻמֵּ֥י מִצְרַ֖יִם בְּלָטֵיהֶ֑ם וַיֶּחֱזַ֤ק לֵב־פַּרְעֹה֙ וְלֹא־שָׁמַ֣ע אֲלֵהֶ֔ם כַּאֲשֶׁ֖ר דִּבֶּ֥ר יְהֹוָֽה׃ \Rashi{\textbf{בלטיהם.}לחש שאומרין אותו בלט ובחשאי; ורבותינו אמרו, "בלטיהם" מעשה שדים, "בלהטיהם" מעשה כשפים (סנהדרין ס"ז): ויחזק לב פרעה. לומר על ידי מכשפות אתם עושים כן; "תבן אתם מכניסים לעפרים", עיר שכלה תבן, אף אתם מביאין מכשפות למצרים שכלה כשפים:}\Peirush{בלטיהם אלו מעשה שדים\Makor{סנהדרין סז:}}%
% ---------- פרק 7 | פסוק 23 ----------
 \textbf{כג} וַיִּ֣פֶן פַּרְעֹ֔ה וַיָּבֹ֖א אֶל־בֵּית֑וֹ וְלֹא־שָׁ֥ת לִבּ֖וֹ גַּם־לָזֹֽאת׃ \Rashi{\textbf{גם לזאת.}למופת המטה שנהפך לתנין, ולא לזה של דם:}%
% ---------- פרק 7 | פסוק 24 ----------
 \textbf{כד} וַיַּחְפְּר֧וּ כׇל־מִצְרַ֛יִם סְבִיבֹ֥ת הַיְאֹ֖ר מַ֣יִם לִשְׁתּ֑וֹת כִּ֣י לֹ֤א יָֽכְלוּ֙ לִשְׁתֹּ֔ת מִמֵּימֵ֖י הַיְאֹֽר׃ %
% ---------- פרק 7 | פסוק 25 ----------
 \textbf{כה} וַיִּמָּלֵ֖א שִׁבְעַ֣ת יָמִ֑ים אַחֲרֵ֥י הַכּוֹת־יְהֹוָ֖ה אֶת־הַיְאֹֽר׃ \{פ\} \Rashi{\textbf{וימלא.}מנין שבעת ימים, שלא שב היאור לקדמותו, שהיתה המכה משמשת רביע חדש, וג' חלקים היה מעיד ומתרה בהם (תנחומא):}%
% ---------- פרק 7 | פסוק 26 ----------
 \textbf{כו} וַיֹּ֤אמֶר יְהֹוָה֙ אֶל־מֹשֶׁ֔ה בֹּ֖א אֶל־פַּרְעֹ֑ה וְאָמַרְתָּ֣ אֵלָ֗יו כֹּ֚ה אָמַ֣ר יְהֹוָ֔ה שַׁלַּ֥ח אֶת־עַמִּ֖י וְיַֽעַבְדֻֽנִי׃ %
% ---------- פרק 7 | פסוק 27 ----------
 \textbf{כז} וְאִם־מָאֵ֥ן אַתָּ֖ה לְשַׁלֵּ֑חַ הִנֵּ֣ה אָנֹכִ֗י נֹגֵ֛ף אֶת־כׇּל־גְּבוּלְךָ֖ בַּֽצְפַרְדְּעִֽים׃ \Rashi{\textbf{ואם מאן אתה.}ואם סרבן אתה: מאן. כמו ממאן, מסרב, אלא כנוי האדם על שם המפעל, כמו "שלו" (איוב ט"ז), "שקט" (שופטים י"ח), "סר וזעף" (מלכים א כ'): נגף את כל גבולך. מכה, וכן כל לשון מגפה אינו לשון מיתה אלא לשון מכה, וכן "ונגפו אשה הרה" (שמות כ"א), אינו לשון מיתה, וכן "ובטרם יתנגפו רגליכם" (ירמיהו י"ג), "פן תגף באבן רגלך" (תהילים צ"א), "ולאבן נגף" (ישעיהו ח'):}%
% ---------- פרק 7 | פסוק 28 ----------
 \textbf{כח} וְשָׁרַ֣ץ הַיְאֹר֮ צְפַרְדְּעִים֒ וְעָלוּ֙ וּבָ֣אוּ בְּבֵיתֶ֔ךָ וּבַחֲדַ֥ר מִשְׁכָּבְךָ֖ וְעַל־מִטָּתֶ֑ךָ וּבְבֵ֤ית עֲבָדֶ֙יךָ֙ וּבְעַמֶּ֔ךָ וּבְתַנּוּרֶ֖יךָ וּבְמִשְׁאֲרוֹתֶֽיךָ׃ \Rashi{\textbf{ועלו.}מן היאור: בביתך. ואחר כך בבית עבדיך. הוא התחיל בעצה תחלה, "ויאמר אל עמו" (שמות א'), וממנו התחילה הפרענות:}\Peirush{מה ראו חנניה מישאל ועזריה שמסרו עצמם על קדושת השם לכבשן האש? קל וחומר מצפרדעים, מה צפרדעים שאין מצווין על קדושת השם, נאמר בהם ובאו ... ובתנוריך ובמשארותיך. אימתי משארות [חתיכות בצק] מצויות אצל תנור? הוי אומר בשעה שהתנור חם. אנו, שמצווין על קדושת השם, על אחת כמה וכמה\Makor{פסחים נג:}}%
% ---------- פרק 7 | פסוק 29 ----------
 \textbf{כט} וּבְכָ֥ה וּֽבְעַמְּךָ֖ וּבְכׇל־עֲבָדֶ֑יךָ יַעֲל֖וּ הַֽצְפַרְדְּעִֽים׃ \Rashi{\textbf{ובכה ובעמך.}בתוך מעיהם נכנסין ומקרקרין (סוטה י"א):}\Peirush{פרעה התנכל לישראל ראשון — לכן נענש ראשון מכל עמו, שנאמר ובכה ובעמך ובכל עבדיך יעלו הצפרדעים\Makor{סוטה יא.}}%
% ---------- פרק 8 | פסוק 1 ----------
 \textbf{א} וַיֹּ֣אמֶר יְהֹוָה֮ אֶל־מֹשֶׁה֒ אֱמֹ֣ר אֶֽל־אַהֲרֹ֗ן נְטֵ֤ה אֶת־יָדְךָ֙ בְּמַטֶּ֔ךָ עַ֨ל־הַנְּהָרֹ֔ת עַל־הַיְאֹרִ֖ים וְעַל־הָאֲגַמִּ֑ים וְהַ֥עַל אֶת־הַֽצְפַרְדְּעִ֖ים עַל־אֶ֥רֶץ מִצְרָֽיִם׃ %
% ---------- פרק 8 | פסוק 2 ----------
 \textbf{ב} וַיֵּ֤ט אַהֲרֹן֙ אֶת־יָד֔וֹ עַ֖ל מֵימֵ֣י מִצְרָ֑יִם וַתַּ֙עַל֙ הַצְּפַרְדֵּ֔עַ וַתְּכַ֖ס אֶת־אֶ֥רֶץ מִצְרָֽיִם׃ \Rashi{\textbf{ותעל הצפרדע.}צפרדע אחת היתה והיו מכין אותה והיא מתזת נחילים נחילים, זהו מדרשו (סנהדרין ס"ז). ופשוטו יש לומר, שרוץ הצפרדעים קורא לשון יחידות; וכן ותהי הכנם – הרחישה, פדולייר"א בלעז, ואף ותעל הצפרדע – גרינולייר"א בלעז:}\Peirush{נאמר ותעל הצפרדע מכאן שצפרדע אחת היתה בתחילה, השריצה ומלאה כל ארץ מצרים צפרדעים\Makor{סנהדרין סז:}}%
% ---------- פרק 8 | פסוק 3 ----------
 \textbf{ג} וַיַּֽעֲשׂוּ־כֵ֥ן הַֽחַרְטֻמִּ֖ים בְּלָטֵיהֶ֑ם וַיַּעֲל֥וּ אֶת־הַֽצְפַרְדְּעִ֖ים עַל־אֶ֥רֶץ מִצְרָֽיִם׃ %
% ---------- פרק 8 | פסוק 4 ----------
 \textbf{ד} וַיִּקְרָ֨א פַרְעֹ֜ה לְמֹשֶׁ֣ה וּֽלְאַהֲרֹ֗ן וַיֹּ֙אמֶר֙ הַעְתִּ֣ירוּ אֶל־יְהֹוָ֔ה וְיָסֵר֙ הַֽצְפַרְדְּעִ֔ים מִמֶּ֖נִּי וּמֵֽעַמִּ֑י וַאֲשַׁלְּחָה֙ אֶת־הָעָ֔ם וְיִזְבְּח֖וּ לַיהֹוָֽה׃ %
% ---------- פרק 8 | פסוק 5 ----------
 \textbf{ה} וַיֹּ֨אמֶר מֹשֶׁ֣ה לְפַרְעֹה֮ הִתְפָּאֵ֣ר עָלַי֒ לְמָתַ֣י ׀ אַעְתִּ֣יר לְךָ֗ וְלַעֲבָדֶ֙יךָ֙ וּֽלְעַמְּךָ֔ לְהַכְרִית֙ הַֽצְפַרְדְּעִ֔ים מִמְּךָ֖ וּמִבָּתֶּ֑יךָ רַ֥ק בַּיְאֹ֖ר תִּשָּׁאַֽרְנָה׃ \Rashi{\textbf{התפאר עלי.}כמו "היתפאר הגרזן על החצב בו" (ישעיהו י') – משתבח לומר אני גדול ממך, ונטי"ר בלעז. וכן התפאר עלי – השתבח להתחכם ולשאל דבר גדול, ולומר שלא אוכל לעשותו: למתי אעתיר לך. את אשר אעתיר לך היום על הכרתת הצפרדעים, למתי תרצה שיכרתו? ותראה אם אשלים דברי למועד שתקבע לי. אלו נאמר "מתי" אעתיר, היה משמע מתי אתפלל, עכשו שנאמר "למתי", משמע אני היום אתפלל עליך, שיכרתו הצפרדעים לזמן שתקבע עלי; אמר לאיזה יום תרצה שיכרתו: אעתיר, העתירו, והעתרתי. ולא נאמר אעתר, עתרו, ועתרתי, מפני שכל לשון עתר הרבות פלל הוא, וכאשר יאמר ארבה, הרבו, והרביתי, לשון מפעיל, כך יאמר אעתיר, העתירו, והעתרתי דברים; ואב לכלם "והעתרתם עלי דבריכם" (יחזקאל ל"ה) – הרביתם:}%
% ---------- פרק 8 | פסוק 6 ----------
 \textbf{ו} וַיֹּ֖אמֶר לְמָחָ֑ר וַיֹּ֙אמֶר֙ כִּדְבָ֣רְךָ֔ לְמַ֣עַן תֵּדַ֔ע כִּי־אֵ֖ין כַּיהֹוָ֥ה אֱלֹהֵֽינוּ׃ \Rashi{\textbf{ויאמר למחר.}התפלל היום, שיכרתו למחר:}%
% ---------- פרק 8 | פסוק 7 ----------
 \textbf{ז} וְסָר֣וּ הַֽצְפַרְדְּעִ֗ים מִמְּךָ֙ וּמִבָּ֣תֶּ֔יךָ וּמֵעֲבָדֶ֖יךָ וּמֵעַמֶּ֑ךָ רַ֥ק בַּיְאֹ֖ר תִּשָּׁאַֽרְנָה׃ %
% ---------- פרק 8 | פסוק 8 ----------
 \textbf{ח} וַיֵּצֵ֥א מֹשֶׁ֛ה וְאַהֲרֹ֖ן מֵעִ֣ם פַּרְעֹ֑ה וַיִּצְעַ֤ק מֹשֶׁה֙ אֶל־יְהֹוָ֔ה עַל־דְּבַ֥ר הַֽצְפַרְדְּעִ֖ים אֲשֶׁר־שָׂ֥ם לְפַרְעֹֽה׃ \Rashi{\textbf{ויצא, ויצעק.}מיד, שיכרתו למחר:}%
% ---------- פרק 8 | פסוק 9 ----------
 \textbf{ט} וַיַּ֥עַשׂ יְהֹוָ֖ה כִּדְבַ֣ר מֹשֶׁ֑ה וַיָּמֻ֙תוּ֙ הַֽצְפַרְדְּעִ֔ים מִן־הַבָּתִּ֥ים מִן־הַחֲצֵרֹ֖ת וּמִן־הַשָּׂדֹֽת׃ %
% ---------- פרק 8 | פסוק 10 ----------
 \textbf{י} וַיִּצְבְּר֥וּ אֹתָ֖ם חֳמָרִ֣ם חֳמָרִ֑ם וַתִּבְאַ֖שׁ הָאָֽרֶץ׃ \Rashi{\textbf{חמרם חמרם.}צבורים צבורים, כתרגומו, דגורין – גלין:}%
% ---------- פרק 8 | פסוק 11 ----------
 \textbf{יא} וַיַּ֣רְא פַּרְעֹ֗ה כִּ֤י הָֽיְתָה֙ הָֽרְוָחָ֔ה וְהַכְבֵּד֙ אֶת־לִבּ֔וֹ וְלֹ֥א שָׁמַ֖ע אֲלֵהֶ֑ם כַּאֲשֶׁ֖ר דִּבֶּ֥ר יְהֹוָֽה׃ \{ס\} \Rashi{\textbf{והכבד את לבו.}לשון פעול הוא, כמו "הלוך ונסוע" (בראשית י"ב), וכן "והכות את מואב" (מלכים ב ג׳:כ״ד), "ושאול לו באלהים" (שמואל א כ״ב:י״ג), "הכה ופצע" (מלכים א כ'): כאשר דבר ה'. והיכן דבר? ולא ישמע אליכם פרעה:}%
% ---------- פרק 8 | פסוק 12 ----------
 \textbf{יב} וַיֹּ֣אמֶר יְהֹוָה֮ אֶל־מֹשֶׁה֒ אֱמֹר֙ אֶֽל־אַהֲרֹ֔ן נְטֵ֣ה אֶֽת־מַטְּךָ֔ וְהַ֖ךְ אֶת־עֲפַ֣ר הָאָ֑רֶץ וְהָיָ֥ה לְכִנִּ֖ם בְּכׇל־אֶ֥רֶץ מִצְרָֽיִם׃ \Rashi{\textbf{אמר אל אהרן.}לא היה העפר כדאי ללקות ע"י משה, לפי שהגן עליו כשהרג את המצרי ויטמנהו בחול, ולקה על ידי אהרן (שמות רבה):}%
% ---------- פרק 8 | פסוק 13 ----------
 \textbf{יג} וַיַּֽעֲשׂוּ־כֵ֗ן וַיֵּט֩ אַהֲרֹ֨ן אֶת־יָד֤וֹ בְמַטֵּ֙הוּ֙ וַיַּךְ֙ אֶת־עֲפַ֣ר הָאָ֔רֶץ וַתְּהִי֙ הַכִּנָּ֔ם בָּאָדָ֖ם וּבַבְּהֵמָ֑ה כׇּל־עֲפַ֥ר הָאָ֛רֶץ הָיָ֥ה כִנִּ֖ים בְּכׇל־אֶ֥רֶץ מִצְרָֽיִם׃ \Rashi{\textbf{ותהי הכנם.}הרחישה, פדוליר"א בלעז:}%
% ---------- פרק 8 | פסוק 14 ----------
 \textbf{יד} וַיַּעֲשׂוּ־כֵ֨ן הַחַרְטֻמִּ֧ים בְּלָטֵיהֶ֛ם לְהוֹצִ֥יא אֶת־הַכִּנִּ֖ים וְלֹ֣א יָכֹ֑לוּ וַתְּהִי֙ הַכִּנָּ֔ם בָּאָדָ֖ם וּבַבְּהֵמָֽה׃ \Rashi{\textbf{להוציא את הכנים.}לבראתם ולהוציאם ממקום אחר: ולא יכלו. שאין השד שולט על בריה פחותה מכשעורה (סנהדרין ס"ז):}%
% ---------- פרק 8 | פסוק 15 ----------
 \textbf{טו} וַיֹּאמְר֤וּ הַֽחַרְטֻמִּם֙ אֶל־פַּרְעֹ֔ה אֶצְבַּ֥ע אֱלֹהִ֖ים הִ֑וא וַיֶּחֱזַ֤ק לֵב־פַּרְעֹה֙ וְלֹֽא־שָׁמַ֣ע אֲלֵהֶ֔ם כַּאֲשֶׁ֖ר דִּבֶּ֥ר יְהֹוָֽה׃ \{ס\} \Rashi{\textbf{אצבע אלהים הוא.}מכה זו אינה על ידי כשפים, מאת המקום היא (שמות רבה): כאשר דבר ה'. ולא ישמע אליכם פרעה:}\Peirush{מזה שהחרטומים לא יכלו לברוא כינים לומדים שאין השד יכול לברוא בריה פחות מכשעורה\Makor{סנהדרין סז:}}\Peirush{אצבע אלהים, מכאן שלקו מצרים באצבע, והיא האצבע שעתידה להפרע מסנחריב\Makor{סנהדרין צה:}}%
% ---------- פרק 8 | פסוק 16 ----------
 \textbf{טז} וַיֹּ֨אמֶר יְהֹוָ֜ה אֶל־מֹשֶׁ֗ה הַשְׁכֵּ֤ם בַּבֹּ֙קֶר֙ וְהִתְיַצֵּב֙ לִפְנֵ֣י פַרְעֹ֔ה הִנֵּ֖ה יוֹצֵ֣א הַמָּ֑יְמָה וְאָמַרְתָּ֣ אֵלָ֗יו כֹּ֚ה אָמַ֣ר יְהֹוָ֔ה שַׁלַּ֥ח עַמִּ֖י וְיַֽעַבְדֻֽנִי׃ %
% ---------- פרק 8 | פסוק 17 ----------
 \textbf{יז} כִּ֣י אִם־אֵינְךָ֮ מְשַׁלֵּ֣חַ אֶת־עַמִּי֒ הִנְנִי֩ מַשְׁלִ֨יחַ בְּךָ֜ וּבַעֲבָדֶ֧יךָ וּֽבְעַמְּךָ֛ וּבְבָתֶּ֖יךָ אֶת־הֶעָרֹ֑ב וּמָ֨לְא֜וּ בָּתֵּ֤י מִצְרַ֙יִם֙ אֶת־הֶ֣עָרֹ֔ב וְגַ֥ם הָאֲדָמָ֖ה אֲשֶׁר־הֵ֥ם עָלֶֽיהָ׃ \Rashi{\textbf{משליח בך.}מגרה בך, וכן "ושן בהמות אשלח בם" (דברים ל"ב), לשון שסוי, אינציט"ר בלעז: את הערב. כל מיני החיות רעות ונחשים ועקרבים בערבוביא, והיו משחיתים בהם; ויש טעם בדבר באגדה בכל מכה ומכה למה זו ולמה זו; בטכסיסי מלחמות מלכים בא עליהם, כסדר מלכות כשצרה על עיר, בתחלה מקלקל מעינותיה, ואחר כך תוקעין עליהם ומריעין בשופרות ליראם ולבהלם, וכן הצפרדעים מקרקרים והומים וכו', כדאיתא במדרש רבי תנחומא (פ' בא):}%
% ---------- פרק 8 | פסוק 18 ----------
 \textbf{יח} וְהִפְלֵיתִי֩ בַיּ֨וֹם הַה֜וּא אֶת־אֶ֣רֶץ גֹּ֗שֶׁן אֲשֶׁ֤ר עַמִּי֙ עֹמֵ֣ד עָלֶ֔יהָ לְבִלְתִּ֥י הֱיֽוֹת־שָׁ֖ם עָרֹ֑ב לְמַ֣עַן תֵּדַ֔ע כִּ֛י אֲנִ֥י יְהֹוָ֖ה בְּקֶ֥רֶב הָאָֽרֶץ׃ \Rashi{\textbf{והפליתי.}והפרשתי, וכן "והפלה ה'" (שמות ט'), וכן "לא נפלאת היא ממך" (דברים ל') – לא מבדלת ומפרשת היא ממך: למען תדע כי אני ה' בקרב הארץ. אף על פי ששכינתי בשמים גזרתי מתקימת בתחתונים:}%
% ---------- פרק 8 | פסוק 19 ----------
 \textbf{יט} וְשַׂמְתִּ֣י פְדֻ֔ת בֵּ֥ין עַמִּ֖י וּבֵ֣ין עַמֶּ֑ךָ לְמָחָ֥ר יִהְיֶ֖ה הָאֹ֥ת הַזֶּֽה׃ \Rashi{\textbf{ושמתי פדת.}שיבדיל בין עמי ובין עמך:}%
% ---------- פרק 8 | פסוק 20 ----------
 \textbf{כ} וַיַּ֤עַשׂ יְהֹוָה֙ כֵּ֔ן וַיָּבֹא֙ עָרֹ֣ב כָּבֵ֔ד בֵּ֥יתָה פַרְעֹ֖ה וּבֵ֣ית עֲבָדָ֑יו וּבְכׇל־אֶ֧רֶץ מִצְרַ֛יִם תִּשָּׁחֵ֥ת הָאָ֖רֶץ מִפְּנֵ֥י הֶעָרֹֽב׃ \Rashi{\textbf{תשחת הארץ.}נשחתת הארץ, אתחבלת ארעא:}%
% ---------- פרק 8 | פסוק 21 ----------
 \textbf{כא} וַיִּקְרָ֣א פַרְעֹ֔ה אֶל־מֹשֶׁ֖ה וּֽלְאַהֲרֹ֑ן וַיֹּ֗אמֶר לְכ֛וּ זִבְח֥וּ לֵאלֹֽהֵיכֶ֖ם בָּאָֽרֶץ׃ \Rashi{\textbf{זבחו לאלהיכם בארץ.}במקומכם ולא תלכו במדבר:}%
% ---------- פרק 8 | פסוק 22 ----------
 \textbf{כב} וַיֹּ֣אמֶר מֹשֶׁ֗ה לֹ֤א נָכוֹן֙ לַעֲשׂ֣וֹת כֵּ֔ן כִּ֚י תּוֹעֲבַ֣ת מִצְרַ֔יִם נִזְבַּ֖ח לַיהֹוָ֣ה אֱלֹהֵ֑ינוּ הֵ֣ן נִזְבַּ֞ח אֶת־תּוֹעֲבַ֥ת מִצְרַ֛יִם לְעֵינֵיהֶ֖ם וְלֹ֥א יִסְקְלֻֽנוּ׃ \Rashi{\textbf{תועבת מצרים.}יראת מצרים, כמו "ולמלכם תועבת בני עמון" (מלכים ב כ״ג:י״ג), ואצל ישראל קורא אותה תועבה; ועוד יש לומר בלשון אחר, תועבת מצרים, דבר שנאוי הוא למצרים זביחה שאנו זובחים, שהרי יראתם אנו זובחים: ולא יסקלנו. בתמיה:}%
% ---------- פרק 8 | פסוק 23 ----------
 \textbf{כג} דֶּ֚רֶךְ שְׁלֹ֣שֶׁת יָמִ֔ים נֵלֵ֖ךְ בַּמִּדְבָּ֑ר וְזָבַ֙חְנוּ֙ לַֽיהֹוָ֣ה אֱלֹהֵ֔ינוּ כַּאֲשֶׁ֖ר יֹאמַ֥ר אֵלֵֽינוּ׃ %
% ---------- פרק 8 | פסוק 24 ----------
 \textbf{כד} וַיֹּ֣אמֶר פַּרְעֹ֗ה אָנֹכִ֞י אֲשַׁלַּ֤ח אֶתְכֶם֙ וּזְבַחְתֶּ֞ם לַיהֹוָ֤ה אֱלֹֽהֵיכֶם֙ בַּמִּדְבָּ֔ר רַ֛ק הַרְחֵ֥ק לֹא־תַרְחִ֖יקוּ לָלֶ֑כֶת הַעְתִּ֖ירוּ בַּעֲדִֽי׃ %
% ---------- פרק 8 | פסוק 25 ----------
 \textbf{כה} וַיֹּ֣אמֶר מֹשֶׁ֗ה הִנֵּ֨ה אָנֹכִ֜י יוֹצֵ֤א מֵֽעִמָּךְ֙ וְהַעְתַּרְתִּ֣י אֶל־יְהֹוָ֔ה וְסָ֣ר הֶעָרֹ֗ב מִפַּרְעֹ֛ה מֵעֲבָדָ֥יו וּמֵעַמּ֖וֹ מָחָ֑ר רַ֗ק אַל־יֹסֵ֤ף פַּרְעֹה֙ הָתֵ֔ל לְבִלְתִּי֙ שַׁלַּ֣ח אֶת־הָעָ֔ם לִזְבֹּ֖חַ לַֽיהֹוָֽה׃ \Rashi{\textbf{התל.}כמו להתל:}%
% ---------- פרק 8 | פסוק 26 ----------
 \textbf{כו} וַיֵּצֵ֥א מֹשֶׁ֖ה מֵעִ֣ם פַּרְעֹ֑ה וַיֶּעְתַּ֖ר אֶל־יְהֹוָֽה׃ \Rashi{\textbf{ויעתר אל ה'.}נתאמץ בתפלה, וכן אם בא לומר ויעתיר, היה יכול לומר, ומשמע וירבה בתפלה, וכשהוא אומר בלשון ויפעל, משמע וירבה להתפלל:}%
% ---------- פרק 8 | פסוק 27 ----------
 \textbf{כז} וַיַּ֤עַשׂ יְהֹוָה֙ כִּדְבַ֣ר מֹשֶׁ֔ה וַיָּ֙סַר֙ הֶעָרֹ֔ב מִפַּרְעֹ֖ה מֵעֲבָדָ֣יו וּמֵעַמּ֑וֹ לֹ֥א נִשְׁאַ֖ר אֶחָֽד׃ \Rashi{\textbf{ויסר הערב.}ולא מתו כמו שמתו הצפרדעים, שאם מתו היה להם הנאה בעורותם (תנחומא):}%
% ---------- פרק 8 | פסוק 28 ----------
 \textbf{כח} וַיַּכְבֵּ֤ד פַּרְעֹה֙ אֶת־לִבּ֔וֹ גַּ֖ם בַּפַּ֣עַם הַזֹּ֑את וְלֹ֥א שִׁלַּ֖ח אֶת־הָעָֽם׃ \{פ\} \Rashi{\textbf{גם בפעם הזאת.}אע"פ שאמר אנכי אשלח אתכם, לא קיים הבטחתו:}%
% ---------- פרק 9 | פסוק 1 ----------
 \textbf{א} וַיֹּ֤אמֶר יְהֹוָה֙ אֶל־מֹשֶׁ֔ה בֹּ֖א אֶל־פַּרְעֹ֑ה וְדִבַּרְתָּ֣ אֵלָ֗יו כֹּֽה־אָמַ֤ר יְהֹוָה֙ אֱלֹהֵ֣י הָֽעִבְרִ֔ים שַׁלַּ֥ח אֶת־עַמִּ֖י וְיַֽעַבְדֻֽנִי׃ %
% ---------- פרק 9 | פסוק 2 ----------
 \textbf{ב} כִּ֛י אִם־מָאֵ֥ן אַתָּ֖ה לְשַׁלֵּ֑חַ וְעוֹדְךָ֖ מַחֲזִ֥יק בָּֽם׃ \Rashi{\textbf{מחזיק בם.}אוחז בם, כמו "והחזיקה במבשיו" (דברים כ"ה):}%
% ---------- פרק 9 | פסוק 3 ----------
 \textbf{ג} הִנֵּ֨ה יַד־יְהֹוָ֜ה הוֹיָ֗ה בְּמִקְנְךָ֙ אֲשֶׁ֣ר בַּשָּׂדֶ֔ה בַּסּוּסִ֤ים בַּֽחֲמֹרִים֙ בַּגְּמַלִּ֔ים בַּבָּקָ֖ר וּבַצֹּ֑אן דֶּ֖בֶר כָּבֵ֥ד מְאֹֽד׃ \Rashi{\textbf{הנה יד ה' הויה.}לשון הוה, כי כן יאמר בלשון נקבה על שעבר היתה, ועל העתיד תהיה, ועל העומד הויה, כמו עושה, רוצה, רועה:}%
% ---------- פרק 9 | פסוק 4 ----------
 \textbf{ד} וְהִפְלָ֣ה יְהֹוָ֔ה בֵּ֚ין מִקְנֵ֣ה יִשְׂרָאֵ֔ל וּבֵ֖ין מִקְנֵ֣ה מִצְרָ֑יִם וְלֹ֥א יָמ֛וּת מִכׇּל־לִבְנֵ֥י יִשְׂרָאֵ֖ל דָּבָֽר׃ \Rashi{\textbf{והפלה.}והבדיל:}%
% ---------- פרק 9 | פסוק 5 ----------
 \textbf{ה} וַיָּ֥שֶׂם יְהֹוָ֖ה מוֹעֵ֣ד לֵאמֹ֑ר מָחָ֗ר יַעֲשֶׂ֧ה יְהֹוָ֛ה הַדָּבָ֥ר הַזֶּ֖ה בָּאָֽרֶץ׃ %
% ---------- פרק 9 | פסוק 6 ----------
 \textbf{ו} וַיַּ֨עַשׂ יְהֹוָ֜ה אֶת־הַדָּבָ֤ר הַזֶּה֙ מִֽמׇּחֳרָ֔ת וַיָּ֕מׇת כֹּ֖ל מִקְנֵ֣ה מִצְרָ֑יִם וּמִמִּקְנֵ֥ה בְנֵֽי־יִשְׂרָאֵ֖ל לֹא־מֵ֥ת אֶחָֽד׃ %
% ---------- פרק 9 | פסוק 7 ----------
 \textbf{ז} וַיִּשְׁלַ֣ח פַּרְעֹ֔ה וְהִנֵּ֗ה לֹא־מֵ֛ת מִמִּקְנֵ֥ה יִשְׂרָאֵ֖ל עַד־אֶחָ֑ד וַיִּכְבַּד֙ לֵ֣ב פַּרְעֹ֔ה וְלֹ֥א שִׁלַּ֖ח אֶת־הָעָֽם׃ \{פ\} %
% ---------- פרק 9 | פסוק 8 ----------
 \textbf{ח} וַיֹּ֣אמֶר יְהֹוָה֮ אֶל־מֹשֶׁ֣ה וְאֶֽל־אַהֲרֹן֒ קְח֤וּ לָכֶם֙ מְלֹ֣א חׇפְנֵיכֶ֔ם פִּ֖יחַ כִּבְשָׁ֑ן וּזְרָק֥וֹ מֹשֶׁ֛ה הַשָּׁמַ֖יְמָה לְעֵינֵ֥י פַרְעֹֽה׃ \Rashi{\textbf{מלא חפניכם.}ילויני"ש בלעז: פיח כבשן. דבר הנפח מן הגחלים עמומים הנשרפים בכבשן, ובלעז אולבי"ש: פיח. לשון הפחה, שהרוח מפיחן ומפריחן: וזרקו משה. וכל דבר הנזרק בכח, אינו נזרק אלא ביד אחת, הרי נסים הרבה – אחד שהחזיק קמצו של משה מלא חפנים שלו ושל אהרן, ואחד שהלך האבק על כל ארץ מצרים (תנחומא):}%
% ---------- פרק 9 | פסוק 9 ----------
 \textbf{ט} וְהָיָ֣ה לְאָבָ֔ק עַ֖ל כׇּל־אֶ֣רֶץ מִצְרָ֑יִם וְהָיָ֨ה עַל־הָאָדָ֜ם וְעַל־הַבְּהֵמָ֗ה לִשְׁחִ֥ין פֹּרֵ֛חַ אֲבַעְבֻּעֹ֖ת בְּכׇל־אֶ֥רֶץ מִצְרָֽיִם׃ \Rashi{\textbf{לשחין פרח אבעבעת.}כתרגומו, לשיחנא סגי אבעבועין, שעל ידו צומחין בהן בועות: שחין. לשון חמימות; והרבה יש בלשון משנה, "שנה שחונה" (יומא נ"ג):}%
% ---------- פרק 9 | פסוק 10 ----------
 \textbf{י} וַיִּקְח֞וּ אֶת־פִּ֣יחַ הַכִּבְשָׁ֗ן וַיַּֽעַמְדוּ֙ לִפְנֵ֣י פַרְעֹ֔ה וַיִּזְרֹ֥ק אֹת֛וֹ מֹשֶׁ֖ה הַשָּׁמָ֑יְמָה וַיְהִ֗י שְׁחִין֙ אֲבַעְבֻּעֹ֔ת פֹּרֵ֕חַ בָּאָדָ֖ם וּבַבְּהֵמָֽה׃ \Rashi{\textbf{באדם ובבהמה.}ואם תאמר מאין היו להם הבהמות, והלא כבר נאמר וימת כל מקנה מצרים? לא נגזרה גזרה אלא על אותן שבשדות בלבד, שנאמר במקנך אשר בשדה, והירא את דבר ה' הכניס את מקנהו אל הבתים, כן שנויה במכילתא אצל "ויקח שש מאות רכב בחור" (שמות י"ד):}\Peirush{השחין היה לח מבחוץ ויבש מבפנים, שחין – יבש מבפנים, אבעבעת פרח – לח מבחוץ\Makor{בבא קמא פ:, בכורות מא.}}%
% ---------- פרק 9 | פסוק 11 ----------
 \textbf{יא} וְלֹֽא־יָכְל֣וּ הַֽחַרְטֻמִּ֗ים לַעֲמֹ֛ד לִפְנֵ֥י מֹשֶׁ֖ה מִפְּנֵ֣י הַשְּׁחִ֑ין כִּֽי־הָיָ֣ה הַשְּׁחִ֔ין בַּֽחַרְטֻמִּ֖ם וּבְכׇל־מִצְרָֽיִם׃ %
% ---------- פרק 9 | פסוק 12 ----------
 \textbf{יב} וַיְחַזֵּ֤ק יְהֹוָה֙ אֶת־לֵ֣ב פַּרְעֹ֔ה וְלֹ֥א שָׁמַ֖ע אֲלֵהֶ֑ם כַּאֲשֶׁ֛ר דִּבֶּ֥ר יְהֹוָ֖ה אֶל־מֹשֶֽׁה׃ \{ס\} %
% ---------- פרק 9 | פסוק 13 ----------
 \textbf{יג} וַיֹּ֤אמֶר יְהֹוָה֙ אֶל־מֹשֶׁ֔ה הַשְׁכֵּ֣ם בַּבֹּ֔קֶר וְהִתְיַצֵּ֖ב לִפְנֵ֣י פַרְעֹ֑ה וְאָמַרְתָּ֣ אֵלָ֗יו כֹּֽה־אָמַ֤ר יְהֹוָה֙ אֱלֹהֵ֣י הָֽעִבְרִ֔ים שַׁלַּ֥ח אֶת־עַמִּ֖י וְיַֽעַבְדֻֽנִי׃ %
% ---------- פרק 9 | פסוק 14 ----------
 \textbf{יד} כִּ֣י ׀ בַּפַּ֣עַם הַזֹּ֗את אֲנִ֨י שֹׁלֵ֜חַ אֶת־כׇּל־מַגֵּפֹתַי֙ אֶֽל־לִבְּךָ֔ וּבַעֲבָדֶ֖יךָ וּבְעַמֶּ֑ךָ בַּעֲב֣וּר תֵּדַ֔ע כִּ֛י אֵ֥ין כָּמֹ֖נִי בְּכׇל־הָאָֽרֶץ׃ \Rashi{\textbf{את כל מגפתי.}למדנו מכאן שמכת בכורות שקולה כנגד כל המכות:}%
% ---------- פרק 9 | פסוק 15 ----------
 \textbf{טו} כִּ֤י עַתָּה֙ שָׁלַ֣חְתִּי אֶת־יָדִ֔י וָאַ֥ךְ אוֹתְךָ֛ וְאֶֽת־עַמְּךָ֖ בַּדָּ֑בֶר וַתִּכָּחֵ֖ד מִן־הָאָֽרֶץ׃ \Rashi{\textbf{כי עתה שלחתי את ידי וגו'.}כי אלו רציתי, כשהיתה ידי במקנך שהכיתים בדבר, שלחתיה והכיתי אותך ואת עמך עם הבהמות, ותכחד מן הארץ, אבל בעבור זאת העמדתיך וגו':}%
% ---------- פרק 9 | פסוק 16 ----------
 \textbf{טז} וְאוּלָ֗ם בַּעֲב֥וּר זֹאת֙ הֶעֱמַדְתִּ֔יךָ בַּעֲב֖וּר הַרְאֹתְךָ֣ אֶת־כֹּחִ֑י וּלְמַ֛עַן סַפֵּ֥ר שְׁמִ֖י בְּכׇל־הָאָֽרֶץ׃ %
% ---------- פרק 9 | פסוק 17 ----------
 \textbf{יז} עוֹדְךָ֖ מִסְתּוֹלֵ֣ל בְּעַמִּ֑י לְבִלְתִּ֖י שַׁלְּחָֽם׃ \Rashi{\textbf{עודך מסתולל בעמי.}כתרגומו, כבישת בה בעמי, והוא מגזרת מסלה, דמתרגמינן ארח כבישא, ובלעז קלקי"ר; וכבר פרשתי בסוף ויהי מקץ, כל תבה שתחלת יסודה סמ"ך והיא באה לדבר בלשון מתפעל, נותן התי"ו של שמוש באמצע אותיות של עקר, כגון זו, וכגון "ויסתבל החגב" (קה' י"ב), מגזרת סבל, "כי תשתרר עלינו" (במדבר ט"ז), מגזרת נגיד ושר, "משתכל הוית" (דניאל ז'):}%
% ---------- פרק 9 | פסוק 18 ----------
 \textbf{יח} הִנְנִ֤י מַמְטִיר֙ כָּעֵ֣ת מָחָ֔ר בָּרָ֖ד כָּבֵ֣ד מְאֹ֑ד אֲשֶׁ֨ר לֹא־הָיָ֤ה כָמֹ֙הוּ֙ בְּמִצְרַ֔יִם לְמִן־הַיּ֥וֹם הִוָּסְדָ֖הֿ וְעַד־עָֽתָּה׃ \Rashi{\textbf{כעת מחר.}כעת הזאת למחר; שרט לו שריטה בכתל – למחר כשתגיע חמה לכאן, ירד הברד: הוסדה. שנתיסדה; וכל תבה שתחלת יסודה יו"ד, כגון יסד, ילד, ידע, יסר, כשהיא מתפעלת, תבא הוי"ו במקום היו"ד, כמו הוסדה, הולדה (הושע ב'), ויודע, "ויולד ליוסף" (בראשית מ"ו), "בדברים לא יוסר עבד" (משלי כ"ט):}%
% ---------- פרק 9 | פסוק 19 ----------
 \textbf{יט} וְעַתָּ֗ה שְׁלַ֤ח הָעֵז֙ אֶֽת־מִקְנְךָ֔ וְאֵ֛ת כׇּל־אֲשֶׁ֥ר לְךָ֖ בַּשָּׂדֶ֑ה כׇּל־הָאָדָ֨ם וְהַבְּהֵמָ֜ה אֲשֶֽׁר־יִמָּצֵ֣א בַשָּׂדֶ֗ה וְלֹ֤א יֵֽאָסֵף֙ הַבַּ֔יְתָה וְיָרַ֧ד עֲלֵהֶ֛ם הַבָּרָ֖ד וָמֵֽתוּ׃ \Rashi{\textbf{שלח העז.}כתרגומו, שלח כנוש, וכן "יושבי הגבים העיזו" (ישעיה י'), "העזו בני בנימן" (ירמיה ו'): ולא יאסף הביתה. לשון הכנסה הוא:}%
% ---------- פרק 9 | פסוק 20 ----------
 \textbf{כ} הַיָּרֵא֙ אֶת־דְּבַ֣ר יְהֹוָ֔ה מֵֽעַבְדֵ֖י פַּרְעֹ֑ה הֵנִ֛יס אֶת־עֲבָדָ֥יו וְאֶת־מִקְנֵ֖הוּ אֶל־הַבָּתִּֽים׃ \Rashi{\textbf{הניס.}הבריח:}%
% ---------- פרק 9 | פסוק 21 ----------
 \textbf{כא} וַאֲשֶׁ֥ר לֹא־שָׂ֛ם לִבּ֖וֹ אֶל־דְּבַ֣ר יְהֹוָ֑ה וַֽיַּעֲזֹ֛ב אֶת־עֲבָדָ֥יו וְאֶת־מִקְנֵ֖הוּ בַּשָּׂדֶֽה׃ \{פ\} %
% ---------- פרק 9 | פסוק 22 ----------
 \textbf{כב} וַיֹּ֨אמֶר יְהֹוָ֜ה אֶל־מֹשֶׁ֗ה נְטֵ֤ה אֶת־יָֽדְךָ֙ עַל־הַשָּׁמַ֔יִם וִיהִ֥י בָרָ֖ד בְּכׇל־אֶ֣רֶץ מִצְרָ֑יִם עַל־הָאָדָ֣ם וְעַל־הַבְּהֵמָ֗ה וְעַ֛ל כׇּל־עֵ֥שֶׂב הַשָּׂדֶ֖ה בְּאֶ֥רֶץ מִצְרָֽיִם׃ \Rashi{\textbf{על השמים.}לצד השמים. ומדרש אגדה: הגביהו הקב"ה למשה למעלה מן השמים:}%
% ---------- פרק 9 | פסוק 23 ----------
 \textbf{כג} וַיֵּ֨ט מֹשֶׁ֣ה אֶת־מַטֵּ֘הוּ֮ עַל־הַשָּׁמַ֒יִם֒ וַֽיהֹוָ֗ה נָתַ֤ן קֹלֹת֙ וּבָרָ֔ד וַתִּ֥הֲלַךְ אֵ֖שׁ אָ֑רְצָה וַיַּמְטֵ֧ר יְהֹוָ֛ה בָּרָ֖ד עַל־אֶ֥רֶץ מִצְרָֽיִם׃ %
% ---------- פרק 9 | פסוק 24 ----------
 \textbf{כד} וַיְהִ֣י בָרָ֔ד וְאֵ֕שׁ מִתְלַקַּ֖חַת בְּת֣וֹךְ הַבָּרָ֑ד כָּבֵ֣ד מְאֹ֔ד אֲ֠שֶׁ֠ר לֹֽא־הָיָ֤ה כָמֹ֙הוּ֙ בְּכׇל־אֶ֣רֶץ מִצְרַ֔יִם מֵאָ֖ז הָיְתָ֥ה לְגֽוֹי׃ \Rashi{\textbf{מתלקחת בתוך הברד.}נס בתוך נס, האש והברד מערבין, והברד מים הוא, ולעשות רצון קונם עשו שלום ביניהם:}%
% ---------- פרק 9 | פסוק 25 ----------
 \textbf{כה} וַיַּ֨ךְ הַבָּרָ֜ד בְּכׇל־אֶ֣רֶץ מִצְרַ֗יִם אֵ֚ת כׇּל־אֲשֶׁ֣ר בַּשָּׂדֶ֔ה מֵאָדָ֖ם וְעַד־בְּהֵמָ֑ה וְאֵ֨ת כׇּל־עֵ֤שֶׂב הַשָּׂדֶה֙ הִכָּ֣ה הַבָּרָ֔ד וְאֶת־כׇּל־עֵ֥ץ הַשָּׂדֶ֖ה שִׁבֵּֽר׃ %
% ---------- פרק 9 | פסוק 26 ----------
 \textbf{כו} רַ֚ק בְּאֶ֣רֶץ גֹּ֔שֶׁן אֲשֶׁר־שָׁ֖ם בְּנֵ֣י יִשְׂרָאֵ֑ל לֹ֥א הָיָ֖ה בָּרָֽד׃ %
% ---------- פרק 9 | פסוק 27 ----------
 \textbf{כז} וַיִּשְׁלַ֣ח פַּרְעֹ֗ה וַיִּקְרָא֙ לְמֹשֶׁ֣ה וּֽלְאַהֲרֹ֔ן וַיֹּ֥אמֶר אֲלֵהֶ֖ם חָטָ֣אתִי הַפָּ֑עַם יְהֹוָה֙ הַצַּדִּ֔יק וַאֲנִ֥י וְעַמִּ֖י הָרְשָׁעִֽים׃ \Peirush{בסוף מסכת ידים מובא הפסוק "מי ה' אשר אשמע בקלו" (ה, ב), ע"מ לסיים בפס' טוב, המשנה אומרת וכשלקה פרעה, אמר ה' הצדיק\Makor{משנה ידים ד, ח}}%
% ---------- פרק 9 | פסוק 28 ----------
 \textbf{כח} הַעְתִּ֙ירוּ֙ אֶל־יְהֹוָ֔ה וְרַ֕ב מִֽהְיֹ֛ת קֹלֹ֥ת אֱלֹהִ֖ים וּבָרָ֑ד וַאֲשַׁלְּחָ֣ה אֶתְכֶ֔ם וְלֹ֥א תֹסִפ֖וּן לַעֲמֹֽד׃ \Rashi{\textbf{ורב.}די לו במה שהוריד כבר:}%
% ---------- פרק 9 | פסוק 29 ----------
 \textbf{כט} וַיֹּ֤אמֶר אֵלָיו֙ מֹשֶׁ֔ה כְּצֵאתִי֙ אֶת־הָעִ֔יר אֶפְרֹ֥שׂ אֶת־כַּפַּ֖י אֶל־יְהֹוָ֑ה הַקֹּל֣וֹת יֶחְדָּל֗וּן וְהַבָּרָד֙ לֹ֣א יִֽהְיֶה־ע֔וֹד לְמַ֣עַן תֵּדַ֔ע כִּ֥י לַיהֹוָ֖ה הָאָֽרֶץ׃ \Rashi{\textbf{כצאתי את העיר.}מן העיר, אבל בתוך העיר לא התפלל לפי שהיתה מלאה גלולים (שמות רבה):}%
% ---------- פרק 9 | פסוק 30 ----------
 \textbf{ל} וְאַתָּ֖ה וַעֲבָדֶ֑יךָ יָדַ֕עְתִּי כִּ֚י טֶ֣רֶם תִּֽירְא֔וּן מִפְּנֵ֖י יְהֹוָ֥ה אֱלֹהִֽים׃ \Rashi{\textbf{טרם תיראון.}עדין לא תיראון; וכן כל טרם שבמקרא "עדין לא" הוא, ואינו לשון קדם, "טרם ישכבו" (בראשית י"ט) – עד לא שכיבו, "טרם יצמח" (שם ב') – עד לא צמח. אף זה כן הוא – ידעתי כי עדין אינכם יראים, ומשתהיה הרוחה תעמדו בקלקולכם:}%
% ---------- פרק 9 | פסוק 31 ----------
 \textbf{לא} וְהַפִּשְׁתָּ֥ה וְהַשְּׂעֹרָ֖ה נֻכָּ֑תָה כִּ֤י הַשְּׂעֹרָה֙ אָבִ֔יב וְהַפִּשְׁתָּ֖ה גִּבְעֹֽל׃ \Rashi{\textbf{והפשתה והשערה נכתה.}נשברה, לשון "פרעה נכה" (מלכים ב כ"ג), "נכאים" (ישעיהו ט"ז), וכן לא נכו; ולא יתכן לפרשו לשון הכאה, שאין נו"ן במקום ה"א לפרש נכתה כמו הכתה, נכו כמו הכו, אלא הנו"ן שרש בתבה והרי הוא מגזרת "ושפו עצמותיו" (איוב ל"ג): כי השערה אביב. כבר בכרה ועומדת בקשיה, ונשתברו ונפלו, וכן הפשתה גדלה כבר, והקשה לעמד בגבעוליה: השערה אביב. עמדה באביה, לשון "באבי הנחל" (שיר ו'):}\Peirush{נאמר אביב כאן ונאמר במנחת הביכורים, מה כאן שעורים אף שם שעורים\Makor{מנחות סח:, פד.}}%
% ---------- פרק 9 | פסוק 32 ----------
 \textbf{לב} וְהַחִטָּ֥ה וְהַכֻּסֶּ֖מֶת לֹ֣א נֻכּ֑וּ כִּ֥י אֲפִילֹ֖ת הֵֽנָּה׃ \Rashi{\textbf{כי אפילת הנה.}מאחרות, ועדין היו רכות ויכולות לעמד בפני קשה; ואע"פ שנאמר "ואת כל עשב השדה הכה הברד", יש לפרש פשוטו של מקרא בעשבים העומדים בקלחם, הראוים ללקות בברד. ובמדרש רבי תנחומא יש מרבותינו שנחלקו על זאת, ודרשו כי אפילת, פלאי פלאות נעשו להם שלא לקו:}%
% ---------- פרק 9 | פסוק 33 ----------
 \textbf{לג} וַיֵּצֵ֨א מֹשֶׁ֜ה מֵעִ֤ם פַּרְעֹה֙ אֶת־הָעִ֔יר וַיִּפְרֹ֥שׂ כַּפָּ֖יו אֶל־יְהֹוָ֑ה וַֽיַּחְדְּל֤וּ הַקֹּלוֹת֙ וְהַבָּרָ֔ד וּמָטָ֖ר לֹא־נִתַּ֥ךְ אָֽרְצָה׃ \Rashi{\textbf{לא נתך.}לא הגיע, ואף אותן שהיו באויר לא הגיעו לארץ; ודומה לו: "ותתך עלינו האלה והשבועה" (דניאל ט') דעזרא – ותגיע עלינו. ומנחם בן סרוק חברו בחלק "כהתוך כסף" (יחזקאל כ"ב), לשון יציקת מתכת, ורואה אני את דבריו, כתרגומו ויצק – ואתיך, לצקת – לאתכא. אף זה לא נתך לארץ – לא הוצק לארץ:}\Peirush{הרואה את המקום שנעשה נס אבני אלגביש חייב להודות לה'. אלגביש - אבנים שעמדו באויר על גב [בגלל] איש, וירדו על גב איש. עמדו על גב איש — זה משה שנאמר "והאיש משה" (יב, ג). ירדו על גב איש – זה יהושע שנאמר "יהושע בן נון איש אשר רוח בו" (במדבר כז, יח) ונאמר "וה' השליך עליהם אבנים גדלות" (יהושע י, יא)\Makor{ברכות נד:}}%
% ---------- פרק 9 | פסוק 34 ----------
 \textbf{לד} וַיַּ֣רְא פַּרְעֹ֗ה כִּֽי־חָדַ֨ל הַמָּטָ֧ר וְהַבָּרָ֛ד וְהַקֹּלֹ֖ת וַיֹּ֣סֶף לַחֲטֹ֑א וַיַּכְבֵּ֥ד לִבּ֖וֹ ה֥וּא וַעֲבָדָֽיו׃ %
% ---------- פרק 9 | פסוק 35 ----------
 \textbf{לה} וַיֶּֽחֱזַק֙ לֵ֣ב פַּרְעֹ֔ה וְלֹ֥א שִׁלַּ֖ח אֶת־בְּנֵ֣י יִשְׂרָאֵ֑ל כַּאֲשֶׁ֛ר דִּבֶּ֥ר יְהֹוָ֖ה בְּיַד־מֹשֶֽׁה׃ \{פ\} %
\addparasha{חומש שמות | פרשת בא}
% ---------- פרק 10 | פסוק 1 ----------
 \textbf{א} וַיֹּ֤אמֶר יְהֹוָה֙ אֶל־מֹשֶׁ֔ה בֹּ֖א אֶל־פַּרְעֹ֑ה כִּֽי־אֲנִ֞י הִכְבַּ֤דְתִּי אֶת־לִבּוֹ֙ וְאֶת־לֵ֣ב עֲבָדָ֔יו לְמַ֗עַן שִׁתִ֛י אֹתֹתַ֥י אֵ֖לֶּה בְּקִרְבּֽוֹ׃ \Rashi{\textbf{ויאמר ה' אל משה בא אל פרעה.}והתרה בו: שיתי. שומי – שאשית אני:}%
% ---------- פרק 10 | פסוק 2 ----------
 \textbf{ב} וּלְמַ֡עַן תְּסַפֵּר֩ בְּאׇזְנֵ֨י בִנְךָ֜ וּבֶן־בִּנְךָ֗ אֵ֣ת אֲשֶׁ֤ר הִתְעַלַּ֙לְתִּי֙ בְּמִצְרַ֔יִם וְאֶת־אֹתֹתַ֖י אֲשֶׁר־שַׂ֣מְתִּי בָ֑ם וִֽידַעְתֶּ֖ם כִּי־אֲנִ֥י יְהֹוָֽה׃ \Rashi{\textbf{התעללתי.}שחקתי, כמו "כי התעללת בי" (במדבר כ"ב), "הלא כאשר התעולל בהם" (שמואל א' ו'), האמור במצרים, ואינו לשון פועל ומעללים, שאם כן היה לו לכתב עוללתי, כמו "ועולל למו כאשר עוללת לי" (איכה א'), "אשר עולל לי" (שם): ולמען תספר. בתורה להודיע לדורות:}%
% ---------- פרק 10 | פסוק 3 ----------
 \textbf{ג} וַיָּבֹ֨א מֹשֶׁ֣ה וְאַהֲרֹן֮ אֶל־פַּרְעֹה֒ וַיֹּאמְר֣וּ אֵלָ֗יו כֹּֽה־אָמַ֤ר יְהֹוָה֙ אֱלֹהֵ֣י הָֽעִבְרִ֔ים עַד־מָתַ֣י מֵאַ֔נְתָּ לֵעָנֹ֖ת מִפָּנָ֑י שַׁלַּ֥ח עַמִּ֖י וְיַֽעַבְדֻֽנִי׃ \Rashi{\textbf{לענת.}כתרגומו, "לאתכנעא", והוא מגזרת עני – מאנת להיות עני ושפל מפני:}%
% ---------- פרק 10 | פסוק 4 ----------
 \textbf{ד} כִּ֛י אִם־מָאֵ֥ן אַתָּ֖ה לְשַׁלֵּ֣חַ אֶת־עַמִּ֑י הִנְנִ֨י מֵבִ֥יא מָחָ֛ר אַרְבֶּ֖ה בִּגְבֻלֶֽךָ׃ %
% ---------- פרק 10 | פסוק 5 ----------
 \textbf{ה} וְכִסָּה֙ אֶת־עֵ֣ין הָאָ֔רֶץ וְלֹ֥א יוּכַ֖ל לִרְאֹ֣ת אֶת־הָאָ֑רֶץ וְאָכַ֣ל ׀ אֶת־יֶ֣תֶר הַפְּלֵטָ֗ה הַנִּשְׁאֶ֤רֶת לָכֶם֙ מִן־הַבָּרָ֔ד וְאָכַל֙ אֶת־כׇּל־הָעֵ֔ץ הַצֹּמֵ֥חַ לָכֶ֖ם מִן־הַשָּׂדֶֽה׃ \Rashi{\textbf{את עין הארץ.}את מראה הארץ: ולא יוכל. הרואה לראת את הארץ, ולשון קצרה דבר:}%
% ---------- פרק 10 | פסוק 6 ----------
 \textbf{ו} וּמָלְא֨וּ בָתֶּ֜יךָ וּבָתֵּ֣י כׇל־עֲבָדֶ֘יךָ֮ וּבָתֵּ֣י כׇל־מִצְרַ֒יִם֒ אֲשֶׁ֨ר לֹֽא־רָא֤וּ אֲבֹתֶ֙יךָ֙ וַאֲב֣וֹת אֲבֹתֶ֔יךָ מִיּ֗וֹם הֱיוֹתָם֙ עַל־הָ֣אֲדָמָ֔ה עַ֖ד הַיּ֣וֹם הַזֶּ֑ה וַיִּ֥פֶן וַיֵּצֵ֖א מֵעִ֥ם פַּרְעֹֽה׃ %
% ---------- פרק 10 | פסוק 7 ----------
 \textbf{ז} וַיֹּאמְרוּ֩ עַבְדֵ֨י פַרְעֹ֜ה אֵלָ֗יו עַד־מָתַי֙ יִהְיֶ֨ה זֶ֥ה לָ֙נוּ֙ לְמוֹקֵ֔שׁ שַׁלַּח֙ אֶת־הָ֣אֲנָשִׁ֔ים וְיַֽעַבְד֖וּ אֶת־יְהֹוָ֣ה אֱלֹהֵיהֶ֑ם הֲטֶ֣רֶם תֵּדַ֔ע כִּ֥י אָבְדָ֖ה מִצְרָֽיִם׃ \Rashi{\textbf{הטרם תדע.}העוד לא ידעת כי אבדה מצרים:}%
% ---------- פרק 10 | פסוק 8 ----------
 \textbf{ח} וַיּוּשַׁ֞ב אֶת־מֹשֶׁ֤ה וְאֶֽת־אַהֲרֹן֙ אֶל־פַּרְעֹ֔ה וַיֹּ֣אמֶר אֲלֵהֶ֔ם לְכ֥וּ עִבְד֖וּ אֶת־יְהֹוָ֣ה אֱלֹהֵיכֶ֑ם מִ֥י וָמִ֖י הַהֹלְכִֽים׃ \Rashi{\textbf{ויושב.}הושבו על ידי שליח, ששלחו אחריהם והשיבום אל פרעה:}%
% ---------- פרק 10 | פסוק 9 ----------
 \textbf{ט} וַיֹּ֣אמֶר מֹשֶׁ֔ה בִּנְעָרֵ֥ינוּ וּבִזְקֵנֵ֖ינוּ נֵלֵ֑ךְ בְּבָנֵ֨ינוּ וּבִבְנוֹתֵ֜נוּ בְּצֹאנֵ֤נוּ וּבִבְקָרֵ֙נוּ֙ נֵלֵ֔ךְ כִּ֥י חַג־יְהֹוָ֖ה לָֽנוּ׃ %
% ---------- פרק 10 | פסוק 10 ----------
 \textbf{י} וַיֹּ֣אמֶר אֲלֵהֶ֗ם יְהִ֨י כֵ֤ן יְהֹוָה֙ עִמָּכֶ֔ם כַּאֲשֶׁ֛ר אֲשַׁלַּ֥ח אֶתְכֶ֖ם וְאֶֽת־טַפְּכֶ֑ם רְא֕וּ כִּ֥י רָעָ֖ה נֶ֥גֶד פְּנֵיכֶֽם׃ \Rashi{\textbf{כאשר אשלח אתכם ואת טפכם.}אף כי אשלח גם את הצאן ואת הבקר כאשר אמרתם: ראו כי רעה נגד פניכם. כתרגומו. ומ"א שמעתי, כוכב אחד יש ששמו רעה, אמר להם פרעה, רואה אני באצטגנינות שלי אותו כוכב עולה לקראתכם במדבר, והוא סימן דם והריגה; וכשחטאו ישראל בעגל ובקש הקב"ה להרגם, אמר משה בתפלתו, למה יאמרו מצרים לאמר "ברעה" הוציאם (שמות ל"ב), זו היא שאמר להם, ראו כי רעה נגד פניכם; מיד וינחם ה' על הרעה והפך את הדם לדם מילה, שמל יהושע אותם, וזהו שנ' "היום גלותי את חרפת מצרים מעליכם" (יהושע ה'), שהיו אומרים לכם, דם אנו רואין עליכם במדבר:}%
% ---------- פרק 10 | פסוק 11 ----------
 \textbf{יא} לֹ֣א כֵ֗ן לְכֽוּ־נָ֤א הַגְּבָרִים֙ וְעִבְד֣וּ אֶת־יְהֹוָ֔ה כִּ֥י אֹתָ֖הּ אַתֶּ֣ם מְבַקְשִׁ֑ים וַיְגָ֣רֶשׁ אֹתָ֔ם מֵאֵ֖ת פְּנֵ֥י פַרְעֹֽה׃ \{ס\} \Rashi{\textbf{לא כן.}כאשר אמרתם להוליך הטף עמכם, אלא לכו הגברים ועבדו את ה': כי אתה אתם מבקשים. כי אותה בקשתם עד הנה – נזבחה לאלהינו, ואין דרך הטף לזבח: ויגרש אותם. הרי זה לשון קצר ולא פרש מי המגרש:}%
% ---------- פרק 10 | פסוק 12 ----------
 \textbf{יב} וַיֹּ֨אמֶר יְהֹוָ֜ה אֶל־מֹשֶׁ֗ה נְטֵ֨ה יָדְךָ֜ עַל־אֶ֤רֶץ מִצְרַ֙יִם֙ בָּֽאַרְבֶּ֔ה וְיַ֖עַל עַל־אֶ֣רֶץ מִצְרָ֑יִם וְיֹאכַל֙ אֶת־כׇּל־עֵ֣שֶׂב הָאָ֔רֶץ אֵ֛ת כׇּל־אֲשֶׁ֥ר הִשְׁאִ֖יר הַבָּרָֽד׃ \Rashi{\textbf{בארבה.}בשביל מכת הארבה:}%
% ---------- פרק 10 | פסוק 13 ----------
 \textbf{יג} וַיֵּ֨ט מֹשֶׁ֣ה אֶת־מַטֵּ֘הוּ֮ עַל־אֶ֣רֶץ מִצְרַ֒יִם֒ וַֽיהֹוָ֗ה נִהַ֤ג רֽוּחַ־קָדִים֙ בָּאָ֔רֶץ כׇּל־הַיּ֥וֹם הַה֖וּא וְכׇל־הַלָּ֑יְלָה הַבֹּ֣קֶר הָיָ֔ה וְר֙וּחַ֙ הַקָּדִ֔ים נָשָׂ֖א אֶת־הָאַרְבֶּֽה׃ \Rashi{\textbf{ורוח הקדים.}רוח מזרחית נשא את הארבה, לפי שבא כנגדו, שמצרים בדרומית מערבית היתה, כמו שמפרש במקום אחר:}%
% ---------- פרק 10 | פסוק 14 ----------
 \textbf{יד} וַיַּ֣עַל הָֽאַרְבֶּ֗ה עַ֚ל כׇּל־אֶ֣רֶץ מִצְרַ֔יִם וַיָּ֕נַח בְּכֹ֖ל גְּב֣וּל מִצְרָ֑יִם כָּבֵ֣ד מְאֹ֔ד לְ֠פָנָ֠יו לֹא־הָ֨יָה כֵ֤ן אַרְבֶּה֙ כָּמֹ֔הוּ וְאַחֲרָ֖יו לֹ֥א יִֽהְיֶה־כֵּֽן׃ \Rashi{\textbf{ואחריו לא יהיה כן.}ואותו שהיה בימי יואל שנ' "כמוהו לא נהיה מן העולם" (יואל ב'), למדנו שהיה כבד משל משה. (אותו של יואל היה) ע"י מינין הרבה, שהיו יחד ארבה, ילק, חסיל, גזם, אבל של משה לא היה אלא מין אחד, וכמוהו לא היה ולא יהיה:}%
% ---------- פרק 10 | פסוק 15 ----------
 \textbf{טו} וַיְכַ֞ס אֶת־עֵ֣ין כׇּל־הָאָ֘רֶץ֮ וַתֶּחְשַׁ֣ךְ הָאָ֒רֶץ֒ וַיֹּ֜אכַל אֶת־כׇּל־עֵ֣שֶׂב הָאָ֗רֶץ וְאֵת֙ כׇּל־פְּרִ֣י הָעֵ֔ץ אֲשֶׁ֥ר הוֹתִ֖יר הַבָּרָ֑ד וְלֹא־נוֹתַ֨ר כׇּל־יֶ֧רֶק בָּעֵ֛ץ וּבְעֵ֥שֶׂב הַשָּׂדֶ֖ה בְּכׇל־אֶ֥רֶץ מִצְרָֽיִם׃ \Rashi{\textbf{כל ירק.}עלה ירק, וירדו"רא בלעז:}%
% ---------- פרק 10 | פסוק 16 ----------
 \textbf{טז} וַיְמַהֵ֣ר פַּרְעֹ֔ה לִקְרֹ֖א לְמֹשֶׁ֣ה וּֽלְאַהֲרֹ֑ן וַיֹּ֗אמֶר חָטָ֛אתִי לַיהֹוָ֥ה אֱלֹֽהֵיכֶ֖ם וְלָכֶֽם׃ %
% ---------- פרק 10 | פסוק 17 ----------
 \textbf{יז} וְעַתָּ֗ה שָׂ֣א נָ֤א חַטָּאתִי֙ אַ֣ךְ הַפַּ֔עַם וְהַעְתִּ֖ירוּ לַיהֹוָ֣ה אֱלֹהֵיכֶ֑ם וְיָסֵר֙ מֵֽעָלַ֔י רַ֖ק אֶת־הַמָּ֥וֶת הַזֶּֽה׃ %
% ---------- פרק 10 | פסוק 18 ----------
 \textbf{יח} וַיֵּצֵ֖א מֵעִ֣ם פַּרְעֹ֑ה וַיֶּעְתַּ֖ר אֶל־יְהֹוָֽה׃ %
% ---------- פרק 10 | פסוק 19 ----------
 \textbf{יט} וַיַּהֲפֹ֨ךְ יְהֹוָ֤ה רֽוּחַ־יָם֙ חָזָ֣ק מְאֹ֔ד וַיִּשָּׂא֙ אֶת־הָ֣אַרְבֶּ֔ה וַיִּתְקָעֵ֖הוּ יָ֣מָּה סּ֑וּף לֹ֤א נִשְׁאַר֙ אַרְבֶּ֣ה אֶחָ֔ד בְּכֹ֖ל גְּב֥וּל מִצְרָֽיִם׃ \Rashi{\textbf{רוח ים.}רוח מערבי: ימה סוף. אומר אני שים סוף היה מקצתו במערב, כנגד כל רוח דרומית, וגם במזרחה של ארץ ישראל, לפיכך רוח ים תקעו לארבה בימה סוף כנגדו. וכן מצינו לענין תחומין שהוא פונה לצד מזרח, שנאמר "מים סוף ועד ים פלשתים" (לקמן כג לא) ממזרח למערב, שים פלשתים במערב היה, שנאמר בפלשתים "יושבי חבל הים גוי כרתים" (צפניה ב'): לא נשאר ארבה אחד. אף המלוחים שמלחו מהם:}%
% ---------- פרק 10 | פסוק 20 ----------
 \textbf{כ} וַיְחַזֵּ֥ק יְהֹוָ֖ה אֶת־לֵ֣ב פַּרְעֹ֑ה וְלֹ֥א שִׁלַּ֖ח אֶת־בְּנֵ֥י יִשְׂרָאֵֽל׃ \{פ\} %
% ---------- פרק 10 | פסוק 21 ----------
 \textbf{כא} וַיֹּ֨אמֶר יְהֹוָ֜ה אֶל־מֹשֶׁ֗ה נְטֵ֤ה יָֽדְךָ֙ עַל־הַשָּׁמַ֔יִם וִ֥יהִי חֹ֖שֶׁךְ עַל־אֶ֣רֶץ מִצְרָ֑יִם וְיָמֵ֖שׁ חֹֽשֶׁךְ׃ \Rashi{\textbf{וימש חשך.}ויחשיך עליהם חשך יותר מחשכו של לילה, חשך של לילה יאמיש ויחשיך עוד: וימש. כמו ויאמש; יש לנו תבות הרבה חסרות אל"ף, לפי שאין הברת האל"ף נכרת כל כך, אין הכתוב מקפיד על חסרונה, כגון "ולא יהל שם ערבי" (ישעיהו י"ג), כמו לא יאהל – לא יטה אהלו, וכן "ותזרני חיל" (שמואל ב' כ"ב), כמו ותאזרני; ואנקלוס תרגם לשון הסרה, כמו לא ימוש; "בתר דיעדי קבל ליליא" – כשיגיע סמוך לאור היום; אבל אין הדבור מישב על הוי"ו של וימש, לפי שהוא כתוב אחר ויהי חשך. ומ"א פותרו לשון "ממשש בצהרים" (דברים כ"ח), שהיה כפול ומכפל ועב עד שהיה בו ממש:}%
% ---------- פרק 10 | פסוק 22 ----------
 \textbf{כב} וַיֵּ֥ט מֹשֶׁ֛ה אֶת־יָד֖וֹ עַל־הַשָּׁמָ֑יִם וַיְהִ֧י חֹֽשֶׁךְ־אֲפֵלָ֛ה בְּכׇל־אֶ֥רֶץ מִצְרַ֖יִם שְׁלֹ֥שֶׁת יָמִֽים׃ \Rashi{\textbf{שלשת ימים.}שלוש של ימים, טרציינ"א בלעז, וכן שבעת ימים בכל מקום שטיי"נא של ימים: ויהי חשך אפלה: שלשת ימים. חשך של אפל, שלא ראו איש את אחיו אותן ג' ימים. ועוד שלושת ימים אחרים חשך מכפל על זה, שלא קמו איש מתחתיו – יושב אין יכול לעמד, ועומד אין יכול לישב; ולמה הביא עליהם חשך? שהיו בישראל באותו הדור רשעים ולא היו רוצים לצאת, ומתו בשלושת ימי אפלה, כדי שלא יראו מצריים במפלתם ויאמרו, אף הם לוקים כמונו. ועוד, שחפשו ישראל וראו את כליהם, וכשיצאו והיו שואלים מהן והיו אומרים אין בידנו כלום, אומר לו, אני ראיתיו בביתך, ובמקום פלוני הוא (שמות רבה):}%
% ---------- פרק 10 | פסוק 23 ----------
 \textbf{כג} לֹֽא־רָא֞וּ אִ֣ישׁ אֶת־אָחִ֗יו וְלֹא־קָ֛מוּ אִ֥ישׁ מִתַּחְתָּ֖יו שְׁלֹ֣שֶׁת יָמִ֑ים וּֽלְכׇל־בְּנֵ֧י יִשְׂרָאֵ֛ל הָ֥יָה א֖וֹר בְּמוֹשְׁבֹתָֽם׃ %
% ---------- פרק 10 | פסוק 24 ----------
 \textbf{כד} וַיִּקְרָ֨א פַרְעֹ֜ה אֶל־מֹשֶׁ֗ה וַיֹּ֙אמֶר֙ לְכוּ֙ עִבְד֣וּ אֶת־יְהֹוָ֔ה רַ֛ק צֹאנְכֶ֥ם וּבְקַרְכֶ֖ם יֻצָּ֑ג גַּֽם־טַפְּכֶ֖ם יֵלֵ֥ךְ עִמָּכֶֽם׃ \Rashi{\textbf{יצג.}יהא מצג במקומו:}%
% ---------- פרק 10 | פסוק 25 ----------
 \textbf{כה} וַיֹּ֣אמֶר מֹשֶׁ֔ה גַּם־אַתָּ֛ה תִּתֵּ֥ן בְּיָדֵ֖נוּ זְבָחִ֣ים וְעֹלֹ֑ת וְעָשִׂ֖ינוּ לַיהֹוָ֥ה אֱלֹהֵֽינוּ׃ \Rashi{\textbf{גם אתה תתן.}לא דיך שמקננו ילך עמנו אלא גם משלך תתן:}\Peirush{נאמר "ויחגו לי במדבר" (ה, א), יכולנו לחשוב חג בלי קרבן, קמ"ל שלא, שנאמר גם אתה תתן בידנו זבחים ועלת\Makor{חגיגה י:}}\Peirush{אסור לקחת קרבנות מהגוים, אך קשה, שנאמר גם אתה תתן בידנו זבחים ועלת\Makor{עבודה זרה כד.}}\Peirush{בני נח לא הקריבו שלמים, אך קשה שהרי עוד לפני מתן תורה הקריבו זבחים – שלמים\Makor{זבחים קטז.}}%
% ---------- פרק 10 | פסוק 26 ----------
 \textbf{כו} וְגַם־מִקְנֵ֜נוּ יֵלֵ֣ךְ עִמָּ֗נוּ לֹ֤א תִשָּׁאֵר֙ פַּרְסָ֔ה כִּ֚י מִמֶּ֣נּוּ נִקַּ֔ח לַעֲבֹ֖ד אֶת־יְהֹוָ֣ה אֱלֹהֵ֑ינוּ וַאֲנַ֣חְנוּ לֹֽא־נֵדַ֗ע מַֽה־נַּעֲבֹד֙ אֶת־יְהֹוָ֔ה עַד־בֹּאֵ֖נוּ שָֽׁמָּה׃ \Rashi{\textbf{פרסה.}פרסת רגל, פלנ"טא בלעז: לא נדע מה נעבד. כמה תכבד העבודה, שמא ישאל יותר ממה שיש בידנו:}%
% ---------- פרק 10 | פסוק 27 ----------
 \textbf{כז} וַיְחַזֵּ֥ק יְהֹוָ֖ה אֶת־לֵ֣ב פַּרְעֹ֑ה וְלֹ֥א אָבָ֖ה לְשַׁלְּחָֽם׃ %
% ---------- פרק 10 | פסוק 28 ----------
 \textbf{כח} וַיֹּֽאמֶר־ל֥וֹ פַרְעֹ֖ה לֵ֣ךְ מֵעָלָ֑י הִשָּׁ֣מֶר לְךָ֗ אַל־תֹּ֙סֶף֙ רְא֣וֹת פָּנַ֔י כִּ֗י בְּי֛וֹם רְאֹתְךָ֥ פָנַ֖י תָּמֽוּת׃ %
% ---------- פרק 10 | פסוק 29 ----------
 \textbf{כט} וַיֹּ֥אמֶר מֹשֶׁ֖ה כֵּ֣ן דִּבַּ֑רְתָּ לֹא־אֹסִ֥ף ע֖וֹד רְא֥וֹת פָּנֶֽיךָ׃ \{פ\} \Rashi{\textbf{כן דברת.}יפה דברת ובזמנו דברת. אמת שלא אוסיף עוד ראות פניך (מכילתא ושמות רבה):}%
% ---------- פרק 11 | פסוק 1 ----------
 \textbf{א} וַיֹּ֨אמֶר יְהֹוָ֜ה אֶל־מֹשֶׁ֗ה ע֣וֹד נֶ֤גַע אֶחָד֙ אָבִ֤יא עַל־פַּרְעֹה֙ וְעַל־מִצְרַ֔יִם אַֽחֲרֵי־כֵ֕ן יְשַׁלַּ֥ח אֶתְכֶ֖ם מִזֶּ֑ה כְּשַׁ֨לְּח֔וֹ כָּלָ֕ה גָּרֵ֛שׁ יְגָרֵ֥שׁ אֶתְכֶ֖ם מִזֶּֽה׃ \Rashi{\textbf{כלה.}גמירא; כליל, כלכם ישלח:}%
% ---------- פרק 11 | פסוק 2 ----------
 \textbf{ב} דַּבֶּר־נָ֖א בְּאׇזְנֵ֣י הָעָ֑ם וְיִשְׁאֲל֞וּ אִ֣ישׁ ׀ מֵאֵ֣ת רֵעֵ֗הוּ וְאִשָּׁה֙ מֵאֵ֣ת רְעוּתָ֔הּ כְּלֵי־כֶ֖סֶף וּכְלֵ֥י זָהָֽב׃ \Rashi{\textbf{דבר נא.}אין נא אלא לשון בקשה, בבקשה ממך הזהירם על כך, שלא יאמר אותו צדיק אברהם "ועבדום וענו אותם" (בראשית ט"ו), קים בהם, "ואחרי כן יצאו ברכש גדול" (שם), לא קים בהם:}\Peirush{אין נא אלא לשון בקשה. אמר ה' למשה בבקשה ממך לך ואמור להם לישראל בבקשה מכם שאלו ממצרים כלי כסף וכלי זהב, שלא יאמר אברהם "ועבדום וענו אותם" (בראשית טו, יג) קיים בהם, "ואחרי כן יצאו ברכוש גדול" (שם) לא קיים בהם. אמרו לו הלואי שנצא בעצנו. משל לאדם שהיה חבוש בבית האסורים, והיו אומרים לו בני אדם מוציאין אותך למחר מבית האסורין ונותנין לך ממון הרבה, ואומר להם: בבקשה מכם הוציאוני היום ואיני מבקש כלום\Makor{ברכות ט.-ט:}}%
% ---------- פרק 11 | פסוק 3 ----------
 \textbf{ג} וַיִּתֵּ֧ן יְהֹוָ֛ה אֶת־חֵ֥ן הָעָ֖ם בְּעֵינֵ֣י מִצְרָ֑יִם גַּ֣ם ׀ הָאִ֣ישׁ מֹשֶׁ֗ה גָּד֤וֹל מְאֹד֙ בְּאֶ֣רֶץ מִצְרַ֔יִם בְּעֵינֵ֥י עַבְדֵֽי־פַרְעֹ֖ה וּבְעֵינֵ֥י הָעָֽם׃ \{ס\} %
% ---------- פרק 11 | פסוק 4 ----------
 \textbf{ד} וַיֹּ֣אמֶר מֹשֶׁ֔ה כֹּ֖ה אָמַ֣ר יְהֹוָ֑ה כַּחֲצֹ֣ת הַלַּ֔יְלָה אֲנִ֥י יוֹצֵ֖א בְּת֥וֹךְ מִצְרָֽיִם׃ \Rashi{\textbf{ויאמר משה כה אמר ה'.}בעמדו לפני פרעה נאמרה לו נבואה זו, שהרי משיצא מלפניו לא הוסיף ראות פניו: כחצת הלילה. כהחלק הלילה: כחצת. כמו "בעלות" (מלכים א י"ח), "בחרות אפם בנו" (תהילים קכ"ד), זהו פשוטו לישבו על אפניו, שאין חצות שם דבר של חצי. ורבותינו דרשוהו כמו כחצי הלילה, ואמרו שאמר משה "כחצות", שמשמע סמוך לו, או לפניו או לאחריו, ולא אמר "בחצות", שמא יטעו אצטגניני פרעה ויאמרו משה בדאי הוא (ברכות ד'):}\Peirush{דוד ידע מתי בדיוק חצות הלילה, אך קשה, שאפילו משה לא ידע, שנאמר כחצות\Makor{ברכות ג:}}\Peirush{משה ידע מתי חצות, הוא אמר כחצות שמא יטעו אצטגניני פרעה בחישוב חצות הלילה ויאמרו משה בדאי הוא\Makor{ברכות ד.}}%
% ---------- פרק 11 | פסוק 5 ----------
 \textbf{ה} וּמֵ֣ת כׇּל־בְּכוֹר֮ בְּאֶ֣רֶץ מִצְרַ֒יִם֒ מִבְּכ֤וֹר פַּרְעֹה֙ הַיֹּשֵׁ֣ב עַל־כִּסְא֔וֹ עַ֚ד בְּכ֣וֹר הַשִּׁפְחָ֔ה אֲשֶׁ֖ר אַחַ֣ר הָרֵחָ֑יִם וְכֹ֖ל בְּכ֥וֹר בְּהֵמָֽה׃ \Rashi{\textbf{עד בכור השבי.}למה לקו השבויים? כדי שלא יאמרו, יראתם תבעה עלבונם והביאה פרענות על מצרים: מבכור פרעה עד בכור השפחה. כל הפחותים מבכור פרעה וחשובים מבכור השפחה היו בכלל; ולמה לקו בני השפחות? שאף הם היו משעבדים בהם ושמחים בצרתם: וכל בכור בהמה. לפי שהיו עובדין לה. כשהקב"ה נפרע מן האמה, נפרע מאלהיה:}%
% ---------- פרק 11 | פסוק 6 ----------
 \textbf{ו} וְהָ֥יְתָ֛ה צְעָקָ֥ה גְדֹלָ֖ה בְּכׇל־אֶ֣רֶץ מִצְרָ֑יִם אֲשֶׁ֤ר כָּמֹ֙הוּ֙ לֹ֣א נִהְיָ֔תָה וְכָמֹ֖הוּ לֹ֥א תֹסִֽף׃ %
% ---------- פרק 11 | פסוק 7 ----------
 \textbf{ז} וּלְכֹ֣ל ׀ בְּנֵ֣י יִשְׂרָאֵ֗ל לֹ֤א יֶֽחֱרַץ־כֶּ֙לֶב֙ לְשֹׁנ֔וֹ לְמֵאִ֖ישׁ וְעַד־בְּהֵמָ֑ה לְמַ֙עַן֙ תֵּֽדְע֔וּן אֲשֶׁר֙ יַפְלֶ֣ה יְהֹוָ֔ה בֵּ֥ין מִצְרַ֖יִם וּבֵ֥ין יִשְׂרָאֵֽל׃ \Rashi{\textbf{לא יחרץ כלב לשונו.}אומר אני שהוא לשון שנון – לא ישנן, וכן "לא חרץ לבני ישראל לאיש את לשונו" (יהושע י') – לא שנן, "אז תחרץ" (שמואל ב ה׳:כ״ד) – תשתנן, "למורג חרוץ" (ישעיהו מ"א) – שנון, "מחשבות חרוץ" (משלי כ"א) – אדם חריף ושנון, "ויד חרוצים תעשיר" (שם י') – חריפים, סוחרים שנונים: אשר יפלה. יבדיל:}\Peirush{הרואה כלב בחלום, ישכים ויאמר ולכל בני ישראל לא יחרץ כלב לשנו לפני שיתגשם הפסוק "והכלבים עזי נפש" (ישעיהו נו, יא)\Makor{ברכות נו:}}%
% ---------- פרק 11 | פסוק 8 ----------
 \textbf{ח} וְיָרְד֣וּ כׇל־עֲבָדֶ֩יךָ֩ אֵ֨לֶּה אֵלַ֜י וְהִשְׁתַּֽחֲווּ־לִ֣י לֵאמֹ֗ר צֵ֤א אַתָּה֙ וְכׇל־הָעָ֣ם אֲשֶׁר־בְּרַגְלֶ֔יךָ וְאַחֲרֵי־כֵ֖ן אֵצֵ֑א וַיֵּצֵ֥א מֵֽעִם־פַּרְעֹ֖ה בׇּחֳרִי־אָֽף׃ \{ס\} \Rashi{\textbf{וירדו כל עבדיך.}חלק כבוד למלכות, שהרי בסוף ירד פרעה בעצמו אליו בלילה "ויאמר קומו צאו מתוך עמי" (שמות י"ב), ולא אמר לו משה מתחלה וירדת אלי והשתחוית לי: אשר ברגליך. ההולכים אחר עצתך והלוכך: ואחרי כן אצא. עם כל העם מארצך: ויצא מעם פרעה. כשגמר דבריו, יצא מלפניו: בחרי אף. על שאמר לו "אל תסף ראות פני" (שמות י'):}\Peirush{לעולם תהא אימת מלכות עליך שמשה אמר לפרעה וירדו כל עבדיך אך לא אתה\Makor{זבחים קב.}}\Peirush{כל מקום שכתוב בתורה חרי אף, כתוב גם איך זה מתבטא, אך קשה שהרי כאן זה לא התבטא, מכאן שמשה סטר לפרעה ויצא\Makor{זבחים קב., מנחות צח.}}%
% ---------- פרק 11 | פסוק 9 ----------
 \textbf{ט} וַיֹּ֤אמֶר יְהֹוָה֙ אֶל־מֹשֶׁ֔ה לֹא־יִשְׁמַ֥ע אֲלֵיכֶ֖ם פַּרְעֹ֑ה לְמַ֛עַן רְב֥וֹת מוֹפְתַ֖י בְּאֶ֥רֶץ מִצְרָֽיִם׃ \Rashi{\textbf{למען רבות מופתי.}מכת בכורות וקריעת ים סוף ולנער את מצרים:}%
% ---------- פרק 11 | פסוק 10 ----------
 \textbf{י} וּמֹשֶׁ֣ה וְאַהֲרֹ֗ן עָשׂ֛וּ אֶת־כׇּל־הַמֹּפְתִ֥ים הָאֵ֖לֶּה לִפְנֵ֣י פַרְעֹ֑ה וַיְחַזֵּ֤ק יְהֹוָה֙ אֶת־לֵ֣ב פַּרְעֹ֔ה וְלֹֽא־שִׁלַּ֥ח אֶת־בְּנֵֽי־יִשְׂרָאֵ֖ל מֵאַרְצֽוֹ׃ \{ס\} \Rashi{\textbf{משה ואהרן עשו וגו'.}כבר כתב לנו זאת בכל המופתים, ולא שנאה כאן אלא בשביל לסמכה לפרשה שלאחריה:}%
% ---------- פרק 12 | פסוק 1 ----------
 \textbf{א} וַיֹּ֤אמֶר יְהֹוָה֙ אֶל־מֹשֶׁ֣ה וְאֶֽל־אַהֲרֹ֔ן בְּאֶ֥רֶץ מִצְרַ֖יִם לֵאמֹֽר׃ \Rashi{\textbf{ויאמר ה' אל משה ואל אהרן.}בשביל שאהרן עשה וטרח במופתים כמשה, חלק לו כבוד זה במצוה ראשונה, שכללו עם משה בדבור: בארץ מצרים. חוץ לכרך, או אינו אלא בתוך הכרך? ת"ל (לעיל ט כט) "כצאתי את העיר וגו'", ומה תפלה קלה לא התפלל בתוך הכרך, דבור חמור לא כל שכן. ומפני מה לא נדבר עמו בתוך הכרך? לפי שהיתה מלאה גלולים:}\Peirush{עדות החודש מותרת בקרובים, שנאמר ויאמר ה' אל משה ואל אהרן... החדש הזה לכם שתהא כשרה במשה ואהרן שהיו אחים\Makor{ראש השנה כב.}}\Peirush{חייבים שלושה דיינים לעדות החודש, שנאמר ויאמר ה' אל משה ואל אהרן... החדש הזה לכם, וכיון שאין עושים בית דין זוגי, מוסיפים עוד אחד\Makor{ראש השנה כה:}}%
% ---------- פרק 12 | פסוק 2 ----------
 \textbf{ב} הַחֹ֧דֶשׁ הַזֶּ֛ה לָכֶ֖ם רֹ֣אשׁ חֳדָשִׁ֑ים רִאשׁ֥וֹן הוּא֙ לָכֶ֔ם לְחׇדְשֵׁ֖י הַשָּׁנָֽה׃ \Rashi{\textbf{החדש הזה.}הראהו לבנה בחדושה ואמר לו כשהירח מתחדש יהיה לך ראש חדש (מכילתא). ואין מקרא יוצא מידי פשוטו, על חדש ניסן אמר לו, זה יהיה ראש לסדר מנין החדשים, שיהא איר קרוי שני, סיון שלישי: הזה. נתקשה משה על מולד הלבנה, באיזו שעור תראה ותהיה ראויה לקדש, והראה לו באצבע את הלבנה ברקיע ואמר לו כזה ראה וקדש (שם). וכיצד הראהו? והלא לא היה נדבר עמו אלא ביום? שנ' "ויהי ביום דבר ה'" (שמות ו'), "ביום צותו" (ויקרא ז'), "מן היום אשר צוה ה' והלאה" (במדבר ט"ו)? אלא סמוך לשקיעת החמה נאמרה לו פרשה זו והראהו עם חשכה:}\Peirush{כל המברך על החדש בזמנו כאילו מקבל פני שכינה, אצלנו כתוב הזה ונאמר על קבלת השכינה בים סוף "זה אלי ואנוהו" (טו, ב)\Makor{סנהדרין מב.}}\Peirush{שלשה דברים היו קשים למשה להבינם בדיוק, עד שהראה לו הקדוש ברוך הוא באצבעו. מעשה המנורה, מולד הירח שנאמר החדש הזה, וזיהוי השרצים האסורים\Makor{מנחות כט.}}\Peirush{חזקיהו עיבר את השנה עם ניסן במקום אדר, אך קשה, הרי כתוב החדש הזה, זה ניסן ואין ניסן אחר\Makor{ברכות י:, סנהדרין יב:}}\Peirush{עם ישראל הגיעו להר סיני בראש חדש סיון, נאמר אצלנו החדש הזה ונאמר "ביום הזה באו מדבר סיני" (יט, א)\Makor{שבת פו:}}\Peirush{ר' אלעזר בן ערך הגיע לפריגיה ודיומסת [שלא למדו תורה], ושכח תלמודו. כשחזר למקומו, בא לקרוא החדש הזה לכם ושכח תלמודו עד כדי כך, שקרא "החרש היה לבם". התפללו עליו חכמים ונזכר\Makor{שבת קמז:}}\Peirush{דורשים בהלכות החג שבועיים לפני החג, שמשה מזהיר על הפסח בראש חדש, נאמר ראש חדשים ונאמר "ויקחו להם איש שה לבית אבת" (ג)\Makor{פסחים ו:}}\Peirush{ניסן זהו החדש הראשון, נאמר בפסח ראש חדשים ראשון הוא לכם לחדשי השנה... בעשור לחדש הזה ויקחו להם, וגם "שמור את חדש האביב" (דברים טז, א) באיזה חדש יש אביב? ניסן\Makor{ראש השנה ז.}}\Peirush{אסור לעבר החדש לצרכי ציבור, שנאמר החדש הזה, כזה ראה וקדש\Makor{ראש השנה כ.}}\Peirush{ראה פס' א'\Makor{ראש השנה כב.}}\Peirush{ראה פס' א'\Makor{ראש השנה כה:}}\Peirush{בשבת הרביעית בחדש אדר קוראים את החדש הזה לכם\Makor{מגילה כט., ל.}}\Peirush{מתחילים לספור שמיטות מתשרי, מגזרה שווה "שנה" (ויקרא כה, ד) "שנה" (דברים יא, יב), אך קשה, למה לא ללמוד מניסן שנאמר בו שנה\Makor{ראש השנה ח:, ט.}}\Peirush{מתחילים לתרום מחצית השקל למקדש בניסן מגזרה שווה "שנה" (במדבר כח, יד) שנה\Makor{ראש השנה ז.}}%
% ---------- פרק 12 | פסוק 3 ----------
 \textbf{ג} דַּבְּר֗וּ אֶֽל־כׇּל־עֲדַ֤ת יִשְׂרָאֵל֙ לֵאמֹ֔ר בֶּעָשֹׂ֖ר לַחֹ֣דֶשׁ הַזֶּ֑ה וְיִקְח֣וּ לָהֶ֗ם אִ֛ישׁ שֶׂ֥ה לְבֵית־אָבֹ֖ת שֶׂ֥ה לַבָּֽיִת׃ \Rashi{\textbf{דברו אל כל עדת.}וכי אהרן מדבר? והלא כבר נאמר אתה תדבר? אלא חולקין כבוד זה לזה, ואומרים זה לזה למדני, והדבור יוצא מבין שניהם כאלו שניהם מדברים (מכילתא): דברו אל כל עדת ישראל לאמר בעשר לחדש. דברו היום, בראש חדש שיקחוהו בעשר לחדש: הזה. פסח מצרים מקחו בעשור, ולא פסח דורות (פסחים צ"ו): שה לבית אבת. למשפחה אחת; הרי שהיו מרבין, יכול שה אחד לכלן? ת"ל "שה לבית" (מכילתא):}\Peirush{בניגוד לפסח מצרים, בפסח לדורות לא צריך לקחת את קרבן הפסח בי' ניסן שנאמר בעשור לחודש הזה ויקחו, הזה ולא אחר\Makor{פסחים צו.}}\Peirush{ראה פס' ב'\Makor{פסחים ו:}}\Peirush{ראה פס' ב'\Makor{ראש השנה ז.}}\Peirush{נאמר ויקחו להם איש שה לבית אבת מכאן שאיש יכול לבצע קנין, ולא קטן\Makor{פסחים צא:}}\Peirush{נאמר לבית אבת, מכאן שניתן למנות שליח לקדשים\Makor{קידושים מב.}}\Peirush{כמו שבקרבן פסח רק שה כשר, כך גם בפטר חמור\Makor{בכורות יב.}}%
% ---------- פרק 12 | פסוק 4 ----------
 \textbf{ד} וְאִם־יִמְעַ֣ט הַבַּ֘יִת֮ מִהְי֣וֹת מִשֶּׂה֒ וְלָקַ֣ח ה֗וּא וּשְׁכֵנ֛וֹ הַקָּרֹ֥ב אֶל־בֵּית֖וֹ בְּמִכְסַ֣ת נְפָשֹׁ֑ת אִ֚ישׁ לְפִ֣י אׇכְל֔וֹ תָּכֹ֖סּוּ עַל־הַשֶּֽׂה׃ \Rashi{\textbf{ואם ימעט הבית מהיות משה.}ואם יהיו מועטים מהיות משה אחד, שאין יכולין לאכלו ויבא לידי נותר: ולקח הוא ושכנו וגו'. זהו משמעו לפי פשוטו. ועוד יש בו מדרש: ללמד, שאחר שנמנו עליו, יכולין להתמעט ולמשך ידיהם הימנו ולהמנות על שה אחר; אך אם באו למשך ידיהם ולהתמעט, מהיות משה, יתמעטו בעוד השה קים, בהיותו בחיים ולא משנשחט (פסחים פ"ט): במכסת. חשבון, וכן "מכסת הערכך" (ויקרא כ"ז): לפי אכלו. הראוי לאכילה, פרט לחולה ולזקן שאינו יכול לאכל כזית (מכילתא): תכסו. "תתמנון":}%
% ---------- פרק 12 | פסוק 5 ----------
 \textbf{ה} שֶׂ֥ה תָמִ֛ים זָכָ֥ר בֶּן־שָׁנָ֖ה יִהְיֶ֣ה לָכֶ֑ם מִן־הַכְּבָשִׂ֥ים וּמִן־הָעִזִּ֖ים תִּקָּֽחוּ׃ \Rashi{\textbf{תמים.}בלא מום: בן שנה. כל שנתו קרוי בן שנה, כלומר שנולד בשנה זו: מן הכבשים ומן העזים. או מזה או מזה, שאף עז קרוי שה, שנאמר "ושה עזים" (דברים י"ד):}%
% ---------- פרק 12 | פסוק 6 ----------
 \textbf{ו} וְהָיָ֤ה לָכֶם֙ לְמִשְׁמֶ֔רֶת עַ֣ד אַרְבָּעָ֥ה עָשָׂ֛ר י֖וֹם לַחֹ֣דֶשׁ הַזֶּ֑ה וְשָׁחֲט֣וּ אֹת֗וֹ כֹּ֛ל קְהַ֥ל עֲדַֽת־יִשְׂרָאֵ֖ל בֵּ֥ין הָעַרְבָּֽיִם׃ \Rashi{\textbf{והיה לכם למשמרת.}זהו לשון בקור, שטעון בקור ממום ארבעה ימים קדם שחיטה (פסחים צ"ו). ומפני מה הקדים לקיחתו לשחיטתו ארבעה ימים, מה שלא צוה כן בפסח דורות? היה ר' מתיא בן חרש אומר, הרי הוא אומר "ואעבר עליך ואראך והנה עתך עת דודים" (יחזקאל ט"ז) – הגיעה שבועה שנשבעתי לאברהם שאגאל את בניו, ולא היו בידם מצוות להתעסק בהם כדי שיגאלו, שנאמר "ואת ערום ועריה" (שם), ונתן להם שתי מצוות, דם פסח ודם מילה, שמלו באותו הלילה, שנאמר "מתבוססת בדמיך" (שם) – בשני דמים, ואומר "גם את בדם בריתך שלחתי אסיריך מבור אין מים בו" (זכריה ט'); ולפי שהיו שטופין באלילים אמר להם משכו וקחו לכם, משכו ידיכם מאלילים, וקחו לכם צאן של מצוה (מכילתא): ושחטו אתו וגו'. וכי כלן שוחטין? אלא מכאן ששלוחו של אדם כמותו (קידושין מ"א): קהל עדת ישראל. קהל ועדה וישראל; מכאן אמרו, פסחי צבור נשחטין בשלוש כתות זו אחר זו, נכנסה כת ראשונה ננעלו דלתות העזרה וכו'. כדאיתא בפסחים (דף ס"ד): בין הערבים. משש שעות ולמעלה קרוי בין הערבים, שהשמש נוטה לבית מבואו לערב; ולשון בין הערבים נראה בעיני אותן שעות שבין עריבת היום לעריבת הלילה, עריבת היום בתחלת ז' שעות מכי ינטו צללי ערב, ועריבת הלילה בתחלת הלילה. ערב לשון נשף וחשך, כמו "ערבה כל שמחה" (ישעיהו כ"ד):}%
% ---------- פרק 12 | פסוק 7 ----------
 \textbf{ז} וְלָֽקְחוּ֙ מִן־הַדָּ֔ם וְנָ֥תְנ֛וּ עַל־שְׁתֵּ֥י הַמְּזוּזֹ֖ת וְעַל־הַמַּשְׁק֑וֹף עַ֚ל הַבָּ֣תִּ֔ים אֲשֶׁר־יֹאכְל֥וּ אֹת֖וֹ בָּהֶֽם׃ \Rashi{\textbf{ולקחו מן הדם.}זו קבלת הדם, יכול ביד? ת"ל אשר בסף (מכילתא): המזוזת. הם הזקופות אחת מכאן לפתח ואחת מכאן: המשקוף. העליון שהדלת שוקף עליו כשסוגרין אותו, לינטי"ל בלעז; ולשון שקיפה חבטה, כמו "קול עלה נדף" (ויקרא כ"ו) – דשקיף, חבורה – משקופי: על הבתים אשר יאכלו אתו בהם. ולא על משקוף ומזוזות שבבית התבן ובית הבקר, שאין דרין בתוכו (ע' מכילתא):}%
% ---------- פרק 12 | פסוק 8 ----------
 \textbf{ח} וְאָכְל֥וּ אֶת־הַבָּשָׂ֖ר בַּלַּ֣יְלָה הַזֶּ֑ה צְלִי־אֵ֣שׁ וּמַצּ֔וֹת עַל־מְרֹרִ֖ים יֹאכְלֻֽהוּ׃ \Rashi{\textbf{את הבשר.}ולא גידים ועצמות (שם): ומצות על מררים. כל עשב מר נקרא מרור; וצום לאכל מר זכר ל"וימררו את חייהם" (שמות א'):}%
% ---------- פרק 12 | פסוק 9 ----------
 \textbf{ט} אַל־תֹּאכְל֤וּ מִמֶּ֙נּוּ֙ נָ֔א וּבָשֵׁ֥ל מְבֻשָּׁ֖ל בַּמָּ֑יִם כִּ֣י אִם־צְלִי־אֵ֔שׁ רֹאשׁ֥וֹ עַל־כְּרָעָ֖יו וְעַל־קִרְבּֽוֹ׃ \Rashi{\textbf{אל תאכלו ממנו נא.}שאינו צלוי כל צרכו קוראו נא בלשון ערבי: ובשל מבשל. כל זה באזהרת אל תאכלו: במים. מנין לשאר משקין? ת"ל ובשל מבשל, מכל מקום (פסחים מ"א): כי אם צלי אש. למעלה גזר עליו במצות עשה, וכאן הוסיף עליו לא תעשה, אל תאכלו ממנו … כי אם צלי אש: ראשו על כרעיו. צולהו כלו כאחד עם ראשו ועם כרעיו ועם קרבו, ובני מעיו נותן לתוכו אחר הדחתן (שם ע"ד); ולשון על קרבו כלשון "על צבאותם" (שמות ו'), כמו בצבאותם – כמות שהן, אף זה כמות שהוא – כל בשרו משלם:}%
% ---------- פרק 12 | פסוק 10 ----------
 \textbf{י} וְלֹא־תוֹתִ֥ירוּ מִמֶּ֖נּוּ עַד־בֹּ֑קֶר וְהַנֹּתָ֥ר מִמֶּ֛נּוּ עַד־בֹּ֖קֶר בָּאֵ֥שׁ תִּשְׂרֹֽפוּ׃ \Rashi{\textbf{והנתר ממנו עד בקר.}מה ת"ל עד בקר פעם שניה? לתן בקר על בקר, שהבקר משמעו משעת הנץ החמה, ובא הכתוב להקדים שאסור באכילה מעלות השחר, זהו לפי משמעו. ועוד מדרש אחר, למד שאינו נשרף ביום טוב אלא ממחרת, וכך תדרשנו: והנותר ממנו בבקר ראשון, עד בקר שני תעמד ותשרפנו (מכילתא):}%
% ---------- פרק 12 | פסוק 11 ----------
 \textbf{יא} וְכָ֘כָה֮ תֹּאכְל֣וּ אֹתוֹ֒ מׇתְנֵיכֶ֣ם חֲגֻרִ֔ים נַֽעֲלֵיכֶם֙ בְּרַגְלֵיכֶ֔ם וּמַקֶּלְכֶ֖ם בְּיֶדְכֶ֑ם וַאֲכַלְתֶּ֤ם אֹתוֹ֙ בְּחִפָּז֔וֹן פֶּ֥סַח ה֖וּא לַיהֹוָֽה׃ \Rashi{\textbf{מתניכם חגרים.}מזמנים לדרך: בחפזון. לשון בהלה ומהירות, כמו "ויהי דוד נחפז ללכת" (שמואל א' כ"ג), "אשר השליכו ארם בחפזם" (מלכים ב' ז'): פסח הוא לה'. הקרבן קרוי פסח על שם הדלוג והפסיחה, שהקב"ה היה מדלג בתי ישראל מבין בתי מצרים וקופץ ממצרי למצרי וישראל אמצעי נמלט; ואתם עשו כל עבודותיו לשם שמים, דרך דלוג וקפיצה, זכר לשמו שקרוי פסח; וגם פשק"א לשון פסיעה:}%
% ---------- פרק 12 | פסוק 12 ----------
 \textbf{יב} וְעָבַרְתִּ֣י בְאֶֽרֶץ־מִצְרַ֘יִם֮ בַּלַּ֣יְלָה הַזֶּה֒ וְהִכֵּיתִ֤י כׇל־בְּכוֹר֙ בְּאֶ֣רֶץ מִצְרַ֔יִם מֵאָדָ֖ם וְעַד־בְּהֵמָ֑ה וּבְכׇל־אֱלֹהֵ֥י מִצְרַ֛יִם אֶֽעֱשֶׂ֥ה שְׁפָטִ֖ים אֲנִ֥י יְהֹוָֽה׃ \Rashi{\textbf{ועברתי.}כמלך העובר ממקום למקום (מכילתא), ובהעברה אחת וברגע אחד כלן לוקים: כל בכור בארץ מצרים. אף בכורות אחרים והם במצרים. ומנין אף בכורי מצרים שבמקומות אחרים? ת"ל "למכה מצרים בבכוריהם" (תהלים קל"ו): מאדם ועד בהמה. מי שהתחיל בעברה ממנו מתחלת הפרענות (מכילתא): ובכל אלהי מצרים. של עץ נרקבת ושל מתכת נמסת ונתכת לארץ (שם): אעשה שפטים אני ה'. אני בעצמי, ולא על ידי שליח:}%
% ---------- פרק 12 | פסוק 13 ----------
 \textbf{יג} וְהָיָה֩ הַדָּ֨ם לָכֶ֜ם לְאֹ֗ת עַ֤ל הַבָּתִּים֙ אֲשֶׁ֣ר אַתֶּ֣ם שָׁ֔ם וְרָאִ֙יתִי֙ אֶת־הַדָּ֔ם וּפָסַחְתִּ֖י עֲלֵכֶ֑ם וְלֹֽא־יִֽהְיֶ֨ה בָכֶ֥ם נֶ֙גֶף֙ לְמַשְׁחִ֔ית בְּהַכֹּתִ֖י בְּאֶ֥רֶץ מִצְרָֽיִם׃ \Rashi{\textbf{והיה הדם לכם לאת.}לכם לאות ולא לאחרים לאות (שם). מכאן שלא נתנו הדם אלא מבפנים: וראיתי את הדם. הכל גלוי לפניו, אלא אמר הקב"ה, נותן אני את עיני לראות שאתם עסוקים במצוותי, ופוסח אני עליכם (שם): ופסחתי. וחמלתי, ודומה לו "פסח והמליט" (ישעיהו ל"א). ואני אומר, כל פסיחה לשון דלוג וקפיצה. ופסחתי, מדלג היה מבתי ישראל לבתי מצריים, שהיו שרויים זה בתוך זה, וכן "פוסחים על שתי הסעפים" (מלכים א י"ח), וכן כל הפסחים – הולכים כקופצים, וכן "פסח והמליט" – מדלגו וממלטו מבין המומתים: ולא יהיה בכם נגף. אבל הוה הוא במצרים. הרי שהיה מצרי בביתו של ישראל, יכול ימלט? ת"ל "ולא יהיה בכם נגף", אבל הוה במצרים שבבתיכם. הרי שהיה ישראל בביתו של מצרי, שומע אני ילקה כמותו, ת"ל "ולא יהיה בכם נגף" (מכילתא):}%
% ---------- פרק 12 | פסוק 14 ----------
 \textbf{יד} וְהָיָה֩ הַיּ֨וֹם הַזֶּ֤ה לָכֶם֙ לְזִכָּר֔וֹן וְחַגֹּתֶ֥ם אֹת֖וֹ חַ֣ג לַֽיהֹוָ֑ה לְדֹרֹ֣תֵיכֶ֔ם חֻקַּ֥ת עוֹלָ֖ם תְּחׇגֻּֽהוּ׃ \Rashi{\textbf{לזכרון.}לדורות: וחגתם אתו. יום שהוא לך לזכרון אתה חוגגו. ועדין לא שמענו איזהו יום הזכרון, ת"ל "זכור את היום הזה אשר יצאתם" (שמות י"ג), למדנו, שיום היציאה הוא יום של זכרון. ואי זה יום יצאו? ת"ל "ממחרת הפסח יצאו" (במדבר ל"ג), הוי אומר יום חמשה עשר בניסן הוא של יום טוב, שהרי ליל חמשה עשר אכלו את הפסח ולבקר יצאו (מכילתא): לדרתיכם וגו'. שומע אני מעוט דורות שנים, ת"ל "חקת עולם תחגהו" (שם):}%
% ---------- פרק 12 | פסוק 15 ----------
 \textbf{טו} שִׁבְעַ֤ת יָמִים֙ מַצּ֣וֹת תֹּאכֵ֔לוּ אַ֚ךְ בַּיּ֣וֹם הָרִאשׁ֔וֹן תַּשְׁבִּ֥יתוּ שְּׂאֹ֖ר מִבָּתֵּיכֶ֑ם כִּ֣י ׀ כׇּל־אֹכֵ֣ל חָמֵ֗ץ וְנִכְרְתָ֞ה הַנֶּ֤פֶשׁ הַהִוא֙ מִיִּשְׂרָאֵ֔ל מִיּ֥וֹם הָרִאשֹׁ֖ן עַד־י֥וֹם הַשְּׁבִעִֽי׃ \Rashi{\textbf{שבעת ימים.}שטיינ"א של ימים: שבעת ימים מצות תאכלו. ובמקום אחר הוא אומר "ששת ימים תאכל מצות" (דברים ט"ז) למד על שביעי של פסח שאינו חובה לאכל מצה, ובלבד שלא יאכל חמץ; מנין אף ששה רשות? ת"ל "ששת ימים". זו מדה בתורה, דבר שהיה בכלל ויצא מן הכלל ללמד, לא ללמד על עצמו בלבד יצא, אלא ללמד על הכלל כלו יצא, מה שביעי רשות אף ששה רשות; יכול אף הלילה הראשון רשות, ת"ל "בערב תאכלו מצת" – הכתוב קבעו חובה (פסחים ק"כ): אך ביום הראשון תשביתו שאור. מערב יום טוב, וקרוי ראשון לפי שהוא לפני השבעה; ומצינו מקדם קרוי ראשון, "הראישון אדם תולד" (איוב ט"ו) – הלפני אדם נולדת; או אינו אלא ראשון של שבעה? ת"ל "לא תשחט על חמץ וגו'" (שמות ל"ד) – לא תשחט הפסח ועדין חמץ קים (פסחים ה'): הנפש ההוא. כשהיא בנפשה ובדעתה – פרט לאנוס (מכילתא): מישראל. שומע אני תכרת מישראל ותלך לה לעם אחר, ת"ל במקום אחר "מלפני" (ויקרא כ"ב) – בכל מקום שהוא רשותי:}%
% ---------- פרק 12 | פסוק 16 ----------
 \textbf{טז} וּבַיּ֤וֹם הָרִאשׁוֹן֙ מִקְרָא־קֹ֔דֶשׁ וּבַיּוֹם֙ הַשְּׁבִיעִ֔י מִקְרָא־קֹ֖דֶשׁ יִהְיֶ֣ה לָכֶ֑ם כׇּל־מְלָאכָה֙ לֹא־יֵעָשֶׂ֣ה בָהֶ֔ם אַ֚ךְ אֲשֶׁ֣ר יֵאָכֵ֣ל לְכׇל־נֶ֔פֶשׁ ה֥וּא לְבַדּ֖וֹ יֵעָשֶׂ֥ה לָכֶֽם׃ \Rashi{\textbf{מקרא קדש.}מקרא שם דבר; קרא אותו קדש לאכילה ושתיה וכסות (מכילתא): לא יעשה בהם. אפלו על ידי אחרים (שם): הוא לבדו. "הוא" ולא מכשיריו שאפשר לעשותן מערב יום טוב (ביצה כ"ח): לכל נפש. אפלו לבהמה; יכול אף לגוים, תלמוד לומר לכם (מכילתא):}%
% ---------- פרק 12 | פסוק 17 ----------
 \textbf{יז} וּשְׁמַרְתֶּם֮ אֶת־הַמַּצּוֹת֒ כִּ֗י בְּעֶ֙צֶם֙ הַיּ֣וֹם הַזֶּ֔ה הוֹצֵ֥אתִי אֶת־צִבְאוֹתֵיכֶ֖ם מֵאֶ֣רֶץ מִצְרָ֑יִם וּשְׁמַרְתֶּ֞ם אֶת־הַיּ֥וֹם הַזֶּ֛ה לְדֹרֹתֵיכֶ֖ם חֻקַּ֥ת עוֹלָֽם׃ \Rashi{\textbf{ושמרתם את המצות.}שלא יבאו לידי חמוץ; מכאן אמרו תפח תלטש בצונן, רבי יאשיה אומר אל תהי קורא את המצות, אלא את המצוות – כדרך שאין מחמיצין את המצה, כך אין מחמיצין את המצוה, אלא אם באה לידך, עשה אותה מיד (שם): ושמרתם את היום הזה. ממלאכה: לדרתיכם חקת עולם. לפי שלא נאמר דורות וחקת עולם על המלאכה אלא על החגיגה, לכך חזר ושנאו כאן, שלא תאמר אזהרת כל מלאכה לא יעשה לא לדורות נאמרה אלא לאותו הדור:}%
% ---------- פרק 12 | פסוק 18 ----------
 \textbf{יח} בָּרִאשֹׁ֡ן בְּאַרְבָּעָה֩ עָשָׂ֨ר י֤וֹם לַחֹ֙דֶשׁ֙ בָּעֶ֔רֶב תֹּאכְל֖וּ מַצֹּ֑ת עַ֠ד י֣וֹם הָאֶחָ֧ד וְעֶשְׂרִ֛ים לַחֹ֖דֶשׁ בָּעָֽרֶב׃ \Rashi{\textbf{עד יום האחד ועשרים.}למה נאמר? והלא כבר נאמר "שבעת ימים"? לפי שנאמר ימים, לילות מנין? ת"ל עד יום האחד ועשרים וגו' (מכילתא):}%
% ---------- פרק 12 | פסוק 19 ----------
 \textbf{יט} שִׁבְעַ֣ת יָמִ֔ים שְׂאֹ֕ר לֹ֥א יִמָּצֵ֖א בְּבָתֵּיכֶ֑ם כִּ֣י ׀ כׇּל־אֹכֵ֣ל מַחְמֶ֗צֶת וְנִכְרְתָ֞ה הַנֶּ֤פֶשׁ הַהִוא֙ מֵעֲדַ֣ת יִשְׂרָאֵ֔ל בַּגֵּ֖ר וּבְאֶזְרַ֥ח הָאָֽרֶץ׃ \Rashi{\textbf{לא ימצא בבתיכם.}מנין לגבולין? תלמוד לומר "בכל גבולך" (שמות י"ג); מה ת"ל בבתיכם? מה ביתך ברשותך, אף גבולך שברשותך, יצא חמצו של נכרי שהוא אצל ישראל ולא קבל עליו אחריות (מכילתא): כי כל אכל מחמצת. לענש כרת על השאור, והלא כבר ענש על החמץ? אלא שלא תאמר חמץ שראוי לאכילה ענש עליו, שאור שאינו ראוי לאכילה לא יענש עליו; ואם ענש על השאור ולא ענש על החמץ, הייתי אומר, שאור שהוא מחמץ אחרים ענש עליו, חמץ שאינו מחמץ אחרים לא יענש עליו, לכך נאמרו שניהם (מכילתא): בגר ובאזרח הארץ. לפי שהנס נעשה לישראל, הצרך לרבות את הגרים (שם):}%
% ---------- פרק 12 | פסוק 20 ----------
 \textbf{כ} כׇּל־מַחְמֶ֖צֶת לֹ֣א תֹאכֵ֑לוּ בְּכֹל֙ מוֹשְׁבֹ֣תֵיכֶ֔ם תֹּאכְל֖וּ מַצּֽוֹת׃ \{פ\} \Rashi{\textbf{מחמצת לא תאכלו.}אזהרה על אכילת שאור: כל מחמצת. להביא את תערבתו: בכל מושבתיכם תאכלו מצות. זה בא ללמד שתהא ראויה להאכל בכל מושבתיכם, פרט למעשר שני וחלות תודה (מכילתא):}%
% ---------- פרק 12 | פסוק 21 ----------
 \textbf{כא} וַיִּקְרָ֥א מֹשֶׁ֛ה לְכׇל־זִקְנֵ֥י יִשְׂרָאֵ֖ל וַיֹּ֣אמֶר אֲלֵהֶ֑ם מִֽשְׁכ֗וּ וּקְח֨וּ לָכֶ֥ם צֹ֛אן לְמִשְׁפְּחֹתֵיכֶ֖ם וְשַׁחֲט֥וּ הַפָּֽסַח׃ \Rashi{\textbf{משכו.}מי שיש לו צאן ימשך משלו: וקחו. מי שאין לו יקח מן השוק (שם): למשפחתיכם. שה לבית אבות:}%
% ---------- פרק 12 | פסוק 22 ----------
 \textbf{כב} וּלְקַחְתֶּ֞ם אֲגֻדַּ֣ת אֵז֗וֹב וּטְבַלְתֶּם֮ בַּדָּ֣ם אֲשֶׁר־בַּסַּף֒ וְהִגַּעְתֶּ֤ם אֶל־הַמַּשְׁקוֹף֙ וְאֶל־שְׁתֵּ֣י הַמְּזוּזֹ֔ת מִן־הַדָּ֖ם אֲשֶׁ֣ר בַּסָּ֑ף וְאַתֶּ֗ם לֹ֥א תֵצְא֛וּ אִ֥ישׁ מִפֶּֽתַח־בֵּית֖וֹ עַד־בֹּֽקֶר׃ \Rashi{\textbf{אזוב.}מין ירק שיש לו גבעולין: אגדת אזוב. שלושה קלחין קרויין אגדה (שבת ק"ט): אשר בסף. בכלי, כמו ספות כסף: מן הדם אשר בסף. למה חזר ושנאו? שלא תאמר טבילה אחת לשלוש המתנות, לכך נאמר עוד אשר בסף, ושתהא כל נתינה ונתינה מן הדם אשר בסף, על כל הגעה טבילה (מכילתא): ואתם לא תצאו וגו'. מגיד, שמאחר שנתנה רשות למשחית לחבל, אינו מבחין בין צדיק לרשע (שם); ולילה רשות למחבלים הוא, שנאמר "בו תרמש כל חיתו יער" (תהלים ק"ד):}%
% ---------- פרק 12 | פסוק 23 ----------
 \textbf{כג} וְעָבַ֣ר יְהֹוָה֮ לִנְגֹּ֣ף אֶת־מִצְרַ֒יִם֒ וְרָאָ֤ה אֶת־הַדָּם֙ עַל־הַמַּשְׁק֔וֹף וְעַ֖ל שְׁתֵּ֣י הַמְּזוּזֹ֑ת וּפָסַ֤ח יְהֹוָה֙ עַל־הַפֶּ֔תַח וְלֹ֤א יִתֵּן֙ הַמַּשְׁחִ֔ית לָבֹ֥א אֶל־בָּתֵּיכֶ֖ם לִנְגֹּֽף׃ \Rashi{\textbf{ופסח.}וחמל, ויש לומר ודלג: ולא יתן המשחית. ולא יתן לו יכלת לבא, כמו "ולא נתנו אלהים להרע עמדי" (בראשית ל"א):}%
% ---------- פרק 12 | פסוק 24 ----------
 \textbf{כד} וּשְׁמַרְתֶּ֖ם אֶת־הַדָּבָ֣ר הַזֶּ֑ה לְחׇק־לְךָ֥ וּלְבָנֶ֖יךָ עַד־עוֹלָֽם׃ %
% ---------- פרק 12 | פסוק 25 ----------
 \textbf{כה} וְהָיָ֞ה כִּֽי־תָבֹ֣אוּ אֶל־הָאָ֗רֶץ אֲשֶׁ֨ר יִתֵּ֧ן יְהֹוָ֛ה לָכֶ֖ם כַּאֲשֶׁ֣ר דִּבֵּ֑ר וּשְׁמַרְתֶּ֖ם אֶת־הָעֲבֹדָ֥ה הַזֹּֽאת׃ \Rashi{\textbf{והיה כי תבאו אל הארץ.}תלה הכתוב עבודה זו בביאתם לארץ, ולא נתחיבו במדבר אלא פסח אחד שעשו בשנה השנית על פי הדבור: כאשר דבר. והיכן דבר? "והבאתי אתכם אל הארץ וגו'" (שמות ו'):}%
% ---------- פרק 12 | פסוק 26 ----------
 \textbf{כו} וְהָיָ֕ה כִּֽי־יֹאמְר֥וּ אֲלֵיכֶ֖ם בְּנֵיכֶ֑ם מָ֛ה הָעֲבֹדָ֥ה הַזֹּ֖את לָכֶֽם׃ %
% ---------- פרק 12 | פסוק 27 ----------
 \textbf{כז} וַאֲמַרְתֶּ֡ם זֶֽבַח־פֶּ֨סַח ה֜וּא לַֽיהֹוָ֗ה אֲשֶׁ֣ר פָּ֠סַ֠ח עַל־בָּתֵּ֤י בְנֵֽי־יִשְׂרָאֵל֙ בְּמִצְרַ֔יִם בְּנׇגְפּ֥וֹ אֶת־מִצְרַ֖יִם וְאֶת־בָּתֵּ֣ינוּ הִצִּ֑יל וַיִּקֹּ֥ד הָעָ֖ם וַיִּֽשְׁתַּחֲוֽוּ׃ \Rashi{\textbf{ויקד העם.}על בשורת הגאלה וביאת הארץ ובשורת הבנים שיהיו להם:}%
% ---------- פרק 12 | פסוק 28 ----------
 \textbf{כח} וַיֵּלְכ֥וּ וַֽיַּעֲשׂ֖וּ בְּנֵ֣י יִשְׂרָאֵ֑ל כַּאֲשֶׁ֨ר צִוָּ֧ה יְהֹוָ֛ה אֶת־מֹשֶׁ֥ה וְאַהֲרֹ֖ן כֵּ֥ן עָשֽׂוּ׃ \{ס\} \Rashi{\textbf{וילכו ויעשו בני ישראל.}וכי כבר עשו, והלא מראש חדש נאמר להם? אלא מכיון שקבלו עליהם מעלה עליהם הכתוב כאלו עשו (מכילתא): וילכו ויעשו. אף ההליכה מנה הכתוב, לתן שכר להליכה ושכר לעשיה (שם): כאשר צוה ה' את משה ואהרן. להגיד שבחן של ישראל שלא הפילו דבר מכל מצוות משה ואהרן, ומהו כן עשו? אף משה ואהרן כן עשו (שם):}%
% ---------- פרק 12 | פסוק 29 ----------
 \textbf{כט} וַיְהִ֣י ׀ בַּחֲצִ֣י הַלַּ֗יְלָה וַֽיהֹוָה֮ הִכָּ֣ה כׇל־בְּכוֹר֮ בְּאֶ֣רֶץ מִצְרַ֒יִם֒ מִבְּכֹ֤ר פַּרְעֹה֙ הַיֹּשֵׁ֣ב עַל־כִּסְא֔וֹ עַ֚ד בְּכ֣וֹר הַשְּׁבִ֔י אֲשֶׁ֖ר בְּבֵ֣ית הַבּ֑וֹר וְכֹ֖ל בְּכ֥וֹר בְּהֵמָֽה׃ \Rashi{\textbf{וה'.}כל מקום שנאמר "וה'" הוא ובית דינו, שהוי"ו לשון תוספת הוא, כמו פלוני ופלוני: הכה כל בכור. אף של אמה אחרת והוא במצרים (מכילתא): מבכר פרעה. אף פרעה בכור היה, ונשתיר מן הבכורות, ועליו הוא אומר "בעבור הראותך את כחי" (שמות ט') – בים סוף (מכילתא): עד בכור השבי. שהיו שמחין לאידם של ישראל; ועוד, שלא יאמרו יראתנו הביאה הפרענות; ובכור השפחה בכלל היה, שהרי מנה מן החשוב שבכלן עד הפחות, ובכור השפחה חשוב מבכור השבי:}%
% ---------- פרק 12 | פסוק 30 ----------
 \textbf{ל} וַיָּ֨קׇם פַּרְעֹ֜ה לַ֗יְלָה ה֤וּא וְכׇל־עֲבָדָיו֙ וְכׇל־מִצְרַ֔יִם וַתְּהִ֛י צְעָקָ֥ה גְדֹלָ֖ה בְּמִצְרָ֑יִם כִּֽי־אֵ֣ין בַּ֔יִת אֲשֶׁ֥ר אֵֽין־שָׁ֖ם מֵֽת׃ \Rashi{\textbf{ויקם פרעה.}ממטתו: לילה. ולא כדרך המלכים בשלוש שעות ביום (שם): הוא. תחלה, ואחר כך עבדיו, מלמד שהיה הוא מחזר על בתי עבדיו ומעמידן (שם): כי אין בית אשר אין שם מת. יש שם בכור, מת, אין שם בכור, גדול שבבית קרוי בכור, שנאמר "אף אני בכור אתנהו" (תהילים פ"ט). דבר אחר, מצריות מזנות תחת בעליהן ויולדות מרוקים פנויים, והיו להם בכורות הרבה, פעמים חמשה לאשה אחת, כל אחד בכור לאביו:}%
% ---------- פרק 12 | פסוק 31 ----------
 \textbf{לא} וַיִּקְרָא֩ לְמֹשֶׁ֨ה וּֽלְאַהֲרֹ֜ן לַ֗יְלָה וַיֹּ֙אמֶר֙ ק֤וּמוּ צְּאוּ֙ מִתּ֣וֹךְ עַמִּ֔י גַּם־אַתֶּ֖ם גַּם־בְּנֵ֣י יִשְׂרָאֵ֑ל וּלְכ֛וּ עִבְד֥וּ אֶת־יְהֹוָ֖ה כְּדַבֶּרְכֶֽם׃ \Rashi{\textbf{ויקרא למשה ולאהרן לילה.}מגיד שהיה מחזר על פתחי העיר וצועק, היכן משה שרוי? היכן אהרן שרוי? (מכילתא): גם אתם. הגברים: גם בני ישראל. הטף: ולכו עבדו את ה' כדברכם. הכל כמו שאמרתם ולא כשאמרתי אני; בטל "לא אשלח", בטל "מי ומי ההולכים", בטל "רק צאנכם ובקרכם יצג":}%
% ---------- פרק 12 | פסוק 32 ----------
 \textbf{לב} גַּם־צֹאנְכֶ֨ם גַּם־בְּקַרְכֶ֥ם קְח֛וּ כַּאֲשֶׁ֥ר דִּבַּרְתֶּ֖ם וָלֵ֑כוּ וּבֵֽרַכְתֶּ֖ם גַּם־אֹתִֽי׃ \Rashi{\textbf{גם צאנכם גם בקרכם קחו.}ומהו כאשר דברתם? "גם אתה תתן בידנו זבחים ועלת" (שמות י'), קחו כאשר דברתם: וברכתם גם אתי. התפללו עלי שלא אמות, שאני בכור:}%
% ---------- פרק 12 | פסוק 33 ----------
 \textbf{לג} וַתֶּחֱזַ֤ק מִצְרַ֙יִם֙ עַל־הָעָ֔ם לְמַהֵ֖ר לְשַׁלְּחָ֣ם מִן־הָאָ֑רֶץ כִּ֥י אָמְר֖וּ כֻּלָּ֥נוּ מֵתִֽים׃ \Rashi{\textbf{כלנו מתים.}אמרו לא כגזרת משה הוא, שהרי אמר "ומת כל בכור", וכאן אף הפשוטים מתים, חמשה או עשרה בבית אחד:}%
% ---------- פרק 12 | פסוק 34 ----------
 \textbf{לד} וַיִּשָּׂ֥א הָעָ֛ם אֶת־בְּצֵק֖וֹ טֶ֣רֶם יֶחְמָ֑ץ מִשְׁאֲרֹתָ֛ם צְרֻרֹ֥ת בְּשִׂמְלֹתָ֖ם עַל־שִׁכְמָֽם׃ \Rashi{\textbf{טרם יחמץ.}המצרים לא הניחום לשהות כדי חמוץ: משארתם. שירי מצה ומרור. על שכמם. אע"פ שבהמות הרבה הוליכו עמהם, מחבבים היו את המצוה:}%
% ---------- פרק 12 | פסוק 35 ----------
 \textbf{לה} וּבְנֵי־יִשְׂרָאֵ֥ל עָשׂ֖וּ כִּדְבַ֣ר מֹשֶׁ֑ה וַֽיִּשְׁאֲלוּ֙ מִמִּצְרַ֔יִם כְּלֵי־כֶ֛סֶף וּכְלֵ֥י זָהָ֖ב וּשְׂמָלֹֽת׃ \Rashi{\textbf{כדבר משה.}שאמר להם במצרים, "וישאלו איש מאת רעהו": ושמלת. אף הן היו חשובות להם מן הכסף ומן הזהב, והמאחר בפסוק חשוב:}%
% ---------- פרק 12 | פסוק 36 ----------
 \textbf{לו} וַֽיהֹוָ֞ה נָתַ֨ן אֶת־חֵ֥ן הָעָ֛ם בְּעֵינֵ֥י מִצְרַ֖יִם וַיַּשְׁאִל֑וּם וַֽיְנַצְּל֖וּ אֶת־מִצְרָֽיִם׃ \{פ\} \Rashi{\textbf{וישאלום.}אף מה שלא היו שואלים מהם היו נותנים להם, אתה אומר אחד? טל שנים ולך. וינצלו. ורוקינו:}%
% ---------- פרק 12 | פסוק 37 ----------
 \textbf{לז} וַיִּסְע֧וּ בְנֵֽי־יִשְׂרָאֵ֛ל מֵרַעְמְסֵ֖ס סֻכֹּ֑תָה כְּשֵׁשׁ־מֵא֨וֹת אֶ֧לֶף רַגְלִ֛י הַגְּבָרִ֖ים לְבַ֥ד מִטָּֽף׃ \Rashi{\textbf{מרעמסס סכתה.}מאה ועשרים מיל היו, ובאו שם לפי שעה, שנאמר "ואשא אתכם על כנפי נשרים" (שמות י"ט): הגברים. מבן עשרים שנה ומעלה (שה"ש רבה ג, ו):}%
% ---------- פרק 12 | פסוק 38 ----------
 \textbf{לח} וְגַם־עֵ֥רֶב רַ֖ב עָלָ֣ה אִתָּ֑ם וְצֹ֣אן וּבָקָ֔ר מִקְנֶ֖ה כָּבֵ֥ד מְאֹֽד׃ \Rashi{\textbf{ערב רב.}תערבת אמות של גרים:}%
% ---------- פרק 12 | פסוק 39 ----------
 \textbf{לט} וַיֹּאפ֨וּ אֶת־הַבָּצֵ֜ק אֲשֶׁ֨ר הוֹצִ֧יאוּ מִמִּצְרַ֛יִם עֻגֹ֥ת מַצּ֖וֹת כִּ֣י לֹ֣א חָמֵ֑ץ כִּֽי־גֹרְשׁ֣וּ מִמִּצְרַ֗יִם וְלֹ֤א יָֽכְלוּ֙ לְהִתְמַהְמֵ֔הַּ וְגַם־צֵדָ֖ה לֹא־עָשׂ֥וּ לָהֶֽם׃ \Rashi{\textbf{עגת מצות.}חררה של מצה; בצק שלא החמיץ קרוי מצה: וגם צדה לא עשו להם. לדרך. מגיד שבחן של ישראל, שלא אמרו היאך נצא למדבר בלא צידה? אלא האמינו והלכו; הוא שמפרש בקבלה, "זכרתי לך חסד נעוריך אהבת כלולתיך לכתך אחרי במדבר בארץ לא זרועה" (ירמיהו ב'), מה שכר מפרש אחריו? "קדש ישראל לה' וגו'":}%
% ---------- פרק 12 | פסוק 40 ----------
 \textbf{מ} וּמוֹשַׁב֙ בְּנֵ֣י יִשְׂרָאֵ֔ל אֲשֶׁ֥ר יָשְׁב֖וּ בְּמִצְרָ֑יִם שְׁלֹשִׁ֣ים שָׁנָ֔ה וְאַרְבַּ֥ע מֵא֖וֹת שָׁנָֽה׃ \Rashi{\textbf{אשר ישבו במצרים.}אחר שאר הישיבות שישבו גרים בארץ לא להם: שלשים שנה וארבע מאות שנה. בין הכל, משנולד יצחק עד עכשו, היו ארבע מאות, משהיה לו זרע לאברהם נתקים כי גר יהיה זרעך, ושלשים שנה היו משנגזרה גזרת בין הבתרים עד שנולד יצחק; ואי אפשר לומר בארץ מצרים לבדה, שהרי קהת מן הבאים עם יעקב היה, צא וחשב כל שנותיו וכל שנות עמרם בנו ושמונים של משה, לא תמצאם כל כך, ועל כרחך הרבה שנים היו לקהת עד שלא ירד למצרים, והרבה משנות עמרם נבלעים בשנות קהת, והרבה משמונים של משה נבלעים בשנות עמרם, הרי שלא תמצא ארבע מאות לביאת מצרים, והזקקתה לומר על כרחך, שאף שאר הישיבות נקראו גרות ואפלו בחברון שנאמר "אשר גר שם אברהם ויצחק" (בראשית ל"ה), ואומר "את ארץ מגריהם אשר גרו בה" (שמות ו'), לפיכך אתה צריך לומר "כי גר יהיה זרעך" משהיה לו זרע, וכשתמנה ת' שנה משנולד יצחק, תמצא מביאתן למצרים עד יציאתן ר"י, וזה אחד מן הדברים ששנו לתלמי המלך (מגילה ט'):}%
% ---------- פרק 12 | פסוק 41 ----------
 \textbf{מא} וַיְהִ֗י מִקֵּץ֙ שְׁלֹשִׁ֣ים שָׁנָ֔ה וְאַרְבַּ֥ע מֵא֖וֹת שָׁנָ֑ה וַיְהִ֗י בְּעֶ֙צֶם֙ הַיּ֣וֹם הַזֶּ֔ה יָ֥צְא֛וּ כׇּל־צִבְא֥וֹת יְהֹוָ֖ה מֵאֶ֥רֶץ מִצְרָֽיִם׃ \Rashi{\textbf{ויהי מקץ שלשים שנה וגו' ויהי בעצם היום הזה.}מגיד, שכיון שהגיע הקץ לא עכבן המקום כהרף עין – בט"ו בניסן באו מלאכי השרת אצל אברהם לבשרו, ובט"ו בניסן נולד יצחק, ובט"ו בניסן נגזרה גזרת בין הבתרים:}%
% ---------- פרק 12 | פסוק 42 ----------
 \textbf{מב} לֵ֣יל שִׁמֻּרִ֥ים הוּא֙ לַֽיהֹוָ֔ה לְהוֹצִיאָ֖ם מֵאֶ֣רֶץ מִצְרָ֑יִם הֽוּא־הַלַּ֤יְלָה הַזֶּה֙ לַֽיהֹוָ֔ה שִׁמֻּרִ֛ים לְכׇל־בְּנֵ֥י יִשְׂרָאֵ֖ל לְדֹרֹתָֽם׃ \{פ\} \Rashi{\textbf{ליל שמרים.}שהיה הקב"ה שומר ומצפה לו, לקים הבטחתו להוציאם מארץ מצרים: הוא הלילה הזה לה'. הוא הלילה שאמר לאברהם, בלילה הזה אני גואל את בניך (מכילתא): שמרים לכל בני ישראל לדרתם. משמר ובא מן המזיקין, כענין שנאמר "ולא יתן המשחית וגו'" (פסחים ק"ט):}%
% ---------- פרק 12 | פסוק 43 ----------
 \textbf{מג} וַיֹּ֤אמֶר יְהֹוָה֙ אֶל־מֹשֶׁ֣ה וְאַהֲרֹ֔ן זֹ֖את חֻקַּ֣ת הַפָּ֑סַח כׇּל־בֶּן־נֵכָ֖ר לֹא־יֹ֥אכַל בּֽוֹ׃ \Rashi{\textbf{זאת חקת הפסח.}בי"ד בניסן נאמרה להם פרשה זו: כל בן נכר. שנתנכרו מעשיו לאביו שבשמים, ואחד הגוי ואחד ישראל משמד במשמע (מכילתא):}%
% ---------- פרק 12 | פסוק 44 ----------
 \textbf{מד} וְכׇל־עֶ֥בֶד אִ֖ישׁ מִקְנַת־כָּ֑סֶף וּמַלְתָּ֣ה אֹת֔וֹ אָ֖ז יֹ֥אכַל בּֽוֹ׃ \Rashi{\textbf{ומלתה אתו אז יאכל בו.}רבו, מגיד שמילת עבדיו מעכבתו מלאכל בפסח, דברי רבי יהושע. רבי אליעזר אומר, אין מילת עבדיו מעכבתו מלאכל בפסח, אם כן מה ת"ל אז יאכל בו? העבד:}%
% ---------- פרק 12 | פסוק 45 ----------
 \textbf{מה} תּוֹשָׁ֥ב וְשָׂכִ֖יר לֹא־יֹ֥אכַל בּֽוֹ׃ \Rashi{\textbf{תושב.}זה גר תושב: ושכיר. זה גוי, ומה ת"ל, והלא ערלים הם, ונאמר וכל ערל לא יאכל בו? אלא כגון ערבי מהול וגבעוני מהול והוא תושב או שכיר (שם):}%
% ---------- פרק 12 | פסוק 46 ----------
 \textbf{מו} בְּבַ֤יִת אֶחָד֙ יֵאָכֵ֔ל לֹא־תוֹצִ֧יא מִן־הַבַּ֛יִת מִן־הַבָּשָׂ֖ר ח֑וּצָה וְעֶ֖צֶם לֹ֥א תִשְׁבְּרוּ־בֽוֹ׃ \Rashi{\textbf{בבית אחד יאכל.}בחבורה אחת, שלא יעשו הנמנין עליו שתי חבורות ויחלקוהו. אתה אומר בחבורה אחת או אינו אלא בבית אחד כמשמעו, וללמד שאם התחילו והיו אוכלים בחצר וירדו גשמים שלא יכנסו לבית? תלמוד לומר, "על הבתים אשר יאכלו אתו בהם", מכאן שהאוכל אוכל בשני מקומות: לא תוציא מן הבית. מן החבורה: ועצם לא תשברו בו. הראוי לאכילה, כגון שיש עליו כזית בשר, יש בו משום שבירת עצם, אין עליו כזית בשר אין בו משום שבירת עצם:}%
% ---------- פרק 12 | פסוק 47 ----------
 \textbf{מז} כׇּל־עֲדַ֥ת יִשְׂרָאֵ֖ל יַעֲשׂ֥וּ אֹתֽוֹ׃ \Rashi{\textbf{כל עדת ישראל יעשו אתו.}למה נאמר, לפי שהוא אומר בפסח מצרים "שה לבית אבות", שנמנו עליו למשפחות, יכול אף פסח דורות כן? ת"ל כל עדת ישראל יעשו אתו:}%
% ---------- פרק 12 | פסוק 48 ----------
 \textbf{מח} וְכִֽי־יָג֨וּר אִתְּךָ֜ גֵּ֗ר וְעָ֣שָׂה פֶ֘סַח֮ לַיהֹוָה֒ הִמּ֧וֹל ל֣וֹ כׇל־זָכָ֗ר וְאָז֙ יִקְרַ֣ב לַעֲשֹׂת֔וֹ וְהָיָ֖ה כְּאֶזְרַ֣ח הָאָ֑רֶץ וְכׇל־עָרֵ֖ל לֹֽא־יֹ֥אכַל בּֽוֹ׃ \Rashi{\textbf{ועשה פסח.}יכול כל המתגיר יעשה פסח מיד, תלמוד לומר והיה כאזרח הארץ, מה אזרח בארבעה עשר אף גר בארבעה עשר (שם): וכל ערל לא יאכל בו. להביא את שמתו אחיו מחמת מילה, שאינו מומר לערלות ואינו למד מ"בן נכר לא יאכל בו" (שם):}%
% ---------- פרק 12 | פסוק 49 ----------
 \textbf{מט} תּוֹרָ֣ה אַחַ֔ת יִהְיֶ֖ה לָֽאֶזְרָ֑ח וְלַגֵּ֖ר הַגָּ֥ר בְּתוֹכְכֶֽם׃ \Rashi{\textbf{תורה אחת וגו'.}להשוות גר לאזרח אף לשאר מצוות שבתורה (שם):}%
% ---------- פרק 12 | פסוק 50 ----------
 \textbf{נ} וַֽיַּעֲשׂ֖וּ כׇּל־בְּנֵ֣י יִשְׂרָאֵ֑ל כַּאֲשֶׁ֨ר צִוָּ֧ה יְהֹוָ֛ה אֶת־מֹשֶׁ֥ה וְאֶֽת־אַהֲרֹ֖ן כֵּ֥ן עָשֽׂוּ׃ \{ס\} %
% ---------- פרק 12 | פסוק 51 ----------
 \textbf{נא} וַיְהִ֕י בְּעֶ֖צֶם הַיּ֣וֹם הַזֶּ֑ה הוֹצִ֨יא יְהֹוָ֜ה אֶת־בְּנֵ֧י יִשְׂרָאֵ֛ל מֵאֶ֥רֶץ מִצְרַ֖יִם עַל־צִבְאֹתָֽם׃ \{פ\} %
% ---------- פרק 13 | פסוק 1 ----------
 \textbf{א} וַיְדַבֵּ֥ר יְהֹוָ֖ה אֶל־מֹשֶׁ֥ה לֵּאמֹֽר׃ %
% ---------- פרק 13 | פסוק 2 ----------
 \textbf{ב} קַדֶּשׁ־לִ֨י כׇל־בְּכ֜וֹר פֶּ֤טֶר כׇּל־רֶ֙חֶם֙ בִּבְנֵ֣י יִשְׂרָאֵ֔ל בָּאָדָ֖ם וּבַבְּהֵמָ֑ה לִ֖י הֽוּא׃ \Rashi{\textbf{פטר כל רחם.}שפתח הרחם תחלה, כמו: "פוטר מים ראשית מדון" (משלי י"ז), וכן "יפטירו בשפה" (תהלים כ"ב) – יפתחו שפתים: לי הוא. לעצמי קניתי, ע"י שהכיתי בכורי מצרים:}%
% ---------- פרק 13 | פסוק 3 ----------
 \textbf{ג} וַיֹּ֨אמֶר מֹשֶׁ֜ה אֶל־הָעָ֗ם זָכ֞וֹר אֶת־הַיּ֤וֹם הַזֶּה֙ אֲשֶׁ֨ר יְצָאתֶ֤ם מִמִּצְרַ֙יִם֙ מִבֵּ֣ית עֲבָדִ֔ים כִּ֚י בְּחֹ֣זֶק יָ֔ד הוֹצִ֧יא יְהֹוָ֛ה אֶתְכֶ֖ם מִזֶּ֑ה וְלֹ֥א יֵאָכֵ֖ל חָמֵֽץ׃ \Rashi{\textbf{זכור את היום הזה.}למד, שמזכירין יציאת מצרים בכל יום:}%
% ---------- פרק 13 | פסוק 4 ----------
 \textbf{ד} הַיּ֖וֹם אַתֶּ֣ם יֹצְאִ֑ים בְּחֹ֖דֶשׁ הָאָבִֽיב׃ \Rashi{\textbf{בחדש האביב.}וכי לא היינו יודעין באיזה חדש? אלא כך אמר להם, ראו חסד שגמלכם, שהוציא אתכם בחדש שהוא כשר לצאת, לא חמה ולא צנה ולא גשמים, וכן הוא אומר, "מוציא אסירים בכושרות" (תהלים ס"ח) – חדש שהוא כשר לצאת (מכילתא):}%
% ---------- פרק 13 | פסוק 5 ----------
 \textbf{ה} וְהָיָ֣ה כִֽי־יְבִיאֲךָ֣ יְהֹוָ֡ה אֶל־אֶ֣רֶץ הַֽ֠כְּנַעֲנִ֠י וְהַחִתִּ֨י וְהָאֱמֹרִ֜י וְהַחִוִּ֣י וְהַיְבוּסִ֗י אֲשֶׁ֨ר נִשְׁבַּ֤ע לַאֲבֹתֶ֙יךָ֙ לָ֣תֶת לָ֔ךְ אֶ֛רֶץ זָבַ֥ת חָלָ֖ב וּדְבָ֑שׁ וְעָבַדְתָּ֛ אֶת־הָעֲבֹדָ֥ה הַזֹּ֖את בַּחֹ֥דֶשׁ הַזֶּֽה׃ \Rashi{\textbf{אל ארץ הכנעני וגו'.}ואע"פ שלא מנה אלא חמשה עממין, כל ז' גוים במשמע, שכלן בכלל כנעני הם ואחת ממשפחות כנען היתה שלא נקרא לה שם אלא כנעני: נשבע לאבתיך. באברהם הוא אומר "ביום ההוא כרת ה' את אברם וגו'" (בראשית ט"ו), וביצחק הוא אומר "גור בארץ הזאת" (שם כ"ו), וביעקב הוא אומר "הארץ אשר אתה שוכב עליה" (שם כ"ח): זבת חלב ודבש. חלב זב מן העזים, והדבש זב מן התמרים ומן התאנים: את העבדה הזאת. של פסח; והלא כבר נאמר למעלה "והיה כי תבאו אל הארץ וגו'", ולמה חזר ושנאה? בשביל דבר שנתחדש בה, בפרשה ראשונה נאמר "והיה כי יאמרו אליכם בניכם מה העבדה הזאת לכם" — בבן רשע הכתוב מדבר שהוציא את עצמו מן הכלל, וכאן "והגדת לבנך" — בבן שאינו יודע לשאל, והכתוב מלמדך, שתפתח לו אתה בדברי אגדה המושכין את הלב (מכילתא):}%
% ---------- פרק 13 | פסוק 6 ----------
 \textbf{ו} שִׁבְעַ֥ת יָמִ֖ים תֹּאכַ֣ל מַצֹּ֑ת וּבַיּוֹם֙ הַשְּׁבִיעִ֔י חַ֖ג לַיהֹוָֽה׃ %
% ---------- פרק 13 | פסוק 7 ----------
 \textbf{ז} מַצּוֹת֙ יֵֽאָכֵ֔ל אֵ֖ת שִׁבְעַ֣ת הַיָּמִ֑ים וְלֹֽא־יֵרָאֶ֨ה לְךָ֜ חָמֵ֗ץ וְלֹֽא־יֵרָאֶ֥ה לְךָ֛ שְׂאֹ֖ר בְּכׇל־גְּבֻלֶֽךָ׃ %
% ---------- פרק 13 | פסוק 8 ----------
 \textbf{ח} וְהִגַּדְתָּ֣ לְבִנְךָ֔ בַּיּ֥וֹם הַה֖וּא לֵאמֹ֑ר בַּעֲב֣וּר זֶ֗ה עָשָׂ֤ה יְהֹוָה֙ לִ֔י בְּצֵאתִ֖י מִמִּצְרָֽיִם׃ \Rashi{\textbf{בעבור זה.}בעבור שאקים מצוותיו, כגון פסח מצה ומרור הללו: עשה ה' לי. רמז תשובה לבן רשע, לומר עשה ה' לי, ולא לך, שאלו היית שם לא היית כדאי לגאל (מכילתא):}%
% ---------- פרק 13 | פסוק 9 ----------
 \textbf{ט} וְהָיָה֩ לְךָ֨ לְא֜וֹת עַל־יָדְךָ֗ וּלְזִכָּרוֹן֙ בֵּ֣ין עֵינֶ֔יךָ לְמַ֗עַן תִּהְיֶ֛ה תּוֹרַ֥ת יְהֹוָ֖ה בְּפִ֑יךָ כִּ֚י בְּיָ֣ד חֲזָקָ֔ה הוֹצִֽאֲךָ֥ יְהֹוָ֖ה מִמִּצְרָֽיִם׃ \Rashi{\textbf{והיה לך לאות.}יציאת מצרים תהיה לך לאות: על ידך ולזכרון בית עיניך. שתכתב הפרשיות הללו ותקשרם בראש ובזרוע: על ידך. על יד שמאל, לפיכך ידכה מלא בפרשה שניה, לדרש בה יד שהיא כהה (מנחות ל"ו):}%
% ---------- פרק 13 | פסוק 10 ----------
 \textbf{י} וְשָׁמַרְתָּ֛ אֶת־הַחֻקָּ֥ה הַזֹּ֖את לְמוֹעֲדָ֑הּ מִיָּמִ֖ים יָמִֽימָה׃ \{פ\} \Rashi{\textbf{מימים ימימה.}משנה לשנה (שם):}%
% ---------- פרק 13 | פסוק 11 ----------
 \textbf{יא} וְהָיָ֞ה כִּֽי־יְבִאֲךָ֤ יְהֹוָה֙ אֶל־אֶ֣רֶץ הַֽכְּנַעֲנִ֔י כַּאֲשֶׁ֛ר נִשְׁבַּ֥ע לְךָ֖ וְלַֽאֲבֹתֶ֑יךָ וּנְתָנָ֖הּ לָֽךְ׃ \Rashi{\textbf{(והיה כי יבאך.}יש מרבותינו שלמדו מכאן שלא קדשו בכורות הנולדים במדבר, והאומר קדשו, מפרש ביאה זו אם תקימוהו במדבר תזכו לכנס לארץ ותקימוהו שם): נשבע לך. והיכן נשבע לך? "והבאתי אתכם אל הארץ אשר נשאתי וגו'" (שמות ו'): ונתנה לך. תהא בעיניך כאלו נתנה לך בו ביום, ואל תהי בעיניך כירשת אבות:}%
% ---------- פרק 13 | פסוק 12 ----------
 \textbf{יב} וְהַעֲבַרְתָּ֥ כׇל־פֶּֽטֶר־רֶ֖חֶם לַֽיהֹוָ֑ה וְכׇל־פֶּ֣טֶר ׀ שֶׁ֣גֶר בְּהֵמָ֗ה אֲשֶׁ֨ר יִהְיֶ֥ה לְךָ֛ הַזְּכָרִ֖ים לַיהֹוָֽה׃ \Rashi{\textbf{והעברת.}אין והעברת אלא לשון הפרשה, וכן הוא אומר "והעברת את נחלתו לבתו" (במדבר כ"ז) (מכילתא): שגר בהמה. נפל ששגרתו אמו ושלחתו בלא עתו; ולמדך הכתוב שהוא קדוש בבכורה לפטר את הבא אחריו ואף שאין נפל קרוי שגר, כמו "שגר אלפיך" (דברים ז'), אבל זה לא בא אלא ללמד על הנפל שהרי כבר כתב כל פטר רחם. וא"ת אף בכור בהמה טמאה במשמע, בא ופרש במקום אחר "בבקרך ובצאנך" (דברים ט"ו). לשון אחר יש לפרש, "והעברת כל פטר רחם" — בבכור אדם הכתוב מדבר:}%
% ---------- פרק 13 | פסוק 13 ----------
 \textbf{יג} וְכׇל־פֶּ֤טֶר חֲמֹר֙ תִּפְדֶּ֣ה בְשֶׂ֔ה וְאִם־לֹ֥א תִפְדֶּ֖ה וַעֲרַפְתּ֑וֹ וְכֹ֨ל בְּכ֥וֹר אָדָ֛ם בְּבָנֶ֖יךָ תִּפְדֶּֽה׃ \Rashi{\textbf{פטר חמור.}ולא פטר שאר בהמה טמאה; וגזרת הכתוב היא לפי שנמשלו בכורי מצרים לחמורים; ועוד, שסיעו את ישראל ביציאתן, שאין לך אחד מישראל שלא נטל הרבה חמורים ממצרים טעונים מכספם ומזהבם של מצרים: תפדה בשה. נותן שה לכהן ופטר חמור מתר בהנאה והשה חלין ביד כהן (בכורות ט'): וערפתו. עורפו בקופיץ מאחוריו והורגו; הוא הפסיד ממונו של כהן, לפיכך יפסיד ממונו (שם י'): וכל בכור אדם בבניך תפדה. חמש סלעים פדיונו קצוב במקום אחר (במדבר י"ח):}%
% ---------- פרק 13 | פסוק 14 ----------
 \textbf{יד} וְהָיָ֞ה כִּֽי־יִשְׁאָלְךָ֥ בִנְךָ֛ מָחָ֖ר לֵאמֹ֣ר מַה־זֹּ֑את וְאָמַרְתָּ֣ אֵלָ֔יו בְּחֹ֣זֶק יָ֗ד הוֹצִיאָ֧נוּ יְהֹוָ֛ה מִמִּצְרַ֖יִם מִבֵּ֥ית עֲבָדִֽים׃ \Rashi{\textbf{כי ישאלך בנך מחר.}יש מחר שהוא עכשיו ויש מחר שהוא לאחר זמן, כגון זה, וכגון "ולא יאמרו בניכם מחר לבנינו" (יהושע כ"ב), דבני גד ודבני ראובן: מה זאת. זה תינוק טפש שאינו יודע להעמיק שאלתו וסותם ושואל "מה זאת?", ובמקום אחר הוא אומר "מה העדת והחקים והמשפטים וגו'" (דברים ו'), הרי זאת שאלת בן חכם; דברה תורה כנגד ארבעה בנים — רשע, ושאינו יודע לשאל, והשואל דרך סתומה, והשואל דרך חכמה (מכילתא):}%
% ---------- פרק 13 | פסוק 15 ----------
 \textbf{טו} וַיְהִ֗י כִּֽי־הִקְשָׁ֣ה פַרְעֹה֮ לְשַׁלְּחֵ֒נוּ֒ וַיַּהֲרֹ֨ג יְהֹוָ֤ה כׇּל־בְּכוֹר֙ בְּאֶ֣רֶץ מִצְרַ֔יִם מִבְּכֹ֥ר אָדָ֖ם וְעַד־בְּכ֣וֹר בְּהֵמָ֑ה עַל־כֵּן֩ אֲנִ֨י זֹבֵ֜חַ לַֽיהֹוָ֗ה כׇּל־פֶּ֤טֶר רֶ֙חֶם֙ הַזְּכָרִ֔ים וְכׇל־בְּכ֥וֹר בָּנַ֖י אֶפְדֶּֽה׃ %
% ---------- פרק 13 | פסוק 16 ----------
 \textbf{טז} וְהָיָ֤ה לְאוֹת֙ עַל־יָ֣דְכָ֔ה וּלְטוֹטָפֹ֖ת בֵּ֣ין עֵינֶ֑יךָ כִּ֚י בְּחֹ֣זֶק יָ֔ד הוֹצִיאָ֥נוּ יְהֹוָ֖ה מִמִּצְרָֽיִם׃ \{ס\} \Rashi{\textbf{ולטוטפת.}תפלין; ועל שם שהם ארבעה בתים קרויין טטפת – "טט" בכתפי שתים, "פת" באפריקי שתים (סנהדרין ד'). ומנחם חברו עם "והטף אל דרום" (יחזקאל כ"א), "אל תטיפו" (מיכה ב'), לשון דבור, כמו ולזכרון, שהרואה אותם קשור בין העינים, יזכר הנס וידבר בו:}%
\addparasha{חומש שמות | פרשת בשלח}
% ---------- פרק 13 | פסוק 17 ----------
 \textbf{יז} וַיְהִ֗י בְּשַׁלַּ֣ח פַּרְעֹה֮ אֶת־הָעָם֒ וְלֹא־נָחָ֣ם אֱלֹהִ֗ים דֶּ֚רֶךְ אֶ֣רֶץ פְּלִשְׁתִּ֔ים כִּ֥י קָר֖וֹב ה֑וּא כִּ֣י ׀ אָמַ֣ר אֱלֹהִ֗ים פֶּֽן־יִנָּחֵ֥ם הָעָ֛ם בִּרְאֹתָ֥ם מִלְחָמָ֖ה וְשָׁ֥בוּ מִצְרָֽיְמָה׃ \Rashi{\textbf{ויהי בשלח פרעה, ולא נחם.}נהגם, כמו "לך נחה את העם" (שמות ל"ב), "בהתהלכך תנחה אותך" (משלי ו'): כי קרוב הוא. ונח לשוב באותו הדרך למצרים; ומדרשי אגדה יש הרבה: בראתם מלחמה. כגון מלחמת "וירד העמלקי והכנעני וגו'" (במדבר י"ד), אם הלכו דרך ישר היו חוזרים, ומה אם כשהקיפם דרך מעקם אמרו "נתנה ראש ונשובה מצרימה" (שם), אם הוליכם בפשוטה על אחת כמה וכמה: פן ינחם. יחשבו מחשבה על שיצאו ויתנו לב לשוב:}%
% ---------- פרק 13 | פסוק 18 ----------
 \textbf{יח} וַיַּסֵּ֨ב אֱלֹהִ֧ים ׀ אֶת־הָעָ֛ם דֶּ֥רֶךְ הַמִּדְבָּ֖ר יַם־ס֑וּף וַחֲמֻשִׁ֛ים עָל֥וּ בְנֵי־יִשְׂרָאֵ֖ל מֵאֶ֥רֶץ מִצְרָֽיִם׃ \Rashi{\textbf{ויסב.}הסבם מן הדרך הפשוטה לדרך העקומה: ים סוף. כמו לים סוף, וסוף הוא לשון אגם שגדלים בו קנים, כמו "ותשם בסוף" (שמות ב'), "קנה וסוף קמלו" (ישעיהו י"ט): וחמשים. אין חמושים אלא מזינים; (לפי שהסבן במדבר הוא גרם להם שעלו חמושים, שאלו הסבן דרך ישוב, לא היו מחמשים להם כל מה שצריכין, אלא כאדם שעובר ממקום למקום ובדעתו לקנות שם מה שיצטרך, אבל כשהוא פורש למדבר צריך לזמן לו כל הצרך; ומקרא זה לא נכתב כי אם לשבר את האזן, שלא תתמה במלחמת עמלק ובמלחמת סיחון ועוג ומדין מהיכן היו להם כלי זין שהכו אותם בחרב) וכן הוא אומר "ואתם תעברו חמשים" (יהושע א'), וכן תרגמו אנקלוס "מזרזין", כמו "וירק את חניכיו" (בראשית י"ד) – וזריז. דבר אחר, חמשים אחד מחמשה יצאו וארבעה חלקים מתו בשלשת ימי אפלה (מכילתא):}%
% ---------- פרק 13 | פסוק 19 ----------
 \textbf{יט} וַיִּקַּ֥ח מֹשֶׁ֛ה אֶת־עַצְמ֥וֹת יוֹסֵ֖ף עִמּ֑וֹ כִּי֩ הַשְׁבֵּ֨עַ הִשְׁבִּ֜יעַ אֶת־בְּנֵ֤י יִשְׂרָאֵל֙ לֵאמֹ֔ר פָּקֹ֨ד יִפְקֹ֤ד אֱלֹהִים֙ אֶתְכֶ֔ם וְהַעֲלִיתֶ֧ם אֶת־עַצְמֹתַ֛י מִזֶּ֖ה אִתְּכֶֽם׃ \Rashi{\textbf{השבע השביע.}השביעם שישביעו לבניהם. ולמה לא השביע בניו שישאוהו לארץ כנען מיד כמו שהשביע יעקב? אמר יוסף, אני שליט הייתי במצרים והיה ספק בידי לעשות, אבל בני לא יניחום מצרים לעשות, לכך השביעם לכשיגאלו ויצאו משם שישאוהו: והעליתם את עצמתי מזה אתכם. לאחיו השביע כן, למדנו שאף עצמות כל השבטים העלו עמהם שנאמר אתכם (מכילתא):}%
% ---------- פרק 13 | פסוק 20 ----------
 \textbf{כ} וַיִּסְע֖וּ מִסֻּכֹּ֑ת וַיַּחֲנ֣וּ בְאֵתָ֔ם בִּקְצֵ֖ה הַמִּדְבָּֽר׃ \Rashi{\textbf{ויסעו מסכת.}ביום השני, שהרי בראשון באו מרעמסס לסכות:}%
% ---------- פרק 13 | פסוק 21 ----------
 \textbf{כא} וַֽיהֹוָ֡ה הֹלֵךְ֩ לִפְנֵיהֶ֨ם יוֹמָ֜ם בְּעַמּ֤וּד עָנָן֙ לַנְחֹתָ֣ם הַדֶּ֔רֶךְ וְלַ֛יְלָה בְּעַמּ֥וּד אֵ֖שׁ לְהָאִ֣יר לָהֶ֑ם לָלֶ֖כֶת יוֹמָ֥ם וָלָֽיְלָה׃ \Rashi{\textbf{לנחתם הדרך.}נקוד פתח שהוא כמו להנחותם, "לראותכם בדרך אשר תלכו בה" (דברים א') שהוא כמו להראותכם, אף כאן להנחתם ע"י שליח, ומי הוא השליח? עמוד הענן; והקדוש ברוך הוא בכבודו מוליכו לפניהם, ומכל מקום את עמוד הענן הכין להנחותם על ידו, שהרי ע"י עמוד הענן הם הולכים; עמוד הענן אינו לאורה אלא להורותם הדרך:}%
% ---------- פרק 13 | פסוק 22 ----------
 \textbf{כב} לֹֽא־יָמִ֞ישׁ עַמּ֤וּד הֶֽעָנָן֙ יוֹמָ֔ם וְעַמּ֥וּד הָאֵ֖שׁ לָ֑יְלָה לִפְנֵ֖י הָעָֽם׃ \{פ\} \Rashi{\textbf{לא ימיש.}הקב"ה את עמוד הענן יומם ועמוד האש לילה; מגיד שעמוד הענן משלים לעמוד האש ועמוד האש משלים לעמוד הענן, שעד שלא ישקע זה עולה זה:}%
% ---------- פרק 14 | פסוק 1 ----------
 \textbf{א} וַיְדַבֵּ֥ר יְהֹוָ֖ה אֶל־מֹשֶׁ֥ה לֵּאמֹֽר׃ %
% ---------- פרק 14 | פסוק 2 ----------
 \textbf{ב} דַּבֵּר֮ אֶל־בְּנֵ֣י יִשְׂרָאֵל֒ וְיָשֻׁ֗בוּ וְיַחֲנוּ֙ לִפְנֵי֙ פִּ֣י הַחִירֹ֔ת בֵּ֥ין מִגְדֹּ֖ל וּבֵ֣ין הַיָּ֑ם לִפְנֵי֙ בַּ֣עַל צְפֹ֔ן נִכְח֥וֹ תַחֲנ֖וּ עַל־הַיָּֽם׃ \Rashi{\textbf{וישבו.}לאחוריהם, לצד מצרים היו מקרבין כל יום השלישי, כדי להטעות את פרעה שיאמר תועים הם בדרך, כמו שנאמר "ואמר פרעה לבני ישראל וגו'": ויחנו לפני פי החירת. הוא פיתום, ועכשו נקרא פי החירת על שם שנעשו שם בני חרין; והם שני סלעים גבוהים זקופים והגיא שביניהם קרוי פי הסלעים: לפני בעל צפן. הוא נשאר מכל אלהי מצרים כדי להטעותן, שיאמרו קשה יראתן, ועליו פרש איוב "משגיא לגוים ויאבדם" (איוב י"ב):}%
% ---------- פרק 14 | פסוק 3 ----------
 \textbf{ג} וְאָמַ֤ר פַּרְעֹה֙ לִבְנֵ֣י יִשְׂרָאֵ֔ל נְבֻכִ֥ים הֵ֖ם בָּאָ֑רֶץ סָגַ֥ר עֲלֵיהֶ֖ם הַמִּדְבָּֽר׃ \Rashi{\textbf{ואמר פרעה.}כשישמע שהם שבים לאחוריהם: לבני ישראל. על בני ישראל. וכן (פסוק יד) "ה' ילחם לכם" – עליכם, (בראשית כ יג) "אמרי לי אחי הוא" – אמרי עלי: נבכים הם. כלואים ומשקעים, ובלעז שירי"ר, כמו "נבכי ים" (איוב ל"ח), "בעמק הבכא" (תהילים פ"ד) "מבכי נהרות" (איוב כ"ח). נבכים הם, כלואים הם במדבר, שאינן יודעין לצאת ממנו ולהיכן ילכו:}%
% ---------- פרק 14 | פסוק 4 ----------
 \textbf{ד} וְחִזַּקְתִּ֣י אֶת־לֵב־פַּרְעֹה֮ וְרָדַ֣ף אַחֲרֵיהֶם֒ וְאִכָּבְדָ֤ה בְּפַרְעֹה֙ וּבְכׇל־חֵיל֔וֹ וְיָדְע֥וּ מִצְרַ֖יִם כִּֽי־אֲנִ֣י יְהֹוָ֑ה וַיַּֽעֲשׂוּ־כֵֽן׃ \Rashi{\textbf{ואכבדה בפרעה.}כשהקב"ה מתנקם ברשעים שמו מתגדל ומתכבד, וכן הוא אומר "ונשפטתי אתו וגו'" ואחר כך "והתגדלתי והתקדשתי ונודעתי וגו'" (יחזקאל ל"ח), ואומר "שמה שבר רשפי קשת", ואחר כך "נודע ביהודה אלהים" (תהלים ע"ו), ואומר "נודע ה' משפט עשה" (שם ט'): בפרעה ובכל חילו. הוא התחיל בעברה וממנו התחילה הפרענות: ויעשו כן. להגיד שבחן, ששמעו לקול משה, ולא אמרו האיך נתקרב אל רודפינו? אנו צריכים לברח אלא אמרו אין לנו לקים אלא דברי בן עמרם (מכילתא):}%
% ---------- פרק 14 | פסוק 5 ----------
 \textbf{ה} וַיֻּגַּד֙ לְמֶ֣לֶךְ מִצְרַ֔יִם כִּ֥י בָרַ֖ח הָעָ֑ם וַ֠יֵּהָפֵ֠ךְ לְבַ֨ב פַּרְעֹ֤ה וַעֲבָדָיו֙ אֶל־הָעָ֔ם וַיֹּֽאמְרוּ֙ מַה־זֹּ֣את עָשִׂ֔ינוּ כִּֽי־שִׁלַּ֥חְנוּ אֶת־יִשְׂרָאֵ֖ל מֵעׇבְדֵֽנוּ׃ \Rashi{\textbf{ויגד למלך מצרים.}איקטורין שלח עמהם, וכיון שהגיעו לשלשת ימים שקבעו לילך ולשוב וראו שאינן חוזרין למצרים, באו והגידו לפרעה ביום הרביעי, בחמישי ובששי רדפו אחריהם, ליל שביעי ירדו לים, בשחרית אמרו שירה והוא יום שביעי של פסח; לכך אנו קורין השירה ביום השביעי: ויהפך. נהפך ממה שהיה, שהרי אמר להם "קומו צאו מתוך עמי" (שמות י"ב), ונהפך לב עבדיו, שהרי לשעבר היו אומרים לו "עד מתי יהיה זה לנו למוקש" (שמות י'), ועכשיו נהפכו לרדף אחריהם בשביל ממונם שהשאילום: מעבדנו. מעבד אותנו:}%
% ---------- פרק 14 | פסוק 6 ----------
 \textbf{ו} וַיֶּאְסֹ֖ר אֶת־רִכְבּ֑וֹ וְאֶת־עַמּ֖וֹ לָקַ֥ח עִמּֽוֹ׃ \Rashi{\textbf{ויאסר את רכבו.}הוא בעצמו (מכילתא): ואת עמו לקח עמו. משכם בדברים, לקינו ונטלו ממוננו ושלחנום, באו עמי ואני לא אתנהג עמכם כשאר מלכים, דרך שאר מלכים עבדיו קודמין לו במלחמה, ואני אקדים לפניכם, שנאמר ופרעה הקריב – הקריב עצמו ומהר לפני חילותיו –, דרך שאר מלכים לטל בזה בראש כמו שיבחר, אני אשוה עמכם בחלק שנאמר "אחלק שלל" (שמות ט"ו):}%
% ---------- פרק 14 | פסוק 7 ----------
 \textbf{ז} וַיִּקַּ֗ח שֵׁשׁ־מֵא֥וֹת רֶ֙כֶב֙ בָּח֔וּר וְכֹ֖ל רֶ֣כֶב מִצְרָ֑יִם וְשָׁלִשִׁ֖ם עַל־כֻּלּֽוֹ׃ \Rashi{\textbf{בחור.}נבחרים, בחור לשון יחיד, כל רכב ורכב שבמנין זה היה בחור: וכל רכב מצרים. ועמהם כל שאר הרכב; ומהיכן היו הבהמות הללו? אם תאמר משל מצרים, הרי נאמר "וימת כל מקנה מצרים" (שמות ט'), ואם משל ישראל, והלא נאמר "וגם מקננו ילך עמנו", משל מי היו? מן הירא את דבר ה'. מכאן היה ר' שמעון אומר, כשר שבמצרים הרג, טוב שבנחשים רצץ את מחו (מכילתא): ושלשם על כלו. שרי צבאות, כתרגומו:}%
% ---------- פרק 14 | פסוק 8 ----------
 \textbf{ח} וַיְחַזֵּ֣ק יְהֹוָ֗ה אֶת־לֵ֤ב פַּרְעֹה֙ מֶ֣לֶךְ מִצְרַ֔יִם וַיִּרְדֹּ֕ף אַחֲרֵ֖י בְּנֵ֣י יִשְׂרָאֵ֑ל וּבְנֵ֣י יִשְׂרָאֵ֔ל יֹצְאִ֖ים בְּיָ֥ד רָמָֽה׃ \Rashi{\textbf{ויחזק ה' את לב פרעה.}שהיה תולה אם לרדף אם לאו, וחזק את לבו לרדף (שם): ביד רמה. בגבורה גבוהה ומפרסמת:}%
% ---------- פרק 14 | פסוק 9 ----------
 \textbf{ט} וַיִּרְדְּפ֨וּ מִצְרַ֜יִם אַחֲרֵיהֶ֗ם וַיַּשִּׂ֤יגוּ אוֹתָם֙ חֹנִ֣ים עַל־הַיָּ֔ם כׇּל־סוּס֙ רֶ֣כֶב פַּרְעֹ֔ה וּפָרָשָׁ֖יו וְחֵיל֑וֹ עַל־פִּי֙ הַֽחִירֹ֔ת לִפְנֵ֖י בַּ֥עַל צְפֹֽן׃ %
% ---------- פרק 14 | פסוק 10 ----------
 \textbf{י} וּפַרְעֹ֖ה הִקְרִ֑יב וַיִּשְׂאוּ֩ בְנֵֽי־יִשְׂרָאֵ֨ל אֶת־עֵינֵיהֶ֜ם וְהִנֵּ֥ה מִצְרַ֣יִם ׀ נֹסֵ֣עַ אַחֲרֵיהֶ֗ם וַיִּֽירְאוּ֙ מְאֹ֔ד וַיִּצְעֲק֥וּ בְנֵֽי־יִשְׂרָאֵ֖ל אֶל־יְהֹוָֽה׃ \Rashi{\textbf{ופרעה הקריב.}היה לו לכתב ופרעה קרב, מהו הקריב? הקריב עצמו ונתאמץ לקדם לפניהם כמו שהתנה עמהם: נסע אחריהם. בלב אחד כאיש אחד. דבר אחר – והנה מצרים נסע אחריהם, ראו שר של מצרים נוסע מן השמים לעזר למצרים; – תנחומא: ויצעקו. תפשו אמנות אבותם – באברהם הוא אומר "אל המקום אשר עמד שם" (בראשית י"ט), ביצחק "לשוח בשדה", (שם כ"ד), ביעקב "ויפגע במקום" (שם כ"ח):}%
% ---------- פרק 14 | פסוק 11 ----------
 \textbf{יא} וַיֹּאמְרוּ֮ אֶל־מֹשֶׁה֒ הֲֽמִבְּלִ֤י אֵין־קְבָרִים֙ בְּמִצְרַ֔יִם לְקַחְתָּ֖נוּ לָמ֣וּת בַּמִּדְבָּ֑ר מַה־זֹּאת֙ עָשִׂ֣יתָ לָּ֔נוּ לְהוֹצִיאָ֖נוּ מִמִּצְרָֽיִם׃ \Rashi{\textbf{המבלי אין קברים.}וכי מחמת חסרון קברים – שאין קברים במצרים לקבר שם – לקחתנו משם. ש"י פו"ר פלינ"צא ד"י נו"ן פוש"יש בלעז:}%
% ---------- פרק 14 | פסוק 12 ----------
 \textbf{יב} הֲלֹא־זֶ֣ה הַדָּבָ֗ר אֲשֶׁר֩ דִּבַּ֨רְנוּ אֵלֶ֤יךָ בְמִצְרַ֙יִם֙ לֵאמֹ֔ר חֲדַ֥ל מִמֶּ֖נּוּ וְנַֽעַבְדָ֣ה אֶת־מִצְרָ֑יִם כִּ֣י ט֥וֹב לָ֙נוּ֙ עֲבֹ֣ד אֶת־מִצְרַ֔יִם מִמֻּתֵ֖נוּ בַּמִּדְבָּֽר׃ \Rashi{\textbf{אשר דברנו אליך בארץ מצרים.}והיכן דברו? "ירא ה' עליכם וישפט" (שמות ה') (מכילתא): ממתנו. מאשר נמות, ואם היה נקוד מלאפום היה נבאר ממיתתנו, עכשיו שנקוד בשורק נבאר מאשר נמות וכן "מי יתן מותנו" (שמות ט"ז) – שנמות, וכן "מי יתן מותי" דאבשלום (שמואל ב' י"ט) – שאמות. כמו "ליום קומי לעד" (צפניה ג'), "עד יום שובי בשלום" (דברי הימים ב' י"ח) – שאקום, שאשוב:}%
% ---------- פרק 14 | פסוק 13 ----------
 \textbf{יג} וַיֹּ֨אמֶר מֹשֶׁ֣ה אֶל־הָעָם֮ אַל־תִּירָ֒אוּ֒ הִֽתְיַצְּב֗וּ וּרְאוּ֙ אֶת־יְשׁוּעַ֣ת יְהֹוָ֔ה אֲשֶׁר־יַעֲשֶׂ֥ה לָכֶ֖ם הַיּ֑וֹם כִּ֗י אֲשֶׁ֨ר רְאִיתֶ֤ם אֶת־מִצְרַ֙יִם֙ הַיּ֔וֹם לֹ֥א תֹסִ֛פוּ לִרְאֹתָ֥ם ע֖וֹד עַד־עוֹלָֽם׃ \Rashi{\textbf{כי אשר ראיתם את מצרים וגו'.}מה שראיתם אותם אינו אלא היום, היום הוא שראיתם אותם ולא תוסיפו עוד:}%
% ---------- פרק 14 | פסוק 14 ----------
 \textbf{יד} יְהֹוָ֖ה יִלָּחֵ֣ם לָכֶ֑ם וְאַתֶּ֖ם תַּחֲרִשֽׁוּן׃ \{פ\} \Rashi{\textbf{ילחם לכם.}בשבילכם; וכן "כי ה' נלחם להם", וכן "אם לאל תריבון" (איוב י"ג), וכן "ואשר דבר לי" (בראשית כ"ד), וכן "האתם תריבון לבעל" (שופטים ו'):}%
% ---------- פרק 14 | פסוק 15 ----------
 \textbf{טו} וַיֹּ֤אמֶר יְהֹוָה֙ אֶל־מֹשֶׁ֔ה מַה־תִּצְעַ֖ק אֵלָ֑י דַּבֵּ֥ר אֶל־בְּנֵי־יִשְׂרָאֵ֖ל וְיִסָּֽעוּ׃ \Rashi{\textbf{מה תצעק אלי.}למדנו שהיה משה עומד ומתפלל, אמר לו הקב"ה, לא עת עתה להאריך בתפלה, שישראל נתונין בצרה (מכילתא). ד"א – מה תצעק אלי, עלי הדבר תלוי ולא עליך, כמו שנאמר להלן "על בני ועל פעל ידי תצוני" (ישעיהו מ"ה): דבר אל בני ישראל ויסעו. אין להם אלא לסע, שאין הים עומד בפניהם, כדאי הוא זכות אבותיהם והאמונה שהאמינו בי ויצאו, לקרע להם הים (מכילתא):}%
% ---------- פרק 14 | פסוק 16 ----------
 \textbf{טז} וְאַתָּ֞ה הָרֵ֣ם אֶֽת־מַטְּךָ֗ וּנְטֵ֧ה אֶת־יָדְךָ֛ עַל־הַיָּ֖ם וּבְקָעֵ֑הוּ וְיָבֹ֧אוּ בְנֵֽי־יִשְׂרָאֵ֛ל בְּת֥וֹךְ הַיָּ֖ם בַּיַּבָּשָֽׁה׃ %
% ---------- פרק 14 | פסוק 17 ----------
 \textbf{יז} וַאֲנִ֗י הִנְנִ֤י מְחַזֵּק֙ אֶת־לֵ֣ב מִצְרַ֔יִם וְיָבֹ֖אוּ אַחֲרֵיהֶ֑ם וְאִכָּבְדָ֤ה בְּפַרְעֹה֙ וּבְכׇל־חֵיל֔וֹ בְּרִכְבּ֖וֹ וּבְפָרָשָֽׁיו׃ %
% ---------- פרק 14 | פסוק 18 ----------
 \textbf{יח} וְיָדְע֥וּ מִצְרַ֖יִם כִּי־אֲנִ֣י יְהֹוָ֑ה בְּהִכָּבְדִ֣י בְּפַרְעֹ֔ה בְּרִכְבּ֖וֹ וּבְפָרָשָֽׁיו׃ %
% ---------- פרק 14 | פסוק 19 ----------
 \textbf{יט} וַיִּסַּ֞ע מַלְאַ֣ךְ הָאֱלֹהִ֗ים הַהֹלֵךְ֙ לִפְנֵי֙ מַחֲנֵ֣ה יִשְׂרָאֵ֔ל וַיֵּ֖לֶךְ מֵאַחֲרֵיהֶ֑ם וַיִּסַּ֞ע עַמּ֤וּד הֶֽעָנָן֙ מִפְּנֵיהֶ֔ם וַֽיַּעֲמֹ֖ד מֵאַחֲרֵיהֶֽם׃ \Rashi{\textbf{וילך מאחריהם.}להבדיל בין מחנה מצרים ובין מחנה ישראל ולקבל חצים ובליסטראות של מצרים. בכל מקום הוא אומר מלאך ה' וכאן מלאך האלהים, אין אלהים בכל מקום אלא דין, מלמד שהיו ישראל נתונין בדין באותה שעה אם להנצל אם להאבד עם מצרים (שם): ויסע עמוד הענן. כשחשכה והשלים עמוד הענן את המחנה לעמוד האש, לא נסתלק הענן כמו שהיה רגיל להסתלק ערבית לגמרי, אלא נסע והלך לו מאחריהם להחשיך למצרים:}%
% ---------- פרק 14 | פסוק 20 ----------
 \textbf{כ} וַיָּבֹ֞א בֵּ֣ין ׀ מַחֲנֵ֣ה מִצְרַ֗יִם וּבֵין֙ מַחֲנֵ֣ה יִשְׂרָאֵ֔ל וַיְהִ֤י הֶֽעָנָן֙ וְהַחֹ֔שֶׁךְ וַיָּ֖אֶר אֶת־הַלָּ֑יְלָה וְלֹא־קָרַ֥ב זֶ֛ה אֶל־זֶ֖ה כׇּל־הַלָּֽיְלָה׃ \Rashi{\textbf{ויבא בין מחנה מצרים.}משל למהלך בדרך ובנו מהלך לפניו, באו לסטים לשבותו, נטלו מלפניו ונתנו לאחריו, בא זאב מאחריו, נתנו לפניו, באו לסטים לפניו וזאבים מאחריו, נתנו על זרועו ונלחם בהם, כך "ואנכי תרגלתי לאפרים קחם על זרועתיו" (הושע י"א) (מכילתא): ויהי הענן והחשך. למצרים: ויאר עמוד האש את הלילה לישראל, והולך לפניהם כדרכו ללכת כל הלילה והחשך של ערפל לצד מצרים: ולא קרב זה אל זה. מחנה אל מחנה (שם):}%
% ---------- פרק 14 | פסוק 21 ----------
 \textbf{כא} וַיֵּ֨ט מֹשֶׁ֣ה אֶת־יָדוֹ֮ עַל־הַיָּם֒ וַיּ֣וֹלֶךְ יְהֹוָ֣ה ׀ אֶת־הַ֠יָּ֠ם בְּר֨וּחַ קָדִ֤ים עַזָּה֙ כׇּל־הַלַּ֔יְלָה וַיָּ֥שֶׂם אֶת־הַיָּ֖ם לֶחָרָבָ֑ה וַיִּבָּקְע֖וּ הַמָּֽיִם׃ \Rashi{\textbf{ברוח קדים עזה.}ברוח קדים שהיא עזה שברוחות, היא הרוח שהקב"ה נפרע בה מן הרשעים, שנאמר, "כרוח קדים אפיצם" (ירמיה י"ח), "יבוא קדים רוח ה'" (הושע י"ג), "רוח הקדים שברך בלב ימים" (יחזקאל כ"ז), "הגה ברוחו הקשה ביום קדים" (ישעיה כ"ז): ויבקעו המים. כל מים שבעולם:}%
% ---------- פרק 14 | פסוק 22 ----------
 \textbf{כב} וַיָּבֹ֧אוּ בְנֵֽי־יִשְׂרָאֵ֛ל בְּת֥וֹךְ הַיָּ֖ם בַּיַּבָּשָׁ֑ה וְהַמַּ֤יִם לָהֶם֙ חוֹמָ֔ה מִֽימִינָ֖ם וּמִשְּׂמֹאלָֽם׃ %
% ---------- פרק 14 | פסוק 23 ----------
 \textbf{כג} וַיִּרְדְּפ֤וּ מִצְרַ֙יִם֙ וַיָּבֹ֣אוּ אַחֲרֵיהֶ֔ם כֹּ֚ל ס֣וּס פַּרְעֹ֔ה רִכְבּ֖וֹ וּפָרָשָׁ֑יו אֶל־תּ֖וֹךְ הַיָּֽם׃ \Rashi{\textbf{כל סוס פרעה.}וכי סוס אחד היה? מגיד שאין כלם חשובין לפני המקום אלא כסוס אחד (מכילתא):}%
% ---------- פרק 14 | פסוק 24 ----------
 \textbf{כד} וַֽיְהִי֙ בְּאַשְׁמֹ֣רֶת הַבֹּ֔קֶר וַיַּשְׁקֵ֤ף יְהֹוָה֙ אֶל־מַחֲנֵ֣ה מִצְרַ֔יִם בְּעַמּ֥וּד אֵ֖שׁ וְעָנָ֑ן וַיָּ֕הׇם אֵ֖ת מַחֲנֵ֥ה מִצְרָֽיִם׃ \Rashi{\textbf{באשמרת הבקר.}שלשת חלקי הלילה קרויין אשמורה, ואותה שלפני הבקר קורא אשמרת הבקר. ואומר אני, לפי שהלילה חלוק למשמרות שיר של מלאכי השרת, כת אחר כת, לשלשה חלקים, לכך קרוי אשמרת, וזהו שתרגם אנקלוס "מטרת": וישקף. ויבט, כלומר פנה אליהם להשחיתם. ותרגומו "ואסתכי" אף הוא לשון הבטה, כמו "שדה צופים", (במדבר כ"ג) – חקל סכותא: בעמוד אש וענן. עמוד ענן יורד ועושה אותו כטיט ועמוד אש מרתיחו, וטלפי סוסיהם משתמטות (מכילתא): ויהם. לשון מהומה, אשדורד"ישון בלעז, ערבבם, נטל סגניות שלהם. ושנינו בפרקי רבי אליעזר בנו של רבי יוסי הגלילי: כל מקום שנאמר בו מהומה הרעמת קול הוא, וזה אב לכלן "וירעם ה' בקול גדול וגו' על פלשתים ויהמם" (שמואל א ז'):}%
% ---------- פרק 14 | פסוק 25 ----------
 \textbf{כה} וַיָּ֗סַר אֵ֚ת אֹפַ֣ן מַרְכְּבֹתָ֔יו וַֽיְנַהֲגֵ֖הוּ בִּכְבֵדֻ֑ת וַיֹּ֣אמֶר מִצְרַ֗יִם אָנ֙וּסָה֙ מִפְּנֵ֣י יִשְׂרָאֵ֔ל כִּ֣י יְהֹוָ֔ה נִלְחָ֥ם לָהֶ֖ם בְּמִצְרָֽיִם׃ \{פ\} \Rashi{\textbf{ויסר את אפן מרכבתיו.}מכח האש נשרפו הגלגלים, והמרכבות נגררות, והיושבים בהם נעים ואבריהן מתפרקין (מכילתא): וינהגהו בכבדת. בהנהגה שהיא כבדה וקשה להם; במדה שמדדו: "ויכבד" לבו הוא ועבדיו (שמות ט'), אף כאן וינהגהו "בכבדת": נלחם להם במצרים. במצריים. דבר אחר במצרים – בארץ מצרים, שכשם שאלו לוקים על הים כך לוקים אותם שנשארו במצרים (מכילתא):}%
% ---------- פרק 14 | פסוק 26 ----------
 \textbf{כו} וַיֹּ֤אמֶר יְהֹוָה֙ אֶל־מֹשֶׁ֔ה נְטֵ֥ה אֶת־יָדְךָ֖ עַל־הַיָּ֑ם וְיָשֻׁ֤בוּ הַמַּ֙יִם֙ עַל־מִצְרַ֔יִם עַל־רִכְבּ֖וֹ וְעַל־פָּרָשָֽׁיו׃ \Rashi{\textbf{וישבו המים.}שזקופין ועומדים כחומה, ישובו למקומם ויכסו על מצרים:}%
% ---------- פרק 14 | פסוק 27 ----------
 \textbf{כז} וַיֵּט֩ מֹשֶׁ֨ה אֶת־יָד֜וֹ עַל־הַיָּ֗ם וַיָּ֨שׇׁב הַיָּ֜ם לִפְנ֥וֹת בֹּ֙קֶר֙ לְאֵ֣יתָנ֔וֹ וּמִצְרַ֖יִם נָסִ֣ים לִקְרָאת֑וֹ וַיְנַעֵ֧ר יְהֹוָ֛ה אֶת־מִצְרַ֖יִם בְּת֥וֹךְ הַיָּֽם׃ \Rashi{\textbf{לפנות בקר.}לעת שהבקר פונה לבא: לאיתנו. לתקפו הראשון (מכילתא): נסים לקראתו. שהיו מהממים ומטרפים ורצין לקראת המים: וינער ה'. כאדם שמנער את הקדרה והופך העליון למטה והתחתון למעלה, כך היו עולין ויורדין ומשתברין בים, ונתן הקב"ה בהם חיות לקבל היסורין (שם): וינער. "ושניק", והוא לשון טרוף בלשון ארמי, והרבה יש במ"א:}%
% ---------- פרק 14 | פסוק 28 ----------
 \textbf{כח} וַיָּשֻׁ֣בוּ הַמַּ֗יִם וַיְכַסּ֤וּ אֶת־הָרֶ֙כֶב֙ וְאֶת־הַפָּ֣רָשִׁ֔ים לְכֹל֙ חֵ֣יל פַּרְעֹ֔ה הַבָּאִ֥ים אַחֲרֵיהֶ֖ם בַּיָּ֑ם לֹֽא־נִשְׁאַ֥ר בָּהֶ֖ם עַד־אֶחָֽד׃ \Rashi{\textbf{ויכסו את הרכב, לכל חיל פרעה.}כך דרך המקראות לכתב למ"ד יתרה, כמו "לכל כליו תעשה נחשת" (שמות כ"ז), וכן "לכל כלי המשכן בכל עבדתו" (במדבר ד'), "ויתדותם ומיתריהם לכל כליהם" (שם), ואינה אלא תקון לשון:}%
% ---------- פרק 14 | פסוק 29 ----------
 \textbf{כט} וּבְנֵ֧י יִשְׂרָאֵ֛ל הָלְכ֥וּ בַיַּבָּשָׁ֖ה בְּת֣וֹךְ הַיָּ֑ם וְהַמַּ֤יִם לָהֶם֙ חֹמָ֔ה מִֽימִינָ֖ם וּמִשְּׂמֹאלָֽם׃ %
% ---------- פרק 14 | פסוק 30 ----------
 \textbf{ל} וַיּ֨וֹשַׁע יְהֹוָ֜ה בַּיּ֥וֹם הַה֛וּא אֶת־יִשְׂרָאֵ֖ל מִיַּ֣ד מִצְרָ֑יִם וַיַּ֤רְא יִשְׂרָאֵל֙ אֶת־מִצְרַ֔יִם מֵ֖ת עַל־שְׂפַ֥ת הַיָּֽם׃ \Rashi{\textbf{וירא ישראל את מצרים מת.}שפלטן הים על שפתו, כדי שלא יאמרו ישראל, כשם שאנו עולים מצד זה כך הם עולין מצד אחר, רחוק ממנו, וירדפו אחרינו (פסחים קי"ח):}%
% ---------- פרק 14 | פסוק 31 ----------
 \textbf{לא} וַיַּ֨רְא יִשְׂרָאֵ֜ל אֶת־הַיָּ֣ד הַגְּדֹלָ֗ה אֲשֶׁ֨ר עָשָׂ֤ה יְהֹוָה֙ בְּמִצְרַ֔יִם וַיִּֽירְא֥וּ הָעָ֖ם אֶת־יְהֹוָ֑ה וַֽיַּאֲמִ֙ינוּ֙ בַּֽיהֹוָ֔ה וּבְמֹשֶׁ֖ה עַבְדּֽוֹ׃ \{פ\} \Rashi{\textbf{את היד הגדולה.}את הגבורה הגדולה שעשתה ידו של הקב"ה. והרבה לשונות נופלין על לשון יד וכלן לשון יד ממש הן, והמפרשו יתקן הלשון אחר ענין הדבור:}%
% ---------- פרק 15 | פסוק 1 ----------
 \textbf{א} אָ֣ז יָשִֽׁיר־מֹשֶׁה֩ וּבְנֵ֨י יִשְׂרָאֵ֜ל אֶת־הַשִּׁירָ֤ה הַזֹּאת֙ לַֽיהֹוָ֔ה וַיֹּאמְר֖וּ לֵאמֹ֑ר        אָשִׁ֤ירָה לַֽיהֹוָה֙ כִּֽי־גָאֹ֣ה גָּאָ֔ה        ס֥וּס וְרֹכְב֖וֹ רָמָ֥ה בַיָּֽם׃ \Rashi{\textbf{אז ישיר משה.}אז כשראה הנס עלה בלבו שישיר שירה. וכן "אז ידבר יהושע" (יהושע י'), וכן "ובית יעשה לבת פרעה" (מלכים א ז') – חשב בלבו שיעשה לה, אף כאן ישיר אמר לו לבו שישיר וכן עשה – ויאמרו לאמר אשירה לה', וכן ביהושע כשראה הנס אמר לו לבו שידבר וכן עשה – "ויאמר לעיני ישראל" (יהושע י'), וכן שירת הבאר, שפתח בה אז ישיר ישראל, פרש אחריו "עלי באר ענו לה" (במדבר י"א), "אז יבנה שלמה במה" (מלכים א י"א), פרשו בו חכמי ישראל שבקש לבנות ולא בנה, למדנו שהיו"ד על שם המחשבה נאמרה, זהו לישב פשוטו. אבל מדרשו אמרו רבותינו ז"ל "מכאן רמז לתחית המתים מן התורה" וכן בכלן, חוץ משל שלמה, שפרשוהו בקש לבנות ולא בנה. ואין לומר ולישב לשון הזה כשאר דברים הנכתבים בלשון עתיד והן מיד, כגון "ככה יעשה איוב" (איוב א'), "ע"פ ה' יחנו" (במדבר ט'), "ויש אשר יהיה הענן" (שם), לפי שהן דבר ההווה תמיד ונופל בו בין לשון עתיד ובין לשון עבר, אבל זה שלא היה אלא לשעה, איני יכול לישבו בלשון הזה: כי גאה גאה. כתרגומו. (ד"א – בא הכפל לומר שעשה דבר שאי אפשר לבשר ודם לעשות; כשהוא נלחם בחברו ומתגבר עליו, מפילו מן הסוס, וכאן הסוס ורכבו רמה בים, וכל שאי אפשר לעשות על ידי זולתו נופל בו לשון גאות, כמו "כי גאות עשה" (ישעיהו י"ב), וכן כל השירה תמצא כפולה, עזי וזמרת יה ויהי לי לישועה, ה' איש מלחמה ה' שמו, וכן כלם). דבר אחר – כי גאה גאה, על כל השירות וכל מה שאקלס בו, עוד יש בו תוספת, ולא כמדת מלך בשר ודם שמקלסין אותו ואין בו: סוס ורכבו. שניהם קשורים זה בזה והמים מעלין אותן ויורדין לעמק ואינן נפרדין (מכילתא): רמה. השליך, וכן "ורמיו לגוא אתון נורא" (דניאל ג'). ומדרש אגדה: כתוב אחד אומר רמה, וכתוב אחד אומר ירה, מלמד שהיו עולין לרום ויורדין לתהום, כמו "מי ירה אבן פנתה" (איוב ל"ח), מלמעלה למטה:}%
% ---------- פרק 15 | פסוק 2 ----------
 \textbf{ב} עׇזִּ֤י וְזִמְרָת֙ יָ֔הּ וַֽיְהִי־לִ֖י לִֽישׁוּעָ֑ה        זֶ֤ה אֵלִי֙ וְאַנְוֵ֔הוּ        אֱלֹהֵ֥י אָבִ֖י וַאֲרֹמְמֶֽנְהוּ׃ \Rashi{\textbf{עזי וזמרת יה.}אנקלוס תרגם "תקפי ותשבחתי", עזי כמו עזי, וזמרת כמו וזמרתי, ואני תמה על לשון המקרא, שאין לך כמוהו בנקדתו במקרא אלא בשלשה מקומות שהוא סמוך אצל וזמרת, וכל שאר מקומות נקוד שור"ק, "ה' עזי ומעזי" (ירמיהו ט"ז), "עזו אליך אשמרה" (תהילים נ"ט). וכן כל תבה בת שתי אותיות הנקודה מלאפום כשהיא מארכת באות שלישית ואין השניה בחטף, הראשונה נקודה בשורק, כגון עז עזי, רק רקי, חק חקי, על עלו – "וסר מעלהם עלו" (ישעיהו י"ד), כל כלו – "ושלשם על כלו" (שמות י"ד), ואלו ג' עזי וזמרת של כאן, ושל ישעיה, ושל תהלים, נקודים בחטף קמץ, ועוד אין באחד מהם כתוב וזמרתי אלא וזמרת, וכלם סמוך להם "ויהי לי לישועה". לכך אני אומר, לישב לשון המקרא, שאין עזי כמו עזי ולא וזמרת כמו וזמרתי, אלא עזי שם דבר הוא כמו "הישבי בשמים" (תהילים קכ"ג), "שכני בחגוי סלע" (עובדיה א'), "שכני סנה" (דברים ל"ג), וזהו השבח עזי וזמרת יה הוא היה לי לישועה. וזמרת דבוק הוא לתבת השם, כמו "לעזרת ה'" (שופטים ה'), "בעברת ה'" (ישעיהו ט'), "על דברת בני האדם" (קהלת ג'). ולשון וזמרת לשון "לא תזמר" (ויקרא כ"ה), "זמיר עריצים" (ישעיהו כ"ה), לשון כסוח וכריתה – עזו ונקמתו של אלהינו היה לנו לישועה. ואל תתמה על לשון ויהי, שלא נאמר היה, שיש לנו מקראות מדברים בלשון זה, וזה דגמתו: "את קירות הבית סביב להיכל ולדביר ויעש צלעות סביב" (מלכים א ו'), היה לו לומר עשה צלעות סביב. וכן בדברי הימים: "ובני ישראל הישבים בערי יהודה וימלך עליהם רחבעם" (דברי הימים ב' י'), היה לו לומר מלך עליהם רחבעם. "מבלתי יכלת ה' וגו' וישחטם" (במדבר י"ד), היה לו לומר שחטם. "והאנשים אשר שלח משה וגו' וימתו" (שם) מתו היה לו לומר, "ואשר לא שם לבו אל דבר ה' ויעזב" (שמות ט'), היה לו לומר עזב: זה אלי. בכבודו נגלה עליהם והיו מראין אותו באצבע, ראתה שפחה על הים מה שלא ראו נביאים (מכילתא): ואנוהו. אנקלוס תרגם לשון נוה – "נוה שאנן" (ישעיהו ל"ג), "לנוה צאן" (שם ס"ה). דבר אחר – ואנוהו לשון נוי, אספר נויו ושבחו לבאי עולם, כגון "מה דודך מדוד וגו' דודי צח ואדום" (שיר השירים ה'), וכל הענין: אלהי אבי. הוא זה וארממנהו: אלהי אבי. לא אני תחלת הקדשה אלא מחזקת ועומדת לי הקדשה ואלהותו עלי מימי אבותי:}%
% ---------- פרק 15 | פסוק 3 ----------
 \textbf{ג} יְהֹוָ֖ה אִ֣ישׁ מִלְחָמָ֑ה יְהֹוָ֖ה שְׁמֽוֹ׃ \Rashi{\textbf{ה' איש מלחמה.}בעל מלחמות, כמו "איש נעמי" (רות א'), וכל איש ואישך מתרגמין בעל, וכן "וחזקת והיית לאיש" (מלכים א ב') – לגבור: ה' שמו. מלחמותיו לא בכלי זין אלא בשמו הוא נלחם; כמו שאמר דוד "ואנכי בא אליך בשם ה' צבאות" (שמואל א י"ז). דבר אחר, ה' שמו – אף בשעה שהוא נלחם ונוקם מאויביו, אוחז הוא במדתו לרחם על ברואיו ולזון את כל באי עולם; ולא כמדת מלכי אדמה כשהוא עוסק במלחמה פונה עצמו מכל עסקים ואין בו כח לעשות זו וזו (מכילתא):}%
% ---------- פרק 15 | פסוק 4 ----------
 \textbf{ד} מַרְכְּבֹ֥ת פַּרְעֹ֛ה וְחֵיל֖וֹ יָרָ֣ה בַיָּ֑ם        וּמִבְחַ֥ר שָֽׁלִשָׁ֖יו טֻבְּע֥וּ בְיַם־סֽוּף׃ \Rashi{\textbf{ירה בים.}"שדי בימא", שדי לשון יריה, וכן הוא אומר "או ירה יירה" (שמות י"ט) – או אשתדאה ישתדי, והתי"ו משמש באלו במקום יתפעל: ומבחר. שם דבר. כמו מרכב, משכב, מקרא קדש: טבעו. אין טביעה אלא במקום טיט, כמו "טבעתי ביון מצולה" (תהילים ס"ט), "ויטבע ירמיהו בטיט" (ירמיהו ל"ח), מלמד, שנעשה הים טיט, לגמל להם כמדתם ששעבדו את ישראל בחמר ובלבנים (מכילתא):}%
% ---------- פרק 15 | פסוק 5 ----------
 \textbf{ה} תְּהֹמֹ֖ת יְכַסְיֻ֑מוּ יָרְד֥וּ בִמְצוֹלֹ֖ת כְּמוֹ־אָֽבֶן׃ \Rashi{\textbf{יכסימו.}כמו יכסום, והיו"ד האמצעית יתרה בו; ודרך מקראות בכך, כמו "וצאנך ירבין" (דברים ח'), "ירוין מדשן ביתך" (תהלים ל"ו), והיו"ד ראשונה שמשמעה לשון עתיד, כך פרשהו: טבעו בים סוף כדי שיחזרו המים ויכסו אותן. יכסימו אין דומה לו במקרא בנקדתו, ודרכו להיות נקוד יכסימו מלאפום: כמו אבן. ובמקום אחר "צללו כעופרת", ובמקום אחר "יאכלמו כקש", הרשעים כקש, הולכים ומטרפין עולין ויורדין, בינונים כאבן, והכשרים כעופרת שנחו מיד (מכילתא):}%
% ---------- פרק 15 | פסוק 6 ----------
 \textbf{ו} יְמִֽינְךָ֣ יְהֹוָ֔ה נֶאְדָּרִ֖י בַּכֹּ֑חַ        יְמִֽינְךָ֥ יְהֹוָ֖ה תִּרְעַ֥ץ אוֹיֵֽב׃ \Rashi{\textbf{ימינך ימינך.}שני פעמים, כשישראל עושין את רצונו של מקום השמאל נעשית ימין (שם): ימינך ה' נאדרי בכח. להציל את ישראל וימינך השנית תרעץ אויב. ולי נראה אותה ימין עצמה תרעץ אויב, מה שאי אפשר לאדם – לעשות שתי מלאכות ביד אחת; ופשוטו של מקרא ימינך הנאדרת בכח מה מלאכתה? ימינך ה' תרעץ אויב; וכמה מקראות דגמתו, "כי הנה אויביך ה' כי הנה אויביך יאבדו" (תהלים צ"ב), ודומיהם: נאדרי. היו"ד יתרה, כמו "רבתי עם, שרתי במדינות" (איכה א'), "גנבתי יום" (בראשית ל"א): תרעץ אויב. תמיד היא רועצת ומשברת האויב; ודומה לו "וירעצו וירוצצו את בני ישראל", בשופטים (י, ח). (ד"א – ימינך הנאדרת בכח היא משברת ומלקה אויב):}%
% ---------- פרק 15 | פסוק 7 ----------
 \textbf{ז} וּבְרֹ֥ב גְּאוֹנְךָ֖ תַּהֲרֹ֣ס קָמֶ֑יךָ        תְּשַׁלַּח֙ חֲרֹ֣נְךָ֔ יֹאכְלֵ֖מוֹ כַּקַּֽשׁ׃ \Rashi{\textbf{וברב גאונך.}זאת היד בלבד רועצת האויב – כשהוא מרימה ברב גאונו אז יהרס קמיו – ואם ברב גאונו לבד אויביו נהרסים קל וחמר כששלח בם חרון אף יאכלמו): תהרס. תמיד אתה הורס קמיך הקמים נגדך, ומי הם הקמים כנגדו? אלו הקמים על ישראל; וכן הוא אומר "כי הנה אויביך יהמיון" (תהלים פ"ג), ומה היא ההמיה? "על עמך יערימו סוד" (שם), ועל זה קורא אותם אויביו של מקום:}%
% ---------- פרק 15 | פסוק 8 ----------
 \textbf{ח} וּבְר֤וּחַ אַפֶּ֙יךָ֙ נֶ֣עֶרְמוּ מַ֔יִם        נִצְּב֥וּ כְמוֹ־נֵ֖ד נֹזְלִ֑ים        קָֽפְא֥וּ תְהֹמֹ֖ת בְּלֶב־יָֽם׃ \Rashi{\textbf{וברוח אפיך.}היוצא משני נחירים של אף. דבר הכתוב כביכול בשכינה דגמת מלך בשר ודם, כדי להשמיע אזן הבריות כפי ההוה, שיוכלו להבין דבר; כשאדם כועס יוצא רוח מנחיריו; וכן "עלה עשן באפו" (תהלים י"ח), וכן "ומרוח אפו יכלו" (איוב ד'), וזהו שאמר "למען שמי אאריך אפי", (ישעיהו מ"ח) – כשזעפו נח, נשימתו ארכה, וכשהוא כועס נשימתו קצרה, "ותהלתי אחטם לך" (שם) – ולמען תהלתי אשים חטם באפי לסתם נחירי בפני האף והרוח שלא יצאו, לך – בשבילך, אחטם, כמו "נאקה בחטם" במסכת שבת, כך נראה בעיני. וכל חרון אף שבמקרא אני אומר כן; חרה אף כמו "ועצמי חרה מני חרב" (איוב ל'), לשון שרפה ומוקד, שהנחירים מתחממים ונחרים בעת הקצף, וחרון מגזרת חרה כמו רצון מגזרת רצה, וכן חמה לשון חמימות, על כן הוא אומר "וחמתו בערה בו" (אסתר א'), ובנוח החמה אומר "נתקררה" דעתו: נערמו מים. אנקלוס תרגם לשון ערמימות, ולשון צחות המקרא כמו "ערמת חטים" (שיר השירים ז'), ונצבו כמו נד יוכיח: נערמו. ממוקד רוח שיצא מאפך יבשו המים, והם נעשו כמין גלים וכריות של ערמה, שהם גבוהים: כמו נד. כתרגומו "כשור" – כחומה: נד. לשון צבור וכנוס, כמו "נד קציר ביום נחלה" (ישעיהו י״ז:י״א), "כונס כנד" (תהילים ל״ג:ז׳), לא כתוב כונס כנאד אלא כנד, ואלו היה כנד כמו כנאד וכונס לשון הכנסה, היה לו לכתב מכניס כבנאד מי הים, אלא כונס לשון אוסף וצובר הוא, וכן "קמו נד אחד" (יהושע ג׳:י״ג), "ויעמדו נד אחד" (שם), ואין לשון קימה ועמידה בנאדות אלא בחומות וצבורים, ולא מצינו נאד נקוד אלא במלאפום כמו "שימה דמעתי בנאדך" (תהילים נ״ו:ט׳), "את נאד החלב" (שופטים ד׳:י״ט): קפאו. כמו "וכגבנה תקפיאני" (איוב י׳:י׳), שהקשו ונעשו כאבנים והמים זורקים את המצריים על האבן בכח ונלחמים בם בכל מיני קשי: בלב ים. בחזק הים; ודרך המקראות לדבר כן, "עד לב השמים" (דברים ד'), "בלב האלה" (שמואל ב י"ח), לשון עקרו ותקפו של דבר:}%
% ---------- פרק 15 | פסוק 9 ----------
 \textbf{ט} אָמַ֥ר אוֹיֵ֛ב אֶרְדֹּ֥ף אַשִּׂ֖יג        אֲחַלֵּ֣ק שָׁלָ֑ל תִּמְלָאֵ֣מוֹ נַפְשִׁ֔י        אָרִ֣יק חַרְבִּ֔י תּוֹרִישֵׁ֖מוֹ יָדִֽי׃ \Rashi{\textbf{אמר אויב.}לעמו, כשפתם בדברים, ארדף ואשיגם ואחלק שלל עם שרי ועבדי: תמלאמו. תתמלא מהם, נפשי רוחי ורצוני. ואל תתמה על תבה המדברת בשתים, תמלאמו – תמלא מהם, יש הרבה בלשון הזה "כי ארץ הנגב נתתני" (שופטים א׳:ט״ו), כמו נתת לי, "ולא יכלו דברו לשלום" (בראשית ל״ז:ד׳), כמו דבר עמו, "בני יצאוני" (ירמיהו י׳:כ׳), כמו יצאו ממני, "מספר צעדי אגידנו" (איוב ל״א:ל״ז), כמו אגיד לו, אף כאן "תמלאמו" – תמלא נפשי מהם: אריק חרבי. אשלף, ועל שם שהוא מריק את התער בשליפתו ונשאר ריק, נופל בו לשון הרקה, כמו "מריקים שקיהם" (בראשית מ״ב:ל״ה), "וכליו יריקו" (ירמיהו מ״ח:י״ב). ואל תאמר אין לשון ריקות נופל על היוצא, אלא על התיק ועל השק ועל הכלי שיצא ממנה, אבל לא על החרב ועל היין, ולדחק ולפרש "אריק חרבי" כלשון "וירק את חניכיו" (בראשית י״ד:י״ד), אזדין בחרבי, מצינו הלשון מוסב אף על היוצא, "שמן תורק" (שיר השירים א׳:ג׳), "ולא הורק מכלי אל כלי" (ירמיהו מ״ח:י״א), לא הורק הכלי אין כתיב כאן אלא לא הורק היין מכלי אל כלי, מצינו הלשון מוסב על היין, וכן "והריקו חרבותם על יפי חכמתך", דחירם (יחזקאל כ״ח:ז׳): תורישמו. לשון רישות ודלות, כמו "מוריש ומעשיר" (שמואל א' ב'):}%
% ---------- פרק 15 | פסוק 10 ----------
 \textbf{י} נָשַׁ֥פְתָּ בְרוּחֲךָ֖ כִּסָּ֣מוֹ יָ֑ם        צָֽלְלוּ֙ כַּֽעוֹפֶ֔רֶת בְּמַ֖יִם אַדִּירִֽים׃ \Rashi{\textbf{נשפת.}לשון הפחה, וכן "וגם נשף בהם" (ישעיה מ'): צללו. שקעו, עמקו, לשון מצולה: כעופרת. אבר, פלו"ם בלעז:}%
% ---------- פרק 15 | פסוק 11 ----------
 \textbf{יא} מִֽי־כָמֹ֤כָה בָּֽאֵלִם֙ יְהֹוָ֔ה        מִ֥י כָּמֹ֖כָה נֶאְדָּ֣ר בַּקֹּ֑דֶשׁ        נוֹרָ֥א תְהִלֹּ֖ת עֹ֥שֵׂה פֶֽלֶא׃ \Rashi{\textbf{באלם.}בחזקים, כמו "ואת אילי הארץ לקח" (יחזקאל י"ז). "אילותי לעזרתי חושה" (תהלים כ"ב): נורא תהלת. יראוי מלהגיד תהלותיך פן ימעטו, כמו שכתוב "לך דומיה תהלה" (תהילים ס״ה:ב׳):}%
% ---------- פרק 15 | פסוק 12 ----------
 \textbf{יב} נָטִ֙יתָ֙ יְמִ֣ינְךָ֔ תִּבְלָעֵ֖מוֹ אָֽרֶץ׃ \Rashi{\textbf{נטית ימינך.}כשהקב"ה נוטה ידו, הרשעים כלים ונופלים, לפי שהכל נתון בידו ונופלים בהטיתה; וכן הוא אומר "וה' יטה ידו וכשל עוזר ונפל עזר" (ישעיהו ל"א), משל לכלי זכוכית הנתונים ביד אדם, מטה ידו מעט והן נופלים ומשתברין (מכילתא): תבלעמו ארץ. מכאן שזכו לקבורה בשכר שאמרו ה' הצדיק (מכילתא):}%
% ---------- פרק 15 | פסוק 13 ----------
 \textbf{יג} נָחִ֥יתָ בְחַסְדְּךָ֖ עַם־ז֣וּ גָּאָ֑לְתָּ        נֵהַ֥לְתָּ בְעׇזְּךָ֖ אֶל־נְוֵ֥ה קׇדְשֶֽׁךָ׃ \Rashi{\textbf{נהלת.}לשון מנהל, ואנקלוס תרגם לשון נושא וסובל, ולא דקדק לפרש אחר לשון העברית:}%
% ---------- פרק 15 | פסוק 14 ----------
 \textbf{יד} שָֽׁמְע֥וּ עַמִּ֖ים יִרְגָּז֑וּן        חִ֣יל אָחַ֔ז יֹשְׁבֵ֖י פְּלָֽשֶׁת׃ \Rashi{\textbf{ירגזון.}מתרגזין: ישבי פלשת. מפני שהרגו את בני אפרים – שמהרו את הקץ ויצאו בחזקה כמפרש בדברי הימים – והרגום אנשי גת (מכילתא):}%
% ---------- פרק 15 | פסוק 15 ----------
 \textbf{טו} אָ֤ז נִבְהֲלוּ֙ אַלּוּפֵ֣י אֱד֔וֹם        אֵילֵ֣י מוֹאָ֔ב יֹֽאחֲזֵ֖מוֹ רָ֑עַד        נָמֹ֕גוּ כֹּ֖ל יֹשְׁבֵ֥י כְנָֽעַן׃ \Rashi{\textbf{אלופי אדום אילי מואב.}והלא לא היה להם לירא כלום, שהרי לא עליהם הולכים? אלא מפני אנינות שהיו מתאוננים ומצטערים על כבודם של ישראל (ילקוט שמעוני): נמגו. נמסו, כמו "ברביבים תמוגגנה" (תהלים ס"ה), אמרו, עלינו הם באים, לכלותינו ולירש את ארצנו (מכילתא):}%
% ---------- פרק 15 | פסוק 16 ----------
 \textbf{טז} תִּפֹּ֨ל עֲלֵיהֶ֤ם אֵימָ֙תָה֙ וָפַ֔חַד        בִּגְדֹ֥ל זְרוֹעֲךָ֖ יִדְּמ֣וּ כָּאָ֑בֶן        עַד־יַעֲבֹ֤ר עַמְּךָ֙ יְהֹוָ֔ה        עַֽד־יַעֲבֹ֖ר עַם־ז֥וּ קָנִֽיתָ׃ \Rashi{\textbf{תפל עליהם אימתה.}על הרחוקים: ופחד. על הקרובים, כענין שנאמר "כי שמענו את אשר הוביש וגו'" (יהושע ב'): עד יעבר, עד יעבר. כתרגומו: קנית. חבבת משאר אמות, כחפץ הקנוי בדמים יקרים שחביב על האדם:}%
% ---------- פרק 15 | פסוק 17 ----------
 \textbf{יז} תְּבִאֵ֗מוֹ וְתִטָּעֵ֙מוֹ֙ בְּהַ֣ר נַחֲלָֽתְךָ֔        מָכ֧וֹן לְשִׁבְתְּךָ֛ פָּעַ֖לְתָּ יְהֹוָ֑ה        מִקְּדָ֕שׁ אֲדֹנָ֖י כּוֹנְנ֥וּ יָדֶֽיךָ׃ \Rashi{\textbf{תבאמו.}נתנבא משה שלא יכנס לארץ לכך לא נאמר "תביאנו" (בבא בתרא קי"ט): מכון לשבתך. מקדש של מטה מכון כנגד כסא של מעלה אשר פעלת (מכילתא): מקדש ה'. הטעם עליו זקף גדול, להפרידו מתבת השם שלאחריו, המקדש אשר כוננו ידיך ה'. חביב בית המקדש, שהעולם נברא ביד אחת, שנאמר "אף ידי יסדה ארץ" (ישעיהו מ"ח), ומקדש בשתי ידים, ואימתי יבנה בשתי ידים? בזמן ש"ה' ימלך לעלם ועד" – לעתיד לבא שכל המלוכה שלו:}%
% ---------- פרק 15 | פסוק 18 ----------
 \textbf{יח} יְהֹוָ֥ה ׀ יִמְלֹ֖ךְ לְעֹלָ֥ם וָעֶֽד׃ \Rashi{\textbf{לעלם ועד.}לשון עולמות הוא והוי"ו בו יסוד, לפיכך היא פתוחה, אבל "אנכי היודע ועד" (ירמיהו כ"ט), שהוי"ו בו שמוש, קמוצה היא:}%
% ---------- פרק 15 | פסוק 19 ----------
 \textbf{יט} כִּ֣י בָא֩ ס֨וּס פַּרְעֹ֜ה בְּרִכְבּ֤וֹ וּבְפָרָשָׁיו֙ בַּיָּ֔ם        וַיָּ֧שֶׁב יְהֹוָ֛ה עֲלֵהֶ֖ם אֶת־מֵ֣י הַיָּ֑ם        וּבְנֵ֧י יִשְׂרָאֵ֛ל הָלְכ֥וּ בַיַּבָּשָׁ֖ה בְּת֥וֹךְ הַיָּֽם׃ \{פ\} \Rashi{\textbf{כי בא סוס פרעה.}כאשר בא:}%
% ---------- פרק 15 | פסוק 20 ----------
 \textbf{כ} וַתִּקַּח֩ מִרְיָ֨ם הַנְּבִיאָ֜ה אֲח֧וֹת אַהֲרֹ֛ן אֶת־הַתֹּ֖ף בְּיָדָ֑הּ וַתֵּצֶ֤אןָ כׇֽל־הַנָּשִׁים֙ אַחֲרֶ֔יהָ בְּתֻפִּ֖ים וּבִמְחֹלֹֽת׃ \Rashi{\textbf{ותקח מרים הנביאה.}היכן נתנבאה? כשהיתה אחות אהרן, קדם שנולד משה, אמרה עתידה אמי שתלד בן וכו' כדאיתא בסוטה. (דף י"ב). ד"א — אחות אהרן, לפי שמסר נפשו עליה כשנצטרעה נקראת על שמו: את התף. כלי של מיני זמר: בתפים ובמחלת. מבטחות היו צדקניות שבדור שהקב"ה עושה להם נסים והוציאו תפים ממצרים (מכילתא):}%
% ---------- פרק 15 | פסוק 21 ----------
 \textbf{כא} וַתַּ֥עַן לָהֶ֖ם מִרְיָ֑ם שִׁ֤ירוּ לַֽיהֹוָה֙ כִּֽי־גָאֹ֣ה גָּאָ֔ה ס֥וּס וְרֹכְב֖וֹ רָמָ֥ה בַיָּֽם׃ \{ס\} \Rashi{\textbf{ותען להם מרים.}משה אמר שירה לאנשים – הוא אומר והם עונין אחריו – ומרים אמרה שירה לנשים (סוטה ל'):}%
% ---------- פרק 15 | פסוק 22 ----------
 \textbf{כב} וַיַּסַּ֨ע מֹשֶׁ֤ה אֶת־יִשְׂרָאֵל֙ מִיַּם־ס֔וּף וַיֵּצְא֖וּ אֶל־מִדְבַּר־שׁ֑וּר וַיֵּלְכ֧וּ שְׁלֹֽשֶׁת־יָמִ֛ים בַּמִּדְבָּ֖ר וְלֹא־מָ֥צְאוּ מָֽיִם׃ \Rashi{\textbf{ויסע משה.}הסיען בעל כרחם, שעטרו מצרים את סוסיהם בתכשיטי זהב וכסף ואבנים טובות, והיו ישראל מוצאין אותן בים – וגדולה היתה בזת הים מבזת מצרים שנאמר "תורי זהב נעשה לך עם נקדות הכסף" (שיר השירים א') – לפיכך הצרך להסיען בעל כרחם:}%
% ---------- פרק 15 | פסוק 23 ----------
 \textbf{כג} וַיָּבֹ֣אוּ מָרָ֔תָה וְלֹ֣א יָֽכְל֗וּ לִשְׁתֹּ֥ת מַ֙יִם֙ מִמָּרָ֔ה כִּ֥י מָרִ֖ים הֵ֑ם עַל־כֵּ֥ן קָרָֽא־שְׁמָ֖הּ מָרָֽה׃ \Rashi{\textbf{ויבאו מרתה.}כמו למרה, ה"א בסוף תבה במקום למ"ד בתחלתה, והתי"ו היא במקום ה"א הנשרשת בתבת מרה, ובסמיכתה, כשהיא נדבקת לה"א שהוא מוסיף במקום הלמ"ד, תהפך הה"א של שרש לתי"ו; וכן כל ה"א שהיא שרש בתבה תתהפך לתי"ו בסמיכתה, כמו "חמה אין לי" (ישעיהו כ"ז), "וחמתו בערה בו" (אסתר א'), הרי ה"א של שרש נהפכת לתי"ו מפני שנסמכת אל הוי"ו הנוספת, וכן "עבד ואמה" (ויקרא כ"ה), "הנה אמתי בלהה" (בראשית ל'), "לנפש חיה" (בראשית ב׳:ז׳), "וזהמתו חיתו לחם" (איוב ל"ג), "בין הרמה" (שופטים ד'), "ותשבתו הרמתה" (שמואל א ז'):}%
% ---------- פרק 15 | פסוק 24 ----------
 \textbf{כד} וַיִּלֹּ֧נוּ הָעָ֛ם עַל־מֹשֶׁ֥ה לֵּאמֹ֖ר מַה־נִּשְׁתֶּֽה׃ \Rashi{\textbf{וילנו.}לשון נפעל הוא, וכן התרגום לשון נפעל הוא ואתרעמו, וכן דרך לשון תלונה להסב הדבור אל האדם – מתלונן, מתרועם ולא אמר לונן, רועם, וכן יאמר הלועז דקומפ"לישנ"ט ש"י בלעז, מסב הדבור אליו באמרו ש"י:}%
% ---------- פרק 15 | פסוק 25 ----------
 \textbf{כה} וַיִּצְעַ֣ק אֶל־יְהֹוָ֗ה וַיּוֹרֵ֤הוּ יְהֹוָה֙ עֵ֔ץ וַיַּשְׁלֵךְ֙ אֶל־הַמַּ֔יִם וַֽיִּמְתְּק֖וּ הַמָּ֑יִם שָׁ֣ם שָׂ֥ם ל֛וֹ חֹ֥ק וּמִשְׁפָּ֖ט וְשָׁ֥ם נִסָּֽהוּ׃ \Rashi{\textbf{שם שם לו.}במרה נתן להם מקצת פרשיות של תורה שיתעסקו בהם, שבת ופרה אדמה ודינין (סנהדרין נ"ו): ושם נסהו. לעם, וראה קשי ערפו, שלא נמלכו במשה בלשון יפה "בקש עלינו רחמים שיהיה לנו מים לשתות" אלא נתלוננו:}%
% ---------- פרק 15 | פסוק 26 ----------
 \textbf{כו} וַיֹּ֩אמֶר֩ אִם־שָׁמ֨וֹעַ תִּשְׁמַ֜ע לְק֣וֹל ׀ יְהֹוָ֣ה אֱלֹהֶ֗יךָ וְהַיָּשָׁ֤ר בְּעֵינָיו֙ תַּעֲשֶׂ֔ה וְהַֽאֲזַנְתָּ֙ לְמִצְוֺתָ֔יו וְשָׁמַרְתָּ֖ כׇּל־חֻקָּ֑יו כׇּֽל־הַמַּחֲלָ֞ה אֲשֶׁר־שַׂ֤מְתִּי בְמִצְרַ֙יִם֙ לֹא־אָשִׂ֣ים עָלֶ֔יךָ כִּ֛י אֲנִ֥י יְהֹוָ֖ה רֹפְאֶֽךָ׃ \{ס\} \Rashi{\textbf{אם שמוע תשמע.}זו קבלה שיקבלו עליהם: תעשה. היא עשיה: והאזנת. תטה אזנים לדקדק בהם: כל חקיו. דברים שאינן אלא גזרת מלך, בלא שום טעם, ויצר הרע מקנטר עליהם: מה אסור באלו? למה נאסרו? כגון לבישת כלאים ואכילת חזיר ופרה אדמה וכיוצא בהם (יומא ס"ז): לא אשים עליך. ואם אשים הרי הוא כלא הושמה, כי אני ה' רפאך – זהו מדרשו. ולפי פשוטו כי אני ה' רפאך ומלמדך תורה ומצוות למען תנצל מהם, כרופא הזה האומר לאדם, אל תאכל דבר זה פן יביאך לידי חלי זה, וכן הוא אומר "רפאות תהי לשרך" (משלי ג'):}%
% ---------- פרק 15 | פסוק 27 ----------
 \textbf{כז} וַיָּבֹ֣אוּ אֵילִ֔מָה וְשָׁ֗ם שְׁתֵּ֥ים עֶשְׂרֵ֛ה עֵינֹ֥ת מַ֖יִם וְשִׁבְעִ֣ים תְּמָרִ֑ים וַיַּחֲנוּ־שָׁ֖ם עַל־הַמָּֽיִם׃ \Rashi{\textbf{שתים עשרה עינת מים.}כנגד י"ב שבטים נזדמנו להם: ושבעים תמרים. כנגד שבעים זקנים (מכילתא):}%
% ---------- פרק 16 | פסוק 1 ----------
 \textbf{א} וַיִּסְעוּ֙ מֵֽאֵילִ֔ם וַיָּבֹ֜אוּ כׇּל־עֲדַ֤ת בְּנֵֽי־יִשְׂרָאֵל֙ אֶל־מִדְבַּר־סִ֔ין אֲשֶׁ֥ר בֵּין־אֵילִ֖ם וּבֵ֣ין סִינָ֑י בַּחֲמִשָּׁ֨ה עָשָׂ֥ר יוֹם֙ לַחֹ֣דֶשׁ הַשֵּׁנִ֔י לְצֵאתָ֖ם מֵאֶ֥רֶץ מִצְרָֽיִם׃ \Rashi{\textbf{בחמשה עשר יום.}נתפרש היום של חניה זו, לפי שבו ביום כלתה החררה שהוציאו ממצרים והצרכו למן, ללמדנו, שאכלו משירי הבצק ששים ואחת סעודות וירד להם מן בששה עשר באיר, ויום א' בשבת היה כדאיתא במסכת שבת (פ"ז):}%
% ---------- פרק 16 | פסוק 2 ----------
 \textbf{ב} (וילינו) [וַיִּלּ֜וֹנוּ] כׇּל־עֲדַ֧ת בְּנֵי־יִשְׂרָאֵ֛ל עַל־מֹשֶׁ֥ה וְעַֽל־אַהֲרֹ֖ן בַּמִּדְבָּֽר׃ \Rashi{\textbf{וילונו.}לפי שכלה הלחם:}%
% ---------- פרק 16 | פסוק 3 ----------
 \textbf{ג} וַיֹּאמְר֨וּ אֲלֵהֶ֜ם בְּנֵ֣י יִשְׂרָאֵ֗ל מִֽי־יִתֵּ֨ן מוּתֵ֤נוּ בְיַד־יְהֹוָה֙ בְּאֶ֣רֶץ מִצְרַ֔יִם בְּשִׁבְתֵּ֙נוּ֙ עַל־סִ֣יר הַבָּשָׂ֔ר בְּאׇכְלֵ֥נוּ לֶ֖חֶם לָשֹׂ֑בַע כִּֽי־הוֹצֵאתֶ֤ם אֹתָ֙נוּ֙ אֶל־הַמִּדְבָּ֣ר הַזֶּ֔ה לְהָמִ֛ית אֶת־כׇּל־הַקָּהָ֥ל הַזֶּ֖ה בָּרָעָֽב׃ \{ס\} \Rashi{\textbf{מי יתן מותנו.}שנמות; ואינו שם דבר כמו מיתתנו, אלא כמו עשותנו חנותנו שובנו, לעשות אנחנו, לחנות אנחנו, למות אנחנו. תרגומו לוי דמיתנא, "לו מתנו" (במדבר יד, ב) — הלואי והיינו מתים:}%
% ---------- פרק 16 | פסוק 4 ----------
 \textbf{ד} וַיֹּ֤אמֶר יְהֹוָה֙ אֶל־מֹשֶׁ֔ה הִנְנִ֨י מַמְטִ֥יר לָכֶ֛ם לֶ֖חֶם מִן־הַשָּׁמָ֑יִם וְיָצָ֨א הָעָ֤ם וְלָֽקְטוּ֙ דְּבַר־י֣וֹם בְּיוֹמ֔וֹ לְמַ֧עַן אֲנַסֶּ֛נּוּ הֲיֵלֵ֥ךְ בְּתוֹרָתִ֖י אִם־לֹֽא׃ \Rashi{\textbf{דבר יום ביומו.}צרך אכילת יום ילקטו ביומו ולא היום לצרך מחר (מכילתא): למען אנסנו הילך בתורתי. אם ישמרו מצוות התלויות בו, שלא יותירו ממנו ולא יצאו בשבת ללקט:}%
% ---------- פרק 16 | פסוק 5 ----------
 \textbf{ה} וְהָיָה֙ בַּיּ֣וֹם הַשִּׁשִּׁ֔י וְהֵכִ֖ינוּ אֵ֣ת אֲשֶׁר־יָבִ֑יאוּ וְהָיָ֣ה מִשְׁנֶ֔ה עַ֥ל אֲשֶֽׁר־יִלְקְט֖וּ י֥וֹם ׀ יֽוֹם׃ \Rashi{\textbf{והיה משנה.}ליום ולמחרת: משנה. על שהיו רגילים ללקט יום יום של שאר ימות השבוע; אומר אני "אשר יביאו והיה משנה" – לאחר שיביאו, ימצאו משנה במדידה על אשר ילקטו וימדו יום יום, וזהו לקטו לחם משנה – בלקיטתו היה נמצא לחם משנה. וזהו על כן הוא נותן לכם ביום הששי לחם יומים – נותן לכם ברכה, פוי"שן, בבית למלאת העמר פעמים ללחם יומים:}%
% ---------- פרק 16 | פסוק 6 ----------
 \textbf{ו} וַיֹּ֤אמֶר מֹשֶׁה֙ וְאַהֲרֹ֔ן אֶֽל־כׇּל־בְּנֵ֖י יִשְׂרָאֵ֑ל עֶ֕רֶב וִֽידַעְתֶּ֕ם כִּ֧י יְהֹוָ֛ה הוֹצִ֥יא אֶתְכֶ֖ם מֵאֶ֥רֶץ מִצְרָֽיִם׃ \Rashi{\textbf{ערב.}כמו לערב: וידעתם כי ה' הוציא אתכם מארץ מצרים. לפי שאמרתם לנו כי הוצאתם אותנו, תדעו כי לא אנחנו המוציאים אלא ה' הוציא אתכם שיגיז לכם את השלו:}%
% ---------- פרק 16 | פסוק 7 ----------
 \textbf{ז} וּבֹ֗קֶר וּרְאִיתֶם֙ אֶת־כְּב֣וֹד יְהֹוָ֔ה בְּשׇׁמְע֥וֹ אֶת־תְּלֻנֹּתֵיכֶ֖ם עַל־יְהֹוָ֑ה וְנַ֣חְנוּ מָ֔ה כִּ֥י (תלונו) [תַלִּ֖ינוּ] עָלֵֽינוּ׃ \Rashi{\textbf{ובקר וראיתם.}לא על הכבוד שנאמר "והנה כבוד ה' נראה בענן" נאמר, אלא כך אמר להם, ערב וידעתם כי היכלת בידו לתן תאותכם, ובשר יתן, אך לא בפנים מאירות יתננה לכם, כי שלא כהגן שאלתם אותו ומכרס מלאה, והלחם ששאלתם לצרך, בירידתו לבקר תראו את כבוד אור פניו – שיורידהו לכם דרך חבה, בבקר, שיש שעות להכינו וטל מלמעלה וטל מלמטה, כמנח בקפסא (יומא ע"ה): את תלנתיכם על ה'. כמו אשר על ה': ונחנו מה. מה אנחנו חשובין? כי תלינו עלינו. שתרעימו עלינו את הכל, ואת בניכם ונשיכם ובנותיכם וערב רב. ועל כרחי אני זקוק לפרש תלינו בלשון תפעילו מפני דגשותו וקריתו, שאלו היה רפה, הייתי מפרשו בלשון תפעלו, כמו "וילן העם על משה" (שמות י"ז), או אם היה דגוש ואין בו יו"ד ונקרא תלונו, הייתי מפרשו לשון תתלוננו, עכשיו הוא משמע תלינו את אחרים, כמו במרגלים "וילינו עליו את כל העדה" (במדבר י"ד):}%
% ---------- פרק 16 | פסוק 8 ----------
 \textbf{ח} וַיֹּ֣אמֶר מֹשֶׁ֗ה בְּתֵ֣ת יְהֹוָה֩ לָכֶ֨ם בָּעֶ֜רֶב בָּשָׂ֣ר לֶאֱכֹ֗ל וְלֶ֤חֶם בַּבֹּ֙קֶר֙ לִשְׂבֹּ֔עַ בִּשְׁמֹ֤עַ יְהֹוָה֙ אֶת־תְּלֻנֹּ֣תֵיכֶ֔ם אֲשֶׁר־אַתֶּ֥ם מַלִּינִ֖ם עָלָ֑יו וְנַ֣חְנוּ מָ֔ה לֹא־עָלֵ֥ינוּ תְלֻנֹּתֵיכֶ֖ם כִּ֥י עַל־יְהֹוָֽה׃ \Rashi{\textbf{בשר לאכל.}ולא לשבע; למדה תורה דרך ארץ שאין אוכלין בשר לשבע. ומה ראה להוריד לחם בבקר ובשר בערב? לפי שהלחם שאלו כהגן, שאי אפשר לו לאדם בלא לחם, אבל בשר שאלו שלא כהגן, שהרבה בהמות היו להם, ועוד שהיה אפשר להם בלא בשר, לפיכך נתן להם בשעת טרח שלא כהגן (יומא ע"ה): אשר אתם מלינם עליו. את האחרים השומעים אתכם מתלוננים:}%
% ---------- פרק 16 | פסוק 9 ----------
 \textbf{ט} וַיֹּ֤אמֶר מֹשֶׁה֙ אֶֽל־אַהֲרֹ֔ן אֱמֹ֗ר אֶֽל־כׇּל־עֲדַת֙ בְּנֵ֣י יִשְׂרָאֵ֔ל קִרְב֖וּ לִפְנֵ֣י יְהֹוָ֑ה כִּ֣י שָׁמַ֔ע אֵ֖ת תְּלֻנֹּתֵיכֶֽם׃ \Rashi{\textbf{קרבו.}למקום שהענן ירד:}%
% ---------- פרק 16 | פסוק 10 ----------
 \textbf{י} וַיְהִ֗י כְּדַבֵּ֤ר אַהֲרֹן֙ אֶל־כׇּל־עֲדַ֣ת בְּנֵֽי־יִשְׂרָאֵ֔ל וַיִּפְנ֖וּ אֶל־הַמִּדְבָּ֑ר וְהִנֵּה֙ כְּב֣וֹד יְהֹוָ֔ה נִרְאָ֖ה בֶּעָנָֽן׃ \{פ\} %
% ---------- פרק 16 | פסוק 11 ----------
 \textbf{יא} וַיְדַבֵּ֥ר יְהֹוָ֖ה אֶל־מֹשֶׁ֥ה לֵּאמֹֽר׃ %
% ---------- פרק 16 | פסוק 12 ----------
 \textbf{יב} שָׁמַ֗עְתִּי אֶת־תְּלוּנֹּת֮ בְּנֵ֣י יִשְׂרָאֵל֒ דַּבֵּ֨ר אֲלֵהֶ֜ם לֵאמֹ֗ר בֵּ֤ין הָֽעַרְבַּ֙יִם֙ תֹּאכְל֣וּ בָשָׂ֔ר וּבַבֹּ֖קֶר תִּשְׂבְּעוּ־לָ֑חֶם וִֽידַעְתֶּ֕ם כִּ֛י אֲנִ֥י יְהֹוָ֖ה אֱלֹהֵיכֶֽם׃ %
% ---------- פרק 16 | פסוק 13 ----------
 \textbf{יג} וַיְהִ֣י בָעֶ֔רֶב וַתַּ֣עַל הַשְּׂלָ֔ו וַתְּכַ֖ס אֶת־הַֽמַּחֲנֶ֑ה וּבַבֹּ֗קֶר הָֽיְתָה֙ שִׁכְבַ֣ת הַטַּ֔ל סָבִ֖יב לַֽמַּחֲנֶֽה׃ \Rashi{\textbf{השלו.}מין עוף ושמן מאד (יומא ע"ה): היתה שכבת הטל. הטל שוכב על המן, ובמקום אחר הוא אומר (במדבר י"א) "וברדת הטל וגו'", הטל יורד על הארץ, והמן יורד עליו, וחוזר ויורד טל עליו, הרי הוא כמנח בקפסא (יומא שם):}%
% ---------- פרק 16 | פסוק 14 ----------
 \textbf{יד} וַתַּ֖עַל שִׁכְבַ֣ת הַטָּ֑ל וְהִנֵּ֞ה עַל־פְּנֵ֤י הַמִּדְבָּר֙ דַּ֣ק מְחֻסְפָּ֔ס דַּ֥ק כַּכְּפֹ֖ר עַל־הָאָֽרֶץ׃ \Rashi{\textbf{ותעל שכבת הטל וגו'.}כשהחמה זורחת, עולה הטל שעל המן לקראת החמה, כדרך טל עולה לקראת החמה; אף אם תמלא שפופרת של ביצה טל ותסתם את פיה ותניחה בחמה היא עולה מאליה באויר. ורבותינו דרשו, שהטל עולה מן הארץ באויר, וכעלות שכבת הטל נתגלה המן וראו והנה על פני המדבר וגו': דק – דבר דק: מחספס. מגלה, ואין דומה לו במקרא. ויש לומר מחספס לשון חפיסה ודלוסקמא שבלשון משנה, כשנתגלה משכבת הטל ראו שהיה דבר דק מחספס בתוכו בין שתי שכבות הטל. ואנקלוס תרגם "מקלף", לשון מחשף הלבן: ככפר. כפור יילי"דא בלעז, דעדק כגיר, "כאבני גר" (ישעיהו כ"ז), והוא מין צבע שחר, כדאמרינן גבי כסוי הדם "הגיר והזרניך" (חולין פ"ח). "דעדק כגיר כגלידא על ארעא" – דק היה כגיר ושוכב מגלד כקרח על הארץ, וכן פרושו: דק ככפר – שטוח, קלוש ומחבר כגליד. דק טינב"ש בלעז, שהיה מגליד גלד דק מלמעלה; "וכגיר" שתרגם אנקלוס תוספת הוא על לשון העברית ואין לו תבה בפסוק:}%
% ---------- פרק 16 | פסוק 15 ----------
 \textbf{טו} וַיִּרְא֣וּ בְנֵֽי־יִשְׂרָאֵ֗ל וַיֹּ֨אמְר֜וּ אִ֤ישׁ אֶל־אָחִיו֙ מָ֣ן ה֔וּא כִּ֛י לֹ֥א יָדְע֖וּ מַה־ה֑וּא וַיֹּ֤אמֶר מֹשֶׁה֙ אֲלֵהֶ֔ם ה֣וּא הַלֶּ֔חֶם אֲשֶׁ֨ר נָתַ֧ן יְהֹוָ֛ה לָכֶ֖ם לְאׇכְלָֽה׃ \Rashi{\textbf{מן הוא.}הכנת מזון הוא, כמו "וימן להם המלך" (דניאל א'): כי לא ידעו מה הוא. שיקראוהו בשמו:}%
% ---------- פרק 16 | פסוק 16 ----------
 \textbf{טז} זֶ֤ה הַדָּבָר֙ אֲשֶׁ֣ר צִוָּ֣ה יְהֹוָ֔ה לִקְט֣וּ מִמֶּ֔נּוּ אִ֖ישׁ לְפִ֣י אׇכְל֑וֹ עֹ֣מֶר לַגֻּלְגֹּ֗לֶת מִסְפַּר֙ נַפְשֹׁ֣תֵיכֶ֔ם אִ֛ישׁ לַאֲשֶׁ֥ר בְּאׇהֳל֖וֹ תִּקָּֽחוּ׃ \Rashi{\textbf{עמר.}שם מדה: מספר נפשתיכם. כפי מנין נפשות שיש לאיש באהלו תקחו עמר לכל גלגלת:}%
% ---------- פרק 16 | פסוק 17 ----------
 \textbf{יז} וַיַּעֲשׂוּ־כֵ֖ן בְּנֵ֣י יִשְׂרָאֵ֑ל וַֽיִּלְקְט֔וּ הַמַּרְבֶּ֖ה וְהַמַּמְעִֽיט׃ \Rashi{\textbf{המרבה והממעיט.}יש שלקטו הרבה ויש שלקטו מעט, וכשבאו לביתם מדדו בעמר איש איש מה שלקטו, ומצאו שהמרבה ללקט לא העדיף על עמר לגלגלת אשר באהלו, והממעיט ללקט לא מצא חסר מעמר לגלגלת; וזהו נס גדול שנעשה בו:}%
% ---------- פרק 16 | פסוק 18 ----------
 \textbf{יח} וַיָּמֹ֣דּוּ בָעֹ֔מֶר וְלֹ֤א הֶעְדִּיף֙ הַמַּרְבֶּ֔ה וְהַמַּמְעִ֖יט לֹ֣א הֶחְסִ֑יר אִ֥ישׁ לְפִֽי־אׇכְל֖וֹ לָקָֽטוּ׃ %
% ---------- פרק 16 | פסוק 19 ----------
 \textbf{יט} וַיֹּ֥אמֶר מֹשֶׁ֖ה אֲלֵהֶ֑ם אִ֕ישׁ אַל־יוֹתֵ֥ר מִמֶּ֖נּוּ עַד־בֹּֽקֶר׃ %
% ---------- פרק 16 | פסוק 20 ----------
 \textbf{כ} וְלֹא־שָׁמְע֣וּ אֶל־מֹשֶׁ֗ה וַיּוֹתִ֨רוּ אֲנָשִׁ֤ים מִמֶּ֙נּוּ֙ עַד־בֹּ֔קֶר וַיָּ֥רֻם תּוֹלָעִ֖ים וַיִּבְאַ֑שׁ וַיִּקְצֹ֥ף עֲלֵהֶ֖ם מֹשֶֽׁה׃ \Rashi{\textbf{ויותרו אנשים.}דתן ואבירם: וירם תולעים. לשון רמה: ויבאש. הרי זה מקרא הפוך, שבתחלה הבאיש ולבסוף התליע, כענין שנאמר "ולא הבאיש ורמה לא היתה בו", וכן דרך כל המתליעים (מכילתא):}%
% ---------- פרק 16 | פסוק 21 ----------
 \textbf{כא} וַיִּלְקְט֤וּ אֹתוֹ֙ בַּבֹּ֣קֶר בַּבֹּ֔קֶר אִ֖ישׁ כְּפִ֣י אׇכְל֑וֹ וְחַ֥ם הַשֶּׁ֖מֶשׁ וְנָמָֽס׃ \Rashi{\textbf{וחם השמש ונמס.}הנשאר בשדה נעשה נחלים ושותין ממנו אילים וצבאים, ואמות העולם צדין מהם וטועמים בהם טעם מן, ויודעים מה שבחן של ישראל: ונמס. פשר, לשון פושרים, על ידי השמש מתחמם ומפשיר, דישטנ"פריר בלעז, ודגמתו בסנהדרין בסוף ד' מיתות:}%
% ---------- פרק 16 | פסוק 22 ----------
 \textbf{כב} וַיְהִ֣י ׀ בַּיּ֣וֹם הַשִּׁשִּׁ֗י לָֽקְט֥וּ לֶ֙חֶם֙ מִשְׁנֶ֔ה שְׁנֵ֥י הָעֹ֖מֶר לָאֶחָ֑ד וַיָּבֹ֙אוּ֙ כׇּל־נְשִׂיאֵ֣י הָֽעֵדָ֔ה וַיַּגִּ֖ידוּ לְמֹשֶֽׁה׃ \Rashi{\textbf{לקטו לחם משנה.}כשמדדו את לקיטתם באהליהם מצאו כפלים, שני העמר לאחד. ומדרש אגדה: לחם משנה – משנה, אותו היום נשתנה לשבח בריחו וטעמו (שם): ויגידו למשה. שאלוהו מה היום מיומים; ומכאן יש ללמד שעדין לא הגיד להם משה פרשת שבת שנצטוה לומר להם והיה ביום הששי והכינו וגו', עד ששאלו מה זאת, אמר להם הוא אשר דבר ה' שנצטויתי לומר לכם, ולכך ענשו הכתוב שאמר לו עד אנה מאנתם, ולא הוציאו מן הכלל:}%
% ---------- פרק 16 | פסוק 23 ----------
 \textbf{כג} וַיֹּ֣אמֶר אֲלֵהֶ֗ם ה֚וּא אֲשֶׁ֣ר דִּבֶּ֣ר יְהֹוָ֔ה שַׁבָּת֧וֹן שַׁבַּת־קֹ֛דֶשׁ לַֽיהֹוָ֖ה מָחָ֑ר אֵ֣ת אֲשֶׁר־תֹּאפ֞וּ אֵפ֗וּ וְאֵ֤ת אֲשֶֽׁר־תְּבַשְּׁלוּ֙ בַּשֵּׁ֔לוּ וְאֵת֙ כׇּל־הָ֣עֹדֵ֔ף הַנִּ֧יחוּ לָכֶ֛ם לְמִשְׁמֶ֖רֶת עַד־הַבֹּֽקֶר׃ \Rashi{\textbf{את אשר תאפו אפו.}מה שאתם רוצים לאפות בתנור, אפו היום הכל לשני ימים, ומה שאתם צריכים לבשל ממנו במים, בשלו היום. לשון אפיה נופל בלחם ולשון בשול בתבשיל: למשמרת. לגניזה:}%
% ---------- פרק 16 | פסוק 24 ----------
 \textbf{כד} וַיַּנִּ֤יחוּ אֹתוֹ֙ עַד־הַבֹּ֔קֶר כַּאֲשֶׁ֖ר צִוָּ֣ה מֹשֶׁ֑ה וְלֹ֣א הִבְאִ֔ישׁ וְרִמָּ֖ה לֹא־הָ֥יְתָה בּֽוֹ׃ %
% ---------- פרק 16 | פסוק 25 ----------
 \textbf{כה} וַיֹּ֤אמֶר מֹשֶׁה֙ אִכְלֻ֣הוּ הַיּ֔וֹם כִּֽי־שַׁבָּ֥ת הַיּ֖וֹם לַיהֹוָ֑ה הַיּ֕וֹם לֹ֥א תִמְצָאֻ֖הוּ בַּשָּׂדֶֽה׃ \Rashi{\textbf{ויאמר משה אכלהו היום וגו'.}שחרית שהיו רגילין לצאת וללקט באו לשאל: "אם נצא אם לאו", אמר להם את שבידכם אכלו, לערב חזרו לפניו ושאלוהו מהו לצאת? אמר להם שבת היום, ראה אותם דואגים שמא פסק המן ולא ירד עוד, אמר להם היום לא תמצאוהו, מה ת"ל היום? היום לא תמצאוהו אבל מחר תמצאוהו:}%
% ---------- פרק 16 | פסוק 26 ----------
 \textbf{כו} שֵׁ֥שֶׁת יָמִ֖ים תִּלְקְטֻ֑הוּ וּבַיּ֧וֹם הַשְּׁבִיעִ֛י שַׁבָּ֖ת לֹ֥א יִֽהְיֶה־בּֽוֹ׃ \Rashi{\textbf{וביום השביעי שבת.}שבת הוא, המן לא יהיה בו; ולא בא הכתוב אלא לרבות יום הכפורים וימים טובים:}%
% ---------- פרק 16 | פסוק 27 ----------
 \textbf{כז} וַֽיְהִי֙ בַּיּ֣וֹם הַשְּׁבִיעִ֔י יָצְא֥וּ מִן־הָעָ֖ם לִלְקֹ֑ט וְלֹ֖א מָצָֽאוּ׃ \{ס\} %
% ---------- פרק 16 | פסוק 28 ----------
 \textbf{כח} וַיֹּ֥אמֶר יְהֹוָ֖ה אֶל־מֹשֶׁ֑ה עַד־אָ֙נָה֙ מֵֽאַנְתֶּ֔ם לִשְׁמֹ֥ר מִצְוֺתַ֖י וְתוֹרֹתָֽי׃ \Rashi{\textbf{עד אנה מאנתם.}משל הדיוט הוא "בהדי הוצא לקי כרבא" – על ידי הרשעים מתגנין הכשרין:}%
% ---------- פרק 16 | פסוק 29 ----------
 \textbf{כט} רְא֗וּ כִּֽי־יְהֹוָה֮ נָתַ֣ן לָכֶ֣ם הַשַּׁבָּת֒ עַל־כֵּ֠ן ה֣וּא נֹתֵ֥ן לָכֶ֛ם בַּיּ֥וֹם הַשִּׁשִּׁ֖י לֶ֣חֶם יוֹמָ֑יִם שְׁב֣וּ ׀ אִ֣ישׁ תַּחְתָּ֗יו אַל־יֵ֥צֵא אִ֛ישׁ מִמְּקֹמ֖וֹ בַּיּ֥וֹם הַשְּׁבִיעִֽי׃ \Rashi{\textbf{ראו.}בעיניכם כי ה' בכבודו מזהיר אתכם על השבת, שהרי נס נעשה בכל ערב שבת לתת לכם לחם יומים: שבו איש תחתיו. מכאן סמכו חכמים ד' אמות ליוצא חוץ לתחום (ערובין נ"א): אל יצא איש ממקמו. אלו אלפים אמה של תחום שבת (מכילתא), ולא במפרש – שאין תחומין אלא מדברי סופרים – ועקרו של מקרא על לוקטי המן נאמר:}%
% ---------- פרק 16 | פסוק 30 ----------
 \textbf{ל} וַיִּשְׁבְּת֥וּ הָעָ֖ם בַּיּ֥וֹם הַשְּׁבִעִֽי׃ %
% ---------- פרק 16 | פסוק 31 ----------
 \textbf{לא} וַיִּקְרְא֧וּ בֵֽית־יִשְׂרָאֵ֛ל אֶת־שְׁמ֖וֹ מָ֑ן וְה֗וּא כְּזֶ֤רַע גַּד֙ לָבָ֔ן וְטַעְמ֖וֹ כְּצַפִּיחִ֥ת בִּדְבָֽשׁ׃ \Rashi{\textbf{והוא כזרע גד לבן.}עשב ששמו קול"יינדר, וזרע שלו עגל ואינו לבן, והמן היה לבן, ואינו נמשל לזרע גד אלא לענין העגול, כזרע גד היה – והוא לבן: כצפיחת. בצק שמטגנין אותו בדבש, וקורין לו אסקריטין בלשון משנה (פסחים ל"ז), והוא תרגום של אנקלוס:}%
% ---------- פרק 16 | פסוק 32 ----------
 \textbf{לב} וַיֹּ֣אמֶר מֹשֶׁ֗ה זֶ֤ה הַדָּבָר֙ אֲשֶׁ֣ר צִוָּ֣ה יְהֹוָ֔ה מְלֹ֤א הָעֹ֙מֶר֙ מִמֶּ֔נּוּ לְמִשְׁמֶ֖רֶת לְדֹרֹתֵיכֶ֑ם לְמַ֣עַן ׀ יִרְא֣וּ אֶת־הַלֶּ֗חֶם אֲשֶׁ֨ר הֶאֱכַ֤לְתִּי אֶתְכֶם֙ בַּמִּדְבָּ֔ר בְּהוֹצִיאִ֥י אֶתְכֶ֖ם מֵאֶ֥רֶץ מִצְרָֽיִם׃ \Rashi{\textbf{למשמרת.}לגניזה: לדרתיכם. בימי ירמיהו; כשהיה ירמיהו מוכיחם למה אין אתם עוסקים בתורה? והם אומרים נניח מלאכתנו ונעסק בתורה, מהיכן נתפרנס? הוציא להם צנצנת המן אמר להם אתם ראו דבר ה', שמעו לא נאמר אלא ראו, בזה נתפרנסו אבותיכם, הרבה שלוחין יש לו למקום להכין מזון ליראיו:}%
% ---------- פרק 16 | פסוק 33 ----------
 \textbf{לג} וַיֹּ֨אמֶר מֹשֶׁ֜ה אֶֽל־אַהֲרֹ֗ן קַ֚ח צִנְצֶ֣נֶת אַחַ֔ת וְתֶן־שָׁ֥מָּה מְלֹֽא־הָעֹ֖מֶר מָ֑ן וְהַנַּ֤ח אֹתוֹ֙ לִפְנֵ֣י יְהֹוָ֔ה לְמִשְׁמֶ֖רֶת לְדֹרֹתֵיכֶֽם׃ \Rashi{\textbf{צנצנת.}צלוחית של חרס, כתרגומו: והנח אותו לפני ה'. לפני הארון; ולא נאמר מקרא זה עד שנבנה אהל מועד, אלא שנכתב כאן בפרשת המן:}%
% ---------- פרק 16 | פסוק 34 ----------
 \textbf{לד} כַּאֲשֶׁ֛ר צִוָּ֥ה יְהֹוָ֖ה אֶל־מֹשֶׁ֑ה וַיַּנִּיחֵ֧הוּ אַהֲרֹ֛ן לִפְנֵ֥י הָעֵדֻ֖ת לְמִשְׁמָֽרֶת׃ %
% ---------- פרק 16 | פסוק 35 ----------
 \textbf{לה} וּבְנֵ֣י יִשְׂרָאֵ֗ל אָֽכְל֤וּ אֶת־הַמָּן֙ אַרְבָּעִ֣ים שָׁנָ֔ה עַד־בֹּאָ֖ם אֶל־אֶ֣רֶץ נוֹשָׁ֑בֶת אֶת־הַמָּן֙ אָֽכְל֔וּ עַד־בֹּאָ֕ם אֶל־קְצֵ֖ה אֶ֥רֶץ כְּנָֽעַן׃ \Rashi{\textbf{ארבעים שנה.}והלא חסר שלשים יום הם? שהרי בט"ו באיר ירד להם המן תחלה ובט"ו בניסן פסק, שנאמר "וישבת המן ממחרת" (יהושע ה')? אלא מגיד, שהעגות שהוציאו ישראל ממצרים טעמו בהם טעם מן (קידושין ל"ח): אל ארץ נושבת. לאחר שעברו את הירדן (שאותה שבעבר הירדן מישבת וטובה שנאמר "אעברה נא ואראה את הארץ הטובה אשר בעבר הירדן", ותרגום של נושבת "יתבתא", רצה לומר — מישבת): אל קצה ארץ כנען. בתחלת הגבול קדם שעברו את הירדן, והוא ערבות מואב. נמצאו מכחישין זה את זה? אלא בערבות מואב כשמת משה בז' באדר פסק המן מלירד, ונסתפקו ממן שלקטו בו ביום עד שהקריבו העמר בששה עשר בניסן, שנאמר "ויאכלו מעבור הארץ ממחרת הפסח" (יהושע ה'):}%
% ---------- פרק 16 | פסוק 36 ----------
 \textbf{לו} וְהָעֹ֕מֶר עֲשִׂרִ֥ית הָאֵיפָ֖ה הֽוּא׃ \{פ\} \Rashi{\textbf{עשרית האיפה.}האיפה שלש סאין והסאה ו' קבין והקב ד' לגין והלג ו' ביצים, נמצא עשירית האיפה מ"ג ביצים וחמש ביצה, והוא שעור לחלה ולמנחות:}%
% ---------- פרק 17 | פסוק 1 ----------
 \textbf{א} וַ֠יִּסְע֠וּ כׇּל־עֲדַ֨ת בְּנֵֽי־יִשְׂרָאֵ֧ל מִמִּדְבַּר־סִ֛ין לְמַסְעֵיהֶ֖ם עַל־פִּ֣י יְהֹוָ֑ה וַֽיַּחֲנוּ֙ בִּרְפִידִ֔ים וְאֵ֥ין מַ֖יִם לִשְׁתֹּ֥ת הָעָֽם׃ %
% ---------- פרק 17 | פסוק 2 ----------
 \textbf{ב} וַיָּ֤רֶב הָעָם֙ עִם־מֹשֶׁ֔ה וַיֹּ֣אמְר֔וּ תְּנוּ־לָ֥נוּ מַ֖יִם וְנִשְׁתֶּ֑ה וַיֹּ֤אמֶר לָהֶם֙ מֹשֶׁ֔ה מַה־תְּרִיבוּן֙ עִמָּדִ֔י מַה־תְּנַסּ֖וּן אֶת־יְהֹוָֽה׃ \Rashi{\textbf{מה תנסון.}לומר היוכל לתת מים בארץ ציה?:}%
% ---------- פרק 17 | פסוק 3 ----------
 \textbf{ג} וַיִּצְמָ֨א שָׁ֤ם הָעָם֙ לַמַּ֔יִם וַיָּ֥לֶן הָעָ֖ם עַל־מֹשֶׁ֑ה וַיֹּ֗אמֶר לָ֤מָּה זֶּה֙ הֶעֱלִיתָ֣נוּ מִמִּצְרַ֔יִם לְהָמִ֥ית אֹתִ֛י וְאֶת־בָּנַ֥י וְאֶת־מִקְנַ֖י בַּצָּמָֽא׃ %
% ---------- פרק 17 | פסוק 4 ----------
 \textbf{ד} וַיִּצְעַ֤ק מֹשֶׁה֙ אֶל־יְהֹוָ֣ה לֵאמֹ֔ר מָ֥ה אֶעֱשֶׂ֖ה לָעָ֣ם הַזֶּ֑ה ע֥וֹד מְעַ֖ט וּסְקָלֻֽנִי׃ \Rashi{\textbf{עוד מעט.}אם אמתין עוד מעט וסקלני:}%
% ---------- פרק 17 | פסוק 5 ----------
 \textbf{ה} וַיֹּ֨אמֶר יְהֹוָ֜ה אֶל־מֹשֶׁ֗ה עֲבֹר֙ לִפְנֵ֣י הָעָ֔ם וְקַ֥ח אִתְּךָ֖ מִזִּקְנֵ֣י יִשְׂרָאֵ֑ל וּמַטְּךָ֗ אֲשֶׁ֨ר הִכִּ֤יתָ בּוֹ֙ אֶת־הַיְאֹ֔ר קַ֥ח בְּיָדְךָ֖ וְהָלָֽכְתָּ׃ \Rashi{\textbf{עבור לפני העם.}וראה אם יסקלוך, למה הוצאת לעז על בני? (שמות רבה): וקח אתך מזקני ישראל. לעדות, שיראו שעל ידך מים יוצאים מן הצור, ולא יאמרו מעינות היו שם מימי קדם: ומטך אשר הכית בו את היאר. מה ת"ל "אשר הכית בו את היאר"? אלא שהיו ישראל אומרים על המטה שאינו מוכן אלא לפרענות – בו לקה פרעה ומצרים כמה מכות במצרים ועל הים – לכך נאמר אשר הכית בו את היאר, יראו עתה שאף לטובה הוא מוכן (מכילתא):}%
% ---------- פרק 17 | פסוק 6 ----------
 \textbf{ו} הִנְנִ֣י עֹמֵד֩ לְפָנֶ֨יךָ שָּׁ֥ם ׀ עַֽל־הַצּוּר֮ בְּחֹרֵב֒ וְהִכִּ֣יתָ בַצּ֗וּר וְיָצְא֥וּ מִמֶּ֛נּוּ מַ֖יִם וְשָׁתָ֣ה הָעָ֑ם וַיַּ֤עַשׂ כֵּן֙ מֹשֶׁ֔ה לְעֵינֵ֖י זִקְנֵ֥י יִשְׂרָאֵֽל׃ \Rashi{\textbf{והכית בצור.}על הצור לא נאמר אלא בצור, מכאן שהמטה היה של מין דבר חזק ושמו סנפירינון, והצור נבקע מפניו (שם):}%
% ---------- פרק 17 | פסוק 7 ----------
 \textbf{ז} וַיִּקְרָא֙ שֵׁ֣ם הַמָּק֔וֹם מַסָּ֖ה וּמְרִיבָ֑ה עַל־רִ֣יב ׀ בְּנֵ֣י יִשְׂרָאֵ֗ל וְעַ֨ל נַסֹּתָ֤ם אֶת־יְהֹוָה֙ לֵאמֹ֔ר הֲיֵ֧שׁ יְהֹוָ֛ה בְּקִרְבֵּ֖נוּ אִם־אָֽיִן׃ \{פ\} %
% ---------- פרק 17 | פסוק 8 ----------
 \textbf{ח} וַיָּבֹ֖א עֲמָלֵ֑ק וַיִּלָּ֥חֶם עִם־יִשְׂרָאֵ֖ל בִּרְפִידִֽם׃ \Rashi{\textbf{ויבא עמלק וגו'.}סמך פרשה זו למקרא זה, לומר, תמיד אני ביניכם ומזמן לכל צרכיכם ואתם אומרים היש ה' בקרבנו אם אין, חייכם שהכלב בא ונושך אתכם ואתם צועקים לי ותדעו היכן אני? משל לאדם שהרכיב בנו על כתפו ויצא לדרך, היה אותו הבן רואה חפץ ואומר, אבא טל חפץ זה ותן לי והוא נותן לו, וכן שניה, וכן שלישית, פגעו באדם אחד אמר לו אותו הבן, ראית את אבא? אמר לו אביו, אינך יודע היכן אני? השליכו מעליו ובא הכלב ונשכו (תנחומא יתרו):}%
% ---------- פרק 17 | פסוק 9 ----------
 \textbf{ט} וַיֹּ֨אמֶר מֹשֶׁ֤ה אֶל־יְהוֹשֻׁ֙עַ֙ בְּחַר־לָ֣נוּ אֲנָשִׁ֔ים וְצֵ֖א הִלָּחֵ֣ם בַּעֲמָלֵ֑ק מָחָ֗ר אָנֹכִ֤י נִצָּב֙ עַל־רֹ֣אשׁ הַגִּבְעָ֔ה וּמַטֵּ֥ה הָאֱלֹהִ֖ים בְּיָדִֽי׃ \Rashi{\textbf{בחר לנו.}לי ולך – השוהו לו; מכאן אמרו חכמים יהי כבוד תלמידך חביב עליך כשלך; וכבוד חברך כמורא רבך מנין? שנאמר "ויאמר אהרן אל משה בי אדני", והלא אהרן גדול מאחיו היה ועושה את חברו כרבו. ומורא רבך כמורא שמים מנין? שנאמר "אדני משה כלאם" – כלם מן העולם, חיבין הם כליה, המורדים בך כאלו מרדו בהקב"ה: וצא הלחם. צא מן הענן והלחם בו: מחר. בעת המלחמה, אנכי נצב: בחר לנו אנשים. גבורים ויראי חטא שתהא זכותן מסיעתן. דבר אחר – בחר לנו אנשים שיודעין לבטל כשפים, לפי שבני עמלק מכשפין היו:}%
% ---------- פרק 17 | פסוק 10 ----------
 \textbf{י} וַיַּ֣עַשׂ יְהוֹשֻׁ֗עַ כַּאֲשֶׁ֤ר אָֽמַר־לוֹ֙ מֹשֶׁ֔ה לְהִלָּחֵ֖ם בַּעֲמָלֵ֑ק וּמֹשֶׁה֙ אַהֲרֹ֣ן וְח֔וּר עָל֖וּ רֹ֥אשׁ הַגִּבְעָֽה׃ \Rashi{\textbf{ומשה אהרן וחור.}מכאן לתענית שצריכים שלשה לעבר לפני התבה, שבתענית היו שרויים (מכילתא): חור. בנה של מרים היה וכלב בעלה (סוטה י"א):}%
% ---------- פרק 17 | פסוק 11 ----------
 \textbf{יא} וְהָיָ֗ה כַּאֲשֶׁ֨ר יָרִ֥ים מֹשֶׁ֛ה יָד֖וֹ וְגָבַ֣ר יִשְׂרָאֵ֑ל וְכַאֲשֶׁ֥ר יָנִ֛יחַ יָד֖וֹ וְגָבַ֥ר עֲמָלֵֽק׃ \Rashi{\textbf{כאשר ירים משה ידו.}וכי ידיו של משה נוצחות היו המלחמה וכו' כדאיתא בר"ה:}%
% ---------- פרק 17 | פסוק 12 ----------
 \textbf{יב} וִידֵ֤י מֹשֶׁה֙ כְּבֵדִ֔ים וַיִּקְחוּ־אֶ֛בֶן וַיָּשִׂ֥ימוּ תַחְתָּ֖יו וַיֵּ֣שֶׁב עָלֶ֑יהָ וְאַהֲרֹ֨ן וְח֜וּר תָּֽמְכ֣וּ בְיָדָ֗יו מִזֶּ֤ה אֶחָד֙ וּמִזֶּ֣ה אֶחָ֔ד וַיְהִ֥י יָדָ֛יו אֱמוּנָ֖ה עַד־בֹּ֥א הַשָּֽׁמֶשׁ׃ \Rashi{\textbf{וידי משה כבדים.}בשביל שנתעצל במצוה ומנה אחר תחתיו, נתיקרו ידיו: ויקחו. אהרן וחור: אבן וישימו תחתיו. ולא ישב לו על כר וכסת. אמר: ישראל שרויין בצער, אף אני אהיה עמהם בצער (תענית י"א): ויהי ידיו אמונה. ויהי משה ידיו באמונה, פרושות השמים בתפלה נאמנה ונכונה: עד בא השמש. שהיו עמלקים מחשבין את השעות באיצטרולוגיאה, באיזו שעה הם נוצחים, והעמיד להם משה חמה וערבב את השעות (תנחומא):}%
% ---------- פרק 17 | פסוק 13 ----------
 \textbf{יג} וַיַּחֲלֹ֧שׁ יְהוֹשֻׁ֛עַ אֶת־עֲמָלֵ֥ק וְאֶת־עַמּ֖וֹ לְפִי־חָֽרֶב׃ \{פ\} \Rashi{\textbf{ויחלש יהושע.}חתך ראשי גבוריו ולא השאיר אלא חלשים שבהם, ולא הרגם כלם, מכאן אנו למדים שעשו ע"פ הדבור של שכינה (שם):}%
% ---------- פרק 17 | פסוק 14 ----------
 \textbf{יד} וַיֹּ֨אמֶר יְהֹוָ֜ה אֶל־מֹשֶׁ֗ה כְּתֹ֨ב זֹ֤את זִכָּרוֹן֙ בַּסֵּ֔פֶר וְשִׂ֖ים בְּאׇזְנֵ֣י יְהוֹשֻׁ֑עַ כִּֽי־מָחֹ֤ה אֶמְחֶה֙ אֶת־זֵ֣כֶר עֲמָלֵ֔ק מִתַּ֖חַת הַשָּׁמָֽיִם׃ \Rashi{\textbf{כתב זאת זכרון.}שבא עמלק להזדוג לישראל קדם לכל האמות: ושים באזני יהושע. המכניס את ישראל לארץ, שיצוה את ישראל לשלם לו את גמולו; כאן נרמז לו למשה שיהושע מכניס את ישראל לארץ (שם): כי מחה אמחה. לכך אני מזהירך כן, כי חפץ אני למחותו:}%
% ---------- פרק 17 | פסוק 15 ----------
 \textbf{טו} וַיִּ֥בֶן מֹשֶׁ֖ה מִזְבֵּ֑חַ וַיִּקְרָ֥א שְׁמ֖וֹ יְהֹוָ֥ה ׀ נִסִּֽי׃ \Rashi{\textbf{ויקרא שמו.}של מזבח: ה' נסי. הקב"ה עשה לנו כאן נס. לא שהמזבח קרוי ה', אלא המזכיר שמו של מזבח זוכר את הנס שעשה המקום – ה' הוא נס שלנו:}%
% ---------- פרק 17 | פסוק 16 ----------
 \textbf{טז} וַיֹּ֗אמֶר כִּֽי־יָד֙ עַל־כֵּ֣ס יָ֔הּ*(בכתר ארם צובה היה כתוב כֵּ֣סְיָ֔הּ בתיבה אחת) מִלְחָמָ֥ה לַיהֹוָ֖ה בַּֽעֲמָלֵ֑ק מִדֹּ֖ר דֹּֽר׃ \{פ\} \Rashi{\textbf{ויאמר.}משה: כי יד על כס יה. ידו של הקב"ה הורמה לשבע בכסאו להיות לו מלחמה ואיבה בעמלק עולמית, ומהו כס ולא נאמר כסא? ואף השם נחלק לחציו? נשבע הקב"ה שאין שמו שלם ואין כסאו שלם עד שימחה שמו של עמלק כלו, וכשימחה שמו יהיה השם שלם והכסא שלם, שנאמר "האויב תמו חרבות לנצח" (תהלים ט'), זהו עמלק שכתוב בו "ועברתו שמרה נצח" (עמוס א'), "וערים נתשת אבד זכרם המה" (תהלים שם), מהו אומר אחריו? "וה' לעולם ישב" — הרי השם שלם, "כונן למשפט כסאו" — הרי כסאו שלם:}%
\addparasha{חומש שמות | פרשת יתרו}
% ---------- פרק 18 | פסוק 1 ----------
 \textbf{א} וַיִּשְׁמַ֞ע יִתְר֨וֹ כֹהֵ֤ן מִדְיָן֙ חֹתֵ֣ן מֹשֶׁ֔ה אֵת֩ כׇּל־אֲשֶׁ֨ר עָשָׂ֤ה אֱלֹהִים֙ לְמֹשֶׁ֔ה וּלְיִשְׂרָאֵ֖ל עַמּ֑וֹ כִּֽי־הוֹצִ֧יא יְהֹוָ֛ה אֶת־יִשְׂרָאֵ֖ל מִמִּצְרָֽיִם׃ \Rashi{\textbf{וישמע יתרו.}מה שמועה שמע ובא? קריעת ים סוף ומלחמת עמלק: יתרו. שבע שמות נקראו לו: רעואל, יתר, יתרו, חובב, חבר, קיני, פוטיאל; יתר, על שם שיתר פרשה אחת בתורה "ואתה תחזה"; יתרו, לכשנתגיר וקים המצוות, הוסיפו לו אות אחד על שמו; חובב, שחבב את התורה; חובב הוא יתרו שנאמר "מבני חובב חתן משה" (שופטים ד'), ויש אומרים רעואל אביו של יתרו, ומהו אומר "ותבאנה אל רעואל אביהן"? שהתינוקות קורין לאבי אביהן אבא. בספרי: חתן משה. כאן היה יתרו מתכבד במשה – אני חותן המלך, ולשעבר היה משה תולה הגדלה בחמיו, שנאמר "וישב אל יתר חותנו" (שמות ד'): למשה ולישראל. שקול משה כנגד כל ישראל (מכילתא): את כל אשר עשה. להם בירידת המן ובבאר ובעמלק: כי הוציא ה' וגו'. זו גדולה על כלם (שם):}%
% ---------- פרק 18 | פסוק 2 ----------
 \textbf{ב} וַיִּקַּ֗ח יִתְרוֹ֙ חֹתֵ֣ן מֹשֶׁ֔ה אֶת־צִפֹּרָ֖ה אֵ֣שֶׁת מֹשֶׁ֑ה אַחַ֖ר שִׁלּוּחֶֽיהָ׃ \Rashi{\textbf{אחר שלוחיה.}כשאמר לו הקב"ה במדין לך שב מצרים ויקח משה את אשתו ואת בניו וגו' ויצא אהרן לקראתו ויפגשהו בהר האלהים (שמות ד'), אמר לו מי הם הללו? אמר לו זו היא אשתי שנשאתי במדין ואלו בני, אמר לו והיכן אתה מוליכן? אמר לו למצרים, אמר לו על הראשונים אנו מצטערים ואתה בא להוסיף עליהם אמר לה לכי לבית אביך, נטלה שני בניה והלכה לה (מכילתא):}%
% ---------- פרק 18 | פסוק 3 ----------
 \textbf{ג} וְאֵ֖ת שְׁנֵ֣י בָנֶ֑יהָ אֲשֶׁ֨ר שֵׁ֤ם הָֽאֶחָד֙ גֵּֽרְשֹׁ֔ם כִּ֣י אָמַ֔ר גֵּ֣ר הָיִ֔יתִי בְּאֶ֖רֶץ נׇכְרִיָּֽה׃ %
% ---------- פרק 18 | פסוק 4 ----------
 \textbf{ד} וְשֵׁ֥ם הָאֶחָ֖ד אֱלִיעֶ֑זֶר כִּֽי־אֱלֹהֵ֤י אָבִי֙ בְּעֶזְרִ֔י וַיַּצִּלֵ֖נִי מֵחֶ֥רֶב פַּרְעֹֽה׃ \Rashi{\textbf{ויצלני מחרב פרעה.}כשגלו דתן ואבירם על דבר המצרי ובקש להרג את משה נעשה צוארו כעמוד של שיש (שם):}%
% ---------- פרק 18 | פסוק 5 ----------
 \textbf{ה} וַיָּבֹ֞א יִתְר֨וֹ חֹתֵ֥ן מֹשֶׁ֛ה וּבָנָ֥יו וְאִשְׁתּ֖וֹ אֶל־מֹשֶׁ֑ה אֶל־הַמִּדְבָּ֗ר אֲשֶׁר־ה֛וּא חֹנֶ֥ה שָׁ֖ם הַ֥ר הָאֱלֹהִֽים׃ \Rashi{\textbf{אל המדבר.}אף אנו יודעים שבמדבר היו אלא בשבחו של יתרו דבר הכתוב, שהיה יושב בכבודו של עולם ונדבו לבו לצאת אל המדבר, מקום תהו לשמע דברי תורה (שם):}%
% ---------- פרק 18 | פסוק 6 ----------
 \textbf{ו} וַיֹּ֙אמֶר֙ אֶל־מֹשֶׁ֔ה אֲנִ֛י חֹתֶנְךָ֥ יִתְר֖וֹ בָּ֣א אֵלֶ֑יךָ וְאִ֨שְׁתְּךָ֔ וּשְׁנֵ֥י בָנֶ֖יהָ עִמָּֽהּ׃ \Rashi{\textbf{ויאמר אל משה.}ע"י שליח (שם): אני חתנך יתרו וגו'. אם אין אתה יוצא בגיני, צא בגין אשתך, ואם אין אתה יוצא בגין אשתך, צא בגין שני בניה (שם):}%
% ---------- פרק 18 | פסוק 7 ----------
 \textbf{ז} וַיֵּצֵ֨א מֹשֶׁ֜ה לִקְרַ֣את חֹֽתְנ֗וֹ וַיִּשְׁתַּ֙חוּ֙ וַיִּשַּׁק־ל֔וֹ וַיִּשְׁאֲל֥וּ אִישׁ־לְרֵעֵ֖הוּ לְשָׁל֑וֹם וַיָּבֹ֖אוּ הָאֹֽהֱלָה׃ \Rashi{\textbf{ויצא משה.}כבוד גדול נתכבד יתרו באותה שעה, כיון שיצא משה, יצא אהרן נדב ואביהוא, ומי הוא שראה את אלו יוצאין ולא יצא? (תנחומא): וישתחו וישק לו. איני יודע מי השתחוה למי, כשהוא אומר איש לרעהו, מי הקרוי "איש" – משה, שנאמר "והאיש משה" (במדבר י"ב):}%
% ---------- פרק 18 | פסוק 8 ----------
 \textbf{ח} וַיְסַפֵּ֤ר מֹשֶׁה֙ לְחֹ֣תְנ֔וֹ אֵת֩ כׇּל־אֲשֶׁ֨ר עָשָׂ֤ה יְהֹוָה֙ לְפַרְעֹ֣ה וּלְמִצְרַ֔יִם עַ֖ל אוֹדֹ֣ת יִשְׂרָאֵ֑ל אֵ֤ת כׇּל־הַתְּלָאָה֙ אֲשֶׁ֣ר מְצָאָ֣תַם בַּדֶּ֔רֶךְ וַיַּצִּלֵ֖ם יְהֹוָֽה׃ \Rashi{\textbf{ויספר משה לחתנו.}למשך את לבו לקרבו לתורה (מכילתא): את כל התלאה. שעל הים ושל עמלק: התלאה. למ"ד אל"ף מן היסוד של תבה, והתי"ו הוא תקון ויסוד הנופל ממנו לפרקים; וכן תרומה, תנופה, תקומה, תנואה:}%
% ---------- פרק 18 | פסוק 9 ----------
 \textbf{ט} וַיִּ֣חַדְּ יִתְר֔וֹ עַ֚ל כׇּל־הַטּוֹבָ֔ה אֲשֶׁר־עָשָׂ֥ה יְהֹוָ֖ה לְיִשְׂרָאֵ֑ל אֲשֶׁ֥ר הִצִּיל֖וֹ מִיַּ֥ד מִצְרָֽיִם׃ \Rashi{\textbf{ויחד יתרו.}וישמח יתרו, זהו פשוטו. ומ"א נעשה בשרו חדודין חדודין, מצר על אבוד מצרים, הינו דאמרי אינשי "גיורא עד עשרה דרי לא תבזי ארמאה באפיה": על כל הטובה. טובת המן והבאר והתורה, ועל כלן אשר הצילו מיד מצרים; עד עכשו לא היה עבד יכול לברח ממצרים, שהיתה הארץ מסגרת, ואלו יצאו ששים רבוא (מכילתא):}%
% ---------- פרק 18 | פסוק 10 ----------
 \textbf{י} וַיֹּ֘אמֶר֮ יִתְרוֹ֒ בָּר֣וּךְ יְהֹוָ֔ה אֲשֶׁ֨ר הִצִּ֥יל אֶתְכֶ֛ם מִיַּ֥ד מִצְרַ֖יִם וּמִיַּ֣ד פַּרְעֹ֑ה אֲשֶׁ֤ר הִצִּיל֙ אֶת־הָעָ֔ם מִתַּ֖חַת יַד־מִצְרָֽיִם׃ \Rashi{\textbf{אשר הציל אתכם מיד מצרים.}אמה קשה: ומיד פרעה. מלך קשה: מתחת יד מצרים. כתרגומו – לשון רדוי ומרות; היד שהיו מכבידים עליהם היא העבודה:}%
% ---------- פרק 18 | פסוק 11 ----------
 \textbf{יא} עַתָּ֣ה יָדַ֔עְתִּי כִּֽי־גָד֥וֹל יְהֹוָ֖ה מִכׇּל־הָאֱלֹהִ֑ים כִּ֣י בַדָּבָ֔ר אֲשֶׁ֥ר זָד֖וּ עֲלֵיהֶֽם׃ \Rashi{\textbf{עתה ידעתי.}מכירו הייתי לשעבר ועכשו ביותר: מכל האלהים. מלמד שהיה מכיר בכל ע"ז שבעולם, שלא הניח ע"ז שלא עבדה (מכילתא): כי בדבר אשר זדו עליהם. כתרגומו – במים דמו לאבדם והם נאבדו במים: אשר זדו. אשר הרשיעו. ורבותינו דרשוהו לשון ויזד יעקב נזיד, בקדרה אשר בשלו בה נתבשלו (סוטה י"א):}%
% ---------- פרק 18 | פסוק 12 ----------
 \textbf{יב} וַיִּקַּ֞ח יִתְר֨וֹ חֹתֵ֥ן מֹשֶׁ֛ה עֹלָ֥ה וּזְבָחִ֖ים לֵֽאלֹהִ֑ים וַיָּבֹ֨א אַהֲרֹ֜ן וְכֹ֣ל ׀ זִקְנֵ֣י יִשְׂרָאֵ֗ל לֶאֱכׇל־לֶ֛חֶם עִם־חֹתֵ֥ן מֹשֶׁ֖ה לִפְנֵ֥י הָאֱלֹהִֽים׃ \Rashi{\textbf{עלה.}כמשמעה, שהיא כלה כליל: וזבחים. שלמים: ויבא אהרן וגו'. ומשה היכן הלך? והלא הוא שיצא לקראתו וגרם לו את כל הכבוד? אלא שהיה עומד ומשמש לפניהם (מכילתא): לפני האלהים. מכאן שהנהנה מסעודה שתלמידי חכמים מסבין בה, כאלו נהנה מזיו השכינה (ברכות ס"ד):}%
% ---------- פרק 18 | פסוק 13 ----------
 \textbf{יג} וַֽיְהִי֙ מִֽמׇּחֳרָ֔ת וַיֵּ֥שֶׁב מֹשֶׁ֖ה לִשְׁפֹּ֣ט אֶת־הָעָ֑ם וַיַּעֲמֹ֤ד הָעָם֙ עַל־מֹשֶׁ֔ה מִן־הַבֹּ֖קֶר עַד־הָעָֽרֶב׃ \Rashi{\textbf{ויהי ממחרת.}מוצאי יום הכפורים היה, כך שנינו בספרי ומהו ממחרת? למחרת רדתו מן ההר; ועל כרחך אי אפשר לומר אלא ממחרת יום הכפורים, שהרי קדם מתן תורה א"א לומר "והודעתי את חקי וגו'", ומשנתנה תורה עד יה"כ לא ישב משה לשפט את העם, שהרי בי"ז בתמוז ירד ושבר את הלחות ולמחר עלה בהשכמה ושהה שמונים יום וירד ביה"כ; ואין פרשה זו כתובה כסדר, שלא נאמר ויהי ממחרת עד שנה שניה – אף לדברי האומר יתרו קדם מתן תורה בא, שלוחו אל ארצו לא היה אלא עד שנה שניה – שהרי נאמר כאן וישלח משה את חתנו, ומצינו במסע הדגלים שאמר לו משה "נסעים אנחנו אל המקום וגו'"(במדבר י כט), אל נא תעזב אתנו" (שם לא), ואם זו קדם מתן תורה, מששלחו והלך היכן מצינו שחזר? ואם תאמר, שם לא נאמר יתרו אלא חובב, ובנו של יתרו היה, הוא חובב הוא יתרו, שהרי כתוב "מבני חבב חתן משה" (שופטים ד'): וישב משה וגו' ויעמד העם. יושב כמלך וכלן עומדים, והקשה הדבר ליתרו שהיה מזלזל בכבודן של ישראל והוכיחו על כך, שנאמר (פסוק יד) מדוע אתה יושב לבדך וכלם נצבים: מן הבקר עד ערב. אפשר לומר כן? אלא כל דין שדן דין אמת לאמתו אפלו שעה אחת, מעלה עליו הכתוב כאלו עוסק בתורה כל היום וכאלו נעשה שתף להקב"ה במעשה בראשית, שנאמר בו ויהי ערב וגו' (שבת י'):}%
% ---------- פרק 18 | פסוק 14 ----------
 \textbf{יד} וַיַּרְא֙ חֹתֵ֣ן מֹשֶׁ֔ה אֵ֛ת כׇּל־אֲשֶׁר־ה֥וּא עֹשֶׂ֖ה לָעָ֑ם וַיֹּ֗אמֶר מָֽה־הַדָּבָ֤ר הַזֶּה֙ אֲשֶׁ֨ר אַתָּ֤ה עֹשֶׂה֙ לָעָ֔ם מַדּ֗וּעַ אַתָּ֤ה יוֹשֵׁב֙ לְבַדֶּ֔ךָ וְכׇל־הָעָ֛ם נִצָּ֥ב עָלֶ֖יךָ מִן־בֹּ֥קֶר עַד־עָֽרֶב׃ %
% ---------- פרק 18 | פסוק 15 ----------
 \textbf{טו} וַיֹּ֥אמֶר מֹשֶׁ֖ה לְחֹתְנ֑וֹ כִּֽי־יָבֹ֥א אֵלַ֛י הָעָ֖ם לִדְרֹ֥שׁ אֱלֹהִֽים׃ \Rashi{\textbf{כי יבא.}כמו כי בא, לשון הוה: לדרש אלהים. כתרגומו "למתבע אלפן" – לשאל תלמוד מפי הגבורה:}%
% ---------- פרק 18 | פסוק 16 ----------
 \textbf{טז} כִּֽי־יִהְיֶ֨ה לָהֶ֤ם דָּבָר֙ בָּ֣א אֵלַ֔י וְשָׁ֣פַטְתִּ֔י בֵּ֥ין אִ֖ישׁ וּבֵ֣ין רֵעֵ֑הוּ וְהוֹדַעְתִּ֛י אֶת־חֻקֵּ֥י הָאֱלֹהִ֖ים וְאֶת־תּוֹרֹתָֽיו׃ \Rashi{\textbf{כי יהיה להם דבר בא.}מי שהיה לו הדבר – בא אלי:}%
% ---------- פרק 18 | פסוק 17 ----------
 \textbf{יז} וַיֹּ֛אמֶר חֹתֵ֥ן מֹשֶׁ֖ה אֵלָ֑יו לֹא־טוֹב֙ הַדָּבָ֔ר אֲשֶׁ֥ר אַתָּ֖ה עֹשֶֽׂה׃ \Rashi{\textbf{ויאמר חתן משה.}דרך כבוד קוראו הכתוב חותנו של מלך:}%
% ---------- פרק 18 | פסוק 18 ----------
 \textbf{יח} נָבֹ֣ל תִּבֹּ֔ל גַּם־אַתָּ֕ה גַּם־הָעָ֥ם הַזֶּ֖ה אֲשֶׁ֣ר עִמָּ֑ךְ כִּֽי־כָבֵ֤ד מִמְּךָ֙ הַדָּבָ֔ר לֹא־תוּכַ֥ל עֲשֹׂ֖הוּ לְבַדֶּֽךָ׃ \Rashi{\textbf{נבל תבל.}כתרגומו, ולשונו לשון כמישה, פלייש"טרא בלעז, כמו "והעלה נבל" (ירמיהו ח'), "כנבל עלה מגפן" (ישעיהו ל"ד), שהוא כמוש על ידי חמה ועל ידי קרח וכחו תש ונלאה: גם אתה. לרבות אהרן וחור ושבעים זקנים: כי כבד ממך. כבדו רב יותר מכחך:}%
% ---------- פרק 18 | פסוק 19 ----------
 \textbf{יט} עַתָּ֞ה שְׁמַ֤ע בְּקֹלִי֙ אִיעָ֣צְךָ֔ וִיהִ֥י אֱלֹהִ֖ים עִמָּ֑ךְ הֱיֵ֧ה אַתָּ֣ה לָעָ֗ם מ֚וּל הָֽאֱלֹהִ֔ים וְהֵבֵאתָ֥ אַתָּ֛ה אֶת־הַדְּבָרִ֖ים אֶל־הָאֱלֹהִֽים׃ \Rashi{\textbf{איעצך ויהי אלהים עמך.}בעצה; אמר לו צא המלך בגבורה (מכילתא): היה אתה לעם מול האלהים. שליח ומליץ בינותם למקום ושואל משפטים מאתו: את הדברים. דברי ריבותם:}%
% ---------- פרק 18 | פסוק 20 ----------
 \textbf{כ} וְהִזְהַרְתָּ֣ה אֶתְהֶ֔ם אֶת־הַחֻקִּ֖ים וְאֶת־הַתּוֹרֹ֑ת וְהוֹדַעְתָּ֣ לָהֶ֗ם אֶת־הַדֶּ֙רֶךְ֙ יֵ֣לְכוּ בָ֔הּ וְאֶת־הַֽמַּעֲשֶׂ֖ה אֲשֶׁ֥ר יַעֲשֽׂוּן׃ %
% ---------- פרק 18 | פסוק 21 ----------
 \textbf{כא} וְאַתָּ֣ה תֶחֱזֶ֣ה מִכׇּל־הָ֠עָ֠ם אַנְשֵׁי־חַ֜יִל יִרְאֵ֧י אֱלֹהִ֛ים אַנְשֵׁ֥י אֱמֶ֖ת שֹׂ֣נְאֵי בָ֑צַע וְשַׂמְתָּ֣ עֲלֵהֶ֗ם שָׂרֵ֤י אֲלָפִים֙ שָׂרֵ֣י מֵא֔וֹת שָׂרֵ֥י חֲמִשִּׁ֖ים וְשָׂרֵ֥י עֲשָׂרֹֽת׃ \Rashi{\textbf{ואתה תחזה.}ברוח הקדש שעליך: אנשי חיל. עשירים, שאין צריכין להחניף ולהכיר פנים (שם): אנשי אמת. אלו בעלי הבטחה, שהם כדאי לסמך על דבריהם, שעל ידי כן יהיו דבריהם נשמעין: שנאי בצע. ששונאין את ממונם בדין, כההיא דאמרינן: כל דינא דמפקין ממונא מיניה בדינא לאו דינא הוא (בבא בתרא נ"ח): שרי אלפים. הם היו שש מאות שרים לשש מאות אלף: שרי מאות. שש אלפים היו: שרי חמשים. י"ב אלף: שרי עשרת. ששים אלף:}%
% ---------- פרק 18 | פסוק 22 ----------
 \textbf{כב} וְשָׁפְט֣וּ אֶת־הָעָם֮ בְּכׇל־עֵת֒ וְהָיָ֞ה כׇּל־הַדָּבָ֤ר הַגָּדֹל֙ יָבִ֣יאוּ אֵלֶ֔יךָ וְכׇל־הַדָּבָ֥ר הַקָּטֹ֖ן יִשְׁפְּטוּ־הֵ֑ם וְהָקֵל֙ מֵֽעָלֶ֔יךָ וְנָשְׂא֖וּ אִתָּֽךְ׃ \Rashi{\textbf{ושפטו.}וידונון, לשון צווי: והקל מעליך. דבר זה להקל מעליך: והקל. כמו "והכבד את לבו" (שמות ח'), "והכות את מואב" (מלכים ב ג'), לשון הוה:}%
% ---------- פרק 18 | פסוק 23 ----------
 \textbf{כג} אִ֣ם אֶת־הַדָּבָ֤ר הַזֶּה֙ תַּעֲשֶׂ֔ה וְצִוְּךָ֣ אֱלֹהִ֔ים וְיָֽכׇלְתָּ֖ עֲמֹ֑ד וְגַם֙ כׇּל־הָעָ֣ם הַזֶּ֔ה עַל־מְקֹמ֖וֹ יָבֹ֥א בְשָׁלֽוֹם׃ \Rashi{\textbf{וצוך אלהים ויכלת עמד.}המלך בגבורה, אם מצוה אותך לעשות כן, תוכל עמד, ואם יעכב על ידך לא תוכל לעמד (מכילתא): וגם כל העם הזה. אהרן נדב ואביהוא ושבעים זקנים הנלוים עתה עמך (שם):}%
% ---------- פרק 18 | פסוק 24 ----------
 \textbf{כד} וַיִּשְׁמַ֥ע מֹשֶׁ֖ה לְק֣וֹל חֹתְנ֑וֹ וַיַּ֕עַשׂ כֹּ֖ל אֲשֶׁ֥ר אָמָֽר׃ %
% ---------- פרק 18 | פסוק 25 ----------
 \textbf{כה} וַיִּבְחַ֨ר מֹשֶׁ֤ה אַנְשֵׁי־חַ֙יִל֙ מִכׇּל־יִשְׂרָאֵ֔ל וַיִּתֵּ֥ן אֹתָ֛ם רָאשִׁ֖ים עַל־הָעָ֑ם שָׂרֵ֤י אֲלָפִים֙ שָׂרֵ֣י מֵא֔וֹת שָׂרֵ֥י חֲמִשִּׁ֖ים וְשָׂרֵ֥י עֲשָׂרֹֽת׃ %
% ---------- פרק 18 | פסוק 26 ----------
 \textbf{כו} וְשָׁפְט֥וּ אֶת־הָעָ֖ם בְּכׇל־עֵ֑ת אֶת־הַדָּבָ֤ר הַקָּשֶׁה֙ יְבִיא֣וּן אֶל־מֹשֶׁ֔ה וְכׇל־הַדָּבָ֥ר הַקָּטֹ֖ן יִשְׁפּוּט֥וּ הֵֽם׃ \Rashi{\textbf{ושפטו.}ודינין ית עמא: יביאון. מיתין: ישפוטו הם. כמו ישפטו, וכן "לא תעבורי" (רות ב'), כמו לא תעברי, ותרגומו דינין אינון; מקראות העליונים הם לשון צווי, לכך מתרגמין וידונון, ייתון, ידונון, ומקראות הללו לשון עשיה:}%
% ---------- פרק 18 | פסוק 27 ----------
 \textbf{כז} וַיְשַׁלַּ֥ח מֹשֶׁ֖ה אֶת־חֹתְנ֑וֹ וַיֵּ֥לֶךְ ל֖וֹ אֶל־אַרְצֽוֹ׃ \{פ\} \Rashi{\textbf{וילך לו אל ארצו.}לגיר בני משפחתו (מכילתא):}%
% ---------- פרק 19 | פסוק 1 ----------
 \textbf{א} בַּחֹ֙דֶשׁ֙ הַשְּׁלִישִׁ֔י לְצֵ֥את בְּנֵי־יִשְׂרָאֵ֖ל מֵאֶ֣רֶץ מִצְרָ֑יִם בַּיּ֣וֹם הַזֶּ֔ה בָּ֖אוּ מִדְבַּ֥ר סִינָֽי׃ \Rashi{\textbf{ביום הזה.}בראש חדש; לא היה צריך לכתב אלא ביום ההוא, מהו ביום הזה? שיהיו דברי תורה חדשים עליך כאלו היום נתנם (ברכות ס"ג):}%
% ---------- פרק 19 | פסוק 2 ----------
 \textbf{ב} וַיִּסְע֣וּ מֵרְפִידִ֗ים וַיָּבֹ֙אוּ֙ מִדְבַּ֣ר סִינַ֔י וַֽיַּחֲנ֖וּ בַּמִּדְבָּ֑ר וַיִּֽחַן־שָׁ֥ם יִשְׂרָאֵ֖ל נֶ֥גֶד הָהָֽר׃ \Rashi{\textbf{ויסעו מרפידים.}מה ת"ל לחזר ולפרש מהיכן נסעו, והלא כבר כתב שברפידים היו חונים, בידוע שמשם נסעו? אלא להקיש נסיעתן מרפידים לביאתן למדבר סיני, מה ביאתן למדבר סיני בתשובה, אף נסיעתן מרפידים בתשובה (מכילתא): ויחן שם ישראל. כאיש אחד בלב אחד, אבל שאר כל החניות בתרעומות ובמחלקת: נגד ההר. למזרחו, וכל מקום שאתה מוצא נגד, פנים למזרח (שם):}%
% ---------- פרק 19 | פסוק 3 ----------
 \textbf{ג} וּמֹשֶׁ֥ה עָלָ֖ה אֶל־הָאֱלֹהִ֑ים וַיִּקְרָ֨א אֵלָ֤יו יְהֹוָה֙ מִן־הָהָ֣ר לֵאמֹ֔ר כֹּ֤ה תֹאמַר֙ לְבֵ֣ית יַעֲקֹ֔ב וְתַגֵּ֖יד לִבְנֵ֥י יִשְׂרָאֵֽל׃ \Rashi{\textbf{ומשה עלה.}ביום השני; וכל עליותיו בהשכמה היו, שנאמר "וישכם משה בבקר" (שמות ל"ד): כה תאמר. בלשון הזה וכסדר הזה: לבית יעקב. אלו הנשים, תאמר להן בלשון רכה (מכילתא): ותגיד לבני ישראל. ענשין ודקדוקים פרש לזכרים, דברים הקשין כגידין (שבת פ"ז):}%
% ---------- פרק 19 | פסוק 4 ----------
 \textbf{ד} אַתֶּ֣ם רְאִיתֶ֔ם אֲשֶׁ֥ר עָשִׂ֖יתִי לְמִצְרָ֑יִם וָאֶשָּׂ֤א אֶתְכֶם֙ עַל־כַּנְפֵ֣י נְשָׁרִ֔ים וָאָבִ֥א אֶתְכֶ֖ם אֵלָֽי׃ \Rashi{\textbf{אתם ראיתם.}לא מסרת היא בידכם, לא בדברים אני משגר לכם, לא בעדים אני מעיד עליכם, אלא אתם ראיתם אשר עשיתי למצרים, על כמה עברות היו חיבין לי קדם שנזדוגו לכם ולא נפרעתי מהם אלא על ידכם (מכילתא): ואשא אתכם. זה היום שבאו ישראל לרעמסס, שהיו ישראל מפזרין בכל ארץ גשן, ולשעה קלה כשבאו לסע ולצאת, נקבצו כלם לרעמסס, ואנקלוס תרגם ואשא כמו ואסיע אתכם – "ואטלית יתכון", תקן את הדבר דרך כבוד למעלה: על כנפי נשרים. כנשר הנושא גוזליו על כנפיו, שכל שאר העופות נותנים את בניהם בין רגליהם, לפי שמתיראין מעוף אחר שפורח על גביהם, אבל הנשר הזה אינו מתירא אלא מן האדם שמא יזרק בו חץ, לפי שאין עוף פורח על גביו, לכך נותנו על כנפיו, אומר מוטב יכנס החץ בי ולא בבני, אף אני עשיתי כן: "ויסע מלאך האלהים וגו', ויבא בין מחנה מצרים וגו'" (שמות י"ד), והיו מצרים זורקים חצים ואבני בליסטראות והענן מקבלם: ואבא אתכם אלי. כתרגומו:}%
% ---------- פרק 19 | פסוק 5 ----------
 \textbf{ה} וְעַתָּ֗ה אִם־שָׁמ֤וֹעַ תִּשְׁמְעוּ֙ בְּקֹלִ֔י וּשְׁמַרְתֶּ֖ם אֶת־בְּרִיתִ֑י וִהְיִ֨יתֶם לִ֤י סְגֻלָּה֙ מִכׇּל־הָ֣עַמִּ֔ים כִּי־לִ֖י כׇּל־הָאָֽרֶץ׃ \Rashi{\textbf{ועתה.}אם עתה תקבלו עליכם, יערב לכם מכאן ואילך, שכל התחלות קשות (מכילתא): ושמרתם את בריתי. שאכרת עמכם על שמירת התורה: סגלה. אוצר חביב, כמו "וסגלת מלכים" (קהלת ב') – כלי יקר ואבנים טובות שהמלכים גונזים אותם, כך אתם תהיו לי סגלה משאר אמות; ולא תאמרו אתם לבדכם שלי ואין לי אחרים עמכם ומה יש לי עוד שתהא חבתכם נכרת: כי לי כל הארץ והם בעיני ולפני לכלום:}%
% ---------- פרק 19 | פסוק 6 ----------
 \textbf{ו} וְאַתֶּ֧ם תִּהְיוּ־לִ֛י מַמְלֶ֥כֶת כֹּהֲנִ֖ים וְג֣וֹי קָד֑וֹשׁ אֵ֚לֶּה הַדְּבָרִ֔ים אֲשֶׁ֥ר תְּדַבֵּ֖ר אֶל־בְּנֵ֥י יִשְׂרָאֵֽל׃ \Rashi{\textbf{ואתם תהיו לי ממלכת כהנים.}שרים, כמה דאת אמר "ובני דוד כהנים" (שמואל ב ח'): אלה הדברים. לא פחות ולא יותר (מכילתא):}%
% ---------- פרק 19 | פסוק 7 ----------
 \textbf{ז} וַיָּבֹ֣א מֹשֶׁ֔ה וַיִּקְרָ֖א לְזִקְנֵ֣י הָעָ֑ם וַיָּ֣שֶׂם לִפְנֵיהֶ֗ם אֵ֚ת כׇּל־הַדְּבָרִ֣ים הָאֵ֔לֶּה אֲשֶׁ֥ר צִוָּ֖הוּ יְהֹוָֽה׃ %
% ---------- פרק 19 | פסוק 8 ----------
 \textbf{ח} וַיַּעֲנ֨וּ כׇל־הָעָ֤ם יַחְדָּו֙ וַיֹּ֣אמְר֔וּ כֹּ֛ל אֲשֶׁר־דִּבֶּ֥ר יְהֹוָ֖ה נַעֲשֶׂ֑ה וַיָּ֧שֶׁב מֹשֶׁ֛ה אֶת־דִּבְרֵ֥י הָעָ֖ם אֶל־יְהֹוָֽה׃ \Rashi{\textbf{וישב משה את דברי העם וגו'.}ביום המחרת, שהוא שלישי, שהרי בהשכמה עלה; וכי צריך היה משה להשיב? אלא בא הכתוב ללמדך דרך ארץ ממשה, שלא אמר הואיל ויודע מי ששלחני איני צריך להשיב (שבת פ"ז):}%
% ---------- פרק 19 | פסוק 9 ----------
 \textbf{ט} וַיֹּ֨אמֶר יְהֹוָ֜ה אֶל־מֹשֶׁ֗ה הִנֵּ֨ה אָנֹכִ֜י בָּ֣א אֵלֶ֘יךָ֮ בְּעַ֣ב הֶֽעָנָן֒ בַּעֲב֞וּר יִשְׁמַ֤ע הָעָם֙ בְּדַבְּרִ֣י עִמָּ֔ךְ וְגַם־בְּךָ֖ יַאֲמִ֣ינוּ לְעוֹלָ֑ם וַיַּגֵּ֥ד מֹשֶׁ֛ה אֶת־דִּבְרֵ֥י הָעָ֖ם אֶל־יְהֹוָֽה׃ \Rashi{\textbf{בעב הענן.}במעבה הענן וזהו ערפל: וגם בך. גם בנביאים הבאים אחריך (מכילתא): ויגד משה את דברי וגו'. ביום המחרת, שהוא רביעי לחדש: את דברי העם וגו'. תשובה על דבר זה; שמעתי מהם שרצונם לשמע ממך, אינו דומה השומע מפי שליח לשומע מפי המלך, רצוננו לראות את מלכנו (מכילתא):}%
% ---------- פרק 19 | פסוק 10 ----------
 \textbf{י} וַיֹּ֨אמֶר יְהֹוָ֤ה אֶל־מֹשֶׁה֙ לֵ֣ךְ אֶל־הָעָ֔ם וְקִדַּשְׁתָּ֥ם הַיּ֖וֹם וּמָחָ֑ר וְכִבְּס֖וּ שִׂמְלֹתָֽם׃ \Rashi{\textbf{ויאמר ה' אל משה.}אם כן שמזקיקין לדבר עמם, לך אל העם:  וקדשתם. וזמנתם, שיכינו עצמם היום ומחר:}%
% ---------- פרק 19 | פסוק 11 ----------
 \textbf{יא} וְהָי֥וּ נְכֹנִ֖ים לַיּ֣וֹם הַשְּׁלִישִׁ֑י כִּ֣י ׀ בַּיּ֣וֹם הַשְּׁלִשִׁ֗י יֵרֵ֧ד יְהֹוָ֛ה לְעֵינֵ֥י כׇל־הָעָ֖ם עַל־הַ֥ר סִינָֽי׃ \Rashi{\textbf{והיו נכנים.}מבדלים מאשה (שבת פ"ז): ליום השלישי. שהוא ששה בחדש, ובחמישי בנה משה את המזבח תחת ההר ושתים עשרה מצבה – כל הענין האמור בפרשת ואלה המשפטים – ואין מקדם ומאחר בתורה: לעיני כל העם. מלמד שלא היה בהם סומא, שנתרפאו כלם (מכילתא):}%
% ---------- פרק 19 | פסוק 12 ----------
 \textbf{יב} וְהִגְבַּלְתָּ֤ אֶת־הָעָם֙ סָבִ֣יב לֵאמֹ֔ר הִשָּׁמְר֥וּ לָכֶ֛ם עֲל֥וֹת בָּהָ֖ר וּנְגֹ֣עַ בְּקָצֵ֑הוּ כׇּל־הַנֹּגֵ֥עַ בָּהָ֖ר מ֥וֹת יוּמָֽת׃ \Rashi{\textbf{והגבלת.}קבע להם תחומין לסימן, שלא יקרבו מן הגבול והלאה: לאמר. הגבול אומר להם השמרו מעלות מכאן ולהלאה, ואתה הזהירם על כך: ונגע בקצהו. אפלו בקצהו:}%
% ---------- פרק 19 | פסוק 13 ----------
 \textbf{יג} לֹא־תִגַּ֨ע בּ֜וֹ יָ֗ד כִּֽי־סָק֤וֹל יִסָּקֵל֙ אוֹ־יָרֹ֣ה יִיָּרֶ֔ה אִם־בְּהֵמָ֥ה אִם־אִ֖ישׁ לֹ֣א יִחְיֶ֑ה בִּמְשֹׁךְ֙ הַיֹּבֵ֔ל הֵ֖מָּה יַעֲל֥וּ בָהָֽר׃ \Rashi{\textbf{ירה יירה.}מכאן לנסקלין שהם נדחין למטה מבית הסקילה שהיה גבוה שתי קומות (סנהדרין מ"ה): יירה. ישלך למטה לארץ, כמו "ירה בים": במשך היבל. כשימשך היובל קול ארך, הוא סימן סלוק שכינה והפסקת הקול, וכיון שאסתלק, הם רשאין לעלות: היבל. הוא שופר של איל, שכן בערביא קורין לדכרא יובלא, ושופר של איל של יצחק היה (פרקי דרבי אליעזר ל"א):}%
% ---------- פרק 19 | פסוק 14 ----------
 \textbf{יד} וַיֵּ֧רֶד מֹשֶׁ֛ה מִן־הָהָ֖ר אֶל־הָעָ֑ם וַיְקַדֵּשׁ֙ אֶת־הָעָ֔ם וַֽיְכַבְּס֖וּ שִׂמְלֹתָֽם׃ \Rashi{\textbf{מן ההר אל העם.}מלמד שלא היה משה פונה לעסקיו, אלא מן ההר – אל העם (מכילתא):}%
% ---------- פרק 19 | פסוק 15 ----------
 \textbf{טו} וַיֹּ֙אמֶר֙ אֶל־הָעָ֔ם הֱי֥וּ נְכֹנִ֖ים לִשְׁלֹ֣שֶׁת יָמִ֑ים אַֽל־תִּגְּשׁ֖וּ אֶל־אִשָּֽׁה׃ \Rashi{\textbf{היו נכנים לשלשת ימים.}לסוף ג' ימים, הוא יום רביעי, שהוסיף משה יום אחד מדעתו כדברי רבי יוסי, ולדברי האומר בששה בחדש נתנו עשרת הדברות, לא הוסיף משה כלום, ולשלשת ימים כמו ליום השלישי: אל תגשו אל אשה. כל ג' ימים הללו, כדי שיהיו הנשים טובלות ליום השלישי ותהיינה טהורות לקבל תורה, שאם ישמשו תוך ג' ימים שמא תפלט האשה שכבת זרע לאחר טבילתה ותחזר ותטמא, אבל מששהתה ג' ימים, כבר הזרע מסריח ואינו ראוי להזריע, וטהור מלטמא את הפולטת (שבת פ"ו):}%
% ---------- פרק 19 | פסוק 16 ----------
 \textbf{טז} וַיְהִי֩ בַיּ֨וֹם הַשְּׁלִישִׁ֜י בִּֽהְיֹ֣ת הַבֹּ֗קֶר וַיְהִי֩ קֹלֹ֨ת וּבְרָקִ֜ים וְעָנָ֤ן כָּבֵד֙ עַל־הָהָ֔ר וְקֹ֥ל שֹׁפָ֖ר חָזָ֣ק מְאֹ֑ד וַיֶּחֱרַ֥ד כׇּל־הָעָ֖ם אֲשֶׁ֥ר בַּֽמַּחֲנֶֽה׃ \Rashi{\textbf{בהית הבקר.}מלמד שהקדים על ידם, מה שאין דרך בשר ודם לעשות כן – שיהא הרב ממתין לתלמיד; וכן מצינו "קום צא אל הבקעה, ואקום ואצא אל הבקעה והנה שם כבוד ה' עמד" (יחזקאל ג'):}%
% ---------- פרק 19 | פסוק 17 ----------
 \textbf{יז} וַיּוֹצֵ֨א מֹשֶׁ֧ה אֶת־הָעָ֛ם לִקְרַ֥את הָֽאֱלֹהִ֖ים מִן־הַֽמַּחֲנֶ֑ה וַיִּֽתְיַצְּב֖וּ בְּתַחְתִּ֥ית הָהָֽר׃ \Rashi{\textbf{לקראת האלהים.}מגיד שהשכינה יצאה לקראתם כחתן היוצא לקראת כלה, וזהו שאמר ה' מסיני בא (דברים ל"ג), ולא נאמר לסיני בא (מכילתא): בתחתית ההר. לפי פשוטו ברגלי ההר; ומדרשו שנתלש ההר ממקומו ונכפה עליהם כגיגית (שבת פ"ח):}%
% ---------- פרק 19 | פסוק 18 ----------
 \textbf{יח} וְהַ֤ר סִינַי֙ עָשַׁ֣ן כֻּלּ֔וֹ מִ֠פְּנֵ֠י אֲשֶׁ֨ר יָרַ֥ד עָלָ֛יו יְהֹוָ֖ה בָּאֵ֑שׁ וַיַּ֤עַל עֲשָׁנוֹ֙ כְּעֶ֣שֶׁן הַכִּבְשָׁ֔ן וַיֶּחֱרַ֥ד כׇּל־הָהָ֖ר מְאֹֽד׃ \Rashi{\textbf{עשן כלו.}אין עשן זה שם דבר, שהרי נקוד השי"ן פתח, אלא לשון פעל כמו אמר, שמר, שמע, לכך תרגומו תנן כלה, ולא תרגם תננא, וכל עשן שבמקרא נקודים קמץ מפני שהם שם דבר: הכבשן. של סיד; יכול ככבשן זה ולא יותר, ת"ל בוער באש עד לב השמים (דברים ד'), ומה ת"ל כבשן? לשבר את האזן מה שהיא יכולה לשמע – נותן לבריות סימן הנכר להם; כיוצא בו כאריה ישאג (הושע י"א), וכי מי נתן כח בארי אלא הוא? והכתוב מושלו כאריה, אלא אנו מכנין ומדמין אותו לבריותיו כדי לשבר את האזן מה שיכולה לשמע; כיוצא בו וקולו כקול מים רבים (יחזקאל מ"ג), וכי מי נתן קול למים אלא הוא? ואתה מכנה אותו, לדמותו לבריותיו – כדי לשבר את האזן (מכילתא):}%
% ---------- פרק 19 | פסוק 19 ----------
 \textbf{יט} וַֽיְהִי֙ ק֣וֹל הַשֹּׁפָ֔ר הוֹלֵ֖ךְ וְחָזֵ֣ק מְאֹ֑ד מֹשֶׁ֣ה יְדַבֵּ֔ר וְהָאֱלֹהִ֖ים יַעֲנֶ֥נּוּ בְקֽוֹל׃ \Rashi{\textbf{הולך וחזק מאד.}מנהג הדיוט כל זמן שהוא מאריך לתקע קולו מחליש וכוהה, אבל כאן הולך וחזק מאד, ולמה כך מתחלה? לשבר אזניהם מה שיכולין לשמע: משה ידבר. כשהיה משה מדבר ומשמיע הדברות לישראל – שהרי לא שמעו מפי הגבורה אלא אנכי ולא יהיה לך – והקב"ה מסיעו לתת בו כח להיות קולו מגביר ונשמע (שם): יעננו בקול. יעננו על דבר הקול, כמו אשר יענה באש (מלכים א י"ח) – על דבר האש, להורידו:}%
% ---------- פרק 19 | פסוק 20 ----------
 \textbf{כ} וַיֵּ֧רֶד יְהֹוָ֛ה עַל־הַ֥ר סִינַ֖י אֶל־רֹ֣אשׁ הָהָ֑ר וַיִּקְרָ֨א יְהֹוָ֧ה לְמֹשֶׁ֛ה אֶל־רֹ֥אשׁ הָהָ֖ר וַיַּ֥עַל מֹשֶֽׁה׃ \Rashi{\textbf{וירד ה' על הר סיני.}יכול ירד עליו ממש, ת"ל כי מן השמים דברתי עמכם (שמות כ'), מלמד שהרכין שמים העליונים והתחתונים והציען על גבי ההר כמצע על המטה וירד כסא הכבוד עליהם (מכילתא):}%
% ---------- פרק 19 | פסוק 21 ----------
 \textbf{כא} וַיֹּ֤אמֶר יְהֹוָה֙ אֶל־מֹשֶׁ֔ה רֵ֖ד הָעֵ֣ד בָּעָ֑ם פֶּן־יֶהֶרְס֤וּ אֶל־יְהֹוָה֙ לִרְא֔וֹת וְנָפַ֥ל מִמֶּ֖נּוּ רָֽב׃ \Rashi{\textbf{העד בעם.}התרה בהם שלא לעלות בהר: פן יהרסו וגו'. שלא יהרסו את מצבם, על ידי שתאותם אל ה', לראות, ויקרבו לצד ההר: ונפל ממנו רב. כל מה שיפל מהם, ואפלו הוא יחידי, חשוב לפני רב (שם): פן יהרסו. כל הריסה מפרדת אסיפת הבנין, אף הנפרדין ממצב אנשים הורסים את המצב:}%
% ---------- פרק 19 | פסוק 22 ----------
 \textbf{כב} וְגַ֧ם הַכֹּהֲנִ֛ים הַנִּגָּשִׁ֥ים אֶל־יְהֹוָ֖ה יִתְקַדָּ֑שׁוּ פֶּן־יִפְרֹ֥ץ בָּהֶ֖ם יְהֹוָֽה׃ \Rashi{\textbf{וגם הכהנים.}אף הבכורות שהעבודה בהם (זבחים קט"ו): הנגשים אל ה'. להקריב קרבנות, אף הם לא יסמכו על חשיבותם לעלות: יתקדשו. יהיו מזמנים להתיצב על עמדן: פן יפרץ. לשון פרצה; יהרג בהם ויעשה בהם פרצה:}%
% ---------- פרק 19 | פסוק 23 ----------
 \textbf{כג} וַיֹּ֤אמֶר מֹשֶׁה֙ אֶל־יְהֹוָ֔ה לֹא־יוּכַ֣ל הָעָ֔ם לַעֲלֹ֖ת אֶל־הַ֣ר סִינָ֑י כִּֽי־אַתָּ֞ה הַעֵדֹ֤תָה בָּ֙נוּ֙ לֵאמֹ֔ר הַגְבֵּ֥ל אֶת־הָהָ֖ר וְקִדַּשְׁתּֽוֹ׃ \Rashi{\textbf{לא יוכל העם.}איני צריך להעיד בהם, שהרי מתרים ועומדין הם היום שלשת ימים, ולא יוכלו לעלות שאין להם רשות:}%
% ---------- פרק 19 | פסוק 24 ----------
 \textbf{כד} וַיֹּ֨אמֶר אֵלָ֤יו יְהֹוָה֙ לֶךְ־רֵ֔ד וְעָלִ֥יתָ אַתָּ֖ה וְאַהֲרֹ֣ן עִמָּ֑ךְ וְהַכֹּהֲנִ֣ים וְהָעָ֗ם אַל־יֶֽהֶרְס֛וּ לַעֲלֹ֥ת אֶל־יְהֹוָ֖ה פֶּן־יִפְרׇץ־בָּֽם׃ \Rashi{\textbf{לך רד.}והעד בהם שנית, שמזרזין את האדם קדם מעשה וחוזרין ומזרזין אותו בשעת מעשה (מכילתא): ועלית אתה ואהרן עמך והכהנים. יכול אף הם עמך, ת"ל ועלית אתה, אמר מעתה, אתה מחצה לעצמך, ואהרן מחצה לעצמו, והכהנים מחצה לעצמם – משה נגש יותר מאהרן, ואהרן יותר מן הכהנים, והעם כל עקר אל יהרסו את מצבם לעלות אל ה' (שם): פן יפרץ בם. אע"פ שהוא נקוד חטף קמץ אינו זז מגזרתו, כך דרך כל תבה שנקדתה מלאפום, כשהיא סמוכה באה במקף, משתנה הנקוד לחטף קמץ:}%
% ---------- פרק 19 | פסוק 25 ----------
 \textbf{כה} וַיֵּ֥רֶד מֹשֶׁ֖ה אֶל־הָעָ֑ם וַיֹּ֖אמֶר אֲלֵהֶֽם׃ \{ס\} \Rashi{\textbf{ויאמר אלהם.}התראה זו:}%
% ---------- פרק 20 | פסוק 1 ----------
 \textbf{א} וַיְדַבֵּ֣ר אֱלֹהִ֔ים אֵ֛ת כׇּל־הַדְּבָרִ֥ים הָאֵ֖לֶּה לֵאמֹֽר׃ \{ס\} \Rashi{\textbf{וידבר אלהים.}אין אלהים אלא דין; לפי שיש פרשיות בתורה שאם עשאן אדם מקבל שכר, ואם לאו אינו מקבל עליהם פרענות, יכול אף עשרת הדברות כן, ת"ל וידבר אלהים – דין לפרע (מכילתא): את כל הדברים האלה. מלמד שאמר הקב"ה עשרת הדברות בדבור אחד, מה שאי אפשר לאדם לומר כן, אם כן מה ת"ל עוד אנכי ולא יהיה לך? שחזר ופרש על כל דבור ודבור בפני עצמו (שם): לאמר. מלמד שהיו עונין על הן הן ועל לאו לאו (שם):}%
% ---------- פרק 20 | פסוק 2 ----------
 \textbf{ב} אָֽנֹכִ֖י֙ יְהֹוָ֣ה אֱלֹהֶ֑֔יךָ אֲשֶׁ֧ר הוֹצֵאתִ֛יךָ מֵאֶ֥רֶץ מִצְרַ֖יִם מִבֵּ֣֥ית עֲבָדִ֑͏ֽים׃ \Rashi{\textbf{אשר הוצאתיך מארץ מצרים.}כדאי היא ההוצאה שתהיו משעבדים לי; ד"א, לפי שנגלה בים כגבור מלחמה ונגלה כאן כזקן מלא רחמים, – שנאמר ותחת רגליו כמעשה לבנת הספיר (שמות כ"ד), זו היתה לפניו בשעת השעבוד, וכעצם השמים, משנגאלו – הואיל ואני משתנה במראות, אל תאמרו שתי רשויות הן, אנכי הוא אשר הוצאתיך ממצרים ועל הים. ד"א לפי שהיו שומעין קולות הרבה, שנאמר את הקולת – קולות באין מארבע רוחות ומן השמים ומן הארץ – אל תאמרו רשויות הרבה הן; ולמה אמר לשון יחיד אלהיך? לתן פתחון פה למשה ללמד סנגוריא במעשה העגל, וזה שאמר למה ה' יחרה אפך בעמך – לא להם צוית לא יהיה לכם אלהים אחרים אלא לי לבדי (שמות רבה): מבית עבדים. מבית פרעה שהייתם עבדים לו, או אינו אומר אלא מבית עבדים שהיו עבדים לעבדים? ת"ל ויפדך מבית עבדים מיד פרעה מלך מצרים (דברים ז'), אמר מעתה עבדים למלך היו ולא עבדים לעבדים (מכילתא):}%
% ---------- פרק 20 | פסוק 3 ----------
 \textbf{ג} לֹֽ֣א־יִהְיֶ֥͏ֽה־לְךָ֛֩ אֱלֹהִ֥֨ים אֲחֵרִ֖֜ים עַל־פָּנָֽ͏ַ֗י׃ \Rashi{\textbf{לא יהיה לך.}למה נאמר? לפי שנ' לא תעשה לך, אין לי אלא שלא יעשה, העשוי כבר מנין שלא יקים? ת"ל לא יהיה לך (שם): אלהים אחרים. שאינן אלוהות אלא אחרים עשאום אלהים עליהם, ולא יתכן לפרש אלהים אחרים זולתי, שגנאי כלפי מעלה לקראותם אלוהות אצלו. ד"א אלהים אחרים, שהם אחרים לעובדיהם – צועקים אליהם ואינן עונין אותם, ודומה כאלו הוא אחר שאינו מכירו מעולם: על פני. כל זמן שאני קים, שלא תאמר לא נצטוו על ע"ז אלא אותו הדור (מכילתא):}%
% ---------- פרק 20 | פסוק 4 ----------
 \textbf{ד} לֹֽ֣א־תַעֲשֶֽׂ֨ה־לְךָ֥֣ פֶ֣֙סֶל֙ ׀ וְכׇל־תְּמוּנָ֔֡ה אֲשֶׁ֤֣ר בַּשָּׁמַ֣֙יִם֙ ׀ מִמַּ֔֡עַל וַֽאֲשֶׁ֥ר֩ בָּאָ֖֨רֶץ מִתָּ֑͏ַ֜חַת וַאֲשֶׁ֥ר בַּמַּ֖֣יִם ׀ מִתַּ֥֣חַת לָאָֽ֗רֶץ׃ \Rashi{\textbf{פסל.}על שם שנפסל: וכל תמונה. תמונת דבר אשר בשמים:}%
% ---------- פרק 20 | פסוק 5 ----------
 \textbf{ה} לֹֽא־תִשְׁתַּחֲוֶ֥֣ה לָהֶ֖ם֮ וְלֹ֣א תׇעׇבְדֵ֑ם֒ כִּ֣י אָֽנֹכִ֞י יְהֹוָ֤ה אֱלֹהֶ֙יךָ֙ אֵ֣ל קַנָּ֔א פֹּ֠קֵ֠ד עֲוֺ֨ן אָבֹ֧ת עַל־בָּנִ֛ים עַל־שִׁלֵּשִׁ֥ים וְעַל־רִבֵּעִ֖ים לְשֹׂנְאָֽ֑י׃ \Rashi{\textbf{אל קנא.}מקנא לפרע ואינו עובר על מדתו למחל על עבודת אלילים; כל לשון קנא אנפרי"מנט בלעז – נותן לב לפרע: לשנאי. כתרגומו, כשאוחזין מעשה אבותיהם בידיהם.}%
% ---------- פרק 20 | פסוק 6 ----------
 \textbf{ו} וְעֹ֥֤שֶׂה חֶ֖֙סֶד֙ לַאֲלָפִ֑֔ים לְאֹהֲבַ֖י וּלְשֹׁמְרֵ֥י מִצְוֺתָֽי׃ \{ס\} \Rashi{\textbf{נוצר חסד.}שאדם עושה, לשלם שכר עד לאלפים דור, נמצאת מדה טובה יתרה על מדת פרעניות אחת על חמש מאות, שזו לארבעה דורות, וזו לאלפים (תוספתא סוטה ד, א):}%
% ---------- פרק 20 | פסוק 7 ----------
 \textbf{ז} לֹ֥א תִשָּׂ֛א אֶת־שֵֽׁם־יְהֹוָ֥ה אֱלֹהֶ֖יךָ לַשָּׁ֑וְא כִּ֣י לֹ֤א יְנַקֶּה֙ יְהֹוָ֔ה אֵ֛ת אֲשֶׁר־יִשָּׂ֥א אֶת־שְׁמ֖וֹ לַשָּֽׁוְא׃ \{פ\} \Rashi{\textbf{לשוא.}חנם, להבל, ואיזהו שבועת שוא? נשבע לשנות את הידוע – על עמוד של אבן שהוא של זהב (שבועות כ"ט):}%
% ---------- פרק 20 | פסוק 8 ----------
 \textbf{ח} זָכ֛וֹר֩ אֶת־י֥֨וֹם הַשַּׁבָּ֖֜ת לְקַדְּשֽׁ֗וֹ׃ \Rashi{\textbf{זכור.}זכור ושמור בדבור אחד נאמרו, וכן מחלליה מות יומת (שמות ל"א), וביום השבת שני כבשים (במדבר כ"ח); וכן לא תלבש שעטנז, גדלים תעשה לך (דברים כ"ב); וכן ערות אשת אחיך (ויקרא י"ח), יבמה יבא עליה (דברים כ"ה); הוא שנאמר אחת דבר אלהים שתים זו שמעתי (תהילים ס"ב): זכור. לשון פעול הוא, כמו אכול ושתו (ישעיהו כ"ב), הלוך ובכה (שמואל ב ג'), וכן פתרונו: תנו לב לזכר תמיד את יום השבת, שאם נזדמן לך חפץ יפה תהא מזמינו לשבת (מכילתא):}%
% ---------- פרק 20 | פסוק 9 ----------
 \textbf{ט} שֵׁ֤֣שֶׁת יָמִ֣ים֙ תַּֽעֲבֹ֔ד֮ וְעָשִׂ֖֣יתָ כׇּֿל־מְלַאכְתֶּֽךָ֒׃ \Rashi{\textbf{ועשית כל מלאכתך.}כשתבא שבת יהא בעיניך כאלו כל מלאכתך עשויה, שלא תהרהר אחר מלאכה (שם):}%
% ---------- פרק 20 | פסוק 10 ----------
 \textbf{י} וְי֨וֹם֙ הַשְּׁבִיעִ֔֜י שַׁבָּ֖֣ת ׀ לַיהֹוָ֣ה אֱלֹהֶ֑֗יךָ לֹֽ֣א־תַעֲשֶׂ֣֨ה כׇל־מְלָאכָ֜֡ה אַתָּ֣ה ׀ וּבִנְךָ֣͏ֽ־וּ֠בִתֶּ֗ךָ עַבְדְּךָ֤֨ וַאֲמָֽתְךָ֜֙ וּבְהֶמְתֶּ֔֗ךָ וְגֵרְךָ֖֙ אֲשֶׁ֥֣ר בִּשְׁעָרֶֽ֔יךָ׃ \Rashi{\textbf{אתה ובנך ובתך.}אלו קטנים; או אינו אלא גדולים? אמרת, הרי כבר מזהרים הם, אלא לא בא אלא להזהיר גדולים על שביתת הקטנים; וזהו ששנינו (שבת קכ"א), קטן שבא לכבות אין שומעין לו מפני ששביתתו עליך:}%
% ---------- פרק 20 | פסוק 11 ----------
 \textbf{יא} כִּ֣י שֵֽׁשֶׁת־יָמִים֩ עָשָׂ֨ה יְהֹוָ֜ה אֶת־הַשָּׁמַ֣יִם וְאֶת־הָאָ֗רֶץ אֶת־הַיָּם֙ וְאֶת־כׇּל־אֲשֶׁר־בָּ֔ם וַיָּ֖נַח בַּיּ֣וֹם הַשְּׁבִיעִ֑י עַל־כֵּ֗ן בֵּרַ֧ךְ יְהֹוָ֛ה אֶת־י֥וֹם הַשַּׁבָּ֖ת וַֽיְקַדְּשֵֽׁהוּ׃ \{ס\} \Rashi{\textbf{וינח ביום השביעי.}כביכול הכתיב בעצמו מנוחה, ללמד הימנו קל וחמר לאדם שמלאכתו בעמל ויגיעה שיהיה נח בשבת (מכילתא): ברך, ויקדשהו. ברכו במן – לכפלו בששי לחם משנה, וקדשו במן – שלא היה יורד בו (מכילתא):}%
% ---------- פרק 20 | פסוק 12 ----------
 \textbf{יב} כַּבֵּ֥ד אֶת־אָבִ֖יךָ וְאֶת־אִמֶּ֑ךָ לְמַ֙עַן֙ יַאֲרִכ֣וּן יָמֶ֔יךָ עַ֚ל הָאֲדָמָ֔ה אֲשֶׁר־יְהֹוָ֥ה אֱלֹהֶ֖יךָ נֹתֵ֥ן לָֽךְ׃ \{ס\} \Rashi{\textbf{למען יארכון ימיך.}אם תכבד יאריכון ואם לאו יקצרון, שדברי תורה נוטריקון הם נדרשים – מכלל הן לאו ומכלל לאו הן (מכילתא):}%
% ---------- פרק 20 | פסוק 13 ----------
 \textbf{יג} לֹ֥֖א תִּֿרְצָ֖͏ֽח׃ \{ס\}        לֹ֣֖א תִּֿנְאָ֑͏ֽף׃ \{ס\}        לֹ֣֖א תִּֿגְנֹֽ֔ב׃ \{ס\}        לֹֽא־תַעֲנֶ֥ה בְרֵעֲךָ֖ עֵ֥ד שָֽׁקֶר׃ \{ס\} \Rashi{\textbf{לא תנאף.}אין נאוף אלא באשת איש, שנאמר מות יומת הנואף והנואפת (ויקרא כ'), ואומר האשה המנאפת תחת אישה תקח את זרים (יחזקאל ט"ז): לא תגנב. בגונב נפשות הכתוב מדבר, לא תגנבו (ויקרא י"ט), בגונב ממון; או אינו אלא זה בגונב ממון ולהלן בגונב נפשות? אמרת דבר למד מענינו – מה לא תרצח, לא תנאף, מדבר בדבר שחיבין עליו מיתת בית דין, אף לא תגנב דבר שחיב עליו מיתת בית דין (סנהדרין פ"ו):}%
% ---------- פרק 20 | פסוק 14 ----------
 \textbf{יד} לֹ֥א תַחְמֹ֖ד בֵּ֣ית רֵעֶ֑ךָ \{ס\}        לֹֽא־תַחְמֹ֞ד אֵ֣שֶׁת רֵעֶ֗ךָ וְעַבְדּ֤וֹ וַאֲמָתוֹ֙ וְשׁוֹר֣וֹ וַחֲמֹר֔וֹ וְכֹ֖ל אֲשֶׁ֥ר לְרֵעֶֽךָ׃ \{פ\} %
% ---------- פרק 20 | פסוק 15 ----------
 \textbf{טו} וְכׇל־הָעָם֩ רֹאִ֨ים אֶת־הַקּוֹלֹ֜ת וְאֶת־הַלַּפִּידִ֗ם וְאֵת֙ ק֣וֹל הַשֹּׁפָ֔ר וְאֶת־הָהָ֖ר עָשֵׁ֑ן וַיַּ֤רְא הָעָם֙ וַיָּנֻ֔עוּ וַיַּֽעַמְד֖וּ מֵֽרָחֹֽק׃ \Rashi{\textbf{וכל העם ראים.}מלמד שלא היה בהם אחד סומא, ומנין שלא היה בהם אלם? ת"ל ויענו כל העם, ומנין שלא היה בהם חרש? ת"ל נעשה ונשמע (מכילתא): ראים את הקולת. רואין את הנשמע, שאי אפשר לראות במקום אחר: את הקולת. היוצאין מפי הגבורה: וינעו. אין נוע אלא זיע: ויעמדו מרחק. היו נרתעין לאחוריהם שנים עשר מיל, כארך מחניהם, ומלאכי השרת באין ומסיעין אותן להחזירם, שנאמר (תהילים ס"ח) מלאכי צבאות ידדון ידדון (שבת פ"ח):}%
% ---------- פרק 20 | פסוק 16 ----------
 \textbf{טז} וַיֹּֽאמְרוּ֙ אֶל־מֹשֶׁ֔ה דַּבֵּר־אַתָּ֥ה עִמָּ֖נוּ וְנִשְׁמָ֑עָה וְאַל־יְדַבֵּ֥ר עִמָּ֛נוּ אֱלֹהִ֖ים פֶּן־נָמֽוּת׃ %
% ---------- פרק 20 | פסוק 17 ----------
 \textbf{יז} וַיֹּ֨אמֶר מֹשֶׁ֣ה אֶל־הָעָם֮ אַל־תִּירָ֒אוּ֒ כִּ֗י לְבַֽעֲבוּר֙ נַסּ֣וֹת אֶתְכֶ֔ם בָּ֖א הָאֱלֹהִ֑ים וּבַעֲב֗וּר תִּהְיֶ֧ה יִרְאָת֛וֹ עַל־פְּנֵיכֶ֖ם לְבִלְתִּ֥י תֶחֱטָֽאוּ׃ \Rashi{\textbf{לבעבור נסות אתכם.}לגדל אתכם בעולם, שיצא לכם שם באמות שהוא בכבודו נגלה עליכם: נסות. לשון הרמה וגדלה, כמו הרימו נס (ישעיהו ס"ב), ארים נסי (שם מ"ט), כנס על הגבעה (שם ל') – שהוא זקוף: ובעבור תהיה יראתו. על ידי שראיתם אותו יראוי ומאים, תדעו כי אין זולתו ותיראו מפניו:}%
% ---------- פרק 20 | פסוק 18 ----------
 \textbf{יח} וַיַּעֲמֹ֥ד הָעָ֖ם מֵרָחֹ֑ק וּמֹשֶׁה֙ נִגַּ֣שׁ אֶל־הָֽעֲרָפֶ֔ל אֲשֶׁר־שָׁ֖ם הָאֱלֹהִֽים׃ \{ס\} \Rashi{\textbf{נגש אל הערפל.}לפנים משלש מחצות: חשך, ענן, וערפל; שנאמר וההר בער באש עד לב השמים חשך ענן וערפל (דברים ד'); ערפל הוא עב הענן שאמר לו הנה אנכי בא אליך בעב הענן (שמות י"ט):}%
% ---------- פרק 20 | פסוק 19 ----------
 \textbf{יט} וַיֹּ֤אמֶר יְהֹוָה֙ אֶל־מֹשֶׁ֔ה כֹּ֥ה תֹאמַ֖ר אֶל־בְּנֵ֣י יִשְׂרָאֵ֑ל אַתֶּ֣ם רְאִיתֶ֔ם כִּ֚י מִן־הַשָּׁמַ֔יִם דִּבַּ֖רְתִּי עִמָּכֶֽם׃ \Rashi{\textbf{כה תאמר.}בלשון הזה: אתם ראיתם. יש הפרש בין מה שאדם רואה למה שאחרים משיחין לו, שמה שאחרים משיחין לו פעמים שלבו חלוק מלהאמין (מכילתא): כי מן השמים דברתי. וכתוב אחד אומר וירד ה' על הר סיני (שמות י"ט), בא הכתוב השלישי והכריע ביניהם – מן השמים השמיעך את קולו ליסרך ועל הארץ הראך את אשו הגדולה (דברים ד'), כבודו בשמים ואשו וגבורתו על הארץ; ד"א: הרכין שמים ושמי השמים והציען על ההר, וכן הוא אומר ויט שמים וירד (תהילים י"ח):}%
% ---------- פרק 20 | פסוק 20 ----------
 \textbf{כ} לֹ֥א תַעֲשׂ֖וּן אִתִּ֑י אֱלֹ֤הֵי כֶ֙סֶף֙ וֵאלֹהֵ֣י זָהָ֔ב לֹ֥א תַעֲשׂ֖וּ לָכֶֽם׃ \Rashi{\textbf{לא תעשון אתי.}לא תעשון דמות שמשי המשמשים לפני במרום (מכילתא): אלהי כסף. בא להזהיר על הכרובים שאתה עושה לעמד אתי, שלא יהיו של כסף, שאם שניתם לעשותם של כסף הרי הן לפני כאלוהות: ואלהי זהב. בא להזהיר שלא יוסיף על שנים, שאם עשית ארבעה הרי הן לפני כאלהי זהב (מכילתא): לא תעשו לכם. לא תאמר הריני עושה כרובים בבתי כנסיות ובבתי מדרשות כדרך שאני עושה בבית עולמים, לכך נאמר לא תעשו לכם (שם):}%
% ---------- פרק 20 | פסוק 21 ----------
 \textbf{כא} מִזְבַּ֣ח אֲדָמָה֮ תַּעֲשֶׂה־לִּי֒ וְזָבַחְתָּ֣ עָלָ֗יו אֶת־עֹלֹתֶ֙יךָ֙ וְאֶת־שְׁלָמֶ֔יךָ אֶת־צֹֽאנְךָ֖ וְאֶת־בְּקָרֶ֑ךָ בְּכׇל־הַמָּקוֹם֙ אֲשֶׁ֣ר אַזְכִּ֣יר אֶת־שְׁמִ֔י אָב֥וֹא אֵלֶ֖יךָ וּבֵרַכְתִּֽיךָ׃ \Rashi{\textbf{מזבח אדמה.}מחבר באדמה, שלא יבננו על גבי עמודים או על גבי בסיס; ד"א שהיה ממלא את חלל מזבח הנחשת אדמה בשעת חניתן (שם): תעשה לי. שתהא תחלת עשיתו לשמי: וזבחת עליו. אצלו, כמו ועליו מטה מנשה (במדבר ב'), או אינו אלא עליו ממש? ת"ל הבשר והדם על מזבח ה' אלהיך (דברים י"ב), ואין שחיטה בראש המזבח: את עלתיך ואת שלמיך. אשר מצאנך ומבקרך: (את צאנך ואת בקרך. פרוש לאת עולתיך ואת שלמיך:) בכל המקום אשר אזכיר את שמי. אשר אתן לך רשות להזכיר שם המפרש שלי, שם אבוא אליך וברכתיך – אשרה שכינתי עליך; מכאן אתה למד שלא נתן רשות להזכיר שם המפרש אלא במקום שהשכינה באה שם, וזהו בית הבחירה, שם נתן רשות לכהנים להזכיר שם המפרש בנשיאת כפים לברך את העם (סוטה ל"ח):}%
% ---------- פרק 20 | פסוק 22 ----------
 \textbf{כב} וְאִם־מִזְבַּ֤ח אֲבָנִים֙ תַּֽעֲשֶׂה־לִּ֔י לֹֽא־תִבְנֶ֥ה אֶתְהֶ֖ן גָּזִ֑ית כִּ֧י חַרְבְּךָ֛ הֵנַ֥פְתָּ עָלֶ֖יהָ וַתְּחַֽלְלֶֽהָ׃ \Rashi{\textbf{ואם מזבח אבנים.}רבי ישמעאל אומר כל אם ואם שבתורה רשות חוץ מג', ואם מזבח אבנים תעשה לי, הרי אם זה משמש בלשון כאשר – וכאשר תעשה לי מזבח אבנים לא תבנה אתהן גזית, שהרי חובה עליך לבנות מזבח אבנים, שנאמר אבנים שלמות תבנה (דברים כ"ז) – וכן אם כסף תלוה (שמות כ"ב), חובה הוא, שנאמר והעבט תעביטנו (דברים ט"ו) – ואף זה משמש בלשון כאשר – וכן ואם תקריב מנחת בכורים (ויקרא ב'), זו מנחת העמר שהיא חובה, ועל כרחך אין אם הללו תלויין אלא ודאין, ובלשון כאשר הם משמשים: גזית. לשון גזיזה, שפוסלן ומכתתן בברזל: כי חרבך הנפת עליה. הרי כי זה משמש בלשון פן, שהוא דילמא – פן תניף חרבך עליה: ותחללה. הא למדת שאם הנפת עליה ברזל חללת, שהמזבח נברא להאריך ימיו של אדם והברזל נברא לקצר ימיו של אדם, אין זה בדין שיונף המקצר על המאריך (מכילתא). ועוד שהמזבח מטיל שלום בין ישראל לאביהם שבשמים, לפיכך לא יבא עליו כורת ומחבל; והרי דברים קל וחמר, ומה אבנים שאינן רואות ולא שומעות ולא מדברות, על ידי שמטילות שלום אמרה תורה ולא תניף עליהם ברזל, המטיל שלום בין איש לאשתו, בין משפחה למשפחה, בין אדם לחברו, על אחת כמה וכמה שלא תבואהו פרענות (שם):}%
% ---------- פרק 20 | פסוק 23 ----------
 \textbf{כג} וְלֹֽא־תַעֲלֶ֥ה בְמַעֲלֹ֖ת עַֽל־מִזְבְּחִ֑י אֲשֶׁ֛ר לֹֽא־תִגָּלֶ֥ה עֶרְוָתְךָ֖ עָלָֽיו׃ \{פ\} \Rashi{\textbf{ולא תעלה במעלת.}כשאתה בונה כבש למזבח לא תעשהו מעלות מעלות, אשקלו"נש בלעז, אלא חלק יהא ומשפע: אשר לא תגלה ערותך. שעל ידי המעלות אתה צריך להרחיב פסיעותיך; ואע"פ שאינו גלוי ערוה ממש, שהרי כתיב ועשה להם מכנסי בד, מכל מקום הרחבת הפסיעות קרוב לגלוי ערוה הוא, ואתה נוהג בם מנהג בזיון; והרי דברים ק"ו, ומה אבנים הללו שאין בהם דעת להקפיד על בזיונן, אמרה תורה הואיל ויש בהם צרך לא תנהג בהם מנהג בזיון, חברך שהוא בדמות יוצרך, ומקפיד על בזיונו, על אחת כמה וכמה:}%
\addparasha{חומש שמות | פרשת משפטים}
% ---------- פרק 21 | פסוק 1 ----------
 \textbf{א} וְאֵ֙לֶּה֙ הַמִּשְׁפָּטִ֔ים אֲשֶׁ֥ר תָּשִׂ֖ים לִפְנֵיהֶֽם׃ \Rashi{\textbf{ואלה המשפטים.}כל מקום שנאמר "אלה" פסל את הראשונים, "ואלה" מוסיף על הראשונים, מה הראשונים מסיני, אף אלו מסיני; ולמה נסמכה פרשת דינין לפרשת מזבח? לומר לך, שתשים סנהדרין אצל המקדש (מכילתא): אשר תשים לפניהם. אמר לו הקב"ה למשה: לא תעלה על דעתך לומר, אשנה להם הפרק וההלכה ב' או ג' פעמים, עד שתהא סדורה בפיהם כמשנתה, ואיני מטריח עצמי להבינם טעמי הדבר ופרושו, לכך נאמר אשר תשים לפניהם – כשלחן הערוך ומוכן לאכל לפני האדם (שם): לפניהם. ולא לפני גוים, ואפלו ידעת בדין אחד שהם דנין אותו כדיני ישראל, אל תביאהו בערכאות שלהם, שהמביא דיני ישראל לפני גוים, מחלל את השם ומיקר שם ע"ז להשביחה, שנאמר כי לא כצורנו צורם ואיבינו פלילים (דברים ל"ב) – כשאויבינו פלילים זהו עדות לעלוי יראתם (תנחומא):}%
% ---------- פרק 21 | פסוק 2 ----------
 \textbf{ב} כִּ֤י תִקְנֶה֙ עֶ֣בֶד עִבְרִ֔י שֵׁ֥שׁ שָׁנִ֖ים יַעֲבֹ֑ד וּבַ֨שְּׁבִעִ֔ת יֵצֵ֥א לַֽחׇפְשִׁ֖י חִנָּֽם׃ \Rashi{\textbf{כי תקנה עבד עברי.}עבד שהוא עברי; או אינו אלא עבדו של עברי – עבד כנעני שלקחתו מישראל – ועליו הוא אומר שש שנים יעבד? ומה אני מקים והתנחלתם אתם (ויקרא כ"ה), בלקוח מן הגוי, אבל בלקוח מישראל יצא בשש? ת"ל כי ימכר לך אחיך העברי (דברים ט"ו), לא אמרתי אלא באחיך: כי תקנה. מיד בית דין שמכרוהו בגנבתו, כמו שנאמר אם אין לו ונמכר בגנבתו, או אינו אלא במוכר עצמו מפני דחקו, אבל מכרוהו בית דין לא יצא בשש? כשהוא אומר וכי ימוך אחיך עמך ונמכר לך (ויקרא כ"ה), הרי מוכר עצמו מפני דחקו אמור, ומה אני מקים כי תקנה? בנמכר בבית דין: לחפשי. לחרות:}%
% ---------- פרק 21 | פסוק 3 ----------
 \textbf{ג} אִם־בְּגַפּ֥וֹ יָבֹ֖א בְּגַפּ֣וֹ יֵצֵ֑א אִם־בַּ֤עַל אִשָּׁה֙ ה֔וּא וְיָצְאָ֥ה אִשְׁתּ֖וֹ עִמּֽוֹ׃ \Rashi{\textbf{אם בגפו יבא.}שלא היה נשוי אשה, כתרגומו אם בלחודוהי, ולשון בגפו – בכנפו, שלא בא אלא כמות שהוא, יחידי, בתוך לבושו – בכנף בגדו: בגפו יצא. מגיד שאם לא היה נשוי מתחלה, אין רבו מוסר לו שפחה כנענית להוליד ממנה עבדים (קידושין כ'): אם בעל אשה הוא. ישראלית: ויצאה אשתו עמו. וכי מי הכניסה שתצא? אלא מגיד הכתוב שהקונה עבד עברי חיב במזונות אשתו ובניו (שם כ"ב):}%
% ---------- פרק 21 | פסוק 4 ----------
 \textbf{ד} אִם־אֲדֹנָיו֙ יִתֶּן־ל֣וֹ אִשָּׁ֔ה וְיָלְדָה־ל֥וֹ בָנִ֖ים א֣וֹ בָנ֑וֹת הָאִשָּׁ֣ה וִילָדֶ֗יהָ תִּהְיֶה֙ לַֽאדֹנֶ֔יהָ וְה֖וּא יֵצֵ֥א בְגַפּֽוֹ׃ \Rashi{\textbf{אם אדניו יתן לו אשה.}מכאן שהרשות ביד רבו למסר לו שפחה כנענית להוליד ממנה עבדים. או אינו אלא בישראלית? ת"ל האשה וילדיה תהיה לאדניה, הא אינו מדבר אלא בכנענית, שהרי העבריה אף היא יוצאה בשש – ואפלו לפני שש אם הביאה סימנין יוצאה – שנאמר אחיך העברי או העבריה (דברים ט"ו), מלמד שאף העבריה יוצאה בשש (מכילתא):}%
% ---------- פרק 21 | פסוק 5 ----------
 \textbf{ה} וְאִם־אָמֹ֤ר יֹאמַר֙ הָעֶ֔בֶד אָהַ֙בְתִּי֙ אֶת־אֲדֹנִ֔י אֶת־אִשְׁתִּ֖י וְאֶת־בָּנָ֑י לֹ֥א אֵצֵ֖א חׇפְשִֽׁי׃ \Rashi{\textbf{את אשתי.}השפחה:}%
% ---------- פרק 21 | פסוק 6 ----------
 \textbf{ו} וְהִגִּישׁ֤וֹ אֲדֹנָיו֙ אֶל־הָ֣אֱלֹהִ֔ים וְהִגִּישׁוֹ֙ אֶל־הַדֶּ֔לֶת א֖וֹ אֶל־הַמְּזוּזָ֑ה וְרָצַ֨ע אֲדֹנָ֤יו אֶת־אׇזְנוֹ֙ בַּמַּרְצֵ֔עַ וַעֲבָד֖וֹ לְעֹלָֽם׃ \{ס\} \Rashi{\textbf{אל האלהים.}לבית דין; צריך שימלך במוכריו שמכרוהו לו (שם): אל הדלת או אל המזוזה. יכול שתהא המזוזה כשרה לרצע עליה, ת"ל ונתתה באזנו ובדלת, בדלת ולא במזוזה, הא מה ת"ל או אל המזוזה? הקיש דלת למזוזה, מה מזוזה מעומד אף דלת מעומד (שם): ורצע אדניו את אזנו במרצע. הימנית, או אינו, אלא של שמאל? ת"ל אזן אזן לגזרה שוה, נאמר כאן ורצע אדניו את אזנו, ונאמר במצרע תנוך אזן המטהר הימנית (ויקרא י"ד), מה להלן הימנית, אף כאן הימנית; ומה ראה אזן לרצע מכל שאר אברים שבגוף? אמר רבן יוחנן בן זכאי: אזן זאת ששמעה על הר סיני לא תגנב, והלך וגנב, תרצע. ואם מוכר עצמו, אזן ששמעה על הר סיני כי לי בני ישראל עבדים, והלך וקנה אדון לעצמו, תרצע. רבי שמעון היה דורש מקרא זה כמין חמר: מה נשתנו דלת ומזוזה מכל כלים שבבית? אמר הקב"ה דלת ומזוזה שהיו עדים במצרים כשפסחתי על המשקוף ועל שתי המזוזות ואמרתי כי לי בני ישראל עבדים – עבדי הם ולא עבדים לעבדים – והלך זה וקנה אדון לעצמו, ירצע בפניהם (קידושין כ"ב): ועבדו לעלם. עד היובל; או אינו אלא לעולם כמשמעו? ת"ל ואיש אל משפחתו תשובו (ויקרא כ"ה), מגיד שחמשים שנה קרויים עולם; ולא שיהא עובדו כל חמשים שנה, אלא עובדו עד היובל, בין סמוך בין מפלג (קידושין ט"ו):}%
% ---------- פרק 21 | פסוק 7 ----------
 \textbf{ז} וְכִֽי־יִמְכֹּ֥ר אִ֛ישׁ אֶת־בִּתּ֖וֹ לְאָמָ֑ה לֹ֥א תֵצֵ֖א כְּצֵ֥את הָעֲבָדִֽים׃ \Rashi{\textbf{וכי ימכר איש את בתו לאמה.}בקטנה הכתוב מדבר; יכול אפלו הביאה סימנים, אמרת קל וחמר, ומה מכורה קדם לכן יוצאה בסימנין – כמו שכתוב ויצאה חנם אין כסף, שאנו דורשין אותו לסימני נערות – שאינה מכורה אינו דין שלא תמכר (ערכין כ"ט): לא תצא כצאת העבדים. כיציאת עבדים כנענים שיוצאים בשן ועין, אבל זו לא תצא בשן ועין אלא עובדת שש, או עד היובל, או עד שתביא סימנין, וכל הקודם קודם לחרותה, ונותן לה דמי עינה או דמי שנה. או אינו אלא לא תצא כצאת העבדים בשש וביובל? ת"ל כי ימכר לך אחיך העברי או העבריה, מקיש עבריה לעברי לכל יציאותיו, מה עברי יוצא בשש וביובל, אף עבריה יוצאה בשש וביובל, ומהו לא תצא כצאת העבדים? לא תצא בראשי אברים כעבדים כנעניים; יכול העברי יוצא בראשי אברים? ת"ל העברי או העבריה, מקיש עברי לעבריה, מה העבריה אינה יוצאה בראשי אברים, אף הוא אינו יוצא בראשי אברים (מכילתא):}%
% ---------- פרק 21 | פסוק 8 ----------
 \textbf{ח} אִם־רָעָ֞ה בְּעֵינֵ֧י אֲדֹנֶ֛יהָ אֲשֶׁר־[ל֥וֹ] (לא) יְעָדָ֖הּ וְהֶפְדָּ֑הּ לְעַ֥ם נׇכְרִ֛י לֹא־יִמְשֹׁ֥ל לְמׇכְרָ֖הּ בְּבִגְדוֹ־בָֽהּ׃ \Rashi{\textbf{אם רעה בעיני אדניה.}שלא נשאה חן בעיניו לכנסה: אשר לא יעדה. שהיה לו ליעדה ולהכניסה לו לאשה, וכסף קניתה הוא כסף קדושיה; וכאן רמז לך הכתוב שמצוה ביעוד, ורמז לך שאינה צריכה קדושין אחרים (קידושין י"ח): והפדה. יתן לה מקום להפדות ולצאת – שאף הוא מסיע בפדיונה; ומה הוא מקום שנותן לה? שמגרע מפדיונה במספר השנים שעשתה אצלו כאלו היא שכורה אצלו. כיצד? הרי שקנאה במנה ועשתה אצלו שתי שנים, אומרים לו, יודע היית שעתידה לצאת לסוף שש, נמצא שקנית עבודת כל שנה ושנה בששית המנה, ועשתה אצלך שתי שנים, הרי שלישית המנה, טל שני שלישיות המנה ותצא מאצלך: לעם נכרי לא ימשל למכרה. שאינו רשאי למכרה לאחר לא האדון ולא האב (שם): בבגדו בה. אם בא לבגד בה, שלא לקים בה מצות יעוד, וכן אביה מאחר שבגד בה ומכרה לזה:}%
% ---------- פרק 21 | פסוק 9 ----------
 \textbf{ט} וְאִם־לִבְנ֖וֹ יִֽיעָדֶ֑נָּה כְּמִשְׁפַּ֥ט הַבָּנ֖וֹת יַעֲשֶׂה־לָּֽהּ׃ \Rashi{\textbf{ואם לבנו ייעדנה.}האדון; מלמד שאף בנו קם תחתיו ליעדה אם ירצה אביו, ואינו צריך לקדשה קדושין אחרים, אלא אומר לה הרי את מיעדת לי בכסף שקבל אביך בדמיך: כמשפט הבנות. שאר כסות ועונה (מכילתא):}%
% ---------- פרק 21 | פסוק 10 ----------
 \textbf{י} אִם־אַחֶ֖רֶת יִֽקַּֽח־ל֑וֹ שְׁאֵרָ֛הּ כְּסוּתָ֥הּ וְעֹנָתָ֖הּ לֹ֥א יִגְרָֽע׃ \Rashi{\textbf{אם אחרת יקח לו.}עליה: שארה כסותה וענתה לא יגרע. מן האמה שיעד לו כבר: שארה. מזונות: כסותה. כמשמעו. ענתה. תשמיש (כתובות מ"ז):}%
% ---------- פרק 21 | פסוק 11 ----------
 \textbf{יא} וְאִ֨ם־שְׁלׇשׁ־אֵ֔לֶּה לֹ֥א יַעֲשֶׂ֖ה לָ֑הּ וְיָצְאָ֥ה חִנָּ֖ם אֵ֥ין כָּֽסֶף׃ \{ס\} \Rashi{\textbf{ואם שלש אלה לא יעשה לה.}אם אחת משלש אלה לא יעשה לה; ומה הן השלש? ייעדנה לו, או לבנו, או יגרע מפדיונה ותצא, וזה לא יעדה, לא לו ולא לבנו, והיא לא היה בידה לפדות את עצמה: ויצאה חנם. רבה לה יציאה לזו יותר ממה שרבה לעבדים, ומה היא היציאה? למדך שתצא בסימנין ותשהה עמו עוד עד שתביא סימנין, – ואם הגיעו שש שנים קדם סימנין, כבר למדנו שתצא, שנאמר העברי או העבריה ועבדך שש שנים (דברים ט"ו) – ומהו האמור כאן ויצאה חנם? שאם קדמו סימנים לשש שנים תצא בהן. או אינו אומר שתצא אלא בבגרות? ת"ל אין כסף לרבות יציאת בגרות; ואם לא נאמרו שניהם הייתי אומר ויצאה חנם זו בגרות, לכך נאמרו שניהם שלא לתן פתחון פה לבעל הדין לחלק (מכילתא):}%
% ---------- פרק 21 | פסוק 12 ----------
 \textbf{יב} מַכֵּ֥ה אִ֛ישׁ וָמֵ֖ת מ֥וֹת יוּמָֽת׃ \Rashi{\textbf{מכה איש ומת.}כמה כתובים נאמרו בפרשת רוצחין, ומה שבידי לפרש למה באו כלם, אפרש: מכה איש ומת. למה נאמר? לפי שנאמר ואיש כי יכה כל נפש אדם (ויקרא כ"ד), שומע אני הכאה בלא מיתה, תלמוד לומר מכה איש ומת – אינו חיב אלא בהכאה של מיתה; ואם נאמר מכה איש ולא נאמר ואיש כי יכה, הייתי אומר אינו חיב עד שיכה איש, הכה את האשה ואת הקטן מנין? ת"ל כי יכה כל נפש אדם – אפלו קטן ואפלו אשה; ועוד, אלו נאמר מכה איש, שומע אני אפלו קטן שהכה והרג יהא חיב, ת"ל ואיש כי יכה, ולא קטן שהכה; ועוד כי יכה כל נפש אדם אפלו נפלים במשמע, ת"ל מכה איש, אינו חיב עד שיכה בן קימא הראוי להיות איש (מכילתא):}%
% ---------- פרק 21 | פסוק 13 ----------
 \textbf{יג} וַאֲשֶׁר֙ לֹ֣א צָדָ֔ה וְהָאֱלֹהִ֖ים אִנָּ֣ה לְיָד֑וֹ וְשַׂמְתִּ֤י לְךָ֙ מָק֔וֹם אֲשֶׁ֥ר יָנ֖וּס שָֽׁמָּה׃ \{ס\} \Rashi{\textbf{ואשר לא צדה.}לא ארב לו ולא נתכון: צדה. לשון ארב, וכן הוא אומר "ואתה צדה את נפשי" (שמואל א כ"ד); ולא יתכן לומר צדה לשון "הצד ציד" (בראשית כ"ז), שצידת חיות אין נופל ה"א בפעל שלה, ושם דבר בה ציד, וזה שם דבר בו צדיה, ופעל שלו צודה, וזה הפעל שלו צד. ואומר אני פתרונו כתרגומו – ודלא כמן לה. ומנחם חברו בחלק צד ציד, ואין אני מודה לו; ואם יש לחברו באחת ממחלקות של צד, נחברנו בחלק "על צד תנשאו" (ישעיהו ס"ו), "צדה אורה" (שמואל א כ'), "ומלין לצד עלאה ימלל" (דניאל ז'), אף כאן אשר לא צדה – לא צדד למצא לו שום צד מיתה; ואף זה יש להרהר עליו, מכל מקום לשון אורב הוא: והאלהים אנה לידו. זמן לידו, לשון "לא תאנה אליך רעה" (תהילים צ"א), "לא יאנה לצדיק כל און" (משלי י"ב), "מתאנה הוא לי" (מלכים ב ה') – מזדמן למצא לי עלה: והאלהים אנה לידו. ולמה תצא זאת מלפניו? הוא שאמר דוד "כאשר יאמר משל הקדמני מרשעים יצא רשע" (שמואל א כ"ד); ומשל הקדמוני היא התורה, שהיא משל הקב"ה שהוא קדמונו של עולם; והיכן אמרה תורה מרשעים יצא רשע? והאלהים אנה לידו. במה הכתוב מדבר? בשני בני אדם, אחד הרג שוגג ואחד הרג מזיד, ולא היו עדים בדבר שיעידו, זה לא נהרג, וזה לא גלה, והקב"ה מזמנן לפנדק אחד, זה שהרג במזיד יושב תחת הסלם וזה שהרג שוגג עולה בסלם ונופל על זה שהרג במזיד והורגו, ועדים מעידים עליו ומחיבים אותו לגלות, נמצא זה שהרג בשוגג גולה, וזה שהרג במזיד נהרג (מכות י'): ושמתי לך מקום. אף במדבר, שינוס שמה, ואי זה מקום קולטו? זה מחנה לויה (מכילתא):}%
% ---------- פרק 21 | פסוק 14 ----------
 \textbf{יד} וְכִֽי־יָזִ֥ד אִ֛ישׁ עַל־רֵעֵ֖הוּ לְהׇרְג֣וֹ בְעׇרְמָ֑ה מֵעִ֣ם מִזְבְּחִ֔י תִּקָּחֶ֖נּוּ לָמֽוּת׃ \{ס\} \Rashi{\textbf{וכי יזד.}למה נאמר? לפי שנאמר מכה איש וגו' שומע אני אפלו גוי, ורופא שהמית ושליח בית דין שהמית במלקות ארבעים, והאב המכה את בנו והרב הרודה את תלמידו, והשוגג, ת"ל וכי יזד – ולא שוגג, על רעהו – ולא על גוי, להרגו בערמה – ולא שליח בית דין והרופא ורודה בנו ותלמידו, שאף על פי שהם מזידין אין מערימין (מכילתא): מעם מזבחי. אם היה כהן ורוצה לעבד עבודה תקחנו למות (יומא פ"ה):}%
% ---------- פרק 21 | פסוק 15 ----------
 \textbf{טו} וּמַכֵּ֥ה אָבִ֛יו וְאִמּ֖וֹ מ֥וֹת יוּמָֽת׃ \{ס\} \Rashi{\textbf{ומכה אביו ואמו.}לפי שלמדנו על החובל בחברו שהוא בתשלומין ולא במיתה, הצרך לומר על החובל באביו שהוא במיתה, ואינו חיב אלא בהכאה שיש בה חבורה (סנהדרין פ"ה): אביו ואמו. או זה או זה: מות יומת. בחנק (מכילתא):}%
% ---------- פרק 21 | פסוק 16 ----------
 \textbf{טז} וְגֹנֵ֨ב אִ֧ישׁ וּמְכָר֛וֹ וְנִמְצָ֥א בְיָד֖וֹ מ֥וֹת יוּמָֽת׃ \{ס\} \Rashi{\textbf{וגנב איש ומכרו.}למה נאמר? לפי שנאמר כי ימצא איש גונב נפש מאחיו (דברים כ"ד), אין לי אלא איש שגנב נפש, אשה או טמטום או אנדרוגינוס שגנבו מנין? ת"ל וגנב איש ומכרו. ולפי שנא' כאן וגנב איש, אין לי אלא גונב איש, גונב אשה מנין? ת"ל גנב נפש, לכך הצרכו שניהם, מה שחסר זה גלה זה (סנהדרין שם): ונמצא בידו. שראוהו עדים שגנבו ומכרו, ונמצא כבר בידו קדם מכירה (שם): מות יומת. בחנק; כל מיתה האמורה בתורה סתם חנק היא (והפסיק הענין וכתב וגנב איש בין מכה אביו ואמו למקלל אביו, ונראה לי הינו פלוגתא, דמר סבר מקשינן הכאה לקללה ומר סבר לא מקשינן (שם)):}%
% ---------- פרק 21 | פסוק 17 ----------
 \textbf{יז} וּמְקַלֵּ֥ל אָבִ֛יו וְאִמּ֖וֹ מ֥וֹת יוּמָֽת׃ \{ס\} \Rashi{\textbf{ומקלל אביו ואמו.}למה נאמר? לפי שהוא אומר איש איש אשר יקלל את אביו (ויקרא כ'), אין לי אלא איש שקלל את אביו, אשה שקללה את אביה מנין? ת"ל ומקלל אביו ואמו סתם, בין איש ובין אשה, אם כן למה נאמר איש אשר יקלל? להוציא את הקטן (מכילתא): מות יומת. בסקילה; וכל מקום שנאמר דמיו בו בסקילה, ובנין אב לכלם באבן ירגמו אתם דמיהם בם (ויקרא כ'), ובמקלל אביו ואמו נאמר דמיו בו:}%
% ---------- פרק 21 | פסוק 18 ----------
 \textbf{יח} וְכִֽי־יְרִיבֻ֣ן אֲנָשִׁ֔ים וְהִכָּה־אִישׁ֙ אֶת־רֵעֵ֔הוּ בְּאֶ֖בֶן א֣וֹ בְאֶגְרֹ֑ף וְלֹ֥א יָמ֖וּת וְנָפַ֥ל לְמִשְׁכָּֽב׃ \Rashi{\textbf{וכי יריבן אנשים.}למה נאמר? לפי שנאמר עין תחת עין, לא למדנו אלא דמי אבריו אבל שבת ורפוי לא למדנו, לכך נאמרה פרשה זו: ונפל למשכב. כתרגומו ויפל לבוטלן – לחלי שמבטלו ממלאכתו:}%
% ---------- פרק 21 | פסוק 19 ----------
 \textbf{יט} אִם־יָק֞וּם וְהִתְהַלֵּ֥ךְ בַּח֛וּץ עַל־מִשְׁעַנְתּ֖וֹ וְנִקָּ֣ה הַמַּכֶּ֑ה רַ֥ק שִׁבְתּ֛וֹ יִתֵּ֖ן וְרַפֹּ֥א יְרַפֵּֽא׃ \{ס\} \Rashi{\textbf{על משענתו.}על בריו וכחו (מכילתא): ונקה המכה. וכי תעלה על דעתך שיהרג זה שלא הרג? אלא למדך כאן שחובשים אותו עד שנראה אם יתרפא זה; וכן משמעו: כשקם זה והולך על משענתו אז נקה המכה, אבל עד שלא יקום זה לא נקה המכה (כתובות ל"ג): רק שבתו. בטול מלאכתו מחמת החלי; אם קטע ידו או רגלו, רואין בטול מלאכתו מחמת החלי כאלו שומר קשואין, שהרי אף לאחר החלי אינו ראוי למלאכת יד ורגל, והוא כבר נתן לו מחמת נזקו דמי ידו ורגלו, שנ' יד תחת יד רגל תחת רגל: ורפא ירפא. כתרגומו – ישלם שכר הרופא:}%
% ---------- פרק 21 | פסוק 20 ----------
 \textbf{כ} וְכִֽי־יַכֶּה֩ אִ֨ישׁ אֶת־עַבְדּ֜וֹ א֤וֹ אֶת־אֲמָתוֹ֙ בַּשֵּׁ֔בֶט וּמֵ֖ת תַּ֣חַת יָד֑וֹ נָקֹ֖ם יִנָּקֵֽם׃ \Rashi{\textbf{וכי יכה איש את עבדו או את אמתו.}בעבד כנעני הכתוב מדבר, או אינו אלא בעברי? ת"ל כי כספו הוא, מה כספו קנוי לו עולמית, אף עבד הקנוי לו עולמית; והרי היה בכלל מכה איש ומת? אלא בא הכתוב והוציאו מן הכלל, להיות נדון בדין יום או יומים, שאם לא מת תחת ידו ושהה מעת לעת פטור (מכילתא): בשבט. כשיש בו כדי להמית הכתוב מדבר. או אפלו אין בו כדי להמית? ת"ל בישראל ואם באבן יד אשר ימות בה הכהו (במדבר ל"ה), והלא דברים קל וחמר: מה ישראל חמור אין חיב עליו אלא אם כן הכהו בדבר שיש בו כדי להמית ועל אבר שהוא כדי למות בהכאה זו, עבד הקל לא כל שכן (מכילתא): נקם ינקם. מיתת סיף, וכן הוא אומר חרב נקמת נקם ברית (ויקרא כ"ו):}%
% ---------- פרק 21 | פסוק 21 ----------
 \textbf{כא} אַ֥ךְ אִם־י֛וֹם א֥וֹ יוֹמַ֖יִם יַעֲמֹ֑ד לֹ֣א יֻקַּ֔ם כִּ֥י כַסְפּ֖וֹ הֽוּא׃ \{ס\} \Rashi{\textbf{אך אם יום או יומים יעמד לא יקם.}אם על יום אחד הוא פטור על יומים לא כל שכן? אלא יום שהוא כיומים, ואיזה? זה מעת לעת (מכילתא): לא יקם כי כספו הוא. הא אחר שהכהו, אע"פ ששהה מעת לעת קדם שמת, חיב:}%
% ---------- פרק 21 | פסוק 22 ----------
 \textbf{כב} וְכִֽי־יִנָּצ֣וּ אֲנָשִׁ֗ים וְנָ֨גְפ֜וּ אִשָּׁ֤ה הָרָה֙ וְיָצְא֣וּ יְלָדֶ֔יהָ וְלֹ֥א יִהְיֶ֖ה אָס֑וֹן עָנ֣וֹשׁ יֵעָנֵ֗שׁ כַּֽאֲשֶׁ֨ר יָשִׁ֤ית עָלָיו֙ בַּ֣עַל הָֽאִשָּׁ֔ה וְנָתַ֖ן בִּפְלִלִֽים׃ \Rashi{\textbf{וכי ינצו אנשים.}זה עם זה, ונתכון להכות את חברו והכה את האשה: ונגפו. אין נגיפה אלא לשון דחיפה והכאה, כמו פן תגף באבן רגלך (תהילים צ"א), בטרם יתנגפו רגליכם (ירמיהו י"ג) ולאבן נגף (ישעיהו ח'): ולא יהיה אסון. באשה: ענוש יענש. לשלם דמי ולדות לבעל; שמין אותה כמה היתה ראויה למכר בשוק להעלות בדמיה בשביל הריונה: ענש יענש. יגבו ממון ממנו, כמו וענשו אתו מאה כסף (דברים כ"ב): כאשר ישית עליו וגו'. כשיתבענו הבעל בבית דין להשית עליו ענש על כך: ונתן. המכה דמי ולדות: בפללים – על פי הדינים (מכילתא):}%
% ---------- פרק 21 | פסוק 23 ----------
 \textbf{כג} וְאִם־אָס֖וֹן יִהְיֶ֑ה וְנָתַתָּ֥ה נֶ֖פֶשׁ תַּ֥חַת נָֽפֶשׁ׃ \Rashi{\textbf{ואם אסון יהיה.}באשה: ונתתה נפש תחת נפש. רבותינו חולקים בדבר, יש אומרים נפש ממש, ויש אומרים ממון אבל לא נפש ממש, שהמתכון להרג את זה והרג את זה פטור ממיתה, ומשלם ליורשיו דמיו כמו שהיה נמכר בשוק (שם):}%
% ---------- פרק 21 | פסוק 24 ----------
 \textbf{כד} עַ֚יִן תַּ֣חַת עַ֔יִן שֵׁ֖ן תַּ֣חַת שֵׁ֑ן יָ֚ד תַּ֣חַת יָ֔ד רֶ֖גֶל תַּ֥חַת רָֽגֶל׃ \Rashi{\textbf{עין תחת עין.}סמא עין חברו נותן לו דמי עינו כמה שפחתו דמיו למכר בשוק, וכן כלם; ולא נטילת אבר ממש, כמו שדרשו רבותינו בפרק החובל (בבא קמא דף פ"ג):}%
% ---------- פרק 21 | פסוק 25 ----------
 \textbf{כה} כְּוִיָּה֙ תַּ֣חַת כְּוִיָּ֔ה פֶּ֖צַע תַּ֣חַת פָּ֑צַע חַבּוּרָ֕ה תַּ֖חַת חַבּוּרָֽה׃ \{ס\} \Rashi{\textbf{כויה תחת כויה.}מכות אש; ועד עכשו דבר בחבלה שיש בה פחת דמים, ועכשיו בשאין בה פחת דמים אלא צער, כגון כואו בשפוד על צפרניו, אומדים כמה אדם כיוצא בזה רוצה לטל להיות מצטער כך (בבא קמא פ"ג): פצע. היא מכה המוציאה דם, שפצע את בשרו, נבר"ורא בלעז; הכל לפי מה שהוא – אם יש בו פחת דמים, נותן נזק, ואם נפל למשכב, נותן שבת ורפוי ובשת וצער; ומקרא זה יתר הוא, ובהחובל דרשוהו רבותינו לחיב על הצער אפלו במקום נזק – שאף על פי שנותן לו דמי ידו, אין פוטרים אותו מן הצער, לומר, הואיל וקנה ידו יש עליו לחתכה בכל מה שירצה, אלא אומרים, יש לו לחתכה בסם, שאינו מצטער כל כך, וזה חתכה בברזל וצערו: חבורה. היא מכה שהדם נצרר בה ואינו יוצא, אלא שמאדים הבשר כנגדו, ולשון חבורה טק"א בלעז, כמו ונמר חברברתיו (ירמיהו י"ג), ותרגומו משקופי, לשון חבטה, בטדו"רא בלעז, וכן שדופות קדים (בראשית מ"א), שקיפן קדום – חבוטות ברוח, וכן על המשקוף, על שם שהדלת נוקש עליו:}%
% ---------- פרק 21 | פסוק 26 ----------
 \textbf{כו} וְכִֽי־יַכֶּ֨ה אִ֜ישׁ אֶת־עֵ֥ין עַבְדּ֛וֹ אֽוֹ־אֶת־עֵ֥ין אֲמָת֖וֹ וְשִֽׁחֲתָ֑הּ לַֽחׇפְשִׁ֥י יְשַׁלְּחֶ֖נּוּ תַּ֥חַת עֵינֽוֹ׃ \Rashi{\textbf{את עין עבדו.}כנעני, אבל עברי אינו יוצא בשן ועין, כמו שאמרנו לא תצא כצאת העבדים: תחת עינו. וכן בכ"ד ראשי אברים: אצבעות הידים והרגלים, ושתי אזנים, והחטם, וראש הגויה שהוא גיד האמה. ולמה נאמר שן ועין? שאם נאמר עין ולא נאמר שן, הייתי אומר מה עין שנברא עמו אף כל שנברא עמו, והרי שן לא נברא עמו; ואם נאמר שן ולא נאמר עין, הייתי אומר אפלו שן תינוק שיש לה חליפין, לכך נאמר עין (מכילתא):}%
% ---------- פרק 21 | פסוק 27 ----------
 \textbf{כז} וְאִם־שֵׁ֥ן עַבְדּ֛וֹ אֽוֹ־שֵׁ֥ן אֲמָת֖וֹ יַפִּ֑יל לַֽחׇפְשִׁ֥י יְשַׁלְּחֶ֖נּוּ תַּ֥חַת שִׁנּֽוֹ׃ \{פ\} %
% ---------- פרק 21 | פסוק 28 ----------
 \textbf{כח} וְכִֽי־יִגַּ֨ח שׁ֥וֹר אֶת־אִ֛ישׁ א֥וֹ אֶת־אִשָּׁ֖ה וָמֵ֑ת סָק֨וֹל יִסָּקֵ֜ל הַשּׁ֗וֹר וְלֹ֤א יֵאָכֵל֙ אֶת־בְּשָׂר֔וֹ וּבַ֥עַל הַשּׁ֖וֹר נָקִֽי׃ \Rashi{\textbf{וכי יגח שור.}אחד שור ואחד כל בהמה וחיה ועוף, אלא שדבר הכתוב בהוה (בבא קמא נ"ד): ולא יאכל את בשרו. ממשמע שנאמר סקול יסקל השור איני יודע שהוא נבלה ונבלה אסורה באכילה? אלא מה תלמוד לומר ולא יאכל את בשרו? שאפלו שחטו לאחר שנגמר דינו אסור באכילה, בהנאה מנין? ת"ל ובעל השור נקי, כאדם האומר לחברו יצא פלוני נקי מנכסיו ואין לו בהם הנאה של כלום, זהו מדרשו (בבא קמא מלכים א). ופשוטו כמשמעו, לפי שנאמר במועד וגם בעליו יומת, הצרך לומר בתם ובעל השור נקי:}%
% ---------- פרק 21 | פסוק 29 ----------
 \textbf{כט} וְאִ֡ם שׁוֹר֩ נַגָּ֨ח ה֜וּא מִתְּמֹ֣ל שִׁלְשֹׁ֗ם וְהוּעַ֤ד בִּבְעָלָיו֙ וְלֹ֣א יִשְׁמְרֶ֔נּוּ וְהֵמִ֥ית אִ֖ישׁ א֣וֹ אִשָּׁ֑ה הַשּׁוֹר֙ יִסָּקֵ֔ל וְגַם־בְּעָלָ֖יו יוּמָֽת׃ \Rashi{\textbf{מתמל שלשם.}הרי שלש נגיחות (מכילתא): והועד בבעליו. לשון התראה בעדים, כמו העד העד בנו האיש (בראשית מ"ג): והמית איש וגו'. לפי שנאמר כי יגח אין לי אלא שהמיתו בנגיחה, המיתו בנשיכה, דחיפה, בעיטה, מנין? תלמוד לומר והמית: וגם בעליו יומת. בידי שמים. יכול בידי אדם? תלמוד לומר מות יומת המכה רוצח הוא, על רציחתו אתה הורגו ואי אתה הורגו על רציחת שורו (סנהדרין ט"ו):}%
% ---------- פרק 21 | פסוק 30 ----------
 \textbf{ל} אִם־כֹּ֖פֶר יוּשַׁ֣ת עָלָ֑יו וְנָתַן֙ פִּדְיֹ֣ן נַפְשׁ֔וֹ כְּכֹ֥ל אֲשֶׁר־יוּשַׁ֖ת עָלָֽיו׃ \Rashi{\textbf{אם כפר יושת עליו.}"אם" זה אינו תלוי, והרי הוא כמו "אם כסף תלוה" – לשון אשר, זה משפטו שישיתו עליו בית דין כפר: ונתן פדין נפשו. דמי נזק דברי ר' ישמעאל, ר' עקיבא אומר דמי מזיק (מכילתא):}%
% ---------- פרק 21 | פסוק 31 ----------
 \textbf{לא} אוֹ־בֵ֥ן יִגָּ֖ח אוֹ־בַ֣ת יִגָּ֑ח כַּמִּשְׁפָּ֥ט הַזֶּ֖ה יֵעָ֥שֶׂה לּֽוֹ׃ \Rashi{\textbf{או בן יגח.}בן שהוא קטן: או בת. שהיא קטנה; לפי שנאמר והמית איש או אשה, יכול אינו חיב אלא על הגדולים, תלמוד לומר או בן יגח וגו', לחיב על הקטנים כגדולים (שם):}%
% ---------- פרק 21 | פסוק 32 ----------
 \textbf{לב} אִם־עֶ֛בֶד יִגַּ֥ח הַשּׁ֖וֹר א֣וֹ אָמָ֑ה כֶּ֣סֶף ׀ שְׁלֹשִׁ֣ים שְׁקָלִ֗ים יִתֵּן֙ לַֽאדֹנָ֔יו וְהַשּׁ֖וֹר יִסָּקֵֽל׃ \{ס\} \Rashi{\textbf{אם עבד, או אמה.}כנעניים: שלשים שקלים יתן. גזרת הכתוב הוא, בין שהוא שוה אלף זוז בין שאינו שוה אלא דינר. והשקל משקלו ארבעה זהובים שהם חצי אנקיא למשקל הישר של קולוני"א:}%
% ---------- פרק 21 | פסוק 33 ----------
 \textbf{לג} וְכִֽי־יִפְתַּ֨ח אִ֜ישׁ בּ֗וֹר א֠וֹ כִּֽי־יִכְרֶ֥ה אִ֛ישׁ בֹּ֖ר וְלֹ֣א יְכַסֶּ֑נּוּ וְנָֽפַל־שָׁ֥מָּה שּׁ֖וֹר א֥וֹ חֲמֽוֹר׃ \Rashi{\textbf{וכי יפתח איש בור.}שהיה מכסה וגלהו: או כי יכרה. למה נאמר? אם על הפתיחה חיב על הכריה לא כל שכן? אלא להביא כורה אחר כורה שהוא חיב (בבא קמא נ"א): ולא יכסנו. הא אם כסהו פטור; ובחופר ברשות הרבים דבר הכתוב: שור או חמור. הוא הדין לכל בהמה וחיה, שבכל מקום שנאמר שור או חמור אנו למדין אותו שור שור משבת, שנאמר למען ינוח שורך וחמרך (שמות כ"ג), מה להלן כל בהמה וחיה כשור – שהרי נאמר במקום אחר וכל בהמתך (דברים ה') – אף כאן כל בהמה וחיה כשור, ולא נאמר שור וחמור אלא שור ולא אדם, חמור ולא כלים (בבא קמא נ"ג):}%
% ---------- פרק 21 | פסוק 34 ----------
 \textbf{לד} בַּ֤עַל הַבּוֹר֙ יְשַׁלֵּ֔ם כֶּ֖סֶף יָשִׁ֣יב לִבְעָלָ֑יו וְהַמֵּ֖ת יִֽהְיֶה־לּֽוֹ׃ \{ס\} \Rashi{\textbf{בעל הבור.}בעל התקלה; אע"פ שאין הבור שלו – שעשאו ברשות הרבים – עשאו הכתוב בעליו להתחיב עליו בנזקיו: כסף ישיב לבעליו. ישיב לרבות שוה כסף ואפלו סבין (שם ז'): והמת יהיה לו. לנזק; שמין את הנבלה ונוטלה בדמים ומשלם לו המזיק עליה תשלומי נזקו (בבא קמא י'):}%
% ---------- פרק 21 | פסוק 35 ----------
 \textbf{לה} וְכִֽי־יִגֹּ֧ף שֽׁוֹר־אִ֛ישׁ אֶת־שׁ֥וֹר רֵעֵ֖הוּ וָמֵ֑ת וּמָ֨כְר֜וּ אֶת־הַשּׁ֤וֹר הַחַי֙ וְחָצ֣וּ אֶת־כַּסְפּ֔וֹ וְגַ֥ם אֶת־הַמֵּ֖ת יֶֽחֱצֽוּן׃ \Rashi{\textbf{וכי יגף.}ידחף; בין בקרניו, בין בגופו, בין ברגלו, בין שנשכו בשניו, כלן בכלל נגיפה הם, שאין נגיפה אלא לשון מכה: שור איש. שור של איש: ומכרו את השור וגו'. בשוים הכתוב מדבר – שור שוה מאתים שהמית שור שוה מאתים – בין שהנבלה שוה הרבה, בין שהיא שוה מעט, כשנוטל זה חצי החי וחצי המת וזה חצי החי וחצי המת, נמצא כל אחד מפסיד חצי נזק שהזיקה המיתה; למדנו שהתם משלם חצי נזק, שמן השוין אתה למד לשאינן שוין, כי דין התם לשלם חצי נזק – לא פחות ולא יותר. או יכול אף בשאינן שוין בדמיהן כשהן חיים אמר הכתוב יחצו את שניהם? אם אמרת כן, פעמים שהמזיק משתכר הרבה, כשהנבלה שוה למכר לנכרים הרבה יותר מדמי שור המזיק, ואי אפשר שיאמר הכתוב שיהא המזיק נשכר; או פעמים שהנזק נוטל הרבה יותר מדמי נזק שלם – שחצי דמי שור המזיק שוין יותר מכל דמי שור הנזק – ואם אמרת כן, הרי תם חמור ממועד. על כרחך לא דבר הכתוב אלא בשוין, ולמדך שהתם משלם חצי נזק, ומן השוין תלמד לשאינן שוין, שהמשתלם חצי נזקו שמין לו את הנבלה, ומה שפחתו דמיו בשביל המיתה, נוטל חצי הפחת והולך. ולמה אמר הכתוב בלשון הזה ולא אמר ישלם חציו? ללמד, שאין התם משלם אלא מגופו, ואם נגח ומת, אין נזק נוטל אלא הנבלה, ואם אינה מגעת לחצי נזקו יפסיד; או שור שוה מנה, שנגח שור שוה חמש מאות זוז, אינו נוטל אלא את השור, שלא נתחיב התם לחיב את בעליו לשלם מן העליה (בבא קמא ל"ג):}%
% ---------- פרק 21 | פסוק 36 ----------
 \textbf{לו} א֣וֹ נוֹדַ֗ע כִּ֠י שׁ֣וֹר נַגָּ֥ח הוּא֙ מִתְּמ֣וֹל שִׁלְשֹׁ֔ם וְלֹ֥א יִשְׁמְרֶ֖נּוּ בְּעָלָ֑יו שַׁלֵּ֨ם יְשַׁלֵּ֥ם שׁוֹר֙ תַּ֣חַת הַשּׁ֔וֹר וְהַמֵּ֖ת יִֽהְיֶה־לּֽוֹ׃ \{ס\} \Rashi{\textbf{או נודע.}או לא היה תם, אלא נודע כי שור נגח הוא, היום ומתמול שלשם הרי שלש נגיחות: שלם ישלם שור. נזק שלם: והמת יהיה לו. לנזק; ועליו ישלים המזיק עד שישתלם נזק כל נזקו:}%
% ---------- פרק 21 | פסוק 37 ----------
 \textbf{לז} כִּ֤י יִגְנֹֽב־אִישׁ֙ שׁ֣וֹר אוֹ־שֶׂ֔ה וּטְבָח֖וֹ א֣וֹ מְכָר֑וֹ חֲמִשָּׁ֣ה בָקָ֗ר יְשַׁלֵּם֙ תַּ֣חַת הַשּׁ֔וֹר וְאַרְבַּע־צֹ֖אן תַּ֥חַת הַשֶּֽׂה׃ \Rashi{\textbf{חמשה בקר וגו'.}אמר ר' יוחנן בן זכאי חס המקום על כבודן של בריות – שור שהולך ברגליו ולא נתבזה בו הגנב לנשאו על כתפו, משלם חמשה, שה שנושאו על כתפו, משלם ארבעה הואיל ונתבזה בו. אמר רבי מאיר בא וראה כמה גדול כחה של מלאכה – שור שבטלו ממלאכתו חמשה, שה שלא בטלו ממלאכתו ארבעה (מכילתא, בבא קמא ע"ט): תחת השור, תחת השה. שנאן הכתוב, לומר שאין מדת תשלומי ארבעה וחמשה נוהגת אלא בשור ושה בלבד (בבא קמא ס"ז):}%
% ---------- פרק 22 | פסוק 1 ----------
 \textbf{א} אִם־בַּמַּחְתֶּ֛רֶת יִמָּצֵ֥א הַגַּנָּ֖ב וְהֻכָּ֣ה וָמֵ֑ת אֵ֥ין ל֖וֹ דָּמִֽים׃ \Rashi{\textbf{אם במחתרת.}כשהיה חותר את הבית: אין לו דמים. אין זו רציחה, הרי הוא כמת מעקרו; כאן למדתך תורה "אם בא להרגך, השכם להרגו" וזה להרגך בא, שהרי יודע הוא שאין אדם מעמיד עצמו ורואה שנוטלין ממונו בפניו ושותק, לפיכך על מנת כן בא שאם יעמד בעל הממון כנגדו יהרגנו (סנהדרין ע"ב):}%
% ---------- פרק 22 | פסוק 2 ----------
 \textbf{ב} אִם־זָרְחָ֥ה הַשֶּׁ֛מֶשׁ עָלָ֖יו דָּמִ֣ים ל֑וֹ שַׁלֵּ֣ם יְשַׁלֵּ֔ם אִם־אֵ֣ין ל֔וֹ וְנִמְכַּ֖ר בִּגְנֵבָתֽוֹ׃ \Rashi{\textbf{אם זרחה השמש עליו.}אין זה אלא כמין משל: אם ברור לך הדבר שיש לו שלום עמך, כשמש הזה שהוא שלום בעולם, כך פשוט לך שאינו בא להרג אפלו יעמד בעל הממון כנגדו, כגון אב החותר לגנב ממון הבן, בידוע שרחמי האב על הבן ואינו בא על עסקי נפשות (שם): דמים לו. כחי הוא חשוב, ורציחה היא אם יהרגנו בעל הבית: שלם ישלם. הגנב ממון שגנב ואינו חיב מיתה; ואנקלוס שתרגם אם עינא דסהדיא נפלת עלוהי, לקח לו שטה אחרת, לומר, שאם מצאוהו עדים קדם שבא בעל הבית, וכשבא בעל הבית נגדו התרו בו שלא יהרגהו, דמים לו – חיב עליו אם הרגו, שמאחר שיש רואים לו, אין הגנב הזה בא על עסקי נפשות, ולא יהרג את בעל הממון:}%
% ---------- פרק 22 | פסוק 3 ----------
 \textbf{ג} אִֽם־הִמָּצֵא֩ תִמָּצֵ֨א בְיָד֜וֹ הַגְּנֵבָ֗ה מִשּׁ֧וֹר עַד־חֲמ֛וֹר עַד־שֶׂ֖ה חַיִּ֑ים שְׁנַ֖יִם יְשַׁלֵּֽם׃ \{ס\} \Rashi{\textbf{המצא תמצא.}ברשותו, שלא טבח ולא מכר: משור עד חמור. כל דבר בכלל תשלומי כפל, בין שיש בו רוח חיים, בין שאין בו רוח חיים, שהרי נאמר במקרא אחר "על שה על שלמה על כל אבדה … ישלם שנים לרעהו" (בבא קמא ס"ב): חיים שנים ישלם. ולא ישלם לו מתים, אלא חיים או דמי חיים (מכילתא):}%
% ---------- פרק 22 | פסוק 4 ----------
 \textbf{ד} כִּ֤י יַבְעֶר־אִישׁ֙ שָׂדֶ֣ה אוֹ־כֶ֔רֶם וְשִׁלַּח֙ אֶת־בְּעִירֹ֔ה וּבִעֵ֖ר בִּשְׂדֵ֣ה אַחֵ֑ר מֵיטַ֥ב שָׂדֵ֛הוּ וּמֵיטַ֥ב כַּרְמ֖וֹ יְשַׁלֵּֽם׃ \{ס\} \Rashi{\textbf{כי יבער, את בעירה ובער.}כלם לשון בהמה, כמו אנחנו ובעירנו (במדבר כ'): כי יבער. יוליך בהמותיו בשדה או בכרם של חברו, ויזיק אותו באחת משתי אלו, או בשלוח בעירו או בבעור, ופרשו רבותינו "ושלח" הוא נזקי מדרך כף רגל, "ובער" הוא נזקי השן, האוכלת ומבערת (בבא קמא ב'): בשדה אחר. בשדה של איש אחר: מיטב שדהו … ישלם. שמין את הנזק, ואם בא לשלם לו קרקע דמי נזקו, ישלם לו ממיטב שדותיו – אם היה נזקו סלע, יתן לו שוה סלע מעדית שיש לו; למדך הכתוב שהנזקין שמין להם בעדית (מכילתא, בבא קמא ו'):}%
% ---------- פרק 22 | פסוק 5 ----------
 \textbf{ה} כִּֽי־תֵצֵ֨א אֵ֜שׁ וּמָצְאָ֤ה קֹצִים֙ וְנֶאֱכַ֣ל גָּדִ֔ישׁ א֥וֹ הַקָּמָ֖ה א֣וֹ הַשָּׂדֶ֑ה שַׁלֵּ֣ם יְשַׁלֵּ֔ם הַמַּבְעִ֖ר אֶת־הַבְּעֵרָֽה׃ \{ס\} \Rashi{\textbf{כי תצא אש.}אפלו מעצמה (בבא קמא כ"ב): ומצאה קצים. קרדו"נש בלעז: ונאכל גדיש. שלחכה בקוצים עד שהגיעה לגדיש או לקמה המחברת בקרקע: או השדה. שלחכה את נירו וצריך לניר אותה פעם שניה (בבא קמא ס'): שלם ישלם המבער. אע"פ שהדליק בתוך שלו, והיא יצאה מעצמה על ידי קוצים שמצאה, חיב לשלם, לפי שלא שמר את גחלתו שלא תצא ותזיק:}%
% ---------- פרק 22 | פסוק 6 ----------
 \textbf{ו} כִּֽי־יִתֵּן֩ אִ֨ישׁ אֶל־רֵעֵ֜הוּ כֶּ֤סֶף אֽוֹ־כֵלִים֙ לִשְׁמֹ֔ר וְגֻנַּ֖ב מִבֵּ֣ית הָאִ֑ישׁ אִם־יִמָּצֵ֥א הַגַּנָּ֖ב יְשַׁלֵּ֥ם שְׁנָֽיִם׃ \Rashi{\textbf{וגנב מבית האיש.}לפי דבריו (שם ס"ג): אם ימצא הגנב ישלם. הגנב שנים לבעלים:}%
% ---------- פרק 22 | פסוק 7 ----------
 \textbf{ז} אִם־לֹ֤א יִמָּצֵא֙ הַגַּנָּ֔ב וְנִקְרַ֥ב בַּֽעַל־הַבַּ֖יִת אֶל־הָֽאֱלֹהִ֑ים אִם־לֹ֥א שָׁלַ֛ח יָד֖וֹ בִּמְלֶ֥אכֶת רֵעֵֽהוּ׃ \Rashi{\textbf{אם לא ימצא הגנב.}ובא השומר הזה, שהוא בעל הבית: ונקרב. אל הדינין, לדון עם זה ולשבע לו שלא שלח ידו בשלו (שם):}%
% ---------- פרק 22 | פסוק 8 ----------
 \textbf{ח} עַֽל־כׇּל־דְּבַר־פֶּ֡שַׁע עַל־שׁ֡וֹר עַל־חֲ֠מ֠וֹר עַל־שֶׂ֨ה עַל־שַׂלְמָ֜ה עַל־כׇּל־אֲבֵדָ֗ה אֲשֶׁ֤ר יֹאמַר֙ כִּי־ה֣וּא זֶ֔ה עַ֚ד הָֽאֱלֹהִ֔ים יָבֹ֖א דְּבַר־שְׁנֵיהֶ֑ם אֲשֶׁ֤ר יַרְשִׁיעֻן֙ אֱלֹהִ֔ים יְשַׁלֵּ֥ם שְׁנַ֖יִם לְרֵעֵֽהוּ׃ \{ס\} \Rashi{\textbf{על כל דבר פשע.}שימצא שקרן בשבועתו – שיעידו עדים שהוא עצמו גנבו – וירשיעוהו אלהים על פי העדים. ישלם שנים לרעהו. למדך הכתוב, שהטוען בפקדון לומר נגנב הימנו, ונמצא שהוא עצמו גנבו, משלם תשלומי כפל; ואימתי? בזמן שנשבע ואח"כ באו עדים, שכך דרשו רבותינו: ונקרב בעל הבית אל האלהים, קריבה זו שבועה היא, אתה אומר לשבועה או אינו אלא לדין, שכיון שבא לדין וכפר לומר נגנבה, מיד יתחיב כפל אם באו עדים שהוא בידו? נאמר כאן שליחות יד ונאמר למטה שליחות יד – שבועת ה' תהיה בין שניהם אם לא שלח ידו – מה להלן שבועה, אף כאן שבועה (שם): אשר יאמר כי הוא זה. לפי פשוטו אשר יאמר העד כי הוא זה שנשבעת עליו, הרי הוא אצלך, עד הדינין יבא דבר שניהם, ויחקרו את העדים, אם כשרים הם וירשיעוהו לשומר זה, ישלם שנים, ואם ירשיעו את העדים, שנמצאו זוממין, ישלמו הם שנים לשומר. ורבותינו ז"ל דרשו כי הוא זה ללמד, שאין מחיבין אותו שבועה אלא אם כן הודה במקצת, לומר כך וכך אני חיב לך, והמותר נגנב ממני (שם ק"ז):}%
% ---------- פרק 22 | פסוק 9 ----------
 \textbf{ט} כִּֽי־יִתֵּן֩ אִ֨ישׁ אֶל־רֵעֵ֜הוּ חֲמ֨וֹר אוֹ־שׁ֥וֹר אוֹ־שֶׂ֛ה וְכׇל־בְּהֵמָ֖ה לִשְׁמֹ֑ר וּמֵ֛ת אוֹ־נִשְׁבַּ֥ר אוֹ־נִשְׁבָּ֖ה אֵ֥ין רֹאֶֽה׃ \Rashi{\textbf{כי יתן איש אל רעהו חמור או שור.}פרשה ראשונה נאמרה בשומר חנם, לפיכך פטר בו את הגנבה, כמו שנאמר (פסוק ו) וגנב מבית האיש, אם לא ימצא הגנב ונקרב בעל הבית לשבועה, למדת שפוטר עצמו בשבועה זו. ופרשה זו אמורה בשומר שכר, לפיכך אינו פטור אם נגנבה, כמו שכתוב (פסוק יא) אם גנב יגנב מעמו ישלם, אבל על האנס, כגון מת מעצמו, או נשבר או נשבה בחזקה על ידי לסטים, ואין רואה שיעיד בדבר:}%
% ---------- פרק 22 | פסוק 10 ----------
 \textbf{י} שְׁבֻעַ֣ת יְהֹוָ֗ה תִּהְיֶה֙ בֵּ֣ין שְׁנֵיהֶ֔ם אִם־לֹ֥א שָׁלַ֛ח יָד֖וֹ בִּמְלֶ֣אכֶת רֵעֵ֑הוּ וְלָקַ֥ח בְּעָלָ֖יו וְלֹ֥א יְשַׁלֵּֽם׃ \Rashi{\textbf{שבעת ה' תהיה.}ישבע שכן הוא כדבריו, והוא לא שלח בה יד להשתמש בה לעצמו, שאם שלח בה יד ואח"כ נאנסה, חיב באנסים (בבא מציעא צ"ד): ולקח בעליו. השבועה: ולא ישלם. לו השומר כלום (בבא קמא ק"ו):}%
% ---------- פרק 22 | פסוק 11 ----------
 \textbf{יא} וְאִם־גָּנֹ֥ב יִגָּנֵ֖ב מֵעִמּ֑וֹ יְשַׁלֵּ֖ם לִבְעָלָֽיו׃ %
% ---------- פרק 22 | פסוק 12 ----------
 \textbf{יב} אִם־טָרֹ֥ף יִטָּרֵ֖ף יְבִאֵ֣הוּ עֵ֑ד הַטְּרֵפָ֖ה לֹ֥א יְשַׁלֵּֽם׃ \{פ\} \Rashi{\textbf{אם טרף יטרף.}ע"י חיה רעה: יביאהו עד. יביא עדים שנטרפה באנס ופטור: הטרפה לא ישלם. אינו אומר טרפה לא ישלם אלא הטרפה, יש טרפה שהוא משלם ויש טרפה שאינו משלם; טרפת חתול ושועל ונמיה משלם, טרפת זאב, ארי, ודב ונחש אינו משלם; ומי לחשך לדון כן? שהרי כתוב ומת או נשבר או נשבה – מה מיתה שאין יכול להציל, אף שבר ושביה שאין יכול להציל (מכילתא):}%
% ---------- פרק 22 | פסוק 13 ----------
 \textbf{יג} וְכִֽי־יִשְׁאַ֥ל אִ֛ישׁ מֵעִ֥ם רֵעֵ֖הוּ וְנִשְׁבַּ֣ר אוֹ־מֵ֑ת בְּעָלָ֥יו אֵין־עִמּ֖וֹ שַׁלֵּ֥ם יְשַׁלֵּֽם׃ \Rashi{\textbf{וכי ישאל.}בא ללמד על השואל שהוא חיב באנסין: בעליו אין עמו. אם בעליו של שור אינו עם השואל במלאכתו (בבא מציעא צ"ה):}%
% ---------- פרק 22 | פסוק 14 ----------
 \textbf{יד} אִם־בְּעָלָ֥יו עִמּ֖וֹ לֹ֣א יְשַׁלֵּ֑ם אִם־שָׂכִ֣יר ה֔וּא בָּ֖א בִּשְׂכָרֽוֹ׃ \{ס\} \Rashi{\textbf{אם בעליו עמו.}בין שהוא באותה מלאכה בין שהוא במלאכה אחרת, היה עמו בשעת שאלה אינו צריך להיות עמו בשעת שבירה ומיתה (שם): אם שכיר הוא. אם השור אינו שאול אלא שכור, בא בשכרו ליד השוכר הזה, ולא בשאלה, ואין כל הנאה שלו, שהרי ע"י שכרו נשתמש, ואין לו משפט שואל להתחיב באנסין; ולא פרש מה דינו, אם כשומר חנם או כשומר שכר, לפיכך נחלקו בו חכמי ישראל שוכר כיצד משלם, רבי מאיר אומר כשומר חנם, ר' יהודה אומר כשומר שכר (בבא מציעא צ"ה):}%
% ---------- פרק 22 | פסוק 15 ----------
 \textbf{טו} וְכִֽי־יְפַתֶּ֣ה אִ֗ישׁ בְּתוּלָ֛ה אֲשֶׁ֥ר לֹא־אֹרָ֖שָׂה וְשָׁכַ֣ב עִמָּ֑הּ מָהֹ֛ר יִמְהָרֶ֥נָּה לּ֖וֹ לְאִשָּֽׁה׃ \Rashi{\textbf{וכי יפתה.}מדבר על לבה עד ששומעת לו, וכן תרגומו וארי ישדל, שדול בלשון ארמי כפתוי בלשון עברי: מהר ימהרנה. יפסק לה מהר כמשפט איש לאשתו – שכותב לה כתבה וישאנה (מכילתא):}%
% ---------- פרק 22 | פסוק 16 ----------
 \textbf{טז} אִם־מָאֵ֧ן יְמָאֵ֛ן אָבִ֖יהָ לְתִתָּ֣הּ ל֑וֹ כֶּ֣סֶף יִשְׁקֹ֔ל כְּמֹ֖הַר הַבְּתוּלֹֽת׃ \{ס\} \Rashi{\textbf{כמהר הבתולת.}שהוא קצוב חמשים כסף אצל התופס את הבתולה ושוכב עמה באנס, שנאמר ונתן האיש השכב עמה לאבי הנערה חמשים כסף (דברים כ"ב):}%
% ---------- פרק 22 | פסוק 17 ----------
 \textbf{יז} מְכַשֵּׁפָ֖ה לֹ֥א תְחַיֶּֽה׃ \Rashi{\textbf{מכשפה לא תחיה.}אלא תומת בבית דין; ואחד זכרים ואחד נקבות, אלא שדבר הכתוב בהווה, שהנשים מצויות מכשפות (סנהדרין ס"ז):}%
% ---------- פרק 22 | פסוק 18 ----------
 \textbf{יח} כׇּל־שֹׁכֵ֥ב עִם־בְּהֵמָ֖ה מ֥וֹת יוּמָֽת׃ \{ס\} \Rashi{\textbf{כל שכב עם בהמה מות יומת.}בסקילה, רובע כנרבעת, שכתוב בהם דמיהם בם (ויקרא כ'):}%
% ---------- פרק 22 | פסוק 19 ----------
 \textbf{יט} זֹבֵ֥חַ לָאֱלֹהִ֖ים יׇֽחֳרָ֑ם בִּלְתִּ֥י לַיהֹוָ֖ה לְבַדּֽוֹ׃ \Rashi{\textbf{לאלהים.}לע"ז; אלו היה נקוד לאלהים, היה צריך לפרש ולכתב "אחרים", עכשו שאמר לאלהים, אין צריך לפרש אחרים, שכל למ"ד ובי"ת המשמשת בראש התבה אם נקודה בחטף, כגון למלך, למדבר, לעיר, צריך לפרש לאיזה מלך, לאיזה מדבר, לאיזה עיר, וכן למלכים ולרגלים צריך לפרש לאיזה, ואם אינו מפרש, כל מלכים במשמע, כן לאלהים כל אלהים במשמע – אפלו קדש, אבל כשהיא נקודה פתח, כמו למלך, למדבר, לעיר, נודע באיזה מלך מדבר, וכן לעיר נודע באיזה עיר מדבר, וכן לאלהים – לאותן שהזהרתם עליהם במקום אחר; כיוצא בו אין כמוך באלהים, לפי שלא פרש, הצרך לנקד פתח: יחרם. יומת; למה נאמר יחרם והלא כבר נאמרה בו מיתה במקום אחר והוצאת את האיש ההוא או את האשה ההוא וגו' (דברים י"ז)? אלא לפי שלא פרש על איזו עבודה חיב מיתה – שלא תאמר כל עבודות במיתה – בא ופרש לך כאן "זבח" לאלהים, לומר לך, מה זביחה עבודה הנעשית בפנים לשמים, אף אני מרבה המקטיר והמנסך שהם עבודות בפנים, וחיבים עליהם לכל ע"ז בין שדרכה לעבדה בכך, בין שאין דרכה לעבדה בכך. אבל שאר עבודות, כגון המכבד והמרבץ והמגפף והמנשק אינו במיתה:}%
% ---------- פרק 22 | פסוק 20 ----------
 \textbf{כ} וְגֵ֥ר לֹא־תוֹנֶ֖ה וְלֹ֣א תִלְחָצֶ֑נּוּ כִּֽי־גֵרִ֥ים הֱיִיתֶ֖ם בְּאֶ֥רֶץ מִצְרָֽיִם׃ \Rashi{\textbf{וגר לא תונה.}אונאת דברים, קונטרארי"ר בלעז, כמו והאכלתי את מוניך את בשרם (ישעיהו מ"ט): ולא תלחצנו. בגזלת ממון: כי גרים הייתם. אם הוניתו, אף הוא יכול להונותך ולומר לך, אף אתה מגרים באת, "מום שבך אל תאמר לחברך"; כל לשון גר אדם שלא נולד באותה מדינה, אלא בא ממדינה אחרת לגור שם:}%
% ---------- פרק 22 | פסוק 21 ----------
 \textbf{כא} כׇּל־אַלְמָנָ֥ה וְיָת֖וֹם לֹ֥א תְעַנּֽוּן׃ \Rashi{\textbf{כל אלמנה ויתום לא תענון.}הוא הדין לכל אדם, אלא שדבר הכתוב בהווה, לפי שהם תשושי כח ודבר מצוי לענותם (מכילתא):}%
% ---------- פרק 22 | פסוק 22 ----------
 \textbf{כב} אִם־עַנֵּ֥ה תְעַנֶּ֖ה אֹת֑וֹ כִּ֣י אִם־צָעֹ֤ק יִצְעַק֙ אֵלַ֔י שָׁמֹ֥עַ אֶשְׁמַ֖ע צַעֲקָתֽוֹ׃ \Rashi{\textbf{אם ענה תענה אתו.}הרי זה מקרא קצר, גזם ולא פרש ענשו, כמו לכן כל הרג קין (בראשית ד'), גזם ולא פרש ענשו, אף כאן אם ענה תענה אותו לשון גזום, כלומר סופך לטל את שלך למה? כי אם צעק יצעק אלי וגו':}%
% ---------- פרק 22 | פסוק 23 ----------
 \textbf{כג} וְחָרָ֣ה אַפִּ֔י וְהָרַגְתִּ֥י אֶתְכֶ֖ם בֶּחָ֑רֶב וְהָי֤וּ נְשֵׁיכֶם֙ אַלְמָנ֔וֹת וּבְנֵיכֶ֖ם יְתֹמִֽים׃ \{פ\} \Rashi{\textbf{והיו נשיכם אלמנות.}ממשמע שנאמר והרגתי אתכם איני יודע שנשיכם אלמנות ובניכם יתומים? אלא הרי זו קללה אחרת, שיהיו הנשים צרורות כאלמנות חיות, שלא יהיו עדים למיתת בעליהן ותהיינה אסורות להנשא, והבנים יהיו יתומים, שלא יניחום בית דין לירד לנכסי אביהם, לפי שאין יודעים אם מתו אם נשבו (בבא מציעא ל"ח):}%
% ---------- פרק 22 | פסוק 24 ----------
 \textbf{כד} אִם־כֶּ֣סֶף ׀ תַּלְוֶ֣ה אֶת־עַמִּ֗י אֶת־הֶֽעָנִי֙ עִמָּ֔ךְ לֹא־תִהְיֶ֥ה ל֖וֹ כְּנֹשֶׁ֑ה לֹֽא־תְשִׂימ֥וּן עָלָ֖יו נֶֽשֶׁךְ׃ \Rashi{\textbf{אם כסף תלוה את עמי.}רבי ישמעאל אומר כל אם ואם שבתורה רשות חוץ משלשה, וזה אחד מהן (מכילתא): את עמי. עמי וגוי עמי קודם, עני ועשיר עני קודם, ענייך ועניי עירך ענייך קודמין, עניי עירך ועניי עיר אחרת עניי עירך קודמין; וזה משמעו: אם כסף תלוה – את עמי תלוהו ולא לגוי, ולאיזה מעמי? את העני, ולאיזה עני? לאותו שעמך. (ד"א, את עמי, שלא תנהג בו בזיון בהלואה שהוא עמי: את העני עמך. הוי מסתכל בעצמך כאלו אתה עני: לא תהיה לו כנשה. לא תתבענו בחזקה. אם אתה יודע שאין לו, אל תהי דומה עליו כאלו הלויתו אלא כאלו לא הלויתו, כלומר לא תכלימהו: נשך. רבית, שהוא כנשיכת נחש שנושך חבורה קטנה ברגלו ואינו מרגיש, ופתאם הוא מבטבט ונופח עד קדקדו, כך רבית אינו מרגיש ואינו נכר עד שהרבית עולה ומחסרו ממון הרבה (תנחומא):}%
% ---------- פרק 22 | פסוק 25 ----------
 \textbf{כה} אִם־חָבֹ֥ל תַּחְבֹּ֖ל שַׂלְמַ֣ת רֵעֶ֑ךָ עַד־בֹּ֥א הַשֶּׁ֖מֶשׁ תְּשִׁיבֶ֥נּוּ לֽוֹ׃ \Rashi{\textbf{אם חבל תחבל.}כל לשון חבלה אינו משכון בשעת הלואה, אלא שממשכנין את הלוה כשמגיע הזמן ואינו פורע; (חבל תחבל – כפל לך בחבלה עד כמה פעמים; אמר הקב"ה כמה אתה חיב לי, והרי נפשך עולה אצלי כל אמש ואמש ונותנת דין ומתחיבת לפני ואני מחזירה לך, אף אתה טל והשב טל והשב): עד בא השמש תשיבנו לו. כל היום תשיבנו לו עד בא השמש, וכבא השמש תחזר ותטלנו עד שיבא בקר של מחר; ובכסות יום הכתוב מדבר, שאין צריך לה בלילה (מכילתא, בבא מציעא קי"ד):}%
% ---------- פרק 22 | פסוק 26 ----------
 \textbf{כו} כִּ֣י הִ֤וא כְסוּתֹה֙ לְבַדָּ֔הּ הִ֥וא שִׂמְלָת֖וֹ לְעֹר֑וֹ בַּמֶּ֣ה יִשְׁכָּ֔ב וְהָיָה֙ כִּֽי־יִצְעַ֣ק אֵלַ֔י וְשָׁמַעְתִּ֖י כִּֽי־חַנּ֥וּן אָֽנִי׃ \{ס\} \Rashi{\textbf{כי הוא כסותה.}זו טלית: שמלתו. זה חלוק: במה ישכב. לרבות את המצע (מכילתא):}%
% ---------- פרק 22 | פסוק 27 ----------
 \textbf{כז} אֱלֹהִ֖ים לֹ֣א תְקַלֵּ֑ל וְנָשִׂ֥יא בְעַמְּךָ֖ לֹ֥א תָאֹֽר׃ \Rashi{\textbf{אלהים לא תקלל.}הרי זו אזהרה לברכת השם ואזהרה לקללת דין (סנהדרין ס"ו):}%
% ---------- פרק 22 | פסוק 28 ----------
 \textbf{כח} מְלֵאָתְךָ֥ וְדִמְעֲךָ֖ לֹ֣א תְאַחֵ֑ר בְּכ֥וֹר בָּנֶ֖יךָ תִּתֶּן־לִֽי׃ \Rashi{\textbf{מלאתך.}חובה המטלת עליך כשתתמלא תבואתך להתבשל, והם בכורים: ודמעך. התרומה – ואיני יודע מהו לשון דמע. לא תאחר. לא תשנה סדר הפרשתן לאחר את המקדם ולהקדים את המאחר, שלא יקדים תרומה לבכורים ומעשר לתרומה (מכילתא): בכור בניך תתן לי. לפדותו בחמש סלעים מן הכהן; והלא כבר צוה עליו במקום אחר? אלא כדי לסמך לו כן תעשה לשרך – מה בכור אדם לאחר שלשים יום פודהו, שנאמר ופדויו מבן חדש תפדה (במדבר י"ח), אף בכור בהמה דקה מטפל בו שלשים יום ואח"כ נותנו לכהן (בכורות כ"ו):}%
% ---------- פרק 22 | פסוק 29 ----------
 \textbf{כט} כֵּֽן־תַּעֲשֶׂ֥ה לְשֹׁרְךָ֖ לְצֹאנֶ֑ךָ שִׁבְעַ֤ת יָמִים֙ יִהְיֶ֣ה עִם־אִמּ֔וֹ בַּיּ֥וֹם הַשְּׁמִינִ֖י תִּתְּנוֹ־לִֽי׃ \Rashi{\textbf{שבעת ימים יהיה עם אמו.}זו אזהרה לכהן, שאם בא למהר את קרבנו לא ימהר קדם שמנה, לפי שהוא מחסר זמן: ביום השמיני תתנו לי. יכול יהא חובה לבו ביום, נאמר כאן שמיני ונאמר להלן ומיום השמיני והלאה ירצה (ויקרא כ"ב), מה שמיני האמור להלן להכשיר משמיני ולהלן, אף שמיני האמור כאן להכשיר משמיני ולהלן; וכן משמעו: וביום השמיני אתה רשאי לתנו לי (מכילתא):}%
% ---------- פרק 22 | פסוק 30 ----------
 \textbf{ל} וְאַנְשֵׁי־קֹ֖דֶשׁ תִּהְי֣וּן לִ֑י וּבָשָׂ֨ר בַּשָּׂדֶ֤ה טְרֵפָה֙ לֹ֣א תֹאכֵ֔לוּ לַכֶּ֖לֶב תַּשְׁלִכ֥וּן אֹתֽוֹ׃ \{ס\} \Rashi{\textbf{ואנשי קדש תהיון לי.}אם אתם קדושים ופרושים משקוצי נבלות וטרפות הרי אתם שלי ואם לאו אינכם שלי: ובשר בשדה טרפה. אף בבית כן, אלא שדבר הכתוב בהווה – מקום שדרך בהמות לטרף; וכן כי בשדה מצאה (דברים כ"ב), וכן אשר לא יהיה טהור מקרה לילה (שם כ"ג), הוא הדין למקרה יום, אלא שדבר הכתוב בהווה. ובשר דתליש מן חיוא חיא, בשר שנתלש ע"י טרפת זאב או ארי מן חיה כשרה או מבהמה כשרה בחייה: לכלב תשלכון אתו. אף הגוי ככלב; או אינו אלא כלב כמשמעו? ת"ל בנבלה או מכר לנכרי (שם י"ד), ק"ו לטרפה שמתרת בכל הנאות. אם כן מה ת"ל לכלב? ללמדך שהכלב נכבד ממנו, ולמדך הכתוב שאין הקב"ה מקפח שכר כל בריה, שנאמר ולכל בני ישראל לא יחרץ כלב לשנו (שמות י"א), אמר הקב"ה תנו לו שכרו (מכילתא):}%
% ---------- פרק 23 | פסוק 1 ----------
 \textbf{א} לֹ֥א תִשָּׂ֖א שֵׁ֣מַע שָׁ֑וְא אַל־תָּ֤שֶׁת יָֽדְךָ֙ עִם־רָשָׁ֔ע לִהְיֹ֖ת עֵ֥ד חָמָֽס׃ \Rashi{\textbf{לא תשא שמע שוא.}כתרגומו, לא תקבל שמע דשקר, אזהרה למקבל לשון הרע ולדין שלא ישמע דברי בעל דין עד שיבא בעל דין חברו (שם): אל תשת ידך עם רשע. הטוען את חברו תביעת שקר, שתבטיחהו להיות לו עד חמס:}%
% ---------- פרק 23 | פסוק 2 ----------
 \textbf{ב} לֹֽא־תִהְיֶ֥ה אַחֲרֵֽי־רַבִּ֖ים לְרָעֹ֑ת וְלֹא־תַעֲנֶ֣ה עַל־רִ֗ב לִנְטֹ֛ת אַחֲרֵ֥י רַבִּ֖ים לְהַטֹּֽת׃ \Rashi{\textbf{לא תהיה אחרי רבים לרעת.}יש במקרא זה מדרשי חכמי ישראל, אבל אין לשון המקרא מישב בהן על אפניו. מכאן דרשו שאין מטין לחובה בהכרעת דין אחד, וסוף המקרא דרשו אחרי רבים להטת, שאם יש שנים מחיבין יותר על המזכין הטה הדין על פיהם לחובה – ובדיני נפשות הכתוב מדבר – ואמצע המקרא דרשו לא תענה על רב – על רב, שאין חולקין על מפלא שבבית דין, לפיכך מתחילין בדיני נפשות מן הצד – לקטנים שבהם שואלין תחלה שיאמרו את דעתם – ולפי דברי רבותינו כך פתרון המקרא: לא תהיה אחרי רבים לרעת. לחיב מיתה בשביל דין אחד שירבו מחיבין על המזכין, ולא תענה על הרב לנטות מדבריו – ולפי שהוא חסר יו"ד דרשו בו כן – אחרי רבים להטת, יש רבים שאתה נוטה אחריהם, ואימתי? בזמן שהן שנים המכריעין במחיבין יותר מן המזכין; וממשמע שנאמר לא תהיה אחרי רבים לרעת שומע אני אבל היה עמהם לטובה, מכאן אמרו דיני נפשות מטין על פי אחד לזכות ועל פי שנים לחובה. ואנקלוס תרגם לא תתמנע מלאלפא מה דמתבעי לך (דבעינך) על דינא, ולשון העברי לפי התרגום כך הוא נדרש: לא תענה על רב לנטת. אם ישאלוך דבר למשפט לא תענה לנטות לצד אחד ולסלק עצמך מן הריב, אלא הוי דן אותו לאמתו. ואני אומר לישבו על אפניו כפשוטו כך פתרונו: לא תהיה אחרי רבים לרעת. אם ראית רשעים מטין משפט, לא תאמר, הואיל ורבים הם הנני נוטה אחריהם: ולא תענה על רב לנטת וגו'. ואם ישאלך הנדון על אותו המשפט, אל תעננו על הריב דבר הנוטה אחרי אותן רבים להטות את המשפט מאמתו, אלא אמר את המשפט כאשר הוא וקולר יהא תלוי בצואר הרבים (סנהדרין ז'):}%
% ---------- פרק 23 | פסוק 3 ----------
 \textbf{ג} וְדָ֕ל לֹ֥א תֶהְדַּ֖ר בְּרִיבֽוֹ׃ \{ס\} \Rashi{\textbf{לא תהדר.}לא תחלק לו כבוד לזכותו בדין ולומר דל הוא אזכנו ואכבדנו:}%
% ---------- פרק 23 | פסוק 4 ----------
 \textbf{ד} כִּ֣י תִפְגַּ֞ע שׁ֧וֹר אֹֽיִבְךָ֛ א֥וֹ חֲמֹר֖וֹ תֹּעֶ֑ה הָשֵׁ֥ב תְּשִׁיבֶ֖נּוּ לֽוֹ׃ \{ס\} %
% ---------- פרק 23 | פסוק 5 ----------
 \textbf{ה} כִּֽי־תִרְאֶ֞ה חֲמ֣וֹר שֹׂנַאֲךָ֗ רֹבֵץ֙ תַּ֣חַת מַשָּׂא֔וֹ וְחָדַלְתָּ֖ מֵעֲזֹ֣ב ל֑וֹ עָזֹ֥ב תַּעֲזֹ֖ב עִמּֽוֹ׃ \{ס\} \Rashi{\textbf{כי תראה חמור שנאך וגו'.}הרי כי משמש לשון דלמא, שהוא מארבע לשונות של שמושי כי, וכה פתרונו: שמא תראה חמורו רבץ תחת משאו וחדלת מעזב לו, בתמיה: עזב תעזב עמו. עזיבה זו לשון עזרה, וכן עצור ועזוב (דברים ל"ב), וכן ויעזבו ירושלים עד החומה (נחמיה ג') – מלאוה עפר לעזר ולסיע את חזק החומה. כיוצא בו כי תאמר בלבבך רבים הגוים האלה ממני וגו' (דברים ז'), שמא תאמר כן, בתמיה? לא תירא מהם. ומדרשו כך דרשו רבותינו: כי תראה וחדלת, פעמים שאתה חודל, ופעמים שאתה עוזר, הא כיצד? זקן ואינו לפי כבודו וחדלת, או בהמת גוי ומשאו של ישראל וחדלת (מכילתא): עזב תעזב עמו. לפרק המשא; מלמשקל ליה – מלטל משאו ממנו:}%
% ---------- פרק 23 | פסוק 6 ----------
 \textbf{ו} לֹ֥א תַטֶּ֛ה מִשְׁפַּ֥ט אֶבְיֹנְךָ֖ בְּרִיבֽוֹ׃ \Rashi{\textbf{אבינך.}לשון אובה, שהוא מדלדל ותאב לכל טובה:}%
% ---------- פרק 23 | פסוק 7 ----------
 \textbf{ז} מִדְּבַר־שֶׁ֖קֶר תִּרְחָ֑ק וְנָקִ֤י וְצַדִּיק֙ אַֽל־תַּהֲרֹ֔ג כִּ֥י לֹא־אַצְדִּ֖יק רָשָֽׁע׃ \Rashi{\textbf{ונקי וצדיק אל תהרג.}מנין ליוצא מבית דין חיב ואמר אחד יש לי ללמד עליו זכות שמחזירין אותו? ת"ל ונקי אל תהרג, ואע"פ שאינו צדיק, שלא נצטדק בבית דין, מכל מקום נקי הוא מדין מיתה שהרי יש לך לזכותו. ומנין ליוצא מבית דין זכאי ואמר אחד יש לי ללמד עליו חובה שאין מחזירין אותו לבית דין? ת"ל וצדיק אל תהרג, וזה צדיק הוא שנצטדק בבית דין (סנהדרין ל"ג): כי לא אצדיק רשע. אין עליך להחזירו, כי אני לא אצדיקנו בדיני, אם יצא מידך זכאי יש לי שלוחים הרבה להמיתו במיתה שנתחיב בה (מכילתא):}%
% ---------- פרק 23 | פסוק 8 ----------
 \textbf{ח} וְשֹׁ֖חַד לֹ֣א תִקָּ֑ח כִּ֤י הַשֹּׁ֙חַד֙ יְעַוֵּ֣ר פִּקְחִ֔ים וִֽיסַלֵּ֖ף דִּבְרֵ֥י צַדִּיקִֽים׃ \Rashi{\textbf{ושחד לא תקח.}אפלו לשפט אמת, וכל שכן כדי להטות את הדין, שהרי להטות את הדין נאמר כבר (דברים ט"ז) לא תטה משפט (כתובות ק"ה): יעור פקחים. אפלו חכם בתורה ונוטל שחד סוף שתטרף דעתו עליו וישתכח תלמודו ויכהה מאור עיניו (שם): ויסלף. כתרגומו ומקלקל: דברי צדיקים. דברים המצדקים, משפטי אמת, וכן תרגומו פתגמין תריצין – ישרים:}%
% ---------- פרק 23 | פסוק 9 ----------
 \textbf{ט} וְגֵ֖ר לֹ֣א תִלְחָ֑ץ וְאַתֶּ֗ם יְדַעְתֶּם֙ אֶת־נֶ֣פֶשׁ הַגֵּ֔ר כִּֽי־גֵרִ֥ים הֱיִיתֶ֖ם בְּאֶ֥רֶץ מִצְרָֽיִם׃ \Rashi{\textbf{וגר לא תלחץ.}בהרבה מקומות הזהירה תורה על הגר מפני שסורו רע (בבא מציעא נ"ט): את נפש הגר. כמה קשה לו כשלוחצים אותו:}%
% ---------- פרק 23 | פסוק 10 ----------
 \textbf{י} וְשֵׁ֥שׁ שָׁנִ֖ים תִּזְרַ֣ע אֶת־אַרְצֶ֑ךָ וְאָסַפְתָּ֖ אֶת־תְּבוּאָתָֽהּ׃ \Rashi{\textbf{ואספת את תבואתה.}לשון הכנסה לבית, כמו ואספתו אל תוך ביתך (דברים כ"ב):}%
% ---------- פרק 23 | פסוק 11 ----------
 \textbf{יא} וְהַשְּׁבִיעִ֞ת תִּשְׁמְטֶ֣נָּה וּנְטַשְׁתָּ֗הּ וְאָֽכְלוּ֙ אֶבְיֹנֵ֣י עַמֶּ֔ךָ וְיִתְרָ֕ם תֹּאכַ֖ל חַיַּ֣ת הַשָּׂדֶ֑ה כֵּֽן־תַּעֲשֶׂ֥ה לְכַרְמְךָ֖ לְזֵיתֶֽךָ׃ \Rashi{\textbf{תשמטנה.}מעבודה: ונטשתה. מאכילה אחר זמן הבעור. ד"א, תשמטנה – מעבודה גמורה, כגון חרישה וזריעה, ונטשתה – מלזבל ומלקשקש (סוכה מ"ד): ויתרם תאכל חית השדה. להקיש מאכל אביון למאכל חיה, מה חיה אוכלת בלא מעשר, אף אביונים אוכלים בלא מעשר; מכאן אמרו אין מעשר בשביעית (מכילתא): כן תעשה לכרמך. ותחלת המקרא מדבר בשדה הלבן, כמו שאמור למעלה הימנו תזרע את ארצך:}%
% ---------- פרק 23 | פסוק 12 ----------
 \textbf{יב} שֵׁ֤שֶׁת יָמִים֙ תַּעֲשֶׂ֣ה מַעֲשֶׂ֔יךָ וּבַיּ֥וֹם הַשְּׁבִיעִ֖י תִּשְׁבֹּ֑ת לְמַ֣עַן יָנ֗וּחַ שֽׁוֹרְךָ֙ וַחֲמֹרֶ֔ךָ וְיִנָּפֵ֥שׁ בֶּן־אֲמָתְךָ֖ וְהַגֵּֽר׃ \Rashi{\textbf{וביום השביעי תשבת.}אף בשנה השביעית לא תעקר שבת בראשית ממקומה, שלא תאמר, הואיל וכל השנה קרויה שבת לא תנהג בה שבת בראשית (שם): למען ינוח שורך וחמרך. תן לו ניח, להתיר שיהא תולש ואוכל עשבים מן הקרקע; או אינו אלא יחבשנו בתוך הבית? אמרת אין זה ניח אלא צער: בן אמתך. בעבד ערל הכתוב מדבר: והגר. גר תושב:}%
% ---------- פרק 23 | פסוק 13 ----------
 \textbf{יג} וּבְכֹ֛ל אֲשֶׁר־אָמַ֥רְתִּי אֲלֵיכֶ֖ם תִּשָּׁמֵ֑רוּ וְשֵׁ֨ם אֱלֹהִ֤ים אֲחֵרִים֙ לֹ֣א תַזְכִּ֔ירוּ לֹ֥א יִשָּׁמַ֖ע עַל־פִּֽיךָ׃ \Rashi{\textbf{ובכל אשר אמרתי אליכם תשמרו.}לעשות כל מצות עשה באזהרה, שכל שמירה שבתורה אזהרה היא במקום לאו (מנחות ל"ו): לא תזכירו. שלא יאמר לו שמר לי בצד ע"ז פלונית, או תעמד עמי ביום ע"ז פלונית; ד"א — ובכל אשר אמרתי אליכם תשמרו ושם אלהים אחרים לא תזכירו, ללמדך ששקולה ע"ז כנגד כל המצוות כלן (הוריות ח'), והנזהר בה כשומר את כלן: לא ישמע מן הגוי על פיך. שלא תעשה שתפות עם גוי וישבע לך בע"ז שלו נמצאת שאתה גורם שיזכר על ידך (סנהדרין ס"ז):}%
% ---------- פרק 23 | פסוק 14 ----------
 \textbf{יד} שָׁלֹ֣שׁ רְגָלִ֔ים תָּחֹ֥ג לִ֖י בַּשָּׁנָֽה׃ \Rashi{\textbf{רגלים.}פעמים; וכן כי הכיתני זה שלש רגלים (במדבר כ"ב):}%
% ---------- פרק 23 | פסוק 15 ----------
 \textbf{טו} אֶת־חַ֣ג הַמַּצּוֹת֮ תִּשְׁמֹר֒ שִׁבְעַ֣ת יָמִים֩ תֹּאכַ֨ל מַצּ֜וֹת כַּֽאֲשֶׁ֣ר צִוִּיתִ֗ךָ לְמוֹעֵד֙ חֹ֣דֶשׁ הָֽאָבִ֔יב כִּי־ב֖וֹ יָצָ֣אתָ מִמִּצְרָ֑יִם וְלֹא־יֵרָא֥וּ פָנַ֖י רֵיקָֽם׃ \Rashi{\textbf{חדש האביב.}שהתבואה מתמלאת בו באביה. אביב לשון אב, בכור וראשון לבשל פרות: ולא יראו פני ריקם. כשתבאו לראות פני ברגלים, הביאו לי עולות (מכילתא):}%
% ---------- פרק 23 | פסוק 16 ----------
 \textbf{טז} וְחַ֤ג הַקָּצִיר֙ בִּכּוּרֵ֣י מַעֲשֶׂ֔יךָ אֲשֶׁ֥ר תִּזְרַ֖ע בַּשָּׂדֶ֑ה וְחַ֤ג הָֽאָסִף֙ בְּצֵ֣את הַשָּׁנָ֔ה בְּאׇסְפְּךָ֥ אֶֽת־מַעֲשֶׂ֖יךָ מִן־הַשָּׂדֶֽה׃ \Rashi{\textbf{וחג הקציר.}הוא חג שבועות: בכורי מעשיך. שהוא זמן הבאת בכורים, ששתי הלחם הבאין בעצרת היו מתירין החדש למנחות ולהביא בכורים למקדש, שנאמר וביום הבכורים וגו' (במדבר כ"ח): וחג האסף. הוא חג הסכות: באספך את מעשיך. שכל ימות החמה התבואה מתיבשת בשדות, ובחג אוספים אותה אל הבית מפני הגשמים:}%
% ---------- פרק 23 | פסוק 17 ----------
 \textbf{יז} שָׁלֹ֥שׁ פְּעָמִ֖ים בַּשָּׁנָ֑ה יֵרָאֶה֙ כׇּל־זְכ֣וּרְךָ֔ אֶל־פְּנֵ֖י הָאָדֹ֥ן ׀ יְהֹוָֽה׃ \Rashi{\textbf{שלש פעמים וגו'.}לפי שהענין מדבר בשביעית, הצרך לומר שלא יסתרסו רגלים ממקומן: כל זכורך. הזכרים שבך:}%
% ---------- פרק 23 | פסוק 18 ----------
 \textbf{יח} לֹֽא־תִזְבַּ֥ח עַל־חָמֵ֖ץ דַּם־זִבְחִ֑י וְלֹֽא־יָלִ֥ין חֵֽלֶב־חַגִּ֖י עַד־בֹּֽקֶר׃ \Rashi{\textbf{לא תזבח על חמץ וגו'.}לא תשחט את הפסח בי"ד בניסן עד שתבער החמץ (מכילתא): ולא ילין חלב חגי חוץ למזבח: עד בקר. יכול אף על המערכה יפסל בלינה, ת"ל על מוקדה על המזבח כל הלילה (ויקרא ו'): ולא ילין. אין לינה אלא בעמוד השחר, שנאמר עד בקר, אבל כל הלילה יכול להעלותו מן הרצפה למזבח (מגילה כ'):}%
% ---------- פרק 23 | פסוק 19 ----------
 \textbf{יט} רֵאשִׁ֗ית בִּכּוּרֵי֙ אַדְמָ֣תְךָ֔ תָּבִ֕יא בֵּ֖ית יְהֹוָ֣ה אֱלֹהֶ֑יךָ לֹֽא־תְבַשֵּׁ֥ל גְּדִ֖י בַּחֲלֵ֥ב אִמּֽוֹ׃ \{פ\} \Rashi{\textbf{ראשית בכורי אדמתך.}אף השביעית חיבת בבכורים, לכך נאמר אף כאן בכורי אדמתך. כיצד? אדם נכנס לתוך שדהו, רואה תאנה שבכרה, כורך עליה גמי לסימן ומקדישה, ואין בכורים אלא בשבעת המינין האמורין במקרא – ארץ חטה ושערה וגו' (בכורים פ"א): לא תבשל גדי. אף עגל וכבש בכלל גדי, שאין גדי אלא לשון ולד רך, ממה שאתה מוצא בכמה מקומות בתורה שכתוב גדי והצרך לפרש אחריו עזים, כגון אנכי אשלח גדי עזים (בראשית ל"ח), את גדי העזים (שם), שני גדיי עזים (שם כ"ז), ללמדך שכל מקום שנאמר גדי סתם אף עגל וכבש במשמע. ובג' מקומות נכתב בתורה, אחד לאסור אכילה, ואחד לאסור הנאה, ואחד לאסור בשול (מכילתא, חולין קט"ו):}%
% ---------- פרק 23 | פסוק 20 ----------
 \textbf{כ} הִנֵּ֨ה אָנֹכִ֜י שֹׁלֵ֤חַ מַלְאָךְ֙ לְפָנֶ֔יךָ לִשְׁמׇרְךָ֖ בַּדָּ֑רֶךְ וְלַהֲבִ֣יאֲךָ֔ אֶל־הַמָּק֖וֹם אֲשֶׁ֥ר הֲכִנֹֽתִי׃ \Rashi{\textbf{הנה אנכי שלח מלאך.}כאן נתבשרו שעתידין לחטא, ושכינה אומרת להם כי לא אעלה בקרבך (שמות ל"ג): אשר הכנתי. אשר זמנתי לתת לכם, זהו פשוטו; ומדרשו אל המקום אשר הכנתי כבר מקומי כנגדו, וזה אחד מן המקראות שאומרים שבית המקדש של מעלה מכון כנגד של מטה (תנחומא):}%
% ---------- פרק 23 | פסוק 21 ----------
 \textbf{כא} הִשָּׁ֧מֶר מִפָּנָ֛יו וּשְׁמַ֥ע בְּקֹל֖וֹ אַל־תַּמֵּ֣ר בּ֑וֹ כִּ֣י לֹ֤א יִשָּׂא֙ לְפִשְׁעֲכֶ֔ם כִּ֥י שְׁמִ֖י בְּקִרְבּֽוֹ׃ \Rashi{\textbf{אל תמר בו.}לשון המראה, כמו אשר ימרה את פיך (יהושע א'): כי לא ישא לפשעכם. אינו מלמד בכך, שהוא מן הכת שאין חוטאין, ועוד שהוא שליח ואינו עושה אלא שליחותו: כי שמי בקרבו. מחבר לראש המקרא – השמר מפניו כי שמי משתף בו. ורבותינו אמרו זה מטטרון ששמו כשם רבו, מטטרון בגימטריא שדי (סנהדרין ל"ח):}%
% ---------- פרק 23 | פסוק 22 ----------
 \textbf{כב} כִּ֣י אִם־שָׁמ֤וֹעַ תִּשְׁמַע֙ בְּקֹל֔וֹ וְעָשִׂ֕יתָ כֹּ֖ל אֲשֶׁ֣ר אֲדַבֵּ֑ר וְאָֽיַבְתִּי֙ אֶת־אֹ֣יְבֶ֔יךָ וְצַרְתִּ֖י אֶת־צֹרְרֶֽיךָ׃ \Rashi{\textbf{וצרתי.}כתרגומו ואעיק:}%
% ---------- פרק 23 | פסוק 23 ----------
 \textbf{כג} כִּֽי־יֵלֵ֣ךְ מַלְאָכִי֮ לְפָנֶ֒יךָ֒ וֶהֱבִֽיאֲךָ֗ אֶל־הָֽאֱמֹרִי֙ וְהַ֣חִתִּ֔י וְהַפְּרִזִּי֙ וְהַֽכְּנַעֲנִ֔י הַחִוִּ֖י וְהַיְבוּסִ֑י וְהִכְחַדְתִּֽיו׃ %
% ---------- פרק 23 | פסוק 24 ----------
 \textbf{כד} לֹֽא־תִשְׁתַּחֲוֶ֤ה לֵאלֹֽהֵיהֶם֙ וְלֹ֣א תׇֽעׇבְדֵ֔ם וְלֹ֥א תַעֲשֶׂ֖ה כְּמַֽעֲשֵׂיהֶ֑ם כִּ֤י הָרֵס֙ תְּהָ֣רְסֵ֔ם וְשַׁבֵּ֥ר תְּשַׁבֵּ֖ר מַצֵּבֹתֵיהֶֽם׃ \Rashi{\textbf{הרס תהרסם.}לאותם אלוהות: מצבתיהם. אבנים שהם מציבין להשתחוות להם:}%
% ---------- פרק 23 | פסוק 25 ----------
 \textbf{כה} וַעֲבַדְתֶּ֗ם אֵ֚ת יְהֹוָ֣ה אֱלֹֽהֵיכֶ֔ם וּבֵרַ֥ךְ אֶֽת־לַחְמְךָ֖ וְאֶת־מֵימֶ֑יךָ וַהֲסִרֹתִ֥י מַחֲלָ֖ה מִקִּרְבֶּֽךָ׃ \{ס\} %
% ---------- פרק 23 | פסוק 26 ----------
 \textbf{כו} לֹ֥א תִהְיֶ֛ה מְשַׁכֵּלָ֥ה וַעֲקָרָ֖ה בְּאַרְצֶ֑ךָ אֶת־מִסְפַּ֥ר יָמֶ֖יךָ אֲמַלֵּֽא׃ \Rashi{\textbf{לא תהיה משכלה.}אם תעשה רצוני: משכלה. מפלת נפלים או קוברת את בניה קרויה משכלה:}%
% ---------- פרק 23 | פסוק 27 ----------
 \textbf{כז} אֶת־אֵֽימָתִי֙ אֲשַׁלַּ֣ח לְפָנֶ֔יךָ וְהַמֹּתִי֙ אֶת־כׇּל־הָעָ֔ם אֲשֶׁ֥ר תָּבֹ֖א בָּהֶ֑ם וְנָתַתִּ֧י אֶת־כׇּל־אֹיְבֶ֛יךָ אֵלֶ֖יךָ עֹֽרֶף׃ \Rashi{\textbf{והמתי.}כמו והממתי, ותרגומו ואשגש. וכן כל תבה שפעל שלה בכפל אות אחרונה, כשתהפך לדבר בלשון פעלתי, יש מקומות שנוטל אות הכפולה ומדגיש את האות ונוקדו במלאפום כגון והמתי מגזרת והמם גלגל עגלתו (ישעיהו כ"ח), וסבותי (קהלת ב'), מגזרת וסבב בית אל (שמואל א ז׳:ט״ז), דלתי (תהילים קט"ז), מגזרת דללו וחרבו (ישעיהו י"ט), על כפים חקתיך (ישעיהו מ״ט:ט״ז), מגזרת חקקי לב (שופטים ה'), את מי רצותי (שמואל א י״ב:ג׳), מגזרת רצץ עזב דלים (איוב כ'). והמתרגם והמתי ואקטל, טועה הוא, שאלו מגזרת מיתה היה, אין ה"א שלה בפתח ולא מ"ם שלה מדגשת ולא נקודה מלאפום, אלא והמתי כגון והמתה את העם הזה (במדבר י״ד:ט״ו), והתי"ו מדגשת לפי שתבא במקום שתי תוי"ן, האחת נשרשת – לפי שאין מיתה בלא תי"ו – והאחרת משמשת, כמו אמרתי, חטאתי, עשיתי, וכן ונתתי התי"ו מדגשת, שהיא באה במקום שתים, לפי שהיה צריך שלש תוי"ן – שתים ליסוד, כמו ביום תת ה' (יהושע י'), מתת אלהים היא (קהלת ג'), והשלישית לשמוש: ערף. שינוסו מפניך ויהפכו לך ערפם:}%
% ---------- פרק 23 | פסוק 28 ----------
 \textbf{כח} וְשָׁלַחְתִּ֥י אֶת־הַצִּרְעָ֖ה לְפָנֶ֑יךָ וְגֵרְשָׁ֗ה אֶת־הַחִוִּ֧י אֶת־הַֽכְּנַעֲנִ֛י וְאֶת־הַחִתִּ֖י מִלְּפָנֶֽיךָ׃ \Rashi{\textbf{הצרעה.}מין שרץ העוף, והיתה מכה אותם בעיניהם ומטילה בהם ארס והם מתים. והצרעה לא עברה את הירדן, והחתי והכנעני הם ארץ סיחון ועוג, לפיכך מכל שבע אמות לא מנה כאן אלא אלו, וחוי אע"פ שהוא מעבר הירדן והלאה, שנו רבותינו במסכת סוטה (דף ל"ו), על שפת הירדן עמדה וזרקה בהם מרה:}%
% ---------- פרק 23 | פסוק 29 ----------
 \textbf{כט} לֹ֧א אֲגָרְשֶׁ֛נּוּ מִפָּנֶ֖יךָ בְּשָׁנָ֣ה אֶחָ֑ת פֶּן־תִּהְיֶ֤ה הָאָ֙רֶץ֙ שְׁמָמָ֔ה וְרַבָּ֥ה עָלֶ֖יךָ חַיַּ֥ת הַשָּׂדֶֽה׃ \Rashi{\textbf{שממה.}רקנית מבני אדם, לפי שאתם מעט ואין בכם כדי למלאות אותה: ורבה עליך. ותרבה עליך:}%
% ---------- פרק 23 | פסוק 30 ----------
 \textbf{ל} מְעַ֥ט מְעַ֛ט אֲגָרְשֶׁ֖נּוּ מִפָּנֶ֑יךָ עַ֚ד אֲשֶׁ֣ר תִּפְרֶ֔ה וְנָחַלְתָּ֖ אֶת־הָאָֽרֶץ׃ \Rashi{\textbf{עד אשר תפרה.}תרבה, לשון פרי, כמו פרו ורבו (בראשית א'):}%
% ---------- פרק 23 | פסוק 31 ----------
 \textbf{לא} וְשַׁתִּ֣י אֶת־גְּבֻלְךָ֗ מִיַּם־סוּף֙ וְעַד־יָ֣ם פְּלִשְׁתִּ֔ים וּמִמִּדְבָּ֖ר עַד־הַנָּהָ֑ר כִּ֣י ׀ אֶתֵּ֣ן בְּיֶדְכֶ֗ם אֵ֚ת יֹשְׁבֵ֣י הָאָ֔רֶץ וְגֵרַשְׁתָּ֖מוֹ מִפָּנֶֽיךָ׃ \Rashi{\textbf{ושתי.}לשון השתה, והתי"ו מדגשת מפני שבאה תחת שתים, שאין שיתה בלא תי"ו, והאחת לשמוש: עד הנהר. פרת: וגרשתמו. ותגרשם:}%
% ---------- פרק 23 | פסוק 32 ----------
 \textbf{לב} לֹֽא־תִכְרֹ֥ת לָהֶ֛ם וְלֵאלֹֽהֵיהֶ֖ם בְּרִֽית׃ %
% ---------- פרק 23 | פסוק 33 ----------
 \textbf{לג} לֹ֤א יֵשְׁבוּ֙ בְּאַרְצְךָ֔ פֶּן־יַחֲטִ֥יאוּ אֹתְךָ֖ לִ֑י כִּ֤י תַעֲבֹד֙ אֶת־אֱלֹ֣הֵיהֶ֔ם כִּֽי־יִהְיֶ֥ה לְךָ֖ לְמוֹקֵֽשׁ׃ \{פ\} \Rashi{\textbf{כי תעבד וגו'.}הרי אלו כי משמשין במקום אשר, וכן בכמה מקומות, וזהו לשון אי שהוא אחד מארבע לשונות שהכי משמש, וגם מצינו בהרבה מקומות אם משמש בלשון אשר, כמו ואם תקריב מנחת בכורים (ויקרא ב'), שהיא חובה:}%
% ---------- פרק 24 | פסוק 1 ----------
 \textbf{א} וְאֶל־מֹשֶׁ֨ה אָמַ֜ר עֲלֵ֣ה אֶל־יְהֹוָ֗ה אַתָּה֙ וְאַהֲרֹן֙ נָדָ֣ב וַאֲבִיה֔וּא וְשִׁבְעִ֖ים מִזִּקְנֵ֣י יִשְׂרָאֵ֑ל וְהִשְׁתַּחֲוִיתֶ֖ם מֵרָחֹֽק׃ \Rashi{\textbf{ואל משה אמר.}פרשה זו נאמרה קדם עשרת הדברות, בד' בסיון נאמרה לו עלה (שבת פ"ח):}%
% ---------- פרק 24 | פסוק 2 ----------
 \textbf{ב} וְנִגַּ֨שׁ מֹשֶׁ֤ה לְבַדּוֹ֙ אֶל־יְהֹוָ֔ה וְהֵ֖ם לֹ֣א יִגָּ֑שׁוּ וְהָעָ֕ם לֹ֥א יַעֲל֖וּ עִמּֽוֹ׃ \Rashi{\textbf{ונגש משה לבדו.}אל הערפל:}%
% ---------- פרק 24 | פסוק 3 ----------
 \textbf{ג} וַיָּבֹ֣א מֹשֶׁ֗ה וַיְסַפֵּ֤ר לָעָם֙ אֵ֚ת כׇּל־דִּבְרֵ֣י יְהֹוָ֔ה וְאֵ֖ת כׇּל־הַמִּשְׁפָּטִ֑ים וַיַּ֨עַן כׇּל־הָעָ֜ם ק֤וֹל אֶחָד֙ וַיֹּ֣אמְר֔וּ כׇּל־הַדְּבָרִ֛ים אֲשֶׁר־דִּבֶּ֥ר יְהֹוָ֖ה נַעֲשֶֽׂה׃ \Rashi{\textbf{ויבא משה ויספר לעם.}בו ביום: את כל דברי ה'. מצות פרישה והגבלה: ואת כל המשפטים. שבע מצוות שנצטוו בני נח, ושבת וכבוד אב ואם ופרה אדמה ודינין שנתנו להם במרה (סנהדרין נ"ו):}%
% ---------- פרק 24 | פסוק 4 ----------
 \textbf{ד} וַיִּכְתֹּ֣ב מֹשֶׁ֗ה אֵ֚ת כׇּל־דִּבְרֵ֣י יְהֹוָ֔ה וַיַּשְׁכֵּ֣ם בַּבֹּ֔קֶר וַיִּ֥בֶן מִזְבֵּ֖חַ תַּ֣חַת הָהָ֑ר וּשְׁתֵּ֤ים עֶשְׂרֵה֙ מַצֵּבָ֔ה לִשְׁנֵ֥ים עָשָׂ֖ר שִׁבְטֵ֥י יִשְׂרָאֵֽל׃ \Rashi{\textbf{ויכתב משה.}מבראשית ועד מתן תורה, וכתב מצוות שנצטוו במרה: וישכם בבקר. בחמשה בסיון:}%
% ---------- פרק 24 | פסוק 5 ----------
 \textbf{ה} וַיִּשְׁלַ֗ח אֶֽת־נַעֲרֵי֙ בְּנֵ֣י יִשְׂרָאֵ֔ל וַֽיַּעֲל֖וּ עֹלֹ֑ת וַֽיִּזְבְּח֞וּ זְבָחִ֧ים שְׁלָמִ֛ים לַיהֹוָ֖ה פָּרִֽים׃ \Rashi{\textbf{את נערי.}הבכורות:}%
% ---------- פרק 24 | פסוק 6 ----------
 \textbf{ו} וַיִּקַּ֤ח מֹשֶׁה֙ חֲצִ֣י הַדָּ֔ם וַיָּ֖שֶׂם בָּאַגָּנֹ֑ת וַחֲצִ֣י הַדָּ֔ם זָרַ֖ק עַל־הַמִּזְבֵּֽחַ׃ \Rashi{\textbf{ויקח משה את חצי הדם.}מי חלקו? מלאך בא וחלקו (ויקרא רבה ו'): באגנת. שתי אגנות, אחד לחצי דם עולה ואחד לחצי דם שלמים להזות אותם על העם. ומכאן למדו רבותינו שנכנסו אבותינו לברית במילה וטבילה והזאת דמים, שאין הזאה בלא טבילה (כריתות ט'):}%
% ---------- פרק 24 | פסוק 7 ----------
 \textbf{ז} וַיִּקַּח֙ סֵ֣פֶר הַבְּרִ֔ית וַיִּקְרָ֖א בְּאׇזְנֵ֣י הָעָ֑ם וַיֹּ֣אמְר֔וּ כֹּ֛ל אֲשֶׁר־דִּבֶּ֥ר יְהֹוָ֖ה נַעֲשֶׂ֥ה וְנִשְׁמָֽע׃ \Rashi{\textbf{ספר הברית.}מבראשית ועד מתן תורה ומצוות שנצטוו במרה:}%
% ---------- פרק 24 | פסוק 8 ----------
 \textbf{ח} וַיִּקַּ֤ח מֹשֶׁה֙ אֶת־הַדָּ֔ם וַיִּזְרֹ֖ק עַל־הָעָ֑ם וַיֹּ֗אמֶר הִנֵּ֤ה דַֽם־הַבְּרִית֙ אֲשֶׁ֨ר כָּרַ֤ת יְהֹוָה֙ עִמָּכֶ֔ם עַ֥ל כׇּל־הַדְּבָרִ֖ים הָאֵֽלֶּה׃ \Rashi{\textbf{ויזרק.}ענין הזאה, ותרגומו וזרק על מדבחא לכפרא על עמא:}%
% ---------- פרק 24 | פסוק 9 ----------
 \textbf{ט} וַיַּ֥עַל מֹשֶׁ֖ה וְאַהֲרֹ֑ן נָדָב֙ וַאֲבִיה֔וּא וְשִׁבְעִ֖ים מִזִּקְנֵ֥י יִשְׂרָאֵֽל׃ %
% ---------- פרק 24 | פסוק 10 ----------
 \textbf{י} וַיִּרְא֕וּ אֵ֖ת אֱלֹהֵ֣י יִשְׂרָאֵ֑ל וְתַ֣חַת רַגְלָ֗יו כְּמַעֲשֵׂה֙ לִבְנַ֣ת הַסַּפִּ֔יר וּכְעֶ֥צֶם הַשָּׁמַ֖יִם לָטֹֽהַר׃ \Rashi{\textbf{ויראו את אלהי ישראל.}נסתכלו והציצו ונתחיבו מיתה, אלא שלא רצה הקב"ה לערבב שמחת התורה והמתין לנדב ואביהוא עד יום חנכת המשכן, ולזקנים עד ויהי העם כמתאוננים, ותבער בם אש ה' ותאכל בקצה המחנה (במדבר י"א) – בקצינים שבמחנה (תנחומא): כמעשה לבנת הספיר. היא היתה לפניו בשעת השעבוד, לזכר צרתן של ישראל שהיו משעבדים במעשה לבנים (ויקרא רבה): וכעצם השמים לטהר. משנגאלו היה אור וחדוה לפניו: וכעצם. כתרגומו, לשון מראה: לטהר. לשון ברור וצלול:}%
% ---------- פרק 24 | פסוק 11 ----------
 \textbf{יא} וְאֶל־אֲצִילֵי֙ בְּנֵ֣י יִשְׂרָאֵ֔ל לֹ֥א שָׁלַ֖ח יָד֑וֹ וַיֶּֽחֱזוּ֙ אֶת־הָ֣אֱלֹהִ֔ים וַיֹּאכְל֖וּ וַיִּשְׁתּֽוּ׃ \{ס\} \Rashi{\textbf{ואל אצילי.}הם נדב ואביהוא והזקנים: לא שלח ידו. מכלל שהיו ראויים להשתלח בהם יד (תנחומא, ויקרא רבה כ'): ויחזו את האלהים. היו מסתכלין בו בלב גס מתוך אכילה ושתיה, כך מדרש תנחומא, ואנקלוס לא תרגם כן. אצילי. לשון גדולים, כמו ומאציליה קראתיך (ישעיהו מ"א), ויאצל מן הרוח (במדבר י"א), שש אמות אצילה (יחזקאל מ"א):}%
% ---------- פרק 24 | פסוק 12 ----------
 \textbf{יב} וַיֹּ֨אמֶר יְהֹוָ֜ה אֶל־מֹשֶׁ֗ה עֲלֵ֥ה אֵלַ֛י הָהָ֖רָה וֶהְיֵה־שָׁ֑ם וְאֶתְּנָ֨ה לְךָ֜ אֶת־לֻחֹ֣ת הָאֶ֗בֶן וְהַתּוֹרָה֙ וְהַמִּצְוָ֔ה אֲשֶׁ֥ר כָּתַ֖בְתִּי לְהוֹרֹתָֽם׃ \Rashi{\textbf{ויאמר ה' אל משה.}לאחר מתן תורה: עלה אלי ההרה והיה שם. ארבעים יום: את לחת האבן והתורה והמצוה אשר כתבתי להורתם. כל שש מאות ושלש עשרה מצוות בכלל עשרת הדברות הן, ורבנו סעדיה פרש באזהרות שיסד לכל דבור ודבור מצוות התלויות בו:}%
% ---------- פרק 24 | פסוק 13 ----------
 \textbf{יג} וַיָּ֣קׇם מֹשֶׁ֔ה וִיהוֹשֻׁ֖עַ מְשָׁרְת֑וֹ וַיַּ֥עַל מֹשֶׁ֖ה אֶל־הַ֥ר הָאֱלֹהִֽים׃ \Rashi{\textbf{ויקם משה ויהושע משרתו.}לא ידעתי מה טיבו של יהושע כאן, ואומר אני, שהיה התלמיד מלוה לרב עד מקום הגבלת תחומי ההר שאינו רשאי לילך משם והלאה, ומשם ויעל משה לבדו אל הר האלהים, ויהושע נטה שם אהלו ונתעכב שם כל ארבעים יום, שכן מצינו כשירד משה וישמע יהושע את קול העם ברעה (שמות ל״ב:י״ז), למדנו שלא היה יהושע עמהם:}%
% ---------- פרק 24 | פסוק 14 ----------
 \textbf{יד} וְאֶל־הַזְּקֵנִ֤ים אָמַר֙ שְׁבוּ־לָ֣נוּ בָזֶ֔ה עַ֥ד אֲשֶׁר־נָשׁ֖וּב אֲלֵיכֶ֑ם וְהִנֵּ֨ה אַהֲרֹ֤ן וְחוּר֙ עִמָּכֶ֔ם מִי־בַ֥עַל דְּבָרִ֖ים יִגַּ֥שׁ אֲלֵהֶֽם׃ \Rashi{\textbf{ואל הזקנים אמר.}בצאתו מן המחנה: שבו לנו בזה. והתעכבו כאן עם שאר העם במחנה, להיות נכונים לשפט לכל איש ריבו: חור. בנה של מרים היה ואביו כלב בן יפנה, שנאמר ויקח לו כלב את אפרת ותלד לו את חור (דהי"א ב'), אפרת זו מרים כדאיתא בסוטה: מי בעל דברים. מי שיש לו דין:}%
% ---------- פרק 24 | פסוק 15 ----------
 \textbf{טו} וַיַּ֥עַל מֹשֶׁ֖ה אֶל־הָהָ֑ר וַיְכַ֥ס הֶעָנָ֖ן אֶת־הָהָֽר׃ %
% ---------- פרק 24 | פסוק 16 ----------
 \textbf{טז} וַיִּשְׁכֹּ֤ן כְּבוֹד־יְהֹוָה֙ עַל־הַ֣ר סִינַ֔י וַיְכַסֵּ֥הוּ הֶעָנָ֖ן שֵׁ֣שֶׁת יָמִ֑ים וַיִּקְרָ֧א אֶל־מֹשֶׁ֛ה בַּיּ֥וֹם הַשְּׁבִיעִ֖י מִתּ֥וֹךְ הֶעָנָֽן׃ \Rashi{\textbf{ויכסהו הענן.}רבותינו חולקים בדבר, יש מהם או' אלו ששה ימים שמר"ח (עד עצרת יום מתן תורה): ויכסהו הענן. להר: ויקרא אל משה ביום השביעי לומר עשרת הדברות, ומשה וכל בני ישראל עומדים, אלא שחלק הכתוב כבוד למשה; וי"א ויכסהו הענן למשה ששת ימים לאחר עשרת הדברות, והם היו בתחלת ארבעים יום שעלה משה לקבל לוחות, ולמדך שכל הנכנס למחנה שכינה טעון פרישה ששה ימים (יומא ג'):}%
% ---------- פרק 24 | פסוק 17 ----------
 \textbf{יז} וּמַרְאֵה֙ כְּב֣וֹד יְהֹוָ֔ה כְּאֵ֥שׁ אֹכֶ֖לֶת בְּרֹ֣אשׁ הָהָ֑ר לְעֵינֵ֖י בְּנֵ֥י יִשְׂרָאֵֽל׃ %
% ---------- פרק 24 | פסוק 18 ----------
 \textbf{יח} וַיָּבֹ֥א מֹשֶׁ֛ה בְּת֥וֹךְ הֶעָנָ֖ן וַיַּ֣עַל אֶל־הָהָ֑ר וַיְהִ֤י מֹשֶׁה֙ בָּהָ֔ר אַרְבָּעִ֣ים י֔וֹם וְאַרְבָּעִ֖ים לָֽיְלָה׃ \{פ\} \Rashi{\textbf{בתוך הענן.}ענן זה כמין עשן הוא, ועשה לו הקב"ה למשה שביל בתוכו (שם ד'):}%
\addparasha{חומש שמות | פרשת תרומה}
% ---------- פרק 25 | פסוק 1 ----------
 \textbf{א} וַיְדַבֵּ֥ר יְהֹוָ֖ה אֶל־מֹשֶׁ֥ה לֵּאמֹֽר׃ %
% ---------- פרק 25 | פסוק 2 ----------
 \textbf{ב} דַּבֵּר֙ אֶל־בְּנֵ֣י יִשְׂרָאֵ֔ל וְיִקְחוּ־לִ֖י תְּרוּמָ֑ה מֵאֵ֤ת כׇּל־אִישׁ֙ אֲשֶׁ֣ר יִדְּבֶ֣נּוּ לִבּ֔וֹ תִּקְח֖וּ אֶת־תְּרוּמָתִֽי׃ \Rashi{\textbf{ויקחו לי תרומה.}לי – לשמי: תרומה. הפרשה, יפרישו לי מממונם נדבה: ידבנו לבו. לשון נדבה והוא לשון רצון טוב, פיישנ"ט בלעז: תקחו את תרומתי. אמרו רבותינו שלש תרומות אמורות כאן, אחת תרומת בקע לגלגלת שנעשו מהם האדנים, כמו שמפרש באלה פקודי, ואחת תרומת המזבח בקע לגלגלת לקפות לקנות מהן קרבנות צבור, ואחת תרומת המשכן נדבת כל אחד ואחד (תלמוד ירושלמי שקלים א'). י"ג דברים האמורים בענין כלם הצרכו למלאכת המשכן או לבגדי כהנה כשתדקדק בהם:}%
% ---------- פרק 25 | פסוק 3 ----------
 \textbf{ג} וְזֹאת֙ הַתְּרוּמָ֔ה אֲשֶׁ֥ר תִּקְח֖וּ מֵאִתָּ֑ם זָהָ֥ב וָכֶ֖סֶף וּנְחֹֽשֶׁת׃ \Rashi{\textbf{זהב וכסף ונחשת וגו'.}כלם באו בנדבה – איש איש מה שנדבו לבו – חוץ מן הכסף שבא בשוה, מחצית השקל לכל אחד, ולא מצינו בכל מלאכת המשכן שהצרך שם כסף יותר, שנאמר וכסף פקודי העדה וגו' בקע לגלגלת וגו', ושאר הכסף הבא שם בנדבה עשאוהו לכלי שרת:}%
% ---------- פרק 25 | פסוק 4 ----------
 \textbf{ד} וּתְכֵ֧לֶת וְאַרְגָּמָ֛ן וְתוֹלַ֥עַת שָׁנִ֖י וְשֵׁ֥שׁ וְעִזִּֽים׃ \Rashi{\textbf{ותכלת.}צמר צבוע בדם חלזון וצבעו ירק (מנחות מ"ד): וארגמן. צמר צבוע ממין צבע ששמו ארגמן: ושש. הוא פשתן. ועזים. נוצה של עזים, לכך תרגם אנקלוס ומעזי – הבא מן העזים – ולא עזים עצמם, שתרגום של עזים עזיא:}%
% ---------- פרק 25 | פסוק 5 ----------
 \textbf{ה} וְעֹרֹ֨ת אֵילִ֧ם מְאׇדָּמִ֛ים וְעֹרֹ֥ת תְּחָשִׁ֖ים וַעֲצֵ֥י שִׁטִּֽים׃ \Rashi{\textbf{מאדמים.}צבועות היו אדם לאחר עבודן: תחשים. מין חיה, ולא היתה אלא לשעה, והרבה גונים היו לה, לכך מתרגם ססגונא ששש ומתפאר בגונין שלו (שבת כ"ח): ועצי שטים. ומאין היו להם במדבר? פרש רבי תנחומא: יעקב אבינו צפה ברוח הקדש שעתידין ישראל לבנות משכן במדבר, והביא ארזים למצרים ונטעם, וצוה לבניו לטלם עמהם כשיצאו ממצרים:}%
% ---------- פרק 25 | פסוק 6 ----------
 \textbf{ו} שֶׁ֖מֶן לַמָּאֹ֑ר בְּשָׂמִים֙ לְשֶׁ֣מֶן הַמִּשְׁחָ֔ה וְלִקְטֹ֖רֶת הַסַּמִּֽים׃ \Rashi{\textbf{שמן למאר.}שמן זית זך להעלות נר תמיד: בשמים לשמן המשחה. שנעשה למשח כלי המשכן והמשכן לקדשו, והצרכו לו בשמים, כמו שמפרש בכי תשא: ולקטרת הסמים. שהיו מקטירין בכל ערב ובקר, כמו שמפרש בואתה תצוה, ולשון קטרת העלאת קיטור ותמרת עשן:}%
% ---------- פרק 25 | פסוק 7 ----------
 \textbf{ז} אַבְנֵי־שֹׁ֕הַם וְאַבְנֵ֖י מִלֻּאִ֑ים לָאֵפֹ֖ד וְלַחֹֽשֶׁן׃ \Rashi{\textbf{אבני שהם.}שתים הצרכו שם לצרך האפוד האמור בואתה תצוה: מלאים. על שם שעושין להם בזהב מושב כמין גמא, ונותנין האבן שם למלאות הגמא, קרויים אבני מלואים, ומקום המושב קרוי משבצת: לאפד ולחשן. אבני השהם לאפוד ואבני המלאים לחשן; וחשן ואפוד מפרשים בואתה תצוה והם מיני תכשיט:}%
% ---------- פרק 25 | פסוק 8 ----------
 \textbf{ח} וְעָ֥שׂוּ לִ֖י מִקְדָּ֑שׁ וְשָׁכַנְתִּ֖י בְּתוֹכָֽם׃ \Rashi{\textbf{ועשו לי מקדש.}ועשו לשמי בית קדשה:}%
% ---------- פרק 25 | פסוק 9 ----------
 \textbf{ט} כְּכֹ֗ל אֲשֶׁ֤ר אֲנִי֙ מַרְאֶ֣ה אוֹתְךָ֔ אֵ֚ת תַּבְנִ֣ית הַמִּשְׁכָּ֔ן וְאֵ֖ת תַּבְנִ֣ית כׇּל־כֵּלָ֑יו וְכֵ֖ן תַּעֲשֽׂוּ׃ \{ס\} \Rashi{\textbf{ככל אשר אני מראה אותך.}כאן את תבנית המשכן, המקרא הזה מחבר למקרא שלמעלה הימנו ועשו לי מקדש ככל אשר אני מראה אותך: וכן תעשו. לדורות (סנהדרין ט"ז), אם יאבד אחד מן הכלים, או כשתעשו לי כלי בית עולמים, כגון שלחנות ומנורות וכיורות ומכונות שעשה שלמה, כתבנית אלו תעשו אותם; ואם לא היה המקרא מחבר למעלה הימנו, לא היה לו לכתב וכן תעשו אלא כן תעשו, והיה מדבר על עשית אהל מועד וכליו:}%
% ---------- פרק 25 | פסוק 10 ----------
 \textbf{י} וְעָשׂ֥וּ אֲר֖וֹן עֲצֵ֣י שִׁטִּ֑ים אַמָּתַ֨יִם וָחֵ֜צִי אׇרְכּ֗וֹ וְאַמָּ֤ה וָחֵ֙צִי֙ רׇחְבּ֔וֹ וְאַמָּ֥ה וָחֵ֖צִי קֹמָתֽוֹ׃ \Rashi{\textbf{ועשו ארון.}כמין ארונות שעושים בלא רגלים, עשוים כמין ארגז שקורין אישקרי"ן בלעז, יושב על שוליו:}%
% ---------- פרק 25 | פסוק 11 ----------
 \textbf{יא} וְצִפִּיתָ֤ אֹתוֹ֙ זָהָ֣ב טָה֔וֹר מִבַּ֥יִת וּמִח֖וּץ תְּצַפֶּ֑נּוּ וְעָשִׂ֧יתָ עָלָ֛יו זֵ֥ר זָהָ֖ב סָבִֽיב׃ \Rashi{\textbf{מבית ומחוץ תצפנו.}שלשה ארונות עשה בצלאל, שנים של זהב ואחד של עץ, ארבע כתלים ושולים לכל אחד ופתוחים מלמעלה, נתן של עץ בתוך של זהב ושל זהב בתוך של עץ וחפה שפתו העליונה בזהב, נמצא מצפה מבית ומחוץ (יומא ע"ב): זר זהב. כמין כתר מקף לו סביב למעלה משפתו, שעשה הארון החיצון גבוה מן הפנימי, עד שעלה למול עבי הכפרת ולמעלה הימנו משהו, וכשהכפרת שוכב על עבי הכתלים, עולה הזר למעלה מכל עבי הכפרת כל שהוא והוא סימן לכתר תורה:}%
% ---------- פרק 25 | פסוק 12 ----------
 \textbf{יב} וְיָצַ֣קְתָּ לּ֗וֹ אַרְבַּע֙ טַבְּעֹ֣ת זָהָ֔ב וְנָ֣תַתָּ֔ה עַ֖ל אַרְבַּ֣ע פַּעֲמֹתָ֑יו וּשְׁתֵּ֣י טַבָּעֹ֗ת עַל־צַלְעוֹ֙ הָֽאֶחָ֔ת וּשְׁתֵּי֙ טַבָּעֹ֔ת עַל־צַלְע֖וֹ הַשֵּׁנִֽית׃ \Rashi{\textbf{ויצקת.}לשון התכה, כתרגומו: פעמותיו. כתרגומו זויתיה; ובזויות העליונות סמוך לכפרת היו נתונות, שתים מכאן ושתים מכאן לרחבו של ארון, והבדים נתונים בהם, וארכו של ארון מפסיק בין הבדים, אמתים וחצי בין בד לבד, שיהיו שני בני אדם הנושאין את הארון מהלכין ביניהם; וכן מפרש במנחות בפ' שתי הלחם (מנחות דף צ"ח): ושתי טבעת על צלעו האחת. הן הן ארבע טבעות שבתחלת המקרא, ופרש לך היכן היו; והוי"ו זו יתרה היא, ופתרונו כמו שתי טבעות, ויש לך לישבה כן: ושתי מן הטבעות האלו על צלעו האחת: צלעו. צדו:}%
% ---------- פרק 25 | פסוק 13 ----------
 \textbf{יג} וְעָשִׂ֥יתָ בַדֵּ֖י עֲצֵ֣י שִׁטִּ֑ים וְצִפִּיתָ֥ אֹתָ֖ם זָהָֽב׃ \Rashi{\textbf{בדי.}מוטות:}%
% ---------- פרק 25 | פסוק 14 ----------
 \textbf{יד} וְהֵבֵאתָ֤ אֶת־הַבַּדִּים֙ בַּטַּבָּעֹ֔ת עַ֖ל צַלְעֹ֣ת הָאָרֹ֑ן לָשֵׂ֥את אֶת־הָאָרֹ֖ן בָּהֶֽם׃ %
% ---------- פרק 25 | פסוק 15 ----------
 \textbf{טו} בְּטַבְּעֹת֙ הָאָרֹ֔ן יִהְי֖וּ הַבַּדִּ֑ים לֹ֥א יָסֻ֖רוּ מִמֶּֽנּוּ׃ \Rashi{\textbf{לא יסרו ממנו.}לעולם (יומא ע"ב):}%
% ---------- פרק 25 | פסוק 16 ----------
 \textbf{טז} וְנָתַתָּ֖ אֶל־הָאָרֹ֑ן אֵ֚ת הָעֵדֻ֔ת אֲשֶׁ֥ר אֶתֵּ֖ן אֵלֶֽיךָ׃ \Rashi{\textbf{ונתת אל הארן.}כמו בארון: העדת. התורה, שהיא לעדות ביני וביניכם שצויתי אתכם מצוות הכתובות בה:}%
% ---------- פרק 25 | פסוק 17 ----------
 \textbf{יז} וְעָשִׂ֥יתָ כַפֹּ֖רֶת זָהָ֣ב טָה֑וֹר אַמָּתַ֤יִם וָחֵ֙צִי֙ אׇרְכָּ֔הּ וְאַמָּ֥ה וָחֵ֖צִי רׇחְבָּֽהּ׃ \Rashi{\textbf{כפרת.}כסוי על הארון, שהיה פתוח מלמעלה ומניחו עליו כמין דף: אמתים וחצי ארכה. כארכו של ארון, ורחבה כרחבו של ארון, ומנחת על עבי הכתלים ארבעתם, ואע"פ שלא נתן שעור לעביה, פרשו רבותינו שהיה עביה טפח (סוכה ה'):}%
% ---------- פרק 25 | פסוק 18 ----------
 \textbf{יח} וְעָשִׂ֛יתָ שְׁנַ֥יִם כְּרֻבִ֖ים זָהָ֑ב מִקְשָׁה֙ תַּעֲשֶׂ֣ה אֹתָ֔ם מִשְּׁנֵ֖י קְצ֥וֹת הַכַּפֹּֽרֶת׃ \Rashi{\textbf{כרבים.}דמות פרצוף תינוק להם: מקשה תעשה. שלא תעשם בפני עצמם ותחברם בראשי הכפרת לאחר עשיתם כמעשה צורפים שקורין שולד"יר בלעז, אלא הטל זהב הרבה בתחלת עשית הכפרת, והכה בפטיש ובקרנס באמצע, וראשין בולטין למעלה, וציר הכרובים בבליטת קצותיו: מקשה. בטדי"ץ בלעז, כמו דא לדא נקשן (דניאל ה'): קצות הכפרת. ראשי הכפרת:}%
% ---------- פרק 25 | פסוק 19 ----------
 \textbf{יט} וַ֠עֲשֵׂ֠ה כְּר֨וּב אֶחָ֤ד מִקָּצָה֙ מִזֶּ֔ה וּכְרוּב־אֶחָ֥ד מִקָּצָ֖ה מִזֶּ֑ה מִן־הַכַּפֹּ֛רֶת תַּעֲשׂ֥וּ אֶת־הַכְּרֻבִ֖ים עַל־שְׁנֵ֥י קְצוֹתָֽיו׃ \Rashi{\textbf{ועשה כרוב אחד מקצה.}שלא תאמר שנים כרובים לכל קצה וקצה, לכך הצרך לפרש כרוב אחד מקצה מזה: מן הכפרת. עצמה תעשה את הכרובים, זהו פרוש של מקשה תעשה אותם, שלא תעשה בפני עצמם ותחברם לכפרת:}%
% ---------- פרק 25 | פסוק 20 ----------
 \textbf{כ} וְהָי֣וּ הַכְּרֻבִים֩ פֹּרְשֵׂ֨י כְנָפַ֜יִם לְמַ֗עְלָה סֹכְכִ֤ים בְּכַנְפֵיהֶם֙ עַל־הַכַּפֹּ֔רֶת וּפְנֵיהֶ֖ם אִ֣ישׁ אֶל־אָחִ֑יו אֶ֨ל־הַכַּפֹּ֔רֶת יִהְי֖וּ פְּנֵ֥י הַכְּרֻבִֽים׃ \Rashi{\textbf{פרשי כנפים.}שלא תעשה כנפיהם שוכבים אלא פרושים וגבוהים למעלה אצל ראשיהם, שיהא עשרה טפחים בחלל שבין הכנפים לכפרת, כדאיתא בסכה:}%
% ---------- פרק 25 | פסוק 21 ----------
 \textbf{כא} וְנָתַתָּ֧ אֶת־הַכַּפֹּ֛רֶת עַל־הָאָרֹ֖ן מִלְמָ֑עְלָה וְאֶל־הָ֣אָרֹ֔ן תִּתֵּן֙ אֶת־הָ֣עֵדֻ֔ת אֲשֶׁ֥ר אֶתֵּ֖ן אֵלֶֽיךָ׃ \Rashi{\textbf{ואל הארן תתן את העדת.}לא ידעתי למה נכפל, שהרי כבר נאמר ונתת אל הארן את העדת, ויש לומר, שבא ללמד שבעודו ארון לבדו, בלא כפרת, יתן תחלה העדות לתוכו ואחר כך יתן את הכפרת עליו, וכן מצינו כשהקים את המשכן, נאמר ויתן את העדת אל הארן ואחר כך ויתן את הכפרת על הארן מלמעלה:}%
% ---------- פרק 25 | פסוק 22 ----------
 \textbf{כב} וְנוֹעַדְתִּ֣י לְךָ֮ שָׁם֒ וְדִבַּרְתִּ֨י אִתְּךָ֜ מֵעַ֣ל הַכַּפֹּ֗רֶת מִבֵּין֙ שְׁנֵ֣י הַכְּרֻבִ֔ים אֲשֶׁ֖ר עַל־אֲר֣וֹן הָעֵדֻ֑ת אֵ֣ת כׇּל־אֲשֶׁ֧ר אֲצַוֶּ֛ה אוֹתְךָ֖ אֶל־בְּנֵ֥י יִשְׂרָאֵֽל׃ \{פ\} \Rashi{\textbf{ונועדתי.}כשאקבע מועד לך לדבר עמך, אותו מקום אקבע למועד שאבא שם לדבר אליך: ודברתי אתך מעל הכפרת. ובמקום אחר הוא אומר וידבר ה' אליו מאהל מועד לאמר (ויקרא א'), זה המשכן, מחוץ לפרכת, נמצאו שני כתובים מכחישים זה את זה, – בא הכתוב השלישי והכריע ביניהם: ובבא משה אל אהל מועד, וישמע את הקול מדבר אליו מעל הכפרת וגו' (במדבר ז'), משה היה נכנס למשכן, וכיון שבא בתוך הפתח, קול יורד מן השמים לבין הכרובים, ומשם יוצא ונשמע למשה באהל מועד (ספרי): ואת כל אשר אצוה אותך אל בני ישראל. הרי וי"ו זו יתרה וטפלה, וכמוה הרבה במקרא, וכה תפתר: ואת אשר אדבר עמך שם את כל אשר אצוה אותך אל בני ישראל הוא:}%
% ---------- פרק 25 | פסוק 23 ----------
 \textbf{כג} וְעָשִׂ֥יתָ שֻׁלְחָ֖ן עֲצֵ֣י שִׁטִּ֑ים אַמָּתַ֤יִם אׇרְכּוֹ֙ וְאַמָּ֣ה רׇחְבּ֔וֹ וְאַמָּ֥ה וָחֵ֖צִי קֹמָתֽוֹ׃ \Rashi{\textbf{קמתו.}גבה רגליו עם עבי השלחן:}%
% ---------- פרק 25 | פסוק 24 ----------
 \textbf{כד} וְצִפִּיתָ֥ אֹת֖וֹ זָהָ֣ב טָה֑וֹר וְעָשִׂ֥יתָ לּ֛וֹ זֵ֥ר זָהָ֖ב סָבִֽיב׃ \Rashi{\textbf{זר זהב.}סימן לכתר מלכות, שהשלחן שם עשר וגדלה, כמו שאומרים שלחן מלכים:}%
% ---------- פרק 25 | פסוק 25 ----------
 \textbf{כה} וְעָשִׂ֨יתָ לּ֥וֹ מִסְגֶּ֛רֶת טֹ֖פַח סָבִ֑יב וְעָשִׂ֧יתָ זֵר־זָהָ֛ב לְמִסְגַּרְתּ֖וֹ סָבִֽיב׃ \Rashi{\textbf{מסגרת.}כתרגומו גדנפא; ונחלקו חכמי ישראל בדבר, י"א למעלה היתה סביב לשלחן כמו לבזבזין שבשפת שלחן שרים, וי"א למטה היתה תקועה מרגל לרגל בארבע רוחות השלחן, ודף השלחן שוכב על אותה מסגרת: ועשית זר זהב למסגרתו. הוא זר האמור למעלה, ופרש לך כאן שעל המסגרת היתה:}%
% ---------- פרק 25 | פסוק 26 ----------
 \textbf{כו} וְעָשִׂ֣יתָ לּ֔וֹ אַרְבַּ֖ע טַבְּעֹ֣ת זָהָ֑ב וְנָתַתָּ֙ אֶת־הַטַּבָּעֹ֔ת עַ֚ל אַרְבַּ֣ע הַפֵּאֹ֔ת אֲשֶׁ֖ר לְאַרְבַּ֥ע רַגְלָֽיו׃ %
% ---------- פרק 25 | פסוק 27 ----------
 \textbf{כז} לְעֻמַּת֙ הַמִּסְגֶּ֔רֶת תִּהְיֶ֖יןָ הַטַּבָּעֹ֑ת לְבָתִּ֣ים לְבַדִּ֔ים לָשֵׂ֖את אֶת־הַשֻּׁלְחָֽן׃ \Rashi{\textbf{לעמת המסגרת תהיין הטבעת.}ברגלים תקועות, כנגד ראשי המסגרת: לבתים לבדים. אותן טבעות יהיו בתים להכניס בהן הבדים: לבתים. לצרך בתים לבדים, כתרגומו – אתרא לאריחיא:}%
% ---------- פרק 25 | פסוק 28 ----------
 \textbf{כח} וְעָשִׂ֤יתָ אֶת־הַבַּדִּים֙ עֲצֵ֣י שִׁטִּ֔ים וְצִפִּיתָ֥ אֹתָ֖ם זָהָ֑ב וְנִשָּׂא־בָ֖ם אֶת־הַשֻּׁלְחָֽן׃ \Rashi{\textbf{ונשא בם.}לשון נפעל, יהיה נשא בם את השלחן:}%
% ---------- פרק 25 | פסוק 29 ----------
 \textbf{כט} וְעָשִׂ֨יתָ קְּעָרֹתָ֜יו וְכַפֹּתָ֗יו וּקְשׂוֹתָיו֙ וּמְנַקִּיֹּתָ֔יו אֲשֶׁ֥ר יֻסַּ֖ךְ בָּהֵ֑ן זָהָ֥ב טָה֖וֹר תַּעֲשֶׂ֥ה אֹתָֽם׃ \Rashi{\textbf{ועשית קערתיו וכפתיו.}קערותיו זה דפוס שהיה עשוי כדפוס הלחם, והלחם היה עשוי כמין תבה פרוצה משתי רוחותיה, שולים לו למטה, וקופל מכאן ומכאן כלפי מעלה כמין כתלים, ולכך קרוי לחם הפנים, שיש לו פנים רואין לכאן ולכאן לצדי הבית מזה ומזה, ונותן ארכו לרחבו של שלחן, וכתליו זקופים כנגד שפת השלחן; והיה עשוי לו דפוס זהב ודפוס ברזל, בשל ברזל הוא נאפה, וכשמוציאו מן התנור נותנו בשל זהב עד למחר בשבת שמסדרו על השלחן, ואותו דפוס קרוי קערה: וכפתיו. בזיכין שנותנין בהם לבונה שתים היו לשני קמצי לבונה שנותנין על שתי המערכות, שנאמר ונתת על המערכת לבנה זכה (ויקרא כ"ד): וקשותיו. הן כמין חצאי קנים חלולים הנסדקין לארכן (מנחות צ"ו), דגמתן עושה של זהב, ומסדר ג' על ראש כל לחם, שישב לחם האחר על גבי אותן הקנים, ומבדילין בין לחם ללחם כדי שתכנס הרוח ביניהם ולא יתעפשו; ובלשון ערבי כל דבר חלול קרוי קס"וא: ומנקיתיו. תרגומו ומכילתיה, הן סניפים כמין יתדות זהב, עומדין בארץ וגבוהים עד למעלה מן השלחן הרבה, כנגד גבה מערכת הלחם, ומפצלים חמשה פצולים זה למעלה מזה, וראשי הקנים שבין לחם ללחם סמוכין על אותן פצולין, כדי שלא יכבד משא הלחם העליונים על התחתונים וישברו; ולשון מכילתיה סבלותיו, כמו נלאיתי הכיל (ירמיהו ו'). אבל לשון מנקיות איני יודע איך נופל על סניפין; ויש מחכמי ישראל אומרים קשתיו אלו סניפין, שמקשין אותו ומחזיקים אותו שלא ישבר, ומנקיתיו אלו הקנים שמנקין אותו שלא יתעפש (מנחות צ"ז), אבל אנקלוס שתרגם ומכילתיה, היה שונה כדברי האומר מנקיות הן סניפין: אשר יסך בהן. אשר יכסה בהן, ועל קשתיו הוא אומר אשר יסך, שהיו עליו כמין סכך וכסוי, וכן במקום אחר הוא אומר ואת קשות הנסך (במדבר ד'), וזה וזה, יסך והנסך, לשון סכך וכסוי הם:}%
% ---------- פרק 25 | פסוק 30 ----------
 \textbf{ל} וְנָתַתָּ֧ עַֽל־הַשֻּׁלְחָ֛ן לֶ֥חֶם פָּנִ֖ים לְפָנַ֥י תָּמִֽיד׃ \{פ\} \Rashi{\textbf{לחם פנים.}שהיו לו פנים כמו שפרשתי, ומנין הלחם וסדר מערכותיו מפרשים באמר אל הכהנים:}%
% ---------- פרק 25 | פסוק 31 ----------
 \textbf{לא} וְעָשִׂ֥יתָ מְנֹרַ֖ת זָהָ֣ב טָה֑וֹר מִקְשָׁ֞ה תֵּעָשֶׂ֤ה*(בספרי ספרד ואשכנז תֵּיעָשֶׂ֤ה) הַמְּנוֹרָה֙ יְרֵכָ֣הּ וְקָנָ֔הּ גְּבִיעֶ֛יהָ כַּפְתֹּרֶ֥יהָ וּפְרָחֶ֖יהָ מִמֶּ֥נָּה יִהְיֽוּ׃ \Rashi{\textbf{מקשה תיעשה המנורה.}שלא יעשנה חליות ולא יעשה קניה ונרותיה אברים אברים ואחר כך ידביקם, כדרך הצורפים שקורין שולד"יר בלעז, אלא כלה באה מחתיכה אחת, ומקיש בקרנס וחותך בכלי האמנות ומפריד הקנים אילך ואילך: מקשה. תרגומו נגיד, לשון המשכה, שממשיך את האברים מן העשת לכאן ולכאן בהקשת הקרנס, ולשון מקשה מכת קרנס, בטדי"ץ בלעז, כמו דא לדא נקשן (דניאל ה'): תיעשה המנורה. מאליה, לפי שהיה משה מתקשה בה, אמר לו הקב"ה השלך את הככר לאור והיא נעשית מאליה, לכך לא נכתב תעשה (תנחומא): ירכה. הוא הרגל שלמטה העשוי כמין תבה, ושלשה רגלים יוצאין הימנו ולמטה: וקנה. הקנה האמצעי שלה העולה באמצע הירך זקוף כלפי מעלה, ועליו נר האמצעי עשוי כמין בזך, לצוק השמן לתוכו ולתת הפתילה: גביעיה. הן כמין כוסות שעושין מזכוכית, ארכים וקצרים, וקורין להם מדיר"נש בלעז, ואלו עשויין של זהב, ובולטין ויוצאין מכל קנה וקנה כמנין שנתן בהם הכתוב, ולא היו בה אלא לנוי: כפתריה. כמין תפוחים היו, עגלין סביב, בולטין סביבות הקנה האמצעי, כדרך שעושין למנורות שלפני השרים, וקורין להם פומיל"ש בלעז, ומנין שלהם כתוב בפרשה, כמה כפתורים בולטין ממנה וכמה חלק בין כפתור לכפתור: ופרחיה. ציורין עשוין בה כמין פרחין: ממנה יהיו. הכל מקשה יוצא מתוך חתיכת העשת, ולא יעשה לבדם וידביקם:}%
% ---------- פרק 25 | פסוק 32 ----------
 \textbf{לב} וְשִׁשָּׁ֣ה קָנִ֔ים יֹצְאִ֖ים מִצִּדֶּ֑יהָ שְׁלֹשָׁ֣ה ׀ קְנֵ֣י מְנֹרָ֗ה מִצִּדָּהּ֙ הָאֶחָ֔ד וּשְׁלֹשָׁה֙ קְנֵ֣י מְנֹרָ֔ה מִצִּדָּ֖הּ הַשֵּׁנִֽי׃ \Rashi{\textbf{יצאים מצדיה.}לכאן ולכאן באלכסון, נמשכים ועולין עד כנגד גבהה של מנורה שהוא קנה האמצעי, ויוצאין מתוך קנה האמצעי זה למעלה מזה, התחתון ארך, ושל מעלה קצר הימנו, והעליון קצר הימנו, לפי שהיה גבה ראשיהן שוה לגבהו של קנה האמצעי, השביעי, שממנו יוצאים הששה קנים:}%
% ---------- פרק 25 | פסוק 33 ----------
 \textbf{לג} שְׁלֹשָׁ֣ה גְ֠בִעִ֠ים מְֽשֻׁקָּדִ֞ים בַּקָּנֶ֣ה הָאֶחָד֮ כַּפְתֹּ֣ר וָפֶ֒רַח֒ וּשְׁלֹשָׁ֣ה גְבִעִ֗ים מְשֻׁקָּדִ֛ים בַּקָּנֶ֥ה הָאֶחָ֖ד כַּפְתֹּ֣ר וָפָ֑רַח כֵּ֚ן לְשֵׁ֣שֶׁת הַקָּנִ֔ים הַיֹּצְאִ֖ים מִן־הַמְּנֹרָֽה׃ \Rashi{\textbf{משקדים.}כתרגומו, מצירים היו, כדרך שעושין לכלי כסף וזהב שקורים נייל"יר בלעז: שלשה גבעים. בולטין מכל קנה וקנה: כפתר ופרח. היה לכל קנה וקנה:}%
% ---------- פרק 25 | פסוק 34 ----------
 \textbf{לד} וּבַמְּנֹרָ֖ה אַרְבָּעָ֣ה גְבִעִ֑ים מְשֻׁ֨קָּדִ֔ים כַּפְתֹּרֶ֖יהָ וּפְרָחֶֽיהָ׃ \Rashi{\textbf{ובמנרה ארבעה גבעים.}בגופה של מנורה היו ארבעה גביעים – אחד בולט בה למטה מן הקנים, והשלשה למעלה מן יציאת הקנים היוצאין מצדיה: משקדים כפתריה ופרחיה. זה אחד מחמשה מקראות שאין להם הכרע (יומא נ"ב), אין ידוע אם גביעים משקדים או משקדים כפתוריה ופרחיה:}%
% ---------- פרק 25 | פסוק 35 ----------
 \textbf{לה} וְכַפְתֹּ֡ר תַּ֩חַת֩ שְׁנֵ֨י הַקָּנִ֜ים מִמֶּ֗נָּה וְכַפְתֹּר֙ תַּ֣חַת שְׁנֵ֤י הַקָּנִים֙ מִמֶּ֔נָּה וְכַפְתֹּ֕ר תַּחַת־שְׁנֵ֥י הַקָּנִ֖ים מִמֶּ֑נָּה לְשֵׁ֙שֶׁת֙ הַקָּנִ֔ים הַיֹּצְאִ֖ים מִן־הַמְּנֹרָֽה׃ \Rashi{\textbf{וכפתר תחת שני הקנים.}מתוך הכפתור היו הקנים נמשכים משני צדיה אילך ואילך, כך שנינו במלאכת המשכן: גבהה של מנורה י"ח טפחים, הרגלים והפרח שלשה טפחים – הוא הפרח האמור בירך, שנאמר עד ירכה עד פרחה (במדבר ח') – וטפחים חלק, וטפח שבו גביע מהארבעה גביעים וכפתור ופרח משני כפתורים ושני פרחים האמורים במנורה עצמה, שנ' משקדים כפתוריה ופרחיה, למדנו שהיו בקנה שני כפתורים ושני פרחים לבד מן השלשה כפתורים שהקנים נמשכין מתוכן, שנ' וכפתר תחת שני הקנים וגו', וטפחים חלק, וטפח כפתור ושני קנים יוצאים ממנו אילך ואילך נמשכים ועולים כנגד גבהה של מנורה, טפח חלק, וטפח כפתור ושני קנים יוצאים ממנו, וטפח חלק, וטפח כפתור ושני קנים יוצאים ממנו ונמשכים ועולין כנגד גבהה של מנורה, וטפחים חלק, נשתירו שם שלשה טפחים שבהם שלשה גביעים וכפתור ופרח, נמצאו גביעים כ"ב – י"ח לששה קנים, שלשה לכל אחד ואחד, וארבעה בגופה של מנורה, הרי כ"ב – ואחד עשר כפתורים, ששה בששת הקנים, ושלשה בגופה של מנורה שהקנים יוצאים מהם, ושנים עוד במנורה שנ' משקדים כפתוריה, ומעוט כפתורים שנים, האחד למטה אצל הירך, והאחד בג' טפחים העליונים עם ג' הגביעים, ותשעה פרחים היו לה, ששה לששת הקנים, שנ' בקנה האחד כפתר ופרח, ושלשה למנורה שנ' משקדים כפתוריה ופרחיה, ומעוט פרחים שנים, ואחד האמור בפרשת בהעלתך עד ירכה עד פרחה, ואם תדקדק במשנה זו הכתובה למעלה, תמצאם כמנינם איש איש במקומו:}%
% ---------- פרק 25 | פסוק 36 ----------
 \textbf{לו} כַּפְתֹּרֵיהֶ֥ם וּקְנֹתָ֖ם מִמֶּ֣נָּה יִהְי֑וּ כֻּלָּ֛הּ מִקְשָׁ֥ה אַחַ֖ת זָהָ֥ב טָהֽוֹר׃ %
% ---------- פרק 25 | פסוק 37 ----------
 \textbf{לז} וְעָשִׂ֥יתָ אֶת־נֵרֹתֶ֖יהָ שִׁבְעָ֑ה וְהֶֽעֱלָה֙ אֶת־נֵ֣רֹתֶ֔יהָ וְהֵאִ֖יר עַל־עֵ֥בֶר פָּנֶֽיהָ׃ \Rashi{\textbf{את נרתיה.}כמין בזיכין שנותנין בתוכן השמן והפתילות: והאיר על עבר פניה. עשה פי ששת הנרות שבראשי הקנים היוצאים מצדיה מסבים כלפי האמצעי, כדי שיהיו הנרות כשתדליקם מאירים אל עבר פניה – מסב אורם אל צד פני הקנה האמצעי, שהוא גוף המנורה:}%
% ---------- פרק 25 | פסוק 38 ----------
 \textbf{לח} וּמַלְקָחֶ֥יהָ וּמַחְתֹּתֶ֖יהָ זָהָ֥ב טָהֽוֹר׃ \Rashi{\textbf{ומלקחיה.}הם הצבתים העשויין לקח בהם הפתילה מתוך השמן, לישבן ולמשכן בפי הנרות, ועל שם שלוקחים בהם קרויים מלקחים; וצבתהא שתרגם אנקלוס לשון צבת, טנייל"ש בלעז: ומחתתיה. הם כמין בזיכין קטנים שחותה בהן את האפר שבנר בבקר בבקר, כשהוא מיטיב את הנרות מאפר הפתילות שדלקו הלילה וכבו, ולשון מחתה פויישר"יא בלעז, כמו לחתות אש מיקוד (ישעיהו ל'):}%
% ---------- פרק 25 | פסוק 39 ----------
 \textbf{לט} כִּכָּ֛ר זָהָ֥ב טָה֖וֹר יַעֲשֶׂ֣ה אֹתָ֑הּ אֵ֥ת כׇּל־הַכֵּלִ֖ים הָאֵֽלֶּה׃ \Rashi{\textbf{ככר זהב טהור.}שלא יהיה משקלה עם כל כליה אלא ככר, לא פחות ולא יותר (מנחות פ"ח). והככר של חל ששים מנה, ושל קדש היה כפול, ק"ך מנה, והמנה הוא ליטרא ששוקלין בה כסף למשקל קולוניא, והם מאה זהובים – עשרים וחמשה סלעים והסלע ארבעה זהובים:}%
% ---------- פרק 25 | פסוק 40 ----------
 \textbf{מ} וּרְאֵ֖ה וַעֲשֵׂ֑ה בְּתַ֨בְנִיתָ֔ם אֲשֶׁר־אַתָּ֥ה מׇרְאֶ֖ה בָּהָֽר׃ \{ס\} \Rashi{\textbf{וראה ועשה.}ראה כאן בהר תבנית שאני מראה אותך, מגיד שנתקשה משה במעשה המנורה עד שהראה לו הקב"ה מנורה של אש (שם כ"ט): אשר אתה מראה. כתרגומו דאת מתחזי בטורא; אלו היה נקוד מראה, בפתח, היה פתרונו אתה מראה לאחרים, עכשו שנקוד חטף קמץ, פתרונו דאת מתחזי – שאחרים מראים לך (שהנקוד מפריד בין עושה לנעשה):}%
% ---------- פרק 26 | פסוק 1 ----------
 \textbf{א} וְאֶת־הַמִּשְׁכָּ֥ן תַּעֲשֶׂ֖ה עֶ֣שֶׂר יְרִיעֹ֑ת שֵׁ֣שׁ מׇשְׁזָ֗ר וּתְכֵ֤לֶת וְאַרְגָּמָן֙ וְתֹלַ֣עַת שָׁנִ֔י כְּרֻבִ֛ים מַעֲשֵׂ֥ה חֹשֵׁ֖ב תַּעֲשֶׂ֥ה אֹתָֽם׃ \Rashi{\textbf{ואת המשכן תעשה עשר יריעת.}להיות לו לגג ולמחצות מחוץ לקרשים, שהיריעות תלויות מאחוריהן לכסותן: שש משזר ותכלת ותלעת שני. הרי ארבעה מינין יחד חוט וחוט – אחד של פשתן ושלשה של צמר – וכל חוט וחוט כפול ו' הרי ארבעה מינין כשהן שזורין יחד כ"ד כפלים לחוט (יומא ע"א): כרבים מעשה חשב. כרובים היו מצירין בהם באריגתן, ולא ברקימה שהוא מעשה מחט, אלא באריגה בשני כתלים, פרצוף אחד מכאן ופרצוף אחד מכאן – ארי מצד זה ונשר מצד זה – כמו שאורגין חגורות של משי שקורין בלעז פיישי"שא:}%
% ---------- פרק 26 | פסוק 2 ----------
 \textbf{ב} אֹ֣רֶךְ ׀ הַיְרִיעָ֣ה הָֽאַחַ֗ת שְׁמֹנֶ֤ה וְעֶשְׂרִים֙ בָּֽאַמָּ֔ה וְרֹ֙חַב֙ אַרְבַּ֣ע בָּאַמָּ֔ה הַיְרִיעָ֖ה הָאֶחָ֑ת מִדָּ֥ה אַחַ֖ת לְכׇל־הַיְרִיעֹֽת׃ %
% ---------- פרק 26 | פסוק 3 ----------
 \textbf{ג} חֲמֵ֣שׁ הַיְרִיעֹ֗ת תִּֽהְיֶ֙יןָ֙ חֹֽבְרֹ֔ת אִשָּׁ֖ה אֶל־אֲחֹתָ֑הּ וְחָמֵ֤שׁ יְרִיעֹת֙ חֹֽבְרֹ֔ת אִשָּׁ֖ה אֶל־אֲחֹתָֽהּ׃ \Rashi{\textbf{תהיין חברת.}תופרן במחט זו בצד זו, חמש לבד וחמש לבד: אשה אל אחתה. כך דרך המקרא לדבר בדבר שהוא לשון נקבה, ובדבר שהוא לשון זכר אומר איש אל אחיו, כמו שנא' בכרובים ופניהם איש אל אחיו (שמות כ"ה):}%
% ---------- פרק 26 | פסוק 4 ----------
 \textbf{ד} וְעָשִׂ֜יתָ לֻֽלְאֹ֣ת תְּכֵ֗לֶת עַ֣ל שְׂפַ֤ת הַיְרִיעָה֙ הָאֶחָ֔ת מִקָּצָ֖ה בַּחֹבָ֑רֶת וְכֵ֤ן תַּעֲשֶׂה֙ בִּשְׂפַ֣ת הַיְרִיעָ֔ה הַקִּ֣יצוֹנָ֔ה בַּמַּחְבֶּ֖רֶת הַשֵּׁנִֽית׃ \Rashi{\textbf{ללאת.}לצול"ש בלעז, וכן תרגם אנקלוס ענובין, לשון עניבה: מקצה בחברת. באותה יריעה שבסוף החבור, קבוצת חמשת היריעות קרויה חוברת: וכן תעשה בשפת היריעה הקיצונה במחברת השנית. באותה יריעה שהיא קיצונה, לשון קצה, כלומר לסוף החוברת:}%
% ---------- פרק 26 | פסוק 5 ----------
 \textbf{ה} חֲמִשִּׁ֣ים לֻֽלָאֹ֗ת תַּעֲשֶׂה֮ בַּיְרִיעָ֣ה הָאֶחָת֒ וַחֲמִשִּׁ֣ים לֻֽלָאֹ֗ת תַּעֲשֶׂה֙ בִּקְצֵ֣ה הַיְרִיעָ֔ה אֲשֶׁ֖ר בַּמַּחְבֶּ֣רֶת הַשֵּׁנִ֑ית מַקְבִּילֹת֙ הַלֻּ֣לָאֹ֔ת אִשָּׁ֖ה אֶל־אֲחֹתָֽהּ׃ \Rashi{\textbf{מקבילת הללאת אשה אל אחתה.}שמר שתעשה הלולאות במדה אחת, מכונת הבדלתן זו מזו, וכמדתן ביריעה זו כן יהא בחברתה, שכשתפרש חוברת אצל חוברת יהיו הלולאות של יריעה זו מכונות כנגד לולאות של זו, וזהו לשון מקבילות – זו כנגד זו, תרגומו של נגד לקבל. היריעות ארכן כ"ח ורחבן ארבע, וכשחבר חמש יריעות יחד נמצא רחבן עשרים, וכן החוברת השנית, והמשכן ארכו שלשים מן המזרח למערב, שנ' עשרים קרשים לפאת נגב תימנה, וכן לצפון, וכל קרש אמה וחצי האמה, הרי שלשים מן המזרח למערב, רחב המשכן מן הצפון לדרום עשר אמות, שנאמר ולירכתי המשכן ימה וגו' ושני קרשים למקצעת הרי עשר – ובמקומם אפרשם למקראות הללו – נותן היריעות ארכן לרחבו של משכן, עשר אמות אמצעיות לגג חלל רחב המשכן, ואמה מכאן ואמה מכאן לעבי ראשי הקרשים שעבים אמה, נשתירו ט"ז אמה, ח' לצפון וח' לדרום, מכסות קומת הקרשים שגבהן עשר, נמצאו שתי אמות התחתונות מגלות; רחבן של יריעות ארבעים אמה כשהן מחברות, עשרים אמה לחוברת. שלשים מהן לגג חלל המשכן לארכו, ואמה כנגד עבי ראשי הקרשים שבמערב, ואמה לכסות עבי העמודים שבמזרח – שלא היו קרשים במזרח, אלא חמשה עמודים שהמסך פרוש ותלוי בווין שבהן כמין וילון – נשתירו שמונה אמות התלויין על אחורי המשכן שבמערב, ושתי אמות התחתונות מגלות, זו מצאתי בבריתא דארבעים ותשע מדות. אבל במסכת שבת אין היריעות מכסות את עמודי המזרח, ותשע אמות תלויות אחורי המשכן, והכתוב בפרשה זו מסיענו, ונתת את הפרכת "תחת" הקרסים, ואם כדברי הבריתא הזאת, נמצאת פרכת משוכה מן הקרסים ולמערב אמה:}%
% ---------- פרק 26 | פסוק 6 ----------
 \textbf{ו} וְעָשִׂ֕יתָ חֲמִשִּׁ֖ים קַרְסֵ֣י זָהָ֑ב וְחִבַּרְתָּ֨ אֶת־הַיְרִיעֹ֜ת אִשָּׁ֤ה אֶל־אֲחֹתָהּ֙ בַּקְּרָסִ֔ים וְהָיָ֥ה הַמִּשְׁכָּ֖ן אֶחָֽד׃ \Rashi{\textbf{קרסי זהב.}פירמי"לש בלעז, ומכניסין ראשן אחד בלולאות שבחוברת זו וראשן אחד בלולאות שבחוברת זו ומחברן בהן:}%
% ---------- פרק 26 | פסוק 7 ----------
 \textbf{ז} וְעָשִׂ֙יתָ֙ יְרִיעֹ֣ת עִזִּ֔ים לְאֹ֖הֶל עַל־הַמִּשְׁכָּ֑ן עַשְׁתֵּי־עֶשְׂרֵ֥ה יְרִיעֹ֖ת תַּעֲשֶׂ֥ה אֹתָֽם׃ \Rashi{\textbf{יריעת עזים.}מנוצה של עזים: לאהל על המשכן. לפרש אותן על היריעות התחתונות:}%
% ---------- פרק 26 | פסוק 8 ----------
 \textbf{ח} אֹ֣רֶךְ ׀ הַיְרִיעָ֣ה הָֽאַחַ֗ת שְׁלֹשִׁים֙ בָּֽאַמָּ֔ה וְרֹ֙חַב֙ אַרְבַּ֣ע בָּאַמָּ֔ה הַיְרִיעָ֖ה הָאֶחָ֑ת מִדָּ֣ה אַחַ֔ת לְעַשְׁתֵּ֥י עֶשְׂרֵ֖ה יְרִיעֹֽת׃ \Rashi{\textbf{שלשים באמה.}שכשנותן ארכן לרחב המשכן, כמו שנתן את הראשונות, נמצאו אלו עודפות אמה מכאן ואמה מכאן לכסות אחת מהשתי אמות שנשארו מגלות מן הקרשים, והאמה התחתונה של קרש שאין היריעה מכסה אותו היא האמה התחובה בנקב האדן, שהאדנים גבהן אמה:}%
% ---------- פרק 26 | פסוק 9 ----------
 \textbf{ט} וְחִבַּרְתָּ֞ אֶת־חֲמֵ֤שׁ הַיְרִיעֹת֙ לְבָ֔ד וְאֶת־שֵׁ֥שׁ הַיְרִיעֹ֖ת לְבָ֑ד וְכָפַלְתָּ֙ אֶת־הַיְרִיעָ֣ה הַשִּׁשִּׁ֔ית אֶל־מ֖וּל פְּנֵ֥י הָאֹֽהֶל׃ \Rashi{\textbf{וכפלת את היריעה הששית.}העודפת באלו העליונות יותר מן התחתונות: אל מול פני האהל. חצי רחבה היה תלוי וכפול על המסך שבמזרח כנגד הפתח, דומה לכלה צנועה המכסה בצעיף על פניה:}%
% ---------- פרק 26 | פסוק 10 ----------
 \textbf{י} וְעָשִׂ֜יתָ חֲמִשִּׁ֣ים לֻֽלָאֹ֗ת עַ֣ל שְׂפַ֤ת הַיְרִיעָה֙ הָֽאֶחָ֔ת הַקִּיצֹנָ֖ה בַּחֹבָ֑רֶת וַחֲמִשִּׁ֣ים לֻֽלָאֹ֗ת עַ֚ל שְׂפַ֣ת הַיְרִיעָ֔ה הַחֹבֶ֖רֶת הַשֵּׁנִֽית׃ %
% ---------- פרק 26 | פסוק 11 ----------
 \textbf{יא} וְעָשִׂ֛יתָ קַרְסֵ֥י נְחֹ֖שֶׁת חֲמִשִּׁ֑ים וְהֵבֵאתָ֤ אֶת־הַקְּרָסִים֙ בַּלֻּ֣לָאֹ֔ת וְחִבַּרְתָּ֥ אֶת־הָאֹ֖הֶל וְהָיָ֥ה אֶחָֽד׃ %
% ---------- פרק 26 | פסוק 12 ----------
 \textbf{יב} וְסֶ֙רַח֙ הָעֹדֵ֔ף בִּירִיעֹ֖ת הָאֹ֑הֶל חֲצִ֤י הַיְרִיעָה֙ הָעֹדֶ֔פֶת תִּסְרַ֕ח עַ֖ל אֲחֹרֵ֥י הַמִּשְׁכָּֽן׃ \Rashi{\textbf{וסרח העדף ביריעת האהל.}על יריעות המשכן; יריעות האהל הן העליונות של עזים, שקרויים אהל, כמו שנאמר בהם לאהל על המשכן, וכל אהל האמור בהן אינו אלא לשון גג, שמאהילות ומסככות על התחתונות, והן היו עודפות על התחתונות חצי היריעה למערב, שהחצי של יריעה אחת עשרה היתרה היה נכפל אל מול פני האהל, נשארו שתי אמות רחב חציה עודף על רחב התחתונות: תסרח על אחרי המשכן. לכסות שתי אמות שהיו מגלות בקרשים: אחרי המשכן. הוא צד מערבי, לפי שהפתח במזרח שהן פניו, וצפון ודרום קרויין צדדין לימין ולשמאל:}%
% ---------- פרק 26 | פסוק 13 ----------
 \textbf{יג} וְהָאַמָּ֨ה מִזֶּ֜ה וְהָאַמָּ֤ה מִזֶּה֙ בָּעֹדֵ֔ף בְּאֹ֖רֶךְ יְרִיעֹ֣ת הָאֹ֑הֶל יִהְיֶ֨ה סָר֜וּחַ עַל־צִדֵּ֧י הַמִּשְׁכָּ֛ן מִזֶּ֥ה וּמִזֶּ֖ה לְכַסֹּתֽוֹ׃ \Rashi{\textbf{והאמה מזה והאמה מזה.}לצפון ולדרום: בעדף בארך יריעת האהל. שהן עודפות על ארך יריעות המשכן שתי אמות: יהיה סרוח על צדי המשכן. לצפון ולדרום, כמו שפרשתי למעלה. למדה תורה דרך ארץ שיהא אדם חס על היפה (ילקוט שמעוני):}%
% ---------- פרק 26 | פסוק 14 ----------
 \textbf{יד} וְעָשִׂ֤יתָ מִכְסֶה֙ לָאֹ֔הֶל עֹרֹ֥ת אֵילִ֖ם מְאׇדָּמִ֑ים וּמִכְסֵ֛ה עֹרֹ֥ת תְּחָשִׁ֖ים מִלְמָֽעְלָה׃ \{פ\} \Rashi{\textbf{מכסה לאהל.}לאותו גג של יריעות עזים עשה עוד מכסה אחר של עורות אילים מאדמים, ועוד למעלה ממנו מכסה עורות תחשים; ואותן מכסאות לא היו מכסין אלא את הגג, ארכן ל' ורחבן י', אלו דברי רבי נחמיה, ולדברי רבי יהודה מכסה אחד היה, חציו של עורות אילים מאדמים וחציו של עורות תחשים (שבת כ"ח):}%
% ---------- פרק 26 | פסוק 15 ----------
 \textbf{טו} וְעָשִׂ֥יתָ אֶת־הַקְּרָשִׁ֖ים לַמִּשְׁכָּ֑ן עֲצֵ֥י שִׁטִּ֖ים עֹמְדִֽים׃ \Rashi{\textbf{ועשית את הקרשים.}היה לו לומר ועשית קרשים, כמו שנאמר בכל דבר ודבר, מהו הקרשים? מאותן העומדין ומיחדין לכך; יעקב אבינו נטע ארזים במצרים, וכשמת צוה לבניו להעלותם עמהם כשיצאו ממצרים, ואמר להם שעתיד הקב"ה לצוות אותן לעשות משכן במדבר מעצי שטים, ראו שיהיו מזמנים בידכם; הוא שיסד הבבלי בפיוט שלו, "טס מטע מזרזים קורות בתינו ארזים", שנזדרזו להיות מוכנים בידם מקדם לכן: עצי שטים עומדים. איתשט"נביש בלעז; שיהא ארך הקרשים זקוף למעלה בקירות המשכן, ולא תעשה הכתלים בקרשים שוכבים, להיות רחב הקרשים לגבה הכתלים קרש על קרש:}%
% ---------- פרק 26 | פסוק 16 ----------
 \textbf{טז} עֶ֥שֶׂר אַמּ֖וֹת אֹ֣רֶךְ הַקָּ֑רֶשׁ וְאַמָּה֙ וַחֲצִ֣י הָֽאַמָּ֔ה רֹ֖חַב הַקֶּ֥רֶשׁ הָאֶחָֽד׃ \Rashi{\textbf{עשר אמות ארך הקרש.}למדנו גבהו של משכן עשר אמות: ואמה וחצי האמה רחב. למדנו ארכו של משכן, לעשרים קרשים שיהיו בצפון ובדרום מן המזרח למערב, שלשים אמה:}%
% ---------- פרק 26 | פסוק 17 ----------
 \textbf{יז} שְׁתֵּ֣י יָד֗וֹת לַקֶּ֙רֶשׁ֙ הָאֶחָ֔ד מְשֻׁ֨לָּבֹ֔ת אִשָּׁ֖ה אֶל־אֲחֹתָ֑הּ כֵּ֣ן תַּעֲשֶׂ֔ה לְכֹ֖ל קַרְשֵׁ֥י הַמִּשְׁכָּֽן׃ \Rashi{\textbf{שתי ידות לקרש האחד.}היה חורץ את הקרש מלמטה באמצעו בגבה אמה, ומניח רביע רחבו מכאן, ורביע רחבו מכאן, והן הן הידות, והחריץ חצי רחב הקרש באמצע, ואותן הידות מכניס באדנים שהיו חלולים, והאדנים גבהן אמה ויושבים רצופים ארבעים זה אצל זה, וידות הקרש הנכנסות בחלל האדנים חרוצות משלשה צדיהן – רחב החריץ כעבי שפת האדן – שיכסה הקרש את כל ראש האדן, שאם לא כן, נמצא רוח בין קרש לקרש כעבי שפת שני האדנים, שיפסיקו ביניהם, וזהו שנאמר ויהיו תאמים מלמטה, שיחרץ את צדי הידות כדי שיתחברו הקרשים זה אצל זה: משלבת. עשויות כמין שליבות סלם, ומבדלות זו מזו, ומשפין ראשיהם לכנס בתוך חלל האדן כשליבה הנכנסת בנקב עמודי הסלם: אשה אל אחתה. מכונות זו כנגד זו, שיהיו חריציהם שוים זו כמדת זו, כדי שלא יהיו שתי ידות זו משוכה לצד פנים וזו משוכה לצד חוץ בעבי הקרש שהוא אמה, ותרגום של ידות צירין, לפי שדומות לצירי הדלת הנכנסים בחרי המפתן:}%
% ---------- פרק 26 | פסוק 18 ----------
 \textbf{יח} וְעָשִׂ֥יתָ אֶת־הַקְּרָשִׁ֖ים לַמִּשְׁכָּ֑ן עֶשְׂרִ֣ים קֶ֔רֶשׁ לִפְאַ֖ת נֶ֥גְבָּה תֵימָֽנָה׃ \Rashi{\textbf{לפאת נגבה תימנה.}אין פאה זו לשון מקצוע, אלא כל הרוח קרויה פאה, כתרגומו: לרוח עבר דרומא:}%
% ---------- פרק 26 | פסוק 19 ----------
 \textbf{יט} וְאַרְבָּעִים֙ אַדְנֵי־כֶ֔סֶף תַּעֲשֶׂ֕ה תַּ֖חַת עֶשְׂרִ֣ים הַקָּ֑רֶשׁ שְׁנֵ֨י אֲדָנִ֜ים תַּֽחַת־הַקֶּ֤רֶשׁ הָאֶחָד֙ לִשְׁתֵּ֣י יְדֹתָ֔יו וּשְׁנֵ֧י אֲדָנִ֛ים תַּֽחַת־הַקֶּ֥רֶשׁ הָאֶחָ֖ד לִשְׁתֵּ֥י יְדֹתָֽיו׃ %
% ---------- פרק 26 | פסוק 20 ----------
 \textbf{כ} וּלְצֶ֧לַע הַמִּשְׁכָּ֛ן הַשֵּׁנִ֖ית לִפְאַ֣ת צָפ֑וֹן עֶשְׂרִ֖ים קָֽרֶשׁ׃ %
% ---------- פרק 26 | פסוק 21 ----------
 \textbf{כא} וְאַרְבָּעִ֥ים אַדְנֵיהֶ֖ם כָּ֑סֶף שְׁנֵ֣י אֲדָנִ֗ים תַּ֚חַת הַקֶּ֣רֶשׁ הָֽאֶחָ֔ד וּשְׁנֵ֣י אֲדָנִ֔ים תַּ֖חַת הַקֶּ֥רֶשׁ הָאֶחָֽד׃ %
% ---------- פרק 26 | פסוק 22 ----------
 \textbf{כב} וּֽלְיַרְכְּתֵ֥י הַמִּשְׁכָּ֖ן יָ֑מָּה תַּעֲשֶׂ֖ה שִׁשָּׁ֥ה קְרָשִֽׁים׃ \Rashi{\textbf{ולירכתי.}לשון סוף, כתרגומו ולסיפי, ולפי שהפתח במזרח, קרוי מזרח פנים והמערב אחורים, וזהו סוף, שהפנים הוא הראש: תעשה ששה קרשים. הרי תשע אמות רחב:}%
% ---------- פרק 26 | פסוק 23 ----------
 \textbf{כג} וּשְׁנֵ֤י קְרָשִׁים֙ תַּעֲשֶׂ֔ה לִמְקֻצְעֹ֖ת הַמִּשְׁכָּ֑ן בַּיַּרְכָתָֽיִם׃ \Rashi{\textbf{ושני קרשים תעשה למקצעת.}אחד למקצוע צפונית מערבית, ואחד למערבית דרומית, כל שמנה קרשים בסדר אחד הן, אלא שאלו השנים אינן בחלל המשכן אלא חצי אמה מזו וחצי אמה מזו נראות בחלל, להשלים רחבו לעשר, והאמה מזה והאמה מזה באות כנגד אמת עבי קרשי המשכן הצפון והדרום, כדי שיהא המקצוע מבחוץ שוה:}%
% ---------- פרק 26 | פסוק 24 ----------
 \textbf{כד} וְיִֽהְי֣וּ תֹֽאֲמִם֮ מִלְּמַ֒טָּה֒ וְיַחְדָּ֗ו יִהְי֤וּ תַמִּים֙ עַל־רֹאשׁ֔וֹ אֶל־הַטַּבַּ֖עַת הָאֶחָ֑ת כֵּ֚ן יִהְיֶ֣ה לִשְׁנֵיהֶ֔ם לִשְׁנֵ֥י הַמִּקְצֹעֹ֖ת יִהְיֽוּ׃ \Rashi{\textbf{ויהיו.}כל הקרשים תואמים זה לזה מלמטה, שלא יפסיק עבי שפת שני האדנים ביניהם להרחיקם זו מזו, זהו שפרשתי שיהיו צירי הידות חרוצים מצדיהן, שיהא רחב הקרש בולט לצדיו חוץ לידי הקרש, לכסות את שפת האדן, וכן הקרש שאצלו, ונמצאו תואמים זה לזה; וקרש המקצוע שבסדר המערב חרוץ לרחבו בעביו כנגד חריץ של צד קרש הצפוני והדרומי, כדי שלא יפרידו האדנים ביניהם: ויחדו יהיו תמים. כמו תואמים: על ראשו. של קרש: אל הטבעת האחת. כל קרש וקרש היה חרוץ מלמעלה ברחבו שני חריצין, בשני צדיו, כמו עבי טבעת, ומכניסו בטבעת אחת, נמצא מתאים לקרש שאצלו, אבל אותן טבעות לא ידעתי אם קבועות הן, אם מטלטלות, ובקרש שבמקצוע היתה טבעת בעבי הקרש הדרומי והצפוני, וראש קרש המקצוע שבסדר מערב נכנס לתוכו, נמצאו שני הכתלים מחברים: כן יהיה לשניהם. לשני הקרשים שבמקצוע, לקרש שבסוף צפון ולקרש המערבי, וכן לשני המקצועות:}%
% ---------- פרק 26 | פסוק 25 ----------
 \textbf{כה} וְהָיוּ֙ שְׁמֹנָ֣ה קְרָשִׁ֔ים וְאַדְנֵיהֶ֣ם כֶּ֔סֶף שִׁשָּׁ֥ה עָשָׂ֖ר אֲדָנִ֑ים שְׁנֵ֣י אֲדָנִ֗ים תַּ֚חַת הַקֶּ֣רֶשׁ הָאֶחָ֔ד וּשְׁנֵ֣י אֲדָנִ֔ים תַּ֖חַת הַקֶּ֥רֶשׁ הָאֶחָֽד׃ \Rashi{\textbf{והיו שמנה קרשים.}הן האמורות למעלה תעשה ששה קרשים ושני קרשים תעשה למקצעת, נמצאו שמנה קרשים בסדר מערבי; כך שנויה במשנה מעשה סדר הקרשים במלאכת המשכן; היה עושה את האדנים חלולים וחורץ את הקרש מלמטה רביע מכאן ורביע מכאן, והחריץ חציו באמצע, ועשה לו שתי ידות כמין שני חמוקין – ולי נראה שהגרסא כמין שני חוקין, כמין שני שליבות סלם המבדלות זו מזו, ומשפות לכנס בחלל האדן כשליבה הנכנסת בנקב עמוד הסלם, והוא לשון משלבות, עשויות כמין שליבה – ומכניסן לתוך שני אדנים, שנאמר שני אדנים שני אדנים, וחורץ את הקרש מלמעלה אצבע מכאן ואצבע מכאן, ונותנן לתוך טבעת אחת של זהב כדי שלא יהיו נפרדין זה מזה, שנאמר ויהיו תאמים מלמטה וגו', כך היא המשנה, והפרוש שלה הצעתי למעלה בסדר המקראות:}%
% ---------- פרק 26 | פסוק 26 ----------
 \textbf{כו} וְעָשִׂ֥יתָ בְרִיחִ֖ם עֲצֵ֣י שִׁטִּ֑ים חֲמִשָּׁ֕ה לְקַרְשֵׁ֥י צֶֽלַע־הַמִּשְׁכָּ֖ן הָאֶחָֽד׃ \Rashi{\textbf{בריחם.}כתרגומו עברין, ובלעז אשפ"ריש: חמשה לקרשי צלע המשכן. אלו חמשה ג' הן, אלא שהבריח העליון והתחתון עשוי משתי חתיכות, זה מבריח עד חצי הכתל וזה מבריח עד חצי הכתל, זה נכנס בטבעות מצד זה, וזה נכנס בטבעות מצד זה, עד שמגיעין זה לזה, נמצא שעליון ותחתון שנים שהן ארבעה, אבל האמצעי ארכו כנגד כל הכתל, ומבריח מקצה הכתל ועד קצהו, שנאמר והבריח התיכן וגו' מברח מן הקצה אל הקצה, העליונים והתחתונים היו להן טבעות בקרשים לכנס בתוכן, שתי טבעות לכל קרש משלשים בתוך עשר אמות של גבה הקרש – חלק אחד מן הטבעת העליונה ולמעלה, וחלק אחד מן התחתונה ולמטה, וכל חלק הוא רביע ארך הקרש, ושני חלקים בין טבעת לטבעת כדי שיהיו כל הטבעות מכונים זו כנגד זו – אבל לבריח התיכון אין טבעות, אלא הקרשים נקובים בעבין, והוא נכנס בהם דרך הנקבים שהם מכונין זה מול זה, וזהו שנאמר "בתוך" הקרשים; הבריחים העליונים והתחתונים שבצפון ושבדרום ארך כל אחד חמש עשרה אמה, והתיכון ארכו שלשים אמה, וזהו מן הקצה אל הקצה – מן המזרח ועד המערב – וחמשה בריחים שבמערב ארך העליונים והתחתונים שש אמות, והתיכון ארכו שתים עשרה, כנגד רחב שמנה קרשים, כך היא מפרשת במלאכת המשכן:}%
% ---------- פרק 26 | פסוק 27 ----------
 \textbf{כז} וַחֲמִשָּׁ֣ה בְרִיחִ֔ם לְקַרְשֵׁ֥י צֶֽלַע־הַמִּשְׁכָּ֖ן הַשֵּׁנִ֑ית וַחֲמִשָּׁ֣ה בְרִיחִ֗ם לְקַרְשֵׁי֙ צֶ֣לַע הַמִּשְׁכָּ֔ן לַיַּרְכָתַ֖יִם יָֽמָּה׃ %
% ---------- פרק 26 | פסוק 28 ----------
 \textbf{כח} וְהַבְּרִ֥יחַ הַתִּיכֹ֖ן בְּת֣וֹךְ הַקְּרָשִׁ֑ים מַבְרִ֕חַ מִן־הַקָּצֶ֖ה אֶל־הַקָּצֶֽה׃ %
% ---------- פרק 26 | פסוק 29 ----------
 \textbf{כט} וְֽאֶת־הַקְּרָשִׁ֞ים תְּצַפֶּ֣ה זָהָ֗ב וְאֶת־טַבְּעֹֽתֵיהֶם֙ תַּעֲשֶׂ֣ה זָהָ֔ב בָּתִּ֖ים לַבְּרִיחִ֑ם וְצִפִּיתָ֥ אֶת־הַבְּרִיחִ֖ם זָהָֽב׃ \Rashi{\textbf{בתים לבריחם.}הטבעות שתעשה בהן יהיו בתים לכנס בהן הבריחים: וצפית את הבריחם זהב. לא שהיה הזהב מדבק על הבריחים – שאין עליהם שום צפוי – אלא בקרש היה קובע כמין שתי פיפיות של זהב, כמין שני סדקי קנה חלול, וקובען אצל הטבעת לכאן ולכאן, ארכן ממלא את רחב הקרש מן הטבעת לכאן וממנה לכאן, והבריח נכנס לתוכו וממנו לטבעת ומן הטבעת לפה השני, נמצאו הבריחים מצפים זהב כשהן תחובין בקרשים; והבריחים הללו מבחוץ היו בולטות, הטבעות והפיפיות לא היו נראות בתוך המשכן, אלא כל הכתל חלק מבפנים:}%
% ---------- פרק 26 | פסוק 30 ----------
 \textbf{ל} וַהֲקֵמֹתָ֖ אֶת־הַמִּשְׁכָּ֑ן כְּמִ֨שְׁפָּט֔וֹ אֲשֶׁ֥ר הׇרְאֵ֖יתָ בָּהָֽר׃ \{ס\} \Rashi{\textbf{והקמת את המשכן.}לאחר שיגמר הקימהו: הראית בהר. קדם לכן, שאני עתיד ללמדך ולהראותך סדר הקמתו:}%
% ---------- פרק 26 | פסוק 31 ----------
 \textbf{לא} וְעָשִׂ֣יתָ פָרֹ֗כֶת תְּכֵ֧לֶת וְאַרְגָּמָ֛ן וְתוֹלַ֥עַת שָׁנִ֖י וְשֵׁ֣שׁ מׇשְׁזָ֑ר מַעֲשֵׂ֥ה חֹשֵׁ֛ב יַעֲשֶׂ֥ה אֹתָ֖הּ כְּרֻבִֽים׃ \Rashi{\textbf{פרכת.}לשון מחצה הוא, ובלשון חכמים פרגוד דבר המבדיל בין המלך ובין העם: תכלת וארגמן. כל מין ומין היה כפול בכל חוט וחוט ששה חוטין: מעשה חשב. כבר פרשתי שזו היא אריגה של שתי קירות, והציורין שמשני עבריה אינן דומין זה לזה: כרבים. ציורין של בריות יעשה בה:}%
% ---------- פרק 26 | פסוק 32 ----------
 \textbf{לב} וְנָתַתָּ֣ה אֹתָ֗הּ עַל־אַרְבָּעָה֙ עַמּוּדֵ֣י שִׁטִּ֔ים מְצֻפִּ֣ים זָהָ֔ב וָוֵיהֶ֖ם זָהָ֑ב עַל־אַרְבָּעָ֖ה אַדְנֵי־כָֽסֶף׃ \Rashi{\textbf{ארבעה עמודי שטים.}תקועים בתוך ארבעה אדנים, ואונקליות קבועין בהן, עקמים למעלה, להושיב עליהן כלונס שראש הפרכת כרוך בה, והאונקליות הן הווין – שהרי כמין ווין הן עשויים – והפרכת ארכה עשר אמות לרחבו של משכן ורחבה עשר אמות כגבהן של קרשים, פרוסה בשלישית של משכן, שיהא הימנה ולפנים עשר אמות והימנה ולחוץ עשרים אמה, נמצא בית קדשי הקדשים עשר על עשר, שנ' ונתתה את הפרכת תחת הקרסים – המחברים את שתי חוברות של יריעות המשכן; רחב החוברת עשרים אמה, וכשפרשה על גג המשכן מן הפתח למערב, כלתה בשני שלישי המשכן, והחוברת השנית כסתה שלישו של משכן, והמותר תלוי לאחוריו לכסות את הקרשים:}%
% ---------- פרק 26 | פסוק 33 ----------
 \textbf{לג} וְנָתַתָּ֣ה אֶת־הַפָּרֹ֘כֶת֮ תַּ֣חַת הַקְּרָסִים֒ וְהֵבֵאתָ֥ שָׁ֙מָּה֙ מִבֵּ֣ית לַפָּרֹ֔כֶת אֵ֖ת אֲר֣וֹן הָעֵד֑וּת וְהִבְדִּילָ֤ה הַפָּרֹ֙כֶת֙ לָכֶ֔ם בֵּ֣ין הַקֹּ֔דֶשׁ וּבֵ֖ין קֹ֥דֶשׁ הַקֳּדָשִֽׁים׃ %
% ---------- פרק 26 | פסוק 34 ----------
 \textbf{לד} וְנָתַתָּ֙ אֶת־הַכַּפֹּ֔רֶת עַ֖ל אֲר֣וֹן הָעֵדֻ֑ת בְּקֹ֖דֶשׁ הַקֳּדָשִֽׁים׃ %
% ---------- פרק 26 | פסוק 35 ----------
 \textbf{לה} וְשַׂמְתָּ֤ אֶת־הַשֻּׁלְחָן֙ מִח֣וּץ לַפָּרֹ֔כֶת וְאֶת־הַמְּנֹרָה֙ נֹ֣כַח הַשֻּׁלְחָ֔ן עַ֛ל צֶ֥לַע הַמִּשְׁכָּ֖ן תֵּימָ֑נָה וְהַ֨שֻּׁלְחָ֔ן תִּתֵּ֖ן עַל־צֶ֥לַע צָפֽוֹן׃ \Rashi{\textbf{ושמת את השלחן.}שלחן בצפון, משוך מן הכתל הצפוני שתי אמות ומחצה, ומנורה בדרום, משוכה מן הכתל הדרומי שתי אמות ומחצה, ומזבח הזהב נתון כנגד אויר שבין שלחן למנורה משוך קמעא כלפי המזרח, וכלם נתונים מן חצי המשכן ולפנים; כיצד? ארך המשכן מן הפתח לפרכת עשרים אמה, המזבח והשלחן והמנורה משוכים מן הפתח לצד מערב עשר אמות (יומא ל"ג):}%
% ---------- פרק 26 | פסוק 36 ----------
 \textbf{לו} וְעָשִׂ֤יתָ מָסָךְ֙ לְפֶ֣תַח הָאֹ֔הֶל תְּכֵ֧לֶת וְאַרְגָּמָ֛ן וְתוֹלַ֥עַת שָׁנִ֖י וְשֵׁ֣שׁ מׇשְׁזָ֑ר מַעֲשֵׂ֖ה רֹקֵֽם׃ \Rashi{\textbf{ועשית מסך.}וילון, הוא מסך כנגד הפתח, כמו שכת בעדו (איוב א'), לשון מגן: מעשה רקם. הצורות עשויות בו מעשה מחט, כפרצוף של עבר זה כך פרצוף של עבר זה: רקם. שם האמן ולא שם האמנות, ותרגומו עובד ציר ולא עובד ציור. מדת המסך כמדת הפרכת – עשר אמות על עשר אמות:}%
% ---------- פרק 26 | פסוק 37 ----------
 \textbf{לז} וְעָשִׂ֣יתָ לַמָּסָ֗ךְ חֲמִשָּׁה֙ עַמּוּדֵ֣י שִׁטִּ֔ים וְצִפִּיתָ֤ אֹתָם֙ זָהָ֔ב וָוֵיהֶ֖ם זָהָ֑ב וְיָצַקְתָּ֣ לָהֶ֔ם חֲמִשָּׁ֖ה אַדְנֵ֥י נְחֹֽשֶׁת׃ \{ס\} %
% ---------- פרק 27 | פסוק 1 ----------
 \textbf{א} וְעָשִׂ֥יתָ אֶת־הַמִּזְבֵּ֖חַ עֲצֵ֣י שִׁטִּ֑ים חָמֵשׁ֩ אַמּ֨וֹת אֹ֜רֶךְ וְחָמֵ֧שׁ אַמּ֣וֹת רֹ֗חַב רָב֤וּעַ יִהְיֶה֙ הַמִּזְבֵּ֔חַ וְשָׁלֹ֥שׁ אַמּ֖וֹת קֹמָתֽוֹ׃ \Rashi{\textbf{ועשית את המזבח וגו' ושלש אמות קמתו.}דברים ככתבן, דברי ר' יהודה, רבי יוסי אומר נאמר כאן רבוע ונאמר בפנימי רבוע, מה להלן גבהו פי שנים כארכו, אף כאן גבהו פי שנים כארכו, ומה אני מקים ושלש אמות קמתו? משפת סובב ולמעלה (זבח' נ"ט):}%
% ---------- פרק 27 | פסוק 2 ----------
 \textbf{ב} וְעָשִׂ֣יתָ קַרְנֹתָ֗יו עַ֚ל אַרְבַּ֣ע פִּנֹּתָ֔יו מִמֶּ֖נּוּ תִּהְיֶ֣יןָ קַרְנֹתָ֑יו וְצִפִּיתָ֥ אֹת֖וֹ נְחֹֽשֶׁת׃ \Rashi{\textbf{ממנו תהיין קרנתיו.}שלא יעשם לבדם ויחברם בו: וצפית אתו נחשת. לכפר על עזות מצח, שנ' ומצחך נחושה (ישעיהו מ"ח):}%
% ---------- פרק 27 | פסוק 3 ----------
 \textbf{ג} וְעָשִׂ֤יתָ סִּֽירֹתָיו֙ לְדַשְּׁנ֔וֹ וְיָעָיו֙ וּמִזְרְקֹתָ֔יו וּמִזְלְגֹתָ֖יו וּמַחְתֹּתָ֑יו לְכׇל־כֵּלָ֖יו תַּעֲשֶׂ֥ה נְחֹֽשֶׁת׃ \Rashi{\textbf{סירתיו.}כמין יורות: לדשנו. להסיר דשנו לתוכם, והוא שתרגם אנקלוס למספי קטמה – לספות הדשן לתוכם; כי יש מלות בלשון עברית מלה אחת מתחלפת בפתרון לשמש בנין וסתירה, כמו ותשרש שרשיה (תהילים פ'), אויל משריש (איוב ה'), וחלופו ובכל תבואתי תשרש (איוב ל״א:י״ב); וכמוהו בסעפיה פריה (ישעיהו י״ז:ו׳), וחלופו מסעף פארה (ישעיהו י׳:ל״ג) – מפשח סעפיה; וכמוהו וזה האחרון עצמו (ירמיהו נ') – שבר עצמיו; וכמוהו ויסקלהו באבנים (מלכים א כ״א:י״ג), וחלופו סקלו מאבן (ישעיהו ס״ב:י׳) – הסירו אבניה, וכן ויעזקהו ויסקלהו (שם ה'), אף כאן לדשנו – להסיר דשנו, ובלעז אדשצנדר"יר: ויעיו. כתרגומו, מגרפות שנוטלים בהם הדשן, והן כמין כסוי הקדרה של מתכת דק ולו בית יד, ובלעז ווד"יל: ומזרקתיו. לקבל בהם דם הזבחים: ומזלגתיו. כמין אנקליות כפופים, ומכה בהם בבשר ונתחבים בו ומהפך בהן על גחלי המערכה שיהא ממהר שרפתן, ובלעז קרוצי"נש, ובלשון חכמים צנוריות (יומא י"ב): ומחתתיו. בית קבול יש להם לטל בהן גחלים מן המזבח, לשאתם על מזבח הפנימי לקטרת, ועל שם חתיתן קרויים מחתות, כמו לחתות אש מיקוד (ישעיהו ל'), לשון שאיבת אש ממקומה, וכן היחתה איש אש בחיקו (משלי ו'): לכל כליו. כמו כל כליו:}%
% ---------- פרק 27 | פסוק 4 ----------
 \textbf{ד} וְעָשִׂ֤יתָ לּוֹ֙ מִכְבָּ֔ר מַעֲשֵׂ֖ה רֶ֣שֶׁת נְחֹ֑שֶׁת וְעָשִׂ֣יתָ עַל־הָרֶ֗שֶׁת אַרְבַּע֙ טַבְּעֹ֣ת נְחֹ֔שֶׁת עַ֖ל אַרְבַּ֥ע קְצוֹתָֽיו׃ \Rashi{\textbf{מכבר.}לשון כברה, שקורין קרי"בלא בלעז, כמין לבוש עשוי לו למזבח, עשוי חרין חרין, כמין רשת; ומקרא זה מסרס וכה פתרונו: ועשית לו מכבר נחשת מעשה רשת:}%
% ---------- פרק 27 | פסוק 5 ----------
 \textbf{ה} וְנָתַתָּ֣ה אֹתָ֗הּ תַּ֛חַת כַּרְכֹּ֥ב הַמִּזְבֵּ֖חַ מִלְּמָ֑טָּה וְהָיְתָ֣ה הָרֶ֔שֶׁת עַ֖ד חֲצִ֥י הַמִּזְבֵּֽחַ׃ \Rashi{\textbf{כרכב המזבח.}סובב; כל דבר המקיף סביב בעגול קרוי כרכב, כמו ששנינו בהכל שוחטין אלו הן גלמי כלי עץ כל שעתיד לשוף ולכרכב, והוא כמו שעושים חריצין עגלין בקרשי דפני התבות וספסלי העץ, אף למזבח עשה חריץ סביבו, והיה רחבו אמה בדפנו לנוי והוא לסוף שלש אמות של גבהו כדברי האומר גבהו פי שנים כארכו, הא מה אני מקים ושלש אמות קומתו? משפת סובב ולמעלה, אבל סובב להלוך הכהנים לא היה למזבח הנחשת, אלא על ראשו לפנים מקרנותיו, וכן שנינו בזבחים: אי זהו כרכוב? בין קרן לקרן; והיה רחב אמה, ולפנים מהן אמה של הלוך רגלי הכהנים, שתי אמות הללו קרויים כרכוב; ודקדקנו שם והכתיב תחת כרכב המזבח מלמטה? למדנו שהכרכוב בדפנו הוא ולבוש המכבר תחתיו? ותרץ המתרץ תרי הוו, חד לנוי וחד לכהנים דלא ישתרגו, זה שבדפן לנוי היה, ומתחתיו הלבישו המכבר והגיע רחבו עד חצי המזבח, והוא היה סימן לחצי גבהו, להבדיל בין דמים העליונים לדמים התחתונים; וכנגדו עשו למזבח בית עולמים דגמת חוט הסקרא באמצעו, וכבש שהיו עולין בו אע"פ שלא פרשו בענין זה, כבר שמענו בפרשת מזבח אדמה תעשה לי ולא תעלה במעלת – לא תעשה לו מעלות בכבש שלו, אלא כבש חלק, למדנו שהיה לו כבש; כך שנינו במכילתא. ומזבח אדמה הוא מזבח הנחשת, שהיו ממלאין חללו אדמה במקום חניתן, והכבש היה בדרום המזבח, מבדל מן המזבח מלא חוט השערה, ורגליו מגיעין עד אמה סמוך לקלעי החצר שבדרום, כדברי האומר עשר אמות קומתו, ולדברי האומר דברים ככתבן – שלש אמות קומתו – לא היה ארך הכבש אלא עשר אמות, כך מצאתי במשנת ארבעים ותשע מדות, וזה שהוא מבדל מן המזבח מלא החוט, במסכת זבחים (דף ס"ב) למדוהו מן המקרא:}%
% ---------- פרק 27 | פסוק 6 ----------
 \textbf{ו} וְעָשִׂ֤יתָ בַדִּים֙ לַמִּזְבֵּ֔חַ בַּדֵּ֖י עֲצֵ֣י שִׁטִּ֑ים וְצִפִּיתָ֥ אֹתָ֖ם נְחֹֽשֶׁת׃ %
% ---------- פרק 27 | פסוק 7 ----------
 \textbf{ז} וְהוּבָ֥א אֶת־בַּדָּ֖יו בַּטַּבָּעֹ֑ת וְהָי֣וּ הַבַּדִּ֗ים עַל־שְׁתֵּ֛י צַלְעֹ֥ת הַמִּזְבֵּ֖חַ בִּשְׂאֵ֥ת אֹתֽוֹ׃ \Rashi{\textbf{בטבעת.}בארבע טבעות שנעשו למכבר:}%
% ---------- פרק 27 | פסוק 8 ----------
 \textbf{ח} נְב֥וּב לֻחֹ֖ת תַּעֲשֶׂ֣ה אֹת֑וֹ כַּאֲשֶׁ֨ר הֶרְאָ֥ה אֹתְךָ֛ בָּהָ֖ר כֵּ֥ן יַעֲשֽׂוּ׃ \{ס\} \Rashi{\textbf{נבוב לחת.}כתרגומו חליל לוחין, לחת עצי שטים מכל צד והחלל באמצע, ולא יהא כלו עץ אחד שיהא עביו חמש אמות על חמש אמות כמין סדן:}%
% ---------- פרק 27 | פסוק 9 ----------
 \textbf{ט} וְעָשִׂ֕יתָ אֵ֖ת חֲצַ֣ר הַמִּשְׁכָּ֑ן לִפְאַ֣ת נֶֽגֶב־תֵּ֠ימָ֠נָה קְלָעִ֨ים לֶחָצֵ֜ר שֵׁ֣שׁ מׇשְׁזָ֗ר מֵאָ֤ה בָֽאַמָּה֙ אֹ֔רֶךְ לַפֵּאָ֖ה הָאֶחָֽת׃ \Rashi{\textbf{קלעים.}עשויין כמין קלעי ספינה, נקבים נקבים, מעשה קליעה ולא מעשה אורג, ותרגום סרדים כתרגומו של מכבר המתרגם סרדה, לפי שהן מנקבין ככברה: לפאה האחת. כל הרוח קרוי פאה:}%
% ---------- פרק 27 | פסוק 10 ----------
 \textbf{י} וְעַמֻּדָ֣יו עֶשְׂרִ֔ים וְאַדְנֵיהֶ֥ם עֶשְׂרִ֖ים נְחֹ֑שֶׁת וָוֵ֧י הָעַמֻּדִ֛ים וַחֲשֻׁקֵיהֶ֖ם כָּֽסֶף׃ \Rashi{\textbf{ועמדיו עשרים.}חמש אמות בין עמוד לעמוד: ואדניהם. של העמודים נחשת, האדנים יושבים על הארץ והעמודים תקועין לתוכן; היה עושה כמין קנדסין שקורין פל"ש בלעז, ארכן ששה טפחים ורחבן שלשה, וטבעת נחשת קבועה בו באמצעו וכורך שפת הקלע סביביו במיתרים כנגד כל עמוד ועמוד, ותולה הקנדס דרך טבעתו באנקליות שבעמוד העשוי כמין וי"ו – ראשו זקוף למעלה וראשו אחד תקוע בעמוד – כאותן שעושין להציב דלתות, שקורין גונ"ש בלעז, ורחב הקלע תלוי מלמטה והיא קומת מחצות החצר: ווי העמדים. הם האנקליות: וחשקיהם. מקפין היו העמודים בחוטי כסף סביב, ואיני יודע אם על פני כלן, אם בראשם, ואם באמצעם, אך יודע אני שחשוק לשון חגורה, שכן מצינו בפילגש בגבעה ועמו צמד חמורים חבושים (שופטים י"ט), תרגומו חשוקים:}%
% ---------- פרק 27 | פסוק 11 ----------
 \textbf{יא} וְכֵ֨ן לִפְאַ֤ת צָפוֹן֙ בָּאֹ֔רֶךְ קְלָעִ֖ים מֵ֣אָה אֹ֑רֶךְ וְעַמֻּדָ֣ו עֶשְׂרִ֗ים וְאַדְנֵיהֶ֤ם עֶשְׂרִים֙ נְחֹ֔שֶׁת וָוֵ֧י הָֽעַמֻּדִ֛ים וַחֲשֻׁקֵיהֶ֖ם כָּֽסֶף׃ %
% ---------- פרק 27 | פסוק 12 ----------
 \textbf{יב} וְרֹ֤חַב הֶֽחָצֵר֙ לִפְאַת־יָ֔ם קְלָעִ֖ים חֲמִשִּׁ֣ים אַמָּ֑ה עַמֻּדֵיהֶ֣ם עֲשָׂרָ֔ה וְאַדְנֵיהֶ֖ם עֲשָׂרָֽה׃ %
% ---------- פרק 27 | פסוק 13 ----------
 \textbf{יג} וְרֹ֣חַב הֶֽחָצֵ֗ר לִפְאַ֛ת קֵ֥דְמָה מִזְרָ֖חָה חֲמִשִּׁ֥ים אַמָּֽה׃ \Rashi{\textbf{לפאת קדמה מזרחה.}פני המזרח; קדם לשון פנים, אחור לשון אחורים, לפיכך המזרח קרוי קדם שהוא פנים, ומערב קרוי אחור, כמה דאת אמר הים האחרון (דברים י"א:כ"ד, דברים ל"ד:ב'), ימא מערבא: חמשים אמה. אותן נ' אמה לא היו סתומים כלם בקלעים, לפי ששם הפתח, אלא ט"ו אמה קלעים לכתף הפתח מכאן, וכן לכתף השנית, נשאר רחב חלל הפתח בינתים כ' אמה, וזהו שנ' ולשער החצר מסך כ' אמה – וילון למסך כנגד הפתח כ' אמה ארך, כרחב הפתח:}%
% ---------- פרק 27 | פסוק 14 ----------
 \textbf{יד} וַחֲמֵ֨שׁ עֶשְׂרֵ֥ה אַמָּ֛ה קְלָעִ֖ים לַכָּתֵ֑ף עַמֻּדֵיהֶ֣ם שְׁלֹשָׁ֔ה וְאַדְנֵיהֶ֖ם שְׁלֹשָֽׁה׃ \Rashi{\textbf{עמדיהם שלשה.}חמש אמות בין עמוד לעמוד. בין עמוד שבראש הדרום העומד במקצוע דרומית מזרחית עד עמוד שהוא מן השלשה שבמזרח חמש אמות, וממנו לשני חמש אמות, ומן השני לשלישי חמש אמות, וכן לכתף השנית, וארבעה עמודים למסך, הרי עשרה עמודים למזרח כנגד י' למערב:}%
% ---------- פרק 27 | פסוק 15 ----------
 \textbf{טו} וְלַכָּתֵף֙ הַשֵּׁנִ֔ית חֲמֵ֥שׁ עֶשְׂרֵ֖ה קְלָעִ֑ים עַמֻּדֵיהֶ֣ם שְׁלֹשָׁ֔ה וְאַדְנֵיהֶ֖ם שְׁלֹשָֽׁה׃ %
% ---------- פרק 27 | פסוק 16 ----------
 \textbf{טז} וּלְשַׁ֨עַר הֶֽחָצֵ֜ר מָסָ֣ךְ ׀ עֶשְׂרִ֣ים אַמָּ֗ה תְּכֵ֨לֶת וְאַרְגָּמָ֜ן וְתוֹלַ֧עַת שָׁנִ֛י וְשֵׁ֥שׁ מׇשְׁזָ֖ר מַעֲשֵׂ֣ה רֹקֵ֑ם עַמֻּֽדֵיהֶם֙ אַרְבָּעָ֔ה וְאַדְנֵיהֶ֖ם אַרְבָּעָֽה׃ %
% ---------- פרק 27 | פסוק 17 ----------
 \textbf{יז} כׇּל־עַמּוּדֵ֨י הֶֽחָצֵ֤ר סָבִיב֙ מְחֻשָּׁקִ֣ים כֶּ֔סֶף וָוֵיהֶ֖ם כָּ֑סֶף וְאַדְנֵיהֶ֖ם נְחֹֽשֶׁת׃ \Rashi{\textbf{כל עמודי החצר סביב וגו'.}לפי שלא פרש ווין וחשוקים ואדני נחשת אלא לצפון ולדרום, אבל למזרח ולמערב לא נאמר ווין וחשוקים ואדני נחשת, לכך בא ולמד כאן:}%
% ---------- פרק 27 | פסוק 18 ----------
 \textbf{יח} אֹ֣רֶךְ הֶֽחָצֵר֩ מֵאָ֨ה בָֽאַמָּ֜ה וְרֹ֣חַב ׀ חֲמִשִּׁ֣ים בַּחֲמִשִּׁ֗ים וְקֹמָ֛ה חָמֵ֥שׁ אַמּ֖וֹת שֵׁ֣שׁ מׇשְׁזָ֑ר וְאַדְנֵיהֶ֖ם נְחֹֽשֶׁת׃ \Rashi{\textbf{ארך החצר.}הצפון והדרום שמן המזרח למערב מאה באמה: ורחב חמשים בחמשים. חצר שבמזרח היתה מרבעת, חמשים על חמשים, שהמשכן ארכו שלשים ורחבו עשר, העמיד מזרח פתחו בשפת נ' החיצונים של ארך החצר, נמצא כלו בחמשים הפנימיים וכלה ארכו לסוף שלשים, נמצאו עשרים אמה רוח לאחוריו בין הקלעים שבמערב ליריעות של אחורי המשכן, ורחב המשכן עשר אמות, באמצע רחב החצר, נמצאו לו עשרים אמה רוח לצפון ולדרום מן קלעי החצר ליריעות המשכן, וכן למערב, וחמשים על חמשים חצר לפניו: וקמה חמש אמות. גבה מחצות החצר והוא רחב הקלעים: ואדניהם נחשת. להביא אדני המסך, שלא תאמר לא נאמרו אדני נחשת אלא לעמודי הקלעים אבל אדני המסך של מין אחר, כן נראה בעיני שלכך חזר ושנאן:}%
% ---------- פרק 27 | פסוק 19 ----------
 \textbf{יט} לְכֹל֙ כְּלֵ֣י הַמִּשְׁכָּ֔ן בְּכֹ֖ל עֲבֹדָת֑וֹ וְכׇל־יְתֵדֹתָ֛יו וְכׇל־יִתְדֹ֥ת הֶחָצֵ֖ר נְחֹֽשֶׁת׃ \{ס\} \Rashi{\textbf{לכל כלי המשכן.}שהיו צריכין להקמתו ולהורדתו, כגון מקבות לתקע יתדות ועמודים: יתדת. כמין נגרי נחשת עשויין ליריעות האהל ולקלעי החצר, קשורים במיתרים סביב סביב בשפוליהן, כדי שלא תהא הרוח מגביהתן, ואיני יודע אם תחובין בארץ או קשורין ותלויין וכבדן מכביד שפולי היריעות שלא ינועו ברוח, ואומר אני ששמן מוכיח עליהם שהם תקועים בארץ, לכך נקראו יתדות, ומקרא זה מסיעני אהל בל יצען בל יסע יתדתיו לנצח (ישעיהו ל"ג):}%
\addparasha{חומש שמות | פרשת תצוה}
% ---------- פרק 27 | פסוק 20 ----------
 \textbf{כ} וְאַתָּ֞ה תְּצַוֶּ֣ה ׀ אֶת־בְּנֵ֣י יִשְׂרָאֵ֗ל וְיִקְח֨וּ אֵלֶ֜יךָ שֶׁ֣מֶן זַ֥יִת זָ֛ךְ כָּתִ֖ית לַמָּא֑וֹר לְהַעֲלֹ֥ת נֵ֖ר תָּמִֽיד׃ \Rashi{\textbf{ואתה תצוה.}זך. בלי שמרים, כמו ששנינו במנחות מגרגרו בראש הזית וכו': כתית. הזיתים כותש במכתשת ואינו טוחנן ברחים, כדי שלא יהיה בו שמרים, ואחר שהוציא טפה ראשונה מכניסן לרחים וטוחנן; והשמן השני פסול למנורה וכשר למנחות שנאמר כתית למאור – ולא כתית למנחות (מנחות פ"ו): להעלות נר תמיד. מדליק עד שתהא שלהבת עולה מאליה (שבת כ"א): תמיד. כל לילה ולילה קרוי תמיד, כמו שאתה אומר עלת תמיד ואינה אלא מיום ליום; וכן במנחת חבתין נאמר תמיד ואינה אלא מחציתה בבקר ומחציתה בערב, אבל תמיד האמור בלחם הפנים משבת לשבת הוא:}%
% ---------- פרק 27 | פסוק 21 ----------
 \textbf{כא} בְּאֹ֣הֶל מוֹעֵד֩ מִח֨וּץ לַפָּרֹ֜כֶת אֲשֶׁ֣ר עַל־הָעֵדֻ֗ת יַעֲרֹךְ֩ אֹת֨וֹ אַהֲרֹ֧ן וּבָנָ֛יו מֵעֶ֥רֶב עַד־בֹּ֖קֶר לִפְנֵ֣י יְהֹוָ֑ה חֻקַּ֤ת עוֹלָם֙ לְדֹ֣רֹתָ֔ם מֵאֵ֖ת בְּנֵ֥י יִשְׂרָאֵֽל׃ \{ס\} \Rashi{\textbf{מערב עד בקר.}תן לה מדתה שתהא דולקת מערב ועד בקר, ושערו חכמים חצי לג ללילי טבת הארכין וכן לכל הלילות, ואם יותר אין בכך כלום (מנחות פ"ט):}%
% ---------- פרק 28 | פסוק 1 ----------
 \textbf{א} וְאַתָּ֡ה הַקְרֵ֣ב אֵלֶ֩יךָ֩ אֶת־אַהֲרֹ֨ן אָחִ֜יךָ וְאֶת־בָּנָ֣יו אִתּ֗וֹ מִתּ֛וֹךְ בְּנֵ֥י יִשְׂרָאֵ֖ל לְכַהֲנוֹ־לִ֑י אַהֲרֹ֕ן נָדָ֧ב וַאֲבִיה֛וּא אֶלְעָזָ֥ר וְאִיתָמָ֖ר בְּנֵ֥י אַהֲרֹֽן׃ \Rashi{\textbf{ואתה הקרב אליך.}לאחר שתגמר מלאכת המשכן:}%
% ---------- פרק 28 | פסוק 2 ----------
 \textbf{ב} וְעָשִׂ֥יתָ בִגְדֵי־קֹ֖דֶשׁ לְאַהֲרֹ֣ן אָחִ֑יךָ לְכָב֖וֹד וּלְתִפְאָֽרֶת׃ %
% ---------- פרק 28 | פסוק 3 ----------
 \textbf{ג} וְאַתָּ֗ה תְּדַבֵּר֙ אֶל־כׇּל־חַכְמֵי־לֵ֔ב אֲשֶׁ֥ר מִלֵּאתִ֖יו ר֣וּחַ חׇכְמָ֑ה וְעָשׂ֞וּ אֶת־בִּגְדֵ֧י אַהֲרֹ֛ן לְקַדְּשׁ֖וֹ לְכַהֲנוֹ־לִֽי׃ \Rashi{\textbf{לקדשו לכהנו לי.}לקדשו להכניסו בכהנה ע"י הבגדים שיהא כהן לי, ולשון כהנה שרות הוא, שנטרי"אה בלעז:}%
% ---------- פרק 28 | פסוק 4 ----------
 \textbf{ד} וְאֵ֨לֶּה הַבְּגָדִ֜ים אֲשֶׁ֣ר יַעֲשׂ֗וּ חֹ֤שֶׁן וְאֵפוֹד֙ וּמְעִ֔יל וּכְתֹ֥נֶת תַּשְׁבֵּ֖ץ מִצְנֶ֣פֶת וְאַבְנֵ֑ט וְעָשׂ֨וּ בִגְדֵי־קֹ֜דֶשׁ לְאַהֲרֹ֥ן אָחִ֛יךָ וּלְבָנָ֖יו לְכַהֲנוֹ־לִֽי׃ \Rashi{\textbf{חשן.}תכשיט כנגד הלב: ואפוד. לא שמעתי ולא מצאתי בבריתא פרוש תבניתו, ולבי אומר לי שהוא חגורה לו מאחוריו, רחבו כרחב גב איש, כמין סינר שקורין פורצי"נט בלעז, שחוגרות השרות כשרוכבות על הסוסים, כך מעשהו מלמטה, שנאמר ודוד חגור אפוד בד, למדנו שהאפוד חגורה היא; ואי אפשר לומר שאין בו אלא חגורה לבדה, שהרי נאמר ויתן עליו את האפד ואח"כ ויחגר אותו בחשב האפד ותרגם אנקלוס בהמין אפודא, למדנו שהחשב הוא החגור והאפוד שם תכשיט לבדו; ואי אפשר לומר שעל שם שתי הכתפות שבו הוא קרוי אפוד, שהרי נאמר שתי כתפות האפוד, למדנו שהאפוד שם לבד והכתפות שם לבד והחשב שם לבד, לכך אני אומר שעל שם הסינר של מטה קרוי אפוד – על שם שאופדו ומקשטו בו – כמו שנ' ויאפד לו בו, והחשב הוא חגור שלמעלה הימנו והכתפות קבועות בו. ועוד אומר לי לבי שיש ראיה שהוא מין לבוש, שתרגם יונתן ודוד חגור אפוד בד (שמואל ב ו') – כרדוט דבוץ, ותרגם כמו כן מעילים כרדוטין, במעשה תמר אחות אבשלום, כי כן תלבשנה בנות המלך הבתולות מעילים (שם י"ג): מעיל. הוא כמין חלוק, וכן הכתנת, אלא שהכתנת סמוך לבשרו ומעיל קרוי חלוק העליון: תשבץ. עשויין משבצות לנוי, והמשבצות הם כמין גמות העשויות בתכשיטי זהב למושב קביעת אבנים טובות ומרגליות, כמו שנאמר באבני האפוד מסבת משבצות זהב, ובלעז קוראין אותו קשטו"נש: מצנפת. כמין כפת כובע, שקורין קופ"יא בלעז, שהרי במקום אחר קורא להם מגבעות ומתרגמינן כובעין: ואבנט. היא חגורה על הכתנת והאפוד חגורה על המעיל, כמו שמצינו בסדר לבישתן ויתן עליו את הכתנת ויחגר אתו באבנט וילבש אתו את המעיל ויתן עליו את האפד: בגדי קדש. מתרומה המקדשת לשמי יעשה אותם:}%
% ---------- פרק 28 | פסוק 5 ----------
 \textbf{ה} וְהֵם֙ יִקְח֣וּ אֶת־הַזָּהָ֔ב וְאֶת־הַתְּכֵ֖לֶת וְאֶת־הָֽאַרְגָּמָ֑ן וְאֶת־תּוֹלַ֥עַת הַשָּׁנִ֖י וְאֶת־הַשֵּֽׁשׁ׃ \{פ\} \Rashi{\textbf{והם יקחו.}אותם חכמי לב שיעשו הבגדים יקבלו מן המתנדבים את הזהב ואת התכלת לעשות מהן את הבגדים:}%
% ---------- פרק 28 | פסוק 6 ----------
 \textbf{ו} וְעָשׂ֖וּ אֶת־הָאֵפֹ֑ד זָ֠הָ֠ב תְּכֵ֨לֶת וְאַרְגָּמָ֜ן תּוֹלַ֧עַת שָׁנִ֛י וְשֵׁ֥שׁ מׇשְׁזָ֖ר מַעֲשֵׂ֥ה חֹשֵֽׁב׃ \Rashi{\textbf{ועשו את האפד.}אם באתי לפרש מעשה האפוד והחשן על סדר המקראות, הרי פרושן פרקים, וישגה הקורא בצרופן, לכך אני כותב מעשיהם כמות שהוא, למען ירוץ הקורא בו, ואחר כך אפרש על סדר המקראות. האפוד עשוי כמין סינר של נשים רוכבות סוסים, וחוגר אותו מאחוריו כנגד לבו למטה מאציליו, רחבו כמדת רחב גבו של אדם ויותר, ומגיע עד עקביו, והחשב מחבר בראשו על פני רחבו מעשה אורג, ומאריך לכאן ולכאן כדי להקיף ולחגר בו, והכתפות מחברות בחשב, אחד לימין ואחד לשמאל, מאחורי הכהן, לשני קצות רחבו של סינר, וכשזוקפן עומדות לו על שתי כתפיו, והן כמין שתי רצועות עשויות ממין האפוד, ארכות כדי שעור לזקפן אצל צוארו מכאן ומכאן ונקפלות לפניו למטה מכתפיו מעט, ואבני השהם קבועות בהם, אחת על כתף ימין ואחת על כתף שמאל, והמשבצות נתונות בראשיהם לפני כתפיו, ושתי עבותות הזהב תחובות בשתי טבעות שבחשן, בשני קצות רחבו העליון, אחת לימין ואחת לשמאל, ושני ראשי השרשרת תקועין במשבצת לימין וכן שני ראשי השרשרת השמאלית תקועין במשבצת שבכתף שמאל, נמצא החשן תלוי במשבצות האפוד על לבו מלפניו, ועוד שתי טבעות בשני קצות החשן בתחתיתו, וכנגדם שתי טבעות בשתי כתפות האפוד מלמטה, בראשו התחתון המחבר בחשב; טבעות החשן אל מול טבעות האפוד שוכבים זה על זה, ומרכסן בפתיל תכלת תחוב בטבעות האפוד והחשן, שיהא תחתית החשן דבוק לחשב האפוד ולא יהא נד ונבדל, הולך וחוזר: זהב תכלת וארגמן תולעת שני ושש משזר. חמשת מינים הללו שזורין בכל חוט וחוט; היו מרדדין את הזהב כמין טסין דקין וקוצצין פתילים מהם, וטווין אותו חוט של זהב עם ששה חוטין של תכלת, וחוט של זהב עם ששה חוטין של ארגמן, וכן בתולעת שני וכן בשש – שכל המינין חוטן כפול ששה וחוט של זהב עם כל אחד ואחד – ואחר כך שוזר את כלם כאחד, נמצא חוטן כפול כ"ח, וכן מפרש במסכת יומא, ולמד מן המקרא הזה וירקעו את פחי הזהב וקצץ פתילם לעשות – את פתילי הזהב – בתוך התכלת ובתוך הארגמן וגו', למדנו שחוט של זהב שזור עם כל מין ומין: מעשה חשב. כבר פרשתי שהוא אריגת שתי קירות, שאין צורות שני עבריה דומות זו לזו:}%
% ---------- פרק 28 | פסוק 7 ----------
 \textbf{ז} שְׁתֵּ֧י כְתֵפֹ֣ת חֹֽבְרֹ֗ת יִֽהְיֶה־לּ֛וֹ אֶל־שְׁנֵ֥י קְצוֹתָ֖יו וְחֻבָּֽר׃ \Rashi{\textbf{שתי כתפת וגו'.}הסינר מלמטה וחשב האפוד היא החגורה, וצמודה לו מלמעלה דגמת סינר הנשים, ומגבו של כהן היו מחברות בחשב שתי חתיכות כמין שתי רצועות רחבות, אחת כנגד כל כתף וכתף, וזוקפן על שתי כתפותיו עד שנקפלות לפניו כנגד החזה, וע"י חבורן לטבעות החשן נאחזין מלפניו כנגד לבו שאין נופלות, כמו שמפרש בענין, והיו זקופות והולכות כנגד כתפיו ושתי אבני השהם קבועות בהן אחת בכל אחת: אל שני קצותיו. אל רחבו של אפוד, שלא היה רחבו אלא כנגד גבו של כהן וגבהו עד כנגד האצילים, שקורין קודי"ש בלעז, שנ' לא יחגרו ביזע (יחזקאל מ"ד) – אין חוגרין במקום זיעה, לא למעלה מאציליהם ולא למטה ממתניהם אלא כנגד אציליהם: וחבר. האפוד עם אותן שתי כתפות האפוד יחבר אותם במחט למטה בחשב, ולא יארגם עמו אלא אורגם לבד ואחר כך מחברם:}%
% ---------- פרק 28 | פסוק 8 ----------
 \textbf{ח} וְחֵ֤שֶׁב אֲפֻדָּתוֹ֙ אֲשֶׁ֣ר עָלָ֔יו כְּמַעֲשֵׂ֖הוּ מִמֶּ֣נּוּ יִהְיֶ֑ה זָהָ֗ב תְּכֵ֧לֶת וְאַרְגָּמָ֛ן וְתוֹלַ֥עַת שָׁנִ֖י וְשֵׁ֥שׁ מׇשְׁזָֽר׃ \Rashi{\textbf{וחשב אפדתו.}וחגור שעל ידו הוא מאפדו ומתקנהו לכהן ומקשטו: אשר עליו. למעלה בשפת הסינר, היא החגורה: כמעשהו. כאריגת הסינר מעשה חושב ומחמשת מינים, כך אריגת החשב מעשה חושב ומחמשת מינין: ממנו יהיה. עמו יהיה ארוג ולא יארגנו לבד ויחברנו:}%
% ---------- פרק 28 | פסוק 9 ----------
 \textbf{ט} וְלָ֣קַחְתָּ֔ אֶת־שְׁתֵּ֖י אַבְנֵי־שֹׁ֑הַם וּפִתַּחְתָּ֣ עֲלֵיהֶ֔ם שְׁמ֖וֹת בְּנֵ֥י יִשְׂרָאֵֽל׃ %
% ---------- פרק 28 | פסוק 10 ----------
 \textbf{י} שִׁשָּׁה֙ מִשְּׁמֹתָ֔ם עַ֖ל הָאֶ֣בֶן הָאֶחָ֑ת וְאֶת־שְׁמ֞וֹת הַשִּׁשָּׁ֧ה הַנּוֹתָרִ֛ים עַל־הָאֶ֥בֶן הַשֵּׁנִ֖ית כְּתוֹלְדֹתָֽם׃ \Rashi{\textbf{כתולדתם.}כסדר שנולדו – ראובן, שמעון, לוי, יהודה, דן, נפתלי על האחת, ועל השניה גד, אשר, יששכר זבולן, יוסף, בנימין מלא, שכן הוא כתוב במקום תולדתו, כ"ה אותיות בכל אחת ואחת:}%
% ---------- פרק 28 | פסוק 11 ----------
 \textbf{יא} מַעֲשֵׂ֣ה חָרַשׁ֮ אֶ֒בֶן֒ פִּתּוּחֵ֣י חֹתָ֗ם תְּפַתַּח֙ אֶת־שְׁתֵּ֣י הָאֲבָנִ֔ים עַל־שְׁמֹ֖ת בְּנֵ֣י יִשְׂרָאֵ֑ל מֻסַבֹּ֛ת מִשְׁבְּצ֥וֹת זָהָ֖ב תַּעֲשֶׂ֥ה אֹתָֽם׃ \Rashi{\textbf{מעשה חרש אבן.}מעשה אמן של אבנים, חרש זה דבוק הוא לתבה שלאחריו ולפיכך הוא נקוד פתח בסופו, וכן חרש עצים נטה קו (ישעיהו מ"ד) – חרש של עצים, וכן חרש ברזל מעצד, (שם), כל אלה דבוקים ופתוחים: פתוחי חתם. כתרגומו כתב מפרש כגלף דעזקא, חרוצות האותיות בתוכן כמו שחורצין חותמי טבעות שהם לחתם אגרות – כתב נכר ומפרש: על שמת. כמו בשמת: מסבת משבצות. מקפות האבנים במשבצות זהב, שעושה מושב האבן בזהב כמין גמא למדת האבן ומשקעה במשבצת, נמצאת המשבצת סובבת את האבן סביב ומחבר המשבצות בכתפות האפוד:}%
% ---------- פרק 28 | פסוק 12 ----------
 \textbf{יב} וְשַׂמְתָּ֞ אֶת־שְׁתֵּ֣י הָאֲבָנִ֗ים עַ֚ל כִּתְפֹ֣ת הָֽאֵפֹ֔ד אַבְנֵ֥י זִכָּרֹ֖ן לִבְנֵ֣י יִשְׂרָאֵ֑ל וְנָשָׂא֩ אַהֲרֹ֨ן אֶת־שְׁמוֹתָ֜ם לִפְנֵ֧י יְהֹוָ֛ה עַל־שְׁתֵּ֥י כְתֵפָ֖יו לְזִכָּרֹֽן׃ \{ס\} \Rashi{\textbf{לזכרן.}שיהא רואה הקב"ה את השבטים כתובים לפניו ויזכר צדקתם:}%
% ---------- פרק 28 | פסוק 13 ----------
 \textbf{יג} וְעָשִׂ֥יתָ מִשְׁבְּצֹ֖ת זָהָֽב׃ \Rashi{\textbf{ועשית משבצת.}מעוט משבצות שתים; ולא פרש לך עתה בפרשה זו אלא מקצת צרכן, ובפרשת החשן גומר לך פרושן:}%
% ---------- פרק 28 | פסוק 14 ----------
 \textbf{יד} וּשְׁתֵּ֤י שַׁרְשְׁרֹת֙ זָהָ֣ב טָה֔וֹר מִגְבָּלֹ֛ת תַּעֲשֶׂ֥ה אֹתָ֖ם מַעֲשֵׂ֣ה עֲבֹ֑ת וְנָתַתָּ֛ה אֶת־שַׁרְשְׁרֹ֥ת הָעֲבֹתֹ֖ת עַל־הַֽמִּשְׁבְּצֹֽת׃ \{ס\} \Rashi{\textbf{שרשרת זהב.}שלשלאות: מגבלת. לסוף גבול החשן תעשה אותם: מעשה עבת. מעשה קליעת חוטין, ולא מעשה נקבים וכפלים כאותן שעושין לבורות, אלא כאותן שעושין לערדסקאות שקורין אנשנשוי"רש בלעז: ונתתה את שרשרת. של עבותות העשויות מעשה עבות על משבצות הללו; ולא זה הוא מקום צואת עשיתן של שרשרות, ולא צואת קביעתן, אין תעשה האמור כאן לשון צווי, ואין ונתתה האמור כאן לשון צווי, אלא לשון עתיד, כי בפרשת החשן חוזר ומצוה על עשיתן ועל קביעתן, ולא נכתב כאן אלא להודיע מקצת צרך המשבצות שצוה לעשות עם האפוד, וכתב לך זאת לומר לך המשבצות הללו יזקקו לך, לכשתעשה שרשרות מגבלות על החשן, תתנם על המשבצות הללו:}%
% ---------- פרק 28 | פסוק 15 ----------
 \textbf{טו} וְעָשִׂ֜יתָ חֹ֤שֶׁן מִשְׁפָּט֙ מַעֲשֵׂ֣ה חֹשֵׁ֔ב כְּמַעֲשֵׂ֥ה אֵפֹ֖ד תַּעֲשֶׂ֑נּוּ זָ֠הָ֠ב תְּכֵ֨לֶת וְאַרְגָּמָ֜ן וְתוֹלַ֧עַת שָׁנִ֛י וְשֵׁ֥שׁ מׇשְׁזָ֖ר תַּעֲשֶׂ֥ה אֹתֽוֹ׃ \Rashi{\textbf{חשן משפט.}שמכפר על קלקול הדין. ד"א, משפט שמברר דבריו והבטחתו אמת, דריש"נטנט בלעז; שהמשפט משמש שלש לשונות, דברי טענות הבעלי דינים וגמר הדין וענש הדין, אם ענש מיתה, אם ענש מכות, אם ענש ממון, וזה משמש לשון ברור דברים, שמפרש ומברר דבריו: כמעשה אפד. מעשה חושב ומחמשת מינין:}%
% ---------- פרק 28 | פסוק 16 ----------
 \textbf{טז} רָב֥וּעַ יִֽהְיֶ֖ה כָּפ֑וּל זֶ֥רֶת אׇרְכּ֖וֹ וְזֶ֥רֶת רׇחְבּֽוֹ׃ \Rashi{\textbf{זרת ארכו וזרת רחבו.}כפול; ומטל לו לפניו כנגד לבו, שנאמר והיו על לב אהרן, תלוי בכתפות האפוד הבאות מאחוריו על כתפיו ונקפלות ויורדות לפניו מעט, והחשן תלוי בהן בשרשרות וטבעות כמו שמפרש בענין:}%
% ---------- פרק 28 | פסוק 17 ----------
 \textbf{יז} וּמִלֵּאתָ֥ בוֹ֙ מִלֻּ֣אַת אֶ֔בֶן אַרְבָּעָ֖ה טוּרִ֣ים אָ֑בֶן ט֗וּר אֹ֤דֶם פִּטְדָה֙ וּבָרֶ֔קֶת הַטּ֖וּר הָאֶחָֽד׃ \Rashi{\textbf{ומלאת בו.}על שם שהאבנים ממלאות גמות המשבצות המתקנות להן, קורא אותן בלשון מלואים:}%
% ---------- פרק 28 | פסוק 18 ----------
 \textbf{יח} וְהַטּ֖וּר הַשֵּׁנִ֑י נֹ֥פֶךְ סַפִּ֖יר וְיָהֲלֹֽם׃ %
% ---------- פרק 28 | פסוק 19 ----------
 \textbf{יט} וְהַטּ֖וּר הַשְּׁלִישִׁ֑י לֶ֥שֶׁם שְׁב֖וֹ וְאַחְלָֽמָה׃ %
% ---------- פרק 28 | פסוק 20 ----------
 \textbf{כ} וְהַטּוּר֙ הָרְבִיעִ֔י תַּרְשִׁ֥ישׁ וְשֹׁ֖הַם וְיָשְׁפֵ֑ה מְשֻׁבָּצִ֥ים זָהָ֛ב יִהְי֖וּ בְּמִלּוּאֹתָֽם׃ \Rashi{\textbf{משבצים זהב יהיו.}הטורים במלואתם – מקפים משבצות זהב בעמק שעור שיתמלא בעבי האבן, זהו לשון במלואתם, כשעור מלוי עבין של אבנים יהיה עמק המשבצות, לא פחות ולא יותר:}%
% ---------- פרק 28 | פסוק 21 ----------
 \textbf{כא} וְ֠הָאֲבָנִ֠ים תִּֽהְיֶ֜יןָ עַל־שְׁמֹ֧ת בְּנֵֽי־יִשְׂרָאֵ֛ל שְׁתֵּ֥ים עֶשְׂרֵ֖ה עַל־שְׁמֹתָ֑ם פִּתּוּחֵ֤י חוֹתָם֙ אִ֣ישׁ עַל־שְׁמ֔וֹ תִּֽהְיֶ֕יןָ לִשְׁנֵ֥י עָשָׂ֖ר שָֽׁבֶט׃ \Rashi{\textbf{איש על שמו.}כסדר תולדותם סדר האבנים, אדם לראובן, פטדה לשמעון, וכן כלם:}%
% ---------- פרק 28 | פסוק 22 ----------
 \textbf{כב} וְעָשִׂ֧יתָ עַל־הַחֹ֛שֶׁן שַֽׁרְשֹׁ֥ת גַּבְלֻ֖ת מַעֲשֵׂ֣ה עֲבֹ֑ת זָהָ֖ב טָהֽוֹר׃ \Rashi{\textbf{על החשן.}בשביל החשן, לקבעם בטבעותיו, כמו שמפרש למטה בענין: שרשת. לשון שרשי אילן, המאחזין לאילן להאחז ולהתקע בארץ, אף אלו יהיו מאחזין לחשן שבהם יהיה תלוי באפוד, והן שתי שרשרות האמורות למעלה בענין המשבצות; ואף שרשרות פתר מנחם בן סרוק לשון שרשים, ואמר שהרי"ש יתרה כמו מ"ם שבשלשום ומ"ם שבריקם, ואיני רואה את דבריו, אלא שרשרת בלשון עברית כשלשלת בלשון משנה: גבלת. הוא מגבלות האמור למעלה, שתתקעם בטבעות שיהיו בגבול החשן, וכל גבול לשון קצה, אשומי"יל בלעז: מעשה עבת. מעשה קליעה:}%
% ---------- פרק 28 | פסוק 23 ----------
 \textbf{כג} וְעָשִׂ֙יתָ֙ עַל־הַחֹ֔שֶׁן שְׁתֵּ֖י טַבְּע֣וֹת זָהָ֑ב וְנָתַתָּ֗ אֶת־שְׁתֵּי֙ הַטַּבָּע֔וֹת עַל־שְׁנֵ֖י קְצ֥וֹת הַחֹֽשֶׁן׃ \Rashi{\textbf{על החשן.}לצרך החשן – כדי לקבעם בו: ולא יתכן לומר שתהא תחלת עשיתן עליו, שאם כן מה הוא שחוזר ואומר ונתת את שתי הטבעות והלא כבר נתונין בו? היה לו לכתב בתחלת המקרא ועשית על קצות החשן שתי טבעות זהב, ואף בשרשרות צריך אתה לפתר כן: על שני קצות החשן. לשתי פאות שכנגד הצואר, לימנית ולשמאלית, הבאים מול כתפות האפוד:}%
% ---------- פרק 28 | פסוק 24 ----------
 \textbf{כד} וְנָתַתָּ֗ה אֶת־שְׁתֵּי֙ עֲבֹתֹ֣ת הַזָּהָ֔ב עַל־שְׁתֵּ֖י הַטַּבָּעֹ֑ת אֶל־קְצ֖וֹת הַחֹֽשֶׁן׃ \Rashi{\textbf{ונתתה את שתי עבתת הזהב.}הן הן שרשות גבלות הכתובות למעלה, ולא פרש מקום קבוען בחשן, עכשו מפרש לך שיהא תוחב אותן בטבעות; ותדע לך שהן הן הראשונות, שהרי בפרשת אלה פקודי לא הכפלו:}%
% ---------- פרק 28 | פסוק 25 ----------
 \textbf{כה} וְאֵ֨ת שְׁתֵּ֤י קְצוֹת֙ שְׁתֵּ֣י הָעֲבֹתֹ֔ת תִּתֵּ֖ן עַל־שְׁתֵּ֣י הַֽמִּשְׁבְּצ֑וֹת וְנָתַתָּ֛ה עַל־כִּתְפ֥וֹת הָאֵפֹ֖ד אֶל־מ֥וּל פָּנָֽיו׃ \Rashi{\textbf{ואת שתי קצות של שתי עבתת, שני ראשיהם של כל אחת.}תתן על שתי המשבצות. הן הן הכתובות למעלה, בין פרשת החשן ופרשת האפוד, ולא פרש את צרכן ואת מקומן, עכשו מפרש שיתקע בהן ראשי העבותות התחובות בטבעות החשן לימין ולשמאל אצל הצואר; שני ראשי שרשרת הימנית תוקע במשבצות של ימין, וכן בשל שמאל שני ראשי השרשרת השמאלית: ונתתה. המשבצות על כתפות האפוד, אחת בזו ואחת בזו, נמצאו כתפות האפוד מחזיקין את החשן שלא יפל, ובהן הוא תלוי, ועדין שפת החשן התחתונה הולכת ובאה ונוקשת על כרסו ואינה דבוקה לו יפה, לכך הצרך עוד שתי טבעות לתחתיתו, כמו שמפרש והולך: אל מול פניו. של אפוד, שלא יתן המשבצות בעבר הכתפות שכלפי המעיל, אלא בעבר העליון שכלפי החוץ, והוא קרוי מול פניו של אפוד, כי אותו עבר שאינו נראה אינו קרוי פנים:}%
% ---------- פרק 28 | פסוק 26 ----------
 \textbf{כו} וְעָשִׂ֗יתָ שְׁתֵּי֙ טַבְּע֣וֹת זָהָ֔ב וְשַׂמְתָּ֣ אֹתָ֔ם עַל־שְׁנֵ֖י קְצ֣וֹת הַחֹ֑שֶׁן עַל־שְׂפָת֕וֹ אֲשֶׁ֛ר אֶל־עֵ֥בֶר הָאֵפֹ֖ד*(בספרי ספרד ואשכנז הָאֵפ֖וֹד) בָּֽיְתָה׃ \Rashi{\textbf{על שני קצות החשן.}הן שתי פאותיו התחתונות לימין ולשמאל: על שפתו אשר אל עבר האפוד ביתה. הרי לך שני סימנין, האחד שיתנם בשני קצוות של תחתיתו שהוא כנגד האפוד, שעליונו אינו כנגד האפוד, שהרי סמוך לצואר הוא והאפוד נתון על מתניו, ועוד נתן סימן, שלא יקבעם בעבר החשן שכלפי החוץ, אלא בעבר שכלפי פנים, שנאמר ביתה, ואותו העבר הוא לצד האפוד, שחשב האפוד חוגרו הכהן ונקפל הסינר לפני הכהן על מתניו, וקצת כרסו מכאן ומכאן עד כנגד קצות החשן וקצותיו שוכבין עליו:}%
% ---------- פרק 28 | פסוק 27 ----------
 \textbf{כז} וְעָשִׂ֘יתָ֮ שְׁתֵּ֣י טַבְּע֣וֹת זָהָב֒ וְנָתַתָּ֣ה אֹתָ֡ם עַל־שְׁתֵּי֩ כִתְפ֨וֹת הָאֵפ֤וֹד מִלְּמַ֙טָּה֙ מִמּ֣וּל פָּנָ֔יו לְעֻמַּ֖ת מַחְבַּרְתּ֑וֹ מִמַּ֕עַל לְחֵ֖שֶׁב הָאֵפֽוֹד׃ \Rashi{\textbf{על שתי כתפות מלמטה.}שהמשבצות נתונות בראשי כתפות האפוד העליונים הבאים על כתפיו, כנגד גרונו, ונקפלות ויורדות לפניו, והטבעות צוה לתן בראשן השני, שהוא מחבר לאפוד, והוא שנ' לעמת מחברתו, סמוך למקום חבורן באפוד, למעלה מן החגורה מעט, שהמחברת לעמת החגורה, ואלו נתונים מעט בגבה זקיפת הכתפות, הוא שנאמר ממעל לחשב האפוד, והן כנגד סוף החשן, ונותן פתיל תכלת באותן הטבעות ובטבעות החשן ורוכסן באותו פתיל לימין ולשמאל, שלא יהא תחתית החשן הולך לפנים וחוזר לאחור, ונוקש על כרסו, ונמצא מישב על המעיל יפה: ממול פניו. בעבר החיצון:}%
% ---------- פרק 28 | פסוק 28 ----------
 \textbf{כח} וְיִרְכְּס֣וּ אֶת־הַ֠חֹ֠שֶׁן מִֽטַּבְּעֹתָ֞ו אֶל־טַבְּעֹ֤ת הָאֵפוֹד֙ בִּפְתִ֣יל תְּכֵ֔לֶת לִֽהְי֖וֹת עַל־חֵ֣שֶׁב הָאֵפ֑וֹד וְלֹֽא־יִזַּ֣ח הַחֹ֔שֶׁן מֵעַ֖ל הָאֵפֽוֹד׃ \Rashi{\textbf{וירכסו.}לשון חבור; וכן מרכסי איש (תהילים ל"א) – חבורי חבלי רשעים, וכן והרכסים לבקעה (ישעיהו מ') – הרים הסמוכים זה לזה, שאי אפשר לירד לגיא שביניהם אלא בקשי גדול, שמתוך סמיכתן הגיא זקופה ועמקה – יהא לבקעת מישור ונוחה לילך: להיות על חשב האפוד. להיות החשן דבוק אל חשב האפוד: ולא יזח. לשון נתוק; ולשון ערבי הוא כדברי דונש בן לברט:}%
% ---------- פרק 28 | פסוק 29 ----------
 \textbf{כט} וְנָשָׂ֣א אַ֠הֲרֹ֠ן אֶת־שְׁמ֨וֹת בְּנֵֽי־יִשְׂרָאֵ֜ל בְּחֹ֧שֶׁן הַמִּשְׁפָּ֛ט עַל־לִבּ֖וֹ בְּבֹא֣וֹ אֶל־הַקֹּ֑דֶשׁ לְזִכָּרֹ֥ן לִפְנֵֽי־יְהֹוָ֖ה תָּמִֽיד׃ %
% ---------- פרק 28 | פסוק 30 ----------
 \textbf{ל} וְנָתַתָּ֞ אֶל־חֹ֣שֶׁן הַמִּשְׁפָּ֗ט אֶת־הָאוּרִים֙ וְאֶת־הַתֻּמִּ֔ים וְהָיוּ֙ עַל־לֵ֣ב אַהֲרֹ֔ן בְּבֹא֖וֹ לִפְנֵ֣י יְהֹוָ֑ה וְנָשָׂ֣א אַ֠הֲרֹ֠ן אֶת־מִשְׁפַּ֨ט בְּנֵי־יִשְׂרָאֵ֧ל עַל־לִבּ֛וֹ לִפְנֵ֥י יְהֹוָ֖ה תָּמִֽיד׃ \{ס\} \Rashi{\textbf{את האורים ואת התמים.}הוא כתב שם המפרש, שהיה נותנו בתוך כפלי החשן, שעל ידו הוא מאיר דבריו ומתמם את דבריו; ובמקדש שני היה החשן – שאי אפשר לכהן גדול להיות מחסר בגדים – אבל אותו השם לא היה בתוכו, ועל שם אותו הכתב הוא קרוי משפט, שנאמר (במדבר כ"ז), ושאל לו במשפט האורים (יומא ע"ג): את משפט בני ישראל. דבר שהם נשפטים ונוכחים על ידו, אם לעשות דבר או לא לעשות; ולפי מ"א שהחשן מכפר על מעותי הדין, נקרא משפט על שם סליחת המשפט:}%
% ---------- פרק 28 | פסוק 31 ----------
 \textbf{לא} וְעָשִׂ֛יתָ אֶת־מְעִ֥יל הָאֵפ֖וֹד כְּלִ֥יל תְּכֵֽלֶת׃ \Rashi{\textbf{את מעיל האפוד.}שהאפוד נתון עליו לחגורה: כליל תכלת. כלו תכלת, שאין מין אחר מערב בו:}%
% ---------- פרק 28 | פסוק 32 ----------
 \textbf{לב} וְהָיָ֥ה פִֽי־רֹאשׁ֖וֹ בְּתוֹכ֑וֹ שָׂפָ֡ה יִֽהְיֶה֩ לְפִ֨יו סָבִ֜יב מַעֲשֵׂ֣ה אֹרֵ֗ג כְּפִ֥י תַחְרָ֛א יִֽהְיֶה־לּ֖וֹ לֹ֥א יִקָּרֵֽעַ׃ \Rashi{\textbf{והיה פי ראשו.}פי המעיל שבגבהו – הוא פתיחת בית הצואר: בתוכו. כתרגומו כפיל לגוה, כפול לתוכו, להיות לו לשפה כפילתו, והיה מעשה אורג ולא במחט: כפי תחרא. למדנו שהשריונים שלהם פיהם כפול: לא יקרע. כדי שלא יקרע; והקורעו עובר בלאו, שזה ממנין לאוין שבתורה, וכן ולא יזח החשן, וכן לא יסרו ממנו הנאמר בבדי הארון:}%
% ---------- פרק 28 | פסוק 33 ----------
 \textbf{לג} וְעָשִׂ֣יתָ עַל־שׁוּלָ֗יו רִמֹּנֵי֙ תְּכֵ֤לֶת וְאַרְגָּמָן֙ וְתוֹלַ֣עַת שָׁנִ֔י עַל־שׁוּלָ֖יו סָבִ֑יב וּפַעֲמֹנֵ֥י זָהָ֛ב בְּתוֹכָ֖ם סָבִֽיב׃ \Rashi{\textbf{רמני.}עגלים וחלולים היו כמין רמונים העשויים כביצת תרנגלת: ופעמני זהב. זגין עם ענבלין שבתוכם: בתוכם סביב. ביניהם סביב – בין שני רמונים פעמון אחד דבוק ותלוי בשולי המעיל:}%
% ---------- פרק 28 | פסוק 34 ----------
 \textbf{לד} פַּעֲמֹ֤ן זָהָב֙ וְרִמּ֔וֹן פַּֽעֲמֹ֥ן זָהָ֖ב וְרִמּ֑וֹן עַל־שׁוּלֵ֥י הַמְּעִ֖יל סָבִֽיב׃ \Rashi{\textbf{פעמן זהב ורמון פעמן זהב ורמון.}אצלו:}%
% ---------- פרק 28 | פסוק 35 ----------
 \textbf{לה} וְהָיָ֥ה עַֽל־אַהֲרֹ֖ן לְשָׁרֵ֑ת וְנִשְׁמַ֣ע ק֠וֹל֠וֹ בְּבֹא֨וֹ אֶל־הַקֹּ֜דֶשׁ לִפְנֵ֧י יְהֹוָ֛ה וּבְצֵאת֖וֹ וְלֹ֥א יָמֽוּת׃ \{ס\} \Rashi{\textbf{ולא ימות.}מכלל לאו אתה שומע הן – אם יהיו לו לא יתחיב מיתה, הא אם יכנס מחסר אחד מן הבגדים הללו, חיב מיתה בידי שמים (סנהדרין פ"ג):}%
% ---------- פרק 28 | פסוק 36 ----------
 \textbf{לו} וְעָשִׂ֥יתָ צִּ֖יץ זָהָ֣ב טָה֑וֹר וּפִתַּחְתָּ֤ עָלָיו֙ פִּתּוּחֵ֣י חֹתָ֔ם קֹ֖דֶשׁ לַֽיהֹוָֽה׃ \Rashi{\textbf{ציץ.}כמין טס של זהב היה, רחב שתי אצבעות, מקיף על המצח מאזן לאזן:}%
% ---------- פרק 28 | פסוק 37 ----------
 \textbf{לז} וְשַׂמְתָּ֤ אֹתוֹ֙ עַל־פְּתִ֣יל תְּכֵ֔לֶת וְהָיָ֖ה עַל־הַמִּצְנָ֑פֶת אֶל־מ֥וּל פְּנֵֽי־הַמִּצְנֶ֖פֶת יִהְיֶֽה׃ \Rashi{\textbf{על פתיל תכלת.}ובמקום אחר הוא אומר ויתנו עליו פתיל תכלת? ועוד כתיב כאן והיה על המצנפת, ולמטה הוא אומר והיה על מצח אהרן? ובשחיטת קדשים שנינו שערו היה נראה בין ציץ למצנפת ששם מניח תפלין? למדנו שהמצנפת למעלה בגבה הראש ואינה עמקה לכנס בה כל הראש עד המצח, והציץ מלמטה, והפתילים היו בנקבים ותלויין בו בשני ראשים ובאמצעו, ששה בשלשה מקומות הללו, פתיל מלמעלה אחד מבחוץ ואחד מבפנים כנגדו, וקושר ראשי הפתילים מאחורי הערף שלשתן, ונמצאו בין ארך הטס ופתילי ראשיו מקיפין את הקדקד, והפתיל האמצעי שבראשו, קשור עם ראשי השנים, והולך על פני רחב הראש מלמעלה, נמצא עשוי כמין כובע; ועל פתיל האמצעי הוא אומר והיה על המצנפת, והיה נותן הציץ על ראשו כמין כובע על המצנפת, והפתיל האמצעי מחזיקו שאינו נופל והטס תלוי כנגד מצחו, ונתקימו כל המקראות – פתיל על הציץ, וציץ על הפתיל, ופתיל על המצנפת מלמעלה:}%
% ---------- פרק 28 | פסוק 38 ----------
 \textbf{לח} וְהָיָה֮ עַל־מֵ֣צַח אַהֲרֹן֒ וְנָשָׂ֨א אַהֲרֹ֜ן אֶת־עֲוֺ֣ן הַקֳּדָשִׁ֗ים אֲשֶׁ֤ר יַקְדִּ֙ישׁוּ֙ בְּנֵ֣י יִשְׂרָאֵ֔ל לְכׇֽל־מַתְּנֹ֖ת קׇדְשֵׁיהֶ֑ם וְהָיָ֤ה עַל־מִצְחוֹ֙ תָּמִ֔יד לְרָצ֥וֹן לָהֶ֖ם לִפְנֵ֥י יְהֹוָֽה׃ \Rashi{\textbf{ונשא אהרן.}לשון סליחה; ואעפ"כ אינו זז ממשמעו – אהרן נושא את המשא של עון – נמצא מסלק העון מן הקדשים (חולין קל"ח): את עון הקדשים. לרצות על הדם ועל החלב שקרבו בטמאה, כמו ששנינו: אי זה עון הוא נושא? אם עון פגול, הרי כבר נאמר לא ירצה (ויקרא י"ט), ואם עון נותר הרי נאמר לא יחשב (שם ז') – ואין לומר שיכפר על עון הכהן שהקריב טמא, שהרי עון הקדשים נאמר ולא עון המקריבים – הא אינו מרצה אלא להכשיר הקרבן (יומא ז'): והיה על מצחו תמיד. אי אפשר לומר שיהא על מצחו תמיד, שהרי אינו עליו אלא בשעת העבודה, אלא תמיד לרצות להם אפלו אינו על מצחו – שלא היה כהן גדול עובד באותה שעה, ולדברי האומר עודהו על מצחו מכפר ומרצה, ואם לאו אינו מרצה, נדרש על מצחו תמיד, מלמד שממשמש בו בעודו על מצחו, שלא יסיח דעתו ממנו (שם):}%
% ---------- פרק 28 | פסוק 39 ----------
 \textbf{לט} וְשִׁבַּצְתָּ֙ הַכְּתֹ֣נֶת שֵׁ֔שׁ וְעָשִׂ֖יתָ מִצְנֶ֣פֶת שֵׁ֑שׁ וְאַבְנֵ֥ט תַּעֲשֶׂ֖ה מַעֲשֵׂ֥ה רֹקֵֽם׃ \Rashi{\textbf{ושבצת.}עשה אותם משבצות משבצות, וכלה של שש:}%
% ---------- פרק 28 | פסוק 40 ----------
 \textbf{מ} וְלִבְנֵ֤י אַהֲרֹן֙ תַּעֲשֶׂ֣ה כֻתֳּנֹ֔ת וְעָשִׂ֥יתָ לָהֶ֖ם אַבְנֵטִ֑ים וּמִגְבָּעוֹת֙ תַּעֲשֶׂ֣ה לָהֶ֔ם לְכָב֖וֹד וּלְתִפְאָֽרֶת׃ \Rashi{\textbf{ולבני אהרן תעשה כתנת.}ארבעה בגדים הללו ולא יותר, כתנת ואבנט ומגבעות – היא מצנפת – ומכנסים כתובים למטה בפרשה:}%
% ---------- פרק 28 | פסוק 41 ----------
 \textbf{מא} וְהִלְבַּשְׁתָּ֤ אֹתָם֙ אֶת־אַהֲרֹ֣ן אָחִ֔יךָ וְאֶת־בָּנָ֖יו אִתּ֑וֹ וּמָשַׁחְתָּ֨ אֹתָ֜ם וּמִלֵּאתָ֧ אֶת־יָדָ֛ם וְקִדַּשְׁתָּ֥ אֹתָ֖ם וְכִהֲנ֥וּ לִֽי׃ \Rashi{\textbf{והלבשת אתם את אהרן.}אותם האמורין באהרן – חשן, ואפוד, ומעיל, וכתנת תשבץ, מצנפת, ואבנט, וציץ, ומכנסים הכתובים למטה בכלם: ואת בניו אתו. אותם הכתובים בהם: ומשחת אתם. את אהרן ואת בניו בשמן המשחה: ומלאת את ידם. כל מלוי ידים לשון חנוך, כשהוא נכנס לדבר להיות מחזק בו מאותו יום והלאה הוא: ובלשון לעז כשממנין אדם על פקידת דבר, נותן השליט בידו בית יד של עור, שקורין גאנ"ט בלעז, ועל ידו הוא מחזיקו בדבר, וקורין לאותו מסירה רווישט"יר בלעז, והוא מלוי ידים:}%
% ---------- פרק 28 | פסוק 42 ----------
 \textbf{מב} וַעֲשֵׂ֤ה לָהֶם֙ מִכְנְסֵי־בָ֔ד לְכַסּ֖וֹת בְּשַׂ֣ר עֶרְוָ֑ה מִמׇּתְנַ֥יִם וְעַד־יְרֵכַ֖יִם יִהְיֽוּ׃ \Rashi{\textbf{ועשה להם.}לאהרן ולבניו: מכנסי בד. הרי שמונה בגדים לכהן גדול וארבעה לכהן הדיוט:}%
% ---------- פרק 28 | פסוק 43 ----------
 \textbf{מג} וְהָיוּ֩ עַל־אַהֲרֹ֨ן וְעַל־בָּנָ֜יו בְּבֹאָ֣ם ׀ אֶל־אֹ֣הֶל מוֹעֵ֗ד א֣וֹ בְגִשְׁתָּ֤ם אֶל־הַמִּזְבֵּ֙חַ֙ לְשָׁרֵ֣ת בַּקֹּ֔דֶשׁ וְלֹא־יִשְׂא֥וּ עָוֺ֖ן וָמֵ֑תוּ חֻקַּ֥ת עוֹלָ֛ם ל֖וֹ וּלְזַרְע֥וֹ אַחֲרָֽיו׃ \{ס\} \Rashi{\textbf{והיו על אהרן.}כל הבגדים האלה, על אהרן הראויין לו: ועל בניו. האמורין בהם: בבאם אל אהל מועד. להיכל, וכן למשכן: ומתו. הא למדת שהמשמש מחסר בגדים במיתה. חקת עולם לו. כל מקום שנאמר חקת עולם הוא גזרה מיד ולדורות, לעכב בו (מנחות י"ט):}%
% ---------- פרק 29 | פסוק 1 ----------
 \textbf{א} וְזֶ֨ה הַדָּבָ֜ר אֲשֶֽׁר־תַּעֲשֶׂ֥ה לָהֶ֛ם לְקַדֵּ֥שׁ אֹתָ֖ם לְכַהֵ֣ן לִ֑י לְ֠קַ֠ח פַּ֣ר אֶחָ֧ד בֶּן־בָּקָ֛ר וְאֵילִ֥ם שְׁנַ֖יִם תְּמִימִֽם׃ \Rashi{\textbf{לקח.}כמו קח. ושתי גזירות הן, אחת של קיחה ואחת של לקיחה, ולהן פתרון אחד: פר אחד. לכפר על מעשה העגל שהוא פר:}%
% ---------- פרק 29 | פסוק 2 ----------
 \textbf{ב} וְלֶ֣חֶם מַצּ֗וֹת וְחַלֹּ֤ת מַצֹּת֙ בְּלוּלֹ֣ת בַּשֶּׁ֔מֶן וּרְקִיקֵ֥י מַצּ֖וֹת מְשֻׁחִ֣ים בַּשָּׁ֑מֶן סֹ֥לֶת חִטִּ֖ים תַּעֲשֶׂ֥ה אֹתָֽם׃ \Rashi{\textbf{ולחם מצות וחלת מצת ורקיקי מצות.}הרי אלו שלשה מינין רבוכה, וחלות ורקיקים לחם מצות, היא הקרויה למטה בענין (פסוק כג) חלת לחם שמן, על שם שנותן שמן ברבוכה כנגד החלות והרקיקין, וכל המינין באים עשר עשר חלות: בלולת בשמן. כשהן קמח יוצק בהן שמן ובוללן: משחים בשמן. אחר אפיתן מושחן כמין כי יונית, שהיא עשויה כנו"ן שלנו:}%
% ---------- פרק 29 | פסוק 3 ----------
 \textbf{ג} וְנָתַתָּ֤ אוֹתָם֙ עַל־סַ֣ל אֶחָ֔ד וְהִקְרַבְתָּ֥ אֹתָ֖ם בַּסָּ֑ל וְאֶ֨ת־הַפָּ֔ר וְאֵ֖ת שְׁנֵ֥י הָאֵילִֽם׃ \Rashi{\textbf{והקרבת אתם.}אל חצר המשכן ביום הקמתו:}%
% ---------- פרק 29 | פסוק 4 ----------
 \textbf{ד} וְאֶת־אַהֲרֹ֤ן וְאֶת־בָּנָיו֙ תַּקְרִ֔יב אֶל־פֶּ֖תַח אֹ֣הֶל מוֹעֵ֑ד וְרָחַצְתָּ֥ אֹתָ֖ם בַּמָּֽיִם׃ \Rashi{\textbf{ורחצת.}זו טבילת כל הגוף:}%
% ---------- פרק 29 | פסוק 5 ----------
 \textbf{ה} וְלָקַחְתָּ֣ אֶת־הַבְּגָדִ֗ים וְהִלְבַּשְׁתָּ֤ אֶֽת־אַהֲרֹן֙ אֶת־הַכֻּתֹּ֔נֶת וְאֵת֙ מְעִ֣יל הָאֵפֹ֔ד וְאֶת־הָאֵפֹ֖ד וְאֶת־הַחֹ֑שֶׁן וְאָפַדְתָּ֣ ל֔וֹ בְּחֵ֖שֶׁב הָאֵפֹֽד׃ \Rashi{\textbf{ואפדת.}קשט ותקן החגורה והסינר סביבותיו:}%
% ---------- פרק 29 | פסוק 6 ----------
 \textbf{ו} וְשַׂמְתָּ֥ הַמִּצְנֶ֖פֶת עַל־רֹאשׁ֑וֹ וְנָתַתָּ֛ אֶת־נֵ֥זֶר הַקֹּ֖דֶשׁ עַל־הַמִּצְנָֽפֶת׃ \Rashi{\textbf{נזר הקדש.}זה הציץ: על המצנפת. כמו שפרשתי למעלה – על ידי הפתיל האמצעי ושני פתילין שבראשו, הקשורין שלשתן מאחורי הערף, הוא נותנו על המצנפת כמין כובע:}%
% ---------- פרק 29 | פסוק 7 ----------
 \textbf{ז} וְלָֽקַחְתָּ֙ אֶת־שֶׁ֣מֶן הַמִּשְׁחָ֔ה וְיָצַקְתָּ֖ עַל־רֹאשׁ֑וֹ וּמָשַׁחְתָּ֖ אֹתֽוֹ׃ \Rashi{\textbf{ומשחת אתו.}אף משיחה זו כמין כי, נותן שמן על ראשו ובין רסי עיניו ומחברן באצבעו (כריתות ה'):}%
% ---------- פרק 29 | פסוק 8 ----------
 \textbf{ח} וְאֶת־בָּנָ֖יו תַּקְרִ֑יב וְהִלְבַּשְׁתָּ֖ם כֻּתֳּנֹֽת׃ %
% ---------- פרק 29 | פסוק 9 ----------
 \textbf{ט} וְחָגַרְתָּ֩ אֹתָ֨ם אַבְנֵ֜ט אַהֲרֹ֣ן וּבָנָ֗יו וְחָבַשְׁתָּ֤ לָהֶם֙ מִגְבָּעֹ֔ת וְהָיְתָ֥ה לָהֶ֛ם כְּהֻנָּ֖ה לְחֻקַּ֣ת עוֹלָ֑ם וּמִלֵּאתָ֥ יַֽד־אַהֲרֹ֖ן וְיַד־בָּנָֽיו׃ \Rashi{\textbf{והיתה להם.}מלוי ידים זה לכהנת עולם: ומלאת. על ידי הדברים האלה: יד אהרן ויד בניו. במנוי ופקדת הכהנה:}%
% ---------- פרק 29 | פסוק 10 ----------
 \textbf{י} וְהִקְרַבְתָּ֙ אֶת־הַפָּ֔ר לִפְנֵ֖י אֹ֣הֶל מוֹעֵ֑ד וְסָמַ֨ךְ אַהֲרֹ֧ן וּבָנָ֛יו אֶת־יְדֵיהֶ֖ם עַל־רֹ֥אשׁ הַפָּֽר׃ %
% ---------- פרק 29 | פסוק 11 ----------
 \textbf{יא} וְשָׁחַטְתָּ֥ אֶת־הַפָּ֖ר לִפְנֵ֣י יְהֹוָ֑ה פֶּ֖תַח אֹ֥הֶל מוֹעֵֽד׃ \Rashi{\textbf{פתח אהל מועד.}בחצר המשכן שלפני הפתח:}%
% ---------- פרק 29 | פסוק 12 ----------
 \textbf{יב} וְלָֽקַחְתָּ֙ מִדַּ֣ם הַפָּ֔ר וְנָתַתָּ֛ה עַל־קַרְנֹ֥ת הַמִּזְבֵּ֖חַ בְּאֶצְבָּעֶ֑ךָ וְאֶת־כׇּל־הַדָּ֣ם תִּשְׁפֹּ֔ךְ אֶל־יְס֖וֹד הַמִּזְבֵּֽחַ׃ \Rashi{\textbf{על קרנת.}למעלה בקרנות ממש: ואת כל הדם. שירי הדם: אל יסוד המזבח. כמין בליטת בית קבול עשוי לו סביב סביב, לאחר שעלה אמה מן הארץ:}%
% ---------- פרק 29 | פסוק 13 ----------
 \textbf{יג} וְלָֽקַחְתָּ֗ אֶֽת־כׇּל־הַחֵ֘לֶב֮ הַֽמְכַסֶּ֣ה אֶת־הַקֶּ֒רֶב֒ וְאֵ֗ת הַיֹּתֶ֙רֶת֙ עַל־הַכָּבֵ֔ד וְאֵת֙ שְׁתֵּ֣י הַכְּלָיֹ֔ת וְאֶת־הַחֵ֖לֶב אֲשֶׁ֣ר עֲלֵיהֶ֑ן וְהִקְטַרְתָּ֖ הַמִּזְבֵּֽחָה׃ \Rashi{\textbf{החלב המכסה את הקרב.}הוא הקרום שעל הכרס, שקורין טי"לא: ואת היתרת. הוא טרפשא דכבדא, שקורין איבר"יש: על הכבד. אף מן הכבד טל עמה (ספרא):}%
% ---------- פרק 29 | פסוק 14 ----------
 \textbf{יד} וְאֶת־בְּשַׂ֤ר הַפָּר֙ וְאֶת־עֹר֣וֹ וְאֶת־פִּרְשׁ֔וֹ תִּשְׂרֹ֣ף בָּאֵ֔שׁ מִח֖וּץ לַֽמַּחֲנֶ֑ה חַטָּ֖את הֽוּא׃ \Rashi{\textbf{תשרף באש.}לא מצינו חטאת חיצונה נשרפת אלא זו:}%
% ---------- פרק 29 | פסוק 15 ----------
 \textbf{טו} וְאֶת־הָאַ֥יִל הָאֶחָ֖ד תִּקָּ֑ח וְסָ֨מְכ֜וּ אַהֲרֹ֧ן וּבָנָ֛יו אֶת־יְדֵיהֶ֖ם עַל־רֹ֥אשׁ הָאָֽיִל׃ %
% ---------- פרק 29 | פסוק 16 ----------
 \textbf{טז} וְשָׁחַטְתָּ֖ אֶת־הָאָ֑יִל וְלָֽקַחְתָּ֙ אֶת־דָּמ֔וֹ וְזָרַקְתָּ֥ עַל־הַמִּזְבֵּ֖חַ סָבִֽיב׃ \Rashi{\textbf{וזרקת.}בכלי; אוחז במזרק וזורק כנגד הקרן כדי שיראה לכאן ולכאן, ואין קרבן טעון מתנה באצבע אלא חטאת בלבד, אבל שאר זבחים אינן טעונין קרן ולא אצבע, שמתן דמם מחצי המזבח ולמטה, ואינו עולה בכבש אלא עומד בארץ וזורק (זבחים נ"ג): סביב. כך מפרש בשחיטת קדשים שאין סביב אלא שתי מתנות שהן ארבע, האחת בקרן זוית זו, והאחת שכנגדה באלכסון, וכל מתנה נראית בשני צדי הקרן אילך ואילך, נמצא הדם נתון בארבע רוחות סביב, לכך קרוי סביב:}%
% ---------- פרק 29 | פסוק 17 ----------
 \textbf{יז} וְאֶ֨ת־הָאַ֔יִל תְּנַתֵּ֖חַ לִנְתָחָ֑יו וְרָחַצְתָּ֤ קִרְבּוֹ֙ וּכְרָעָ֔יו וְנָתַתָּ֥ עַל־נְתָחָ֖יו וְעַל־רֹאשֽׁוֹ׃ \Rashi{\textbf{על נתחיו.}עם נתחיו – מוסף על שאר הנתחים:}%
% ---------- פרק 29 | פסוק 18 ----------
 \textbf{יח} וְהִקְטַרְתָּ֤ אֶת־כׇּל־הָאַ֙יִל֙ הַמִּזְבֵּ֔חָה עֹלָ֥ה ה֖וּא לַֽיהֹוָ֑ה רֵ֣יחַ נִיח֔וֹחַ אִשֶּׁ֥ה לַיהֹוָ֖ה הֽוּא׃ \Rashi{\textbf{ריח ניחוח.}נחת רוח לפני, שאמרתי ונעשה רצוני: אשה. לשון אש, והיא הקטרת אברים שעל האש:}%
% ---------- פרק 29 | פסוק 19 ----------
 \textbf{יט} וְלָ֣קַחְתָּ֔ אֵ֖ת הָאַ֣יִל הַשֵּׁנִ֑י וְסָמַ֨ךְ אַהֲרֹ֧ן וּבָנָ֛יו אֶת־יְדֵיהֶ֖ם עַל־רֹ֥אשׁ הָאָֽיִל׃ %
% ---------- פרק 29 | פסוק 20 ----------
 \textbf{כ} וְשָׁחַטְתָּ֣ אֶת־הָאַ֗יִל וְלָקַחְתָּ֤ מִדָּמוֹ֙ וְנָֽתַתָּ֡ה עַל־תְּנוּךְ֩ אֹ֨זֶן אַהֲרֹ֜ן וְעַל־תְּנ֨וּךְ אֹ֤זֶן בָּנָיו֙ הַיְמָנִ֔ית וְעַל־בֹּ֤הֶן יָדָם֙ הַיְמָנִ֔ית וְעַל־בֹּ֥הֶן רַגְלָ֖ם הַיְמָנִ֑ית וְזָרַקְתָּ֧ אֶת־הַדָּ֛ם עַל־הַמִּזְבֵּ֖חַ סָבִֽיב׃ \Rashi{\textbf{תנוך.}הוא הסחוס האמצעי שבתוך האזן, שקורין טנדרו"ס בלעז: בהן ידם. הגודל ובפרק האמצעי (ספרא):}%
% ---------- פרק 29 | פסוק 21 ----------
 \textbf{כא} וְלָקַחְתָּ֞ מִן־הַדָּ֨ם אֲשֶׁ֥ר עַֽל־הַמִּזְבֵּ֘חַ֮ וּמִשֶּׁ֣מֶן הַמִּשְׁחָה֒ וְהִזֵּיתָ֤ עַֽל־אַהֲרֹן֙ וְעַל־בְּגָדָ֔יו וְעַל־בָּנָ֛יו וְעַל־בִּגְדֵ֥י בָנָ֖יו אִתּ֑וֹ וְקָדַ֥שׁ הוּא֙ וּבְגָדָ֔יו וּבָנָ֛יו וּבִגְדֵ֥י בָנָ֖יו אִתּֽוֹ׃ %
% ---------- פרק 29 | פסוק 22 ----------
 \textbf{כב} וְלָקַחְתָּ֣ מִן־הָ֠אַ֠יִל הַחֵ֨לֶב וְהָֽאַלְיָ֜ה וְאֶת־הַחֵ֣לֶב ׀ הַֽמְכַסֶּ֣ה אֶת־הַקֶּ֗רֶב וְאֵ֨ת יֹתֶ֤רֶת הַכָּבֵד֙ וְאֵ֣ת ׀ שְׁתֵּ֣י הַכְּלָיֹ֗ת וְאֶת־הַחֵ֙לֶב֙ אֲשֶׁ֣ר עֲלֵיהֶ֔ן וְאֵ֖ת שׁ֣וֹק הַיָּמִ֑ין כִּ֛י אֵ֥יל מִלֻּאִ֖ים הֽוּא׃ \Rashi{\textbf{החלב.}זה חלב הדקים או הקבה: והאליה. מן הכליות ולמטה, כמו שמפרש בויקרא שנ' (ויקרא ג'), לעמת העצה יסירנה – מקום שהכליות יועצות (חולין י"א), ובאמורי הפר לא נאמר אליה, שאין אליה קרבה אלא בכבש וכבשה ואיל, אבל שור ועז אין טעונין אליה: ואת שוק הימין. לא מצינו הקטרה בשוק הימין עם האמורים אלא זו בלבד: כי איל מלאים הוא. שלמים; לשון שלמות – שהשלם בכל; מגיד הכתוב שהמלואים שלמים, שמשימים שלום למזבח ולעובד העבודה ולבעלים, לכך אני מצריכו החזה, להיות לו לעובד העבודה למנה, וזהו משה ששמש במלואים, והשאר אכלו אהרן ובניו שהם בעלים, כמפרש בענין:}%
% ---------- פרק 29 | פסוק 23 ----------
 \textbf{כג} וְכִכַּ֨ר לֶ֜חֶם אַחַ֗ת וְֽחַלַּ֨ת לֶ֥חֶם שֶׁ֛מֶן אַחַ֖ת וְרָקִ֣יק אֶחָ֑ד מִסַּל֙ הַמַּצּ֔וֹת אֲשֶׁ֖ר לִפְנֵ֥י יְהֹוָֽה׃ \Rashi{\textbf{וככר לחם.}מן החלות: וחלת לחם שמן. ממין הרבוכה: ורקיק. מן הרקיקין, אחד מעשרה שבכל מין ומין. ולא מצינו תרומת לחם הבא עם זבח נקטרת אלא זו בלבד, שתרומת לחמי תודה ואיל נזיר נתונה לכהנים עם חזה ושוק, ומזה לא היה למשה למנה אלא חזה בלבד:}%
% ---------- פרק 29 | פסוק 24 ----------
 \textbf{כד} וְשַׂמְתָּ֣ הַכֹּ֔ל עַ֚ל כַּפֵּ֣י אַהֲרֹ֔ן וְעַ֖ל כַּפֵּ֣י בָנָ֑יו וְהֵנַפְתָּ֥ אֹתָ֛ם תְּנוּפָ֖ה לִפְנֵ֥י יְהֹוָֽה׃ \Rashi{\textbf{על כפי אהרן: והנפת.}שניהם עסוקים בתנופה – הבעלים והכהן, הא כיצד? כהן מניח ידו תחת יד הבעלים ומניף; ובזה היו אהרן ובניו בעלים, ומשה כהן: תנופה. מוליך ומביא למי שארבע רוחות העולם שלו; תנופה מעכבת ומבטלת פרעניות ורוחות רעות, מעלה ומוריד למי שהשמים והארץ שלו, ומעכבת טללים רעים (מנחות ס"ב):}%
% ---------- פרק 29 | פסוק 25 ----------
 \textbf{כה} וְלָקַחְתָּ֤ אֹתָם֙ מִיָּדָ֔ם וְהִקְטַרְתָּ֥ הַמִּזְבֵּ֖חָה עַל־הָעֹלָ֑ה לְרֵ֤יחַ נִיח֙וֹחַ֙ לִפְנֵ֣י יְהֹוָ֔ה אִשֶּׁ֥ה ה֖וּא לַיהֹוָֽה׃ \Rashi{\textbf{על העלה.}על האיל הראשון שהעלית עולה: לריח ניחוח. לנחת רוח למי שאמר ונעשה רצונו: אשה. לאש נתן: לה'. לשמו של מקום:}%
% ---------- פרק 29 | פסוק 26 ----------
 \textbf{כו} וְלָקַחְתָּ֣ אֶת־הֶֽחָזֶ֗ה מֵאֵ֤יל הַמִּלֻּאִים֙ אֲשֶׁ֣ר לְאַהֲרֹ֔ן וְהֵנַפְתָּ֥ אֹת֛וֹ תְּנוּפָ֖ה לִפְנֵ֣י יְהֹוָ֑ה וְהָיָ֥ה לְךָ֖ לְמָנָֽה׃ %
% ---------- פרק 29 | פסוק 27 ----------
 \textbf{כז} וְקִדַּשְׁתָּ֞ אֵ֣ת ׀ חֲזֵ֣ה הַתְּנוּפָ֗ה וְאֵת֙ שׁ֣וֹק הַתְּרוּמָ֔ה אֲשֶׁ֥ר הוּנַ֖ף וַאֲשֶׁ֣ר הוּרָ֑ם מֵאֵיל֙ הַמִּלֻּאִ֔ים מֵאֲשֶׁ֥ר לְאַהֲרֹ֖ן וּמֵאֲשֶׁ֥ר לְבָנָֽיו׃ \Rashi{\textbf{וקדשת את חזה התנופה ואת שוק התרומה וגו'.}קדשם לדורות, להיות נוהגת תרומתם והנפתם בחזה ושוק של שלמים, אבל לא להקטרה, אלא והיה לאהרן ולבניו לאכל: תנופה. לשון הולכה והבאה, ונטלי"ה בלעז: הורם. לשון מעלה ומוריד:}%
% ---------- פרק 29 | פסוק 28 ----------
 \textbf{כח} וְהָיָה֩ לְאַהֲרֹ֨ן וּלְבָנָ֜יו לְחׇק־עוֹלָ֗ם מֵאֵת֙ בְּנֵ֣י יִשְׂרָאֵ֔ל כִּ֥י תְרוּמָ֖ה ה֑וּא וּתְרוּמָ֞ה יִהְיֶ֨ה מֵאֵ֤ת בְּנֵֽי־יִשְׂרָאֵל֙ מִזִּבְחֵ֣י שַׁלְמֵיהֶ֔ם תְּרוּמָתָ֖ם לַיהֹוָֽה׃ \Rashi{\textbf{לחק עולם מאת בני ישראל.}שהשלמים לבעלים, ואת החזה ואת השוק יתנו לכהן: כי תרומה הוא. החזה ושוק הזה:}%
% ---------- פרק 29 | פסוק 29 ----------
 \textbf{כט} וּבִגְדֵ֤י הַקֹּ֙דֶשׁ֙ אֲשֶׁ֣ר לְאַהֲרֹ֔ן יִהְי֥וּ לְבָנָ֖יו אַחֲרָ֑יו לְמׇשְׁחָ֣ה בָהֶ֔ם וּלְמַלֵּא־בָ֖ם אֶת־יָדָֽם׃ \Rashi{\textbf{לבניו אחריו.}למי שבא בגדלה אחריו: למשחה. להתגדל בהם, שיש משיחה שהיא לשון שררה, כמו לך נתתים למשחה (במדבר י"ח), אל תגעו במשיחי (דהי"א טז): ולמלא בם את ידם. על ידי הבגדים הוא מתלבש בכהנה גדולה:}%
% ---------- פרק 29 | פסוק 30 ----------
 \textbf{ל} שִׁבְעַ֣ת יָמִ֗ים יִלְבָּשָׁ֧ם הַכֹּהֵ֛ן תַּחְתָּ֖יו מִבָּנָ֑יו אֲשֶׁ֥ר יָבֹ֛א אֶל־אֹ֥הֶל מוֹעֵ֖ד לְשָׁרֵ֥ת בַּקֹּֽדֶשׁ׃ \Rashi{\textbf{שבעת ימים.}רצופין: ילבשם הכהן. אשר יקום מבניו תחתיו לכהנה גדולה, כשימנוהו להיות כהן גדול: אשר יבא אל אהל מועד. אותו כהן המוכן לכנס לפני ולפנים ביום הכפורים, וזהו כהן גדול, שאין עבודת יום הכפורים כשרה אלא בו: תחתיו מבניו. מלמד שאם יש לו לכהן גדול בן ממלא את מקומו, ימנוהו כהן גדול תחתיו: (הכהן תחתיו מבניו. מכאן ראיה כל לשון כהן פועל עובד ממש, לפיכך נגון תביר נמשך לפניו):}%
% ---------- פרק 29 | פסוק 31 ----------
 \textbf{לא} וְאֵ֛ת אֵ֥יל הַמִּלֻּאִ֖ים תִּקָּ֑ח וּבִשַּׁלְתָּ֥ אֶת־בְּשָׂר֖וֹ בְּמָקֹ֥ם קָדֹֽשׁ׃ \Rashi{\textbf{במקם קדש.}בחצר אהל מועד, שהשלמים הללו קדשי קדשים היו:}%
% ---------- פרק 29 | פסוק 32 ----------
 \textbf{לב} וְאָכַ֨ל אַהֲרֹ֤ן וּבָנָיו֙ אֶת־בְּשַׂ֣ר הָאַ֔יִל וְאֶת־הַלֶּ֖חֶם אֲשֶׁ֣ר בַּסָּ֑ל פֶּ֖תַח אֹ֥הֶל מוֹעֵֽד׃ \Rashi{\textbf{פתח אהל מועד.}כל החצר קרוי כן:}%
% ---------- פרק 29 | פסוק 33 ----------
 \textbf{לג} וְאָכְל֤וּ אֹתָם֙ אֲשֶׁ֣ר כֻּפַּ֣ר בָּהֶ֔ם לְמַלֵּ֥א אֶת־יָדָ֖ם לְקַדֵּ֣שׁ אֹתָ֑ם וְזָ֥ר לֹא־יֹאכַ֖ל כִּי־קֹ֥דֶשׁ הֵֽם׃ \Rashi{\textbf{ואכלו אתם.}אהרן ובניו, לפי שהם בעליהם: אשר כפר בהם. כל זרות ותעוב: למלא את ידם. באיל ולחם הללו: לקדש אתם. שעל ידי המלואים הללו נתמלאו ידיהם ונתקדשו לכהנה: כי קדש הם. קדשי קדשים; ומכאן למדנו אזהרה לזר האוכל קדשי קדשים, שנתן המקרא טעם לדבר משום דקדש הם (פסח' כ"ד):}%
% ---------- פרק 29 | פסוק 34 ----------
 \textbf{לד} וְֽאִם־יִוָּתֵ֞ר מִבְּשַׂ֧ר הַמִּלֻּאִ֛ים וּמִן־הַלֶּ֖חֶם עַד־הַבֹּ֑קֶר וְשָׂרַפְתָּ֤ אֶת־הַנּוֹתָר֙ בָּאֵ֔שׁ לֹ֥א יֵאָכֵ֖ל כִּי־קֹ֥דֶשׁ הֽוּא׃ %
% ---------- פרק 29 | פסוק 35 ----------
 \textbf{לה} וְעָשִׂ֜יתָ לְאַהֲרֹ֤ן וּלְבָנָיו֙ כָּ֔כָה כְּכֹ֥ל אֲשֶׁר־צִוִּ֖יתִי אֹתָ֑כָה שִׁבְעַ֥ת יָמִ֖ים תְּמַלֵּ֥א יָדָֽם׃ \Rashi{\textbf{ועשית לאהרן ולבניו ככה.}שנה הכתוב וכפל לעכב – שאם חסר דבר אחד מכל האמור בענין, לא נתמלאו ידיהם להיות כהנים ועבודתם פסולה: אתכה. כמו אותך: שבעת ימים תמלא וגו'. בענין הזה ובקרבנות הללו בכל יום:}%
% ---------- פרק 29 | פסוק 36 ----------
 \textbf{לו} וּפַ֨ר חַטָּ֜את תַּעֲשֶׂ֤ה לַיּוֹם֙ עַל־הַכִּפֻּרִ֔ים וְחִטֵּאתָ֙ עַל־הַמִּזְבֵּ֔חַ בְּכַפֶּרְךָ֖ עָלָ֑יו וּמָֽשַׁחְתָּ֥ אֹת֖וֹ לְקַדְּשֽׁוֹ׃ \Rashi{\textbf{על הכפרים.}בשביל הכפורים, לכפר על המזבח מכל זרות ותעוב; ולפי שנ' שבעת ימים תמלא ידם, אין לי אלא דבר הבא בשבילם, כגון האילים והלחם, אבל הבא בשביל המזבח, כגון פר שהוא לחטוי המזבח, לא שמענו, לכך הצרך מקרא זה; ומדרש ת"כ אומר כפרת המזבח הצרכה שמא התנדב איש דבר גזל במלאכת המשכן והמזבח: וחטאת. ותדכי, לשון מתנת דמים הנתונים באצבע קרוי חטוי: ומשחת אתו. בשמן המשחה, וכל המשיחות כמין כי:}%
% ---------- פרק 29 | פסוק 37 ----------
 \textbf{לז} שִׁבְעַ֣ת יָמִ֗ים תְּכַפֵּר֙ עַל־הַמִּזְבֵּ֔חַ וְקִדַּשְׁתָּ֖ אֹת֑וֹ וְהָיָ֤ה הַמִּזְבֵּ֙חַ֙ קֹ֣דֶשׁ קׇֽדָשִׁ֔ים כׇּל־הַנֹּגֵ֥עַ בַּמִּזְבֵּ֖חַ יִקְדָּֽשׁ׃ \{ס\} \Rashi{\textbf{והיה המזבח קדש.}ומה היא קדשתו? כל הנגע במזבח יקדש. אפלו קרבן פסול שעלה עליו קדשו המזבח להכשירו שלא ירד; מתוך שנאמר כל הנגע במזבח יקדש, שומע אני בין ראוי בין שאינו ראוי – כגון דבר שלא היה פסולו בקדש, כגון הרובע והנרבע ומקצה ונעבד והטרפה וכיוצא בהן – ת"ל וזה אשר תעשה הסמוך אחריו, מה עולה ראויה, אף כל ראוי שנראה לו כבר ונפסל משבא לעזרה, כגון הלן והיוצא והטמא ושנשחט במחשבת חוץ לזמנו וחוץ למקומו וכיוצא בהן (זבחים פ"ג):}%
% ---------- פרק 29 | פסוק 38 ----------
 \textbf{לח} וְזֶ֕ה אֲשֶׁ֥ר תַּעֲשֶׂ֖ה עַל־הַמִּזְבֵּ֑חַ כְּבָשִׂ֧ים בְּנֵֽי־שָׁנָ֛ה שְׁנַ֥יִם לַיּ֖וֹם תָּמִֽיד׃ %
% ---------- פרק 29 | פסוק 39 ----------
 \textbf{לט} אֶת־הַכֶּ֥בֶשׂ הָאֶחָ֖ד תַּעֲשֶׂ֣ה בַבֹּ֑קֶר וְאֵת֙ הַכֶּ֣בֶשׂ הַשֵּׁנִ֔י תַּעֲשֶׂ֖ה בֵּ֥ין הָעַרְבָּֽיִם׃ %
% ---------- פרק 29 | פסוק 40 ----------
 \textbf{מ} וְעִשָּׂרֹ֨ן סֹ֜לֶת בָּל֨וּל בְּשֶׁ֤מֶן כָּתִית֙ רֶ֣בַע הַהִ֔ין וְנֵ֕סֶךְ רְבִיעִ֥ת הַהִ֖ין יָ֑יִן לַכֶּ֖בֶשׂ הָאֶחָֽד׃ \Rashi{\textbf{ועשרן סלת.}עשירית האיפה – מ"ג ביצים וחמש ביצה: בשמן כתית. לא לחובה נאמר כתית, אלא להכשיר, לפי שנאמר כתית למאור – ומשמע למאור ולא למנחות – יכול לפסלו למנחות, ת"ל כאן כתית ולא נאמר כתית למאור, אלא למעט מנחות שאין צריך כתית, שאף הטחון בריחים כשר בהן (מנחות פ"ו): רבע ההין. שלשה לגין: ונסך. לספלים, כמו ששנינו במסכת סכה (מ"ח) – שני ספלים של כסף היו בראש המזבח ומנקבים כמין שני חטמין דקים, נותן היין לתוכו והוא מקלח ויוצא דרך החטם ונופל על גג המזבח, ומשם יורד לשיתין, במזבח בית עולמים, ובמזבח הנחשת יורד מן המזבח לארץ:}%
% ---------- פרק 29 | פסוק 41 ----------
 \textbf{מא} וְאֵת֙ הַכֶּ֣בֶשׂ הַשֵּׁנִ֔י תַּעֲשֶׂ֖ה בֵּ֣ין הָעַרְבָּ֑יִם כְּמִנְחַ֨ת הַבֹּ֤קֶר וּכְנִסְכָּהּ֙ תַּֽעֲשֶׂה־לָּ֔הּ לְרֵ֣יחַ נִיחֹ֔חַ אִשֶּׁ֖ה לַיהֹוָֽה׃ \Rashi{\textbf{לריח ניחח.}על המנחה נאמר, שמנחת נסכים כלה כליל, וסדר הקרבתם האברים בתחלה ואחר כך המנחה, שנאמר עלה ומנחה:}%
% ---------- פרק 29 | פסוק 42 ----------
 \textbf{מב} עֹלַ֤ת תָּמִיד֙ לְדֹרֹ֣תֵיכֶ֔ם פֶּ֥תַח אֹֽהֶל־מוֹעֵ֖ד לִפְנֵ֣י יְהֹוָ֑ה אֲשֶׁ֨ר אִוָּעֵ֤ד לָכֶם֙ שָׁ֔מָּה לְדַבֵּ֥ר אֵלֶ֖יךָ שָֽׁם׃ \Rashi{\textbf{תמיד.}מיום אל יום – ולא יפסיק יום בינתים: אשר אועד לכם. כשאקבע מועד לדבר אליך, שם אקבענו לבא. ויש מרבותינו למדים מכאן שמעל מזבח הנחשת היה הקב"ה מדבר עם משה משהוקם המשכן, ויש אומרים מעל הכפרת, כמו שנאמר ודברתי אתך מעל הכפרת (שמות כ"ו), ואשר אועד לכם האמור כאן אינו אמור על המזבח אלא על אהל מועד הנזכר במקרא:}%
% ---------- פרק 29 | פסוק 43 ----------
 \textbf{מג} וְנֹעַדְתִּ֥י שָׁ֖מָּה לִבְנֵ֣י יִשְׂרָאֵ֑ל וְנִקְדַּ֖שׁ בִּכְבֹדִֽי׃ \Rashi{\textbf{ונעדתי שמה.}אתועד עמם בדבור, כמלך הקובע מקום מועד לדבר עם עבדיו שם: ונקדש. המשכן: בכבדי. – שתשרה שכינתי בו. ומדרש אגדה אל תקרי בכבדי אלא בכבודי – במכבדים שלי – כאן רמז לו מיתת בני אהרן ביום הקמתו, וזהו שאמר משה הוא אשר דבר ה' לאמר בקרבי אקדש, והיכן דבר? ונקדש בכבדי (זבחים קט"ו):}%
% ---------- פרק 29 | פסוק 44 ----------
 \textbf{מד} וְקִדַּשְׁתִּ֛י אֶת־אֹ֥הֶל מוֹעֵ֖ד וְאֶת־הַמִּזְבֵּ֑חַ וְאֶת־אַהֲרֹ֧ן וְאֶת־בָּנָ֛יו אֲקַדֵּ֖שׁ לְכַהֵ֥ן לִֽי׃ %
% ---------- פרק 29 | פסוק 45 ----------
 \textbf{מה} וְשָׁ֣כַנְתִּ֔י בְּת֖וֹךְ בְּנֵ֣י יִשְׂרָאֵ֑ל וְהָיִ֥יתִי לָהֶ֖ם לֵאלֹהִֽים׃ %
% ---------- פרק 29 | פסוק 46 ----------
 \textbf{מו} וְיָדְע֗וּ כִּ֣י אֲנִ֤י יְהֹוָה֙ אֱלֹ֣הֵיהֶ֔ם אֲשֶׁ֨ר הוֹצֵ֧אתִי אֹתָ֛ם מֵאֶ֥רֶץ מִצְרַ֖יִם לְשׇׁכְנִ֣י בְתוֹכָ֑ם אֲנִ֖י יְהֹוָ֥ה אֱלֹהֵיהֶֽם׃ \{פ\} \Rashi{\textbf{לשכני בתוכם.}על מנת לשכן אני בתוכם:}%
% ---------- פרק 30 | פסוק 1 ----------
 \textbf{א} וְעָשִׂ֥יתָ מִזְבֵּ֖חַ מִקְטַ֣ר קְטֹ֑רֶת עֲצֵ֥י שִׁטִּ֖ים תַּעֲשֶׂ֥ה אֹתֽוֹ׃ \Rashi{\textbf{מקטר קטרת.}להעלות עליו קטור עשן סמים:}%
% ---------- פרק 30 | פסוק 2 ----------
 \textbf{ב} אַמָּ֨ה אׇרְכּ֜וֹ וְאַמָּ֤ה רׇחְבּוֹ֙ רָב֣וּעַ יִהְיֶ֔ה וְאַמָּתַ֖יִם קֹמָת֑וֹ מִמֶּ֖נּוּ קַרְנֹתָֽיו׃ %
% ---------- פרק 30 | פסוק 3 ----------
 \textbf{ג} וְצִפִּיתָ֨ אֹת֜וֹ זָהָ֣ב טָה֗וֹר אֶת־גַּגּ֧וֹ וְאֶת־קִירֹתָ֛יו סָבִ֖יב וְאֶת־קַרְנֹתָ֑יו וְעָשִׂ֥יתָ לּ֛וֹ זֵ֥ר זָהָ֖ב סָבִֽיב׃ \Rashi{\textbf{את גגו.}זה היה לו גג, אבל מזבח העולה לא היה לו גג אלא ממלאים חללו אדמה בכל חניתן: זר זהב. סימן הוא לכתר כהנה:}%
% ---------- פרק 30 | פסוק 4 ----------
 \textbf{ד} וּשְׁתֵּי֩ טַבְּעֹ֨ת זָהָ֜ב תַּֽעֲשֶׂה־לּ֣וֹ ׀ מִתַּ֣חַת לְזֵר֗וֹ עַ֚ל שְׁתֵּ֣י צַלְעֹתָ֔יו תַּעֲשֶׂ֖ה עַל־שְׁנֵ֣י צִדָּ֑יו וְהָיָה֙ לְבָתִּ֣ים לְבַדִּ֔ים לָשֵׂ֥את אֹת֖וֹ בָּהֵֽמָּה׃ \Rashi{\textbf{צלעתיו.}כאן הוא לשון זויות, כתרגומו, לפי שנאמר על שני צדיו – על שתי זויות שבשני צדיו: והיה. מעשה הטבעות האלה: לבתים לבדים. בית תהיה הטבעת לבד:}%
% ---------- פרק 30 | פסוק 5 ----------
 \textbf{ה} וְעָשִׂ֥יתָ אֶת־הַבַּדִּ֖ים עֲצֵ֣י שִׁטִּ֑ים וְצִפִּיתָ֥ אֹתָ֖ם זָהָֽב׃ %
% ---------- פרק 30 | פסוק 6 ----------
 \textbf{ו} וְנָתַתָּ֤ה אֹתוֹ֙ לִפְנֵ֣י הַפָּרֹ֔כֶת אֲשֶׁ֖ר עַל־אֲרֹ֣ן הָעֵדֻ֑ת לִפְנֵ֣י הַכַּפֹּ֗רֶת אֲשֶׁר֙ עַל־הָ֣עֵדֻ֔ת אֲשֶׁ֛ר אִוָּעֵ֥ד לְךָ֖ שָֽׁמָּה׃ \Rashi{\textbf{לפני הפרכת.}שמא תאמר משוך מכנגד הארון לצפון או לדרום, ת"ל לפני הכפרת מכון כנגד הארון מבחוץ:}%
% ---------- פרק 30 | פסוק 7 ----------
 \textbf{ז} וְהִקְטִ֥יר עָלָ֛יו אַהֲרֹ֖ן קְטֹ֣רֶת סַמִּ֑ים בַּבֹּ֣קֶר בַּבֹּ֗קֶר בְּהֵיטִיב֛וֹ אֶת־הַנֵּרֹ֖ת יַקְטִירֶֽנָּה׃ \Rashi{\textbf{בהיטיבו.}לשון נקוי הבזיכין של המנורה מדשן הפתילות שנשרפו בלילה, והיה מטיבן בכל בקר ובקר: הנרת. לו"שיש בלעז, וכן כל נרות האמורות במנורה, חוץ ממקום שנאמר שם העלאה שהוא לשון הדלקה:}%
% ---------- פרק 30 | פסוק 8 ----------
 \textbf{ח} וּבְהַעֲלֹ֨ת אַהֲרֹ֧ן אֶת־הַנֵּרֹ֛ת בֵּ֥ין הָעַרְבַּ֖יִם יַקְטִירֶ֑נָּה קְטֹ֧רֶת תָּמִ֛יד לִפְנֵ֥י יְהֹוָ֖ה לְדֹרֹתֵיכֶֽם׃ \Rashi{\textbf{ובהעלת.}כשידליקם להעלות להבתן: יקטירנה. בכל יום, פרס מקטיר שחרית ופרס מקטיר בין הערבים (כריתות ו'):}%
% ---------- פרק 30 | פסוק 9 ----------
 \textbf{ט} לֹא־תַעֲל֥וּ עָלָ֛יו קְטֹ֥רֶת זָרָ֖ה וְעֹלָ֣ה וּמִנְחָ֑ה וְנֵ֕סֶךְ לֹ֥א תִסְּכ֖וּ עָלָֽיו׃ \Rashi{\textbf{לא תעלו עליו.}על מזבח זה: קטרת זרה. שום קטרת של נדבה – כלן זרות לו חוץ מזו: ועלה ומנחה. ולא עולה ומנחה, עולה של בהמה ועוף, מנחה היא של לחם:}%
% ---------- פרק 30 | פסוק 10 ----------
 \textbf{י} וְכִפֶּ֤ר אַהֲרֹן֙ עַל־קַרְנֹתָ֔יו אַחַ֖ת בַּשָּׁנָ֑ה מִדַּ֞ם חַטַּ֣את הַכִּפֻּרִ֗ים אַחַ֤ת בַּשָּׁנָה֙ יְכַפֵּ֤ר עָלָיו֙ לְדֹרֹ֣תֵיכֶ֔ם קֹֽדֶשׁ־קׇדָשִׁ֥ים ה֖וּא לַיהֹוָֽה׃ \{פ\} \Rashi{\textbf{וכפר אהרן.}מתן דמים: אחת בשנה. ביום הכפורים, הוא שנאמר באחרי מות ויצא אל המזבח אשר לפני ה' וכפר עליו (ויקרא ט"ז): חטאת הכפרים. הם פר ושעיר של יום הכפורים המכפרים על טמאת מקדש וקדשיו: קדש קדשים. המזבח מקדש לדברים הללו בלבד ולא לעבודה אחרת:}%
\addparasha{חומש שמות | פרשת כי תשא}
% ---------- פרק 30 | פסוק 11 ----------
 \textbf{יא} וַיְדַבֵּ֥ר יְהֹוָ֖ה אֶל־מֹשֶׁ֥ה לֵּאמֹֽר׃ %
% ---------- פרק 30 | פסוק 12 ----------
 \textbf{יב} כִּ֣י תִשָּׂ֞א אֶת־רֹ֥אשׁ בְּנֵֽי־יִשְׂרָאֵל֮ לִפְקֻדֵיהֶם֒ וְנָ֨תְנ֜וּ אִ֣ישׁ כֹּ֧פֶר נַפְשׁ֛וֹ לַיהֹוָ֖ה בִּפְקֹ֣ד אֹתָ֑ם וְלֹא־יִהְיֶ֥ה בָהֶ֛ם נֶ֖גֶף בִּפְקֹ֥ד אֹתָֽם׃ \Rashi{\textbf{כי תשא.}לשון קבלה, כתרגומו; כשתחפץ לקבל סכום מנינם לדעת כמה הם, אל תמנם לגלגלת, אלא יתנו כל אחד מחצית השקל ותמנה את השקלים ותדע מנינם: ולא יהיה בהם נגף. שהמנין שולט בו עין הרע, והדבר בא עליהם, כמו שמצינו בימי דוד (שמואל ב כ"ד):}%
% ---------- פרק 30 | פסוק 13 ----------
 \textbf{יג} זֶ֣ה ׀ יִתְּנ֗וּ כׇּל־הָעֹבֵר֙ עַל־הַפְּקֻדִ֔ים מַחֲצִ֥ית הַשֶּׁ֖קֶל בְּשֶׁ֣קֶל הַקֹּ֑דֶשׁ עֶשְׂרִ֤ים גֵּרָה֙ הַשֶּׁ֔קֶל מַחֲצִ֣ית הַשֶּׁ֔קֶל תְּרוּמָ֖ה לַֽיהֹוָֽה׃ \Rashi{\textbf{זה יתנו.}הראה לו כמין מטבע של אש ומשקלה מחצית השקל ואומר לו כזה יתנו (תלמוד ירושלמי שק' א'): העבר על הפקדים. דרך המונין מעבירין את הנמנין זה אחר זה, וכן כל אשר יעבר תחת השבט (ויקרא כ"ז), וכן תעברנה הצאן על ידי מונה (ירמיהו ל"ג): מחצית השקל בשקל הקדש. במשקל השקל שקצבתי לך לשקל בו שקלי הקדש, כגון שקלים האמורין בפרשת ערכין ושדה אחזה: עשרים גרה השקל. עכשו פרש לך כמה הוא: גרה. לשון מעה; וכן בשמואל (א' ב'), יבוא להשתחות לו לאגורת כסף וככר לחם: עשרים גרה השקל. שהשקל השלם ארבעה זוזים, והזוז מתחלתו חמש מעות, אלא באו והוסיפו עליו שתות והעלוהו לשש מעה כסף, ומחצית השקל הזה שאמרתי לך יתנו תרומה לה' (בכורות ה'):}%
% ---------- פרק 30 | פסוק 14 ----------
 \textbf{יד} כֹּ֗ל הָעֹבֵר֙ עַל־הַפְּקֻדִ֔ים מִבֶּ֛ן עֶשְׂרִ֥ים שָׁנָ֖ה וָמָ֑עְלָה יִתֵּ֖ן תְּרוּמַ֥ת יְהֹוָֽה׃ \Rashi{\textbf{מבן עשרים שנה ומעלה.}למדך כאן שאין פחות מבן עשרים יוצא לצבא ונמנה בכלל אנשים:}%
% ---------- פרק 30 | פסוק 15 ----------
 \textbf{טו} הֶֽעָשִׁ֣יר לֹֽא־יַרְבֶּ֗ה וְהַדַּל֙ לֹ֣א יַמְעִ֔יט מִֽמַּחֲצִ֖ית הַשָּׁ֑קֶל לָתֵת֙ אֶת־תְּרוּמַ֣ת יְהֹוָ֔ה לְכַפֵּ֖ר עַל־נַפְשֹׁתֵיכֶֽם׃ \Rashi{\textbf{לכפר על נפשתיכם.}שלא תנגפו על ידי מנין; ד"א לכפר על נפשותיכם, לפי שרמז להם כאן ג' תרומות – שנכתב כאן תרומת ה' שלש פעמים, אחת תרומת אדנים שמנאן כשהתחילו בנדבת המשכן, שנתנו כל אחד ואחד מחצית השקל, ועלה למאת הככר, שנאמר וכסף פקודי העדה מאת ככר (שמות ל"ח), ומהם נעשו האדנים שנאמר ויהי מאת ככר הכסף וגו' (שם); והשנית אף היא על ידי מנין, שמנאן משהוקם המשכן, הוא המנין האמור בתחלת חמש הפקודים באחד לחדש השני בשנה השנית (במדבר א'), ונתנו כל אחד מחצית השקל, והן לקנות מהן קרבנות צבור של כל שנה ושנה, והשוו בהם עניים ועשירים, ועל אותה תרומה נאמר לכפר על נפשותיכם, שהקרבנות לכפרה הם באים; והשלישית היא תרומת המשכן כמו שנאמר כל מרים תרומת כסף ונחשת (שמות ל"ה), ולא היתה יד כלם שוה בה אלא איש מה שנדבו לבו:}%
% ---------- פרק 30 | פסוק 16 ----------
 \textbf{טז} וְלָקַחְתָּ֞ אֶת־כֶּ֣סֶף הַכִּפֻּרִ֗ים מֵאֵת֙ בְּנֵ֣י יִשְׂרָאֵ֔ל וְנָתַתָּ֣ אֹת֔וֹ עַל־עֲבֹדַ֖ת אֹ֣הֶל מוֹעֵ֑ד וְהָיָה֩ לִבְנֵ֨י יִשְׂרָאֵ֤ל לְזִכָּרוֹן֙ לִפְנֵ֣י יְהֹוָ֔ה לְכַפֵּ֖ר עַל־נַפְשֹׁתֵיכֶֽם׃ \{פ\} \Rashi{\textbf{ונתת אתו על עבדת אהל מועד.}למדת שנצטוו למנותם בתחלת נדבת המשכן, אחר מעשה העגל, מפני שנכנס בהם מגפה, כמו שנאמר ויגף ה' את העם; משל לצאן החביבה על בעליה שנפל בה דבר ומשפסק אמר לו לרועה בבקשה ממך מנה את צאני ודע כמה נותרו בהם, להודיע שהיא חביבה עליו; ואי אפשר לומר שהמנין הזה הוא האמור בחמש הפקודים, שהרי נאמר בו באחד לחדש השני והמשכן הוקם באחד לחדש הראשון, שנאמר ביום החדש הראשון באחד לחדש תקים וגו' (שמות מ'), ומהמנין הזה נעשו האדנים – משקלים שלו – שנאמר ויהי מאת ככר הכסף לצקת וגו' (שם ל"ח), הא למדת שתים היו, אחד בתחלת נדבתן אחר יום הכפורים בשנה ראשונה, ואחת בשנה שניה באיר משהוקם המשכן. ואם תאמר וכי אפשר שבשניהם היו ישראל שוים – ו' מאות אלף וג' אלפים וחמש מאות וחמשים – שהרי בכסף פקודי העדה נאמר כן, ובחמש הפקודים אף בו נאמר כן ויהיו כל הפקודים שש מאות אלף ושלשת אלפים וחמש מאות וחמשים (במדבר א'), והלא בשתי שנים היו, וא"א שלא היו בשעת מנין הראשון בני י"ט שנה שלא נמנו ובשניה נעשו בני כ'! תשובה לדבר: אצל שנות האנשים בשנה אחת נמנו, אבל למנין יציאת מצרים היו שתי שנים, לפי שליציאת מצרים מונין מניסן, כמו ששנינו במס' ראש השנה (דף ב'), ונבנה המשכן בראשונה והוקם בשניה, שנתחדשה שנה באחד בניסן, אבל שנות האנשים מנויין למנין שנות עולם, המתחילין מתשרי, נמצאו שני המנינים בשנה אחת – המנין הראשון היה בתשרי לאחר י"הכ שנתרצה המקום לישראל לסלח להם ונצטוו על המשכן, והשני באחד באיר: על עבדת אהל מועד. הן האדנים שנעשו בו:}%
% ---------- פרק 30 | פסוק 17 ----------
 \textbf{יז} וַיְדַבֵּ֥ר יְהֹוָ֖ה אֶל־מֹשֶׁ֥ה לֵּאמֹֽר׃ %
% ---------- פרק 30 | פסוק 18 ----------
 \textbf{יח} וְעָשִׂ֜יתָ כִּיּ֥וֹר נְחֹ֛שֶׁת וְכַנּ֥וֹ נְחֹ֖שֶׁת לְרׇחְצָ֑ה וְנָתַתָּ֣ אֹת֗וֹ בֵּֽין־אֹ֤הֶל מוֹעֵד֙ וּבֵ֣ין הַמִּזְבֵּ֔חַ וְנָתַתָּ֥ שָׁ֖מָּה מָֽיִם׃ \Rashi{\textbf{כיור.}כמין דוד גדולה ולה דדים המריקים בפיהם מים: וכנו. כתרגומו בסיסה, מושב מתקן לכיור: לרחצה. מוסב על הכיור: ובין המזבח. מזבח העולה, שכתוב בו שהוא לפני פתח משכן אהל מועד, והיה הכיור משוך קמעא, ועומד כנגד אויר שבין המזבח והמשכן, ואינו מפסיק כלל בינתים, משום שנאמר ואת מזבח העלה שם פתח משכן אהל מועד (שמות מ'), – כלומר מזבח לפני אהל מועד ואין כיור לפני אהל מועד, הא כיצד? משוך קמעא כלפי הדרום; כך שנויה בזבחים:}%
% ---------- פרק 30 | פסוק 19 ----------
 \textbf{יט} וְרָחֲצ֛וּ אַהֲרֹ֥ן וּבָנָ֖יו מִמֶּ֑נּוּ אֶת־יְדֵיהֶ֖ם וְאֶת־רַגְלֵיהֶֽם׃ \Rashi{\textbf{את ידיהם ואת רגליהם.}בבת אחת היה מקדש ידיו ורגליו; וכך שנינו בזבחים (דף י"ט); כיצד קדוש ידים ורגלים? מניח ידו הימנית על גבי רגלו הימנית וידו השמאלית על גבי רגלו השמאלית ומקדש:}%
% ---------- פרק 30 | פסוק 20 ----------
 \textbf{כ} בְּבֹאָ֞ם אֶל־אֹ֧הֶל מוֹעֵ֛ד יִרְחֲצוּ־מַ֖יִם וְלֹ֣א יָמֻ֑תוּ א֣וֹ בְגִשְׁתָּ֤ם אֶל־הַמִּזְבֵּ֙חַ֙ לְשָׁרֵ֔ת לְהַקְטִ֥יר אִשֶּׁ֖ה לַֽיהֹוָֽה׃ \Rashi{\textbf{בבאם אל אהל מועד.}להקטיר שחרית ובין הערבים קטרת, או להזות מדם פר כהן המשיח ושעירי ע"ז: ולא ימתו. הא אם לא ירחצו ימותו, שבתורה נאמרו כללות, ומכלל לאו אתה שומע הן: אל המזבח. החיצון, שאין כאן ביאת אהל מועד אלא בחצר:}%
% ---------- פרק 30 | פסוק 21 ----------
 \textbf{כא} וְרָחֲצ֛וּ יְדֵיהֶ֥ם וְרַגְלֵיהֶ֖ם וְלֹ֣א יָמֻ֑תוּ וְהָיְתָ֨ה לָהֶ֧ם חׇק־עוֹלָ֛ם ל֥וֹ וּלְזַרְע֖וֹ לְדֹרֹתָֽם׃ \{פ\} \Rashi{\textbf{ולא ימתו.}לחיב מיתה על המשמש במזבח ואינו רחוץ ידים ורגלים, שהמיתה הראשונה לא שמענו אלא על הנכנס להיכל:}%
% ---------- פרק 30 | פסוק 22 ----------
 \textbf{כב} וַיְדַבֵּ֥ר יְהֹוָ֖ה אֶל־מֹשֶׁ֥ה לֵּאמֹֽר׃ %
% ---------- פרק 30 | פסוק 23 ----------
 \textbf{כג} וְאַתָּ֣ה קַח־לְךָ֮ בְּשָׂמִ֣ים רֹאשׁ֒ מׇר־דְּרוֹר֙ חֲמֵ֣שׁ מֵא֔וֹת וְקִנְּמׇן־בֶּ֥שֶׂם מַחֲצִית֖וֹ חֲמִשִּׁ֣ים וּמָאתָ֑יִם וּקְנֵה־בֹ֖שֶׂם חֲמִשִּׁ֥ים וּמָאתָֽיִם׃ \Rashi{\textbf{בשמים ראש.}חשובים: וקנמן בשם. לפי שהקנמון קלפת עץ הוא, יש שהוא טוב ויש בו ריח טוב וטעם, ויש שאינו אלא כעץ, לכך הצרך לומר קנמן בשם – מן הטוב: מחציתו חמשים ומאתים. מחצית הבאתו תהא חמשים ומאתים, נמצא כלו חמש מאות, כמו שעור מר דרור, אם כן למה נאמר בו חצאין, גזרת הכתוב היא להביאו לחצאין, להרבות בו ב' הכרעות, שאין שוקלין עין בעין; וכך שנויה בכרתות: וקנה בשם. קנה של בשם; לפי שיש קנים שאינן של בשם הצרך לומר קנה בשם: חמשים ומאתים. סך משקל כלו:}%
% ---------- פרק 30 | פסוק 24 ----------
 \textbf{כד} וְקִדָּ֕ה חֲמֵ֥שׁ מֵא֖וֹת בְּשֶׁ֣קֶל הַקֹּ֑דֶשׁ וְשֶׁ֥מֶן זַ֖יִת הִֽין׃ \Rashi{\textbf{וקדה.}שם שרש עשב, ובלשון חכמים קציעה: הין. י"ב לגין; ונחלקו בו חכמי ישראל — רבי מאיר אומר בו שלקו את העקרין, אמר לו ר' יהודה והלא לסוך את העקרין אינו ספק אלא שראום במים שלא יבלעו את השמן, ואחר כך הציף עליהם השמן עד שקלט הריח, וקפחו לשמן מעל העקרין (הוריות י"א):}%
% ---------- פרק 30 | פסוק 25 ----------
 \textbf{כה} וְעָשִׂ֣יתָ אֹת֗וֹ שֶׁ֚מֶן מִשְׁחַת־קֹ֔דֶשׁ רֹ֥קַח מִרְקַ֖חַת מַעֲשֵׂ֣ה רֹקֵ֑חַ שֶׁ֥מֶן מִשְׁחַת־קֹ֖דֶשׁ יִהְיֶֽה׃ \Rashi{\textbf{רקח מרקחת.}רקח שם דבר הוא, והטעם מוכיח, שהוא למעלה, והרי הוא כמו רקח, רגע; ואינו כמו רגע הים (ישעיהו נ"א), וכמו רקע הארץ (שם מלכים ב), שהטעם למטה; וכל דבר המערב בחברו עד שזה קופח מזה או ריח או טעם קרוי מרקחת: רקח מרקחת. רקח העשוי ע"י אמנות ותערובות: מעשה רקח. שם האמן בדבר:}%
% ---------- פרק 30 | פסוק 26 ----------
 \textbf{כו} וּמָשַׁחְתָּ֥ ב֖וֹ אֶת־אֹ֣הֶל מוֹעֵ֑ד וְאֵ֖ת אֲר֥וֹן הָעֵדֻֽת׃ \Rashi{\textbf{ומשחת בו.}כל המשיחות כמין כי, חוץ משל מלכים שהן כמין נזר (כריתות ה'):}%
% ---------- פרק 30 | פסוק 27 ----------
 \textbf{כז} וְאֶת־הַשֻּׁלְחָן֙ וְאֶת־כׇּל־כֵּלָ֔יו וְאֶת־הַמְּנֹרָ֖ה וְאֶת־כֵּלֶ֑יהָ וְאֵ֖ת מִזְבַּ֥ח הַקְּטֹֽרֶת׃ %
% ---------- פרק 30 | פסוק 28 ----------
 \textbf{כח} וְאֶת־מִזְבַּ֥ח הָעֹלָ֖ה וְאֶת־כׇּל־כֵּלָ֑יו וְאֶת־הַכִּיֹּ֖ר וְאֶת־כַּנּֽוֹ׃ %
% ---------- פרק 30 | פסוק 29 ----------
 \textbf{כט} וְקִדַּשְׁתָּ֣ אֹתָ֔ם וְהָי֖וּ קֹ֣דֶשׁ קׇֽדָשִׁ֑ים כׇּל־הַנֹּגֵ֥עַ בָּהֶ֖ם יִקְדָּֽשׁ׃ \Rashi{\textbf{וקדשת אתם.}משיחה זו מקדשתם להיות קדש קדשים, ומה היא קדשתם? כל הנגע וגו': כל הראוי לכלי שרת, משנכנס לתוכו, קדוש קדשת הגוף – לפסל ביוצא ולינה וטבול יום ואינו נפדה לצאת לחלין, אבל דבר שאינו ראוי להם אין מקדשין; ושנויה היא משנה שלמה אצל מזבח: מתוך שנ' כל הנגע במזבח יקדש (שמות כ"ט), שומע אני בין ראוי בין שאינו ראוי, ת"ל כבשים, מה כבשים ראויים אף כל ראוי; כל משיחת משכן וכהנים ומלכים מתרגם לשון רבוי, לפי שאין צרך משיחתן אלא לגדלה, כי כן יסד המלך שזה חנוך גדלתן, ושאר משיחות – כמו רקיקין משוחין (שם), וראשית שמנים ימשחו (עמוס ו') – לשון ארמית בהן כלשון עברית:}%
% ---------- פרק 30 | פסוק 30 ----------
 \textbf{ל} וְאֶת־אַהֲרֹ֥ן וְאֶת־בָּנָ֖יו תִּמְשָׁ֑ח וְקִדַּשְׁתָּ֥ אֹתָ֖ם לְכַהֵ֥ן לִֽי׃ %
% ---------- פרק 30 | פסוק 31 ----------
 \textbf{לא} וְאֶל־בְּנֵ֥י יִשְׂרָאֵ֖ל תְּדַבֵּ֣ר לֵאמֹ֑ר שֶׁ֠מֶן מִשְׁחַת־קֹ֨דֶשׁ יִהְיֶ֥ה זֶ֛ה לִ֖י לְדֹרֹתֵיכֶֽם׃ \Rashi{\textbf{לדרתיכם.}מכאן למדו רבותינו לומר שכלו קים לעתיד לבא: זה. בגימטריא תריסר לגין הוו (כריתות ה'):}%
% ---------- פרק 30 | פסוק 32 ----------
 \textbf{לב} עַל־בְּשַׂ֤ר אָדָם֙ לֹ֣א יִיסָ֔ךְ וּ֨בְמַתְכֻּנְתּ֔וֹ לֹ֥א תַעֲשׂ֖וּ כָּמֹ֑הוּ קֹ֣דֶשׁ ה֔וּא קֹ֖דֶשׁ יִהְיֶ֥ה לָכֶֽם׃ \Rashi{\textbf{לא ייסך.}בשני יודי"ן, לשון לא יפעל, כמו (דברים ה טו) למען ייטב לך: על בשר אדם לא ייסך. מן השמן הזה עצמו: ובמתכנתו לא תעשו כמהו. בסכום סמניו לא תעשו אחר כמוהו במשקל סמנין הללו לפי מדת הין שמן, אבל אם פחת או רבה סממנין לפי מדת הין שמן מתר, ואף העשוי במתכנתו של זה, אין הסך ממנו חיב אלא הרוקחו: ובמתכנתו. לשון חשבון, כמו (שמות ה ח) מתכנת הלבנים, וכן במתכנתה של קטרת:}%
% ---------- פרק 30 | פסוק 33 ----------
 \textbf{לג} אִ֚ישׁ אֲשֶׁ֣ר יִרְקַ֣ח כָּמֹ֔הוּ וַאֲשֶׁ֥ר יִתֵּ֛ן מִמֶּ֖נּוּ עַל־זָ֑ר וְנִכְרַ֖ת מֵעַמָּֽיו׃ \{ס\} \Rashi{\textbf{ואשר יתן ממנו.}מאותו של משה: על זר. שאינו צרך כהנה ומלכות:}%
% ---------- פרק 30 | פסוק 34 ----------
 \textbf{לד} וַיֹּ֩אמֶר֩ יְהֹוָ֨ה אֶל־מֹשֶׁ֜ה קַח־לְךָ֣ סַמִּ֗ים נָטָ֤ף ׀ וּשְׁחֵ֙לֶת֙ וְחֶלְבְּנָ֔ה סַמִּ֖ים וּלְבֹנָ֣ה זַכָּ֑ה בַּ֥ד בְּבַ֖ד יִהְיֶֽה׃ \Rashi{\textbf{נטף.}הוא צרי, ועל שאינו אלא שרף הנוטף מעצי הקטף קרוי נטף, ובלעז גומ"א, והצרי קורין לו טרי"אקה: ושחלת. שרש בשם חלק ומצהיר כצפרן, ובלשון המשנה קרוי צפרן, וזהו שתרגם אנקלוס וטופרא: וחלבנה. בשם שריחו רע, וקורין לו גלבנא, ומנאה הכתוב בין סמני הקטרת ללמדנו שלא יקל בעינינו לצרף עמנו באגדת תעניותינו ותפלותנו את פושעי ישראל שיהיו נמנין עמנו: סמים. אחרים: ולבנה זכה. מכאן למדו רבותינו י"א סמנין נאמרו לו למשה בסיני – מעוט סמים שנים, נטף ושחלת וחלבנה ג', הרי חמשה, סמים לרבות עוד כמו אלו הרי עשרה, ולבונה הרי י"א; ואלו הן הצרי והצפרן החלבנה והלבונה מר וקציעה שבלת נרד וכרכם הרי ח' – שהשבלת ונרד אחד, שהנרד דומה לשבלת – הקשט והקלופה וקנמון הרי י"א; ברית כרשינה אינו נקטר אלא בו שפין את הצפרן ללבנה שתהא נאה (כריתות ו'): בד בבד יהיה. אלו הארבעה הנזכרים כאן יהיו שוין משקל במשקל – כמשקלו של זה כך משקלו של זה – וכן שנינו: הצרי והצפרן החלבנה והלבונה משקל שבעים שבעים מנה (שם); ולשון בד נראה בעיני שהוא לשון יחיד, אחד באחד יהיו – זה כמות זה:}%
% ---------- פרק 30 | פסוק 35 ----------
 \textbf{לה} וְעָשִׂ֤יתָ אֹתָהּ֙ קְטֹ֔רֶת רֹ֖קַח מַעֲשֵׂ֣ה רוֹקֵ֑חַ מְמֻלָּ֖ח טָה֥וֹר קֹֽדֶשׁ׃ \Rashi{\textbf{ממלח.}כתרגומו מערב, שיערב שחיקתן יפה יפה זה עם זה; ואומר אני שדומה לו וייראו המלחים (יונה א'), מלחיך וחבליך (יחזקאל כ"ז) – על שם שמהפכין את המים במשוטות כשמנהיגים את הספינה, כאדם המהפך בכף ביצים טרופות לערבן עם המים, וכל דבר שאדם רוצה לערב יפה יפה מהפכו באצבע או בבזך: ממלח טהור קדש. ממלח יהיה, וטהור יהיה וקדש יהיה:}%
% ---------- פרק 30 | פסוק 36 ----------
 \textbf{לו} וְשָֽׁחַקְתָּ֣ מִמֶּ֘נָּה֮ הָדֵק֒ וְנָתַתָּ֨ה מִמֶּ֜נָּה לִפְנֵ֤י הָעֵדֻת֙ בְּאֹ֣הֶל מוֹעֵ֔ד אֲשֶׁ֛ר אִוָּעֵ֥ד לְךָ֖ שָׁ֑מָּה קֹ֥דֶשׁ קׇֽדָשִׁ֖ים תִּהְיֶ֥ה לָכֶֽם׃ \Rashi{\textbf{ונתתה ממנה וגו'.}היא קטרת שבכל יום ויום שעל מזבח הפנימי שהוא באהל מועד: אשר אועד לך שמה. כל מועדי דבור שאקבע לך, אני קובעם לאותו מקום:}%
% ---------- פרק 30 | פסוק 37 ----------
 \textbf{לז} וְהַקְּטֹ֙רֶת֙ אֲשֶׁ֣ר תַּעֲשֶׂ֔ה בְּמַ֨תְכֻּנְתָּ֔הּ לֹ֥א תַעֲשׂ֖וּ לָכֶ֑ם קֹ֛דֶשׁ תִּהְיֶ֥ה לְךָ֖ לַיהֹוָֽה׃ \Rashi{\textbf{במתכנתה.}במנין סממניה: קדש תהיה לך לה'. שלא תעשנה אלא לשמי:}%
% ---------- פרק 30 | פסוק 38 ----------
 \textbf{לח} אִ֛ישׁ אֲשֶׁר־יַעֲשֶׂ֥ה כָמ֖וֹהָ לְהָרִ֣יחַ בָּ֑הּ וְנִכְרַ֖ת מֵעַמָּֽיו׃ \{ס\} \Rashi{\textbf{להריח בה.}אבל עושה אתה במתכנתה משלך, כדי למכרה לצבור (כריתות ה'):}%
% ---------- פרק 31 | פסוק 1 ----------
 \textbf{א} וַיְדַבֵּ֥ר יְהֹוָ֖ה אֶל־מֹשֶׁ֥ה לֵּאמֹֽר׃ %
% ---------- פרק 31 | פסוק 2 ----------
 \textbf{ב} רְאֵ֖ה קָרָ֣אתִֽי בְשֵׁ֑ם בְּצַלְאֵ֛ל בֶּן־אוּרִ֥י בֶן־ח֖וּר לְמַטֵּ֥ה יְהוּדָֽה׃ \Rashi{\textbf{קראתי בשם –.}לעשות מלאכתי – את בצלאל:}%
% ---------- פרק 31 | פסוק 3 ----------
 \textbf{ג} וָאֲמַלֵּ֥א אֹת֖וֹ ר֣וּחַ אֱלֹהִ֑ים בְּחׇכְמָ֛ה וּבִתְבוּנָ֥ה וּבְדַ֖עַת וּבְכׇל־מְלָאכָֽה׃ \Rashi{\textbf{חכמה.}מה שאדם שומע מאחרים ולמד: תבונה. מבין דבר מלבו מתוך דברים שלמד: דעת. רוח הקדש:}%
% ---------- פרק 31 | פסוק 4 ----------
 \textbf{ד} לַחְשֹׁ֖ב מַחֲשָׁבֹ֑ת לַעֲשׂ֛וֹת בַּזָּהָ֥ב וּבַכֶּ֖סֶף וּבַנְּחֹֽשֶׁת׃ \Rashi{\textbf{לחשב מחשבת.}אריגת מעשה חשב:}%
% ---------- פרק 31 | פסוק 5 ----------
 \textbf{ה} וּבַחֲרֹ֥שֶׁת אֶ֛בֶן לְמַלֹּ֖את וּבַחֲרֹ֣שֶׁת עֵ֑ץ לַעֲשׂ֖וֹת בְּכׇל־מְלָאכָֽה׃ \Rashi{\textbf{ובחרשת.}לשון אמנות, כמו חרש חכם (ישעיהו מ'), ואנקלוס פרש, ושנה בפרושן, שחרש אבנים קרוי אמן וחרש עץ קרוי נגר: למלאת. להושיבה במשבצת שלה במלואה – לעשות המשבצת למדת מושב האבן ועביה:}%
% ---------- פרק 31 | פסוק 6 ----------
 \textbf{ו} וַאֲנִ֞י הִנֵּ֧ה נָתַ֣תִּי אִתּ֗וֹ אֵ֣ת אׇהֳלִיאָ֞ב בֶּן־אֲחִֽיסָמָךְ֙ לְמַטֵּה־דָ֔ן וּבְלֵ֥ב כׇּל־חֲכַם־לֵ֖ב נָתַ֣תִּי חׇכְמָ֑ה וְעָשׂ֕וּ אֵ֖ת כׇּל־אֲשֶׁ֥ר צִוִּיתִֽךָ׃ \Rashi{\textbf{ובלב כל חכם לב וגו'.}ועוד שאר חכמי לב יש בכם, וכל אשר נתתי בו חכמה ועשו את כל אשר צויתיך:}%
% ---------- פרק 31 | פסוק 7 ----------
 \textbf{ז} אֵ֣ת ׀ אֹ֣הֶל מוֹעֵ֗ד וְאֶת־הָֽאָרֹן֙ לָֽעֵדֻ֔ת וְאֶת־הַכַּפֹּ֖רֶת אֲשֶׁ֣ר עָלָ֑יו וְאֵ֖ת כׇּל־כְּלֵ֥י הָאֹֽהֶל׃ \Rashi{\textbf{ואת הארן לעדת.}לצרך לחות העדות:}%
% ---------- פרק 31 | פסוק 8 ----------
 \textbf{ח} וְאֶת־הַשֻּׁלְחָן֙ וְאֶת־כֵּלָ֔יו וְאֶת־הַמְּנֹרָ֥ה הַטְּהֹרָ֖ה וְאֶת־כׇּל־כֵּלֶ֑יהָ וְאֵ֖ת מִזְבַּ֥ח הַקְּטֹֽרֶת׃ \Rashi{\textbf{הטהרה.}על שם זהב טהור:}%
% ---------- פרק 31 | פסוק 9 ----------
 \textbf{ט} וְאֶת־מִזְבַּ֥ח הָעֹלָ֖ה וְאֶת־כׇּל־כֵּלָ֑יו וְאֶת־הַכִּיּ֖וֹר וְאֶת־כַּנּֽוֹ׃ %
% ---------- פרק 31 | פסוק 10 ----------
 \textbf{י} וְאֵ֖ת בִּגְדֵ֣י הַשְּׂרָ֑ד וְאֶת־בִּגְדֵ֤י הַקֹּ֙דֶשׁ֙ לְאַהֲרֹ֣ן הַכֹּהֵ֔ן וְאֶת־בִּגְדֵ֥י בָנָ֖יו לְכַהֵֽן׃ \Rashi{\textbf{ואת בגדי השרד.}אומר אני – לפי פשוטו של מקרא – שאי אפשר לומר שבבגדי כהנה מדבר, לפי שנ' אצלם ואת בגדי הקדש לאהרן הכהן ואת בגדי בניו לכהן, אלא אלו בגדי השרד הם בגדי התכלת והארגמן ותולעת שני האמורין בפרשת מסעות, ונתנו אל בגד תכלת ופרשו עליו בגד ארגמן ופרשו עליהם בגד תולעת שני (במדבר ד'); ונראין דברי, שנ' ומן התכלת והארגמן ותולעת השני עשו בגדי שרד לשרת בקדש (שמות ל"ט), ולא הזכר שש עמהם, ואם בבגדי כהנה מדבר, לא מצינו באחד מהם ארגמן או תולעת שני בלא שש: בגדי השרד. יש מפרשים לשון עבודה ושרות, כתרגומו לבושי שמושא, ואין לו דמיון במקרא; ואני אומר שהוא לשון ארמי, כתרגום של קלעים ותרגום של מכבר, שהיו ארוגים במחט עשויים נקבים נקבים לצי"דץ בלעז:}%
% ---------- פרק 31 | פסוק 11 ----------
 \textbf{יא} וְאֵ֨ת שֶׁ֧מֶן הַמִּשְׁחָ֛ה וְאֶת־קְטֹ֥רֶת הַסַּמִּ֖ים לַקֹּ֑דֶשׁ כְּכֹ֥ל אֲשֶׁר־צִוִּיתִ֖ךָ יַעֲשֽׂוּ׃ \{פ\} \Rashi{\textbf{ואת קטרת הסמים לקדש.}לצרך הקטרת ההיכל שהוא קדש:}%
% ---------- פרק 31 | פסוק 12 ----------
 \textbf{יב} וַיֹּ֥אמֶר יְהֹוָ֖ה אֶל־מֹשֶׁ֥ה לֵּאמֹֽר׃ %
% ---------- פרק 31 | פסוק 13 ----------
 \textbf{יג} וְאַתָּ֞ה דַּבֵּ֨ר אֶל־בְּנֵ֤י יִשְׂרָאֵל֙ לֵאמֹ֔ר אַ֥ךְ אֶת־שַׁבְּתֹתַ֖י תִּשְׁמֹ֑רוּ כִּי֩ א֨וֹת הִ֜וא בֵּינִ֤י וּבֵֽינֵיכֶם֙ לְדֹרֹ֣תֵיכֶ֔ם לָדַ֕עַת כִּ֛י אֲנִ֥י יְהֹוָ֖ה מְקַדִּשְׁכֶֽם׃ \Rashi{\textbf{ואתה דבר אל בני ישראל.}ואתה אע"פ שהפקדתיך לצוותם על מלאכת המשכן, אל יקל בעיניך לדחות את השבת מפני אותה מלאכה: אך את שבתתי תשמרו. אע"פ שתהיו רדופין וזריזין בזריזות המלאכה, שבת אל תדחה מפניה, כל אכין ורקין מעוטין – למעט שבת ממלאכת המשכן: כי אות הוא ביני וביניכם. אות גדלה היא בינינו, שבחרתי בכם בהנחילי לכם את יום מנוחתי למנוחה: לדעת. האמות בה כי אני ה' מקדשכם:}%
% ---------- פרק 31 | פסוק 14 ----------
 \textbf{יד} וּשְׁמַרְתֶּם֙ אֶת־הַשַּׁבָּ֔ת כִּ֛י קֹ֥דֶשׁ הִ֖וא לָכֶ֑ם מְחַֽלְלֶ֙יהָ֙ מ֣וֹת יוּמָ֔ת כִּ֗י כׇּל־הָעֹשֶׂ֥ה בָהּ֙ מְלָאכָ֔ה וְנִכְרְתָ֛ה הַנֶּ֥פֶשׁ הַהִ֖וא מִקֶּ֥רֶב עַמֶּֽיהָ׃ \Rashi{\textbf{מות יומת.}אם יש עדים והתראה: ונכרתה. בלא התראה (מכילתא): מחלליה. הנוהג בה חל בקדשתה:}%
% ---------- פרק 31 | פסוק 15 ----------
 \textbf{טו} שֵׁ֣שֶׁת יָמִים֮ יֵעָשֶׂ֣ה מְלָאכָה֒ וּבַיּ֣וֹם הַשְּׁבִיעִ֗י שַׁבַּ֧ת שַׁבָּת֛וֹן קֹ֖דֶשׁ לַיהֹוָ֑ה כׇּל־הָעֹשֶׂ֧ה מְלָאכָ֛ה בְּי֥וֹם הַשַּׁבָּ֖ת מ֥וֹת יוּמָֽת׃ \Rashi{\textbf{שבת שבתון.}מנוחת מרגוע ולא מנוחת עראי: קדש לה'. שמירת קדשתה לשמי ובמצותי:}%
% ---------- פרק 31 | פסוק 16 ----------
 \textbf{טז} וְשָׁמְר֥וּ בְנֵֽי־יִשְׂרָאֵ֖ל אֶת־הַשַּׁבָּ֑ת לַעֲשׂ֧וֹת אֶת־הַשַּׁבָּ֛ת לְדֹרֹתָ֖ם בְּרִ֥ית עוֹלָֽם׃ %
% ---------- פרק 31 | פסוק 17 ----------
 \textbf{יז} בֵּינִ֗י וּבֵין֙ בְּנֵ֣י יִשְׂרָאֵ֔ל א֥וֹת הִ֖וא לְעֹלָ֑ם כִּי־שֵׁ֣שֶׁת יָמִ֗ים עָשָׂ֤ה יְהֹוָה֙ אֶת־הַשָּׁמַ֣יִם וְאֶת־הָאָ֔רֶץ וּבַיּוֹם֙ הַשְּׁבִיעִ֔י שָׁבַ֖ת וַיִּנָּפַֽשׁ׃ \{ס\} \Rashi{\textbf{וינפש.}כתרגומו ונח, וכל לשון נפש הוא לשון נפש, שמשיב נפשו ונשימתו בהרגיעו מטרח המלאכה; ומי שכתוב בו לא ייעף ולא ייגע (ישעיה מ'), וכל פעלו במאמר, הכתיב מנוחה בעצמו? לשבר האזן מה שהיא יכולה לשמע:}%
% ---------- פרק 31 | פסוק 18 ----------
 \textbf{יח} וַיִּתֵּ֣ן אֶל־מֹשֶׁ֗ה כְּכַלֹּתוֹ֙ לְדַבֵּ֤ר אִתּוֹ֙ בְּהַ֣ר סִינַ֔י שְׁנֵ֖י לֻחֹ֣ת הָעֵדֻ֑ת לֻחֹ֣ת אֶ֔בֶן כְּתֻבִ֖ים בְּאֶצְבַּ֥ע אֱלֹהִֽים׃ \Rashi{\textbf{ויתן אל משה וגו'.}אין מקדם ומאחר בתורה – מעשה העגל קדם לצווי מלאכת המשכן ימים רבים היה – שהרי בי"ז בתמוז נשתברו הלוחות, וביום הכפורים נתרצה הקב"ה לישראל, ולמחרת התחילו בנדבת המשכן והוקם באחד בניסן (תנחומא): ככלתו. ככלתו כתיב – חסר – שנמסרה לו תורה במתנה ככלה לחתן, שלא היה יכול ללמד כלה בזמן מעט כזה. דבר אחר, מה כלה מתקשטת בכ"ד קשוטין – הן האמורים בספר ישעיה – אף תלמיד חכם צריך להיות בקי בכ"ד ספרים: לדבר אתו. החקים והמשפטים שבואלה המשפטים: לדבר אתו. מלמד שהיה משה שומע מפי הגבורה וחוזרין ושונין את ההלכה שניהם יחד: לחת. לחת כתיב שהיו שתיהן שוות (שם):}%
% ---------- פרק 32 | פסוק 1 ----------
 \textbf{א} וַיַּ֣רְא הָעָ֔ם כִּֽי־בֹשֵׁ֥שׁ מֹשֶׁ֖ה לָרֶ֣דֶת מִן־הָהָ֑ר וַיִּקָּהֵ֨ל הָעָ֜ם עַֽל־אַהֲרֹ֗ן וַיֹּאמְר֤וּ אֵלָיו֙ ק֣וּם ׀ עֲשֵׂה־לָ֣נוּ אֱלֹהִ֗ים אֲשֶׁ֤ר יֵֽלְכוּ֙ לְפָנֵ֔ינוּ כִּי־זֶ֣ה ׀ מֹשֶׁ֣ה הָאִ֗ישׁ אֲשֶׁ֤ר הֶֽעֱלָ֙נוּ֙ מֵאֶ֣רֶץ מִצְרַ֔יִם לֹ֥א יָדַ֖עְנוּ מֶה־הָ֥יָה לֽוֹ׃ \Rashi{\textbf{כי בשש משה.}כתרגומו, לשון אחור, וכן בשש רכבו (שופטים ה'), ויחילו עד בוש (שם ג'); כי כשעלה משה להר אמר להם לסוף ארבעים יום אני בא בתוך שש שעות, כסבורים הם שאותו יום שעלה מן המנין הוא, והוא אמר להם שלמים – ארבעים יום ולילו עמו – ויום עליתו אין לילו עמו, שהרי בז' בסיון עלה, נמצא יום ארבעים בשבעה עשר בתמוז. בי"ו בא שטן וערבב את העולם, והראה דמות חשך ואפלה וערבוביה, לומר ודאי מת משה לכך בא ערבוביא לעולם, אמר להם מת משה, שכבר באו שש שעות ולא בא וכו' כדאיתא במסכת שבת (דף פ"ט); ואי אפשר לומר שלא טעו אלא ביום המענן בין קדם חצות בין לאחר חצות, שהרי לא ירד משה עד יום המחרת, שנאמר וישכימו ממחרת ויעלו עלת: אשר ילכו לפנינו. אלוהות הרבה אוו להם (סנהדרין ס"ג): כי זה משה האיש. כמין דמות משה הראה להם השטן, שנושאים אותו באויר רקיע השמים (שבת פ"ט): אשר העלנו מארץ מצרים. והיה מורה לנו דרך אשר נעלה בה, עתה צריכין אנו לאלוהות אשר ילכו לפנינו:}%
% ---------- פרק 32 | פסוק 2 ----------
 \textbf{ב} וַיֹּ֤אמֶר אֲלֵהֶם֙ אַהֲרֹ֔ן פָּֽרְקוּ֙ נִזְמֵ֣י הַזָּהָ֔ב אֲשֶׁר֙ בְּאׇזְנֵ֣י נְשֵׁיכֶ֔ם בְּנֵיכֶ֖ם וּבְנֹתֵיכֶ֑ם וְהָבִ֖יאוּ אֵלָֽי׃ \Rashi{\textbf{באזני נשיכם.}אמר אהרן בלבו: הנשים והילדים חסים על תכשיטיהן, שמא יתעכב הדבר ובתוך כך יבא משה, והם לא המתינו ופרקו מעל עצמן: פרקו. לשון צווי, מגזרת פרק ליחיד, כמו ברכו מגזרת ברך:}%
% ---------- פרק 32 | פסוק 3 ----------
 \textbf{ג} וַיִּתְפָּֽרְקוּ֙ כׇּל־הָעָ֔ם אֶת־נִזְמֵ֥י הַזָּהָ֖ב אֲשֶׁ֣ר בְּאׇזְנֵיהֶ֑ם וַיָּבִ֖יאוּ אֶֽל־אַהֲרֹֽן׃ \Rashi{\textbf{ויתפרקו.}לשון פריקת משא, כשנטלום מאזניהם נמצאו הם מפרקים מנזמיהם, דישקריי"ר בלעז: את נזמי. כמו מנזמי, כמו כצאתי את העיר (שמות ט') – מן העיר:}%
% ---------- פרק 32 | פסוק 4 ----------
 \textbf{ד} וַיִּקַּ֣ח מִיָּדָ֗ם וַיָּ֤צַר אֹתוֹ֙ בַּחֶ֔רֶט וַֽיַּעֲשֵׂ֖הוּ עֵ֣גֶל מַסֵּכָ֑ה וַיֹּ֣אמְר֔וּ אֵ֤לֶּה אֱלֹהֶ֙יךָ֙ יִשְׂרָאֵ֔ל אֲשֶׁ֥ר הֶעֱל֖וּךָ מֵאֶ֥רֶץ מִצְרָֽיִם׃ \Rashi{\textbf{ויצר אתו בחרט.}יש לתרגמו בשני פנים, האחד ויצר – לשון קשירה, בחרט – לשון סודר, כמו והמטפחות והחריטים (ישעיהו ג'), ויצר ככרים כסף בשני חרטים (מלכים ב ה'), והשני ויצר – לשון צורה, בחרט – כלי אמנות הצורפין שחורצין וחורתין בו צורות בזהב, כעט סופר החורת אותיות בלוחות ופנקסין, כמו וכתב עליו בחרט אנוש (ישעיהו ח'), וזהו שתרגם אנקלוס וצר יתיה בזיפא, לשון זיוף, הוא כלי אמנות שחורצין בו בזהב אותיות ושקדים, שקורין בלעז ניי"ל, ומזיפין על ידו חותמות: עגל מסכה. כיון שהשליכו לאור בכור, באו מכשפי ערב רב שעלו עמהם ממצרים ועשאוהו בכשפים; ויש אומרים מיכה היה שם, שיצא מתוך דמוסי בנין שנתמעך בו במצרים, והיה בידו שם וטס שכתב בו משה "עלה שור" "עלה שור" להעלות ארונו של יוסף מתוך נילוס – והשליכו לתוך הכור ויצא העגל (תנחומא): מסכה. לשון מתכת; ד"א קכ"ה קנטרין זהב היו בו כגימטריא של מסכה (שם): אלה אלהיך. ולא נאמר אלה אלהינו, מכאן שערב רב שעלו ממצרים הם שנקהלו על אהרן והם שעשאוהו, ואחר כך הטעו את ישראל אחריו (שם):}%
% ---------- פרק 32 | פסוק 5 ----------
 \textbf{ה} וַיַּ֣רְא אַהֲרֹ֔ן וַיִּ֥בֶן מִזְבֵּ֖חַ לְפָנָ֑יו וַיִּקְרָ֤א אַֽהֲרֹן֙ וַיֹּאמַ֔ר חַ֥ג לַיהֹוָ֖ה מָחָֽר׃ \Rashi{\textbf{וירא אהרן.}שהיה בו רוח חיים, שנאמר בתבנית שור אכל עשב (תהלים ק"ו), וראה שהצליח מעשה שטן ולא היה לו פה לדחותם לגמרי: ויבן מזבח. לדחותם: ויקרא חג לה' מחר. ולא היום, שמא יבא משה קדם שיעבדוהו, זהו פשוטו; ומדרשו, בויקרא רבה, דברים הרבה ראה אהרן – ראה חור בן אחותו שהיה מוכיחם והרגוהו, וזהו ויבן מזבח לפניו – ויבן מזבוח לפניו – ועוד ראה ואמר מוטב שיתלה בי הסרחון ולא בהם, ועוד ראה ואמר אם הם בונים את המזבח, זה מביא צרור וזה מביא אבן ונמצאת מלאכתן עשוי בבת אחת, מתוך שאני בונה אותו ואני מתעצל במלאכתי, בין כך ובין כך משה בא: חג לה'. בלבו היה לשמים, בטוח היה שיבא משה ויעבדו את המקום:}%
% ---------- פרק 32 | פסוק 6 ----------
 \textbf{ו} וַיַּשְׁכִּ֙ימוּ֙ מִֽמׇּחֳרָ֔ת וַיַּעֲל֣וּ עֹלֹ֔ת וַיַּגִּ֖שׁוּ שְׁלָמִ֑ים וַיֵּ֤שֶׁב הָעָם֙ לֶֽאֱכֹ֣ל וְשָׁת֔וֹ וַיָּקֻ֖מוּ לְצַחֵֽק׃ \{פ\} \Rashi{\textbf{וישכימו.}השטן זרזם, כדי שיחטאו: לצחק. יש במשמע הזה גלוי עריות, כמו שנ' לצחק בי (בראשית ל"ט), ושפיכות דמים, כמו שנ' יקומו נא הנערים וישחקו לפנינו (שמואל ב ב'), אף כן נהרג חור (שמות רבה):}%
% ---------- פרק 32 | פסוק 7 ----------
 \textbf{ז} וַיְדַבֵּ֥ר יְהֹוָ֖ה אֶל־מֹשֶׁ֑ה לֶךְ־רֵ֕ד כִּ֚י שִׁחֵ֣ת עַמְּךָ֔ אֲשֶׁ֥ר הֶעֱלֵ֖יתָ מֵאֶ֥רֶץ מִצְרָֽיִם׃ \Rashi{\textbf{וידבר.}לשון קשי הוא, כמו וידבר אתם קשות (בראשית מ"ב): לך רד. מגדלתך – לא נתתי לך גדלה אלא בשבילם (ברכות ל"ב), באותה שעה נתנדה משה מפי בית דין של מעלה (תנחומא): שחת עמך. שחת העם לא נאמר, אלא עמך – ערב רב שקבלת מעצמך וגירתם ולא נמלכת בי, ואמרת טוב שידבקו גרים בשכינה – הם שחתו והשחיתו (שמות רבה):}%
% ---------- פרק 32 | פסוק 8 ----------
 \textbf{ח} סָ֣רוּ מַהֵ֗ר מִן־הַדֶּ֙רֶךְ֙ אֲשֶׁ֣ר צִוִּיתִ֔ם עָשׂ֣וּ לָהֶ֔ם עֵ֖גֶל מַסֵּכָ֑ה וַיִּשְׁתַּֽחֲווּ־לוֹ֙ וַיִּזְבְּחוּ־ל֔וֹ וַיֹּ֣אמְר֔וּ אֵ֤לֶּה אֱלֹהֶ֙יךָ֙ יִשְׂרָאֵ֔ל אֲשֶׁ֥ר הֶֽעֱל֖וּךָ מֵאֶ֥רֶץ מִצְרָֽיִם׃ %
% ---------- פרק 32 | פסוק 9 ----------
 \textbf{ט} וַיֹּ֥אמֶר יְהֹוָ֖ה אֶל־מֹשֶׁ֑ה רָאִ֙יתִי֙ אֶת־הָעָ֣ם הַזֶּ֔ה וְהִנֵּ֥ה עַם־קְשֵׁה־עֹ֖רֶף הֽוּא׃ \Rashi{\textbf{קשה ערף.}מחזרין קשי ערפם לנגד מוכיחיהם וממאנים לשמע:}%
% ---------- פרק 32 | פסוק 10 ----------
 \textbf{י} וְעַתָּה֙ הַנִּ֣יחָה לִּ֔י וְיִֽחַר־אַפִּ֥י בָהֶ֖ם וַאֲכַלֵּ֑ם וְאֶֽעֱשֶׂ֥ה אוֹתְךָ֖ לְג֥וֹי גָּדֽוֹל׃ \Rashi{\textbf{הניחה לי.}עדין לא שמענו שהתפלל משה עליהם והוא אומר הניחה לי? אלא כאן פתח לו פתח והודיעו שהדבר תלוי בו – שאם יתפלל עליהם לא יכלם (ברכות ל"ב):}%
% ---------- פרק 32 | פסוק 11 ----------
 \textbf{יא} וַיְחַ֣ל מֹשֶׁ֔ה אֶת־פְּנֵ֖י יְהֹוָ֣ה אֱלֹהָ֑יו וַיֹּ֗אמֶר לָמָ֤ה יְהֹוָה֙ יֶחֱרֶ֤ה אַפְּךָ֙ בְּעַמֶּ֔ךָ אֲשֶׁ֤ר הוֹצֵ֙אתָ֙ מֵאֶ֣רֶץ מִצְרַ֔יִם בְּכֹ֥חַ גָּד֖וֹל וּבְיָ֥ד חֲזָקָֽה׃ \Rashi{\textbf{למה ה' יחרה אפך.}כלום מתקנא אלא חכם בחכם, גבור בגבור:}%
% ---------- פרק 32 | פסוק 12 ----------
 \textbf{יב} לָ֩מָּה֩ יֹאמְר֨וּ מִצְרַ֜יִם לֵאמֹ֗ר בְּרָעָ֤ה הֽוֹצִיאָם֙ לַהֲרֹ֤ג אֹתָם֙ בֶּֽהָרִ֔ים וּ֨לְכַלֹּתָ֔ם מֵעַ֖ל פְּנֵ֣י הָֽאֲדָמָ֑ה שׁ֚וּב מֵחֲר֣וֹן אַפֶּ֔ךָ וְהִנָּחֵ֥ם עַל־הָרָעָ֖ה לְעַמֶּֽךָ׃ \Rashi{\textbf{והנחם.}התעשת מחשבה אחרת – להיטיב להם על הרעה. אשר חשבת להם:}%
% ---------- פרק 32 | פסוק 13 ----------
 \textbf{יג} זְכֹ֡ר לְאַבְרָהָם֩ לְיִצְחָ֨ק וּלְיִשְׂרָאֵ֜ל עֲבָדֶ֗יךָ אֲשֶׁ֨ר נִשְׁבַּ֣עְתָּ לָהֶם֮ בָּךְ֒ וַתְּדַבֵּ֣ר אֲלֵהֶ֔ם אַרְבֶּה֙ אֶֽת־זַרְעֲכֶ֔ם כְּכוֹכְבֵ֖י הַשָּׁמָ֑יִם וְכׇל־הָאָ֨רֶץ הַזֹּ֜את אֲשֶׁ֣ר אָמַ֗רְתִּי אֶתֵּן֙ לְזַרְעֲכֶ֔ם וְנָחֲל֖וּ לְעֹלָֽם׃ \Rashi{\textbf{זכר לאברהם.}אם עברו על עשרת הדברות, אברהם אביהם נתנסה בעשרה נסיונות ועדין לא קבל שכרו, תנהו לו ויצאו עשרה בעשרה (תנחומא): לאברהם ליצחק ולישראל. אם לשרפה הם, זכר לאברהם שמסר עצמו לשרף עליך באור כשדים, אם להריגה, זכר ליצחק שפשט צוארו לעקדה, אם לגלות, זכר ליעקב שגלה לחרן, ואם אינן נצולין בזכותן, מה אתה אומר לי ואעשה אותך לגוי גדול? ואם כסא של ג' רגלים אינו עומד לפניך בשעת כעסך, קל וחמר לכסא של רגל אחת (ברכות ל"ב): אשר נשבעת להם בך. לא נשבעת להם בדבר שהוא כלה – לא בשמים ולא בארץ ולא בהרים ולא בגבעות – אלא בך, שאתה קים ושבועתך קימת לעולם, שנ' בי נשבעתי נאם ה' (בראשית כ"ב), וליצחק נאמר והקמתי את השבועה אשר נשבעתי לאברהם אביך (שם), וליעקב נאמר אני אל שדי פרה ורבה (שם ל"ה) – נשבע לו באל שדי:}%
% ---------- פרק 32 | פסוק 14 ----------
 \textbf{יד} וַיִּנָּ֖חֶם יְהֹוָ֑ה עַל־הָ֣רָעָ֔ה אֲשֶׁ֥ר דִּבֶּ֖ר לַעֲשׂ֥וֹת לְעַמּֽוֹ׃ \{פ\} %
% ---------- פרק 32 | פסוק 15 ----------
 \textbf{טו} וַיִּ֜פֶן וַיֵּ֤רֶד מֹשֶׁה֙ מִן־הָהָ֔ר וּשְׁנֵ֛י לֻחֹ֥ת הָעֵדֻ֖ת בְּיָד֑וֹ לֻחֹ֗ת כְּתֻבִים֙ מִשְּׁנֵ֣י עֶבְרֵיהֶ֔ם מִזֶּ֥ה וּמִזֶּ֖ה הֵ֥ם כְּתֻבִֽים׃ \Rashi{\textbf{משני עבריהם.}היו האותיות נקראות, ומעשה נסים הוא (שבת ק"ד):}%
% ---------- פרק 32 | פסוק 16 ----------
 \textbf{טז} וְהַ֨לֻּחֹ֔ת מַעֲשֵׂ֥ה אֱלֹהִ֖ים הֵ֑מָּה וְהַמִּכְתָּ֗ב מִכְתַּ֤ב אֱלֹהִים֙ ה֔וּא חָר֖וּת עַל־הַלֻּחֹֽת׃ \Rashi{\textbf{מעשה אלהים המה.}כמשמעו – הוא בכבודו עשאן; ד"א, כאדם האומר לחברו כל עסקיו של פלוני במלאכה פלונית, כך כל שעשועיו של הקב"ה בתורה: חרות. לשון חרת וחרט אחד הוא, שניהם לשון חקוק, אנטלי"יר בלעז:}%
% ---------- פרק 32 | פסוק 17 ----------
 \textbf{יז} וַיִּשְׁמַ֧ע יְהוֹשֻׁ֛עַ אֶת־ק֥וֹל הָעָ֖ם בְּרֵעֹ֑ה וַיֹּ֙אמֶר֙ אֶל־מֹשֶׁ֔ה ק֥וֹל מִלְחָמָ֖ה בַּֽמַּחֲנֶֽה׃ \Rashi{\textbf{ברעה.}בהריעו, שהיו מריעים ושמחים וצוחקים:}%
% ---------- פרק 32 | פסוק 18 ----------
 \textbf{יח} וַיֹּ֗אמֶר אֵ֥ין קוֹל֙ עֲנ֣וֹת גְּבוּרָ֔ה וְאֵ֥ין ק֖וֹל עֲנ֣וֹת חֲלוּשָׁ֑ה ק֣וֹל עַנּ֔וֹת אָנֹכִ֖י שֹׁמֵֽעַ׃ \Rashi{\textbf{אין קול ענות גבורה.}אין קול הזה נראה קול ענית גבורים הצועקים "נצחון" ולא קול חלשים שצועקים "וי אנוסה": קול ענות. קול חרופין וגדופין, המענין את נפש שומען כשנאמרין לו:}%
% ---------- פרק 32 | פסוק 19 ----------
 \textbf{יט} וַֽיְהִ֗י כַּאֲשֶׁ֤ר קָרַב֙ אֶל־הַֽמַּחֲנֶ֔ה וַיַּ֥רְא אֶת־הָעֵ֖גֶל וּמְחֹלֹ֑ת וַיִּֽחַר־אַ֣ף מֹשֶׁ֗ה וַיַּשְׁלֵ֤ךְ מִיָּדָו֙ אֶת־הַלֻּחֹ֔ת וַיְשַׁבֵּ֥ר אֹתָ֖ם תַּ֥חַת הָהָֽר׃ \Rashi{\textbf{וישלך מידו וגו'.}אמר: מה פסח שהוא אחד מן המצוות, אמרה תורה כל בן נכר לא יאכל בו (שמות י"ב), התורה כלה כאן וכל ישראל משמדים ואתננה להם? (שבת פ"ז): תחת ההר. לרגלי ההר:}%
% ---------- פרק 32 | פסוק 20 ----------
 \textbf{כ} וַיִּקַּ֞ח אֶת־הָעֵ֨גֶל אֲשֶׁ֤ר עָשׂוּ֙ וַיִּשְׂרֹ֣ף בָּאֵ֔שׁ וַיִּטְחַ֖ן עַ֣ד אֲשֶׁר־דָּ֑ק וַיִּ֙זֶר֙ עַל־פְּנֵ֣י הַמַּ֔יִם וַיַּ֖שְׁקְ אֶת־בְּנֵ֥י יִשְׂרָאֵֽל׃ \Rashi{\textbf{ויזר.}לשון נפוץ, וכן יזרה על נוהו גפרית (איוב י"ח), וכן כי חנם מזרה הרשת (משלי א'), שזורין בה דגן וקטניות: וישק את בני ישראל. נתכון לבדקן כסוטות; שלש מיתות נדונו שם, אם יש עדים והתראה בסיף – כמשפט אנשי עיר הנדחת שהן מרבין – עדים בלא התראה במגפה, שנאמר ויגף ה' את העם, לא עדים ולא התראה בהדרוקן, שבדקום המים וצבו בטניהם (יומא ס"ו):}%
% ---------- פרק 32 | פסוק 21 ----------
 \textbf{כא} וַיֹּ֤אמֶר מֹשֶׁה֙ אֶֽל־אַהֲרֹ֔ן מֶֽה־עָשָׂ֥ה לְךָ֖ הָעָ֣ם הַזֶּ֑ה כִּֽי־הֵבֵ֥אתָ עָלָ֖יו חֲטָאָ֥ה גְדֹלָֽה׃ \Rashi{\textbf{מה עשה לך העם.}כמה יסורים סבלת שיסרוך עד שלא תביא עליהם חטא זה:}%
% ---------- פרק 32 | פסוק 22 ----------
 \textbf{כב} וַיֹּ֣אמֶר אַהֲרֹ֔ן אַל־יִ֥חַר אַ֖ף אֲדֹנִ֑י אַתָּה֙ יָדַ֣עְתָּ אֶת־הָעָ֔ם כִּ֥י בְרָ֖ע הֽוּא׃ \Rashi{\textbf{כי ברע הוא.}בדרך רע הם הולכין תמיד, ובנסיונות לפני המקום:}%
% ---------- פרק 32 | פסוק 23 ----------
 \textbf{כג} וַיֹּ֣אמְרוּ לִ֔י עֲשֵׂה־לָ֣נוּ אֱלֹהִ֔ים אֲשֶׁ֥ר יֵלְכ֖וּ לְפָנֵ֑ינוּ כִּי־זֶ֣ה ׀ מֹשֶׁ֣ה הָאִ֗ישׁ אֲשֶׁ֤ר הֶֽעֱלָ֙נוּ֙ מֵאֶ֣רֶץ מִצְרַ֔יִם לֹ֥א יָדַ֖עְנוּ מֶה־הָ֥יָה לֽוֹ׃ %
% ---------- פרק 32 | פסוק 24 ----------
 \textbf{כד} וָאֹמַ֤ר לָהֶם֙ לְמִ֣י זָהָ֔ב הִתְפָּרָ֖קוּ וַיִּתְּנוּ־לִ֑י וָאַשְׁלִכֵ֣הוּ בָאֵ֔שׁ וַיֵּצֵ֖א הָעֵ֥גֶל הַזֶּֽה׃ \Rashi{\textbf{ואמר להם.}דבר אחד – "למי זהב" לבד – והם מהרו והתפרקו ויתנו לי: ואשלכהו באש. ולא ידעתי שיצא העגל הזה, ויצא:}%
% ---------- פרק 32 | פסוק 25 ----------
 \textbf{כה} וַיַּ֤רְא מֹשֶׁה֙ אֶת־הָעָ֔ם כִּ֥י פָרֻ֖עַ ה֑וּא כִּֽי־פְרָעֹ֣ה אַהֲרֹ֔ן לְשִׁמְצָ֖ה בְּקָמֵיהֶֽם׃ \Rashi{\textbf{פרע.}מגלה; נתגלה שמצו וקלונו, כמו ופרע את ראש האשה (במדבר ה'): לשמצה בקמיהם. להיות להם הדבר הזה לגנות בפי כל הקמים עליהם:}%
% ---------- פרק 32 | פסוק 26 ----------
 \textbf{כו} וַיַּעֲמֹ֤ד מֹשֶׁה֙ בְּשַׁ֣עַר הַֽמַּחֲנֶ֔ה וַיֹּ֕אמֶר מִ֥י לַיהֹוָ֖ה אֵלָ֑י וַיֵּאָסְפ֥וּ אֵלָ֖יו כׇּל־בְּנֵ֥י לֵוִֽי׃ \Rashi{\textbf{מי לה' אלי.}יבא אלי: כל בני לוי. מכאן שכל השבט כשר (יומא ס"ו):}%
% ---------- פרק 32 | פסוק 27 ----------
 \textbf{כז} וַיֹּ֣אמֶר לָהֶ֗ם כֹּֽה־אָמַ֤ר יְהֹוָה֙ אֱלֹהֵ֣י יִשְׂרָאֵ֔ל שִׂ֥ימוּ אִישׁ־חַרְבּ֖וֹ עַל־יְרֵכ֑וֹ עִבְר֨וּ וָשׁ֜וּבוּ מִשַּׁ֤עַר לָשַׁ֙עַר֙ בַּֽמַּחֲנֶ֔ה וְהִרְג֧וּ אִֽישׁ־אֶת־אָחִ֛יו וְאִ֥ישׁ אֶת־רֵעֵ֖הוּ וְאִ֥ישׁ אֶת־קְרֹבֽוֹ׃ \Rashi{\textbf{כה אמר וגו'.}והיכן אמר? זבח לאלהים יחרם (שמות כ"ב) – כך שנויה במכילתא: אחיו. מאמו והוא ישראל (יומא ס"ו):}%
% ---------- פרק 32 | פסוק 28 ----------
 \textbf{כח} וַיַּֽעֲשׂ֥וּ בְנֵֽי־לֵוִ֖י כִּדְבַ֣ר מֹשֶׁ֑ה וַיִּפֹּ֤ל מִן־הָעָם֙ בַּיּ֣וֹם הַה֔וּא כִּשְׁלֹ֥שֶׁת אַלְפֵ֖י אִֽישׁ׃ %
% ---------- פרק 32 | פסוק 29 ----------
 \textbf{כט} וַיֹּ֣אמֶר מֹשֶׁ֗ה מִלְא֨וּ יֶדְכֶ֤ם הַיּוֹם֙ לַֽיהֹוָ֔ה כִּ֛י אִ֥ישׁ בִּבְנ֖וֹ וּבְאָחִ֑יו וְלָתֵ֧ת עֲלֵיכֶ֛ם הַיּ֖וֹם בְּרָכָֽה׃ \Rashi{\textbf{מלאו ידכם.}אתם ההורגים אותם – בדבר זה תתחנכו להיות כהנים למקום: כי איש. מכם ימלא ידו בבנו ובאחיו:}%
% ---------- פרק 32 | פסוק 30 ----------
 \textbf{ל} וַֽיְהִי֙ מִֽמׇּחֳרָ֔ת וַיֹּ֤אמֶר מֹשֶׁה֙ אֶל־הָעָ֔ם אַתֶּ֥ם חֲטָאתֶ֖ם חֲטָאָ֣ה גְדֹלָ֑ה וְעַתָּה֙ אֶֽעֱלֶ֣ה אֶל־יְהֹוָ֔ה אוּלַ֥י אֲכַפְּרָ֖ה בְּעַ֥ד חַטַּאתְכֶֽם׃ \Rashi{\textbf{אכפרה בעד חטאתכם.}אשים כפר וקנוח וסתימה לנגד חטאתכם, להבדיל ביניכם ובין החטא:}%
% ---------- פרק 32 | פסוק 31 ----------
 \textbf{לא} וַיָּ֧שׇׁב מֹשֶׁ֛ה אֶל־יְהֹוָ֖ה וַיֹּאמַ֑ר אָ֣נָּ֗א חָטָ֞א הָעָ֤ם הַזֶּה֙ חֲטָאָ֣ה גְדֹלָ֔ה וַיַּֽעֲשׂ֥וּ לָהֶ֖ם אֱלֹהֵ֥י זָהָֽב׃ \Rashi{\textbf{אלהי זהב.}אתה הוא שגרמת להם, שהשפעת להם זהב וכל חפצם, מה יעשו שלא יחטאו? משל למלך שהיה מאכיל ומשקה את בנו ומקשטו ותולה לו כיס בצוארו ומעמידו בפתח בית זונות, מה יעשה הבן שלא יחטא:}%
% ---------- פרק 32 | פסוק 32 ----------
 \textbf{לב} וְעַתָּ֖ה אִם־תִּשָּׂ֣א חַטָּאתָ֑ם וְאִם־אַ֕יִן מְחֵ֣נִי נָ֔א מִֽסִּפְרְךָ֖ אֲשֶׁ֥ר כָּתָֽבְתָּ׃ \Rashi{\textbf{ועתה אם תשא חטאתם.}הרי טוב – איני אומר לך מחני, ואם אין מחני; וזה מקרא קצר, וכן הרבה: מספרך. מכל התורה כלה, שלא יאמרו עלי שלא הייתי כדאי לבקש עליהם רחמים:}%
% ---------- פרק 32 | פסוק 33 ----------
 \textbf{לג} וַיֹּ֥אמֶר יְהֹוָ֖ה אֶל־מֹשֶׁ֑ה מִ֚י אֲשֶׁ֣ר חָֽטָא־לִ֔י אֶמְחֶ֖נּוּ מִסִּפְרִֽי׃ %
% ---------- פרק 32 | פסוק 34 ----------
 \textbf{לד} וְעַתָּ֞ה לֵ֣ךְ ׀ נְחֵ֣ה אֶת־הָעָ֗ם אֶ֤ל אֲשֶׁר־דִּבַּ֙רְתִּי֙ לָ֔ךְ הִנֵּ֥ה מַלְאָכִ֖י יֵלֵ֣ךְ לְפָנֶ֑יךָ וּבְי֣וֹם פׇּקְדִ֔י וּפָקַדְתִּ֥י עֲלֵהֶ֖ם חַטָּאתָֽם׃ \Rashi{\textbf{אל אשר דברתי לך.}יש כאן לך אצל דבור במקום אליך, וכן לדבר לו על אדניהו (מלכים א ב׳:י״ט): הנה מלאכי. ולא אני: וביום פקדי וגו'. עתה שמעתי אליך מלכלותם יחד, ותמיד תמיד כשאפקד עליהם עונותיהם ופקדתי עליהם מעט מן העון הזה עם שאר העונות; ואין פרענות באה על ישראל שאין בה קצת מפרעון עון העגל (סנהדרין ק"ב):}%
% ---------- פרק 32 | פסוק 35 ----------
 \textbf{לה} וַיִּגֹּ֥ף יְהֹוָ֖ה אֶת־הָעָ֑ם עַ֚ל אֲשֶׁ֣ר עָשׂ֣וּ אֶת־הָעֵ֔גֶל אֲשֶׁ֥ר עָשָׂ֖ה אַהֲרֹֽן׃ \{ס\} \Rashi{\textbf{ויגף ה' את העם.}מיתה בידי שמים לעדים בלא התראה:}%
% ---------- פרק 33 | פסוק 1 ----------
 \textbf{א} וַיְדַבֵּ֨ר יְהֹוָ֤ה אֶל־מֹשֶׁה֙ לֵ֣ךְ עֲלֵ֣ה מִזֶּ֔ה אַתָּ֣ה וְהָעָ֔ם אֲשֶׁ֥ר הֶֽעֱלִ֖יתָ מֵאֶ֣רֶץ מִצְרָ֑יִם אֶל־הָאָ֗רֶץ אֲשֶׁ֣ר נִ֠שְׁבַּ֠עְתִּי לְאַבְרָהָ֨ם לְיִצְחָ֤ק וּֽלְיַעֲקֹב֙ לֵאמֹ֔ר לְזַרְעֲךָ֖ אֶתְּנֶֽנָּה׃ \Rashi{\textbf{לך עלה מזה.}ארץ ישראל גבוהה מכל הארצות, לכך נאמר עלה; ד"א, כלפי שאמר לו בשעת הכעס "לך רד" אמר לו בשעת רצון "לך עלה" (תנחומא): אתה והעם. כאן לא נאמר ועמך:}%
% ---------- פרק 33 | פסוק 2 ----------
 \textbf{ב} וְשָׁלַחְתִּ֥י לְפָנֶ֖יךָ מַלְאָ֑ךְ וְגֵֽרַשְׁתִּ֗י אֶת־הַֽכְּנַעֲנִי֙ הָֽאֱמֹרִ֔י וְהַֽחִתִּי֙ וְהַפְּרִזִּ֔י הַחִוִּ֖י וְהַיְבוּסִֽי׃ \Rashi{\textbf{וגרשתי את הכנעני וגו'.}ו' אמות הם והגרגשי עמד ופנה מפניהם מאליו (תלמוד ירושלמי שביעית ו'):}%
% ---------- פרק 33 | פסוק 3 ----------
 \textbf{ג} אֶל־אֶ֛רֶץ זָבַ֥ת חָלָ֖ב וּדְבָ֑שׁ כִּי֩ לֹ֨א אֶֽעֱלֶ֜ה בְּקִרְבְּךָ֗ כִּ֤י עַם־קְשֵׁה־עֹ֙רֶף֙ אַ֔תָּה פֶּן־אֲכֶלְךָ֖ בַּדָּֽרֶךְ׃ \Rashi{\textbf{אל ארץ זבת חלב ודבש.}אני אומר לך להעלותם: כי לא אעלה בקרבך. לכך אני אומר לך ושלחתי לפניך מלאך: כי עם קשה ערף אתה. וכששכינתי בקרבכם ואתם ממרים בי, מרבה אני עליכם זעם: אכלך. לשון כליון:}%
% ---------- פרק 33 | פסוק 4 ----------
 \textbf{ד} וַיִּשְׁמַ֣ע הָעָ֗ם אֶת־הַדָּבָ֥ר הָרָ֛ע הַזֶּ֖ה וַיִּתְאַבָּ֑לוּ וְלֹא־שָׁ֛תוּ אִ֥ישׁ עֶדְי֖וֹ עָלָֽיו׃ \Rashi{\textbf{הדבר הרע.}שאין השכינה שורה ומהלכת עמם: איש עדיו. כתרים שנתנו להם בחורב כשאמרו נעשה ונשמע (שבת פ"ח):}%
% ---------- פרק 33 | פסוק 5 ----------
 \textbf{ה} וַיֹּ֨אמֶר יְהֹוָ֜ה אֶל־מֹשֶׁ֗ה אֱמֹ֤ר אֶל־בְּנֵֽי־יִשְׂרָאֵל֙ אַתֶּ֣ם עַם־קְשֵׁה־עֹ֔רֶף רֶ֧גַע אֶחָ֛ד אֶֽעֱלֶ֥ה בְקִרְבְּךָ֖ וְכִלִּיתִ֑יךָ וְעַתָּ֗ה הוֹרֵ֤ד עֶדְיְךָ֙ מֵֽעָלֶ֔יךָ וְאֵדְעָ֖ה מָ֥ה אֶֽעֱשֶׂה־לָּֽךְ׃ \Rashi{\textbf{רגע אחד אעלה בקרבך וכליתיך.}אם אעלה בקרבך ואתם ממרים בי בקשיות ערפכם אזעם עליכם רגע אחד – שהוא שעור זעמו, שנ' חבי כמעט רגע עד יעבר זעם (ישעיהו כ"ו) – ואכלה אתכם, לפיכך טוב לכם שאשלח מלאך: ועתה. פרענות זו תלקו מיד – שתורידו עדיכם מעליכם: ואדעה מה אעשה לך. בפקדת שאר העון אני יודע מה שבלבי לעשות לך:}%
% ---------- פרק 33 | פסוק 6 ----------
 \textbf{ו} וַיִּֽתְנַצְּל֧וּ בְנֵֽי־יִשְׂרָאֵ֛ל אֶת־עֶדְיָ֖ם מֵהַ֥ר חוֹרֵֽב׃ \Rashi{\textbf{את עדים מהר חורב.}את העדי שהיה בידם מהר חורב:}%
% ---------- פרק 33 | פסוק 7 ----------
 \textbf{ז} וּמֹשֶׁה֩ יִקַּ֨ח אֶת־הָאֹ֜הֶל וְנָֽטָה־ל֣וֹ ׀ מִח֣וּץ לַֽמַּחֲנֶ֗ה הַרְחֵק֙ מִן־הַֽמַּחֲנֶ֔ה וְקָ֥רָא ל֖וֹ אֹ֣הֶל מוֹעֵ֑ד וְהָיָה֙ כׇּל־מְבַקֵּ֣שׁ יְהֹוָ֔ה יֵצֵא֙ אֶל־אֹ֣הֶל מוֹעֵ֔ד אֲשֶׁ֖ר מִח֥וּץ לַֽמַּחֲנֶֽה׃ \Rashi{\textbf{ומשה.}מאותו עון והלאה, יקח את האהל. לשון הווה הוא – לוקח אהלו ונוטהו מחוץ למחנה; אמר מנדה לרב מנדה לתלמיד (תנחומא): הרחק. אלפים אמה, כענין שנ' אך רחוק יהיה ביניכם ובניו כאלפים אמה במדה (יהושע ג'): וקרא לו. והיה קורא לו אהל מועד, הוא בית ועד למבקשי תורה: כל מבקש ה'. מכאן למבקש פני זקן כמקבל פני שכינה (תנחומא): יצא אל אהל מועד. כמו יוצא. ד"א והיה כל מבקש ה' – אפלו מלאכי השרת כשהיו שואלים מקום שכינה, חבריהם אומרים להם הרי הוא באהלו של משה:}%
% ---------- פרק 33 | פסוק 8 ----------
 \textbf{ח} וְהָיָ֗ה כְּצֵ֤את מֹשֶׁה֙ אֶל־הָאֹ֔הֶל יָק֙וּמוּ֙ כׇּל־הָעָ֔ם וְנִ֨צְּב֔וּ אִ֖ישׁ פֶּ֣תַח אׇהֳל֑וֹ וְהִבִּ֙יטוּ֙ אַחֲרֵ֣י מֹשֶׁ֔ה עַד־בֹּא֖וֹ הָאֹֽהֱלָה׃ \Rashi{\textbf{והיה.}לשון הווה: כצאת. משה מן המחנה ללכת אל האהל: יקומו כל העם. עומדים מפניו, ואין יושבין עד שנתכסה מהם: והביטו אחרי משה. לשבח – אשרי ילוד אשה שכך מבטח שהשכינה תכנס אחריו לפתח אהלו:}%
% ---------- פרק 33 | פסוק 9 ----------
 \textbf{ט} וְהָיָ֗ה כְּבֹ֤א מֹשֶׁה֙ הָאֹ֔הֱלָה יֵרֵד֙ עַמּ֣וּד הֶֽעָנָ֔ן וְעָמַ֖ד פֶּ֣תַח הָאֹ֑הֶל וְדִבֶּ֖ר עִם־מֹשֶֽׁה׃ \Rashi{\textbf{ודבר עם משה.}כמו ומדבר עם משה; תרגומו ומתמלל עם משה שהוא כבוד שכינה, כמו וישמע את הקול מדבר אליו (במדבר ז'), ואינו קורא מדבר אליו; כשהוא קורא מדבר, פתרונו הקול מדבר בינו לבין עצמו, וההדיוט שומע מאליו, וכשהוא קורא מדבר משמע שהמלך מדבר עם ההדיוט:}%
% ---------- פרק 33 | פסוק 10 ----------
 \textbf{י} וְרָאָ֤ה כׇל־הָעָם֙ אֶת־עַמּ֣וּד הֶֽעָנָ֔ן עֹמֵ֖ד פֶּ֣תַח הָאֹ֑הֶל וְקָ֤ם כׇּל־הָעָם֙ וְהִֽשְׁתַּחֲו֔וּ אִ֖ישׁ פֶּ֥תַח אׇהֳלֽוֹ׃ \Rashi{\textbf{והשתחוו.}לשכינה:}%
% ---------- פרק 33 | פסוק 11 ----------
 \textbf{יא} וְדִבֶּ֨ר יְהֹוָ֤ה אֶל־מֹשֶׁה֙ פָּנִ֣ים אֶל־פָּנִ֔ים כַּאֲשֶׁ֛ר יְדַבֵּ֥ר אִ֖ישׁ אֶל־רֵעֵ֑הוּ וְשָׁב֙ אֶל־הַֽמַּחֲנֶ֔ה וּמְשָׁ֨רְת֜וֹ יְהוֹשֻׁ֤עַ בִּן־נוּן֙ נַ֔עַר לֹ֥א יָמִ֖ישׁ מִתּ֥וֹךְ הָאֹֽהֶל׃ \{פ\} \Rashi{\textbf{ודבר ה' אל משה פנים אל פנים.}ומתמלל עם משה: ושב אל המחנה. לאחר שנדבר עמו היה משה שב אל המחנה ומלמד לזקנים מה שלמד; והדבר הזה נהג משה מיום הכפורים עד שהוקם המשכן, ולא יותר, שהרי בשבעה עשר בתמוז נשתברו הלוחות, ובי"ח שרף את העגל ודן את החוטאים, ובי"ט עלה שנאמר ויהי ממחרת ויאמר משה אל העם וגו' (שמות ל"ב), ועשה שם ארבעים יום ובקש רחמים, שנ' ואתנפל לפני ה' וגו' (דברים ט'), ובראש חדש אלול נאמר לו ועלית בבקר אל הר סיני (שמות ל"ד), לקבל לוחות האחרונות – ועשה שם מ' יום, שנאמר בהם ואנכי עמדתי בהר כימים הראשונים וגו' – מה הראשונים ברצון אף האחרונים ברצון, אמר מעתה אמצעיים היו בכעס – בעשרה בתשרי נתרצה הקב"ה לישראל בשמחה ובלב שלם, ואמר לו למשה סלחתי, ומסר לו לוחות אחרונות, וירד והתחיל לצוותם על מלאכת המשכן, ועשאוהו עד אחד בניסן, ומשהוקם לא נדבר עמו עוד אלא מאהל מועד (תנחומא): ושב אל המחנה. תרגומו ותיב למשריתא, לפי שהוא לשון הווה, וכן כל הענין, וראה כל העם – וחזן, ונצבו – וקימין, והביטו – ומסתכלין, והשתחוו – וסגדין. ומדרשו: ודבר ה' אל משה – שישוב אל המחנה, אמר לו אני בכעס ואתה בכעס, אם כן מי יקריבם? (תנחומא):}%
% ---------- פרק 33 | פסוק 12 ----------
 \textbf{יב} וַיֹּ֨אמֶר מֹשֶׁ֜ה אֶל־יְהֹוָ֗ה רְ֠אֵ֠ה אַתָּ֞ה אֹמֵ֤ר אֵלַי֙ הַ֚עַל אֶת־הָעָ֣ם הַזֶּ֔ה וְאַתָּה֙ לֹ֣א הֽוֹדַעְתַּ֔נִי אֵ֥ת אֲשֶׁר־תִּשְׁלַ֖ח עִמִּ֑י וְאַתָּ֤ה אָמַ֙רְתָּ֙ יְדַעְתִּ֣יךָֽ בְשֵׁ֔ם וְגַם־מָצָ֥אתָ חֵ֖ן בְּעֵינָֽי׃ \Rashi{\textbf{ראה אתה אמר אלי.}ראה – תן עיניך ולבך על דבריך: אתה אמר אלי וגו' ואתה לא הודעתני וגו'. ואשר אמרת לי הנה אנכי שלח מלאך, אין זו הודעה, שאין אני חפץ בה: ואתה אמרת ידעתיך בשם. הכרתיך משאר בני אדם בשם חשיבות, שהרי אמרת לי הנה אנכי בא אליך בעב הענן וגו' וגם בך יאמינו לעולם (שמות י"ט):}%
% ---------- פרק 33 | פסוק 13 ----------
 \textbf{יג} וְעַתָּ֡ה אִם־נָא֩ מָצָ֨אתִי חֵ֜ן בְּעֵינֶ֗יךָ הוֹדִעֵ֤נִי נָא֙ אֶת־דְּרָכֶ֔ךָ וְאֵדָ֣עֲךָ֔ לְמַ֥עַן אֶמְצָא־חֵ֖ן בְּעֵינֶ֑יךָ וּרְאֵ֕ה כִּ֥י עַמְּךָ֖ הַגּ֥וֹי הַזֶּֽה׃ \Rashi{\textbf{ועתה וגו'.}אם אמת שמצאתי חן בעיניך הודעני נא את דרכך. מה שכר אתה נותן למוצאי חן בעיניך: ואדעך למען אמצא חן בעיניך. ואדע בזו מדת תגמוליך – מה היא מציאת חן שמצאתי בעיניך; ופתרון למען אמצא חן – למען אכיר כמה שכר מציאת החן: וראה כי עמך הגוי הזה. שלא תאמר ואעשה אותך לגוי גדול ואת אלה תעזב, ראה כי עמך הם מקדם, ואם בהם תמאס, איני סומך על היוצאים מחלצי שיתקימו, ואת תשלום השכר שלי בעם הזה תודיעני; ורבותינו דרשו במסכת ברכות (דף ז'), ואני לישב המקראות על אפניהם ועל סדרם באתי:}%
% ---------- פרק 33 | פסוק 14 ----------
 \textbf{יד} וַיֹּאמַ֑ר פָּנַ֥י יֵלֵ֖כוּ וַהֲנִחֹ֥תִי לָֽךְ׃ \Rashi{\textbf{ויאמר פני ילכו.}כתרגומו – לא אשלח עוד מלאך, אני בעצמי אלך, כמו ופניך הולכים בקרב (שמואל ב י״ז:י״א):}%
% ---------- פרק 33 | פסוק 15 ----------
 \textbf{טו} וַיֹּ֖אמֶר אֵלָ֑יו אִם־אֵ֤ין פָּנֶ֙יךָ֙ הֹלְכִ֔ים אַֽל־תַּעֲלֵ֖נוּ מִזֶּֽה׃ \Rashi{\textbf{ויאמר אליו.}בזו אני חפץ, כי על ידי מלאך אל תעלנו מזה:}%
% ---------- פרק 33 | פסוק 16 ----------
 \textbf{טז} וּבַמֶּ֣ה ׀ יִוָּדַ֣ע אֵפ֗וֹא כִּֽי־מָצָ֨אתִי חֵ֤ן בְּעֵינֶ֙יךָ֙ אֲנִ֣י וְעַמֶּ֔ךָ הֲל֖וֹא בְּלֶכְתְּךָ֣ עִמָּ֑נוּ וְנִפְלִ֙ינוּ֙ אֲנִ֣י וְעַמְּךָ֔ מִכׇּ֨ל־הָעָ֔ם אֲשֶׁ֖ר עַל־פְּנֵ֥י הָאֲדָמָֽה׃ \{פ\} \Rashi{\textbf{ובמה יודע אפוא.}יודע מציאת החן, הלוא בלכתך עמנו. ועוד דבר אני שואל ממך, שלא תשרה שכינתך עוד על אמות העולם. ונפלינו אני ועמך. ונהיה מבדלים בדבר הזה מכל העם (ברכות ז'), כמו והפלה ה' בין מקנה ישראל וגו' (שמות ט'):}%
% ---------- פרק 33 | פסוק 17 ----------
 \textbf{יז} וַיֹּ֤אמֶר יְהֹוָה֙ אֶל־מֹשֶׁ֔ה גַּ֣ם אֶת־הַדָּבָ֥ר הַזֶּ֛ה אֲשֶׁ֥ר דִּבַּ֖רְתָּ אֶֽעֱשֶׂ֑ה כִּֽי־מָצָ֤אתָ חֵן֙ בְּעֵינַ֔י וָאֵדָעֲךָ֖ בְּשֵֽׁם׃ \Rashi{\textbf{גם את הדבר הזה.}שלא תשרה שכינתי עוד על אמות העולם אעשה, ואין דבריו של בלעם על ידי שרית השכינה, אלא "נופל וגלוי עינים" כגון ואלי דבר יגנב (איוב ד') – שומעין ע"י שליח:}%
% ---------- פרק 33 | פסוק 18 ----------
 \textbf{יח} וַיֹּאמַ֑ר הַרְאֵ֥נִי נָ֖א אֶת־כְּבֹדֶֽךָ׃ \Rashi{\textbf{ויאמר הראני נא את כבדך.}ראה משה שהיה עת רצון ודבריו מקבלים, והוסיף לשאל להראותו מראית כבודו:}%
% ---------- פרק 33 | פסוק 19 ----------
 \textbf{יט} וַיֹּ֗אמֶר אֲנִ֨י אַעֲבִ֤יר כׇּל־טוּבִי֙ עַל־פָּנֶ֔יךָ וְקָרָ֧אתִֽי בְשֵׁ֛ם יְהֹוָ֖ה לְפָנֶ֑יךָ וְחַנֹּתִי֙ אֶת־אֲשֶׁ֣ר אָחֹ֔ן וְרִחַמְתִּ֖י אֶת־אֲשֶׁ֥ר אֲרַחֵֽם׃ \Rashi{\textbf{ויאמר אני אעביר וגו'.}הגיעה שעה שתראה בכבודי מה שארשה אותך לראות, לפי שאני רוצה וצריך ללמדך סדר תפלה; שכשנצרכת לבקש רחמים על ישראל הזכרת לי זכות אבות, כסבור אתה שאם תמה זכות אבות אין עוד תקוה, אני אעביר כל מדת טובי לפניך על הצור ואתה נתון במערה: וקראתי בשם ה' לפניך. ללמדך סדר בקשת רחמים אף אם תכלה זכות אבות, וכסדר זה שאתה רואה אותי – מעטף וקורא י"ג מדות – הוי מלמד את ישראל לעשות כן, וע"י שיזכירו לפני "רחום וחנון" יהיו נענין – כי רחמי לא כלים: וחנתי את אשר אחון. אותן פעמים שארצה לחן: ורחמתי. עת שאחפץ לרחם; עד כאן לא הבטיחו אלא עתים אענה ועתים לא אענה, אבל בשעת מעשה אמר לו הנה אנכי כרת ברית (שמות ל"ד) – הבטיחו שאינן חוזרות ריקם:}%
% ---------- פרק 33 | פסוק 20 ----------
 \textbf{כ} וַיֹּ֕אמֶר לֹ֥א תוּכַ֖ל לִרְאֹ֣ת אֶת־פָּנָ֑י כִּ֛י לֹֽא־יִרְאַ֥נִי הָאָדָ֖ם וָחָֽי׃ \Rashi{\textbf{ויאמר לא תוכל וגו'.}אף כשאעביר כל טובי על פניך איני נותן לך רשות לראות את פני:}%
% ---------- פרק 33 | פסוק 21 ----------
 \textbf{כא} וַיֹּ֣אמֶר יְהֹוָ֔ה הִנֵּ֥ה מָק֖וֹם אִתִּ֑י וְנִצַּבְתָּ֖ עַל־הַצּֽוּר׃ \Rashi{\textbf{הנה מקום אתי.}בהר אשר אני מדבר עמך תמיד, יש מקום מוכן לי לצרכך שאטמינך שם שלא תזק, ומשם תראה מה שתראה, זהו פשוטו; ומדרשו, על מקום שהשכינה שם מדבר, ואומר המקום אתי, ואינו אומר אני במקום, שהקב"ה מקומו של עולם ואין עולמו מקומו (תנחומא):}%
% ---------- פרק 33 | פסוק 22 ----------
 \textbf{כב} וְהָיָה֙ בַּעֲבֹ֣ר כְּבֹדִ֔י וְשַׂמְתִּ֖יךָ בְּנִקְרַ֣ת הַצּ֑וּר וְשַׂכֹּתִ֥י כַפִּ֛י עָלֶ֖יךָ עַד־עׇבְרִֽי׃ \Rashi{\textbf{בעבר כבדי.}כשאעבר לפניך: בנקרת הצור. כמו העיני האנשים ההם תנקר (במדבר ט"ז), יקרוה ערבי נחל (משלי ל'), אני קרתי ושתיתי מים (מלכים ב' יט, כד; ישעיה לז, כה), גזרה אחת להם, נקרת הצור: כרית הצור: ושכתי כפי. מכאן שנתנה רשות למחבלים לחבל, ותרגומו "ואגין במימרי" כנוי הוא לדרך כבוד של מעלה, שאינו צריך לסוכך עליו בכף ממש:}%
% ---------- פרק 33 | פסוק 23 ----------
 \textbf{כג} וַהֲסִרֹתִי֙ אֶת־כַּפִּ֔י וְרָאִ֖יתָ אֶת־אֲחֹרָ֑י וּפָנַ֖י לֹ֥א יֵרָאֽוּ׃ \{פ\}*(ספק פרשה סתומה בכתר ארם צובה) \Rashi{\textbf{והסרתי את כפי.}ואעדי ית דברת יקרי – כשאסלק הנהגת כבודי מנגד פניך, ללכת משם ולהלן: וראית את אחרי. הראהו קשר של תפלין (ברכות ז'):}%
% ---------- פרק 34 | פסוק 1 ----------
 \textbf{א} וַיֹּ֤אמֶר יְהֹוָה֙ אֶל־מֹשֶׁ֔ה פְּסׇל־לְךָ֛ שְׁנֵֽי־לֻחֹ֥ת אֲבָנִ֖ים כָּרִאשֹׁנִ֑ים וְכָתַבְתִּי֙ עַל־הַלֻּחֹ֔ת אֶ֨ת־הַדְּבָרִ֔ים אֲשֶׁ֥ר הָי֛וּ עַל־הַלֻּחֹ֥ת הָרִאשֹׁנִ֖ים אֲשֶׁ֥ר שִׁבַּֽרְתָּ׃ \Rashi{\textbf{פסל לך.}הראהו מחצב סנפירינון מתוך אהלו ואמר לו הפסלת יהיה שלך; משם נתעשר משה הרבה (תנחומא): פסל לך. אתה שברת הראשונות, אתה פסל לך אחרות, משל למלך שהלך למדינת הים והניח ארוסתו עם השפחות, מתוך קלקול השפחות יצא עליה שם רע, עמד שושבינה וקרע כתבתה, אמר אם יאמר המלך להרגה, אמר לו עדין אינה אשתך; בדק המלך ומצא שלא היה הקלקול אלא מן השפחות, נתרצה לה, אמר לו שושבינה כתב לה כתבה אחרת, שנקרעה הראשונה, אמר לו המלך אתה קרעת אותה אתה קנה לה ניר אחר ואני אכתב לה בכתב ידי; כן המלך זה הקב"ה, השפחות אלו ערב רב, והשושבין זה משה, ארוסתו של הקב"ה ישראל, לכך נאמר פסל לך (שם):}%
% ---------- פרק 34 | פסוק 2 ----------
 \textbf{ב} וֶהְיֵ֥ה נָכ֖וֹן לַבֹּ֑קֶר וְעָלִ֤יתָ בַבֹּ֙קֶר֙ אֶל־הַ֣ר סִינַ֔י וְנִצַּבְתָּ֥ לִ֛י שָׁ֖ם עַל־רֹ֥אשׁ הָהָֽר׃ \Rashi{\textbf{נכון.}מזמן:}%
% ---------- פרק 34 | פסוק 3 ----------
 \textbf{ג} וְאִישׁ֙ לֹֽא־יַעֲלֶ֣ה עִמָּ֔ךְ וְגַם־אִ֥ישׁ אַל־יֵרָ֖א בְּכׇל־הָהָ֑ר גַּם־הַצֹּ֤אן וְהַבָּקָר֙ אַל־יִרְע֔וּ אֶל־מ֖וּל הָהָ֥ר הַהֽוּא׃ \Rashi{\textbf{ואיש לא יעלה עמך.}הראשונות ע"י שהיו בתשואות וקולות וקהלות, שלטה בהן עין רעה – אין לך יפה מן הצניעות (שם):}%
% ---------- פרק 34 | פסוק 4 ----------
 \textbf{ד} וַיִּפְסֹ֡ל שְׁנֵֽי־לֻחֹ֨ת אֲבָנִ֜ים כָּרִאשֹׁנִ֗ים וַיַּשְׁכֵּ֨ם מֹשֶׁ֤ה בַבֹּ֙קֶר֙ וַיַּ֙עַל֙ אֶל־הַ֣ר סִינַ֔י כַּאֲשֶׁ֛ר צִוָּ֥ה יְהֹוָ֖ה אֹת֑וֹ וַיִּקַּ֣ח בְּיָד֔וֹ שְׁנֵ֖י לֻחֹ֥ת אֲבָנִֽים׃ %
% ---------- פרק 34 | פסוק 5 ----------
 \textbf{ה} וַיֵּ֤רֶד יְהֹוָה֙ בֶּֽעָנָ֔ן וַיִּתְיַצֵּ֥ב עִמּ֖וֹ שָׁ֑ם וַיִּקְרָ֥א בְשֵׁ֖ם יְהֹוָֽה׃ \Rashi{\textbf{ויקרא בשם ה'.}מתרגמינן וקרא בשמא דה':}%
% ---------- פרק 34 | פסוק 6 ----------
 \textbf{ו} וַיַּעֲבֹ֨ר יְהֹוָ֥ה ׀ עַל־פָּנָיו֮ וַיִּקְרָא֒ יְהֹוָ֣ה ׀ יְהֹוָ֔ה אֵ֥ל רַח֖וּם וְחַנּ֑וּן אֶ֥רֶךְ אַפַּ֖יִם וְרַב־חֶ֥סֶד וֶאֱמֶֽת׃ \Rashi{\textbf{ה' ה'.}מדת רחמים היא, אחת קדם שיחטא, ואחת אחר שיחטא וישוב (ראש השנה י"ז): אל. אף זו מדת רחמים, וכן הוא אומר אלי אלי למה עזבתני (תהילים כ"ב) – ואין לומר למדת הדין למה עזבתני, כך מצאתי במכילתא: ארך אפים. מאריך אפו ואינו ממהר לפרע, שמא יעשה תשובה: ורב חסד. לצריכים חסד – שאין להם זכיות כל כך: ואמת. לשלם שכר טוב לעושי רצונו:}%
% ---------- פרק 34 | פסוק 7 ----------
 \textbf{ז} נֹצֵ֥ר*(בספרי תימן נֹצֵ֥ר בנו״ן רגילה) חֶ֙סֶד֙ לָאֲלָפִ֔ים נֹשֵׂ֥א עָוֺ֛ן וָפֶ֖שַׁע וְחַטָּאָ֑ה וְנַקֵּה֙ לֹ֣א יְנַקֶּ֔ה פֹּקֵ֣ד ׀ עֲוֺ֣ן אָב֗וֹת עַל־בָּנִים֙ וְעַל־בְּנֵ֣י בָנִ֔ים עַל־שִׁלֵּשִׁ֖ים וְעַל־רִבֵּעִֽים׃ \Rashi{\textbf{נצר חסד.}שהאדם עושה לפניו: לאלפים. לשני אלפים דורות: עון ופשע. עונות — אלו הזדונות, פשעים — אלו המרדים שאדם עושה להכעיס: ונקה לא ינקה. לפי פשוטו משמע שאינו מותר על העון לגמרי, אלא נפרע ממנו מעט מעט, ורבותינו דרשו מנקה הוא לשבים ולא ינקה לשאינן שבים (יומא פ"ו): פקד עון אבות על בנים. כשאוחזים מעשה אבותיהם בידיהם, שכבר פרש במקרא אחר לשנאי (שמות כ'): ועל רבעים. דור רביעי; נמצאת מדה טובה מרבה על מדת פרענות אחת לחמש מאות, שבמדה טובה הוא אומר נצר חסד לאלפים (תוספתא סוטה ד'):}%
% ---------- פרק 34 | פסוק 8 ----------
 \textbf{ח} וַיְמַהֵ֖ר מֹשֶׁ֑ה וַיִּקֹּ֥ד אַ֖רְצָה וַיִּשְׁתָּֽחוּ׃ \Rashi{\textbf{וימהר משה.}כשראה משה שכינה עוברת, ושמע קול הקריאה, מיד וישתחו:}%
% ---------- פרק 34 | פסוק 9 ----------
 \textbf{ט} וַיֹּ֡אמֶר אִם־נָא֩ מָצָ֨אתִי חֵ֤ן בְּעֵינֶ֙יךָ֙ אֲדֹנָ֔י יֵֽלֶךְ־נָ֥א אֲדֹנָ֖י בְּקִרְבֵּ֑נוּ כִּ֤י עַם־קְשֵׁה־עֹ֙רֶף֙ ה֔וּא וְסָלַחְתָּ֛ לַעֲוֺנֵ֥נוּ וּלְחַטָּאתֵ֖נוּ וּנְחַלְתָּֽנוּ׃ \Rashi{\textbf{ילך נא אדני בקרבנו.}כמו שהבטחת, מאחר שאתה נושא עון; ואם עם קשה ערף הוא וימרו בך ואמרת על זאת פן אכלך בדרך, אתה תסלח לעונינו וגומר; יש כי במקום אם: ונחלתנו. ותתננו לך לנחלה מיחדת, זו היא בקשת ונפלינו אני ועמך – שלא תשרה שכינתך על האמות:}%
% ---------- פרק 34 | פסוק 10 ----------
 \textbf{י} וַיֹּ֗אמֶר הִנֵּ֣ה אָנֹכִי֮ כֹּרֵ֣ת בְּרִית֒ נֶ֤גֶד כׇּֽל־עַמְּךָ֙ אֶעֱשֶׂ֣ה נִפְלָאֹ֔ת אֲשֶׁ֛ר לֹֽא־נִבְרְא֥וּ בְכׇל־הָאָ֖רֶץ וּבְכׇל־הַגּוֹיִ֑ם וְרָאָ֣ה כׇל־הָ֠עָ֠ם אֲשֶׁר־אַתָּ֨ה בְקִרְבּ֜וֹ אֶת־מַעֲשֵׂ֤ה יְהֹוָה֙ כִּֽי־נוֹרָ֣א ה֔וּא אֲשֶׁ֥ר אֲנִ֖י עֹשֶׂ֥ה עִמָּֽךְ׃ \Rashi{\textbf{כרת ברית.}על זאת: אעשה נפלאת. לשון ונפלינו, שתהיו מבדלים בזו מכל האמות, שלא תשרה שכינתי עליהם:}%
% ---------- פרק 34 | פסוק 11 ----------
 \textbf{יא} שְׁמׇ֨ר־לְךָ֔ אֵ֛ת אֲשֶׁ֥ר אָנֹכִ֖י מְצַוְּךָ֣ הַיּ֑וֹם הִנְנִ֧י גֹרֵ֣שׁ מִפָּנֶ֗יךָ אֶת־הָאֱמֹרִי֙ וְהַֽכְּנַעֲנִ֔י וְהַחִתִּי֙ וְהַפְּרִזִּ֔י וְהַחִוִּ֖י וְהַיְבוּסִֽי׃ \Rashi{\textbf{את האמרי וגו'.}ו' אמות יש כאן, כי הגרגשי עמד ופנה מפניהם:}%
% ---------- פרק 34 | פסוק 12 ----------
 \textbf{יב} הִשָּׁ֣מֶר לְךָ֗ פֶּן־תִּכְרֹ֤ת בְּרִית֙ לְיוֹשֵׁ֣ב הָאָ֔רֶץ אֲשֶׁ֥ר אַתָּ֖ה בָּ֣א עָלֶ֑יהָ פֶּן־יִהְיֶ֥ה לְמוֹקֵ֖שׁ בְּקִרְבֶּֽךָ׃ %
% ---------- פרק 34 | פסוק 13 ----------
 \textbf{יג} כִּ֤י אֶת־מִזְבְּחֹתָם֙ תִּתֹּצ֔וּן וְאֶת־מַצֵּבֹתָ֖ם תְּשַׁבֵּר֑וּן וְאֶת־אֲשֵׁרָ֖יו תִּכְרֹתֽוּן׃ \Rashi{\textbf{אשריו.}הוא אילן שעובדים אותו:}%
% ---------- פרק 34 | פסוק 14 ----------
 \textbf{יד} כִּ֛י לֹ֥א תִֽשְׁתַּחֲוֶ֖ה לְאֵ֣ל אַחֵ֑ר כִּ֤י יְהֹוָה֙ קַנָּ֣א שְׁמ֔וֹ אֵ֥ל קַנָּ֖א הֽוּא׃ \Rashi{\textbf{קנא שמו.}מקנא לפרע ואינו מותר, וזהו כל לשון קנאה – אוחז בנצחונו ופורע מאויביו:}%
% ---------- פרק 34 | פסוק 15 ----------
 \textbf{טו} פֶּן־תִּכְרֹ֥ת בְּרִ֖ית לְיוֹשֵׁ֣ב הָאָ֑רֶץ וְזָנ֣וּ ׀ אַחֲרֵ֣י אֱלֹֽהֵיהֶ֗ם וְזָבְחוּ֙ לֵאלֹ֣הֵיהֶ֔ם וְקָרָ֣א לְךָ֔ וְאָכַלְתָּ֖ מִזִּבְחֽוֹ׃ \Rashi{\textbf{ואכלת מזבחו.}כסבור אתה שאין ענש באכילתו, ואני מעלה עליך כמודה בעבודתם, שמתוך כך אתה בא ולוקח מבנותיו לבניך (עבודה זרה ח'):}%
% ---------- פרק 34 | פסוק 16 ----------
 \textbf{טז} וְלָקַחְתָּ֥ מִבְּנֹתָ֖יו לְבָנֶ֑יךָ וְזָנ֣וּ בְנֹתָ֗יו אַחֲרֵי֙ אֱלֹ֣הֵיהֶ֔ן וְהִזְנוּ֙ אֶת־בָּנֶ֔יךָ אַחֲרֵ֖י אֱלֹהֵיהֶֽן׃ %
% ---------- פרק 34 | פסוק 17 ----------
 \textbf{יז} אֱלֹהֵ֥י מַסֵּכָ֖ה לֹ֥א תַעֲשֶׂה־לָּֽךְ׃ %
% ---------- פרק 34 | פסוק 18 ----------
 \textbf{יח} אֶת־חַ֣ג הַמַּצּוֹת֮ תִּשְׁמֹר֒ שִׁבְעַ֨ת יָמִ֜ים תֹּאכַ֤ל מַצּוֹת֙ אֲשֶׁ֣ר צִוִּיתִ֔ךָ לְמוֹעֵ֖ד חֹ֣דֶשׁ הָאָבִ֑יב כִּ֚י בְּחֹ֣דֶשׁ הָֽאָבִ֔יב יָצָ֖אתָ מִמִּצְרָֽיִם׃ \Rashi{\textbf{חדש האביב.}חדש הבכור, שהתבואה מתבכרת בבשולה:}%
% ---------- פרק 34 | פסוק 19 ----------
 \textbf{יט} כׇּל־פֶּ֥טֶר רֶ֖חֶם לִ֑י וְכׇֽל־מִקְנְךָ֙ תִּזָּכָ֔ר פֶּ֖טֶר שׁ֥וֹר וָשֶֽׂה׃ \Rashi{\textbf{כל פטר רחם לי.}באדם: וכל מקנך וגו'. וכל מקנך אשר תזכר בפטר שור ושה – אשר יפטר זכר את רחמה: פטר. לשון פתיחה, וכן פוטר מים ראשית מדון (משלי י"ז); תי"ו של תזכר לשון נקבה היא, מוסב על היולדת:}%
% ---------- פרק 34 | פסוק 20 ----------
 \textbf{כ} וּפֶ֤טֶר חֲמוֹר֙ תִּפְדֶּ֣ה בְשֶׂ֔ה וְאִם־לֹ֥א תִפְדֶּ֖ה וַעֲרַפְתּ֑וֹ כֹּ֣ל בְּכ֤וֹר בָּנֶ֙יךָ֙ תִּפְדֶּ֔ה וְלֹֽא־יֵרָא֥וּ פָנַ֖י רֵיקָֽם׃ \Rashi{\textbf{ופטר חמור.}ולא שאר בהמה טמאה: תפדה בשה. נותן שה לכהן והוא חלין ביד כהן, ופטר חמור מתר בעבודה לבעלים: וערפתו. עורפו בקופיץ; הוא הפסיד ממון כהן, לפיכך יפסד ממונו: כל בכור בניך תפדה. חמשה סלעים פדיונו קצוב, שנאמר ופדויו מבן חדש תפדה (במדבר י"ח): ולא יראו פני רקם. לפי פשוטו של מקרא דבר בפני עצמו הוא, ואינו מוסב על הבכור – שאין במצות בכור ראית פנים – אלא אזהרה אחרת היא, וכשתעלו לרגל לראות לא יראו פני ריקם – מצוה עליכם להביא עולת ראית פנים. ולפי מדרש בריתא מקרא יתר הוא, ומפנה לגזרה שוה, ללמד על הענקתו של עבד עברי שהוא חמשה סלעים מכל מין ומין, כפדיון בכור; במסכת קדושין (דף י"ז):}%
% ---------- פרק 34 | פסוק 21 ----------
 \textbf{כא} שֵׁ֤שֶׁת יָמִים֙ תַּעֲבֹ֔ד וּבַיּ֥וֹם הַשְּׁבִיעִ֖י תִּשְׁבֹּ֑ת בֶּחָרִ֥ישׁ וּבַקָּצִ֖יר תִּשְׁבֹּֽת׃ \Rashi{\textbf{בחריש ובקציר.}למה נזכר חריש וקציר? יש מרבותינו אומרים על חריש של ערב שביעית הנכנס לשביעית, וקציר של שביעית היוצא למוצאי שביעית, ללמדך שמוסיפין מחל על הקדש, וכך משמעו: ששת ימים תעבד וביום השביעי תשבת, ועבודת ו' הימים שהתרתי לך, יש שנה שהחריש והקציר אסור, ואין צרך לומר חריש וקציר של שביעית, שהרי כבר נאמר שדך לא תזרע וגו' (ויקרא כ"ה), ויש אומרים שאינו מדבר אלא בשבת, וחריש וקציר שהזכר בו לומר לך, מה חריש רשות אף קציר רשות, יצא קציר העמר שהוא מצוה ודוחה את השבת (ראש השנה ט'):}%
% ---------- פרק 34 | פסוק 22 ----------
 \textbf{כב} וְחַ֤ג שָׁבֻעֹת֙ תַּעֲשֶׂ֣ה לְךָ֔ בִּכּוּרֵ֖י קְצִ֣יר חִטִּ֑ים וְחַג֙ הָֽאָסִ֔יף תְּקוּפַ֖ת הַשָּׁנָֽה׃ \Rashi{\textbf{בכורי קציר חטים.}שאתה מביא בו שתי הלחם מן החטים: בכורי. שהיא מנחה ראשונה הבאה מן החדש של חטים למקדש, כי מנחת העמר הבאה בפסח, מן השעורים היא (מנחות פ"ד): וחג האסיף. בזמן שאתה אוסף תבואתך מן השדה לבית; אסיפה זו לשון הכנסה לבית, כמו ואספתו אל תוך ביתך (דברים כ"ב): תקופת השנה. שהיא בחזרת השנה – בתחלת השנה הבאה: תקופת. לשון מסבה והקפה:}%
% ---------- פרק 34 | פסוק 23 ----------
 \textbf{כג} שָׁלֹ֥שׁ פְּעָמִ֖ים בַּשָּׁנָ֑ה יֵרָאֶה֙ כׇּל־זְכ֣וּרְךָ֔ אֶת־פְּנֵ֛י הָֽאָדֹ֥ן ׀ יְהֹוָ֖ה אֱלֹהֵ֥י יִשְׂרָאֵֽל׃ \Rashi{\textbf{כל זכורך.}כל הזכרים שבך; הרבה מצות בתורה נאמרו ונכפלו – ויש מהם שלש פעמים וארבע – לחיב ולענש על מנין לאוין שבהם ועל מנין עשה שבהם:}%
% ---------- פרק 34 | פסוק 24 ----------
 \textbf{כד} כִּֽי־אוֹרִ֤ישׁ גּוֹיִם֙ מִפָּנֶ֔יךָ וְהִרְחַבְתִּ֖י אֶת־גְּבֻלֶ֑ךָ וְלֹא־יַחְמֹ֥ד אִישׁ֙ אֶֽת־אַרְצְךָ֔ בַּעֲלֹֽתְךָ֗ לֵרָאוֹת֙ אֶת־פְּנֵי֙ יְהֹוָ֣ה אֱלֹהֶ֔יךָ שָׁלֹ֥שׁ פְּעָמִ֖ים בַּשָּׁנָֽה׃ \Rashi{\textbf{אוריש.}כתרגומו אתריך, וכן החל רש (דברים ב'), וכן ויורש את האמרי (במדבר כ"א) – לשון גרושין: והרחבתי את גבלך. ואתה רחוק מבית הבחירה, ואינך יכול לראות לפני תמיד, לכך אני קובע לך שלש רגלים הללו:}%
% ---------- פרק 34 | פסוק 25 ----------
 \textbf{כה} לֹֽא־תִשְׁחַ֥ט עַל־חָמֵ֖ץ דַּם־זִבְחִ֑י וְלֹא־יָלִ֣ין לַבֹּ֔קֶר זֶ֖בַח חַ֥ג הַפָּֽסַח׃ \Rashi{\textbf{לא תשחט וגו'.}לא תשחט את הפסח ועדין חמץ קים; אזהרה לשוחט או לזורק או לאחד מבני חבורה (פסח' ס"ג): ולא ילין. כתרגומו; אין לינה מועלת בראש המזבח ואין לינה אלא בעמוד השחר (זבחים פ"ז): זבח חג הפסח. אמוריו, ומכאן אתה למד לכל הקטר חלבים ואברים:}%
% ---------- פרק 34 | פסוק 26 ----------
 \textbf{כו} רֵאשִׁ֗ית בִּכּוּרֵי֙ אַדְמָ֣תְךָ֔ תָּבִ֕יא בֵּ֖ית יְהֹוָ֣ה אֱלֹהֶ֑יךָ לֹא־תְבַשֵּׁ֥ל גְּדִ֖י בַּחֲלֵ֥ב אִמּֽוֹ׃ \{פ\} \Rashi{\textbf{ראשית בכורי אדמתך.}משבעת המינין האמורים בשבח ארצך, ארץ חטה ושערה וגפן וגו' (דברים ח') – ודבש הוא דבש תמרים: לא תבשל גדי. אזהרה לבשר בחלב, ושלש פעמים כתוב בתורה, אחד לאכילה ואחד להנאה ואחד לאסור בשול (חולין קט"ו): גדי. כל ולד רך במשמע, ואף עגל וכבש; ממה שהצרך לפרש בכמה מקומות "גדי עזים" למדת שגדי סתם כל יונקים במשמע (שם קי"ג): בחלב אמו. פרט לעוף שאין לו חלב, שאין אסורו מן התורה אלא מדברי סופרים (שם):}%
% ---------- פרק 34 | פסוק 27 ----------
 \textbf{כז} וַיֹּ֤אמֶר יְהֹוָה֙ אֶל־מֹשֶׁ֔ה כְּתׇב־לְךָ֖ אֶת־הַדְּבָרִ֣ים הָאֵ֑לֶּה כִּ֞י עַל־פִּ֣י ׀ הַדְּבָרִ֣ים הָאֵ֗לֶּה כָּרַ֧תִּי אִתְּךָ֛ בְּרִ֖ית וְאֶת־יִשְׂרָאֵֽל׃ \Rashi{\textbf{את הדברים האלה.}ולא אתה רשאי לכתב תורה שבעל פה (גיטין ס'):}%
% ---------- פרק 34 | פסוק 28 ----------
 \textbf{כח} וַֽיְהִי־שָׁ֣ם עִם־יְהֹוָ֗ה אַרְבָּעִ֥ים יוֹם֙ וְאַרְבָּעִ֣ים לַ֔יְלָה לֶ֚חֶם לֹ֣א אָכַ֔ל וּמַ֖יִם לֹ֣א שָׁתָ֑ה וַיִּכְתֹּ֣ב עַל־הַלֻּחֹ֗ת אֵ֚ת דִּבְרֵ֣י הַבְּרִ֔ית עֲשֶׂ֖רֶת הַדְּבָרִֽים׃ %
% ---------- פרק 34 | פסוק 29 ----------
 \textbf{כט} וַיְהִ֗י בְּרֶ֤דֶת מֹשֶׁה֙ מֵהַ֣ר סִינַ֔י וּשְׁנֵ֨י לֻחֹ֤ת הָֽעֵדֻת֙ בְּיַד־מֹשֶׁ֔ה בְּרִדְתּ֖וֹ מִן־הָהָ֑ר וּמֹשֶׁ֣ה לֹֽא־יָדַ֗ע כִּ֥י קָרַ֛ן ע֥וֹר פָּנָ֖יו בְּדַבְּר֥וֹ אִתּֽוֹ׃ \Rashi{\textbf{ויהי ברדת משה.}כשהביא לוחות אחרונות ביום הכפורים: כי קרן. לשון קרנים, שהאור מבהיק ובולט כמין קרן; ומהיכן זכה משה לקרני ההוד? רבותינו אמרו מן המערה, שנתן הקדוש ברוך הוא ידו על פניו, שנאמר ושכתי כפי (תנחומא):}%
% ---------- פרק 34 | פסוק 30 ----------
 \textbf{ל} וַיַּ֨רְא אַהֲרֹ֜ן וְכׇל־בְּנֵ֤י יִשְׂרָאֵל֙ אֶת־מֹשֶׁ֔ה וְהִנֵּ֥ה קָרַ֖ן ע֣וֹר פָּנָ֑יו וַיִּֽירְא֖וּ מִגֶּ֥שֶׁת אֵלָֽיו׃ \Rashi{\textbf{וייראו מגשת אליו.}בא וראה כמה גדול כחה של עברה שעד שלא פשטו ידיהם בעברה מהו אומר? ומראה כבוד ה' כאש אכלת בראש ההר לעיני בני ישראל (שמות כ"ד) – ולא יראים ולא מזדעזעים, ומשעשו את העגל אף מקרני הודו של משה היו מרתיעים ומזדעזעים (ספרי):}%
% ---------- פרק 34 | פסוק 31 ----------
 \textbf{לא} וַיִּקְרָ֤א אֲלֵהֶם֙ מֹשֶׁ֔ה וַיָּשֻׁ֧בוּ אֵלָ֛יו אַהֲרֹ֥ן וְכׇל־הַנְּשִׂאִ֖ים בָּעֵדָ֑ה וַיְדַבֵּ֥ר מֹשֶׁ֖ה אֲלֵהֶֽם׃ \Rashi{\textbf{הנשאים בעדה.}כמו נשיאי העדה: וידבר משה אלהם. שליחותו של מקום, ולשון הוה הוא כל הענין הזה:}%
% ---------- פרק 34 | פסוק 32 ----------
 \textbf{לב} וְאַחֲרֵי־כֵ֥ן נִגְּשׁ֖וּ כׇּל־בְּנֵ֣י יִשְׂרָאֵ֑ל וַיְצַוֵּ֕ם אֵת֩ כׇּל־אֲשֶׁ֨ר דִּבֶּ֧ר יְהֹוָ֛ה אִתּ֖וֹ בְּהַ֥ר סִינָֽי׃ \Rashi{\textbf{ואחרי כן נגשו.}אחר שלמד לזקנים, חוזר ומלמד הפרשה או ההלכה לישראל; תנו רבנן כיצד סדר המשנה? משה היה למד מפי הגבורה, נכנס אהרן שנה לו משה פרקו, נסתלק אהרן וישב לו לשמאל משה, נכנסו בניו, שנה להם משה פרקם נסתלקו הן, ישב אלעזר לימין משה ואיתמר לשמאל אהרן, נכנסו זקנים, שנה להם משה פרקם, נסתלקו זקנים ישבו לצדדין, נכנסו כל העם, שנה להם משה פרקם; נמצא ביד כל העם אחד, ביד הזקנים שנים, ביד בני אהרן שלשה, ביד אהרן ארבעה וכו'. כדאיתא בערובין (דף נ"ד):}%
% ---------- פרק 34 | פסוק 33 ----------
 \textbf{לג} וַיְכַ֣ל מֹשֶׁ֔ה מִדַּבֵּ֖ר אִתָּ֑ם וַיִּתֵּ֥ן עַל־פָּנָ֖יו מַסְוֶֽה׃ \Rashi{\textbf{ויתן על פניו מסוה.}כתרגומו בית אפי, לשון ארמי הוא, בתלמוד (כתובות ס"ב) סוי לבה, ועוד בכתבות (דף ס'), הוה קא מסוה לאפה, לשון הבטה – היה מסתכל בה, אף כאן מסוה בגד הנתן כנגד הפרצוף ובית העינים; ולכבוד קרני ההוד – שלא יזונו הכל מהם – היה נותן המסוה כנגדן, ונוטלו בשעה שהיה מדבר עם ישראל, ובשעה שהמקום נדבר עמו עד צאתו, ובצאתו יצא בלא מסוה:}%
% ---------- פרק 34 | פסוק 34 ----------
 \textbf{לד} וּבְבֹ֨א מֹשֶׁ֜ה לִפְנֵ֤י יְהֹוָה֙ לְדַבֵּ֣ר אִתּ֔וֹ יָסִ֥יר אֶת־הַמַּסְוֶ֖ה עַד־צֵאת֑וֹ וְיָצָ֗א וְדִבֶּר֙ אֶל־בְּנֵ֣י יִשְׂרָאֵ֔ל אֵ֖ת אֲשֶׁ֥ר יְצֻוֶּֽה׃ \Rashi{\textbf{ודבר אל בני ישראל.}וראו קרני ההוד בפניו; וכשהוא מסתלק מהם …}%
% ---------- פרק 34 | פסוק 35 ----------
 \textbf{לה} וְרָא֤וּ בְנֵֽי־יִשְׂרָאֵל֙ אֶת־פְּנֵ֣י מֹשֶׁ֔ה כִּ֣י קָרַ֔ן ע֖וֹר פְּנֵ֣י מֹשֶׁ֑ה וְהֵשִׁ֨יב מֹשֶׁ֤ה אֶת־הַמַּסְוֶה֙ עַל־פָּנָ֔יו עַד־בֹּא֖וֹ לְדַבֵּ֥ר אִתּֽוֹ׃ \{ס\} \Rashi{\textbf{והשיב משה את המסוה על פניו עד באו לדבר אתו.}וכשבא לדבר אתו נוטלו מעל פניו:}%
\addparasha{חומש שמות | פרשת ויקהל}
% ---------- פרק 35 | פסוק 1 ----------
 \textbf{א} וַיַּקְהֵ֣ל מֹשֶׁ֗ה אֶֽת־כׇּל־עֲדַ֛ת בְּנֵ֥י יִשְׂרָאֵ֖ל וַיֹּ֣אמֶר אֲלֵהֶ֑ם אֵ֚לֶּה הַדְּבָרִ֔ים אֲשֶׁר־צִוָּ֥ה יְהֹוָ֖ה לַעֲשֹׂ֥ת אֹתָֽם׃ \Rashi{\textbf{ויקהל משה.}למחרת יום הכפורים כשירד מן ההר, והוא לשון הפעיל, שאינו אוסף אנשים בידים, אלא הן נאספין על פי דבורו, ותרגומו ואכניש:}%
% ---------- פרק 35 | פסוק 2 ----------
 \textbf{ב} שֵׁ֣שֶׁת יָמִים֮ תֵּעָשֶׂ֣ה מְלָאכָה֒ וּבַיּ֣וֹם הַשְּׁבִיעִ֗י יִהְיֶ֨ה לָכֶ֥ם קֹ֛דֶשׁ שַׁבַּ֥ת שַׁבָּת֖וֹן לַיהֹוָ֑ה כׇּל־הָעֹשֶׂ֥ה ב֛וֹ מְלָאכָ֖ה יוּמָֽת׃ \Rashi{\textbf{ששת ימים.}הקדים להם אזהרת שבת לצווי מלאכת המשכן, לומר שאינו דוחה את השבת (מכילתא):}%
% ---------- פרק 35 | פסוק 3 ----------
 \textbf{ג} לֹא־תְבַעֲר֣וּ אֵ֔שׁ בְּכֹ֖ל מֹשְׁבֹֽתֵיכֶ֑ם בְּי֖וֹם הַשַּׁבָּֽת׃ \{פ\} \Rashi{\textbf{לא תבערו אש.}יש מרבותינו אומרים הבערה ללאו יצאת, ויש אומרים לחלק יצאת (שבת ע'):}%
% ---------- פרק 35 | פסוק 4 ----------
 \textbf{ד} וַיֹּ֣אמֶר מֹשֶׁ֔ה אֶל־כׇּל־עֲדַ֥ת בְּנֵֽי־יִשְׂרָאֵ֖ל לֵאמֹ֑ר זֶ֣ה הַדָּבָ֔ר אֲשֶׁר־צִוָּ֥ה יְהֹוָ֖ה לֵאמֹֽר׃ \Rashi{\textbf{זה הדבר אשר צוה ה'.}לי לאמר לכם:}%
% ---------- פרק 35 | פסוק 5 ----------
 \textbf{ה} קְח֨וּ מֵֽאִתְּכֶ֤ם תְּרוּמָה֙ לַֽיהֹוָ֔ה כֹּ֚ל נְדִ֣יב לִבּ֔וֹ יְבִיאֶ֕הָ אֵ֖ת תְּרוּמַ֣ת יְהֹוָ֑ה זָהָ֥ב וָכֶ֖סֶף וּנְחֹֽשֶׁת׃ \Rashi{\textbf{נדיב לבו.}על שם שלבו נודבו, קרוי נדיב לב. כבר פרשתי נדבת המשכן ומלאכתו במקום צואתן:}%
% ---------- פרק 35 | פסוק 6 ----------
 \textbf{ו} וּתְכֵ֧לֶת וְאַרְגָּמָ֛ן וְתוֹלַ֥עַת שָׁנִ֖י וְשֵׁ֥שׁ וְעִזִּֽים׃ %
% ---------- פרק 35 | פסוק 7 ----------
 \textbf{ז} וְעֹרֹ֨ת אֵילִ֧ם מְאׇדָּמִ֛ים וְעֹרֹ֥ת תְּחָשִׁ֖ים וַעֲצֵ֥י שִׁטִּֽים׃ %
% ---------- פרק 35 | פסוק 8 ----------
 \textbf{ח} וְשֶׁ֖מֶן לַמָּא֑וֹר וּבְשָׂמִים֙ לְשֶׁ֣מֶן הַמִּשְׁחָ֔ה וְלִקְטֹ֖רֶת הַסַּמִּֽים׃ %
% ---------- פרק 35 | פסוק 9 ----------
 \textbf{ט} וְאַ֨בְנֵי־שֹׁ֔הַם וְאַבְנֵ֖י מִלֻּאִ֑ים לָאֵפ֖וֹד וְלַחֹֽשֶׁן׃ %
% ---------- פרק 35 | פסוק 10 ----------
 \textbf{י} וְכׇל־חֲכַם־לֵ֖ב בָּכֶ֑ם יָבֹ֣אוּ וְיַעֲשׂ֔וּ אֵ֛ת כׇּל־אֲשֶׁ֥ר צִוָּ֖ה יְהֹוָֽה׃ %
% ---------- פרק 35 | פסוק 11 ----------
 \textbf{יא} אֶ֨ת־הַמִּשְׁכָּ֔ן אֶֽת־אׇהֳל֖וֹ וְאֶת־מִכְסֵ֑הוּ אֶת־קְרָסָיו֙ וְאֶת־קְרָשָׁ֔יו אֶת־בְּרִיחָ֕ו אֶת־עַמֻּדָ֖יו וְאֶת־אֲדָנָֽיו׃ \Rashi{\textbf{את המשכן.}יריעות התחתונות הנראות בתוכו קרויים משכן: את אהלו. היא אהל יריעות עזים, העשוי לגג: ואת מכסהו. מכסה עורות אילים והתחשים:}%
% ---------- פרק 35 | פסוק 12 ----------
 \textbf{יב} אֶת־הָאָרֹ֥ן וְאֶת־בַּדָּ֖יו אֶת־הַכַּפֹּ֑רֶת וְאֵ֖ת פָּרֹ֥כֶת הַמָּסָֽךְ׃ \Rashi{\textbf{ואת פרכת המסך.}פרכת המחיצה; כל דבר המגן, בין מלמעלה בין מכנגד, קרוי מסך וסכך, וכן סכת בעדו (איוב א'), הנני סך את דרכך (הושע ב'):}%
% ---------- פרק 35 | פסוק 13 ----------
 \textbf{יג} אֶת־הַשֻּׁלְחָ֥ן וְאֶת־בַּדָּ֖יו וְאֶת־כׇּל־כֵּלָ֑יו וְאֵ֖ת לֶ֥חֶם הַפָּנִֽים׃ \Rashi{\textbf{לחם הפנים.}כבר פרשתי, על שם שהיו לו פנים לכאן ולכאן, שהוא עשוי כמין תבה פרוצה:}%
% ---------- פרק 35 | פסוק 14 ----------
 \textbf{יד} וְאֶת־מְנֹרַ֧ת הַמָּא֛וֹר וְאֶת־כֵּלֶ֖יהָ וְאֶת־נֵרֹתֶ֑יהָ וְאֵ֖ת שֶׁ֥מֶן הַמָּאֽוֹר׃ \Rashi{\textbf{ואת כליה.}מלקחים ומחתות: נרתיה. לוצי"יש בלעז, בזיכים שהשמן והפתילות נתונין בהן: ואת שמן המאור. אף הוא צריך חכמי לב, שהוא משנה משאר שמנים, כמו שמפרש במנחות (דף פ"ו), מגרגרו בראש הזית, והוא כתית וזך:}%
% ---------- פרק 35 | פסוק 15 ----------
 \textbf{טו} וְאֶת־מִזְבַּ֤ח הַקְּטֹ֙רֶת֙ וְאֶת־בַּדָּ֔יו וְאֵת֙ שֶׁ֣מֶן הַמִּשְׁחָ֔ה וְאֵ֖ת קְטֹ֣רֶת הַסַּמִּ֑ים וְאֶת־מָסַ֥ךְ הַפֶּ֖תַח לְפֶ֥תַח הַמִּשְׁכָּֽן׃ \Rashi{\textbf{מסך הפתח.}וילון שלפני המזרח, שלא היו שם קרשים ולא יריעות:}%
% ---------- פרק 35 | פסוק 16 ----------
 \textbf{טז} אֵ֣ת ׀ מִזְבַּ֣ח הָעֹלָ֗ה וְאֶת־מִכְבַּ֤ר הַנְּחֹ֙שֶׁת֙ אֲשֶׁר־ל֔וֹ אֶת־בַּדָּ֖יו וְאֶת־כׇּל־כֵּלָ֑יו אֶת־הַכִּיֹּ֖ר וְאֶת־כַּנּֽוֹ׃ %
% ---------- פרק 35 | פסוק 17 ----------
 \textbf{יז} אֵ֚ת קַלְעֵ֣י הֶחָצֵ֔ר אֶת־עַמֻּדָ֖יו וְאֶת־אֲדָנֶ֑יהָ וְאֵ֕ת מָסַ֖ךְ שַׁ֥עַר הֶחָצֵֽר׃ \Rashi{\textbf{את עמדיו ואת אדניה.}הרי חצר קרוי כאן לשון זכר ולשון נקבה, וכן דברים הרבה: ואת מסך שער החצר. וילון פרוש לצד המזרח עשרים אמה אמצעיות של רחב החצר – שהיה חמשים רחב – וסתומין הימנו לצד צפון ט"ו אמה, וכן לדרום, שנאמר, וחמש עשרה אמה קלעים לכתף:}%
% ---------- פרק 35 | פסוק 18 ----------
 \textbf{יח} אֶת־יִתְדֹ֧ת הַמִּשְׁכָּ֛ן וְאֶת־יִתְדֹ֥ת הֶחָצֵ֖ר וְאֶת־מֵיתְרֵיהֶֽם׃ \Rashi{\textbf{יתדת.}לתקע ולקשור בהם סופי היריעות בארץ, שלא ינועו ברוח: מיתריהם. חבלים לקשר:}%
% ---------- פרק 35 | פסוק 19 ----------
 \textbf{יט} אֶת־בִּגְדֵ֥י הַשְּׂרָ֖ד לְשָׁרֵ֣ת בַּקֹּ֑דֶשׁ אֶת־בִּגְדֵ֤י הַקֹּ֙דֶשׁ֙ לְאַהֲרֹ֣ן הַכֹּהֵ֔ן וְאֶת־בִּגְדֵ֥י בָנָ֖יו לְכַהֵֽן׃ \Rashi{\textbf{בגדי השרד.}לכסות הארון והשלחן, המנורה והמזבחות בשעת סלוק מסעות:}%
% ---------- פרק 35 | פסוק 20 ----------
 \textbf{כ} וַיֵּ֥צְא֛וּ כׇּל־עֲדַ֥ת בְּנֵֽי־יִשְׂרָאֵ֖ל מִלִּפְנֵ֥י מֹשֶֽׁה׃ %
% ---------- פרק 35 | פסוק 21 ----------
 \textbf{כא} וַיָּבֹ֕אוּ כׇּל־אִ֖ישׁ אֲשֶׁר־נְשָׂא֣וֹ לִבּ֑וֹ וְכֹ֡ל אֲשֶׁר֩ נָדְבָ֨ה רוּח֜וֹ אֹת֗וֹ הֵ֠בִ֠יאוּ אֶת־תְּרוּמַ֨ת יְהֹוָ֜ה לִמְלֶ֨אכֶת אֹ֤הֶל מוֹעֵד֙ וּלְכׇל־עֲבֹ֣דָת֔וֹ וּלְבִגְדֵ֖י הַקֹּֽדֶשׁ׃ %
% ---------- פרק 35 | פסוק 22 ----------
 \textbf{כב} וַיָּבֹ֥אוּ הָאֲנָשִׁ֖ים עַל־הַנָּשִׁ֑ים כֹּ֣ל ׀ נְדִ֣יב לֵ֗ב הֵ֠בִ֠יאוּ חָ֣ח וָנֶ֜זֶם וְטַבַּ֤עַת וְכוּמָז֙ כׇּל־כְּלִ֣י זָהָ֔ב וְכׇל־אִ֕ישׁ אֲשֶׁ֥ר הֵנִ֛יף תְּנוּפַ֥ת זָהָ֖ב לַיהֹוָֽה׃ \Rashi{\textbf{(על הנשים.}עם הנשים וסמוכין אליהן): חח. הוא תכשיט של זהב עגל, נתון על הזרוע, והוא הצמיד: וכומז. כלי זהב הוא נתון כנגד אותו מקום לאשה, ורבותינו פרשו שם כומז, כאן מקום זמה (שבת ס"ד):}%
% ---------- פרק 35 | פסוק 23 ----------
 \textbf{כג} וְכׇל־אִ֞ישׁ אֲשֶׁר־נִמְצָ֣א אִתּ֗וֹ תְּכֵ֧לֶת וְאַרְגָּמָ֛ן וְתוֹלַ֥עַת שָׁנִ֖י וְשֵׁ֣שׁ וְעִזִּ֑ים וְעֹרֹ֨ת אֵילִ֧ם מְאׇדָּמִ֛ים וְעֹרֹ֥ת תְּחָשִׁ֖ים הֵבִֽיאוּ׃ \Rashi{\textbf{וכל איש אשר נמצא אתו.}תכלת או ארגמן או תולעת שני או עורות אילים או תחשים כלם הביאו:}%
% ---------- פרק 35 | פסוק 24 ----------
 \textbf{כד} כׇּל־מֵרִ֗ים תְּר֤וּמַת כֶּ֙סֶף֙ וּנְחֹ֔שֶׁת הֵבִ֕יאוּ אֵ֖ת תְּרוּמַ֣ת יְהֹוָ֑ה וְכֹ֡ל אֲשֶׁר֩ נִמְצָ֨א אִתּ֜וֹ עֲצֵ֥י שִׁטִּ֛ים לְכׇל־מְלֶ֥אכֶת הָעֲבֹדָ֖ה הֵבִֽיאוּ׃ %
% ---------- פרק 35 | פסוק 25 ----------
 \textbf{כה} וְכׇל־אִשָּׁ֥ה חַכְמַת־לֵ֖ב בְּיָדֶ֣יהָ טָו֑וּ וַיָּבִ֣יאוּ מַטְוֶ֗ה אֶֽת־הַתְּכֵ֙לֶת֙ וְאֶת־הָֽאַרְגָּמָ֔ן אֶת־תּוֹלַ֥עַת הַשָּׁנִ֖י וְאֶת־הַשֵּֽׁשׁ׃ %
% ---------- פרק 35 | פסוק 26 ----------
 \textbf{כו} וְכׇ֨ל־הַנָּשִׁ֔ים אֲשֶׁ֨ר נָשָׂ֥א לִבָּ֛ן אֹתָ֖נָה בְּחׇכְמָ֑ה טָו֖וּ אֶת־הָעִזִּֽים׃ \Rashi{\textbf{טוו את העזים.}היא היתה אמנות יתרה, שמעל גבי העזים טווין אותם (שבת ע"ד):}%
% ---------- פרק 35 | פסוק 27 ----------
 \textbf{כז} וְהַנְּשִׂאִ֣ם הֵבִ֔יאוּ אֵ֚ת אַבְנֵ֣י הַשֹּׁ֔הַם וְאֵ֖ת אַבְנֵ֣י הַמִּלֻּאִ֑ים לָאֵפ֖וֹד וְלַחֹֽשֶׁן׃ \Rashi{\textbf{והנשאם הביאו.}אמר ר' נתן: מה ראו נשיאים להתנדב בחנכת המזבח בתחלה ובמלאכת המשכן לא התנדבו בתחלה? אלא כך אמרו נשיאים, יתנדבו צבור מה שמתנדבין, ומה שמחסרים, אנו משלימין אותו, כיון שהשלימו צבור את הכל – שנ' והמלאכה היתה דים – אמרו נשיאים מה עלינו לעשות? הביאו את אבני השהם וגו', לכך התנדבו בחנכת המזבח תחלה, ולפי שנתעצלו מתחלה נחסרה אות משמם, והנשאם כתיב (ספרי במדבר מ"ה):}%
% ---------- פרק 35 | פסוק 28 ----------
 \textbf{כח} וְאֶת־הַבֹּ֖שֶׂם וְאֶת־הַשָּׁ֑מֶן לְמָא֕וֹר וּלְשֶׁ֙מֶן֙ הַמִּשְׁחָ֔ה וְלִקְטֹ֖רֶת הַסַּמִּֽים׃ %
% ---------- פרק 35 | פסוק 29 ----------
 \textbf{כט} כׇּל־אִ֣ישׁ וְאִשָּׁ֗ה אֲשֶׁ֨ר נָדַ֣ב לִבָּם֮ אֹתָם֒ לְהָבִיא֙ לְכׇל־הַמְּלָאכָ֔ה אֲשֶׁ֨ר צִוָּ֧ה יְהֹוָ֛ה לַעֲשׂ֖וֹת בְּיַד־מֹשֶׁ֑ה הֵבִ֧יאוּ בְנֵי־יִשְׂרָאֵ֛ל נְדָבָ֖ה לַיהֹוָֽה׃ \{פ\} %
% ---------- פרק 35 | פסוק 30 ----------
 \textbf{ל} וַיֹּ֤אמֶר מֹשֶׁה֙ אֶל־בְּנֵ֣י יִשְׂרָאֵ֔ל רְא֛וּ קָרָ֥א יְהֹוָ֖ה בְּשֵׁ֑ם בְּצַלְאֵ֛ל בֶּן־אוּרִ֥י בֶן־ח֖וּר לְמַטֵּ֥ה יְהוּדָֽה׃ \Rashi{\textbf{חור.}בנה של מרים היה:}%
% ---------- פרק 35 | פסוק 31 ----------
 \textbf{לא} וַיְמַלֵּ֥א אֹת֖וֹ ר֣וּחַ אֱלֹהִ֑ים בְּחׇכְמָ֛ה בִּתְבוּנָ֥ה וּבְדַ֖עַת וּבְכׇל־מְלָאכָֽה׃ %
% ---------- פרק 35 | פסוק 32 ----------
 \textbf{לב} וְלַחְשֹׁ֖ב מַֽחֲשָׁבֹ֑ת לַעֲשֹׂ֛ת בַּזָּהָ֥ב וּבַכֶּ֖סֶף וּבַנְּחֹֽשֶׁת׃ %
% ---------- פרק 35 | פסוק 33 ----------
 \textbf{לג} וּבַחֲרֹ֥שֶׁת אֶ֛בֶן לְמַלֹּ֖את וּבַחֲרֹ֣שֶׁת עֵ֑ץ לַעֲשׂ֖וֹת בְּכׇל־מְלֶ֥אכֶת מַחֲשָֽׁבֶת׃ %
% ---------- פרק 35 | פסוק 34 ----------
 \textbf{לד} וּלְהוֹרֹ֖ת נָתַ֣ן בְּלִבּ֑וֹ ה֕וּא וְאׇֽהֳלִיאָ֥ב בֶּן־אֲחִיסָמָ֖ךְ לְמַטֵּה־דָֽן׃ \Rashi{\textbf{ואהליאב.}משבט דן מן הירודין שבשבטים – מבני השפחות – והשוהו המקום לבצלאל למלאכת המשכן והוא מגדולי השבטים, לקים מה שנ' ולא נכר שוע לפני דל (איוב ל"ד):}%
% ---------- פרק 35 | פסוק 35 ----------
 \textbf{לה} מִלֵּ֨א אֹתָ֜ם חׇכְמַת־לֵ֗ב לַעֲשׂוֹת֮ כׇּל־מְלֶ֣אכֶת חָרָ֣שׁ ׀ וְחֹשֵׁב֒ וְרֹקֵ֞ם בַּתְּכֵ֣לֶת וּבָֽאַרְגָּמָ֗ן בְּתוֹלַ֧עַת הַשָּׁנִ֛י וּבַשֵּׁ֖שׁ וְאֹרֵ֑ג עֹשֵׂי֙ כׇּל־מְלָאכָ֔ה וְחֹשְׁבֵ֖י מַחֲשָׁבֹֽת׃ %
% ---------- פרק 36 | פסוק 1 ----------
 \textbf{א} וְעָשָׂה֩ בְצַלְאֵ֨ל וְאׇהֳלִיאָ֜ב וְכֹ֣ל ׀ אִ֣ישׁ חֲכַם־לֵ֗ב אֲשֶׁר֩ נָתַ֨ן יְהֹוָ֜ה חׇכְמָ֤ה וּתְבוּנָה֙ בָּהֵ֔מָּה לָדַ֣עַת לַעֲשֹׂ֔ת אֶֽת־כׇּל־מְלֶ֖אכֶת עֲבֹדַ֣ת הַקֹּ֑דֶשׁ לְכֹ֥ל אֲשֶׁר־צִוָּ֖ה יְהֹוָֽה׃ %
% ---------- פרק 36 | פסוק 2 ----------
 \textbf{ב} וַיִּקְרָ֣א מֹשֶׁ֗ה אֶל־בְּצַלְאֵל֮ וְאֶל־אׇֽהֳלִיאָב֒ וְאֶל֙ כׇּל־אִ֣ישׁ חֲכַם־לֵ֔ב אֲשֶׁ֨ר נָתַ֧ן יְהֹוָ֛ה חׇכְמָ֖ה בְּלִבּ֑וֹ כֹּ֚ל אֲשֶׁ֣ר נְשָׂא֣וֹ לִבּ֔וֹ לְקׇרְבָ֥ה אֶל־הַמְּלָאכָ֖ה לַעֲשֹׂ֥ת אֹתָֽהּ׃ %
% ---------- פרק 36 | פסוק 3 ----------
 \textbf{ג} וַיִּקְח֞וּ מִלִּפְנֵ֣י מֹשֶׁ֗ה אֵ֤ת כׇּל־הַתְּרוּמָה֙ אֲשֶׁ֨ר הֵבִ֜יאוּ בְּנֵ֣י יִשְׂרָאֵ֗ל לִמְלֶ֛אכֶת עֲבֹדַ֥ת הַקֹּ֖דֶשׁ לַעֲשֹׂ֣ת אֹתָ֑הּ וְ֠הֵ֠ם הֵבִ֨יאוּ אֵלָ֥יו ע֛וֹד נְדָבָ֖ה בַּבֹּ֥קֶר בַּבֹּֽקֶר׃ %
% ---------- פרק 36 | פסוק 4 ----------
 \textbf{ד} וַיָּבֹ֙אוּ֙ כׇּל־הַ֣חֲכָמִ֔ים הָעֹשִׂ֕ים אֵ֖ת כׇּל־מְלֶ֣אכֶת הַקֹּ֑דֶשׁ אִֽישׁ־אִ֥ישׁ מִמְּלַאכְתּ֖וֹ אֲשֶׁר־הֵ֥מָּה עֹשִֽׂים׃ %
% ---------- פרק 36 | פסוק 5 ----------
 \textbf{ה} וַיֹּאמְרוּ֙ אֶל־מֹשֶׁ֣ה לֵּאמֹ֔ר מַרְבִּ֥ים הָעָ֖ם לְהָבִ֑יא מִדֵּ֤י הָֽעֲבֹדָה֙ לַמְּלָאכָ֔ה אֲשֶׁר־צִוָּ֥ה יְהֹוָ֖ה לַעֲשֹׂ֥ת אֹתָֽהּ׃ \Rashi{\textbf{מדי העבדה.}יותר מכדי צרך העבודה:}%
% ---------- פרק 36 | פסוק 6 ----------
 \textbf{ו} וַיְצַ֣ו מֹשֶׁ֗ה וַיַּעֲבִ֨ירוּ ק֥וֹל בַּֽמַּחֲנֶה֮ לֵאמֹר֒ אִ֣ישׁ וְאִשָּׁ֗ה אַל־יַעֲשׂוּ־ע֛וֹד מְלָאכָ֖ה לִתְרוּמַ֣ת הַקֹּ֑דֶשׁ וַיִּכָּלֵ֥א הָעָ֖ם מֵהָבִֽיא׃ \Rashi{\textbf{ויכלא.}לשון מניעה:}%
% ---------- פרק 36 | פסוק 7 ----------
 \textbf{ז} וְהַמְּלָאכָ֗ה הָיְתָ֥ה דַיָּ֛ם לְכׇל־הַמְּלָאכָ֖ה לַעֲשׂ֣וֹת אֹתָ֑הּ וְהוֹתֵֽר׃ \{ס\} \Rashi{\textbf{והמלאכה היתה דים לכל המלאכה.}ומלאכת ההבאה היתה דים של עושי המשכן, לכל המלאכה של משכן, לעשות אותה ולהותר: והותר. כמו והכבד את לבו (שמות ח'), והכות את מואב (מלכים ב ג'):}%
% ---------- פרק 36 | פסוק 8 ----------
 \textbf{ח} וַיַּעֲשׂ֨וּ כׇל־חֲכַם־לֵ֜ב בְּעֹשֵׂ֧י הַמְּלָאכָ֛ה אֶת־הַמִּשְׁכָּ֖ן עֶ֣שֶׂר יְרִיעֹ֑ת שֵׁ֣שׁ מׇשְׁזָ֗ר וּתְכֵ֤לֶת וְאַרְגָּמָן֙ וְתוֹלַ֣עַת שָׁנִ֔י כְּרֻבִ֛ים מַעֲשֵׂ֥ה חֹשֵׁ֖ב עָשָׂ֥ה אֹתָֽם׃ %
% ---------- פרק 36 | פסוק 9 ----------
 \textbf{ט} אֹ֜רֶךְ הַיְרִיעָ֣ה הָֽאַחַ֗ת שְׁמֹנֶ֤ה וְעֶשְׂרִים֙ בָּֽאַמָּ֔ה וְרֹ֙חַב֙ אַרְבַּ֣ע בָּֽאַמָּ֔ה הַיְרִיעָ֖ה הָאֶחָ֑ת מִדָּ֥ה אַחַ֖ת לְכׇל־הַיְרִיעֹֽת׃ %
% ---------- פרק 36 | פסוק 10 ----------
 \textbf{י} וַיְחַבֵּר֙ אֶת־חֲמֵ֣שׁ הַיְרִיעֹ֔ת אַחַ֖ת אֶל־אֶחָ֑ת וְחָמֵ֤שׁ יְרִיעֹת֙ חִבַּ֔ר אַחַ֖ת אֶל־אֶחָֽת׃ %
% ---------- פרק 36 | פסוק 11 ----------
 \textbf{יא} וַיַּ֜עַשׂ לֻֽלְאֹ֣ת תְּכֵ֗לֶת עַ֣ל שְׂפַ֤ת הַיְרִיעָה֙ הָֽאֶחָ֔ת מִקָּצָ֖ה בַּמַּחְבָּ֑רֶת כֵּ֤ן עָשָׂה֙ בִּשְׂפַ֣ת הַיְרִיעָ֔ה הַקִּ֣יצוֹנָ֔ה בַּמַּחְבֶּ֖רֶת הַשֵּׁנִֽית׃ %
% ---------- פרק 36 | פסוק 12 ----------
 \textbf{יב} חֲמִשִּׁ֣ים לֻלָאֹ֗ת עָשָׂה֮ בַּיְרִיעָ֣ה הָאֶחָת֒ וַחֲמִשִּׁ֣ים לֻלָאֹ֗ת עָשָׂה֙ בִּקְצֵ֣ה הַיְרִיעָ֔ה אֲשֶׁ֖ר בַּמַּחְבֶּ֣רֶת הַשֵּׁנִ֑ית מַקְבִּילֹת֙ הַלֻּ֣לָאֹ֔ת אַחַ֖ת אֶל־אֶחָֽת׃ %
% ---------- פרק 36 | פסוק 13 ----------
 \textbf{יג} וַיַּ֕עַשׂ חֲמִשִּׁ֖ים קַרְסֵ֣י זָהָ֑ב וַיְחַבֵּ֨ר אֶת־הַיְרִיעֹ֜ת אַחַ֤ת אֶל־אַחַת֙ בַּקְּרָסִ֔ים וַיְהִ֥י הַמִּשְׁכָּ֖ן אֶחָֽד׃ \{פ\} %
% ---------- פרק 36 | פסוק 14 ----------
 \textbf{יד} וַיַּ֙עַשׂ֙ יְרִיעֹ֣ת עִזִּ֔ים לְאֹ֖הֶל עַל־הַמִּשְׁכָּ֑ן עַשְׁתֵּֽי־עֶשְׂרֵ֥ה יְרִיעֹ֖ת עָשָׂ֥ה אֹתָֽם׃ %
% ---------- פרק 36 | פסוק 15 ----------
 \textbf{טו} אֹ֜רֶךְ הַיְרִיעָ֣ה הָאַחַ֗ת שְׁלֹשִׁים֙ בָּֽאַמָּ֔ה וְאַרְבַּ֣ע אַמּ֔וֹת רֹ֖חַב הַיְרִיעָ֣ה הָאֶחָ֑ת מִדָּ֣ה אַחַ֔ת לְעַשְׁתֵּ֥י עֶשְׂרֵ֖ה יְרִיעֹֽת׃ %
% ---------- פרק 36 | פסוק 16 ----------
 \textbf{טז} וַיְחַבֵּ֛ר אֶת־חֲמֵ֥שׁ הַיְרִיעֹ֖ת לְבָ֑ד וְאֶת־שֵׁ֥שׁ הַיְרִיעֹ֖ת לְבָֽד׃ %
% ---------- פרק 36 | פסוק 17 ----------
 \textbf{יז} וַיַּ֜עַשׂ לֻֽלָאֹ֣ת חֲמִשִּׁ֗ים עַ֚ל שְׂפַ֣ת הַיְרִיעָ֔ה הַקִּיצֹנָ֖ה בַּמַּחְבָּ֑רֶת וַחֲמִשִּׁ֣ים לֻלָאֹ֗ת עָשָׂה֙ עַל־שְׂפַ֣ת הַיְרִיעָ֔ה הַחֹבֶ֖רֶת הַשֵּׁנִֽית׃ %
% ---------- פרק 36 | פסוק 18 ----------
 \textbf{יח} וַיַּ֛עַשׂ קַרְסֵ֥י נְחֹ֖שֶׁת חֲמִשִּׁ֑ים לְחַבֵּ֥ר אֶת־הָאֹ֖הֶל לִהְיֹ֥ת אֶחָֽד׃ %
% ---------- פרק 36 | פסוק 19 ----------
 \textbf{יט} וַיַּ֤עַשׂ מִכְסֶה֙ לָאֹ֔הֶל עֹרֹ֥ת אֵילִ֖ם מְאׇדָּמִ֑ים וּמִכְסֵ֛ה עֹרֹ֥ת תְּחָשִׁ֖ים מִלְמָֽעְלָה׃ \{ס\} %
% ---------- פרק 36 | פסוק 20 ----------
 \textbf{כ} וַיַּ֥עַשׂ אֶת־הַקְּרָשִׁ֖ים לַמִּשְׁכָּ֑ן עֲצֵ֥י שִׁטִּ֖ים עֹמְדִֽים׃ %
% ---------- פרק 36 | פסוק 21 ----------
 \textbf{כא} עֶ֥שֶׂר אַמֹּ֖ת אֹ֣רֶךְ הַקָּ֑רֶשׁ וְאַמָּה֙ וַחֲצִ֣י הָֽאַמָּ֔ה רֹ֖חַב הַקֶּ֥רֶשׁ הָאֶחָֽד׃ %
% ---------- פרק 36 | פסוק 22 ----------
 \textbf{כב} שְׁתֵּ֣י יָדֹ֗ת לַקֶּ֙רֶשׁ֙ הָֽאֶחָ֔ד מְשֻׁ֨לָּבֹ֔ת אַחַ֖ת אֶל־אֶחָ֑ת כֵּ֣ן עָשָׂ֔ה לְכֹ֖ל קַרְשֵׁ֥י הַמִּשְׁכָּֽן׃ %
% ---------- פרק 36 | פסוק 23 ----------
 \textbf{כג} וַיַּ֥עַשׂ אֶת־הַקְּרָשִׁ֖ים לַמִּשְׁכָּ֑ן עֶשְׂרִ֣ים קְרָשִׁ֔ים לִפְאַ֖ת נֶ֥גֶב תֵּימָֽנָה׃ %
% ---------- פרק 36 | פסוק 24 ----------
 \textbf{כד} וְאַרְבָּעִים֙ אַדְנֵי־כֶ֔סֶף עָשָׂ֕ה תַּ֖חַת עֶשְׂרִ֣ים הַקְּרָשִׁ֑ים שְׁנֵ֨י אֲדָנִ֜ים תַּֽחַת־הַקֶּ֤רֶשׁ הָאֶחָד֙ לִשְׁתֵּ֣י יְדֹתָ֔יו וּשְׁנֵ֧י אֲדָנִ֛ים תַּֽחַת־הַקֶּ֥רֶשׁ הָאֶחָ֖ד לִשְׁתֵּ֥י יְדֹתָֽיו׃ %
% ---------- פרק 36 | פסוק 25 ----------
 \textbf{כה} וּלְצֶ֧לַע הַמִּשְׁכָּ֛ן הַשֵּׁנִ֖ית לִפְאַ֣ת צָפ֑וֹן עָשָׂ֖ה עֶשְׂרִ֥ים קְרָשִֽׁים׃ %
% ---------- פרק 36 | פסוק 26 ----------
 \textbf{כו} וְאַרְבָּעִ֥ים אַדְנֵיהֶ֖ם כָּ֑סֶף שְׁנֵ֣י אֲדָנִ֗ים תַּ֚חַת הַקֶּ֣רֶשׁ הָאֶחָ֔ד וּשְׁנֵ֣י אֲדָנִ֔ים תַּ֖חַת הַקֶּ֥רֶשׁ הָאֶחָֽד׃ %
% ---------- פרק 36 | פסוק 27 ----------
 \textbf{כז} וּֽלְיַרְכְּתֵ֥י הַמִּשְׁכָּ֖ן יָ֑מָּה עָשָׂ֖ה שִׁשָּׁ֥ה קְרָשִֽׁים׃ %
% ---------- פרק 36 | פסוק 28 ----------
 \textbf{כח} וּשְׁנֵ֤י קְרָשִׁים֙ עָשָׂ֔ה לִמְקֻצְעֹ֖ת הַמִּשְׁכָּ֑ן בַּיַּרְכָתָֽיִם׃ %
% ---------- פרק 36 | פסוק 29 ----------
 \textbf{כט} וְהָי֣וּ תוֹאֲמִם֮ מִלְּמַ֒טָּה֒ וְיַחְדָּ֗ו יִהְי֤וּ תַמִּים֙ אֶל־רֹאשׁ֔וֹ אֶל־הַטַּבַּ֖עַת הָאֶחָ֑ת כֵּ֚ן עָשָׂ֣ה לִשְׁנֵיהֶ֔ם לִשְׁנֵ֖י הַמִּקְצֹעֹֽת׃ %
% ---------- פרק 36 | פסוק 30 ----------
 \textbf{ל} וְהָיוּ֙ שְׁמֹנָ֣ה קְרָשִׁ֔ים וְאַדְנֵיהֶ֣ם כֶּ֔סֶף שִׁשָּׁ֥ה עָשָׂ֖ר אֲדָנִ֑ים שְׁנֵ֤י אֲדָנִים֙ שְׁנֵ֣י אֲדָנִ֔ים תַּ֖חַת הַקֶּ֥רֶשׁ הָאֶחָֽד׃ %
% ---------- פרק 36 | פסוק 31 ----------
 \textbf{לא} וַיַּ֥עַשׂ בְּרִיחֵ֖י עֲצֵ֣י שִׁטִּ֑ים חֲמִשָּׁ֕ה לְקַרְשֵׁ֥י צֶֽלַע־הַמִּשְׁכָּ֖ן הָאֶחָֽת׃ %
% ---------- פרק 36 | פסוק 32 ----------
 \textbf{לב} וַחֲמִשָּׁ֣ה בְרִיחִ֔ם לְקַרְשֵׁ֥י צֶֽלַע־הַמִּשְׁכָּ֖ן הַשֵּׁנִ֑ית וַחֲמִשָּׁ֤ה בְרִיחִם֙ לְקַרְשֵׁ֣י הַמִּשְׁכָּ֔ן לַיַּרְכָתַ֖יִם יָֽמָּה׃ %
% ---------- פרק 36 | פסוק 33 ----------
 \textbf{לג} וַיַּ֖עַשׂ אֶת־הַבְּרִ֣יחַ הַתִּיכֹ֑ן לִבְרֹ֙חַ֙ בְּת֣וֹךְ הַקְּרָשִׁ֔ים מִן־הַקָּצֶ֖ה אֶל־הַקָּצֶֽה׃ %
% ---------- פרק 36 | פסוק 34 ----------
 \textbf{לד} וְֽאֶת־הַקְּרָשִׁ֞ים צִפָּ֣ה זָהָ֗ב וְאֶת־טַבְּעֹתָם֙ עָשָׂ֣ה זָהָ֔ב בָּתִּ֖ים לַבְּרִיחִ֑ם וַיְצַ֥ף אֶת־הַבְּרִיחִ֖ם זָהָֽב׃ %
% ---------- פרק 36 | פסוק 35 ----------
 \textbf{לה} וַיַּ֙עַשׂ֙ אֶת־הַפָּרֹ֔כֶת תְּכֵ֧לֶת וְאַרְגָּמָ֛ן וְתוֹלַ֥עַת שָׁנִ֖י וְשֵׁ֣שׁ מׇשְׁזָ֑ר מַעֲשֵׂ֥ה חֹשֵׁ֛ב עָשָׂ֥ה אֹתָ֖הּ כְּרֻבִֽים׃ %
% ---------- פרק 36 | פסוק 36 ----------
 \textbf{לו} וַיַּ֣עַשׂ לָ֗הּ אַרְבָּעָה֙ עַמּוּדֵ֣י שִׁטִּ֔ים וַיְצַפֵּ֣ם זָהָ֔ב וָוֵיהֶ֖ם זָהָ֑ב וַיִּצֹ֣ק לָהֶ֔ם אַרְבָּעָ֖ה אַדְנֵי־כָֽסֶף׃ %
% ---------- פרק 36 | פסוק 37 ----------
 \textbf{לז} וַיַּ֤עַשׂ מָסָךְ֙ לְפֶ֣תַח הָאֹ֔הֶל תְּכֵ֧לֶת וְאַרְגָּמָ֛ן וְתוֹלַ֥עַת שָׁנִ֖י וְשֵׁ֣שׁ מׇשְׁזָ֑ר מַעֲשֵׂ֖ה רֹקֵֽם׃ %
% ---------- פרק 36 | פסוק 38 ----------
 \textbf{לח} וְאֶת־עַמּוּדָ֤יו חֲמִשָּׁה֙ וְאֶת־וָ֣וֵיהֶ֔ם וְצִפָּ֧ה רָאשֵׁיהֶ֛ם וַחֲשֻׁקֵיהֶ֖ם זָהָ֑ב וְאַדְנֵיהֶ֥ם חֲמִשָּׁ֖ה נְחֹֽשֶׁת׃ \{פ\} %
% ---------- פרק 37 | פסוק 1 ----------
 \textbf{א} וַיַּ֧עַשׂ בְּצַלְאֵ֛ל אֶת־הָאָרֹ֖ן עֲצֵ֣י שִׁטִּ֑ים אַמָּתַ֨יִם וָחֵ֜צִי אׇרְכּ֗וֹ וְאַמָּ֤ה וָחֵ֙צִי֙ רׇחְבּ֔וֹ וְאַמָּ֥ה וָחֵ֖צִי קֹמָתֽוֹ׃ \Rashi{\textbf{ויעש בצלאל.}לפי שנתן נפשו על המלאכה יותר משאר חכמים, נקראת על שמו (שמות רבה):}%
% ---------- פרק 37 | פסוק 2 ----------
 \textbf{ב} וַיְצַפֵּ֛הוּ זָהָ֥ב טָה֖וֹר מִבַּ֣יִת וּמִח֑וּץ וַיַּ֥עַשׂ ל֛וֹ זֵ֥ר זָהָ֖ב סָבִֽיב׃ %
% ---------- פרק 37 | פסוק 3 ----------
 \textbf{ג} וַיִּצֹ֣ק ל֗וֹ אַרְבַּע֙ טַבְּעֹ֣ת זָהָ֔ב עַ֖ל אַרְבַּ֣ע פַּעֲמֹתָ֑יו וּשְׁתֵּ֣י טַבָּעֹ֗ת עַל־צַלְעוֹ֙ הָֽאֶחָ֔ת וּשְׁתֵּי֙ טַבָּעֹ֔ת עַל־צַלְע֖וֹ הַשֵּׁנִֽית׃ %
% ---------- פרק 37 | פסוק 4 ----------
 \textbf{ד} וַיַּ֥עַשׂ בַּדֵּ֖י עֲצֵ֣י שִׁטִּ֑ים וַיְצַ֥ף אֹתָ֖ם זָהָֽב׃ %
% ---------- פרק 37 | פסוק 5 ----------
 \textbf{ה} וַיָּבֵ֤א אֶת־הַבַּדִּים֙ בַּטַּבָּעֹ֔ת עַ֖ל צַלְעֹ֣ת הָאָרֹ֑ן לָשֵׂ֖את אֶת־הָאָרֹֽן׃ %
% ---------- פרק 37 | פסוק 6 ----------
 \textbf{ו} וַיַּ֥עַשׂ כַּפֹּ֖רֶת זָהָ֣ב טָה֑וֹר אַמָּתַ֤יִם וָחֵ֙צִי֙ אׇרְכָּ֔הּ וְאַמָּ֥ה וָחֵ֖צִי רׇחְבָּֽהּ׃ %
% ---------- פרק 37 | פסוק 7 ----------
 \textbf{ז} וַיַּ֛עַשׂ שְׁנֵ֥י כְרֻבִ֖ים זָהָ֑ב מִקְשָׁה֙ עָשָׂ֣ה אֹתָ֔ם מִשְּׁנֵ֖י קְצ֥וֹת הַכַּפֹּֽרֶת׃ %
% ---------- פרק 37 | פסוק 8 ----------
 \textbf{ח} כְּרוּב־אֶחָ֤ד מִקָּצָה֙ מִזֶּ֔ה וּכְרוּב־אֶחָ֥ד מִקָּצָ֖ה מִזֶּ֑ה מִן־הַכַּפֹּ֛רֶת עָשָׂ֥ה אֶת־הַכְּרֻבִ֖ים מִשְּׁנֵ֥י (קצוותו) [קְצוֹתָֽיו]׃ %
% ---------- פרק 37 | פסוק 9 ----------
 \textbf{ט} וַיִּהְי֣וּ הַכְּרֻבִים֩ פֹּרְשֵׂ֨י כְנָפַ֜יִם לְמַ֗עְלָה סֹֽכְכִ֤ים בְּכַנְפֵיהֶם֙ עַל־הַכַּפֹּ֔רֶת וּפְנֵיהֶ֖ם אִ֣ישׁ אֶל־אָחִ֑יו אֶ֨ל־הַכַּפֹּ֔רֶת הָי֖וּ פְּנֵ֥י הַכְּרֻבִֽים׃ \{פ\} %
% ---------- פרק 37 | פסוק 10 ----------
 \textbf{י} וַיַּ֥עַשׂ אֶת־הַשֻּׁלְחָ֖ן עֲצֵ֣י שִׁטִּ֑ים אַמָּתַ֤יִם אׇרְכּוֹ֙ וְאַמָּ֣ה רׇחְבּ֔וֹ וְאַמָּ֥ה וָחֵ֖צִי קֹמָתֽוֹ׃ %
% ---------- פרק 37 | פסוק 11 ----------
 \textbf{יא} וַיְצַ֥ף אֹת֖וֹ זָהָ֣ב טָה֑וֹר וַיַּ֥עַשׂ ל֛וֹ זֵ֥ר זָהָ֖ב סָבִֽיב׃ %
% ---------- פרק 37 | פסוק 12 ----------
 \textbf{יב} וַיַּ֨עַשׂ ל֥וֹ מִסְגֶּ֛רֶת טֹ֖פַח סָבִ֑יב וַיַּ֧עַשׂ זֵר־זָהָ֛ב לְמִסְגַּרְתּ֖וֹ סָבִֽיב׃ %
% ---------- פרק 37 | פסוק 13 ----------
 \textbf{יג} וַיִּצֹ֣ק ל֔וֹ אַרְבַּ֖ע טַבְּעֹ֣ת זָהָ֑ב וַיִּתֵּן֙ אֶת־הַטַּבָּעֹ֔ת עַ֚ל אַרְבַּ֣ע הַפֵּאֹ֔ת אֲשֶׁ֖ר לְאַרְבַּ֥ע רַגְלָֽיו׃ %
% ---------- פרק 37 | פסוק 14 ----------
 \textbf{יד} לְעֻמַּת֙ הַמִּסְגֶּ֔רֶת הָי֖וּ הַטַּבָּעֹ֑ת בָּתִּים֙ לַבַּדִּ֔ים לָשֵׂ֖את אֶת־הַשֻּׁלְחָֽן׃ %
% ---------- פרק 37 | פסוק 15 ----------
 \textbf{טו} וַיַּ֤עַשׂ אֶת־הַבַּדִּים֙ עֲצֵ֣י שִׁטִּ֔ים וַיְצַ֥ף אֹתָ֖ם זָהָ֑ב לָשֵׂ֖את אֶת־הַשֻּׁלְחָֽן׃ %
% ---------- פרק 37 | פסוק 16 ----------
 \textbf{טז} וַיַּ֜עַשׂ אֶֽת־הַכֵּלִ֣ים ׀ אֲשֶׁ֣ר עַל־הַשֻּׁלְחָ֗ן אֶת־קְעָרֹתָ֤יו וְאֶת־כַּפֹּתָיו֙ וְאֵת֙ מְנַקִּיֹּתָ֔יו וְאֶ֨ת־הַקְּשָׂוֺ֔ת אֲשֶׁ֥ר יֻסַּ֖ךְ בָּהֵ֑ן זָהָ֖ב טָהֽוֹר׃ \{פ\} %
% ---------- פרק 37 | פסוק 17 ----------
 \textbf{יז} וַיַּ֥עַשׂ אֶת־הַמְּנֹרָ֖ה זָהָ֣ב טָה֑וֹר מִקְשָׁ֞ה עָשָׂ֤ה אֶת־הַמְּנֹרָה֙ יְרֵכָ֣הּ וְקָנָ֔הּ גְּבִיעֶ֛יהָ כַּפְתֹּרֶ֥יהָ וּפְרָחֶ֖יהָ מִמֶּ֥נָּה הָיֽוּ׃ %
% ---------- פרק 37 | פסוק 18 ----------
 \textbf{יח} וְשִׁשָּׁ֣ה קָנִ֔ים יֹצְאִ֖ים מִצִּדֶּ֑יהָ שְׁלֹשָׁ֣ה ׀ קְנֵ֣י מְנֹרָ֗ה מִצִּדָּהּ֙ הָֽאֶחָ֔ד וּשְׁלֹשָׁה֙ קְנֵ֣י מְנֹרָ֔ה מִצִּדָּ֖הּ הַשֵּׁנִֽי׃ %
% ---------- פרק 37 | פסוק 19 ----------
 \textbf{יט} שְׁלֹשָׁ֣ה גְ֠בִעִ֠ים מְֽשֻׁקָּדִ֞ים בַּקָּנֶ֣ה הָאֶחָד֮ כַּפְתֹּ֣ר וָפֶ֒רַח֒ וּשְׁלֹשָׁ֣ה גְבִעִ֗ים מְשֻׁקָּדִ֛ים בְּקָנֶ֥ה אֶחָ֖ד כַּפְתֹּ֣ר וָפָ֑רַח כֵּ֚ן לְשֵׁ֣שֶׁת הַקָּנִ֔ים הַיֹּצְאִ֖ים מִן־הַמְּנֹרָֽה׃ %
% ---------- פרק 37 | פסוק 20 ----------
 \textbf{כ} וּבַמְּנֹרָ֖ה אַרְבָּעָ֣ה גְבִעִ֑ים מְשֻׁ֨קָּדִ֔ים כַּפְתֹּרֶ֖יהָ וּפְרָחֶֽיהָ׃ %
% ---------- פרק 37 | פסוק 21 ----------
 \textbf{כא} וְכַפְתֹּ֡ר תַּ֩חַת֩ שְׁנֵ֨י הַקָּנִ֜ים מִמֶּ֗נָּה וְכַפְתֹּר֙ תַּ֣חַת שְׁנֵ֤י הַקָּנִים֙ מִמֶּ֔נָּה וְכַפְתֹּ֕ר תַּֽחַת־שְׁנֵ֥י הַקָּנִ֖ים מִמֶּ֑נָּה לְשֵׁ֙שֶׁת֙ הַקָּנִ֔ים הַיֹּצְאִ֖ים מִמֶּֽנָּה׃ %
% ---------- פרק 37 | פסוק 22 ----------
 \textbf{כב} כַּפְתֹּרֵיהֶ֥ם וּקְנֹתָ֖ם מִמֶּ֣נָּה הָי֑וּ כֻּלָּ֛הּ מִקְשָׁ֥ה אַחַ֖ת זָהָ֥ב טָהֽוֹר׃ %
% ---------- פרק 37 | פסוק 23 ----------
 \textbf{כג} וַיַּ֥עַשׂ אֶת־נֵרֹתֶ֖יהָ שִׁבְעָ֑ה וּמַלְקָחֶ֥יהָ וּמַחְתֹּתֶ֖יהָ זָהָ֥ב טָהֽוֹר׃ %
% ---------- פרק 37 | פסוק 24 ----------
 \textbf{כד} כִּכָּ֛ר זָהָ֥ב טָה֖וֹר עָשָׂ֣ה אֹתָ֑הּ וְאֵ֖ת כׇּל־כֵּלֶֽיהָ׃ \{פ\} %
% ---------- פרק 37 | פסוק 25 ----------
 \textbf{כה} וַיַּ֛עַשׂ אֶת־מִזְבַּ֥ח הַקְּטֹ֖רֶת עֲצֵ֣י שִׁטִּ֑ים אַמָּ֣ה אׇרְכּוֹ֩ וְאַמָּ֨ה רׇחְבּ֜וֹ רָב֗וּעַ וְאַמָּתַ֙יִם֙ קֹֽמָת֔וֹ מִמֶּ֖נּוּ הָי֥וּ קַרְנֹתָֽיו׃ %
% ---------- פרק 37 | פסוק 26 ----------
 \textbf{כו} וַיְצַ֨ף אֹת֜וֹ זָהָ֣ב טָה֗וֹר אֶת־גַּגּ֧וֹ וְאֶת־קִירֹתָ֛יו סָבִ֖יב וְאֶת־קַרְנֹתָ֑יו וַיַּ֥עַשׂ ל֛וֹ זֵ֥ר זָהָ֖ב סָבִֽיב׃ %
% ---------- פרק 37 | פסוק 27 ----------
 \textbf{כז} וּשְׁתֵּי֩ טַבְּעֹ֨ת זָהָ֜ב עָֽשָׂה־ל֣וֹ ׀ מִתַּ֣חַת לְזֵר֗וֹ עַ֚ל שְׁתֵּ֣י צַלְעֹתָ֔יו עַ֖ל שְׁנֵ֣י צִדָּ֑יו לְבָתִּ֣ים לְבַדִּ֔ים לָשֵׂ֥את אֹת֖וֹ בָּהֶֽם׃ %
% ---------- פרק 37 | פסוק 28 ----------
 \textbf{כח} וַיַּ֥עַשׂ אֶת־הַבַּדִּ֖ים עֲצֵ֣י שִׁטִּ֑ים וַיְצַ֥ף אֹתָ֖ם זָהָֽב׃ %
% ---------- פרק 37 | פסוק 29 ----------
 \textbf{כט} וַיַּ֜עַשׂ אֶת־שֶׁ֤מֶן הַמִּשְׁחָה֙ קֹ֔דֶשׁ וְאֶת־קְטֹ֥רֶת הַסַּמִּ֖ים טָה֑וֹר מַעֲשֵׂ֖ה רֹקֵֽחַ׃ \{ס\} %
% ---------- פרק 38 | פסוק 1 ----------
 \textbf{א} וַיַּ֛עַשׂ אֶת־מִזְבַּ֥ח הָעֹלָ֖ה עֲצֵ֣י שִׁטִּ֑ים חָמֵשׁ֩ אַמּ֨וֹת אׇרְכּ֜וֹ וְחָֽמֵשׁ־אַמּ֤וֹת רׇחְבּוֹ֙ רָב֔וּעַ וְשָׁלֹ֥שׁ אַמּ֖וֹת קֹמָתֽוֹ׃ %
% ---------- פרק 38 | פסוק 2 ----------
 \textbf{ב} וַיַּ֣עַשׂ קַרְנֹתָ֗יו עַ֚ל אַרְבַּ֣ע פִּנֹּתָ֔יו מִמֶּ֖נּוּ הָי֣וּ קַרְנֹתָ֑יו וַיְצַ֥ף אֹת֖וֹ נְחֹֽשֶׁת׃ %
% ---------- פרק 38 | פסוק 3 ----------
 \textbf{ג} וַיַּ֜עַשׂ אֶֽת־כׇּל־כְּלֵ֣י הַמִּזְבֵּ֗חַ אֶת־הַסִּירֹ֤ת וְאֶת־הַיָּעִים֙ וְאֶת־הַמִּזְרָקֹ֔ת אֶת־הַמִּזְלָגֹ֖ת וְאֶת־הַמַּחְתֹּ֑ת כׇּל־כֵּלָ֖יו עָשָׂ֥ה נְחֹֽשֶׁת׃ %
% ---------- פרק 38 | פסוק 4 ----------
 \textbf{ד} וַיַּ֤עַשׂ לַמִּזְבֵּ֙חַ֙ מִכְבָּ֔ר מַעֲשֵׂ֖ה רֶ֣שֶׁת נְחֹ֑שֶׁת תַּ֧חַת כַּרְכֻּבּ֛וֹ מִלְּמַ֖טָּה עַד־חֶצְיֽוֹ׃ %
% ---------- פרק 38 | פסוק 5 ----------
 \textbf{ה} וַיִּצֹ֞ק אַרְבַּ֧ע טַבָּעֹ֛ת בְּאַרְבַּ֥ע הַקְּצָוֺ֖ת לְמִכְבַּ֣ר הַנְּחֹ֑שֶׁת בָּתִּ֖ים לַבַּדִּֽים׃ %
% ---------- פרק 38 | פסוק 6 ----------
 \textbf{ו} וַיַּ֥עַשׂ אֶת־הַבַּדִּ֖ים עֲצֵ֣י שִׁטִּ֑ים וַיְצַ֥ף אֹתָ֖ם נְחֹֽשֶׁת׃ %
% ---------- פרק 38 | פסוק 7 ----------
 \textbf{ז} וַיָּבֵ֨א אֶת־הַבַּדִּ֜ים בַּטַּבָּעֹ֗ת עַ֚ל צַלְעֹ֣ת הַמִּזְבֵּ֔חַ לָשֵׂ֥את אֹת֖וֹ בָּהֶ֑ם נְב֥וּב לֻחֹ֖ת עָשָׂ֥ה אֹתֽוֹ׃ \{ס\} \Rashi{\textbf{נבוב לחת.}נבוב הוא חלול, וכן ועביו ארבע אצבעות נבוב (ירמיהו נ"ב): נבוב לחת. הלוחות של עצי שטים לכל רוח, והחלל באמצע:}%
% ---------- פרק 38 | פסוק 8 ----------
 \textbf{ח} וַיַּ֗עַשׂ אֵ֚ת הַכִּיּ֣וֹר נְחֹ֔שֶׁת וְאֵ֖ת כַּנּ֣וֹ נְחֹ֑שֶׁת בְּמַרְאֹת֙ הַצֹּ֣בְאֹ֔ת אֲשֶׁ֣ר צָֽבְא֔וּ פֶּ֖תַח אֹ֥הֶל מוֹעֵֽד׃ \{ס\} \Rashi{\textbf{במראת הצבאת.}בנות ישראל היו בידן מראות שרואות בהן כשהן מתקשטות, ואף אותן לא עכבו מלהביא לנדבת המשכן, והיה מואס משה בהן מפני שעשוים ליצר הרע, אמר לו הקב"ה קבל, כי אלו חביבין עלי מן הכל, שעל ידיהם העמידו הנשים צבאות רבות במצרים; כשהיו בעליהם יגעים בעבודת פרך, היו הולכות ומוליכות להם מאכל ומשתה, ומאכילות אותם ונוטלות המראות, וכל אחת רואה עצמה עם בעלה במראה, ומשדלתו בדברים, לומר אני נאה ממך, ומתוך כך מביאות לבעליהן לידי תאוה ונזקקות להם ומתעברות ויולדות שם, שנאמר תחת התפוח עוררתיך (שיר השירים ח'), וזה שנאמר במראת הצובאות. ונעשה הכיור מהם, שהוא לשום שלום בין איש לאשתו – להשקות ממים שבתוכו למי שקנא לה בעלה ונסתרה; ותדע לך שהן מראות ממש, שהרי נאמר ונחשת התנופה שבעים ככר וגו', ויעש בה וגומר, וכיור וכנו לא הזכרו שם, למדת שלא היה נחשת של כיור מנחשת התנופה, כך דרש רבי תנחומא, וכן תרגם אנקלוס במחזית נשיא והוא תרגום של מראות, מירוריי"ש בלעז. וכן מצינו והגליונים (ישעיה ג') – מתרגמינן מחזיתא: אשר צבאו. להביא נדבתן:}%
% ---------- פרק 38 | פסוק 9 ----------
 \textbf{ט} וַיַּ֖עַשׂ אֶת־הֶחָצֵ֑ר לִפְאַ֣ת ׀ נֶ֣גֶב תֵּימָ֗נָה קַלְעֵ֤י הֶֽחָצֵר֙ שֵׁ֣שׁ מׇשְׁזָ֔ר מֵאָ֖ה בָּאַמָּֽה׃ %
% ---------- פרק 38 | פסוק 10 ----------
 \textbf{י} עַמּוּדֵיהֶ֣ם עֶשְׂרִ֔ים וְאַדְנֵיהֶ֥ם עֶשְׂרִ֖ים נְחֹ֑שֶׁת וָוֵ֧י הָעַמּוּדִ֛ים וַחֲשֻׁקֵיהֶ֖ם כָּֽסֶף׃ %
% ---------- פרק 38 | פסוק 11 ----------
 \textbf{יא} וְלִפְאַ֤ת צָפוֹן֙ מֵאָ֣ה בָֽאַמָּ֔ה עַמּוּדֵיהֶ֣ם עֶשְׂרִ֔ים וְאַדְנֵיהֶ֥ם עֶשְׂרִ֖ים נְחֹ֑שֶׁת וָוֵ֧י הָֽעַמּוּדִ֛ים וַחֲשֻׁקֵיהֶ֖ם כָּֽסֶף׃ %
% ---------- פרק 38 | פסוק 12 ----------
 \textbf{יב} וְלִפְאַת־יָ֗ם קְלָעִים֙ חֲמִשִּׁ֣ים בָּֽאַמָּ֔ה עַמּוּדֵיהֶ֣ם עֲשָׂרָ֔ה וְאַדְנֵיהֶ֖ם עֲשָׂרָ֑ה וָוֵ֧י הָעַמֻּדִ֛ים וַחֲשׁוּקֵיהֶ֖ם כָּֽסֶף׃ %
% ---------- פרק 38 | פסוק 13 ----------
 \textbf{יג} וְלִפְאַ֛ת קֵ֥דְמָה מִזְרָ֖חָה חֲמִשִּׁ֥ים אַמָּֽה׃ %
% ---------- פרק 38 | פסוק 14 ----------
 \textbf{יד} קְלָעִ֛ים חֲמֵשׁ־עֶשְׂרֵ֥ה אַמָּ֖ה אֶל־הַכָּתֵ֑ף עַמּוּדֵיהֶ֣ם שְׁלֹשָׁ֔ה וְאַדְנֵיהֶ֖ם שְׁלֹשָֽׁה׃ %
% ---------- פרק 38 | פסוק 15 ----------
 \textbf{טו} וְלַכָּתֵ֣ף הַשֵּׁנִ֗ית מִזֶּ֤ה וּמִזֶּה֙ לְשַׁ֣עַר הֶֽחָצֵ֔ר קְלָעִ֕ים חֲמֵ֥שׁ עֶשְׂרֵ֖ה אַמָּ֑ה עַמֻּדֵיהֶ֣ם שְׁלֹשָׁ֔ה וְאַדְנֵיהֶ֖ם שְׁלֹשָֽׁה׃ %
% ---------- פרק 38 | פסוק 16 ----------
 \textbf{טז} כׇּל־קַלְעֵ֧י הֶחָצֵ֛ר סָבִ֖יב שֵׁ֥שׁ מׇשְׁזָֽר׃ %
% ---------- פרק 38 | פסוק 17 ----------
 \textbf{יז} וְהָאֲדָנִ֣ים לָֽעַמֻּדִים֮ נְחֹ֒שֶׁת֒ וָוֵ֨י הָֽעַמּוּדִ֜ים וַחֲשׁוּקֵיהֶם֙ כֶּ֔סֶף וְצִפּ֥וּי רָאשֵׁיהֶ֖ם כָּ֑סֶף וְהֵם֙ מְחֻשָּׁקִ֣ים כֶּ֔סֶף כֹּ֖ל עַמֻּדֵ֥י הֶחָצֵֽר׃ %
% ---------- פרק 38 | פסוק 18 ----------
 \textbf{יח} וּמָסַ֞ךְ שַׁ֤עַר הֶחָצֵר֙ מַעֲשֵׂ֣ה רֹקֵ֔ם תְּכֵ֧לֶת וְאַרְגָּמָ֛ן וְתוֹלַ֥עַת שָׁנִ֖י וְשֵׁ֣שׁ מׇשְׁזָ֑ר וְעֶשְׂרִ֤ים אַמָּה֙ אֹ֔רֶךְ וְקוֹמָ֤ה בְרֹ֙חַב֙ חָמֵ֣שׁ אַמּ֔וֹת לְעֻמַּ֖ת קַלְעֵ֥י הֶחָצֵֽר׃ \Rashi{\textbf{לעמת קלעי החצר.}כמדת קלעי החצר:}%
% ---------- פרק 38 | פסוק 19 ----------
 \textbf{יט} וְעַמֻּֽדֵיהֶם֙ אַרְבָּעָ֔ה וְאַדְנֵיהֶ֥ם אַרְבָּעָ֖ה נְחֹ֑שֶׁת וָוֵיהֶ֣ם כֶּ֔סֶף וְצִפּ֧וּי רָאשֵׁיהֶ֛ם וַחֲשֻׁקֵיהֶ֖ם כָּֽסֶף׃ %
% ---------- פרק 38 | פסוק 20 ----------
 \textbf{כ} וְֽכׇל־הַיְתֵדֹ֞ת לַמִּשְׁכָּ֧ן וְלֶחָצֵ֛ר סָבִ֖יב נְחֹֽשֶׁת׃ \{ס\} %
\addparasha{חומש שמות | פרשת פקודי}
% ---------- פרק 38 | פסוק 21 ----------
 \textbf{כא} אֵ֣לֶּה פְקוּדֵ֤י הַמִּשְׁכָּן֙ מִשְׁכַּ֣ן הָעֵדֻ֔ת אֲשֶׁ֥ר פֻּקַּ֖ד עַל־פִּ֣י מֹשֶׁ֑ה עֲבֹדַת֙ הַלְוִיִּ֔ם בְּיַד֙ אִֽיתָמָ֔ר בֶּֽן־אַהֲרֹ֖ן הַכֹּהֵֽן׃ \Rashi{\textbf{אלה פקודי.}בפרשה זו נמנו כל משקלי נדבת המשכן, לכסף ולזהב ולנחשת, ונמנו כל כליו לכל עבודתו: המשכן משכן. שני פעמים, רמז למקדש שנתמשכן בשני חרבנין על עונותיהן של ישראל (תנחומא): משכן העדת. עדות לישראל שותר להם הקב"ה על מעשה העגל, שהרי השרה שכינתו ביניהם: עבדת הלוים. פקודי המשכן וכליו היא עבודה המסורה ללוים במדבר – לשאת ולהוריד ולהקים, איש איש למשאו המפקד עליו, כמו שאמור בפ' נשא: ביד איתמר. הוא היה פקיד עליהם למסר לכל בית אב עבודה שעליו:}%
% ---------- פרק 38 | פסוק 22 ----------
 \textbf{כב} וּבְצַלְאֵ֛ל בֶּן־אוּרִ֥י בֶן־ח֖וּר לְמַטֵּ֣ה יְהוּדָ֑ה עָשָׂ֕ה אֵ֛ת כׇּל־אֲשֶׁר־צִוָּ֥ה יְהֹוָ֖ה אֶת־מֹשֶֽׁה׃ \Rashi{\textbf{ובצלאל בן אורי וגו' עשה את כל אשר צוה ה' את משה.}אשר צוה משה אין כתיב כאן, אלא כל אשר צוה ה' את משה, אפלו דברים שלא אמר לו רבו, הסכימה דעתו למה שנאמר למשה בסיני, כי משה צוה לבצלאל לעשות תחלה כלים ואחר כך משכן, אמר לו בצלאל מנהג עולם לעשות תחלה בית ואחר כך משים כלים בתוכו. אמר לו כך שמעתי מפי הקב"ה. אמר לו משה בצל אל היית, כי בודאי כך צוה לי הקב"ה, וכן עשה המשכן תחלה ואחר כך עשה הכלים:}%
% ---------- פרק 38 | פסוק 23 ----------
 \textbf{כג} וְאִתּ֗וֹ אׇהֳלִיאָ֞ב בֶּן־אֲחִיסָמָ֛ךְ לְמַטֵּה־דָ֖ן חָרָ֣שׁ וְחֹשֵׁ֑ב וְרֹקֵ֗ם בַּתְּכֵ֙לֶת֙ וּבָֽאַרְגָּמָ֔ן וּבְתוֹלַ֥עַת הַשָּׁנִ֖י וּבַשֵּֽׁשׁ׃ \{ס\} %
% ---------- פרק 38 | פסוק 24 ----------
 \textbf{כד} כׇּל־הַזָּהָ֗ב הֶֽעָשׂוּי֙ לַמְּלָאכָ֔ה בְּכֹ֖ל מְלֶ֣אכֶת הַקֹּ֑דֶשׁ וַיְהִ֣י ׀ זְהַ֣ב הַתְּנוּפָ֗ה תֵּ֤שַׁע וְעֶשְׂרִים֙ כִּכָּ֔ר וּשְׁבַ֨ע מֵא֧וֹת וּשְׁלֹשִׁ֛ים שֶׁ֖קֶל בְּשֶׁ֥קֶל הַקֹּֽדֶשׁ׃ \Rashi{\textbf{ככר.}ששים מנה, ומנה של קדש כפול היה, הרי הככר ק"כ מנה, והמנה כ"ה סלעים, הרי ככר של קדש שלשת אלפים שקלים, לפיכך מנה בפרוטרוט כל השקלים שפחותין במנינם מג' אלפים, שאין מגיעין לככר:}%
% ---------- פרק 38 | פסוק 25 ----------
 \textbf{כה} וְכֶ֛סֶף פְּקוּדֵ֥י הָעֵדָ֖ה מְאַ֣ת כִּכָּ֑ר וְאֶ֩לֶף֩ וּשְׁבַ֨ע מֵא֜וֹת וַחֲמִשָּׁ֧ה וְשִׁבְעִ֛ים שֶׁ֖קֶל בְּשֶׁ֥קֶל הַקֹּֽדֶשׁ׃ %
% ---------- פרק 38 | פסוק 26 ----------
 \textbf{כו} בֶּ֚קַע לַגֻּלְגֹּ֔לֶת מַחֲצִ֥ית הַשֶּׁ֖קֶל בְּשֶׁ֣קֶל הַקֹּ֑דֶשׁ לְכֹ֨ל הָעֹבֵ֜ר עַל־הַפְּקֻדִ֗ים מִבֶּ֨ן עֶשְׂרִ֤ים שָׁנָה֙ וָמַ֔עְלָה לְשֵׁשׁ־מֵא֥וֹת אֶ֙לֶף֙ וּשְׁלֹ֣שֶׁת אֲלָפִ֔ים וַחֲמֵ֥שׁ מֵא֖וֹת וַחֲמִשִּֽׁים׃ \Rashi{\textbf{בקע.}הוא שם משקל של מחצית השקל: לשש מאות אלף וגו'. כך היו ישראל, וכך עלה מנינם אחר שהוקם המשכן בספר וידבר, ואף עתה בנדבת המשכן כך היו, ומנין חצאי השקלים של שש מאות אלף עולה מאת ככר, כל אחד של שלשת אלפים שקלים, כיצד? שש מאות אלף חצאין הרי הן שלש מאות אלף שלמים, הרי מאת ככר, והשלשת אלפים וה' מאות וחמשים חצאים עולים אלף וז' מאות וע"ה שקלים:}%
% ---------- פרק 38 | פסוק 27 ----------
 \textbf{כז} וַיְהִ֗י מְאַת֙ כִּכַּ֣ר הַכֶּ֔סֶף לָצֶ֗קֶת אֵ֚ת אַדְנֵ֣י הַקֹּ֔דֶשׁ וְאֵ֖ת אַדְנֵ֣י הַפָּרֹ֑כֶת מְאַ֧ת אֲדָנִ֛ים לִמְאַ֥ת הַכִּכָּ֖ר כִּכָּ֥ר לָאָֽדֶן׃ \Rashi{\textbf{לצקת.}כתרגומו לאתכא: את אדני הקדש. של קרשי המשכן, שהם מ"ח קרשים ולהן צ"ו אדנים, ואדני הפרכת ארבעה, הרי מאה, וכל שאר האדנים נחשת כתוב בהם:}%
% ---------- פרק 38 | פסוק 28 ----------
 \textbf{כח} וְאֶת־הָאֶ֜לֶף וּשְׁבַ֤ע הַמֵּאוֹת֙ וַחֲמִשָּׁ֣ה וְשִׁבְעִ֔ים עָשָׂ֥ה וָוִ֖ים לָעַמּוּדִ֑ים וְצִפָּ֥ה רָאשֵׁיהֶ֖ם וְחִשַּׁ֥ק אֹתָֽם׃ \Rashi{\textbf{וצפה ראשיהם.}של עמודים מהן, שבכלן כתיב וצפוי ראשיהם וחשקיהם כסף:}%
% ---------- פרק 38 | פסוק 29 ----------
 \textbf{כט} וּנְחֹ֥שֶׁת הַתְּנוּפָ֖ה שִׁבְעִ֣ים כִּכָּ֑ר וְאַלְפַּ֥יִם וְאַרְבַּע־מֵא֖וֹת שָֽׁקֶל׃ %
% ---------- פרק 38 | פסוק 30 ----------
 \textbf{ל} וַיַּ֣עַשׂ בָּ֗הּ אֶת־אַדְנֵי֙ פֶּ֚תַח אֹ֣הֶל מוֹעֵ֔ד וְאֵת֙ מִזְבַּ֣ח הַנְּחֹ֔שֶׁת וְאֶת־מִכְבַּ֥ר הַנְּחֹ֖שֶׁת אֲשֶׁר־ל֑וֹ וְאֵ֖ת כׇּל־כְּלֵ֥י הַמִּזְבֵּֽחַ׃ %
% ---------- פרק 38 | פסוק 31 ----------
 \textbf{לא} וְאֶת־אַדְנֵ֤י הֶֽחָצֵר֙ סָבִ֔יב וְאֶת־אַדְנֵ֖י שַׁ֣עַר הֶחָצֵ֑ר וְאֵ֨ת כׇּל־יִתְדֹ֧ת הַמִּשְׁכָּ֛ן וְאֶת־כׇּל־יִתְדֹ֥ת הֶחָצֵ֖ר סָבִֽיב׃ %
% ---------- פרק 39 | פסוק 1 ----------
 \textbf{א} וּמִן־הַתְּכֵ֤לֶת וְהָֽאַרְגָּמָן֙ וְתוֹלַ֣עַת הַשָּׁנִ֔י עָשׂ֥וּ בִגְדֵי־שְׂרָ֖ד לְשָׁרֵ֣ת בַּקֹּ֑דֶשׁ וַֽיַּעֲשׂ֞וּ אֶת־בִּגְדֵ֤י הַקֹּ֙דֶשׁ֙ אֲשֶׁ֣ר לְאַהֲרֹ֔ן כַּאֲשֶׁ֛ר צִוָּ֥ה יְהֹוָ֖ה אֶת־מֹשֶֽׁה׃ \{פ\} \Rashi{\textbf{ומן התכלת והארגמן וגו'.}שש לא נאמר כאן, ומכאן אני אומר שאין בגדי שרד הללו בגדי כהנה, שבבגדי כהנה היה שש, אלא הם בגדים שמכסים בהם כלי הקדש בשעת סלוק מסעות, שלא היה בהם שש:}%
% ---------- פרק 39 | פסוק 2 ----------
 \textbf{ב} וַיַּ֖עַשׂ אֶת־הָאֵפֹ֑ד זָהָ֗ב תְּכֵ֧לֶת וְאַרְגָּמָ֛ן וְתוֹלַ֥עַת שָׁנִ֖י וְשֵׁ֥שׁ מׇשְׁזָֽר׃ %
% ---------- פרק 39 | פסוק 3 ----------
 \textbf{ג} וַֽיְרַקְּע֞וּ אֶת־פַּחֵ֣י הַזָּהָב֮ וְקִצֵּ֣ץ פְּתִילִם֒ לַעֲשׂ֗וֹת בְּת֤וֹךְ הַתְּכֵ֙לֶת֙ וּבְת֣וֹךְ הָֽאַרְגָּמָ֔ן וּבְת֛וֹךְ תּוֹלַ֥עַת הַשָּׁנִ֖י וּבְת֣וֹךְ הַשֵּׁ֑שׁ מַעֲשֵׂ֖ה חֹשֵֽׁב׃ \Rashi{\textbf{וירקעו.}כמו לרקע הארץ (תהילים קל"ו), כתרגומו ורדידו, טסין היו מרדדין מן הזהב, אשטנ"דרא בלעז, טסין דקות. כאן הוא מלמדך היאך היו טווין את הזהב עם החוטין, מרדדים טסין דקין, וקוצצין מהן פתילים לארך הטס – לעשות אותן פתילים תערבת עם כל מין ומין בחשן ואפוד שנ' בהן זהב, חוט אחד של זהב עם ו' חוטין של תכלת, וכן עם כל מין ומין, שכל המינים חוטן כפול ו', והזהב חוט שביעי עם כל אחד ואחד:}%
% ---------- פרק 39 | פסוק 4 ----------
 \textbf{ד} כְּתֵפֹ֥ת עָֽשׂוּ־ל֖וֹ חֹבְרֹ֑ת עַל־שְׁנֵ֥י (קצוותו) [קְצוֹתָ֖יו] חֻבָּֽר׃ %
% ---------- פרק 39 | פסוק 5 ----------
 \textbf{ה} וְחֵ֨שֶׁב אֲפֻדָּת֜וֹ אֲשֶׁ֣ר עָלָ֗יו מִמֶּ֣נּוּ הוּא֮ כְּמַעֲשֵׂ֒הוּ֒ זָהָ֗ב תְּכֵ֧לֶת וְאַרְגָּמָ֛ן וְתוֹלַ֥עַת שָׁנִ֖י וְשֵׁ֣שׁ מׇשְׁזָ֑ר כַּאֲשֶׁ֛ר צִוָּ֥ה יְהֹוָ֖ה אֶת־מֹשֶֽׁה׃ \{ס\} %
% ---------- פרק 39 | פסוק 6 ----------
 \textbf{ו} וַֽיַּעֲשׂוּ֙ אֶת־אַבְנֵ֣י הַשֹּׁ֔הַם מֻֽסַבֹּ֖ת מִשְׁבְּצֹ֣ת זָהָ֑ב מְפֻתָּחֹת֙ פִּתּוּחֵ֣י חוֹתָ֔ם עַל־שְׁמ֖וֹת בְּנֵ֥י יִשְׂרָאֵֽל׃ %
% ---------- פרק 39 | פסוק 7 ----------
 \textbf{ז} וַיָּ֣שֶׂם אֹתָ֗ם עַ֚ל כִּתְפֹ֣ת הָאֵפֹ֔ד אַבְנֵ֥י זִכָּר֖וֹן לִבְנֵ֣י יִשְׂרָאֵ֑ל כַּאֲשֶׁ֛ר צִוָּ֥ה יְהֹוָ֖ה אֶת־מֹשֶֽׁה׃ \{פ\} %
% ---------- פרק 39 | פסוק 8 ----------
 \textbf{ח} וַיַּ֧עַשׂ אֶת־הַחֹ֛שֶׁן מַעֲשֵׂ֥ה חֹשֵׁ֖ב כְּמַעֲשֵׂ֣ה אֵפֹ֑ד זָהָ֗ב תְּכֵ֧לֶת וְאַרְגָּמָ֛ן וְתוֹלַ֥עַת שָׁנִ֖י וְשֵׁ֥שׁ מׇשְׁזָֽר׃ %
% ---------- פרק 39 | פסוק 9 ----------
 \textbf{ט} רָב֧וּעַ הָיָ֛ה כָּפ֖וּל עָשׂ֣וּ אֶת־הַחֹ֑שֶׁן זֶ֧רֶת אׇרְכּ֛וֹ וְזֶ֥רֶת רׇחְבּ֖וֹ כָּפֽוּל׃ %
% ---------- פרק 39 | פסוק 10 ----------
 \textbf{י} וַיְמַ֨לְאוּ־ב֔וֹ אַרְבָּעָ֖ה ט֣וּרֵי אָ֑בֶן ט֗וּר אֹ֤דֶם פִּטְדָה֙ וּבָרֶ֔קֶת הַטּ֖וּר הָאֶחָֽד׃ %
% ---------- פרק 39 | פסוק 11 ----------
 \textbf{יא} וְהַטּ֖וּר הַשֵּׁנִ֑י נֹ֥פֶךְ סַפִּ֖יר וְיָהֲלֹֽם׃ %
% ---------- פרק 39 | פסוק 12 ----------
 \textbf{יב} וְהַטּ֖וּר הַשְּׁלִישִׁ֑י לֶ֥שֶׁם שְׁב֖וֹ וְאַחְלָֽמָה׃ %
% ---------- פרק 39 | פסוק 13 ----------
 \textbf{יג} וְהַטּוּר֙ הָֽרְבִיעִ֔י תַּרְשִׁ֥ישׁ שֹׁ֖הַם וְיָשְׁפֵ֑ה מֽוּסַבֹּ֛ת מִשְׁבְּצֹ֥ת זָהָ֖ב בְּמִלֻּאֹתָֽם׃ %
% ---------- פרק 39 | פסוק 14 ----------
 \textbf{יד} וְ֠הָאֲבָנִ֠ים עַל־שְׁמֹ֨ת בְּנֵי־יִשְׂרָאֵ֥ל הֵ֛נָּה שְׁתֵּ֥ים עֶשְׂרֵ֖ה עַל־שְׁמֹתָ֑ם פִּתּוּחֵ֤י חֹתָם֙ אִ֣ישׁ עַל־שְׁמ֔וֹ לִשְׁנֵ֥ים עָשָׂ֖ר שָֽׁבֶט׃ %
% ---------- פרק 39 | פסוק 15 ----------
 \textbf{טו} וַיַּעֲשׂ֧וּ עַל־הַחֹ֛שֶׁן שַׁרְשְׁרֹ֥ת גַּבְלֻ֖ת מַעֲשֵׂ֣ה עֲבֹ֑ת זָהָ֖ב טָהֽוֹר׃ %
% ---------- פרק 39 | פסוק 16 ----------
 \textbf{טז} וַֽיַּעֲשׂ֗וּ שְׁתֵּי֙ מִשְׁבְּצֹ֣ת זָהָ֔ב וּשְׁתֵּ֖י טַבְּעֹ֣ת זָהָ֑ב וַֽיִּתְּנ֗וּ אֶת־שְׁתֵּי֙ הַטַּבָּעֹ֔ת עַל־שְׁנֵ֖י קְצ֥וֹת הַחֹֽשֶׁן׃ %
% ---------- פרק 39 | פסוק 17 ----------
 \textbf{יז} וַֽיִּתְּנ֗וּ שְׁתֵּי֙ הָעֲבֹתֹ֣ת הַזָּהָ֔ב עַל־שְׁתֵּ֖י הַטַּבָּעֹ֑ת עַל־קְצ֖וֹת הַחֹֽשֶׁן׃ %
% ---------- פרק 39 | פסוק 18 ----------
 \textbf{יח} וְאֵ֨ת שְׁתֵּ֤י קְצוֹת֙ שְׁתֵּ֣י הָֽעֲבֹתֹ֔ת נָתְנ֖וּ עַל־שְׁתֵּ֣י הַֽמִּשְׁבְּצֹ֑ת וַֽיִּתְּנֻ֛ם עַל־כִּתְפֹ֥ת הָאֵפֹ֖ד אֶל־מ֥וּל פָּנָֽיו׃ %
% ---------- פרק 39 | פסוק 19 ----------
 \textbf{יט} וַֽיַּעֲשׂ֗וּ שְׁתֵּי֙ טַבְּעֹ֣ת זָהָ֔ב וַיָּשִׂ֕ימוּ עַל־שְׁנֵ֖י קְצ֣וֹת הַחֹ֑שֶׁן עַל־שְׂפָת֕וֹ אֲשֶׁ֛ר אֶל־עֵ֥בֶר הָאֵפֹ֖ד בָּֽיְתָה׃ %
% ---------- פרק 39 | פסוק 20 ----------
 \textbf{כ} וַֽיַּעֲשׂוּ֮ שְׁתֵּ֣י טַבְּעֹ֣ת זָהָב֒ וַֽיִּתְּנֻ֡ם עַל־שְׁתֵּי֩ כִתְפֹ֨ת הָאֵפֹ֤ד מִלְּמַ֙טָּה֙ מִמּ֣וּל פָּנָ֔יו לְעֻמַּ֖ת מַחְבַּרְתּ֑וֹ מִמַּ֕עַל לְחֵ֖שֶׁב הָאֵפֹֽד׃ %
% ---------- פרק 39 | פסוק 21 ----------
 \textbf{כא} וַיִּרְכְּס֣וּ אֶת־הַחֹ֡שֶׁן מִטַּבְּעֹתָיו֩ אֶל־טַבְּעֹ֨ת הָאֵפֹ֜ד בִּפְתִ֣יל תְּכֵ֗לֶת לִֽהְיֹת֙ עַל־חֵ֣שֶׁב הָאֵפֹ֔ד וְלֹֽא־יִזַּ֣ח הַחֹ֔שֶׁן מֵעַ֖ל הָאֵפֹ֑ד כַּאֲשֶׁ֛ר צִוָּ֥ה יְהֹוָ֖ה אֶת־מֹשֶֽׁה׃ \{פ\} %
% ---------- פרק 39 | פסוק 22 ----------
 \textbf{כב} וַיַּ֛עַשׂ אֶת־מְעִ֥יל הָאֵפֹ֖ד מַעֲשֵׂ֣ה אֹרֵ֑ג כְּלִ֖יל תְּכֵֽלֶת׃ %
% ---------- פרק 39 | פסוק 23 ----------
 \textbf{כג} וּפִֽי־הַמְּעִ֥יל בְּתוֹכ֖וֹ כְּפִ֣י תַחְרָ֑א שָׂפָ֥ה לְפִ֛יו סָבִ֖יב לֹ֥א יִקָּרֵֽעַ׃ %
% ---------- פרק 39 | פסוק 24 ----------
 \textbf{כד} וַֽיַּעֲשׂוּ֙ עַל־שׁוּלֵ֣י הַמְּעִ֔יל רִמּוֹנֵ֕י תְּכֵ֥לֶת וְאַרְגָּמָ֖ן וְתוֹלַ֣עַת שָׁנִ֑י מׇשְׁזָֽר׃ %
% ---------- פרק 39 | פסוק 25 ----------
 \textbf{כה} וַיַּעֲשׂ֥וּ פַעֲמֹנֵ֖י זָהָ֣ב טָה֑וֹר וַיִּתְּנ֨וּ אֶת־הַפַּֽעֲמֹנִ֜ים בְּת֣וֹךְ הָרִמֹּנִ֗ים עַל־שׁוּלֵ֤י הַמְּעִיל֙ סָבִ֔יב בְּת֖וֹךְ הָרִמֹּנִֽים׃ %
% ---------- פרק 39 | פסוק 26 ----------
 \textbf{כו} פַּעֲמֹ֤ן וְרִמֹּן֙ פַּעֲמֹ֣ן וְרִמֹּ֔ן עַל־שׁוּלֵ֥י הַמְּעִ֖יל סָבִ֑יב לְשָׁרֵ֕ת כַּאֲשֶׁ֛ר צִוָּ֥ה יְהֹוָ֖ה אֶת־מֹשֶֽׁה׃ \{ס\} %
% ---------- פרק 39 | פסוק 27 ----------
 \textbf{כז} וַֽיַּעֲשׂ֛וּ אֶת־הַכׇּתְנֹ֥ת שֵׁ֖שׁ מַעֲשֵׂ֣ה אֹרֵ֑ג לְאַהֲרֹ֖ן וּלְבָנָֽיו׃ %
% ---------- פרק 39 | פסוק 28 ----------
 \textbf{כח} וְאֵת֙ הַמִּצְנֶ֣פֶת שֵׁ֔שׁ וְאֶת־פַּאֲרֵ֥י הַמִּגְבָּעֹ֖ת שֵׁ֑שׁ וְאֶת־מִכְנְסֵ֥י הַבָּ֖ד שֵׁ֥שׁ מׇשְׁזָֽר׃ \Rashi{\textbf{ואת פארי המגבעת.}תפארת המגבעות – המגבעות המפארות:}%
% ---------- פרק 39 | פסוק 29 ----------
 \textbf{כט} וְֽאֶת־הָאַבְנֵ֞ט שֵׁ֣שׁ מׇשְׁזָ֗ר וּתְכֵ֧לֶת וְאַרְגָּמָ֛ן וְתוֹלַ֥עַת שָׁנִ֖י מַעֲשֵׂ֣ה רֹקֵ֑ם כַּאֲשֶׁ֛ר צִוָּ֥ה יְהֹוָ֖ה אֶת־מֹשֶֽׁה׃ \{ס\} %
% ---------- פרק 39 | פסוק 30 ----------
 \textbf{ל} וַֽיַּעֲשׂ֛וּ אֶת־צִ֥יץ נֵֽזֶר־הַקֹּ֖דֶשׁ זָהָ֣ב טָה֑וֹר וַיִּכְתְּב֣וּ עָלָ֗יו מִכְתַּב֙ פִּתּוּחֵ֣י חוֹתָ֔ם קֹ֖דֶשׁ לַיהֹוָֽה׃ %
% ---------- פרק 39 | פסוק 31 ----------
 \textbf{לא} וַיִּתְּנ֤וּ עָלָיו֙ פְּתִ֣יל תְּכֵ֔לֶת לָתֵ֥ת עַל־הַמִּצְנֶ֖פֶת מִלְמָ֑עְלָה כַּאֲשֶׁ֛ר צִוָּ֥ה יְהֹוָ֖ה אֶת־מֹשֶֽׁה׃ \{ס\} \Rashi{\textbf{לתת על המצנפת מלמעלה.}וע"י הפתילים היה מושיבו על המצנפת כמין כתר; וא"א לומר הציץ על המצנפת, שהרי בשחיטת קדשים שנינו שערו היה נראה בין ציץ למצנפת ששם מניח תפלין, והציץ היה נתון על המצח, הרי המצנפת למעלה והציץ למטה, ומהו על המצנפת מלמעלה? ועוד הקשיתי בה – כאן הוא אומר ויתנו עליו פתיל תכלת, ובענין הצואה הוא אומר ושמת אותו על פתיל תכלת? ואומר אני, פתיל תכלת זה חוטין הן, לקשרו בהן במצנפת – לפי שהציץ אינו אלא מאזן לאזן ובמה יקשרנו במצחו? – והיו קבועין בו חוטי תכלת לשני ראשיו ובאמצעיתו, שבהן קושרו ותולהו במצנפת כשהוא בראשו, ושני חוטין היו בכל קצה וקצה, אחד ממעל ואחד מתחת לצד מצחו, וכן באמצעו, שכך הוא נח לקשר, ואין דרך קשירה בפחות משני חוטין, לכך נאמר על פתיל תכלת, ועליו פתיל תכלת, וקושר ראשיהם השנים כלן יחד מאחוריו למול ערפו, ומושיבו על המצנפת; ואל תתמה שלא נאמר פתילי תכלת הואיל ומרבין הן, שהרי מצינו בחשן ואפוד וירכסו את החשן … בפתיל תכלת, ועל כרחך פחות משנים לא היו שהרי בשני קצות החשן היו ב' טבעות החשן, ובשתי כתפות האפוד היו ב' טבעות האפוד, שכנגדן, ולפי דרך קשירה ד' חוטין היו ומכל מקום פחות משנים א"א:}%
% ---------- פרק 39 | פסוק 32 ----------
 \textbf{לב} וַתֵּ֕כֶל כׇּל־עֲבֹדַ֕ת מִשְׁכַּ֖ן אֹ֣הֶל מוֹעֵ֑ד וַֽיַּעֲשׂוּ֙ בְּנֵ֣י יִשְׂרָאֵ֔ל כְּ֠כֹ֠ל אֲשֶׁ֨ר צִוָּ֧ה יְהֹוָ֛ה אֶת־מֹשֶׁ֖ה כֵּ֥ן עָשֽׂוּ׃ \{פ\} \Rashi{\textbf{ויעשו בני ישראל.}את המלאכה ככל אשר צוה ה' וגו':}%
% ---------- פרק 39 | פסוק 33 ----------
 \textbf{לג} וַיָּבִ֤יאוּ אֶת־הַמִּשְׁכָּן֙ אֶל־מֹשֶׁ֔ה אֶת־הָאֹ֖הֶל וְאֶת־כׇּל־כֵּלָ֑יו קְרָסָ֣יו קְרָשָׁ֔יו בְּרִיחָ֖ו וְעַמֻּדָ֥יו וַאֲדָנָֽיו׃ \Rashi{\textbf{ויביאו את המשכן וגו'.}שלא היו יכולין להקימו; ולפי שלא עשה משה שום מלאכה במשכן, הניח לו הקב"ה הקמתו, שלא היה יכול להקימו שום אדם מחמת כבד הקרשים, שאין כח באדם לזקפן, ומשה העמידו; אמר משה לפני הקב"ה איך אפשר הקמתו ע"י אדם? אמר לו עסק אתה בידך, ונראה כמקימו והוא נזקף וקם מאליו, וזהו שנ' הוקם המשכן (שמות מ') – הוקם מאליו; מדרש ר' תנחומא:}%
% ---------- פרק 39 | פסוק 34 ----------
 \textbf{לד} וְאֶת־מִכְסֵ֞ה עוֹרֹ֤ת הָֽאֵילִם֙ הַמְאׇדָּמִ֔ים וְאֶת־מִכְסֵ֖ה עֹרֹ֣ת הַתְּחָשִׁ֑ים וְאֵ֖ת פָּרֹ֥כֶת הַמָּסָֽךְ׃ %
% ---------- פרק 39 | פסוק 35 ----------
 \textbf{לה} אֶת־אֲר֥וֹן הָעֵדֻ֖ת וְאֶת־בַּדָּ֑יו וְאֵ֖ת הַכַּפֹּֽרֶת׃ %
% ---------- פרק 39 | פסוק 36 ----------
 \textbf{לו} אֶת־הַשֻּׁלְחָן֙ אֶת־כׇּל־כֵּלָ֔יו וְאֵ֖ת לֶ֥חֶם הַפָּנִֽים׃ %
% ---------- פרק 39 | פסוק 37 ----------
 \textbf{לז} אֶת־הַמְּנֹרָ֨ה הַטְּהֹרָ֜ה אֶת־נֵרֹתֶ֗יהָ נֵרֹ֛ת הַמַּֽעֲרָכָ֖ה וְאֶת־כׇּל־כֵּלֶ֑יהָ וְאֵ֖ת שֶׁ֥מֶן הַמָּאֽוֹר׃ %
% ---------- פרק 39 | פסוק 38 ----------
 \textbf{לח} וְאֵת֙ מִזְבַּ֣ח הַזָּהָ֔ב וְאֵת֙ שֶׁ֣מֶן הַמִּשְׁחָ֔ה וְאֵ֖ת קְטֹ֣רֶת הַסַּמִּ֑ים וְאֵ֕ת מָסַ֖ךְ פֶּ֥תַח הָאֹֽהֶל׃ %
% ---------- פרק 39 | פסוק 39 ----------
 \textbf{לט} אֵ֣ת ׀ מִזְבַּ֣ח הַנְּחֹ֗שֶׁת וְאֶת־מִכְבַּ֤ר הַנְּחֹ֙שֶׁת֙ אֲשֶׁר־ל֔וֹ אֶת־בַּדָּ֖יו וְאֶת־כׇּל־כֵּלָ֑יו אֶת־הַכִּיֹּ֖ר וְאֶת־כַּנּֽוֹ׃ %
% ---------- פרק 39 | פסוק 40 ----------
 \textbf{מ} אֵת֩ קַלְעֵ֨י הֶחָצֵ֜ר אֶת־עַמֻּדֶ֣יהָ וְאֶת־אֲדָנֶ֗יהָ וְאֶת־הַמָּסָךְ֙ לְשַׁ֣עַר הֶֽחָצֵ֔ר אֶת־מֵיתָרָ֖יו וִיתֵדֹתֶ֑יהָ וְאֵ֗ת כׇּל־כְּלֵ֛י עֲבֹדַ֥ת הַמִּשְׁכָּ֖ן לְאֹ֥הֶל מוֹעֵֽד׃ %
% ---------- פרק 39 | פסוק 41 ----------
 \textbf{מא} אֶת־בִּגְדֵ֥י הַשְּׂרָ֖ד לְשָׁרֵ֣ת בַּקֹּ֑דֶשׁ אֶת־בִּגְדֵ֤י הַקֹּ֙דֶשׁ֙ לְאַהֲרֹ֣ן הַכֹּהֵ֔ן וְאֶת־בִּגְדֵ֥י בָנָ֖יו לְכַהֵֽן׃ %
% ---------- פרק 39 | פסוק 42 ----------
 \textbf{מב} כְּכֹ֛ל אֲשֶׁר־צִוָּ֥ה יְהֹוָ֖ה אֶת־מֹשֶׁ֑ה כֵּ֤ן עָשׂוּ֙ בְּנֵ֣י יִשְׂרָאֵ֔ל אֵ֖ת כׇּל־הָעֲבֹדָֽה׃ %
% ---------- פרק 39 | פסוק 43 ----------
 \textbf{מג} וַיַּ֨רְא מֹשֶׁ֜ה אֶת־כׇּל־הַמְּלָאכָ֗ה וְהִנֵּה֙ עָשׂ֣וּ אֹתָ֔הּ כַּאֲשֶׁ֛ר צִוָּ֥ה יְהֹוָ֖ה כֵּ֣ן עָשׂ֑וּ וַיְבָ֥רֶךְ אֹתָ֖ם מֹשֶֽׁה׃ \{פ\} \Rashi{\textbf{ויברך אתם משה.}אמר להם יהי רצון שתשרה שכינה במעשה ידיכם, ויהי נעם ה' אלהינו עלינו וגו', והוא אחד מי"א מזמורים שבתפלה למשה (ספרא):}%
% ---------- פרק 40 | פסוק 1 ----------
 \textbf{א} וַיְדַבֵּ֥ר יְהֹוָ֖ה אֶל־מֹשֶׁ֥ה לֵּאמֹֽר׃ %
% ---------- פרק 40 | פסוק 2 ----------
 \textbf{ב} בְּיוֹם־הַחֹ֥דֶשׁ הָרִאשׁ֖וֹן בְּאֶחָ֣ד לַחֹ֑דֶשׁ תָּקִ֕ים אֶת־מִשְׁכַּ֖ן אֹ֥הֶל מוֹעֵֽד׃ %
% ---------- פרק 40 | פסוק 3 ----------
 \textbf{ג} וְשַׂמְתָּ֣ שָׁ֔ם אֵ֖ת אֲר֣וֹן הָעֵד֑וּת וְסַכֹּתָ֥ עַל־הָאָרֹ֖ן אֶת־הַפָּרֹֽכֶת׃ \Rashi{\textbf{וסכת על הארן.}לשון הגנה, שהרי מחצה היתה:}%
% ---------- פרק 40 | פסוק 4 ----------
 \textbf{ד} וְהֵבֵאתָ֙ אֶת־הַשֻּׁלְחָ֔ן וְעָרַכְתָּ֖ אֶת־עֶרְכּ֑וֹ וְהֵבֵאתָ֙ אֶת־הַמְּנֹרָ֔ה וְהַעֲלֵיתָ֖ אֶת־נֵרֹתֶֽיהָ׃ \Rashi{\textbf{וערכת את ערכו.}שתים מערכות של לחם הפנים:}%
% ---------- פרק 40 | פסוק 5 ----------
 \textbf{ה} וְנָתַתָּ֞ה אֶת־מִזְבַּ֤ח הַזָּהָב֙ לִקְטֹ֔רֶת לִפְנֵ֖י אֲר֣וֹן הָעֵדֻ֑ת וְשַׂמְתָּ֛ אֶת־מָסַ֥ךְ הַפֶּ֖תַח לַמִּשְׁכָּֽן׃ %
% ---------- פרק 40 | פסוק 6 ----------
 \textbf{ו} וְנָ֣תַתָּ֔ה אֵ֖ת מִזְבַּ֣ח הָעֹלָ֑ה לִפְנֵ֕י פֶּ֖תַח מִשְׁכַּ֥ן אֹֽהֶל־מוֹעֵֽד׃ %
% ---------- פרק 40 | פסוק 7 ----------
 \textbf{ז} וְנָֽתַתָּ֙ אֶת־הַכִּיֹּ֔ר בֵּֽין־אֹ֥הֶל מוֹעֵ֖ד וּבֵ֣ין הַמִּזְבֵּ֑חַ וְנָתַתָּ֥ שָׁ֖ם מָֽיִם׃ %
% ---------- פרק 40 | פסוק 8 ----------
 \textbf{ח} וְשַׂמְתָּ֥ אֶת־הֶחָצֵ֖ר סָבִ֑יב וְנָ֣תַתָּ֔ אֶת־מָסַ֖ךְ שַׁ֥עַר הֶחָצֵֽר׃ %
% ---------- פרק 40 | פסוק 9 ----------
 \textbf{ט} וְלָקַחְתָּ֙ אֶת־שֶׁ֣מֶן הַמִּשְׁחָ֔ה וּמָשַׁחְתָּ֥ אֶת־הַמִּשְׁכָּ֖ן וְאֶת־כׇּל־אֲשֶׁר־בּ֑וֹ וְקִדַּשְׁתָּ֥ אֹת֛וֹ וְאֶת־כׇּל־כֵּלָ֖יו וְהָ֥יָה קֹֽדֶשׁ׃ %
% ---------- פרק 40 | פסוק 10 ----------
 \textbf{י} וּמָשַׁחְתָּ֛ אֶת־מִזְבַּ֥ח הָעֹלָ֖ה וְאֶת־כׇּל־כֵּלָ֑יו וְקִדַּשְׁתָּ֙ אֶת־הַמִּזְבֵּ֔חַ וְהָיָ֥ה הַמִּזְבֵּ֖חַ קֹ֥דֶשׁ קׇֽדָשִֽׁים׃ %
% ---------- פרק 40 | פסוק 11 ----------
 \textbf{יא} וּמָשַׁחְתָּ֥ אֶת־הַכִּיֹּ֖ר וְאֶת־כַּנּ֑וֹ וְקִדַּשְׁתָּ֖ אֹתֽוֹ׃ %
% ---------- פרק 40 | פסוק 12 ----------
 \textbf{יב} וְהִקְרַבְתָּ֤ אֶֽת־אַהֲרֹן֙ וְאֶת־בָּנָ֔יו אֶל־פֶּ֖תַח אֹ֣הֶל מוֹעֵ֑ד וְרָחַצְתָּ֥ אֹתָ֖ם בַּמָּֽיִם׃ %
% ---------- פרק 40 | פסוק 13 ----------
 \textbf{יג} וְהִלְבַּשְׁתָּ֙ אֶֽת־אַהֲרֹ֔ן אֵ֖ת בִּגְדֵ֣י הַקֹּ֑דֶשׁ וּמָשַׁחְתָּ֥ אֹת֛וֹ וְקִדַּשְׁתָּ֥ אֹת֖וֹ וְכִהֵ֥ן לִֽי׃ %
% ---------- פרק 40 | פסוק 14 ----------
 \textbf{יד} וְאֶת־בָּנָ֖יו תַּקְרִ֑יב וְהִלְבַּשְׁתָּ֥ אֹתָ֖ם כֻּתֳּנֹֽת׃ %
% ---------- פרק 40 | פסוק 15 ----------
 \textbf{טו} וּמָשַׁחְתָּ֣ אֹתָ֗ם כַּאֲשֶׁ֤ר מָשַׁ֙חְתָּ֙ אֶת־אֲבִיהֶ֔ם וְכִהֲנ֖וּ לִ֑י וְ֠הָיְתָ֠ה לִהְיֹ֨ת לָהֶ֧ם מׇשְׁחָתָ֛ם לִכְהֻנַּ֥ת עוֹלָ֖ם לְדֹרֹתָֽם׃ %
% ---------- פרק 40 | פסוק 16 ----------
 \textbf{טז} וַיַּ֖עַשׂ מֹשֶׁ֑ה כְּ֠כֹ֠ל אֲשֶׁ֨ר צִוָּ֧ה יְהֹוָ֛ה אֹת֖וֹ כֵּ֥ן עָשָֽׂה׃ \{ס\} %
% ---------- פרק 40 | פסוק 17 ----------
 \textbf{יז} וַיְהִ֞י בַּחֹ֧דֶשׁ הָרִאשׁ֛וֹן בַּשָּׁנָ֥ה הַשֵּׁנִ֖ית בְּאֶחָ֣ד לַחֹ֑דֶשׁ הוּקַ֖ם הַמִּשְׁכָּֽן׃ %
% ---------- פרק 40 | פסוק 18 ----------
 \textbf{יח} וַיָּ֨קֶם מֹשֶׁ֜ה אֶת־הַמִּשְׁכָּ֗ן וַיִּתֵּן֙ אֶת־אֲדָנָ֔יו וַיָּ֙שֶׂם֙ אֶת־קְרָשָׁ֔יו וַיִּתֵּ֖ן אֶת־בְּרִיחָ֑יו וַיָּ֖קֶם אֶת־עַמּוּדָֽיו׃ %
% ---------- פרק 40 | פסוק 19 ----------
 \textbf{יט} וַיִּפְרֹ֤שׂ אֶת־הָאֹ֙הֶל֙ עַל־הַמִּשְׁכָּ֔ן וַיָּ֜שֶׂם אֶת־מִכְסֵ֥ה הָאֹ֛הֶל עָלָ֖יו מִלְמָ֑עְלָה כַּאֲשֶׁ֛ר צִוָּ֥ה יְהֹוָ֖ה אֶת־מֹשֶֽׁה׃ \{ס\} \Rashi{\textbf{ויפרש את האהל.}הן יריעות העזים:}%
% ---------- פרק 40 | פסוק 20 ----------
 \textbf{כ} וַיִּקַּ֞ח וַיִּתֵּ֤ן אֶת־הָעֵדֻת֙ אֶל־הָ֣אָרֹ֔ן וַיָּ֥שֶׂם אֶת־הַבַּדִּ֖ים עַל־הָאָרֹ֑ן וַיִּתֵּ֧ן אֶת־הַכַּפֹּ֛רֶת עַל־הָאָרֹ֖ן מִלְמָֽעְלָה׃ \Rashi{\textbf{את העדת.}הלוחות:}%
% ---------- פרק 40 | פסוק 21 ----------
 \textbf{כא} וַיָּבֵ֣א אֶת־הָאָרֹן֮ אֶל־הַמִּשְׁכָּן֒ וַיָּ֗שֶׂם אֵ֚ת פָּרֹ֣כֶת הַמָּסָ֔ךְ וַיָּ֕סֶךְ עַ֖ל אֲר֣וֹן הָעֵד֑וּת כַּאֲשֶׁ֛ר צִוָּ֥ה יְהֹוָ֖ה אֶת־מֹשֶֽׁה׃ \{ס\} %
% ---------- פרק 40 | פסוק 22 ----------
 \textbf{כב} וַיִּתֵּ֤ן אֶת־הַשֻּׁלְחָן֙ בְּאֹ֣הֶל מוֹעֵ֔ד עַ֛ל יֶ֥רֶךְ הַמִּשְׁכָּ֖ן צָפֹ֑נָה מִח֖וּץ לַפָּרֹֽכֶת׃ \Rashi{\textbf{על ירך המשכן צפנה.}בחצי הצפוני של רחב הבית: ירך. כתרגומו, צדא, כירך הזה שהוא בצדו של אדם:}%
% ---------- פרק 40 | פסוק 23 ----------
 \textbf{כג} וַיַּעֲרֹ֥ךְ עָלָ֛יו עֵ֥רֶךְ לֶ֖חֶם לִפְנֵ֣י יְהֹוָ֑ה כַּאֲשֶׁ֛ר צִוָּ֥ה יְהֹוָ֖ה אֶת־מֹשֶֽׁה׃ \{ס\} %
% ---------- פרק 40 | פסוק 24 ----------
 \textbf{כד} וַיָּ֤שֶׂם אֶת־הַמְּנֹרָה֙ בְּאֹ֣הֶל מוֹעֵ֔ד נֹ֖כַח הַשֻּׁלְחָ֑ן עַ֛ל יֶ֥רֶךְ הַמִּשְׁכָּ֖ן נֶֽגְבָּה׃ %
% ---------- פרק 40 | פסוק 25 ----------
 \textbf{כה} וַיַּ֥עַל הַנֵּרֹ֖ת לִפְנֵ֣י יְהֹוָ֑ה כַּאֲשֶׁ֛ר צִוָּ֥ה יְהֹוָ֖ה אֶת־מֹשֶֽׁה׃ \{ס\} %
% ---------- פרק 40 | פסוק 26 ----------
 \textbf{כו} וַיָּ֛שֶׂם אֶת־מִזְבַּ֥ח הַזָּהָ֖ב בְּאֹ֣הֶל מוֹעֵ֑ד לִפְנֵ֖י הַפָּרֹֽכֶת׃ %
% ---------- פרק 40 | פסוק 27 ----------
 \textbf{כז} וַיַּקְטֵ֥ר עָלָ֖יו קְטֹ֣רֶת סַמִּ֑ים כַּאֲשֶׁ֛ר צִוָּ֥ה יְהֹוָ֖ה אֶת־מֹשֶֽׁה׃ \{ס\} \Rashi{\textbf{ויקטר עליו קטרת.}שחרית וערבית, כמו שנאמר בבקר בבקר בהיטיבו את הנרות וגומר (שמות ל'):}%
% ---------- פרק 40 | פסוק 28 ----------
 \textbf{כח} וַיָּ֛שֶׂם אֶת־מָסַ֥ךְ הַפֶּ֖תַח לַמִּשְׁכָּֽן׃ %
% ---------- פרק 40 | פסוק 29 ----------
 \textbf{כט} וְאֵת֙ מִזְבַּ֣ח הָעֹלָ֔ה שָׂ֕ם פֶּ֖תַח מִשְׁכַּ֣ן אֹֽהֶל־מוֹעֵ֑ד וַיַּ֣עַל עָלָ֗יו אֶת־הָעֹלָה֙ וְאֶת־הַמִּנְחָ֔ה כַּאֲשֶׁ֛ר צִוָּ֥ה יְהֹוָ֖ה אֶת־מֹשֶֽׁה׃ \{ס\} \Rashi{\textbf{ויעל עליו וגו'.}אף ביום השמיני למלואים, שהוא יום הקמת המשכן, שמש משה והקריב קרבנות צבור, חוץ מאותן שנצטוו בו ביום, שנאמר קרב אל המזבח וגומר (ויקרא ט'): את העלה. עולת התמיד: ואת המנחה. מנחת נסכים של תמיד, כמו שנאמר ועשרן סלת בלול בשמן וגומר (שמות כ"ט):}%
% ---------- פרק 40 | פסוק 30 ----------
 \textbf{ל} וַיָּ֙שֶׂם֙ אֶת־הַכִּיֹּ֔ר בֵּֽין־אֹ֥הֶל מוֹעֵ֖ד וּבֵ֣ין הַמִּזְבֵּ֑חַ וַיִּתֵּ֥ן שָׁ֛מָּה מַ֖יִם לְרׇחְצָֽה׃ %
% ---------- פרק 40 | פסוק 31 ----------
 \textbf{לא} וְרָחֲצ֣וּ מִמֶּ֔נּוּ מֹשֶׁ֖ה וְאַהֲרֹ֣ן וּבָנָ֑יו אֶת־יְדֵיהֶ֖ם וְאֶת־רַגְלֵיהֶֽם׃ \Rashi{\textbf{ורחצו ממנו משה ואהרן.}יום שמיני למלואים השוו כלם לכהנה, ותרגומו ומקדשין מנה, בו ביום קדש משה עמהם:}%
% ---------- פרק 40 | פסוק 32 ----------
 \textbf{לב} בְּבֹאָ֞ם אֶל־אֹ֣הֶל מוֹעֵ֗ד וּבְקׇרְבָתָ֛ם אֶל־הַמִּזְבֵּ֖חַ יִרְחָ֑צוּ כַּאֲשֶׁ֛ר צִוָּ֥ה יְהֹוָ֖ה אֶת־מֹשֶֽׁה׃ \{ס\} \Rashi{\textbf{ובקרבתם.}כמו ובקרבם – כשיקרבו:}%
% ---------- פרק 40 | פסוק 33 ----------
 \textbf{לג} וַיָּ֣קֶם אֶת־הֶחָצֵ֗ר סָבִיב֙ לַמִּשְׁכָּ֣ן וְלַמִּזְבֵּ֔חַ וַיִּתֵּ֕ן אֶת־מָסַ֖ךְ שַׁ֣עַר הֶחָצֵ֑ר וַיְכַ֥ל מֹשֶׁ֖ה אֶת־הַמְּלָאכָֽה׃ \{פ\} %
% ---------- פרק 40 | פסוק 34 ----------
 \textbf{לד} וַיְכַ֥ס הֶעָנָ֖ן אֶת־אֹ֣הֶל מוֹעֵ֑ד וּכְב֣וֹד יְהֹוָ֔ה מָלֵ֖א אֶת־הַמִּשְׁכָּֽן׃ %
% ---------- פרק 40 | פסוק 35 ----------
 \textbf{לה} וְלֹא־יָכֹ֣ל מֹשֶׁ֗ה לָבוֹא֙ אֶל־אֹ֣הֶל מוֹעֵ֔ד כִּֽי־שָׁכַ֥ן עָלָ֖יו הֶעָנָ֑ן וּכְב֣וֹד יְהֹוָ֔ה מָלֵ֖א אֶת־הַמִּשְׁכָּֽן׃ \Rashi{\textbf{ולא יכל משה לבוא אל אהל מועד.}וכתוב אחד אומר ובבא משה אל אהל מועד (במדבר ז')? בא הכתוב השלישי והכריע ביניהם – כי שכן עליו הענן, אמר מעתה כל זמן שהיה עליו הענן לא היה יכול לבא, נסתלק הענן נכנס ומדבר עמו (ספרא):}%
% ---------- פרק 40 | פסוק 36 ----------
 \textbf{לו} וּבְהֵעָל֤וֹת הֶֽעָנָן֙ מֵעַ֣ל הַמִּשְׁכָּ֔ן יִסְע֖וּ בְּנֵ֣י יִשְׂרָאֵ֑ל בְּכֹ֖ל מַסְעֵיהֶֽם׃ %
% ---------- פרק 40 | פסוק 37 ----------
 \textbf{לז} וְאִם־לֹ֥א יֵעָלֶ֖ה הֶעָנָ֑ן וְלֹ֣א יִסְע֔וּ עַד־י֖וֹם הֵעָלֹתֽוֹ׃ %
% ---------- פרק 40 | פסוק 38 ----------
 \textbf{לח} כִּי֩ עֲנַ֨ן יְהֹוָ֤ה עַֽל־הַמִּשְׁכָּן֙ יוֹמָ֔ם וְאֵ֕שׁ תִּהְיֶ֥ה לַ֖יְלָה בּ֑וֹ לְעֵינֵ֥י כׇל־בֵּֽית־יִשְׂרָאֵ֖ל בְּכׇל־מַסְעֵיהֶֽם׃ \Rashi{\textbf{לעיני כל בית ישראל בכל מסעיהם.}בכל מסע שהיו נוסעים, היה הענן שוכן במקום אשר יחנו שם; מקום חניתן אף הוא קרוי מסע, וכן וילך למסעיו (בראשית י"ג), וכן אלה מסעי (במדבר ל"ג), לפי שממקום החניה חזרו ונסעו, לכך נקראו כלן מסעות:}%
\addparasha{חומש שמות | פרשת שמות}
% ---------- פרק 1 | פסוק 1 ----------
 \textbf{א} וְאֵ֗לֶּה שְׁמוֹת֙ בְּנֵ֣י יִשְׂרָאֵ֔ל הַבָּאִ֖ים מִצְרָ֑יְמָה אֵ֣ת יַעֲקֹ֔ב אִ֥ישׁ וּבֵית֖וֹ בָּֽאוּ׃ \Rashi{\textbf{ואלה שמות בני ישראל.}אע"פ שמנאן בחייהם בשמותם, חזר ומנאם במיתתם, להודיע חבתם, שנמשלו לכוכבים, שמוציאם ומכניסם במספר ובשמותם (שמות רבה), שנ' "המוציא במספר צבאם לכלם בשם יקרא" (ישעיהו מ'):}%
% ---------- פרק 1 | פסוק 2 ----------
 \textbf{ב} רְאוּבֵ֣ן שִׁמְע֔וֹן לֵוִ֖י וִיהוּדָֽה׃ %
% ---------- פרק 1 | פסוק 3 ----------
 \textbf{ג} יִשָּׂשכָ֥ר זְבוּלֻ֖ן וּבִנְיָמִֽן׃ %
% ---------- פרק 1 | פסוק 4 ----------
 \textbf{ד} דָּ֥ן וְנַפְתָּלִ֖י גָּ֥ד וְאָשֵֽׁר׃ %
% ---------- פרק 1 | פסוק 5 ----------
 \textbf{ה} וַֽיְהִ֗י כׇּל־נֶ֛פֶשׁ יֹצְאֵ֥י יֶֽרֶךְ־יַעֲקֹ֖ב שִׁבְעִ֣ים נָ֑פֶשׁ וְיוֹסֵ֖ף הָיָ֥ה בְמִצְרָֽיִם׃ \Rashi{\textbf{ויוסף היה במצרים.}והלא הוא ובניו היו בכלל שבעים, ומה בא ללמדנו? וכי לא היינו יודעים שהוא היה במצרים? אלא להודיעך צדקתו של יוסף, הוא יוסף הרועה את צאן אביו, הוא יוסף שהיה במצרים, ונעשה מלך, ועומד בצדקו (שמות רבה):}%
% ---------- פרק 1 | פסוק 6 ----------
 \textbf{ו} וַיָּ֤מׇת יוֹסֵף֙ וְכׇל־אֶחָ֔יו וְכֹ֖ל הַדּ֥וֹר הַהֽוּא׃ %
% ---------- פרק 1 | פסוק 7 ----------
 \textbf{ז} וּבְנֵ֣י יִשְׂרָאֵ֗ל פָּר֧וּ וַֽיִּשְׁרְצ֛וּ וַיִּרְבּ֥וּ וַיַּֽעַצְמ֖וּ בִּמְאֹ֣ד מְאֹ֑ד וַתִּמָּלֵ֥א הָאָ֖רֶץ אֹתָֽם׃ \{פ\} \Rashi{\textbf{פרו.}שלא הפילו נשותיהם ולא מתו כשהם קטנים:  וישרצו. שהיו יולדות ששה בכרס אחד (שם):}%
% ---------- פרק 1 | פסוק 8 ----------
 \textbf{ח} וַיָּ֥קׇם מֶֽלֶךְ־חָדָ֖שׁ עַל־מִצְרָ֑יִם אֲשֶׁ֥ר לֹֽא־יָדַ֖ע אֶת־יוֹסֵֽף׃ \Rashi{\textbf{ויקם מלך חדש.}רב ושמואל, חד אמר חדש ממש, וחד אמר שנתחדשו גזרותיו: אשר לא ידע. עשה עצמו כאלו לא ידעו (סוטה י"א):}%
% ---------- פרק 1 | פסוק 9 ----------
 \textbf{ט} וַיֹּ֖אמֶר אֶל־עַמּ֑וֹ הִנֵּ֗ה עַ֚ם בְּנֵ֣י יִשְׂרָאֵ֔ל רַ֥ב וְעָצ֖וּם מִמֶּֽנּוּ׃ %
% ---------- פרק 1 | פסוק 10 ----------
 \textbf{י} הָ֥בָה נִֽתְחַכְּמָ֖ה ל֑וֹ פֶּן־יִרְבֶּ֗ה וְהָיָ֞ה כִּֽי־תִקְרֶ֤אנָה מִלְחָמָה֙ וְנוֹסַ֤ף גַּם־הוּא֙ עַל־שֹׂ֣נְאֵ֔ינוּ וְנִלְחַם־בָּ֖נוּ וְעָלָ֥ה מִן־הָאָֽרֶץ׃ \Rashi{\textbf{הבה נתחכמה.}כל הבה לשון הכנה והזמנה לדבר הוא, כלומר, הזמינו עצמכם לכך: נתחכמה לו. לעם; נתחכמה מה לעשות לו. ורבותינו דרשו נתחכם למושיען של ישראל לדונם במים, שכבר נשבע שלא יביא מבול לעולם (שם): ועלה מן הארץ. על כרחנו. ורבותינו דרשו כאדם שמקלל עצמו ותולה קללתו באחרים (שם), והרי הוא כאלו כתב ועלינו מן הארץ – והם יירשוה:}%
% ---------- פרק 1 | פסוק 11 ----------
 \textbf{יא} וַיָּשִׂ֤ימוּ עָלָיו֙ שָׂרֵ֣י מִסִּ֔ים לְמַ֥עַן עַנֹּת֖וֹ בְּסִבְלֹתָ֑ם וַיִּ֜בֶן עָרֵ֤י מִסְכְּנוֹת֙ לְפַרְעֹ֔ה אֶת־פִּתֹ֖ם וְאֶת־רַעַמְסֵֽס׃ \Rashi{\textbf{עליו.}על העם: מסים. לשון מס, שרים שגובין מהם המס, ומהו המס? שיבנו ערי מסכנות לפרעה: למען ענתו בסבלתם. של מצרים: ערי מסכנות. כתרגומו, וכן "לך בא אל הסוכן הזה" (ישעיהו כ"ב) – גזבר הממנה על האוצרות: את פתם ואת רעמסס. שלא היו ראויות מתחלה לכך, ועשאום חזקות ובצורות לאוצר:}%
% ---------- פרק 1 | פסוק 12 ----------
 \textbf{יב} וְכַאֲשֶׁר֙ יְעַנּ֣וּ אֹת֔וֹ כֵּ֥ן יִרְבֶּ֖ה וְכֵ֣ן יִפְרֹ֑ץ וַיָּקֻ֕צוּ מִפְּנֵ֖י בְּנֵ֥י יִשְׂרָאֵֽל׃ \Rashi{\textbf{וכאשר יענו אתו.}בכל מה שהם נותנין לב לענות, כן לב הקב"ה להרבות ולהפריץ: כן ירבה. כן רבה וכן פרץ. ומדרשו, רוח הקדש אומרת כן, אתם אומרים פן ירבה ואני אומר כן ירבה (סוטה י"א): ויקצו. קצו בחייהם. ורבותינו דרשו כקוצים היו בעיניהם (שם):}%
% ---------- פרק 1 | פסוק 13 ----------
 \textbf{יג} וַיַּעֲבִ֧דוּ מִצְרַ֛יִם אֶת־בְּנֵ֥י יִשְׂרָאֵ֖ל בְּפָֽרֶךְ׃ \Rashi{\textbf{בפרך.}בעבודה קשה, המפרכת את הגוף ומשברתו (שם):}%
% ---------- פרק 1 | פסוק 14 ----------
 \textbf{יד} וַיְמָרְר֨וּ אֶת־חַיֵּיהֶ֜ם בַּעֲבֹדָ֣ה קָשָׁ֗ה בְּחֹ֙מֶר֙ וּבִלְבֵנִ֔ים וּבְכׇל־עֲבֹדָ֖ה בַּשָּׂדֶ֑ה אֵ֚ת כׇּל־עֲבֹ֣דָתָ֔ם אֲשֶׁר־עָבְד֥וּ בָהֶ֖ם בְּפָֽרֶךְ׃ %
% ---------- פרק 1 | פסוק 15 ----------
 \textbf{טו} וַיֹּ֙אמֶר֙ מֶ֣לֶךְ מִצְרַ֔יִם לַֽמְיַלְּדֹ֖ת הָֽעִבְרִיֹּ֑ת אֲשֶׁ֨ר שֵׁ֤ם הָֽאַחַת֙ שִׁפְרָ֔ה וְשֵׁ֥ם הַשֵּׁנִ֖ית פּוּעָֽה׃ \Rashi{\textbf{למילדת.}הוא לשון מולידות, אלא שיש ל' קל ויש ל' כבד, כמו שובר ומשבר, דובר ומדבר, כך מוליד ומילד: שפרה. יוכבד, על שם שמשפרת את הולד (שם): פועה. זו מרים, שפועה ומדברת והוגה לולד (שם), כדרך הנשים המפיסות תינוק הבוכה. פועה. לשון צעקה, כמו "כיולדה אפעה" (ישעיה מ"ב):}%
% ---------- פרק 1 | פסוק 16 ----------
 \textbf{טז} וַיֹּ֗אמֶר בְּיַלֶּדְכֶן֙ אֶת־הָֽעִבְרִיּ֔וֹת וּרְאִיתֶ֖ן עַל־הָאׇבְנָ֑יִם אִם־בֵּ֥ן הוּא֙ וַהֲמִתֶּ֣ן אֹת֔וֹ וְאִם־בַּ֥ת הִ֖וא וָחָֽיָה׃ \Rashi{\textbf{בילדכן.}כמו בהולידכן: על האבנים. מושב האשה היולדת, במקום אחר (ישעיהו ל"ז), קוראו משבר; וכמוהו "עושה מלאכה על האבנים" (ירמיהו י"ח), – מושב כלי אמנות יוצר חרש: אם בן הוא וגו'. לא היה מקפיד אלא על הזכרים, שאמרו לו אצטגניניו שעתיד להולד בן המושיע אותם: וחיה. ותחיה:}%
% ---------- פרק 1 | פסוק 17 ----------
 \textbf{יז} וַתִּירֶ֤אןָ הַֽמְיַלְּדֹת֙ אֶת־הָ֣אֱלֹהִ֔ים וְלֹ֣א עָשׂ֔וּ כַּאֲשֶׁ֛ר דִּבֶּ֥ר אֲלֵיהֶ֖ן מֶ֣לֶךְ מִצְרָ֑יִם וַתְּחַיֶּ֖יןָ אֶת־הַיְלָדִֽים׃ \Rashi{\textbf{ותחיין את הילדים.}מספקות להם מים ומזון. תרגום הראשון וקימא, והשני וקימתן, לפי שלשון עבר לנקבות רבות תבה זו וכיוצא בה משמשת לשון פעלו ולשון פעלתם, כגון "ותאמרן איש מצרי" (שמות ב'), (לשון עבר) כמו ויאמרו לזכרים, "ותדברנה בפיכם" (ירמיהו מ"ד), לשון דברתן, כמו ותדברו לזכרים, וכן "ותחללנה אתי אל עמי" (יחזקאל י"ג), לשון עבר חללתן, כמו ותחללו לזכרים:}%
% ---------- פרק 1 | פסוק 18 ----------
 \textbf{יח} וַיִּקְרָ֤א מֶֽלֶךְ־מִצְרַ֙יִם֙ לַֽמְיַלְּדֹ֔ת וַיֹּ֣אמֶר לָהֶ֔ן מַדּ֥וּעַ עֲשִׂיתֶ֖ן הַדָּבָ֣ר הַזֶּ֑ה וַתְּחַיֶּ֖יןָ אֶת־הַיְלָדִֽים׃ %
% ---------- פרק 1 | פסוק 19 ----------
 \textbf{יט} וַתֹּאמַ֤רְןָ הַֽמְיַלְּדֹת֙ אֶל־פַּרְעֹ֔ה כִּ֣י לֹ֧א כַנָּשִׁ֛ים הַמִּצְרִיֹּ֖ת הָֽעִבְרִיֹּ֑ת כִּֽי־חָי֣וֹת הֵ֔נָּה בְּטֶ֨רֶם תָּב֧וֹא אֲלֵהֶ֛ן*(בכתר ארם צובה היה כתוב אֲלֵיהֶ֛ן) הַמְיַלֶּ֖דֶת וְיָלָֽדוּ׃ \Rashi{\textbf{כי חיות הנה.}בקיאות כמילדות, תרגום מילדות חיתא. ורבותינו דרשו הרי הן משולות כחיות השדה (סוטה י"א), שאינן צריכות מילדות. והיכן משולות לחיות? גור אריה, זאב יטרף (בראשית מ"ט), בכור שורו (דברים ל"ג), אילה שלחה (בראשית מ"ט), ומי שלא נכתב בו, הרי הכתוב כללן ויברך אותם (שם), ועוד כתיב מה אמך לביא (יחזקאל י"ט):}%
% ---------- פרק 1 | פסוק 20 ----------
 \textbf{כ} וַיֵּ֥יטֶב אֱלֹהִ֖ים לַֽמְיַלְּדֹ֑ת וַיִּ֧רֶב הָעָ֛ם וַיַּֽעַצְמ֖וּ מְאֹֽד׃ \Rashi{\textbf{וייטב.}הטיב להן. וזה חלוק בתבה שיסודה שתי אותיות ונתן לה וי"ו יו"ד בראשה, כשהיא באה לדבר לשון ויפעיל הוא נקוד היו"ד בצירי, שהוא קמץ קטן, כגון וייטב אלהים למילדת, "וירב בבת יהודה" (איכה ב'), – הרבה תאניה. וכן "ויגל השארית" דנבוזראדן (דברי הימים ב' ל"ו), – הגלה את השארית, "ויפן זנב אל זנב" (שופטים ט"ו) – הפנה הזנבות זו לזו, כל אלו לשון הפעיל את אחרים; וכשהוא מדבר בלשון ויפעל, הוא נקוד היו"ד בחירק, כגון "וייטב בעיניו" (ויקרא י'), לשון הוטב, וכן וירב העם, – נתרבה העם, "ויגל יהודה" (מלכים ב כ"ה) – הגלה יהודה, "ויפן כה וכה" (שמות ב׳:י״ב) – הפנה לכאן ולכאן. ואל תשיבני וילך, וישב, וירד, ויצא, לפי שאינן מגזרתן של אלו, שהרי היו"ד יסוד בהן, ילך, ישב, ירד, יצא – י' אות שלישית בו: וייטב אלהים למילדת. מהו הטובה?:}%
% ---------- פרק 1 | פסוק 21 ----------
 \textbf{כא} וַיְהִ֕י כִּֽי־יָרְא֥וּ הַֽמְיַלְּדֹ֖ת אֶת־הָאֱלֹהִ֑ים וַיַּ֥עַשׂ לָהֶ֖ם בָּתִּֽים׃ \Rashi{\textbf{ויעש להם בתים.}בתי כהנה ולויה ומלכות שקרויין בתים, כמו שכתוב: "לבנות את בית ה' ואת בית המלך" (מלכים א ט׳:א׳), כהנה ולויה מיוכבד ומלכות ממרים. כדאיתא במסכת סוטה:}%
% ---------- פרק 1 | פסוק 22 ----------
 \textbf{כב} וַיְצַ֣ו פַּרְעֹ֔ה לְכׇל־עַמּ֖וֹ לֵאמֹ֑ר כׇּל־הַבֵּ֣ן הַיִּלּ֗וֹד הַיְאֹ֙רָה֙ תַּשְׁלִיכֻ֔הוּ וְכׇל־הַבַּ֖ת תְּחַיּֽוּן׃ \{פ\} \Rashi{\textbf{לכל עמו.}אף עליהם גזר (סוטה י"ב), יום שנולד משה אמרו לו אצטגניניו, היום נולד מושיען, ואין אנו יודעים אם ממצרים אם מישראל, ורואין אנו שסופו ללקות במים, לפיכך גזר אותו היום אף על המצריים, שנאמר כל הבן הילוד, ולא נאמר הילוד לעברים; והם לא היו יודעים שסופו ללקות על מי מריבה:}%
% ---------- פרק 2 | פסוק 1 ----------
 \textbf{א} וַיֵּ֥לֶךְ אִ֖ישׁ מִבֵּ֣ית לֵוִ֑י וַיִּקַּ֖ח אֶת־בַּת־לֵוִֽי׃ \Rashi{\textbf{ויקח את בת לוי.}פרוש היה ממנה מפני גזרת פרעה, והחזירה ועשה בה לקוחין שניים, ואף היא נהפכה להיות נערה; ובת ק"ל שנה היתה, שנולדה בבואם למצרים בין החומות, ומאתים ועשר נשתהו שם, וכשיצאו היה משה בן שמונים שנה, אם כן כשנתעברה ממנו היתה בת מאה ושלושים וקורא אותה בת לוי (עי' סוטה י"ב, בבא בתרא קי"ט ושמות רבה):}%
% ---------- פרק 2 | פסוק 2 ----------
 \textbf{ב} וַתַּ֥הַר הָאִשָּׁ֖ה וַתֵּ֣לֶד בֵּ֑ן וַתֵּ֤רֶא אֹתוֹ֙ כִּי־ט֣וֹב ה֔וּא וַֽתִּצְפְּנֵ֖הוּ שְׁלֹשָׁ֥ה יְרָחִֽים׃ \Rashi{\textbf{כי טוב הוא.}כשנולד נתמלא הבית כלו אורה (סוטה י"ב):}%
% ---------- פרק 2 | פסוק 3 ----------
 \textbf{ג} וְלֹא־יָכְלָ֣ה עוֹד֮ הַצְּפִינוֹ֒ וַתִּֽקַּֽח־לוֹ֙ תֵּ֣בַת גֹּ֔מֶא וַתַּחְמְרָ֥הֿ בַחֵמָ֖ר וּבַזָּ֑פֶת וַתָּ֤שֶׂם בָּהּ֙ אֶת־הַיֶּ֔לֶד וַתָּ֥שֶׂם בַּסּ֖וּף עַל־שְׂפַ֥ת הַיְאֹֽר׃ \Rashi{\textbf{ולא יכלה עוד הצפינו.}שמנו לה המצריים מיום שהחזירה, והיא ילדתו לששה חדשים ויום אחד (שם), שהיולדת לשבעה יולדת למקטעין (נדה ל"ח), והם בדקו אחריה לסוף תשעה: גמא. גמי בלשון משנה (שבת ע"ח), ובלעז יונק"ו, ודבר רך הוא ועומד בפני רך ובפני קשה (סוטה י"ב): בחמר ובזפת. זפת מבחוץ וטיט מבפנים, כדי שלא יריח אותו צדיק ריח רע של זפת (שם): ותשם בסוף. הוא לשון אגם, רושיי"ל בלעז, ודומה לו "קנה וסוף קמלו" (ישעיהו י"ט):}%
% ---------- פרק 2 | פסוק 4 ----------
 \textbf{ד} וַתֵּתַצַּ֥ב אֲחֹת֖וֹ מֵרָחֹ֑ק לְדֵעָ֕ה מַה־יֵּעָשֶׂ֖ה לֽוֹ׃ %
% ---------- פרק 2 | פסוק 5 ----------
 \textbf{ה} וַתֵּ֤רֶד בַּת־פַּרְעֹה֙ לִרְחֹ֣ץ עַל־הַיְאֹ֔ר וְנַעֲרֹתֶ֥יהָ הֹלְכֹ֖ת עַל־יַ֣ד הַיְאֹ֑ר וַתֵּ֤רֶא אֶת־הַתֵּבָה֙ בְּת֣וֹךְ הַסּ֔וּף וַתִּשְׁלַ֥ח אֶת־אֲמָתָ֖הּ וַתִּקָּחֶֽהָ׃ \Rashi{\textbf{לרחץ על היאר.}סרס המקרא ופרשהו: ותרד בת פרעה על היאר לרחץ בו: על יד היאר. אצל היאור, כמו "ראו חלקת יואב אל ידי" (שמואל ב י"ד), והוא לשון יד ממש, שיד האדם סמוכה לו. ורבותינו אמרו, "הולכות" לשון מיתה (סוטה י"ב), כמו "הנה אנכי הולך למות" (בראשית כ"ה) – הולכות למות לפי שמחו בה; והכתוב מסיען, כי למה לנו לכתב ונערותיה הולכות: את אמתה. את שפחתה. ורבותינו דרשו (סוטה שם), לשון יד, אבל לפי דקדוק לשון הקדש היה לו להנקד אמתה, דגושה, והם דרשו את אמתה – את ידה, ונשתרבבה אמתה אמות הרבה:}%
% ---------- פרק 2 | פסוק 6 ----------
 \textbf{ו} וַתִּפְתַּח֙ וַתִּרְאֵ֣הוּ אֶת־הַיֶּ֔לֶד וְהִנֵּה־נַ֖עַר בֹּכֶ֑ה וַתַּחְמֹ֣ל עָלָ֔יו וַתֹּ֕אמֶר מִיַּלְדֵ֥י הָֽעִבְרִ֖ים זֶֽה׃ \Rashi{\textbf{ותפתח ותראהו.}את מי ראתה? את הילד, זהו פשוטו. ומדרשו (שם), שראתה עמו שכינה: והנה נער בכה. קולו כנער (שם):}%
% ---------- פרק 2 | פסוק 7 ----------
 \textbf{ז} וַתֹּ֣אמֶר אֲחֹתוֹ֮ אֶל־בַּת־פַּרְעֹה֒ הַאֵלֵ֗ךְ וְקָרָ֤אתִי לָךְ֙ אִשָּׁ֣ה מֵינֶ֔קֶת מִ֖ן הָעִבְרִיֹּ֑ת וְתֵינִ֥ק לָ֖ךְ אֶת־הַיָּֽלֶד׃ \Rashi{\textbf{מן העברית.}שהחזירתו על מצריות הרבה לינק ולא ינק, לפי שהיה עתיד לדבר עם השכינה (שם):}%
% ---------- פרק 2 | פסוק 8 ----------
 \textbf{ח} וַתֹּֽאמֶר־לָ֥הּ בַּת־פַּרְעֹ֖ה לֵ֑כִי וַתֵּ֙לֶךְ֙ הָֽעַלְמָ֔ה וַתִּקְרָ֖א אֶת־אֵ֥ם הַיָּֽלֶד׃ \Rashi{\textbf{ותלך העלמה.}הלכה בזריזות ועלמות כעלם (שם):}%
% ---------- פרק 2 | פסוק 9 ----------
 \textbf{ט} וַתֹּ֧אמֶר לָ֣הּ בַּת־פַּרְעֹ֗ה הֵילִ֜יכִי אֶת־הַיֶּ֤לֶד הַזֶּה֙ וְהֵינִקִ֣הוּ לִ֔י וַאֲנִ֖י אֶתֵּ֣ן אֶת־שְׂכָרֵ֑ךְ וַתִּקַּ֧ח הָאִשָּׁ֛ה הַיֶּ֖לֶד וַתְּנִיקֵֽהוּ׃ \Rashi{\textbf{היליכי.}נתנבאה ולא ידעה מה נתנבאה – הי שליכי:}%
% ---------- פרק 2 | פסוק 10 ----------
 \textbf{י} וַיִּגְדַּ֣ל הַיֶּ֗לֶד וַתְּבִאֵ֙הוּ֙ לְבַת־פַּרְעֹ֔ה וַֽיְהִי־לָ֖הּ לְבֵ֑ן וַתִּקְרָ֤א שְׁמוֹ֙ מֹשֶׁ֔ה וַתֹּ֕אמֶר כִּ֥י מִן־הַמַּ֖יִם מְשִׁיתִֽהוּ׃ \Rashi{\textbf{משיתהו.}שחליתה, הוא לשון הוצאה בלשון ארמי, "כמשחל בינתא מחלבא", ובלשון עברי משיתיהו לשון הסירותיו, כמו לא ימוש, לא משו, כך חברו מנחם. ואני אומר שאינו ממחברת מש וימוש אלא מגזרת משה, ולשון הוצאה הוא, וכן "ימשני ממים רבים", (שמואל ב כ"ב), שאלו היה ממחברת מש לא יתכן לומר משיתיהו אלא המישותיהו, כאשר יאמר מן קם הקימותי, ומן שב השיבותי, ומן בא הביאותי, או משתיהו, כמו "ומשתי את עון הארץ" (זכריה ג'), אבל משיתי אינו אלא מגזרת תבה שפעל שלה מיסד בה"א בסוף התבה, כגון משה, בנה, עשה, צוה, פנה, כשיבא לומר בהם פעלתי, תבא היו"ד במקום ה"א, כמו בניתי, עשיתי, צויתי:}%
% ---------- פרק 2 | פסוק 11 ----------
 \textbf{יא} וַיְהִ֣י ׀ בַּיָּמִ֣ים הָהֵ֗ם וַיִּגְדַּ֤ל מֹשֶׁה֙ וַיֵּצֵ֣א אֶל־אֶחָ֔יו וַיַּ֖רְא בְּסִבְלֹתָ֑ם וַיַּרְא֙ אִ֣ישׁ מִצְרִ֔י מַכֶּ֥ה אִישׁ־עִבְרִ֖י מֵאֶחָֽיו׃ \Rashi{\textbf{ויגדל משה.}והלא כבר כתב ויגדל הילד? אמר רבי יהודה ברבי אלעאי, הראשון לקומה והשני לגדלה, שמנהו פרעה על ביתו (ילקוט שמעוני): וירא בסבלתם. נתן עיניו ולבו להיות מצר עליהם (שמות רבה א'): איש מצרי. נוגש היה ממנה על שוטרי ישראל והיה מעמידם מקרות הגבר למלאכתם (שם): מכה איש עברי. מלקהו ורודהו. ובעלה של שלומית בת דברי היה ונתן עיניו בה, ובלילה העמידו והוציאו מביתו, והוא חזר ונכנס לבית ובא על אשתו, כסבורה שהוא בעלה, וחזר האיש לביתו והרגיש בדבר, וכשראה אותו מצרי שהרגיש בדבר, היה מכהו ורודהו כל היום (שם):}%
% ---------- פרק 2 | פסוק 12 ----------
 \textbf{יב} וַיִּ֤פֶן כֹּה֙ וָכֹ֔ה וַיַּ֖רְא כִּ֣י אֵ֣ין אִ֑ישׁ וַיַּךְ֙ אֶת־הַמִּצְרִ֔י וַֽיִּטְמְנֵ֖הוּ בַּחֽוֹל׃ \Rashi{\textbf{ויפן כה וכה.}ראה מה עשה לו בבית ומה עשה לו בשדה (שם). ולפי פשוטו כמשמעו: וירא כי אין איש. עתיד לצאת ממנו שיתגיר (ת"י):}%
% ---------- פרק 2 | פסוק 13 ----------
 \textbf{יג} וַיֵּצֵא֙ בַּיּ֣וֹם הַשֵּׁנִ֔י וְהִנֵּ֛ה שְׁנֵֽי־אֲנָשִׁ֥ים עִבְרִ֖ים נִצִּ֑ים וַיֹּ֙אמֶר֙ לָֽרָשָׁ֔ע לָ֥מָּה תַכֶּ֖ה רֵעֶֽךָ׃ \Rashi{\textbf{שני אנשים עברים.}דתן ואבירם, הם שהותירו מן המן (נדרים ס"ד): נצים. מריבים: למה תכה. אע"פ שלא הכהו נקרא רשע בהרמת יד (סנהדרין נ"ח): רעך. רשע כמותך (תנחומא):}%
% ---------- פרק 2 | פסוק 14 ----------
 \textbf{יד} וַ֠יֹּ֠אמֶר מִ֣י שָֽׂמְךָ֞ לְאִ֨ישׁ שַׂ֤ר וְשֹׁפֵט֙ עָלֵ֔ינוּ הַלְהׇרְגֵ֙נִי֙ אַתָּ֣ה אֹמֵ֔ר כַּאֲשֶׁ֥ר הָרַ֖גְתָּ אֶת־הַמִּצְרִ֑י וַיִּירָ֤א מֹשֶׁה֙ וַיֹּאמַ֔ר אָכֵ֖ן נוֹדַ֥ע הַדָּבָֽר׃ \Rashi{\textbf{מי שמך לאיש.}והנה עודך נער: הלהרגני אתה אמר. מכאן אנו למדים שהרגו בשם המפרש (שמות רבה א'): ויירא משה. כפשוטו. ומדרשו: דאג לו על שראה בישראל רשעים דלטורין, אמר, מעתה שמא אינם ראויין להגאל (שמות רבה א'): אכן נודע הדבר. כמשמעו. ומדרשו, נודע לי הדבר שהייתי תמה עליו, מה חטאו ישראל מכל שבעים אמות להיות נרדים בעבודת פרך, אבל רואה אני שהם ראויים לכך (שמות רבה א'):}%
% ---------- פרק 2 | פסוק 15 ----------
 \textbf{טו} וַיִּשְׁמַ֤ע פַּרְעֹה֙ אֶת־הַדָּבָ֣ר הַזֶּ֔ה וַיְבַקֵּ֖שׁ לַהֲרֹ֣ג אֶת־מֹשֶׁ֑ה וַיִּבְרַ֤ח מֹשֶׁה֙ מִפְּנֵ֣י פַרְעֹ֔ה וַיֵּ֥שֶׁב בְּאֶֽרֶץ־מִדְיָ֖ן וַיֵּ֥שֶׁב עַֽל־הַבְּאֵֽר׃ \Rashi{\textbf{וישמע פרעה.}הם הלשינו עליו (שמות רבה א'): ויבקש להרג את משה. מסרו לקוסטינר להרגו ולא שלטה בו החרב (שם), הוא שאמר משה "ויצלני מחרב פרעה" (שמות י״ח:ד׳): (וישב בארץ מדין. נתעכב שם; כמו: "וישב יעקב" (בראשית ל"ז):) וישב על הבאר. לשון ישיבה. למד מיעקב שנזדוג לו זווגו מן הבאר (שמות רבה א'):}%
% ---------- פרק 2 | פסוק 16 ----------
 \textbf{טז} וּלְכֹהֵ֥ן מִדְיָ֖ן שֶׁ֣בַע בָּנ֑וֹת וַתָּבֹ֣אנָה וַתִּדְלֶ֗נָה וַתְּמַלֶּ֙אנָה֙ אֶת־הָ֣רְהָטִ֔ים לְהַשְׁק֖וֹת צֹ֥אן אֲבִיהֶֽן׃ \Rashi{\textbf{ולכהן מדין.}רב שבהן; ופרש לו מע"ז ונדוהו מאצלם (שם): את הרהטים. את ברכות מרוצות המים העשויות בארץ:}%
% ---------- פרק 2 | פסוק 17 ----------
 \textbf{יז} וַיָּבֹ֥אוּ הָרֹעִ֖ים וַיְגָרְשׁ֑וּם וַיָּ֤קׇם מֹשֶׁה֙ וַיּ֣וֹשִׁעָ֔ן וַיַּ֖שְׁקְ אֶת־צֹאנָֽם׃ \Rashi{\textbf{ויגרשום.}מפני הנדוי (שם):}%
% ---------- פרק 2 | פסוק 18 ----------
 \textbf{יח} וַתָּבֹ֕אנָה אֶל־רְעוּאֵ֖ל אֲבִיהֶ֑ן וַיֹּ֕אמֶר מַדּ֛וּעַ מִהַרְתֶּ֥ן בֹּ֖א הַיּֽוֹם׃ %
% ---------- פרק 2 | פסוק 19 ----------
 \textbf{יט} וַתֹּאמַ֕רְןָ אִ֣ישׁ מִצְרִ֔י הִצִּילָ֖נוּ מִיַּ֣ד הָרֹעִ֑ים וְגַם־דָּלֹ֤ה דָלָה֙ לָ֔נוּ וַיַּ֖שְׁקְ אֶת־הַצֹּֽאן׃ %
% ---------- פרק 2 | פסוק 20 ----------
 \textbf{כ} וַיֹּ֥אמֶר אֶל־בְּנֹתָ֖יו וְאַיּ֑וֹ לָ֤מָּה זֶּה֙ עֲזַבְתֶּ֣ן אֶת־הָאִ֔ישׁ קִרְאֶ֥ן ל֖וֹ וְיֹ֥אכַל לָֽחֶם׃ \Rashi{\textbf{למה זה עזבתן.}הכיר בו שהוא מזרעו של יעקב, שהמים עולים לקראתו (שמות רבה א'): ויאכל לחם. שמא ישא אחת מכם; כמה דאת אמר: "כי אם הלחם אשר הוא אוכל" (בראשית ל"ט):}%
% ---------- פרק 2 | פסוק 21 ----------
 \textbf{כא} וַיּ֥וֹאֶל מֹשֶׁ֖ה לָשֶׁ֣בֶת אֶת־הָאִ֑ישׁ וַיִּתֵּ֛ן אֶת־צִפֹּרָ֥ה בִתּ֖וֹ לְמֹשֶֽׁה׃ \Rashi{\textbf{ויואל.}כתרגומו, ודומה לו "הואל נא ולין" (שופטים י"ט), "ולו הואלנו" (יהושע ז'), "הואלתי לדבר" (בראשית י"ח). ומדרשו: לשון אלה – נשבע לו שלא יזוז ממדין כי אם ברשותו (נדרים ס"ה):}%
% ---------- פרק 2 | פסוק 22 ----------
 \textbf{כב} וַתֵּ֣לֶד בֵּ֔ן וַיִּקְרָ֥א אֶת־שְׁמ֖וֹ גֵּרְשֹׁ֑ם כִּ֣י אָמַ֔ר גֵּ֣ר הָיִ֔יתִי בְּאֶ֖רֶץ נׇכְרִיָּֽה׃ \{פ\} %
% ---------- פרק 2 | פסוק 23 ----------
 \textbf{כג} וַיְהִי֩ בַיָּמִ֨ים הָֽרַבִּ֜ים הָהֵ֗ם וַיָּ֙מׇת֙ מֶ֣לֶךְ מִצְרַ֔יִם וַיֵּאָנְח֧וּ בְנֵֽי־יִשְׂרָאֵ֛ל מִן־הָעֲבֹדָ֖ה וַיִּזְעָ֑קוּ וַתַּ֧עַל שַׁוְעָתָ֛ם אֶל־הָאֱלֹהִ֖ים מִן־הָעֲבֹדָֽה׃ \Rashi{\textbf{ויהי בימים הרבים ההם.}שהיה משה גר במדין, וימת מלך מצרים, והצרכו ישראל לתשועה, ומשה היה רעה (שמות ג׳:א׳), ובאת תשועה על ידו, לכך נסמכו פרשיות הללו: וימת מלך מצרים. נצטרע והיה שוחט תינוקות ישראל ורוחץ בדמם (שמות רבה א'):}%
% ---------- פרק 2 | פסוק 24 ----------
 \textbf{כד} וַיִּשְׁמַ֥ע אֱלֹהִ֖ים אֶת־נַאֲקָתָ֑ם וַיִּזְכֹּ֤ר אֱלֹהִים֙ אֶת־בְּרִית֔וֹ אֶת־אַבְרָהָ֖ם אֶת־יִצְחָ֥ק וְאֶֽת־יַעֲקֹֽב׃ \Rashi{\textbf{נאקתם.}צעקתם, וכן "מעיר מתים ינאקו" (איוב כ"ד): את בריתו את אברהם. עם אברהם:}%
% ---------- פרק 2 | פסוק 25 ----------
 \textbf{כה} וַיַּ֥רְא אֱלֹהִ֖ים אֶת־בְּנֵ֣י יִשְׂרָאֵ֑ל וַיֵּ֖דַע אֱלֹהִֽים׃ \{ס\} \Rashi{\textbf{וידע אלהים.}נתן עליהם לב ולא העלים עיניו:}%
% ---------- פרק 3 | פסוק 1 ----------
 \textbf{א} וּמֹשֶׁ֗ה הָיָ֥ה רֹעֶ֛ה אֶת־צֹ֛אן יִתְר֥וֹ חֹתְנ֖וֹ כֹּהֵ֣ן מִדְיָ֑ן וַיִּנְהַ֤ג אֶת־הַצֹּאן֙ אַחַ֣ר הַמִּדְבָּ֔ר וַיָּבֹ֛א אֶל־הַ֥ר הָאֱלֹהִ֖ים חֹרֵֽבָה׃ \Rashi{\textbf{אחר המדבר.}להתרחק מן הגזל, שלא ירעו בשדות אחרים (שמות רבה ב'): אל הר האלהים. על שם העתיד:}%
% ---------- פרק 3 | פסוק 2 ----------
 \textbf{ב} וַ֠יֵּרָ֠א מַלְאַ֨ךְ יְהֹוָ֥ה אֵלָ֛יו בְּלַבַּת־אֵ֖שׁ מִתּ֣וֹךְ הַסְּנֶ֑ה וַיַּ֗רְא וְהִנֵּ֤ה הַסְּנֶה֙ בֹּעֵ֣ר בָּאֵ֔שׁ וְהַסְּנֶ֖ה אֵינֶ֥נּוּ אֻכָּֽל׃ \Rashi{\textbf{בלבת אש.}בשלהבת אש – לבו של אש, כמו "לב השמים" (דברים ד'), "בלב האלה" (שמואל ב י"ח); ואל תתמה על התי"ו, שיש לנו כיוצא בו, "מה אמלה לבתך" (יחזקאל ט"ז): מתוך הסנה. ולא אילן אחר, משום "עמו אנכי בצרה" (תהילים צא טו): אכל. נאכל, כמו "לא עבד בה" (דברים כ"א), "אשר לקח משם" (בראשית ג'):}%
% ---------- פרק 3 | פסוק 3 ----------
 \textbf{ג} וַיֹּ֣אמֶר מֹשֶׁ֔ה אָסֻֽרָה־נָּ֣א וְאֶרְאֶ֔ה אֶת־הַמַּרְאֶ֥ה הַגָּדֹ֖ל הַזֶּ֑ה מַדּ֖וּעַ לֹא־יִבְעַ֥ר הַסְּנֶֽה׃ \Rashi{\textbf{אסרה נא.}אסורה מכאן להתקרב שם:}%
% ---------- פרק 3 | פסוק 4 ----------
 \textbf{ד} וַיַּ֥רְא יְהֹוָ֖ה כִּ֣י סָ֣ר לִרְא֑וֹת וַיִּקְרָא֩ אֵלָ֨יו אֱלֹהִ֜ים מִתּ֣וֹךְ הַסְּנֶ֗ה וַיֹּ֛אמֶר מֹשֶׁ֥ה מֹשֶׁ֖ה וַיֹּ֥אמֶר הִנֵּֽנִי׃ %
% ---------- פרק 3 | פסוק 5 ----------
 \textbf{ה} וַיֹּ֖אמֶר אַל־תִּקְרַ֣ב הֲלֹ֑ם שַׁל־נְעָלֶ֙יךָ֙ מֵעַ֣ל רַגְלֶ֔יךָ כִּ֣י הַמָּק֗וֹם אֲשֶׁ֤ר אַתָּה֙ עוֹמֵ֣ד עָלָ֔יו אַדְמַת־קֹ֖דֶשׁ הֽוּא׃ \Rashi{\textbf{של.}שלף והוצא, כמו "ונשל הברזל" (דברים י"ט), "כי ישל זיתך" (שם כ"ח): אדמת קדש הוא. המקום:}%
% ---------- פרק 3 | פסוק 6 ----------
 \textbf{ו} וַיֹּ֗אמֶר אָנֹכִי֙ אֱלֹהֵ֣י אָבִ֔יךָ אֱלֹהֵ֧י אַבְרָהָ֛ם אֱלֹהֵ֥י יִצְחָ֖ק וֵאלֹהֵ֣י יַעֲקֹ֑ב וַיַּסְתֵּ֤ר מֹשֶׁה֙ פָּנָ֔יו כִּ֣י יָרֵ֔א מֵהַבִּ֖יט אֶל־הָאֱלֹהִֽים׃ %
% ---------- פרק 3 | פסוק 7 ----------
 \textbf{ז} וַיֹּ֣אמֶר יְהֹוָ֔ה רָאֹ֥ה רָאִ֛יתִי אֶת־עֳנִ֥י עַמִּ֖י אֲשֶׁ֣ר בְּמִצְרָ֑יִם וְאֶת־צַעֲקָתָ֤ם שָׁמַ֙עְתִּי֙ מִפְּנֵ֣י נֹֽגְשָׂ֔יו כִּ֥י יָדַ֖עְתִּי אֶת־מַכְאֹבָֽיו׃ \Rashi{\textbf{כי ידעתי את מכאביו.}כמו "וידע אלהים" (שמות ב'), כלומר כי שמתי לב להתבונן ולדעת את מכאוביו ולא העלמתי עיני לאטם אזני מצעקתם:}%
% ---------- פרק 3 | פסוק 8 ----------
 \textbf{ח} וָאֵרֵ֞ד לְהַצִּיל֣וֹ ׀ מִיַּ֣ד מִצְרַ֗יִם וּֽלְהַעֲלֹתוֹ֮ מִן־הָאָ֣רֶץ הַהִוא֒ אֶל־אֶ֤רֶץ טוֹבָה֙ וּרְחָבָ֔ה אֶל־אֶ֛רֶץ זָבַ֥ת חָלָ֖ב וּדְבָ֑שׁ אֶל־מְק֤וֹם הַֽכְּנַעֲנִי֙ וְהַ֣חִתִּ֔י וְהָֽאֱמֹרִי֙ וְהַפְּרִזִּ֔י וְהַחִוִּ֖י וְהַיְבוּסִֽי׃ %
% ---------- פרק 3 | פסוק 9 ----------
 \textbf{ט} וְעַתָּ֕ה הִנֵּ֛ה צַעֲקַ֥ת בְּנֵי־יִשְׂרָאֵ֖ל בָּ֣אָה אֵלָ֑י וְגַם־רָאִ֙יתִי֙ אֶת־הַלַּ֔חַץ אֲשֶׁ֥ר מִצְרַ֖יִם לֹחֲצִ֥ים אֹתָֽם׃ %
% ---------- פרק 3 | פסוק 10 ----------
 \textbf{י} וְעַתָּ֣ה לְכָ֔ה וְאֶֽשְׁלָחֲךָ֖ אֶל־פַּרְעֹ֑ה וְהוֹצֵ֛א אֶת־עַמִּ֥י בְנֵֽי־יִשְׂרָאֵ֖ל מִמִּצְרָֽיִם׃ \Rashi{\textbf{ועתה לכה ואשלחך אל פרעה.}ואם תאמר מה תועיל? והוצא את עמי. יועילו דבריך ותוציאם משם:}%
% ---------- פרק 3 | פסוק 11 ----------
 \textbf{יא} וַיֹּ֤אמֶר מֹשֶׁה֙ אֶל־הָ֣אֱלֹהִ֔ים מִ֣י אָנֹ֔כִי כִּ֥י אֵלֵ֖ךְ אֶל־פַּרְעֹ֑ה וְכִ֥י אוֹצִ֛יא אֶת־בְּנֵ֥י יִשְׂרָאֵ֖ל מִמִּצְרָֽיִם׃ \Rashi{\textbf{מי אנכי.}מה אני חשוב לדבר עם המלכים? וכי אוציא את בני ישראל. ואף אם חשוב אני, מה זכו ישראל שיעשה להם נס ואוציאם ממצרים?:}%
% ---------- פרק 3 | פסוק 12 ----------
 \textbf{יב} וַיֹּ֙אמֶר֙ כִּֽי־אֶֽהְיֶ֣ה עִמָּ֔ךְ וְזֶה־לְּךָ֣ הָא֔וֹת כִּ֥י אָנֹכִ֖י שְׁלַחְתִּ֑יךָ בְּהוֹצִֽיאֲךָ֤ אֶת־הָעָם֙ מִמִּצְרַ֔יִם תַּֽעַבְדוּן֙ אֶת־הָ֣אֱלֹהִ֔ים עַ֖ל הָהָ֥ר הַזֶּֽה׃ \Rashi{\textbf{ויאמר כי אהיה עמך.}השיבו על ראשון ראשון ועל אחרון אחרון; שאמרת מי אנכי כי אלך אל פרעה – לא שלך היא כי אם משלי, כי אהיה עמך – וזה המראה אשר ראית בסנה לך האות כי אנכי שלחתיך, וכדאי אני להציל, כאשר ראית הסנה עושה שליחותי ואיננו אכל, כך תלך בשליחותי ואינך נזוק; וששאלת מה זכות יש לישראל שיצאו ממצרים? דבר גדול יש לי על הוצאה זו, שהרי עתידים לקבל התורה על ההר הזה לסוף ג' חדשים שיצאו ממצרים. ד"א: כי אהיה עמך, וזה שתצליח בשליחותך לך האות על הבטחה אחרת, שאני מבטיחך שכשתוציאם ממצרים תעבדון אותי על ההר הזה, שתקבלו התורה עליו, והיא הזכות העומדת לישראל. ודגמת לשון זה מצינו: "וזה לך האות אכול השנה ספיח וגו'" (ישעיהו ל"ז), מפלת סנחריב תהיה לך אות על הבטחה אחרת, שארצכם חרבה מפרות ואני אברך הספיחים:}%
% ---------- פרק 3 | פסוק 13 ----------
 \textbf{יג} וַיֹּ֨אמֶר מֹשֶׁ֜ה אֶל־הָֽאֱלֹהִ֗ים הִנֵּ֨ה אָנֹכִ֣י בָא֮ אֶל־בְּנֵ֣י יִשְׂרָאֵל֒ וְאָמַרְתִּ֣י לָהֶ֔ם אֱלֹהֵ֥י אֲבוֹתֵיכֶ֖ם שְׁלָחַ֣נִי אֲלֵיכֶ֑ם וְאָֽמְרוּ־לִ֣י מַה־שְּׁמ֔וֹ מָ֥ה אֹמַ֖ר אֲלֵהֶֽם׃ %
% ---------- פרק 3 | פסוק 14 ----------
 \textbf{יד} וַיֹּ֤אמֶר אֱלֹהִים֙ אֶל־מֹשֶׁ֔ה אֶֽהְיֶ֖ה אֲשֶׁ֣ר אֶֽהְיֶ֑ה וַיֹּ֗אמֶר כֹּ֤ה תֹאמַר֙ לִבְנֵ֣י יִשְׂרָאֵ֔ל אֶֽהְיֶ֖ה שְׁלָחַ֥נִי אֲלֵיכֶֽם׃ \Rashi{\textbf{אהיה אשר אהיה.}אהיה עמם בצרה זו אשר אהיה עמם בשעבוד שאר מלכיות. אמר לפניו, רבונו של עולם, מה אני מזכיר להם צרה אחרת? דים בצרה זו, אמר לו יפה אמרת, כה תאמר וגו' (ברכות ט'):}%
% ---------- פרק 3 | פסוק 15 ----------
 \textbf{טו} וַיֹּ֩אמֶר֩ ע֨וֹד אֱלֹהִ֜ים אֶל־מֹשֶׁ֗ה כֹּֽה־תֹאמַר֮ אֶל־בְּנֵ֣י יִשְׂרָאֵל֒ יְהֹוָ֞ה אֱלֹהֵ֣י אֲבֹתֵיכֶ֗ם אֱלֹהֵ֨י אַבְרָהָ֜ם אֱלֹהֵ֥י יִצְחָ֛ק וֵאלֹהֵ֥י יַעֲקֹ֖ב שְׁלָחַ֣נִי אֲלֵיכֶ֑ם זֶה־שְּׁמִ֣י לְעֹלָ֔ם וְזֶ֥ה זִכְרִ֖י לְדֹ֥ר דֹּֽר׃ \Rashi{\textbf{זה שמי לעלם.}חסר וי"ו, לומר, העלימהו – שלא יקרא ככתבו (פסחים נ'): וזה זכרי. למדו היאך נקרא. וכן דוד הוא אומר, "ה' שמך לעולם ה' זכרך לדר ודר" (תהלים קל"ה):}%
% ---------- פרק 3 | פסוק 16 ----------
 \textbf{טז} לֵ֣ךְ וְאָֽסַפְתָּ֞ אֶת־זִקְנֵ֣י יִשְׂרָאֵ֗ל וְאָמַרְתָּ֤ אֲלֵהֶם֙ יְהֹוָ֞ה אֱלֹהֵ֤י אֲבֹֽתֵיכֶם֙ נִרְאָ֣ה אֵלַ֔י אֱלֹהֵ֧י אַבְרָהָ֛ם יִצְחָ֥ק וְיַעֲקֹ֖ב לֵאמֹ֑ר פָּקֹ֤ד פָּקַ֙דְתִּי֙ אֶתְכֶ֔ם וְאֶת־הֶעָשׂ֥וּי לָכֶ֖ם בְּמִצְרָֽיִם׃ \Rashi{\textbf{את זקני ישראל.}מיחדים לישיבה. ואם תאמר זקנים סתם, היאך אפשר לו לאסף זקנים של ששים רבוא? (יומא כ"ח)}%
% ---------- פרק 3 | פסוק 17 ----------
 \textbf{יז} וָאֹמַ֗ר אַעֲלֶ֣ה אֶתְכֶם֮ מֵעֳנִ֣י מִצְרַ֒יִם֒ אֶל־אֶ֤רֶץ הַֽכְּנַעֲנִי֙ וְהַ֣חִתִּ֔י וְהָֽאֱמֹרִי֙ וְהַפְּרִזִּ֔י וְהַחִוִּ֖י וְהַיְבוּסִ֑י אֶל־אֶ֛רֶץ זָבַ֥ת חָלָ֖ב וּדְבָֽשׁ׃ %
% ---------- פרק 3 | פסוק 18 ----------
 \textbf{יח} וְשָׁמְע֖וּ לְקֹלֶ֑ךָ וּבָאתָ֡ אַתָּה֩ וְזִקְנֵ֨י יִשְׂרָאֵ֜ל אֶל־מֶ֣לֶךְ מִצְרַ֗יִם וַאֲמַרְתֶּ֤ם אֵלָיו֙ יְהֹוָ֞ה אֱלֹהֵ֤י הָֽעִבְרִיִּים֙ נִקְרָ֣ה עָלֵ֔ינוּ וְעַתָּ֗ה נֵֽלְכָה־נָּ֞א דֶּ֣רֶךְ שְׁלֹ֤שֶׁת יָמִים֙ בַּמִּדְבָּ֔ר וְנִזְבְּחָ֖ה לַֽיהֹוָ֥ה אֱלֹהֵֽינוּ׃ \Rashi{\textbf{ושמעו לקלך.}מכיון שתאמר להם לשון זה, ישמעו לקולך, שכבר סימן זה מסור בידם מיעקב ומיוסף, שבלשון זה הם נגאלים, יעקב אמר להם "ואלהים פקד יפקד אתכם" (בראשית נ'), יוסף אמר להם "פקד יפקד אלהים אתכם" (שם): נקרה עלינו. לשון מקרה; וכן "ויקר אלהים", "ואנכי אקרה כה" (במדבר כ"ג) – אהא נקרה מאתו הלם: (אלהי העבריים. יו"ד יתירה, רמז לי' מכות. ברש"י ישן):}%
% ---------- פרק 3 | פסוק 19 ----------
 \textbf{יט} וַאֲנִ֣י יָדַ֔עְתִּי כִּ֠י לֹֽא־יִתֵּ֥ן אֶתְכֶ֛ם מֶ֥לֶךְ מִצְרַ֖יִם לַהֲלֹ֑ךְ וְלֹ֖א בְּיָ֥ד חֲזָקָֽה׃ \Rashi{\textbf{לא יתן אתכם מלך מצרים להלך.}אם אין אני מראה לו ידי החזקה; כלומר, כל עוד שאין אני מודיעו ידי החזקה, לא יתן אתכם להלך: לא יתן. לא ישבק כמו: "על כן לא נתתיך" (בראשית כ'), "ולא נתנו אלהים להרע עמדי" (שם ל"א), וכלן לשון נתינה הם. וי"מ ולא ביד חזקה — ולא בשביל שידו חזקה, "כי מאז אשלח את ידי והכיתי את מצרים וגו'", ומתרגמין אותו "ולא מן קדם דחילה תקיף". משמו של רבי יעקב ברבי מנחם נאמר לי:}%
% ---------- פרק 3 | פסוק 20 ----------
 \textbf{כ} וְשָׁלַחְתִּ֤י אֶת־יָדִי֙ וְהִכֵּיתִ֣י אֶת־מִצְרַ֔יִם בְּכֹל֙ נִפְלְאֹתַ֔י אֲשֶׁ֥ר אֶֽעֱשֶׂ֖ה בְּקִרְבּ֑וֹ וְאַחֲרֵי־כֵ֖ן יְשַׁלַּ֥ח אֶתְכֶֽם׃ %
% ---------- פרק 3 | פסוק 21 ----------
 \textbf{כא} וְנָתַתִּ֛י אֶת־חֵ֥ן הָֽעָם־הַזֶּ֖ה בְּעֵינֵ֣י מִצְרָ֑יִם וְהָיָה֙ כִּ֣י תֵֽלֵכ֔וּן לֹ֥א תֵלְכ֖וּ רֵיקָֽם׃ %
% ---------- פרק 3 | פסוק 22 ----------
 \textbf{כב} וְשָׁאֲלָ֨ה אִשָּׁ֤ה מִשְּׁכֶנְתָּהּ֙ וּמִגָּרַ֣ת בֵּיתָ֔הּ כְּלֵי־כֶ֛סֶף וּכְלֵ֥י זָהָ֖ב וּשְׂמָלֹ֑ת וְשַׂמְתֶּ֗ם עַל־בְּנֵיכֶם֙ וְעַל־בְּנֹ֣תֵיכֶ֔ם וְנִצַּלְתֶּ֖ם אֶת־מִצְרָֽיִם׃ \Rashi{\textbf{ומגרת ביתה.}מאותה שהיא גרה אתה בבית: ונצלתם. כתרגומו "ותרוקנון", וכן "וינצלו את מצרים" (שמות י"ב), "ויתנצלו בני ישראל את עדים" (שם ל"ג), והנו"ן בו יסוד; ומנחם חברו במחברת צד"י עם "ויצל אלהים את מקנה אביכם", "אשר הציל אלהים מאבינו" (בראשית ל"א), ולא יאמנו דבריו, כי אם לא היתה הנו"ן יסוד והיא נקודה בחירק, לא תהא משמשת בלשון ופעלתם אלא בלשון ונפעלתם, כמו "ונסחתם מעל האדמה" (דברים כ"ח), "ונתתם ביד אויב", ונגפתם לפני אויביכם (ויקרא כ"ו), "ונתכתם בתוכה" (יחזקאל כ"ב), "ואמרתם נצלנו" (ירמיהו ז'), לשון נפעלנו, וכל נו"ן שהיא באה בתבה לפרקים ונופלת ממנה – כנו"ן של נוגף, נושא, נותן, נושך – כשהיא מדברת לשון ופעלתם תנקד בחטף, כגון "ונשאתם את אביכם" (בראשית מ"ה), "ונתתם להם את ארץ הגלעד" (במדבר ל"ב), "ונמלתם את בשר ערלתכם" (בראשית י"ז), לכך אני אומר שזאת הנקודה בחירק מן היסוד היא, ויסוד שם דבר נצול, והוא מן הלשונות הכבדים, כמו דבור, כפור, למוד, כשידבר בלשון ופעלתם, ינקד בחירק כמו "ודברתם אל הסלע" (במדבר כ'), "וכפרתם את הבית" (יחזקאל מ"ה), "ולמדתם אותם את בניכם" (דברים י"א):}%
% ---------- פרק 4 | פסוק 1 ----------
 \textbf{א} וַיַּ֤עַן מֹשֶׁה֙ וַיֹּ֔אמֶר וְהֵן֙ לֹֽא־יַאֲמִ֣ינוּ לִ֔י וְלֹ֥א יִשְׁמְע֖וּ בְּקֹלִ֑י כִּ֣י יֹֽאמְר֔וּ לֹֽא־נִרְאָ֥ה אֵלֶ֖יךָ יְהֹוָֽה׃ %
% ---------- פרק 4 | פסוק 2 ----------
 \textbf{ב} וַיֹּ֧אמֶר אֵלָ֛יו יְהֹוָ֖ה (מזה) [מַה־זֶּ֣ה] בְיָדֶ֑ךָ וַיֹּ֖אמֶר מַטֶּֽה׃ \Rashi{\textbf{מזה בידך.}לכך נכתב תבה אחת, לדרש: מזה שבידך אתה חיב ללקות, שחשדת בכשרים; ופשוטו: כאדם שאומר לחברו מודה אתה שזו שלפניך אבן היא? אומר לו הן; אמר לו הריני עושה אותה עץ:}%
% ---------- פרק 4 | פסוק 3 ----------
 \textbf{ג} וַיֹּ֙אמֶר֙ הַשְׁלִיכֵ֣הוּ אַ֔רְצָה וַיַּשְׁלִכֵ֥הוּ אַ֖רְצָה וַיְהִ֣י לְנָחָ֑שׁ וַיָּ֥נׇס מֹשֶׁ֖ה מִפָּנָֽיו׃ \Rashi{\textbf{ויהי לנחש.}רמז לו שספר לשון הרע על ישראל ותפש אמנותו של נחש (שמות רבה):}%
% ---------- פרק 4 | פסוק 4 ----------
 \textbf{ד} וַיֹּ֤אמֶר יְהֹוָה֙ אֶל־מֹשֶׁ֔ה שְׁלַח֙ יָֽדְךָ֔ וֶאֱחֹ֖ז בִּזְנָב֑וֹ וַיִּשְׁלַ֤ח יָדוֹ֙ וַיַּ֣חֲזֶק בּ֔וֹ וַיְהִ֥י לְמַטֶּ֖ה בְּכַפּֽוֹ׃ \Rashi{\textbf{ויחזק בו.}לשון אחיזה הוא; והרבה יש במקרא, "ויחזיקו האנשים בידו" (בראשית י"ט), "והחזיקה במבשיו" (דברים כ"ה), "והחזקתי בזקנו" (שמואל א י"ז), (כל לשון חזוק הדבוק לבי"ת לשון אחיזה הוא):}%
% ---------- פרק 4 | פסוק 5 ----------
 \textbf{ה} לְמַ֣עַן יַאֲמִ֔ינוּ כִּֽי־נִרְאָ֥ה אֵלֶ֛יךָ יְהֹוָ֖ה אֱלֹהֵ֣י אֲבֹתָ֑ם אֱלֹהֵ֧י אַבְרָהָ֛ם אֱלֹהֵ֥י יִצְחָ֖ק וֵאלֹהֵ֥י יַעֲקֹֽב׃ %
% ---------- פרק 4 | פסוק 6 ----------
 \textbf{ו} וַיֹּ֩אמֶר֩ יְהֹוָ֨ה ל֜וֹ ע֗וֹד הָֽבֵא־נָ֤א יָֽדְךָ֙ בְּחֵיקֶ֔ךָ וַיָּבֵ֥א יָד֖וֹ בְּחֵיק֑וֹ וַיּ֣וֹצִאָ֔הּ וְהִנֵּ֥ה יָד֖וֹ מְצֹרַ֥עַת כַּשָּֽׁלֶג׃ \Rashi{\textbf{מצרעת כשלג.}דרך צרעת להיות לבנה – "ואם בהרת לבנה היא" (ויקרא י"ג); אף באות זה רמז לו שלשון הרע ספר באמרו לא יאמינו לי, לפיכך הלקהו בצרעת, כמו שלקתה מרים על לשון הרע (שבת צ"ז):}%
% ---------- פרק 4 | פסוק 7 ----------
 \textbf{ז} וַיֹּ֗אמֶר הָשֵׁ֤ב יָֽדְךָ֙ אֶל־חֵיקֶ֔ךָ וַיָּ֥שֶׁב יָד֖וֹ אֶל־חֵיק֑וֹ וַיּֽוֹצִאָהּ֙ מֵֽחֵיק֔וֹ וְהִנֵּה־שָׁ֖בָה כִּבְשָׂרֽוֹ׃ \Rashi{\textbf{ויוצאה מחיקו והנה שבה כבשרו.}מכאן שמדה טובה ממהרת לבא ממדת פרענות, שהרי בראשונה לא נאמר "מחיקו" (שם):}%
% ---------- פרק 4 | פסוק 8 ----------
 \textbf{ח} וְהָיָה֙ אִם־לֹ֣א יַאֲמִ֣ינוּ לָ֔ךְ וְלֹ֣א יִשְׁמְע֔וּ לְקֹ֖ל הָאֹ֣ת הָרִאשׁ֑וֹן וְהֶֽאֱמִ֔ינוּ לְקֹ֖ל הָאֹ֥ת הָאַחֲרֽוֹן׃ \Rashi{\textbf{והאמינו לקל האת האחרון.}משתאמר להם בשבילכם לקיתי, על שספרתי עליכם לשון הרע, יאמינו לך, שכבר למדו בכך שהמזדוגין להרע להם לוקים בנגעים, כגון פרעה ואבימלך בשביל שרה:}%
% ---------- פרק 4 | פסוק 9 ----------
 \textbf{ט} וְהָיָ֡ה אִם־לֹ֣א יַאֲמִ֡ינוּ גַּם֩ לִשְׁנֵ֨י הָאֹת֜וֹת הָאֵ֗לֶּה וְלֹ֤א יִשְׁמְעוּן֙ לְקֹלֶ֔ךָ וְלָקַחְתָּ֙ מִמֵּימֵ֣י הַיְאֹ֔ר וְשָׁפַכְתָּ֖ הַיַּבָּשָׁ֑ה וְהָי֤וּ הַמַּ֙יִם֙ אֲשֶׁ֣ר תִּקַּ֣ח מִן־הַיְאֹ֔ר וְהָי֥וּ לְדָ֖ם בַּיַּבָּֽשֶׁת׃ \Rashi{\textbf{ולקחת ממימי היאור.}רמז להם שבמכה ראשונה נפרע מאלהותם, (פירוש, כשהקב"ה נפרע מן האומות, נפרע מאלהותם תחלה, שהיו עובדים לנילוס המחיה אותם, והפכם לדם. ברש"י ישן): והיו המים וגו'. "והיו" "והיו", ב' פעמים נראה בעיני, אלו נאמר והיו המים אשר תקח מן היאור לדם ביבשת, שומע אני, שבידו הם נהפכים לדם – ואף כשירדו לארץ יהיו בהויתן, אבל עכשו מלמדנו שלא יהיו דם עד שיהיו ביבשת:}%
% ---------- פרק 4 | פסוק 10 ----------
 \textbf{י} וַיֹּ֨אמֶר מֹשֶׁ֣ה אֶל־יְהֹוָה֮ בִּ֣י אֲדֹנָי֒ לֹא֩ אִ֨ישׁ דְּבָרִ֜ים אָנֹ֗כִי גַּ֤ם מִתְּמוֹל֙ גַּ֣ם מִשִּׁלְשֹׁ֔ם גַּ֛ם מֵאָ֥ז דַּבֶּרְךָ֖ אֶל־עַבְדֶּ֑ךָ כִּ֧י כְבַד־פֶּ֛ה וּכְבַ֥ד לָשׁ֖וֹן אָנֹֽכִי׃ \Rashi{\textbf{גם מתמול וגו'.}למדנו שכל ז' ימים היה הקב"ה מפתה את משה בסנה לילך בשליחותו, "מתמול", "שלשם", "מאז דברך" הרי ג' ושלושה גמין רבויין הם, הרי ששה, והוא היה עומד ביום השביעי כשאמר לו זאת עוד שלח נא ביד תשלח, עד שחרה בו וקבל עליו (שמות רבה); וכל זה שלא היה רוצה לטל גדלה על אהרן אחיו שהיה גדול הימנו ונביא היה, שנאמר "הנגלה נגליתי אל בית אביך בהיותם במצרים" (שמואל א ב') – הוא אהרן, וכן "ואודע להם בארץ מצרים … ואמר אליהם איש שקוצי עיניו השליכו" (יחזקאל כ'), ואותה נבואה לאהרן נאמרה (תנחומא): כבד פה. בכבדות אני מדבר. ובלשון לעז בלב"ו:}%
% ---------- פרק 4 | פסוק 11 ----------
 \textbf{יא} וַיֹּ֨אמֶר יְהֹוָ֜ה אֵלָ֗יו מִ֣י שָׂ֣ם פֶּה֮ לָֽאָדָם֒ א֚וֹ מִֽי־יָשׂ֣וּם אִלֵּ֔ם א֣וֹ חֵרֵ֔שׁ א֥וֹ פִקֵּ֖חַ א֣וֹ עִוֵּ֑ר הֲלֹ֥א אָנֹכִ֖י יְהֹוָֽה׃ \Rashi{\textbf{מי שם פה וגו'.}מי למדך לדבר, כשהיית נדון לפני פרעה על המצרי? או מי ישום אלם. מי עשה פרעה אלם, שלא נתאמץ במצות הריגתך, ואת משרתיו חרשים שלא שמעו בצוותו עליך, ואת אספקלטורין ההורגים מי עשאם עורים שלא ראו כשברחת מן הבימה ונמלטת? הלא אנכי. ששמי ה' עשיתי כל זאת (שם):}%
% ---------- פרק 4 | פסוק 12 ----------
 \textbf{יב} וְעַתָּ֖ה לֵ֑ךְ וְאָנֹכִי֙ אֶֽהְיֶ֣ה עִם־פִּ֔יךָ וְהוֹרֵיתִ֖יךָ אֲשֶׁ֥ר תְּדַבֵּֽר׃ %
% ---------- פרק 4 | פסוק 13 ----------
 \textbf{יג} וַיֹּ֖אמֶר בִּ֣י אֲדֹנָ֑י שְֽׁלַֽח־נָ֖א בְּיַד־תִּשְׁלָֽח׃ \Rashi{\textbf{ביד תשלח.}ביד מי שאתה רגיל לשלח, והוא אהרן. ד"א: ביד אחר שתרצה לשלח, אין סופי להכניסם לארץ ולהיות גואלם לעתיד, יש לך שלוחים הרבה:}%
% ---------- פרק 4 | פסוק 14 ----------
 \textbf{יד} וַיִּֽחַר־אַ֨ף יְהֹוָ֜ה בְּמֹשֶׁ֗ה וַיֹּ֙אמֶר֙ הֲלֹ֨א אַהֲרֹ֤ן אָחִ֙יךָ֙ הַלֵּוִ֔י יָדַ֕עְתִּי כִּֽי־דַבֵּ֥ר יְדַבֵּ֖ר ה֑וּא וְגַ֤ם הִנֵּה־הוּא֙ יֹצֵ֣א לִקְרָאתֶ֔ךָ וְרָאֲךָ֖ וְשָׂמַ֥ח בְּלִבּֽוֹ׃ \Rashi{\textbf{ויחר אף.}רבי יהושע בן קרחה אומר כל חרון אף שבתורה עושה רשם וזה לא נאמר בו רשם ולא מצינו שבא ענש ע"י אותו חרון; אמר לו רבי יוסי אף בזו נאמר בו רשם – הלא אהרן אחיך הלוי – שהיה עתיד להיות לוי ולא כהן, והכהנה הייתי אומר לצאת ממך, מעתה לא יהיה כן אלא הוא יהיה כהן ואתה הלוי, שנאמר "ומשה איש האלהים בניו יקראו על שבט הלוי" (דהי"א כ"ג) (זבחים ק"ב): הנה הוא יצא לקראתך. כשתלך למצרים: וראך ושמח בלבו. לא כשאתה סבור, שיהא מקפיד עליך שאתה עולה לגדלה; ומשם זכה אהרן לעדי החשן הנתון על הלב (שבת קל"ט):}%
% ---------- פרק 4 | פסוק 15 ----------
 \textbf{טו} וְדִבַּרְתָּ֣ אֵלָ֔יו וְשַׂמְתָּ֥ אֶת־הַדְּבָרִ֖ים בְּפִ֑יו וְאָנֹכִ֗י אֶֽהְיֶ֤ה עִם־פִּ֙יךָ֙ וְעִם־פִּ֔יהוּ וְהוֹרֵיתִ֣י אֶתְכֶ֔ם אֵ֖ת אֲשֶׁ֥ר תַּעֲשֽׂוּן׃ %
% ---------- פרק 4 | פסוק 16 ----------
 \textbf{טז} וְדִבֶּר־ה֥וּא לְךָ֖ אֶל־הָעָ֑ם וְהָ֤יָה הוּא֙ יִֽהְיֶה־לְּךָ֣ לְפֶ֔ה וְאַתָּ֖ה תִּֽהְיֶה־לּ֥וֹ לֵֽאלֹהִֽים׃ \Rashi{\textbf{ודבר הוא לך.}בשבילך ידבר אל העם. וזה יוכיח על כל לך ולי ולו ולכם ולהם הסמוכים לדבור שכלם לשון על הם: יהיה לך לפה. למליץ, לפי שאתה כבד פה: לאלהים. לרב ולשר:}%
% ---------- פרק 4 | פסוק 17 ----------
 \textbf{יז} וְאֶת־הַמַּטֶּ֥ה הַזֶּ֖ה תִּקַּ֣ח בְּיָדֶ֑ךָ אֲשֶׁ֥ר תַּעֲשֶׂה־בּ֖וֹ אֶת־הָאֹתֹֽת׃ \{פ\} %
% ---------- פרק 4 | פסוק 18 ----------
 \textbf{יח} וַיֵּ֨לֶךְ מֹשֶׁ֜ה וַיָּ֣שׇׁב ׀ אֶל־יֶ֣תֶר חֹֽתְנ֗וֹ וַיֹּ֤אמֶר לוֹ֙ אֵ֣לְכָה נָּ֗א וְאָשׁ֙וּבָה֙ אֶל־אַחַ֣י אֲשֶׁר־בְּמִצְרַ֔יִם וְאֶרְאֶ֖ה הַעוֹדָ֣ם חַיִּ֑ים וַיֹּ֧אמֶר יִתְר֛וֹ לְמֹשֶׁ֖ה לֵ֥ךְ לְשָׁלֽוֹם׃ \Rashi{\textbf{וישב אל יתר חתנו.}לטל רשות, שהרי נשבע לו (נדרים ס"ה). וז' שמות היו לו, רעואל, יתר, יתרו, קיני, חובב, חבר, פוטיאל:}%
% ---------- פרק 4 | פסוק 19 ----------
 \textbf{יט} וַיֹּ֨אמֶר יְהֹוָ֤ה אֶל־מֹשֶׁה֙ בְּמִדְיָ֔ן לֵ֖ךְ שֻׁ֣ב מִצְרָ֑יִם כִּי־מֵ֙תוּ֙ כׇּל־הָ֣אֲנָשִׁ֔ים הַֽמְבַקְשִׁ֖ים אֶת־נַפְשֶֽׁךָ׃ \Rashi{\textbf{כי מתו כל האנשים.}מי הם? דתן ואבירם; חיים היו, אלא שירדו מנכסיהם והעני חשוב כמת:}%
% ---------- פרק 4 | פסוק 20 ----------
 \textbf{כ} וַיִּקַּ֨ח מֹשֶׁ֜ה אֶת־אִשְׁתּ֣וֹ וְאֶת־בָּנָ֗יו וַיַּרְכִּבֵם֙ עַֽל־הַחֲמֹ֔ר וַיָּ֖שׇׁב אַ֣רְצָה מִצְרָ֑יִם וַיִּקַּ֥ח מֹשֶׁ֛ה אֶת־מַטֵּ֥ה הָאֱלֹהִ֖ים בְּיָדֽוֹ׃ \Rashi{\textbf{על החמר.}חמור המיחד, הוא החמור שחבש אברהם לעקדת יצחק, והוא שעתיד מלך המשיח להגלות עליו, שנאמר "עני ורכב על חמור" (זכריה ט') (פרקי דרבי אליעזר ל"א): וישב ארצה מצרים, ויקח משה את מטה. אין מקדם ומאחר מדקדקים במקרא:}%
% ---------- פרק 4 | פסוק 21 ----------
 \textbf{כא} וַיֹּ֣אמֶר יְהֹוָה֮ אֶל־מֹשֶׁה֒ בְּלֶכְתְּךָ֙ לָשׁ֣וּב מִצְרַ֔יְמָה רְאֵ֗ה כׇּל־הַמֹּֽפְתִים֙ אֲשֶׁר־שַׂ֣מְתִּי בְיָדֶ֔ךָ וַעֲשִׂיתָ֖ם לִפְנֵ֣י פַרְעֹ֑ה וַאֲנִי֙ אֲחַזֵּ֣ק אֶת־לִבּ֔וֹ וְלֹ֥א יְשַׁלַּ֖ח אֶת־הָעָֽם׃ \Rashi{\textbf{בלכתך לשוב מצרימה וגו'.}דע שעל מנת כן תלך שתהא גבור בשליחותי לעשות כל מופתי לפני פרעה ולא תירא ממנו: אשר שמתי בידך. לא על שלוש אותות האמורות למעלה, שהרי לא לפני פרעה צוה לעשותם אלא לפני ישראל שיאמינו לו, ולא מצינו שעשאם לפניו, אלא מופתים שאני עתיד לשום בידך במצרים, כמו "כי ידבר אליכם פרעה וגו'" (שמות ז'); ואל תתמה על אשר כתיב "אשר שמתי", שכן משמעו, כשתדבר עמו כבר שמתים בידך:}%
% ---------- פרק 4 | פסוק 22 ----------
 \textbf{כב} וְאָמַרְתָּ֖ אֶל־פַּרְעֹ֑ה כֹּ֚ה אָמַ֣ר יְהֹוָ֔ה בְּנִ֥י בְכֹרִ֖י יִשְׂרָאֵֽל׃ \Rashi{\textbf{ואמרת אל פרעה.}כשתשמע שלבו חזק וימאן לשלח אמר לו כן: בני בכרי. לשון גדלה; כמו: "אף אני בכור אתנהו" (תהלים פ"ט), זהו פשוטו. ומדרשו: כאן חתם הקב"ה על מכירת הבכורה שלקח יעקב מעשו (בראשית רבה ס"ג):}%
% ---------- פרק 4 | פסוק 23 ----------
 \textbf{כג} וָאֹמַ֣ר אֵלֶ֗יךָ שַׁלַּ֤ח אֶת־בְּנִי֙ וְיַֽעַבְדֵ֔נִי וַתְּמָאֵ֖ן לְשַׁלְּח֑וֹ הִנֵּה֙ אָנֹכִ֣י הֹרֵ֔ג אֶת־בִּנְךָ֖ בְּכֹרֶֽךָ׃ \Rashi{\textbf{ואמר אליך.}בשליחותו של מקום שלח את בני: הנה אנכי הרג וגו'. היא מכה אחרונה, ובה התרהו תחלה מפני שהיא קשה; וזה הוא שאמר "הן אל ישגיב בכחו" – לפיכך "מי כמהו מורה" (איוב ל"ו); בשר ודם המבקש להנקם מחברו, מעלים את דבריו, שלא יבקש הצלה, אבל הקב"ה ישגיב בכחו ואין יכלת להמלט מידו כי אם בשובו אליו, לפיכך הוא מורהו ומתרה בו לשוב:}%
% ---------- פרק 4 | פסוק 24 ----------
 \textbf{כד} וַיְהִ֥י בַדֶּ֖רֶךְ בַּמָּל֑וֹן וַיִּפְגְּשֵׁ֣הוּ יְהֹוָ֔ה וַיְבַקֵּ֖שׁ הֲמִיתֽוֹ׃ \Rashi{\textbf{ויהי משה בדרך במלון: ויבקש המיתו.}(המלאך ל)משה: לפי שלא מל את אליעזר בנו, ועל שנתרשל נענש מיתה. תניא אמר רבי יוסי: חס ושלום, לא נתרשל אלא אמר, אמול ואצא לדרך, סכנה היא לתינוק עד שלושת ימים, אמול ואשהה ג' ימים, הקב"ה צוני לך שב מצרים ומפני מה נענש? לפי שנתעסק במלון תחלה (מכילתא). במסכת נדרים: והיה המלאך נעשה כמין נחש ובולעו מראשו ועד ירכיו וחוזר ובולעו מרגליו ועד אותו מקום, הבינה צפורה שבשביל המילה הוא (נדרים ל"ב):}%
% ---------- פרק 4 | פסוק 25 ----------
 \textbf{כה} וַתִּקַּ֨ח צִפֹּרָ֜ה צֹ֗ר וַתִּכְרֹת֙ אֶת־עׇרְלַ֣ת בְּנָ֔הּ וַתַּגַּ֖ע לְרַגְלָ֑יו וַתֹּ֕אמֶר כִּ֧י חֲתַן־דָּמִ֛ים אַתָּ֖ה לִֽי׃ \Rashi{\textbf{ותגע לרגליו.}השליכתו לפני רגליו של משה: ותאמר. על בנה: כי חתן דמים אתה לי. אתה הייתה גורם להיות החתן שלי נרצח עליך – הורג אישי אתה לי:}%
% ---------- פרק 4 | פסוק 26 ----------
 \textbf{כו} וַיִּ֖רֶף מִמֶּ֑נּוּ אָ֚ז אָֽמְרָ֔ה חֲתַ֥ן דָּמִ֖ים לַמּוּלֹֽת׃ \{פ\} \Rashi{\textbf{וירף.}המלאך ממנו, אז הבינה שעל המילה בא להרגו: אמרה חתן דמים למולת. חתני היה נרצח על דבר המילה: למולת. על דבר המולות. שם דבר הוא, והלמ"ד משמשת בלשון על, כמו "ואמר פרעה לבני ישראל" (שמות י"ד), ואנקלוס תרגם דמים על דם המילה:}%
% ---------- פרק 4 | פסוק 27 ----------
 \textbf{כז} וַיֹּ֤אמֶר יְהֹוָה֙ אֶֽל־אַהֲרֹ֔ן לֵ֛ךְ לִקְרַ֥את מֹשֶׁ֖ה הַמִּדְבָּ֑רָה וַיֵּ֗לֶךְ וַֽיִּפְגְּשֵׁ֛הוּ בְּהַ֥ר הָאֱלֹהִ֖ים וַיִּשַּׁק־לֽוֹ׃ %
% ---------- פרק 4 | פסוק 28 ----------
 \textbf{כח} וַיַּגֵּ֤ד מֹשֶׁה֙ לְאַֽהֲרֹ֔ן אֵ֛ת כׇּל־דִּבְרֵ֥י יְהֹוָ֖ה אֲשֶׁ֣ר שְׁלָח֑וֹ וְאֵ֥ת כׇּל־הָאֹתֹ֖ת אֲשֶׁ֥ר צִוָּֽהוּ׃ %
% ---------- פרק 4 | פסוק 29 ----------
 \textbf{כט} וַיֵּ֥לֶךְ מֹשֶׁ֖ה וְאַהֲרֹ֑ן וַיַּ֣אַסְפ֔וּ אֶת־כׇּל־זִקְנֵ֖י בְּנֵ֥י יִשְׂרָאֵֽל׃ %
% ---------- פרק 4 | פסוק 30 ----------
 \textbf{ל} וַיְדַבֵּ֣ר אַהֲרֹ֔ן אֵ֚ת כׇּל־הַדְּבָרִ֔ים אֲשֶׁר־דִּבֶּ֥ר יְהֹוָ֖ה אֶל־מֹשֶׁ֑ה וַיַּ֥עַשׂ הָאֹתֹ֖ת לְעֵינֵ֥י הָעָֽם׃ %
% ---------- פרק 4 | פסוק 31 ----------
 \textbf{לא} וַֽיַּאֲמֵ֖ן הָעָ֑ם וַֽיִּשְׁמְע֡וּ כִּֽי־פָקַ֨ד יְהֹוָ֜ה אֶת־בְּנֵ֣י יִשְׂרָאֵ֗ל וְכִ֤י רָאָה֙ אֶת־עׇנְיָ֔ם וַֽיִּקְּד֖וּ וַיִּֽשְׁתַּחֲוֽוּ׃ %
% ---------- פרק 5 | פסוק 1 ----------
 \textbf{א} וְאַחַ֗ר בָּ֚אוּ מֹשֶׁ֣ה וְאַהֲרֹ֔ן וַיֹּאמְר֖וּ אֶל־פַּרְעֹ֑ה כֹּֽה־אָמַ֤ר יְהֹוָה֙ אֱלֹהֵ֣י יִשְׂרָאֵ֔ל שַׁלַּח֙ אֶת־עַמִּ֔י וְיָחֹ֥גּוּ לִ֖י בַּמִּדְבָּֽר׃ \Rashi{\textbf{ואחר באו משה ואהרן וגו'.}אבל הזקנים נשמטו אחד אחד מאחר משה ואהרן, עד שנשמטו כלם קדם שהגיעו לפלטין, לפי שיראו ללכת; ובסיני נפרע להם, "ונגש משה לבדו אל ה' והם לא יגשו" (שמות כ"ד) – החזירם לאחוריהם (שמות רבה):}%
% ---------- פרק 5 | פסוק 2 ----------
 \textbf{ב} וַיֹּ֣אמֶר פַּרְעֹ֔ה מִ֤י יְהֹוָה֙ אֲשֶׁ֣ר אֶשְׁמַ֣ע בְּקֹל֔וֹ לְשַׁלַּ֖ח אֶת־יִשְׂרָאֵ֑ל לֹ֤א יָדַ֙עְתִּי֙ אֶת־יְהֹוָ֔ה וְגַ֥ם אֶת־יִשְׂרָאֵ֖ל לֹ֥א אֲשַׁלֵּֽחַ׃ %
% ---------- פרק 5 | פסוק 3 ----------
 \textbf{ג} וַיֹּ֣אמְר֔וּ אֱלֹהֵ֥י הָעִבְרִ֖ים נִקְרָ֣א עָלֵ֑ינוּ נֵ֣לְכָה נָּ֡א דֶּ֩רֶךְ֩ שְׁלֹ֨שֶׁת יָמִ֜ים בַּמִּדְבָּ֗ר וְנִזְבְּחָה֙ לַֽיהֹוָ֣ה אֱלֹהֵ֔ינוּ פֶּ֨ן־יִפְגָּעֵ֔נוּ בַּדֶּ֖בֶר א֥וֹ בֶחָֽרֶב׃ \Rashi{\textbf{פן יפגענו.}פן יפגעך היו צריכים לומר, אלא שחלקו כבוד למלכות. פגיעה זו לשון מקרה מות היא (שם):}%
% ---------- פרק 5 | פסוק 4 ----------
 \textbf{ד} וַיֹּ֤אמֶר אֲלֵהֶם֙ מֶ֣לֶךְ מִצְרַ֔יִם לָ֚מָּה מֹשֶׁ֣ה וְאַהֲרֹ֔ן תַּפְרִ֥יעוּ אֶת־הָעָ֖ם מִֽמַּעֲשָׂ֑יו לְכ֖וּ לְסִבְלֹתֵיכֶֽם׃ \Rashi{\textbf{תפריעו את העם ממעשיו.}תבדילו ותרחיקו אותם ממלאכתם, ששומעין לכם וסבורים לנוח מן המלאכה; וכן "פרעהו אל תעבר בו" (משלי ד') – רחקהו, וכן "ותפרעו כל עצתי" (שם א'), "כי פרע הוא" (שמות ל"ב) – נרחק ונתעב: לכו לסבלתיכם. לכו למלאכתכם שיש לכם לעשות בבתיכם, אבל מלאכת שעבוד מצרים לא היתה על שבטו של לוי; תדע לך, שהרי משה ואהרן יוצאים ובאים (שמות רבה):}%
% ---------- פרק 5 | פסוק 5 ----------
 \textbf{ה} וַיֹּ֣אמֶר פַּרְעֹ֔ה הֵן־רַבִּ֥ים עַתָּ֖ה עַ֣ם הָאָ֑רֶץ וְהִשְׁבַּתֶּ֥ם אֹתָ֖ם מִסִּבְלֹתָֽם׃ \Rashi{\textbf{הן רבים עתה עם הארץ.}שהעבודה מטלת עליהם, ואתם משביתים אותם מסבלותם, הפסד גדול הוא זה:}%
% ---------- פרק 5 | פסוק 6 ----------
 \textbf{ו} וַיְצַ֥ו פַּרְעֹ֖ה בַּיּ֣וֹם הַה֑וּא אֶת־הַנֹּגְשִׂ֣ים בָּעָ֔ם וְאֶת־שֹׁטְרָ֖יו לֵאמֹֽר׃ \Rashi{\textbf{הנגשים.}מצריים היו והשוטרים היו ישראלים; הנוגש ממנה על כמה שוטרים והשוטר ממנה לרדות בעושי המלאכה (שמות רבה):}%
% ---------- פרק 5 | פסוק 7 ----------
 \textbf{ז} לֹ֣א תֹאסִפ֞וּן לָתֵ֨ת תֶּ֧בֶן לָעָ֛ם לִלְבֹּ֥ן הַלְּבֵנִ֖ים כִּתְמ֣וֹל שִׁלְשֹׁ֑ם הֵ֚ם יֵֽלְכ֔וּ וְקֹשְׁשׁ֥וּ לָהֶ֖ם תֶּֽבֶן׃ \Rashi{\textbf{תבן.}אשטו"בלא, היו גובלין אותו עם הטיט: לבנים. טיוול"ש בלעז, שעושים מטיט ומיבשין אותן בחמה; ויש ששורפין אותן בכבשן: כתמול שלשום. כאשר הייתם עושים עד הנה; וקששו. ולקטו:}%
% ---------- פרק 5 | פסוק 8 ----------
 \textbf{ח} וְאֶת־מַתְכֹּ֨נֶת הַלְּבֵנִ֜ים אֲשֶׁ֣ר הֵם֩ עֹשִׂ֨ים תְּמ֤וֹל שִׁלְשֹׁם֙ תָּשִׂ֣ימוּ עֲלֵיהֶ֔ם לֹ֥א תִגְרְע֖וּ מִמֶּ֑נּוּ כִּֽי־נִרְפִּ֣ים הֵ֔ם עַל־כֵּ֗ן הֵ֤ם צֹֽעֲקִים֙ לֵאמֹ֔ר נֵלְכָ֖ה נִזְבְּחָ֥ה לֵאלֹהֵֽינוּ׃ \Rashi{\textbf{ואת מתכנת הלבנים.}סכום חשבון הלבנים שהיה כל אחד עושה ליום כשהיה התבן נתן להם, אותו סכום תשימו עליהם גם עתה, למען תכבד העבודה עליהם: כי נרפים. מן העבודה הם, לכך לבם פונה אל הבטלה וצועקים לאמר נלכה וגו': מתכנת. ותכן לבנים, "ולו נתכנו עלילות" (שמואל א ב'), "את הכסף המתכן" (מלכים ב י"ב), כלם לשון חשבון הם: נרפים. המלאכה רפויה בידם ועזובה מהם והם נרפים ממנה; רטריי"ט בלעז:}%
% ---------- פרק 5 | פסוק 9 ----------
 \textbf{ט} תִּכְבַּ֧ד הָעֲבֹדָ֛ה עַל־הָאֲנָשִׁ֖ים וְיַעֲשׂוּ־בָ֑הּ וְאַל־יִשְׁע֖וּ בְּדִבְרֵי־שָֽׁקֶר׃ \Rashi{\textbf{ואל ישעו בדברי שקר.}ואל יהגו וידברו תמיד בדברי רוח לאמר, נלכה נזבחה; ודומה לו "ואשעה בחקיך תמיד" (תהילים קי"ט); "למשל ולשנינה" (דברים כ"ח), מתרגמינן "ולשועי"; ויספר – "ואשתעי"; ואי אפשר לומר ישעו לשון "וישע ה' אל הבל", "ואל קין ואל מנחתו לא שעה" (בראשית ד'), ולפרש אל ישעו – אל יפנו, שאם כן היה לו לכתב ואל ישעו אל דברי שקר או לדברי שקר, כי כן גזרת כלם, "ישעה האדם על עושהו" (ישעיהו י"ז), "ולא ישעה אל המזבחות", (שם), "ולא שעו על קדוש ישראל" (שם ל"א), ולא מצאתי שמוש של בי"ת סמוכה לאחריהם, אבל אחר לשון דבור – כמתעסק לדבר בדבר – נופל לשון שמוש בי"ת, כגון: "הנדברים בך" (יחזקאל ל"ג), "ותדבר מרים ואהרן במשה" (במדבר י"ב), "המלאך הדובר בי" (זכריה ד'), "לדבר בם" (דברים י"א), "ואדברה בעדותיך" (תהילים קי"ט), אף כאן – אל ישעו בדברי שקר – אל יהיו נדברים בדברי שוא והבאי:}%
% ---------- פרק 5 | פסוק 10 ----------
 \textbf{י} וַיֵּ֨צְא֜וּ נֹגְשֵׂ֤י הָעָם֙ וְשֹׁ֣טְרָ֔יו וַיֹּאמְר֥וּ אֶל־הָעָ֖ם לֵאמֹ֑ר כֹּ֚ה אָמַ֣ר פַּרְעֹ֔ה אֵינֶ֛נִּי נֹתֵ֥ן לָכֶ֖ם תֶּֽבֶן׃ %
% ---------- פרק 5 | פסוק 11 ----------
 \textbf{יא} אַתֶּ֗ם לְכ֨וּ קְח֤וּ לָכֶם֙ תֶּ֔בֶן מֵאֲשֶׁ֖ר תִּמְצָ֑אוּ כִּ֣י אֵ֥ין נִגְרָ֛ע מֵעֲבֹדַתְכֶ֖ם דָּבָֽר׃ \Rashi{\textbf{אתם לכו קחו לכם תבן.}וצריכים אתם לילך בזריזות, כי אין נגרע מעבדתכם דבר, מכל סכום לבנים שהייתם עושים ליום בהיות התבן נתן לכם מזמן מבית המלך:}%
% ---------- פרק 5 | פסוק 12 ----------
 \textbf{יב} וַיָּ֥פֶץ הָעָ֖ם בְּכׇל־אֶ֣רֶץ מִצְרָ֑יִם לְקֹשֵׁ֥שׁ קַ֖שׁ לַתֶּֽבֶן׃ \Rashi{\textbf{לקשש קש לתבן.}לאסף אסיפה, ללקט לקט, לצרך תבן הטיט: קש. לשון לקוט, על שם שדבר המתפזר הוא וצריך לקוששו קרוי קש בשאר מקומות:}%
% ---------- פרק 5 | פסוק 13 ----------
 \textbf{יג} וְהַנֹּגְשִׂ֖ים אָצִ֣ים לֵאמֹ֑ר כַּלּ֤וּ מַעֲשֵׂיכֶם֙ דְּבַר־י֣וֹם בְּיוֹמ֔וֹ כַּאֲשֶׁ֖ר בִּהְי֥וֹת הַתֶּֽבֶן׃ \Rashi{\textbf{אצים.}דוחקים: דבר יום ביומו. חשבון של כל יום כלו ביומו, כאשר עשיתם בהיות התבן מוכן:}%
% ---------- פרק 5 | פסוק 14 ----------
 \textbf{יד} וַיֻּכּ֗וּ שֹֽׁטְרֵי֙ בְּנֵ֣י יִשְׂרָאֵ֔ל אֲשֶׁר־שָׂ֣מוּ עֲלֵהֶ֔ם נֹגְשֵׂ֥י פַרְעֹ֖ה לֵאמֹ֑ר מַדּ֡וּעַ לֹא֩ כִלִּיתֶ֨ם חׇקְכֶ֤ם לִלְבֹּן֙ כִּתְמ֣וֹל שִׁלְשֹׁ֔ם גַּם־תְּמ֖וֹל גַּם־הַיּֽוֹם׃ \Rashi{\textbf{ויכו שטרי בני ישראל.}השוטרים ישראלים היו וחסים על חבריהם מלדחקם, וכשהיו משלימים הלבנים לנוגשים שהם מצריים, והיה חסר מן הסכום, היו מלקין אותם על שלא דחקו את עושי המלאכה; לפיכך זכו אותם שוטרים להיות סנהדרין ונאצל מן הרוח אשר על משה והושם עליהם, שנאמר: "אספה לי שבעים איש מזקני ישראל וגו'" – מאותן שידעת הטובה שעשו במצרים – כי הם זקני העם ושוטריו (במדבר י"א), (שמות רבה): ויכו שטרי בני ישראל אשר שמו נגשי פרעה אותם לשוטרים עליהם, לאמר מדוע וגו'. למה ויכו? שהיו אומרים להם מדוע לא כליתם גם תמול גם היום חק הקצוב עליכם ללבן כתמול השלישי, שהוא יום שלפני אתמול; והוא היה בהיות התבן נתן להם: ויכו. לשון ויפעלו, הכו מיד אחרים – הנוגשים הכום:}%
% ---------- פרק 5 | פסוק 15 ----------
 \textbf{טו} וַיָּבֹ֗אוּ שֹֽׁטְרֵי֙ בְּנֵ֣י יִשְׂרָאֵ֔ל וַיִּצְעֲק֥וּ אֶל־פַּרְעֹ֖ה לֵאמֹ֑ר לָ֧מָּה תַעֲשֶׂ֦ה כֹ֖ה לַעֲבָדֶֽיךָ׃ %
% ---------- פרק 5 | פסוק 16 ----------
 \textbf{טז} תֶּ֗בֶן אֵ֤ין נִתָּן֙ לַעֲבָדֶ֔יךָ וּלְבֵנִ֛ים אֹמְרִ֥ים לָ֖נוּ עֲשׂ֑וּ וְהִנֵּ֧ה עֲבָדֶ֛יךָ מֻכִּ֖ים וְחָטָ֥את עַמֶּֽךָ׃ \Rashi{\textbf{ולבנים אומרים לנו.}הנוגשים עשו, כמנין הראשון: וחטאת עמך. אלו היה נקוד פתח, הייתי אומר שהוא דבוק, ודבר זה חטאת עמך הוא, עכשו שהוא קמץ, שם דבר הוא, וכך פרושו: ודבר זה מביא חטאת על עמך, כאלו כתוב וחטאת לעמך; כמו "כבאנה בית לחם" (רות א'), שהוא כמו לבית לחם, וכן הרבה:}%
% ---------- פרק 5 | פסוק 17 ----------
 \textbf{יז} וַיֹּ֛אמֶר נִרְפִּ֥ים אַתֶּ֖ם נִרְפִּ֑ים עַל־כֵּן֙ אַתֶּ֣ם אֹֽמְרִ֔ים נֵלְכָ֖ה נִזְבְּחָ֥ה לַֽיהֹוָֽה׃ %
% ---------- פרק 5 | פסוק 18 ----------
 \textbf{יח} וְעַתָּה֙ לְכ֣וּ עִבְד֔וּ וְתֶ֖בֶן לֹא־יִנָּתֵ֣ן לָכֶ֑ם וְתֹ֥כֶן לְבֵנִ֖ים תִּתֵּֽנוּ׃ \Rashi{\textbf{ותכן לבנים.}חשבון הלבנים; וכן "את הכסף המתכן" (מלכים ב י"ב) – המנוי, כמו שאמור בענין "ויצרו וימנו את הכסף":}%
% ---------- פרק 5 | פסוק 19 ----------
 \textbf{יט} וַיִּרְא֞וּ שֹֽׁטְרֵ֧י בְנֵֽי־יִשְׂרָאֵ֛ל אֹתָ֖ם בְּרָ֣ע לֵאמֹ֑ר לֹא־תִגְרְע֥וּ מִלִּבְנֵיכֶ֖ם דְּבַר־י֥וֹם בְּיוֹמֽוֹ׃ \Rashi{\textbf{ויראו שטרי בני ישראל.}את חבריהם הנרדים על ידם, ברע, ראו אותם ברעה וצרה המוצאת אותם, בהכבידם העבודה עליהם לאמר לא תגרעו וגו':}%
% ---------- פרק 5 | פסוק 20 ----------
 \textbf{כ} וַֽיִּפְגְּעוּ֙ אֶת־מֹשֶׁ֣ה וְאֶֽת־אַהֲרֹ֔ן נִצָּבִ֖ים לִקְרָאתָ֑ם בְּצֵאתָ֖ם מֵאֵ֥ת פַּרְעֹֽה׃ \Rashi{\textbf{ויפגעו.}אנשים מישראל את משה ואת אהרן וגו'. ורבותינו דרשו: כל "נצים" "ונצבים" דתן ואבירם היו, שנ' בהם "יצאו נצבים" (במדבר ט״ז:כ״ד), (נדרים ס"ד):}%
% ---------- פרק 5 | פסוק 21 ----------
 \textbf{כא} וַיֹּאמְר֣וּ אֲלֵהֶ֔ם יֵ֧רֶא יְהֹוָ֛ה עֲלֵיכֶ֖ם וְיִשְׁפֹּ֑ט אֲשֶׁ֧ר הִבְאַשְׁתֶּ֣ם אֶת־רֵיחֵ֗נוּ בְּעֵינֵ֤י פַרְעֹה֙ וּבְעֵינֵ֣י עֲבָדָ֔יו לָֽתֶת־חֶ֥רֶב בְּיָדָ֖ם לְהׇרְגֵֽנוּ׃ %
% ---------- פרק 5 | פסוק 22 ----------
 \textbf{כב} וַיָּ֧שׇׁב מֹשֶׁ֛ה אֶל־יְהֹוָ֖ה וַיֹּאמַ֑ר אֲדֹנָ֗י לָמָ֤ה הֲרֵעֹ֙תָה֙ לָעָ֣ם הַזֶּ֔ה לָ֥מָּה זֶּ֖ה שְׁלַחְתָּֽנִי׃ \Rashi{\textbf{למה הרעתה לעם הזה.}ואם תאמר מה אכפת לך? קובל אני על ששלחתני (שמות רבה):}%
% ---------- פרק 5 | פסוק 23 ----------
 \textbf{כג} וּמֵאָ֞ז בָּ֤אתִי אֶל־פַּרְעֹה֙ לְדַבֵּ֣ר בִּשְׁמֶ֔ךָ הֵרַ֖ע לָעָ֣ם הַזֶּ֑ה וְהַצֵּ֥ל לֹא־הִצַּ֖לְתָּ אֶת־עַמֶּֽךָ׃ \Rashi{\textbf{הרע.}לשון הפעיל הוא – הרבה רעה עליהם, ותרגומו "אבאיש":}%
% ---------- פרק 6 | פסוק 1 ----------
 \textbf{א} וַיֹּ֤אמֶר יְהֹוָה֙ אֶל־מֹשֶׁ֔ה עַתָּ֣ה תִרְאֶ֔ה אֲשֶׁ֥ר אֶֽעֱשֶׂ֖ה לְפַרְעֹ֑ה כִּ֣י בְיָ֤ד חֲזָקָה֙ יְשַׁלְּחֵ֔ם וּבְיָ֣ד חֲזָקָ֔ה יְגָרְשֵׁ֖ם מֵאַרְצֽוֹ׃ \{ס\} \Rashi{\textbf{עתה תראה וגו'.}הרהרת על מדותי, לא כאברהם שאמרתי לו "כי ביצחק יקרא לך זרע" (בראשית כ"א), ואחר כך אמרתי לו "העלהו לעולה" ולא הרהר אחרי; לפיכך עתה תראה. העשוי לפרעה תראה, ולא העשוי למלכי שבעה אמות כשאביאם לארץ (סנהדרין קי"א): כי ביד חזקה ישלחם. מפני ידי החזקה שתחזק עליו ישלחם: וביד חזקה יגרשם מארצו. על כרחם של ישראל יגרשם, לא יספיקו לעשות להם צדה; וכן הוא אומר "ותחזק מצרים על העם וגו'" (שמות י"ב):}%
\addparasha{חומש שמות | פרשת וארא}
% ---------- פרק 6 | פסוק 2 ----------
 \textbf{ב} וַיְדַבֵּ֥ר אֱלֹהִ֖ים אֶל־מֹשֶׁ֑ה וַיֹּ֥אמֶר אֵלָ֖יו אֲנִ֥י יְהֹוָֽה׃ \Rashi{\textbf{וידבר אלהים אל משה.}דבר אתו משפט, על שהקשה לדבר ולומר "למה הרעתה לעם הזה": ויאמר אליו אני ה'. נאמן לשלם שכר טוב למתהלכים לפני, ולא לחנם שלחתיך כי אם לקים דברי שדברתי לאבות הראשונים. ובלשון הזה מצינו שהוא נדרש בכמה מקומות אני ה' נאמן להפרע – כשהוא אמור אצל ענש – כגון "וחללת את שם אלהיך אני ה'" (ויקרא י״ט:י״ב), וכשהוא אומר אצל קיום מצוות – כגון "ושמרתם מצותי ועשיתם אותם אני ה'" (שם כ"ב) – נאמן לתן שכר:}%
% ---------- פרק 6 | פסוק 3 ----------
 \textbf{ג} וָאֵרָ֗א אֶל־אַבְרָהָ֛ם אֶל־יִצְחָ֥ק וְאֶֽל־יַעֲקֹ֖ב בְּאֵ֣ל שַׁדָּ֑י וּשְׁמִ֣י יְהֹוָ֔ה לֹ֥א נוֹדַ֖עְתִּי לָהֶֽם׃ \Rashi{\textbf{וארא.}אל האבות: באל שדי. הבטחתים הבטחות ובכלן אמרתי להם אני אל שדי: ושמי ה' לא נודעתי להם. לא הודעתי אין כתיב כאן אלא לא נודעתי, לא נכרתי להם במדת אמתות שלי, שעליה נקרא שמי ה', נאמן לאמת דברי, שהרי הבטחתים ולא קימתי:}%
% ---------- פרק 6 | פסוק 4 ----------
 \textbf{ד} וְגַ֨ם הֲקִמֹ֤תִי אֶת־בְּרִיתִי֙ אִתָּ֔ם לָתֵ֥ת לָהֶ֖ם אֶת־אֶ֣רֶץ כְּנָ֑עַן אֵ֛ת אֶ֥רֶץ מְגֻרֵיהֶ֖ם אֲשֶׁר־גָּ֥רוּ בָֽהּ׃ \Rashi{\textbf{וגם הקמתי את בריתי וגו'.}וגם כשנראיתי להם באל שדי הצבתי והעמדתי בריתי ביני וביניהם, לתת להם ארץ כנען. לאברהם בפרשת מילה נאמר "אני אל שדי וגו' ונתתי לך ולזרעך אחריך את ארץ מגריך" (בראשית י"ז), ליצחק "כי לך ולזרעך אתן את כל הארצת האל והקמתי את השבעה אשר נשבעתי לאברהם" (שם כ"ו), ואותה שבועה שנשבעתי לאברהם באל שדי אמרתי ליעקב – "אני אל שדי פרה ורבה וגו' ואת הארץ אשר וגו'" (שם ל"ה), הרי שנדרתי להם ולא קימתי:}\Peirush{לכם לא נאמר אלא להם מכאן לתחיית המתים מן התורה\Makor{סנהדרין צ:.}}%
% ---------- פרק 6 | פסוק 5 ----------
 \textbf{ה} וְגַ֣ם ׀ אֲנִ֣י שָׁמַ֗עְתִּי אֶֽת־נַאֲקַת֙ בְּנֵ֣י יִשְׂרָאֵ֔ל אֲשֶׁ֥ר מִצְרַ֖יִם מַעֲבִדִ֣ים אֹתָ֑ם וָאֶזְכֹּ֖ר אֶת־בְּרִיתִֽי׃ \Rashi{\textbf{וגם אני.}כמו שהצבתי והעמדתי הברית יש עלי לקים, לפיכך שמעתי את נאקת בני ישראל הנואקים, אשר מצרים מעבדים אתם ואזכר אותו הברית, כי בברית בין הבתרים אמרתי לו "וגם את הגוי אשר יעבדו דן אנכי" (בראשית ט"ו):}%
% ---------- פרק 6 | פסוק 6 ----------
 \textbf{ו} לָכֵ֞ן אֱמֹ֥ר לִבְנֵֽי־יִשְׂרָאֵל֮ אֲנִ֣י יְהֹוָה֒ וְהוֹצֵאתִ֣י אֶתְכֶ֗ם מִתַּ֙חַת֙ סִבְלֹ֣ת מִצְרַ֔יִם וְהִצַּלְתִּ֥י אֶתְכֶ֖ם מֵעֲבֹדָתָ֑ם וְגָאַלְתִּ֤י אֶתְכֶם֙ בִּזְר֣וֹעַ נְטוּיָ֔ה וּבִשְׁפָטִ֖ים גְּדֹלִֽים׃ \Rashi{\textbf{לכן.}על פי אותה השבועה: אמר לבני ישראל אני ה'. הנאמן בהבטחתי: והוצאתי אתכם. כי כן הבטחתי: "ואחרי כן יצאו ברכש גדול" (שם): סבלת. טרח משא מצרים:}\Peirush{נאמר כאן סבלות מצרים, ועל יוסף נאמר "הסירותי מסבל שכמו" (תהלים פא, ז), כמו שיוסף יצא מבית האסורים בראש השנה כך יצאנו ממצרים בראש השנה\Makor{ראש השנה יא:}}%
% ---------- פרק 6 | פסוק 7 ----------
 \textbf{ז} וְלָקַחְתִּ֨י אֶתְכֶ֥ם לִי֙ לְעָ֔ם וְהָיִ֥יתִי לָכֶ֖ם לֵֽאלֹהִ֑ים וִֽידַעְתֶּ֗ם כִּ֣י אֲנִ֤י יְהֹוָה֙ אֱלֹ֣הֵיכֶ֔ם הַמּוֹצִ֣יא אֶתְכֶ֔ם מִתַּ֖חַת סִבְל֥וֹת מִצְרָֽיִם׃ \Peirush{נאמר ולקחתי אתכם ונאמר "והבאתי אתכם" (ח), מקיש יציאת מצרים לכניסת הארץ מה כניסתם לארץ נכנסו שנים מששים ריבוא [יהושע וכלב] אף ביציאת מצרים יצאו שנים משישים ריבוא\Makor{סנהדרין קיא.}}\Peirush{ר' נחמיה סובר שמברכים מוציא לחם מן הארץ, משום שהמוציא זהו לשון הווה שנאמר המוציא אתכם\Makor{ברכות לח.}}%
% ---------- פרק 6 | פסוק 8 ----------
 \textbf{ח} וְהֵבֵאתִ֤י אֶתְכֶם֙ אֶל־הָאָ֔רֶץ אֲשֶׁ֤ר נָשָׂ֙אתִי֙ אֶת־יָדִ֔י לָתֵ֣ת אֹתָ֔הּ לְאַבְרָהָ֥ם לְיִצְחָ֖ק וּֽלְיַעֲקֹ֑ב וְנָתַתִּ֨י אֹתָ֥הּ לָכֶ֛ם מוֹרָשָׁ֖ה אֲנִ֥י יְהֹוָֽה׃ \Rashi{\textbf{נשאתי את ידי.}הרימותיה להשבע בכסאי:}\Peirush{ראה פס' ז'\Makor{סנהדרין קיא.}}\Peirush{נאמר לכם, מכאן שחלוקת הארץ היא באחוזים לפי יוצאי מצרים ולא לפי השבטים\Makor{בבא בתרא קיז: שם קיט:}}%
% ---------- פרק 6 | פסוק 9 ----------
 \textbf{ט} וַיְדַבֵּ֥ר מֹשֶׁ֛ה כֵּ֖ן אֶל־בְּנֵ֣י יִשְׂרָאֵ֑ל וְלֹ֤א שָֽׁמְעוּ֙ אֶל־מֹשֶׁ֔ה מִקֹּ֣צֶר ר֔וּחַ וּמֵעֲבֹדָ֖ה קָשָֽׁה׃ \{פ\} \Rashi{\textbf{ולא שמעו אל משה.}לא קבלו תנחומין: מקצר רוח. כל מי שהוא מצר, רוחו ונשימתו קצרה, ואינו יכול להאריך בנשימתו: קרוב לענין זה שמעתי בפרשה זו מרבי ברוך ברבי אליעזר, והביא לי ראיה ממקרא זה, "בפעם הזאת אודיעם את ידי ואת גבורתי וידעו כי שמי ה'" (ירמיהו ט"ז) – למדנו כשהקב"ה מאמן את דבריו אפלו לפרענות מודיע ששמו ה', וכל שכן האמנה לטובה. ורבותינו דרשוהו לענין של מעלה, שאמר משה למה הרעתה, אמר לו הקב"ה "חבל על דאבדין ולא משתכחין" – יש לי להתאונן על מיתת האבות, הרבה פעמים נגליתי אליהם באל שדי ולא אמרו לי מה שמך, ואתה אמרת מה שמו, מה אמר אליהם: וגם הקמתי וגו'. וכשבקש אברהם לקבר את שרה לא מצא קבר עד שקנה בדמים מרבים, וכן ביצחק ערערו עליו על הבארות אשר חפר, וכן יעקב ויקן את חלקת השדה, לנטות אהלו, ולא הרהרו אחר מדותי, ואתה אמרת למה הרעתה? ואין המדרש מתישב אחר המקרא מפני כמה דברים, אחת, שלא נאמר ושמי ה' לא שאלו לי, ואם תאמר לא הודיעם שכך שמו, הרי תחלה כשנגלה לאברהם בין הבתרים נאמר "אני ה' אשר הוצאתיך מאור כשדים" (בראשית ט"ו), ועוד היאך הסמיכה נמשכת בדברים שהוא סומך לכאן וגם אני שמעתי וגו' לכן אמר לבני ישראל? לכך אני אומר יתישב המקרא על פשוטו, דבור דבור על אפניו, והדרשה תדרש, שנאמר "הלוא כה דברי כאש נאם ה' וכפטיש יפוצץ סלע" (ירמיהו כ"ג) – מתחלק לכמה ניצוצות (סנהדרין ל"ד):}%
% ---------- פרק 6 | פסוק 10 ----------
 \textbf{י} וַיְדַבֵּ֥ר יְהֹוָ֖ה אֶל־מֹשֶׁ֥ה לֵּאמֹֽר׃ %
% ---------- פרק 6 | פסוק 11 ----------
 \textbf{יא} בֹּ֣א דַבֵּ֔ר אֶל־פַּרְעֹ֖ה מֶ֣לֶךְ מִצְרָ֑יִם וִֽישַׁלַּ֥ח אֶת־בְּנֵֽי־יִשְׂרָאֵ֖ל מֵאַרְצֽוֹ׃ %
% ---------- פרק 6 | פסוק 12 ----------
 \textbf{יב} וַיְדַבֵּ֣ר מֹשֶׁ֔ה לִפְנֵ֥י יְהֹוָ֖ה לֵאמֹ֑ר הֵ֤ן בְּנֵֽי־יִשְׂרָאֵל֙ לֹֽא־שָׁמְע֣וּ אֵלַ֔י וְאֵיךְ֙ יִשְׁמָעֵ֣נִי פַרְעֹ֔ה וַאֲנִ֖י עֲרַ֥ל שְׂפָתָֽיִם׃ \{פ\} \Rashi{\textbf{ואיך ישמעני פרעה.}זה אחד מעשרה קל וחמר שבתורה (בראשית רבה צ"ב): ערל שפתים. אטום שפתים; וכן כל לשון ערלה אני אומר שהוא אטום, "ערלה אזנם" (ירמיהו ו') – אטומה משמע, "ערלי לב" (שם ט') – אטומים מהבין, "שתה גם אתה והערל" (חבקוק ב') – והאטם משכרות כוס הקללה, "ערלת בשר" – שהגיד אטום ומכסה בה, "וערלתם ערלתו" (ויקרא י"ט) – עשו לו אטם וכסוי אסור שיבדיל בפני אכילתו, "שלש שנים יהיה לכם ערלים" (שם) – אטום ומכסה ומבדל מלאכלו:}%
% ---------- פרק 6 | פסוק 13 ----------
 \textbf{יג} וַיְדַבֵּ֣ר יְהֹוָה֮ אֶל־מֹשֶׁ֣ה וְאֶֽל־אַהֲרֹן֒ וַיְצַוֵּם֙ אֶל־בְּנֵ֣י יִשְׂרָאֵ֔ל וְאֶל־פַּרְעֹ֖ה מֶ֣לֶךְ מִצְרָ֑יִם לְהוֹצִ֥יא אֶת־בְּנֵֽי־יִשְׂרָאֵ֖ל מֵאֶ֥רֶץ מִצְרָֽיִם׃ \{ס\} \Rashi{\textbf{וידבר ה' אל משה ואל אהרן.}לפי שאמר משה אני ערל שפתים, צרף לו הקב"ה את אהרן להיות לו לפה ולמליץ: ויצום אל בני ישראל. צום עליהם להנהיגם בנחת ולסבל אותם (שמות רבה): ואל פרעה מלך מצרים. צום עליו לחלק לו כבוד בדבריהם, זה מדרשו (שם); ופשוטו, צום על דבר ישראל ועל שליחותו אל פרעה; ודבר הצווי מהו, מפרש בפרשה שניה לאחר סדר היחס, אלא מתוך שהזכיר משה ואהרן, הפסיק הענין באלה ראשי בית אבתם ללמדנו היאך נולדו משה ואהרן ובמי נתיחסו:}%
% ---------- פרק 6 | פסוק 14 ----------
 \textbf{יד} אֵ֖לֶּה רָאשֵׁ֣י בֵית־אֲבֹתָ֑ם בְּנֵ֨י רְאוּבֵ֜ן בְּכֹ֣ר יִשְׂרָאֵ֗ל חֲנ֤וֹךְ וּפַלּוּא֙ חֶצְרֹ֣ן וְכַרְמִ֔י אֵ֖לֶּה מִשְׁפְּחֹ֥ת רְאוּבֵֽן׃ \Rashi{\textbf{אלה ראשי בית אבתם.}מתוך שהזקק ליחס שבטו של לוי עד משה ואהרן בשביל משה ואהרן, התחיל ליחסם דרך תולדותם מראובן. ובפסיקתא רבתי ראיתי, לפי שקנטרם יעקב אבינו לשלושה שבטים הללו בשעת מותו, חזר הכתוב ויחסם כאן לבדם, לומר שחשובים הם:}%
% ---------- פרק 6 | פסוק 15 ----------
 \textbf{טו} וּבְנֵ֣י שִׁמְע֗וֹן יְמוּאֵ֨ל וְיָמִ֤ין וְאֹ֙הַד֙ וְיָכִ֣ין וְצֹ֔חַר וְשָׁא֖וּל בֶּן־הַֽכְּנַעֲנִ֑ית אֵ֖לֶּה מִשְׁפְּחֹ֥ת שִׁמְעֽוֹן׃ %
% ---------- פרק 6 | פסוק 16 ----------
 \textbf{טז} וְאֵ֨לֶּה שְׁמ֤וֹת בְּנֵֽי־לֵוִי֙ לְתֹ֣לְדֹתָ֔ם גֵּרְשׁ֕וֹן וּקְהָ֖ת וּמְרָרִ֑י וּשְׁנֵי֙ חַיֵּ֣י לֵוִ֔י שֶׁ֧בַע וּשְׁלֹשִׁ֛ים וּמְאַ֖ת שָׁנָֽה׃ \Rashi{\textbf{ושני חיי לוי וגו'.}למה נמנו שנותיו של לוי? להודיע כמה ימי השעבוד; שכל זמן שאחד מן השבטים קים לא היה שעבוד (שמות רבה א'), שנאמר "וימת יוסף וכל אחיו" (שמות א), ואחר כך "ויקם מלך חדש", ולוי האריך ימים על כלם:}%
% ---------- פרק 6 | פסוק 17 ----------
 \textbf{יז} בְּנֵ֥י גֵרְשׁ֛וֹן לִבְנִ֥י וְשִׁמְעִ֖י לְמִשְׁפְּחֹתָֽם׃ %
% ---------- פרק 6 | פסוק 18 ----------
 \textbf{יח} וּבְנֵ֣י קְהָ֔ת עַמְרָ֣ם וְיִצְהָ֔ר וְחֶבְר֖וֹן וְעֻזִּיאֵ֑ל וּשְׁנֵי֙ חַיֵּ֣י קְהָ֔ת שָׁלֹ֧שׁ וּשְׁלֹשִׁ֛ים וּמְאַ֖ת שָׁנָֽה׃ \Rashi{\textbf{ושני חיי קהת, ושני חיי עמרם וגו'.}מחשבון זה אנו למדים על מושב בני ישראל ארבע מאות שנה שאמר הכתוב (בראשית ט״ו:י״ג), שלא בארץ מצרים לבדה היו, אלא מיום שנולד יצחק; שהרי קהת מיורדי מצרים היה, חשב כל שנותיו ושנות עמרם ושמונים של משה, לא תמצאם ארבע מאות שנה, והרבה שנים נבלעים לבנים בשני האבות:}%
% ---------- פרק 6 | פסוק 19 ----------
 \textbf{יט} וּבְנֵ֥י מְרָרִ֖י מַחְלִ֣י וּמוּשִׁ֑י אֵ֛לֶּה מִשְׁפְּחֹ֥ת הַלֵּוִ֖י לְתֹלְדֹתָֽם׃ %
% ---------- פרק 6 | פסוק 20 ----------
 \textbf{כ} וַיִּקַּ֨ח עַמְרָ֜ם אֶת־יוֹכֶ֤בֶד דֹּֽדָתוֹ֙ ל֣וֹ לְאִשָּׁ֔ה וַתֵּ֣לֶד ל֔וֹ אֶֽת־אַהֲרֹ֖ן וְאֶת־מֹשֶׁ֑ה וּשְׁנֵי֙ חַיֵּ֣י עַמְרָ֔ם שֶׁ֧בַע וּשְׁלֹשִׁ֛ים וּמְאַ֖ת שָׁנָֽה׃ \Rashi{\textbf{יוכבד דודתו.}אחת אבוהי, בת לוי אחות קהת:}\Peirush{נאמר דודתו, אולי משמע דודתו מן האם, ומותר לשאת את אחות האם. לא, הכוונה לדודתו מן האב\Makor{סנהדרין נח:}}\Peirush{נאמר באחשוורוש שמלך על "שבע ועשרים ומאה מדינה" (אסתר א, א) ולומדים שבתחילה מלך על שבע ולבסוף מלך על עשרים ולבסוף מלך על מאה. ומקשים, נאמר בעמרם שבע ושלשים ומאת שנה, אך לא לומדים משם דרשה\Makor{מגילה יא.}}%
% ---------- פרק 6 | פסוק 21 ----------
 \textbf{כא} וּבְנֵ֖י יִצְהָ֑ר קֹ֥רַח וָנֶ֖פֶג וְזִכְרִֽי׃ %
% ---------- פרק 6 | פסוק 22 ----------
 \textbf{כב} וּבְנֵ֖י עֻזִּיאֵ֑ל מִֽישָׁאֵ֥ל וְאֶלְצָפָ֖ן וְסִתְרִֽי׃ %
% ---------- פרק 6 | פסוק 23 ----------
 \textbf{כג} וַיִּקַּ֨ח אַהֲרֹ֜ן אֶת־אֱלִישֶׁ֧בַע בַּת־עַמִּינָדָ֛ב אֲח֥וֹת נַחְשׁ֖וֹן ל֣וֹ לְאִשָּׁ֑ה וַתֵּ֣לֶד ל֗וֹ אֶת־נָדָב֙ וְאֶת־אֲבִיה֔וּא אֶת־אֶלְעָזָ֖ר וְאֶת־אִֽיתָמָֽר׃ \Rashi{\textbf{אחות נחשון.}מכאן למדו, הנושא אשה צריך לבדק באחיה (בבא בתרא ק"י):}\Peirush{הנושא אשה צריך שיבדוק באחיה שנאמר אחות נחשון\Makor{בבא בתרא קי.}}%
% ---------- פרק 6 | פסוק 24 ----------
 \textbf{כד} וּבְנֵ֣י קֹ֔רַח אַסִּ֥יר וְאֶלְקָנָ֖ה וַאֲבִיאָסָ֑ף אֵ֖לֶּה מִשְׁפְּחֹ֥ת הַקׇּרְחִֽי׃ %
% ---------- פרק 6 | פסוק 25 ----------
 \textbf{כה} וְאֶלְעָזָ֨ר בֶּֽן־אַהֲרֹ֜ן לָקַֽח־ל֨וֹ מִבְּנ֤וֹת פּֽוּטִיאֵל֙ ל֣וֹ לְאִשָּׁ֔ה וַתֵּ֥לֶד ל֖וֹ אֶת־פִּֽינְחָ֑ס אֵ֗לֶּה רָאשֵׁ֛י אֲב֥וֹת הַלְוִיִּ֖ם לְמִשְׁפְּחֹתָֽם׃ \Rashi{\textbf{מבנות פוטיאל.}מזרע יתרו שפטם עגלים לע"ז, ומזרע יוסף שפטפט ביצרו (סוטה מ"ג):}\Peirush{נאמר מבנות פוטיאל, מכאן שפנחס יצא מיתרו שפיטם עגלים לעבודת כוכבים, ומיוסף שפיטפט ביצרו [כבש אותו]\Makor{סוטה מג. בבא בתרא קט:}}%
% ---------- פרק 6 | פסוק 26 ----------
 \textbf{כו} ה֥וּא אַהֲרֹ֖ן וּמֹשֶׁ֑ה אֲשֶׁ֨ר אָמַ֤ר יְהֹוָה֙ לָהֶ֔ם הוֹצִ֜יאוּ אֶת־בְּנֵ֧י יִשְׂרָאֵ֛ל מֵאֶ֥רֶץ מִצְרַ֖יִם עַל־צִבְאֹתָֽם׃ \Rashi{\textbf{הוא אהרן ומשה.}אלו שהזכרו למעלה, שילדה יוכבד לעמרם. הוא אהרן ומשה. יש מקומות שמקדים אהרן למשה ויש מקומות שמקדים משה לאהרן, לומר לך ששקולין כאחד: על צבאתם. בצבאותם, כל צבאם לשבטיהם; יש על שאינו אלא במקום אות אחת, "ועל חרבך תחיה" (בראשית כ"ז), כמו בחרבך, "עמדתם על חרבכם" (יחזקאל ל"ג) – בחרבכם:}%
% ---------- פרק 6 | פסוק 27 ----------
 \textbf{כז} הֵ֗ם הַֽמְדַבְּרִים֙ אֶל־פַּרְעֹ֣ה מֶֽלֶךְ־מִצְרַ֔יִם לְהוֹצִ֥יא אֶת־בְּנֵֽי־יִשְׂרָאֵ֖ל מִמִּצְרָ֑יִם ה֥וּא מֹשֶׁ֖ה וְאַהֲרֹֽן׃ \Rashi{\textbf{הם המדברים וגו'.}הם שנצטוו, הם שקימו: הוא משה ואהרן. הם בשליחותם ובצדקתם מתחלה ועד סוף (מגילה י"א):}\Peirush{נאמר הוא אהרן ומשה, הן בצדקן מתחילתן ועד סופן\Makor{מגילה יא.}}%
% ---------- פרק 6 | פסוק 28 ----------
 \textbf{כח} וַיְהִ֗י בְּי֨וֹם דִּבֶּ֧ר יְהֹוָ֛ה אֶל־מֹשֶׁ֖ה בְּאֶ֥רֶץ מִצְרָֽיִם׃ \{ס\} \Rashi{\textbf{ויהי ביום דבר וגו'.}מחבר למקרא שלאחריו:}%
% ---------- פרק 6 | פסוק 29 ----------
 \textbf{כט} וַיְדַבֵּ֧ר יְהֹוָ֛ה אֶל־מֹשֶׁ֥ה לֵּאמֹ֖ר אֲנִ֣י יְהֹוָ֑ה דַּבֵּ֗ר אֶל־פַּרְעֹה֙ מֶ֣לֶךְ מִצְרַ֔יִם אֵ֛ת כׇּל־אֲשֶׁ֥ר אֲנִ֖י דֹּבֵ֥ר אֵלֶֽיךָ׃ \Rashi{\textbf{וידבר ה'.}הוא הדבור עצמו האמור למעלה בא דבר אל פרעה מלך מצרים, אלא מתוך שהפסיק הענין כדי ליחסם, חזר עליו להתחיל בו: אני ה'. כדאי אני לשלחך ולקים דברי שליחותי:}%
% ---------- פרק 6 | פסוק 30 ----------
 \textbf{ל} וַיֹּ֥אמֶר מֹשֶׁ֖ה לִפְנֵ֣י יְהֹוָ֑ה הֵ֤ן אֲנִי֙ עֲרַ֣ל שְׂפָתַ֔יִם וְאֵ֕יךְ יִשְׁמַ֥ע אֵלַ֖י פַּרְעֹֽה׃ \{פ\} \Rashi{\textbf{ויאמר משה לפני ה'.}היא האמירה שאמר למעלה הן בני ישראל לא שמעו אלי, ושנה הכתוב כאן, כיון שהפסיק הענין; וכן היא השיטה, כאדם האומר "נחזר על הראשונות":}%
% ---------- פרק 7 | פסוק 1 ----------
 \textbf{א} וַיֹּ֤אמֶר יְהֹוָה֙ אֶל־מֹשֶׁ֔ה רְאֵ֛ה נְתַתִּ֥יךָ אֱלֹהִ֖ים לְפַרְעֹ֑ה וְאַהֲרֹ֥ן אָחִ֖יךָ יִהְיֶ֥ה נְבִיאֶֽךָ׃ \Rashi{\textbf{נתתיך אלהים לפרעה.}שופט ורודה, לרדותו במכות ויסורין: יהיה נביאך. כתרגומו – "מתרגמנך". וכן כל לשון נבואה אדם המכריז ומשמיע לעם דברי תוכחות, והוא מגזרת "ניב שפתים" (ישעיהו נ"ז), "ינוב חכמה" (משלי י'), "ויכל מהתנבות" דשמואל (א' י'); ובלעז קוראין לו פרידיכ"ר:}%
% ---------- פרק 7 | פסוק 2 ----------
 \textbf{ב} אַתָּ֣ה תְדַבֵּ֔ר אֵ֖ת כׇּל־אֲשֶׁ֣ר אֲצַוֶּ֑ךָּ וְאַהֲרֹ֤ן אָחִ֙יךָ֙ יְדַבֵּ֣ר אֶל־פַּרְעֹ֔ה וְשִׁלַּ֥ח אֶת־בְּנֵֽי־יִשְׂרָאֵ֖ל מֵאַרְצֽוֹ׃ \Rashi{\textbf{אתה תדבר.}פעם אחת כל שליחות ושליחות כפי ששמעתו מפי, ואהרן אחיך ימליצנו ויטעימנו באזני פרעה:}%
% ---------- פרק 7 | פסוק 3 ----------
 \textbf{ג} וַאֲנִ֥י אַקְשֶׁ֖ה אֶת־לֵ֣ב פַּרְעֹ֑ה וְהִרְבֵּיתִ֧י אֶת־אֹתֹתַ֛י וְאֶת־מוֹפְתַ֖י בְּאֶ֥רֶץ מִצְרָֽיִם׃ \Rashi{\textbf{ואני אקשה.}מאחר שהרשיע והתריס כנגדי, וגלוי לפני שאין נחת רוח באמות לתת לב שלם לשוב, טוב שיתקשה לבו, למען הרבות בו אותותי, ותכירו אתם את גבורתי. וכן מדתו של הקב"ה, מביא פרענות על האמות כדי שישמעו ישראל וייראו, שנא' "הכרתי גוים נשמו פנותם … אמרתי אך תיראי אותי תקחי מוסר" (צפניה ג'); ואעפ"כ בחמש מכות הראשונות לא נאמר ויחזק ה' את לב פרעה, אלא ויחזק לב פרעה (תנחומא):}%
% ---------- פרק 7 | פסוק 4 ----------
 \textbf{ד} וְלֹֽא־יִשְׁמַ֤ע אֲלֵכֶם֙ פַּרְעֹ֔ה וְנָתַתִּ֥י אֶת־יָדִ֖י בְּמִצְרָ֑יִם וְהוֹצֵאתִ֨י אֶת־צִבְאֹתַ֜י אֶת־עַמִּ֤י בְנֵֽי־יִשְׂרָאֵל֙ מֵאֶ֣רֶץ מִצְרַ֔יִם בִּשְׁפָטִ֖ים גְּדֹלִֽים׃ \Rashi{\textbf{את ידי.}יד ממש, להכות בהם:}\Peirush{ניתן למחוק את השם 'צבאות' שלא נקרא אלא על שם ישראל שנאמר את צבאותי את עמי בני ישראל\Makor{שבועות לה:}}%
% ---------- פרק 7 | פסוק 5 ----------
 \textbf{ה} וְיָדְע֤וּ מִצְרַ֙יִם֙ כִּֽי־אֲנִ֣י יְהֹוָ֔ה בִּנְטֹתִ֥י אֶת־יָדִ֖י עַל־מִצְרָ֑יִם וְהוֹצֵאתִ֥י אֶת־בְּנֵֽי־יִשְׂרָאֵ֖ל מִתּוֹכָֽם׃ %
% ---------- פרק 7 | פסוק 6 ----------
 \textbf{ו} וַיַּ֥עַשׂ מֹשֶׁ֖ה וְאַהֲרֹ֑ן כַּאֲשֶׁ֨ר צִוָּ֧ה יְהֹוָ֛ה אֹתָ֖ם כֵּ֥ן עָשֽׂוּ׃ %
% ---------- פרק 7 | פסוק 7 ----------
 \textbf{ז} וּמֹשֶׁה֙ בֶּן־שְׁמֹנִ֣ים שָׁנָ֔ה וְאַֽהֲרֹ֔ן בֶּן־שָׁלֹ֥שׁ וּשְׁמֹנִ֖ים שָׁנָ֑ה בְּדַבְּרָ֖ם אֶל־פַּרְעֹֽה׃ \{פ\} %
% ---------- פרק 7 | פסוק 8 ----------
 \textbf{ח} וַיֹּ֣אמֶר יְהֹוָ֔ה אֶל־מֹשֶׁ֥ה וְאֶֽל־אַהֲרֹ֖ן לֵאמֹֽר׃ %
% ---------- פרק 7 | פסוק 9 ----------
 \textbf{ט} כִּי֩ יְדַבֵּ֨ר אֲלֵכֶ֤ם פַּרְעֹה֙ לֵאמֹ֔ר תְּנ֥וּ לָכֶ֖ם מוֹפֵ֑ת וְאָמַרְתָּ֣ אֶֽל־אַהֲרֹ֗ן קַ֧ח אֶֽת־מַטְּךָ֛ וְהַשְׁלֵ֥ךְ לִפְנֵֽי־פַרְעֹ֖ה יְהִ֥י לְתַנִּֽין׃ \Rashi{\textbf{מופת.}אות, להודיע שיש צרוך במי ששולח אתכם:}%
% ---------- פרק 7 | פסוק 10 ----------
 \textbf{י} וַיָּבֹ֨א מֹשֶׁ֤ה וְאַהֲרֹן֙ אֶל־פַּרְעֹ֔ה וַיַּ֣עֲשׂוּ כֵ֔ן כַּאֲשֶׁ֖ר צִוָּ֣ה יְהֹוָ֑ה וַיַּשְׁלֵ֨ךְ אַהֲרֹ֜ן אֶת־מַטֵּ֗הוּ לִפְנֵ֥י פַרְעֹ֛ה וְלִפְנֵ֥י עֲבָדָ֖יו וַיְהִ֥י לְתַנִּֽין׃ \Rashi{\textbf{לתנין.}נחש:}%
% ---------- פרק 7 | פסוק 11 ----------
 \textbf{יא} וַיִּקְרָא֙ גַּם־פַּרְעֹ֔ה לַֽחֲכָמִ֖ים וְלַֽמְכַשְּׁפִ֑ים וַיַּֽעֲשׂ֨וּ גַם־הֵ֜ם חַרְטֻמֵּ֥י מִצְרַ֛יִם בְּלַהֲטֵיהֶ֖ם כֵּֽן׃ \Rashi{\textbf{בלהטיהם.}בלחשיהון; ואין לו דמיון במקרא; ויש לדמות לו "להט החרב המתהפכת" (בראשית ג') – דומה שהיא מתהפכת על ידי לחש:}\Peirush{בלהטיהם אלו מעשה כשפים\Makor{סנהדרין סז:}}%
% ---------- פרק 7 | פסוק 12 ----------
 \textbf{יב} וַיַּשְׁלִ֙יכוּ֙ אִ֣ישׁ מַטֵּ֔הוּ וַיִּהְי֖וּ לְתַנִּינִ֑ם וַיִּבְלַ֥ע מַטֵּֽה־אַהֲרֹ֖ן אֶת־מַטֹּתָֽם׃ \Rashi{\textbf{ויבלע מטה אהרן.}מאחר שחזר ונעשה מטה, בלע את כלן (שבת צ"ז):}\Peirush{נס בתוך נס, לאחר שחזר ונעשה מטה בלע את המטות, ולא כשהוא תנין\Makor{שבת צז.}}%
% ---------- פרק 7 | פסוק 13 ----------
 \textbf{יג} וַיֶּֽחֱזַק֙ לֵ֣ב פַּרְעֹ֔ה וְלֹ֥א שָׁמַ֖ע אֲלֵהֶ֑ם כַּאֲשֶׁ֖ר דִּבֶּ֥ר יְהֹוָֽה׃ \{ס\} %
% ---------- פרק 7 | פסוק 14 ----------
 \textbf{יד} וַיֹּ֤אמֶר יְהֹוָה֙ אֶל־מֹשֶׁ֔ה כָּבֵ֖ד לֵ֣ב פַּרְעֹ֑ה מֵאֵ֖ן לְשַׁלַּ֥ח הָעָֽם׃ \Rashi{\textbf{כבד.}תרגומו יקיר ולא אתיקר, מפני שהוא שם דבר, כמו "כי כבד ממך הדבר" (שמות י"ח):}%
% ---------- פרק 7 | פסוק 15 ----------
 \textbf{טו} לֵ֣ךְ אֶל־פַּרְעֹ֞ה בַּבֹּ֗קֶר הִנֵּה֙ יֹצֵ֣א הַמַּ֔יְמָה וְנִצַּבְתָּ֥ לִקְרָאת֖וֹ עַל־שְׂפַ֣ת הַיְאֹ֑ר וְהַמַּטֶּ֛ה אֲשֶׁר־נֶהְפַּ֥ךְ לְנָחָ֖שׁ תִּקַּ֥ח בְּיָדֶֽךָ׃ \Rashi{\textbf{הנה יצא המימה.}לנקביו; שהיה עושה עצמו אלוה ואומר שאינו צריך לנקביו, ומשכים ויוצא לנילוס ועושה שם צרכיו (תנחומא):}\Peirush{אמר ר"ל מלך הוא והסבר לו פנים ורבי יוחנן אמר רשע הוא והעז פניך בו\Makor{זבחים קב.}}\Peirush{פרעה היה אמגושי\Makor{מועד קטן יח:}}%
% ---------- פרק 7 | פסוק 16 ----------
 \textbf{טז} וְאָמַרְתָּ֣ אֵלָ֗יו יְהֹוָ֞ה אֱלֹהֵ֤י הָעִבְרִים֙ שְׁלָחַ֤נִי אֵלֶ֙יךָ֙ לֵאמֹ֔ר שַׁלַּח֙ אֶת־עַמִּ֔י וְיַֽעַבְדֻ֖נִי בַּמִּדְבָּ֑ר וְהִנֵּ֥ה לֹא־שָׁמַ֖עְתָּ עַד־כֹּֽה׃ \Rashi{\textbf{עד כה.}עד הנה; ומדרשו: עד שתשמע ממני מכת בכורות שאפתח בה ב"כה" אמר ה' כחצות הלילה (שמות י״א:ד׳):}%
% ---------- פרק 7 | פסוק 17 ----------
 \textbf{יז} כֹּ֚ה אָמַ֣ר יְהֹוָ֔ה בְּזֹ֣את תֵּדַ֔ע כִּ֖י אֲנִ֣י יְהֹוָ֑ה הִנֵּ֨ה אָנֹכִ֜י מַכֶּ֣ה ׀ בַּמַּטֶּ֣ה אֲשֶׁר־בְּיָדִ֗י עַל־הַמַּ֛יִם אֲשֶׁ֥ר בַּיְאֹ֖ר וְנֶהֶפְכ֥וּ לְדָֽם׃ \Rashi{\textbf{ונהפכו לדם.}לפי שאין גשמים יורדים במצרים ונילוס עולה ומשקה את הארץ ומצרים עובדים לנילוס, לפיכך הלקה את יראתם ואח"כ הלקה אותם (תנחומא):}%
% ---------- פרק 7 | פסוק 18 ----------
 \textbf{יח} וְהַדָּגָ֧ה אֲשֶׁר־בַּיְאֹ֛ר תָּמ֖וּת וּבָאַ֣שׁ הַיְאֹ֑ר וְנִלְא֣וּ מִצְרַ֔יִם לִשְׁתּ֥וֹת מַ֖יִם מִן־הַיְאֹֽר׃ \{ס\} \Rashi{\textbf{ונלאו מצרים.}לבקש רפואה למי היאור שיהיו ראויין לשתות:}%
% ---------- פרק 7 | פסוק 19 ----------
 \textbf{יט} וַיֹּ֨אמֶר יְהֹוָ֜ה אֶל־מֹשֶׁ֗ה אֱמֹ֣ר אֶֽל־אַהֲרֹ֡ן קַ֣ח מַטְּךָ֣ וּנְטֵֽה־יָדְךָ֩ עַל־מֵימֵ֨י מִצְרַ֜יִם עַֽל־נַהֲרֹתָ֣ם ׀ עַל־יְאֹרֵיהֶ֣ם וְעַל־אַגְמֵיהֶ֗ם וְעַ֛ל כׇּל־מִקְוֵ֥ה מֵימֵיהֶ֖ם וְיִֽהְיוּ־דָ֑ם וְהָ֤יָה דָם֙ בְּכׇל־אֶ֣רֶץ מִצְרַ֔יִם וּבָעֵצִ֖ים וּבָאֲבָנִֽים׃ \Rashi{\textbf{אמר אל אהרן.}לפי שהגן היאור על משה כשנשלך לתוכו, לפיכך לא לקה על ידו לא בדם ולא בצפרדעים, ולקה על ידי אהרן (שמות רבה): נהרתם. הם נהרות המושכים, כעין נהרות שלנו: יאריהם. הם ברכות נגרים, העשויות בידי אדם משפת הנהר לשדות, ונילוס מימיו מתברכים ועולה דרך היאורים ומשקה השדות: אגמיהם. קבוצת מים שאינן נובעין ואין מושכין, אלא עומדין במקום אחד, וקורין לו אשטנ"ק: בכל ארץ מצרים. אף במרחצאות ובאמבטאות שבבתים: ובעצים ובאבנים. מים שבכלי עץ ובכלי אבן:}%
% ---------- פרק 7 | פסוק 20 ----------
 \textbf{כ} וַיַּֽעֲשׂוּ־כֵן֩ מֹשֶׁ֨ה וְאַהֲרֹ֜ן כַּאֲשֶׁ֣ר ׀ צִוָּ֣ה יְהֹוָ֗ה וַיָּ֤רֶם בַּמַּטֶּה֙ וַיַּ֤ךְ אֶת־הַמַּ֙יִם֙ אֲשֶׁ֣ר בַּיְאֹ֔ר לְעֵינֵ֣י פַרְעֹ֔ה וּלְעֵינֵ֖י עֲבָדָ֑יו וַיֵּהָ֥פְכ֛וּ כׇּל־הַמַּ֥יִם אֲשֶׁר־בַּיְאֹ֖ר לְדָֽם׃ %
% ---------- פרק 7 | פסוק 21 ----------
 \textbf{כא} וְהַדָּגָ֨ה אֲשֶׁר־בַּיְאֹ֥ר מֵ֙תָה֙ וַיִּבְאַ֣שׁ הַיְאֹ֔ר וְלֹא־יָכְל֣וּ מִצְרַ֔יִם לִשְׁתּ֥וֹת מַ֖יִם מִן־הַיְאֹ֑ר וַיְהִ֥י הַדָּ֖ם בְּכׇל־אֶ֥רֶץ מִצְרָֽיִם׃ \Peirush{הנודר לא לאכול דגה אסור בדגים קטנים וחייב בגדולים. אך קשה, הרי כתוב והדגה אשר ביאור מתה והכוונה גם לגדולים וגם לקטנים\Makor{נדרים נא:}}%
% ---------- פרק 7 | פסוק 22 ----------
 \textbf{כב} וַיַּֽעֲשׂוּ־כֵ֛ן חַרְטֻמֵּ֥י מִצְרַ֖יִם בְּלָטֵיהֶ֑ם וַיֶּחֱזַ֤ק לֵב־פַּרְעֹה֙ וְלֹא־שָׁמַ֣ע אֲלֵהֶ֔ם כַּאֲשֶׁ֖ר דִּבֶּ֥ר יְהֹוָֽה׃ \Rashi{\textbf{בלטיהם.}לחש שאומרין אותו בלט ובחשאי; ורבותינו אמרו, "בלטיהם" מעשה שדים, "בלהטיהם" מעשה כשפים (סנהדרין ס"ז): ויחזק לב פרעה. לומר על ידי מכשפות אתם עושים כן; "תבן אתם מכניסים לעפרים", עיר שכלה תבן, אף אתם מביאין מכשפות למצרים שכלה כשפים:}\Peirush{בלטיהם אלו מעשה שדים\Makor{סנהדרין סז:}}%
% ---------- פרק 7 | פסוק 23 ----------
 \textbf{כג} וַיִּ֣פֶן פַּרְעֹ֔ה וַיָּבֹ֖א אֶל־בֵּית֑וֹ וְלֹא־שָׁ֥ת לִבּ֖וֹ גַּם־לָזֹֽאת׃ \Rashi{\textbf{גם לזאת.}למופת המטה שנהפך לתנין, ולא לזה של דם:}%
% ---------- פרק 7 | פסוק 24 ----------
 \textbf{כד} וַיַּחְפְּר֧וּ כׇל־מִצְרַ֛יִם סְבִיבֹ֥ת הַיְאֹ֖ר מַ֣יִם לִשְׁתּ֑וֹת כִּ֣י לֹ֤א יָֽכְלוּ֙ לִשְׁתֹּ֔ת מִמֵּימֵ֖י הַיְאֹֽר׃ %
% ---------- פרק 7 | פסוק 25 ----------
 \textbf{כה} וַיִּמָּלֵ֖א שִׁבְעַ֣ת יָמִ֑ים אַחֲרֵ֥י הַכּוֹת־יְהֹוָ֖ה אֶת־הַיְאֹֽר׃ \{פ\} \Rashi{\textbf{וימלא.}מנין שבעת ימים, שלא שב היאור לקדמותו, שהיתה המכה משמשת רביע חדש, וג' חלקים היה מעיד ומתרה בהם (תנחומא):}%
% ---------- פרק 7 | פסוק 26 ----------
 \textbf{כו} וַיֹּ֤אמֶר יְהֹוָה֙ אֶל־מֹשֶׁ֔ה בֹּ֖א אֶל־פַּרְעֹ֑ה וְאָמַרְתָּ֣ אֵלָ֗יו כֹּ֚ה אָמַ֣ר יְהֹוָ֔ה שַׁלַּ֥ח אֶת־עַמִּ֖י וְיַֽעַבְדֻֽנִי׃ %
% ---------- פרק 7 | פסוק 27 ----------
 \textbf{כז} וְאִם־מָאֵ֥ן אַתָּ֖ה לְשַׁלֵּ֑חַ הִנֵּ֣ה אָנֹכִ֗י נֹגֵ֛ף אֶת־כׇּל־גְּבוּלְךָ֖ בַּֽצְפַרְדְּעִֽים׃ \Rashi{\textbf{ואם מאן אתה.}ואם סרבן אתה: מאן. כמו ממאן, מסרב, אלא כנוי האדם על שם המפעל, כמו "שלו" (איוב ט"ז), "שקט" (שופטים י"ח), "סר וזעף" (מלכים א כ'): נגף את כל גבולך. מכה, וכן כל לשון מגפה אינו לשון מיתה אלא לשון מכה, וכן "ונגפו אשה הרה" (שמות כ"א), אינו לשון מיתה, וכן "ובטרם יתנגפו רגליכם" (ירמיהו י"ג), "פן תגף באבן רגלך" (תהילים צ"א), "ולאבן נגף" (ישעיהו ח'):}%
% ---------- פרק 7 | פסוק 28 ----------
 \textbf{כח} וְשָׁרַ֣ץ הַיְאֹר֮ צְפַרְדְּעִים֒ וְעָלוּ֙ וּבָ֣אוּ בְּבֵיתֶ֔ךָ וּבַחֲדַ֥ר מִשְׁכָּבְךָ֖ וְעַל־מִטָּתֶ֑ךָ וּבְבֵ֤ית עֲבָדֶ֙יךָ֙ וּבְעַמֶּ֔ךָ וּבְתַנּוּרֶ֖יךָ וּבְמִשְׁאֲרוֹתֶֽיךָ׃ \Rashi{\textbf{ועלו.}מן היאור: בביתך. ואחר כך בבית עבדיך. הוא התחיל בעצה תחלה, "ויאמר אל עמו" (שמות א'), וממנו התחילה הפרענות:}\Peirush{מה ראו חנניה מישאל ועזריה שמסרו עצמם על קדושת השם לכבשן האש? קל וחומר מצפרדעים, מה צפרדעים שאין מצווין על קדושת השם, נאמר בהם ובאו ... ובתנוריך ובמשארותיך. אימתי משארות [חתיכות בצק] מצויות אצל תנור? הוי אומר בשעה שהתנור חם. אנו, שמצווין על קדושת השם, על אחת כמה וכמה\Makor{פסחים נג:}}%
% ---------- פרק 7 | פסוק 29 ----------
 \textbf{כט} וּבְכָ֥ה וּֽבְעַמְּךָ֖ וּבְכׇל־עֲבָדֶ֑יךָ יַעֲל֖וּ הַֽצְפַרְדְּעִֽים׃ \Rashi{\textbf{ובכה ובעמך.}בתוך מעיהם נכנסין ומקרקרין (סוטה י"א):}\Peirush{פרעה התנכל לישראל ראשון — לכן נענש ראשון מכל עמו, שנאמר ובכה ובעמך ובכל עבדיך יעלו הצפרדעים\Makor{סוטה יא.}}%
% ---------- פרק 8 | פסוק 1 ----------
 \textbf{א} וַיֹּ֣אמֶר יְהֹוָה֮ אֶל־מֹשֶׁה֒ אֱמֹ֣ר אֶֽל־אַהֲרֹ֗ן נְטֵ֤ה אֶת־יָדְךָ֙ בְּמַטֶּ֔ךָ עַ֨ל־הַנְּהָרֹ֔ת עַל־הַיְאֹרִ֖ים וְעַל־הָאֲגַמִּ֑ים וְהַ֥עַל אֶת־הַֽצְפַרְדְּעִ֖ים עַל־אֶ֥רֶץ מִצְרָֽיִם׃ %
% ---------- פרק 8 | פסוק 2 ----------
 \textbf{ב} וַיֵּ֤ט אַהֲרֹן֙ אֶת־יָד֔וֹ עַ֖ל מֵימֵ֣י מִצְרָ֑יִם וַתַּ֙עַל֙ הַצְּפַרְדֵּ֔עַ וַתְּכַ֖ס אֶת־אֶ֥רֶץ מִצְרָֽיִם׃ \Rashi{\textbf{ותעל הצפרדע.}צפרדע אחת היתה והיו מכין אותה והיא מתזת נחילים נחילים, זהו מדרשו (סנהדרין ס"ז). ופשוטו יש לומר, שרוץ הצפרדעים קורא לשון יחידות; וכן ותהי הכנם – הרחישה, פדולייר"א בלעז, ואף ותעל הצפרדע – גרינולייר"א בלעז:}\Peirush{נאמר ותעל הצפרדע מכאן שצפרדע אחת היתה בתחילה, השריצה ומלאה כל ארץ מצרים צפרדעים\Makor{סנהדרין סז:}}%
% ---------- פרק 8 | פסוק 3 ----------
 \textbf{ג} וַיַּֽעֲשׂוּ־כֵ֥ן הַֽחַרְטֻמִּ֖ים בְּלָטֵיהֶ֑ם וַיַּעֲל֥וּ אֶת־הַֽצְפַרְדְּעִ֖ים עַל־אֶ֥רֶץ מִצְרָֽיִם׃ %
% ---------- פרק 8 | פסוק 4 ----------
 \textbf{ד} וַיִּקְרָ֨א פַרְעֹ֜ה לְמֹשֶׁ֣ה וּֽלְאַהֲרֹ֗ן וַיֹּ֙אמֶר֙ הַעְתִּ֣ירוּ אֶל־יְהֹוָ֔ה וְיָסֵר֙ הַֽצְפַרְדְּעִ֔ים מִמֶּ֖נִּי וּמֵֽעַמִּ֑י וַאֲשַׁלְּחָה֙ אֶת־הָעָ֔ם וְיִזְבְּח֖וּ לַיהֹוָֽה׃ %
% ---------- פרק 8 | פסוק 5 ----------
 \textbf{ה} וַיֹּ֨אמֶר מֹשֶׁ֣ה לְפַרְעֹה֮ הִתְפָּאֵ֣ר עָלַי֒ לְמָתַ֣י ׀ אַעְתִּ֣יר לְךָ֗ וְלַעֲבָדֶ֙יךָ֙ וּֽלְעַמְּךָ֔ לְהַכְרִית֙ הַֽצְפַרְדְּעִ֔ים מִמְּךָ֖ וּמִבָּתֶּ֑יךָ רַ֥ק בַּיְאֹ֖ר תִּשָּׁאַֽרְנָה׃ \Rashi{\textbf{התפאר עלי.}כמו "היתפאר הגרזן על החצב בו" (ישעיהו י') – משתבח לומר אני גדול ממך, ונטי"ר בלעז. וכן התפאר עלי – השתבח להתחכם ולשאל דבר גדול, ולומר שלא אוכל לעשותו: למתי אעתיר לך. את אשר אעתיר לך היום על הכרתת הצפרדעים, למתי תרצה שיכרתו? ותראה אם אשלים דברי למועד שתקבע לי. אלו נאמר "מתי" אעתיר, היה משמע מתי אתפלל, עכשו שנאמר "למתי", משמע אני היום אתפלל עליך, שיכרתו הצפרדעים לזמן שתקבע עלי; אמר לאיזה יום תרצה שיכרתו: אעתיר, העתירו, והעתרתי. ולא נאמר אעתר, עתרו, ועתרתי, מפני שכל לשון עתר הרבות פלל הוא, וכאשר יאמר ארבה, הרבו, והרביתי, לשון מפעיל, כך יאמר אעתיר, העתירו, והעתרתי דברים; ואב לכלם "והעתרתם עלי דבריכם" (יחזקאל ל"ה) – הרביתם:}%
% ---------- פרק 8 | פסוק 6 ----------
 \textbf{ו} וַיֹּ֖אמֶר לְמָחָ֑ר וַיֹּ֙אמֶר֙ כִּדְבָ֣רְךָ֔ לְמַ֣עַן תֵּדַ֔ע כִּי־אֵ֖ין כַּיהֹוָ֥ה אֱלֹהֵֽינוּ׃ \Rashi{\textbf{ויאמר למחר.}התפלל היום, שיכרתו למחר:}%
% ---------- פרק 8 | פסוק 7 ----------
 \textbf{ז} וְסָר֣וּ הַֽצְפַרְדְּעִ֗ים מִמְּךָ֙ וּמִבָּ֣תֶּ֔יךָ וּמֵעֲבָדֶ֖יךָ וּמֵעַמֶּ֑ךָ רַ֥ק בַּיְאֹ֖ר תִּשָּׁאַֽרְנָה׃ %
% ---------- פרק 8 | פסוק 8 ----------
 \textbf{ח} וַיֵּצֵ֥א מֹשֶׁ֛ה וְאַהֲרֹ֖ן מֵעִ֣ם פַּרְעֹ֑ה וַיִּצְעַ֤ק מֹשֶׁה֙ אֶל־יְהֹוָ֔ה עַל־דְּבַ֥ר הַֽצְפַרְדְּעִ֖ים אֲשֶׁר־שָׂ֥ם לְפַרְעֹֽה׃ \Rashi{\textbf{ויצא, ויצעק.}מיד, שיכרתו למחר:}%
% ---------- פרק 8 | פסוק 9 ----------
 \textbf{ט} וַיַּ֥עַשׂ יְהֹוָ֖ה כִּדְבַ֣ר מֹשֶׁ֑ה וַיָּמֻ֙תוּ֙ הַֽצְפַרְדְּעִ֔ים מִן־הַבָּתִּ֥ים מִן־הַחֲצֵרֹ֖ת וּמִן־הַשָּׂדֹֽת׃ %
% ---------- פרק 8 | פסוק 10 ----------
 \textbf{י} וַיִּצְבְּר֥וּ אֹתָ֖ם חֳמָרִ֣ם חֳמָרִ֑ם וַתִּבְאַ֖שׁ הָאָֽרֶץ׃ \Rashi{\textbf{חמרם חמרם.}צבורים צבורים, כתרגומו, דגורין – גלין:}%
% ---------- פרק 8 | פסוק 11 ----------
 \textbf{יא} וַיַּ֣רְא פַּרְעֹ֗ה כִּ֤י הָֽיְתָה֙ הָֽרְוָחָ֔ה וְהַכְבֵּד֙ אֶת־לִבּ֔וֹ וְלֹ֥א שָׁמַ֖ע אֲלֵהֶ֑ם כַּאֲשֶׁ֖ר דִּבֶּ֥ר יְהֹוָֽה׃ \{ס\} \Rashi{\textbf{והכבד את לבו.}לשון פעול הוא, כמו "הלוך ונסוע" (בראשית י"ב), וכן "והכות את מואב" (מלכים ב ג׳:כ״ד), "ושאול לו באלהים" (שמואל א כ״ב:י״ג), "הכה ופצע" (מלכים א כ'): כאשר דבר ה'. והיכן דבר? ולא ישמע אליכם פרעה:}%
% ---------- פרק 8 | פסוק 12 ----------
 \textbf{יב} וַיֹּ֣אמֶר יְהֹוָה֮ אֶל־מֹשֶׁה֒ אֱמֹר֙ אֶֽל־אַהֲרֹ֔ן נְטֵ֣ה אֶֽת־מַטְּךָ֔ וְהַ֖ךְ אֶת־עֲפַ֣ר הָאָ֑רֶץ וְהָיָ֥ה לְכִנִּ֖ם בְּכׇל־אֶ֥רֶץ מִצְרָֽיִם׃ \Rashi{\textbf{אמר אל אהרן.}לא היה העפר כדאי ללקות ע"י משה, לפי שהגן עליו כשהרג את המצרי ויטמנהו בחול, ולקה על ידי אהרן (שמות רבה):}%
% ---------- פרק 8 | פסוק 13 ----------
 \textbf{יג} וַיַּֽעֲשׂוּ־כֵ֗ן וַיֵּט֩ אַהֲרֹ֨ן אֶת־יָד֤וֹ בְמַטֵּ֙הוּ֙ וַיַּךְ֙ אֶת־עֲפַ֣ר הָאָ֔רֶץ וַתְּהִי֙ הַכִּנָּ֔ם בָּאָדָ֖ם וּבַבְּהֵמָ֑ה כׇּל־עֲפַ֥ר הָאָ֛רֶץ הָיָ֥ה כִנִּ֖ים בְּכׇל־אֶ֥רֶץ מִצְרָֽיִם׃ \Rashi{\textbf{ותהי הכנם.}הרחישה, פדוליר"א בלעז:}%
% ---------- פרק 8 | פסוק 14 ----------
 \textbf{יד} וַיַּעֲשׂוּ־כֵ֨ן הַחַרְטֻמִּ֧ים בְּלָטֵיהֶ֛ם לְהוֹצִ֥יא אֶת־הַכִּנִּ֖ים וְלֹ֣א יָכֹ֑לוּ וַתְּהִי֙ הַכִּנָּ֔ם בָּאָדָ֖ם וּבַבְּהֵמָֽה׃ \Rashi{\textbf{להוציא את הכנים.}לבראתם ולהוציאם ממקום אחר: ולא יכלו. שאין השד שולט על בריה פחותה מכשעורה (סנהדרין ס"ז):}%
% ---------- פרק 8 | פסוק 15 ----------
 \textbf{טו} וַיֹּאמְר֤וּ הַֽחַרְטֻמִּם֙ אֶל־פַּרְעֹ֔ה אֶצְבַּ֥ע אֱלֹהִ֖ים הִ֑וא וַיֶּחֱזַ֤ק לֵב־פַּרְעֹה֙ וְלֹֽא־שָׁמַ֣ע אֲלֵהֶ֔ם כַּאֲשֶׁ֖ר דִּבֶּ֥ר יְהֹוָֽה׃ \{ס\} \Rashi{\textbf{אצבע אלהים הוא.}מכה זו אינה על ידי כשפים, מאת המקום היא (שמות רבה): כאשר דבר ה'. ולא ישמע אליכם פרעה:}\Peirush{מזה שהחרטומים לא יכלו לברוא כינים לומדים שאין השד יכול לברוא בריה פחות מכשעורה\Makor{סנהדרין סז:}}\Peirush{אצבע אלהים, מכאן שלקו מצרים באצבע, והיא האצבע שעתידה להפרע מסנחריב\Makor{סנהדרין צה:}}%
% ---------- פרק 8 | פסוק 16 ----------
 \textbf{טז} וַיֹּ֨אמֶר יְהֹוָ֜ה אֶל־מֹשֶׁ֗ה הַשְׁכֵּ֤ם בַּבֹּ֙קֶר֙ וְהִתְיַצֵּב֙ לִפְנֵ֣י פַרְעֹ֔ה הִנֵּ֖ה יוֹצֵ֣א הַמָּ֑יְמָה וְאָמַרְתָּ֣ אֵלָ֗יו כֹּ֚ה אָמַ֣ר יְהֹוָ֔ה שַׁלַּ֥ח עַמִּ֖י וְיַֽעַבְדֻֽנִי׃ %
% ---------- פרק 8 | פסוק 17 ----------
 \textbf{יז} כִּ֣י אִם־אֵינְךָ֮ מְשַׁלֵּ֣חַ אֶת־עַמִּי֒ הִנְנִי֩ מַשְׁלִ֨יחַ בְּךָ֜ וּבַעֲבָדֶ֧יךָ וּֽבְעַמְּךָ֛ וּבְבָתֶּ֖יךָ אֶת־הֶעָרֹ֑ב וּמָ֨לְא֜וּ בָּתֵּ֤י מִצְרַ֙יִם֙ אֶת־הֶ֣עָרֹ֔ב וְגַ֥ם הָאֲדָמָ֖ה אֲשֶׁר־הֵ֥ם עָלֶֽיהָ׃ \Rashi{\textbf{משליח בך.}מגרה בך, וכן "ושן בהמות אשלח בם" (דברים ל"ב), לשון שסוי, אינציט"ר בלעז: את הערב. כל מיני החיות רעות ונחשים ועקרבים בערבוביא, והיו משחיתים בהם; ויש טעם בדבר באגדה בכל מכה ומכה למה זו ולמה זו; בטכסיסי מלחמות מלכים בא עליהם, כסדר מלכות כשצרה על עיר, בתחלה מקלקל מעינותיה, ואחר כך תוקעין עליהם ומריעין בשופרות ליראם ולבהלם, וכן הצפרדעים מקרקרים והומים וכו', כדאיתא במדרש רבי תנחומא (פ' בא):}%
% ---------- פרק 8 | פסוק 18 ----------
 \textbf{יח} וְהִפְלֵיתִי֩ בַיּ֨וֹם הַה֜וּא אֶת־אֶ֣רֶץ גֹּ֗שֶׁן אֲשֶׁ֤ר עַמִּי֙ עֹמֵ֣ד עָלֶ֔יהָ לְבִלְתִּ֥י הֱיֽוֹת־שָׁ֖ם עָרֹ֑ב לְמַ֣עַן תֵּדַ֔ע כִּ֛י אֲנִ֥י יְהֹוָ֖ה בְּקֶ֥רֶב הָאָֽרֶץ׃ \Rashi{\textbf{והפליתי.}והפרשתי, וכן "והפלה ה'" (שמות ט'), וכן "לא נפלאת היא ממך" (דברים ל') – לא מבדלת ומפרשת היא ממך: למען תדע כי אני ה' בקרב הארץ. אף על פי ששכינתי בשמים גזרתי מתקימת בתחתונים:}%
% ---------- פרק 8 | פסוק 19 ----------
 \textbf{יט} וְשַׂמְתִּ֣י פְדֻ֔ת בֵּ֥ין עַמִּ֖י וּבֵ֣ין עַמֶּ֑ךָ לְמָחָ֥ר יִהְיֶ֖ה הָאֹ֥ת הַזֶּֽה׃ \Rashi{\textbf{ושמתי פדת.}שיבדיל בין עמי ובין עמך:}%
% ---------- פרק 8 | פסוק 20 ----------
 \textbf{כ} וַיַּ֤עַשׂ יְהֹוָה֙ כֵּ֔ן וַיָּבֹא֙ עָרֹ֣ב כָּבֵ֔ד בֵּ֥יתָה פַרְעֹ֖ה וּבֵ֣ית עֲבָדָ֑יו וּבְכׇל־אֶ֧רֶץ מִצְרַ֛יִם תִּשָּׁחֵ֥ת הָאָ֖רֶץ מִפְּנֵ֥י הֶעָרֹֽב׃ \Rashi{\textbf{תשחת הארץ.}נשחתת הארץ, אתחבלת ארעא:}%
% ---------- פרק 8 | פסוק 21 ----------
 \textbf{כא} וַיִּקְרָ֣א פַרְעֹ֔ה אֶל־מֹשֶׁ֖ה וּֽלְאַהֲרֹ֑ן וַיֹּ֗אמֶר לְכ֛וּ זִבְח֥וּ לֵאלֹֽהֵיכֶ֖ם בָּאָֽרֶץ׃ \Rashi{\textbf{זבחו לאלהיכם בארץ.}במקומכם ולא תלכו במדבר:}%
% ---------- פרק 8 | פסוק 22 ----------
 \textbf{כב} וַיֹּ֣אמֶר מֹשֶׁ֗ה לֹ֤א נָכוֹן֙ לַעֲשׂ֣וֹת כֵּ֔ן כִּ֚י תּוֹעֲבַ֣ת מִצְרַ֔יִם נִזְבַּ֖ח לַיהֹוָ֣ה אֱלֹהֵ֑ינוּ הֵ֣ן נִזְבַּ֞ח אֶת־תּוֹעֲבַ֥ת מִצְרַ֛יִם לְעֵינֵיהֶ֖ם וְלֹ֥א יִסְקְלֻֽנוּ׃ \Rashi{\textbf{תועבת מצרים.}יראת מצרים, כמו "ולמלכם תועבת בני עמון" (מלכים ב כ״ג:י״ג), ואצל ישראל קורא אותה תועבה; ועוד יש לומר בלשון אחר, תועבת מצרים, דבר שנאוי הוא למצרים זביחה שאנו זובחים, שהרי יראתם אנו זובחים: ולא יסקלנו. בתמיה:}%
% ---------- פרק 8 | פסוק 23 ----------
 \textbf{כג} דֶּ֚רֶךְ שְׁלֹ֣שֶׁת יָמִ֔ים נֵלֵ֖ךְ בַּמִּדְבָּ֑ר וְזָבַ֙חְנוּ֙ לַֽיהֹוָ֣ה אֱלֹהֵ֔ינוּ כַּאֲשֶׁ֖ר יֹאמַ֥ר אֵלֵֽינוּ׃ %
% ---------- פרק 8 | פסוק 24 ----------
 \textbf{כד} וַיֹּ֣אמֶר פַּרְעֹ֗ה אָנֹכִ֞י אֲשַׁלַּ֤ח אֶתְכֶם֙ וּזְבַחְתֶּ֞ם לַיהֹוָ֤ה אֱלֹֽהֵיכֶם֙ בַּמִּדְבָּ֔ר רַ֛ק הַרְחֵ֥ק לֹא־תַרְחִ֖יקוּ לָלֶ֑כֶת הַעְתִּ֖ירוּ בַּעֲדִֽי׃ %
% ---------- פרק 8 | פסוק 25 ----------
 \textbf{כה} וַיֹּ֣אמֶר מֹשֶׁ֗ה הִנֵּ֨ה אָנֹכִ֜י יוֹצֵ֤א מֵֽעִמָּךְ֙ וְהַעְתַּרְתִּ֣י אֶל־יְהֹוָ֔ה וְסָ֣ר הֶעָרֹ֗ב מִפַּרְעֹ֛ה מֵעֲבָדָ֥יו וּמֵעַמּ֖וֹ מָחָ֑ר רַ֗ק אַל־יֹסֵ֤ף פַּרְעֹה֙ הָתֵ֔ל לְבִלְתִּי֙ שַׁלַּ֣ח אֶת־הָעָ֔ם לִזְבֹּ֖חַ לַֽיהֹוָֽה׃ \Rashi{\textbf{התל.}כמו להתל:}%
% ---------- פרק 8 | פסוק 26 ----------
 \textbf{כו} וַיֵּצֵ֥א מֹשֶׁ֖ה מֵעִ֣ם פַּרְעֹ֑ה וַיֶּעְתַּ֖ר אֶל־יְהֹוָֽה׃ \Rashi{\textbf{ויעתר אל ה'.}נתאמץ בתפלה, וכן אם בא לומר ויעתיר, היה יכול לומר, ומשמע וירבה בתפלה, וכשהוא אומר בלשון ויפעל, משמע וירבה להתפלל:}%
% ---------- פרק 8 | פסוק 27 ----------
 \textbf{כז} וַיַּ֤עַשׂ יְהֹוָה֙ כִּדְבַ֣ר מֹשֶׁ֔ה וַיָּ֙סַר֙ הֶעָרֹ֔ב מִפַּרְעֹ֖ה מֵעֲבָדָ֣יו וּמֵעַמּ֑וֹ לֹ֥א נִשְׁאַ֖ר אֶחָֽד׃ \Rashi{\textbf{ויסר הערב.}ולא מתו כמו שמתו הצפרדעים, שאם מתו היה להם הנאה בעורותם (תנחומא):}%
% ---------- פרק 8 | פסוק 28 ----------
 \textbf{כח} וַיַּכְבֵּ֤ד פַּרְעֹה֙ אֶת־לִבּ֔וֹ גַּ֖ם בַּפַּ֣עַם הַזֹּ֑את וְלֹ֥א שִׁלַּ֖ח אֶת־הָעָֽם׃ \{פ\} \Rashi{\textbf{גם בפעם הזאת.}אע"פ שאמר אנכי אשלח אתכם, לא קיים הבטחתו:}%
% ---------- פרק 9 | פסוק 1 ----------
 \textbf{א} וַיֹּ֤אמֶר יְהֹוָה֙ אֶל־מֹשֶׁ֔ה בֹּ֖א אֶל־פַּרְעֹ֑ה וְדִבַּרְתָּ֣ אֵלָ֗יו כֹּֽה־אָמַ֤ר יְהֹוָה֙ אֱלֹהֵ֣י הָֽעִבְרִ֔ים שַׁלַּ֥ח אֶת־עַמִּ֖י וְיַֽעַבְדֻֽנִי׃ %
% ---------- פרק 9 | פסוק 2 ----------
 \textbf{ב} כִּ֛י אִם־מָאֵ֥ן אַתָּ֖ה לְשַׁלֵּ֑חַ וְעוֹדְךָ֖ מַחֲזִ֥יק בָּֽם׃ \Rashi{\textbf{מחזיק בם.}אוחז בם, כמו "והחזיקה במבשיו" (דברים כ"ה):}%
% ---------- פרק 9 | פסוק 3 ----------
 \textbf{ג} הִנֵּ֨ה יַד־יְהֹוָ֜ה הוֹיָ֗ה בְּמִקְנְךָ֙ אֲשֶׁ֣ר בַּשָּׂדֶ֔ה בַּסּוּסִ֤ים בַּֽחֲמֹרִים֙ בַּגְּמַלִּ֔ים בַּבָּקָ֖ר וּבַצֹּ֑אן דֶּ֖בֶר כָּבֵ֥ד מְאֹֽד׃ \Rashi{\textbf{הנה יד ה' הויה.}לשון הוה, כי כן יאמר בלשון נקבה על שעבר היתה, ועל העתיד תהיה, ועל העומד הויה, כמו עושה, רוצה, רועה:}%
% ---------- פרק 9 | פסוק 4 ----------
 \textbf{ד} וְהִפְלָ֣ה יְהֹוָ֔ה בֵּ֚ין מִקְנֵ֣ה יִשְׂרָאֵ֔ל וּבֵ֖ין מִקְנֵ֣ה מִצְרָ֑יִם וְלֹ֥א יָמ֛וּת מִכׇּל־לִבְנֵ֥י יִשְׂרָאֵ֖ל דָּבָֽר׃ \Rashi{\textbf{והפלה.}והבדיל:}%
% ---------- פרק 9 | פסוק 5 ----------
 \textbf{ה} וַיָּ֥שֶׂם יְהֹוָ֖ה מוֹעֵ֣ד לֵאמֹ֑ר מָחָ֗ר יַעֲשֶׂ֧ה יְהֹוָ֛ה הַדָּבָ֥ר הַזֶּ֖ה בָּאָֽרֶץ׃ %
% ---------- פרק 9 | פסוק 6 ----------
 \textbf{ו} וַיַּ֨עַשׂ יְהֹוָ֜ה אֶת־הַדָּבָ֤ר הַזֶּה֙ מִֽמׇּחֳרָ֔ת וַיָּ֕מׇת כֹּ֖ל מִקְנֵ֣ה מִצְרָ֑יִם וּמִמִּקְנֵ֥ה בְנֵֽי־יִשְׂרָאֵ֖ל לֹא־מֵ֥ת אֶחָֽד׃ %
% ---------- פרק 9 | פסוק 7 ----------
 \textbf{ז} וַיִּשְׁלַ֣ח פַּרְעֹ֔ה וְהִנֵּ֗ה לֹא־מֵ֛ת מִמִּקְנֵ֥ה יִשְׂרָאֵ֖ל עַד־אֶחָ֑ד וַיִּכְבַּד֙ לֵ֣ב פַּרְעֹ֔ה וְלֹ֥א שִׁלַּ֖ח אֶת־הָעָֽם׃ \{פ\} %
% ---------- פרק 9 | פסוק 8 ----------
 \textbf{ח} וַיֹּ֣אמֶר יְהֹוָה֮ אֶל־מֹשֶׁ֣ה וְאֶֽל־אַהֲרֹן֒ קְח֤וּ לָכֶם֙ מְלֹ֣א חׇפְנֵיכֶ֔ם פִּ֖יחַ כִּבְשָׁ֑ן וּזְרָק֥וֹ מֹשֶׁ֛ה הַשָּׁמַ֖יְמָה לְעֵינֵ֥י פַרְעֹֽה׃ \Rashi{\textbf{מלא חפניכם.}ילויני"ש בלעז: פיח כבשן. דבר הנפח מן הגחלים עמומים הנשרפים בכבשן, ובלעז אולבי"ש: פיח. לשון הפחה, שהרוח מפיחן ומפריחן: וזרקו משה. וכל דבר הנזרק בכח, אינו נזרק אלא ביד אחת, הרי נסים הרבה – אחד שהחזיק קמצו של משה מלא חפנים שלו ושל אהרן, ואחד שהלך האבק על כל ארץ מצרים (תנחומא):}%
% ---------- פרק 9 | פסוק 9 ----------
 \textbf{ט} וְהָיָ֣ה לְאָבָ֔ק עַ֖ל כׇּל־אֶ֣רֶץ מִצְרָ֑יִם וְהָיָ֨ה עַל־הָאָדָ֜ם וְעַל־הַבְּהֵמָ֗ה לִשְׁחִ֥ין פֹּרֵ֛חַ אֲבַעְבֻּעֹ֖ת בְּכׇל־אֶ֥רֶץ מִצְרָֽיִם׃ \Rashi{\textbf{לשחין פרח אבעבעת.}כתרגומו, לשיחנא סגי אבעבועין, שעל ידו צומחין בהן בועות: שחין. לשון חמימות; והרבה יש בלשון משנה, "שנה שחונה" (יומא נ"ג):}%
% ---------- פרק 9 | פסוק 10 ----------
 \textbf{י} וַיִּקְח֞וּ אֶת־פִּ֣יחַ הַכִּבְשָׁ֗ן וַיַּֽעַמְדוּ֙ לִפְנֵ֣י פַרְעֹ֔ה וַיִּזְרֹ֥ק אֹת֛וֹ מֹשֶׁ֖ה הַשָּׁמָ֑יְמָה וַיְהִ֗י שְׁחִין֙ אֲבַעְבֻּעֹ֔ת פֹּרֵ֕חַ בָּאָדָ֖ם וּבַבְּהֵמָֽה׃ \Rashi{\textbf{באדם ובבהמה.}ואם תאמר מאין היו להם הבהמות, והלא כבר נאמר וימת כל מקנה מצרים? לא נגזרה גזרה אלא על אותן שבשדות בלבד, שנאמר במקנך אשר בשדה, והירא את דבר ה' הכניס את מקנהו אל הבתים, כן שנויה במכילתא אצל "ויקח שש מאות רכב בחור" (שמות י"ד):}\Peirush{השחין היה לח מבחוץ ויבש מבפנים, שחין – יבש מבפנים, אבעבעת פרח – לח מבחוץ\Makor{בבא קמא פ:, בכורות מא.}}%
% ---------- פרק 9 | פסוק 11 ----------
 \textbf{יא} וְלֹֽא־יָכְל֣וּ הַֽחַרְטֻמִּ֗ים לַעֲמֹ֛ד לִפְנֵ֥י מֹשֶׁ֖ה מִפְּנֵ֣י הַשְּׁחִ֑ין כִּֽי־הָיָ֣ה הַשְּׁחִ֔ין בַּֽחַרְטֻמִּ֖ם וּבְכׇל־מִצְרָֽיִם׃ %
% ---------- פרק 9 | פסוק 12 ----------
 \textbf{יב} וַיְחַזֵּ֤ק יְהֹוָה֙ אֶת־לֵ֣ב פַּרְעֹ֔ה וְלֹ֥א שָׁמַ֖ע אֲלֵהֶ֑ם כַּאֲשֶׁ֛ר דִּבֶּ֥ר יְהֹוָ֖ה אֶל־מֹשֶֽׁה׃ \{ס\} %
% ---------- פרק 9 | פסוק 13 ----------
 \textbf{יג} וַיֹּ֤אמֶר יְהֹוָה֙ אֶל־מֹשֶׁ֔ה הַשְׁכֵּ֣ם בַּבֹּ֔קֶר וְהִתְיַצֵּ֖ב לִפְנֵ֣י פַרְעֹ֑ה וְאָמַרְתָּ֣ אֵלָ֗יו כֹּֽה־אָמַ֤ר יְהֹוָה֙ אֱלֹהֵ֣י הָֽעִבְרִ֔ים שַׁלַּ֥ח אֶת־עַמִּ֖י וְיַֽעַבְדֻֽנִי׃ %
% ---------- פרק 9 | פסוק 14 ----------
 \textbf{יד} כִּ֣י ׀ בַּפַּ֣עַם הַזֹּ֗את אֲנִ֨י שֹׁלֵ֜חַ אֶת־כׇּל־מַגֵּפֹתַי֙ אֶֽל־לִבְּךָ֔ וּבַעֲבָדֶ֖יךָ וּבְעַמֶּ֑ךָ בַּעֲב֣וּר תֵּדַ֔ע כִּ֛י אֵ֥ין כָּמֹ֖נִי בְּכׇל־הָאָֽרֶץ׃ \Rashi{\textbf{את כל מגפתי.}למדנו מכאן שמכת בכורות שקולה כנגד כל המכות:}%
% ---------- פרק 9 | פסוק 15 ----------
 \textbf{טו} כִּ֤י עַתָּה֙ שָׁלַ֣חְתִּי אֶת־יָדִ֔י וָאַ֥ךְ אוֹתְךָ֛ וְאֶֽת־עַמְּךָ֖ בַּדָּ֑בֶר וַתִּכָּחֵ֖ד מִן־הָאָֽרֶץ׃ \Rashi{\textbf{כי עתה שלחתי את ידי וגו'.}כי אלו רציתי, כשהיתה ידי במקנך שהכיתים בדבר, שלחתיה והכיתי אותך ואת עמך עם הבהמות, ותכחד מן הארץ, אבל בעבור זאת העמדתיך וגו':}%
% ---------- פרק 9 | פסוק 16 ----------
 \textbf{טז} וְאוּלָ֗ם בַּעֲב֥וּר זֹאת֙ הֶעֱמַדְתִּ֔יךָ בַּעֲב֖וּר הַרְאֹתְךָ֣ אֶת־כֹּחִ֑י וּלְמַ֛עַן סַפֵּ֥ר שְׁמִ֖י בְּכׇל־הָאָֽרֶץ׃ %
% ---------- פרק 9 | פסוק 17 ----------
 \textbf{יז} עוֹדְךָ֖ מִסְתּוֹלֵ֣ל בְּעַמִּ֑י לְבִלְתִּ֖י שַׁלְּחָֽם׃ \Rashi{\textbf{עודך מסתולל בעמי.}כתרגומו, כבישת בה בעמי, והוא מגזרת מסלה, דמתרגמינן ארח כבישא, ובלעז קלקי"ר; וכבר פרשתי בסוף ויהי מקץ, כל תבה שתחלת יסודה סמ"ך והיא באה לדבר בלשון מתפעל, נותן התי"ו של שמוש באמצע אותיות של עקר, כגון זו, וכגון "ויסתבל החגב" (קה' י"ב), מגזרת סבל, "כי תשתרר עלינו" (במדבר ט"ז), מגזרת נגיד ושר, "משתכל הוית" (דניאל ז'):}%
% ---------- פרק 9 | פסוק 18 ----------
 \textbf{יח} הִנְנִ֤י מַמְטִיר֙ כָּעֵ֣ת מָחָ֔ר בָּרָ֖ד כָּבֵ֣ד מְאֹ֑ד אֲשֶׁ֨ר לֹא־הָיָ֤ה כָמֹ֙הוּ֙ בְּמִצְרַ֔יִם לְמִן־הַיּ֥וֹם הִוָּסְדָ֖הֿ וְעַד־עָֽתָּה׃ \Rashi{\textbf{כעת מחר.}כעת הזאת למחר; שרט לו שריטה בכתל – למחר כשתגיע חמה לכאן, ירד הברד: הוסדה. שנתיסדה; וכל תבה שתחלת יסודה יו"ד, כגון יסד, ילד, ידע, יסר, כשהיא מתפעלת, תבא הוי"ו במקום היו"ד, כמו הוסדה, הולדה (הושע ב'), ויודע, "ויולד ליוסף" (בראשית מ"ו), "בדברים לא יוסר עבד" (משלי כ"ט):}%
% ---------- פרק 9 | פסוק 19 ----------
 \textbf{יט} וְעַתָּ֗ה שְׁלַ֤ח הָעֵז֙ אֶֽת־מִקְנְךָ֔ וְאֵ֛ת כׇּל־אֲשֶׁ֥ר לְךָ֖ בַּשָּׂדֶ֑ה כׇּל־הָאָדָ֨ם וְהַבְּהֵמָ֜ה אֲשֶֽׁר־יִמָּצֵ֣א בַשָּׂדֶ֗ה וְלֹ֤א יֵֽאָסֵף֙ הַבַּ֔יְתָה וְיָרַ֧ד עֲלֵהֶ֛ם הַבָּרָ֖ד וָמֵֽתוּ׃ \Rashi{\textbf{שלח העז.}כתרגומו, שלח כנוש, וכן "יושבי הגבים העיזו" (ישעיה י'), "העזו בני בנימן" (ירמיה ו'): ולא יאסף הביתה. לשון הכנסה הוא:}%
% ---------- פרק 9 | פסוק 20 ----------
 \textbf{כ} הַיָּרֵא֙ אֶת־דְּבַ֣ר יְהֹוָ֔ה מֵֽעַבְדֵ֖י פַּרְעֹ֑ה הֵנִ֛יס אֶת־עֲבָדָ֥יו וְאֶת־מִקְנֵ֖הוּ אֶל־הַבָּתִּֽים׃ \Rashi{\textbf{הניס.}הבריח:}%
% ---------- פרק 9 | פסוק 21 ----------
 \textbf{כא} וַאֲשֶׁ֥ר לֹא־שָׂ֛ם לִבּ֖וֹ אֶל־דְּבַ֣ר יְהֹוָ֑ה וַֽיַּעֲזֹ֛ב אֶת־עֲבָדָ֥יו וְאֶת־מִקְנֵ֖הוּ בַּשָּׂדֶֽה׃ \{פ\} %
% ---------- פרק 9 | פסוק 22 ----------
 \textbf{כב} וַיֹּ֨אמֶר יְהֹוָ֜ה אֶל־מֹשֶׁ֗ה נְטֵ֤ה אֶת־יָֽדְךָ֙ עַל־הַשָּׁמַ֔יִם וִיהִ֥י בָרָ֖ד בְּכׇל־אֶ֣רֶץ מִצְרָ֑יִם עַל־הָאָדָ֣ם וְעַל־הַבְּהֵמָ֗ה וְעַ֛ל כׇּל־עֵ֥שֶׂב הַשָּׂדֶ֖ה בְּאֶ֥רֶץ מִצְרָֽיִם׃ \Rashi{\textbf{על השמים.}לצד השמים. ומדרש אגדה: הגביהו הקב"ה למשה למעלה מן השמים:}%
% ---------- פרק 9 | פסוק 23 ----------
 \textbf{כג} וַיֵּ֨ט מֹשֶׁ֣ה אֶת־מַטֵּ֘הוּ֮ עַל־הַשָּׁמַ֒יִם֒ וַֽיהֹוָ֗ה נָתַ֤ן קֹלֹת֙ וּבָרָ֔ד וַתִּ֥הֲלַךְ אֵ֖שׁ אָ֑רְצָה וַיַּמְטֵ֧ר יְהֹוָ֛ה בָּרָ֖ד עַל־אֶ֥רֶץ מִצְרָֽיִם׃ %
% ---------- פרק 9 | פסוק 24 ----------
 \textbf{כד} וַיְהִ֣י בָרָ֔ד וְאֵ֕שׁ מִתְלַקַּ֖חַת בְּת֣וֹךְ הַבָּרָ֑ד כָּבֵ֣ד מְאֹ֔ד אֲ֠שֶׁ֠ר לֹֽא־הָיָ֤ה כָמֹ֙הוּ֙ בְּכׇל־אֶ֣רֶץ מִצְרַ֔יִם מֵאָ֖ז הָיְתָ֥ה לְגֽוֹי׃ \Rashi{\textbf{מתלקחת בתוך הברד.}נס בתוך נס, האש והברד מערבין, והברד מים הוא, ולעשות רצון קונם עשו שלום ביניהם:}%
% ---------- פרק 9 | פסוק 25 ----------
 \textbf{כה} וַיַּ֨ךְ הַבָּרָ֜ד בְּכׇל־אֶ֣רֶץ מִצְרַ֗יִם אֵ֚ת כׇּל־אֲשֶׁ֣ר בַּשָּׂדֶ֔ה מֵאָדָ֖ם וְעַד־בְּהֵמָ֑ה וְאֵ֨ת כׇּל־עֵ֤שֶׂב הַשָּׂדֶה֙ הִכָּ֣ה הַבָּרָ֔ד וְאֶת־כׇּל־עֵ֥ץ הַשָּׂדֶ֖ה שִׁבֵּֽר׃ %
% ---------- פרק 9 | פסוק 26 ----------
 \textbf{כו} רַ֚ק בְּאֶ֣רֶץ גֹּ֔שֶׁן אֲשֶׁר־שָׁ֖ם בְּנֵ֣י יִשְׂרָאֵ֑ל לֹ֥א הָיָ֖ה בָּרָֽד׃ %
% ---------- פרק 9 | פסוק 27 ----------
 \textbf{כז} וַיִּשְׁלַ֣ח פַּרְעֹ֗ה וַיִּקְרָא֙ לְמֹשֶׁ֣ה וּֽלְאַהֲרֹ֔ן וַיֹּ֥אמֶר אֲלֵהֶ֖ם חָטָ֣אתִי הַפָּ֑עַם יְהֹוָה֙ הַצַּדִּ֔יק וַאֲנִ֥י וְעַמִּ֖י הָרְשָׁעִֽים׃ \Peirush{בסוף מסכת ידים מובא הפסוק "מי ה' אשר אשמע בקלו" (ה, ב), ע"מ לסיים בפס' טוב, המשנה אומרת וכשלקה פרעה, אמר ה' הצדיק\Makor{משנה ידים ד, ח}}%
% ---------- פרק 9 | פסוק 28 ----------
 \textbf{כח} הַעְתִּ֙ירוּ֙ אֶל־יְהֹוָ֔ה וְרַ֕ב מִֽהְיֹ֛ת קֹלֹ֥ת אֱלֹהִ֖ים וּבָרָ֑ד וַאֲשַׁלְּחָ֣ה אֶתְכֶ֔ם וְלֹ֥א תֹסִפ֖וּן לַעֲמֹֽד׃ \Rashi{\textbf{ורב.}די לו במה שהוריד כבר:}%
% ---------- פרק 9 | פסוק 29 ----------
 \textbf{כט} וַיֹּ֤אמֶר אֵלָיו֙ מֹשֶׁ֔ה כְּצֵאתִי֙ אֶת־הָעִ֔יר אֶפְרֹ֥שׂ אֶת־כַּפַּ֖י אֶל־יְהֹוָ֑ה הַקֹּל֣וֹת יֶחְדָּל֗וּן וְהַבָּרָד֙ לֹ֣א יִֽהְיֶה־ע֔וֹד לְמַ֣עַן תֵּדַ֔ע כִּ֥י לַיהֹוָ֖ה הָאָֽרֶץ׃ \Rashi{\textbf{כצאתי את העיר.}מן העיר, אבל בתוך העיר לא התפלל לפי שהיתה מלאה גלולים (שמות רבה):}%
% ---------- פרק 9 | פסוק 30 ----------
 \textbf{ל} וְאַתָּ֖ה וַעֲבָדֶ֑יךָ יָדַ֕עְתִּי כִּ֚י טֶ֣רֶם תִּֽירְא֔וּן מִפְּנֵ֖י יְהֹוָ֥ה אֱלֹהִֽים׃ \Rashi{\textbf{טרם תיראון.}עדין לא תיראון; וכן כל טרם שבמקרא "עדין לא" הוא, ואינו לשון קדם, "טרם ישכבו" (בראשית י"ט) – עד לא שכיבו, "טרם יצמח" (שם ב') – עד לא צמח. אף זה כן הוא – ידעתי כי עדין אינכם יראים, ומשתהיה הרוחה תעמדו בקלקולכם:}%
% ---------- פרק 9 | פסוק 31 ----------
 \textbf{לא} וְהַפִּשְׁתָּ֥ה וְהַשְּׂעֹרָ֖ה נֻכָּ֑תָה כִּ֤י הַשְּׂעֹרָה֙ אָבִ֔יב וְהַפִּשְׁתָּ֖ה גִּבְעֹֽל׃ \Rashi{\textbf{והפשתה והשערה נכתה.}נשברה, לשון "פרעה נכה" (מלכים ב כ"ג), "נכאים" (ישעיהו ט"ז), וכן לא נכו; ולא יתכן לפרשו לשון הכאה, שאין נו"ן במקום ה"א לפרש נכתה כמו הכתה, נכו כמו הכו, אלא הנו"ן שרש בתבה והרי הוא מגזרת "ושפו עצמותיו" (איוב ל"ג): כי השערה אביב. כבר בכרה ועומדת בקשיה, ונשתברו ונפלו, וכן הפשתה גדלה כבר, והקשה לעמד בגבעוליה: השערה אביב. עמדה באביה, לשון "באבי הנחל" (שיר ו'):}\Peirush{נאמר אביב כאן ונאמר במנחת הביכורים, מה כאן שעורים אף שם שעורים\Makor{מנחות סח:, פד.}}%
% ---------- פרק 9 | פסוק 32 ----------
 \textbf{לב} וְהַחִטָּ֥ה וְהַכֻּסֶּ֖מֶת לֹ֣א נֻכּ֑וּ כִּ֥י אֲפִילֹ֖ת הֵֽנָּה׃ \Rashi{\textbf{כי אפילת הנה.}מאחרות, ועדין היו רכות ויכולות לעמד בפני קשה; ואע"פ שנאמר "ואת כל עשב השדה הכה הברד", יש לפרש פשוטו של מקרא בעשבים העומדים בקלחם, הראוים ללקות בברד. ובמדרש רבי תנחומא יש מרבותינו שנחלקו על זאת, ודרשו כי אפילת, פלאי פלאות נעשו להם שלא לקו:}%
% ---------- פרק 9 | פסוק 33 ----------
 \textbf{לג} וַיֵּצֵ֨א מֹשֶׁ֜ה מֵעִ֤ם פַּרְעֹה֙ אֶת־הָעִ֔יר וַיִּפְרֹ֥שׂ כַּפָּ֖יו אֶל־יְהֹוָ֑ה וַֽיַּחְדְּל֤וּ הַקֹּלוֹת֙ וְהַבָּרָ֔ד וּמָטָ֖ר לֹא־נִתַּ֥ךְ אָֽרְצָה׃ \Rashi{\textbf{לא נתך.}לא הגיע, ואף אותן שהיו באויר לא הגיעו לארץ; ודומה לו: "ותתך עלינו האלה והשבועה" (דניאל ט') דעזרא – ותגיע עלינו. ומנחם בן סרוק חברו בחלק "כהתוך כסף" (יחזקאל כ"ב), לשון יציקת מתכת, ורואה אני את דבריו, כתרגומו ויצק – ואתיך, לצקת – לאתכא. אף זה לא נתך לארץ – לא הוצק לארץ:}\Peirush{הרואה את המקום שנעשה נס אבני אלגביש חייב להודות לה'. אלגביש - אבנים שעמדו באויר על גב [בגלל] איש, וירדו על גב איש. עמדו על גב איש — זה משה שנאמר "והאיש משה" (יב, ג). ירדו על גב איש – זה יהושע שנאמר "יהושע בן נון איש אשר רוח בו" (במדבר כז, יח) ונאמר "וה' השליך עליהם אבנים גדלות" (יהושע י, יא)\Makor{ברכות נד:}}%
% ---------- פרק 9 | פסוק 34 ----------
 \textbf{לד} וַיַּ֣רְא פַּרְעֹ֗ה כִּֽי־חָדַ֨ל הַמָּטָ֧ר וְהַבָּרָ֛ד וְהַקֹּלֹ֖ת וַיֹּ֣סֶף לַחֲטֹ֑א וַיַּכְבֵּ֥ד לִבּ֖וֹ ה֥וּא וַעֲבָדָֽיו׃ %
% ---------- פרק 9 | פסוק 35 ----------
 \textbf{לה} וַיֶּֽחֱזַק֙ לֵ֣ב פַּרְעֹ֔ה וְלֹ֥א שִׁלַּ֖ח אֶת־בְּנֵ֣י יִשְׂרָאֵ֑ל כַּאֲשֶׁ֛ר דִּבֶּ֥ר יְהֹוָ֖ה בְּיַד־מֹשֶֽׁה׃ \{פ\} %
\addparasha{חומש שמות | פרשת בא}
% ---------- פרק 10 | פסוק 1 ----------
 \textbf{א} וַיֹּ֤אמֶר יְהֹוָה֙ אֶל־מֹשֶׁ֔ה בֹּ֖א אֶל־פַּרְעֹ֑ה כִּֽי־אֲנִ֞י הִכְבַּ֤דְתִּי אֶת־לִבּוֹ֙ וְאֶת־לֵ֣ב עֲבָדָ֔יו לְמַ֗עַן שִׁתִ֛י אֹתֹתַ֥י אֵ֖לֶּה בְּקִרְבּֽוֹ׃ \Rashi{\textbf{ויאמר ה' אל משה בא אל פרעה.}והתרה בו: שיתי. שומי – שאשית אני:}%
% ---------- פרק 10 | פסוק 2 ----------
 \textbf{ב} וּלְמַ֡עַן תְּסַפֵּר֩ בְּאׇזְנֵ֨י בִנְךָ֜ וּבֶן־בִּנְךָ֗ אֵ֣ת אֲשֶׁ֤ר הִתְעַלַּ֙לְתִּי֙ בְּמִצְרַ֔יִם וְאֶת־אֹתֹתַ֖י אֲשֶׁר־שַׂ֣מְתִּי בָ֑ם וִֽידַעְתֶּ֖ם כִּי־אֲנִ֥י יְהֹוָֽה׃ \Rashi{\textbf{התעללתי.}שחקתי, כמו "כי התעללת בי" (במדבר כ"ב), "הלא כאשר התעולל בהם" (שמואל א' ו'), האמור במצרים, ואינו לשון פועל ומעללים, שאם כן היה לו לכתב עוללתי, כמו "ועולל למו כאשר עוללת לי" (איכה א'), "אשר עולל לי" (שם): ולמען תספר. בתורה להודיע לדורות:}%
% ---------- פרק 10 | פסוק 3 ----------
 \textbf{ג} וַיָּבֹ֨א מֹשֶׁ֣ה וְאַהֲרֹן֮ אֶל־פַּרְעֹה֒ וַיֹּאמְר֣וּ אֵלָ֗יו כֹּֽה־אָמַ֤ר יְהֹוָה֙ אֱלֹהֵ֣י הָֽעִבְרִ֔ים עַד־מָתַ֣י מֵאַ֔נְתָּ לֵעָנֹ֖ת מִפָּנָ֑י שַׁלַּ֥ח עַמִּ֖י וְיַֽעַבְדֻֽנִי׃ \Rashi{\textbf{לענת.}כתרגומו, "לאתכנעא", והוא מגזרת עני – מאנת להיות עני ושפל מפני:}%
% ---------- פרק 10 | פסוק 4 ----------
 \textbf{ד} כִּ֛י אִם־מָאֵ֥ן אַתָּ֖ה לְשַׁלֵּ֣חַ אֶת־עַמִּ֑י הִנְנִ֨י מֵבִ֥יא מָחָ֛ר אַרְבֶּ֖ה בִּגְבֻלֶֽךָ׃ %
% ---------- פרק 10 | פסוק 5 ----------
 \textbf{ה} וְכִסָּה֙ אֶת־עֵ֣ין הָאָ֔רֶץ וְלֹ֥א יוּכַ֖ל לִרְאֹ֣ת אֶת־הָאָ֑רֶץ וְאָכַ֣ל ׀ אֶת־יֶ֣תֶר הַפְּלֵטָ֗ה הַנִּשְׁאֶ֤רֶת לָכֶם֙ מִן־הַבָּרָ֔ד וְאָכַל֙ אֶת־כׇּל־הָעֵ֔ץ הַצֹּמֵ֥חַ לָכֶ֖ם מִן־הַשָּׂדֶֽה׃ \Rashi{\textbf{את עין הארץ.}את מראה הארץ: ולא יוכל. הרואה לראת את הארץ, ולשון קצרה דבר:}%
% ---------- פרק 10 | פסוק 6 ----------
 \textbf{ו} וּמָלְא֨וּ בָתֶּ֜יךָ וּבָתֵּ֣י כׇל־עֲבָדֶ֘יךָ֮ וּבָתֵּ֣י כׇל־מִצְרַ֒יִם֒ אֲשֶׁ֨ר לֹֽא־רָא֤וּ אֲבֹתֶ֙יךָ֙ וַאֲב֣וֹת אֲבֹתֶ֔יךָ מִיּ֗וֹם הֱיוֹתָם֙ עַל־הָ֣אֲדָמָ֔ה עַ֖ד הַיּ֣וֹם הַזֶּ֑ה וַיִּ֥פֶן וַיֵּצֵ֖א מֵעִ֥ם פַּרְעֹֽה׃ %
% ---------- פרק 10 | פסוק 7 ----------
 \textbf{ז} וַיֹּאמְרוּ֩ עַבְדֵ֨י פַרְעֹ֜ה אֵלָ֗יו עַד־מָתַי֙ יִהְיֶ֨ה זֶ֥ה לָ֙נוּ֙ לְמוֹקֵ֔שׁ שַׁלַּח֙ אֶת־הָ֣אֲנָשִׁ֔ים וְיַֽעַבְד֖וּ אֶת־יְהֹוָ֣ה אֱלֹהֵיהֶ֑ם הֲטֶ֣רֶם תֵּדַ֔ע כִּ֥י אָבְדָ֖ה מִצְרָֽיִם׃ \Rashi{\textbf{הטרם תדע.}העוד לא ידעת כי אבדה מצרים:}%
% ---------- פרק 10 | פסוק 8 ----------
 \textbf{ח} וַיּוּשַׁ֞ב אֶת־מֹשֶׁ֤ה וְאֶֽת־אַהֲרֹן֙ אֶל־פַּרְעֹ֔ה וַיֹּ֣אמֶר אֲלֵהֶ֔ם לְכ֥וּ עִבְד֖וּ אֶת־יְהֹוָ֣ה אֱלֹהֵיכֶ֑ם מִ֥י וָמִ֖י הַהֹלְכִֽים׃ \Rashi{\textbf{ויושב.}הושבו על ידי שליח, ששלחו אחריהם והשיבום אל פרעה:}%
% ---------- פרק 10 | פסוק 9 ----------
 \textbf{ט} וַיֹּ֣אמֶר מֹשֶׁ֔ה בִּנְעָרֵ֥ינוּ וּבִזְקֵנֵ֖ינוּ נֵלֵ֑ךְ בְּבָנֵ֨ינוּ וּבִבְנוֹתֵ֜נוּ בְּצֹאנֵ֤נוּ וּבִבְקָרֵ֙נוּ֙ נֵלֵ֔ךְ כִּ֥י חַג־יְהֹוָ֖ה לָֽנוּ׃ %
% ---------- פרק 10 | פסוק 10 ----------
 \textbf{י} וַיֹּ֣אמֶר אֲלֵהֶ֗ם יְהִ֨י כֵ֤ן יְהֹוָה֙ עִמָּכֶ֔ם כַּאֲשֶׁ֛ר אֲשַׁלַּ֥ח אֶתְכֶ֖ם וְאֶֽת־טַפְּכֶ֑ם רְא֕וּ כִּ֥י רָעָ֖ה נֶ֥גֶד פְּנֵיכֶֽם׃ \Rashi{\textbf{כאשר אשלח אתכם ואת טפכם.}אף כי אשלח גם את הצאן ואת הבקר כאשר אמרתם: ראו כי רעה נגד פניכם. כתרגומו. ומ"א שמעתי, כוכב אחד יש ששמו רעה, אמר להם פרעה, רואה אני באצטגנינות שלי אותו כוכב עולה לקראתכם במדבר, והוא סימן דם והריגה; וכשחטאו ישראל בעגל ובקש הקב"ה להרגם, אמר משה בתפלתו, למה יאמרו מצרים לאמר "ברעה" הוציאם (שמות ל"ב), זו היא שאמר להם, ראו כי רעה נגד פניכם; מיד וינחם ה' על הרעה והפך את הדם לדם מילה, שמל יהושע אותם, וזהו שנ' "היום גלותי את חרפת מצרים מעליכם" (יהושע ה'), שהיו אומרים לכם, דם אנו רואין עליכם במדבר:}%
% ---------- פרק 10 | פסוק 11 ----------
 \textbf{יא} לֹ֣א כֵ֗ן לְכֽוּ־נָ֤א הַגְּבָרִים֙ וְעִבְד֣וּ אֶת־יְהֹוָ֔ה כִּ֥י אֹתָ֖הּ אַתֶּ֣ם מְבַקְשִׁ֑ים וַיְגָ֣רֶשׁ אֹתָ֔ם מֵאֵ֖ת פְּנֵ֥י פַרְעֹֽה׃ \{ס\} \Rashi{\textbf{לא כן.}כאשר אמרתם להוליך הטף עמכם, אלא לכו הגברים ועבדו את ה': כי אתה אתם מבקשים. כי אותה בקשתם עד הנה – נזבחה לאלהינו, ואין דרך הטף לזבח: ויגרש אותם. הרי זה לשון קצר ולא פרש מי המגרש:}%
% ---------- פרק 10 | פסוק 12 ----------
 \textbf{יב} וַיֹּ֨אמֶר יְהֹוָ֜ה אֶל־מֹשֶׁ֗ה נְטֵ֨ה יָדְךָ֜ עַל־אֶ֤רֶץ מִצְרַ֙יִם֙ בָּֽאַרְבֶּ֔ה וְיַ֖עַל עַל־אֶ֣רֶץ מִצְרָ֑יִם וְיֹאכַל֙ אֶת־כׇּל־עֵ֣שֶׂב הָאָ֔רֶץ אֵ֛ת כׇּל־אֲשֶׁ֥ר הִשְׁאִ֖יר הַבָּרָֽד׃ \Rashi{\textbf{בארבה.}בשביל מכת הארבה:}%
% ---------- פרק 10 | פסוק 13 ----------
 \textbf{יג} וַיֵּ֨ט מֹשֶׁ֣ה אֶת־מַטֵּ֘הוּ֮ עַל־אֶ֣רֶץ מִצְרַ֒יִם֒ וַֽיהֹוָ֗ה נִהַ֤ג רֽוּחַ־קָדִים֙ בָּאָ֔רֶץ כׇּל־הַיּ֥וֹם הַה֖וּא וְכׇל־הַלָּ֑יְלָה הַבֹּ֣קֶר הָיָ֔ה וְר֙וּחַ֙ הַקָּדִ֔ים נָשָׂ֖א אֶת־הָאַרְבֶּֽה׃ \Rashi{\textbf{ורוח הקדים.}רוח מזרחית נשא את הארבה, לפי שבא כנגדו, שמצרים בדרומית מערבית היתה, כמו שמפרש במקום אחר:}%
% ---------- פרק 10 | פסוק 14 ----------
 \textbf{יד} וַיַּ֣עַל הָֽאַרְבֶּ֗ה עַ֚ל כׇּל־אֶ֣רֶץ מִצְרַ֔יִם וַיָּ֕נַח בְּכֹ֖ל גְּב֣וּל מִצְרָ֑יִם כָּבֵ֣ד מְאֹ֔ד לְ֠פָנָ֠יו לֹא־הָ֨יָה כֵ֤ן אַרְבֶּה֙ כָּמֹ֔הוּ וְאַחֲרָ֖יו לֹ֥א יִֽהְיֶה־כֵּֽן׃ \Rashi{\textbf{ואחריו לא יהיה כן.}ואותו שהיה בימי יואל שנ' "כמוהו לא נהיה מן העולם" (יואל ב'), למדנו שהיה כבד משל משה. (אותו של יואל היה) ע"י מינין הרבה, שהיו יחד ארבה, ילק, חסיל, גזם, אבל של משה לא היה אלא מין אחד, וכמוהו לא היה ולא יהיה:}%
% ---------- פרק 10 | פסוק 15 ----------
 \textbf{טו} וַיְכַ֞ס אֶת־עֵ֣ין כׇּל־הָאָ֘רֶץ֮ וַתֶּחְשַׁ֣ךְ הָאָ֒רֶץ֒ וַיֹּ֜אכַל אֶת־כׇּל־עֵ֣שֶׂב הָאָ֗רֶץ וְאֵת֙ כׇּל־פְּרִ֣י הָעֵ֔ץ אֲשֶׁ֥ר הוֹתִ֖יר הַבָּרָ֑ד וְלֹא־נוֹתַ֨ר כׇּל־יֶ֧רֶק בָּעֵ֛ץ וּבְעֵ֥שֶׂב הַשָּׂדֶ֖ה בְּכׇל־אֶ֥רֶץ מִצְרָֽיִם׃ \Rashi{\textbf{כל ירק.}עלה ירק, וירדו"רא בלעז:}%
% ---------- פרק 10 | פסוק 16 ----------
 \textbf{טז} וַיְמַהֵ֣ר פַּרְעֹ֔ה לִקְרֹ֖א לְמֹשֶׁ֣ה וּֽלְאַהֲרֹ֑ן וַיֹּ֗אמֶר חָטָ֛אתִי לַיהֹוָ֥ה אֱלֹֽהֵיכֶ֖ם וְלָכֶֽם׃ %
% ---------- פרק 10 | פסוק 17 ----------
 \textbf{יז} וְעַתָּ֗ה שָׂ֣א נָ֤א חַטָּאתִי֙ אַ֣ךְ הַפַּ֔עַם וְהַעְתִּ֖ירוּ לַיהֹוָ֣ה אֱלֹהֵיכֶ֑ם וְיָסֵר֙ מֵֽעָלַ֔י רַ֖ק אֶת־הַמָּ֥וֶת הַזֶּֽה׃ %
% ---------- פרק 10 | פסוק 18 ----------
 \textbf{יח} וַיֵּצֵ֖א מֵעִ֣ם פַּרְעֹ֑ה וַיֶּעְתַּ֖ר אֶל־יְהֹוָֽה׃ %
% ---------- פרק 10 | פסוק 19 ----------
 \textbf{יט} וַיַּהֲפֹ֨ךְ יְהֹוָ֤ה רֽוּחַ־יָם֙ חָזָ֣ק מְאֹ֔ד וַיִּשָּׂא֙ אֶת־הָ֣אַרְבֶּ֔ה וַיִּתְקָעֵ֖הוּ יָ֣מָּה סּ֑וּף לֹ֤א נִשְׁאַר֙ אַרְבֶּ֣ה אֶחָ֔ד בְּכֹ֖ל גְּב֥וּל מִצְרָֽיִם׃ \Rashi{\textbf{רוח ים.}רוח מערבי: ימה סוף. אומר אני שים סוף היה מקצתו במערב, כנגד כל רוח דרומית, וגם במזרחה של ארץ ישראל, לפיכך רוח ים תקעו לארבה בימה סוף כנגדו. וכן מצינו לענין תחומין שהוא פונה לצד מזרח, שנאמר "מים סוף ועד ים פלשתים" (לקמן כג לא) ממזרח למערב, שים פלשתים במערב היה, שנאמר בפלשתים "יושבי חבל הים גוי כרתים" (צפניה ב'): לא נשאר ארבה אחד. אף המלוחים שמלחו מהם:}%
% ---------- פרק 10 | פסוק 20 ----------
 \textbf{כ} וַיְחַזֵּ֥ק יְהֹוָ֖ה אֶת־לֵ֣ב פַּרְעֹ֑ה וְלֹ֥א שִׁלַּ֖ח אֶת־בְּנֵ֥י יִשְׂרָאֵֽל׃ \{פ\} %
% ---------- פרק 10 | פסוק 21 ----------
 \textbf{כא} וַיֹּ֨אמֶר יְהֹוָ֜ה אֶל־מֹשֶׁ֗ה נְטֵ֤ה יָֽדְךָ֙ עַל־הַשָּׁמַ֔יִם וִ֥יהִי חֹ֖שֶׁךְ עַל־אֶ֣רֶץ מִצְרָ֑יִם וְיָמֵ֖שׁ חֹֽשֶׁךְ׃ \Rashi{\textbf{וימש חשך.}ויחשיך עליהם חשך יותר מחשכו של לילה, חשך של לילה יאמיש ויחשיך עוד: וימש. כמו ויאמש; יש לנו תבות הרבה חסרות אל"ף, לפי שאין הברת האל"ף נכרת כל כך, אין הכתוב מקפיד על חסרונה, כגון "ולא יהל שם ערבי" (ישעיהו י"ג), כמו לא יאהל – לא יטה אהלו, וכן "ותזרני חיל" (שמואל ב' כ"ב), כמו ותאזרני; ואנקלוס תרגם לשון הסרה, כמו לא ימוש; "בתר דיעדי קבל ליליא" – כשיגיע סמוך לאור היום; אבל אין הדבור מישב על הוי"ו של וימש, לפי שהוא כתוב אחר ויהי חשך. ומ"א פותרו לשון "ממשש בצהרים" (דברים כ"ח), שהיה כפול ומכפל ועב עד שהיה בו ממש:}%
% ---------- פרק 10 | פסוק 22 ----------
 \textbf{כב} וַיֵּ֥ט מֹשֶׁ֛ה אֶת־יָד֖וֹ עַל־הַשָּׁמָ֑יִם וַיְהִ֧י חֹֽשֶׁךְ־אֲפֵלָ֛ה בְּכׇל־אֶ֥רֶץ מִצְרַ֖יִם שְׁלֹ֥שֶׁת יָמִֽים׃ \Rashi{\textbf{שלשת ימים.}שלוש של ימים, טרציינ"א בלעז, וכן שבעת ימים בכל מקום שטיי"נא של ימים: ויהי חשך אפלה: שלשת ימים. חשך של אפל, שלא ראו איש את אחיו אותן ג' ימים. ועוד שלושת ימים אחרים חשך מכפל על זה, שלא קמו איש מתחתיו – יושב אין יכול לעמד, ועומד אין יכול לישב; ולמה הביא עליהם חשך? שהיו בישראל באותו הדור רשעים ולא היו רוצים לצאת, ומתו בשלושת ימי אפלה, כדי שלא יראו מצריים במפלתם ויאמרו, אף הם לוקים כמונו. ועוד, שחפשו ישראל וראו את כליהם, וכשיצאו והיו שואלים מהן והיו אומרים אין בידנו כלום, אומר לו, אני ראיתיו בביתך, ובמקום פלוני הוא (שמות רבה):}%
% ---------- פרק 10 | פסוק 23 ----------
 \textbf{כג} לֹֽא־רָא֞וּ אִ֣ישׁ אֶת־אָחִ֗יו וְלֹא־קָ֛מוּ אִ֥ישׁ מִתַּחְתָּ֖יו שְׁלֹ֣שֶׁת יָמִ֑ים וּֽלְכׇל־בְּנֵ֧י יִשְׂרָאֵ֛ל הָ֥יָה א֖וֹר בְּמוֹשְׁבֹתָֽם׃ %
% ---------- פרק 10 | פסוק 24 ----------
 \textbf{כד} וַיִּקְרָ֨א פַרְעֹ֜ה אֶל־מֹשֶׁ֗ה וַיֹּ֙אמֶר֙ לְכוּ֙ עִבְד֣וּ אֶת־יְהֹוָ֔ה רַ֛ק צֹאנְכֶ֥ם וּבְקַרְכֶ֖ם יֻצָּ֑ג גַּֽם־טַפְּכֶ֖ם יֵלֵ֥ךְ עִמָּכֶֽם׃ \Rashi{\textbf{יצג.}יהא מצג במקומו:}%
% ---------- פרק 10 | פסוק 25 ----------
 \textbf{כה} וַיֹּ֣אמֶר מֹשֶׁ֔ה גַּם־אַתָּ֛ה תִּתֵּ֥ן בְּיָדֵ֖נוּ זְבָחִ֣ים וְעֹלֹ֑ת וְעָשִׂ֖ינוּ לַיהֹוָ֥ה אֱלֹהֵֽינוּ׃ \Rashi{\textbf{גם אתה תתן.}לא דיך שמקננו ילך עמנו אלא גם משלך תתן:}\Peirush{נאמר "ויחגו לי במדבר" (ה, א), יכולנו לחשוב חג בלי קרבן, קמ"ל שלא, שנאמר גם אתה תתן בידנו זבחים ועלת\Makor{חגיגה י:}}\Peirush{אסור לקחת קרבנות מהגוים, אך קשה, שנאמר גם אתה תתן בידנו זבחים ועלת\Makor{עבודה זרה כד.}}\Peirush{בני נח לא הקריבו שלמים, אך קשה שהרי עוד לפני מתן תורה הקריבו זבחים – שלמים\Makor{זבחים קטז.}}%
% ---------- פרק 10 | פסוק 26 ----------
 \textbf{כו} וְגַם־מִקְנֵ֜נוּ יֵלֵ֣ךְ עִמָּ֗נוּ לֹ֤א תִשָּׁאֵר֙ פַּרְסָ֔ה כִּ֚י מִמֶּ֣נּוּ נִקַּ֔ח לַעֲבֹ֖ד אֶת־יְהֹוָ֣ה אֱלֹהֵ֑ינוּ וַאֲנַ֣חְנוּ לֹֽא־נֵדַ֗ע מַֽה־נַּעֲבֹד֙ אֶת־יְהֹוָ֔ה עַד־בֹּאֵ֖נוּ שָֽׁמָּה׃ \Rashi{\textbf{פרסה.}פרסת רגל, פלנ"טא בלעז: לא נדע מה נעבד. כמה תכבד העבודה, שמא ישאל יותר ממה שיש בידנו:}%
% ---------- פרק 10 | פסוק 27 ----------
 \textbf{כז} וַיְחַזֵּ֥ק יְהֹוָ֖ה אֶת־לֵ֣ב פַּרְעֹ֑ה וְלֹ֥א אָבָ֖ה לְשַׁלְּחָֽם׃ %
% ---------- פרק 10 | פסוק 28 ----------
 \textbf{כח} וַיֹּֽאמֶר־ל֥וֹ פַרְעֹ֖ה לֵ֣ךְ מֵעָלָ֑י הִשָּׁ֣מֶר לְךָ֗ אַל־תֹּ֙סֶף֙ רְא֣וֹת פָּנַ֔י כִּ֗י בְּי֛וֹם רְאֹתְךָ֥ פָנַ֖י תָּמֽוּת׃ %
% ---------- פרק 10 | פסוק 29 ----------
 \textbf{כט} וַיֹּ֥אמֶר מֹשֶׁ֖ה כֵּ֣ן דִּבַּ֑רְתָּ לֹא־אֹסִ֥ף ע֖וֹד רְא֥וֹת פָּנֶֽיךָ׃ \{פ\} \Rashi{\textbf{כן דברת.}יפה דברת ובזמנו דברת. אמת שלא אוסיף עוד ראות פניך (מכילתא ושמות רבה):}%
% ---------- פרק 11 | פסוק 1 ----------
 \textbf{א} וַיֹּ֨אמֶר יְהֹוָ֜ה אֶל־מֹשֶׁ֗ה ע֣וֹד נֶ֤גַע אֶחָד֙ אָבִ֤יא עַל־פַּרְעֹה֙ וְעַל־מִצְרַ֔יִם אַֽחֲרֵי־כֵ֕ן יְשַׁלַּ֥ח אֶתְכֶ֖ם מִזֶּ֑ה כְּשַׁ֨לְּח֔וֹ כָּלָ֕ה גָּרֵ֛שׁ יְגָרֵ֥שׁ אֶתְכֶ֖ם מִזֶּֽה׃ \Rashi{\textbf{כלה.}גמירא; כליל, כלכם ישלח:}%
% ---------- פרק 11 | פסוק 2 ----------
 \textbf{ב} דַּבֶּר־נָ֖א בְּאׇזְנֵ֣י הָעָ֑ם וְיִשְׁאֲל֞וּ אִ֣ישׁ ׀ מֵאֵ֣ת רֵעֵ֗הוּ וְאִשָּׁה֙ מֵאֵ֣ת רְעוּתָ֔הּ כְּלֵי־כֶ֖סֶף וּכְלֵ֥י זָהָֽב׃ \Rashi{\textbf{דבר נא.}אין נא אלא לשון בקשה, בבקשה ממך הזהירם על כך, שלא יאמר אותו צדיק אברהם "ועבדום וענו אותם" (בראשית ט"ו), קים בהם, "ואחרי כן יצאו ברכש גדול" (שם), לא קים בהם:}\Peirush{אין נא אלא לשון בקשה. אמר ה' למשה בבקשה ממך לך ואמור להם לישראל בבקשה מכם שאלו ממצרים כלי כסף וכלי זהב, שלא יאמר אברהם "ועבדום וענו אותם" (בראשית טו, יג) קיים בהם, "ואחרי כן יצאו ברכוש גדול" (שם) לא קיים בהם. אמרו לו הלואי שנצא בעצנו. משל לאדם שהיה חבוש בבית האסורים, והיו אומרים לו בני אדם מוציאין אותך למחר מבית האסורין ונותנין לך ממון הרבה, ואומר להם: בבקשה מכם הוציאוני היום ואיני מבקש כלום\Makor{ברכות ט.-ט:}}%
% ---------- פרק 11 | פסוק 3 ----------
 \textbf{ג} וַיִּתֵּ֧ן יְהֹוָ֛ה אֶת־חֵ֥ן הָעָ֖ם בְּעֵינֵ֣י מִצְרָ֑יִם גַּ֣ם ׀ הָאִ֣ישׁ מֹשֶׁ֗ה גָּד֤וֹל מְאֹד֙ בְּאֶ֣רֶץ מִצְרַ֔יִם בְּעֵינֵ֥י עַבְדֵֽי־פַרְעֹ֖ה וּבְעֵינֵ֥י הָעָֽם׃ \{ס\} %
% ---------- פרק 11 | פסוק 4 ----------
 \textbf{ד} וַיֹּ֣אמֶר מֹשֶׁ֔ה כֹּ֖ה אָמַ֣ר יְהֹוָ֑ה כַּחֲצֹ֣ת הַלַּ֔יְלָה אֲנִ֥י יוֹצֵ֖א בְּת֥וֹךְ מִצְרָֽיִם׃ \Rashi{\textbf{ויאמר משה כה אמר ה'.}בעמדו לפני פרעה נאמרה לו נבואה זו, שהרי משיצא מלפניו לא הוסיף ראות פניו: כחצת הלילה. כהחלק הלילה: כחצת. כמו "בעלות" (מלכים א י"ח), "בחרות אפם בנו" (תהילים קכ"ד), זהו פשוטו לישבו על אפניו, שאין חצות שם דבר של חצי. ורבותינו דרשוהו כמו כחצי הלילה, ואמרו שאמר משה "כחצות", שמשמע סמוך לו, או לפניו או לאחריו, ולא אמר "בחצות", שמא יטעו אצטגניני פרעה ויאמרו משה בדאי הוא (ברכות ד'):}\Peirush{דוד ידע מתי בדיוק חצות הלילה, אך קשה, שאפילו משה לא ידע, שנאמר כחצות\Makor{ברכות ג:}}\Peirush{משה ידע מתי חצות, הוא אמר כחצות שמא יטעו אצטגניני פרעה בחישוב חצות הלילה ויאמרו משה בדאי הוא\Makor{ברכות ד.}}%
% ---------- פרק 11 | פסוק 5 ----------
 \textbf{ה} וּמֵ֣ת כׇּל־בְּכוֹר֮ בְּאֶ֣רֶץ מִצְרַ֒יִם֒ מִבְּכ֤וֹר פַּרְעֹה֙ הַיֹּשֵׁ֣ב עַל־כִּסְא֔וֹ עַ֚ד בְּכ֣וֹר הַשִּׁפְחָ֔ה אֲשֶׁ֖ר אַחַ֣ר הָרֵחָ֑יִם וְכֹ֖ל בְּכ֥וֹר בְּהֵמָֽה׃ \Rashi{\textbf{עד בכור השבי.}למה לקו השבויים? כדי שלא יאמרו, יראתם תבעה עלבונם והביאה פרענות על מצרים: מבכור פרעה עד בכור השפחה. כל הפחותים מבכור פרעה וחשובים מבכור השפחה היו בכלל; ולמה לקו בני השפחות? שאף הם היו משעבדים בהם ושמחים בצרתם: וכל בכור בהמה. לפי שהיו עובדין לה. כשהקב"ה נפרע מן האמה, נפרע מאלהיה:}%
% ---------- פרק 11 | פסוק 6 ----------
 \textbf{ו} וְהָ֥יְתָ֛ה צְעָקָ֥ה גְדֹלָ֖ה בְּכׇל־אֶ֣רֶץ מִצְרָ֑יִם אֲשֶׁ֤ר כָּמֹ֙הוּ֙ לֹ֣א נִהְיָ֔תָה וְכָמֹ֖הוּ לֹ֥א תֹסִֽף׃ %
% ---------- פרק 11 | פסוק 7 ----------
 \textbf{ז} וּלְכֹ֣ל ׀ בְּנֵ֣י יִשְׂרָאֵ֗ל לֹ֤א יֶֽחֱרַץ־כֶּ֙לֶב֙ לְשֹׁנ֔וֹ לְמֵאִ֖ישׁ וְעַד־בְּהֵמָ֑ה לְמַ֙עַן֙ תֵּֽדְע֔וּן אֲשֶׁר֙ יַפְלֶ֣ה יְהֹוָ֔ה בֵּ֥ין מִצְרַ֖יִם וּבֵ֥ין יִשְׂרָאֵֽל׃ \Rashi{\textbf{לא יחרץ כלב לשונו.}אומר אני שהוא לשון שנון – לא ישנן, וכן "לא חרץ לבני ישראל לאיש את לשונו" (יהושע י') – לא שנן, "אז תחרץ" (שמואל ב ה׳:כ״ד) – תשתנן, "למורג חרוץ" (ישעיהו מ"א) – שנון, "מחשבות חרוץ" (משלי כ"א) – אדם חריף ושנון, "ויד חרוצים תעשיר" (שם י') – חריפים, סוחרים שנונים: אשר יפלה. יבדיל:}\Peirush{הרואה כלב בחלום, ישכים ויאמר ולכל בני ישראל לא יחרץ כלב לשנו לפני שיתגשם הפסוק "והכלבים עזי נפש" (ישעיהו נו, יא)\Makor{ברכות נו:}}%
% ---------- פרק 11 | פסוק 8 ----------
 \textbf{ח} וְיָרְד֣וּ כׇל־עֲבָדֶ֩יךָ֩ אֵ֨לֶּה אֵלַ֜י וְהִשְׁתַּֽחֲווּ־לִ֣י לֵאמֹ֗ר צֵ֤א אַתָּה֙ וְכׇל־הָעָ֣ם אֲשֶׁר־בְּרַגְלֶ֔יךָ וְאַחֲרֵי־כֵ֖ן אֵצֵ֑א וַיֵּצֵ֥א מֵֽעִם־פַּרְעֹ֖ה בׇּחֳרִי־אָֽף׃ \{ס\} \Rashi{\textbf{וירדו כל עבדיך.}חלק כבוד למלכות, שהרי בסוף ירד פרעה בעצמו אליו בלילה "ויאמר קומו צאו מתוך עמי" (שמות י"ב), ולא אמר לו משה מתחלה וירדת אלי והשתחוית לי: אשר ברגליך. ההולכים אחר עצתך והלוכך: ואחרי כן אצא. עם כל העם מארצך: ויצא מעם פרעה. כשגמר דבריו, יצא מלפניו: בחרי אף. על שאמר לו "אל תסף ראות פני" (שמות י'):}\Peirush{לעולם תהא אימת מלכות עליך שמשה אמר לפרעה וירדו כל עבדיך אך לא אתה\Makor{זבחים קב.}}\Peirush{כל מקום שכתוב בתורה חרי אף, כתוב גם איך זה מתבטא, אך קשה שהרי כאן זה לא התבטא, מכאן שמשה סטר לפרעה ויצא\Makor{זבחים קב., מנחות צח.}}%
% ---------- פרק 11 | פסוק 9 ----------
 \textbf{ט} וַיֹּ֤אמֶר יְהֹוָה֙ אֶל־מֹשֶׁ֔ה לֹא־יִשְׁמַ֥ע אֲלֵיכֶ֖ם פַּרְעֹ֑ה לְמַ֛עַן רְב֥וֹת מוֹפְתַ֖י בְּאֶ֥רֶץ מִצְרָֽיִם׃ \Rashi{\textbf{למען רבות מופתי.}מכת בכורות וקריעת ים סוף ולנער את מצרים:}%
% ---------- פרק 11 | פסוק 10 ----------
 \textbf{י} וּמֹשֶׁ֣ה וְאַהֲרֹ֗ן עָשׂ֛וּ אֶת־כׇּל־הַמֹּפְתִ֥ים הָאֵ֖לֶּה לִפְנֵ֣י פַרְעֹ֑ה וַיְחַזֵּ֤ק יְהֹוָה֙ אֶת־לֵ֣ב פַּרְעֹ֔ה וְלֹֽא־שִׁלַּ֥ח אֶת־בְּנֵֽי־יִשְׂרָאֵ֖ל מֵאַרְצֽוֹ׃ \{ס\} \Rashi{\textbf{משה ואהרן עשו וגו'.}כבר כתב לנו זאת בכל המופתים, ולא שנאה כאן אלא בשביל לסמכה לפרשה שלאחריה:}%
% ---------- פרק 12 | פסוק 1 ----------
 \textbf{א} וַיֹּ֤אמֶר יְהֹוָה֙ אֶל־מֹשֶׁ֣ה וְאֶֽל־אַהֲרֹ֔ן בְּאֶ֥רֶץ מִצְרַ֖יִם לֵאמֹֽר׃ \Rashi{\textbf{ויאמר ה' אל משה ואל אהרן.}בשביל שאהרן עשה וטרח במופתים כמשה, חלק לו כבוד זה במצוה ראשונה, שכללו עם משה בדבור: בארץ מצרים. חוץ לכרך, או אינו אלא בתוך הכרך? ת"ל (לעיל ט כט) "כצאתי את העיר וגו'", ומה תפלה קלה לא התפלל בתוך הכרך, דבור חמור לא כל שכן. ומפני מה לא נדבר עמו בתוך הכרך? לפי שהיתה מלאה גלולים:}\Peirush{עדות החודש מותרת בקרובים, שנאמר ויאמר ה' אל משה ואל אהרן... החדש הזה לכם שתהא כשרה במשה ואהרן שהיו אחים\Makor{ראש השנה כב.}}\Peirush{חייבים שלושה דיינים לעדות החודש, שנאמר ויאמר ה' אל משה ואל אהרן... החדש הזה לכם, וכיון שאין עושים בית דין זוגי, מוסיפים עוד אחד\Makor{ראש השנה כה:}}%
% ---------- פרק 12 | פסוק 2 ----------
 \textbf{ב} הַחֹ֧דֶשׁ הַזֶּ֛ה לָכֶ֖ם רֹ֣אשׁ חֳדָשִׁ֑ים רִאשׁ֥וֹן הוּא֙ לָכֶ֔ם לְחׇדְשֵׁ֖י הַשָּׁנָֽה׃ \Rashi{\textbf{החדש הזה.}הראהו לבנה בחדושה ואמר לו כשהירח מתחדש יהיה לך ראש חדש (מכילתא). ואין מקרא יוצא מידי פשוטו, על חדש ניסן אמר לו, זה יהיה ראש לסדר מנין החדשים, שיהא איר קרוי שני, סיון שלישי: הזה. נתקשה משה על מולד הלבנה, באיזו שעור תראה ותהיה ראויה לקדש, והראה לו באצבע את הלבנה ברקיע ואמר לו כזה ראה וקדש (שם). וכיצד הראהו? והלא לא היה נדבר עמו אלא ביום? שנ' "ויהי ביום דבר ה'" (שמות ו'), "ביום צותו" (ויקרא ז'), "מן היום אשר צוה ה' והלאה" (במדבר ט"ו)? אלא סמוך לשקיעת החמה נאמרה לו פרשה זו והראהו עם חשכה:}\Peirush{כל המברך על החדש בזמנו כאילו מקבל פני שכינה, אצלנו כתוב הזה ונאמר על קבלת השכינה בים סוף "זה אלי ואנוהו" (טו, ב)\Makor{סנהדרין מב.}}\Peirush{שלשה דברים היו קשים למשה להבינם בדיוק, עד שהראה לו הקדוש ברוך הוא באצבעו. מעשה המנורה, מולד הירח שנאמר החדש הזה, וזיהוי השרצים האסורים\Makor{מנחות כט.}}\Peirush{חזקיהו עיבר את השנה עם ניסן במקום אדר, אך קשה, הרי כתוב החדש הזה, זה ניסן ואין ניסן אחר\Makor{ברכות י:, סנהדרין יב:}}\Peirush{עם ישראל הגיעו להר סיני בראש חדש סיון, נאמר אצלנו החדש הזה ונאמר "ביום הזה באו מדבר סיני" (יט, א)\Makor{שבת פו:}}\Peirush{ר' אלעזר בן ערך הגיע לפריגיה ודיומסת [שלא למדו תורה], ושכח תלמודו. כשחזר למקומו, בא לקרוא החדש הזה לכם ושכח תלמודו עד כדי כך, שקרא "החרש היה לבם". התפללו עליו חכמים ונזכר\Makor{שבת קמז:}}\Peirush{דורשים בהלכות החג שבועיים לפני החג, שמשה מזהיר על הפסח בראש חדש, נאמר ראש חדשים ונאמר "ויקחו להם איש שה לבית אבת" (ג)\Makor{פסחים ו:}}\Peirush{ניסן זהו החדש הראשון, נאמר בפסח ראש חדשים ראשון הוא לכם לחדשי השנה... בעשור לחדש הזה ויקחו להם, וגם "שמור את חדש האביב" (דברים טז, א) באיזה חדש יש אביב? ניסן\Makor{ראש השנה ז.}}\Peirush{אסור לעבר החדש לצרכי ציבור, שנאמר החדש הזה, כזה ראה וקדש\Makor{ראש השנה כ.}}\Peirush{ראה פס' א'\Makor{ראש השנה כב.}}\Peirush{ראה פס' א'\Makor{ראש השנה כה:}}\Peirush{בשבת הרביעית בחדש אדר קוראים את החדש הזה לכם\Makor{מגילה כט., ל.}}\Peirush{מתחילים לספור שמיטות מתשרי, מגזרה שווה "שנה" (ויקרא כה, ד) "שנה" (דברים יא, יב), אך קשה, למה לא ללמוד מניסן שנאמר בו שנה\Makor{ראש השנה ח:, ט.}}\Peirush{מתחילים לתרום מחצית השקל למקדש בניסן מגזרה שווה "שנה" (במדבר כח, יד) שנה\Makor{ראש השנה ז.}}%
% ---------- פרק 12 | פסוק 3 ----------
 \textbf{ג} דַּבְּר֗וּ אֶֽל־כׇּל־עֲדַ֤ת יִשְׂרָאֵל֙ לֵאמֹ֔ר בֶּעָשֹׂ֖ר לַחֹ֣דֶשׁ הַזֶּ֑ה וְיִקְח֣וּ לָהֶ֗ם אִ֛ישׁ שֶׂ֥ה לְבֵית־אָבֹ֖ת שֶׂ֥ה לַבָּֽיִת׃ \Rashi{\textbf{דברו אל כל עדת.}וכי אהרן מדבר? והלא כבר נאמר אתה תדבר? אלא חולקין כבוד זה לזה, ואומרים זה לזה למדני, והדבור יוצא מבין שניהם כאלו שניהם מדברים (מכילתא): דברו אל כל עדת ישראל לאמר בעשר לחדש. דברו היום, בראש חדש שיקחוהו בעשר לחדש: הזה. פסח מצרים מקחו בעשור, ולא פסח דורות (פסחים צ"ו): שה לבית אבת. למשפחה אחת; הרי שהיו מרבין, יכול שה אחד לכלן? ת"ל "שה לבית" (מכילתא):}\Peirush{בניגוד לפסח מצרים, בפסח לדורות לא צריך לקחת את קרבן הפסח בי' ניסן שנאמר בעשור לחודש הזה ויקחו, הזה ולא אחר\Makor{פסחים צו.}}\Peirush{ראה פס' ב'\Makor{פסחים ו:}}\Peirush{ראה פס' ב'\Makor{ראש השנה ז.}}\Peirush{נאמר ויקחו להם איש שה לבית אבת מכאן שאיש יכול לבצע קנין, ולא קטן\Makor{פסחים צא:}}\Peirush{נאמר לבית אבת, מכאן שניתן למנות שליח לקדשים\Makor{קידושים מב.}}\Peirush{כמו שבקרבן פסח רק שה כשר, כך גם בפטר חמור\Makor{בכורות יב.}}%
% ---------- פרק 12 | פסוק 4 ----------
 \textbf{ד} וְאִם־יִמְעַ֣ט הַבַּ֘יִת֮ מִהְי֣וֹת מִשֶּׂה֒ וְלָקַ֣ח ה֗וּא וּשְׁכֵנ֛וֹ הַקָּרֹ֥ב אֶל־בֵּית֖וֹ בְּמִכְסַ֣ת נְפָשֹׁ֑ת אִ֚ישׁ לְפִ֣י אׇכְל֔וֹ תָּכֹ֖סּוּ עַל־הַשֶּֽׂה׃ \Rashi{\textbf{ואם ימעט הבית מהיות משה.}ואם יהיו מועטים מהיות משה אחד, שאין יכולין לאכלו ויבא לידי נותר: ולקח הוא ושכנו וגו'. זהו משמעו לפי פשוטו. ועוד יש בו מדרש: ללמד, שאחר שנמנו עליו, יכולין להתמעט ולמשך ידיהם הימנו ולהמנות על שה אחר; אך אם באו למשך ידיהם ולהתמעט, מהיות משה, יתמעטו בעוד השה קים, בהיותו בחיים ולא משנשחט (פסחים פ"ט): במכסת. חשבון, וכן "מכסת הערכך" (ויקרא כ"ז): לפי אכלו. הראוי לאכילה, פרט לחולה ולזקן שאינו יכול לאכל כזית (מכילתא): תכסו. "תתמנון":}%
% ---------- פרק 12 | פסוק 5 ----------
 \textbf{ה} שֶׂ֥ה תָמִ֛ים זָכָ֥ר בֶּן־שָׁנָ֖ה יִהְיֶ֣ה לָכֶ֑ם מִן־הַכְּבָשִׂ֥ים וּמִן־הָעִזִּ֖ים תִּקָּֽחוּ׃ \Rashi{\textbf{תמים.}בלא מום: בן שנה. כל שנתו קרוי בן שנה, כלומר שנולד בשנה זו: מן הכבשים ומן העזים. או מזה או מזה, שאף עז קרוי שה, שנאמר "ושה עזים" (דברים י"ד):}%
% ---------- פרק 12 | פסוק 6 ----------
 \textbf{ו} וְהָיָ֤ה לָכֶם֙ לְמִשְׁמֶ֔רֶת עַ֣ד אַרְבָּעָ֥ה עָשָׂ֛ר י֖וֹם לַחֹ֣דֶשׁ הַזֶּ֑ה וְשָׁחֲט֣וּ אֹת֗וֹ כֹּ֛ל קְהַ֥ל עֲדַֽת־יִשְׂרָאֵ֖ל בֵּ֥ין הָעַרְבָּֽיִם׃ \Rashi{\textbf{והיה לכם למשמרת.}זהו לשון בקור, שטעון בקור ממום ארבעה ימים קדם שחיטה (פסחים צ"ו). ומפני מה הקדים לקיחתו לשחיטתו ארבעה ימים, מה שלא צוה כן בפסח דורות? היה ר' מתיא בן חרש אומר, הרי הוא אומר "ואעבר עליך ואראך והנה עתך עת דודים" (יחזקאל ט"ז) – הגיעה שבועה שנשבעתי לאברהם שאגאל את בניו, ולא היו בידם מצוות להתעסק בהם כדי שיגאלו, שנאמר "ואת ערום ועריה" (שם), ונתן להם שתי מצוות, דם פסח ודם מילה, שמלו באותו הלילה, שנאמר "מתבוססת בדמיך" (שם) – בשני דמים, ואומר "גם את בדם בריתך שלחתי אסיריך מבור אין מים בו" (זכריה ט'); ולפי שהיו שטופין באלילים אמר להם משכו וקחו לכם, משכו ידיכם מאלילים, וקחו לכם צאן של מצוה (מכילתא): ושחטו אתו וגו'. וכי כלן שוחטין? אלא מכאן ששלוחו של אדם כמותו (קידושין מ"א): קהל עדת ישראל. קהל ועדה וישראל; מכאן אמרו, פסחי צבור נשחטין בשלוש כתות זו אחר זו, נכנסה כת ראשונה ננעלו דלתות העזרה וכו'. כדאיתא בפסחים (דף ס"ד): בין הערבים. משש שעות ולמעלה קרוי בין הערבים, שהשמש נוטה לבית מבואו לערב; ולשון בין הערבים נראה בעיני אותן שעות שבין עריבת היום לעריבת הלילה, עריבת היום בתחלת ז' שעות מכי ינטו צללי ערב, ועריבת הלילה בתחלת הלילה. ערב לשון נשף וחשך, כמו "ערבה כל שמחה" (ישעיהו כ"ד):}%
% ---------- פרק 12 | פסוק 7 ----------
 \textbf{ז} וְלָֽקְחוּ֙ מִן־הַדָּ֔ם וְנָ֥תְנ֛וּ עַל־שְׁתֵּ֥י הַמְּזוּזֹ֖ת וְעַל־הַמַּשְׁק֑וֹף עַ֚ל הַבָּ֣תִּ֔ים אֲשֶׁר־יֹאכְל֥וּ אֹת֖וֹ בָּהֶֽם׃ \Rashi{\textbf{ולקחו מן הדם.}זו קבלת הדם, יכול ביד? ת"ל אשר בסף (מכילתא): המזוזת. הם הזקופות אחת מכאן לפתח ואחת מכאן: המשקוף. העליון שהדלת שוקף עליו כשסוגרין אותו, לינטי"ל בלעז; ולשון שקיפה חבטה, כמו "קול עלה נדף" (ויקרא כ"ו) – דשקיף, חבורה – משקופי: על הבתים אשר יאכלו אתו בהם. ולא על משקוף ומזוזות שבבית התבן ובית הבקר, שאין דרין בתוכו (ע' מכילתא):}%
% ---------- פרק 12 | פסוק 8 ----------
 \textbf{ח} וְאָכְל֥וּ אֶת־הַבָּשָׂ֖ר בַּלַּ֣יְלָה הַזֶּ֑ה צְלִי־אֵ֣שׁ וּמַצּ֔וֹת עַל־מְרֹרִ֖ים יֹאכְלֻֽהוּ׃ \Rashi{\textbf{את הבשר.}ולא גידים ועצמות (שם): ומצות על מררים. כל עשב מר נקרא מרור; וצום לאכל מר זכר ל"וימררו את חייהם" (שמות א'):}%
% ---------- פרק 12 | פסוק 9 ----------
 \textbf{ט} אַל־תֹּאכְל֤וּ מִמֶּ֙נּוּ֙ נָ֔א וּבָשֵׁ֥ל מְבֻשָּׁ֖ל בַּמָּ֑יִם כִּ֣י אִם־צְלִי־אֵ֔שׁ רֹאשׁ֥וֹ עַל־כְּרָעָ֖יו וְעַל־קִרְבּֽוֹ׃ \Rashi{\textbf{אל תאכלו ממנו נא.}שאינו צלוי כל צרכו קוראו נא בלשון ערבי: ובשל מבשל. כל זה באזהרת אל תאכלו: במים. מנין לשאר משקין? ת"ל ובשל מבשל, מכל מקום (פסחים מ"א): כי אם צלי אש. למעלה גזר עליו במצות עשה, וכאן הוסיף עליו לא תעשה, אל תאכלו ממנו … כי אם צלי אש: ראשו על כרעיו. צולהו כלו כאחד עם ראשו ועם כרעיו ועם קרבו, ובני מעיו נותן לתוכו אחר הדחתן (שם ע"ד); ולשון על קרבו כלשון "על צבאותם" (שמות ו'), כמו בצבאותם – כמות שהן, אף זה כמות שהוא – כל בשרו משלם:}%
% ---------- פרק 12 | פסוק 10 ----------
 \textbf{י} וְלֹא־תוֹתִ֥ירוּ מִמֶּ֖נּוּ עַד־בֹּ֑קֶר וְהַנֹּתָ֥ר מִמֶּ֛נּוּ עַד־בֹּ֖קֶר בָּאֵ֥שׁ תִּשְׂרֹֽפוּ׃ \Rashi{\textbf{והנתר ממנו עד בקר.}מה ת"ל עד בקר פעם שניה? לתן בקר על בקר, שהבקר משמעו משעת הנץ החמה, ובא הכתוב להקדים שאסור באכילה מעלות השחר, זהו לפי משמעו. ועוד מדרש אחר, למד שאינו נשרף ביום טוב אלא ממחרת, וכך תדרשנו: והנותר ממנו בבקר ראשון, עד בקר שני תעמד ותשרפנו (מכילתא):}%
% ---------- פרק 12 | פסוק 11 ----------
 \textbf{יא} וְכָ֘כָה֮ תֹּאכְל֣וּ אֹתוֹ֒ מׇתְנֵיכֶ֣ם חֲגֻרִ֔ים נַֽעֲלֵיכֶם֙ בְּרַגְלֵיכֶ֔ם וּמַקֶּלְכֶ֖ם בְּיֶדְכֶ֑ם וַאֲכַלְתֶּ֤ם אֹתוֹ֙ בְּחִפָּז֔וֹן פֶּ֥סַח ה֖וּא לַיהֹוָֽה׃ \Rashi{\textbf{מתניכם חגרים.}מזמנים לדרך: בחפזון. לשון בהלה ומהירות, כמו "ויהי דוד נחפז ללכת" (שמואל א' כ"ג), "אשר השליכו ארם בחפזם" (מלכים ב' ז'): פסח הוא לה'. הקרבן קרוי פסח על שם הדלוג והפסיחה, שהקב"ה היה מדלג בתי ישראל מבין בתי מצרים וקופץ ממצרי למצרי וישראל אמצעי נמלט; ואתם עשו כל עבודותיו לשם שמים, דרך דלוג וקפיצה, זכר לשמו שקרוי פסח; וגם פשק"א לשון פסיעה:}%
% ---------- פרק 12 | פסוק 12 ----------
 \textbf{יב} וְעָבַרְתִּ֣י בְאֶֽרֶץ־מִצְרַ֘יִם֮ בַּלַּ֣יְלָה הַזֶּה֒ וְהִכֵּיתִ֤י כׇל־בְּכוֹר֙ בְּאֶ֣רֶץ מִצְרַ֔יִם מֵאָדָ֖ם וְעַד־בְּהֵמָ֑ה וּבְכׇל־אֱלֹהֵ֥י מִצְרַ֛יִם אֶֽעֱשֶׂ֥ה שְׁפָטִ֖ים אֲנִ֥י יְהֹוָֽה׃ \Rashi{\textbf{ועברתי.}כמלך העובר ממקום למקום (מכילתא), ובהעברה אחת וברגע אחד כלן לוקים: כל בכור בארץ מצרים. אף בכורות אחרים והם במצרים. ומנין אף בכורי מצרים שבמקומות אחרים? ת"ל "למכה מצרים בבכוריהם" (תהלים קל"ו): מאדם ועד בהמה. מי שהתחיל בעברה ממנו מתחלת הפרענות (מכילתא): ובכל אלהי מצרים. של עץ נרקבת ושל מתכת נמסת ונתכת לארץ (שם): אעשה שפטים אני ה'. אני בעצמי, ולא על ידי שליח:}%
% ---------- פרק 12 | פסוק 13 ----------
 \textbf{יג} וְהָיָה֩ הַדָּ֨ם לָכֶ֜ם לְאֹ֗ת עַ֤ל הַבָּתִּים֙ אֲשֶׁ֣ר אַתֶּ֣ם שָׁ֔ם וְרָאִ֙יתִי֙ אֶת־הַדָּ֔ם וּפָסַחְתִּ֖י עֲלֵכֶ֑ם וְלֹֽא־יִֽהְיֶ֨ה בָכֶ֥ם נֶ֙גֶף֙ לְמַשְׁחִ֔ית בְּהַכֹּתִ֖י בְּאֶ֥רֶץ מִצְרָֽיִם׃ \Rashi{\textbf{והיה הדם לכם לאת.}לכם לאות ולא לאחרים לאות (שם). מכאן שלא נתנו הדם אלא מבפנים: וראיתי את הדם. הכל גלוי לפניו, אלא אמר הקב"ה, נותן אני את עיני לראות שאתם עסוקים במצוותי, ופוסח אני עליכם (שם): ופסחתי. וחמלתי, ודומה לו "פסח והמליט" (ישעיהו ל"א). ואני אומר, כל פסיחה לשון דלוג וקפיצה. ופסחתי, מדלג היה מבתי ישראל לבתי מצריים, שהיו שרויים זה בתוך זה, וכן "פוסחים על שתי הסעפים" (מלכים א י"ח), וכן כל הפסחים – הולכים כקופצים, וכן "פסח והמליט" – מדלגו וממלטו מבין המומתים: ולא יהיה בכם נגף. אבל הוה הוא במצרים. הרי שהיה מצרי בביתו של ישראל, יכול ימלט? ת"ל "ולא יהיה בכם נגף", אבל הוה במצרים שבבתיכם. הרי שהיה ישראל בביתו של מצרי, שומע אני ילקה כמותו, ת"ל "ולא יהיה בכם נגף" (מכילתא):}%
% ---------- פרק 12 | פסוק 14 ----------
 \textbf{יד} וְהָיָה֩ הַיּ֨וֹם הַזֶּ֤ה לָכֶם֙ לְזִכָּר֔וֹן וְחַגֹּתֶ֥ם אֹת֖וֹ חַ֣ג לַֽיהֹוָ֑ה לְדֹרֹ֣תֵיכֶ֔ם חֻקַּ֥ת עוֹלָ֖ם תְּחׇגֻּֽהוּ׃ \Rashi{\textbf{לזכרון.}לדורות: וחגתם אתו. יום שהוא לך לזכרון אתה חוגגו. ועדין לא שמענו איזהו יום הזכרון, ת"ל "זכור את היום הזה אשר יצאתם" (שמות י"ג), למדנו, שיום היציאה הוא יום של זכרון. ואי זה יום יצאו? ת"ל "ממחרת הפסח יצאו" (במדבר ל"ג), הוי אומר יום חמשה עשר בניסן הוא של יום טוב, שהרי ליל חמשה עשר אכלו את הפסח ולבקר יצאו (מכילתא): לדרתיכם וגו'. שומע אני מעוט דורות שנים, ת"ל "חקת עולם תחגהו" (שם):}%
% ---------- פרק 12 | פסוק 15 ----------
 \textbf{טו} שִׁבְעַ֤ת יָמִים֙ מַצּ֣וֹת תֹּאכֵ֔לוּ אַ֚ךְ בַּיּ֣וֹם הָרִאשׁ֔וֹן תַּשְׁבִּ֥יתוּ שְּׂאֹ֖ר מִבָּתֵּיכֶ֑ם כִּ֣י ׀ כׇּל־אֹכֵ֣ל חָמֵ֗ץ וְנִכְרְתָ֞ה הַנֶּ֤פֶשׁ הַהִוא֙ מִיִּשְׂרָאֵ֔ל מִיּ֥וֹם הָרִאשֹׁ֖ן עַד־י֥וֹם הַשְּׁבִעִֽי׃ \Rashi{\textbf{שבעת ימים.}שטיינ"א של ימים: שבעת ימים מצות תאכלו. ובמקום אחר הוא אומר "ששת ימים תאכל מצות" (דברים ט"ז) למד על שביעי של פסח שאינו חובה לאכל מצה, ובלבד שלא יאכל חמץ; מנין אף ששה רשות? ת"ל "ששת ימים". זו מדה בתורה, דבר שהיה בכלל ויצא מן הכלל ללמד, לא ללמד על עצמו בלבד יצא, אלא ללמד על הכלל כלו יצא, מה שביעי רשות אף ששה רשות; יכול אף הלילה הראשון רשות, ת"ל "בערב תאכלו מצת" – הכתוב קבעו חובה (פסחים ק"כ): אך ביום הראשון תשביתו שאור. מערב יום טוב, וקרוי ראשון לפי שהוא לפני השבעה; ומצינו מקדם קרוי ראשון, "הראישון אדם תולד" (איוב ט"ו) – הלפני אדם נולדת; או אינו אלא ראשון של שבעה? ת"ל "לא תשחט על חמץ וגו'" (שמות ל"ד) – לא תשחט הפסח ועדין חמץ קים (פסחים ה'): הנפש ההוא. כשהיא בנפשה ובדעתה – פרט לאנוס (מכילתא): מישראל. שומע אני תכרת מישראל ותלך לה לעם אחר, ת"ל במקום אחר "מלפני" (ויקרא כ"ב) – בכל מקום שהוא רשותי:}%
% ---------- פרק 12 | פסוק 16 ----------
 \textbf{טז} וּבַיּ֤וֹם הָרִאשׁוֹן֙ מִקְרָא־קֹ֔דֶשׁ וּבַיּוֹם֙ הַשְּׁבִיעִ֔י מִקְרָא־קֹ֖דֶשׁ יִהְיֶ֣ה לָכֶ֑ם כׇּל־מְלָאכָה֙ לֹא־יֵעָשֶׂ֣ה בָהֶ֔ם אַ֚ךְ אֲשֶׁ֣ר יֵאָכֵ֣ל לְכׇל־נֶ֔פֶשׁ ה֥וּא לְבַדּ֖וֹ יֵעָשֶׂ֥ה לָכֶֽם׃ \Rashi{\textbf{מקרא קדש.}מקרא שם דבר; קרא אותו קדש לאכילה ושתיה וכסות (מכילתא): לא יעשה בהם. אפלו על ידי אחרים (שם): הוא לבדו. "הוא" ולא מכשיריו שאפשר לעשותן מערב יום טוב (ביצה כ"ח): לכל נפש. אפלו לבהמה; יכול אף לגוים, תלמוד לומר לכם (מכילתא):}%
% ---------- פרק 12 | פסוק 17 ----------
 \textbf{יז} וּשְׁמַרְתֶּם֮ אֶת־הַמַּצּוֹת֒ כִּ֗י בְּעֶ֙צֶם֙ הַיּ֣וֹם הַזֶּ֔ה הוֹצֵ֥אתִי אֶת־צִבְאוֹתֵיכֶ֖ם מֵאֶ֣רֶץ מִצְרָ֑יִם וּשְׁמַרְתֶּ֞ם אֶת־הַיּ֥וֹם הַזֶּ֛ה לְדֹרֹתֵיכֶ֖ם חֻקַּ֥ת עוֹלָֽם׃ \Rashi{\textbf{ושמרתם את המצות.}שלא יבאו לידי חמוץ; מכאן אמרו תפח תלטש בצונן, רבי יאשיה אומר אל תהי קורא את המצות, אלא את המצוות – כדרך שאין מחמיצין את המצה, כך אין מחמיצין את המצוה, אלא אם באה לידך, עשה אותה מיד (שם): ושמרתם את היום הזה. ממלאכה: לדרתיכם חקת עולם. לפי שלא נאמר דורות וחקת עולם על המלאכה אלא על החגיגה, לכך חזר ושנאו כאן, שלא תאמר אזהרת כל מלאכה לא יעשה לא לדורות נאמרה אלא לאותו הדור:}%
% ---------- פרק 12 | פסוק 18 ----------
 \textbf{יח} בָּרִאשֹׁ֡ן בְּאַרְבָּעָה֩ עָשָׂ֨ר י֤וֹם לַחֹ֙דֶשׁ֙ בָּעֶ֔רֶב תֹּאכְל֖וּ מַצֹּ֑ת עַ֠ד י֣וֹם הָאֶחָ֧ד וְעֶשְׂרִ֛ים לַחֹ֖דֶשׁ בָּעָֽרֶב׃ \Rashi{\textbf{עד יום האחד ועשרים.}למה נאמר? והלא כבר נאמר "שבעת ימים"? לפי שנאמר ימים, לילות מנין? ת"ל עד יום האחד ועשרים וגו' (מכילתא):}%
% ---------- פרק 12 | פסוק 19 ----------
 \textbf{יט} שִׁבְעַ֣ת יָמִ֔ים שְׂאֹ֕ר לֹ֥א יִמָּצֵ֖א בְּבָתֵּיכֶ֑ם כִּ֣י ׀ כׇּל־אֹכֵ֣ל מַחְמֶ֗צֶת וְנִכְרְתָ֞ה הַנֶּ֤פֶשׁ הַהִוא֙ מֵעֲדַ֣ת יִשְׂרָאֵ֔ל בַּגֵּ֖ר וּבְאֶזְרַ֥ח הָאָֽרֶץ׃ \Rashi{\textbf{לא ימצא בבתיכם.}מנין לגבולין? תלמוד לומר "בכל גבולך" (שמות י"ג); מה ת"ל בבתיכם? מה ביתך ברשותך, אף גבולך שברשותך, יצא חמצו של נכרי שהוא אצל ישראל ולא קבל עליו אחריות (מכילתא): כי כל אכל מחמצת. לענש כרת על השאור, והלא כבר ענש על החמץ? אלא שלא תאמר חמץ שראוי לאכילה ענש עליו, שאור שאינו ראוי לאכילה לא יענש עליו; ואם ענש על השאור ולא ענש על החמץ, הייתי אומר, שאור שהוא מחמץ אחרים ענש עליו, חמץ שאינו מחמץ אחרים לא יענש עליו, לכך נאמרו שניהם (מכילתא): בגר ובאזרח הארץ. לפי שהנס נעשה לישראל, הצרך לרבות את הגרים (שם):}%
% ---------- פרק 12 | פסוק 20 ----------
 \textbf{כ} כׇּל־מַחְמֶ֖צֶת לֹ֣א תֹאכֵ֑לוּ בְּכֹל֙ מוֹשְׁבֹ֣תֵיכֶ֔ם תֹּאכְל֖וּ מַצּֽוֹת׃ \{פ\} \Rashi{\textbf{מחמצת לא תאכלו.}אזהרה על אכילת שאור: כל מחמצת. להביא את תערבתו: בכל מושבתיכם תאכלו מצות. זה בא ללמד שתהא ראויה להאכל בכל מושבתיכם, פרט למעשר שני וחלות תודה (מכילתא):}%
% ---------- פרק 12 | פסוק 21 ----------
 \textbf{כא} וַיִּקְרָ֥א מֹשֶׁ֛ה לְכׇל־זִקְנֵ֥י יִשְׂרָאֵ֖ל וַיֹּ֣אמֶר אֲלֵהֶ֑ם מִֽשְׁכ֗וּ וּקְח֨וּ לָכֶ֥ם צֹ֛אן לְמִשְׁפְּחֹתֵיכֶ֖ם וְשַׁחֲט֥וּ הַפָּֽסַח׃ \Rashi{\textbf{משכו.}מי שיש לו צאן ימשך משלו: וקחו. מי שאין לו יקח מן השוק (שם): למשפחתיכם. שה לבית אבות:}%
% ---------- פרק 12 | פסוק 22 ----------
 \textbf{כב} וּלְקַחְתֶּ֞ם אֲגֻדַּ֣ת אֵז֗וֹב וּטְבַלְתֶּם֮ בַּדָּ֣ם אֲשֶׁר־בַּסַּף֒ וְהִגַּעְתֶּ֤ם אֶל־הַמַּשְׁקוֹף֙ וְאֶל־שְׁתֵּ֣י הַמְּזוּזֹ֔ת מִן־הַדָּ֖ם אֲשֶׁ֣ר בַּסָּ֑ף וְאַתֶּ֗ם לֹ֥א תֵצְא֛וּ אִ֥ישׁ מִפֶּֽתַח־בֵּית֖וֹ עַד־בֹּֽקֶר׃ \Rashi{\textbf{אזוב.}מין ירק שיש לו גבעולין: אגדת אזוב. שלושה קלחין קרויין אגדה (שבת ק"ט): אשר בסף. בכלי, כמו ספות כסף: מן הדם אשר בסף. למה חזר ושנאו? שלא תאמר טבילה אחת לשלוש המתנות, לכך נאמר עוד אשר בסף, ושתהא כל נתינה ונתינה מן הדם אשר בסף, על כל הגעה טבילה (מכילתא): ואתם לא תצאו וגו'. מגיד, שמאחר שנתנה רשות למשחית לחבל, אינו מבחין בין צדיק לרשע (שם); ולילה רשות למחבלים הוא, שנאמר "בו תרמש כל חיתו יער" (תהלים ק"ד):}%
% ---------- פרק 12 | פסוק 23 ----------
 \textbf{כג} וְעָבַ֣ר יְהֹוָה֮ לִנְגֹּ֣ף אֶת־מִצְרַ֒יִם֒ וְרָאָ֤ה אֶת־הַדָּם֙ עַל־הַמַּשְׁק֔וֹף וְעַ֖ל שְׁתֵּ֣י הַמְּזוּזֹ֑ת וּפָסַ֤ח יְהֹוָה֙ עַל־הַפֶּ֔תַח וְלֹ֤א יִתֵּן֙ הַמַּשְׁחִ֔ית לָבֹ֥א אֶל־בָּתֵּיכֶ֖ם לִנְגֹּֽף׃ \Rashi{\textbf{ופסח.}וחמל, ויש לומר ודלג: ולא יתן המשחית. ולא יתן לו יכלת לבא, כמו "ולא נתנו אלהים להרע עמדי" (בראשית ל"א):}%
% ---------- פרק 12 | פסוק 24 ----------
 \textbf{כד} וּשְׁמַרְתֶּ֖ם אֶת־הַדָּבָ֣ר הַזֶּ֑ה לְחׇק־לְךָ֥ וּלְבָנֶ֖יךָ עַד־עוֹלָֽם׃ %
% ---------- פרק 12 | פסוק 25 ----------
 \textbf{כה} וְהָיָ֞ה כִּֽי־תָבֹ֣אוּ אֶל־הָאָ֗רֶץ אֲשֶׁ֨ר יִתֵּ֧ן יְהֹוָ֛ה לָכֶ֖ם כַּאֲשֶׁ֣ר דִּבֵּ֑ר וּשְׁמַרְתֶּ֖ם אֶת־הָעֲבֹדָ֥ה הַזֹּֽאת׃ \Rashi{\textbf{והיה כי תבאו אל הארץ.}תלה הכתוב עבודה זו בביאתם לארץ, ולא נתחיבו במדבר אלא פסח אחד שעשו בשנה השנית על פי הדבור: כאשר דבר. והיכן דבר? "והבאתי אתכם אל הארץ וגו'" (שמות ו'):}%
% ---------- פרק 12 | פסוק 26 ----------
 \textbf{כו} וְהָיָ֕ה כִּֽי־יֹאמְר֥וּ אֲלֵיכֶ֖ם בְּנֵיכֶ֑ם מָ֛ה הָעֲבֹדָ֥ה הַזֹּ֖את לָכֶֽם׃ %
% ---------- פרק 12 | פסוק 27 ----------
 \textbf{כז} וַאֲמַרְתֶּ֡ם זֶֽבַח־פֶּ֨סַח ה֜וּא לַֽיהֹוָ֗ה אֲשֶׁ֣ר פָּ֠סַ֠ח עַל־בָּתֵּ֤י בְנֵֽי־יִשְׂרָאֵל֙ בְּמִצְרַ֔יִם בְּנׇגְפּ֥וֹ אֶת־מִצְרַ֖יִם וְאֶת־בָּתֵּ֣ינוּ הִצִּ֑יל וַיִּקֹּ֥ד הָעָ֖ם וַיִּֽשְׁתַּחֲוֽוּ׃ \Rashi{\textbf{ויקד העם.}על בשורת הגאלה וביאת הארץ ובשורת הבנים שיהיו להם:}%
% ---------- פרק 12 | פסוק 28 ----------
 \textbf{כח} וַיֵּלְכ֥וּ וַֽיַּעֲשׂ֖וּ בְּנֵ֣י יִשְׂרָאֵ֑ל כַּאֲשֶׁ֨ר צִוָּ֧ה יְהֹוָ֛ה אֶת־מֹשֶׁ֥ה וְאַהֲרֹ֖ן כֵּ֥ן עָשֽׂוּ׃ \{ס\} \Rashi{\textbf{וילכו ויעשו בני ישראל.}וכי כבר עשו, והלא מראש חדש נאמר להם? אלא מכיון שקבלו עליהם מעלה עליהם הכתוב כאלו עשו (מכילתא): וילכו ויעשו. אף ההליכה מנה הכתוב, לתן שכר להליכה ושכר לעשיה (שם): כאשר צוה ה' את משה ואהרן. להגיד שבחן של ישראל שלא הפילו דבר מכל מצוות משה ואהרן, ומהו כן עשו? אף משה ואהרן כן עשו (שם):}%
% ---------- פרק 12 | פסוק 29 ----------
 \textbf{כט} וַיְהִ֣י ׀ בַּחֲצִ֣י הַלַּ֗יְלָה וַֽיהֹוָה֮ הִכָּ֣ה כׇל־בְּכוֹר֮ בְּאֶ֣רֶץ מִצְרַ֒יִם֒ מִבְּכֹ֤ר פַּרְעֹה֙ הַיֹּשֵׁ֣ב עַל־כִּסְא֔וֹ עַ֚ד בְּכ֣וֹר הַשְּׁבִ֔י אֲשֶׁ֖ר בְּבֵ֣ית הַבּ֑וֹר וְכֹ֖ל בְּכ֥וֹר בְּהֵמָֽה׃ \Rashi{\textbf{וה'.}כל מקום שנאמר "וה'" הוא ובית דינו, שהוי"ו לשון תוספת הוא, כמו פלוני ופלוני: הכה כל בכור. אף של אמה אחרת והוא במצרים (מכילתא): מבכר פרעה. אף פרעה בכור היה, ונשתיר מן הבכורות, ועליו הוא אומר "בעבור הראותך את כחי" (שמות ט') – בים סוף (מכילתא): עד בכור השבי. שהיו שמחין לאידם של ישראל; ועוד, שלא יאמרו יראתנו הביאה הפרענות; ובכור השפחה בכלל היה, שהרי מנה מן החשוב שבכלן עד הפחות, ובכור השפחה חשוב מבכור השבי:}%
% ---------- פרק 12 | פסוק 30 ----------
 \textbf{ל} וַיָּ֨קׇם פַּרְעֹ֜ה לַ֗יְלָה ה֤וּא וְכׇל־עֲבָדָיו֙ וְכׇל־מִצְרַ֔יִם וַתְּהִ֛י צְעָקָ֥ה גְדֹלָ֖ה בְּמִצְרָ֑יִם כִּֽי־אֵ֣ין בַּ֔יִת אֲשֶׁ֥ר אֵֽין־שָׁ֖ם מֵֽת׃ \Rashi{\textbf{ויקם פרעה.}ממטתו: לילה. ולא כדרך המלכים בשלוש שעות ביום (שם): הוא. תחלה, ואחר כך עבדיו, מלמד שהיה הוא מחזר על בתי עבדיו ומעמידן (שם): כי אין בית אשר אין שם מת. יש שם בכור, מת, אין שם בכור, גדול שבבית קרוי בכור, שנאמר "אף אני בכור אתנהו" (תהילים פ"ט). דבר אחר, מצריות מזנות תחת בעליהן ויולדות מרוקים פנויים, והיו להם בכורות הרבה, פעמים חמשה לאשה אחת, כל אחד בכור לאביו:}%
% ---------- פרק 12 | פסוק 31 ----------
 \textbf{לא} וַיִּקְרָא֩ לְמֹשֶׁ֨ה וּֽלְאַהֲרֹ֜ן לַ֗יְלָה וַיֹּ֙אמֶר֙ ק֤וּמוּ צְּאוּ֙ מִתּ֣וֹךְ עַמִּ֔י גַּם־אַתֶּ֖ם גַּם־בְּנֵ֣י יִשְׂרָאֵ֑ל וּלְכ֛וּ עִבְד֥וּ אֶת־יְהֹוָ֖ה כְּדַבֶּרְכֶֽם׃ \Rashi{\textbf{ויקרא למשה ולאהרן לילה.}מגיד שהיה מחזר על פתחי העיר וצועק, היכן משה שרוי? היכן אהרן שרוי? (מכילתא): גם אתם. הגברים: גם בני ישראל. הטף: ולכו עבדו את ה' כדברכם. הכל כמו שאמרתם ולא כשאמרתי אני; בטל "לא אשלח", בטל "מי ומי ההולכים", בטל "רק צאנכם ובקרכם יצג":}%
% ---------- פרק 12 | פסוק 32 ----------
 \textbf{לב} גַּם־צֹאנְכֶ֨ם גַּם־בְּקַרְכֶ֥ם קְח֛וּ כַּאֲשֶׁ֥ר דִּבַּרְתֶּ֖ם וָלֵ֑כוּ וּבֵֽרַכְתֶּ֖ם גַּם־אֹתִֽי׃ \Rashi{\textbf{גם צאנכם גם בקרכם קחו.}ומהו כאשר דברתם? "גם אתה תתן בידנו זבחים ועלת" (שמות י'), קחו כאשר דברתם: וברכתם גם אתי. התפללו עלי שלא אמות, שאני בכור:}%
% ---------- פרק 12 | פסוק 33 ----------
 \textbf{לג} וַתֶּחֱזַ֤ק מִצְרַ֙יִם֙ עַל־הָעָ֔ם לְמַהֵ֖ר לְשַׁלְּחָ֣ם מִן־הָאָ֑רֶץ כִּ֥י אָמְר֖וּ כֻּלָּ֥נוּ מֵתִֽים׃ \Rashi{\textbf{כלנו מתים.}אמרו לא כגזרת משה הוא, שהרי אמר "ומת כל בכור", וכאן אף הפשוטים מתים, חמשה או עשרה בבית אחד:}%
% ---------- פרק 12 | פסוק 34 ----------
 \textbf{לד} וַיִּשָּׂ֥א הָעָ֛ם אֶת־בְּצֵק֖וֹ טֶ֣רֶם יֶחְמָ֑ץ מִשְׁאֲרֹתָ֛ם צְרֻרֹ֥ת בְּשִׂמְלֹתָ֖ם עַל־שִׁכְמָֽם׃ \Rashi{\textbf{טרם יחמץ.}המצרים לא הניחום לשהות כדי חמוץ: משארתם. שירי מצה ומרור. על שכמם. אע"פ שבהמות הרבה הוליכו עמהם, מחבבים היו את המצוה:}%
% ---------- פרק 12 | פסוק 35 ----------
 \textbf{לה} וּבְנֵי־יִשְׂרָאֵ֥ל עָשׂ֖וּ כִּדְבַ֣ר מֹשֶׁ֑ה וַֽיִּשְׁאֲלוּ֙ מִמִּצְרַ֔יִם כְּלֵי־כֶ֛סֶף וּכְלֵ֥י זָהָ֖ב וּשְׂמָלֹֽת׃ \Rashi{\textbf{כדבר משה.}שאמר להם במצרים, "וישאלו איש מאת רעהו": ושמלת. אף הן היו חשובות להם מן הכסף ומן הזהב, והמאחר בפסוק חשוב:}%
% ---------- פרק 12 | פסוק 36 ----------
 \textbf{לו} וַֽיהֹוָ֞ה נָתַ֨ן אֶת־חֵ֥ן הָעָ֛ם בְּעֵינֵ֥י מִצְרַ֖יִם וַיַּשְׁאִל֑וּם וַֽיְנַצְּל֖וּ אֶת־מִצְרָֽיִם׃ \{פ\} \Rashi{\textbf{וישאלום.}אף מה שלא היו שואלים מהם היו נותנים להם, אתה אומר אחד? טל שנים ולך. וינצלו. ורוקינו:}%
% ---------- פרק 12 | פסוק 37 ----------
 \textbf{לז} וַיִּסְע֧וּ בְנֵֽי־יִשְׂרָאֵ֛ל מֵרַעְמְסֵ֖ס סֻכֹּ֑תָה כְּשֵׁשׁ־מֵא֨וֹת אֶ֧לֶף רַגְלִ֛י הַגְּבָרִ֖ים לְבַ֥ד מִטָּֽף׃ \Rashi{\textbf{מרעמסס סכתה.}מאה ועשרים מיל היו, ובאו שם לפי שעה, שנאמר "ואשא אתכם על כנפי נשרים" (שמות י"ט): הגברים. מבן עשרים שנה ומעלה (שה"ש רבה ג, ו):}%
% ---------- פרק 12 | פסוק 38 ----------
 \textbf{לח} וְגַם־עֵ֥רֶב רַ֖ב עָלָ֣ה אִתָּ֑ם וְצֹ֣אן וּבָקָ֔ר מִקְנֶ֖ה כָּבֵ֥ד מְאֹֽד׃ \Rashi{\textbf{ערב רב.}תערבת אמות של גרים:}%
% ---------- פרק 12 | פסוק 39 ----------
 \textbf{לט} וַיֹּאפ֨וּ אֶת־הַבָּצֵ֜ק אֲשֶׁ֨ר הוֹצִ֧יאוּ מִמִּצְרַ֛יִם עֻגֹ֥ת מַצּ֖וֹת כִּ֣י לֹ֣א חָמֵ֑ץ כִּֽי־גֹרְשׁ֣וּ מִמִּצְרַ֗יִם וְלֹ֤א יָֽכְלוּ֙ לְהִתְמַהְמֵ֔הַּ וְגַם־צֵדָ֖ה לֹא־עָשׂ֥וּ לָהֶֽם׃ \Rashi{\textbf{עגת מצות.}חררה של מצה; בצק שלא החמיץ קרוי מצה: וגם צדה לא עשו להם. לדרך. מגיד שבחן של ישראל, שלא אמרו היאך נצא למדבר בלא צידה? אלא האמינו והלכו; הוא שמפרש בקבלה, "זכרתי לך חסד נעוריך אהבת כלולתיך לכתך אחרי במדבר בארץ לא זרועה" (ירמיהו ב'), מה שכר מפרש אחריו? "קדש ישראל לה' וגו'":}%
% ---------- פרק 12 | פסוק 40 ----------
 \textbf{מ} וּמוֹשַׁב֙ בְּנֵ֣י יִשְׂרָאֵ֔ל אֲשֶׁ֥ר יָשְׁב֖וּ בְּמִצְרָ֑יִם שְׁלֹשִׁ֣ים שָׁנָ֔ה וְאַרְבַּ֥ע מֵא֖וֹת שָׁנָֽה׃ \Rashi{\textbf{אשר ישבו במצרים.}אחר שאר הישיבות שישבו גרים בארץ לא להם: שלשים שנה וארבע מאות שנה. בין הכל, משנולד יצחק עד עכשו, היו ארבע מאות, משהיה לו זרע לאברהם נתקים כי גר יהיה זרעך, ושלשים שנה היו משנגזרה גזרת בין הבתרים עד שנולד יצחק; ואי אפשר לומר בארץ מצרים לבדה, שהרי קהת מן הבאים עם יעקב היה, צא וחשב כל שנותיו וכל שנות עמרם בנו ושמונים של משה, לא תמצאם כל כך, ועל כרחך הרבה שנים היו לקהת עד שלא ירד למצרים, והרבה משנות עמרם נבלעים בשנות קהת, והרבה משמונים של משה נבלעים בשנות עמרם, הרי שלא תמצא ארבע מאות לביאת מצרים, והזקקתה לומר על כרחך, שאף שאר הישיבות נקראו גרות ואפלו בחברון שנאמר "אשר גר שם אברהם ויצחק" (בראשית ל"ה), ואומר "את ארץ מגריהם אשר גרו בה" (שמות ו'), לפיכך אתה צריך לומר "כי גר יהיה זרעך" משהיה לו זרע, וכשתמנה ת' שנה משנולד יצחק, תמצא מביאתן למצרים עד יציאתן ר"י, וזה אחד מן הדברים ששנו לתלמי המלך (מגילה ט'):}%
% ---------- פרק 12 | פסוק 41 ----------
 \textbf{מא} וַיְהִ֗י מִקֵּץ֙ שְׁלֹשִׁ֣ים שָׁנָ֔ה וְאַרְבַּ֥ע מֵא֖וֹת שָׁנָ֑ה וַיְהִ֗י בְּעֶ֙צֶם֙ הַיּ֣וֹם הַזֶּ֔ה יָ֥צְא֛וּ כׇּל־צִבְא֥וֹת יְהֹוָ֖ה מֵאֶ֥רֶץ מִצְרָֽיִם׃ \Rashi{\textbf{ויהי מקץ שלשים שנה וגו' ויהי בעצם היום הזה.}מגיד, שכיון שהגיע הקץ לא עכבן המקום כהרף עין – בט"ו בניסן באו מלאכי השרת אצל אברהם לבשרו, ובט"ו בניסן נולד יצחק, ובט"ו בניסן נגזרה גזרת בין הבתרים:}%
% ---------- פרק 12 | פסוק 42 ----------
 \textbf{מב} לֵ֣יל שִׁמֻּרִ֥ים הוּא֙ לַֽיהֹוָ֔ה לְהוֹצִיאָ֖ם מֵאֶ֣רֶץ מִצְרָ֑יִם הֽוּא־הַלַּ֤יְלָה הַזֶּה֙ לַֽיהֹוָ֔ה שִׁמֻּרִ֛ים לְכׇל־בְּנֵ֥י יִשְׂרָאֵ֖ל לְדֹרֹתָֽם׃ \{פ\} \Rashi{\textbf{ליל שמרים.}שהיה הקב"ה שומר ומצפה לו, לקים הבטחתו להוציאם מארץ מצרים: הוא הלילה הזה לה'. הוא הלילה שאמר לאברהם, בלילה הזה אני גואל את בניך (מכילתא): שמרים לכל בני ישראל לדרתם. משמר ובא מן המזיקין, כענין שנאמר "ולא יתן המשחית וגו'" (פסחים ק"ט):}%
% ---------- פרק 12 | פסוק 43 ----------
 \textbf{מג} וַיֹּ֤אמֶר יְהֹוָה֙ אֶל־מֹשֶׁ֣ה וְאַהֲרֹ֔ן זֹ֖את חֻקַּ֣ת הַפָּ֑סַח כׇּל־בֶּן־נֵכָ֖ר לֹא־יֹ֥אכַל בּֽוֹ׃ \Rashi{\textbf{זאת חקת הפסח.}בי"ד בניסן נאמרה להם פרשה זו: כל בן נכר. שנתנכרו מעשיו לאביו שבשמים, ואחד הגוי ואחד ישראל משמד במשמע (מכילתא):}%
% ---------- פרק 12 | פסוק 44 ----------
 \textbf{מד} וְכׇל־עֶ֥בֶד אִ֖ישׁ מִקְנַת־כָּ֑סֶף וּמַלְתָּ֣ה אֹת֔וֹ אָ֖ז יֹ֥אכַל בּֽוֹ׃ \Rashi{\textbf{ומלתה אתו אז יאכל בו.}רבו, מגיד שמילת עבדיו מעכבתו מלאכל בפסח, דברי רבי יהושע. רבי אליעזר אומר, אין מילת עבדיו מעכבתו מלאכל בפסח, אם כן מה ת"ל אז יאכל בו? העבד:}%
% ---------- פרק 12 | פסוק 45 ----------
 \textbf{מה} תּוֹשָׁ֥ב וְשָׂכִ֖יר לֹא־יֹ֥אכַל בּֽוֹ׃ \Rashi{\textbf{תושב.}זה גר תושב: ושכיר. זה גוי, ומה ת"ל, והלא ערלים הם, ונאמר וכל ערל לא יאכל בו? אלא כגון ערבי מהול וגבעוני מהול והוא תושב או שכיר (שם):}%
% ---------- פרק 12 | פסוק 46 ----------
 \textbf{מו} בְּבַ֤יִת אֶחָד֙ יֵאָכֵ֔ל לֹא־תוֹצִ֧יא מִן־הַבַּ֛יִת מִן־הַבָּשָׂ֖ר ח֑וּצָה וְעֶ֖צֶם לֹ֥א תִשְׁבְּרוּ־בֽוֹ׃ \Rashi{\textbf{בבית אחד יאכל.}בחבורה אחת, שלא יעשו הנמנין עליו שתי חבורות ויחלקוהו. אתה אומר בחבורה אחת או אינו אלא בבית אחד כמשמעו, וללמד שאם התחילו והיו אוכלים בחצר וירדו גשמים שלא יכנסו לבית? תלמוד לומר, "על הבתים אשר יאכלו אתו בהם", מכאן שהאוכל אוכל בשני מקומות: לא תוציא מן הבית. מן החבורה: ועצם לא תשברו בו. הראוי לאכילה, כגון שיש עליו כזית בשר, יש בו משום שבירת עצם, אין עליו כזית בשר אין בו משום שבירת עצם:}%
% ---------- פרק 12 | פסוק 47 ----------
 \textbf{מז} כׇּל־עֲדַ֥ת יִשְׂרָאֵ֖ל יַעֲשׂ֥וּ אֹתֽוֹ׃ \Rashi{\textbf{כל עדת ישראל יעשו אתו.}למה נאמר, לפי שהוא אומר בפסח מצרים "שה לבית אבות", שנמנו עליו למשפחות, יכול אף פסח דורות כן? ת"ל כל עדת ישראל יעשו אתו:}%
% ---------- פרק 12 | פסוק 48 ----------
 \textbf{מח} וְכִֽי־יָג֨וּר אִתְּךָ֜ גֵּ֗ר וְעָ֣שָׂה פֶ֘סַח֮ לַיהֹוָה֒ הִמּ֧וֹל ל֣וֹ כׇל־זָכָ֗ר וְאָז֙ יִקְרַ֣ב לַעֲשֹׂת֔וֹ וְהָיָ֖ה כְּאֶזְרַ֣ח הָאָ֑רֶץ וְכׇל־עָרֵ֖ל לֹֽא־יֹ֥אכַל בּֽוֹ׃ \Rashi{\textbf{ועשה פסח.}יכול כל המתגיר יעשה פסח מיד, תלמוד לומר והיה כאזרח הארץ, מה אזרח בארבעה עשר אף גר בארבעה עשר (שם): וכל ערל לא יאכל בו. להביא את שמתו אחיו מחמת מילה, שאינו מומר לערלות ואינו למד מ"בן נכר לא יאכל בו" (שם):}%
% ---------- פרק 12 | פסוק 49 ----------
 \textbf{מט} תּוֹרָ֣ה אַחַ֔ת יִהְיֶ֖ה לָֽאֶזְרָ֑ח וְלַגֵּ֖ר הַגָּ֥ר בְּתוֹכְכֶֽם׃ \Rashi{\textbf{תורה אחת וגו'.}להשוות גר לאזרח אף לשאר מצוות שבתורה (שם):}%
% ---------- פרק 12 | פסוק 50 ----------
 \textbf{נ} וַֽיַּעֲשׂ֖וּ כׇּל־בְּנֵ֣י יִשְׂרָאֵ֑ל כַּאֲשֶׁ֨ר צִוָּ֧ה יְהֹוָ֛ה אֶת־מֹשֶׁ֥ה וְאֶֽת־אַהֲרֹ֖ן כֵּ֥ן עָשֽׂוּ׃ \{ס\} %
% ---------- פרק 12 | פסוק 51 ----------
 \textbf{נא} וַיְהִ֕י בְּעֶ֖צֶם הַיּ֣וֹם הַזֶּ֑ה הוֹצִ֨יא יְהֹוָ֜ה אֶת־בְּנֵ֧י יִשְׂרָאֵ֛ל מֵאֶ֥רֶץ מִצְרַ֖יִם עַל־צִבְאֹתָֽם׃ \{פ\} %
% ---------- פרק 13 | פסוק 1 ----------
 \textbf{א} וַיְדַבֵּ֥ר יְהֹוָ֖ה אֶל־מֹשֶׁ֥ה לֵּאמֹֽר׃ %
% ---------- פרק 13 | פסוק 2 ----------
 \textbf{ב} קַדֶּשׁ־לִ֨י כׇל־בְּכ֜וֹר פֶּ֤טֶר כׇּל־רֶ֙חֶם֙ בִּבְנֵ֣י יִשְׂרָאֵ֔ל בָּאָדָ֖ם וּבַבְּהֵמָ֑ה לִ֖י הֽוּא׃ \Rashi{\textbf{פטר כל רחם.}שפתח הרחם תחלה, כמו: "פוטר מים ראשית מדון" (משלי י"ז), וכן "יפטירו בשפה" (תהלים כ"ב) – יפתחו שפתים: לי הוא. לעצמי קניתי, ע"י שהכיתי בכורי מצרים:}%
% ---------- פרק 13 | פסוק 3 ----------
 \textbf{ג} וַיֹּ֨אמֶר מֹשֶׁ֜ה אֶל־הָעָ֗ם זָכ֞וֹר אֶת־הַיּ֤וֹם הַזֶּה֙ אֲשֶׁ֨ר יְצָאתֶ֤ם מִמִּצְרַ֙יִם֙ מִבֵּ֣ית עֲבָדִ֔ים כִּ֚י בְּחֹ֣זֶק יָ֔ד הוֹצִ֧יא יְהֹוָ֛ה אֶתְכֶ֖ם מִזֶּ֑ה וְלֹ֥א יֵאָכֵ֖ל חָמֵֽץ׃ \Rashi{\textbf{זכור את היום הזה.}למד, שמזכירין יציאת מצרים בכל יום:}%
% ---------- פרק 13 | פסוק 4 ----------
 \textbf{ד} הַיּ֖וֹם אַתֶּ֣ם יֹצְאִ֑ים בְּחֹ֖דֶשׁ הָאָבִֽיב׃ \Rashi{\textbf{בחדש האביב.}וכי לא היינו יודעין באיזה חדש? אלא כך אמר להם, ראו חסד שגמלכם, שהוציא אתכם בחדש שהוא כשר לצאת, לא חמה ולא צנה ולא גשמים, וכן הוא אומר, "מוציא אסירים בכושרות" (תהלים ס"ח) – חדש שהוא כשר לצאת (מכילתא):}%
% ---------- פרק 13 | פסוק 5 ----------
 \textbf{ה} וְהָיָ֣ה כִֽי־יְבִיאֲךָ֣ יְהֹוָ֡ה אֶל־אֶ֣רֶץ הַֽ֠כְּנַעֲנִ֠י וְהַחִתִּ֨י וְהָאֱמֹרִ֜י וְהַחִוִּ֣י וְהַיְבוּסִ֗י אֲשֶׁ֨ר נִשְׁבַּ֤ע לַאֲבֹתֶ֙יךָ֙ לָ֣תֶת לָ֔ךְ אֶ֛רֶץ זָבַ֥ת חָלָ֖ב וּדְבָ֑שׁ וְעָבַדְתָּ֛ אֶת־הָעֲבֹדָ֥ה הַזֹּ֖את בַּחֹ֥דֶשׁ הַזֶּֽה׃ \Rashi{\textbf{אל ארץ הכנעני וגו'.}ואע"פ שלא מנה אלא חמשה עממין, כל ז' גוים במשמע, שכלן בכלל כנעני הם ואחת ממשפחות כנען היתה שלא נקרא לה שם אלא כנעני: נשבע לאבתיך. באברהם הוא אומר "ביום ההוא כרת ה' את אברם וגו'" (בראשית ט"ו), וביצחק הוא אומר "גור בארץ הזאת" (שם כ"ו), וביעקב הוא אומר "הארץ אשר אתה שוכב עליה" (שם כ"ח): זבת חלב ודבש. חלב זב מן העזים, והדבש זב מן התמרים ומן התאנים: את העבדה הזאת. של פסח; והלא כבר נאמר למעלה "והיה כי תבאו אל הארץ וגו'", ולמה חזר ושנאה? בשביל דבר שנתחדש בה, בפרשה ראשונה נאמר "והיה כי יאמרו אליכם בניכם מה העבדה הזאת לכם" — בבן רשע הכתוב מדבר שהוציא את עצמו מן הכלל, וכאן "והגדת לבנך" — בבן שאינו יודע לשאל, והכתוב מלמדך, שתפתח לו אתה בדברי אגדה המושכין את הלב (מכילתא):}%
% ---------- פרק 13 | פסוק 6 ----------
 \textbf{ו} שִׁבְעַ֥ת יָמִ֖ים תֹּאכַ֣ל מַצֹּ֑ת וּבַיּוֹם֙ הַשְּׁבִיעִ֔י חַ֖ג לַיהֹוָֽה׃ %
% ---------- פרק 13 | פסוק 7 ----------
 \textbf{ז} מַצּוֹת֙ יֵֽאָכֵ֔ל אֵ֖ת שִׁבְעַ֣ת הַיָּמִ֑ים וְלֹֽא־יֵרָאֶ֨ה לְךָ֜ חָמֵ֗ץ וְלֹֽא־יֵרָאֶ֥ה לְךָ֛ שְׂאֹ֖ר בְּכׇל־גְּבֻלֶֽךָ׃ %
% ---------- פרק 13 | פסוק 8 ----------
 \textbf{ח} וְהִגַּדְתָּ֣ לְבִנְךָ֔ בַּיּ֥וֹם הַה֖וּא לֵאמֹ֑ר בַּעֲב֣וּר זֶ֗ה עָשָׂ֤ה יְהֹוָה֙ לִ֔י בְּצֵאתִ֖י מִמִּצְרָֽיִם׃ \Rashi{\textbf{בעבור זה.}בעבור שאקים מצוותיו, כגון פסח מצה ומרור הללו: עשה ה' לי. רמז תשובה לבן רשע, לומר עשה ה' לי, ולא לך, שאלו היית שם לא היית כדאי לגאל (מכילתא):}%
% ---------- פרק 13 | פסוק 9 ----------
 \textbf{ט} וְהָיָה֩ לְךָ֨ לְא֜וֹת עַל־יָדְךָ֗ וּלְזִכָּרוֹן֙ בֵּ֣ין עֵינֶ֔יךָ לְמַ֗עַן תִּהְיֶ֛ה תּוֹרַ֥ת יְהֹוָ֖ה בְּפִ֑יךָ כִּ֚י בְּיָ֣ד חֲזָקָ֔ה הוֹצִֽאֲךָ֥ יְהֹוָ֖ה מִמִּצְרָֽיִם׃ \Rashi{\textbf{והיה לך לאות.}יציאת מצרים תהיה לך לאות: על ידך ולזכרון בית עיניך. שתכתב הפרשיות הללו ותקשרם בראש ובזרוע: על ידך. על יד שמאל, לפיכך ידכה מלא בפרשה שניה, לדרש בה יד שהיא כהה (מנחות ל"ו):}%
% ---------- פרק 13 | פסוק 10 ----------
 \textbf{י} וְשָׁמַרְתָּ֛ אֶת־הַחֻקָּ֥ה הַזֹּ֖את לְמוֹעֲדָ֑הּ מִיָּמִ֖ים יָמִֽימָה׃ \{פ\} \Rashi{\textbf{מימים ימימה.}משנה לשנה (שם):}%
% ---------- פרק 13 | פסוק 11 ----------
 \textbf{יא} וְהָיָ֞ה כִּֽי־יְבִאֲךָ֤ יְהֹוָה֙ אֶל־אֶ֣רֶץ הַֽכְּנַעֲנִ֔י כַּאֲשֶׁ֛ר נִשְׁבַּ֥ע לְךָ֖ וְלַֽאֲבֹתֶ֑יךָ וּנְתָנָ֖הּ לָֽךְ׃ \Rashi{\textbf{(והיה כי יבאך.}יש מרבותינו שלמדו מכאן שלא קדשו בכורות הנולדים במדבר, והאומר קדשו, מפרש ביאה זו אם תקימוהו במדבר תזכו לכנס לארץ ותקימוהו שם): נשבע לך. והיכן נשבע לך? "והבאתי אתכם אל הארץ אשר נשאתי וגו'" (שמות ו'): ונתנה לך. תהא בעיניך כאלו נתנה לך בו ביום, ואל תהי בעיניך כירשת אבות:}%
% ---------- פרק 13 | פסוק 12 ----------
 \textbf{יב} וְהַעֲבַרְתָּ֥ כׇל־פֶּֽטֶר־רֶ֖חֶם לַֽיהֹוָ֑ה וְכׇל־פֶּ֣טֶר ׀ שֶׁ֣גֶר בְּהֵמָ֗ה אֲשֶׁ֨ר יִהְיֶ֥ה לְךָ֛ הַזְּכָרִ֖ים לַיהֹוָֽה׃ \Rashi{\textbf{והעברת.}אין והעברת אלא לשון הפרשה, וכן הוא אומר "והעברת את נחלתו לבתו" (במדבר כ"ז) (מכילתא): שגר בהמה. נפל ששגרתו אמו ושלחתו בלא עתו; ולמדך הכתוב שהוא קדוש בבכורה לפטר את הבא אחריו ואף שאין נפל קרוי שגר, כמו "שגר אלפיך" (דברים ז'), אבל זה לא בא אלא ללמד על הנפל שהרי כבר כתב כל פטר רחם. וא"ת אף בכור בהמה טמאה במשמע, בא ופרש במקום אחר "בבקרך ובצאנך" (דברים ט"ו). לשון אחר יש לפרש, "והעברת כל פטר רחם" — בבכור אדם הכתוב מדבר:}%
% ---------- פרק 13 | פסוק 13 ----------
 \textbf{יג} וְכׇל־פֶּ֤טֶר חֲמֹר֙ תִּפְדֶּ֣ה בְשֶׂ֔ה וְאִם־לֹ֥א תִפְדֶּ֖ה וַעֲרַפְתּ֑וֹ וְכֹ֨ל בְּכ֥וֹר אָדָ֛ם בְּבָנֶ֖יךָ תִּפְדֶּֽה׃ \Rashi{\textbf{פטר חמור.}ולא פטר שאר בהמה טמאה; וגזרת הכתוב היא לפי שנמשלו בכורי מצרים לחמורים; ועוד, שסיעו את ישראל ביציאתן, שאין לך אחד מישראל שלא נטל הרבה חמורים ממצרים טעונים מכספם ומזהבם של מצרים: תפדה בשה. נותן שה לכהן ופטר חמור מתר בהנאה והשה חלין ביד כהן (בכורות ט'): וערפתו. עורפו בקופיץ מאחוריו והורגו; הוא הפסיד ממונו של כהן, לפיכך יפסיד ממונו (שם י'): וכל בכור אדם בבניך תפדה. חמש סלעים פדיונו קצוב במקום אחר (במדבר י"ח):}%
% ---------- פרק 13 | פסוק 14 ----------
 \textbf{יד} וְהָיָ֞ה כִּֽי־יִשְׁאָלְךָ֥ בִנְךָ֛ מָחָ֖ר לֵאמֹ֣ר מַה־זֹּ֑את וְאָמַרְתָּ֣ אֵלָ֔יו בְּחֹ֣זֶק יָ֗ד הוֹצִיאָ֧נוּ יְהֹוָ֛ה מִמִּצְרַ֖יִם מִבֵּ֥ית עֲבָדִֽים׃ \Rashi{\textbf{כי ישאלך בנך מחר.}יש מחר שהוא עכשיו ויש מחר שהוא לאחר זמן, כגון זה, וכגון "ולא יאמרו בניכם מחר לבנינו" (יהושע כ"ב), דבני גד ודבני ראובן: מה זאת. זה תינוק טפש שאינו יודע להעמיק שאלתו וסותם ושואל "מה זאת?", ובמקום אחר הוא אומר "מה העדת והחקים והמשפטים וגו'" (דברים ו'), הרי זאת שאלת בן חכם; דברה תורה כנגד ארבעה בנים — רשע, ושאינו יודע לשאל, והשואל דרך סתומה, והשואל דרך חכמה (מכילתא):}%
% ---------- פרק 13 | פסוק 15 ----------
 \textbf{טו} וַיְהִ֗י כִּֽי־הִקְשָׁ֣ה פַרְעֹה֮ לְשַׁלְּחֵ֒נוּ֒ וַיַּהֲרֹ֨ג יְהֹוָ֤ה כׇּל־בְּכוֹר֙ בְּאֶ֣רֶץ מִצְרַ֔יִם מִבְּכֹ֥ר אָדָ֖ם וְעַד־בְּכ֣וֹר בְּהֵמָ֑ה עַל־כֵּן֩ אֲנִ֨י זֹבֵ֜חַ לַֽיהֹוָ֗ה כׇּל־פֶּ֤טֶר רֶ֙חֶם֙ הַזְּכָרִ֔ים וְכׇל־בְּכ֥וֹר בָּנַ֖י אֶפְדֶּֽה׃ %
% ---------- פרק 13 | פסוק 16 ----------
 \textbf{טז} וְהָיָ֤ה לְאוֹת֙ עַל־יָ֣דְכָ֔ה וּלְטוֹטָפֹ֖ת בֵּ֣ין עֵינֶ֑יךָ כִּ֚י בְּחֹ֣זֶק יָ֔ד הוֹצִיאָ֥נוּ יְהֹוָ֖ה מִמִּצְרָֽיִם׃ \{ס\} \Rashi{\textbf{ולטוטפת.}תפלין; ועל שם שהם ארבעה בתים קרויין טטפת – "טט" בכתפי שתים, "פת" באפריקי שתים (סנהדרין ד'). ומנחם חברו עם "והטף אל דרום" (יחזקאל כ"א), "אל תטיפו" (מיכה ב'), לשון דבור, כמו ולזכרון, שהרואה אותם קשור בין העינים, יזכר הנס וידבר בו:}%
\addparasha{חומש שמות | פרשת בשלח}
% ---------- פרק 13 | פסוק 17 ----------
 \textbf{יז} וַיְהִ֗י בְּשַׁלַּ֣ח פַּרְעֹה֮ אֶת־הָעָם֒ וְלֹא־נָחָ֣ם אֱלֹהִ֗ים דֶּ֚רֶךְ אֶ֣רֶץ פְּלִשְׁתִּ֔ים כִּ֥י קָר֖וֹב ה֑וּא כִּ֣י ׀ אָמַ֣ר אֱלֹהִ֗ים פֶּֽן־יִנָּחֵ֥ם הָעָ֛ם בִּרְאֹתָ֥ם מִלְחָמָ֖ה וְשָׁ֥בוּ מִצְרָֽיְמָה׃ \Rashi{\textbf{ויהי בשלח פרעה, ולא נחם.}נהגם, כמו "לך נחה את העם" (שמות ל"ב), "בהתהלכך תנחה אותך" (משלי ו'): כי קרוב הוא. ונח לשוב באותו הדרך למצרים; ומדרשי אגדה יש הרבה: בראתם מלחמה. כגון מלחמת "וירד העמלקי והכנעני וגו'" (במדבר י"ד), אם הלכו דרך ישר היו חוזרים, ומה אם כשהקיפם דרך מעקם אמרו "נתנה ראש ונשובה מצרימה" (שם), אם הוליכם בפשוטה על אחת כמה וכמה: פן ינחם. יחשבו מחשבה על שיצאו ויתנו לב לשוב:}%
% ---------- פרק 13 | פסוק 18 ----------
 \textbf{יח} וַיַּסֵּ֨ב אֱלֹהִ֧ים ׀ אֶת־הָעָ֛ם דֶּ֥רֶךְ הַמִּדְבָּ֖ר יַם־ס֑וּף וַחֲמֻשִׁ֛ים עָל֥וּ בְנֵי־יִשְׂרָאֵ֖ל מֵאֶ֥רֶץ מִצְרָֽיִם׃ \Rashi{\textbf{ויסב.}הסבם מן הדרך הפשוטה לדרך העקומה: ים סוף. כמו לים סוף, וסוף הוא לשון אגם שגדלים בו קנים, כמו "ותשם בסוף" (שמות ב'), "קנה וסוף קמלו" (ישעיהו י"ט): וחמשים. אין חמושים אלא מזינים; (לפי שהסבן במדבר הוא גרם להם שעלו חמושים, שאלו הסבן דרך ישוב, לא היו מחמשים להם כל מה שצריכין, אלא כאדם שעובר ממקום למקום ובדעתו לקנות שם מה שיצטרך, אבל כשהוא פורש למדבר צריך לזמן לו כל הצרך; ומקרא זה לא נכתב כי אם לשבר את האזן, שלא תתמה במלחמת עמלק ובמלחמת סיחון ועוג ומדין מהיכן היו להם כלי זין שהכו אותם בחרב) וכן הוא אומר "ואתם תעברו חמשים" (יהושע א'), וכן תרגמו אנקלוס "מזרזין", כמו "וירק את חניכיו" (בראשית י"ד) – וזריז. דבר אחר, חמשים אחד מחמשה יצאו וארבעה חלקים מתו בשלשת ימי אפלה (מכילתא):}%
% ---------- פרק 13 | פסוק 19 ----------
 \textbf{יט} וַיִּקַּ֥ח מֹשֶׁ֛ה אֶת־עַצְמ֥וֹת יוֹסֵ֖ף עִמּ֑וֹ כִּי֩ הַשְׁבֵּ֨עַ הִשְׁבִּ֜יעַ אֶת־בְּנֵ֤י יִשְׂרָאֵל֙ לֵאמֹ֔ר פָּקֹ֨ד יִפְקֹ֤ד אֱלֹהִים֙ אֶתְכֶ֔ם וְהַעֲלִיתֶ֧ם אֶת־עַצְמֹתַ֛י מִזֶּ֖ה אִתְּכֶֽם׃ \Rashi{\textbf{השבע השביע.}השביעם שישביעו לבניהם. ולמה לא השביע בניו שישאוהו לארץ כנען מיד כמו שהשביע יעקב? אמר יוסף, אני שליט הייתי במצרים והיה ספק בידי לעשות, אבל בני לא יניחום מצרים לעשות, לכך השביעם לכשיגאלו ויצאו משם שישאוהו: והעליתם את עצמתי מזה אתכם. לאחיו השביע כן, למדנו שאף עצמות כל השבטים העלו עמהם שנאמר אתכם (מכילתא):}%
% ---------- פרק 13 | פסוק 20 ----------
 \textbf{כ} וַיִּסְע֖וּ מִסֻּכֹּ֑ת וַיַּחֲנ֣וּ בְאֵתָ֔ם בִּקְצֵ֖ה הַמִּדְבָּֽר׃ \Rashi{\textbf{ויסעו מסכת.}ביום השני, שהרי בראשון באו מרעמסס לסכות:}%
% ---------- פרק 13 | פסוק 21 ----------
 \textbf{כא} וַֽיהֹוָ֡ה הֹלֵךְ֩ לִפְנֵיהֶ֨ם יוֹמָ֜ם בְּעַמּ֤וּד עָנָן֙ לַנְחֹתָ֣ם הַדֶּ֔רֶךְ וְלַ֛יְלָה בְּעַמּ֥וּד אֵ֖שׁ לְהָאִ֣יר לָהֶ֑ם לָלֶ֖כֶת יוֹמָ֥ם וָלָֽיְלָה׃ \Rashi{\textbf{לנחתם הדרך.}נקוד פתח שהוא כמו להנחותם, "לראותכם בדרך אשר תלכו בה" (דברים א') שהוא כמו להראותכם, אף כאן להנחתם ע"י שליח, ומי הוא השליח? עמוד הענן; והקדוש ברוך הוא בכבודו מוליכו לפניהם, ומכל מקום את עמוד הענן הכין להנחותם על ידו, שהרי ע"י עמוד הענן הם הולכים; עמוד הענן אינו לאורה אלא להורותם הדרך:}%
% ---------- פרק 13 | פסוק 22 ----------
 \textbf{כב} לֹֽא־יָמִ֞ישׁ עַמּ֤וּד הֶֽעָנָן֙ יוֹמָ֔ם וְעַמּ֥וּד הָאֵ֖שׁ לָ֑יְלָה לִפְנֵ֖י הָעָֽם׃ \{פ\} \Rashi{\textbf{לא ימיש.}הקב"ה את עמוד הענן יומם ועמוד האש לילה; מגיד שעמוד הענן משלים לעמוד האש ועמוד האש משלים לעמוד הענן, שעד שלא ישקע זה עולה זה:}%
% ---------- פרק 14 | פסוק 1 ----------
 \textbf{א} וַיְדַבֵּ֥ר יְהֹוָ֖ה אֶל־מֹשֶׁ֥ה לֵּאמֹֽר׃ %
% ---------- פרק 14 | פסוק 2 ----------
 \textbf{ב} דַּבֵּר֮ אֶל־בְּנֵ֣י יִשְׂרָאֵל֒ וְיָשֻׁ֗בוּ וְיַחֲנוּ֙ לִפְנֵי֙ פִּ֣י הַחִירֹ֔ת בֵּ֥ין מִגְדֹּ֖ל וּבֵ֣ין הַיָּ֑ם לִפְנֵי֙ בַּ֣עַל צְפֹ֔ן נִכְח֥וֹ תַחֲנ֖וּ עַל־הַיָּֽם׃ \Rashi{\textbf{וישבו.}לאחוריהם, לצד מצרים היו מקרבין כל יום השלישי, כדי להטעות את פרעה שיאמר תועים הם בדרך, כמו שנאמר "ואמר פרעה לבני ישראל וגו'": ויחנו לפני פי החירת. הוא פיתום, ועכשו נקרא פי החירת על שם שנעשו שם בני חרין; והם שני סלעים גבוהים זקופים והגיא שביניהם קרוי פי הסלעים: לפני בעל צפן. הוא נשאר מכל אלהי מצרים כדי להטעותן, שיאמרו קשה יראתן, ועליו פרש איוב "משגיא לגוים ויאבדם" (איוב י"ב):}%
% ---------- פרק 14 | פסוק 3 ----------
 \textbf{ג} וְאָמַ֤ר פַּרְעֹה֙ לִבְנֵ֣י יִשְׂרָאֵ֔ל נְבֻכִ֥ים הֵ֖ם בָּאָ֑רֶץ סָגַ֥ר עֲלֵיהֶ֖ם הַמִּדְבָּֽר׃ \Rashi{\textbf{ואמר פרעה.}כשישמע שהם שבים לאחוריהם: לבני ישראל. על בני ישראל. וכן (פסוק יד) "ה' ילחם לכם" – עליכם, (בראשית כ יג) "אמרי לי אחי הוא" – אמרי עלי: נבכים הם. כלואים ומשקעים, ובלעז שירי"ר, כמו "נבכי ים" (איוב ל"ח), "בעמק הבכא" (תהילים פ"ד) "מבכי נהרות" (איוב כ"ח). נבכים הם, כלואים הם במדבר, שאינן יודעין לצאת ממנו ולהיכן ילכו:}%
% ---------- פרק 14 | פסוק 4 ----------
 \textbf{ד} וְחִזַּקְתִּ֣י אֶת־לֵב־פַּרְעֹה֮ וְרָדַ֣ף אַחֲרֵיהֶם֒ וְאִכָּבְדָ֤ה בְּפַרְעֹה֙ וּבְכׇל־חֵיל֔וֹ וְיָדְע֥וּ מִצְרַ֖יִם כִּֽי־אֲנִ֣י יְהֹוָ֑ה וַיַּֽעֲשׂוּ־כֵֽן׃ \Rashi{\textbf{ואכבדה בפרעה.}כשהקב"ה מתנקם ברשעים שמו מתגדל ומתכבד, וכן הוא אומר "ונשפטתי אתו וגו'" ואחר כך "והתגדלתי והתקדשתי ונודעתי וגו'" (יחזקאל ל"ח), ואומר "שמה שבר רשפי קשת", ואחר כך "נודע ביהודה אלהים" (תהלים ע"ו), ואומר "נודע ה' משפט עשה" (שם ט'): בפרעה ובכל חילו. הוא התחיל בעברה וממנו התחילה הפרענות: ויעשו כן. להגיד שבחן, ששמעו לקול משה, ולא אמרו האיך נתקרב אל רודפינו? אנו צריכים לברח אלא אמרו אין לנו לקים אלא דברי בן עמרם (מכילתא):}%
% ---------- פרק 14 | פסוק 5 ----------
 \textbf{ה} וַיֻּגַּד֙ לְמֶ֣לֶךְ מִצְרַ֔יִם כִּ֥י בָרַ֖ח הָעָ֑ם וַ֠יֵּהָפֵ֠ךְ לְבַ֨ב פַּרְעֹ֤ה וַעֲבָדָיו֙ אֶל־הָעָ֔ם וַיֹּֽאמְרוּ֙ מַה־זֹּ֣את עָשִׂ֔ינוּ כִּֽי־שִׁלַּ֥חְנוּ אֶת־יִשְׂרָאֵ֖ל מֵעׇבְדֵֽנוּ׃ \Rashi{\textbf{ויגד למלך מצרים.}איקטורין שלח עמהם, וכיון שהגיעו לשלשת ימים שקבעו לילך ולשוב וראו שאינן חוזרין למצרים, באו והגידו לפרעה ביום הרביעי, בחמישי ובששי רדפו אחריהם, ליל שביעי ירדו לים, בשחרית אמרו שירה והוא יום שביעי של פסח; לכך אנו קורין השירה ביום השביעי: ויהפך. נהפך ממה שהיה, שהרי אמר להם "קומו צאו מתוך עמי" (שמות י"ב), ונהפך לב עבדיו, שהרי לשעבר היו אומרים לו "עד מתי יהיה זה לנו למוקש" (שמות י'), ועכשיו נהפכו לרדף אחריהם בשביל ממונם שהשאילום: מעבדנו. מעבד אותנו:}%
% ---------- פרק 14 | פסוק 6 ----------
 \textbf{ו} וַיֶּאְסֹ֖ר אֶת־רִכְבּ֑וֹ וְאֶת־עַמּ֖וֹ לָקַ֥ח עִמּֽוֹ׃ \Rashi{\textbf{ויאסר את רכבו.}הוא בעצמו (מכילתא): ואת עמו לקח עמו. משכם בדברים, לקינו ונטלו ממוננו ושלחנום, באו עמי ואני לא אתנהג עמכם כשאר מלכים, דרך שאר מלכים עבדיו קודמין לו במלחמה, ואני אקדים לפניכם, שנאמר ופרעה הקריב – הקריב עצמו ומהר לפני חילותיו –, דרך שאר מלכים לטל בזה בראש כמו שיבחר, אני אשוה עמכם בחלק שנאמר "אחלק שלל" (שמות ט"ו):}%
% ---------- פרק 14 | פסוק 7 ----------
 \textbf{ז} וַיִּקַּ֗ח שֵׁשׁ־מֵא֥וֹת רֶ֙כֶב֙ בָּח֔וּר וְכֹ֖ל רֶ֣כֶב מִצְרָ֑יִם וְשָׁלִשִׁ֖ם עַל־כֻּלּֽוֹ׃ \Rashi{\textbf{בחור.}נבחרים, בחור לשון יחיד, כל רכב ורכב שבמנין זה היה בחור: וכל רכב מצרים. ועמהם כל שאר הרכב; ומהיכן היו הבהמות הללו? אם תאמר משל מצרים, הרי נאמר "וימת כל מקנה מצרים" (שמות ט'), ואם משל ישראל, והלא נאמר "וגם מקננו ילך עמנו", משל מי היו? מן הירא את דבר ה'. מכאן היה ר' שמעון אומר, כשר שבמצרים הרג, טוב שבנחשים רצץ את מחו (מכילתא): ושלשם על כלו. שרי צבאות, כתרגומו:}%
% ---------- פרק 14 | פסוק 8 ----------
 \textbf{ח} וַיְחַזֵּ֣ק יְהֹוָ֗ה אֶת־לֵ֤ב פַּרְעֹה֙ מֶ֣לֶךְ מִצְרַ֔יִם וַיִּרְדֹּ֕ף אַחֲרֵ֖י בְּנֵ֣י יִשְׂרָאֵ֑ל וּבְנֵ֣י יִשְׂרָאֵ֔ל יֹצְאִ֖ים בְּיָ֥ד רָמָֽה׃ \Rashi{\textbf{ויחזק ה' את לב פרעה.}שהיה תולה אם לרדף אם לאו, וחזק את לבו לרדף (שם): ביד רמה. בגבורה גבוהה ומפרסמת:}%
% ---------- פרק 14 | פסוק 9 ----------
 \textbf{ט} וַיִּרְדְּפ֨וּ מִצְרַ֜יִם אַחֲרֵיהֶ֗ם וַיַּשִּׂ֤יגוּ אוֹתָם֙ חֹנִ֣ים עַל־הַיָּ֔ם כׇּל־סוּס֙ רֶ֣כֶב פַּרְעֹ֔ה וּפָרָשָׁ֖יו וְחֵיל֑וֹ עַל־פִּי֙ הַֽחִירֹ֔ת לִפְנֵ֖י בַּ֥עַל צְפֹֽן׃ %
% ---------- פרק 14 | פסוק 10 ----------
 \textbf{י} וּפַרְעֹ֖ה הִקְרִ֑יב וַיִּשְׂאוּ֩ בְנֵֽי־יִשְׂרָאֵ֨ל אֶת־עֵינֵיהֶ֜ם וְהִנֵּ֥ה מִצְרַ֣יִם ׀ נֹסֵ֣עַ אַחֲרֵיהֶ֗ם וַיִּֽירְאוּ֙ מְאֹ֔ד וַיִּצְעֲק֥וּ בְנֵֽי־יִשְׂרָאֵ֖ל אֶל־יְהֹוָֽה׃ \Rashi{\textbf{ופרעה הקריב.}היה לו לכתב ופרעה קרב, מהו הקריב? הקריב עצמו ונתאמץ לקדם לפניהם כמו שהתנה עמהם: נסע אחריהם. בלב אחד כאיש אחד. דבר אחר – והנה מצרים נסע אחריהם, ראו שר של מצרים נוסע מן השמים לעזר למצרים; – תנחומא: ויצעקו. תפשו אמנות אבותם – באברהם הוא אומר "אל המקום אשר עמד שם" (בראשית י"ט), ביצחק "לשוח בשדה", (שם כ"ד), ביעקב "ויפגע במקום" (שם כ"ח):}%
% ---------- פרק 14 | פסוק 11 ----------
 \textbf{יא} וַיֹּאמְרוּ֮ אֶל־מֹשֶׁה֒ הֲֽמִבְּלִ֤י אֵין־קְבָרִים֙ בְּמִצְרַ֔יִם לְקַחְתָּ֖נוּ לָמ֣וּת בַּמִּדְבָּ֑ר מַה־זֹּאת֙ עָשִׂ֣יתָ לָּ֔נוּ לְהוֹצִיאָ֖נוּ מִמִּצְרָֽיִם׃ \Rashi{\textbf{המבלי אין קברים.}וכי מחמת חסרון קברים – שאין קברים במצרים לקבר שם – לקחתנו משם. ש"י פו"ר פלינ"צא ד"י נו"ן פוש"יש בלעז:}%
% ---------- פרק 14 | פסוק 12 ----------
 \textbf{יב} הֲלֹא־זֶ֣ה הַדָּבָ֗ר אֲשֶׁר֩ דִּבַּ֨רְנוּ אֵלֶ֤יךָ בְמִצְרַ֙יִם֙ לֵאמֹ֔ר חֲדַ֥ל מִמֶּ֖נּוּ וְנַֽעַבְדָ֣ה אֶת־מִצְרָ֑יִם כִּ֣י ט֥וֹב לָ֙נוּ֙ עֲבֹ֣ד אֶת־מִצְרַ֔יִם מִמֻּתֵ֖נוּ בַּמִּדְבָּֽר׃ \Rashi{\textbf{אשר דברנו אליך בארץ מצרים.}והיכן דברו? "ירא ה' עליכם וישפט" (שמות ה') (מכילתא): ממתנו. מאשר נמות, ואם היה נקוד מלאפום היה נבאר ממיתתנו, עכשיו שנקוד בשורק נבאר מאשר נמות וכן "מי יתן מותנו" (שמות ט"ז) – שנמות, וכן "מי יתן מותי" דאבשלום (שמואל ב' י"ט) – שאמות. כמו "ליום קומי לעד" (צפניה ג'), "עד יום שובי בשלום" (דברי הימים ב' י"ח) – שאקום, שאשוב:}%
% ---------- פרק 14 | פסוק 13 ----------
 \textbf{יג} וַיֹּ֨אמֶר מֹשֶׁ֣ה אֶל־הָעָם֮ אַל־תִּירָ֒אוּ֒ הִֽתְיַצְּב֗וּ וּרְאוּ֙ אֶת־יְשׁוּעַ֣ת יְהֹוָ֔ה אֲשֶׁר־יַעֲשֶׂ֥ה לָכֶ֖ם הַיּ֑וֹם כִּ֗י אֲשֶׁ֨ר רְאִיתֶ֤ם אֶת־מִצְרַ֙יִם֙ הַיּ֔וֹם לֹ֥א תֹסִ֛פוּ לִרְאֹתָ֥ם ע֖וֹד עַד־עוֹלָֽם׃ \Rashi{\textbf{כי אשר ראיתם את מצרים וגו'.}מה שראיתם אותם אינו אלא היום, היום הוא שראיתם אותם ולא תוסיפו עוד:}%
% ---------- פרק 14 | פסוק 14 ----------
 \textbf{יד} יְהֹוָ֖ה יִלָּחֵ֣ם לָכֶ֑ם וְאַתֶּ֖ם תַּחֲרִשֽׁוּן׃ \{פ\} \Rashi{\textbf{ילחם לכם.}בשבילכם; וכן "כי ה' נלחם להם", וכן "אם לאל תריבון" (איוב י"ג), וכן "ואשר דבר לי" (בראשית כ"ד), וכן "האתם תריבון לבעל" (שופטים ו'):}%
% ---------- פרק 14 | פסוק 15 ----------
 \textbf{טו} וַיֹּ֤אמֶר יְהֹוָה֙ אֶל־מֹשֶׁ֔ה מַה־תִּצְעַ֖ק אֵלָ֑י דַּבֵּ֥ר אֶל־בְּנֵי־יִשְׂרָאֵ֖ל וְיִסָּֽעוּ׃ \Rashi{\textbf{מה תצעק אלי.}למדנו שהיה משה עומד ומתפלל, אמר לו הקב"ה, לא עת עתה להאריך בתפלה, שישראל נתונין בצרה (מכילתא). ד"א – מה תצעק אלי, עלי הדבר תלוי ולא עליך, כמו שנאמר להלן "על בני ועל פעל ידי תצוני" (ישעיהו מ"ה): דבר אל בני ישראל ויסעו. אין להם אלא לסע, שאין הים עומד בפניהם, כדאי הוא זכות אבותיהם והאמונה שהאמינו בי ויצאו, לקרע להם הים (מכילתא):}%
% ---------- פרק 14 | פסוק 16 ----------
 \textbf{טז} וְאַתָּ֞ה הָרֵ֣ם אֶֽת־מַטְּךָ֗ וּנְטֵ֧ה אֶת־יָדְךָ֛ עַל־הַיָּ֖ם וּבְקָעֵ֑הוּ וְיָבֹ֧אוּ בְנֵֽי־יִשְׂרָאֵ֛ל בְּת֥וֹךְ הַיָּ֖ם בַּיַּבָּשָֽׁה׃ %
% ---------- פרק 14 | פסוק 17 ----------
 \textbf{יז} וַאֲנִ֗י הִנְנִ֤י מְחַזֵּק֙ אֶת־לֵ֣ב מִצְרַ֔יִם וְיָבֹ֖אוּ אַחֲרֵיהֶ֑ם וְאִכָּבְדָ֤ה בְּפַרְעֹה֙ וּבְכׇל־חֵיל֔וֹ בְּרִכְבּ֖וֹ וּבְפָרָשָֽׁיו׃ %
% ---------- פרק 14 | פסוק 18 ----------
 \textbf{יח} וְיָדְע֥וּ מִצְרַ֖יִם כִּי־אֲנִ֣י יְהֹוָ֑ה בְּהִכָּבְדִ֣י בְּפַרְעֹ֔ה בְּרִכְבּ֖וֹ וּבְפָרָשָֽׁיו׃ %
% ---------- פרק 14 | פסוק 19 ----------
 \textbf{יט} וַיִּסַּ֞ע מַלְאַ֣ךְ הָאֱלֹהִ֗ים הַהֹלֵךְ֙ לִפְנֵי֙ מַחֲנֵ֣ה יִשְׂרָאֵ֔ל וַיֵּ֖לֶךְ מֵאַחֲרֵיהֶ֑ם וַיִּסַּ֞ע עַמּ֤וּד הֶֽעָנָן֙ מִפְּנֵיהֶ֔ם וַֽיַּעֲמֹ֖ד מֵאַחֲרֵיהֶֽם׃ \Rashi{\textbf{וילך מאחריהם.}להבדיל בין מחנה מצרים ובין מחנה ישראל ולקבל חצים ובליסטראות של מצרים. בכל מקום הוא אומר מלאך ה' וכאן מלאך האלהים, אין אלהים בכל מקום אלא דין, מלמד שהיו ישראל נתונין בדין באותה שעה אם להנצל אם להאבד עם מצרים (שם): ויסע עמוד הענן. כשחשכה והשלים עמוד הענן את המחנה לעמוד האש, לא נסתלק הענן כמו שהיה רגיל להסתלק ערבית לגמרי, אלא נסע והלך לו מאחריהם להחשיך למצרים:}%
% ---------- פרק 14 | פסוק 20 ----------
 \textbf{כ} וַיָּבֹ֞א בֵּ֣ין ׀ מַחֲנֵ֣ה מִצְרַ֗יִם וּבֵין֙ מַחֲנֵ֣ה יִשְׂרָאֵ֔ל וַיְהִ֤י הֶֽעָנָן֙ וְהַחֹ֔שֶׁךְ וַיָּ֖אֶר אֶת־הַלָּ֑יְלָה וְלֹא־קָרַ֥ב זֶ֛ה אֶל־זֶ֖ה כׇּל־הַלָּֽיְלָה׃ \Rashi{\textbf{ויבא בין מחנה מצרים.}משל למהלך בדרך ובנו מהלך לפניו, באו לסטים לשבותו, נטלו מלפניו ונתנו לאחריו, בא זאב מאחריו, נתנו לפניו, באו לסטים לפניו וזאבים מאחריו, נתנו על זרועו ונלחם בהם, כך "ואנכי תרגלתי לאפרים קחם על זרועתיו" (הושע י"א) (מכילתא): ויהי הענן והחשך. למצרים: ויאר עמוד האש את הלילה לישראל, והולך לפניהם כדרכו ללכת כל הלילה והחשך של ערפל לצד מצרים: ולא קרב זה אל זה. מחנה אל מחנה (שם):}%
% ---------- פרק 14 | פסוק 21 ----------
 \textbf{כא} וַיֵּ֨ט מֹשֶׁ֣ה אֶת־יָדוֹ֮ עַל־הַיָּם֒ וַיּ֣וֹלֶךְ יְהֹוָ֣ה ׀ אֶת־הַ֠יָּ֠ם בְּר֨וּחַ קָדִ֤ים עַזָּה֙ כׇּל־הַלַּ֔יְלָה וַיָּ֥שֶׂם אֶת־הַיָּ֖ם לֶחָרָבָ֑ה וַיִּבָּקְע֖וּ הַמָּֽיִם׃ \Rashi{\textbf{ברוח קדים עזה.}ברוח קדים שהיא עזה שברוחות, היא הרוח שהקב"ה נפרע בה מן הרשעים, שנאמר, "כרוח קדים אפיצם" (ירמיה י"ח), "יבוא קדים רוח ה'" (הושע י"ג), "רוח הקדים שברך בלב ימים" (יחזקאל כ"ז), "הגה ברוחו הקשה ביום קדים" (ישעיה כ"ז): ויבקעו המים. כל מים שבעולם:}%
% ---------- פרק 14 | פסוק 22 ----------
 \textbf{כב} וַיָּבֹ֧אוּ בְנֵֽי־יִשְׂרָאֵ֛ל בְּת֥וֹךְ הַיָּ֖ם בַּיַּבָּשָׁ֑ה וְהַמַּ֤יִם לָהֶם֙ חוֹמָ֔ה מִֽימִינָ֖ם וּמִשְּׂמֹאלָֽם׃ %
% ---------- פרק 14 | פסוק 23 ----------
 \textbf{כג} וַיִּרְדְּפ֤וּ מִצְרַ֙יִם֙ וַיָּבֹ֣אוּ אַחֲרֵיהֶ֔ם כֹּ֚ל ס֣וּס פַּרְעֹ֔ה רִכְבּ֖וֹ וּפָרָשָׁ֑יו אֶל־תּ֖וֹךְ הַיָּֽם׃ \Rashi{\textbf{כל סוס פרעה.}וכי סוס אחד היה? מגיד שאין כלם חשובין לפני המקום אלא כסוס אחד (מכילתא):}%
% ---------- פרק 14 | פסוק 24 ----------
 \textbf{כד} וַֽיְהִי֙ בְּאַשְׁמֹ֣רֶת הַבֹּ֔קֶר וַיַּשְׁקֵ֤ף יְהֹוָה֙ אֶל־מַחֲנֵ֣ה מִצְרַ֔יִם בְּעַמּ֥וּד אֵ֖שׁ וְעָנָ֑ן וַיָּ֕הׇם אֵ֖ת מַחֲנֵ֥ה מִצְרָֽיִם׃ \Rashi{\textbf{באשמרת הבקר.}שלשת חלקי הלילה קרויין אשמורה, ואותה שלפני הבקר קורא אשמרת הבקר. ואומר אני, לפי שהלילה חלוק למשמרות שיר של מלאכי השרת, כת אחר כת, לשלשה חלקים, לכך קרוי אשמרת, וזהו שתרגם אנקלוס "מטרת": וישקף. ויבט, כלומר פנה אליהם להשחיתם. ותרגומו "ואסתכי" אף הוא לשון הבטה, כמו "שדה צופים", (במדבר כ"ג) – חקל סכותא: בעמוד אש וענן. עמוד ענן יורד ועושה אותו כטיט ועמוד אש מרתיחו, וטלפי סוסיהם משתמטות (מכילתא): ויהם. לשון מהומה, אשדורד"ישון בלעז, ערבבם, נטל סגניות שלהם. ושנינו בפרקי רבי אליעזר בנו של רבי יוסי הגלילי: כל מקום שנאמר בו מהומה הרעמת קול הוא, וזה אב לכלן "וירעם ה' בקול גדול וגו' על פלשתים ויהמם" (שמואל א ז'):}%
% ---------- פרק 14 | פסוק 25 ----------
 \textbf{כה} וַיָּ֗סַר אֵ֚ת אֹפַ֣ן מַרְכְּבֹתָ֔יו וַֽיְנַהֲגֵ֖הוּ בִּכְבֵדֻ֑ת וַיֹּ֣אמֶר מִצְרַ֗יִם אָנ֙וּסָה֙ מִפְּנֵ֣י יִשְׂרָאֵ֔ל כִּ֣י יְהֹוָ֔ה נִלְחָ֥ם לָהֶ֖ם בְּמִצְרָֽיִם׃ \{פ\} \Rashi{\textbf{ויסר את אפן מרכבתיו.}מכח האש נשרפו הגלגלים, והמרכבות נגררות, והיושבים בהם נעים ואבריהן מתפרקין (מכילתא): וינהגהו בכבדת. בהנהגה שהיא כבדה וקשה להם; במדה שמדדו: "ויכבד" לבו הוא ועבדיו (שמות ט'), אף כאן וינהגהו "בכבדת": נלחם להם במצרים. במצריים. דבר אחר במצרים – בארץ מצרים, שכשם שאלו לוקים על הים כך לוקים אותם שנשארו במצרים (מכילתא):}%
% ---------- פרק 14 | פסוק 26 ----------
 \textbf{כו} וַיֹּ֤אמֶר יְהֹוָה֙ אֶל־מֹשֶׁ֔ה נְטֵ֥ה אֶת־יָדְךָ֖ עַל־הַיָּ֑ם וְיָשֻׁ֤בוּ הַמַּ֙יִם֙ עַל־מִצְרַ֔יִם עַל־רִכְבּ֖וֹ וְעַל־פָּרָשָֽׁיו׃ \Rashi{\textbf{וישבו המים.}שזקופין ועומדים כחומה, ישובו למקומם ויכסו על מצרים:}%
% ---------- פרק 14 | פסוק 27 ----------
 \textbf{כז} וַיֵּט֩ מֹשֶׁ֨ה אֶת־יָד֜וֹ עַל־הַיָּ֗ם וַיָּ֨שׇׁב הַיָּ֜ם לִפְנ֥וֹת בֹּ֙קֶר֙ לְאֵ֣יתָנ֔וֹ וּמִצְרַ֖יִם נָסִ֣ים לִקְרָאת֑וֹ וַיְנַעֵ֧ר יְהֹוָ֛ה אֶת־מִצְרַ֖יִם בְּת֥וֹךְ הַיָּֽם׃ \Rashi{\textbf{לפנות בקר.}לעת שהבקר פונה לבא: לאיתנו. לתקפו הראשון (מכילתא): נסים לקראתו. שהיו מהממים ומטרפים ורצין לקראת המים: וינער ה'. כאדם שמנער את הקדרה והופך העליון למטה והתחתון למעלה, כך היו עולין ויורדין ומשתברין בים, ונתן הקב"ה בהם חיות לקבל היסורין (שם): וינער. "ושניק", והוא לשון טרוף בלשון ארמי, והרבה יש במ"א:}%
% ---------- פרק 14 | פסוק 28 ----------
 \textbf{כח} וַיָּשֻׁ֣בוּ הַמַּ֗יִם וַיְכַסּ֤וּ אֶת־הָרֶ֙כֶב֙ וְאֶת־הַפָּ֣רָשִׁ֔ים לְכֹל֙ חֵ֣יל פַּרְעֹ֔ה הַבָּאִ֥ים אַחֲרֵיהֶ֖ם בַּיָּ֑ם לֹֽא־נִשְׁאַ֥ר בָּהֶ֖ם עַד־אֶחָֽד׃ \Rashi{\textbf{ויכסו את הרכב, לכל חיל פרעה.}כך דרך המקראות לכתב למ"ד יתרה, כמו "לכל כליו תעשה נחשת" (שמות כ"ז), וכן "לכל כלי המשכן בכל עבדתו" (במדבר ד'), "ויתדותם ומיתריהם לכל כליהם" (שם), ואינה אלא תקון לשון:}%
% ---------- פרק 14 | פסוק 29 ----------
 \textbf{כט} וּבְנֵ֧י יִשְׂרָאֵ֛ל הָלְכ֥וּ בַיַּבָּשָׁ֖ה בְּת֣וֹךְ הַיָּ֑ם וְהַמַּ֤יִם לָהֶם֙ חֹמָ֔ה מִֽימִינָ֖ם וּמִשְּׂמֹאלָֽם׃ %
% ---------- פרק 14 | פסוק 30 ----------
 \textbf{ל} וַיּ֨וֹשַׁע יְהֹוָ֜ה בַּיּ֥וֹם הַה֛וּא אֶת־יִשְׂרָאֵ֖ל מִיַּ֣ד מִצְרָ֑יִם וַיַּ֤רְא יִשְׂרָאֵל֙ אֶת־מִצְרַ֔יִם מֵ֖ת עַל־שְׂפַ֥ת הַיָּֽם׃ \Rashi{\textbf{וירא ישראל את מצרים מת.}שפלטן הים על שפתו, כדי שלא יאמרו ישראל, כשם שאנו עולים מצד זה כך הם עולין מצד אחר, רחוק ממנו, וירדפו אחרינו (פסחים קי"ח):}%
% ---------- פרק 14 | פסוק 31 ----------
 \textbf{לא} וַיַּ֨רְא יִשְׂרָאֵ֜ל אֶת־הַיָּ֣ד הַגְּדֹלָ֗ה אֲשֶׁ֨ר עָשָׂ֤ה יְהֹוָה֙ בְּמִצְרַ֔יִם וַיִּֽירְא֥וּ הָעָ֖ם אֶת־יְהֹוָ֑ה וַֽיַּאֲמִ֙ינוּ֙ בַּֽיהֹוָ֔ה וּבְמֹשֶׁ֖ה עַבְדּֽוֹ׃ \{פ\} \Rashi{\textbf{את היד הגדולה.}את הגבורה הגדולה שעשתה ידו של הקב"ה. והרבה לשונות נופלין על לשון יד וכלן לשון יד ממש הן, והמפרשו יתקן הלשון אחר ענין הדבור:}%
% ---------- פרק 15 | פסוק 1 ----------
 \textbf{א} אָ֣ז יָשִֽׁיר־מֹשֶׁה֩ וּבְנֵ֨י יִשְׂרָאֵ֜ל אֶת־הַשִּׁירָ֤ה הַזֹּאת֙ לַֽיהֹוָ֔ה וַיֹּאמְר֖וּ לֵאמֹ֑ר        אָשִׁ֤ירָה לַֽיהֹוָה֙ כִּֽי־גָאֹ֣ה גָּאָ֔ה        ס֥וּס וְרֹכְב֖וֹ רָמָ֥ה בַיָּֽם׃ \Rashi{\textbf{אז ישיר משה.}אז כשראה הנס עלה בלבו שישיר שירה. וכן "אז ידבר יהושע" (יהושע י'), וכן "ובית יעשה לבת פרעה" (מלכים א ז') – חשב בלבו שיעשה לה, אף כאן ישיר אמר לו לבו שישיר וכן עשה – ויאמרו לאמר אשירה לה', וכן ביהושע כשראה הנס אמר לו לבו שידבר וכן עשה – "ויאמר לעיני ישראל" (יהושע י'), וכן שירת הבאר, שפתח בה אז ישיר ישראל, פרש אחריו "עלי באר ענו לה" (במדבר י"א), "אז יבנה שלמה במה" (מלכים א י"א), פרשו בו חכמי ישראל שבקש לבנות ולא בנה, למדנו שהיו"ד על שם המחשבה נאמרה, זהו לישב פשוטו. אבל מדרשו אמרו רבותינו ז"ל "מכאן רמז לתחית המתים מן התורה" וכן בכלן, חוץ משל שלמה, שפרשוהו בקש לבנות ולא בנה. ואין לומר ולישב לשון הזה כשאר דברים הנכתבים בלשון עתיד והן מיד, כגון "ככה יעשה איוב" (איוב א'), "ע"פ ה' יחנו" (במדבר ט'), "ויש אשר יהיה הענן" (שם), לפי שהן דבר ההווה תמיד ונופל בו בין לשון עתיד ובין לשון עבר, אבל זה שלא היה אלא לשעה, איני יכול לישבו בלשון הזה: כי גאה גאה. כתרגומו. (ד"א – בא הכפל לומר שעשה דבר שאי אפשר לבשר ודם לעשות; כשהוא נלחם בחברו ומתגבר עליו, מפילו מן הסוס, וכאן הסוס ורכבו רמה בים, וכל שאי אפשר לעשות על ידי זולתו נופל בו לשון גאות, כמו "כי גאות עשה" (ישעיהו י"ב), וכן כל השירה תמצא כפולה, עזי וזמרת יה ויהי לי לישועה, ה' איש מלחמה ה' שמו, וכן כלם). דבר אחר – כי גאה גאה, על כל השירות וכל מה שאקלס בו, עוד יש בו תוספת, ולא כמדת מלך בשר ודם שמקלסין אותו ואין בו: סוס ורכבו. שניהם קשורים זה בזה והמים מעלין אותן ויורדין לעמק ואינן נפרדין (מכילתא): רמה. השליך, וכן "ורמיו לגוא אתון נורא" (דניאל ג'). ומדרש אגדה: כתוב אחד אומר רמה, וכתוב אחד אומר ירה, מלמד שהיו עולין לרום ויורדין לתהום, כמו "מי ירה אבן פנתה" (איוב ל"ח), מלמעלה למטה:}%
% ---------- פרק 15 | פסוק 2 ----------
 \textbf{ב} עׇזִּ֤י וְזִמְרָת֙ יָ֔הּ וַֽיְהִי־לִ֖י לִֽישׁוּעָ֑ה        זֶ֤ה אֵלִי֙ וְאַנְוֵ֔הוּ        אֱלֹהֵ֥י אָבִ֖י וַאֲרֹמְמֶֽנְהוּ׃ \Rashi{\textbf{עזי וזמרת יה.}אנקלוס תרגם "תקפי ותשבחתי", עזי כמו עזי, וזמרת כמו וזמרתי, ואני תמה על לשון המקרא, שאין לך כמוהו בנקדתו במקרא אלא בשלשה מקומות שהוא סמוך אצל וזמרת, וכל שאר מקומות נקוד שור"ק, "ה' עזי ומעזי" (ירמיהו ט"ז), "עזו אליך אשמרה" (תהילים נ"ט). וכן כל תבה בת שתי אותיות הנקודה מלאפום כשהיא מארכת באות שלישית ואין השניה בחטף, הראשונה נקודה בשורק, כגון עז עזי, רק רקי, חק חקי, על עלו – "וסר מעלהם עלו" (ישעיהו י"ד), כל כלו – "ושלשם על כלו" (שמות י"ד), ואלו ג' עזי וזמרת של כאן, ושל ישעיה, ושל תהלים, נקודים בחטף קמץ, ועוד אין באחד מהם כתוב וזמרתי אלא וזמרת, וכלם סמוך להם "ויהי לי לישועה". לכך אני אומר, לישב לשון המקרא, שאין עזי כמו עזי ולא וזמרת כמו וזמרתי, אלא עזי שם דבר הוא כמו "הישבי בשמים" (תהילים קכ"ג), "שכני בחגוי סלע" (עובדיה א'), "שכני סנה" (דברים ל"ג), וזהו השבח עזי וזמרת יה הוא היה לי לישועה. וזמרת דבוק הוא לתבת השם, כמו "לעזרת ה'" (שופטים ה'), "בעברת ה'" (ישעיהו ט'), "על דברת בני האדם" (קהלת ג'). ולשון וזמרת לשון "לא תזמר" (ויקרא כ"ה), "זמיר עריצים" (ישעיהו כ"ה), לשון כסוח וכריתה – עזו ונקמתו של אלהינו היה לנו לישועה. ואל תתמה על לשון ויהי, שלא נאמר היה, שיש לנו מקראות מדברים בלשון זה, וזה דגמתו: "את קירות הבית סביב להיכל ולדביר ויעש צלעות סביב" (מלכים א ו'), היה לו לומר עשה צלעות סביב. וכן בדברי הימים: "ובני ישראל הישבים בערי יהודה וימלך עליהם רחבעם" (דברי הימים ב' י'), היה לו לומר מלך עליהם רחבעם. "מבלתי יכלת ה' וגו' וישחטם" (במדבר י"ד), היה לו לומר שחטם. "והאנשים אשר שלח משה וגו' וימתו" (שם) מתו היה לו לומר, "ואשר לא שם לבו אל דבר ה' ויעזב" (שמות ט'), היה לו לומר עזב: זה אלי. בכבודו נגלה עליהם והיו מראין אותו באצבע, ראתה שפחה על הים מה שלא ראו נביאים (מכילתא): ואנוהו. אנקלוס תרגם לשון נוה – "נוה שאנן" (ישעיהו ל"ג), "לנוה צאן" (שם ס"ה). דבר אחר – ואנוהו לשון נוי, אספר נויו ושבחו לבאי עולם, כגון "מה דודך מדוד וגו' דודי צח ואדום" (שיר השירים ה'), וכל הענין: אלהי אבי. הוא זה וארממנהו: אלהי אבי. לא אני תחלת הקדשה אלא מחזקת ועומדת לי הקדשה ואלהותו עלי מימי אבותי:}%
% ---------- פרק 15 | פסוק 3 ----------
 \textbf{ג} יְהֹוָ֖ה אִ֣ישׁ מִלְחָמָ֑ה יְהֹוָ֖ה שְׁמֽוֹ׃ \Rashi{\textbf{ה' איש מלחמה.}בעל מלחמות, כמו "איש נעמי" (רות א'), וכל איש ואישך מתרגמין בעל, וכן "וחזקת והיית לאיש" (מלכים א ב') – לגבור: ה' שמו. מלחמותיו לא בכלי זין אלא בשמו הוא נלחם; כמו שאמר דוד "ואנכי בא אליך בשם ה' צבאות" (שמואל א י"ז). דבר אחר, ה' שמו – אף בשעה שהוא נלחם ונוקם מאויביו, אוחז הוא במדתו לרחם על ברואיו ולזון את כל באי עולם; ולא כמדת מלכי אדמה כשהוא עוסק במלחמה פונה עצמו מכל עסקים ואין בו כח לעשות זו וזו (מכילתא):}%
% ---------- פרק 15 | פסוק 4 ----------
 \textbf{ד} מַרְכְּבֹ֥ת פַּרְעֹ֛ה וְחֵיל֖וֹ יָרָ֣ה בַיָּ֑ם        וּמִבְחַ֥ר שָֽׁלִשָׁ֖יו טֻבְּע֥וּ בְיַם־סֽוּף׃ \Rashi{\textbf{ירה בים.}"שדי בימא", שדי לשון יריה, וכן הוא אומר "או ירה יירה" (שמות י"ט) – או אשתדאה ישתדי, והתי"ו משמש באלו במקום יתפעל: ומבחר. שם דבר. כמו מרכב, משכב, מקרא קדש: טבעו. אין טביעה אלא במקום טיט, כמו "טבעתי ביון מצולה" (תהילים ס"ט), "ויטבע ירמיהו בטיט" (ירמיהו ל"ח), מלמד, שנעשה הים טיט, לגמל להם כמדתם ששעבדו את ישראל בחמר ובלבנים (מכילתא):}%
% ---------- פרק 15 | פסוק 5 ----------
 \textbf{ה} תְּהֹמֹ֖ת יְכַסְיֻ֑מוּ יָרְד֥וּ בִמְצוֹלֹ֖ת כְּמוֹ־אָֽבֶן׃ \Rashi{\textbf{יכסימו.}כמו יכסום, והיו"ד האמצעית יתרה בו; ודרך מקראות בכך, כמו "וצאנך ירבין" (דברים ח'), "ירוין מדשן ביתך" (תהלים ל"ו), והיו"ד ראשונה שמשמעה לשון עתיד, כך פרשהו: טבעו בים סוף כדי שיחזרו המים ויכסו אותן. יכסימו אין דומה לו במקרא בנקדתו, ודרכו להיות נקוד יכסימו מלאפום: כמו אבן. ובמקום אחר "צללו כעופרת", ובמקום אחר "יאכלמו כקש", הרשעים כקש, הולכים ומטרפין עולין ויורדין, בינונים כאבן, והכשרים כעופרת שנחו מיד (מכילתא):}%
% ---------- פרק 15 | פסוק 6 ----------
 \textbf{ו} יְמִֽינְךָ֣ יְהֹוָ֔ה נֶאְדָּרִ֖י בַּכֹּ֑חַ        יְמִֽינְךָ֥ יְהֹוָ֖ה תִּרְעַ֥ץ אוֹיֵֽב׃ \Rashi{\textbf{ימינך ימינך.}שני פעמים, כשישראל עושין את רצונו של מקום השמאל נעשית ימין (שם): ימינך ה' נאדרי בכח. להציל את ישראל וימינך השנית תרעץ אויב. ולי נראה אותה ימין עצמה תרעץ אויב, מה שאי אפשר לאדם – לעשות שתי מלאכות ביד אחת; ופשוטו של מקרא ימינך הנאדרת בכח מה מלאכתה? ימינך ה' תרעץ אויב; וכמה מקראות דגמתו, "כי הנה אויביך ה' כי הנה אויביך יאבדו" (תהלים צ"ב), ודומיהם: נאדרי. היו"ד יתרה, כמו "רבתי עם, שרתי במדינות" (איכה א'), "גנבתי יום" (בראשית ל"א): תרעץ אויב. תמיד היא רועצת ומשברת האויב; ודומה לו "וירעצו וירוצצו את בני ישראל", בשופטים (י, ח). (ד"א – ימינך הנאדרת בכח היא משברת ומלקה אויב):}%
% ---------- פרק 15 | פסוק 7 ----------
 \textbf{ז} וּבְרֹ֥ב גְּאוֹנְךָ֖ תַּהֲרֹ֣ס קָמֶ֑יךָ        תְּשַׁלַּח֙ חֲרֹ֣נְךָ֔ יֹאכְלֵ֖מוֹ כַּקַּֽשׁ׃ \Rashi{\textbf{וברב גאונך.}זאת היד בלבד רועצת האויב – כשהוא מרימה ברב גאונו אז יהרס קמיו – ואם ברב גאונו לבד אויביו נהרסים קל וחמר כששלח בם חרון אף יאכלמו): תהרס. תמיד אתה הורס קמיך הקמים נגדך, ומי הם הקמים כנגדו? אלו הקמים על ישראל; וכן הוא אומר "כי הנה אויביך יהמיון" (תהלים פ"ג), ומה היא ההמיה? "על עמך יערימו סוד" (שם), ועל זה קורא אותם אויביו של מקום:}%
% ---------- פרק 15 | פסוק 8 ----------
 \textbf{ח} וּבְר֤וּחַ אַפֶּ֙יךָ֙ נֶ֣עֶרְמוּ מַ֔יִם        נִצְּב֥וּ כְמוֹ־נֵ֖ד נֹזְלִ֑ים        קָֽפְא֥וּ תְהֹמֹ֖ת בְּלֶב־יָֽם׃ \Rashi{\textbf{וברוח אפיך.}היוצא משני נחירים של אף. דבר הכתוב כביכול בשכינה דגמת מלך בשר ודם, כדי להשמיע אזן הבריות כפי ההוה, שיוכלו להבין דבר; כשאדם כועס יוצא רוח מנחיריו; וכן "עלה עשן באפו" (תהלים י"ח), וכן "ומרוח אפו יכלו" (איוב ד'), וזהו שאמר "למען שמי אאריך אפי", (ישעיהו מ"ח) – כשזעפו נח, נשימתו ארכה, וכשהוא כועס נשימתו קצרה, "ותהלתי אחטם לך" (שם) – ולמען תהלתי אשים חטם באפי לסתם נחירי בפני האף והרוח שלא יצאו, לך – בשבילך, אחטם, כמו "נאקה בחטם" במסכת שבת, כך נראה בעיני. וכל חרון אף שבמקרא אני אומר כן; חרה אף כמו "ועצמי חרה מני חרב" (איוב ל'), לשון שרפה ומוקד, שהנחירים מתחממים ונחרים בעת הקצף, וחרון מגזרת חרה כמו רצון מגזרת רצה, וכן חמה לשון חמימות, על כן הוא אומר "וחמתו בערה בו" (אסתר א'), ובנוח החמה אומר "נתקררה" דעתו: נערמו מים. אנקלוס תרגם לשון ערמימות, ולשון צחות המקרא כמו "ערמת חטים" (שיר השירים ז'), ונצבו כמו נד יוכיח: נערמו. ממוקד רוח שיצא מאפך יבשו המים, והם נעשו כמין גלים וכריות של ערמה, שהם גבוהים: כמו נד. כתרגומו "כשור" – כחומה: נד. לשון צבור וכנוס, כמו "נד קציר ביום נחלה" (ישעיהו י״ז:י״א), "כונס כנד" (תהילים ל״ג:ז׳), לא כתוב כונס כנאד אלא כנד, ואלו היה כנד כמו כנאד וכונס לשון הכנסה, היה לו לכתב מכניס כבנאד מי הים, אלא כונס לשון אוסף וצובר הוא, וכן "קמו נד אחד" (יהושע ג׳:י״ג), "ויעמדו נד אחד" (שם), ואין לשון קימה ועמידה בנאדות אלא בחומות וצבורים, ולא מצינו נאד נקוד אלא במלאפום כמו "שימה דמעתי בנאדך" (תהילים נ״ו:ט׳), "את נאד החלב" (שופטים ד׳:י״ט): קפאו. כמו "וכגבנה תקפיאני" (איוב י׳:י׳), שהקשו ונעשו כאבנים והמים זורקים את המצריים על האבן בכח ונלחמים בם בכל מיני קשי: בלב ים. בחזק הים; ודרך המקראות לדבר כן, "עד לב השמים" (דברים ד'), "בלב האלה" (שמואל ב י"ח), לשון עקרו ותקפו של דבר:}%
% ---------- פרק 15 | פסוק 9 ----------
 \textbf{ט} אָמַ֥ר אוֹיֵ֛ב אֶרְדֹּ֥ף אַשִּׂ֖יג        אֲחַלֵּ֣ק שָׁלָ֑ל תִּמְלָאֵ֣מוֹ נַפְשִׁ֔י        אָרִ֣יק חַרְבִּ֔י תּוֹרִישֵׁ֖מוֹ יָדִֽי׃ \Rashi{\textbf{אמר אויב.}לעמו, כשפתם בדברים, ארדף ואשיגם ואחלק שלל עם שרי ועבדי: תמלאמו. תתמלא מהם, נפשי רוחי ורצוני. ואל תתמה על תבה המדברת בשתים, תמלאמו – תמלא מהם, יש הרבה בלשון הזה "כי ארץ הנגב נתתני" (שופטים א׳:ט״ו), כמו נתת לי, "ולא יכלו דברו לשלום" (בראשית ל״ז:ד׳), כמו דבר עמו, "בני יצאוני" (ירמיהו י׳:כ׳), כמו יצאו ממני, "מספר צעדי אגידנו" (איוב ל״א:ל״ז), כמו אגיד לו, אף כאן "תמלאמו" – תמלא נפשי מהם: אריק חרבי. אשלף, ועל שם שהוא מריק את התער בשליפתו ונשאר ריק, נופל בו לשון הרקה, כמו "מריקים שקיהם" (בראשית מ״ב:ל״ה), "וכליו יריקו" (ירמיהו מ״ח:י״ב). ואל תאמר אין לשון ריקות נופל על היוצא, אלא על התיק ועל השק ועל הכלי שיצא ממנה, אבל לא על החרב ועל היין, ולדחק ולפרש "אריק חרבי" כלשון "וירק את חניכיו" (בראשית י״ד:י״ד), אזדין בחרבי, מצינו הלשון מוסב אף על היוצא, "שמן תורק" (שיר השירים א׳:ג׳), "ולא הורק מכלי אל כלי" (ירמיהו מ״ח:י״א), לא הורק הכלי אין כתיב כאן אלא לא הורק היין מכלי אל כלי, מצינו הלשון מוסב על היין, וכן "והריקו חרבותם על יפי חכמתך", דחירם (יחזקאל כ״ח:ז׳): תורישמו. לשון רישות ודלות, כמו "מוריש ומעשיר" (שמואל א' ב'):}%
% ---------- פרק 15 | פסוק 10 ----------
 \textbf{י} נָשַׁ֥פְתָּ בְרוּחֲךָ֖ כִּסָּ֣מוֹ יָ֑ם        צָֽלְלוּ֙ כַּֽעוֹפֶ֔רֶת בְּמַ֖יִם אַדִּירִֽים׃ \Rashi{\textbf{נשפת.}לשון הפחה, וכן "וגם נשף בהם" (ישעיה מ'): צללו. שקעו, עמקו, לשון מצולה: כעופרת. אבר, פלו"ם בלעז:}%
% ---------- פרק 15 | פסוק 11 ----------
 \textbf{יא} מִֽי־כָמֹ֤כָה בָּֽאֵלִם֙ יְהֹוָ֔ה        מִ֥י כָּמֹ֖כָה נֶאְדָּ֣ר בַּקֹּ֑דֶשׁ        נוֹרָ֥א תְהִלֹּ֖ת עֹ֥שֵׂה פֶֽלֶא׃ \Rashi{\textbf{באלם.}בחזקים, כמו "ואת אילי הארץ לקח" (יחזקאל י"ז). "אילותי לעזרתי חושה" (תהלים כ"ב): נורא תהלת. יראוי מלהגיד תהלותיך פן ימעטו, כמו שכתוב "לך דומיה תהלה" (תהילים ס״ה:ב׳):}%
% ---------- פרק 15 | פסוק 12 ----------
 \textbf{יב} נָטִ֙יתָ֙ יְמִ֣ינְךָ֔ תִּבְלָעֵ֖מוֹ אָֽרֶץ׃ \Rashi{\textbf{נטית ימינך.}כשהקב"ה נוטה ידו, הרשעים כלים ונופלים, לפי שהכל נתון בידו ונופלים בהטיתה; וכן הוא אומר "וה' יטה ידו וכשל עוזר ונפל עזר" (ישעיהו ל"א), משל לכלי זכוכית הנתונים ביד אדם, מטה ידו מעט והן נופלים ומשתברין (מכילתא): תבלעמו ארץ. מכאן שזכו לקבורה בשכר שאמרו ה' הצדיק (מכילתא):}%
% ---------- פרק 15 | פסוק 13 ----------
 \textbf{יג} נָחִ֥יתָ בְחַסְדְּךָ֖ עַם־ז֣וּ גָּאָ֑לְתָּ        נֵהַ֥לְתָּ בְעׇזְּךָ֖ אֶל־נְוֵ֥ה קׇדְשֶֽׁךָ׃ \Rashi{\textbf{נהלת.}לשון מנהל, ואנקלוס תרגם לשון נושא וסובל, ולא דקדק לפרש אחר לשון העברית:}%
% ---------- פרק 15 | פסוק 14 ----------
 \textbf{יד} שָֽׁמְע֥וּ עַמִּ֖ים יִרְגָּז֑וּן        חִ֣יל אָחַ֔ז יֹשְׁבֵ֖י פְּלָֽשֶׁת׃ \Rashi{\textbf{ירגזון.}מתרגזין: ישבי פלשת. מפני שהרגו את בני אפרים – שמהרו את הקץ ויצאו בחזקה כמפרש בדברי הימים – והרגום אנשי גת (מכילתא):}%
% ---------- פרק 15 | פסוק 15 ----------
 \textbf{טו} אָ֤ז נִבְהֲלוּ֙ אַלּוּפֵ֣י אֱד֔וֹם        אֵילֵ֣י מוֹאָ֔ב יֹֽאחֲזֵ֖מוֹ רָ֑עַד        נָמֹ֕גוּ כֹּ֖ל יֹשְׁבֵ֥י כְנָֽעַן׃ \Rashi{\textbf{אלופי אדום אילי מואב.}והלא לא היה להם לירא כלום, שהרי לא עליהם הולכים? אלא מפני אנינות שהיו מתאוננים ומצטערים על כבודם של ישראל (ילקוט שמעוני): נמגו. נמסו, כמו "ברביבים תמוגגנה" (תהלים ס"ה), אמרו, עלינו הם באים, לכלותינו ולירש את ארצנו (מכילתא):}%
% ---------- פרק 15 | פסוק 16 ----------
 \textbf{טז} תִּפֹּ֨ל עֲלֵיהֶ֤ם אֵימָ֙תָה֙ וָפַ֔חַד        בִּגְדֹ֥ל זְרוֹעֲךָ֖ יִדְּמ֣וּ כָּאָ֑בֶן        עַד־יַעֲבֹ֤ר עַמְּךָ֙ יְהֹוָ֔ה        עַֽד־יַעֲבֹ֖ר עַם־ז֥וּ קָנִֽיתָ׃ \Rashi{\textbf{תפל עליהם אימתה.}על הרחוקים: ופחד. על הקרובים, כענין שנאמר "כי שמענו את אשר הוביש וגו'" (יהושע ב'): עד יעבר, עד יעבר. כתרגומו: קנית. חבבת משאר אמות, כחפץ הקנוי בדמים יקרים שחביב על האדם:}%
% ---------- פרק 15 | פסוק 17 ----------
 \textbf{יז} תְּבִאֵ֗מוֹ וְתִטָּעֵ֙מוֹ֙ בְּהַ֣ר נַחֲלָֽתְךָ֔        מָכ֧וֹן לְשִׁבְתְּךָ֛ פָּעַ֖לְתָּ יְהֹוָ֑ה        מִקְּדָ֕שׁ אֲדֹנָ֖י כּוֹנְנ֥וּ יָדֶֽיךָ׃ \Rashi{\textbf{תבאמו.}נתנבא משה שלא יכנס לארץ לכך לא נאמר "תביאנו" (בבא בתרא קי"ט): מכון לשבתך. מקדש של מטה מכון כנגד כסא של מעלה אשר פעלת (מכילתא): מקדש ה'. הטעם עליו זקף גדול, להפרידו מתבת השם שלאחריו, המקדש אשר כוננו ידיך ה'. חביב בית המקדש, שהעולם נברא ביד אחת, שנאמר "אף ידי יסדה ארץ" (ישעיהו מ"ח), ומקדש בשתי ידים, ואימתי יבנה בשתי ידים? בזמן ש"ה' ימלך לעלם ועד" – לעתיד לבא שכל המלוכה שלו:}%
% ---------- פרק 15 | פסוק 18 ----------
 \textbf{יח} יְהֹוָ֥ה ׀ יִמְלֹ֖ךְ לְעֹלָ֥ם וָעֶֽד׃ \Rashi{\textbf{לעלם ועד.}לשון עולמות הוא והוי"ו בו יסוד, לפיכך היא פתוחה, אבל "אנכי היודע ועד" (ירמיהו כ"ט), שהוי"ו בו שמוש, קמוצה היא:}%
% ---------- פרק 15 | פסוק 19 ----------
 \textbf{יט} כִּ֣י בָא֩ ס֨וּס פַּרְעֹ֜ה בְּרִכְבּ֤וֹ וּבְפָרָשָׁיו֙ בַּיָּ֔ם        וַיָּ֧שֶׁב יְהֹוָ֛ה עֲלֵהֶ֖ם אֶת־מֵ֣י הַיָּ֑ם        וּבְנֵ֧י יִשְׂרָאֵ֛ל הָלְכ֥וּ בַיַּבָּשָׁ֖ה בְּת֥וֹךְ הַיָּֽם׃ \{פ\} \Rashi{\textbf{כי בא סוס פרעה.}כאשר בא:}%
% ---------- פרק 15 | פסוק 20 ----------
 \textbf{כ} וַתִּקַּח֩ מִרְיָ֨ם הַנְּבִיאָ֜ה אֲח֧וֹת אַהֲרֹ֛ן אֶת־הַתֹּ֖ף בְּיָדָ֑הּ וַתֵּצֶ֤אןָ כׇֽל־הַנָּשִׁים֙ אַחֲרֶ֔יהָ בְּתֻפִּ֖ים וּבִמְחֹלֹֽת׃ \Rashi{\textbf{ותקח מרים הנביאה.}היכן נתנבאה? כשהיתה אחות אהרן, קדם שנולד משה, אמרה עתידה אמי שתלד בן וכו' כדאיתא בסוטה. (דף י"ב). ד"א — אחות אהרן, לפי שמסר נפשו עליה כשנצטרעה נקראת על שמו: את התף. כלי של מיני זמר: בתפים ובמחלת. מבטחות היו צדקניות שבדור שהקב"ה עושה להם נסים והוציאו תפים ממצרים (מכילתא):}%
% ---------- פרק 15 | פסוק 21 ----------
 \textbf{כא} וַתַּ֥עַן לָהֶ֖ם מִרְיָ֑ם שִׁ֤ירוּ לַֽיהֹוָה֙ כִּֽי־גָאֹ֣ה גָּאָ֔ה ס֥וּס וְרֹכְב֖וֹ רָמָ֥ה בַיָּֽם׃ \{ס\} \Rashi{\textbf{ותען להם מרים.}משה אמר שירה לאנשים – הוא אומר והם עונין אחריו – ומרים אמרה שירה לנשים (סוטה ל'):}%
% ---------- פרק 15 | פסוק 22 ----------
 \textbf{כב} וַיַּסַּ֨ע מֹשֶׁ֤ה אֶת־יִשְׂרָאֵל֙ מִיַּם־ס֔וּף וַיֵּצְא֖וּ אֶל־מִדְבַּר־שׁ֑וּר וַיֵּלְכ֧וּ שְׁלֹֽשֶׁת־יָמִ֛ים בַּמִּדְבָּ֖ר וְלֹא־מָ֥צְאוּ מָֽיִם׃ \Rashi{\textbf{ויסע משה.}הסיען בעל כרחם, שעטרו מצרים את סוסיהם בתכשיטי זהב וכסף ואבנים טובות, והיו ישראל מוצאין אותן בים – וגדולה היתה בזת הים מבזת מצרים שנאמר "תורי זהב נעשה לך עם נקדות הכסף" (שיר השירים א') – לפיכך הצרך להסיען בעל כרחם:}%
% ---------- פרק 15 | פסוק 23 ----------
 \textbf{כג} וַיָּבֹ֣אוּ מָרָ֔תָה וְלֹ֣א יָֽכְל֗וּ לִשְׁתֹּ֥ת מַ֙יִם֙ מִמָּרָ֔ה כִּ֥י מָרִ֖ים הֵ֑ם עַל־כֵּ֥ן קָרָֽא־שְׁמָ֖הּ מָרָֽה׃ \Rashi{\textbf{ויבאו מרתה.}כמו למרה, ה"א בסוף תבה במקום למ"ד בתחלתה, והתי"ו היא במקום ה"א הנשרשת בתבת מרה, ובסמיכתה, כשהיא נדבקת לה"א שהוא מוסיף במקום הלמ"ד, תהפך הה"א של שרש לתי"ו; וכן כל ה"א שהיא שרש בתבה תתהפך לתי"ו בסמיכתה, כמו "חמה אין לי" (ישעיהו כ"ז), "וחמתו בערה בו" (אסתר א'), הרי ה"א של שרש נהפכת לתי"ו מפני שנסמכת אל הוי"ו הנוספת, וכן "עבד ואמה" (ויקרא כ"ה), "הנה אמתי בלהה" (בראשית ל'), "לנפש חיה" (בראשית ב׳:ז׳), "וזהמתו חיתו לחם" (איוב ל"ג), "בין הרמה" (שופטים ד'), "ותשבתו הרמתה" (שמואל א ז'):}%
% ---------- פרק 15 | פסוק 24 ----------
 \textbf{כד} וַיִּלֹּ֧נוּ הָעָ֛ם עַל־מֹשֶׁ֥ה לֵּאמֹ֖ר מַה־נִּשְׁתֶּֽה׃ \Rashi{\textbf{וילנו.}לשון נפעל הוא, וכן התרגום לשון נפעל הוא ואתרעמו, וכן דרך לשון תלונה להסב הדבור אל האדם – מתלונן, מתרועם ולא אמר לונן, רועם, וכן יאמר הלועז דקומפ"לישנ"ט ש"י בלעז, מסב הדבור אליו באמרו ש"י:}%
% ---------- פרק 15 | פסוק 25 ----------
 \textbf{כה} וַיִּצְעַ֣ק אֶל־יְהֹוָ֗ה וַיּוֹרֵ֤הוּ יְהֹוָה֙ עֵ֔ץ וַיַּשְׁלֵךְ֙ אֶל־הַמַּ֔יִם וַֽיִּמְתְּק֖וּ הַמָּ֑יִם שָׁ֣ם שָׂ֥ם ל֛וֹ חֹ֥ק וּמִשְׁפָּ֖ט וְשָׁ֥ם נִסָּֽהוּ׃ \Rashi{\textbf{שם שם לו.}במרה נתן להם מקצת פרשיות של תורה שיתעסקו בהם, שבת ופרה אדמה ודינין (סנהדרין נ"ו): ושם נסהו. לעם, וראה קשי ערפו, שלא נמלכו במשה בלשון יפה "בקש עלינו רחמים שיהיה לנו מים לשתות" אלא נתלוננו:}%
% ---------- פרק 15 | פסוק 26 ----------
 \textbf{כו} וַיֹּ֩אמֶר֩ אִם־שָׁמ֨וֹעַ תִּשְׁמַ֜ע לְק֣וֹל ׀ יְהֹוָ֣ה אֱלֹהֶ֗יךָ וְהַיָּשָׁ֤ר בְּעֵינָיו֙ תַּעֲשֶׂ֔ה וְהַֽאֲזַנְתָּ֙ לְמִצְוֺתָ֔יו וְשָׁמַרְתָּ֖ כׇּל־חֻקָּ֑יו כׇּֽל־הַמַּחֲלָ֞ה אֲשֶׁר־שַׂ֤מְתִּי בְמִצְרַ֙יִם֙ לֹא־אָשִׂ֣ים עָלֶ֔יךָ כִּ֛י אֲנִ֥י יְהֹוָ֖ה רֹפְאֶֽךָ׃ \{ס\} \Rashi{\textbf{אם שמוע תשמע.}זו קבלה שיקבלו עליהם: תעשה. היא עשיה: והאזנת. תטה אזנים לדקדק בהם: כל חקיו. דברים שאינן אלא גזרת מלך, בלא שום טעם, ויצר הרע מקנטר עליהם: מה אסור באלו? למה נאסרו? כגון לבישת כלאים ואכילת חזיר ופרה אדמה וכיוצא בהם (יומא ס"ז): לא אשים עליך. ואם אשים הרי הוא כלא הושמה, כי אני ה' רפאך – זהו מדרשו. ולפי פשוטו כי אני ה' רפאך ומלמדך תורה ומצוות למען תנצל מהם, כרופא הזה האומר לאדם, אל תאכל דבר זה פן יביאך לידי חלי זה, וכן הוא אומר "רפאות תהי לשרך" (משלי ג'):}%
% ---------- פרק 15 | פסוק 27 ----------
 \textbf{כז} וַיָּבֹ֣אוּ אֵילִ֔מָה וְשָׁ֗ם שְׁתֵּ֥ים עֶשְׂרֵ֛ה עֵינֹ֥ת מַ֖יִם וְשִׁבְעִ֣ים תְּמָרִ֑ים וַיַּחֲנוּ־שָׁ֖ם עַל־הַמָּֽיִם׃ \Rashi{\textbf{שתים עשרה עינת מים.}כנגד י"ב שבטים נזדמנו להם: ושבעים תמרים. כנגד שבעים זקנים (מכילתא):}%
% ---------- פרק 16 | פסוק 1 ----------
 \textbf{א} וַיִּסְעוּ֙ מֵֽאֵילִ֔ם וַיָּבֹ֜אוּ כׇּל־עֲדַ֤ת בְּנֵֽי־יִשְׂרָאֵל֙ אֶל־מִדְבַּר־סִ֔ין אֲשֶׁ֥ר בֵּין־אֵילִ֖ם וּבֵ֣ין סִינָ֑י בַּחֲמִשָּׁ֨ה עָשָׂ֥ר יוֹם֙ לַחֹ֣דֶשׁ הַשֵּׁנִ֔י לְצֵאתָ֖ם מֵאֶ֥רֶץ מִצְרָֽיִם׃ \Rashi{\textbf{בחמשה עשר יום.}נתפרש היום של חניה זו, לפי שבו ביום כלתה החררה שהוציאו ממצרים והצרכו למן, ללמדנו, שאכלו משירי הבצק ששים ואחת סעודות וירד להם מן בששה עשר באיר, ויום א' בשבת היה כדאיתא במסכת שבת (פ"ז):}%
% ---------- פרק 16 | פסוק 2 ----------
 \textbf{ב} (וילינו) [וַיִּלּ֜וֹנוּ] כׇּל־עֲדַ֧ת בְּנֵי־יִשְׂרָאֵ֛ל עַל־מֹשֶׁ֥ה וְעַֽל־אַהֲרֹ֖ן בַּמִּדְבָּֽר׃ \Rashi{\textbf{וילונו.}לפי שכלה הלחם:}%
% ---------- פרק 16 | פסוק 3 ----------
 \textbf{ג} וַיֹּאמְר֨וּ אֲלֵהֶ֜ם בְּנֵ֣י יִשְׂרָאֵ֗ל מִֽי־יִתֵּ֨ן מוּתֵ֤נוּ בְיַד־יְהֹוָה֙ בְּאֶ֣רֶץ מִצְרַ֔יִם בְּשִׁבְתֵּ֙נוּ֙ עַל־סִ֣יר הַבָּשָׂ֔ר בְּאׇכְלֵ֥נוּ לֶ֖חֶם לָשֹׂ֑בַע כִּֽי־הוֹצֵאתֶ֤ם אֹתָ֙נוּ֙ אֶל־הַמִּדְבָּ֣ר הַזֶּ֔ה לְהָמִ֛ית אֶת־כׇּל־הַקָּהָ֥ל הַזֶּ֖ה בָּרָעָֽב׃ \{ס\} \Rashi{\textbf{מי יתן מותנו.}שנמות; ואינו שם דבר כמו מיתתנו, אלא כמו עשותנו חנותנו שובנו, לעשות אנחנו, לחנות אנחנו, למות אנחנו. תרגומו לוי דמיתנא, "לו מתנו" (במדבר יד, ב) — הלואי והיינו מתים:}%
% ---------- פרק 16 | פסוק 4 ----------
 \textbf{ד} וַיֹּ֤אמֶר יְהֹוָה֙ אֶל־מֹשֶׁ֔ה הִנְנִ֨י מַמְטִ֥יר לָכֶ֛ם לֶ֖חֶם מִן־הַשָּׁמָ֑יִם וְיָצָ֨א הָעָ֤ם וְלָֽקְטוּ֙ דְּבַר־י֣וֹם בְּיוֹמ֔וֹ לְמַ֧עַן אֲנַסֶּ֛נּוּ הֲיֵלֵ֥ךְ בְּתוֹרָתִ֖י אִם־לֹֽא׃ \Rashi{\textbf{דבר יום ביומו.}צרך אכילת יום ילקטו ביומו ולא היום לצרך מחר (מכילתא): למען אנסנו הילך בתורתי. אם ישמרו מצוות התלויות בו, שלא יותירו ממנו ולא יצאו בשבת ללקט:}%
% ---------- פרק 16 | פסוק 5 ----------
 \textbf{ה} וְהָיָה֙ בַּיּ֣וֹם הַשִּׁשִּׁ֔י וְהֵכִ֖ינוּ אֵ֣ת אֲשֶׁר־יָבִ֑יאוּ וְהָיָ֣ה מִשְׁנֶ֔ה עַ֥ל אֲשֶֽׁר־יִלְקְט֖וּ י֥וֹם ׀ יֽוֹם׃ \Rashi{\textbf{והיה משנה.}ליום ולמחרת: משנה. על שהיו רגילים ללקט יום יום של שאר ימות השבוע; אומר אני "אשר יביאו והיה משנה" – לאחר שיביאו, ימצאו משנה במדידה על אשר ילקטו וימדו יום יום, וזהו לקטו לחם משנה – בלקיטתו היה נמצא לחם משנה. וזהו על כן הוא נותן לכם ביום הששי לחם יומים – נותן לכם ברכה, פוי"שן, בבית למלאת העמר פעמים ללחם יומים:}%
% ---------- פרק 16 | פסוק 6 ----------
 \textbf{ו} וַיֹּ֤אמֶר מֹשֶׁה֙ וְאַהֲרֹ֔ן אֶֽל־כׇּל־בְּנֵ֖י יִשְׂרָאֵ֑ל עֶ֕רֶב וִֽידַעְתֶּ֕ם כִּ֧י יְהֹוָ֛ה הוֹצִ֥יא אֶתְכֶ֖ם מֵאֶ֥רֶץ מִצְרָֽיִם׃ \Rashi{\textbf{ערב.}כמו לערב: וידעתם כי ה' הוציא אתכם מארץ מצרים. לפי שאמרתם לנו כי הוצאתם אותנו, תדעו כי לא אנחנו המוציאים אלא ה' הוציא אתכם שיגיז לכם את השלו:}%
% ---------- פרק 16 | פסוק 7 ----------
 \textbf{ז} וּבֹ֗קֶר וּרְאִיתֶם֙ אֶת־כְּב֣וֹד יְהֹוָ֔ה בְּשׇׁמְע֥וֹ אֶת־תְּלֻנֹּתֵיכֶ֖ם עַל־יְהֹוָ֑ה וְנַ֣חְנוּ מָ֔ה כִּ֥י (תלונו) [תַלִּ֖ינוּ] עָלֵֽינוּ׃ \Rashi{\textbf{ובקר וראיתם.}לא על הכבוד שנאמר "והנה כבוד ה' נראה בענן" נאמר, אלא כך אמר להם, ערב וידעתם כי היכלת בידו לתן תאותכם, ובשר יתן, אך לא בפנים מאירות יתננה לכם, כי שלא כהגן שאלתם אותו ומכרס מלאה, והלחם ששאלתם לצרך, בירידתו לבקר תראו את כבוד אור פניו – שיורידהו לכם דרך חבה, בבקר, שיש שעות להכינו וטל מלמעלה וטל מלמטה, כמנח בקפסא (יומא ע"ה): את תלנתיכם על ה'. כמו אשר על ה': ונחנו מה. מה אנחנו חשובין? כי תלינו עלינו. שתרעימו עלינו את הכל, ואת בניכם ונשיכם ובנותיכם וערב רב. ועל כרחי אני זקוק לפרש תלינו בלשון תפעילו מפני דגשותו וקריתו, שאלו היה רפה, הייתי מפרשו בלשון תפעלו, כמו "וילן העם על משה" (שמות י"ז), או אם היה דגוש ואין בו יו"ד ונקרא תלונו, הייתי מפרשו לשון תתלוננו, עכשיו הוא משמע תלינו את אחרים, כמו במרגלים "וילינו עליו את כל העדה" (במדבר י"ד):}%
% ---------- פרק 16 | פסוק 8 ----------
 \textbf{ח} וַיֹּ֣אמֶר מֹשֶׁ֗ה בְּתֵ֣ת יְהֹוָה֩ לָכֶ֨ם בָּעֶ֜רֶב בָּשָׂ֣ר לֶאֱכֹ֗ל וְלֶ֤חֶם בַּבֹּ֙קֶר֙ לִשְׂבֹּ֔עַ בִּשְׁמֹ֤עַ יְהֹוָה֙ אֶת־תְּלֻנֹּ֣תֵיכֶ֔ם אֲשֶׁר־אַתֶּ֥ם מַלִּינִ֖ם עָלָ֑יו וְנַ֣חְנוּ מָ֔ה לֹא־עָלֵ֥ינוּ תְלֻנֹּתֵיכֶ֖ם כִּ֥י עַל־יְהֹוָֽה׃ \Rashi{\textbf{בשר לאכל.}ולא לשבע; למדה תורה דרך ארץ שאין אוכלין בשר לשבע. ומה ראה להוריד לחם בבקר ובשר בערב? לפי שהלחם שאלו כהגן, שאי אפשר לו לאדם בלא לחם, אבל בשר שאלו שלא כהגן, שהרבה בהמות היו להם, ועוד שהיה אפשר להם בלא בשר, לפיכך נתן להם בשעת טרח שלא כהגן (יומא ע"ה): אשר אתם מלינם עליו. את האחרים השומעים אתכם מתלוננים:}%
% ---------- פרק 16 | פסוק 9 ----------
 \textbf{ט} וַיֹּ֤אמֶר מֹשֶׁה֙ אֶֽל־אַהֲרֹ֔ן אֱמֹ֗ר אֶֽל־כׇּל־עֲדַת֙ בְּנֵ֣י יִשְׂרָאֵ֔ל קִרְב֖וּ לִפְנֵ֣י יְהֹוָ֑ה כִּ֣י שָׁמַ֔ע אֵ֖ת תְּלֻנֹּתֵיכֶֽם׃ \Rashi{\textbf{קרבו.}למקום שהענן ירד:}%
% ---------- פרק 16 | פסוק 10 ----------
 \textbf{י} וַיְהִ֗י כְּדַבֵּ֤ר אַהֲרֹן֙ אֶל־כׇּל־עֲדַ֣ת בְּנֵֽי־יִשְׂרָאֵ֔ל וַיִּפְנ֖וּ אֶל־הַמִּדְבָּ֑ר וְהִנֵּה֙ כְּב֣וֹד יְהֹוָ֔ה נִרְאָ֖ה בֶּעָנָֽן׃ \{פ\} %
% ---------- פרק 16 | פסוק 11 ----------
 \textbf{יא} וַיְדַבֵּ֥ר יְהֹוָ֖ה אֶל־מֹשֶׁ֥ה לֵּאמֹֽר׃ %
% ---------- פרק 16 | פסוק 12 ----------
 \textbf{יב} שָׁמַ֗עְתִּי אֶת־תְּלוּנֹּת֮ בְּנֵ֣י יִשְׂרָאֵל֒ דַּבֵּ֨ר אֲלֵהֶ֜ם לֵאמֹ֗ר בֵּ֤ין הָֽעַרְבַּ֙יִם֙ תֹּאכְל֣וּ בָשָׂ֔ר וּבַבֹּ֖קֶר תִּשְׂבְּעוּ־לָ֑חֶם וִֽידַעְתֶּ֕ם כִּ֛י אֲנִ֥י יְהֹוָ֖ה אֱלֹהֵיכֶֽם׃ %
% ---------- פרק 16 | פסוק 13 ----------
 \textbf{יג} וַיְהִ֣י בָעֶ֔רֶב וַתַּ֣עַל הַשְּׂלָ֔ו וַתְּכַ֖ס אֶת־הַֽמַּחֲנֶ֑ה וּבַבֹּ֗קֶר הָֽיְתָה֙ שִׁכְבַ֣ת הַטַּ֔ל סָבִ֖יב לַֽמַּחֲנֶֽה׃ \Rashi{\textbf{השלו.}מין עוף ושמן מאד (יומא ע"ה): היתה שכבת הטל. הטל שוכב על המן, ובמקום אחר הוא אומר (במדבר י"א) "וברדת הטל וגו'", הטל יורד על הארץ, והמן יורד עליו, וחוזר ויורד טל עליו, הרי הוא כמנח בקפסא (יומא שם):}%
% ---------- פרק 16 | פסוק 14 ----------
 \textbf{יד} וַתַּ֖עַל שִׁכְבַ֣ת הַטָּ֑ל וְהִנֵּ֞ה עַל־פְּנֵ֤י הַמִּדְבָּר֙ דַּ֣ק מְחֻסְפָּ֔ס דַּ֥ק כַּכְּפֹ֖ר עַל־הָאָֽרֶץ׃ \Rashi{\textbf{ותעל שכבת הטל וגו'.}כשהחמה זורחת, עולה הטל שעל המן לקראת החמה, כדרך טל עולה לקראת החמה; אף אם תמלא שפופרת של ביצה טל ותסתם את פיה ותניחה בחמה היא עולה מאליה באויר. ורבותינו דרשו, שהטל עולה מן הארץ באויר, וכעלות שכבת הטל נתגלה המן וראו והנה על פני המדבר וגו': דק – דבר דק: מחספס. מגלה, ואין דומה לו במקרא. ויש לומר מחספס לשון חפיסה ודלוסקמא שבלשון משנה, כשנתגלה משכבת הטל ראו שהיה דבר דק מחספס בתוכו בין שתי שכבות הטל. ואנקלוס תרגם "מקלף", לשון מחשף הלבן: ככפר. כפור יילי"דא בלעז, דעדק כגיר, "כאבני גר" (ישעיהו כ"ז), והוא מין צבע שחר, כדאמרינן גבי כסוי הדם "הגיר והזרניך" (חולין פ"ח). "דעדק כגיר כגלידא על ארעא" – דק היה כגיר ושוכב מגלד כקרח על הארץ, וכן פרושו: דק ככפר – שטוח, קלוש ומחבר כגליד. דק טינב"ש בלעז, שהיה מגליד גלד דק מלמעלה; "וכגיר" שתרגם אנקלוס תוספת הוא על לשון העברית ואין לו תבה בפסוק:}%
% ---------- פרק 16 | פסוק 15 ----------
 \textbf{טו} וַיִּרְא֣וּ בְנֵֽי־יִשְׂרָאֵ֗ל וַיֹּ֨אמְר֜וּ אִ֤ישׁ אֶל־אָחִיו֙ מָ֣ן ה֔וּא כִּ֛י לֹ֥א יָדְע֖וּ מַה־ה֑וּא וַיֹּ֤אמֶר מֹשֶׁה֙ אֲלֵהֶ֔ם ה֣וּא הַלֶּ֔חֶם אֲשֶׁ֨ר נָתַ֧ן יְהֹוָ֛ה לָכֶ֖ם לְאׇכְלָֽה׃ \Rashi{\textbf{מן הוא.}הכנת מזון הוא, כמו "וימן להם המלך" (דניאל א'): כי לא ידעו מה הוא. שיקראוהו בשמו:}%
% ---------- פרק 16 | פסוק 16 ----------
 \textbf{טז} זֶ֤ה הַדָּבָר֙ אֲשֶׁ֣ר צִוָּ֣ה יְהֹוָ֔ה לִקְט֣וּ מִמֶּ֔נּוּ אִ֖ישׁ לְפִ֣י אׇכְל֑וֹ עֹ֣מֶר לַגֻּלְגֹּ֗לֶת מִסְפַּר֙ נַפְשֹׁ֣תֵיכֶ֔ם אִ֛ישׁ לַאֲשֶׁ֥ר בְּאׇהֳל֖וֹ תִּקָּֽחוּ׃ \Rashi{\textbf{עמר.}שם מדה: מספר נפשתיכם. כפי מנין נפשות שיש לאיש באהלו תקחו עמר לכל גלגלת:}%
% ---------- פרק 16 | פסוק 17 ----------
 \textbf{יז} וַיַּעֲשׂוּ־כֵ֖ן בְּנֵ֣י יִשְׂרָאֵ֑ל וַֽיִּלְקְט֔וּ הַמַּרְבֶּ֖ה וְהַמַּמְעִֽיט׃ \Rashi{\textbf{המרבה והממעיט.}יש שלקטו הרבה ויש שלקטו מעט, וכשבאו לביתם מדדו בעמר איש איש מה שלקטו, ומצאו שהמרבה ללקט לא העדיף על עמר לגלגלת אשר באהלו, והממעיט ללקט לא מצא חסר מעמר לגלגלת; וזהו נס גדול שנעשה בו:}%
% ---------- פרק 16 | פסוק 18 ----------
 \textbf{יח} וַיָּמֹ֣דּוּ בָעֹ֔מֶר וְלֹ֤א הֶעְדִּיף֙ הַמַּרְבֶּ֔ה וְהַמַּמְעִ֖יט לֹ֣א הֶחְסִ֑יר אִ֥ישׁ לְפִֽי־אׇכְל֖וֹ לָקָֽטוּ׃ %
% ---------- פרק 16 | פסוק 19 ----------
 \textbf{יט} וַיֹּ֥אמֶר מֹשֶׁ֖ה אֲלֵהֶ֑ם אִ֕ישׁ אַל־יוֹתֵ֥ר מִמֶּ֖נּוּ עַד־בֹּֽקֶר׃ %
% ---------- פרק 16 | פסוק 20 ----------
 \textbf{כ} וְלֹא־שָׁמְע֣וּ אֶל־מֹשֶׁ֗ה וַיּוֹתִ֨רוּ אֲנָשִׁ֤ים מִמֶּ֙נּוּ֙ עַד־בֹּ֔קֶר וַיָּ֥רֻם תּוֹלָעִ֖ים וַיִּבְאַ֑שׁ וַיִּקְצֹ֥ף עֲלֵהֶ֖ם מֹשֶֽׁה׃ \Rashi{\textbf{ויותרו אנשים.}דתן ואבירם: וירם תולעים. לשון רמה: ויבאש. הרי זה מקרא הפוך, שבתחלה הבאיש ולבסוף התליע, כענין שנאמר "ולא הבאיש ורמה לא היתה בו", וכן דרך כל המתליעים (מכילתא):}%
% ---------- פרק 16 | פסוק 21 ----------
 \textbf{כא} וַיִּלְקְט֤וּ אֹתוֹ֙ בַּבֹּ֣קֶר בַּבֹּ֔קֶר אִ֖ישׁ כְּפִ֣י אׇכְל֑וֹ וְחַ֥ם הַשֶּׁ֖מֶשׁ וְנָמָֽס׃ \Rashi{\textbf{וחם השמש ונמס.}הנשאר בשדה נעשה נחלים ושותין ממנו אילים וצבאים, ואמות העולם צדין מהם וטועמים בהם טעם מן, ויודעים מה שבחן של ישראל: ונמס. פשר, לשון פושרים, על ידי השמש מתחמם ומפשיר, דישטנ"פריר בלעז, ודגמתו בסנהדרין בסוף ד' מיתות:}%
% ---------- פרק 16 | פסוק 22 ----------
 \textbf{כב} וַיְהִ֣י ׀ בַּיּ֣וֹם הַשִּׁשִּׁ֗י לָֽקְט֥וּ לֶ֙חֶם֙ מִשְׁנֶ֔ה שְׁנֵ֥י הָעֹ֖מֶר לָאֶחָ֑ד וַיָּבֹ֙אוּ֙ כׇּל־נְשִׂיאֵ֣י הָֽעֵדָ֔ה וַיַּגִּ֖ידוּ לְמֹשֶֽׁה׃ \Rashi{\textbf{לקטו לחם משנה.}כשמדדו את לקיטתם באהליהם מצאו כפלים, שני העמר לאחד. ומדרש אגדה: לחם משנה – משנה, אותו היום נשתנה לשבח בריחו וטעמו (שם): ויגידו למשה. שאלוהו מה היום מיומים; ומכאן יש ללמד שעדין לא הגיד להם משה פרשת שבת שנצטוה לומר להם והיה ביום הששי והכינו וגו', עד ששאלו מה זאת, אמר להם הוא אשר דבר ה' שנצטויתי לומר לכם, ולכך ענשו הכתוב שאמר לו עד אנה מאנתם, ולא הוציאו מן הכלל:}%
% ---------- פרק 16 | פסוק 23 ----------
 \textbf{כג} וַיֹּ֣אמֶר אֲלֵהֶ֗ם ה֚וּא אֲשֶׁ֣ר דִּבֶּ֣ר יְהֹוָ֔ה שַׁבָּת֧וֹן שַׁבַּת־קֹ֛דֶשׁ לַֽיהֹוָ֖ה מָחָ֑ר אֵ֣ת אֲשֶׁר־תֹּאפ֞וּ אֵפ֗וּ וְאֵ֤ת אֲשֶֽׁר־תְּבַשְּׁלוּ֙ בַּשֵּׁ֔לוּ וְאֵת֙ כׇּל־הָ֣עֹדֵ֔ף הַנִּ֧יחוּ לָכֶ֛ם לְמִשְׁמֶ֖רֶת עַד־הַבֹּֽקֶר׃ \Rashi{\textbf{את אשר תאפו אפו.}מה שאתם רוצים לאפות בתנור, אפו היום הכל לשני ימים, ומה שאתם צריכים לבשל ממנו במים, בשלו היום. לשון אפיה נופל בלחם ולשון בשול בתבשיל: למשמרת. לגניזה:}%
% ---------- פרק 16 | פסוק 24 ----------
 \textbf{כד} וַיַּנִּ֤יחוּ אֹתוֹ֙ עַד־הַבֹּ֔קֶר כַּאֲשֶׁ֖ר צִוָּ֣ה מֹשֶׁ֑ה וְלֹ֣א הִבְאִ֔ישׁ וְרִמָּ֖ה לֹא־הָ֥יְתָה בּֽוֹ׃ %
% ---------- פרק 16 | פסוק 25 ----------
 \textbf{כה} וַיֹּ֤אמֶר מֹשֶׁה֙ אִכְלֻ֣הוּ הַיּ֔וֹם כִּֽי־שַׁבָּ֥ת הַיּ֖וֹם לַיהֹוָ֑ה הַיּ֕וֹם לֹ֥א תִמְצָאֻ֖הוּ בַּשָּׂדֶֽה׃ \Rashi{\textbf{ויאמר משה אכלהו היום וגו'.}שחרית שהיו רגילין לצאת וללקט באו לשאל: "אם נצא אם לאו", אמר להם את שבידכם אכלו, לערב חזרו לפניו ושאלוהו מהו לצאת? אמר להם שבת היום, ראה אותם דואגים שמא פסק המן ולא ירד עוד, אמר להם היום לא תמצאוהו, מה ת"ל היום? היום לא תמצאוהו אבל מחר תמצאוהו:}%
% ---------- פרק 16 | פסוק 26 ----------
 \textbf{כו} שֵׁ֥שֶׁת יָמִ֖ים תִּלְקְטֻ֑הוּ וּבַיּ֧וֹם הַשְּׁבִיעִ֛י שַׁבָּ֖ת לֹ֥א יִֽהְיֶה־בּֽוֹ׃ \Rashi{\textbf{וביום השביעי שבת.}שבת הוא, המן לא יהיה בו; ולא בא הכתוב אלא לרבות יום הכפורים וימים טובים:}%
% ---------- פרק 16 | פסוק 27 ----------
 \textbf{כז} וַֽיְהִי֙ בַּיּ֣וֹם הַשְּׁבִיעִ֔י יָצְא֥וּ מִן־הָעָ֖ם לִלְקֹ֑ט וְלֹ֖א מָצָֽאוּ׃ \{ס\} %
% ---------- פרק 16 | פסוק 28 ----------
 \textbf{כח} וַיֹּ֥אמֶר יְהֹוָ֖ה אֶל־מֹשֶׁ֑ה עַד־אָ֙נָה֙ מֵֽאַנְתֶּ֔ם לִשְׁמֹ֥ר מִצְוֺתַ֖י וְתוֹרֹתָֽי׃ \Rashi{\textbf{עד אנה מאנתם.}משל הדיוט הוא "בהדי הוצא לקי כרבא" – על ידי הרשעים מתגנין הכשרין:}%
% ---------- פרק 16 | פסוק 29 ----------
 \textbf{כט} רְא֗וּ כִּֽי־יְהֹוָה֮ נָתַ֣ן לָכֶ֣ם הַשַּׁבָּת֒ עַל־כֵּ֠ן ה֣וּא נֹתֵ֥ן לָכֶ֛ם בַּיּ֥וֹם הַשִּׁשִּׁ֖י לֶ֣חֶם יוֹמָ֑יִם שְׁב֣וּ ׀ אִ֣ישׁ תַּחְתָּ֗יו אַל־יֵ֥צֵא אִ֛ישׁ מִמְּקֹמ֖וֹ בַּיּ֥וֹם הַשְּׁבִיעִֽי׃ \Rashi{\textbf{ראו.}בעיניכם כי ה' בכבודו מזהיר אתכם על השבת, שהרי נס נעשה בכל ערב שבת לתת לכם לחם יומים: שבו איש תחתיו. מכאן סמכו חכמים ד' אמות ליוצא חוץ לתחום (ערובין נ"א): אל יצא איש ממקמו. אלו אלפים אמה של תחום שבת (מכילתא), ולא במפרש – שאין תחומין אלא מדברי סופרים – ועקרו של מקרא על לוקטי המן נאמר:}%
% ---------- פרק 16 | פסוק 30 ----------
 \textbf{ל} וַיִּשְׁבְּת֥וּ הָעָ֖ם בַּיּ֥וֹם הַשְּׁבִעִֽי׃ %
% ---------- פרק 16 | פסוק 31 ----------
 \textbf{לא} וַיִּקְרְא֧וּ בֵֽית־יִשְׂרָאֵ֛ל אֶת־שְׁמ֖וֹ מָ֑ן וְה֗וּא כְּזֶ֤רַע גַּד֙ לָבָ֔ן וְטַעְמ֖וֹ כְּצַפִּיחִ֥ת בִּדְבָֽשׁ׃ \Rashi{\textbf{והוא כזרע גד לבן.}עשב ששמו קול"יינדר, וזרע שלו עגל ואינו לבן, והמן היה לבן, ואינו נמשל לזרע גד אלא לענין העגול, כזרע גד היה – והוא לבן: כצפיחת. בצק שמטגנין אותו בדבש, וקורין לו אסקריטין בלשון משנה (פסחים ל"ז), והוא תרגום של אנקלוס:}%
% ---------- פרק 16 | פסוק 32 ----------
 \textbf{לב} וַיֹּ֣אמֶר מֹשֶׁ֗ה זֶ֤ה הַדָּבָר֙ אֲשֶׁ֣ר צִוָּ֣ה יְהֹוָ֔ה מְלֹ֤א הָעֹ֙מֶר֙ מִמֶּ֔נּוּ לְמִשְׁמֶ֖רֶת לְדֹרֹתֵיכֶ֑ם לְמַ֣עַן ׀ יִרְא֣וּ אֶת־הַלֶּ֗חֶם אֲשֶׁ֨ר הֶאֱכַ֤לְתִּי אֶתְכֶם֙ בַּמִּדְבָּ֔ר בְּהוֹצִיאִ֥י אֶתְכֶ֖ם מֵאֶ֥רֶץ מִצְרָֽיִם׃ \Rashi{\textbf{למשמרת.}לגניזה: לדרתיכם. בימי ירמיהו; כשהיה ירמיהו מוכיחם למה אין אתם עוסקים בתורה? והם אומרים נניח מלאכתנו ונעסק בתורה, מהיכן נתפרנס? הוציא להם צנצנת המן אמר להם אתם ראו דבר ה', שמעו לא נאמר אלא ראו, בזה נתפרנסו אבותיכם, הרבה שלוחין יש לו למקום להכין מזון ליראיו:}%
% ---------- פרק 16 | פסוק 33 ----------
 \textbf{לג} וַיֹּ֨אמֶר מֹשֶׁ֜ה אֶֽל־אַהֲרֹ֗ן קַ֚ח צִנְצֶ֣נֶת אַחַ֔ת וְתֶן־שָׁ֥מָּה מְלֹֽא־הָעֹ֖מֶר מָ֑ן וְהַנַּ֤ח אֹתוֹ֙ לִפְנֵ֣י יְהֹוָ֔ה לְמִשְׁמֶ֖רֶת לְדֹרֹתֵיכֶֽם׃ \Rashi{\textbf{צנצנת.}צלוחית של חרס, כתרגומו: והנח אותו לפני ה'. לפני הארון; ולא נאמר מקרא זה עד שנבנה אהל מועד, אלא שנכתב כאן בפרשת המן:}%
% ---------- פרק 16 | פסוק 34 ----------
 \textbf{לד} כַּאֲשֶׁ֛ר צִוָּ֥ה יְהֹוָ֖ה אֶל־מֹשֶׁ֑ה וַיַּנִּיחֵ֧הוּ אַהֲרֹ֛ן לִפְנֵ֥י הָעֵדֻ֖ת לְמִשְׁמָֽרֶת׃ %
% ---------- פרק 16 | פסוק 35 ----------
 \textbf{לה} וּבְנֵ֣י יִשְׂרָאֵ֗ל אָֽכְל֤וּ אֶת־הַמָּן֙ אַרְבָּעִ֣ים שָׁנָ֔ה עַד־בֹּאָ֖ם אֶל־אֶ֣רֶץ נוֹשָׁ֑בֶת אֶת־הַמָּן֙ אָֽכְל֔וּ עַד־בֹּאָ֕ם אֶל־קְצֵ֖ה אֶ֥רֶץ כְּנָֽעַן׃ \Rashi{\textbf{ארבעים שנה.}והלא חסר שלשים יום הם? שהרי בט"ו באיר ירד להם המן תחלה ובט"ו בניסן פסק, שנאמר "וישבת המן ממחרת" (יהושע ה')? אלא מגיד, שהעגות שהוציאו ישראל ממצרים טעמו בהם טעם מן (קידושין ל"ח): אל ארץ נושבת. לאחר שעברו את הירדן (שאותה שבעבר הירדן מישבת וטובה שנאמר "אעברה נא ואראה את הארץ הטובה אשר בעבר הירדן", ותרגום של נושבת "יתבתא", רצה לומר — מישבת): אל קצה ארץ כנען. בתחלת הגבול קדם שעברו את הירדן, והוא ערבות מואב. נמצאו מכחישין זה את זה? אלא בערבות מואב כשמת משה בז' באדר פסק המן מלירד, ונסתפקו ממן שלקטו בו ביום עד שהקריבו העמר בששה עשר בניסן, שנאמר "ויאכלו מעבור הארץ ממחרת הפסח" (יהושע ה'):}%
% ---------- פרק 16 | פסוק 36 ----------
 \textbf{לו} וְהָעֹ֕מֶר עֲשִׂרִ֥ית הָאֵיפָ֖ה הֽוּא׃ \{פ\} \Rashi{\textbf{עשרית האיפה.}האיפה שלש סאין והסאה ו' קבין והקב ד' לגין והלג ו' ביצים, נמצא עשירית האיפה מ"ג ביצים וחמש ביצה, והוא שעור לחלה ולמנחות:}%
% ---------- פרק 17 | פסוק 1 ----------
 \textbf{א} וַ֠יִּסְע֠וּ כׇּל־עֲדַ֨ת בְּנֵֽי־יִשְׂרָאֵ֧ל מִמִּדְבַּר־סִ֛ין לְמַסְעֵיהֶ֖ם עַל־פִּ֣י יְהֹוָ֑ה וַֽיַּחֲנוּ֙ בִּרְפִידִ֔ים וְאֵ֥ין מַ֖יִם לִשְׁתֹּ֥ת הָעָֽם׃ %
% ---------- פרק 17 | פסוק 2 ----------
 \textbf{ב} וַיָּ֤רֶב הָעָם֙ עִם־מֹשֶׁ֔ה וַיֹּ֣אמְר֔וּ תְּנוּ־לָ֥נוּ מַ֖יִם וְנִשְׁתֶּ֑ה וַיֹּ֤אמֶר לָהֶם֙ מֹשֶׁ֔ה מַה־תְּרִיבוּן֙ עִמָּדִ֔י מַה־תְּנַסּ֖וּן אֶת־יְהֹוָֽה׃ \Rashi{\textbf{מה תנסון.}לומר היוכל לתת מים בארץ ציה?:}%
% ---------- פרק 17 | פסוק 3 ----------
 \textbf{ג} וַיִּצְמָ֨א שָׁ֤ם הָעָם֙ לַמַּ֔יִם וַיָּ֥לֶן הָעָ֖ם עַל־מֹשֶׁ֑ה וַיֹּ֗אמֶר לָ֤מָּה זֶּה֙ הֶעֱלִיתָ֣נוּ מִמִּצְרַ֔יִם לְהָמִ֥ית אֹתִ֛י וְאֶת־בָּנַ֥י וְאֶת־מִקְנַ֖י בַּצָּמָֽא׃ %
% ---------- פרק 17 | פסוק 4 ----------
 \textbf{ד} וַיִּצְעַ֤ק מֹשֶׁה֙ אֶל־יְהֹוָ֣ה לֵאמֹ֔ר מָ֥ה אֶעֱשֶׂ֖ה לָעָ֣ם הַזֶּ֑ה ע֥וֹד מְעַ֖ט וּסְקָלֻֽנִי׃ \Rashi{\textbf{עוד מעט.}אם אמתין עוד מעט וסקלני:}%
% ---------- פרק 17 | פסוק 5 ----------
 \textbf{ה} וַיֹּ֨אמֶר יְהֹוָ֜ה אֶל־מֹשֶׁ֗ה עֲבֹר֙ לִפְנֵ֣י הָעָ֔ם וְקַ֥ח אִתְּךָ֖ מִזִּקְנֵ֣י יִשְׂרָאֵ֑ל וּמַטְּךָ֗ אֲשֶׁ֨ר הִכִּ֤יתָ בּוֹ֙ אֶת־הַיְאֹ֔ר קַ֥ח בְּיָדְךָ֖ וְהָלָֽכְתָּ׃ \Rashi{\textbf{עבור לפני העם.}וראה אם יסקלוך, למה הוצאת לעז על בני? (שמות רבה): וקח אתך מזקני ישראל. לעדות, שיראו שעל ידך מים יוצאים מן הצור, ולא יאמרו מעינות היו שם מימי קדם: ומטך אשר הכית בו את היאר. מה ת"ל "אשר הכית בו את היאר"? אלא שהיו ישראל אומרים על המטה שאינו מוכן אלא לפרענות – בו לקה פרעה ומצרים כמה מכות במצרים ועל הים – לכך נאמר אשר הכית בו את היאר, יראו עתה שאף לטובה הוא מוכן (מכילתא):}%
% ---------- פרק 17 | פסוק 6 ----------
 \textbf{ו} הִנְנִ֣י עֹמֵד֩ לְפָנֶ֨יךָ שָּׁ֥ם ׀ עַֽל־הַצּוּר֮ בְּחֹרֵב֒ וְהִכִּ֣יתָ בַצּ֗וּר וְיָצְא֥וּ מִמֶּ֛נּוּ מַ֖יִם וְשָׁתָ֣ה הָעָ֑ם וַיַּ֤עַשׂ כֵּן֙ מֹשֶׁ֔ה לְעֵינֵ֖י זִקְנֵ֥י יִשְׂרָאֵֽל׃ \Rashi{\textbf{והכית בצור.}על הצור לא נאמר אלא בצור, מכאן שהמטה היה של מין דבר חזק ושמו סנפירינון, והצור נבקע מפניו (שם):}%
% ---------- פרק 17 | פסוק 7 ----------
 \textbf{ז} וַיִּקְרָא֙ שֵׁ֣ם הַמָּק֔וֹם מַסָּ֖ה וּמְרִיבָ֑ה עַל־רִ֣יב ׀ בְּנֵ֣י יִשְׂרָאֵ֗ל וְעַ֨ל נַסֹּתָ֤ם אֶת־יְהֹוָה֙ לֵאמֹ֔ר הֲיֵ֧שׁ יְהֹוָ֛ה בְּקִרְבֵּ֖נוּ אִם־אָֽיִן׃ \{פ\} %
% ---------- פרק 17 | פסוק 8 ----------
 \textbf{ח} וַיָּבֹ֖א עֲמָלֵ֑ק וַיִּלָּ֥חֶם עִם־יִשְׂרָאֵ֖ל בִּרְפִידִֽם׃ \Rashi{\textbf{ויבא עמלק וגו'.}סמך פרשה זו למקרא זה, לומר, תמיד אני ביניכם ומזמן לכל צרכיכם ואתם אומרים היש ה' בקרבנו אם אין, חייכם שהכלב בא ונושך אתכם ואתם צועקים לי ותדעו היכן אני? משל לאדם שהרכיב בנו על כתפו ויצא לדרך, היה אותו הבן רואה חפץ ואומר, אבא טל חפץ זה ותן לי והוא נותן לו, וכן שניה, וכן שלישית, פגעו באדם אחד אמר לו אותו הבן, ראית את אבא? אמר לו אביו, אינך יודע היכן אני? השליכו מעליו ובא הכלב ונשכו (תנחומא יתרו):}%
% ---------- פרק 17 | פסוק 9 ----------
 \textbf{ט} וַיֹּ֨אמֶר מֹשֶׁ֤ה אֶל־יְהוֹשֻׁ֙עַ֙ בְּחַר־לָ֣נוּ אֲנָשִׁ֔ים וְצֵ֖א הִלָּחֵ֣ם בַּעֲמָלֵ֑ק מָחָ֗ר אָנֹכִ֤י נִצָּב֙ עַל־רֹ֣אשׁ הַגִּבְעָ֔ה וּמַטֵּ֥ה הָאֱלֹהִ֖ים בְּיָדִֽי׃ \Rashi{\textbf{בחר לנו.}לי ולך – השוהו לו; מכאן אמרו חכמים יהי כבוד תלמידך חביב עליך כשלך; וכבוד חברך כמורא רבך מנין? שנאמר "ויאמר אהרן אל משה בי אדני", והלא אהרן גדול מאחיו היה ועושה את חברו כרבו. ומורא רבך כמורא שמים מנין? שנאמר "אדני משה כלאם" – כלם מן העולם, חיבין הם כליה, המורדים בך כאלו מרדו בהקב"ה: וצא הלחם. צא מן הענן והלחם בו: מחר. בעת המלחמה, אנכי נצב: בחר לנו אנשים. גבורים ויראי חטא שתהא זכותן מסיעתן. דבר אחר – בחר לנו אנשים שיודעין לבטל כשפים, לפי שבני עמלק מכשפין היו:}%
% ---------- פרק 17 | פסוק 10 ----------
 \textbf{י} וַיַּ֣עַשׂ יְהוֹשֻׁ֗עַ כַּאֲשֶׁ֤ר אָֽמַר־לוֹ֙ מֹשֶׁ֔ה לְהִלָּחֵ֖ם בַּעֲמָלֵ֑ק וּמֹשֶׁה֙ אַהֲרֹ֣ן וְח֔וּר עָל֖וּ רֹ֥אשׁ הַגִּבְעָֽה׃ \Rashi{\textbf{ומשה אהרן וחור.}מכאן לתענית שצריכים שלשה לעבר לפני התבה, שבתענית היו שרויים (מכילתא): חור. בנה של מרים היה וכלב בעלה (סוטה י"א):}%
% ---------- פרק 17 | פסוק 11 ----------
 \textbf{יא} וְהָיָ֗ה כַּאֲשֶׁ֨ר יָרִ֥ים מֹשֶׁ֛ה יָד֖וֹ וְגָבַ֣ר יִשְׂרָאֵ֑ל וְכַאֲשֶׁ֥ר יָנִ֛יחַ יָד֖וֹ וְגָבַ֥ר עֲמָלֵֽק׃ \Rashi{\textbf{כאשר ירים משה ידו.}וכי ידיו של משה נוצחות היו המלחמה וכו' כדאיתא בר"ה:}%
% ---------- פרק 17 | פסוק 12 ----------
 \textbf{יב} וִידֵ֤י מֹשֶׁה֙ כְּבֵדִ֔ים וַיִּקְחוּ־אֶ֛בֶן וַיָּשִׂ֥ימוּ תַחְתָּ֖יו וַיֵּ֣שֶׁב עָלֶ֑יהָ וְאַהֲרֹ֨ן וְח֜וּר תָּֽמְכ֣וּ בְיָדָ֗יו מִזֶּ֤ה אֶחָד֙ וּמִזֶּ֣ה אֶחָ֔ד וַיְהִ֥י יָדָ֛יו אֱמוּנָ֖ה עַד־בֹּ֥א הַשָּֽׁמֶשׁ׃ \Rashi{\textbf{וידי משה כבדים.}בשביל שנתעצל במצוה ומנה אחר תחתיו, נתיקרו ידיו: ויקחו. אהרן וחור: אבן וישימו תחתיו. ולא ישב לו על כר וכסת. אמר: ישראל שרויין בצער, אף אני אהיה עמהם בצער (תענית י"א): ויהי ידיו אמונה. ויהי משה ידיו באמונה, פרושות השמים בתפלה נאמנה ונכונה: עד בא השמש. שהיו עמלקים מחשבין את השעות באיצטרולוגיאה, באיזו שעה הם נוצחים, והעמיד להם משה חמה וערבב את השעות (תנחומא):}%
% ---------- פרק 17 | פסוק 13 ----------
 \textbf{יג} וַיַּחֲלֹ֧שׁ יְהוֹשֻׁ֛עַ אֶת־עֲמָלֵ֥ק וְאֶת־עַמּ֖וֹ לְפִי־חָֽרֶב׃ \{פ\} \Rashi{\textbf{ויחלש יהושע.}חתך ראשי גבוריו ולא השאיר אלא חלשים שבהם, ולא הרגם כלם, מכאן אנו למדים שעשו ע"פ הדבור של שכינה (שם):}%
% ---------- פרק 17 | פסוק 14 ----------
 \textbf{יד} וַיֹּ֨אמֶר יְהֹוָ֜ה אֶל־מֹשֶׁ֗ה כְּתֹ֨ב זֹ֤את זִכָּרוֹן֙ בַּסֵּ֔פֶר וְשִׂ֖ים בְּאׇזְנֵ֣י יְהוֹשֻׁ֑עַ כִּֽי־מָחֹ֤ה אֶמְחֶה֙ אֶת־זֵ֣כֶר עֲמָלֵ֔ק מִתַּ֖חַת הַשָּׁמָֽיִם׃ \Rashi{\textbf{כתב זאת זכרון.}שבא עמלק להזדוג לישראל קדם לכל האמות: ושים באזני יהושע. המכניס את ישראל לארץ, שיצוה את ישראל לשלם לו את גמולו; כאן נרמז לו למשה שיהושע מכניס את ישראל לארץ (שם): כי מחה אמחה. לכך אני מזהירך כן, כי חפץ אני למחותו:}%
% ---------- פרק 17 | פסוק 15 ----------
 \textbf{טו} וַיִּ֥בֶן מֹשֶׁ֖ה מִזְבֵּ֑חַ וַיִּקְרָ֥א שְׁמ֖וֹ יְהֹוָ֥ה ׀ נִסִּֽי׃ \Rashi{\textbf{ויקרא שמו.}של מזבח: ה' נסי. הקב"ה עשה לנו כאן נס. לא שהמזבח קרוי ה', אלא המזכיר שמו של מזבח זוכר את הנס שעשה המקום – ה' הוא נס שלנו:}%
% ---------- פרק 17 | פסוק 16 ----------
 \textbf{טז} וַיֹּ֗אמֶר כִּֽי־יָד֙ עַל־כֵּ֣ס יָ֔הּ*(בכתר ארם צובה היה כתוב כֵּ֣סְיָ֔הּ בתיבה אחת) מִלְחָמָ֥ה לַיהֹוָ֖ה בַּֽעֲמָלֵ֑ק מִדֹּ֖ר דֹּֽר׃ \{פ\} \Rashi{\textbf{ויאמר.}משה: כי יד על כס יה. ידו של הקב"ה הורמה לשבע בכסאו להיות לו מלחמה ואיבה בעמלק עולמית, ומהו כס ולא נאמר כסא? ואף השם נחלק לחציו? נשבע הקב"ה שאין שמו שלם ואין כסאו שלם עד שימחה שמו של עמלק כלו, וכשימחה שמו יהיה השם שלם והכסא שלם, שנאמר "האויב תמו חרבות לנצח" (תהלים ט'), זהו עמלק שכתוב בו "ועברתו שמרה נצח" (עמוס א'), "וערים נתשת אבד זכרם המה" (תהלים שם), מהו אומר אחריו? "וה' לעולם ישב" — הרי השם שלם, "כונן למשפט כסאו" — הרי כסאו שלם:}%
\addparasha{חומש שמות | פרשת יתרו}
% ---------- פרק 18 | פסוק 1 ----------
 \textbf{א} וַיִּשְׁמַ֞ע יִתְר֨וֹ כֹהֵ֤ן מִדְיָן֙ חֹתֵ֣ן מֹשֶׁ֔ה אֵת֩ כׇּל־אֲשֶׁ֨ר עָשָׂ֤ה אֱלֹהִים֙ לְמֹשֶׁ֔ה וּלְיִשְׂרָאֵ֖ל עַמּ֑וֹ כִּֽי־הוֹצִ֧יא יְהֹוָ֛ה אֶת־יִשְׂרָאֵ֖ל מִמִּצְרָֽיִם׃ \Rashi{\textbf{וישמע יתרו.}מה שמועה שמע ובא? קריעת ים סוף ומלחמת עמלק: יתרו. שבע שמות נקראו לו: רעואל, יתר, יתרו, חובב, חבר, קיני, פוטיאל; יתר, על שם שיתר פרשה אחת בתורה "ואתה תחזה"; יתרו, לכשנתגיר וקים המצוות, הוסיפו לו אות אחד על שמו; חובב, שחבב את התורה; חובב הוא יתרו שנאמר "מבני חובב חתן משה" (שופטים ד'), ויש אומרים רעואל אביו של יתרו, ומהו אומר "ותבאנה אל רעואל אביהן"? שהתינוקות קורין לאבי אביהן אבא. בספרי: חתן משה. כאן היה יתרו מתכבד במשה – אני חותן המלך, ולשעבר היה משה תולה הגדלה בחמיו, שנאמר "וישב אל יתר חותנו" (שמות ד'): למשה ולישראל. שקול משה כנגד כל ישראל (מכילתא): את כל אשר עשה. להם בירידת המן ובבאר ובעמלק: כי הוציא ה' וגו'. זו גדולה על כלם (שם):}%
% ---------- פרק 18 | פסוק 2 ----------
 \textbf{ב} וַיִּקַּ֗ח יִתְרוֹ֙ חֹתֵ֣ן מֹשֶׁ֔ה אֶת־צִפֹּרָ֖ה אֵ֣שֶׁת מֹשֶׁ֑ה אַחַ֖ר שִׁלּוּחֶֽיהָ׃ \Rashi{\textbf{אחר שלוחיה.}כשאמר לו הקב"ה במדין לך שב מצרים ויקח משה את אשתו ואת בניו וגו' ויצא אהרן לקראתו ויפגשהו בהר האלהים (שמות ד'), אמר לו מי הם הללו? אמר לו זו היא אשתי שנשאתי במדין ואלו בני, אמר לו והיכן אתה מוליכן? אמר לו למצרים, אמר לו על הראשונים אנו מצטערים ואתה בא להוסיף עליהם אמר לה לכי לבית אביך, נטלה שני בניה והלכה לה (מכילתא):}%
% ---------- פרק 18 | פסוק 3 ----------
 \textbf{ג} וְאֵ֖ת שְׁנֵ֣י בָנֶ֑יהָ אֲשֶׁ֨ר שֵׁ֤ם הָֽאֶחָד֙ גֵּֽרְשֹׁ֔ם כִּ֣י אָמַ֔ר גֵּ֣ר הָיִ֔יתִי בְּאֶ֖רֶץ נׇכְרִיָּֽה׃ %
% ---------- פרק 18 | פסוק 4 ----------
 \textbf{ד} וְשֵׁ֥ם הָאֶחָ֖ד אֱלִיעֶ֑זֶר כִּֽי־אֱלֹהֵ֤י אָבִי֙ בְּעֶזְרִ֔י וַיַּצִּלֵ֖נִי מֵחֶ֥רֶב פַּרְעֹֽה׃ \Rashi{\textbf{ויצלני מחרב פרעה.}כשגלו דתן ואבירם על דבר המצרי ובקש להרג את משה נעשה צוארו כעמוד של שיש (שם):}%
% ---------- פרק 18 | פסוק 5 ----------
 \textbf{ה} וַיָּבֹ֞א יִתְר֨וֹ חֹתֵ֥ן מֹשֶׁ֛ה וּבָנָ֥יו וְאִשְׁתּ֖וֹ אֶל־מֹשֶׁ֑ה אֶל־הַמִּדְבָּ֗ר אֲשֶׁר־ה֛וּא חֹנֶ֥ה שָׁ֖ם הַ֥ר הָאֱלֹהִֽים׃ \Rashi{\textbf{אל המדבר.}אף אנו יודעים שבמדבר היו אלא בשבחו של יתרו דבר הכתוב, שהיה יושב בכבודו של עולם ונדבו לבו לצאת אל המדבר, מקום תהו לשמע דברי תורה (שם):}%
% ---------- פרק 18 | פסוק 6 ----------
 \textbf{ו} וַיֹּ֙אמֶר֙ אֶל־מֹשֶׁ֔ה אֲנִ֛י חֹתֶנְךָ֥ יִתְר֖וֹ בָּ֣א אֵלֶ֑יךָ וְאִ֨שְׁתְּךָ֔ וּשְׁנֵ֥י בָנֶ֖יהָ עִמָּֽהּ׃ \Rashi{\textbf{ויאמר אל משה.}ע"י שליח (שם): אני חתנך יתרו וגו'. אם אין אתה יוצא בגיני, צא בגין אשתך, ואם אין אתה יוצא בגין אשתך, צא בגין שני בניה (שם):}%
% ---------- פרק 18 | פסוק 7 ----------
 \textbf{ז} וַיֵּצֵ֨א מֹשֶׁ֜ה לִקְרַ֣את חֹֽתְנ֗וֹ וַיִּשְׁתַּ֙חוּ֙ וַיִּשַּׁק־ל֔וֹ וַיִּשְׁאֲל֥וּ אִישׁ־לְרֵעֵ֖הוּ לְשָׁל֑וֹם וַיָּבֹ֖אוּ הָאֹֽהֱלָה׃ \Rashi{\textbf{ויצא משה.}כבוד גדול נתכבד יתרו באותה שעה, כיון שיצא משה, יצא אהרן נדב ואביהוא, ומי הוא שראה את אלו יוצאין ולא יצא? (תנחומא): וישתחו וישק לו. איני יודע מי השתחוה למי, כשהוא אומר איש לרעהו, מי הקרוי "איש" – משה, שנאמר "והאיש משה" (במדבר י"ב):}%
% ---------- פרק 18 | פסוק 8 ----------
 \textbf{ח} וַיְסַפֵּ֤ר מֹשֶׁה֙ לְחֹ֣תְנ֔וֹ אֵת֩ כׇּל־אֲשֶׁ֨ר עָשָׂ֤ה יְהֹוָה֙ לְפַרְעֹ֣ה וּלְמִצְרַ֔יִם עַ֖ל אוֹדֹ֣ת יִשְׂרָאֵ֑ל אֵ֤ת כׇּל־הַתְּלָאָה֙ אֲשֶׁ֣ר מְצָאָ֣תַם בַּדֶּ֔רֶךְ וַיַּצִּלֵ֖ם יְהֹוָֽה׃ \Rashi{\textbf{ויספר משה לחתנו.}למשך את לבו לקרבו לתורה (מכילתא): את כל התלאה. שעל הים ושל עמלק: התלאה. למ"ד אל"ף מן היסוד של תבה, והתי"ו הוא תקון ויסוד הנופל ממנו לפרקים; וכן תרומה, תנופה, תקומה, תנואה:}%
% ---------- פרק 18 | פסוק 9 ----------
 \textbf{ט} וַיִּ֣חַדְּ יִתְר֔וֹ עַ֚ל כׇּל־הַטּוֹבָ֔ה אֲשֶׁר־עָשָׂ֥ה יְהֹוָ֖ה לְיִשְׂרָאֵ֑ל אֲשֶׁ֥ר הִצִּיל֖וֹ מִיַּ֥ד מִצְרָֽיִם׃ \Rashi{\textbf{ויחד יתרו.}וישמח יתרו, זהו פשוטו. ומ"א נעשה בשרו חדודין חדודין, מצר על אבוד מצרים, הינו דאמרי אינשי "גיורא עד עשרה דרי לא תבזי ארמאה באפיה": על כל הטובה. טובת המן והבאר והתורה, ועל כלן אשר הצילו מיד מצרים; עד עכשו לא היה עבד יכול לברח ממצרים, שהיתה הארץ מסגרת, ואלו יצאו ששים רבוא (מכילתא):}%
% ---------- פרק 18 | פסוק 10 ----------
 \textbf{י} וַיֹּ֘אמֶר֮ יִתְרוֹ֒ בָּר֣וּךְ יְהֹוָ֔ה אֲשֶׁ֨ר הִצִּ֥יל אֶתְכֶ֛ם מִיַּ֥ד מִצְרַ֖יִם וּמִיַּ֣ד פַּרְעֹ֑ה אֲשֶׁ֤ר הִצִּיל֙ אֶת־הָעָ֔ם מִתַּ֖חַת יַד־מִצְרָֽיִם׃ \Rashi{\textbf{אשר הציל אתכם מיד מצרים.}אמה קשה: ומיד פרעה. מלך קשה: מתחת יד מצרים. כתרגומו – לשון רדוי ומרות; היד שהיו מכבידים עליהם היא העבודה:}%
% ---------- פרק 18 | פסוק 11 ----------
 \textbf{יא} עַתָּ֣ה יָדַ֔עְתִּי כִּֽי־גָד֥וֹל יְהֹוָ֖ה מִכׇּל־הָאֱלֹהִ֑ים כִּ֣י בַדָּבָ֔ר אֲשֶׁ֥ר זָד֖וּ עֲלֵיהֶֽם׃ \Rashi{\textbf{עתה ידעתי.}מכירו הייתי לשעבר ועכשו ביותר: מכל האלהים. מלמד שהיה מכיר בכל ע"ז שבעולם, שלא הניח ע"ז שלא עבדה (מכילתא): כי בדבר אשר זדו עליהם. כתרגומו – במים דמו לאבדם והם נאבדו במים: אשר זדו. אשר הרשיעו. ורבותינו דרשוהו לשון ויזד יעקב נזיד, בקדרה אשר בשלו בה נתבשלו (סוטה י"א):}%
% ---------- פרק 18 | פסוק 12 ----------
 \textbf{יב} וַיִּקַּ֞ח יִתְר֨וֹ חֹתֵ֥ן מֹשֶׁ֛ה עֹלָ֥ה וּזְבָחִ֖ים לֵֽאלֹהִ֑ים וַיָּבֹ֨א אַהֲרֹ֜ן וְכֹ֣ל ׀ זִקְנֵ֣י יִשְׂרָאֵ֗ל לֶאֱכׇל־לֶ֛חֶם עִם־חֹתֵ֥ן מֹשֶׁ֖ה לִפְנֵ֥י הָאֱלֹהִֽים׃ \Rashi{\textbf{עלה.}כמשמעה, שהיא כלה כליל: וזבחים. שלמים: ויבא אהרן וגו'. ומשה היכן הלך? והלא הוא שיצא לקראתו וגרם לו את כל הכבוד? אלא שהיה עומד ומשמש לפניהם (מכילתא): לפני האלהים. מכאן שהנהנה מסעודה שתלמידי חכמים מסבין בה, כאלו נהנה מזיו השכינה (ברכות ס"ד):}%
% ---------- פרק 18 | פסוק 13 ----------
 \textbf{יג} וַֽיְהִי֙ מִֽמׇּחֳרָ֔ת וַיֵּ֥שֶׁב מֹשֶׁ֖ה לִשְׁפֹּ֣ט אֶת־הָעָ֑ם וַיַּעֲמֹ֤ד הָעָם֙ עַל־מֹשֶׁ֔ה מִן־הַבֹּ֖קֶר עַד־הָעָֽרֶב׃ \Rashi{\textbf{ויהי ממחרת.}מוצאי יום הכפורים היה, כך שנינו בספרי ומהו ממחרת? למחרת רדתו מן ההר; ועל כרחך אי אפשר לומר אלא ממחרת יום הכפורים, שהרי קדם מתן תורה א"א לומר "והודעתי את חקי וגו'", ומשנתנה תורה עד יה"כ לא ישב משה לשפט את העם, שהרי בי"ז בתמוז ירד ושבר את הלחות ולמחר עלה בהשכמה ושהה שמונים יום וירד ביה"כ; ואין פרשה זו כתובה כסדר, שלא נאמר ויהי ממחרת עד שנה שניה – אף לדברי האומר יתרו קדם מתן תורה בא, שלוחו אל ארצו לא היה אלא עד שנה שניה – שהרי נאמר כאן וישלח משה את חתנו, ומצינו במסע הדגלים שאמר לו משה "נסעים אנחנו אל המקום וגו'"(במדבר י כט), אל נא תעזב אתנו" (שם לא), ואם זו קדם מתן תורה, מששלחו והלך היכן מצינו שחזר? ואם תאמר, שם לא נאמר יתרו אלא חובב, ובנו של יתרו היה, הוא חובב הוא יתרו, שהרי כתוב "מבני חבב חתן משה" (שופטים ד'): וישב משה וגו' ויעמד העם. יושב כמלך וכלן עומדים, והקשה הדבר ליתרו שהיה מזלזל בכבודן של ישראל והוכיחו על כך, שנאמר (פסוק יד) מדוע אתה יושב לבדך וכלם נצבים: מן הבקר עד ערב. אפשר לומר כן? אלא כל דין שדן דין אמת לאמתו אפלו שעה אחת, מעלה עליו הכתוב כאלו עוסק בתורה כל היום וכאלו נעשה שתף להקב"ה במעשה בראשית, שנאמר בו ויהי ערב וגו' (שבת י'):}%
% ---------- פרק 18 | פסוק 14 ----------
 \textbf{יד} וַיַּרְא֙ חֹתֵ֣ן מֹשֶׁ֔ה אֵ֛ת כׇּל־אֲשֶׁר־ה֥וּא עֹשֶׂ֖ה לָעָ֑ם וַיֹּ֗אמֶר מָֽה־הַדָּבָ֤ר הַזֶּה֙ אֲשֶׁ֨ר אַתָּ֤ה עֹשֶׂה֙ לָעָ֔ם מַדּ֗וּעַ אַתָּ֤ה יוֹשֵׁב֙ לְבַדֶּ֔ךָ וְכׇל־הָעָ֛ם נִצָּ֥ב עָלֶ֖יךָ מִן־בֹּ֥קֶר עַד־עָֽרֶב׃ %
% ---------- פרק 18 | פסוק 15 ----------
 \textbf{טו} וַיֹּ֥אמֶר מֹשֶׁ֖ה לְחֹתְנ֑וֹ כִּֽי־יָבֹ֥א אֵלַ֛י הָעָ֖ם לִדְרֹ֥שׁ אֱלֹהִֽים׃ \Rashi{\textbf{כי יבא.}כמו כי בא, לשון הוה: לדרש אלהים. כתרגומו "למתבע אלפן" – לשאל תלמוד מפי הגבורה:}%
% ---------- פרק 18 | פסוק 16 ----------
 \textbf{טז} כִּֽי־יִהְיֶ֨ה לָהֶ֤ם דָּבָר֙ בָּ֣א אֵלַ֔י וְשָׁ֣פַטְתִּ֔י בֵּ֥ין אִ֖ישׁ וּבֵ֣ין רֵעֵ֑הוּ וְהוֹדַעְתִּ֛י אֶת־חֻקֵּ֥י הָאֱלֹהִ֖ים וְאֶת־תּוֹרֹתָֽיו׃ \Rashi{\textbf{כי יהיה להם דבר בא.}מי שהיה לו הדבר – בא אלי:}%
% ---------- פרק 18 | פסוק 17 ----------
 \textbf{יז} וַיֹּ֛אמֶר חֹתֵ֥ן מֹשֶׁ֖ה אֵלָ֑יו לֹא־טוֹב֙ הַדָּבָ֔ר אֲשֶׁ֥ר אַתָּ֖ה עֹשֶֽׂה׃ \Rashi{\textbf{ויאמר חתן משה.}דרך כבוד קוראו הכתוב חותנו של מלך:}%
% ---------- פרק 18 | פסוק 18 ----------
 \textbf{יח} נָבֹ֣ל תִּבֹּ֔ל גַּם־אַתָּ֕ה גַּם־הָעָ֥ם הַזֶּ֖ה אֲשֶׁ֣ר עִמָּ֑ךְ כִּֽי־כָבֵ֤ד מִמְּךָ֙ הַדָּבָ֔ר לֹא־תוּכַ֥ל עֲשֹׂ֖הוּ לְבַדֶּֽךָ׃ \Rashi{\textbf{נבל תבל.}כתרגומו, ולשונו לשון כמישה, פלייש"טרא בלעז, כמו "והעלה נבל" (ירמיהו ח'), "כנבל עלה מגפן" (ישעיהו ל"ד), שהוא כמוש על ידי חמה ועל ידי קרח וכחו תש ונלאה: גם אתה. לרבות אהרן וחור ושבעים זקנים: כי כבד ממך. כבדו רב יותר מכחך:}%
% ---------- פרק 18 | פסוק 19 ----------
 \textbf{יט} עַתָּ֞ה שְׁמַ֤ע בְּקֹלִי֙ אִיעָ֣צְךָ֔ וִיהִ֥י אֱלֹהִ֖ים עִמָּ֑ךְ הֱיֵ֧ה אַתָּ֣ה לָעָ֗ם מ֚וּל הָֽאֱלֹהִ֔ים וְהֵבֵאתָ֥ אַתָּ֛ה אֶת־הַדְּבָרִ֖ים אֶל־הָאֱלֹהִֽים׃ \Rashi{\textbf{איעצך ויהי אלהים עמך.}בעצה; אמר לו צא המלך בגבורה (מכילתא): היה אתה לעם מול האלהים. שליח ומליץ בינותם למקום ושואל משפטים מאתו: את הדברים. דברי ריבותם:}%
% ---------- פרק 18 | פסוק 20 ----------
 \textbf{כ} וְהִזְהַרְתָּ֣ה אֶתְהֶ֔ם אֶת־הַחֻקִּ֖ים וְאֶת־הַתּוֹרֹ֑ת וְהוֹדַעְתָּ֣ לָהֶ֗ם אֶת־הַדֶּ֙רֶךְ֙ יֵ֣לְכוּ בָ֔הּ וְאֶת־הַֽמַּעֲשֶׂ֖ה אֲשֶׁ֥ר יַעֲשֽׂוּן׃ %
% ---------- פרק 18 | פסוק 21 ----------
 \textbf{כא} וְאַתָּ֣ה תֶחֱזֶ֣ה מִכׇּל־הָ֠עָ֠ם אַנְשֵׁי־חַ֜יִל יִרְאֵ֧י אֱלֹהִ֛ים אַנְשֵׁ֥י אֱמֶ֖ת שֹׂ֣נְאֵי בָ֑צַע וְשַׂמְתָּ֣ עֲלֵהֶ֗ם שָׂרֵ֤י אֲלָפִים֙ שָׂרֵ֣י מֵא֔וֹת שָׂרֵ֥י חֲמִשִּׁ֖ים וְשָׂרֵ֥י עֲשָׂרֹֽת׃ \Rashi{\textbf{ואתה תחזה.}ברוח הקדש שעליך: אנשי חיל. עשירים, שאין צריכין להחניף ולהכיר פנים (שם): אנשי אמת. אלו בעלי הבטחה, שהם כדאי לסמך על דבריהם, שעל ידי כן יהיו דבריהם נשמעין: שנאי בצע. ששונאין את ממונם בדין, כההיא דאמרינן: כל דינא דמפקין ממונא מיניה בדינא לאו דינא הוא (בבא בתרא נ"ח): שרי אלפים. הם היו שש מאות שרים לשש מאות אלף: שרי מאות. שש אלפים היו: שרי חמשים. י"ב אלף: שרי עשרת. ששים אלף:}%
% ---------- פרק 18 | פסוק 22 ----------
 \textbf{כב} וְשָׁפְט֣וּ אֶת־הָעָם֮ בְּכׇל־עֵת֒ וְהָיָ֞ה כׇּל־הַדָּבָ֤ר הַגָּדֹל֙ יָבִ֣יאוּ אֵלֶ֔יךָ וְכׇל־הַדָּבָ֥ר הַקָּטֹ֖ן יִשְׁפְּטוּ־הֵ֑ם וְהָקֵל֙ מֵֽעָלֶ֔יךָ וְנָשְׂא֖וּ אִתָּֽךְ׃ \Rashi{\textbf{ושפטו.}וידונון, לשון צווי: והקל מעליך. דבר זה להקל מעליך: והקל. כמו "והכבד את לבו" (שמות ח'), "והכות את מואב" (מלכים ב ג'), לשון הוה:}%
% ---------- פרק 18 | פסוק 23 ----------
 \textbf{כג} אִ֣ם אֶת־הַדָּבָ֤ר הַזֶּה֙ תַּעֲשֶׂ֔ה וְצִוְּךָ֣ אֱלֹהִ֔ים וְיָֽכׇלְתָּ֖ עֲמֹ֑ד וְגַם֙ כׇּל־הָעָ֣ם הַזֶּ֔ה עַל־מְקֹמ֖וֹ יָבֹ֥א בְשָׁלֽוֹם׃ \Rashi{\textbf{וצוך אלהים ויכלת עמד.}המלך בגבורה, אם מצוה אותך לעשות כן, תוכל עמד, ואם יעכב על ידך לא תוכל לעמד (מכילתא): וגם כל העם הזה. אהרן נדב ואביהוא ושבעים זקנים הנלוים עתה עמך (שם):}%
% ---------- פרק 18 | פסוק 24 ----------
 \textbf{כד} וַיִּשְׁמַ֥ע מֹשֶׁ֖ה לְק֣וֹל חֹתְנ֑וֹ וַיַּ֕עַשׂ כֹּ֖ל אֲשֶׁ֥ר אָמָֽר׃ %
% ---------- פרק 18 | פסוק 25 ----------
 \textbf{כה} וַיִּבְחַ֨ר מֹשֶׁ֤ה אַנְשֵׁי־חַ֙יִל֙ מִכׇּל־יִשְׂרָאֵ֔ל וַיִּתֵּ֥ן אֹתָ֛ם רָאשִׁ֖ים עַל־הָעָ֑ם שָׂרֵ֤י אֲלָפִים֙ שָׂרֵ֣י מֵא֔וֹת שָׂרֵ֥י חֲמִשִּׁ֖ים וְשָׂרֵ֥י עֲשָׂרֹֽת׃ %
% ---------- פרק 18 | פסוק 26 ----------
 \textbf{כו} וְשָׁפְט֥וּ אֶת־הָעָ֖ם בְּכׇל־עֵ֑ת אֶת־הַדָּבָ֤ר הַקָּשֶׁה֙ יְבִיא֣וּן אֶל־מֹשֶׁ֔ה וְכׇל־הַדָּבָ֥ר הַקָּטֹ֖ן יִשְׁפּוּט֥וּ הֵֽם׃ \Rashi{\textbf{ושפטו.}ודינין ית עמא: יביאון. מיתין: ישפוטו הם. כמו ישפטו, וכן "לא תעבורי" (רות ב'), כמו לא תעברי, ותרגומו דינין אינון; מקראות העליונים הם לשון צווי, לכך מתרגמין וידונון, ייתון, ידונון, ומקראות הללו לשון עשיה:}%
% ---------- פרק 18 | פסוק 27 ----------
 \textbf{כז} וַיְשַׁלַּ֥ח מֹשֶׁ֖ה אֶת־חֹתְנ֑וֹ וַיֵּ֥לֶךְ ל֖וֹ אֶל־אַרְצֽוֹ׃ \{פ\} \Rashi{\textbf{וילך לו אל ארצו.}לגיר בני משפחתו (מכילתא):}%
% ---------- פרק 19 | פסוק 1 ----------
 \textbf{א} בַּחֹ֙דֶשׁ֙ הַשְּׁלִישִׁ֔י לְצֵ֥את בְּנֵי־יִשְׂרָאֵ֖ל מֵאֶ֣רֶץ מִצְרָ֑יִם בַּיּ֣וֹם הַזֶּ֔ה בָּ֖אוּ מִדְבַּ֥ר סִינָֽי׃ \Rashi{\textbf{ביום הזה.}בראש חדש; לא היה צריך לכתב אלא ביום ההוא, מהו ביום הזה? שיהיו דברי תורה חדשים עליך כאלו היום נתנם (ברכות ס"ג):}%
% ---------- פרק 19 | פסוק 2 ----------
 \textbf{ב} וַיִּסְע֣וּ מֵרְפִידִ֗ים וַיָּבֹ֙אוּ֙ מִדְבַּ֣ר סִינַ֔י וַֽיַּחֲנ֖וּ בַּמִּדְבָּ֑ר וַיִּֽחַן־שָׁ֥ם יִשְׂרָאֵ֖ל נֶ֥גֶד הָהָֽר׃ \Rashi{\textbf{ויסעו מרפידים.}מה ת"ל לחזר ולפרש מהיכן נסעו, והלא כבר כתב שברפידים היו חונים, בידוע שמשם נסעו? אלא להקיש נסיעתן מרפידים לביאתן למדבר סיני, מה ביאתן למדבר סיני בתשובה, אף נסיעתן מרפידים בתשובה (מכילתא): ויחן שם ישראל. כאיש אחד בלב אחד, אבל שאר כל החניות בתרעומות ובמחלקת: נגד ההר. למזרחו, וכל מקום שאתה מוצא נגד, פנים למזרח (שם):}%
% ---------- פרק 19 | פסוק 3 ----------
 \textbf{ג} וּמֹשֶׁ֥ה עָלָ֖ה אֶל־הָאֱלֹהִ֑ים וַיִּקְרָ֨א אֵלָ֤יו יְהֹוָה֙ מִן־הָהָ֣ר לֵאמֹ֔ר כֹּ֤ה תֹאמַר֙ לְבֵ֣ית יַעֲקֹ֔ב וְתַגֵּ֖יד לִבְנֵ֥י יִשְׂרָאֵֽל׃ \Rashi{\textbf{ומשה עלה.}ביום השני; וכל עליותיו בהשכמה היו, שנאמר "וישכם משה בבקר" (שמות ל"ד): כה תאמר. בלשון הזה וכסדר הזה: לבית יעקב. אלו הנשים, תאמר להן בלשון רכה (מכילתא): ותגיד לבני ישראל. ענשין ודקדוקים פרש לזכרים, דברים הקשין כגידין (שבת פ"ז):}%
% ---------- פרק 19 | פסוק 4 ----------
 \textbf{ד} אַתֶּ֣ם רְאִיתֶ֔ם אֲשֶׁ֥ר עָשִׂ֖יתִי לְמִצְרָ֑יִם וָאֶשָּׂ֤א אֶתְכֶם֙ עַל־כַּנְפֵ֣י נְשָׁרִ֔ים וָאָבִ֥א אֶתְכֶ֖ם אֵלָֽי׃ \Rashi{\textbf{אתם ראיתם.}לא מסרת היא בידכם, לא בדברים אני משגר לכם, לא בעדים אני מעיד עליכם, אלא אתם ראיתם אשר עשיתי למצרים, על כמה עברות היו חיבין לי קדם שנזדוגו לכם ולא נפרעתי מהם אלא על ידכם (מכילתא): ואשא אתכם. זה היום שבאו ישראל לרעמסס, שהיו ישראל מפזרין בכל ארץ גשן, ולשעה קלה כשבאו לסע ולצאת, נקבצו כלם לרעמסס, ואנקלוס תרגם ואשא כמו ואסיע אתכם – "ואטלית יתכון", תקן את הדבר דרך כבוד למעלה: על כנפי נשרים. כנשר הנושא גוזליו על כנפיו, שכל שאר העופות נותנים את בניהם בין רגליהם, לפי שמתיראין מעוף אחר שפורח על גביהם, אבל הנשר הזה אינו מתירא אלא מן האדם שמא יזרק בו חץ, לפי שאין עוף פורח על גביו, לכך נותנו על כנפיו, אומר מוטב יכנס החץ בי ולא בבני, אף אני עשיתי כן: "ויסע מלאך האלהים וגו', ויבא בין מחנה מצרים וגו'" (שמות י"ד), והיו מצרים זורקים חצים ואבני בליסטראות והענן מקבלם: ואבא אתכם אלי. כתרגומו:}%
% ---------- פרק 19 | פסוק 5 ----------
 \textbf{ה} וְעַתָּ֗ה אִם־שָׁמ֤וֹעַ תִּשְׁמְעוּ֙ בְּקֹלִ֔י וּשְׁמַרְתֶּ֖ם אֶת־בְּרִיתִ֑י וִהְיִ֨יתֶם לִ֤י סְגֻלָּה֙ מִכׇּל־הָ֣עַמִּ֔ים כִּי־לִ֖י כׇּל־הָאָֽרֶץ׃ \Rashi{\textbf{ועתה.}אם עתה תקבלו עליכם, יערב לכם מכאן ואילך, שכל התחלות קשות (מכילתא): ושמרתם את בריתי. שאכרת עמכם על שמירת התורה: סגלה. אוצר חביב, כמו "וסגלת מלכים" (קהלת ב') – כלי יקר ואבנים טובות שהמלכים גונזים אותם, כך אתם תהיו לי סגלה משאר אמות; ולא תאמרו אתם לבדכם שלי ואין לי אחרים עמכם ומה יש לי עוד שתהא חבתכם נכרת: כי לי כל הארץ והם בעיני ולפני לכלום:}%
% ---------- פרק 19 | פסוק 6 ----------
 \textbf{ו} וְאַתֶּ֧ם תִּהְיוּ־לִ֛י מַמְלֶ֥כֶת כֹּהֲנִ֖ים וְג֣וֹי קָד֑וֹשׁ אֵ֚לֶּה הַדְּבָרִ֔ים אֲשֶׁ֥ר תְּדַבֵּ֖ר אֶל־בְּנֵ֥י יִשְׂרָאֵֽל׃ \Rashi{\textbf{ואתם תהיו לי ממלכת כהנים.}שרים, כמה דאת אמר "ובני דוד כהנים" (שמואל ב ח'): אלה הדברים. לא פחות ולא יותר (מכילתא):}%
% ---------- פרק 19 | פסוק 7 ----------
 \textbf{ז} וַיָּבֹ֣א מֹשֶׁ֔ה וַיִּקְרָ֖א לְזִקְנֵ֣י הָעָ֑ם וַיָּ֣שֶׂם לִפְנֵיהֶ֗ם אֵ֚ת כׇּל־הַדְּבָרִ֣ים הָאֵ֔לֶּה אֲשֶׁ֥ר צִוָּ֖הוּ יְהֹוָֽה׃ %
% ---------- פרק 19 | פסוק 8 ----------
 \textbf{ח} וַיַּעֲנ֨וּ כׇל־הָעָ֤ם יַחְדָּו֙ וַיֹּ֣אמְר֔וּ כֹּ֛ל אֲשֶׁר־דִּבֶּ֥ר יְהֹוָ֖ה נַעֲשֶׂ֑ה וַיָּ֧שֶׁב מֹשֶׁ֛ה אֶת־דִּבְרֵ֥י הָעָ֖ם אֶל־יְהֹוָֽה׃ \Rashi{\textbf{וישב משה את דברי העם וגו'.}ביום המחרת, שהוא שלישי, שהרי בהשכמה עלה; וכי צריך היה משה להשיב? אלא בא הכתוב ללמדך דרך ארץ ממשה, שלא אמר הואיל ויודע מי ששלחני איני צריך להשיב (שבת פ"ז):}%
% ---------- פרק 19 | פסוק 9 ----------
 \textbf{ט} וַיֹּ֨אמֶר יְהֹוָ֜ה אֶל־מֹשֶׁ֗ה הִנֵּ֨ה אָנֹכִ֜י בָּ֣א אֵלֶ֘יךָ֮ בְּעַ֣ב הֶֽעָנָן֒ בַּעֲב֞וּר יִשְׁמַ֤ע הָעָם֙ בְּדַבְּרִ֣י עִמָּ֔ךְ וְגַם־בְּךָ֖ יַאֲמִ֣ינוּ לְעוֹלָ֑ם וַיַּגֵּ֥ד מֹשֶׁ֛ה אֶת־דִּבְרֵ֥י הָעָ֖ם אֶל־יְהֹוָֽה׃ \Rashi{\textbf{בעב הענן.}במעבה הענן וזהו ערפל: וגם בך. גם בנביאים הבאים אחריך (מכילתא): ויגד משה את דברי וגו'. ביום המחרת, שהוא רביעי לחדש: את דברי העם וגו'. תשובה על דבר זה; שמעתי מהם שרצונם לשמע ממך, אינו דומה השומע מפי שליח לשומע מפי המלך, רצוננו לראות את מלכנו (מכילתא):}%
% ---------- פרק 19 | פסוק 10 ----------
 \textbf{י} וַיֹּ֨אמֶר יְהֹוָ֤ה אֶל־מֹשֶׁה֙ לֵ֣ךְ אֶל־הָעָ֔ם וְקִדַּשְׁתָּ֥ם הַיּ֖וֹם וּמָחָ֑ר וְכִבְּס֖וּ שִׂמְלֹתָֽם׃ \Rashi{\textbf{ויאמר ה' אל משה.}אם כן שמזקיקין לדבר עמם, לך אל העם:  וקדשתם. וזמנתם, שיכינו עצמם היום ומחר:}%
% ---------- פרק 19 | פסוק 11 ----------
 \textbf{יא} וְהָי֥וּ נְכֹנִ֖ים לַיּ֣וֹם הַשְּׁלִישִׁ֑י כִּ֣י ׀ בַּיּ֣וֹם הַשְּׁלִשִׁ֗י יֵרֵ֧ד יְהֹוָ֛ה לְעֵינֵ֥י כׇל־הָעָ֖ם עַל־הַ֥ר סִינָֽי׃ \Rashi{\textbf{והיו נכנים.}מבדלים מאשה (שבת פ"ז): ליום השלישי. שהוא ששה בחדש, ובחמישי בנה משה את המזבח תחת ההר ושתים עשרה מצבה – כל הענין האמור בפרשת ואלה המשפטים – ואין מקדם ומאחר בתורה: לעיני כל העם. מלמד שלא היה בהם סומא, שנתרפאו כלם (מכילתא):}%
% ---------- פרק 19 | פסוק 12 ----------
 \textbf{יב} וְהִגְבַּלְתָּ֤ אֶת־הָעָם֙ סָבִ֣יב לֵאמֹ֔ר הִשָּׁמְר֥וּ לָכֶ֛ם עֲל֥וֹת בָּהָ֖ר וּנְגֹ֣עַ בְּקָצֵ֑הוּ כׇּל־הַנֹּגֵ֥עַ בָּהָ֖ר מ֥וֹת יוּמָֽת׃ \Rashi{\textbf{והגבלת.}קבע להם תחומין לסימן, שלא יקרבו מן הגבול והלאה: לאמר. הגבול אומר להם השמרו מעלות מכאן ולהלאה, ואתה הזהירם על כך: ונגע בקצהו. אפלו בקצהו:}%
% ---------- פרק 19 | פסוק 13 ----------
 \textbf{יג} לֹא־תִגַּ֨ע בּ֜וֹ יָ֗ד כִּֽי־סָק֤וֹל יִסָּקֵל֙ אוֹ־יָרֹ֣ה יִיָּרֶ֔ה אִם־בְּהֵמָ֥ה אִם־אִ֖ישׁ לֹ֣א יִחְיֶ֑ה בִּמְשֹׁךְ֙ הַיֹּבֵ֔ל הֵ֖מָּה יַעֲל֥וּ בָהָֽר׃ \Rashi{\textbf{ירה יירה.}מכאן לנסקלין שהם נדחין למטה מבית הסקילה שהיה גבוה שתי קומות (סנהדרין מ"ה): יירה. ישלך למטה לארץ, כמו "ירה בים": במשך היבל. כשימשך היובל קול ארך, הוא סימן סלוק שכינה והפסקת הקול, וכיון שאסתלק, הם רשאין לעלות: היבל. הוא שופר של איל, שכן בערביא קורין לדכרא יובלא, ושופר של איל של יצחק היה (פרקי דרבי אליעזר ל"א):}%
% ---------- פרק 19 | פסוק 14 ----------
 \textbf{יד} וַיֵּ֧רֶד מֹשֶׁ֛ה מִן־הָהָ֖ר אֶל־הָעָ֑ם וַיְקַדֵּשׁ֙ אֶת־הָעָ֔ם וַֽיְכַבְּס֖וּ שִׂמְלֹתָֽם׃ \Rashi{\textbf{מן ההר אל העם.}מלמד שלא היה משה פונה לעסקיו, אלא מן ההר – אל העם (מכילתא):}%
% ---------- פרק 19 | פסוק 15 ----------
 \textbf{טו} וַיֹּ֙אמֶר֙ אֶל־הָעָ֔ם הֱי֥וּ נְכֹנִ֖ים לִשְׁלֹ֣שֶׁת יָמִ֑ים אַֽל־תִּגְּשׁ֖וּ אֶל־אִשָּֽׁה׃ \Rashi{\textbf{היו נכנים לשלשת ימים.}לסוף ג' ימים, הוא יום רביעי, שהוסיף משה יום אחד מדעתו כדברי רבי יוסי, ולדברי האומר בששה בחדש נתנו עשרת הדברות, לא הוסיף משה כלום, ולשלשת ימים כמו ליום השלישי: אל תגשו אל אשה. כל ג' ימים הללו, כדי שיהיו הנשים טובלות ליום השלישי ותהיינה טהורות לקבל תורה, שאם ישמשו תוך ג' ימים שמא תפלט האשה שכבת זרע לאחר טבילתה ותחזר ותטמא, אבל מששהתה ג' ימים, כבר הזרע מסריח ואינו ראוי להזריע, וטהור מלטמא את הפולטת (שבת פ"ו):}%
% ---------- פרק 19 | פסוק 16 ----------
 \textbf{טז} וַיְהִי֩ בַיּ֨וֹם הַשְּׁלִישִׁ֜י בִּֽהְיֹ֣ת הַבֹּ֗קֶר וַיְהִי֩ קֹלֹ֨ת וּבְרָקִ֜ים וְעָנָ֤ן כָּבֵד֙ עַל־הָהָ֔ר וְקֹ֥ל שֹׁפָ֖ר חָזָ֣ק מְאֹ֑ד וַיֶּחֱרַ֥ד כׇּל־הָעָ֖ם אֲשֶׁ֥ר בַּֽמַּחֲנֶֽה׃ \Rashi{\textbf{בהית הבקר.}מלמד שהקדים על ידם, מה שאין דרך בשר ודם לעשות כן – שיהא הרב ממתין לתלמיד; וכן מצינו "קום צא אל הבקעה, ואקום ואצא אל הבקעה והנה שם כבוד ה' עמד" (יחזקאל ג'):}%
% ---------- פרק 19 | פסוק 17 ----------
 \textbf{יז} וַיּוֹצֵ֨א מֹשֶׁ֧ה אֶת־הָעָ֛ם לִקְרַ֥את הָֽאֱלֹהִ֖ים מִן־הַֽמַּחֲנֶ֑ה וַיִּֽתְיַצְּב֖וּ בְּתַחְתִּ֥ית הָהָֽר׃ \Rashi{\textbf{לקראת האלהים.}מגיד שהשכינה יצאה לקראתם כחתן היוצא לקראת כלה, וזהו שאמר ה' מסיני בא (דברים ל"ג), ולא נאמר לסיני בא (מכילתא): בתחתית ההר. לפי פשוטו ברגלי ההר; ומדרשו שנתלש ההר ממקומו ונכפה עליהם כגיגית (שבת פ"ח):}%
% ---------- פרק 19 | פסוק 18 ----------
 \textbf{יח} וְהַ֤ר סִינַי֙ עָשַׁ֣ן כֻּלּ֔וֹ מִ֠פְּנֵ֠י אֲשֶׁ֨ר יָרַ֥ד עָלָ֛יו יְהֹוָ֖ה בָּאֵ֑שׁ וַיַּ֤עַל עֲשָׁנוֹ֙ כְּעֶ֣שֶׁן הַכִּבְשָׁ֔ן וַיֶּחֱרַ֥ד כׇּל־הָהָ֖ר מְאֹֽד׃ \Rashi{\textbf{עשן כלו.}אין עשן זה שם דבר, שהרי נקוד השי"ן פתח, אלא לשון פעל כמו אמר, שמר, שמע, לכך תרגומו תנן כלה, ולא תרגם תננא, וכל עשן שבמקרא נקודים קמץ מפני שהם שם דבר: הכבשן. של סיד; יכול ככבשן זה ולא יותר, ת"ל בוער באש עד לב השמים (דברים ד'), ומה ת"ל כבשן? לשבר את האזן מה שהיא יכולה לשמע – נותן לבריות סימן הנכר להם; כיוצא בו כאריה ישאג (הושע י"א), וכי מי נתן כח בארי אלא הוא? והכתוב מושלו כאריה, אלא אנו מכנין ומדמין אותו לבריותיו כדי לשבר את האזן מה שיכולה לשמע; כיוצא בו וקולו כקול מים רבים (יחזקאל מ"ג), וכי מי נתן קול למים אלא הוא? ואתה מכנה אותו, לדמותו לבריותיו – כדי לשבר את האזן (מכילתא):}%
% ---------- פרק 19 | פסוק 19 ----------
 \textbf{יט} וַֽיְהִי֙ ק֣וֹל הַשֹּׁפָ֔ר הוֹלֵ֖ךְ וְחָזֵ֣ק מְאֹ֑ד מֹשֶׁ֣ה יְדַבֵּ֔ר וְהָאֱלֹהִ֖ים יַעֲנֶ֥נּוּ בְקֽוֹל׃ \Rashi{\textbf{הולך וחזק מאד.}מנהג הדיוט כל זמן שהוא מאריך לתקע קולו מחליש וכוהה, אבל כאן הולך וחזק מאד, ולמה כך מתחלה? לשבר אזניהם מה שיכולין לשמע: משה ידבר. כשהיה משה מדבר ומשמיע הדברות לישראל – שהרי לא שמעו מפי הגבורה אלא אנכי ולא יהיה לך – והקב"ה מסיעו לתת בו כח להיות קולו מגביר ונשמע (שם): יעננו בקול. יעננו על דבר הקול, כמו אשר יענה באש (מלכים א י"ח) – על דבר האש, להורידו:}%
% ---------- פרק 19 | פסוק 20 ----------
 \textbf{כ} וַיֵּ֧רֶד יְהֹוָ֛ה עַל־הַ֥ר סִינַ֖י אֶל־רֹ֣אשׁ הָהָ֑ר וַיִּקְרָ֨א יְהֹוָ֧ה לְמֹשֶׁ֛ה אֶל־רֹ֥אשׁ הָהָ֖ר וַיַּ֥עַל מֹשֶֽׁה׃ \Rashi{\textbf{וירד ה' על הר סיני.}יכול ירד עליו ממש, ת"ל כי מן השמים דברתי עמכם (שמות כ'), מלמד שהרכין שמים העליונים והתחתונים והציען על גבי ההר כמצע על המטה וירד כסא הכבוד עליהם (מכילתא):}%
% ---------- פרק 19 | פסוק 21 ----------
 \textbf{כא} וַיֹּ֤אמֶר יְהֹוָה֙ אֶל־מֹשֶׁ֔ה רֵ֖ד הָעֵ֣ד בָּעָ֑ם פֶּן־יֶהֶרְס֤וּ אֶל־יְהֹוָה֙ לִרְא֔וֹת וְנָפַ֥ל מִמֶּ֖נּוּ רָֽב׃ \Rashi{\textbf{העד בעם.}התרה בהם שלא לעלות בהר: פן יהרסו וגו'. שלא יהרסו את מצבם, על ידי שתאותם אל ה', לראות, ויקרבו לצד ההר: ונפל ממנו רב. כל מה שיפל מהם, ואפלו הוא יחידי, חשוב לפני רב (שם): פן יהרסו. כל הריסה מפרדת אסיפת הבנין, אף הנפרדין ממצב אנשים הורסים את המצב:}%
% ---------- פרק 19 | פסוק 22 ----------
 \textbf{כב} וְגַ֧ם הַכֹּהֲנִ֛ים הַנִּגָּשִׁ֥ים אֶל־יְהֹוָ֖ה יִתְקַדָּ֑שׁוּ פֶּן־יִפְרֹ֥ץ בָּהֶ֖ם יְהֹוָֽה׃ \Rashi{\textbf{וגם הכהנים.}אף הבכורות שהעבודה בהם (זבחים קט"ו): הנגשים אל ה'. להקריב קרבנות, אף הם לא יסמכו על חשיבותם לעלות: יתקדשו. יהיו מזמנים להתיצב על עמדן: פן יפרץ. לשון פרצה; יהרג בהם ויעשה בהם פרצה:}%
% ---------- פרק 19 | פסוק 23 ----------
 \textbf{כג} וַיֹּ֤אמֶר מֹשֶׁה֙ אֶל־יְהֹוָ֔ה לֹא־יוּכַ֣ל הָעָ֔ם לַעֲלֹ֖ת אֶל־הַ֣ר סִינָ֑י כִּֽי־אַתָּ֞ה הַעֵדֹ֤תָה בָּ֙נוּ֙ לֵאמֹ֔ר הַגְבֵּ֥ל אֶת־הָהָ֖ר וְקִדַּשְׁתּֽוֹ׃ \Rashi{\textbf{לא יוכל העם.}איני צריך להעיד בהם, שהרי מתרים ועומדין הם היום שלשת ימים, ולא יוכלו לעלות שאין להם רשות:}%
% ---------- פרק 19 | פסוק 24 ----------
 \textbf{כד} וַיֹּ֨אמֶר אֵלָ֤יו יְהֹוָה֙ לֶךְ־רֵ֔ד וְעָלִ֥יתָ אַתָּ֖ה וְאַהֲרֹ֣ן עִמָּ֑ךְ וְהַכֹּהֲנִ֣ים וְהָעָ֗ם אַל־יֶֽהֶרְס֛וּ לַעֲלֹ֥ת אֶל־יְהֹוָ֖ה פֶּן־יִפְרׇץ־בָּֽם׃ \Rashi{\textbf{לך רד.}והעד בהם שנית, שמזרזין את האדם קדם מעשה וחוזרין ומזרזין אותו בשעת מעשה (מכילתא): ועלית אתה ואהרן עמך והכהנים. יכול אף הם עמך, ת"ל ועלית אתה, אמר מעתה, אתה מחצה לעצמך, ואהרן מחצה לעצמו, והכהנים מחצה לעצמם – משה נגש יותר מאהרן, ואהרן יותר מן הכהנים, והעם כל עקר אל יהרסו את מצבם לעלות אל ה' (שם): פן יפרץ בם. אע"פ שהוא נקוד חטף קמץ אינו זז מגזרתו, כך דרך כל תבה שנקדתה מלאפום, כשהיא סמוכה באה במקף, משתנה הנקוד לחטף קמץ:}%
% ---------- פרק 19 | פסוק 25 ----------
 \textbf{כה} וַיֵּ֥רֶד מֹשֶׁ֖ה אֶל־הָעָ֑ם וַיֹּ֖אמֶר אֲלֵהֶֽם׃ \{ס\} \Rashi{\textbf{ויאמר אלהם.}התראה זו:}%
% ---------- פרק 20 | פסוק 1 ----------
 \textbf{א} וַיְדַבֵּ֣ר אֱלֹהִ֔ים אֵ֛ת כׇּל־הַדְּבָרִ֥ים הָאֵ֖לֶּה לֵאמֹֽר׃ \{ס\} \Rashi{\textbf{וידבר אלהים.}אין אלהים אלא דין; לפי שיש פרשיות בתורה שאם עשאן אדם מקבל שכר, ואם לאו אינו מקבל עליהם פרענות, יכול אף עשרת הדברות כן, ת"ל וידבר אלהים – דין לפרע (מכילתא): את כל הדברים האלה. מלמד שאמר הקב"ה עשרת הדברות בדבור אחד, מה שאי אפשר לאדם לומר כן, אם כן מה ת"ל עוד אנכי ולא יהיה לך? שחזר ופרש על כל דבור ודבור בפני עצמו (שם): לאמר. מלמד שהיו עונין על הן הן ועל לאו לאו (שם):}%
% ---------- פרק 20 | פסוק 2 ----------
 \textbf{ב} אָֽנֹכִ֖י֙ יְהֹוָ֣ה אֱלֹהֶ֑֔יךָ אֲשֶׁ֧ר הוֹצֵאתִ֛יךָ מֵאֶ֥רֶץ מִצְרַ֖יִם מִבֵּ֣֥ית עֲבָדִ֑͏ֽים׃ \Rashi{\textbf{אשר הוצאתיך מארץ מצרים.}כדאי היא ההוצאה שתהיו משעבדים לי; ד"א, לפי שנגלה בים כגבור מלחמה ונגלה כאן כזקן מלא רחמים, – שנאמר ותחת רגליו כמעשה לבנת הספיר (שמות כ"ד), זו היתה לפניו בשעת השעבוד, וכעצם השמים, משנגאלו – הואיל ואני משתנה במראות, אל תאמרו שתי רשויות הן, אנכי הוא אשר הוצאתיך ממצרים ועל הים. ד"א לפי שהיו שומעין קולות הרבה, שנאמר את הקולת – קולות באין מארבע רוחות ומן השמים ומן הארץ – אל תאמרו רשויות הרבה הן; ולמה אמר לשון יחיד אלהיך? לתן פתחון פה למשה ללמד סנגוריא במעשה העגל, וזה שאמר למה ה' יחרה אפך בעמך – לא להם צוית לא יהיה לכם אלהים אחרים אלא לי לבדי (שמות רבה): מבית עבדים. מבית פרעה שהייתם עבדים לו, או אינו אומר אלא מבית עבדים שהיו עבדים לעבדים? ת"ל ויפדך מבית עבדים מיד פרעה מלך מצרים (דברים ז'), אמר מעתה עבדים למלך היו ולא עבדים לעבדים (מכילתא):}%
% ---------- פרק 20 | פסוק 3 ----------
 \textbf{ג} לֹֽ֣א־יִהְיֶ֥͏ֽה־לְךָ֛֩ אֱלֹהִ֥֨ים אֲחֵרִ֖֜ים עַל־פָּנָֽ͏ַ֗י׃ \Rashi{\textbf{לא יהיה לך.}למה נאמר? לפי שנ' לא תעשה לך, אין לי אלא שלא יעשה, העשוי כבר מנין שלא יקים? ת"ל לא יהיה לך (שם): אלהים אחרים. שאינן אלוהות אלא אחרים עשאום אלהים עליהם, ולא יתכן לפרש אלהים אחרים זולתי, שגנאי כלפי מעלה לקראותם אלוהות אצלו. ד"א אלהים אחרים, שהם אחרים לעובדיהם – צועקים אליהם ואינן עונין אותם, ודומה כאלו הוא אחר שאינו מכירו מעולם: על פני. כל זמן שאני קים, שלא תאמר לא נצטוו על ע"ז אלא אותו הדור (מכילתא):}%
% ---------- פרק 20 | פסוק 4 ----------
 \textbf{ד} לֹֽ֣א־תַעֲשֶֽׂ֨ה־לְךָ֥֣ פֶ֣֙סֶל֙ ׀ וְכׇל־תְּמוּנָ֔֡ה אֲשֶׁ֤֣ר בַּשָּׁמַ֣֙יִם֙ ׀ מִמַּ֔֡עַל וַֽאֲשֶׁ֥ר֩ בָּאָ֖֨רֶץ מִתָּ֑͏ַ֜חַת וַאֲשֶׁ֥ר בַּמַּ֖֣יִם ׀ מִתַּ֥֣חַת לָאָֽ֗רֶץ׃ \Rashi{\textbf{פסל.}על שם שנפסל: וכל תמונה. תמונת דבר אשר בשמים:}%
% ---------- פרק 20 | פסוק 5 ----------
 \textbf{ה} לֹֽא־תִשְׁתַּחֲוֶ֥֣ה לָהֶ֖ם֮ וְלֹ֣א תׇעׇבְדֵ֑ם֒ כִּ֣י אָֽנֹכִ֞י יְהֹוָ֤ה אֱלֹהֶ֙יךָ֙ אֵ֣ל קַנָּ֔א פֹּ֠קֵ֠ד עֲוֺ֨ן אָבֹ֧ת עַל־בָּנִ֛ים עַל־שִׁלֵּשִׁ֥ים וְעַל־רִבֵּעִ֖ים לְשֹׂנְאָֽ֑י׃ \Rashi{\textbf{אל קנא.}מקנא לפרע ואינו עובר על מדתו למחל על עבודת אלילים; כל לשון קנא אנפרי"מנט בלעז – נותן לב לפרע: לשנאי. כתרגומו, כשאוחזין מעשה אבותיהם בידיהם.}%
% ---------- פרק 20 | פסוק 6 ----------
 \textbf{ו} וְעֹ֥֤שֶׂה חֶ֖֙סֶד֙ לַאֲלָפִ֑֔ים לְאֹהֲבַ֖י וּלְשֹׁמְרֵ֥י מִצְוֺתָֽי׃ \{ס\} \Rashi{\textbf{נוצר חסד.}שאדם עושה, לשלם שכר עד לאלפים דור, נמצאת מדה טובה יתרה על מדת פרעניות אחת על חמש מאות, שזו לארבעה דורות, וזו לאלפים (תוספתא סוטה ד, א):}%
% ---------- פרק 20 | פסוק 7 ----------
 \textbf{ז} לֹ֥א תִשָּׂ֛א אֶת־שֵֽׁם־יְהֹוָ֥ה אֱלֹהֶ֖יךָ לַשָּׁ֑וְא כִּ֣י לֹ֤א יְנַקֶּה֙ יְהֹוָ֔ה אֵ֛ת אֲשֶׁר־יִשָּׂ֥א אֶת־שְׁמ֖וֹ לַשָּֽׁוְא׃ \{פ\} \Rashi{\textbf{לשוא.}חנם, להבל, ואיזהו שבועת שוא? נשבע לשנות את הידוע – על עמוד של אבן שהוא של זהב (שבועות כ"ט):}%
% ---------- פרק 20 | פסוק 8 ----------
 \textbf{ח} זָכ֛וֹר֩ אֶת־י֥֨וֹם הַשַּׁבָּ֖֜ת לְקַדְּשֽׁ֗וֹ׃ \Rashi{\textbf{זכור.}זכור ושמור בדבור אחד נאמרו, וכן מחלליה מות יומת (שמות ל"א), וביום השבת שני כבשים (במדבר כ"ח); וכן לא תלבש שעטנז, גדלים תעשה לך (דברים כ"ב); וכן ערות אשת אחיך (ויקרא י"ח), יבמה יבא עליה (דברים כ"ה); הוא שנאמר אחת דבר אלהים שתים זו שמעתי (תהילים ס"ב): זכור. לשון פעול הוא, כמו אכול ושתו (ישעיהו כ"ב), הלוך ובכה (שמואל ב ג'), וכן פתרונו: תנו לב לזכר תמיד את יום השבת, שאם נזדמן לך חפץ יפה תהא מזמינו לשבת (מכילתא):}%
% ---------- פרק 20 | פסוק 9 ----------
 \textbf{ט} שֵׁ֤֣שֶׁת יָמִ֣ים֙ תַּֽעֲבֹ֔ד֮ וְעָשִׂ֖֣יתָ כׇּֿל־מְלַאכְתֶּֽךָ֒׃ \Rashi{\textbf{ועשית כל מלאכתך.}כשתבא שבת יהא בעיניך כאלו כל מלאכתך עשויה, שלא תהרהר אחר מלאכה (שם):}%
% ---------- פרק 20 | פסוק 10 ----------
 \textbf{י} וְי֨וֹם֙ הַשְּׁבִיעִ֔֜י שַׁבָּ֖֣ת ׀ לַיהֹוָ֣ה אֱלֹהֶ֑֗יךָ לֹֽ֣א־תַעֲשֶׂ֣֨ה כׇל־מְלָאכָ֜֡ה אַתָּ֣ה ׀ וּבִנְךָ֣͏ֽ־וּ֠בִתֶּ֗ךָ עַבְדְּךָ֤֨ וַאֲמָֽתְךָ֜֙ וּבְהֶמְתֶּ֔֗ךָ וְגֵרְךָ֖֙ אֲשֶׁ֥֣ר בִּשְׁעָרֶֽ֔יךָ׃ \Rashi{\textbf{אתה ובנך ובתך.}אלו קטנים; או אינו אלא גדולים? אמרת, הרי כבר מזהרים הם, אלא לא בא אלא להזהיר גדולים על שביתת הקטנים; וזהו ששנינו (שבת קכ"א), קטן שבא לכבות אין שומעין לו מפני ששביתתו עליך:}%
% ---------- פרק 20 | פסוק 11 ----------
 \textbf{יא} כִּ֣י שֵֽׁשֶׁת־יָמִים֩ עָשָׂ֨ה יְהֹוָ֜ה אֶת־הַשָּׁמַ֣יִם וְאֶת־הָאָ֗רֶץ אֶת־הַיָּם֙ וְאֶת־כׇּל־אֲשֶׁר־בָּ֔ם וַיָּ֖נַח בַּיּ֣וֹם הַשְּׁבִיעִ֑י עַל־כֵּ֗ן בֵּרַ֧ךְ יְהֹוָ֛ה אֶת־י֥וֹם הַשַּׁבָּ֖ת וַֽיְקַדְּשֵֽׁהוּ׃ \{ס\} \Rashi{\textbf{וינח ביום השביעי.}כביכול הכתיב בעצמו מנוחה, ללמד הימנו קל וחמר לאדם שמלאכתו בעמל ויגיעה שיהיה נח בשבת (מכילתא): ברך, ויקדשהו. ברכו במן – לכפלו בששי לחם משנה, וקדשו במן – שלא היה יורד בו (מכילתא):}%
% ---------- פרק 20 | פסוק 12 ----------
 \textbf{יב} כַּבֵּ֥ד אֶת־אָבִ֖יךָ וְאֶת־אִמֶּ֑ךָ לְמַ֙עַן֙ יַאֲרִכ֣וּן יָמֶ֔יךָ עַ֚ל הָאֲדָמָ֔ה אֲשֶׁר־יְהֹוָ֥ה אֱלֹהֶ֖יךָ נֹתֵ֥ן לָֽךְ׃ \{ס\} \Rashi{\textbf{למען יארכון ימיך.}אם תכבד יאריכון ואם לאו יקצרון, שדברי תורה נוטריקון הם נדרשים – מכלל הן לאו ומכלל לאו הן (מכילתא):}%
% ---------- פרק 20 | פסוק 13 ----------
 \textbf{יג} לֹ֥֖א תִּֿרְצָ֖͏ֽח׃ \{ס\}        לֹ֣֖א תִּֿנְאָ֑͏ֽף׃ \{ס\}        לֹ֣֖א תִּֿגְנֹֽ֔ב׃ \{ס\}        לֹֽא־תַעֲנֶ֥ה בְרֵעֲךָ֖ עֵ֥ד שָֽׁקֶר׃ \{ס\} \Rashi{\textbf{לא תנאף.}אין נאוף אלא באשת איש, שנאמר מות יומת הנואף והנואפת (ויקרא כ'), ואומר האשה המנאפת תחת אישה תקח את זרים (יחזקאל ט"ז): לא תגנב. בגונב נפשות הכתוב מדבר, לא תגנבו (ויקרא י"ט), בגונב ממון; או אינו אלא זה בגונב ממון ולהלן בגונב נפשות? אמרת דבר למד מענינו – מה לא תרצח, לא תנאף, מדבר בדבר שחיבין עליו מיתת בית דין, אף לא תגנב דבר שחיב עליו מיתת בית דין (סנהדרין פ"ו):}%
% ---------- פרק 20 | פסוק 14 ----------
 \textbf{יד} לֹ֥א תַחְמֹ֖ד בֵּ֣ית רֵעֶ֑ךָ \{ס\}        לֹֽא־תַחְמֹ֞ד אֵ֣שֶׁת רֵעֶ֗ךָ וְעַבְדּ֤וֹ וַאֲמָתוֹ֙ וְשׁוֹר֣וֹ וַחֲמֹר֔וֹ וְכֹ֖ל אֲשֶׁ֥ר לְרֵעֶֽךָ׃ \{פ\} %
% ---------- פרק 20 | פסוק 15 ----------
 \textbf{טו} וְכׇל־הָעָם֩ רֹאִ֨ים אֶת־הַקּוֹלֹ֜ת וְאֶת־הַלַּפִּידִ֗ם וְאֵת֙ ק֣וֹל הַשֹּׁפָ֔ר וְאֶת־הָהָ֖ר עָשֵׁ֑ן וַיַּ֤רְא הָעָם֙ וַיָּנֻ֔עוּ וַיַּֽעַמְד֖וּ מֵֽרָחֹֽק׃ \Rashi{\textbf{וכל העם ראים.}מלמד שלא היה בהם אחד סומא, ומנין שלא היה בהם אלם? ת"ל ויענו כל העם, ומנין שלא היה בהם חרש? ת"ל נעשה ונשמע (מכילתא): ראים את הקולת. רואין את הנשמע, שאי אפשר לראות במקום אחר: את הקולת. היוצאין מפי הגבורה: וינעו. אין נוע אלא זיע: ויעמדו מרחק. היו נרתעין לאחוריהם שנים עשר מיל, כארך מחניהם, ומלאכי השרת באין ומסיעין אותן להחזירם, שנאמר (תהילים ס"ח) מלאכי צבאות ידדון ידדון (שבת פ"ח):}%
% ---------- פרק 20 | פסוק 16 ----------
 \textbf{טז} וַיֹּֽאמְרוּ֙ אֶל־מֹשֶׁ֔ה דַּבֵּר־אַתָּ֥ה עִמָּ֖נוּ וְנִשְׁמָ֑עָה וְאַל־יְדַבֵּ֥ר עִמָּ֛נוּ אֱלֹהִ֖ים פֶּן־נָמֽוּת׃ %
% ---------- פרק 20 | פסוק 17 ----------
 \textbf{יז} וַיֹּ֨אמֶר מֹשֶׁ֣ה אֶל־הָעָם֮ אַל־תִּירָ֒אוּ֒ כִּ֗י לְבַֽעֲבוּר֙ נַסּ֣וֹת אֶתְכֶ֔ם בָּ֖א הָאֱלֹהִ֑ים וּבַעֲב֗וּר תִּהְיֶ֧ה יִרְאָת֛וֹ עַל־פְּנֵיכֶ֖ם לְבִלְתִּ֥י תֶחֱטָֽאוּ׃ \Rashi{\textbf{לבעבור נסות אתכם.}לגדל אתכם בעולם, שיצא לכם שם באמות שהוא בכבודו נגלה עליכם: נסות. לשון הרמה וגדלה, כמו הרימו נס (ישעיהו ס"ב), ארים נסי (שם מ"ט), כנס על הגבעה (שם ל') – שהוא זקוף: ובעבור תהיה יראתו. על ידי שראיתם אותו יראוי ומאים, תדעו כי אין זולתו ותיראו מפניו:}%
% ---------- פרק 20 | פסוק 18 ----------
 \textbf{יח} וַיַּעֲמֹ֥ד הָעָ֖ם מֵרָחֹ֑ק וּמֹשֶׁה֙ נִגַּ֣שׁ אֶל־הָֽעֲרָפֶ֔ל אֲשֶׁר־שָׁ֖ם הָאֱלֹהִֽים׃ \{ס\} \Rashi{\textbf{נגש אל הערפל.}לפנים משלש מחצות: חשך, ענן, וערפל; שנאמר וההר בער באש עד לב השמים חשך ענן וערפל (דברים ד'); ערפל הוא עב הענן שאמר לו הנה אנכי בא אליך בעב הענן (שמות י"ט):}%
% ---------- פרק 20 | פסוק 19 ----------
 \textbf{יט} וַיֹּ֤אמֶר יְהֹוָה֙ אֶל־מֹשֶׁ֔ה כֹּ֥ה תֹאמַ֖ר אֶל־בְּנֵ֣י יִשְׂרָאֵ֑ל אַתֶּ֣ם רְאִיתֶ֔ם כִּ֚י מִן־הַשָּׁמַ֔יִם דִּבַּ֖רְתִּי עִמָּכֶֽם׃ \Rashi{\textbf{כה תאמר.}בלשון הזה: אתם ראיתם. יש הפרש בין מה שאדם רואה למה שאחרים משיחין לו, שמה שאחרים משיחין לו פעמים שלבו חלוק מלהאמין (מכילתא): כי מן השמים דברתי. וכתוב אחד אומר וירד ה' על הר סיני (שמות י"ט), בא הכתוב השלישי והכריע ביניהם – מן השמים השמיעך את קולו ליסרך ועל הארץ הראך את אשו הגדולה (דברים ד'), כבודו בשמים ואשו וגבורתו על הארץ; ד"א: הרכין שמים ושמי השמים והציען על ההר, וכן הוא אומר ויט שמים וירד (תהילים י"ח):}%
% ---------- פרק 20 | פסוק 20 ----------
 \textbf{כ} לֹ֥א תַעֲשׂ֖וּן אִתִּ֑י אֱלֹ֤הֵי כֶ֙סֶף֙ וֵאלֹהֵ֣י זָהָ֔ב לֹ֥א תַעֲשׂ֖וּ לָכֶֽם׃ \Rashi{\textbf{לא תעשון אתי.}לא תעשון דמות שמשי המשמשים לפני במרום (מכילתא): אלהי כסף. בא להזהיר על הכרובים שאתה עושה לעמד אתי, שלא יהיו של כסף, שאם שניתם לעשותם של כסף הרי הן לפני כאלוהות: ואלהי זהב. בא להזהיר שלא יוסיף על שנים, שאם עשית ארבעה הרי הן לפני כאלהי זהב (מכילתא): לא תעשו לכם. לא תאמר הריני עושה כרובים בבתי כנסיות ובבתי מדרשות כדרך שאני עושה בבית עולמים, לכך נאמר לא תעשו לכם (שם):}%
% ---------- פרק 20 | פסוק 21 ----------
 \textbf{כא} מִזְבַּ֣ח אֲדָמָה֮ תַּעֲשֶׂה־לִּי֒ וְזָבַחְתָּ֣ עָלָ֗יו אֶת־עֹלֹתֶ֙יךָ֙ וְאֶת־שְׁלָמֶ֔יךָ אֶת־צֹֽאנְךָ֖ וְאֶת־בְּקָרֶ֑ךָ בְּכׇל־הַמָּקוֹם֙ אֲשֶׁ֣ר אַזְכִּ֣יר אֶת־שְׁמִ֔י אָב֥וֹא אֵלֶ֖יךָ וּבֵרַכְתִּֽיךָ׃ \Rashi{\textbf{מזבח אדמה.}מחבר באדמה, שלא יבננו על גבי עמודים או על גבי בסיס; ד"א שהיה ממלא את חלל מזבח הנחשת אדמה בשעת חניתן (שם): תעשה לי. שתהא תחלת עשיתו לשמי: וזבחת עליו. אצלו, כמו ועליו מטה מנשה (במדבר ב'), או אינו אלא עליו ממש? ת"ל הבשר והדם על מזבח ה' אלהיך (דברים י"ב), ואין שחיטה בראש המזבח: את עלתיך ואת שלמיך. אשר מצאנך ומבקרך: (את צאנך ואת בקרך. פרוש לאת עולתיך ואת שלמיך:) בכל המקום אשר אזכיר את שמי. אשר אתן לך רשות להזכיר שם המפרש שלי, שם אבוא אליך וברכתיך – אשרה שכינתי עליך; מכאן אתה למד שלא נתן רשות להזכיר שם המפרש אלא במקום שהשכינה באה שם, וזהו בית הבחירה, שם נתן רשות לכהנים להזכיר שם המפרש בנשיאת כפים לברך את העם (סוטה ל"ח):}%
% ---------- פרק 20 | פסוק 22 ----------
 \textbf{כב} וְאִם־מִזְבַּ֤ח אֲבָנִים֙ תַּֽעֲשֶׂה־לִּ֔י לֹֽא־תִבְנֶ֥ה אֶתְהֶ֖ן גָּזִ֑ית כִּ֧י חַרְבְּךָ֛ הֵנַ֥פְתָּ עָלֶ֖יהָ וַתְּחַֽלְלֶֽהָ׃ \Rashi{\textbf{ואם מזבח אבנים.}רבי ישמעאל אומר כל אם ואם שבתורה רשות חוץ מג', ואם מזבח אבנים תעשה לי, הרי אם זה משמש בלשון כאשר – וכאשר תעשה לי מזבח אבנים לא תבנה אתהן גזית, שהרי חובה עליך לבנות מזבח אבנים, שנאמר אבנים שלמות תבנה (דברים כ"ז) – וכן אם כסף תלוה (שמות כ"ב), חובה הוא, שנאמר והעבט תעביטנו (דברים ט"ו) – ואף זה משמש בלשון כאשר – וכן ואם תקריב מנחת בכורים (ויקרא ב'), זו מנחת העמר שהיא חובה, ועל כרחך אין אם הללו תלויין אלא ודאין, ובלשון כאשר הם משמשים: גזית. לשון גזיזה, שפוסלן ומכתתן בברזל: כי חרבך הנפת עליה. הרי כי זה משמש בלשון פן, שהוא דילמא – פן תניף חרבך עליה: ותחללה. הא למדת שאם הנפת עליה ברזל חללת, שהמזבח נברא להאריך ימיו של אדם והברזל נברא לקצר ימיו של אדם, אין זה בדין שיונף המקצר על המאריך (מכילתא). ועוד שהמזבח מטיל שלום בין ישראל לאביהם שבשמים, לפיכך לא יבא עליו כורת ומחבל; והרי דברים קל וחמר, ומה אבנים שאינן רואות ולא שומעות ולא מדברות, על ידי שמטילות שלום אמרה תורה ולא תניף עליהם ברזל, המטיל שלום בין איש לאשתו, בין משפחה למשפחה, בין אדם לחברו, על אחת כמה וכמה שלא תבואהו פרענות (שם):}%
% ---------- פרק 20 | פסוק 23 ----------
 \textbf{כג} וְלֹֽא־תַעֲלֶ֥ה בְמַעֲלֹ֖ת עַֽל־מִזְבְּחִ֑י אֲשֶׁ֛ר לֹֽא־תִגָּלֶ֥ה עֶרְוָתְךָ֖ עָלָֽיו׃ \{פ\} \Rashi{\textbf{ולא תעלה במעלת.}כשאתה בונה כבש למזבח לא תעשהו מעלות מעלות, אשקלו"נש בלעז, אלא חלק יהא ומשפע: אשר לא תגלה ערותך. שעל ידי המעלות אתה צריך להרחיב פסיעותיך; ואע"פ שאינו גלוי ערוה ממש, שהרי כתיב ועשה להם מכנסי בד, מכל מקום הרחבת הפסיעות קרוב לגלוי ערוה הוא, ואתה נוהג בם מנהג בזיון; והרי דברים ק"ו, ומה אבנים הללו שאין בהם דעת להקפיד על בזיונן, אמרה תורה הואיל ויש בהם צרך לא תנהג בהם מנהג בזיון, חברך שהוא בדמות יוצרך, ומקפיד על בזיונו, על אחת כמה וכמה:}%
\addparasha{חומש שמות | פרשת משפטים}
% ---------- פרק 21 | פסוק 1 ----------
 \textbf{א} וְאֵ֙לֶּה֙ הַמִּשְׁפָּטִ֔ים אֲשֶׁ֥ר תָּשִׂ֖ים לִפְנֵיהֶֽם׃ \Rashi{\textbf{ואלה המשפטים.}כל מקום שנאמר "אלה" פסל את הראשונים, "ואלה" מוסיף על הראשונים, מה הראשונים מסיני, אף אלו מסיני; ולמה נסמכה פרשת דינין לפרשת מזבח? לומר לך, שתשים סנהדרין אצל המקדש (מכילתא): אשר תשים לפניהם. אמר לו הקב"ה למשה: לא תעלה על דעתך לומר, אשנה להם הפרק וההלכה ב' או ג' פעמים, עד שתהא סדורה בפיהם כמשנתה, ואיני מטריח עצמי להבינם טעמי הדבר ופרושו, לכך נאמר אשר תשים לפניהם – כשלחן הערוך ומוכן לאכל לפני האדם (שם): לפניהם. ולא לפני גוים, ואפלו ידעת בדין אחד שהם דנין אותו כדיני ישראל, אל תביאהו בערכאות שלהם, שהמביא דיני ישראל לפני גוים, מחלל את השם ומיקר שם ע"ז להשביחה, שנאמר כי לא כצורנו צורם ואיבינו פלילים (דברים ל"ב) – כשאויבינו פלילים זהו עדות לעלוי יראתם (תנחומא):}%
% ---------- פרק 21 | פסוק 2 ----------
 \textbf{ב} כִּ֤י תִקְנֶה֙ עֶ֣בֶד עִבְרִ֔י שֵׁ֥שׁ שָׁנִ֖ים יַעֲבֹ֑ד וּבַ֨שְּׁבִעִ֔ת יֵצֵ֥א לַֽחׇפְשִׁ֖י חִנָּֽם׃ \Rashi{\textbf{כי תקנה עבד עברי.}עבד שהוא עברי; או אינו אלא עבדו של עברי – עבד כנעני שלקחתו מישראל – ועליו הוא אומר שש שנים יעבד? ומה אני מקים והתנחלתם אתם (ויקרא כ"ה), בלקוח מן הגוי, אבל בלקוח מישראל יצא בשש? ת"ל כי ימכר לך אחיך העברי (דברים ט"ו), לא אמרתי אלא באחיך: כי תקנה. מיד בית דין שמכרוהו בגנבתו, כמו שנאמר אם אין לו ונמכר בגנבתו, או אינו אלא במוכר עצמו מפני דחקו, אבל מכרוהו בית דין לא יצא בשש? כשהוא אומר וכי ימוך אחיך עמך ונמכר לך (ויקרא כ"ה), הרי מוכר עצמו מפני דחקו אמור, ומה אני מקים כי תקנה? בנמכר בבית דין: לחפשי. לחרות:}%
% ---------- פרק 21 | פסוק 3 ----------
 \textbf{ג} אִם־בְּגַפּ֥וֹ יָבֹ֖א בְּגַפּ֣וֹ יֵצֵ֑א אִם־בַּ֤עַל אִשָּׁה֙ ה֔וּא וְיָצְאָ֥ה אִשְׁתּ֖וֹ עִמּֽוֹ׃ \Rashi{\textbf{אם בגפו יבא.}שלא היה נשוי אשה, כתרגומו אם בלחודוהי, ולשון בגפו – בכנפו, שלא בא אלא כמות שהוא, יחידי, בתוך לבושו – בכנף בגדו: בגפו יצא. מגיד שאם לא היה נשוי מתחלה, אין רבו מוסר לו שפחה כנענית להוליד ממנה עבדים (קידושין כ'): אם בעל אשה הוא. ישראלית: ויצאה אשתו עמו. וכי מי הכניסה שתצא? אלא מגיד הכתוב שהקונה עבד עברי חיב במזונות אשתו ובניו (שם כ"ב):}%
% ---------- פרק 21 | פסוק 4 ----------
 \textbf{ד} אִם־אֲדֹנָיו֙ יִתֶּן־ל֣וֹ אִשָּׁ֔ה וְיָלְדָה־ל֥וֹ בָנִ֖ים א֣וֹ בָנ֑וֹת הָאִשָּׁ֣ה וִילָדֶ֗יהָ תִּהְיֶה֙ לַֽאדֹנֶ֔יהָ וְה֖וּא יֵצֵ֥א בְגַפּֽוֹ׃ \Rashi{\textbf{אם אדניו יתן לו אשה.}מכאן שהרשות ביד רבו למסר לו שפחה כנענית להוליד ממנה עבדים. או אינו אלא בישראלית? ת"ל האשה וילדיה תהיה לאדניה, הא אינו מדבר אלא בכנענית, שהרי העבריה אף היא יוצאה בשש – ואפלו לפני שש אם הביאה סימנין יוצאה – שנאמר אחיך העברי או העבריה (דברים ט"ו), מלמד שאף העבריה יוצאה בשש (מכילתא):}%
% ---------- פרק 21 | פסוק 5 ----------
 \textbf{ה} וְאִם־אָמֹ֤ר יֹאמַר֙ הָעֶ֔בֶד אָהַ֙בְתִּי֙ אֶת־אֲדֹנִ֔י אֶת־אִשְׁתִּ֖י וְאֶת־בָּנָ֑י לֹ֥א אֵצֵ֖א חׇפְשִֽׁי׃ \Rashi{\textbf{את אשתי.}השפחה:}%
% ---------- פרק 21 | פסוק 6 ----------
 \textbf{ו} וְהִגִּישׁ֤וֹ אֲדֹנָיו֙ אֶל־הָ֣אֱלֹהִ֔ים וְהִגִּישׁוֹ֙ אֶל־הַדֶּ֔לֶת א֖וֹ אֶל־הַמְּזוּזָ֑ה וְרָצַ֨ע אֲדֹנָ֤יו אֶת־אׇזְנוֹ֙ בַּמַּרְצֵ֔עַ וַעֲבָד֖וֹ לְעֹלָֽם׃ \{ס\} \Rashi{\textbf{אל האלהים.}לבית דין; צריך שימלך במוכריו שמכרוהו לו (שם): אל הדלת או אל המזוזה. יכול שתהא המזוזה כשרה לרצע עליה, ת"ל ונתתה באזנו ובדלת, בדלת ולא במזוזה, הא מה ת"ל או אל המזוזה? הקיש דלת למזוזה, מה מזוזה מעומד אף דלת מעומד (שם): ורצע אדניו את אזנו במרצע. הימנית, או אינו, אלא של שמאל? ת"ל אזן אזן לגזרה שוה, נאמר כאן ורצע אדניו את אזנו, ונאמר במצרע תנוך אזן המטהר הימנית (ויקרא י"ד), מה להלן הימנית, אף כאן הימנית; ומה ראה אזן לרצע מכל שאר אברים שבגוף? אמר רבן יוחנן בן זכאי: אזן זאת ששמעה על הר סיני לא תגנב, והלך וגנב, תרצע. ואם מוכר עצמו, אזן ששמעה על הר סיני כי לי בני ישראל עבדים, והלך וקנה אדון לעצמו, תרצע. רבי שמעון היה דורש מקרא זה כמין חמר: מה נשתנו דלת ומזוזה מכל כלים שבבית? אמר הקב"ה דלת ומזוזה שהיו עדים במצרים כשפסחתי על המשקוף ועל שתי המזוזות ואמרתי כי לי בני ישראל עבדים – עבדי הם ולא עבדים לעבדים – והלך זה וקנה אדון לעצמו, ירצע בפניהם (קידושין כ"ב): ועבדו לעלם. עד היובל; או אינו אלא לעולם כמשמעו? ת"ל ואיש אל משפחתו תשובו (ויקרא כ"ה), מגיד שחמשים שנה קרויים עולם; ולא שיהא עובדו כל חמשים שנה, אלא עובדו עד היובל, בין סמוך בין מפלג (קידושין ט"ו):}%
% ---------- פרק 21 | פסוק 7 ----------
 \textbf{ז} וְכִֽי־יִמְכֹּ֥ר אִ֛ישׁ אֶת־בִּתּ֖וֹ לְאָמָ֑ה לֹ֥א תֵצֵ֖א כְּצֵ֥את הָעֲבָדִֽים׃ \Rashi{\textbf{וכי ימכר איש את בתו לאמה.}בקטנה הכתוב מדבר; יכול אפלו הביאה סימנים, אמרת קל וחמר, ומה מכורה קדם לכן יוצאה בסימנין – כמו שכתוב ויצאה חנם אין כסף, שאנו דורשין אותו לסימני נערות – שאינה מכורה אינו דין שלא תמכר (ערכין כ"ט): לא תצא כצאת העבדים. כיציאת עבדים כנענים שיוצאים בשן ועין, אבל זו לא תצא בשן ועין אלא עובדת שש, או עד היובל, או עד שתביא סימנין, וכל הקודם קודם לחרותה, ונותן לה דמי עינה או דמי שנה. או אינו אלא לא תצא כצאת העבדים בשש וביובל? ת"ל כי ימכר לך אחיך העברי או העבריה, מקיש עבריה לעברי לכל יציאותיו, מה עברי יוצא בשש וביובל, אף עבריה יוצאה בשש וביובל, ומהו לא תצא כצאת העבדים? לא תצא בראשי אברים כעבדים כנעניים; יכול העברי יוצא בראשי אברים? ת"ל העברי או העבריה, מקיש עברי לעבריה, מה העבריה אינה יוצאה בראשי אברים, אף הוא אינו יוצא בראשי אברים (מכילתא):}%
% ---------- פרק 21 | פסוק 8 ----------
 \textbf{ח} אִם־רָעָ֞ה בְּעֵינֵ֧י אֲדֹנֶ֛יהָ אֲשֶׁר־[ל֥וֹ] (לא) יְעָדָ֖הּ וְהֶפְדָּ֑הּ לְעַ֥ם נׇכְרִ֛י לֹא־יִמְשֹׁ֥ל לְמׇכְרָ֖הּ בְּבִגְדוֹ־בָֽהּ׃ \Rashi{\textbf{אם רעה בעיני אדניה.}שלא נשאה חן בעיניו לכנסה: אשר לא יעדה. שהיה לו ליעדה ולהכניסה לו לאשה, וכסף קניתה הוא כסף קדושיה; וכאן רמז לך הכתוב שמצוה ביעוד, ורמז לך שאינה צריכה קדושין אחרים (קידושין י"ח): והפדה. יתן לה מקום להפדות ולצאת – שאף הוא מסיע בפדיונה; ומה הוא מקום שנותן לה? שמגרע מפדיונה במספר השנים שעשתה אצלו כאלו היא שכורה אצלו. כיצד? הרי שקנאה במנה ועשתה אצלו שתי שנים, אומרים לו, יודע היית שעתידה לצאת לסוף שש, נמצא שקנית עבודת כל שנה ושנה בששית המנה, ועשתה אצלך שתי שנים, הרי שלישית המנה, טל שני שלישיות המנה ותצא מאצלך: לעם נכרי לא ימשל למכרה. שאינו רשאי למכרה לאחר לא האדון ולא האב (שם): בבגדו בה. אם בא לבגד בה, שלא לקים בה מצות יעוד, וכן אביה מאחר שבגד בה ומכרה לזה:}%
% ---------- פרק 21 | פסוק 9 ----------
 \textbf{ט} וְאִם־לִבְנ֖וֹ יִֽיעָדֶ֑נָּה כְּמִשְׁפַּ֥ט הַבָּנ֖וֹת יַעֲשֶׂה־לָּֽהּ׃ \Rashi{\textbf{ואם לבנו ייעדנה.}האדון; מלמד שאף בנו קם תחתיו ליעדה אם ירצה אביו, ואינו צריך לקדשה קדושין אחרים, אלא אומר לה הרי את מיעדת לי בכסף שקבל אביך בדמיך: כמשפט הבנות. שאר כסות ועונה (מכילתא):}%
% ---------- פרק 21 | פסוק 10 ----------
 \textbf{י} אִם־אַחֶ֖רֶת יִֽקַּֽח־ל֑וֹ שְׁאֵרָ֛הּ כְּסוּתָ֥הּ וְעֹנָתָ֖הּ לֹ֥א יִגְרָֽע׃ \Rashi{\textbf{אם אחרת יקח לו.}עליה: שארה כסותה וענתה לא יגרע. מן האמה שיעד לו כבר: שארה. מזונות: כסותה. כמשמעו. ענתה. תשמיש (כתובות מ"ז):}%
% ---------- פרק 21 | פסוק 11 ----------
 \textbf{יא} וְאִ֨ם־שְׁלׇשׁ־אֵ֔לֶּה לֹ֥א יַעֲשֶׂ֖ה לָ֑הּ וְיָצְאָ֥ה חִנָּ֖ם אֵ֥ין כָּֽסֶף׃ \{ס\} \Rashi{\textbf{ואם שלש אלה לא יעשה לה.}אם אחת משלש אלה לא יעשה לה; ומה הן השלש? ייעדנה לו, או לבנו, או יגרע מפדיונה ותצא, וזה לא יעדה, לא לו ולא לבנו, והיא לא היה בידה לפדות את עצמה: ויצאה חנם. רבה לה יציאה לזו יותר ממה שרבה לעבדים, ומה היא היציאה? למדך שתצא בסימנין ותשהה עמו עוד עד שתביא סימנין, – ואם הגיעו שש שנים קדם סימנין, כבר למדנו שתצא, שנאמר העברי או העבריה ועבדך שש שנים (דברים ט"ו) – ומהו האמור כאן ויצאה חנם? שאם קדמו סימנים לשש שנים תצא בהן. או אינו אומר שתצא אלא בבגרות? ת"ל אין כסף לרבות יציאת בגרות; ואם לא נאמרו שניהם הייתי אומר ויצאה חנם זו בגרות, לכך נאמרו שניהם שלא לתן פתחון פה לבעל הדין לחלק (מכילתא):}%
% ---------- פרק 21 | פסוק 12 ----------
 \textbf{יב} מַכֵּ֥ה אִ֛ישׁ וָמֵ֖ת מ֥וֹת יוּמָֽת׃ \Rashi{\textbf{מכה איש ומת.}כמה כתובים נאמרו בפרשת רוצחין, ומה שבידי לפרש למה באו כלם, אפרש: מכה איש ומת. למה נאמר? לפי שנאמר ואיש כי יכה כל נפש אדם (ויקרא כ"ד), שומע אני הכאה בלא מיתה, תלמוד לומר מכה איש ומת – אינו חיב אלא בהכאה של מיתה; ואם נאמר מכה איש ולא נאמר ואיש כי יכה, הייתי אומר אינו חיב עד שיכה איש, הכה את האשה ואת הקטן מנין? ת"ל כי יכה כל נפש אדם – אפלו קטן ואפלו אשה; ועוד, אלו נאמר מכה איש, שומע אני אפלו קטן שהכה והרג יהא חיב, ת"ל ואיש כי יכה, ולא קטן שהכה; ועוד כי יכה כל נפש אדם אפלו נפלים במשמע, ת"ל מכה איש, אינו חיב עד שיכה בן קימא הראוי להיות איש (מכילתא):}%
% ---------- פרק 21 | פסוק 13 ----------
 \textbf{יג} וַאֲשֶׁר֙ לֹ֣א צָדָ֔ה וְהָאֱלֹהִ֖ים אִנָּ֣ה לְיָד֑וֹ וְשַׂמְתִּ֤י לְךָ֙ מָק֔וֹם אֲשֶׁ֥ר יָנ֖וּס שָֽׁמָּה׃ \{ס\} \Rashi{\textbf{ואשר לא צדה.}לא ארב לו ולא נתכון: צדה. לשון ארב, וכן הוא אומר "ואתה צדה את נפשי" (שמואל א כ"ד); ולא יתכן לומר צדה לשון "הצד ציד" (בראשית כ"ז), שצידת חיות אין נופל ה"א בפעל שלה, ושם דבר בה ציד, וזה שם דבר בו צדיה, ופעל שלו צודה, וזה הפעל שלו צד. ואומר אני פתרונו כתרגומו – ודלא כמן לה. ומנחם חברו בחלק צד ציד, ואין אני מודה לו; ואם יש לחברו באחת ממחלקות של צד, נחברנו בחלק "על צד תנשאו" (ישעיהו ס"ו), "צדה אורה" (שמואל א כ'), "ומלין לצד עלאה ימלל" (דניאל ז'), אף כאן אשר לא צדה – לא צדד למצא לו שום צד מיתה; ואף זה יש להרהר עליו, מכל מקום לשון אורב הוא: והאלהים אנה לידו. זמן לידו, לשון "לא תאנה אליך רעה" (תהילים צ"א), "לא יאנה לצדיק כל און" (משלי י"ב), "מתאנה הוא לי" (מלכים ב ה') – מזדמן למצא לי עלה: והאלהים אנה לידו. ולמה תצא זאת מלפניו? הוא שאמר דוד "כאשר יאמר משל הקדמני מרשעים יצא רשע" (שמואל א כ"ד); ומשל הקדמוני היא התורה, שהיא משל הקב"ה שהוא קדמונו של עולם; והיכן אמרה תורה מרשעים יצא רשע? והאלהים אנה לידו. במה הכתוב מדבר? בשני בני אדם, אחד הרג שוגג ואחד הרג מזיד, ולא היו עדים בדבר שיעידו, זה לא נהרג, וזה לא גלה, והקב"ה מזמנן לפנדק אחד, זה שהרג במזיד יושב תחת הסלם וזה שהרג שוגג עולה בסלם ונופל על זה שהרג במזיד והורגו, ועדים מעידים עליו ומחיבים אותו לגלות, נמצא זה שהרג בשוגג גולה, וזה שהרג במזיד נהרג (מכות י'): ושמתי לך מקום. אף במדבר, שינוס שמה, ואי זה מקום קולטו? זה מחנה לויה (מכילתא):}%
% ---------- פרק 21 | פסוק 14 ----------
 \textbf{יד} וְכִֽי־יָזִ֥ד אִ֛ישׁ עַל־רֵעֵ֖הוּ לְהׇרְג֣וֹ בְעׇרְמָ֑ה מֵעִ֣ם מִזְבְּחִ֔י תִּקָּחֶ֖נּוּ לָמֽוּת׃ \{ס\} \Rashi{\textbf{וכי יזד.}למה נאמר? לפי שנאמר מכה איש וגו' שומע אני אפלו גוי, ורופא שהמית ושליח בית דין שהמית במלקות ארבעים, והאב המכה את בנו והרב הרודה את תלמידו, והשוגג, ת"ל וכי יזד – ולא שוגג, על רעהו – ולא על גוי, להרגו בערמה – ולא שליח בית דין והרופא ורודה בנו ותלמידו, שאף על פי שהם מזידין אין מערימין (מכילתא): מעם מזבחי. אם היה כהן ורוצה לעבד עבודה תקחנו למות (יומא פ"ה):}%
% ---------- פרק 21 | פסוק 15 ----------
 \textbf{טו} וּמַכֵּ֥ה אָבִ֛יו וְאִמּ֖וֹ מ֥וֹת יוּמָֽת׃ \{ס\} \Rashi{\textbf{ומכה אביו ואמו.}לפי שלמדנו על החובל בחברו שהוא בתשלומין ולא במיתה, הצרך לומר על החובל באביו שהוא במיתה, ואינו חיב אלא בהכאה שיש בה חבורה (סנהדרין פ"ה): אביו ואמו. או זה או זה: מות יומת. בחנק (מכילתא):}%
% ---------- פרק 21 | פסוק 16 ----------
 \textbf{טז} וְגֹנֵ֨ב אִ֧ישׁ וּמְכָר֛וֹ וְנִמְצָ֥א בְיָד֖וֹ מ֥וֹת יוּמָֽת׃ \{ס\} \Rashi{\textbf{וגנב איש ומכרו.}למה נאמר? לפי שנאמר כי ימצא איש גונב נפש מאחיו (דברים כ"ד), אין לי אלא איש שגנב נפש, אשה או טמטום או אנדרוגינוס שגנבו מנין? ת"ל וגנב איש ומכרו. ולפי שנא' כאן וגנב איש, אין לי אלא גונב איש, גונב אשה מנין? ת"ל גנב נפש, לכך הצרכו שניהם, מה שחסר זה גלה זה (סנהדרין שם): ונמצא בידו. שראוהו עדים שגנבו ומכרו, ונמצא כבר בידו קדם מכירה (שם): מות יומת. בחנק; כל מיתה האמורה בתורה סתם חנק היא (והפסיק הענין וכתב וגנב איש בין מכה אביו ואמו למקלל אביו, ונראה לי הינו פלוגתא, דמר סבר מקשינן הכאה לקללה ומר סבר לא מקשינן (שם)):}%
% ---------- פרק 21 | פסוק 17 ----------
 \textbf{יז} וּמְקַלֵּ֥ל אָבִ֛יו וְאִמּ֖וֹ מ֥וֹת יוּמָֽת׃ \{ס\} \Rashi{\textbf{ומקלל אביו ואמו.}למה נאמר? לפי שהוא אומר איש איש אשר יקלל את אביו (ויקרא כ'), אין לי אלא איש שקלל את אביו, אשה שקללה את אביה מנין? ת"ל ומקלל אביו ואמו סתם, בין איש ובין אשה, אם כן למה נאמר איש אשר יקלל? להוציא את הקטן (מכילתא): מות יומת. בסקילה; וכל מקום שנאמר דמיו בו בסקילה, ובנין אב לכלם באבן ירגמו אתם דמיהם בם (ויקרא כ'), ובמקלל אביו ואמו נאמר דמיו בו:}%
% ---------- פרק 21 | פסוק 18 ----------
 \textbf{יח} וְכִֽי־יְרִיבֻ֣ן אֲנָשִׁ֔ים וְהִכָּה־אִישׁ֙ אֶת־רֵעֵ֔הוּ בְּאֶ֖בֶן א֣וֹ בְאֶגְרֹ֑ף וְלֹ֥א יָמ֖וּת וְנָפַ֥ל לְמִשְׁכָּֽב׃ \Rashi{\textbf{וכי יריבן אנשים.}למה נאמר? לפי שנאמר עין תחת עין, לא למדנו אלא דמי אבריו אבל שבת ורפוי לא למדנו, לכך נאמרה פרשה זו: ונפל למשכב. כתרגומו ויפל לבוטלן – לחלי שמבטלו ממלאכתו:}%
% ---------- פרק 21 | פסוק 19 ----------
 \textbf{יט} אִם־יָק֞וּם וְהִתְהַלֵּ֥ךְ בַּח֛וּץ עַל־מִשְׁעַנְתּ֖וֹ וְנִקָּ֣ה הַמַּכֶּ֑ה רַ֥ק שִׁבְתּ֛וֹ יִתֵּ֖ן וְרַפֹּ֥א יְרַפֵּֽא׃ \{ס\} \Rashi{\textbf{על משענתו.}על בריו וכחו (מכילתא): ונקה המכה. וכי תעלה על דעתך שיהרג זה שלא הרג? אלא למדך כאן שחובשים אותו עד שנראה אם יתרפא זה; וכן משמעו: כשקם זה והולך על משענתו אז נקה המכה, אבל עד שלא יקום זה לא נקה המכה (כתובות ל"ג): רק שבתו. בטול מלאכתו מחמת החלי; אם קטע ידו או רגלו, רואין בטול מלאכתו מחמת החלי כאלו שומר קשואין, שהרי אף לאחר החלי אינו ראוי למלאכת יד ורגל, והוא כבר נתן לו מחמת נזקו דמי ידו ורגלו, שנ' יד תחת יד רגל תחת רגל: ורפא ירפא. כתרגומו – ישלם שכר הרופא:}%
% ---------- פרק 21 | פסוק 20 ----------
 \textbf{כ} וְכִֽי־יַכֶּה֩ אִ֨ישׁ אֶת־עַבְדּ֜וֹ א֤וֹ אֶת־אֲמָתוֹ֙ בַּשֵּׁ֔בֶט וּמֵ֖ת תַּ֣חַת יָד֑וֹ נָקֹ֖ם יִנָּקֵֽם׃ \Rashi{\textbf{וכי יכה איש את עבדו או את אמתו.}בעבד כנעני הכתוב מדבר, או אינו אלא בעברי? ת"ל כי כספו הוא, מה כספו קנוי לו עולמית, אף עבד הקנוי לו עולמית; והרי היה בכלל מכה איש ומת? אלא בא הכתוב והוציאו מן הכלל, להיות נדון בדין יום או יומים, שאם לא מת תחת ידו ושהה מעת לעת פטור (מכילתא): בשבט. כשיש בו כדי להמית הכתוב מדבר. או אפלו אין בו כדי להמית? ת"ל בישראל ואם באבן יד אשר ימות בה הכהו (במדבר ל"ה), והלא דברים קל וחמר: מה ישראל חמור אין חיב עליו אלא אם כן הכהו בדבר שיש בו כדי להמית ועל אבר שהוא כדי למות בהכאה זו, עבד הקל לא כל שכן (מכילתא): נקם ינקם. מיתת סיף, וכן הוא אומר חרב נקמת נקם ברית (ויקרא כ"ו):}%
% ---------- פרק 21 | פסוק 21 ----------
 \textbf{כא} אַ֥ךְ אִם־י֛וֹם א֥וֹ יוֹמַ֖יִם יַעֲמֹ֑ד לֹ֣א יֻקַּ֔ם כִּ֥י כַסְפּ֖וֹ הֽוּא׃ \{ס\} \Rashi{\textbf{אך אם יום או יומים יעמד לא יקם.}אם על יום אחד הוא פטור על יומים לא כל שכן? אלא יום שהוא כיומים, ואיזה? זה מעת לעת (מכילתא): לא יקם כי כספו הוא. הא אחר שהכהו, אע"פ ששהה מעת לעת קדם שמת, חיב:}%
% ---------- פרק 21 | פסוק 22 ----------
 \textbf{כב} וְכִֽי־יִנָּצ֣וּ אֲנָשִׁ֗ים וְנָ֨גְפ֜וּ אִשָּׁ֤ה הָרָה֙ וְיָצְא֣וּ יְלָדֶ֔יהָ וְלֹ֥א יִהְיֶ֖ה אָס֑וֹן עָנ֣וֹשׁ יֵעָנֵ֗שׁ כַּֽאֲשֶׁ֨ר יָשִׁ֤ית עָלָיו֙ בַּ֣עַל הָֽאִשָּׁ֔ה וְנָתַ֖ן בִּפְלִלִֽים׃ \Rashi{\textbf{וכי ינצו אנשים.}זה עם זה, ונתכון להכות את חברו והכה את האשה: ונגפו. אין נגיפה אלא לשון דחיפה והכאה, כמו פן תגף באבן רגלך (תהילים צ"א), בטרם יתנגפו רגליכם (ירמיהו י"ג) ולאבן נגף (ישעיהו ח'): ולא יהיה אסון. באשה: ענוש יענש. לשלם דמי ולדות לבעל; שמין אותה כמה היתה ראויה למכר בשוק להעלות בדמיה בשביל הריונה: ענש יענש. יגבו ממון ממנו, כמו וענשו אתו מאה כסף (דברים כ"ב): כאשר ישית עליו וגו'. כשיתבענו הבעל בבית דין להשית עליו ענש על כך: ונתן. המכה דמי ולדות: בפללים – על פי הדינים (מכילתא):}%
% ---------- פרק 21 | פסוק 23 ----------
 \textbf{כג} וְאִם־אָס֖וֹן יִהְיֶ֑ה וְנָתַתָּ֥ה נֶ֖פֶשׁ תַּ֥חַת נָֽפֶשׁ׃ \Rashi{\textbf{ואם אסון יהיה.}באשה: ונתתה נפש תחת נפש. רבותינו חולקים בדבר, יש אומרים נפש ממש, ויש אומרים ממון אבל לא נפש ממש, שהמתכון להרג את זה והרג את זה פטור ממיתה, ומשלם ליורשיו דמיו כמו שהיה נמכר בשוק (שם):}%
% ---------- פרק 21 | פסוק 24 ----------
 \textbf{כד} עַ֚יִן תַּ֣חַת עַ֔יִן שֵׁ֖ן תַּ֣חַת שֵׁ֑ן יָ֚ד תַּ֣חַת יָ֔ד רֶ֖גֶל תַּ֥חַת רָֽגֶל׃ \Rashi{\textbf{עין תחת עין.}סמא עין חברו נותן לו דמי עינו כמה שפחתו דמיו למכר בשוק, וכן כלם; ולא נטילת אבר ממש, כמו שדרשו רבותינו בפרק החובל (בבא קמא דף פ"ג):}%
% ---------- פרק 21 | פסוק 25 ----------
 \textbf{כה} כְּוִיָּה֙ תַּ֣חַת כְּוִיָּ֔ה פֶּ֖צַע תַּ֣חַת פָּ֑צַע חַבּוּרָ֕ה תַּ֖חַת חַבּוּרָֽה׃ \{ס\} \Rashi{\textbf{כויה תחת כויה.}מכות אש; ועד עכשו דבר בחבלה שיש בה פחת דמים, ועכשיו בשאין בה פחת דמים אלא צער, כגון כואו בשפוד על צפרניו, אומדים כמה אדם כיוצא בזה רוצה לטל להיות מצטער כך (בבא קמא פ"ג): פצע. היא מכה המוציאה דם, שפצע את בשרו, נבר"ורא בלעז; הכל לפי מה שהוא – אם יש בו פחת דמים, נותן נזק, ואם נפל למשכב, נותן שבת ורפוי ובשת וצער; ומקרא זה יתר הוא, ובהחובל דרשוהו רבותינו לחיב על הצער אפלו במקום נזק – שאף על פי שנותן לו דמי ידו, אין פוטרים אותו מן הצער, לומר, הואיל וקנה ידו יש עליו לחתכה בכל מה שירצה, אלא אומרים, יש לו לחתכה בסם, שאינו מצטער כל כך, וזה חתכה בברזל וצערו: חבורה. היא מכה שהדם נצרר בה ואינו יוצא, אלא שמאדים הבשר כנגדו, ולשון חבורה טק"א בלעז, כמו ונמר חברברתיו (ירמיהו י"ג), ותרגומו משקופי, לשון חבטה, בטדו"רא בלעז, וכן שדופות קדים (בראשית מ"א), שקיפן קדום – חבוטות ברוח, וכן על המשקוף, על שם שהדלת נוקש עליו:}%
% ---------- פרק 21 | פסוק 26 ----------
 \textbf{כו} וְכִֽי־יַכֶּ֨ה אִ֜ישׁ אֶת־עֵ֥ין עַבְדּ֛וֹ אֽוֹ־אֶת־עֵ֥ין אֲמָת֖וֹ וְשִֽׁחֲתָ֑הּ לַֽחׇפְשִׁ֥י יְשַׁלְּחֶ֖נּוּ תַּ֥חַת עֵינֽוֹ׃ \Rashi{\textbf{את עין עבדו.}כנעני, אבל עברי אינו יוצא בשן ועין, כמו שאמרנו לא תצא כצאת העבדים: תחת עינו. וכן בכ"ד ראשי אברים: אצבעות הידים והרגלים, ושתי אזנים, והחטם, וראש הגויה שהוא גיד האמה. ולמה נאמר שן ועין? שאם נאמר עין ולא נאמר שן, הייתי אומר מה עין שנברא עמו אף כל שנברא עמו, והרי שן לא נברא עמו; ואם נאמר שן ולא נאמר עין, הייתי אומר אפלו שן תינוק שיש לה חליפין, לכך נאמר עין (מכילתא):}%
% ---------- פרק 21 | פסוק 27 ----------
 \textbf{כז} וְאִם־שֵׁ֥ן עַבְדּ֛וֹ אֽוֹ־שֵׁ֥ן אֲמָת֖וֹ יַפִּ֑יל לַֽחׇפְשִׁ֥י יְשַׁלְּחֶ֖נּוּ תַּ֥חַת שִׁנּֽוֹ׃ \{פ\} %
% ---------- פרק 21 | פסוק 28 ----------
 \textbf{כח} וְכִֽי־יִגַּ֨ח שׁ֥וֹר אֶת־אִ֛ישׁ א֥וֹ אֶת־אִשָּׁ֖ה וָמֵ֑ת סָק֨וֹל יִסָּקֵ֜ל הַשּׁ֗וֹר וְלֹ֤א יֵאָכֵל֙ אֶת־בְּשָׂר֔וֹ וּבַ֥עַל הַשּׁ֖וֹר נָקִֽי׃ \Rashi{\textbf{וכי יגח שור.}אחד שור ואחד כל בהמה וחיה ועוף, אלא שדבר הכתוב בהוה (בבא קמא נ"ד): ולא יאכל את בשרו. ממשמע שנאמר סקול יסקל השור איני יודע שהוא נבלה ונבלה אסורה באכילה? אלא מה תלמוד לומר ולא יאכל את בשרו? שאפלו שחטו לאחר שנגמר דינו אסור באכילה, בהנאה מנין? ת"ל ובעל השור נקי, כאדם האומר לחברו יצא פלוני נקי מנכסיו ואין לו בהם הנאה של כלום, זהו מדרשו (בבא קמא מלכים א). ופשוטו כמשמעו, לפי שנאמר במועד וגם בעליו יומת, הצרך לומר בתם ובעל השור נקי:}%
% ---------- פרק 21 | פסוק 29 ----------
 \textbf{כט} וְאִ֡ם שׁוֹר֩ נַגָּ֨ח ה֜וּא מִתְּמֹ֣ל שִׁלְשֹׁ֗ם וְהוּעַ֤ד בִּבְעָלָיו֙ וְלֹ֣א יִשְׁמְרֶ֔נּוּ וְהֵמִ֥ית אִ֖ישׁ א֣וֹ אִשָּׁ֑ה הַשּׁוֹר֙ יִסָּקֵ֔ל וְגַם־בְּעָלָ֖יו יוּמָֽת׃ \Rashi{\textbf{מתמל שלשם.}הרי שלש נגיחות (מכילתא): והועד בבעליו. לשון התראה בעדים, כמו העד העד בנו האיש (בראשית מ"ג): והמית איש וגו'. לפי שנאמר כי יגח אין לי אלא שהמיתו בנגיחה, המיתו בנשיכה, דחיפה, בעיטה, מנין? תלמוד לומר והמית: וגם בעליו יומת. בידי שמים. יכול בידי אדם? תלמוד לומר מות יומת המכה רוצח הוא, על רציחתו אתה הורגו ואי אתה הורגו על רציחת שורו (סנהדרין ט"ו):}%
% ---------- פרק 21 | פסוק 30 ----------
 \textbf{ל} אִם־כֹּ֖פֶר יוּשַׁ֣ת עָלָ֑יו וְנָתַן֙ פִּדְיֹ֣ן נַפְשׁ֔וֹ כְּכֹ֥ל אֲשֶׁר־יוּשַׁ֖ת עָלָֽיו׃ \Rashi{\textbf{אם כפר יושת עליו.}"אם" זה אינו תלוי, והרי הוא כמו "אם כסף תלוה" – לשון אשר, זה משפטו שישיתו עליו בית דין כפר: ונתן פדין נפשו. דמי נזק דברי ר' ישמעאל, ר' עקיבא אומר דמי מזיק (מכילתא):}%
% ---------- פרק 21 | פסוק 31 ----------
 \textbf{לא} אוֹ־בֵ֥ן יִגָּ֖ח אוֹ־בַ֣ת יִגָּ֑ח כַּמִּשְׁפָּ֥ט הַזֶּ֖ה יֵעָ֥שֶׂה לּֽוֹ׃ \Rashi{\textbf{או בן יגח.}בן שהוא קטן: או בת. שהיא קטנה; לפי שנאמר והמית איש או אשה, יכול אינו חיב אלא על הגדולים, תלמוד לומר או בן יגח וגו', לחיב על הקטנים כגדולים (שם):}%
% ---------- פרק 21 | פסוק 32 ----------
 \textbf{לב} אִם־עֶ֛בֶד יִגַּ֥ח הַשּׁ֖וֹר א֣וֹ אָמָ֑ה כֶּ֣סֶף ׀ שְׁלֹשִׁ֣ים שְׁקָלִ֗ים יִתֵּן֙ לַֽאדֹנָ֔יו וְהַשּׁ֖וֹר יִסָּקֵֽל׃ \{ס\} \Rashi{\textbf{אם עבד, או אמה.}כנעניים: שלשים שקלים יתן. גזרת הכתוב הוא, בין שהוא שוה אלף זוז בין שאינו שוה אלא דינר. והשקל משקלו ארבעה זהובים שהם חצי אנקיא למשקל הישר של קולוני"א:}%
% ---------- פרק 21 | פסוק 33 ----------
 \textbf{לג} וְכִֽי־יִפְתַּ֨ח אִ֜ישׁ בּ֗וֹר א֠וֹ כִּֽי־יִכְרֶ֥ה אִ֛ישׁ בֹּ֖ר וְלֹ֣א יְכַסֶּ֑נּוּ וְנָֽפַל־שָׁ֥מָּה שּׁ֖וֹר א֥וֹ חֲמֽוֹר׃ \Rashi{\textbf{וכי יפתח איש בור.}שהיה מכסה וגלהו: או כי יכרה. למה נאמר? אם על הפתיחה חיב על הכריה לא כל שכן? אלא להביא כורה אחר כורה שהוא חיב (בבא קמא נ"א): ולא יכסנו. הא אם כסהו פטור; ובחופר ברשות הרבים דבר הכתוב: שור או חמור. הוא הדין לכל בהמה וחיה, שבכל מקום שנאמר שור או חמור אנו למדין אותו שור שור משבת, שנאמר למען ינוח שורך וחמרך (שמות כ"ג), מה להלן כל בהמה וחיה כשור – שהרי נאמר במקום אחר וכל בהמתך (דברים ה') – אף כאן כל בהמה וחיה כשור, ולא נאמר שור וחמור אלא שור ולא אדם, חמור ולא כלים (בבא קמא נ"ג):}%
% ---------- פרק 21 | פסוק 34 ----------
 \textbf{לד} בַּ֤עַל הַבּוֹר֙ יְשַׁלֵּ֔ם כֶּ֖סֶף יָשִׁ֣יב לִבְעָלָ֑יו וְהַמֵּ֖ת יִֽהְיֶה־לּֽוֹ׃ \{ס\} \Rashi{\textbf{בעל הבור.}בעל התקלה; אע"פ שאין הבור שלו – שעשאו ברשות הרבים – עשאו הכתוב בעליו להתחיב עליו בנזקיו: כסף ישיב לבעליו. ישיב לרבות שוה כסף ואפלו סבין (שם ז'): והמת יהיה לו. לנזק; שמין את הנבלה ונוטלה בדמים ומשלם לו המזיק עליה תשלומי נזקו (בבא קמא י'):}%
% ---------- פרק 21 | פסוק 35 ----------
 \textbf{לה} וְכִֽי־יִגֹּ֧ף שֽׁוֹר־אִ֛ישׁ אֶת־שׁ֥וֹר רֵעֵ֖הוּ וָמֵ֑ת וּמָ֨כְר֜וּ אֶת־הַשּׁ֤וֹר הַחַי֙ וְחָצ֣וּ אֶת־כַּסְפּ֔וֹ וְגַ֥ם אֶת־הַמֵּ֖ת יֶֽחֱצֽוּן׃ \Rashi{\textbf{וכי יגף.}ידחף; בין בקרניו, בין בגופו, בין ברגלו, בין שנשכו בשניו, כלן בכלל נגיפה הם, שאין נגיפה אלא לשון מכה: שור איש. שור של איש: ומכרו את השור וגו'. בשוים הכתוב מדבר – שור שוה מאתים שהמית שור שוה מאתים – בין שהנבלה שוה הרבה, בין שהיא שוה מעט, כשנוטל זה חצי החי וחצי המת וזה חצי החי וחצי המת, נמצא כל אחד מפסיד חצי נזק שהזיקה המיתה; למדנו שהתם משלם חצי נזק, שמן השוין אתה למד לשאינן שוין, כי דין התם לשלם חצי נזק – לא פחות ולא יותר. או יכול אף בשאינן שוין בדמיהן כשהן חיים אמר הכתוב יחצו את שניהם? אם אמרת כן, פעמים שהמזיק משתכר הרבה, כשהנבלה שוה למכר לנכרים הרבה יותר מדמי שור המזיק, ואי אפשר שיאמר הכתוב שיהא המזיק נשכר; או פעמים שהנזק נוטל הרבה יותר מדמי נזק שלם – שחצי דמי שור המזיק שוין יותר מכל דמי שור הנזק – ואם אמרת כן, הרי תם חמור ממועד. על כרחך לא דבר הכתוב אלא בשוין, ולמדך שהתם משלם חצי נזק, ומן השוין תלמד לשאינן שוין, שהמשתלם חצי נזקו שמין לו את הנבלה, ומה שפחתו דמיו בשביל המיתה, נוטל חצי הפחת והולך. ולמה אמר הכתוב בלשון הזה ולא אמר ישלם חציו? ללמד, שאין התם משלם אלא מגופו, ואם נגח ומת, אין נזק נוטל אלא הנבלה, ואם אינה מגעת לחצי נזקו יפסיד; או שור שוה מנה, שנגח שור שוה חמש מאות זוז, אינו נוטל אלא את השור, שלא נתחיב התם לחיב את בעליו לשלם מן העליה (בבא קמא ל"ג):}%
% ---------- פרק 21 | פסוק 36 ----------
 \textbf{לו} א֣וֹ נוֹדַ֗ע כִּ֠י שׁ֣וֹר נַגָּ֥ח הוּא֙ מִתְּמ֣וֹל שִׁלְשֹׁ֔ם וְלֹ֥א יִשְׁמְרֶ֖נּוּ בְּעָלָ֑יו שַׁלֵּ֨ם יְשַׁלֵּ֥ם שׁוֹר֙ תַּ֣חַת הַשּׁ֔וֹר וְהַמֵּ֖ת יִֽהְיֶה־לּֽוֹ׃ \{ס\} \Rashi{\textbf{או נודע.}או לא היה תם, אלא נודע כי שור נגח הוא, היום ומתמול שלשם הרי שלש נגיחות: שלם ישלם שור. נזק שלם: והמת יהיה לו. לנזק; ועליו ישלים המזיק עד שישתלם נזק כל נזקו:}%
% ---------- פרק 21 | פסוק 37 ----------
 \textbf{לז} כִּ֤י יִגְנֹֽב־אִישׁ֙ שׁ֣וֹר אוֹ־שֶׂ֔ה וּטְבָח֖וֹ א֣וֹ מְכָר֑וֹ חֲמִשָּׁ֣ה בָקָ֗ר יְשַׁלֵּם֙ תַּ֣חַת הַשּׁ֔וֹר וְאַרְבַּע־צֹ֖אן תַּ֥חַת הַשֶּֽׂה׃ \Rashi{\textbf{חמשה בקר וגו'.}אמר ר' יוחנן בן זכאי חס המקום על כבודן של בריות – שור שהולך ברגליו ולא נתבזה בו הגנב לנשאו על כתפו, משלם חמשה, שה שנושאו על כתפו, משלם ארבעה הואיל ונתבזה בו. אמר רבי מאיר בא וראה כמה גדול כחה של מלאכה – שור שבטלו ממלאכתו חמשה, שה שלא בטלו ממלאכתו ארבעה (מכילתא, בבא קמא ע"ט): תחת השור, תחת השה. שנאן הכתוב, לומר שאין מדת תשלומי ארבעה וחמשה נוהגת אלא בשור ושה בלבד (בבא קמא ס"ז):}%
% ---------- פרק 22 | פסוק 1 ----------
 \textbf{א} אִם־בַּמַּחְתֶּ֛רֶת יִמָּצֵ֥א הַגַּנָּ֖ב וְהֻכָּ֣ה וָמֵ֑ת אֵ֥ין ל֖וֹ דָּמִֽים׃ \Rashi{\textbf{אם במחתרת.}כשהיה חותר את הבית: אין לו דמים. אין זו רציחה, הרי הוא כמת מעקרו; כאן למדתך תורה "אם בא להרגך, השכם להרגו" וזה להרגך בא, שהרי יודע הוא שאין אדם מעמיד עצמו ורואה שנוטלין ממונו בפניו ושותק, לפיכך על מנת כן בא שאם יעמד בעל הממון כנגדו יהרגנו (סנהדרין ע"ב):}%
% ---------- פרק 22 | פסוק 2 ----------
 \textbf{ב} אִם־זָרְחָ֥ה הַשֶּׁ֛מֶשׁ עָלָ֖יו דָּמִ֣ים ל֑וֹ שַׁלֵּ֣ם יְשַׁלֵּ֔ם אִם־אֵ֣ין ל֔וֹ וְנִמְכַּ֖ר בִּגְנֵבָתֽוֹ׃ \Rashi{\textbf{אם זרחה השמש עליו.}אין זה אלא כמין משל: אם ברור לך הדבר שיש לו שלום עמך, כשמש הזה שהוא שלום בעולם, כך פשוט לך שאינו בא להרג אפלו יעמד בעל הממון כנגדו, כגון אב החותר לגנב ממון הבן, בידוע שרחמי האב על הבן ואינו בא על עסקי נפשות (שם): דמים לו. כחי הוא חשוב, ורציחה היא אם יהרגנו בעל הבית: שלם ישלם. הגנב ממון שגנב ואינו חיב מיתה; ואנקלוס שתרגם אם עינא דסהדיא נפלת עלוהי, לקח לו שטה אחרת, לומר, שאם מצאוהו עדים קדם שבא בעל הבית, וכשבא בעל הבית נגדו התרו בו שלא יהרגהו, דמים לו – חיב עליו אם הרגו, שמאחר שיש רואים לו, אין הגנב הזה בא על עסקי נפשות, ולא יהרג את בעל הממון:}%
% ---------- פרק 22 | פסוק 3 ----------
 \textbf{ג} אִֽם־הִמָּצֵא֩ תִמָּצֵ֨א בְיָד֜וֹ הַגְּנֵבָ֗ה מִשּׁ֧וֹר עַד־חֲמ֛וֹר עַד־שֶׂ֖ה חַיִּ֑ים שְׁנַ֖יִם יְשַׁלֵּֽם׃ \{ס\} \Rashi{\textbf{המצא תמצא.}ברשותו, שלא טבח ולא מכר: משור עד חמור. כל דבר בכלל תשלומי כפל, בין שיש בו רוח חיים, בין שאין בו רוח חיים, שהרי נאמר במקרא אחר "על שה על שלמה על כל אבדה … ישלם שנים לרעהו" (בבא קמא ס"ב): חיים שנים ישלם. ולא ישלם לו מתים, אלא חיים או דמי חיים (מכילתא):}%
% ---------- פרק 22 | פסוק 4 ----------
 \textbf{ד} כִּ֤י יַבְעֶר־אִישׁ֙ שָׂדֶ֣ה אוֹ־כֶ֔רֶם וְשִׁלַּח֙ אֶת־בְּעִירֹ֔ה וּבִעֵ֖ר בִּשְׂדֵ֣ה אַחֵ֑ר מֵיטַ֥ב שָׂדֵ֛הוּ וּמֵיטַ֥ב כַּרְמ֖וֹ יְשַׁלֵּֽם׃ \{ס\} \Rashi{\textbf{כי יבער, את בעירה ובער.}כלם לשון בהמה, כמו אנחנו ובעירנו (במדבר כ'): כי יבער. יוליך בהמותיו בשדה או בכרם של חברו, ויזיק אותו באחת משתי אלו, או בשלוח בעירו או בבעור, ופרשו רבותינו "ושלח" הוא נזקי מדרך כף רגל, "ובער" הוא נזקי השן, האוכלת ומבערת (בבא קמא ב'): בשדה אחר. בשדה של איש אחר: מיטב שדהו … ישלם. שמין את הנזק, ואם בא לשלם לו קרקע דמי נזקו, ישלם לו ממיטב שדותיו – אם היה נזקו סלע, יתן לו שוה סלע מעדית שיש לו; למדך הכתוב שהנזקין שמין להם בעדית (מכילתא, בבא קמא ו'):}%
% ---------- פרק 22 | פסוק 5 ----------
 \textbf{ה} כִּֽי־תֵצֵ֨א אֵ֜שׁ וּמָצְאָ֤ה קֹצִים֙ וְנֶאֱכַ֣ל גָּדִ֔ישׁ א֥וֹ הַקָּמָ֖ה א֣וֹ הַשָּׂדֶ֑ה שַׁלֵּ֣ם יְשַׁלֵּ֔ם הַמַּבְעִ֖ר אֶת־הַבְּעֵרָֽה׃ \{ס\} \Rashi{\textbf{כי תצא אש.}אפלו מעצמה (בבא קמא כ"ב): ומצאה קצים. קרדו"נש בלעז: ונאכל גדיש. שלחכה בקוצים עד שהגיעה לגדיש או לקמה המחברת בקרקע: או השדה. שלחכה את נירו וצריך לניר אותה פעם שניה (בבא קמא ס'): שלם ישלם המבער. אע"פ שהדליק בתוך שלו, והיא יצאה מעצמה על ידי קוצים שמצאה, חיב לשלם, לפי שלא שמר את גחלתו שלא תצא ותזיק:}%
% ---------- פרק 22 | פסוק 6 ----------
 \textbf{ו} כִּֽי־יִתֵּן֩ אִ֨ישׁ אֶל־רֵעֵ֜הוּ כֶּ֤סֶף אֽוֹ־כֵלִים֙ לִשְׁמֹ֔ר וְגֻנַּ֖ב מִבֵּ֣ית הָאִ֑ישׁ אִם־יִמָּצֵ֥א הַגַּנָּ֖ב יְשַׁלֵּ֥ם שְׁנָֽיִם׃ \Rashi{\textbf{וגנב מבית האיש.}לפי דבריו (שם ס"ג): אם ימצא הגנב ישלם. הגנב שנים לבעלים:}%
% ---------- פרק 22 | פסוק 7 ----------
 \textbf{ז} אִם־לֹ֤א יִמָּצֵא֙ הַגַּנָּ֔ב וְנִקְרַ֥ב בַּֽעַל־הַבַּ֖יִת אֶל־הָֽאֱלֹהִ֑ים אִם־לֹ֥א שָׁלַ֛ח יָד֖וֹ בִּמְלֶ֥אכֶת רֵעֵֽהוּ׃ \Rashi{\textbf{אם לא ימצא הגנב.}ובא השומר הזה, שהוא בעל הבית: ונקרב. אל הדינין, לדון עם זה ולשבע לו שלא שלח ידו בשלו (שם):}%
% ---------- פרק 22 | פסוק 8 ----------
 \textbf{ח} עַֽל־כׇּל־דְּבַר־פֶּ֡שַׁע עַל־שׁ֡וֹר עַל־חֲ֠מ֠וֹר עַל־שֶׂ֨ה עַל־שַׂלְמָ֜ה עַל־כׇּל־אֲבֵדָ֗ה אֲשֶׁ֤ר יֹאמַר֙ כִּי־ה֣וּא זֶ֔ה עַ֚ד הָֽאֱלֹהִ֔ים יָבֹ֖א דְּבַר־שְׁנֵיהֶ֑ם אֲשֶׁ֤ר יַרְשִׁיעֻן֙ אֱלֹהִ֔ים יְשַׁלֵּ֥ם שְׁנַ֖יִם לְרֵעֵֽהוּ׃ \{ס\} \Rashi{\textbf{על כל דבר פשע.}שימצא שקרן בשבועתו – שיעידו עדים שהוא עצמו גנבו – וירשיעוהו אלהים על פי העדים. ישלם שנים לרעהו. למדך הכתוב, שהטוען בפקדון לומר נגנב הימנו, ונמצא שהוא עצמו גנבו, משלם תשלומי כפל; ואימתי? בזמן שנשבע ואח"כ באו עדים, שכך דרשו רבותינו: ונקרב בעל הבית אל האלהים, קריבה זו שבועה היא, אתה אומר לשבועה או אינו אלא לדין, שכיון שבא לדין וכפר לומר נגנבה, מיד יתחיב כפל אם באו עדים שהוא בידו? נאמר כאן שליחות יד ונאמר למטה שליחות יד – שבועת ה' תהיה בין שניהם אם לא שלח ידו – מה להלן שבועה, אף כאן שבועה (שם): אשר יאמר כי הוא זה. לפי פשוטו אשר יאמר העד כי הוא זה שנשבעת עליו, הרי הוא אצלך, עד הדינין יבא דבר שניהם, ויחקרו את העדים, אם כשרים הם וירשיעוהו לשומר זה, ישלם שנים, ואם ירשיעו את העדים, שנמצאו זוממין, ישלמו הם שנים לשומר. ורבותינו ז"ל דרשו כי הוא זה ללמד, שאין מחיבין אותו שבועה אלא אם כן הודה במקצת, לומר כך וכך אני חיב לך, והמותר נגנב ממני (שם ק"ז):}%
% ---------- פרק 22 | פסוק 9 ----------
 \textbf{ט} כִּֽי־יִתֵּן֩ אִ֨ישׁ אֶל־רֵעֵ֜הוּ חֲמ֨וֹר אוֹ־שׁ֥וֹר אוֹ־שֶׂ֛ה וְכׇל־בְּהֵמָ֖ה לִשְׁמֹ֑ר וּמֵ֛ת אוֹ־נִשְׁבַּ֥ר אוֹ־נִשְׁבָּ֖ה אֵ֥ין רֹאֶֽה׃ \Rashi{\textbf{כי יתן איש אל רעהו חמור או שור.}פרשה ראשונה נאמרה בשומר חנם, לפיכך פטר בו את הגנבה, כמו שנאמר (פסוק ו) וגנב מבית האיש, אם לא ימצא הגנב ונקרב בעל הבית לשבועה, למדת שפוטר עצמו בשבועה זו. ופרשה זו אמורה בשומר שכר, לפיכך אינו פטור אם נגנבה, כמו שכתוב (פסוק יא) אם גנב יגנב מעמו ישלם, אבל על האנס, כגון מת מעצמו, או נשבר או נשבה בחזקה על ידי לסטים, ואין רואה שיעיד בדבר:}%
% ---------- פרק 22 | פסוק 10 ----------
 \textbf{י} שְׁבֻעַ֣ת יְהֹוָ֗ה תִּהְיֶה֙ בֵּ֣ין שְׁנֵיהֶ֔ם אִם־לֹ֥א שָׁלַ֛ח יָד֖וֹ בִּמְלֶ֣אכֶת רֵעֵ֑הוּ וְלָקַ֥ח בְּעָלָ֖יו וְלֹ֥א יְשַׁלֵּֽם׃ \Rashi{\textbf{שבעת ה' תהיה.}ישבע שכן הוא כדבריו, והוא לא שלח בה יד להשתמש בה לעצמו, שאם שלח בה יד ואח"כ נאנסה, חיב באנסים (בבא מציעא צ"ד): ולקח בעליו. השבועה: ולא ישלם. לו השומר כלום (בבא קמא ק"ו):}%
% ---------- פרק 22 | פסוק 11 ----------
 \textbf{יא} וְאִם־גָּנֹ֥ב יִגָּנֵ֖ב מֵעִמּ֑וֹ יְשַׁלֵּ֖ם לִבְעָלָֽיו׃ %
% ---------- פרק 22 | פסוק 12 ----------
 \textbf{יב} אִם־טָרֹ֥ף יִטָּרֵ֖ף יְבִאֵ֣הוּ עֵ֑ד הַטְּרֵפָ֖ה לֹ֥א יְשַׁלֵּֽם׃ \{פ\} \Rashi{\textbf{אם טרף יטרף.}ע"י חיה רעה: יביאהו עד. יביא עדים שנטרפה באנס ופטור: הטרפה לא ישלם. אינו אומר טרפה לא ישלם אלא הטרפה, יש טרפה שהוא משלם ויש טרפה שאינו משלם; טרפת חתול ושועל ונמיה משלם, טרפת זאב, ארי, ודב ונחש אינו משלם; ומי לחשך לדון כן? שהרי כתוב ומת או נשבר או נשבה – מה מיתה שאין יכול להציל, אף שבר ושביה שאין יכול להציל (מכילתא):}%
% ---------- פרק 22 | פסוק 13 ----------
 \textbf{יג} וְכִֽי־יִשְׁאַ֥ל אִ֛ישׁ מֵעִ֥ם רֵעֵ֖הוּ וְנִשְׁבַּ֣ר אוֹ־מֵ֑ת בְּעָלָ֥יו אֵין־עִמּ֖וֹ שַׁלֵּ֥ם יְשַׁלֵּֽם׃ \Rashi{\textbf{וכי ישאל.}בא ללמד על השואל שהוא חיב באנסין: בעליו אין עמו. אם בעליו של שור אינו עם השואל במלאכתו (בבא מציעא צ"ה):}%
% ---------- פרק 22 | פסוק 14 ----------
 \textbf{יד} אִם־בְּעָלָ֥יו עִמּ֖וֹ לֹ֣א יְשַׁלֵּ֑ם אִם־שָׂכִ֣יר ה֔וּא בָּ֖א בִּשְׂכָרֽוֹ׃ \{ס\} \Rashi{\textbf{אם בעליו עמו.}בין שהוא באותה מלאכה בין שהוא במלאכה אחרת, היה עמו בשעת שאלה אינו צריך להיות עמו בשעת שבירה ומיתה (שם): אם שכיר הוא. אם השור אינו שאול אלא שכור, בא בשכרו ליד השוכר הזה, ולא בשאלה, ואין כל הנאה שלו, שהרי ע"י שכרו נשתמש, ואין לו משפט שואל להתחיב באנסין; ולא פרש מה דינו, אם כשומר חנם או כשומר שכר, לפיכך נחלקו בו חכמי ישראל שוכר כיצד משלם, רבי מאיר אומר כשומר חנם, ר' יהודה אומר כשומר שכר (בבא מציעא צ"ה):}%
% ---------- פרק 22 | פסוק 15 ----------
 \textbf{טו} וְכִֽי־יְפַתֶּ֣ה אִ֗ישׁ בְּתוּלָ֛ה אֲשֶׁ֥ר לֹא־אֹרָ֖שָׂה וְשָׁכַ֣ב עִמָּ֑הּ מָהֹ֛ר יִמְהָרֶ֥נָּה לּ֖וֹ לְאִשָּֽׁה׃ \Rashi{\textbf{וכי יפתה.}מדבר על לבה עד ששומעת לו, וכן תרגומו וארי ישדל, שדול בלשון ארמי כפתוי בלשון עברי: מהר ימהרנה. יפסק לה מהר כמשפט איש לאשתו – שכותב לה כתבה וישאנה (מכילתא):}%
% ---------- פרק 22 | פסוק 16 ----------
 \textbf{טז} אִם־מָאֵ֧ן יְמָאֵ֛ן אָבִ֖יהָ לְתִתָּ֣הּ ל֑וֹ כֶּ֣סֶף יִשְׁקֹ֔ל כְּמֹ֖הַר הַבְּתוּלֹֽת׃ \{ס\} \Rashi{\textbf{כמהר הבתולת.}שהוא קצוב חמשים כסף אצל התופס את הבתולה ושוכב עמה באנס, שנאמר ונתן האיש השכב עמה לאבי הנערה חמשים כסף (דברים כ"ב):}%
% ---------- פרק 22 | פסוק 17 ----------
 \textbf{יז} מְכַשֵּׁפָ֖ה לֹ֥א תְחַיֶּֽה׃ \Rashi{\textbf{מכשפה לא תחיה.}אלא תומת בבית דין; ואחד זכרים ואחד נקבות, אלא שדבר הכתוב בהווה, שהנשים מצויות מכשפות (סנהדרין ס"ז):}%
% ---------- פרק 22 | פסוק 18 ----------
 \textbf{יח} כׇּל־שֹׁכֵ֥ב עִם־בְּהֵמָ֖ה מ֥וֹת יוּמָֽת׃ \{ס\} \Rashi{\textbf{כל שכב עם בהמה מות יומת.}בסקילה, רובע כנרבעת, שכתוב בהם דמיהם בם (ויקרא כ'):}%
% ---------- פרק 22 | פסוק 19 ----------
 \textbf{יט} זֹבֵ֥חַ לָאֱלֹהִ֖ים יׇֽחֳרָ֑ם בִּלְתִּ֥י לַיהֹוָ֖ה לְבַדּֽוֹ׃ \Rashi{\textbf{לאלהים.}לע"ז; אלו היה נקוד לאלהים, היה צריך לפרש ולכתב "אחרים", עכשו שאמר לאלהים, אין צריך לפרש אחרים, שכל למ"ד ובי"ת המשמשת בראש התבה אם נקודה בחטף, כגון למלך, למדבר, לעיר, צריך לפרש לאיזה מלך, לאיזה מדבר, לאיזה עיר, וכן למלכים ולרגלים צריך לפרש לאיזה, ואם אינו מפרש, כל מלכים במשמע, כן לאלהים כל אלהים במשמע – אפלו קדש, אבל כשהיא נקודה פתח, כמו למלך, למדבר, לעיר, נודע באיזה מלך מדבר, וכן לעיר נודע באיזה עיר מדבר, וכן לאלהים – לאותן שהזהרתם עליהם במקום אחר; כיוצא בו אין כמוך באלהים, לפי שלא פרש, הצרך לנקד פתח: יחרם. יומת; למה נאמר יחרם והלא כבר נאמרה בו מיתה במקום אחר והוצאת את האיש ההוא או את האשה ההוא וגו' (דברים י"ז)? אלא לפי שלא פרש על איזו עבודה חיב מיתה – שלא תאמר כל עבודות במיתה – בא ופרש לך כאן "זבח" לאלהים, לומר לך, מה זביחה עבודה הנעשית בפנים לשמים, אף אני מרבה המקטיר והמנסך שהם עבודות בפנים, וחיבים עליהם לכל ע"ז בין שדרכה לעבדה בכך, בין שאין דרכה לעבדה בכך. אבל שאר עבודות, כגון המכבד והמרבץ והמגפף והמנשק אינו במיתה:}%
% ---------- פרק 22 | פסוק 20 ----------
 \textbf{כ} וְגֵ֥ר לֹא־תוֹנֶ֖ה וְלֹ֣א תִלְחָצֶ֑נּוּ כִּֽי־גֵרִ֥ים הֱיִיתֶ֖ם בְּאֶ֥רֶץ מִצְרָֽיִם׃ \Rashi{\textbf{וגר לא תונה.}אונאת דברים, קונטרארי"ר בלעז, כמו והאכלתי את מוניך את בשרם (ישעיהו מ"ט): ולא תלחצנו. בגזלת ממון: כי גרים הייתם. אם הוניתו, אף הוא יכול להונותך ולומר לך, אף אתה מגרים באת, "מום שבך אל תאמר לחברך"; כל לשון גר אדם שלא נולד באותה מדינה, אלא בא ממדינה אחרת לגור שם:}%
% ---------- פרק 22 | פסוק 21 ----------
 \textbf{כא} כׇּל־אַלְמָנָ֥ה וְיָת֖וֹם לֹ֥א תְעַנּֽוּן׃ \Rashi{\textbf{כל אלמנה ויתום לא תענון.}הוא הדין לכל אדם, אלא שדבר הכתוב בהווה, לפי שהם תשושי כח ודבר מצוי לענותם (מכילתא):}%
% ---------- פרק 22 | פסוק 22 ----------
 \textbf{כב} אִם־עַנֵּ֥ה תְעַנֶּ֖ה אֹת֑וֹ כִּ֣י אִם־צָעֹ֤ק יִצְעַק֙ אֵלַ֔י שָׁמֹ֥עַ אֶשְׁמַ֖ע צַעֲקָתֽוֹ׃ \Rashi{\textbf{אם ענה תענה אתו.}הרי זה מקרא קצר, גזם ולא פרש ענשו, כמו לכן כל הרג קין (בראשית ד'), גזם ולא פרש ענשו, אף כאן אם ענה תענה אותו לשון גזום, כלומר סופך לטל את שלך למה? כי אם צעק יצעק אלי וגו':}%
% ---------- פרק 22 | פסוק 23 ----------
 \textbf{כג} וְחָרָ֣ה אַפִּ֔י וְהָרַגְתִּ֥י אֶתְכֶ֖ם בֶּחָ֑רֶב וְהָי֤וּ נְשֵׁיכֶם֙ אַלְמָנ֔וֹת וּבְנֵיכֶ֖ם יְתֹמִֽים׃ \{פ\} \Rashi{\textbf{והיו נשיכם אלמנות.}ממשמע שנאמר והרגתי אתכם איני יודע שנשיכם אלמנות ובניכם יתומים? אלא הרי זו קללה אחרת, שיהיו הנשים צרורות כאלמנות חיות, שלא יהיו עדים למיתת בעליהן ותהיינה אסורות להנשא, והבנים יהיו יתומים, שלא יניחום בית דין לירד לנכסי אביהם, לפי שאין יודעים אם מתו אם נשבו (בבא מציעא ל"ח):}%
% ---------- פרק 22 | פסוק 24 ----------
 \textbf{כד} אִם־כֶּ֣סֶף ׀ תַּלְוֶ֣ה אֶת־עַמִּ֗י אֶת־הֶֽעָנִי֙ עִמָּ֔ךְ לֹא־תִהְיֶ֥ה ל֖וֹ כְּנֹשֶׁ֑ה לֹֽא־תְשִׂימ֥וּן עָלָ֖יו נֶֽשֶׁךְ׃ \Rashi{\textbf{אם כסף תלוה את עמי.}רבי ישמעאל אומר כל אם ואם שבתורה רשות חוץ משלשה, וזה אחד מהן (מכילתא): את עמי. עמי וגוי עמי קודם, עני ועשיר עני קודם, ענייך ועניי עירך ענייך קודמין, עניי עירך ועניי עיר אחרת עניי עירך קודמין; וזה משמעו: אם כסף תלוה – את עמי תלוהו ולא לגוי, ולאיזה מעמי? את העני, ולאיזה עני? לאותו שעמך. (ד"א, את עמי, שלא תנהג בו בזיון בהלואה שהוא עמי: את העני עמך. הוי מסתכל בעצמך כאלו אתה עני: לא תהיה לו כנשה. לא תתבענו בחזקה. אם אתה יודע שאין לו, אל תהי דומה עליו כאלו הלויתו אלא כאלו לא הלויתו, כלומר לא תכלימהו: נשך. רבית, שהוא כנשיכת נחש שנושך חבורה קטנה ברגלו ואינו מרגיש, ופתאם הוא מבטבט ונופח עד קדקדו, כך רבית אינו מרגיש ואינו נכר עד שהרבית עולה ומחסרו ממון הרבה (תנחומא):}%
% ---------- פרק 22 | פסוק 25 ----------
 \textbf{כה} אִם־חָבֹ֥ל תַּחְבֹּ֖ל שַׂלְמַ֣ת רֵעֶ֑ךָ עַד־בֹּ֥א הַשֶּׁ֖מֶשׁ תְּשִׁיבֶ֥נּוּ לֽוֹ׃ \Rashi{\textbf{אם חבל תחבל.}כל לשון חבלה אינו משכון בשעת הלואה, אלא שממשכנין את הלוה כשמגיע הזמן ואינו פורע; (חבל תחבל – כפל לך בחבלה עד כמה פעמים; אמר הקב"ה כמה אתה חיב לי, והרי נפשך עולה אצלי כל אמש ואמש ונותנת דין ומתחיבת לפני ואני מחזירה לך, אף אתה טל והשב טל והשב): עד בא השמש תשיבנו לו. כל היום תשיבנו לו עד בא השמש, וכבא השמש תחזר ותטלנו עד שיבא בקר של מחר; ובכסות יום הכתוב מדבר, שאין צריך לה בלילה (מכילתא, בבא מציעא קי"ד):}%
% ---------- פרק 22 | פסוק 26 ----------
 \textbf{כו} כִּ֣י הִ֤וא כְסוּתֹה֙ לְבַדָּ֔הּ הִ֥וא שִׂמְלָת֖וֹ לְעֹר֑וֹ בַּמֶּ֣ה יִשְׁכָּ֔ב וְהָיָה֙ כִּֽי־יִצְעַ֣ק אֵלַ֔י וְשָׁמַעְתִּ֖י כִּֽי־חַנּ֥וּן אָֽנִי׃ \{ס\} \Rashi{\textbf{כי הוא כסותה.}זו טלית: שמלתו. זה חלוק: במה ישכב. לרבות את המצע (מכילתא):}%
% ---------- פרק 22 | פסוק 27 ----------
 \textbf{כז} אֱלֹהִ֖ים לֹ֣א תְקַלֵּ֑ל וְנָשִׂ֥יא בְעַמְּךָ֖ לֹ֥א תָאֹֽר׃ \Rashi{\textbf{אלהים לא תקלל.}הרי זו אזהרה לברכת השם ואזהרה לקללת דין (סנהדרין ס"ו):}%
% ---------- פרק 22 | פסוק 28 ----------
 \textbf{כח} מְלֵאָתְךָ֥ וְדִמְעֲךָ֖ לֹ֣א תְאַחֵ֑ר בְּכ֥וֹר בָּנֶ֖יךָ תִּתֶּן־לִֽי׃ \Rashi{\textbf{מלאתך.}חובה המטלת עליך כשתתמלא תבואתך להתבשל, והם בכורים: ודמעך. התרומה – ואיני יודע מהו לשון דמע. לא תאחר. לא תשנה סדר הפרשתן לאחר את המקדם ולהקדים את המאחר, שלא יקדים תרומה לבכורים ומעשר לתרומה (מכילתא): בכור בניך תתן לי. לפדותו בחמש סלעים מן הכהן; והלא כבר צוה עליו במקום אחר? אלא כדי לסמך לו כן תעשה לשרך – מה בכור אדם לאחר שלשים יום פודהו, שנאמר ופדויו מבן חדש תפדה (במדבר י"ח), אף בכור בהמה דקה מטפל בו שלשים יום ואח"כ נותנו לכהן (בכורות כ"ו):}%
% ---------- פרק 22 | פסוק 29 ----------
 \textbf{כט} כֵּֽן־תַּעֲשֶׂ֥ה לְשֹׁרְךָ֖ לְצֹאנֶ֑ךָ שִׁבְעַ֤ת יָמִים֙ יִהְיֶ֣ה עִם־אִמּ֔וֹ בַּיּ֥וֹם הַשְּׁמִינִ֖י תִּתְּנוֹ־לִֽי׃ \Rashi{\textbf{שבעת ימים יהיה עם אמו.}זו אזהרה לכהן, שאם בא למהר את קרבנו לא ימהר קדם שמנה, לפי שהוא מחסר זמן: ביום השמיני תתנו לי. יכול יהא חובה לבו ביום, נאמר כאן שמיני ונאמר להלן ומיום השמיני והלאה ירצה (ויקרא כ"ב), מה שמיני האמור להלן להכשיר משמיני ולהלן, אף שמיני האמור כאן להכשיר משמיני ולהלן; וכן משמעו: וביום השמיני אתה רשאי לתנו לי (מכילתא):}%
% ---------- פרק 22 | פסוק 30 ----------
 \textbf{ל} וְאַנְשֵׁי־קֹ֖דֶשׁ תִּהְי֣וּן לִ֑י וּבָשָׂ֨ר בַּשָּׂדֶ֤ה טְרֵפָה֙ לֹ֣א תֹאכֵ֔לוּ לַכֶּ֖לֶב תַּשְׁלִכ֥וּן אֹתֽוֹ׃ \{ס\} \Rashi{\textbf{ואנשי קדש תהיון לי.}אם אתם קדושים ופרושים משקוצי נבלות וטרפות הרי אתם שלי ואם לאו אינכם שלי: ובשר בשדה טרפה. אף בבית כן, אלא שדבר הכתוב בהווה – מקום שדרך בהמות לטרף; וכן כי בשדה מצאה (דברים כ"ב), וכן אשר לא יהיה טהור מקרה לילה (שם כ"ג), הוא הדין למקרה יום, אלא שדבר הכתוב בהווה. ובשר דתליש מן חיוא חיא, בשר שנתלש ע"י טרפת זאב או ארי מן חיה כשרה או מבהמה כשרה בחייה: לכלב תשלכון אתו. אף הגוי ככלב; או אינו אלא כלב כמשמעו? ת"ל בנבלה או מכר לנכרי (שם י"ד), ק"ו לטרפה שמתרת בכל הנאות. אם כן מה ת"ל לכלב? ללמדך שהכלב נכבד ממנו, ולמדך הכתוב שאין הקב"ה מקפח שכר כל בריה, שנאמר ולכל בני ישראל לא יחרץ כלב לשנו (שמות י"א), אמר הקב"ה תנו לו שכרו (מכילתא):}%
% ---------- פרק 23 | פסוק 1 ----------
 \textbf{א} לֹ֥א תִשָּׂ֖א שֵׁ֣מַע שָׁ֑וְא אַל־תָּ֤שֶׁת יָֽדְךָ֙ עִם־רָשָׁ֔ע לִהְיֹ֖ת עֵ֥ד חָמָֽס׃ \Rashi{\textbf{לא תשא שמע שוא.}כתרגומו, לא תקבל שמע דשקר, אזהרה למקבל לשון הרע ולדין שלא ישמע דברי בעל דין עד שיבא בעל דין חברו (שם): אל תשת ידך עם רשע. הטוען את חברו תביעת שקר, שתבטיחהו להיות לו עד חמס:}%
% ---------- פרק 23 | פסוק 2 ----------
 \textbf{ב} לֹֽא־תִהְיֶ֥ה אַחֲרֵֽי־רַבִּ֖ים לְרָעֹ֑ת וְלֹא־תַעֲנֶ֣ה עַל־רִ֗ב לִנְטֹ֛ת אַחֲרֵ֥י רַבִּ֖ים לְהַטֹּֽת׃ \Rashi{\textbf{לא תהיה אחרי רבים לרעת.}יש במקרא זה מדרשי חכמי ישראל, אבל אין לשון המקרא מישב בהן על אפניו. מכאן דרשו שאין מטין לחובה בהכרעת דין אחד, וסוף המקרא דרשו אחרי רבים להטת, שאם יש שנים מחיבין יותר על המזכין הטה הדין על פיהם לחובה – ובדיני נפשות הכתוב מדבר – ואמצע המקרא דרשו לא תענה על רב – על רב, שאין חולקין על מפלא שבבית דין, לפיכך מתחילין בדיני נפשות מן הצד – לקטנים שבהם שואלין תחלה שיאמרו את דעתם – ולפי דברי רבותינו כך פתרון המקרא: לא תהיה אחרי רבים לרעת. לחיב מיתה בשביל דין אחד שירבו מחיבין על המזכין, ולא תענה על הרב לנטות מדבריו – ולפי שהוא חסר יו"ד דרשו בו כן – אחרי רבים להטת, יש רבים שאתה נוטה אחריהם, ואימתי? בזמן שהן שנים המכריעין במחיבין יותר מן המזכין; וממשמע שנאמר לא תהיה אחרי רבים לרעת שומע אני אבל היה עמהם לטובה, מכאן אמרו דיני נפשות מטין על פי אחד לזכות ועל פי שנים לחובה. ואנקלוס תרגם לא תתמנע מלאלפא מה דמתבעי לך (דבעינך) על דינא, ולשון העברי לפי התרגום כך הוא נדרש: לא תענה על רב לנטת. אם ישאלוך דבר למשפט לא תענה לנטות לצד אחד ולסלק עצמך מן הריב, אלא הוי דן אותו לאמתו. ואני אומר לישבו על אפניו כפשוטו כך פתרונו: לא תהיה אחרי רבים לרעת. אם ראית רשעים מטין משפט, לא תאמר, הואיל ורבים הם הנני נוטה אחריהם: ולא תענה על רב לנטת וגו'. ואם ישאלך הנדון על אותו המשפט, אל תעננו על הריב דבר הנוטה אחרי אותן רבים להטות את המשפט מאמתו, אלא אמר את המשפט כאשר הוא וקולר יהא תלוי בצואר הרבים (סנהדרין ז'):}%
% ---------- פרק 23 | פסוק 3 ----------
 \textbf{ג} וְדָ֕ל לֹ֥א תֶהְדַּ֖ר בְּרִיבֽוֹ׃ \{ס\} \Rashi{\textbf{לא תהדר.}לא תחלק לו כבוד לזכותו בדין ולומר דל הוא אזכנו ואכבדנו:}%
% ---------- פרק 23 | פסוק 4 ----------
 \textbf{ד} כִּ֣י תִפְגַּ֞ע שׁ֧וֹר אֹֽיִבְךָ֛ א֥וֹ חֲמֹר֖וֹ תֹּעֶ֑ה הָשֵׁ֥ב תְּשִׁיבֶ֖נּוּ לֽוֹ׃ \{ס\} %
% ---------- פרק 23 | פסוק 5 ----------
 \textbf{ה} כִּֽי־תִרְאֶ֞ה חֲמ֣וֹר שֹׂנַאֲךָ֗ רֹבֵץ֙ תַּ֣חַת מַשָּׂא֔וֹ וְחָדַלְתָּ֖ מֵעֲזֹ֣ב ל֑וֹ עָזֹ֥ב תַּעֲזֹ֖ב עִמּֽוֹ׃ \{ס\} \Rashi{\textbf{כי תראה חמור שנאך וגו'.}הרי כי משמש לשון דלמא, שהוא מארבע לשונות של שמושי כי, וכה פתרונו: שמא תראה חמורו רבץ תחת משאו וחדלת מעזב לו, בתמיה: עזב תעזב עמו. עזיבה זו לשון עזרה, וכן עצור ועזוב (דברים ל"ב), וכן ויעזבו ירושלים עד החומה (נחמיה ג') – מלאוה עפר לעזר ולסיע את חזק החומה. כיוצא בו כי תאמר בלבבך רבים הגוים האלה ממני וגו' (דברים ז'), שמא תאמר כן, בתמיה? לא תירא מהם. ומדרשו כך דרשו רבותינו: כי תראה וחדלת, פעמים שאתה חודל, ופעמים שאתה עוזר, הא כיצד? זקן ואינו לפי כבודו וחדלת, או בהמת גוי ומשאו של ישראל וחדלת (מכילתא): עזב תעזב עמו. לפרק המשא; מלמשקל ליה – מלטל משאו ממנו:}%
% ---------- פרק 23 | פסוק 6 ----------
 \textbf{ו} לֹ֥א תַטֶּ֛ה מִשְׁפַּ֥ט אֶבְיֹנְךָ֖ בְּרִיבֽוֹ׃ \Rashi{\textbf{אבינך.}לשון אובה, שהוא מדלדל ותאב לכל טובה:}%
% ---------- פרק 23 | פסוק 7 ----------
 \textbf{ז} מִדְּבַר־שֶׁ֖קֶר תִּרְחָ֑ק וְנָקִ֤י וְצַדִּיק֙ אַֽל־תַּהֲרֹ֔ג כִּ֥י לֹא־אַצְדִּ֖יק רָשָֽׁע׃ \Rashi{\textbf{ונקי וצדיק אל תהרג.}מנין ליוצא מבית דין חיב ואמר אחד יש לי ללמד עליו זכות שמחזירין אותו? ת"ל ונקי אל תהרג, ואע"פ שאינו צדיק, שלא נצטדק בבית דין, מכל מקום נקי הוא מדין מיתה שהרי יש לך לזכותו. ומנין ליוצא מבית דין זכאי ואמר אחד יש לי ללמד עליו חובה שאין מחזירין אותו לבית דין? ת"ל וצדיק אל תהרג, וזה צדיק הוא שנצטדק בבית דין (סנהדרין ל"ג): כי לא אצדיק רשע. אין עליך להחזירו, כי אני לא אצדיקנו בדיני, אם יצא מידך זכאי יש לי שלוחים הרבה להמיתו במיתה שנתחיב בה (מכילתא):}%
% ---------- פרק 23 | פסוק 8 ----------
 \textbf{ח} וְשֹׁ֖חַד לֹ֣א תִקָּ֑ח כִּ֤י הַשֹּׁ֙חַד֙ יְעַוֵּ֣ר פִּקְחִ֔ים וִֽיסַלֵּ֖ף דִּבְרֵ֥י צַדִּיקִֽים׃ \Rashi{\textbf{ושחד לא תקח.}אפלו לשפט אמת, וכל שכן כדי להטות את הדין, שהרי להטות את הדין נאמר כבר (דברים ט"ז) לא תטה משפט (כתובות ק"ה): יעור פקחים. אפלו חכם בתורה ונוטל שחד סוף שתטרף דעתו עליו וישתכח תלמודו ויכהה מאור עיניו (שם): ויסלף. כתרגומו ומקלקל: דברי צדיקים. דברים המצדקים, משפטי אמת, וכן תרגומו פתגמין תריצין – ישרים:}%
% ---------- פרק 23 | פסוק 9 ----------
 \textbf{ט} וְגֵ֖ר לֹ֣א תִלְחָ֑ץ וְאַתֶּ֗ם יְדַעְתֶּם֙ אֶת־נֶ֣פֶשׁ הַגֵּ֔ר כִּֽי־גֵרִ֥ים הֱיִיתֶ֖ם בְּאֶ֥רֶץ מִצְרָֽיִם׃ \Rashi{\textbf{וגר לא תלחץ.}בהרבה מקומות הזהירה תורה על הגר מפני שסורו רע (בבא מציעא נ"ט): את נפש הגר. כמה קשה לו כשלוחצים אותו:}%
% ---------- פרק 23 | פסוק 10 ----------
 \textbf{י} וְשֵׁ֥שׁ שָׁנִ֖ים תִּזְרַ֣ע אֶת־אַרְצֶ֑ךָ וְאָסַפְתָּ֖ אֶת־תְּבוּאָתָֽהּ׃ \Rashi{\textbf{ואספת את תבואתה.}לשון הכנסה לבית, כמו ואספתו אל תוך ביתך (דברים כ"ב):}%
% ---------- פרק 23 | פסוק 11 ----------
 \textbf{יא} וְהַשְּׁבִיעִ֞ת תִּשְׁמְטֶ֣נָּה וּנְטַשְׁתָּ֗הּ וְאָֽכְלוּ֙ אֶבְיֹנֵ֣י עַמֶּ֔ךָ וְיִתְרָ֕ם תֹּאכַ֖ל חַיַּ֣ת הַשָּׂדֶ֑ה כֵּֽן־תַּעֲשֶׂ֥ה לְכַרְמְךָ֖ לְזֵיתֶֽךָ׃ \Rashi{\textbf{תשמטנה.}מעבודה: ונטשתה. מאכילה אחר זמן הבעור. ד"א, תשמטנה – מעבודה גמורה, כגון חרישה וזריעה, ונטשתה – מלזבל ומלקשקש (סוכה מ"ד): ויתרם תאכל חית השדה. להקיש מאכל אביון למאכל חיה, מה חיה אוכלת בלא מעשר, אף אביונים אוכלים בלא מעשר; מכאן אמרו אין מעשר בשביעית (מכילתא): כן תעשה לכרמך. ותחלת המקרא מדבר בשדה הלבן, כמו שאמור למעלה הימנו תזרע את ארצך:}%
% ---------- פרק 23 | פסוק 12 ----------
 \textbf{יב} שֵׁ֤שֶׁת יָמִים֙ תַּעֲשֶׂ֣ה מַעֲשֶׂ֔יךָ וּבַיּ֥וֹם הַשְּׁבִיעִ֖י תִּשְׁבֹּ֑ת לְמַ֣עַן יָנ֗וּחַ שֽׁוֹרְךָ֙ וַחֲמֹרֶ֔ךָ וְיִנָּפֵ֥שׁ בֶּן־אֲמָתְךָ֖ וְהַגֵּֽר׃ \Rashi{\textbf{וביום השביעי תשבת.}אף בשנה השביעית לא תעקר שבת בראשית ממקומה, שלא תאמר, הואיל וכל השנה קרויה שבת לא תנהג בה שבת בראשית (שם): למען ינוח שורך וחמרך. תן לו ניח, להתיר שיהא תולש ואוכל עשבים מן הקרקע; או אינו אלא יחבשנו בתוך הבית? אמרת אין זה ניח אלא צער: בן אמתך. בעבד ערל הכתוב מדבר: והגר. גר תושב:}%
% ---------- פרק 23 | פסוק 13 ----------
 \textbf{יג} וּבְכֹ֛ל אֲשֶׁר־אָמַ֥רְתִּי אֲלֵיכֶ֖ם תִּשָּׁמֵ֑רוּ וְשֵׁ֨ם אֱלֹהִ֤ים אֲחֵרִים֙ לֹ֣א תַזְכִּ֔ירוּ לֹ֥א יִשָּׁמַ֖ע עַל־פִּֽיךָ׃ \Rashi{\textbf{ובכל אשר אמרתי אליכם תשמרו.}לעשות כל מצות עשה באזהרה, שכל שמירה שבתורה אזהרה היא במקום לאו (מנחות ל"ו): לא תזכירו. שלא יאמר לו שמר לי בצד ע"ז פלונית, או תעמד עמי ביום ע"ז פלונית; ד"א — ובכל אשר אמרתי אליכם תשמרו ושם אלהים אחרים לא תזכירו, ללמדך ששקולה ע"ז כנגד כל המצוות כלן (הוריות ח'), והנזהר בה כשומר את כלן: לא ישמע מן הגוי על פיך. שלא תעשה שתפות עם גוי וישבע לך בע"ז שלו נמצאת שאתה גורם שיזכר על ידך (סנהדרין ס"ז):}%
% ---------- פרק 23 | פסוק 14 ----------
 \textbf{יד} שָׁלֹ֣שׁ רְגָלִ֔ים תָּחֹ֥ג לִ֖י בַּשָּׁנָֽה׃ \Rashi{\textbf{רגלים.}פעמים; וכן כי הכיתני זה שלש רגלים (במדבר כ"ב):}%
% ---------- פרק 23 | פסוק 15 ----------
 \textbf{טו} אֶת־חַ֣ג הַמַּצּוֹת֮ תִּשְׁמֹר֒ שִׁבְעַ֣ת יָמִים֩ תֹּאכַ֨ל מַצּ֜וֹת כַּֽאֲשֶׁ֣ר צִוִּיתִ֗ךָ לְמוֹעֵד֙ חֹ֣דֶשׁ הָֽאָבִ֔יב כִּי־ב֖וֹ יָצָ֣אתָ מִמִּצְרָ֑יִם וְלֹא־יֵרָא֥וּ פָנַ֖י רֵיקָֽם׃ \Rashi{\textbf{חדש האביב.}שהתבואה מתמלאת בו באביה. אביב לשון אב, בכור וראשון לבשל פרות: ולא יראו פני ריקם. כשתבאו לראות פני ברגלים, הביאו לי עולות (מכילתא):}%
% ---------- פרק 23 | פסוק 16 ----------
 \textbf{טז} וְחַ֤ג הַקָּצִיר֙ בִּכּוּרֵ֣י מַעֲשֶׂ֔יךָ אֲשֶׁ֥ר תִּזְרַ֖ע בַּשָּׂדֶ֑ה וְחַ֤ג הָֽאָסִף֙ בְּצֵ֣את הַשָּׁנָ֔ה בְּאׇסְפְּךָ֥ אֶֽת־מַעֲשֶׂ֖יךָ מִן־הַשָּׂדֶֽה׃ \Rashi{\textbf{וחג הקציר.}הוא חג שבועות: בכורי מעשיך. שהוא זמן הבאת בכורים, ששתי הלחם הבאין בעצרת היו מתירין החדש למנחות ולהביא בכורים למקדש, שנאמר וביום הבכורים וגו' (במדבר כ"ח): וחג האסף. הוא חג הסכות: באספך את מעשיך. שכל ימות החמה התבואה מתיבשת בשדות, ובחג אוספים אותה אל הבית מפני הגשמים:}%
% ---------- פרק 23 | פסוק 17 ----------
 \textbf{יז} שָׁלֹ֥שׁ פְּעָמִ֖ים בַּשָּׁנָ֑ה יֵרָאֶה֙ כׇּל־זְכ֣וּרְךָ֔ אֶל־פְּנֵ֖י הָאָדֹ֥ן ׀ יְהֹוָֽה׃ \Rashi{\textbf{שלש פעמים וגו'.}לפי שהענין מדבר בשביעית, הצרך לומר שלא יסתרסו רגלים ממקומן: כל זכורך. הזכרים שבך:}%
% ---------- פרק 23 | פסוק 18 ----------
 \textbf{יח} לֹֽא־תִזְבַּ֥ח עַל־חָמֵ֖ץ דַּם־זִבְחִ֑י וְלֹֽא־יָלִ֥ין חֵֽלֶב־חַגִּ֖י עַד־בֹּֽקֶר׃ \Rashi{\textbf{לא תזבח על חמץ וגו'.}לא תשחט את הפסח בי"ד בניסן עד שתבער החמץ (מכילתא): ולא ילין חלב חגי חוץ למזבח: עד בקר. יכול אף על המערכה יפסל בלינה, ת"ל על מוקדה על המזבח כל הלילה (ויקרא ו'): ולא ילין. אין לינה אלא בעמוד השחר, שנאמר עד בקר, אבל כל הלילה יכול להעלותו מן הרצפה למזבח (מגילה כ'):}%
% ---------- פרק 23 | פסוק 19 ----------
 \textbf{יט} רֵאשִׁ֗ית בִּכּוּרֵי֙ אַדְמָ֣תְךָ֔ תָּבִ֕יא בֵּ֖ית יְהֹוָ֣ה אֱלֹהֶ֑יךָ לֹֽא־תְבַשֵּׁ֥ל גְּדִ֖י בַּחֲלֵ֥ב אִמּֽוֹ׃ \{פ\} \Rashi{\textbf{ראשית בכורי אדמתך.}אף השביעית חיבת בבכורים, לכך נאמר אף כאן בכורי אדמתך. כיצד? אדם נכנס לתוך שדהו, רואה תאנה שבכרה, כורך עליה גמי לסימן ומקדישה, ואין בכורים אלא בשבעת המינין האמורין במקרא – ארץ חטה ושערה וגו' (בכורים פ"א): לא תבשל גדי. אף עגל וכבש בכלל גדי, שאין גדי אלא לשון ולד רך, ממה שאתה מוצא בכמה מקומות בתורה שכתוב גדי והצרך לפרש אחריו עזים, כגון אנכי אשלח גדי עזים (בראשית ל"ח), את גדי העזים (שם), שני גדיי עזים (שם כ"ז), ללמדך שכל מקום שנאמר גדי סתם אף עגל וכבש במשמע. ובג' מקומות נכתב בתורה, אחד לאסור אכילה, ואחד לאסור הנאה, ואחד לאסור בשול (מכילתא, חולין קט"ו):}%
% ---------- פרק 23 | פסוק 20 ----------
 \textbf{כ} הִנֵּ֨ה אָנֹכִ֜י שֹׁלֵ֤חַ מַלְאָךְ֙ לְפָנֶ֔יךָ לִשְׁמׇרְךָ֖ בַּדָּ֑רֶךְ וְלַהֲבִ֣יאֲךָ֔ אֶל־הַמָּק֖וֹם אֲשֶׁ֥ר הֲכִנֹֽתִי׃ \Rashi{\textbf{הנה אנכי שלח מלאך.}כאן נתבשרו שעתידין לחטא, ושכינה אומרת להם כי לא אעלה בקרבך (שמות ל"ג): אשר הכנתי. אשר זמנתי לתת לכם, זהו פשוטו; ומדרשו אל המקום אשר הכנתי כבר מקומי כנגדו, וזה אחד מן המקראות שאומרים שבית המקדש של מעלה מכון כנגד של מטה (תנחומא):}%
% ---------- פרק 23 | פסוק 21 ----------
 \textbf{כא} הִשָּׁ֧מֶר מִפָּנָ֛יו וּשְׁמַ֥ע בְּקֹל֖וֹ אַל־תַּמֵּ֣ר בּ֑וֹ כִּ֣י לֹ֤א יִשָּׂא֙ לְפִשְׁעֲכֶ֔ם כִּ֥י שְׁמִ֖י בְּקִרְבּֽוֹ׃ \Rashi{\textbf{אל תמר בו.}לשון המראה, כמו אשר ימרה את פיך (יהושע א'): כי לא ישא לפשעכם. אינו מלמד בכך, שהוא מן הכת שאין חוטאין, ועוד שהוא שליח ואינו עושה אלא שליחותו: כי שמי בקרבו. מחבר לראש המקרא – השמר מפניו כי שמי משתף בו. ורבותינו אמרו זה מטטרון ששמו כשם רבו, מטטרון בגימטריא שדי (סנהדרין ל"ח):}%
% ---------- פרק 23 | פסוק 22 ----------
 \textbf{כב} כִּ֣י אִם־שָׁמ֤וֹעַ תִּשְׁמַע֙ בְּקֹל֔וֹ וְעָשִׂ֕יתָ כֹּ֖ל אֲשֶׁ֣ר אֲדַבֵּ֑ר וְאָֽיַבְתִּי֙ אֶת־אֹ֣יְבֶ֔יךָ וְצַרְתִּ֖י אֶת־צֹרְרֶֽיךָ׃ \Rashi{\textbf{וצרתי.}כתרגומו ואעיק:}%
% ---------- פרק 23 | פסוק 23 ----------
 \textbf{כג} כִּֽי־יֵלֵ֣ךְ מַלְאָכִי֮ לְפָנֶ֒יךָ֒ וֶהֱבִֽיאֲךָ֗ אֶל־הָֽאֱמֹרִי֙ וְהַ֣חִתִּ֔י וְהַפְּרִזִּי֙ וְהַֽכְּנַעֲנִ֔י הַחִוִּ֖י וְהַיְבוּסִ֑י וְהִכְחַדְתִּֽיו׃ %
% ---------- פרק 23 | פסוק 24 ----------
 \textbf{כד} לֹֽא־תִשְׁתַּחֲוֶ֤ה לֵאלֹֽהֵיהֶם֙ וְלֹ֣א תׇֽעׇבְדֵ֔ם וְלֹ֥א תַעֲשֶׂ֖ה כְּמַֽעֲשֵׂיהֶ֑ם כִּ֤י הָרֵס֙ תְּהָ֣רְסֵ֔ם וְשַׁבֵּ֥ר תְּשַׁבֵּ֖ר מַצֵּבֹתֵיהֶֽם׃ \Rashi{\textbf{הרס תהרסם.}לאותם אלוהות: מצבתיהם. אבנים שהם מציבין להשתחוות להם:}%
% ---------- פרק 23 | פסוק 25 ----------
 \textbf{כה} וַעֲבַדְתֶּ֗ם אֵ֚ת יְהֹוָ֣ה אֱלֹֽהֵיכֶ֔ם וּבֵרַ֥ךְ אֶֽת־לַחְמְךָ֖ וְאֶת־מֵימֶ֑יךָ וַהֲסִרֹתִ֥י מַחֲלָ֖ה מִקִּרְבֶּֽךָ׃ \{ס\} %
% ---------- פרק 23 | פסוק 26 ----------
 \textbf{כו} לֹ֥א תִהְיֶ֛ה מְשַׁכֵּלָ֥ה וַעֲקָרָ֖ה בְּאַרְצֶ֑ךָ אֶת־מִסְפַּ֥ר יָמֶ֖יךָ אֲמַלֵּֽא׃ \Rashi{\textbf{לא תהיה משכלה.}אם תעשה רצוני: משכלה. מפלת נפלים או קוברת את בניה קרויה משכלה:}%
% ---------- פרק 23 | פסוק 27 ----------
 \textbf{כז} אֶת־אֵֽימָתִי֙ אֲשַׁלַּ֣ח לְפָנֶ֔יךָ וְהַמֹּתִי֙ אֶת־כׇּל־הָעָ֔ם אֲשֶׁ֥ר תָּבֹ֖א בָּהֶ֑ם וְנָתַתִּ֧י אֶת־כׇּל־אֹיְבֶ֛יךָ אֵלֶ֖יךָ עֹֽרֶף׃ \Rashi{\textbf{והמתי.}כמו והממתי, ותרגומו ואשגש. וכן כל תבה שפעל שלה בכפל אות אחרונה, כשתהפך לדבר בלשון פעלתי, יש מקומות שנוטל אות הכפולה ומדגיש את האות ונוקדו במלאפום כגון והמתי מגזרת והמם גלגל עגלתו (ישעיהו כ"ח), וסבותי (קהלת ב'), מגזרת וסבב בית אל (שמואל א ז׳:ט״ז), דלתי (תהילים קט"ז), מגזרת דללו וחרבו (ישעיהו י"ט), על כפים חקתיך (ישעיהו מ״ט:ט״ז), מגזרת חקקי לב (שופטים ה'), את מי רצותי (שמואל א י״ב:ג׳), מגזרת רצץ עזב דלים (איוב כ'). והמתרגם והמתי ואקטל, טועה הוא, שאלו מגזרת מיתה היה, אין ה"א שלה בפתח ולא מ"ם שלה מדגשת ולא נקודה מלאפום, אלא והמתי כגון והמתה את העם הזה (במדבר י״ד:ט״ו), והתי"ו מדגשת לפי שתבא במקום שתי תוי"ן, האחת נשרשת – לפי שאין מיתה בלא תי"ו – והאחרת משמשת, כמו אמרתי, חטאתי, עשיתי, וכן ונתתי התי"ו מדגשת, שהיא באה במקום שתים, לפי שהיה צריך שלש תוי"ן – שתים ליסוד, כמו ביום תת ה' (יהושע י'), מתת אלהים היא (קהלת ג'), והשלישית לשמוש: ערף. שינוסו מפניך ויהפכו לך ערפם:}%
% ---------- פרק 23 | פסוק 28 ----------
 \textbf{כח} וְשָׁלַחְתִּ֥י אֶת־הַצִּרְעָ֖ה לְפָנֶ֑יךָ וְגֵרְשָׁ֗ה אֶת־הַחִוִּ֧י אֶת־הַֽכְּנַעֲנִ֛י וְאֶת־הַחִתִּ֖י מִלְּפָנֶֽיךָ׃ \Rashi{\textbf{הצרעה.}מין שרץ העוף, והיתה מכה אותם בעיניהם ומטילה בהם ארס והם מתים. והצרעה לא עברה את הירדן, והחתי והכנעני הם ארץ סיחון ועוג, לפיכך מכל שבע אמות לא מנה כאן אלא אלו, וחוי אע"פ שהוא מעבר הירדן והלאה, שנו רבותינו במסכת סוטה (דף ל"ו), על שפת הירדן עמדה וזרקה בהם מרה:}%
% ---------- פרק 23 | פסוק 29 ----------
 \textbf{כט} לֹ֧א אֲגָרְשֶׁ֛נּוּ מִפָּנֶ֖יךָ בְּשָׁנָ֣ה אֶחָ֑ת פֶּן־תִּהְיֶ֤ה הָאָ֙רֶץ֙ שְׁמָמָ֔ה וְרַבָּ֥ה עָלֶ֖יךָ חַיַּ֥ת הַשָּׂדֶֽה׃ \Rashi{\textbf{שממה.}רקנית מבני אדם, לפי שאתם מעט ואין בכם כדי למלאות אותה: ורבה עליך. ותרבה עליך:}%
% ---------- פרק 23 | פסוק 30 ----------
 \textbf{ל} מְעַ֥ט מְעַ֛ט אֲגָרְשֶׁ֖נּוּ מִפָּנֶ֑יךָ עַ֚ד אֲשֶׁ֣ר תִּפְרֶ֔ה וְנָחַלְתָּ֖ אֶת־הָאָֽרֶץ׃ \Rashi{\textbf{עד אשר תפרה.}תרבה, לשון פרי, כמו פרו ורבו (בראשית א'):}%
% ---------- פרק 23 | פסוק 31 ----------
 \textbf{לא} וְשַׁתִּ֣י אֶת־גְּבֻלְךָ֗ מִיַּם־סוּף֙ וְעַד־יָ֣ם פְּלִשְׁתִּ֔ים וּמִמִּדְבָּ֖ר עַד־הַנָּהָ֑ר כִּ֣י ׀ אֶתֵּ֣ן בְּיֶדְכֶ֗ם אֵ֚ת יֹשְׁבֵ֣י הָאָ֔רֶץ וְגֵרַשְׁתָּ֖מוֹ מִפָּנֶֽיךָ׃ \Rashi{\textbf{ושתי.}לשון השתה, והתי"ו מדגשת מפני שבאה תחת שתים, שאין שיתה בלא תי"ו, והאחת לשמוש: עד הנהר. פרת: וגרשתמו. ותגרשם:}%
% ---------- פרק 23 | פסוק 32 ----------
 \textbf{לב} לֹֽא־תִכְרֹ֥ת לָהֶ֛ם וְלֵאלֹֽהֵיהֶ֖ם בְּרִֽית׃ %
% ---------- פרק 23 | פסוק 33 ----------
 \textbf{לג} לֹ֤א יֵשְׁבוּ֙ בְּאַרְצְךָ֔ פֶּן־יַחֲטִ֥יאוּ אֹתְךָ֖ לִ֑י כִּ֤י תַעֲבֹד֙ אֶת־אֱלֹ֣הֵיהֶ֔ם כִּֽי־יִהְיֶ֥ה לְךָ֖ לְמוֹקֵֽשׁ׃ \{פ\} \Rashi{\textbf{כי תעבד וגו'.}הרי אלו כי משמשין במקום אשר, וכן בכמה מקומות, וזהו לשון אי שהוא אחד מארבע לשונות שהכי משמש, וגם מצינו בהרבה מקומות אם משמש בלשון אשר, כמו ואם תקריב מנחת בכורים (ויקרא ב'), שהיא חובה:}%
% ---------- פרק 24 | פסוק 1 ----------
 \textbf{א} וְאֶל־מֹשֶׁ֨ה אָמַ֜ר עֲלֵ֣ה אֶל־יְהֹוָ֗ה אַתָּה֙ וְאַהֲרֹן֙ נָדָ֣ב וַאֲבִיה֔וּא וְשִׁבְעִ֖ים מִזִּקְנֵ֣י יִשְׂרָאֵ֑ל וְהִשְׁתַּחֲוִיתֶ֖ם מֵרָחֹֽק׃ \Rashi{\textbf{ואל משה אמר.}פרשה זו נאמרה קדם עשרת הדברות, בד' בסיון נאמרה לו עלה (שבת פ"ח):}%
% ---------- פרק 24 | פסוק 2 ----------
 \textbf{ב} וְנִגַּ֨שׁ מֹשֶׁ֤ה לְבַדּוֹ֙ אֶל־יְהֹוָ֔ה וְהֵ֖ם לֹ֣א יִגָּ֑שׁוּ וְהָעָ֕ם לֹ֥א יַעֲל֖וּ עִמּֽוֹ׃ \Rashi{\textbf{ונגש משה לבדו.}אל הערפל:}%
% ---------- פרק 24 | פסוק 3 ----------
 \textbf{ג} וַיָּבֹ֣א מֹשֶׁ֗ה וַיְסַפֵּ֤ר לָעָם֙ אֵ֚ת כׇּל־דִּבְרֵ֣י יְהֹוָ֔ה וְאֵ֖ת כׇּל־הַמִּשְׁפָּטִ֑ים וַיַּ֨עַן כׇּל־הָעָ֜ם ק֤וֹל אֶחָד֙ וַיֹּ֣אמְר֔וּ כׇּל־הַדְּבָרִ֛ים אֲשֶׁר־דִּבֶּ֥ר יְהֹוָ֖ה נַעֲשֶֽׂה׃ \Rashi{\textbf{ויבא משה ויספר לעם.}בו ביום: את כל דברי ה'. מצות פרישה והגבלה: ואת כל המשפטים. שבע מצוות שנצטוו בני נח, ושבת וכבוד אב ואם ופרה אדמה ודינין שנתנו להם במרה (סנהדרין נ"ו):}%
% ---------- פרק 24 | פסוק 4 ----------
 \textbf{ד} וַיִּכְתֹּ֣ב מֹשֶׁ֗ה אֵ֚ת כׇּל־דִּבְרֵ֣י יְהֹוָ֔ה וַיַּשְׁכֵּ֣ם בַּבֹּ֔קֶר וַיִּ֥בֶן מִזְבֵּ֖חַ תַּ֣חַת הָהָ֑ר וּשְׁתֵּ֤ים עֶשְׂרֵה֙ מַצֵּבָ֔ה לִשְׁנֵ֥ים עָשָׂ֖ר שִׁבְטֵ֥י יִשְׂרָאֵֽל׃ \Rashi{\textbf{ויכתב משה.}מבראשית ועד מתן תורה, וכתב מצוות שנצטוו במרה: וישכם בבקר. בחמשה בסיון:}%
% ---------- פרק 24 | פסוק 5 ----------
 \textbf{ה} וַיִּשְׁלַ֗ח אֶֽת־נַעֲרֵי֙ בְּנֵ֣י יִשְׂרָאֵ֔ל וַֽיַּעֲל֖וּ עֹלֹ֑ת וַֽיִּזְבְּח֞וּ זְבָחִ֧ים שְׁלָמִ֛ים לַיהֹוָ֖ה פָּרִֽים׃ \Rashi{\textbf{את נערי.}הבכורות:}%
% ---------- פרק 24 | פסוק 6 ----------
 \textbf{ו} וַיִּקַּ֤ח מֹשֶׁה֙ חֲצִ֣י הַדָּ֔ם וַיָּ֖שֶׂם בָּאַגָּנֹ֑ת וַחֲצִ֣י הַדָּ֔ם זָרַ֖ק עַל־הַמִּזְבֵּֽחַ׃ \Rashi{\textbf{ויקח משה את חצי הדם.}מי חלקו? מלאך בא וחלקו (ויקרא רבה ו'): באגנת. שתי אגנות, אחד לחצי דם עולה ואחד לחצי דם שלמים להזות אותם על העם. ומכאן למדו רבותינו שנכנסו אבותינו לברית במילה וטבילה והזאת דמים, שאין הזאה בלא טבילה (כריתות ט'):}%
% ---------- פרק 24 | פסוק 7 ----------
 \textbf{ז} וַיִּקַּח֙ סֵ֣פֶר הַבְּרִ֔ית וַיִּקְרָ֖א בְּאׇזְנֵ֣י הָעָ֑ם וַיֹּ֣אמְר֔וּ כֹּ֛ל אֲשֶׁר־דִּבֶּ֥ר יְהֹוָ֖ה נַעֲשֶׂ֥ה וְנִשְׁמָֽע׃ \Rashi{\textbf{ספר הברית.}מבראשית ועד מתן תורה ומצוות שנצטוו במרה:}%
% ---------- פרק 24 | פסוק 8 ----------
 \textbf{ח} וַיִּקַּ֤ח מֹשֶׁה֙ אֶת־הַדָּ֔ם וַיִּזְרֹ֖ק עַל־הָעָ֑ם וַיֹּ֗אמֶר הִנֵּ֤ה דַֽם־הַבְּרִית֙ אֲשֶׁ֨ר כָּרַ֤ת יְהֹוָה֙ עִמָּכֶ֔ם עַ֥ל כׇּל־הַדְּבָרִ֖ים הָאֵֽלֶּה׃ \Rashi{\textbf{ויזרק.}ענין הזאה, ותרגומו וזרק על מדבחא לכפרא על עמא:}%
% ---------- פרק 24 | פסוק 9 ----------
 \textbf{ט} וַיַּ֥עַל מֹשֶׁ֖ה וְאַהֲרֹ֑ן נָדָב֙ וַאֲבִיה֔וּא וְשִׁבְעִ֖ים מִזִּקְנֵ֥י יִשְׂרָאֵֽל׃ %
% ---------- פרק 24 | פסוק 10 ----------
 \textbf{י} וַיִּרְא֕וּ אֵ֖ת אֱלֹהֵ֣י יִשְׂרָאֵ֑ל וְתַ֣חַת רַגְלָ֗יו כְּמַעֲשֵׂה֙ לִבְנַ֣ת הַסַּפִּ֔יר וּכְעֶ֥צֶם הַשָּׁמַ֖יִם לָטֹֽהַר׃ \Rashi{\textbf{ויראו את אלהי ישראל.}נסתכלו והציצו ונתחיבו מיתה, אלא שלא רצה הקב"ה לערבב שמחת התורה והמתין לנדב ואביהוא עד יום חנכת המשכן, ולזקנים עד ויהי העם כמתאוננים, ותבער בם אש ה' ותאכל בקצה המחנה (במדבר י"א) – בקצינים שבמחנה (תנחומא): כמעשה לבנת הספיר. היא היתה לפניו בשעת השעבוד, לזכר צרתן של ישראל שהיו משעבדים במעשה לבנים (ויקרא רבה): וכעצם השמים לטהר. משנגאלו היה אור וחדוה לפניו: וכעצם. כתרגומו, לשון מראה: לטהר. לשון ברור וצלול:}%
% ---------- פרק 24 | פסוק 11 ----------
 \textbf{יא} וְאֶל־אֲצִילֵי֙ בְּנֵ֣י יִשְׂרָאֵ֔ל לֹ֥א שָׁלַ֖ח יָד֑וֹ וַיֶּֽחֱזוּ֙ אֶת־הָ֣אֱלֹהִ֔ים וַיֹּאכְל֖וּ וַיִּשְׁתּֽוּ׃ \{ס\} \Rashi{\textbf{ואל אצילי.}הם נדב ואביהוא והזקנים: לא שלח ידו. מכלל שהיו ראויים להשתלח בהם יד (תנחומא, ויקרא רבה כ'): ויחזו את האלהים. היו מסתכלין בו בלב גס מתוך אכילה ושתיה, כך מדרש תנחומא, ואנקלוס לא תרגם כן. אצילי. לשון גדולים, כמו ומאציליה קראתיך (ישעיהו מ"א), ויאצל מן הרוח (במדבר י"א), שש אמות אצילה (יחזקאל מ"א):}%
% ---------- פרק 24 | פסוק 12 ----------
 \textbf{יב} וַיֹּ֨אמֶר יְהֹוָ֜ה אֶל־מֹשֶׁ֗ה עֲלֵ֥ה אֵלַ֛י הָהָ֖רָה וֶהְיֵה־שָׁ֑ם וְאֶתְּנָ֨ה לְךָ֜ אֶת־לֻחֹ֣ת הָאֶ֗בֶן וְהַתּוֹרָה֙ וְהַמִּצְוָ֔ה אֲשֶׁ֥ר כָּתַ֖בְתִּי לְהוֹרֹתָֽם׃ \Rashi{\textbf{ויאמר ה' אל משה.}לאחר מתן תורה: עלה אלי ההרה והיה שם. ארבעים יום: את לחת האבן והתורה והמצוה אשר כתבתי להורתם. כל שש מאות ושלש עשרה מצוות בכלל עשרת הדברות הן, ורבנו סעדיה פרש באזהרות שיסד לכל דבור ודבור מצוות התלויות בו:}%
% ---------- פרק 24 | פסוק 13 ----------
 \textbf{יג} וַיָּ֣קׇם מֹשֶׁ֔ה וִיהוֹשֻׁ֖עַ מְשָׁרְת֑וֹ וַיַּ֥עַל מֹשֶׁ֖ה אֶל־הַ֥ר הָאֱלֹהִֽים׃ \Rashi{\textbf{ויקם משה ויהושע משרתו.}לא ידעתי מה טיבו של יהושע כאן, ואומר אני, שהיה התלמיד מלוה לרב עד מקום הגבלת תחומי ההר שאינו רשאי לילך משם והלאה, ומשם ויעל משה לבדו אל הר האלהים, ויהושע נטה שם אהלו ונתעכב שם כל ארבעים יום, שכן מצינו כשירד משה וישמע יהושע את קול העם ברעה (שמות ל״ב:י״ז), למדנו שלא היה יהושע עמהם:}%
% ---------- פרק 24 | פסוק 14 ----------
 \textbf{יד} וְאֶל־הַזְּקֵנִ֤ים אָמַר֙ שְׁבוּ־לָ֣נוּ בָזֶ֔ה עַ֥ד אֲשֶׁר־נָשׁ֖וּב אֲלֵיכֶ֑ם וְהִנֵּ֨ה אַהֲרֹ֤ן וְחוּר֙ עִמָּכֶ֔ם מִי־בַ֥עַל דְּבָרִ֖ים יִגַּ֥שׁ אֲלֵהֶֽם׃ \Rashi{\textbf{ואל הזקנים אמר.}בצאתו מן המחנה: שבו לנו בזה. והתעכבו כאן עם שאר העם במחנה, להיות נכונים לשפט לכל איש ריבו: חור. בנה של מרים היה ואביו כלב בן יפנה, שנאמר ויקח לו כלב את אפרת ותלד לו את חור (דהי"א ב'), אפרת זו מרים כדאיתא בסוטה: מי בעל דברים. מי שיש לו דין:}%
% ---------- פרק 24 | פסוק 15 ----------
 \textbf{טו} וַיַּ֥עַל מֹשֶׁ֖ה אֶל־הָהָ֑ר וַיְכַ֥ס הֶעָנָ֖ן אֶת־הָהָֽר׃ %
% ---------- פרק 24 | פסוק 16 ----------
 \textbf{טז} וַיִּשְׁכֹּ֤ן כְּבוֹד־יְהֹוָה֙ עַל־הַ֣ר סִינַ֔י וַיְכַסֵּ֥הוּ הֶעָנָ֖ן שֵׁ֣שֶׁת יָמִ֑ים וַיִּקְרָ֧א אֶל־מֹשֶׁ֛ה בַּיּ֥וֹם הַשְּׁבִיעִ֖י מִתּ֥וֹךְ הֶעָנָֽן׃ \Rashi{\textbf{ויכסהו הענן.}רבותינו חולקים בדבר, יש מהם או' אלו ששה ימים שמר"ח (עד עצרת יום מתן תורה): ויכסהו הענן. להר: ויקרא אל משה ביום השביעי לומר עשרת הדברות, ומשה וכל בני ישראל עומדים, אלא שחלק הכתוב כבוד למשה; וי"א ויכסהו הענן למשה ששת ימים לאחר עשרת הדברות, והם היו בתחלת ארבעים יום שעלה משה לקבל לוחות, ולמדך שכל הנכנס למחנה שכינה טעון פרישה ששה ימים (יומא ג'):}%
% ---------- פרק 24 | פסוק 17 ----------
 \textbf{יז} וּמַרְאֵה֙ כְּב֣וֹד יְהֹוָ֔ה כְּאֵ֥שׁ אֹכֶ֖לֶת בְּרֹ֣אשׁ הָהָ֑ר לְעֵינֵ֖י בְּנֵ֥י יִשְׂרָאֵֽל׃ %
% ---------- פרק 24 | פסוק 18 ----------
 \textbf{יח} וַיָּבֹ֥א מֹשֶׁ֛ה בְּת֥וֹךְ הֶעָנָ֖ן וַיַּ֣עַל אֶל־הָהָ֑ר וַיְהִ֤י מֹשֶׁה֙ בָּהָ֔ר אַרְבָּעִ֣ים י֔וֹם וְאַרְבָּעִ֖ים לָֽיְלָה׃ \{פ\} \Rashi{\textbf{בתוך הענן.}ענן זה כמין עשן הוא, ועשה לו הקב"ה למשה שביל בתוכו (שם ד'):}%
\addparasha{חומש שמות | פרשת תרומה}
% ---------- פרק 25 | פסוק 1 ----------
 \textbf{א} וַיְדַבֵּ֥ר יְהֹוָ֖ה אֶל־מֹשֶׁ֥ה לֵּאמֹֽר׃ %
% ---------- פרק 25 | פסוק 2 ----------
 \textbf{ב} דַּבֵּר֙ אֶל־בְּנֵ֣י יִשְׂרָאֵ֔ל וְיִקְחוּ־לִ֖י תְּרוּמָ֑ה מֵאֵ֤ת כׇּל־אִישׁ֙ אֲשֶׁ֣ר יִדְּבֶ֣נּוּ לִבּ֔וֹ תִּקְח֖וּ אֶת־תְּרוּמָתִֽי׃ \Rashi{\textbf{ויקחו לי תרומה.}לי – לשמי: תרומה. הפרשה, יפרישו לי מממונם נדבה: ידבנו לבו. לשון נדבה והוא לשון רצון טוב, פיישנ"ט בלעז: תקחו את תרומתי. אמרו רבותינו שלש תרומות אמורות כאן, אחת תרומת בקע לגלגלת שנעשו מהם האדנים, כמו שמפרש באלה פקודי, ואחת תרומת המזבח בקע לגלגלת לקפות לקנות מהן קרבנות צבור, ואחת תרומת המשכן נדבת כל אחד ואחד (תלמוד ירושלמי שקלים א'). י"ג דברים האמורים בענין כלם הצרכו למלאכת המשכן או לבגדי כהנה כשתדקדק בהם:}%
% ---------- פרק 25 | פסוק 3 ----------
 \textbf{ג} וְזֹאת֙ הַתְּרוּמָ֔ה אֲשֶׁ֥ר תִּקְח֖וּ מֵאִתָּ֑ם זָהָ֥ב וָכֶ֖סֶף וּנְחֹֽשֶׁת׃ \Rashi{\textbf{זהב וכסף ונחשת וגו'.}כלם באו בנדבה – איש איש מה שנדבו לבו – חוץ מן הכסף שבא בשוה, מחצית השקל לכל אחד, ולא מצינו בכל מלאכת המשכן שהצרך שם כסף יותר, שנאמר וכסף פקודי העדה וגו' בקע לגלגלת וגו', ושאר הכסף הבא שם בנדבה עשאוהו לכלי שרת:}%
% ---------- פרק 25 | פסוק 4 ----------
 \textbf{ד} וּתְכֵ֧לֶת וְאַרְגָּמָ֛ן וְתוֹלַ֥עַת שָׁנִ֖י וְשֵׁ֥שׁ וְעִזִּֽים׃ \Rashi{\textbf{ותכלת.}צמר צבוע בדם חלזון וצבעו ירק (מנחות מ"ד): וארגמן. צמר צבוע ממין צבע ששמו ארגמן: ושש. הוא פשתן. ועזים. נוצה של עזים, לכך תרגם אנקלוס ומעזי – הבא מן העזים – ולא עזים עצמם, שתרגום של עזים עזיא:}%
% ---------- פרק 25 | פסוק 5 ----------
 \textbf{ה} וְעֹרֹ֨ת אֵילִ֧ם מְאׇדָּמִ֛ים וְעֹרֹ֥ת תְּחָשִׁ֖ים וַעֲצֵ֥י שִׁטִּֽים׃ \Rashi{\textbf{מאדמים.}צבועות היו אדם לאחר עבודן: תחשים. מין חיה, ולא היתה אלא לשעה, והרבה גונים היו לה, לכך מתרגם ססגונא ששש ומתפאר בגונין שלו (שבת כ"ח): ועצי שטים. ומאין היו להם במדבר? פרש רבי תנחומא: יעקב אבינו צפה ברוח הקדש שעתידין ישראל לבנות משכן במדבר, והביא ארזים למצרים ונטעם, וצוה לבניו לטלם עמהם כשיצאו ממצרים:}%
% ---------- פרק 25 | פסוק 6 ----------
 \textbf{ו} שֶׁ֖מֶן לַמָּאֹ֑ר בְּשָׂמִים֙ לְשֶׁ֣מֶן הַמִּשְׁחָ֔ה וְלִקְטֹ֖רֶת הַסַּמִּֽים׃ \Rashi{\textbf{שמן למאר.}שמן זית זך להעלות נר תמיד: בשמים לשמן המשחה. שנעשה למשח כלי המשכן והמשכן לקדשו, והצרכו לו בשמים, כמו שמפרש בכי תשא: ולקטרת הסמים. שהיו מקטירין בכל ערב ובקר, כמו שמפרש בואתה תצוה, ולשון קטרת העלאת קיטור ותמרת עשן:}%
% ---------- פרק 25 | פסוק 7 ----------
 \textbf{ז} אַבְנֵי־שֹׁ֕הַם וְאַבְנֵ֖י מִלֻּאִ֑ים לָאֵפֹ֖ד וְלַחֹֽשֶׁן׃ \Rashi{\textbf{אבני שהם.}שתים הצרכו שם לצרך האפוד האמור בואתה תצוה: מלאים. על שם שעושין להם בזהב מושב כמין גמא, ונותנין האבן שם למלאות הגמא, קרויים אבני מלואים, ומקום המושב קרוי משבצת: לאפד ולחשן. אבני השהם לאפוד ואבני המלאים לחשן; וחשן ואפוד מפרשים בואתה תצוה והם מיני תכשיט:}%
% ---------- פרק 25 | פסוק 8 ----------
 \textbf{ח} וְעָ֥שׂוּ לִ֖י מִקְדָּ֑שׁ וְשָׁכַנְתִּ֖י בְּתוֹכָֽם׃ \Rashi{\textbf{ועשו לי מקדש.}ועשו לשמי בית קדשה:}%
% ---------- פרק 25 | פסוק 9 ----------
 \textbf{ט} כְּכֹ֗ל אֲשֶׁ֤ר אֲנִי֙ מַרְאֶ֣ה אוֹתְךָ֔ אֵ֚ת תַּבְנִ֣ית הַמִּשְׁכָּ֔ן וְאֵ֖ת תַּבְנִ֣ית כׇּל־כֵּלָ֑יו וְכֵ֖ן תַּעֲשֽׂוּ׃ \{ס\} \Rashi{\textbf{ככל אשר אני מראה אותך.}כאן את תבנית המשכן, המקרא הזה מחבר למקרא שלמעלה הימנו ועשו לי מקדש ככל אשר אני מראה אותך: וכן תעשו. לדורות (סנהדרין ט"ז), אם יאבד אחד מן הכלים, או כשתעשו לי כלי בית עולמים, כגון שלחנות ומנורות וכיורות ומכונות שעשה שלמה, כתבנית אלו תעשו אותם; ואם לא היה המקרא מחבר למעלה הימנו, לא היה לו לכתב וכן תעשו אלא כן תעשו, והיה מדבר על עשית אהל מועד וכליו:}%
% ---------- פרק 25 | פסוק 10 ----------
 \textbf{י} וְעָשׂ֥וּ אֲר֖וֹן עֲצֵ֣י שִׁטִּ֑ים אַמָּתַ֨יִם וָחֵ֜צִי אׇרְכּ֗וֹ וְאַמָּ֤ה וָחֵ֙צִי֙ רׇחְבּ֔וֹ וְאַמָּ֥ה וָחֵ֖צִי קֹמָתֽוֹ׃ \Rashi{\textbf{ועשו ארון.}כמין ארונות שעושים בלא רגלים, עשוים כמין ארגז שקורין אישקרי"ן בלעז, יושב על שוליו:}%
% ---------- פרק 25 | פסוק 11 ----------
 \textbf{יא} וְצִפִּיתָ֤ אֹתוֹ֙ זָהָ֣ב טָה֔וֹר מִבַּ֥יִת וּמִח֖וּץ תְּצַפֶּ֑נּוּ וְעָשִׂ֧יתָ עָלָ֛יו זֵ֥ר זָהָ֖ב סָבִֽיב׃ \Rashi{\textbf{מבית ומחוץ תצפנו.}שלשה ארונות עשה בצלאל, שנים של זהב ואחד של עץ, ארבע כתלים ושולים לכל אחד ופתוחים מלמעלה, נתן של עץ בתוך של זהב ושל זהב בתוך של עץ וחפה שפתו העליונה בזהב, נמצא מצפה מבית ומחוץ (יומא ע"ב): זר זהב. כמין כתר מקף לו סביב למעלה משפתו, שעשה הארון החיצון גבוה מן הפנימי, עד שעלה למול עבי הכפרת ולמעלה הימנו משהו, וכשהכפרת שוכב על עבי הכתלים, עולה הזר למעלה מכל עבי הכפרת כל שהוא והוא סימן לכתר תורה:}%
% ---------- פרק 25 | פסוק 12 ----------
 \textbf{יב} וְיָצַ֣קְתָּ לּ֗וֹ אַרְבַּע֙ טַבְּעֹ֣ת זָהָ֔ב וְנָ֣תַתָּ֔ה עַ֖ל אַרְבַּ֣ע פַּעֲמֹתָ֑יו וּשְׁתֵּ֣י טַבָּעֹ֗ת עַל־צַלְעוֹ֙ הָֽאֶחָ֔ת וּשְׁתֵּי֙ טַבָּעֹ֔ת עַל־צַלְע֖וֹ הַשֵּׁנִֽית׃ \Rashi{\textbf{ויצקת.}לשון התכה, כתרגומו: פעמותיו. כתרגומו זויתיה; ובזויות העליונות סמוך לכפרת היו נתונות, שתים מכאן ושתים מכאן לרחבו של ארון, והבדים נתונים בהם, וארכו של ארון מפסיק בין הבדים, אמתים וחצי בין בד לבד, שיהיו שני בני אדם הנושאין את הארון מהלכין ביניהם; וכן מפרש במנחות בפ' שתי הלחם (מנחות דף צ"ח): ושתי טבעת על צלעו האחת. הן הן ארבע טבעות שבתחלת המקרא, ופרש לך היכן היו; והוי"ו זו יתרה היא, ופתרונו כמו שתי טבעות, ויש לך לישבה כן: ושתי מן הטבעות האלו על צלעו האחת: צלעו. צדו:}%
% ---------- פרק 25 | פסוק 13 ----------
 \textbf{יג} וְעָשִׂ֥יתָ בַדֵּ֖י עֲצֵ֣י שִׁטִּ֑ים וְצִפִּיתָ֥ אֹתָ֖ם זָהָֽב׃ \Rashi{\textbf{בדי.}מוטות:}%
% ---------- פרק 25 | פסוק 14 ----------
 \textbf{יד} וְהֵבֵאתָ֤ אֶת־הַבַּדִּים֙ בַּטַּבָּעֹ֔ת עַ֖ל צַלְעֹ֣ת הָאָרֹ֑ן לָשֵׂ֥את אֶת־הָאָרֹ֖ן בָּהֶֽם׃ %
% ---------- פרק 25 | פסוק 15 ----------
 \textbf{טו} בְּטַבְּעֹת֙ הָאָרֹ֔ן יִהְי֖וּ הַבַּדִּ֑ים לֹ֥א יָסֻ֖רוּ מִמֶּֽנּוּ׃ \Rashi{\textbf{לא יסרו ממנו.}לעולם (יומא ע"ב):}%
% ---------- פרק 25 | פסוק 16 ----------
 \textbf{טז} וְנָתַתָּ֖ אֶל־הָאָרֹ֑ן אֵ֚ת הָעֵדֻ֔ת אֲשֶׁ֥ר אֶתֵּ֖ן אֵלֶֽיךָ׃ \Rashi{\textbf{ונתת אל הארן.}כמו בארון: העדת. התורה, שהיא לעדות ביני וביניכם שצויתי אתכם מצוות הכתובות בה:}%
% ---------- פרק 25 | פסוק 17 ----------
 \textbf{יז} וְעָשִׂ֥יתָ כַפֹּ֖רֶת זָהָ֣ב טָה֑וֹר אַמָּתַ֤יִם וָחֵ֙צִי֙ אׇרְכָּ֔הּ וְאַמָּ֥ה וָחֵ֖צִי רׇחְבָּֽהּ׃ \Rashi{\textbf{כפרת.}כסוי על הארון, שהיה פתוח מלמעלה ומניחו עליו כמין דף: אמתים וחצי ארכה. כארכו של ארון, ורחבה כרחבו של ארון, ומנחת על עבי הכתלים ארבעתם, ואע"פ שלא נתן שעור לעביה, פרשו רבותינו שהיה עביה טפח (סוכה ה'):}%
% ---------- פרק 25 | פסוק 18 ----------
 \textbf{יח} וְעָשִׂ֛יתָ שְׁנַ֥יִם כְּרֻבִ֖ים זָהָ֑ב מִקְשָׁה֙ תַּעֲשֶׂ֣ה אֹתָ֔ם מִשְּׁנֵ֖י קְצ֥וֹת הַכַּפֹּֽרֶת׃ \Rashi{\textbf{כרבים.}דמות פרצוף תינוק להם: מקשה תעשה. שלא תעשם בפני עצמם ותחברם בראשי הכפרת לאחר עשיתם כמעשה צורפים שקורין שולד"יר בלעז, אלא הטל זהב הרבה בתחלת עשית הכפרת, והכה בפטיש ובקרנס באמצע, וראשין בולטין למעלה, וציר הכרובים בבליטת קצותיו: מקשה. בטדי"ץ בלעז, כמו דא לדא נקשן (דניאל ה'): קצות הכפרת. ראשי הכפרת:}%
% ---------- פרק 25 | פסוק 19 ----------
 \textbf{יט} וַ֠עֲשֵׂ֠ה כְּר֨וּב אֶחָ֤ד מִקָּצָה֙ מִזֶּ֔ה וּכְרוּב־אֶחָ֥ד מִקָּצָ֖ה מִזֶּ֑ה מִן־הַכַּפֹּ֛רֶת תַּעֲשׂ֥וּ אֶת־הַכְּרֻבִ֖ים עַל־שְׁנֵ֥י קְצוֹתָֽיו׃ \Rashi{\textbf{ועשה כרוב אחד מקצה.}שלא תאמר שנים כרובים לכל קצה וקצה, לכך הצרך לפרש כרוב אחד מקצה מזה: מן הכפרת. עצמה תעשה את הכרובים, זהו פרוש של מקשה תעשה אותם, שלא תעשה בפני עצמם ותחברם לכפרת:}%
% ---------- פרק 25 | פסוק 20 ----------
 \textbf{כ} וְהָי֣וּ הַכְּרֻבִים֩ פֹּרְשֵׂ֨י כְנָפַ֜יִם לְמַ֗עְלָה סֹכְכִ֤ים בְּכַנְפֵיהֶם֙ עַל־הַכַּפֹּ֔רֶת וּפְנֵיהֶ֖ם אִ֣ישׁ אֶל־אָחִ֑יו אֶ֨ל־הַכַּפֹּ֔רֶת יִהְי֖וּ פְּנֵ֥י הַכְּרֻבִֽים׃ \Rashi{\textbf{פרשי כנפים.}שלא תעשה כנפיהם שוכבים אלא פרושים וגבוהים למעלה אצל ראשיהם, שיהא עשרה טפחים בחלל שבין הכנפים לכפרת, כדאיתא בסכה:}%
% ---------- פרק 25 | פסוק 21 ----------
 \textbf{כא} וְנָתַתָּ֧ אֶת־הַכַּפֹּ֛רֶת עַל־הָאָרֹ֖ן מִלְמָ֑עְלָה וְאֶל־הָ֣אָרֹ֔ן תִּתֵּן֙ אֶת־הָ֣עֵדֻ֔ת אֲשֶׁ֥ר אֶתֵּ֖ן אֵלֶֽיךָ׃ \Rashi{\textbf{ואל הארן תתן את העדת.}לא ידעתי למה נכפל, שהרי כבר נאמר ונתת אל הארן את העדת, ויש לומר, שבא ללמד שבעודו ארון לבדו, בלא כפרת, יתן תחלה העדות לתוכו ואחר כך יתן את הכפרת עליו, וכן מצינו כשהקים את המשכן, נאמר ויתן את העדת אל הארן ואחר כך ויתן את הכפרת על הארן מלמעלה:}%
% ---------- פרק 25 | פסוק 22 ----------
 \textbf{כב} וְנוֹעַדְתִּ֣י לְךָ֮ שָׁם֒ וְדִבַּרְתִּ֨י אִתְּךָ֜ מֵעַ֣ל הַכַּפֹּ֗רֶת מִבֵּין֙ שְׁנֵ֣י הַכְּרֻבִ֔ים אֲשֶׁ֖ר עַל־אֲר֣וֹן הָעֵדֻ֑ת אֵ֣ת כׇּל־אֲשֶׁ֧ר אֲצַוֶּ֛ה אוֹתְךָ֖ אֶל־בְּנֵ֥י יִשְׂרָאֵֽל׃ \{פ\} \Rashi{\textbf{ונועדתי.}כשאקבע מועד לך לדבר עמך, אותו מקום אקבע למועד שאבא שם לדבר אליך: ודברתי אתך מעל הכפרת. ובמקום אחר הוא אומר וידבר ה' אליו מאהל מועד לאמר (ויקרא א'), זה המשכן, מחוץ לפרכת, נמצאו שני כתובים מכחישים זה את זה, – בא הכתוב השלישי והכריע ביניהם: ובבא משה אל אהל מועד, וישמע את הקול מדבר אליו מעל הכפרת וגו' (במדבר ז'), משה היה נכנס למשכן, וכיון שבא בתוך הפתח, קול יורד מן השמים לבין הכרובים, ומשם יוצא ונשמע למשה באהל מועד (ספרי): ואת כל אשר אצוה אותך אל בני ישראל. הרי וי"ו זו יתרה וטפלה, וכמוה הרבה במקרא, וכה תפתר: ואת אשר אדבר עמך שם את כל אשר אצוה אותך אל בני ישראל הוא:}%
% ---------- פרק 25 | פסוק 23 ----------
 \textbf{כג} וְעָשִׂ֥יתָ שֻׁלְחָ֖ן עֲצֵ֣י שִׁטִּ֑ים אַמָּתַ֤יִם אׇרְכּוֹ֙ וְאַמָּ֣ה רׇחְבּ֔וֹ וְאַמָּ֥ה וָחֵ֖צִי קֹמָתֽוֹ׃ \Rashi{\textbf{קמתו.}גבה רגליו עם עבי השלחן:}%
% ---------- פרק 25 | פסוק 24 ----------
 \textbf{כד} וְצִפִּיתָ֥ אֹת֖וֹ זָהָ֣ב טָה֑וֹר וְעָשִׂ֥יתָ לּ֛וֹ זֵ֥ר זָהָ֖ב סָבִֽיב׃ \Rashi{\textbf{זר זהב.}סימן לכתר מלכות, שהשלחן שם עשר וגדלה, כמו שאומרים שלחן מלכים:}%
% ---------- פרק 25 | פסוק 25 ----------
 \textbf{כה} וְעָשִׂ֨יתָ לּ֥וֹ מִסְגֶּ֛רֶת טֹ֖פַח סָבִ֑יב וְעָשִׂ֧יתָ זֵר־זָהָ֛ב לְמִסְגַּרְתּ֖וֹ סָבִֽיב׃ \Rashi{\textbf{מסגרת.}כתרגומו גדנפא; ונחלקו חכמי ישראל בדבר, י"א למעלה היתה סביב לשלחן כמו לבזבזין שבשפת שלחן שרים, וי"א למטה היתה תקועה מרגל לרגל בארבע רוחות השלחן, ודף השלחן שוכב על אותה מסגרת: ועשית זר זהב למסגרתו. הוא זר האמור למעלה, ופרש לך כאן שעל המסגרת היתה:}%
% ---------- פרק 25 | פסוק 26 ----------
 \textbf{כו} וְעָשִׂ֣יתָ לּ֔וֹ אַרְבַּ֖ע טַבְּעֹ֣ת זָהָ֑ב וְנָתַתָּ֙ אֶת־הַטַּבָּעֹ֔ת עַ֚ל אַרְבַּ֣ע הַפֵּאֹ֔ת אֲשֶׁ֖ר לְאַרְבַּ֥ע רַגְלָֽיו׃ %
% ---------- פרק 25 | פסוק 27 ----------
 \textbf{כז} לְעֻמַּת֙ הַמִּסְגֶּ֔רֶת תִּהְיֶ֖יןָ הַטַּבָּעֹ֑ת לְבָתִּ֣ים לְבַדִּ֔ים לָשֵׂ֖את אֶת־הַשֻּׁלְחָֽן׃ \Rashi{\textbf{לעמת המסגרת תהיין הטבעת.}ברגלים תקועות, כנגד ראשי המסגרת: לבתים לבדים. אותן טבעות יהיו בתים להכניס בהן הבדים: לבתים. לצרך בתים לבדים, כתרגומו – אתרא לאריחיא:}%
% ---------- פרק 25 | פסוק 28 ----------
 \textbf{כח} וְעָשִׂ֤יתָ אֶת־הַבַּדִּים֙ עֲצֵ֣י שִׁטִּ֔ים וְצִפִּיתָ֥ אֹתָ֖ם זָהָ֑ב וְנִשָּׂא־בָ֖ם אֶת־הַשֻּׁלְחָֽן׃ \Rashi{\textbf{ונשא בם.}לשון נפעל, יהיה נשא בם את השלחן:}%
% ---------- פרק 25 | פסוק 29 ----------
 \textbf{כט} וְעָשִׂ֨יתָ קְּעָרֹתָ֜יו וְכַפֹּתָ֗יו וּקְשׂוֹתָיו֙ וּמְנַקִּיֹּתָ֔יו אֲשֶׁ֥ר יֻסַּ֖ךְ בָּהֵ֑ן זָהָ֥ב טָה֖וֹר תַּעֲשֶׂ֥ה אֹתָֽם׃ \Rashi{\textbf{ועשית קערתיו וכפתיו.}קערותיו זה דפוס שהיה עשוי כדפוס הלחם, והלחם היה עשוי כמין תבה פרוצה משתי רוחותיה, שולים לו למטה, וקופל מכאן ומכאן כלפי מעלה כמין כתלים, ולכך קרוי לחם הפנים, שיש לו פנים רואין לכאן ולכאן לצדי הבית מזה ומזה, ונותן ארכו לרחבו של שלחן, וכתליו זקופים כנגד שפת השלחן; והיה עשוי לו דפוס זהב ודפוס ברזל, בשל ברזל הוא נאפה, וכשמוציאו מן התנור נותנו בשל זהב עד למחר בשבת שמסדרו על השלחן, ואותו דפוס קרוי קערה: וכפתיו. בזיכין שנותנין בהם לבונה שתים היו לשני קמצי לבונה שנותנין על שתי המערכות, שנאמר ונתת על המערכת לבנה זכה (ויקרא כ"ד): וקשותיו. הן כמין חצאי קנים חלולים הנסדקין לארכן (מנחות צ"ו), דגמתן עושה של זהב, ומסדר ג' על ראש כל לחם, שישב לחם האחר על גבי אותן הקנים, ומבדילין בין לחם ללחם כדי שתכנס הרוח ביניהם ולא יתעפשו; ובלשון ערבי כל דבר חלול קרוי קס"וא: ומנקיתיו. תרגומו ומכילתיה, הן סניפים כמין יתדות זהב, עומדין בארץ וגבוהים עד למעלה מן השלחן הרבה, כנגד גבה מערכת הלחם, ומפצלים חמשה פצולים זה למעלה מזה, וראשי הקנים שבין לחם ללחם סמוכין על אותן פצולין, כדי שלא יכבד משא הלחם העליונים על התחתונים וישברו; ולשון מכילתיה סבלותיו, כמו נלאיתי הכיל (ירמיהו ו'). אבל לשון מנקיות איני יודע איך נופל על סניפין; ויש מחכמי ישראל אומרים קשתיו אלו סניפין, שמקשין אותו ומחזיקים אותו שלא ישבר, ומנקיתיו אלו הקנים שמנקין אותו שלא יתעפש (מנחות צ"ז), אבל אנקלוס שתרגם ומכילתיה, היה שונה כדברי האומר מנקיות הן סניפין: אשר יסך בהן. אשר יכסה בהן, ועל קשתיו הוא אומר אשר יסך, שהיו עליו כמין סכך וכסוי, וכן במקום אחר הוא אומר ואת קשות הנסך (במדבר ד'), וזה וזה, יסך והנסך, לשון סכך וכסוי הם:}%
% ---------- פרק 25 | פסוק 30 ----------
 \textbf{ל} וְנָתַתָּ֧ עַֽל־הַשֻּׁלְחָ֛ן לֶ֥חֶם פָּנִ֖ים לְפָנַ֥י תָּמִֽיד׃ \{פ\} \Rashi{\textbf{לחם פנים.}שהיו לו פנים כמו שפרשתי, ומנין הלחם וסדר מערכותיו מפרשים באמר אל הכהנים:}%
% ---------- פרק 25 | פסוק 31 ----------
 \textbf{לא} וְעָשִׂ֥יתָ מְנֹרַ֖ת זָהָ֣ב טָה֑וֹר מִקְשָׁ֞ה תֵּעָשֶׂ֤ה*(בספרי ספרד ואשכנז תֵּיעָשֶׂ֤ה) הַמְּנוֹרָה֙ יְרֵכָ֣הּ וְקָנָ֔הּ גְּבִיעֶ֛יהָ כַּפְתֹּרֶ֥יהָ וּפְרָחֶ֖יהָ מִמֶּ֥נָּה יִהְיֽוּ׃ \Rashi{\textbf{מקשה תיעשה המנורה.}שלא יעשנה חליות ולא יעשה קניה ונרותיה אברים אברים ואחר כך ידביקם, כדרך הצורפים שקורין שולד"יר בלעז, אלא כלה באה מחתיכה אחת, ומקיש בקרנס וחותך בכלי האמנות ומפריד הקנים אילך ואילך: מקשה. תרגומו נגיד, לשון המשכה, שממשיך את האברים מן העשת לכאן ולכאן בהקשת הקרנס, ולשון מקשה מכת קרנס, בטדי"ץ בלעז, כמו דא לדא נקשן (דניאל ה'): תיעשה המנורה. מאליה, לפי שהיה משה מתקשה בה, אמר לו הקב"ה השלך את הככר לאור והיא נעשית מאליה, לכך לא נכתב תעשה (תנחומא): ירכה. הוא הרגל שלמטה העשוי כמין תבה, ושלשה רגלים יוצאין הימנו ולמטה: וקנה. הקנה האמצעי שלה העולה באמצע הירך זקוף כלפי מעלה, ועליו נר האמצעי עשוי כמין בזך, לצוק השמן לתוכו ולתת הפתילה: גביעיה. הן כמין כוסות שעושין מזכוכית, ארכים וקצרים, וקורין להם מדיר"נש בלעז, ואלו עשויין של זהב, ובולטין ויוצאין מכל קנה וקנה כמנין שנתן בהם הכתוב, ולא היו בה אלא לנוי: כפתריה. כמין תפוחים היו, עגלין סביב, בולטין סביבות הקנה האמצעי, כדרך שעושין למנורות שלפני השרים, וקורין להם פומיל"ש בלעז, ומנין שלהם כתוב בפרשה, כמה כפתורים בולטין ממנה וכמה חלק בין כפתור לכפתור: ופרחיה. ציורין עשוין בה כמין פרחין: ממנה יהיו. הכל מקשה יוצא מתוך חתיכת העשת, ולא יעשה לבדם וידביקם:}%
% ---------- פרק 25 | פסוק 32 ----------
 \textbf{לב} וְשִׁשָּׁ֣ה קָנִ֔ים יֹצְאִ֖ים מִצִּדֶּ֑יהָ שְׁלֹשָׁ֣ה ׀ קְנֵ֣י מְנֹרָ֗ה מִצִּדָּהּ֙ הָאֶחָ֔ד וּשְׁלֹשָׁה֙ קְנֵ֣י מְנֹרָ֔ה מִצִּדָּ֖הּ הַשֵּׁנִֽי׃ \Rashi{\textbf{יצאים מצדיה.}לכאן ולכאן באלכסון, נמשכים ועולין עד כנגד גבהה של מנורה שהוא קנה האמצעי, ויוצאין מתוך קנה האמצעי זה למעלה מזה, התחתון ארך, ושל מעלה קצר הימנו, והעליון קצר הימנו, לפי שהיה גבה ראשיהן שוה לגבהו של קנה האמצעי, השביעי, שממנו יוצאים הששה קנים:}%
% ---------- פרק 25 | פסוק 33 ----------
 \textbf{לג} שְׁלֹשָׁ֣ה גְ֠בִעִ֠ים מְֽשֻׁקָּדִ֞ים בַּקָּנֶ֣ה הָאֶחָד֮ כַּפְתֹּ֣ר וָפֶ֒רַח֒ וּשְׁלֹשָׁ֣ה גְבִעִ֗ים מְשֻׁקָּדִ֛ים בַּקָּנֶ֥ה הָאֶחָ֖ד כַּפְתֹּ֣ר וָפָ֑רַח כֵּ֚ן לְשֵׁ֣שֶׁת הַקָּנִ֔ים הַיֹּצְאִ֖ים מִן־הַמְּנֹרָֽה׃ \Rashi{\textbf{משקדים.}כתרגומו, מצירים היו, כדרך שעושין לכלי כסף וזהב שקורים נייל"יר בלעז: שלשה גבעים. בולטין מכל קנה וקנה: כפתר ופרח. היה לכל קנה וקנה:}%
% ---------- פרק 25 | פסוק 34 ----------
 \textbf{לד} וּבַמְּנֹרָ֖ה אַרְבָּעָ֣ה גְבִעִ֑ים מְשֻׁ֨קָּדִ֔ים כַּפְתֹּרֶ֖יהָ וּפְרָחֶֽיהָ׃ \Rashi{\textbf{ובמנרה ארבעה גבעים.}בגופה של מנורה היו ארבעה גביעים – אחד בולט בה למטה מן הקנים, והשלשה למעלה מן יציאת הקנים היוצאין מצדיה: משקדים כפתריה ופרחיה. זה אחד מחמשה מקראות שאין להם הכרע (יומא נ"ב), אין ידוע אם גביעים משקדים או משקדים כפתוריה ופרחיה:}%
% ---------- פרק 25 | פסוק 35 ----------
 \textbf{לה} וְכַפְתֹּ֡ר תַּ֩חַת֩ שְׁנֵ֨י הַקָּנִ֜ים מִמֶּ֗נָּה וְכַפְתֹּר֙ תַּ֣חַת שְׁנֵ֤י הַקָּנִים֙ מִמֶּ֔נָּה וְכַפְתֹּ֕ר תַּחַת־שְׁנֵ֥י הַקָּנִ֖ים מִמֶּ֑נָּה לְשֵׁ֙שֶׁת֙ הַקָּנִ֔ים הַיֹּצְאִ֖ים מִן־הַמְּנֹרָֽה׃ \Rashi{\textbf{וכפתר תחת שני הקנים.}מתוך הכפתור היו הקנים נמשכים משני צדיה אילך ואילך, כך שנינו במלאכת המשכן: גבהה של מנורה י"ח טפחים, הרגלים והפרח שלשה טפחים – הוא הפרח האמור בירך, שנאמר עד ירכה עד פרחה (במדבר ח') – וטפחים חלק, וטפח שבו גביע מהארבעה גביעים וכפתור ופרח משני כפתורים ושני פרחים האמורים במנורה עצמה, שנ' משקדים כפתוריה ופרחיה, למדנו שהיו בקנה שני כפתורים ושני פרחים לבד מן השלשה כפתורים שהקנים נמשכין מתוכן, שנ' וכפתר תחת שני הקנים וגו', וטפחים חלק, וטפח כפתור ושני קנים יוצאים ממנו אילך ואילך נמשכים ועולים כנגד גבהה של מנורה, טפח חלק, וטפח כפתור ושני קנים יוצאים ממנו, וטפח חלק, וטפח כפתור ושני קנים יוצאים ממנו ונמשכים ועולין כנגד גבהה של מנורה, וטפחים חלק, נשתירו שם שלשה טפחים שבהם שלשה גביעים וכפתור ופרח, נמצאו גביעים כ"ב – י"ח לששה קנים, שלשה לכל אחד ואחד, וארבעה בגופה של מנורה, הרי כ"ב – ואחד עשר כפתורים, ששה בששת הקנים, ושלשה בגופה של מנורה שהקנים יוצאים מהם, ושנים עוד במנורה שנ' משקדים כפתוריה, ומעוט כפתורים שנים, האחד למטה אצל הירך, והאחד בג' טפחים העליונים עם ג' הגביעים, ותשעה פרחים היו לה, ששה לששת הקנים, שנ' בקנה האחד כפתר ופרח, ושלשה למנורה שנ' משקדים כפתוריה ופרחיה, ומעוט פרחים שנים, ואחד האמור בפרשת בהעלתך עד ירכה עד פרחה, ואם תדקדק במשנה זו הכתובה למעלה, תמצאם כמנינם איש איש במקומו:}%
% ---------- פרק 25 | פסוק 36 ----------
 \textbf{לו} כַּפְתֹּרֵיהֶ֥ם וּקְנֹתָ֖ם מִמֶּ֣נָּה יִהְי֑וּ כֻּלָּ֛הּ מִקְשָׁ֥ה אַחַ֖ת זָהָ֥ב טָהֽוֹר׃ %
% ---------- פרק 25 | פסוק 37 ----------
 \textbf{לז} וְעָשִׂ֥יתָ אֶת־נֵרֹתֶ֖יהָ שִׁבְעָ֑ה וְהֶֽעֱלָה֙ אֶת־נֵ֣רֹתֶ֔יהָ וְהֵאִ֖יר עַל־עֵ֥בֶר פָּנֶֽיהָ׃ \Rashi{\textbf{את נרתיה.}כמין בזיכין שנותנין בתוכן השמן והפתילות: והאיר על עבר פניה. עשה פי ששת הנרות שבראשי הקנים היוצאים מצדיה מסבים כלפי האמצעי, כדי שיהיו הנרות כשתדליקם מאירים אל עבר פניה – מסב אורם אל צד פני הקנה האמצעי, שהוא גוף המנורה:}%
% ---------- פרק 25 | פסוק 38 ----------
 \textbf{לח} וּמַלְקָחֶ֥יהָ וּמַחְתֹּתֶ֖יהָ זָהָ֥ב טָהֽוֹר׃ \Rashi{\textbf{ומלקחיה.}הם הצבתים העשויין לקח בהם הפתילה מתוך השמן, לישבן ולמשכן בפי הנרות, ועל שם שלוקחים בהם קרויים מלקחים; וצבתהא שתרגם אנקלוס לשון צבת, טנייל"ש בלעז: ומחתתיה. הם כמין בזיכין קטנים שחותה בהן את האפר שבנר בבקר בבקר, כשהוא מיטיב את הנרות מאפר הפתילות שדלקו הלילה וכבו, ולשון מחתה פויישר"יא בלעז, כמו לחתות אש מיקוד (ישעיהו ל'):}%
% ---------- פרק 25 | פסוק 39 ----------
 \textbf{לט} כִּכָּ֛ר זָהָ֥ב טָה֖וֹר יַעֲשֶׂ֣ה אֹתָ֑הּ אֵ֥ת כׇּל־הַכֵּלִ֖ים הָאֵֽלֶּה׃ \Rashi{\textbf{ככר זהב טהור.}שלא יהיה משקלה עם כל כליה אלא ככר, לא פחות ולא יותר (מנחות פ"ח). והככר של חל ששים מנה, ושל קדש היה כפול, ק"ך מנה, והמנה הוא ליטרא ששוקלין בה כסף למשקל קולוניא, והם מאה זהובים – עשרים וחמשה סלעים והסלע ארבעה זהובים:}%
% ---------- פרק 25 | פסוק 40 ----------
 \textbf{מ} וּרְאֵ֖ה וַעֲשֵׂ֑ה בְּתַ֨בְנִיתָ֔ם אֲשֶׁר־אַתָּ֥ה מׇרְאֶ֖ה בָּהָֽר׃ \{ס\} \Rashi{\textbf{וראה ועשה.}ראה כאן בהר תבנית שאני מראה אותך, מגיד שנתקשה משה במעשה המנורה עד שהראה לו הקב"ה מנורה של אש (שם כ"ט): אשר אתה מראה. כתרגומו דאת מתחזי בטורא; אלו היה נקוד מראה, בפתח, היה פתרונו אתה מראה לאחרים, עכשו שנקוד חטף קמץ, פתרונו דאת מתחזי – שאחרים מראים לך (שהנקוד מפריד בין עושה לנעשה):}%
% ---------- פרק 26 | פסוק 1 ----------
 \textbf{א} וְאֶת־הַמִּשְׁכָּ֥ן תַּעֲשֶׂ֖ה עֶ֣שֶׂר יְרִיעֹ֑ת שֵׁ֣שׁ מׇשְׁזָ֗ר וּתְכֵ֤לֶת וְאַרְגָּמָן֙ וְתֹלַ֣עַת שָׁנִ֔י כְּרֻבִ֛ים מַעֲשֵׂ֥ה חֹשֵׁ֖ב תַּעֲשֶׂ֥ה אֹתָֽם׃ \Rashi{\textbf{ואת המשכן תעשה עשר יריעת.}להיות לו לגג ולמחצות מחוץ לקרשים, שהיריעות תלויות מאחוריהן לכסותן: שש משזר ותכלת ותלעת שני. הרי ארבעה מינין יחד חוט וחוט – אחד של פשתן ושלשה של צמר – וכל חוט וחוט כפול ו' הרי ארבעה מינין כשהן שזורין יחד כ"ד כפלים לחוט (יומא ע"א): כרבים מעשה חשב. כרובים היו מצירין בהם באריגתן, ולא ברקימה שהוא מעשה מחט, אלא באריגה בשני כתלים, פרצוף אחד מכאן ופרצוף אחד מכאן – ארי מצד זה ונשר מצד זה – כמו שאורגין חגורות של משי שקורין בלעז פיישי"שא:}%
% ---------- פרק 26 | פסוק 2 ----------
 \textbf{ב} אֹ֣רֶךְ ׀ הַיְרִיעָ֣ה הָֽאַחַ֗ת שְׁמֹנֶ֤ה וְעֶשְׂרִים֙ בָּֽאַמָּ֔ה וְרֹ֙חַב֙ אַרְבַּ֣ע בָּאַמָּ֔ה הַיְרִיעָ֖ה הָאֶחָ֑ת מִדָּ֥ה אַחַ֖ת לְכׇל־הַיְרִיעֹֽת׃ %
% ---------- פרק 26 | פסוק 3 ----------
 \textbf{ג} חֲמֵ֣שׁ הַיְרִיעֹ֗ת תִּֽהְיֶ֙יןָ֙ חֹֽבְרֹ֔ת אִשָּׁ֖ה אֶל־אֲחֹתָ֑הּ וְחָמֵ֤שׁ יְרִיעֹת֙ חֹֽבְרֹ֔ת אִשָּׁ֖ה אֶל־אֲחֹתָֽהּ׃ \Rashi{\textbf{תהיין חברת.}תופרן במחט זו בצד זו, חמש לבד וחמש לבד: אשה אל אחתה. כך דרך המקרא לדבר בדבר שהוא לשון נקבה, ובדבר שהוא לשון זכר אומר איש אל אחיו, כמו שנא' בכרובים ופניהם איש אל אחיו (שמות כ"ה):}%
% ---------- פרק 26 | פסוק 4 ----------
 \textbf{ד} וְעָשִׂ֜יתָ לֻֽלְאֹ֣ת תְּכֵ֗לֶת עַ֣ל שְׂפַ֤ת הַיְרִיעָה֙ הָאֶחָ֔ת מִקָּצָ֖ה בַּחֹבָ֑רֶת וְכֵ֤ן תַּעֲשֶׂה֙ בִּשְׂפַ֣ת הַיְרִיעָ֔ה הַקִּ֣יצוֹנָ֔ה בַּמַּחְבֶּ֖רֶת הַשֵּׁנִֽית׃ \Rashi{\textbf{ללאת.}לצול"ש בלעז, וכן תרגם אנקלוס ענובין, לשון עניבה: מקצה בחברת. באותה יריעה שבסוף החבור, קבוצת חמשת היריעות קרויה חוברת: וכן תעשה בשפת היריעה הקיצונה במחברת השנית. באותה יריעה שהיא קיצונה, לשון קצה, כלומר לסוף החוברת:}%
% ---------- פרק 26 | פסוק 5 ----------
 \textbf{ה} חֲמִשִּׁ֣ים לֻֽלָאֹ֗ת תַּעֲשֶׂה֮ בַּיְרִיעָ֣ה הָאֶחָת֒ וַחֲמִשִּׁ֣ים לֻֽלָאֹ֗ת תַּעֲשֶׂה֙ בִּקְצֵ֣ה הַיְרִיעָ֔ה אֲשֶׁ֖ר בַּמַּחְבֶּ֣רֶת הַשֵּׁנִ֑ית מַקְבִּילֹת֙ הַלֻּ֣לָאֹ֔ת אִשָּׁ֖ה אֶל־אֲחֹתָֽהּ׃ \Rashi{\textbf{מקבילת הללאת אשה אל אחתה.}שמר שתעשה הלולאות במדה אחת, מכונת הבדלתן זו מזו, וכמדתן ביריעה זו כן יהא בחברתה, שכשתפרש חוברת אצל חוברת יהיו הלולאות של יריעה זו מכונות כנגד לולאות של זו, וזהו לשון מקבילות – זו כנגד זו, תרגומו של נגד לקבל. היריעות ארכן כ"ח ורחבן ארבע, וכשחבר חמש יריעות יחד נמצא רחבן עשרים, וכן החוברת השנית, והמשכן ארכו שלשים מן המזרח למערב, שנ' עשרים קרשים לפאת נגב תימנה, וכן לצפון, וכל קרש אמה וחצי האמה, הרי שלשים מן המזרח למערב, רחב המשכן מן הצפון לדרום עשר אמות, שנאמר ולירכתי המשכן ימה וגו' ושני קרשים למקצעת הרי עשר – ובמקומם אפרשם למקראות הללו – נותן היריעות ארכן לרחבו של משכן, עשר אמות אמצעיות לגג חלל רחב המשכן, ואמה מכאן ואמה מכאן לעבי ראשי הקרשים שעבים אמה, נשתירו ט"ז אמה, ח' לצפון וח' לדרום, מכסות קומת הקרשים שגבהן עשר, נמצאו שתי אמות התחתונות מגלות; רחבן של יריעות ארבעים אמה כשהן מחברות, עשרים אמה לחוברת. שלשים מהן לגג חלל המשכן לארכו, ואמה כנגד עבי ראשי הקרשים שבמערב, ואמה לכסות עבי העמודים שבמזרח – שלא היו קרשים במזרח, אלא חמשה עמודים שהמסך פרוש ותלוי בווין שבהן כמין וילון – נשתירו שמונה אמות התלויין על אחורי המשכן שבמערב, ושתי אמות התחתונות מגלות, זו מצאתי בבריתא דארבעים ותשע מדות. אבל במסכת שבת אין היריעות מכסות את עמודי המזרח, ותשע אמות תלויות אחורי המשכן, והכתוב בפרשה זו מסיענו, ונתת את הפרכת "תחת" הקרסים, ואם כדברי הבריתא הזאת, נמצאת פרכת משוכה מן הקרסים ולמערב אמה:}%
% ---------- פרק 26 | פסוק 6 ----------
 \textbf{ו} וְעָשִׂ֕יתָ חֲמִשִּׁ֖ים קַרְסֵ֣י זָהָ֑ב וְחִבַּרְתָּ֨ אֶת־הַיְרִיעֹ֜ת אִשָּׁ֤ה אֶל־אֲחֹתָהּ֙ בַּקְּרָסִ֔ים וְהָיָ֥ה הַמִּשְׁכָּ֖ן אֶחָֽד׃ \Rashi{\textbf{קרסי זהב.}פירמי"לש בלעז, ומכניסין ראשן אחד בלולאות שבחוברת זו וראשן אחד בלולאות שבחוברת זו ומחברן בהן:}%
% ---------- פרק 26 | פסוק 7 ----------
 \textbf{ז} וְעָשִׂ֙יתָ֙ יְרִיעֹ֣ת עִזִּ֔ים לְאֹ֖הֶל עַל־הַמִּשְׁכָּ֑ן עַשְׁתֵּי־עֶשְׂרֵ֥ה יְרִיעֹ֖ת תַּעֲשֶׂ֥ה אֹתָֽם׃ \Rashi{\textbf{יריעת עזים.}מנוצה של עזים: לאהל על המשכן. לפרש אותן על היריעות התחתונות:}%
% ---------- פרק 26 | פסוק 8 ----------
 \textbf{ח} אֹ֣רֶךְ ׀ הַיְרִיעָ֣ה הָֽאַחַ֗ת שְׁלֹשִׁים֙ בָּֽאַמָּ֔ה וְרֹ֙חַב֙ אַרְבַּ֣ע בָּאַמָּ֔ה הַיְרִיעָ֖ה הָאֶחָ֑ת מִדָּ֣ה אַחַ֔ת לְעַשְׁתֵּ֥י עֶשְׂרֵ֖ה יְרִיעֹֽת׃ \Rashi{\textbf{שלשים באמה.}שכשנותן ארכן לרחב המשכן, כמו שנתן את הראשונות, נמצאו אלו עודפות אמה מכאן ואמה מכאן לכסות אחת מהשתי אמות שנשארו מגלות מן הקרשים, והאמה התחתונה של קרש שאין היריעה מכסה אותו היא האמה התחובה בנקב האדן, שהאדנים גבהן אמה:}%
% ---------- פרק 26 | פסוק 9 ----------
 \textbf{ט} וְחִבַּרְתָּ֞ אֶת־חֲמֵ֤שׁ הַיְרִיעֹת֙ לְבָ֔ד וְאֶת־שֵׁ֥שׁ הַיְרִיעֹ֖ת לְבָ֑ד וְכָפַלְתָּ֙ אֶת־הַיְרִיעָ֣ה הַשִּׁשִּׁ֔ית אֶל־מ֖וּל פְּנֵ֥י הָאֹֽהֶל׃ \Rashi{\textbf{וכפלת את היריעה הששית.}העודפת באלו העליונות יותר מן התחתונות: אל מול פני האהל. חצי רחבה היה תלוי וכפול על המסך שבמזרח כנגד הפתח, דומה לכלה צנועה המכסה בצעיף על פניה:}%
% ---------- פרק 26 | פסוק 10 ----------
 \textbf{י} וְעָשִׂ֜יתָ חֲמִשִּׁ֣ים לֻֽלָאֹ֗ת עַ֣ל שְׂפַ֤ת הַיְרִיעָה֙ הָֽאֶחָ֔ת הַקִּיצֹנָ֖ה בַּחֹבָ֑רֶת וַחֲמִשִּׁ֣ים לֻֽלָאֹ֗ת עַ֚ל שְׂפַ֣ת הַיְרִיעָ֔ה הַחֹבֶ֖רֶת הַשֵּׁנִֽית׃ %
% ---------- פרק 26 | פסוק 11 ----------
 \textbf{יא} וְעָשִׂ֛יתָ קַרְסֵ֥י נְחֹ֖שֶׁת חֲמִשִּׁ֑ים וְהֵבֵאתָ֤ אֶת־הַקְּרָסִים֙ בַּלֻּ֣לָאֹ֔ת וְחִבַּרְתָּ֥ אֶת־הָאֹ֖הֶל וְהָיָ֥ה אֶחָֽד׃ %
% ---------- פרק 26 | פסוק 12 ----------
 \textbf{יב} וְסֶ֙רַח֙ הָעֹדֵ֔ף בִּירִיעֹ֖ת הָאֹ֑הֶל חֲצִ֤י הַיְרִיעָה֙ הָעֹדֶ֔פֶת תִּסְרַ֕ח עַ֖ל אֲחֹרֵ֥י הַמִּשְׁכָּֽן׃ \Rashi{\textbf{וסרח העדף ביריעת האהל.}על יריעות המשכן; יריעות האהל הן העליונות של עזים, שקרויים אהל, כמו שנאמר בהם לאהל על המשכן, וכל אהל האמור בהן אינו אלא לשון גג, שמאהילות ומסככות על התחתונות, והן היו עודפות על התחתונות חצי היריעה למערב, שהחצי של יריעה אחת עשרה היתרה היה נכפל אל מול פני האהל, נשארו שתי אמות רחב חציה עודף על רחב התחתונות: תסרח על אחרי המשכן. לכסות שתי אמות שהיו מגלות בקרשים: אחרי המשכן. הוא צד מערבי, לפי שהפתח במזרח שהן פניו, וצפון ודרום קרויין צדדין לימין ולשמאל:}%
% ---------- פרק 26 | פסוק 13 ----------
 \textbf{יג} וְהָאַמָּ֨ה מִזֶּ֜ה וְהָאַמָּ֤ה מִזֶּה֙ בָּעֹדֵ֔ף בְּאֹ֖רֶךְ יְרִיעֹ֣ת הָאֹ֑הֶל יִהְיֶ֨ה סָר֜וּחַ עַל־צִדֵּ֧י הַמִּשְׁכָּ֛ן מִזֶּ֥ה וּמִזֶּ֖ה לְכַסֹּתֽוֹ׃ \Rashi{\textbf{והאמה מזה והאמה מזה.}לצפון ולדרום: בעדף בארך יריעת האהל. שהן עודפות על ארך יריעות המשכן שתי אמות: יהיה סרוח על צדי המשכן. לצפון ולדרום, כמו שפרשתי למעלה. למדה תורה דרך ארץ שיהא אדם חס על היפה (ילקוט שמעוני):}%
% ---------- פרק 26 | פסוק 14 ----------
 \textbf{יד} וְעָשִׂ֤יתָ מִכְסֶה֙ לָאֹ֔הֶל עֹרֹ֥ת אֵילִ֖ם מְאׇדָּמִ֑ים וּמִכְסֵ֛ה עֹרֹ֥ת תְּחָשִׁ֖ים מִלְמָֽעְלָה׃ \{פ\} \Rashi{\textbf{מכסה לאהל.}לאותו גג של יריעות עזים עשה עוד מכסה אחר של עורות אילים מאדמים, ועוד למעלה ממנו מכסה עורות תחשים; ואותן מכסאות לא היו מכסין אלא את הגג, ארכן ל' ורחבן י', אלו דברי רבי נחמיה, ולדברי רבי יהודה מכסה אחד היה, חציו של עורות אילים מאדמים וחציו של עורות תחשים (שבת כ"ח):}%
% ---------- פרק 26 | פסוק 15 ----------
 \textbf{טו} וְעָשִׂ֥יתָ אֶת־הַקְּרָשִׁ֖ים לַמִּשְׁכָּ֑ן עֲצֵ֥י שִׁטִּ֖ים עֹמְדִֽים׃ \Rashi{\textbf{ועשית את הקרשים.}היה לו לומר ועשית קרשים, כמו שנאמר בכל דבר ודבר, מהו הקרשים? מאותן העומדין ומיחדין לכך; יעקב אבינו נטע ארזים במצרים, וכשמת צוה לבניו להעלותם עמהם כשיצאו ממצרים, ואמר להם שעתיד הקב"ה לצוות אותן לעשות משכן במדבר מעצי שטים, ראו שיהיו מזמנים בידכם; הוא שיסד הבבלי בפיוט שלו, "טס מטע מזרזים קורות בתינו ארזים", שנזדרזו להיות מוכנים בידם מקדם לכן: עצי שטים עומדים. איתשט"נביש בלעז; שיהא ארך הקרשים זקוף למעלה בקירות המשכן, ולא תעשה הכתלים בקרשים שוכבים, להיות רחב הקרשים לגבה הכתלים קרש על קרש:}%
% ---------- פרק 26 | פסוק 16 ----------
 \textbf{טז} עֶ֥שֶׂר אַמּ֖וֹת אֹ֣רֶךְ הַקָּ֑רֶשׁ וְאַמָּה֙ וַחֲצִ֣י הָֽאַמָּ֔ה רֹ֖חַב הַקֶּ֥רֶשׁ הָאֶחָֽד׃ \Rashi{\textbf{עשר אמות ארך הקרש.}למדנו גבהו של משכן עשר אמות: ואמה וחצי האמה רחב. למדנו ארכו של משכן, לעשרים קרשים שיהיו בצפון ובדרום מן המזרח למערב, שלשים אמה:}%
% ---------- פרק 26 | פסוק 17 ----------
 \textbf{יז} שְׁתֵּ֣י יָד֗וֹת לַקֶּ֙רֶשׁ֙ הָאֶחָ֔ד מְשֻׁ֨לָּבֹ֔ת אִשָּׁ֖ה אֶל־אֲחֹתָ֑הּ כֵּ֣ן תַּעֲשֶׂ֔ה לְכֹ֖ל קַרְשֵׁ֥י הַמִּשְׁכָּֽן׃ \Rashi{\textbf{שתי ידות לקרש האחד.}היה חורץ את הקרש מלמטה באמצעו בגבה אמה, ומניח רביע רחבו מכאן, ורביע רחבו מכאן, והן הן הידות, והחריץ חצי רחב הקרש באמצע, ואותן הידות מכניס באדנים שהיו חלולים, והאדנים גבהן אמה ויושבים רצופים ארבעים זה אצל זה, וידות הקרש הנכנסות בחלל האדנים חרוצות משלשה צדיהן – רחב החריץ כעבי שפת האדן – שיכסה הקרש את כל ראש האדן, שאם לא כן, נמצא רוח בין קרש לקרש כעבי שפת שני האדנים, שיפסיקו ביניהם, וזהו שנאמר ויהיו תאמים מלמטה, שיחרץ את צדי הידות כדי שיתחברו הקרשים זה אצל זה: משלבת. עשויות כמין שליבות סלם, ומבדלות זו מזו, ומשפין ראשיהם לכנס בתוך חלל האדן כשליבה הנכנסת בנקב עמודי הסלם: אשה אל אחתה. מכונות זו כנגד זו, שיהיו חריציהם שוים זו כמדת זו, כדי שלא יהיו שתי ידות זו משוכה לצד פנים וזו משוכה לצד חוץ בעבי הקרש שהוא אמה, ותרגום של ידות צירין, לפי שדומות לצירי הדלת הנכנסים בחרי המפתן:}%
% ---------- פרק 26 | פסוק 18 ----------
 \textbf{יח} וְעָשִׂ֥יתָ אֶת־הַקְּרָשִׁ֖ים לַמִּשְׁכָּ֑ן עֶשְׂרִ֣ים קֶ֔רֶשׁ לִפְאַ֖ת נֶ֥גְבָּה תֵימָֽנָה׃ \Rashi{\textbf{לפאת נגבה תימנה.}אין פאה זו לשון מקצוע, אלא כל הרוח קרויה פאה, כתרגומו: לרוח עבר דרומא:}%
% ---------- פרק 26 | פסוק 19 ----------
 \textbf{יט} וְאַרְבָּעִים֙ אַדְנֵי־כֶ֔סֶף תַּעֲשֶׂ֕ה תַּ֖חַת עֶשְׂרִ֣ים הַקָּ֑רֶשׁ שְׁנֵ֨י אֲדָנִ֜ים תַּֽחַת־הַקֶּ֤רֶשׁ הָאֶחָד֙ לִשְׁתֵּ֣י יְדֹתָ֔יו וּשְׁנֵ֧י אֲדָנִ֛ים תַּֽחַת־הַקֶּ֥רֶשׁ הָאֶחָ֖ד לִשְׁתֵּ֥י יְדֹתָֽיו׃ %
% ---------- פרק 26 | פסוק 20 ----------
 \textbf{כ} וּלְצֶ֧לַע הַמִּשְׁכָּ֛ן הַשֵּׁנִ֖ית לִפְאַ֣ת צָפ֑וֹן עֶשְׂרִ֖ים קָֽרֶשׁ׃ %
% ---------- פרק 26 | פסוק 21 ----------
 \textbf{כא} וְאַרְבָּעִ֥ים אַדְנֵיהֶ֖ם כָּ֑סֶף שְׁנֵ֣י אֲדָנִ֗ים תַּ֚חַת הַקֶּ֣רֶשׁ הָֽאֶחָ֔ד וּשְׁנֵ֣י אֲדָנִ֔ים תַּ֖חַת הַקֶּ֥רֶשׁ הָאֶחָֽד׃ %
% ---------- פרק 26 | פסוק 22 ----------
 \textbf{כב} וּֽלְיַרְכְּתֵ֥י הַמִּשְׁכָּ֖ן יָ֑מָּה תַּעֲשֶׂ֖ה שִׁשָּׁ֥ה קְרָשִֽׁים׃ \Rashi{\textbf{ולירכתי.}לשון סוף, כתרגומו ולסיפי, ולפי שהפתח במזרח, קרוי מזרח פנים והמערב אחורים, וזהו סוף, שהפנים הוא הראש: תעשה ששה קרשים. הרי תשע אמות רחב:}%
% ---------- פרק 26 | פסוק 23 ----------
 \textbf{כג} וּשְׁנֵ֤י קְרָשִׁים֙ תַּעֲשֶׂ֔ה לִמְקֻצְעֹ֖ת הַמִּשְׁכָּ֑ן בַּיַּרְכָתָֽיִם׃ \Rashi{\textbf{ושני קרשים תעשה למקצעת.}אחד למקצוע צפונית מערבית, ואחד למערבית דרומית, כל שמנה קרשים בסדר אחד הן, אלא שאלו השנים אינן בחלל המשכן אלא חצי אמה מזו וחצי אמה מזו נראות בחלל, להשלים רחבו לעשר, והאמה מזה והאמה מזה באות כנגד אמת עבי קרשי המשכן הצפון והדרום, כדי שיהא המקצוע מבחוץ שוה:}%
% ---------- פרק 26 | פסוק 24 ----------
 \textbf{כד} וְיִֽהְי֣וּ תֹֽאֲמִם֮ מִלְּמַ֒טָּה֒ וְיַחְדָּ֗ו יִהְי֤וּ תַמִּים֙ עַל־רֹאשׁ֔וֹ אֶל־הַטַּבַּ֖עַת הָאֶחָ֑ת כֵּ֚ן יִהְיֶ֣ה לִשְׁנֵיהֶ֔ם לִשְׁנֵ֥י הַמִּקְצֹעֹ֖ת יִהְיֽוּ׃ \Rashi{\textbf{ויהיו.}כל הקרשים תואמים זה לזה מלמטה, שלא יפסיק עבי שפת שני האדנים ביניהם להרחיקם זו מזו, זהו שפרשתי שיהיו צירי הידות חרוצים מצדיהן, שיהא רחב הקרש בולט לצדיו חוץ לידי הקרש, לכסות את שפת האדן, וכן הקרש שאצלו, ונמצאו תואמים זה לזה; וקרש המקצוע שבסדר המערב חרוץ לרחבו בעביו כנגד חריץ של צד קרש הצפוני והדרומי, כדי שלא יפרידו האדנים ביניהם: ויחדו יהיו תמים. כמו תואמים: על ראשו. של קרש: אל הטבעת האחת. כל קרש וקרש היה חרוץ מלמעלה ברחבו שני חריצין, בשני צדיו, כמו עבי טבעת, ומכניסו בטבעת אחת, נמצא מתאים לקרש שאצלו, אבל אותן טבעות לא ידעתי אם קבועות הן, אם מטלטלות, ובקרש שבמקצוע היתה טבעת בעבי הקרש הדרומי והצפוני, וראש קרש המקצוע שבסדר מערב נכנס לתוכו, נמצאו שני הכתלים מחברים: כן יהיה לשניהם. לשני הקרשים שבמקצוע, לקרש שבסוף צפון ולקרש המערבי, וכן לשני המקצועות:}%
% ---------- פרק 26 | פסוק 25 ----------
 \textbf{כה} וְהָיוּ֙ שְׁמֹנָ֣ה קְרָשִׁ֔ים וְאַדְנֵיהֶ֣ם כֶּ֔סֶף שִׁשָּׁ֥ה עָשָׂ֖ר אֲדָנִ֑ים שְׁנֵ֣י אֲדָנִ֗ים תַּ֚חַת הַקֶּ֣רֶשׁ הָאֶחָ֔ד וּשְׁנֵ֣י אֲדָנִ֔ים תַּ֖חַת הַקֶּ֥רֶשׁ הָאֶחָֽד׃ \Rashi{\textbf{והיו שמנה קרשים.}הן האמורות למעלה תעשה ששה קרשים ושני קרשים תעשה למקצעת, נמצאו שמנה קרשים בסדר מערבי; כך שנויה במשנה מעשה סדר הקרשים במלאכת המשכן; היה עושה את האדנים חלולים וחורץ את הקרש מלמטה רביע מכאן ורביע מכאן, והחריץ חציו באמצע, ועשה לו שתי ידות כמין שני חמוקין – ולי נראה שהגרסא כמין שני חוקין, כמין שני שליבות סלם המבדלות זו מזו, ומשפות לכנס בחלל האדן כשליבה הנכנסת בנקב עמוד הסלם, והוא לשון משלבות, עשויות כמין שליבה – ומכניסן לתוך שני אדנים, שנאמר שני אדנים שני אדנים, וחורץ את הקרש מלמעלה אצבע מכאן ואצבע מכאן, ונותנן לתוך טבעת אחת של זהב כדי שלא יהיו נפרדין זה מזה, שנאמר ויהיו תאמים מלמטה וגו', כך היא המשנה, והפרוש שלה הצעתי למעלה בסדר המקראות:}%
% ---------- פרק 26 | פסוק 26 ----------
 \textbf{כו} וְעָשִׂ֥יתָ בְרִיחִ֖ם עֲצֵ֣י שִׁטִּ֑ים חֲמִשָּׁ֕ה לְקַרְשֵׁ֥י צֶֽלַע־הַמִּשְׁכָּ֖ן הָאֶחָֽד׃ \Rashi{\textbf{בריחם.}כתרגומו עברין, ובלעז אשפ"ריש: חמשה לקרשי צלע המשכן. אלו חמשה ג' הן, אלא שהבריח העליון והתחתון עשוי משתי חתיכות, זה מבריח עד חצי הכתל וזה מבריח עד חצי הכתל, זה נכנס בטבעות מצד זה, וזה נכנס בטבעות מצד זה, עד שמגיעין זה לזה, נמצא שעליון ותחתון שנים שהן ארבעה, אבל האמצעי ארכו כנגד כל הכתל, ומבריח מקצה הכתל ועד קצהו, שנאמר והבריח התיכן וגו' מברח מן הקצה אל הקצה, העליונים והתחתונים היו להן טבעות בקרשים לכנס בתוכן, שתי טבעות לכל קרש משלשים בתוך עשר אמות של גבה הקרש – חלק אחד מן הטבעת העליונה ולמעלה, וחלק אחד מן התחתונה ולמטה, וכל חלק הוא רביע ארך הקרש, ושני חלקים בין טבעת לטבעת כדי שיהיו כל הטבעות מכונים זו כנגד זו – אבל לבריח התיכון אין טבעות, אלא הקרשים נקובים בעבין, והוא נכנס בהם דרך הנקבים שהם מכונין זה מול זה, וזהו שנאמר "בתוך" הקרשים; הבריחים העליונים והתחתונים שבצפון ושבדרום ארך כל אחד חמש עשרה אמה, והתיכון ארכו שלשים אמה, וזהו מן הקצה אל הקצה – מן המזרח ועד המערב – וחמשה בריחים שבמערב ארך העליונים והתחתונים שש אמות, והתיכון ארכו שתים עשרה, כנגד רחב שמנה קרשים, כך היא מפרשת במלאכת המשכן:}%
% ---------- פרק 26 | פסוק 27 ----------
 \textbf{כז} וַחֲמִשָּׁ֣ה בְרִיחִ֔ם לְקַרְשֵׁ֥י צֶֽלַע־הַמִּשְׁכָּ֖ן הַשֵּׁנִ֑ית וַחֲמִשָּׁ֣ה בְרִיחִ֗ם לְקַרְשֵׁי֙ צֶ֣לַע הַמִּשְׁכָּ֔ן לַיַּרְכָתַ֖יִם יָֽמָּה׃ %
% ---------- פרק 26 | פסוק 28 ----------
 \textbf{כח} וְהַבְּרִ֥יחַ הַתִּיכֹ֖ן בְּת֣וֹךְ הַקְּרָשִׁ֑ים מַבְרִ֕חַ מִן־הַקָּצֶ֖ה אֶל־הַקָּצֶֽה׃ %
% ---------- פרק 26 | פסוק 29 ----------
 \textbf{כט} וְֽאֶת־הַקְּרָשִׁ֞ים תְּצַפֶּ֣ה זָהָ֗ב וְאֶת־טַבְּעֹֽתֵיהֶם֙ תַּעֲשֶׂ֣ה זָהָ֔ב בָּתִּ֖ים לַבְּרִיחִ֑ם וְצִפִּיתָ֥ אֶת־הַבְּרִיחִ֖ם זָהָֽב׃ \Rashi{\textbf{בתים לבריחם.}הטבעות שתעשה בהן יהיו בתים לכנס בהן הבריחים: וצפית את הבריחם זהב. לא שהיה הזהב מדבק על הבריחים – שאין עליהם שום צפוי – אלא בקרש היה קובע כמין שתי פיפיות של זהב, כמין שני סדקי קנה חלול, וקובען אצל הטבעת לכאן ולכאן, ארכן ממלא את רחב הקרש מן הטבעת לכאן וממנה לכאן, והבריח נכנס לתוכו וממנו לטבעת ומן הטבעת לפה השני, נמצאו הבריחים מצפים זהב כשהן תחובין בקרשים; והבריחים הללו מבחוץ היו בולטות, הטבעות והפיפיות לא היו נראות בתוך המשכן, אלא כל הכתל חלק מבפנים:}%
% ---------- פרק 26 | פסוק 30 ----------
 \textbf{ל} וַהֲקֵמֹתָ֖ אֶת־הַמִּשְׁכָּ֑ן כְּמִ֨שְׁפָּט֔וֹ אֲשֶׁ֥ר הׇרְאֵ֖יתָ בָּהָֽר׃ \{ס\} \Rashi{\textbf{והקמת את המשכן.}לאחר שיגמר הקימהו: הראית בהר. קדם לכן, שאני עתיד ללמדך ולהראותך סדר הקמתו:}%
% ---------- פרק 26 | פסוק 31 ----------
 \textbf{לא} וְעָשִׂ֣יתָ פָרֹ֗כֶת תְּכֵ֧לֶת וְאַרְגָּמָ֛ן וְתוֹלַ֥עַת שָׁנִ֖י וְשֵׁ֣שׁ מׇשְׁזָ֑ר מַעֲשֵׂ֥ה חֹשֵׁ֛ב יַעֲשֶׂ֥ה אֹתָ֖הּ כְּרֻבִֽים׃ \Rashi{\textbf{פרכת.}לשון מחצה הוא, ובלשון חכמים פרגוד דבר המבדיל בין המלך ובין העם: תכלת וארגמן. כל מין ומין היה כפול בכל חוט וחוט ששה חוטין: מעשה חשב. כבר פרשתי שזו היא אריגה של שתי קירות, והציורין שמשני עבריה אינן דומין זה לזה: כרבים. ציורין של בריות יעשה בה:}%
% ---------- פרק 26 | פסוק 32 ----------
 \textbf{לב} וְנָתַתָּ֣ה אֹתָ֗הּ עַל־אַרְבָּעָה֙ עַמּוּדֵ֣י שִׁטִּ֔ים מְצֻפִּ֣ים זָהָ֔ב וָוֵיהֶ֖ם זָהָ֑ב עַל־אַרְבָּעָ֖ה אַדְנֵי־כָֽסֶף׃ \Rashi{\textbf{ארבעה עמודי שטים.}תקועים בתוך ארבעה אדנים, ואונקליות קבועין בהן, עקמים למעלה, להושיב עליהן כלונס שראש הפרכת כרוך בה, והאונקליות הן הווין – שהרי כמין ווין הן עשויים – והפרכת ארכה עשר אמות לרחבו של משכן ורחבה עשר אמות כגבהן של קרשים, פרוסה בשלישית של משכן, שיהא הימנה ולפנים עשר אמות והימנה ולחוץ עשרים אמה, נמצא בית קדשי הקדשים עשר על עשר, שנ' ונתתה את הפרכת תחת הקרסים – המחברים את שתי חוברות של יריעות המשכן; רחב החוברת עשרים אמה, וכשפרשה על גג המשכן מן הפתח למערב, כלתה בשני שלישי המשכן, והחוברת השנית כסתה שלישו של משכן, והמותר תלוי לאחוריו לכסות את הקרשים:}%
% ---------- פרק 26 | פסוק 33 ----------
 \textbf{לג} וְנָתַתָּ֣ה אֶת־הַפָּרֹ֘כֶת֮ תַּ֣חַת הַקְּרָסִים֒ וְהֵבֵאתָ֥ שָׁ֙מָּה֙ מִבֵּ֣ית לַפָּרֹ֔כֶת אֵ֖ת אֲר֣וֹן הָעֵד֑וּת וְהִבְדִּילָ֤ה הַפָּרֹ֙כֶת֙ לָכֶ֔ם בֵּ֣ין הַקֹּ֔דֶשׁ וּבֵ֖ין קֹ֥דֶשׁ הַקֳּדָשִֽׁים׃ %
% ---------- פרק 26 | פסוק 34 ----------
 \textbf{לד} וְנָתַתָּ֙ אֶת־הַכַּפֹּ֔רֶת עַ֖ל אֲר֣וֹן הָעֵדֻ֑ת בְּקֹ֖דֶשׁ הַקֳּדָשִֽׁים׃ %
% ---------- פרק 26 | פסוק 35 ----------
 \textbf{לה} וְשַׂמְתָּ֤ אֶת־הַשֻּׁלְחָן֙ מִח֣וּץ לַפָּרֹ֔כֶת וְאֶת־הַמְּנֹרָה֙ נֹ֣כַח הַשֻּׁלְחָ֔ן עַ֛ל צֶ֥לַע הַמִּשְׁכָּ֖ן תֵּימָ֑נָה וְהַ֨שֻּׁלְחָ֔ן תִּתֵּ֖ן עַל־צֶ֥לַע צָפֽוֹן׃ \Rashi{\textbf{ושמת את השלחן.}שלחן בצפון, משוך מן הכתל הצפוני שתי אמות ומחצה, ומנורה בדרום, משוכה מן הכתל הדרומי שתי אמות ומחצה, ומזבח הזהב נתון כנגד אויר שבין שלחן למנורה משוך קמעא כלפי המזרח, וכלם נתונים מן חצי המשכן ולפנים; כיצד? ארך המשכן מן הפתח לפרכת עשרים אמה, המזבח והשלחן והמנורה משוכים מן הפתח לצד מערב עשר אמות (יומא ל"ג):}%
% ---------- פרק 26 | פסוק 36 ----------
 \textbf{לו} וְעָשִׂ֤יתָ מָסָךְ֙ לְפֶ֣תַח הָאֹ֔הֶל תְּכֵ֧לֶת וְאַרְגָּמָ֛ן וְתוֹלַ֥עַת שָׁנִ֖י וְשֵׁ֣שׁ מׇשְׁזָ֑ר מַעֲשֵׂ֖ה רֹקֵֽם׃ \Rashi{\textbf{ועשית מסך.}וילון, הוא מסך כנגד הפתח, כמו שכת בעדו (איוב א'), לשון מגן: מעשה רקם. הצורות עשויות בו מעשה מחט, כפרצוף של עבר זה כך פרצוף של עבר זה: רקם. שם האמן ולא שם האמנות, ותרגומו עובד ציר ולא עובד ציור. מדת המסך כמדת הפרכת – עשר אמות על עשר אמות:}%
% ---------- פרק 26 | פסוק 37 ----------
 \textbf{לז} וְעָשִׂ֣יתָ לַמָּסָ֗ךְ חֲמִשָּׁה֙ עַמּוּדֵ֣י שִׁטִּ֔ים וְצִפִּיתָ֤ אֹתָם֙ זָהָ֔ב וָוֵיהֶ֖ם זָהָ֑ב וְיָצַקְתָּ֣ לָהֶ֔ם חֲמִשָּׁ֖ה אַדְנֵ֥י נְחֹֽשֶׁת׃ \{ס\} %
% ---------- פרק 27 | פסוק 1 ----------
 \textbf{א} וְעָשִׂ֥יתָ אֶת־הַמִּזְבֵּ֖חַ עֲצֵ֣י שִׁטִּ֑ים חָמֵשׁ֩ אַמּ֨וֹת אֹ֜רֶךְ וְחָמֵ֧שׁ אַמּ֣וֹת רֹ֗חַב רָב֤וּעַ יִהְיֶה֙ הַמִּזְבֵּ֔חַ וְשָׁלֹ֥שׁ אַמּ֖וֹת קֹמָתֽוֹ׃ \Rashi{\textbf{ועשית את המזבח וגו' ושלש אמות קמתו.}דברים ככתבן, דברי ר' יהודה, רבי יוסי אומר נאמר כאן רבוע ונאמר בפנימי רבוע, מה להלן גבהו פי שנים כארכו, אף כאן גבהו פי שנים כארכו, ומה אני מקים ושלש אמות קמתו? משפת סובב ולמעלה (זבח' נ"ט):}%
% ---------- פרק 27 | פסוק 2 ----------
 \textbf{ב} וְעָשִׂ֣יתָ קַרְנֹתָ֗יו עַ֚ל אַרְבַּ֣ע פִּנֹּתָ֔יו מִמֶּ֖נּוּ תִּהְיֶ֣יןָ קַרְנֹתָ֑יו וְצִפִּיתָ֥ אֹת֖וֹ נְחֹֽשֶׁת׃ \Rashi{\textbf{ממנו תהיין קרנתיו.}שלא יעשם לבדם ויחברם בו: וצפית אתו נחשת. לכפר על עזות מצח, שנ' ומצחך נחושה (ישעיהו מ"ח):}%
% ---------- פרק 27 | פסוק 3 ----------
 \textbf{ג} וְעָשִׂ֤יתָ סִּֽירֹתָיו֙ לְדַשְּׁנ֔וֹ וְיָעָיו֙ וּמִזְרְקֹתָ֔יו וּמִזְלְגֹתָ֖יו וּמַחְתֹּתָ֑יו לְכׇל־כֵּלָ֖יו תַּעֲשֶׂ֥ה נְחֹֽשֶׁת׃ \Rashi{\textbf{סירתיו.}כמין יורות: לדשנו. להסיר דשנו לתוכם, והוא שתרגם אנקלוס למספי קטמה – לספות הדשן לתוכם; כי יש מלות בלשון עברית מלה אחת מתחלפת בפתרון לשמש בנין וסתירה, כמו ותשרש שרשיה (תהילים פ'), אויל משריש (איוב ה'), וחלופו ובכל תבואתי תשרש (איוב ל״א:י״ב); וכמוהו בסעפיה פריה (ישעיהו י״ז:ו׳), וחלופו מסעף פארה (ישעיהו י׳:ל״ג) – מפשח סעפיה; וכמוהו וזה האחרון עצמו (ירמיהו נ') – שבר עצמיו; וכמוהו ויסקלהו באבנים (מלכים א כ״א:י״ג), וחלופו סקלו מאבן (ישעיהו ס״ב:י׳) – הסירו אבניה, וכן ויעזקהו ויסקלהו (שם ה'), אף כאן לדשנו – להסיר דשנו, ובלעז אדשצנדר"יר: ויעיו. כתרגומו, מגרפות שנוטלים בהם הדשן, והן כמין כסוי הקדרה של מתכת דק ולו בית יד, ובלעז ווד"יל: ומזרקתיו. לקבל בהם דם הזבחים: ומזלגתיו. כמין אנקליות כפופים, ומכה בהם בבשר ונתחבים בו ומהפך בהן על גחלי המערכה שיהא ממהר שרפתן, ובלעז קרוצי"נש, ובלשון חכמים צנוריות (יומא י"ב): ומחתתיו. בית קבול יש להם לטל בהן גחלים מן המזבח, לשאתם על מזבח הפנימי לקטרת, ועל שם חתיתן קרויים מחתות, כמו לחתות אש מיקוד (ישעיהו ל'), לשון שאיבת אש ממקומה, וכן היחתה איש אש בחיקו (משלי ו'): לכל כליו. כמו כל כליו:}%
% ---------- פרק 27 | פסוק 4 ----------
 \textbf{ד} וְעָשִׂ֤יתָ לּוֹ֙ מִכְבָּ֔ר מַעֲשֵׂ֖ה רֶ֣שֶׁת נְחֹ֑שֶׁת וְעָשִׂ֣יתָ עַל־הָרֶ֗שֶׁת אַרְבַּע֙ טַבְּעֹ֣ת נְחֹ֔שֶׁת עַ֖ל אַרְבַּ֥ע קְצוֹתָֽיו׃ \Rashi{\textbf{מכבר.}לשון כברה, שקורין קרי"בלא בלעז, כמין לבוש עשוי לו למזבח, עשוי חרין חרין, כמין רשת; ומקרא זה מסרס וכה פתרונו: ועשית לו מכבר נחשת מעשה רשת:}%
% ---------- פרק 27 | פסוק 5 ----------
 \textbf{ה} וְנָתַתָּ֣ה אֹתָ֗הּ תַּ֛חַת כַּרְכֹּ֥ב הַמִּזְבֵּ֖חַ מִלְּמָ֑טָּה וְהָיְתָ֣ה הָרֶ֔שֶׁת עַ֖ד חֲצִ֥י הַמִּזְבֵּֽחַ׃ \Rashi{\textbf{כרכב המזבח.}סובב; כל דבר המקיף סביב בעגול קרוי כרכב, כמו ששנינו בהכל שוחטין אלו הן גלמי כלי עץ כל שעתיד לשוף ולכרכב, והוא כמו שעושים חריצין עגלין בקרשי דפני התבות וספסלי העץ, אף למזבח עשה חריץ סביבו, והיה רחבו אמה בדפנו לנוי והוא לסוף שלש אמות של גבהו כדברי האומר גבהו פי שנים כארכו, הא מה אני מקים ושלש אמות קומתו? משפת סובב ולמעלה, אבל סובב להלוך הכהנים לא היה למזבח הנחשת, אלא על ראשו לפנים מקרנותיו, וכן שנינו בזבחים: אי זהו כרכוב? בין קרן לקרן; והיה רחב אמה, ולפנים מהן אמה של הלוך רגלי הכהנים, שתי אמות הללו קרויים כרכוב; ודקדקנו שם והכתיב תחת כרכב המזבח מלמטה? למדנו שהכרכוב בדפנו הוא ולבוש המכבר תחתיו? ותרץ המתרץ תרי הוו, חד לנוי וחד לכהנים דלא ישתרגו, זה שבדפן לנוי היה, ומתחתיו הלבישו המכבר והגיע רחבו עד חצי המזבח, והוא היה סימן לחצי גבהו, להבדיל בין דמים העליונים לדמים התחתונים; וכנגדו עשו למזבח בית עולמים דגמת חוט הסקרא באמצעו, וכבש שהיו עולין בו אע"פ שלא פרשו בענין זה, כבר שמענו בפרשת מזבח אדמה תעשה לי ולא תעלה במעלת – לא תעשה לו מעלות בכבש שלו, אלא כבש חלק, למדנו שהיה לו כבש; כך שנינו במכילתא. ומזבח אדמה הוא מזבח הנחשת, שהיו ממלאין חללו אדמה במקום חניתן, והכבש היה בדרום המזבח, מבדל מן המזבח מלא חוט השערה, ורגליו מגיעין עד אמה סמוך לקלעי החצר שבדרום, כדברי האומר עשר אמות קומתו, ולדברי האומר דברים ככתבן – שלש אמות קומתו – לא היה ארך הכבש אלא עשר אמות, כך מצאתי במשנת ארבעים ותשע מדות, וזה שהוא מבדל מן המזבח מלא החוט, במסכת זבחים (דף ס"ב) למדוהו מן המקרא:}%
% ---------- פרק 27 | פסוק 6 ----------
 \textbf{ו} וְעָשִׂ֤יתָ בַדִּים֙ לַמִּזְבֵּ֔חַ בַּדֵּ֖י עֲצֵ֣י שִׁטִּ֑ים וְצִפִּיתָ֥ אֹתָ֖ם נְחֹֽשֶׁת׃ %
% ---------- פרק 27 | פסוק 7 ----------
 \textbf{ז} וְהוּבָ֥א אֶת־בַּדָּ֖יו בַּטַּבָּעֹ֑ת וְהָי֣וּ הַבַּדִּ֗ים עַל־שְׁתֵּ֛י צַלְעֹ֥ת הַמִּזְבֵּ֖חַ בִּשְׂאֵ֥ת אֹתֽוֹ׃ \Rashi{\textbf{בטבעת.}בארבע טבעות שנעשו למכבר:}%
% ---------- פרק 27 | פסוק 8 ----------
 \textbf{ח} נְב֥וּב לֻחֹ֖ת תַּעֲשֶׂ֣ה אֹת֑וֹ כַּאֲשֶׁ֨ר הֶרְאָ֥ה אֹתְךָ֛ בָּהָ֖ר כֵּ֥ן יַעֲשֽׂוּ׃ \{ס\} \Rashi{\textbf{נבוב לחת.}כתרגומו חליל לוחין, לחת עצי שטים מכל צד והחלל באמצע, ולא יהא כלו עץ אחד שיהא עביו חמש אמות על חמש אמות כמין סדן:}%
% ---------- פרק 27 | פסוק 9 ----------
 \textbf{ט} וְעָשִׂ֕יתָ אֵ֖ת חֲצַ֣ר הַמִּשְׁכָּ֑ן לִפְאַ֣ת נֶֽגֶב־תֵּ֠ימָ֠נָה קְלָעִ֨ים לֶחָצֵ֜ר שֵׁ֣שׁ מׇשְׁזָ֗ר מֵאָ֤ה בָֽאַמָּה֙ אֹ֔רֶךְ לַפֵּאָ֖ה הָאֶחָֽת׃ \Rashi{\textbf{קלעים.}עשויין כמין קלעי ספינה, נקבים נקבים, מעשה קליעה ולא מעשה אורג, ותרגום סרדים כתרגומו של מכבר המתרגם סרדה, לפי שהן מנקבין ככברה: לפאה האחת. כל הרוח קרוי פאה:}%
% ---------- פרק 27 | פסוק 10 ----------
 \textbf{י} וְעַמֻּדָ֣יו עֶשְׂרִ֔ים וְאַדְנֵיהֶ֥ם עֶשְׂרִ֖ים נְחֹ֑שֶׁת וָוֵ֧י הָעַמֻּדִ֛ים וַחֲשֻׁקֵיהֶ֖ם כָּֽסֶף׃ \Rashi{\textbf{ועמדיו עשרים.}חמש אמות בין עמוד לעמוד: ואדניהם. של העמודים נחשת, האדנים יושבים על הארץ והעמודים תקועין לתוכן; היה עושה כמין קנדסין שקורין פל"ש בלעז, ארכן ששה טפחים ורחבן שלשה, וטבעת נחשת קבועה בו באמצעו וכורך שפת הקלע סביביו במיתרים כנגד כל עמוד ועמוד, ותולה הקנדס דרך טבעתו באנקליות שבעמוד העשוי כמין וי"ו – ראשו זקוף למעלה וראשו אחד תקוע בעמוד – כאותן שעושין להציב דלתות, שקורין גונ"ש בלעז, ורחב הקלע תלוי מלמטה והיא קומת מחצות החצר: ווי העמדים. הם האנקליות: וחשקיהם. מקפין היו העמודים בחוטי כסף סביב, ואיני יודע אם על פני כלן, אם בראשם, ואם באמצעם, אך יודע אני שחשוק לשון חגורה, שכן מצינו בפילגש בגבעה ועמו צמד חמורים חבושים (שופטים י"ט), תרגומו חשוקים:}%
% ---------- פרק 27 | פסוק 11 ----------
 \textbf{יא} וְכֵ֨ן לִפְאַ֤ת צָפוֹן֙ בָּאֹ֔רֶךְ קְלָעִ֖ים מֵ֣אָה אֹ֑רֶךְ וְעַמֻּדָ֣ו עֶשְׂרִ֗ים וְאַדְנֵיהֶ֤ם עֶשְׂרִים֙ נְחֹ֔שֶׁת וָוֵ֧י הָֽעַמֻּדִ֛ים וַחֲשֻׁקֵיהֶ֖ם כָּֽסֶף׃ %
% ---------- פרק 27 | פסוק 12 ----------
 \textbf{יב} וְרֹ֤חַב הֶֽחָצֵר֙ לִפְאַת־יָ֔ם קְלָעִ֖ים חֲמִשִּׁ֣ים אַמָּ֑ה עַמֻּדֵיהֶ֣ם עֲשָׂרָ֔ה וְאַדְנֵיהֶ֖ם עֲשָׂרָֽה׃ %
% ---------- פרק 27 | פסוק 13 ----------
 \textbf{יג} וְרֹ֣חַב הֶֽחָצֵ֗ר לִפְאַ֛ת קֵ֥דְמָה מִזְרָ֖חָה חֲמִשִּׁ֥ים אַמָּֽה׃ \Rashi{\textbf{לפאת קדמה מזרחה.}פני המזרח; קדם לשון פנים, אחור לשון אחורים, לפיכך המזרח קרוי קדם שהוא פנים, ומערב קרוי אחור, כמה דאת אמר הים האחרון (דברים י"א:כ"ד, דברים ל"ד:ב'), ימא מערבא: חמשים אמה. אותן נ' אמה לא היו סתומים כלם בקלעים, לפי ששם הפתח, אלא ט"ו אמה קלעים לכתף הפתח מכאן, וכן לכתף השנית, נשאר רחב חלל הפתח בינתים כ' אמה, וזהו שנ' ולשער החצר מסך כ' אמה – וילון למסך כנגד הפתח כ' אמה ארך, כרחב הפתח:}%
% ---------- פרק 27 | פסוק 14 ----------
 \textbf{יד} וַחֲמֵ֨שׁ עֶשְׂרֵ֥ה אַמָּ֛ה קְלָעִ֖ים לַכָּתֵ֑ף עַמֻּדֵיהֶ֣ם שְׁלֹשָׁ֔ה וְאַדְנֵיהֶ֖ם שְׁלֹשָֽׁה׃ \Rashi{\textbf{עמדיהם שלשה.}חמש אמות בין עמוד לעמוד. בין עמוד שבראש הדרום העומד במקצוע דרומית מזרחית עד עמוד שהוא מן השלשה שבמזרח חמש אמות, וממנו לשני חמש אמות, ומן השני לשלישי חמש אמות, וכן לכתף השנית, וארבעה עמודים למסך, הרי עשרה עמודים למזרח כנגד י' למערב:}%
% ---------- פרק 27 | פסוק 15 ----------
 \textbf{טו} וְלַכָּתֵף֙ הַשֵּׁנִ֔ית חֲמֵ֥שׁ עֶשְׂרֵ֖ה קְלָעִ֑ים עַמֻּדֵיהֶ֣ם שְׁלֹשָׁ֔ה וְאַדְנֵיהֶ֖ם שְׁלֹשָֽׁה׃ %
% ---------- פרק 27 | פסוק 16 ----------
 \textbf{טז} וּלְשַׁ֨עַר הֶֽחָצֵ֜ר מָסָ֣ךְ ׀ עֶשְׂרִ֣ים אַמָּ֗ה תְּכֵ֨לֶת וְאַרְגָּמָ֜ן וְתוֹלַ֧עַת שָׁנִ֛י וְשֵׁ֥שׁ מׇשְׁזָ֖ר מַעֲשֵׂ֣ה רֹקֵ֑ם עַמֻּֽדֵיהֶם֙ אַרְבָּעָ֔ה וְאַדְנֵיהֶ֖ם אַרְבָּעָֽה׃ %
% ---------- פרק 27 | פסוק 17 ----------
 \textbf{יז} כׇּל־עַמּוּדֵ֨י הֶֽחָצֵ֤ר סָבִיב֙ מְחֻשָּׁקִ֣ים כֶּ֔סֶף וָוֵיהֶ֖ם כָּ֑סֶף וְאַדְנֵיהֶ֖ם נְחֹֽשֶׁת׃ \Rashi{\textbf{כל עמודי החצר סביב וגו'.}לפי שלא פרש ווין וחשוקים ואדני נחשת אלא לצפון ולדרום, אבל למזרח ולמערב לא נאמר ווין וחשוקים ואדני נחשת, לכך בא ולמד כאן:}%
% ---------- פרק 27 | פסוק 18 ----------
 \textbf{יח} אֹ֣רֶךְ הֶֽחָצֵר֩ מֵאָ֨ה בָֽאַמָּ֜ה וְרֹ֣חַב ׀ חֲמִשִּׁ֣ים בַּחֲמִשִּׁ֗ים וְקֹמָ֛ה חָמֵ֥שׁ אַמּ֖וֹת שֵׁ֣שׁ מׇשְׁזָ֑ר וְאַדְנֵיהֶ֖ם נְחֹֽשֶׁת׃ \Rashi{\textbf{ארך החצר.}הצפון והדרום שמן המזרח למערב מאה באמה: ורחב חמשים בחמשים. חצר שבמזרח היתה מרבעת, חמשים על חמשים, שהמשכן ארכו שלשים ורחבו עשר, העמיד מזרח פתחו בשפת נ' החיצונים של ארך החצר, נמצא כלו בחמשים הפנימיים וכלה ארכו לסוף שלשים, נמצאו עשרים אמה רוח לאחוריו בין הקלעים שבמערב ליריעות של אחורי המשכן, ורחב המשכן עשר אמות, באמצע רחב החצר, נמצאו לו עשרים אמה רוח לצפון ולדרום מן קלעי החצר ליריעות המשכן, וכן למערב, וחמשים על חמשים חצר לפניו: וקמה חמש אמות. גבה מחצות החצר והוא רחב הקלעים: ואדניהם נחשת. להביא אדני המסך, שלא תאמר לא נאמרו אדני נחשת אלא לעמודי הקלעים אבל אדני המסך של מין אחר, כן נראה בעיני שלכך חזר ושנאן:}%
% ---------- פרק 27 | פסוק 19 ----------
 \textbf{יט} לְכֹל֙ כְּלֵ֣י הַמִּשְׁכָּ֔ן בְּכֹ֖ל עֲבֹדָת֑וֹ וְכׇל־יְתֵדֹתָ֛יו וְכׇל־יִתְדֹ֥ת הֶחָצֵ֖ר נְחֹֽשֶׁת׃ \{ס\} \Rashi{\textbf{לכל כלי המשכן.}שהיו צריכין להקמתו ולהורדתו, כגון מקבות לתקע יתדות ועמודים: יתדת. כמין נגרי נחשת עשויין ליריעות האהל ולקלעי החצר, קשורים במיתרים סביב סביב בשפוליהן, כדי שלא תהא הרוח מגביהתן, ואיני יודע אם תחובין בארץ או קשורין ותלויין וכבדן מכביד שפולי היריעות שלא ינועו ברוח, ואומר אני ששמן מוכיח עליהם שהם תקועים בארץ, לכך נקראו יתדות, ומקרא זה מסיעני אהל בל יצען בל יסע יתדתיו לנצח (ישעיהו ל"ג):}%
\addparasha{חומש שמות | פרשת תצוה}
% ---------- פרק 27 | פסוק 20 ----------
 \textbf{כ} וְאַתָּ֞ה תְּצַוֶּ֣ה ׀ אֶת־בְּנֵ֣י יִשְׂרָאֵ֗ל וְיִקְח֨וּ אֵלֶ֜יךָ שֶׁ֣מֶן זַ֥יִת זָ֛ךְ כָּתִ֖ית לַמָּא֑וֹר לְהַעֲלֹ֥ת נֵ֖ר תָּמִֽיד׃ \Rashi{\textbf{ואתה תצוה.}זך. בלי שמרים, כמו ששנינו במנחות מגרגרו בראש הזית וכו': כתית. הזיתים כותש במכתשת ואינו טוחנן ברחים, כדי שלא יהיה בו שמרים, ואחר שהוציא טפה ראשונה מכניסן לרחים וטוחנן; והשמן השני פסול למנורה וכשר למנחות שנאמר כתית למאור – ולא כתית למנחות (מנחות פ"ו): להעלות נר תמיד. מדליק עד שתהא שלהבת עולה מאליה (שבת כ"א): תמיד. כל לילה ולילה קרוי תמיד, כמו שאתה אומר עלת תמיד ואינה אלא מיום ליום; וכן במנחת חבתין נאמר תמיד ואינה אלא מחציתה בבקר ומחציתה בערב, אבל תמיד האמור בלחם הפנים משבת לשבת הוא:}%
% ---------- פרק 27 | פסוק 21 ----------
 \textbf{כא} בְּאֹ֣הֶל מוֹעֵד֩ מִח֨וּץ לַפָּרֹ֜כֶת אֲשֶׁ֣ר עַל־הָעֵדֻ֗ת יַעֲרֹךְ֩ אֹת֨וֹ אַהֲרֹ֧ן וּבָנָ֛יו מֵעֶ֥רֶב עַד־בֹּ֖קֶר לִפְנֵ֣י יְהֹוָ֑ה חֻקַּ֤ת עוֹלָם֙ לְדֹ֣רֹתָ֔ם מֵאֵ֖ת בְּנֵ֥י יִשְׂרָאֵֽל׃ \{ס\} \Rashi{\textbf{מערב עד בקר.}תן לה מדתה שתהא דולקת מערב ועד בקר, ושערו חכמים חצי לג ללילי טבת הארכין וכן לכל הלילות, ואם יותר אין בכך כלום (מנחות פ"ט):}%
% ---------- פרק 28 | פסוק 1 ----------
 \textbf{א} וְאַתָּ֡ה הַקְרֵ֣ב אֵלֶ֩יךָ֩ אֶת־אַהֲרֹ֨ן אָחִ֜יךָ וְאֶת־בָּנָ֣יו אִתּ֗וֹ מִתּ֛וֹךְ בְּנֵ֥י יִשְׂרָאֵ֖ל לְכַהֲנוֹ־לִ֑י אַהֲרֹ֕ן נָדָ֧ב וַאֲבִיה֛וּא אֶלְעָזָ֥ר וְאִיתָמָ֖ר בְּנֵ֥י אַהֲרֹֽן׃ \Rashi{\textbf{ואתה הקרב אליך.}לאחר שתגמר מלאכת המשכן:}%
% ---------- פרק 28 | פסוק 2 ----------
 \textbf{ב} וְעָשִׂ֥יתָ בִגְדֵי־קֹ֖דֶשׁ לְאַהֲרֹ֣ן אָחִ֑יךָ לְכָב֖וֹד וּלְתִפְאָֽרֶת׃ %
% ---------- פרק 28 | פסוק 3 ----------
 \textbf{ג} וְאַתָּ֗ה תְּדַבֵּר֙ אֶל־כׇּל־חַכְמֵי־לֵ֔ב אֲשֶׁ֥ר מִלֵּאתִ֖יו ר֣וּחַ חׇכְמָ֑ה וְעָשׂ֞וּ אֶת־בִּגְדֵ֧י אַהֲרֹ֛ן לְקַדְּשׁ֖וֹ לְכַהֲנוֹ־לִֽי׃ \Rashi{\textbf{לקדשו לכהנו לי.}לקדשו להכניסו בכהנה ע"י הבגדים שיהא כהן לי, ולשון כהנה שרות הוא, שנטרי"אה בלעז:}%
% ---------- פרק 28 | פסוק 4 ----------
 \textbf{ד} וְאֵ֨לֶּה הַבְּגָדִ֜ים אֲשֶׁ֣ר יַעֲשׂ֗וּ חֹ֤שֶׁן וְאֵפוֹד֙ וּמְעִ֔יל וּכְתֹ֥נֶת תַּשְׁבֵּ֖ץ מִצְנֶ֣פֶת וְאַבְנֵ֑ט וְעָשׂ֨וּ בִגְדֵי־קֹ֜דֶשׁ לְאַהֲרֹ֥ן אָחִ֛יךָ וּלְבָנָ֖יו לְכַהֲנוֹ־לִֽי׃ \Rashi{\textbf{חשן.}תכשיט כנגד הלב: ואפוד. לא שמעתי ולא מצאתי בבריתא פרוש תבניתו, ולבי אומר לי שהוא חגורה לו מאחוריו, רחבו כרחב גב איש, כמין סינר שקורין פורצי"נט בלעז, שחוגרות השרות כשרוכבות על הסוסים, כך מעשהו מלמטה, שנאמר ודוד חגור אפוד בד, למדנו שהאפוד חגורה היא; ואי אפשר לומר שאין בו אלא חגורה לבדה, שהרי נאמר ויתן עליו את האפד ואח"כ ויחגר אותו בחשב האפד ותרגם אנקלוס בהמין אפודא, למדנו שהחשב הוא החגור והאפוד שם תכשיט לבדו; ואי אפשר לומר שעל שם שתי הכתפות שבו הוא קרוי אפוד, שהרי נאמר שתי כתפות האפוד, למדנו שהאפוד שם לבד והכתפות שם לבד והחשב שם לבד, לכך אני אומר שעל שם הסינר של מטה קרוי אפוד – על שם שאופדו ומקשטו בו – כמו שנ' ויאפד לו בו, והחשב הוא חגור שלמעלה הימנו והכתפות קבועות בו. ועוד אומר לי לבי שיש ראיה שהוא מין לבוש, שתרגם יונתן ודוד חגור אפוד בד (שמואל ב ו') – כרדוט דבוץ, ותרגם כמו כן מעילים כרדוטין, במעשה תמר אחות אבשלום, כי כן תלבשנה בנות המלך הבתולות מעילים (שם י"ג): מעיל. הוא כמין חלוק, וכן הכתנת, אלא שהכתנת סמוך לבשרו ומעיל קרוי חלוק העליון: תשבץ. עשויין משבצות לנוי, והמשבצות הם כמין גמות העשויות בתכשיטי זהב למושב קביעת אבנים טובות ומרגליות, כמו שנאמר באבני האפוד מסבת משבצות זהב, ובלעז קוראין אותו קשטו"נש: מצנפת. כמין כפת כובע, שקורין קופ"יא בלעז, שהרי במקום אחר קורא להם מגבעות ומתרגמינן כובעין: ואבנט. היא חגורה על הכתנת והאפוד חגורה על המעיל, כמו שמצינו בסדר לבישתן ויתן עליו את הכתנת ויחגר אתו באבנט וילבש אתו את המעיל ויתן עליו את האפד: בגדי קדש. מתרומה המקדשת לשמי יעשה אותם:}%
% ---------- פרק 28 | פסוק 5 ----------
 \textbf{ה} וְהֵם֙ יִקְח֣וּ אֶת־הַזָּהָ֔ב וְאֶת־הַתְּכֵ֖לֶת וְאֶת־הָֽאַרְגָּמָ֑ן וְאֶת־תּוֹלַ֥עַת הַשָּׁנִ֖י וְאֶת־הַשֵּֽׁשׁ׃ \{פ\} \Rashi{\textbf{והם יקחו.}אותם חכמי לב שיעשו הבגדים יקבלו מן המתנדבים את הזהב ואת התכלת לעשות מהן את הבגדים:}%
% ---------- פרק 28 | פסוק 6 ----------
 \textbf{ו} וְעָשׂ֖וּ אֶת־הָאֵפֹ֑ד זָ֠הָ֠ב תְּכֵ֨לֶת וְאַרְגָּמָ֜ן תּוֹלַ֧עַת שָׁנִ֛י וְשֵׁ֥שׁ מׇשְׁזָ֖ר מַעֲשֵׂ֥ה חֹשֵֽׁב׃ \Rashi{\textbf{ועשו את האפד.}אם באתי לפרש מעשה האפוד והחשן על סדר המקראות, הרי פרושן פרקים, וישגה הקורא בצרופן, לכך אני כותב מעשיהם כמות שהוא, למען ירוץ הקורא בו, ואחר כך אפרש על סדר המקראות. האפוד עשוי כמין סינר של נשים רוכבות סוסים, וחוגר אותו מאחוריו כנגד לבו למטה מאציליו, רחבו כמדת רחב גבו של אדם ויותר, ומגיע עד עקביו, והחשב מחבר בראשו על פני רחבו מעשה אורג, ומאריך לכאן ולכאן כדי להקיף ולחגר בו, והכתפות מחברות בחשב, אחד לימין ואחד לשמאל, מאחורי הכהן, לשני קצות רחבו של סינר, וכשזוקפן עומדות לו על שתי כתפיו, והן כמין שתי רצועות עשויות ממין האפוד, ארכות כדי שעור לזקפן אצל צוארו מכאן ומכאן ונקפלות לפניו למטה מכתפיו מעט, ואבני השהם קבועות בהם, אחת על כתף ימין ואחת על כתף שמאל, והמשבצות נתונות בראשיהם לפני כתפיו, ושתי עבותות הזהב תחובות בשתי טבעות שבחשן, בשני קצות רחבו העליון, אחת לימין ואחת לשמאל, ושני ראשי השרשרת תקועין במשבצת לימין וכן שני ראשי השרשרת השמאלית תקועין במשבצת שבכתף שמאל, נמצא החשן תלוי במשבצות האפוד על לבו מלפניו, ועוד שתי טבעות בשני קצות החשן בתחתיתו, וכנגדם שתי טבעות בשתי כתפות האפוד מלמטה, בראשו התחתון המחבר בחשב; טבעות החשן אל מול טבעות האפוד שוכבים זה על זה, ומרכסן בפתיל תכלת תחוב בטבעות האפוד והחשן, שיהא תחתית החשן דבוק לחשב האפוד ולא יהא נד ונבדל, הולך וחוזר: זהב תכלת וארגמן תולעת שני ושש משזר. חמשת מינים הללו שזורין בכל חוט וחוט; היו מרדדין את הזהב כמין טסין דקין וקוצצין פתילים מהם, וטווין אותו חוט של זהב עם ששה חוטין של תכלת, וחוט של זהב עם ששה חוטין של ארגמן, וכן בתולעת שני וכן בשש – שכל המינין חוטן כפול ששה וחוט של זהב עם כל אחד ואחד – ואחר כך שוזר את כלם כאחד, נמצא חוטן כפול כ"ח, וכן מפרש במסכת יומא, ולמד מן המקרא הזה וירקעו את פחי הזהב וקצץ פתילם לעשות – את פתילי הזהב – בתוך התכלת ובתוך הארגמן וגו', למדנו שחוט של זהב שזור עם כל מין ומין: מעשה חשב. כבר פרשתי שהוא אריגת שתי קירות, שאין צורות שני עבריה דומות זו לזו:}%
% ---------- פרק 28 | פסוק 7 ----------
 \textbf{ז} שְׁתֵּ֧י כְתֵפֹ֣ת חֹֽבְרֹ֗ת יִֽהְיֶה־לּ֛וֹ אֶל־שְׁנֵ֥י קְצוֹתָ֖יו וְחֻבָּֽר׃ \Rashi{\textbf{שתי כתפת וגו'.}הסינר מלמטה וחשב האפוד היא החגורה, וצמודה לו מלמעלה דגמת סינר הנשים, ומגבו של כהן היו מחברות בחשב שתי חתיכות כמין שתי רצועות רחבות, אחת כנגד כל כתף וכתף, וזוקפן על שתי כתפותיו עד שנקפלות לפניו כנגד החזה, וע"י חבורן לטבעות החשן נאחזין מלפניו כנגד לבו שאין נופלות, כמו שמפרש בענין, והיו זקופות והולכות כנגד כתפיו ושתי אבני השהם קבועות בהן אחת בכל אחת: אל שני קצותיו. אל רחבו של אפוד, שלא היה רחבו אלא כנגד גבו של כהן וגבהו עד כנגד האצילים, שקורין קודי"ש בלעז, שנ' לא יחגרו ביזע (יחזקאל מ"ד) – אין חוגרין במקום זיעה, לא למעלה מאציליהם ולא למטה ממתניהם אלא כנגד אציליהם: וחבר. האפוד עם אותן שתי כתפות האפוד יחבר אותם במחט למטה בחשב, ולא יארגם עמו אלא אורגם לבד ואחר כך מחברם:}%
% ---------- פרק 28 | פסוק 8 ----------
 \textbf{ח} וְחֵ֤שֶׁב אֲפֻדָּתוֹ֙ אֲשֶׁ֣ר עָלָ֔יו כְּמַעֲשֵׂ֖הוּ מִמֶּ֣נּוּ יִהְיֶ֑ה זָהָ֗ב תְּכֵ֧לֶת וְאַרְגָּמָ֛ן וְתוֹלַ֥עַת שָׁנִ֖י וְשֵׁ֥שׁ מׇשְׁזָֽר׃ \Rashi{\textbf{וחשב אפדתו.}וחגור שעל ידו הוא מאפדו ומתקנהו לכהן ומקשטו: אשר עליו. למעלה בשפת הסינר, היא החגורה: כמעשהו. כאריגת הסינר מעשה חושב ומחמשת מינים, כך אריגת החשב מעשה חושב ומחמשת מינין: ממנו יהיה. עמו יהיה ארוג ולא יארגנו לבד ויחברנו:}%
% ---------- פרק 28 | פסוק 9 ----------
 \textbf{ט} וְלָ֣קַחְתָּ֔ אֶת־שְׁתֵּ֖י אַבְנֵי־שֹׁ֑הַם וּפִתַּחְתָּ֣ עֲלֵיהֶ֔ם שְׁמ֖וֹת בְּנֵ֥י יִשְׂרָאֵֽל׃ %
% ---------- פרק 28 | פסוק 10 ----------
 \textbf{י} שִׁשָּׁה֙ מִשְּׁמֹתָ֔ם עַ֖ל הָאֶ֣בֶן הָאֶחָ֑ת וְאֶת־שְׁמ֞וֹת הַשִּׁשָּׁ֧ה הַנּוֹתָרִ֛ים עַל־הָאֶ֥בֶן הַשֵּׁנִ֖ית כְּתוֹלְדֹתָֽם׃ \Rashi{\textbf{כתולדתם.}כסדר שנולדו – ראובן, שמעון, לוי, יהודה, דן, נפתלי על האחת, ועל השניה גד, אשר, יששכר זבולן, יוסף, בנימין מלא, שכן הוא כתוב במקום תולדתו, כ"ה אותיות בכל אחת ואחת:}%
% ---------- פרק 28 | פסוק 11 ----------
 \textbf{יא} מַעֲשֵׂ֣ה חָרַשׁ֮ אֶ֒בֶן֒ פִּתּוּחֵ֣י חֹתָ֗ם תְּפַתַּח֙ אֶת־שְׁתֵּ֣י הָאֲבָנִ֔ים עַל־שְׁמֹ֖ת בְּנֵ֣י יִשְׂרָאֵ֑ל מֻסַבֹּ֛ת מִשְׁבְּצ֥וֹת זָהָ֖ב תַּעֲשֶׂ֥ה אֹתָֽם׃ \Rashi{\textbf{מעשה חרש אבן.}מעשה אמן של אבנים, חרש זה דבוק הוא לתבה שלאחריו ולפיכך הוא נקוד פתח בסופו, וכן חרש עצים נטה קו (ישעיהו מ"ד) – חרש של עצים, וכן חרש ברזל מעצד, (שם), כל אלה דבוקים ופתוחים: פתוחי חתם. כתרגומו כתב מפרש כגלף דעזקא, חרוצות האותיות בתוכן כמו שחורצין חותמי טבעות שהם לחתם אגרות – כתב נכר ומפרש: על שמת. כמו בשמת: מסבת משבצות. מקפות האבנים במשבצות זהב, שעושה מושב האבן בזהב כמין גמא למדת האבן ומשקעה במשבצת, נמצאת המשבצת סובבת את האבן סביב ומחבר המשבצות בכתפות האפוד:}%
% ---------- פרק 28 | פסוק 12 ----------
 \textbf{יב} וְשַׂמְתָּ֞ אֶת־שְׁתֵּ֣י הָאֲבָנִ֗ים עַ֚ל כִּתְפֹ֣ת הָֽאֵפֹ֔ד אַבְנֵ֥י זִכָּרֹ֖ן לִבְנֵ֣י יִשְׂרָאֵ֑ל וְנָשָׂא֩ אַהֲרֹ֨ן אֶת־שְׁמוֹתָ֜ם לִפְנֵ֧י יְהֹוָ֛ה עַל־שְׁתֵּ֥י כְתֵפָ֖יו לְזִכָּרֹֽן׃ \{ס\} \Rashi{\textbf{לזכרן.}שיהא רואה הקב"ה את השבטים כתובים לפניו ויזכר צדקתם:}%
% ---------- פרק 28 | פסוק 13 ----------
 \textbf{יג} וְעָשִׂ֥יתָ מִשְׁבְּצֹ֖ת זָהָֽב׃ \Rashi{\textbf{ועשית משבצת.}מעוט משבצות שתים; ולא פרש לך עתה בפרשה זו אלא מקצת צרכן, ובפרשת החשן גומר לך פרושן:}%
% ---------- פרק 28 | פסוק 14 ----------
 \textbf{יד} וּשְׁתֵּ֤י שַׁרְשְׁרֹת֙ זָהָ֣ב טָה֔וֹר מִגְבָּלֹ֛ת תַּעֲשֶׂ֥ה אֹתָ֖ם מַעֲשֵׂ֣ה עֲבֹ֑ת וְנָתַתָּ֛ה אֶת־שַׁרְשְׁרֹ֥ת הָעֲבֹתֹ֖ת עַל־הַֽמִּשְׁבְּצֹֽת׃ \{ס\} \Rashi{\textbf{שרשרת זהב.}שלשלאות: מגבלת. לסוף גבול החשן תעשה אותם: מעשה עבת. מעשה קליעת חוטין, ולא מעשה נקבים וכפלים כאותן שעושין לבורות, אלא כאותן שעושין לערדסקאות שקורין אנשנשוי"רש בלעז: ונתתה את שרשרת. של עבותות העשויות מעשה עבות על משבצות הללו; ולא זה הוא מקום צואת עשיתן של שרשרות, ולא צואת קביעתן, אין תעשה האמור כאן לשון צווי, ואין ונתתה האמור כאן לשון צווי, אלא לשון עתיד, כי בפרשת החשן חוזר ומצוה על עשיתן ועל קביעתן, ולא נכתב כאן אלא להודיע מקצת צרך המשבצות שצוה לעשות עם האפוד, וכתב לך זאת לומר לך המשבצות הללו יזקקו לך, לכשתעשה שרשרות מגבלות על החשן, תתנם על המשבצות הללו:}%
% ---------- פרק 28 | פסוק 15 ----------
 \textbf{טו} וְעָשִׂ֜יתָ חֹ֤שֶׁן מִשְׁפָּט֙ מַעֲשֵׂ֣ה חֹשֵׁ֔ב כְּמַעֲשֵׂ֥ה אֵפֹ֖ד תַּעֲשֶׂ֑נּוּ זָ֠הָ֠ב תְּכֵ֨לֶת וְאַרְגָּמָ֜ן וְתוֹלַ֧עַת שָׁנִ֛י וְשֵׁ֥שׁ מׇשְׁזָ֖ר תַּעֲשֶׂ֥ה אֹתֽוֹ׃ \Rashi{\textbf{חשן משפט.}שמכפר על קלקול הדין. ד"א, משפט שמברר דבריו והבטחתו אמת, דריש"נטנט בלעז; שהמשפט משמש שלש לשונות, דברי טענות הבעלי דינים וגמר הדין וענש הדין, אם ענש מיתה, אם ענש מכות, אם ענש ממון, וזה משמש לשון ברור דברים, שמפרש ומברר דבריו: כמעשה אפד. מעשה חושב ומחמשת מינין:}%
% ---------- פרק 28 | פסוק 16 ----------
 \textbf{טז} רָב֥וּעַ יִֽהְיֶ֖ה כָּפ֑וּל זֶ֥רֶת אׇרְכּ֖וֹ וְזֶ֥רֶת רׇחְבּֽוֹ׃ \Rashi{\textbf{זרת ארכו וזרת רחבו.}כפול; ומטל לו לפניו כנגד לבו, שנאמר והיו על לב אהרן, תלוי בכתפות האפוד הבאות מאחוריו על כתפיו ונקפלות ויורדות לפניו מעט, והחשן תלוי בהן בשרשרות וטבעות כמו שמפרש בענין:}%
% ---------- פרק 28 | פסוק 17 ----------
 \textbf{יז} וּמִלֵּאתָ֥ בוֹ֙ מִלֻּ֣אַת אֶ֔בֶן אַרְבָּעָ֖ה טוּרִ֣ים אָ֑בֶן ט֗וּר אֹ֤דֶם פִּטְדָה֙ וּבָרֶ֔קֶת הַטּ֖וּר הָאֶחָֽד׃ \Rashi{\textbf{ומלאת בו.}על שם שהאבנים ממלאות גמות המשבצות המתקנות להן, קורא אותן בלשון מלואים:}%
% ---------- פרק 28 | פסוק 18 ----------
 \textbf{יח} וְהַטּ֖וּר הַשֵּׁנִ֑י נֹ֥פֶךְ סַפִּ֖יר וְיָהֲלֹֽם׃ %
% ---------- פרק 28 | פסוק 19 ----------
 \textbf{יט} וְהַטּ֖וּר הַשְּׁלִישִׁ֑י לֶ֥שֶׁם שְׁב֖וֹ וְאַחְלָֽמָה׃ %
% ---------- פרק 28 | פסוק 20 ----------
 \textbf{כ} וְהַטּוּר֙ הָרְבִיעִ֔י תַּרְשִׁ֥ישׁ וְשֹׁ֖הַם וְיָשְׁפֵ֑ה מְשֻׁבָּצִ֥ים זָהָ֛ב יִהְי֖וּ בְּמִלּוּאֹתָֽם׃ \Rashi{\textbf{משבצים זהב יהיו.}הטורים במלואתם – מקפים משבצות זהב בעמק שעור שיתמלא בעבי האבן, זהו לשון במלואתם, כשעור מלוי עבין של אבנים יהיה עמק המשבצות, לא פחות ולא יותר:}%
% ---------- פרק 28 | פסוק 21 ----------
 \textbf{כא} וְ֠הָאֲבָנִ֠ים תִּֽהְיֶ֜יןָ עַל־שְׁמֹ֧ת בְּנֵֽי־יִשְׂרָאֵ֛ל שְׁתֵּ֥ים עֶשְׂרֵ֖ה עַל־שְׁמֹתָ֑ם פִּתּוּחֵ֤י חוֹתָם֙ אִ֣ישׁ עַל־שְׁמ֔וֹ תִּֽהְיֶ֕יןָ לִשְׁנֵ֥י עָשָׂ֖ר שָֽׁבֶט׃ \Rashi{\textbf{איש על שמו.}כסדר תולדותם סדר האבנים, אדם לראובן, פטדה לשמעון, וכן כלם:}%
% ---------- פרק 28 | פסוק 22 ----------
 \textbf{כב} וְעָשִׂ֧יתָ עַל־הַחֹ֛שֶׁן שַֽׁרְשֹׁ֥ת גַּבְלֻ֖ת מַעֲשֵׂ֣ה עֲבֹ֑ת זָהָ֖ב טָהֽוֹר׃ \Rashi{\textbf{על החשן.}בשביל החשן, לקבעם בטבעותיו, כמו שמפרש למטה בענין: שרשת. לשון שרשי אילן, המאחזין לאילן להאחז ולהתקע בארץ, אף אלו יהיו מאחזין לחשן שבהם יהיה תלוי באפוד, והן שתי שרשרות האמורות למעלה בענין המשבצות; ואף שרשרות פתר מנחם בן סרוק לשון שרשים, ואמר שהרי"ש יתרה כמו מ"ם שבשלשום ומ"ם שבריקם, ואיני רואה את דבריו, אלא שרשרת בלשון עברית כשלשלת בלשון משנה: גבלת. הוא מגבלות האמור למעלה, שתתקעם בטבעות שיהיו בגבול החשן, וכל גבול לשון קצה, אשומי"יל בלעז: מעשה עבת. מעשה קליעה:}%
% ---------- פרק 28 | פסוק 23 ----------
 \textbf{כג} וְעָשִׂ֙יתָ֙ עַל־הַחֹ֔שֶׁן שְׁתֵּ֖י טַבְּע֣וֹת זָהָ֑ב וְנָתַתָּ֗ אֶת־שְׁתֵּי֙ הַטַּבָּע֔וֹת עַל־שְׁנֵ֖י קְצ֥וֹת הַחֹֽשֶׁן׃ \Rashi{\textbf{על החשן.}לצרך החשן – כדי לקבעם בו: ולא יתכן לומר שתהא תחלת עשיתן עליו, שאם כן מה הוא שחוזר ואומר ונתת את שתי הטבעות והלא כבר נתונין בו? היה לו לכתב בתחלת המקרא ועשית על קצות החשן שתי טבעות זהב, ואף בשרשרות צריך אתה לפתר כן: על שני קצות החשן. לשתי פאות שכנגד הצואר, לימנית ולשמאלית, הבאים מול כתפות האפוד:}%
% ---------- פרק 28 | פסוק 24 ----------
 \textbf{כד} וְנָתַתָּ֗ה אֶת־שְׁתֵּי֙ עֲבֹתֹ֣ת הַזָּהָ֔ב עַל־שְׁתֵּ֖י הַטַּבָּעֹ֑ת אֶל־קְצ֖וֹת הַחֹֽשֶׁן׃ \Rashi{\textbf{ונתתה את שתי עבתת הזהב.}הן הן שרשות גבלות הכתובות למעלה, ולא פרש מקום קבוען בחשן, עכשו מפרש לך שיהא תוחב אותן בטבעות; ותדע לך שהן הן הראשונות, שהרי בפרשת אלה פקודי לא הכפלו:}%
% ---------- פרק 28 | פסוק 25 ----------
 \textbf{כה} וְאֵ֨ת שְׁתֵּ֤י קְצוֹת֙ שְׁתֵּ֣י הָעֲבֹתֹ֔ת תִּתֵּ֖ן עַל־שְׁתֵּ֣י הַֽמִּשְׁבְּצ֑וֹת וְנָתַתָּ֛ה עַל־כִּתְפ֥וֹת הָאֵפֹ֖ד אֶל־מ֥וּל פָּנָֽיו׃ \Rashi{\textbf{ואת שתי קצות של שתי עבתת, שני ראשיהם של כל אחת.}תתן על שתי המשבצות. הן הן הכתובות למעלה, בין פרשת החשן ופרשת האפוד, ולא פרש את צרכן ואת מקומן, עכשו מפרש שיתקע בהן ראשי העבותות התחובות בטבעות החשן לימין ולשמאל אצל הצואר; שני ראשי שרשרת הימנית תוקע במשבצות של ימין, וכן בשל שמאל שני ראשי השרשרת השמאלית: ונתתה. המשבצות על כתפות האפוד, אחת בזו ואחת בזו, נמצאו כתפות האפוד מחזיקין את החשן שלא יפל, ובהן הוא תלוי, ועדין שפת החשן התחתונה הולכת ובאה ונוקשת על כרסו ואינה דבוקה לו יפה, לכך הצרך עוד שתי טבעות לתחתיתו, כמו שמפרש והולך: אל מול פניו. של אפוד, שלא יתן המשבצות בעבר הכתפות שכלפי המעיל, אלא בעבר העליון שכלפי החוץ, והוא קרוי מול פניו של אפוד, כי אותו עבר שאינו נראה אינו קרוי פנים:}%
% ---------- פרק 28 | פסוק 26 ----------
 \textbf{כו} וְעָשִׂ֗יתָ שְׁתֵּי֙ טַבְּע֣וֹת זָהָ֔ב וְשַׂמְתָּ֣ אֹתָ֔ם עַל־שְׁנֵ֖י קְצ֣וֹת הַחֹ֑שֶׁן עַל־שְׂפָת֕וֹ אֲשֶׁ֛ר אֶל־עֵ֥בֶר הָאֵפֹ֖ד*(בספרי ספרד ואשכנז הָאֵפ֖וֹד) בָּֽיְתָה׃ \Rashi{\textbf{על שני קצות החשן.}הן שתי פאותיו התחתונות לימין ולשמאל: על שפתו אשר אל עבר האפוד ביתה. הרי לך שני סימנין, האחד שיתנם בשני קצוות של תחתיתו שהוא כנגד האפוד, שעליונו אינו כנגד האפוד, שהרי סמוך לצואר הוא והאפוד נתון על מתניו, ועוד נתן סימן, שלא יקבעם בעבר החשן שכלפי החוץ, אלא בעבר שכלפי פנים, שנאמר ביתה, ואותו העבר הוא לצד האפוד, שחשב האפוד חוגרו הכהן ונקפל הסינר לפני הכהן על מתניו, וקצת כרסו מכאן ומכאן עד כנגד קצות החשן וקצותיו שוכבין עליו:}%
% ---------- פרק 28 | פסוק 27 ----------
 \textbf{כז} וְעָשִׂ֘יתָ֮ שְׁתֵּ֣י טַבְּע֣וֹת זָהָב֒ וְנָתַתָּ֣ה אֹתָ֡ם עַל־שְׁתֵּי֩ כִתְפ֨וֹת הָאֵפ֤וֹד מִלְּמַ֙טָּה֙ מִמּ֣וּל פָּנָ֔יו לְעֻמַּ֖ת מַחְבַּרְתּ֑וֹ מִמַּ֕עַל לְחֵ֖שֶׁב הָאֵפֽוֹד׃ \Rashi{\textbf{על שתי כתפות מלמטה.}שהמשבצות נתונות בראשי כתפות האפוד העליונים הבאים על כתפיו, כנגד גרונו, ונקפלות ויורדות לפניו, והטבעות צוה לתן בראשן השני, שהוא מחבר לאפוד, והוא שנ' לעמת מחברתו, סמוך למקום חבורן באפוד, למעלה מן החגורה מעט, שהמחברת לעמת החגורה, ואלו נתונים מעט בגבה זקיפת הכתפות, הוא שנאמר ממעל לחשב האפוד, והן כנגד סוף החשן, ונותן פתיל תכלת באותן הטבעות ובטבעות החשן ורוכסן באותו פתיל לימין ולשמאל, שלא יהא תחתית החשן הולך לפנים וחוזר לאחור, ונוקש על כרסו, ונמצא מישב על המעיל יפה: ממול פניו. בעבר החיצון:}%
% ---------- פרק 28 | פסוק 28 ----------
 \textbf{כח} וְיִרְכְּס֣וּ אֶת־הַ֠חֹ֠שֶׁן מִֽטַּבְּעֹתָ֞ו אֶל־טַבְּעֹ֤ת הָאֵפוֹד֙ בִּפְתִ֣יל תְּכֵ֔לֶת לִֽהְי֖וֹת עַל־חֵ֣שֶׁב הָאֵפ֑וֹד וְלֹֽא־יִזַּ֣ח הַחֹ֔שֶׁן מֵעַ֖ל הָאֵפֽוֹד׃ \Rashi{\textbf{וירכסו.}לשון חבור; וכן מרכסי איש (תהילים ל"א) – חבורי חבלי רשעים, וכן והרכסים לבקעה (ישעיהו מ') – הרים הסמוכים זה לזה, שאי אפשר לירד לגיא שביניהם אלא בקשי גדול, שמתוך סמיכתן הגיא זקופה ועמקה – יהא לבקעת מישור ונוחה לילך: להיות על חשב האפוד. להיות החשן דבוק אל חשב האפוד: ולא יזח. לשון נתוק; ולשון ערבי הוא כדברי דונש בן לברט:}%
% ---------- פרק 28 | פסוק 29 ----------
 \textbf{כט} וְנָשָׂ֣א אַ֠הֲרֹ֠ן אֶת־שְׁמ֨וֹת בְּנֵֽי־יִשְׂרָאֵ֜ל בְּחֹ֧שֶׁן הַמִּשְׁפָּ֛ט עַל־לִבּ֖וֹ בְּבֹא֣וֹ אֶל־הַקֹּ֑דֶשׁ לְזִכָּרֹ֥ן לִפְנֵֽי־יְהֹוָ֖ה תָּמִֽיד׃ %
% ---------- פרק 28 | פסוק 30 ----------
 \textbf{ל} וְנָתַתָּ֞ אֶל־חֹ֣שֶׁן הַמִּשְׁפָּ֗ט אֶת־הָאוּרִים֙ וְאֶת־הַתֻּמִּ֔ים וְהָיוּ֙ עַל־לֵ֣ב אַהֲרֹ֔ן בְּבֹא֖וֹ לִפְנֵ֣י יְהֹוָ֑ה וְנָשָׂ֣א אַ֠הֲרֹ֠ן אֶת־מִשְׁפַּ֨ט בְּנֵי־יִשְׂרָאֵ֧ל עַל־לִבּ֛וֹ לִפְנֵ֥י יְהֹוָ֖ה תָּמִֽיד׃ \{ס\} \Rashi{\textbf{את האורים ואת התמים.}הוא כתב שם המפרש, שהיה נותנו בתוך כפלי החשן, שעל ידו הוא מאיר דבריו ומתמם את דבריו; ובמקדש שני היה החשן – שאי אפשר לכהן גדול להיות מחסר בגדים – אבל אותו השם לא היה בתוכו, ועל שם אותו הכתב הוא קרוי משפט, שנאמר (במדבר כ"ז), ושאל לו במשפט האורים (יומא ע"ג): את משפט בני ישראל. דבר שהם נשפטים ונוכחים על ידו, אם לעשות דבר או לא לעשות; ולפי מ"א שהחשן מכפר על מעותי הדין, נקרא משפט על שם סליחת המשפט:}%
% ---------- פרק 28 | פסוק 31 ----------
 \textbf{לא} וְעָשִׂ֛יתָ אֶת־מְעִ֥יל הָאֵפ֖וֹד כְּלִ֥יל תְּכֵֽלֶת׃ \Rashi{\textbf{את מעיל האפוד.}שהאפוד נתון עליו לחגורה: כליל תכלת. כלו תכלת, שאין מין אחר מערב בו:}%
% ---------- פרק 28 | פסוק 32 ----------
 \textbf{לב} וְהָיָ֥ה פִֽי־רֹאשׁ֖וֹ בְּתוֹכ֑וֹ שָׂפָ֡ה יִֽהְיֶה֩ לְפִ֨יו סָבִ֜יב מַעֲשֵׂ֣ה אֹרֵ֗ג כְּפִ֥י תַחְרָ֛א יִֽהְיֶה־לּ֖וֹ לֹ֥א יִקָּרֵֽעַ׃ \Rashi{\textbf{והיה פי ראשו.}פי המעיל שבגבהו – הוא פתיחת בית הצואר: בתוכו. כתרגומו כפיל לגוה, כפול לתוכו, להיות לו לשפה כפילתו, והיה מעשה אורג ולא במחט: כפי תחרא. למדנו שהשריונים שלהם פיהם כפול: לא יקרע. כדי שלא יקרע; והקורעו עובר בלאו, שזה ממנין לאוין שבתורה, וכן ולא יזח החשן, וכן לא יסרו ממנו הנאמר בבדי הארון:}%
% ---------- פרק 28 | פסוק 33 ----------
 \textbf{לג} וְעָשִׂ֣יתָ עַל־שׁוּלָ֗יו רִמֹּנֵי֙ תְּכֵ֤לֶת וְאַרְגָּמָן֙ וְתוֹלַ֣עַת שָׁנִ֔י עַל־שׁוּלָ֖יו סָבִ֑יב וּפַעֲמֹנֵ֥י זָהָ֛ב בְּתוֹכָ֖ם סָבִֽיב׃ \Rashi{\textbf{רמני.}עגלים וחלולים היו כמין רמונים העשויים כביצת תרנגלת: ופעמני זהב. זגין עם ענבלין שבתוכם: בתוכם סביב. ביניהם סביב – בין שני רמונים פעמון אחד דבוק ותלוי בשולי המעיל:}%
% ---------- פרק 28 | פסוק 34 ----------
 \textbf{לד} פַּעֲמֹ֤ן זָהָב֙ וְרִמּ֔וֹן פַּֽעֲמֹ֥ן זָהָ֖ב וְרִמּ֑וֹן עַל־שׁוּלֵ֥י הַמְּעִ֖יל סָבִֽיב׃ \Rashi{\textbf{פעמן זהב ורמון פעמן זהב ורמון.}אצלו:}%
% ---------- פרק 28 | פסוק 35 ----------
 \textbf{לה} וְהָיָ֥ה עַֽל־אַהֲרֹ֖ן לְשָׁרֵ֑ת וְנִשְׁמַ֣ע ק֠וֹל֠וֹ בְּבֹא֨וֹ אֶל־הַקֹּ֜דֶשׁ לִפְנֵ֧י יְהֹוָ֛ה וּבְצֵאת֖וֹ וְלֹ֥א יָמֽוּת׃ \{ס\} \Rashi{\textbf{ולא ימות.}מכלל לאו אתה שומע הן – אם יהיו לו לא יתחיב מיתה, הא אם יכנס מחסר אחד מן הבגדים הללו, חיב מיתה בידי שמים (סנהדרין פ"ג):}%
% ---------- פרק 28 | פסוק 36 ----------
 \textbf{לו} וְעָשִׂ֥יתָ צִּ֖יץ זָהָ֣ב טָה֑וֹר וּפִתַּחְתָּ֤ עָלָיו֙ פִּתּוּחֵ֣י חֹתָ֔ם קֹ֖דֶשׁ לַֽיהֹוָֽה׃ \Rashi{\textbf{ציץ.}כמין טס של זהב היה, רחב שתי אצבעות, מקיף על המצח מאזן לאזן:}%
% ---------- פרק 28 | פסוק 37 ----------
 \textbf{לז} וְשַׂמְתָּ֤ אֹתוֹ֙ עַל־פְּתִ֣יל תְּכֵ֔לֶת וְהָיָ֖ה עַל־הַמִּצְנָ֑פֶת אֶל־מ֥וּל פְּנֵֽי־הַמִּצְנֶ֖פֶת יִהְיֶֽה׃ \Rashi{\textbf{על פתיל תכלת.}ובמקום אחר הוא אומר ויתנו עליו פתיל תכלת? ועוד כתיב כאן והיה על המצנפת, ולמטה הוא אומר והיה על מצח אהרן? ובשחיטת קדשים שנינו שערו היה נראה בין ציץ למצנפת ששם מניח תפלין? למדנו שהמצנפת למעלה בגבה הראש ואינה עמקה לכנס בה כל הראש עד המצח, והציץ מלמטה, והפתילים היו בנקבים ותלויין בו בשני ראשים ובאמצעו, ששה בשלשה מקומות הללו, פתיל מלמעלה אחד מבחוץ ואחד מבפנים כנגדו, וקושר ראשי הפתילים מאחורי הערף שלשתן, ונמצאו בין ארך הטס ופתילי ראשיו מקיפין את הקדקד, והפתיל האמצעי שבראשו, קשור עם ראשי השנים, והולך על פני רחב הראש מלמעלה, נמצא עשוי כמין כובע; ועל פתיל האמצעי הוא אומר והיה על המצנפת, והיה נותן הציץ על ראשו כמין כובע על המצנפת, והפתיל האמצעי מחזיקו שאינו נופל והטס תלוי כנגד מצחו, ונתקימו כל המקראות – פתיל על הציץ, וציץ על הפתיל, ופתיל על המצנפת מלמעלה:}%
% ---------- פרק 28 | פסוק 38 ----------
 \textbf{לח} וְהָיָה֮ עַל־מֵ֣צַח אַהֲרֹן֒ וְנָשָׂ֨א אַהֲרֹ֜ן אֶת־עֲוֺ֣ן הַקֳּדָשִׁ֗ים אֲשֶׁ֤ר יַקְדִּ֙ישׁוּ֙ בְּנֵ֣י יִשְׂרָאֵ֔ל לְכׇֽל־מַתְּנֹ֖ת קׇדְשֵׁיהֶ֑ם וְהָיָ֤ה עַל־מִצְחוֹ֙ תָּמִ֔יד לְרָצ֥וֹן לָהֶ֖ם לִפְנֵ֥י יְהֹוָֽה׃ \Rashi{\textbf{ונשא אהרן.}לשון סליחה; ואעפ"כ אינו זז ממשמעו – אהרן נושא את המשא של עון – נמצא מסלק העון מן הקדשים (חולין קל"ח): את עון הקדשים. לרצות על הדם ועל החלב שקרבו בטמאה, כמו ששנינו: אי זה עון הוא נושא? אם עון פגול, הרי כבר נאמר לא ירצה (ויקרא י"ט), ואם עון נותר הרי נאמר לא יחשב (שם ז') – ואין לומר שיכפר על עון הכהן שהקריב טמא, שהרי עון הקדשים נאמר ולא עון המקריבים – הא אינו מרצה אלא להכשיר הקרבן (יומא ז'): והיה על מצחו תמיד. אי אפשר לומר שיהא על מצחו תמיד, שהרי אינו עליו אלא בשעת העבודה, אלא תמיד לרצות להם אפלו אינו על מצחו – שלא היה כהן גדול עובד באותה שעה, ולדברי האומר עודהו על מצחו מכפר ומרצה, ואם לאו אינו מרצה, נדרש על מצחו תמיד, מלמד שממשמש בו בעודו על מצחו, שלא יסיח דעתו ממנו (שם):}%
% ---------- פרק 28 | פסוק 39 ----------
 \textbf{לט} וְשִׁבַּצְתָּ֙ הַכְּתֹ֣נֶת שֵׁ֔שׁ וְעָשִׂ֖יתָ מִצְנֶ֣פֶת שֵׁ֑שׁ וְאַבְנֵ֥ט תַּעֲשֶׂ֖ה מַעֲשֵׂ֥ה רֹקֵֽם׃ \Rashi{\textbf{ושבצת.}עשה אותם משבצות משבצות, וכלה של שש:}%
% ---------- פרק 28 | פסוק 40 ----------
 \textbf{מ} וְלִבְנֵ֤י אַהֲרֹן֙ תַּעֲשֶׂ֣ה כֻתֳּנֹ֔ת וְעָשִׂ֥יתָ לָהֶ֖ם אַבְנֵטִ֑ים וּמִגְבָּעוֹת֙ תַּעֲשֶׂ֣ה לָהֶ֔ם לְכָב֖וֹד וּלְתִפְאָֽרֶת׃ \Rashi{\textbf{ולבני אהרן תעשה כתנת.}ארבעה בגדים הללו ולא יותר, כתנת ואבנט ומגבעות – היא מצנפת – ומכנסים כתובים למטה בפרשה:}%
% ---------- פרק 28 | פסוק 41 ----------
 \textbf{מא} וְהִלְבַּשְׁתָּ֤ אֹתָם֙ אֶת־אַהֲרֹ֣ן אָחִ֔יךָ וְאֶת־בָּנָ֖יו אִתּ֑וֹ וּמָשַׁחְתָּ֨ אֹתָ֜ם וּמִלֵּאתָ֧ אֶת־יָדָ֛ם וְקִדַּשְׁתָּ֥ אֹתָ֖ם וְכִהֲנ֥וּ לִֽי׃ \Rashi{\textbf{והלבשת אתם את אהרן.}אותם האמורין באהרן – חשן, ואפוד, ומעיל, וכתנת תשבץ, מצנפת, ואבנט, וציץ, ומכנסים הכתובים למטה בכלם: ואת בניו אתו. אותם הכתובים בהם: ומשחת אתם. את אהרן ואת בניו בשמן המשחה: ומלאת את ידם. כל מלוי ידים לשון חנוך, כשהוא נכנס לדבר להיות מחזק בו מאותו יום והלאה הוא: ובלשון לעז כשממנין אדם על פקידת דבר, נותן השליט בידו בית יד של עור, שקורין גאנ"ט בלעז, ועל ידו הוא מחזיקו בדבר, וקורין לאותו מסירה רווישט"יר בלעז, והוא מלוי ידים:}%
% ---------- פרק 28 | פסוק 42 ----------
 \textbf{מב} וַעֲשֵׂ֤ה לָהֶם֙ מִכְנְסֵי־בָ֔ד לְכַסּ֖וֹת בְּשַׂ֣ר עֶרְוָ֑ה מִמׇּתְנַ֥יִם וְעַד־יְרֵכַ֖יִם יִהְיֽוּ׃ \Rashi{\textbf{ועשה להם.}לאהרן ולבניו: מכנסי בד. הרי שמונה בגדים לכהן גדול וארבעה לכהן הדיוט:}%
% ---------- פרק 28 | פסוק 43 ----------
 \textbf{מג} וְהָיוּ֩ עַל־אַהֲרֹ֨ן וְעַל־בָּנָ֜יו בְּבֹאָ֣ם ׀ אֶל־אֹ֣הֶל מוֹעֵ֗ד א֣וֹ בְגִשְׁתָּ֤ם אֶל־הַמִּזְבֵּ֙חַ֙ לְשָׁרֵ֣ת בַּקֹּ֔דֶשׁ וְלֹא־יִשְׂא֥וּ עָוֺ֖ן וָמֵ֑תוּ חֻקַּ֥ת עוֹלָ֛ם ל֖וֹ וּלְזַרְע֥וֹ אַחֲרָֽיו׃ \{ס\} \Rashi{\textbf{והיו על אהרן.}כל הבגדים האלה, על אהרן הראויין לו: ועל בניו. האמורין בהם: בבאם אל אהל מועד. להיכל, וכן למשכן: ומתו. הא למדת שהמשמש מחסר בגדים במיתה. חקת עולם לו. כל מקום שנאמר חקת עולם הוא גזרה מיד ולדורות, לעכב בו (מנחות י"ט):}%
% ---------- פרק 29 | פסוק 1 ----------
 \textbf{א} וְזֶ֨ה הַדָּבָ֜ר אֲשֶֽׁר־תַּעֲשֶׂ֥ה לָהֶ֛ם לְקַדֵּ֥שׁ אֹתָ֖ם לְכַהֵ֣ן לִ֑י לְ֠קַ֠ח פַּ֣ר אֶחָ֧ד בֶּן־בָּקָ֛ר וְאֵילִ֥ם שְׁנַ֖יִם תְּמִימִֽם׃ \Rashi{\textbf{לקח.}כמו קח. ושתי גזירות הן, אחת של קיחה ואחת של לקיחה, ולהן פתרון אחד: פר אחד. לכפר על מעשה העגל שהוא פר:}%
% ---------- פרק 29 | פסוק 2 ----------
 \textbf{ב} וְלֶ֣חֶם מַצּ֗וֹת וְחַלֹּ֤ת מַצֹּת֙ בְּלוּלֹ֣ת בַּשֶּׁ֔מֶן וּרְקִיקֵ֥י מַצּ֖וֹת מְשֻׁחִ֣ים בַּשָּׁ֑מֶן סֹ֥לֶת חִטִּ֖ים תַּעֲשֶׂ֥ה אֹתָֽם׃ \Rashi{\textbf{ולחם מצות וחלת מצת ורקיקי מצות.}הרי אלו שלשה מינין רבוכה, וחלות ורקיקים לחם מצות, היא הקרויה למטה בענין (פסוק כג) חלת לחם שמן, על שם שנותן שמן ברבוכה כנגד החלות והרקיקין, וכל המינין באים עשר עשר חלות: בלולת בשמן. כשהן קמח יוצק בהן שמן ובוללן: משחים בשמן. אחר אפיתן מושחן כמין כי יונית, שהיא עשויה כנו"ן שלנו:}%
% ---------- פרק 29 | פסוק 3 ----------
 \textbf{ג} וְנָתַתָּ֤ אוֹתָם֙ עַל־סַ֣ל אֶחָ֔ד וְהִקְרַבְתָּ֥ אֹתָ֖ם בַּסָּ֑ל וְאֶ֨ת־הַפָּ֔ר וְאֵ֖ת שְׁנֵ֥י הָאֵילִֽם׃ \Rashi{\textbf{והקרבת אתם.}אל חצר המשכן ביום הקמתו:}%
% ---------- פרק 29 | פסוק 4 ----------
 \textbf{ד} וְאֶת־אַהֲרֹ֤ן וְאֶת־בָּנָיו֙ תַּקְרִ֔יב אֶל־פֶּ֖תַח אֹ֣הֶל מוֹעֵ֑ד וְרָחַצְתָּ֥ אֹתָ֖ם בַּמָּֽיִם׃ \Rashi{\textbf{ורחצת.}זו טבילת כל הגוף:}%
% ---------- פרק 29 | פסוק 5 ----------
 \textbf{ה} וְלָקַחְתָּ֣ אֶת־הַבְּגָדִ֗ים וְהִלְבַּשְׁתָּ֤ אֶֽת־אַהֲרֹן֙ אֶת־הַכֻּתֹּ֔נֶת וְאֵת֙ מְעִ֣יל הָאֵפֹ֔ד וְאֶת־הָאֵפֹ֖ד וְאֶת־הַחֹ֑שֶׁן וְאָפַדְתָּ֣ ל֔וֹ בְּחֵ֖שֶׁב הָאֵפֹֽד׃ \Rashi{\textbf{ואפדת.}קשט ותקן החגורה והסינר סביבותיו:}%
% ---------- פרק 29 | פסוק 6 ----------
 \textbf{ו} וְשַׂמְתָּ֥ הַמִּצְנֶ֖פֶת עַל־רֹאשׁ֑וֹ וְנָתַתָּ֛ אֶת־נֵ֥זֶר הַקֹּ֖דֶשׁ עַל־הַמִּצְנָֽפֶת׃ \Rashi{\textbf{נזר הקדש.}זה הציץ: על המצנפת. כמו שפרשתי למעלה – על ידי הפתיל האמצעי ושני פתילין שבראשו, הקשורין שלשתן מאחורי הערף, הוא נותנו על המצנפת כמין כובע:}%
% ---------- פרק 29 | פסוק 7 ----------
 \textbf{ז} וְלָֽקַחְתָּ֙ אֶת־שֶׁ֣מֶן הַמִּשְׁחָ֔ה וְיָצַקְתָּ֖ עַל־רֹאשׁ֑וֹ וּמָשַׁחְתָּ֖ אֹתֽוֹ׃ \Rashi{\textbf{ומשחת אתו.}אף משיחה זו כמין כי, נותן שמן על ראשו ובין רסי עיניו ומחברן באצבעו (כריתות ה'):}%
% ---------- פרק 29 | פסוק 8 ----------
 \textbf{ח} וְאֶת־בָּנָ֖יו תַּקְרִ֑יב וְהִלְבַּשְׁתָּ֖ם כֻּתֳּנֹֽת׃ %
% ---------- פרק 29 | פסוק 9 ----------
 \textbf{ט} וְחָגַרְתָּ֩ אֹתָ֨ם אַבְנֵ֜ט אַהֲרֹ֣ן וּבָנָ֗יו וְחָבַשְׁתָּ֤ לָהֶם֙ מִגְבָּעֹ֔ת וְהָיְתָ֥ה לָהֶ֛ם כְּהֻנָּ֖ה לְחֻקַּ֣ת עוֹלָ֑ם וּמִלֵּאתָ֥ יַֽד־אַהֲרֹ֖ן וְיַד־בָּנָֽיו׃ \Rashi{\textbf{והיתה להם.}מלוי ידים זה לכהנת עולם: ומלאת. על ידי הדברים האלה: יד אהרן ויד בניו. במנוי ופקדת הכהנה:}%
% ---------- פרק 29 | פסוק 10 ----------
 \textbf{י} וְהִקְרַבְתָּ֙ אֶת־הַפָּ֔ר לִפְנֵ֖י אֹ֣הֶל מוֹעֵ֑ד וְסָמַ֨ךְ אַהֲרֹ֧ן וּבָנָ֛יו אֶת־יְדֵיהֶ֖ם עַל־רֹ֥אשׁ הַפָּֽר׃ %
% ---------- פרק 29 | פסוק 11 ----------
 \textbf{יא} וְשָׁחַטְתָּ֥ אֶת־הַפָּ֖ר לִפְנֵ֣י יְהֹוָ֑ה פֶּ֖תַח אֹ֥הֶל מוֹעֵֽד׃ \Rashi{\textbf{פתח אהל מועד.}בחצר המשכן שלפני הפתח:}%
% ---------- פרק 29 | פסוק 12 ----------
 \textbf{יב} וְלָֽקַחְתָּ֙ מִדַּ֣ם הַפָּ֔ר וְנָתַתָּ֛ה עַל־קַרְנֹ֥ת הַמִּזְבֵּ֖חַ בְּאֶצְבָּעֶ֑ךָ וְאֶת־כׇּל־הַדָּ֣ם תִּשְׁפֹּ֔ךְ אֶל־יְס֖וֹד הַמִּזְבֵּֽחַ׃ \Rashi{\textbf{על קרנת.}למעלה בקרנות ממש: ואת כל הדם. שירי הדם: אל יסוד המזבח. כמין בליטת בית קבול עשוי לו סביב סביב, לאחר שעלה אמה מן הארץ:}%
% ---------- פרק 29 | פסוק 13 ----------
 \textbf{יג} וְלָֽקַחְתָּ֗ אֶֽת־כׇּל־הַחֵ֘לֶב֮ הַֽמְכַסֶּ֣ה אֶת־הַקֶּ֒רֶב֒ וְאֵ֗ת הַיֹּתֶ֙רֶת֙ עַל־הַכָּבֵ֔ד וְאֵת֙ שְׁתֵּ֣י הַכְּלָיֹ֔ת וְאֶת־הַחֵ֖לֶב אֲשֶׁ֣ר עֲלֵיהֶ֑ן וְהִקְטַרְתָּ֖ הַמִּזְבֵּֽחָה׃ \Rashi{\textbf{החלב המכסה את הקרב.}הוא הקרום שעל הכרס, שקורין טי"לא: ואת היתרת. הוא טרפשא דכבדא, שקורין איבר"יש: על הכבד. אף מן הכבד טל עמה (ספרא):}%
% ---------- פרק 29 | פסוק 14 ----------
 \textbf{יד} וְאֶת־בְּשַׂ֤ר הַפָּר֙ וְאֶת־עֹר֣וֹ וְאֶת־פִּרְשׁ֔וֹ תִּשְׂרֹ֣ף בָּאֵ֔שׁ מִח֖וּץ לַֽמַּחֲנֶ֑ה חַטָּ֖את הֽוּא׃ \Rashi{\textbf{תשרף באש.}לא מצינו חטאת חיצונה נשרפת אלא זו:}%
% ---------- פרק 29 | פסוק 15 ----------
 \textbf{טו} וְאֶת־הָאַ֥יִל הָאֶחָ֖ד תִּקָּ֑ח וְסָ֨מְכ֜וּ אַהֲרֹ֧ן וּבָנָ֛יו אֶת־יְדֵיהֶ֖ם עַל־רֹ֥אשׁ הָאָֽיִל׃ %
% ---------- פרק 29 | פסוק 16 ----------
 \textbf{טז} וְשָׁחַטְתָּ֖ אֶת־הָאָ֑יִל וְלָֽקַחְתָּ֙ אֶת־דָּמ֔וֹ וְזָרַקְתָּ֥ עַל־הַמִּזְבֵּ֖חַ סָבִֽיב׃ \Rashi{\textbf{וזרקת.}בכלי; אוחז במזרק וזורק כנגד הקרן כדי שיראה לכאן ולכאן, ואין קרבן טעון מתנה באצבע אלא חטאת בלבד, אבל שאר זבחים אינן טעונין קרן ולא אצבע, שמתן דמם מחצי המזבח ולמטה, ואינו עולה בכבש אלא עומד בארץ וזורק (זבחים נ"ג): סביב. כך מפרש בשחיטת קדשים שאין סביב אלא שתי מתנות שהן ארבע, האחת בקרן זוית זו, והאחת שכנגדה באלכסון, וכל מתנה נראית בשני צדי הקרן אילך ואילך, נמצא הדם נתון בארבע רוחות סביב, לכך קרוי סביב:}%
% ---------- פרק 29 | פסוק 17 ----------
 \textbf{יז} וְאֶ֨ת־הָאַ֔יִל תְּנַתֵּ֖חַ לִנְתָחָ֑יו וְרָחַצְתָּ֤ קִרְבּוֹ֙ וּכְרָעָ֔יו וְנָתַתָּ֥ עַל־נְתָחָ֖יו וְעַל־רֹאשֽׁוֹ׃ \Rashi{\textbf{על נתחיו.}עם נתחיו – מוסף על שאר הנתחים:}%
% ---------- פרק 29 | פסוק 18 ----------
 \textbf{יח} וְהִקְטַרְתָּ֤ אֶת־כׇּל־הָאַ֙יִל֙ הַמִּזְבֵּ֔חָה עֹלָ֥ה ה֖וּא לַֽיהֹוָ֑ה רֵ֣יחַ נִיח֔וֹחַ אִשֶּׁ֥ה לַיהֹוָ֖ה הֽוּא׃ \Rashi{\textbf{ריח ניחוח.}נחת רוח לפני, שאמרתי ונעשה רצוני: אשה. לשון אש, והיא הקטרת אברים שעל האש:}%
% ---------- פרק 29 | פסוק 19 ----------
 \textbf{יט} וְלָ֣קַחְתָּ֔ אֵ֖ת הָאַ֣יִל הַשֵּׁנִ֑י וְסָמַ֨ךְ אַהֲרֹ֧ן וּבָנָ֛יו אֶת־יְדֵיהֶ֖ם עַל־רֹ֥אשׁ הָאָֽיִל׃ %
% ---------- פרק 29 | פסוק 20 ----------
 \textbf{כ} וְשָׁחַטְתָּ֣ אֶת־הָאַ֗יִל וְלָקַחְתָּ֤ מִדָּמוֹ֙ וְנָֽתַתָּ֡ה עַל־תְּנוּךְ֩ אֹ֨זֶן אַהֲרֹ֜ן וְעַל־תְּנ֨וּךְ אֹ֤זֶן בָּנָיו֙ הַיְמָנִ֔ית וְעַל־בֹּ֤הֶן יָדָם֙ הַיְמָנִ֔ית וְעַל־בֹּ֥הֶן רַגְלָ֖ם הַיְמָנִ֑ית וְזָרַקְתָּ֧ אֶת־הַדָּ֛ם עַל־הַמִּזְבֵּ֖חַ סָבִֽיב׃ \Rashi{\textbf{תנוך.}הוא הסחוס האמצעי שבתוך האזן, שקורין טנדרו"ס בלעז: בהן ידם. הגודל ובפרק האמצעי (ספרא):}%
% ---------- פרק 29 | פסוק 21 ----------
 \textbf{כא} וְלָקַחְתָּ֞ מִן־הַדָּ֨ם אֲשֶׁ֥ר עַֽל־הַמִּזְבֵּ֘חַ֮ וּמִשֶּׁ֣מֶן הַמִּשְׁחָה֒ וְהִזֵּיתָ֤ עַֽל־אַהֲרֹן֙ וְעַל־בְּגָדָ֔יו וְעַל־בָּנָ֛יו וְעַל־בִּגְדֵ֥י בָנָ֖יו אִתּ֑וֹ וְקָדַ֥שׁ הוּא֙ וּבְגָדָ֔יו וּבָנָ֛יו וּבִגְדֵ֥י בָנָ֖יו אִתּֽוֹ׃ %
% ---------- פרק 29 | פסוק 22 ----------
 \textbf{כב} וְלָקַחְתָּ֣ מִן־הָ֠אַ֠יִל הַחֵ֨לֶב וְהָֽאַלְיָ֜ה וְאֶת־הַחֵ֣לֶב ׀ הַֽמְכַסֶּ֣ה אֶת־הַקֶּ֗רֶב וְאֵ֨ת יֹתֶ֤רֶת הַכָּבֵד֙ וְאֵ֣ת ׀ שְׁתֵּ֣י הַכְּלָיֹ֗ת וְאֶת־הַחֵ֙לֶב֙ אֲשֶׁ֣ר עֲלֵיהֶ֔ן וְאֵ֖ת שׁ֣וֹק הַיָּמִ֑ין כִּ֛י אֵ֥יל מִלֻּאִ֖ים הֽוּא׃ \Rashi{\textbf{החלב.}זה חלב הדקים או הקבה: והאליה. מן הכליות ולמטה, כמו שמפרש בויקרא שנ' (ויקרא ג'), לעמת העצה יסירנה – מקום שהכליות יועצות (חולין י"א), ובאמורי הפר לא נאמר אליה, שאין אליה קרבה אלא בכבש וכבשה ואיל, אבל שור ועז אין טעונין אליה: ואת שוק הימין. לא מצינו הקטרה בשוק הימין עם האמורים אלא זו בלבד: כי איל מלאים הוא. שלמים; לשון שלמות – שהשלם בכל; מגיד הכתוב שהמלואים שלמים, שמשימים שלום למזבח ולעובד העבודה ולבעלים, לכך אני מצריכו החזה, להיות לו לעובד העבודה למנה, וזהו משה ששמש במלואים, והשאר אכלו אהרן ובניו שהם בעלים, כמפרש בענין:}%
% ---------- פרק 29 | פסוק 23 ----------
 \textbf{כג} וְכִכַּ֨ר לֶ֜חֶם אַחַ֗ת וְֽחַלַּ֨ת לֶ֥חֶם שֶׁ֛מֶן אַחַ֖ת וְרָקִ֣יק אֶחָ֑ד מִסַּל֙ הַמַּצּ֔וֹת אֲשֶׁ֖ר לִפְנֵ֥י יְהֹוָֽה׃ \Rashi{\textbf{וככר לחם.}מן החלות: וחלת לחם שמן. ממין הרבוכה: ורקיק. מן הרקיקין, אחד מעשרה שבכל מין ומין. ולא מצינו תרומת לחם הבא עם זבח נקטרת אלא זו בלבד, שתרומת לחמי תודה ואיל נזיר נתונה לכהנים עם חזה ושוק, ומזה לא היה למשה למנה אלא חזה בלבד:}%
% ---------- פרק 29 | פסוק 24 ----------
 \textbf{כד} וְשַׂמְתָּ֣ הַכֹּ֔ל עַ֚ל כַּפֵּ֣י אַהֲרֹ֔ן וְעַ֖ל כַּפֵּ֣י בָנָ֑יו וְהֵנַפְתָּ֥ אֹתָ֛ם תְּנוּפָ֖ה לִפְנֵ֥י יְהֹוָֽה׃ \Rashi{\textbf{על כפי אהרן: והנפת.}שניהם עסוקים בתנופה – הבעלים והכהן, הא כיצד? כהן מניח ידו תחת יד הבעלים ומניף; ובזה היו אהרן ובניו בעלים, ומשה כהן: תנופה. מוליך ומביא למי שארבע רוחות העולם שלו; תנופה מעכבת ומבטלת פרעניות ורוחות רעות, מעלה ומוריד למי שהשמים והארץ שלו, ומעכבת טללים רעים (מנחות ס"ב):}%
% ---------- פרק 29 | פסוק 25 ----------
 \textbf{כה} וְלָקַחְתָּ֤ אֹתָם֙ מִיָּדָ֔ם וְהִקְטַרְתָּ֥ הַמִּזְבֵּ֖חָה עַל־הָעֹלָ֑ה לְרֵ֤יחַ נִיח֙וֹחַ֙ לִפְנֵ֣י יְהֹוָ֔ה אִשֶּׁ֥ה ה֖וּא לַיהֹוָֽה׃ \Rashi{\textbf{על העלה.}על האיל הראשון שהעלית עולה: לריח ניחוח. לנחת רוח למי שאמר ונעשה רצונו: אשה. לאש נתן: לה'. לשמו של מקום:}%
% ---------- פרק 29 | פסוק 26 ----------
 \textbf{כו} וְלָקַחְתָּ֣ אֶת־הֶֽחָזֶ֗ה מֵאֵ֤יל הַמִּלֻּאִים֙ אֲשֶׁ֣ר לְאַהֲרֹ֔ן וְהֵנַפְתָּ֥ אֹת֛וֹ תְּנוּפָ֖ה לִפְנֵ֣י יְהֹוָ֑ה וְהָיָ֥ה לְךָ֖ לְמָנָֽה׃ %
% ---------- פרק 29 | פסוק 27 ----------
 \textbf{כז} וְקִדַּשְׁתָּ֞ אֵ֣ת ׀ חֲזֵ֣ה הַתְּנוּפָ֗ה וְאֵת֙ שׁ֣וֹק הַתְּרוּמָ֔ה אֲשֶׁ֥ר הוּנַ֖ף וַאֲשֶׁ֣ר הוּרָ֑ם מֵאֵיל֙ הַמִּלֻּאִ֔ים מֵאֲשֶׁ֥ר לְאַהֲרֹ֖ן וּמֵאֲשֶׁ֥ר לְבָנָֽיו׃ \Rashi{\textbf{וקדשת את חזה התנופה ואת שוק התרומה וגו'.}קדשם לדורות, להיות נוהגת תרומתם והנפתם בחזה ושוק של שלמים, אבל לא להקטרה, אלא והיה לאהרן ולבניו לאכל: תנופה. לשון הולכה והבאה, ונטלי"ה בלעז: הורם. לשון מעלה ומוריד:}%
% ---------- פרק 29 | פסוק 28 ----------
 \textbf{כח} וְהָיָה֩ לְאַהֲרֹ֨ן וּלְבָנָ֜יו לְחׇק־עוֹלָ֗ם מֵאֵת֙ בְּנֵ֣י יִשְׂרָאֵ֔ל כִּ֥י תְרוּמָ֖ה ה֑וּא וּתְרוּמָ֞ה יִהְיֶ֨ה מֵאֵ֤ת בְּנֵֽי־יִשְׂרָאֵל֙ מִזִּבְחֵ֣י שַׁלְמֵיהֶ֔ם תְּרוּמָתָ֖ם לַיהֹוָֽה׃ \Rashi{\textbf{לחק עולם מאת בני ישראל.}שהשלמים לבעלים, ואת החזה ואת השוק יתנו לכהן: כי תרומה הוא. החזה ושוק הזה:}%
% ---------- פרק 29 | פסוק 29 ----------
 \textbf{כט} וּבִגְדֵ֤י הַקֹּ֙דֶשׁ֙ אֲשֶׁ֣ר לְאַהֲרֹ֔ן יִהְי֥וּ לְבָנָ֖יו אַחֲרָ֑יו לְמׇשְׁחָ֣ה בָהֶ֔ם וּלְמַלֵּא־בָ֖ם אֶת־יָדָֽם׃ \Rashi{\textbf{לבניו אחריו.}למי שבא בגדלה אחריו: למשחה. להתגדל בהם, שיש משיחה שהיא לשון שררה, כמו לך נתתים למשחה (במדבר י"ח), אל תגעו במשיחי (דהי"א טז): ולמלא בם את ידם. על ידי הבגדים הוא מתלבש בכהנה גדולה:}%
% ---------- פרק 29 | פסוק 30 ----------
 \textbf{ל} שִׁבְעַ֣ת יָמִ֗ים יִלְבָּשָׁ֧ם הַכֹּהֵ֛ן תַּחְתָּ֖יו מִבָּנָ֑יו אֲשֶׁ֥ר יָבֹ֛א אֶל־אֹ֥הֶל מוֹעֵ֖ד לְשָׁרֵ֥ת בַּקֹּֽדֶשׁ׃ \Rashi{\textbf{שבעת ימים.}רצופין: ילבשם הכהן. אשר יקום מבניו תחתיו לכהנה גדולה, כשימנוהו להיות כהן גדול: אשר יבא אל אהל מועד. אותו כהן המוכן לכנס לפני ולפנים ביום הכפורים, וזהו כהן גדול, שאין עבודת יום הכפורים כשרה אלא בו: תחתיו מבניו. מלמד שאם יש לו לכהן גדול בן ממלא את מקומו, ימנוהו כהן גדול תחתיו: (הכהן תחתיו מבניו. מכאן ראיה כל לשון כהן פועל עובד ממש, לפיכך נגון תביר נמשך לפניו):}%
% ---------- פרק 29 | פסוק 31 ----------
 \textbf{לא} וְאֵ֛ת אֵ֥יל הַמִּלֻּאִ֖ים תִּקָּ֑ח וּבִשַּׁלְתָּ֥ אֶת־בְּשָׂר֖וֹ בְּמָקֹ֥ם קָדֹֽשׁ׃ \Rashi{\textbf{במקם קדש.}בחצר אהל מועד, שהשלמים הללו קדשי קדשים היו:}%
% ---------- פרק 29 | פסוק 32 ----------
 \textbf{לב} וְאָכַ֨ל אַהֲרֹ֤ן וּבָנָיו֙ אֶת־בְּשַׂ֣ר הָאַ֔יִל וְאֶת־הַלֶּ֖חֶם אֲשֶׁ֣ר בַּסָּ֑ל פֶּ֖תַח אֹ֥הֶל מוֹעֵֽד׃ \Rashi{\textbf{פתח אהל מועד.}כל החצר קרוי כן:}%
% ---------- פרק 29 | פסוק 33 ----------
 \textbf{לג} וְאָכְל֤וּ אֹתָם֙ אֲשֶׁ֣ר כֻּפַּ֣ר בָּהֶ֔ם לְמַלֵּ֥א אֶת־יָדָ֖ם לְקַדֵּ֣שׁ אֹתָ֑ם וְזָ֥ר לֹא־יֹאכַ֖ל כִּי־קֹ֥דֶשׁ הֵֽם׃ \Rashi{\textbf{ואכלו אתם.}אהרן ובניו, לפי שהם בעליהם: אשר כפר בהם. כל זרות ותעוב: למלא את ידם. באיל ולחם הללו: לקדש אתם. שעל ידי המלואים הללו נתמלאו ידיהם ונתקדשו לכהנה: כי קדש הם. קדשי קדשים; ומכאן למדנו אזהרה לזר האוכל קדשי קדשים, שנתן המקרא טעם לדבר משום דקדש הם (פסח' כ"ד):}%
% ---------- פרק 29 | פסוק 34 ----------
 \textbf{לד} וְֽאִם־יִוָּתֵ֞ר מִבְּשַׂ֧ר הַמִּלֻּאִ֛ים וּמִן־הַלֶּ֖חֶם עַד־הַבֹּ֑קֶר וְשָׂרַפְתָּ֤ אֶת־הַנּוֹתָר֙ בָּאֵ֔שׁ לֹ֥א יֵאָכֵ֖ל כִּי־קֹ֥דֶשׁ הֽוּא׃ %
% ---------- פרק 29 | פסוק 35 ----------
 \textbf{לה} וְעָשִׂ֜יתָ לְאַהֲרֹ֤ן וּלְבָנָיו֙ כָּ֔כָה כְּכֹ֥ל אֲשֶׁר־צִוִּ֖יתִי אֹתָ֑כָה שִׁבְעַ֥ת יָמִ֖ים תְּמַלֵּ֥א יָדָֽם׃ \Rashi{\textbf{ועשית לאהרן ולבניו ככה.}שנה הכתוב וכפל לעכב – שאם חסר דבר אחד מכל האמור בענין, לא נתמלאו ידיהם להיות כהנים ועבודתם פסולה: אתכה. כמו אותך: שבעת ימים תמלא וגו'. בענין הזה ובקרבנות הללו בכל יום:}%
% ---------- פרק 29 | פסוק 36 ----------
 \textbf{לו} וּפַ֨ר חַטָּ֜את תַּעֲשֶׂ֤ה לַיּוֹם֙ עַל־הַכִּפֻּרִ֔ים וְחִטֵּאתָ֙ עַל־הַמִּזְבֵּ֔חַ בְּכַפֶּרְךָ֖ עָלָ֑יו וּמָֽשַׁחְתָּ֥ אֹת֖וֹ לְקַדְּשֽׁוֹ׃ \Rashi{\textbf{על הכפרים.}בשביל הכפורים, לכפר על המזבח מכל זרות ותעוב; ולפי שנ' שבעת ימים תמלא ידם, אין לי אלא דבר הבא בשבילם, כגון האילים והלחם, אבל הבא בשביל המזבח, כגון פר שהוא לחטוי המזבח, לא שמענו, לכך הצרך מקרא זה; ומדרש ת"כ אומר כפרת המזבח הצרכה שמא התנדב איש דבר גזל במלאכת המשכן והמזבח: וחטאת. ותדכי, לשון מתנת דמים הנתונים באצבע קרוי חטוי: ומשחת אתו. בשמן המשחה, וכל המשיחות כמין כי:}%
% ---------- פרק 29 | פסוק 37 ----------
 \textbf{לז} שִׁבְעַ֣ת יָמִ֗ים תְּכַפֵּר֙ עַל־הַמִּזְבֵּ֔חַ וְקִדַּשְׁתָּ֖ אֹת֑וֹ וְהָיָ֤ה הַמִּזְבֵּ֙חַ֙ קֹ֣דֶשׁ קׇֽדָשִׁ֔ים כׇּל־הַנֹּגֵ֥עַ בַּמִּזְבֵּ֖חַ יִקְדָּֽשׁ׃ \{ס\} \Rashi{\textbf{והיה המזבח קדש.}ומה היא קדשתו? כל הנגע במזבח יקדש. אפלו קרבן פסול שעלה עליו קדשו המזבח להכשירו שלא ירד; מתוך שנאמר כל הנגע במזבח יקדש, שומע אני בין ראוי בין שאינו ראוי – כגון דבר שלא היה פסולו בקדש, כגון הרובע והנרבע ומקצה ונעבד והטרפה וכיוצא בהן – ת"ל וזה אשר תעשה הסמוך אחריו, מה עולה ראויה, אף כל ראוי שנראה לו כבר ונפסל משבא לעזרה, כגון הלן והיוצא והטמא ושנשחט במחשבת חוץ לזמנו וחוץ למקומו וכיוצא בהן (זבחים פ"ג):}%
% ---------- פרק 29 | פסוק 38 ----------
 \textbf{לח} וְזֶ֕ה אֲשֶׁ֥ר תַּעֲשֶׂ֖ה עַל־הַמִּזְבֵּ֑חַ כְּבָשִׂ֧ים בְּנֵֽי־שָׁנָ֛ה שְׁנַ֥יִם לַיּ֖וֹם תָּמִֽיד׃ %
% ---------- פרק 29 | פסוק 39 ----------
 \textbf{לט} אֶת־הַכֶּ֥בֶשׂ הָאֶחָ֖ד תַּעֲשֶׂ֣ה בַבֹּ֑קֶר וְאֵת֙ הַכֶּ֣בֶשׂ הַשֵּׁנִ֔י תַּעֲשֶׂ֖ה בֵּ֥ין הָעַרְבָּֽיִם׃ %
% ---------- פרק 29 | פסוק 40 ----------
 \textbf{מ} וְעִשָּׂרֹ֨ן סֹ֜לֶת בָּל֨וּל בְּשֶׁ֤מֶן כָּתִית֙ רֶ֣בַע הַהִ֔ין וְנֵ֕סֶךְ רְבִיעִ֥ת הַהִ֖ין יָ֑יִן לַכֶּ֖בֶשׂ הָאֶחָֽד׃ \Rashi{\textbf{ועשרן סלת.}עשירית האיפה – מ"ג ביצים וחמש ביצה: בשמן כתית. לא לחובה נאמר כתית, אלא להכשיר, לפי שנאמר כתית למאור – ומשמע למאור ולא למנחות – יכול לפסלו למנחות, ת"ל כאן כתית ולא נאמר כתית למאור, אלא למעט מנחות שאין צריך כתית, שאף הטחון בריחים כשר בהן (מנחות פ"ו): רבע ההין. שלשה לגין: ונסך. לספלים, כמו ששנינו במסכת סכה (מ"ח) – שני ספלים של כסף היו בראש המזבח ומנקבים כמין שני חטמין דקים, נותן היין לתוכו והוא מקלח ויוצא דרך החטם ונופל על גג המזבח, ומשם יורד לשיתין, במזבח בית עולמים, ובמזבח הנחשת יורד מן המזבח לארץ:}%
% ---------- פרק 29 | פסוק 41 ----------
 \textbf{מא} וְאֵת֙ הַכֶּ֣בֶשׂ הַשֵּׁנִ֔י תַּעֲשֶׂ֖ה בֵּ֣ין הָעַרְבָּ֑יִם כְּמִנְחַ֨ת הַבֹּ֤קֶר וּכְנִסְכָּהּ֙ תַּֽעֲשֶׂה־לָּ֔הּ לְרֵ֣יחַ נִיחֹ֔חַ אִשֶּׁ֖ה לַיהֹוָֽה׃ \Rashi{\textbf{לריח ניחח.}על המנחה נאמר, שמנחת נסכים כלה כליל, וסדר הקרבתם האברים בתחלה ואחר כך המנחה, שנאמר עלה ומנחה:}%
% ---------- פרק 29 | פסוק 42 ----------
 \textbf{מב} עֹלַ֤ת תָּמִיד֙ לְדֹרֹ֣תֵיכֶ֔ם פֶּ֥תַח אֹֽהֶל־מוֹעֵ֖ד לִפְנֵ֣י יְהֹוָ֑ה אֲשֶׁ֨ר אִוָּעֵ֤ד לָכֶם֙ שָׁ֔מָּה לְדַבֵּ֥ר אֵלֶ֖יךָ שָֽׁם׃ \Rashi{\textbf{תמיד.}מיום אל יום – ולא יפסיק יום בינתים: אשר אועד לכם. כשאקבע מועד לדבר אליך, שם אקבענו לבא. ויש מרבותינו למדים מכאן שמעל מזבח הנחשת היה הקב"ה מדבר עם משה משהוקם המשכן, ויש אומרים מעל הכפרת, כמו שנאמר ודברתי אתך מעל הכפרת (שמות כ"ו), ואשר אועד לכם האמור כאן אינו אמור על המזבח אלא על אהל מועד הנזכר במקרא:}%
% ---------- פרק 29 | פסוק 43 ----------
 \textbf{מג} וְנֹעַדְתִּ֥י שָׁ֖מָּה לִבְנֵ֣י יִשְׂרָאֵ֑ל וְנִקְדַּ֖שׁ בִּכְבֹדִֽי׃ \Rashi{\textbf{ונעדתי שמה.}אתועד עמם בדבור, כמלך הקובע מקום מועד לדבר עם עבדיו שם: ונקדש. המשכן: בכבדי. – שתשרה שכינתי בו. ומדרש אגדה אל תקרי בכבדי אלא בכבודי – במכבדים שלי – כאן רמז לו מיתת בני אהרן ביום הקמתו, וזהו שאמר משה הוא אשר דבר ה' לאמר בקרבי אקדש, והיכן דבר? ונקדש בכבדי (זבחים קט"ו):}%
% ---------- פרק 29 | פסוק 44 ----------
 \textbf{מד} וְקִדַּשְׁתִּ֛י אֶת־אֹ֥הֶל מוֹעֵ֖ד וְאֶת־הַמִּזְבֵּ֑חַ וְאֶת־אַהֲרֹ֧ן וְאֶת־בָּנָ֛יו אֲקַדֵּ֖שׁ לְכַהֵ֥ן לִֽי׃ %
% ---------- פרק 29 | פסוק 45 ----------
 \textbf{מה} וְשָׁ֣כַנְתִּ֔י בְּת֖וֹךְ בְּנֵ֣י יִשְׂרָאֵ֑ל וְהָיִ֥יתִי לָהֶ֖ם לֵאלֹהִֽים׃ %
% ---------- פרק 29 | פסוק 46 ----------
 \textbf{מו} וְיָדְע֗וּ כִּ֣י אֲנִ֤י יְהֹוָה֙ אֱלֹ֣הֵיהֶ֔ם אֲשֶׁ֨ר הוֹצֵ֧אתִי אֹתָ֛ם מֵאֶ֥רֶץ מִצְרַ֖יִם לְשׇׁכְנִ֣י בְתוֹכָ֑ם אֲנִ֖י יְהֹוָ֥ה אֱלֹהֵיהֶֽם׃ \{פ\} \Rashi{\textbf{לשכני בתוכם.}על מנת לשכן אני בתוכם:}%
% ---------- פרק 30 | פסוק 1 ----------
 \textbf{א} וְעָשִׂ֥יתָ מִזְבֵּ֖חַ מִקְטַ֣ר קְטֹ֑רֶת עֲצֵ֥י שִׁטִּ֖ים תַּעֲשֶׂ֥ה אֹתֽוֹ׃ \Rashi{\textbf{מקטר קטרת.}להעלות עליו קטור עשן סמים:}%
% ---------- פרק 30 | פסוק 2 ----------
 \textbf{ב} אַמָּ֨ה אׇרְכּ֜וֹ וְאַמָּ֤ה רׇחְבּוֹ֙ רָב֣וּעַ יִהְיֶ֔ה וְאַמָּתַ֖יִם קֹמָת֑וֹ מִמֶּ֖נּוּ קַרְנֹתָֽיו׃ %
% ---------- פרק 30 | פסוק 3 ----------
 \textbf{ג} וְצִפִּיתָ֨ אֹת֜וֹ זָהָ֣ב טָה֗וֹר אֶת־גַּגּ֧וֹ וְאֶת־קִירֹתָ֛יו סָבִ֖יב וְאֶת־קַרְנֹתָ֑יו וְעָשִׂ֥יתָ לּ֛וֹ זֵ֥ר זָהָ֖ב סָבִֽיב׃ \Rashi{\textbf{את גגו.}זה היה לו גג, אבל מזבח העולה לא היה לו גג אלא ממלאים חללו אדמה בכל חניתן: זר זהב. סימן הוא לכתר כהנה:}%
% ---------- פרק 30 | פסוק 4 ----------
 \textbf{ד} וּשְׁתֵּי֩ טַבְּעֹ֨ת זָהָ֜ב תַּֽעֲשֶׂה־לּ֣וֹ ׀ מִתַּ֣חַת לְזֵר֗וֹ עַ֚ל שְׁתֵּ֣י צַלְעֹתָ֔יו תַּעֲשֶׂ֖ה עַל־שְׁנֵ֣י צִדָּ֑יו וְהָיָה֙ לְבָתִּ֣ים לְבַדִּ֔ים לָשֵׂ֥את אֹת֖וֹ בָּהֵֽמָּה׃ \Rashi{\textbf{צלעתיו.}כאן הוא לשון זויות, כתרגומו, לפי שנאמר על שני צדיו – על שתי זויות שבשני צדיו: והיה. מעשה הטבעות האלה: לבתים לבדים. בית תהיה הטבעת לבד:}%
% ---------- פרק 30 | פסוק 5 ----------
 \textbf{ה} וְעָשִׂ֥יתָ אֶת־הַבַּדִּ֖ים עֲצֵ֣י שִׁטִּ֑ים וְצִפִּיתָ֥ אֹתָ֖ם זָהָֽב׃ %
% ---------- פרק 30 | פסוק 6 ----------
 \textbf{ו} וְנָתַתָּ֤ה אֹתוֹ֙ לִפְנֵ֣י הַפָּרֹ֔כֶת אֲשֶׁ֖ר עַל־אֲרֹ֣ן הָעֵדֻ֑ת לִפְנֵ֣י הַכַּפֹּ֗רֶת אֲשֶׁר֙ עַל־הָ֣עֵדֻ֔ת אֲשֶׁ֛ר אִוָּעֵ֥ד לְךָ֖ שָֽׁמָּה׃ \Rashi{\textbf{לפני הפרכת.}שמא תאמר משוך מכנגד הארון לצפון או לדרום, ת"ל לפני הכפרת מכון כנגד הארון מבחוץ:}%
% ---------- פרק 30 | פסוק 7 ----------
 \textbf{ז} וְהִקְטִ֥יר עָלָ֛יו אַהֲרֹ֖ן קְטֹ֣רֶת סַמִּ֑ים בַּבֹּ֣קֶר בַּבֹּ֗קֶר בְּהֵיטִיב֛וֹ אֶת־הַנֵּרֹ֖ת יַקְטִירֶֽנָּה׃ \Rashi{\textbf{בהיטיבו.}לשון נקוי הבזיכין של המנורה מדשן הפתילות שנשרפו בלילה, והיה מטיבן בכל בקר ובקר: הנרת. לו"שיש בלעז, וכן כל נרות האמורות במנורה, חוץ ממקום שנאמר שם העלאה שהוא לשון הדלקה:}%
% ---------- פרק 30 | פסוק 8 ----------
 \textbf{ח} וּבְהַעֲלֹ֨ת אַהֲרֹ֧ן אֶת־הַנֵּרֹ֛ת בֵּ֥ין הָעַרְבַּ֖יִם יַקְטִירֶ֑נָּה קְטֹ֧רֶת תָּמִ֛יד לִפְנֵ֥י יְהֹוָ֖ה לְדֹרֹתֵיכֶֽם׃ \Rashi{\textbf{ובהעלת.}כשידליקם להעלות להבתן: יקטירנה. בכל יום, פרס מקטיר שחרית ופרס מקטיר בין הערבים (כריתות ו'):}%
% ---------- פרק 30 | פסוק 9 ----------
 \textbf{ט} לֹא־תַעֲל֥וּ עָלָ֛יו קְטֹ֥רֶת זָרָ֖ה וְעֹלָ֣ה וּמִנְחָ֑ה וְנֵ֕סֶךְ לֹ֥א תִסְּכ֖וּ עָלָֽיו׃ \Rashi{\textbf{לא תעלו עליו.}על מזבח זה: קטרת זרה. שום קטרת של נדבה – כלן זרות לו חוץ מזו: ועלה ומנחה. ולא עולה ומנחה, עולה של בהמה ועוף, מנחה היא של לחם:}%
% ---------- פרק 30 | פסוק 10 ----------
 \textbf{י} וְכִפֶּ֤ר אַהֲרֹן֙ עַל־קַרְנֹתָ֔יו אַחַ֖ת בַּשָּׁנָ֑ה מִדַּ֞ם חַטַּ֣את הַכִּפֻּרִ֗ים אַחַ֤ת בַּשָּׁנָה֙ יְכַפֵּ֤ר עָלָיו֙ לְדֹרֹ֣תֵיכֶ֔ם קֹֽדֶשׁ־קׇדָשִׁ֥ים ה֖וּא לַיהֹוָֽה׃ \{פ\} \Rashi{\textbf{וכפר אהרן.}מתן דמים: אחת בשנה. ביום הכפורים, הוא שנאמר באחרי מות ויצא אל המזבח אשר לפני ה' וכפר עליו (ויקרא ט"ז): חטאת הכפרים. הם פר ושעיר של יום הכפורים המכפרים על טמאת מקדש וקדשיו: קדש קדשים. המזבח מקדש לדברים הללו בלבד ולא לעבודה אחרת:}%
\addparasha{חומש שמות | פרשת כי תשא}
% ---------- פרק 30 | פסוק 11 ----------
 \textbf{יא} וַיְדַבֵּ֥ר יְהֹוָ֖ה אֶל־מֹשֶׁ֥ה לֵּאמֹֽר׃ %
% ---------- פרק 30 | פסוק 12 ----------
 \textbf{יב} כִּ֣י תִשָּׂ֞א אֶת־רֹ֥אשׁ בְּנֵֽי־יִשְׂרָאֵל֮ לִפְקֻדֵיהֶם֒ וְנָ֨תְנ֜וּ אִ֣ישׁ כֹּ֧פֶר נַפְשׁ֛וֹ לַיהֹוָ֖ה בִּפְקֹ֣ד אֹתָ֑ם וְלֹא־יִהְיֶ֥ה בָהֶ֛ם נֶ֖גֶף בִּפְקֹ֥ד אֹתָֽם׃ \Rashi{\textbf{כי תשא.}לשון קבלה, כתרגומו; כשתחפץ לקבל סכום מנינם לדעת כמה הם, אל תמנם לגלגלת, אלא יתנו כל אחד מחצית השקל ותמנה את השקלים ותדע מנינם: ולא יהיה בהם נגף. שהמנין שולט בו עין הרע, והדבר בא עליהם, כמו שמצינו בימי דוד (שמואל ב כ"ד):}%
% ---------- פרק 30 | פסוק 13 ----------
 \textbf{יג} זֶ֣ה ׀ יִתְּנ֗וּ כׇּל־הָעֹבֵר֙ עַל־הַפְּקֻדִ֔ים מַחֲצִ֥ית הַשֶּׁ֖קֶל בְּשֶׁ֣קֶל הַקֹּ֑דֶשׁ עֶשְׂרִ֤ים גֵּרָה֙ הַשֶּׁ֔קֶל מַחֲצִ֣ית הַשֶּׁ֔קֶל תְּרוּמָ֖ה לַֽיהֹוָֽה׃ \Rashi{\textbf{זה יתנו.}הראה לו כמין מטבע של אש ומשקלה מחצית השקל ואומר לו כזה יתנו (תלמוד ירושלמי שק' א'): העבר על הפקדים. דרך המונין מעבירין את הנמנין זה אחר זה, וכן כל אשר יעבר תחת השבט (ויקרא כ"ז), וכן תעברנה הצאן על ידי מונה (ירמיהו ל"ג): מחצית השקל בשקל הקדש. במשקל השקל שקצבתי לך לשקל בו שקלי הקדש, כגון שקלים האמורין בפרשת ערכין ושדה אחזה: עשרים גרה השקל. עכשו פרש לך כמה הוא: גרה. לשון מעה; וכן בשמואל (א' ב'), יבוא להשתחות לו לאגורת כסף וככר לחם: עשרים גרה השקל. שהשקל השלם ארבעה זוזים, והזוז מתחלתו חמש מעות, אלא באו והוסיפו עליו שתות והעלוהו לשש מעה כסף, ומחצית השקל הזה שאמרתי לך יתנו תרומה לה' (בכורות ה'):}%
% ---------- פרק 30 | פסוק 14 ----------
 \textbf{יד} כֹּ֗ל הָעֹבֵר֙ עַל־הַפְּקֻדִ֔ים מִבֶּ֛ן עֶשְׂרִ֥ים שָׁנָ֖ה וָמָ֑עְלָה יִתֵּ֖ן תְּרוּמַ֥ת יְהֹוָֽה׃ \Rashi{\textbf{מבן עשרים שנה ומעלה.}למדך כאן שאין פחות מבן עשרים יוצא לצבא ונמנה בכלל אנשים:}%
% ---------- פרק 30 | פסוק 15 ----------
 \textbf{טו} הֶֽעָשִׁ֣יר לֹֽא־יַרְבֶּ֗ה וְהַדַּל֙ לֹ֣א יַמְעִ֔יט מִֽמַּחֲצִ֖ית הַשָּׁ֑קֶל לָתֵת֙ אֶת־תְּרוּמַ֣ת יְהֹוָ֔ה לְכַפֵּ֖ר עַל־נַפְשֹׁתֵיכֶֽם׃ \Rashi{\textbf{לכפר על נפשתיכם.}שלא תנגפו על ידי מנין; ד"א לכפר על נפשותיכם, לפי שרמז להם כאן ג' תרומות – שנכתב כאן תרומת ה' שלש פעמים, אחת תרומת אדנים שמנאן כשהתחילו בנדבת המשכן, שנתנו כל אחד ואחד מחצית השקל, ועלה למאת הככר, שנאמר וכסף פקודי העדה מאת ככר (שמות ל"ח), ומהם נעשו האדנים שנאמר ויהי מאת ככר הכסף וגו' (שם); והשנית אף היא על ידי מנין, שמנאן משהוקם המשכן, הוא המנין האמור בתחלת חמש הפקודים באחד לחדש השני בשנה השנית (במדבר א'), ונתנו כל אחד מחצית השקל, והן לקנות מהן קרבנות צבור של כל שנה ושנה, והשוו בהם עניים ועשירים, ועל אותה תרומה נאמר לכפר על נפשותיכם, שהקרבנות לכפרה הם באים; והשלישית היא תרומת המשכן כמו שנאמר כל מרים תרומת כסף ונחשת (שמות ל"ה), ולא היתה יד כלם שוה בה אלא איש מה שנדבו לבו:}%
% ---------- פרק 30 | פסוק 16 ----------
 \textbf{טז} וְלָקַחְתָּ֞ אֶת־כֶּ֣סֶף הַכִּפֻּרִ֗ים מֵאֵת֙ בְּנֵ֣י יִשְׂרָאֵ֔ל וְנָתַתָּ֣ אֹת֔וֹ עַל־עֲבֹדַ֖ת אֹ֣הֶל מוֹעֵ֑ד וְהָיָה֩ לִבְנֵ֨י יִשְׂרָאֵ֤ל לְזִכָּרוֹן֙ לִפְנֵ֣י יְהֹוָ֔ה לְכַפֵּ֖ר עַל־נַפְשֹׁתֵיכֶֽם׃ \{פ\} \Rashi{\textbf{ונתת אתו על עבדת אהל מועד.}למדת שנצטוו למנותם בתחלת נדבת המשכן, אחר מעשה העגל, מפני שנכנס בהם מגפה, כמו שנאמר ויגף ה' את העם; משל לצאן החביבה על בעליה שנפל בה דבר ומשפסק אמר לו לרועה בבקשה ממך מנה את צאני ודע כמה נותרו בהם, להודיע שהיא חביבה עליו; ואי אפשר לומר שהמנין הזה הוא האמור בחמש הפקודים, שהרי נאמר בו באחד לחדש השני והמשכן הוקם באחד לחדש הראשון, שנאמר ביום החדש הראשון באחד לחדש תקים וגו' (שמות מ'), ומהמנין הזה נעשו האדנים – משקלים שלו – שנאמר ויהי מאת ככר הכסף לצקת וגו' (שם ל"ח), הא למדת שתים היו, אחד בתחלת נדבתן אחר יום הכפורים בשנה ראשונה, ואחת בשנה שניה באיר משהוקם המשכן. ואם תאמר וכי אפשר שבשניהם היו ישראל שוים – ו' מאות אלף וג' אלפים וחמש מאות וחמשים – שהרי בכסף פקודי העדה נאמר כן, ובחמש הפקודים אף בו נאמר כן ויהיו כל הפקודים שש מאות אלף ושלשת אלפים וחמש מאות וחמשים (במדבר א'), והלא בשתי שנים היו, וא"א שלא היו בשעת מנין הראשון בני י"ט שנה שלא נמנו ובשניה נעשו בני כ'! תשובה לדבר: אצל שנות האנשים בשנה אחת נמנו, אבל למנין יציאת מצרים היו שתי שנים, לפי שליציאת מצרים מונין מניסן, כמו ששנינו במס' ראש השנה (דף ב'), ונבנה המשכן בראשונה והוקם בשניה, שנתחדשה שנה באחד בניסן, אבל שנות האנשים מנויין למנין שנות עולם, המתחילין מתשרי, נמצאו שני המנינים בשנה אחת – המנין הראשון היה בתשרי לאחר י"הכ שנתרצה המקום לישראל לסלח להם ונצטוו על המשכן, והשני באחד באיר: על עבדת אהל מועד. הן האדנים שנעשו בו:}%
% ---------- פרק 30 | פסוק 17 ----------
 \textbf{יז} וַיְדַבֵּ֥ר יְהֹוָ֖ה אֶל־מֹשֶׁ֥ה לֵּאמֹֽר׃ %
% ---------- פרק 30 | פסוק 18 ----------
 \textbf{יח} וְעָשִׂ֜יתָ כִּיּ֥וֹר נְחֹ֛שֶׁת וְכַנּ֥וֹ נְחֹ֖שֶׁת לְרׇחְצָ֑ה וְנָתַתָּ֣ אֹת֗וֹ בֵּֽין־אֹ֤הֶל מוֹעֵד֙ וּבֵ֣ין הַמִּזְבֵּ֔חַ וְנָתַתָּ֥ שָׁ֖מָּה מָֽיִם׃ \Rashi{\textbf{כיור.}כמין דוד גדולה ולה דדים המריקים בפיהם מים: וכנו. כתרגומו בסיסה, מושב מתקן לכיור: לרחצה. מוסב על הכיור: ובין המזבח. מזבח העולה, שכתוב בו שהוא לפני פתח משכן אהל מועד, והיה הכיור משוך קמעא, ועומד כנגד אויר שבין המזבח והמשכן, ואינו מפסיק כלל בינתים, משום שנאמר ואת מזבח העלה שם פתח משכן אהל מועד (שמות מ'), – כלומר מזבח לפני אהל מועד ואין כיור לפני אהל מועד, הא כיצד? משוך קמעא כלפי הדרום; כך שנויה בזבחים:}%
% ---------- פרק 30 | פסוק 19 ----------
 \textbf{יט} וְרָחֲצ֛וּ אַהֲרֹ֥ן וּבָנָ֖יו מִמֶּ֑נּוּ אֶת־יְדֵיהֶ֖ם וְאֶת־רַגְלֵיהֶֽם׃ \Rashi{\textbf{את ידיהם ואת רגליהם.}בבת אחת היה מקדש ידיו ורגליו; וכך שנינו בזבחים (דף י"ט); כיצד קדוש ידים ורגלים? מניח ידו הימנית על גבי רגלו הימנית וידו השמאלית על גבי רגלו השמאלית ומקדש:}%
% ---------- פרק 30 | פסוק 20 ----------
 \textbf{כ} בְּבֹאָ֞ם אֶל־אֹ֧הֶל מוֹעֵ֛ד יִרְחֲצוּ־מַ֖יִם וְלֹ֣א יָמֻ֑תוּ א֣וֹ בְגִשְׁתָּ֤ם אֶל־הַמִּזְבֵּ֙חַ֙ לְשָׁרֵ֔ת לְהַקְטִ֥יר אִשֶּׁ֖ה לַֽיהֹוָֽה׃ \Rashi{\textbf{בבאם אל אהל מועד.}להקטיר שחרית ובין הערבים קטרת, או להזות מדם פר כהן המשיח ושעירי ע"ז: ולא ימתו. הא אם לא ירחצו ימותו, שבתורה נאמרו כללות, ומכלל לאו אתה שומע הן: אל המזבח. החיצון, שאין כאן ביאת אהל מועד אלא בחצר:}%
% ---------- פרק 30 | פסוק 21 ----------
 \textbf{כא} וְרָחֲצ֛וּ יְדֵיהֶ֥ם וְרַגְלֵיהֶ֖ם וְלֹ֣א יָמֻ֑תוּ וְהָיְתָ֨ה לָהֶ֧ם חׇק־עוֹלָ֛ם ל֥וֹ וּלְזַרְע֖וֹ לְדֹרֹתָֽם׃ \{פ\} \Rashi{\textbf{ולא ימתו.}לחיב מיתה על המשמש במזבח ואינו רחוץ ידים ורגלים, שהמיתה הראשונה לא שמענו אלא על הנכנס להיכל:}%
% ---------- פרק 30 | פסוק 22 ----------
 \textbf{כב} וַיְדַבֵּ֥ר יְהֹוָ֖ה אֶל־מֹשֶׁ֥ה לֵּאמֹֽר׃ %
% ---------- פרק 30 | פסוק 23 ----------
 \textbf{כג} וְאַתָּ֣ה קַח־לְךָ֮ בְּשָׂמִ֣ים רֹאשׁ֒ מׇר־דְּרוֹר֙ חֲמֵ֣שׁ מֵא֔וֹת וְקִנְּמׇן־בֶּ֥שֶׂם מַחֲצִית֖וֹ חֲמִשִּׁ֣ים וּמָאתָ֑יִם וּקְנֵה־בֹ֖שֶׂם חֲמִשִּׁ֥ים וּמָאתָֽיִם׃ \Rashi{\textbf{בשמים ראש.}חשובים: וקנמן בשם. לפי שהקנמון קלפת עץ הוא, יש שהוא טוב ויש בו ריח טוב וטעם, ויש שאינו אלא כעץ, לכך הצרך לומר קנמן בשם – מן הטוב: מחציתו חמשים ומאתים. מחצית הבאתו תהא חמשים ומאתים, נמצא כלו חמש מאות, כמו שעור מר דרור, אם כן למה נאמר בו חצאין, גזרת הכתוב היא להביאו לחצאין, להרבות בו ב' הכרעות, שאין שוקלין עין בעין; וכך שנויה בכרתות: וקנה בשם. קנה של בשם; לפי שיש קנים שאינן של בשם הצרך לומר קנה בשם: חמשים ומאתים. סך משקל כלו:}%
% ---------- פרק 30 | פסוק 24 ----------
 \textbf{כד} וְקִדָּ֕ה חֲמֵ֥שׁ מֵא֖וֹת בְּשֶׁ֣קֶל הַקֹּ֑דֶשׁ וְשֶׁ֥מֶן זַ֖יִת הִֽין׃ \Rashi{\textbf{וקדה.}שם שרש עשב, ובלשון חכמים קציעה: הין. י"ב לגין; ונחלקו בו חכמי ישראל — רבי מאיר אומר בו שלקו את העקרין, אמר לו ר' יהודה והלא לסוך את העקרין אינו ספק אלא שראום במים שלא יבלעו את השמן, ואחר כך הציף עליהם השמן עד שקלט הריח, וקפחו לשמן מעל העקרין (הוריות י"א):}%
% ---------- פרק 30 | פסוק 25 ----------
 \textbf{כה} וְעָשִׂ֣יתָ אֹת֗וֹ שֶׁ֚מֶן מִשְׁחַת־קֹ֔דֶשׁ רֹ֥קַח מִרְקַ֖חַת מַעֲשֵׂ֣ה רֹקֵ֑חַ שֶׁ֥מֶן מִשְׁחַת־קֹ֖דֶשׁ יִהְיֶֽה׃ \Rashi{\textbf{רקח מרקחת.}רקח שם דבר הוא, והטעם מוכיח, שהוא למעלה, והרי הוא כמו רקח, רגע; ואינו כמו רגע הים (ישעיהו נ"א), וכמו רקע הארץ (שם מלכים ב), שהטעם למטה; וכל דבר המערב בחברו עד שזה קופח מזה או ריח או טעם קרוי מרקחת: רקח מרקחת. רקח העשוי ע"י אמנות ותערובות: מעשה רקח. שם האמן בדבר:}%
% ---------- פרק 30 | פסוק 26 ----------
 \textbf{כו} וּמָשַׁחְתָּ֥ ב֖וֹ אֶת־אֹ֣הֶל מוֹעֵ֑ד וְאֵ֖ת אֲר֥וֹן הָעֵדֻֽת׃ \Rashi{\textbf{ומשחת בו.}כל המשיחות כמין כי, חוץ משל מלכים שהן כמין נזר (כריתות ה'):}%
% ---------- פרק 30 | פסוק 27 ----------
 \textbf{כז} וְאֶת־הַשֻּׁלְחָן֙ וְאֶת־כׇּל־כֵּלָ֔יו וְאֶת־הַמְּנֹרָ֖ה וְאֶת־כֵּלֶ֑יהָ וְאֵ֖ת מִזְבַּ֥ח הַקְּטֹֽרֶת׃ %
% ---------- פרק 30 | פסוק 28 ----------
 \textbf{כח} וְאֶת־מִזְבַּ֥ח הָעֹלָ֖ה וְאֶת־כׇּל־כֵּלָ֑יו וְאֶת־הַכִּיֹּ֖ר וְאֶת־כַּנּֽוֹ׃ %
% ---------- פרק 30 | פסוק 29 ----------
 \textbf{כט} וְקִדַּשְׁתָּ֣ אֹתָ֔ם וְהָי֖וּ קֹ֣דֶשׁ קׇֽדָשִׁ֑ים כׇּל־הַנֹּגֵ֥עַ בָּהֶ֖ם יִקְדָּֽשׁ׃ \Rashi{\textbf{וקדשת אתם.}משיחה זו מקדשתם להיות קדש קדשים, ומה היא קדשתם? כל הנגע וגו': כל הראוי לכלי שרת, משנכנס לתוכו, קדוש קדשת הגוף – לפסל ביוצא ולינה וטבול יום ואינו נפדה לצאת לחלין, אבל דבר שאינו ראוי להם אין מקדשין; ושנויה היא משנה שלמה אצל מזבח: מתוך שנ' כל הנגע במזבח יקדש (שמות כ"ט), שומע אני בין ראוי בין שאינו ראוי, ת"ל כבשים, מה כבשים ראויים אף כל ראוי; כל משיחת משכן וכהנים ומלכים מתרגם לשון רבוי, לפי שאין צרך משיחתן אלא לגדלה, כי כן יסד המלך שזה חנוך גדלתן, ושאר משיחות – כמו רקיקין משוחין (שם), וראשית שמנים ימשחו (עמוס ו') – לשון ארמית בהן כלשון עברית:}%
% ---------- פרק 30 | פסוק 30 ----------
 \textbf{ל} וְאֶת־אַהֲרֹ֥ן וְאֶת־בָּנָ֖יו תִּמְשָׁ֑ח וְקִדַּשְׁתָּ֥ אֹתָ֖ם לְכַהֵ֥ן לִֽי׃ %
% ---------- פרק 30 | פסוק 31 ----------
 \textbf{לא} וְאֶל־בְּנֵ֥י יִשְׂרָאֵ֖ל תְּדַבֵּ֣ר לֵאמֹ֑ר שֶׁ֠מֶן מִשְׁחַת־קֹ֨דֶשׁ יִהְיֶ֥ה זֶ֛ה לִ֖י לְדֹרֹתֵיכֶֽם׃ \Rashi{\textbf{לדרתיכם.}מכאן למדו רבותינו לומר שכלו קים לעתיד לבא: זה. בגימטריא תריסר לגין הוו (כריתות ה'):}%
% ---------- פרק 30 | פסוק 32 ----------
 \textbf{לב} עַל־בְּשַׂ֤ר אָדָם֙ לֹ֣א יִיסָ֔ךְ וּ֨בְמַתְכֻּנְתּ֔וֹ לֹ֥א תַעֲשׂ֖וּ כָּמֹ֑הוּ קֹ֣דֶשׁ ה֔וּא קֹ֖דֶשׁ יִהְיֶ֥ה לָכֶֽם׃ \Rashi{\textbf{לא ייסך.}בשני יודי"ן, לשון לא יפעל, כמו (דברים ה טו) למען ייטב לך: על בשר אדם לא ייסך. מן השמן הזה עצמו: ובמתכנתו לא תעשו כמהו. בסכום סמניו לא תעשו אחר כמוהו במשקל סמנין הללו לפי מדת הין שמן, אבל אם פחת או רבה סממנין לפי מדת הין שמן מתר, ואף העשוי במתכנתו של זה, אין הסך ממנו חיב אלא הרוקחו: ובמתכנתו. לשון חשבון, כמו (שמות ה ח) מתכנת הלבנים, וכן במתכנתה של קטרת:}%
% ---------- פרק 30 | פסוק 33 ----------
 \textbf{לג} אִ֚ישׁ אֲשֶׁ֣ר יִרְקַ֣ח כָּמֹ֔הוּ וַאֲשֶׁ֥ר יִתֵּ֛ן מִמֶּ֖נּוּ עַל־זָ֑ר וְנִכְרַ֖ת מֵעַמָּֽיו׃ \{ס\} \Rashi{\textbf{ואשר יתן ממנו.}מאותו של משה: על זר. שאינו צרך כהנה ומלכות:}%
% ---------- פרק 30 | פסוק 34 ----------
 \textbf{לד} וַיֹּ֩אמֶר֩ יְהֹוָ֨ה אֶל־מֹשֶׁ֜ה קַח־לְךָ֣ סַמִּ֗ים נָטָ֤ף ׀ וּשְׁחֵ֙לֶת֙ וְחֶלְבְּנָ֔ה סַמִּ֖ים וּלְבֹנָ֣ה זַכָּ֑ה בַּ֥ד בְּבַ֖ד יִהְיֶֽה׃ \Rashi{\textbf{נטף.}הוא צרי, ועל שאינו אלא שרף הנוטף מעצי הקטף קרוי נטף, ובלעז גומ"א, והצרי קורין לו טרי"אקה: ושחלת. שרש בשם חלק ומצהיר כצפרן, ובלשון המשנה קרוי צפרן, וזהו שתרגם אנקלוס וטופרא: וחלבנה. בשם שריחו רע, וקורין לו גלבנא, ומנאה הכתוב בין סמני הקטרת ללמדנו שלא יקל בעינינו לצרף עמנו באגדת תעניותינו ותפלותנו את פושעי ישראל שיהיו נמנין עמנו: סמים. אחרים: ולבנה זכה. מכאן למדו רבותינו י"א סמנין נאמרו לו למשה בסיני – מעוט סמים שנים, נטף ושחלת וחלבנה ג', הרי חמשה, סמים לרבות עוד כמו אלו הרי עשרה, ולבונה הרי י"א; ואלו הן הצרי והצפרן החלבנה והלבונה מר וקציעה שבלת נרד וכרכם הרי ח' – שהשבלת ונרד אחד, שהנרד דומה לשבלת – הקשט והקלופה וקנמון הרי י"א; ברית כרשינה אינו נקטר אלא בו שפין את הצפרן ללבנה שתהא נאה (כריתות ו'): בד בבד יהיה. אלו הארבעה הנזכרים כאן יהיו שוין משקל במשקל – כמשקלו של זה כך משקלו של זה – וכן שנינו: הצרי והצפרן החלבנה והלבונה משקל שבעים שבעים מנה (שם); ולשון בד נראה בעיני שהוא לשון יחיד, אחד באחד יהיו – זה כמות זה:}%
% ---------- פרק 30 | פסוק 35 ----------
 \textbf{לה} וְעָשִׂ֤יתָ אֹתָהּ֙ קְטֹ֔רֶת רֹ֖קַח מַעֲשֵׂ֣ה רוֹקֵ֑חַ מְמֻלָּ֖ח טָה֥וֹר קֹֽדֶשׁ׃ \Rashi{\textbf{ממלח.}כתרגומו מערב, שיערב שחיקתן יפה יפה זה עם זה; ואומר אני שדומה לו וייראו המלחים (יונה א'), מלחיך וחבליך (יחזקאל כ"ז) – על שם שמהפכין את המים במשוטות כשמנהיגים את הספינה, כאדם המהפך בכף ביצים טרופות לערבן עם המים, וכל דבר שאדם רוצה לערב יפה יפה מהפכו באצבע או בבזך: ממלח טהור קדש. ממלח יהיה, וטהור יהיה וקדש יהיה:}%
% ---------- פרק 30 | פסוק 36 ----------
 \textbf{לו} וְשָֽׁחַקְתָּ֣ מִמֶּ֘נָּה֮ הָדֵק֒ וְנָתַתָּ֨ה מִמֶּ֜נָּה לִפְנֵ֤י הָעֵדֻת֙ בְּאֹ֣הֶל מוֹעֵ֔ד אֲשֶׁ֛ר אִוָּעֵ֥ד לְךָ֖ שָׁ֑מָּה קֹ֥דֶשׁ קׇֽדָשִׁ֖ים תִּהְיֶ֥ה לָכֶֽם׃ \Rashi{\textbf{ונתתה ממנה וגו'.}היא קטרת שבכל יום ויום שעל מזבח הפנימי שהוא באהל מועד: אשר אועד לך שמה. כל מועדי דבור שאקבע לך, אני קובעם לאותו מקום:}%
% ---------- פרק 30 | פסוק 37 ----------
 \textbf{לז} וְהַקְּטֹ֙רֶת֙ אֲשֶׁ֣ר תַּעֲשֶׂ֔ה בְּמַ֨תְכֻּנְתָּ֔הּ לֹ֥א תַעֲשׂ֖וּ לָכֶ֑ם קֹ֛דֶשׁ תִּהְיֶ֥ה לְךָ֖ לַיהֹוָֽה׃ \Rashi{\textbf{במתכנתה.}במנין סממניה: קדש תהיה לך לה'. שלא תעשנה אלא לשמי:}%
% ---------- פרק 30 | פסוק 38 ----------
 \textbf{לח} אִ֛ישׁ אֲשֶׁר־יַעֲשֶׂ֥ה כָמ֖וֹהָ לְהָרִ֣יחַ בָּ֑הּ וְנִכְרַ֖ת מֵעַמָּֽיו׃ \{ס\} \Rashi{\textbf{להריח בה.}אבל עושה אתה במתכנתה משלך, כדי למכרה לצבור (כריתות ה'):}%
% ---------- פרק 31 | פסוק 1 ----------
 \textbf{א} וַיְדַבֵּ֥ר יְהֹוָ֖ה אֶל־מֹשֶׁ֥ה לֵּאמֹֽר׃ %
% ---------- פרק 31 | פסוק 2 ----------
 \textbf{ב} רְאֵ֖ה קָרָ֣אתִֽי בְשֵׁ֑ם בְּצַלְאֵ֛ל בֶּן־אוּרִ֥י בֶן־ח֖וּר לְמַטֵּ֥ה יְהוּדָֽה׃ \Rashi{\textbf{קראתי בשם –.}לעשות מלאכתי – את בצלאל:}%
% ---------- פרק 31 | פסוק 3 ----------
 \textbf{ג} וָאֲמַלֵּ֥א אֹת֖וֹ ר֣וּחַ אֱלֹהִ֑ים בְּחׇכְמָ֛ה וּבִתְבוּנָ֥ה וּבְדַ֖עַת וּבְכׇל־מְלָאכָֽה׃ \Rashi{\textbf{חכמה.}מה שאדם שומע מאחרים ולמד: תבונה. מבין דבר מלבו מתוך דברים שלמד: דעת. רוח הקדש:}%
% ---------- פרק 31 | פסוק 4 ----------
 \textbf{ד} לַחְשֹׁ֖ב מַחֲשָׁבֹ֑ת לַעֲשׂ֛וֹת בַּזָּהָ֥ב וּבַכֶּ֖סֶף וּבַנְּחֹֽשֶׁת׃ \Rashi{\textbf{לחשב מחשבת.}אריגת מעשה חשב:}%
% ---------- פרק 31 | פסוק 5 ----------
 \textbf{ה} וּבַחֲרֹ֥שֶׁת אֶ֛בֶן לְמַלֹּ֖את וּבַחֲרֹ֣שֶׁת עֵ֑ץ לַעֲשׂ֖וֹת בְּכׇל־מְלָאכָֽה׃ \Rashi{\textbf{ובחרשת.}לשון אמנות, כמו חרש חכם (ישעיהו מ'), ואנקלוס פרש, ושנה בפרושן, שחרש אבנים קרוי אמן וחרש עץ קרוי נגר: למלאת. להושיבה במשבצת שלה במלואה – לעשות המשבצת למדת מושב האבן ועביה:}%
% ---------- פרק 31 | פסוק 6 ----------
 \textbf{ו} וַאֲנִ֞י הִנֵּ֧ה נָתַ֣תִּי אִתּ֗וֹ אֵ֣ת אׇהֳלִיאָ֞ב בֶּן־אֲחִֽיסָמָךְ֙ לְמַטֵּה־דָ֔ן וּבְלֵ֥ב כׇּל־חֲכַם־לֵ֖ב נָתַ֣תִּי חׇכְמָ֑ה וְעָשׂ֕וּ אֵ֖ת כׇּל־אֲשֶׁ֥ר צִוִּיתִֽךָ׃ \Rashi{\textbf{ובלב כל חכם לב וגו'.}ועוד שאר חכמי לב יש בכם, וכל אשר נתתי בו חכמה ועשו את כל אשר צויתיך:}%
% ---------- פרק 31 | פסוק 7 ----------
 \textbf{ז} אֵ֣ת ׀ אֹ֣הֶל מוֹעֵ֗ד וְאֶת־הָֽאָרֹן֙ לָֽעֵדֻ֔ת וְאֶת־הַכַּפֹּ֖רֶת אֲשֶׁ֣ר עָלָ֑יו וְאֵ֖ת כׇּל־כְּלֵ֥י הָאֹֽהֶל׃ \Rashi{\textbf{ואת הארן לעדת.}לצרך לחות העדות:}%
% ---------- פרק 31 | פסוק 8 ----------
 \textbf{ח} וְאֶת־הַשֻּׁלְחָן֙ וְאֶת־כֵּלָ֔יו וְאֶת־הַמְּנֹרָ֥ה הַטְּהֹרָ֖ה וְאֶת־כׇּל־כֵּלֶ֑יהָ וְאֵ֖ת מִזְבַּ֥ח הַקְּטֹֽרֶת׃ \Rashi{\textbf{הטהרה.}על שם זהב טהור:}%
% ---------- פרק 31 | פסוק 9 ----------
 \textbf{ט} וְאֶת־מִזְבַּ֥ח הָעֹלָ֖ה וְאֶת־כׇּל־כֵּלָ֑יו וְאֶת־הַכִּיּ֖וֹר וְאֶת־כַּנּֽוֹ׃ %
% ---------- פרק 31 | פסוק 10 ----------
 \textbf{י} וְאֵ֖ת בִּגְדֵ֣י הַשְּׂרָ֑ד וְאֶת־בִּגְדֵ֤י הַקֹּ֙דֶשׁ֙ לְאַהֲרֹ֣ן הַכֹּהֵ֔ן וְאֶת־בִּגְדֵ֥י בָנָ֖יו לְכַהֵֽן׃ \Rashi{\textbf{ואת בגדי השרד.}אומר אני – לפי פשוטו של מקרא – שאי אפשר לומר שבבגדי כהנה מדבר, לפי שנ' אצלם ואת בגדי הקדש לאהרן הכהן ואת בגדי בניו לכהן, אלא אלו בגדי השרד הם בגדי התכלת והארגמן ותולעת שני האמורין בפרשת מסעות, ונתנו אל בגד תכלת ופרשו עליו בגד ארגמן ופרשו עליהם בגד תולעת שני (במדבר ד'); ונראין דברי, שנ' ומן התכלת והארגמן ותולעת השני עשו בגדי שרד לשרת בקדש (שמות ל"ט), ולא הזכר שש עמהם, ואם בבגדי כהנה מדבר, לא מצינו באחד מהם ארגמן או תולעת שני בלא שש: בגדי השרד. יש מפרשים לשון עבודה ושרות, כתרגומו לבושי שמושא, ואין לו דמיון במקרא; ואני אומר שהוא לשון ארמי, כתרגום של קלעים ותרגום של מכבר, שהיו ארוגים במחט עשויים נקבים נקבים לצי"דץ בלעז:}%
% ---------- פרק 31 | פסוק 11 ----------
 \textbf{יא} וְאֵ֨ת שֶׁ֧מֶן הַמִּשְׁחָ֛ה וְאֶת־קְטֹ֥רֶת הַסַּמִּ֖ים לַקֹּ֑דֶשׁ כְּכֹ֥ל אֲשֶׁר־צִוִּיתִ֖ךָ יַעֲשֽׂוּ׃ \{פ\} \Rashi{\textbf{ואת קטרת הסמים לקדש.}לצרך הקטרת ההיכל שהוא קדש:}%
% ---------- פרק 31 | פסוק 12 ----------
 \textbf{יב} וַיֹּ֥אמֶר יְהֹוָ֖ה אֶל־מֹשֶׁ֥ה לֵּאמֹֽר׃ %
% ---------- פרק 31 | פסוק 13 ----------
 \textbf{יג} וְאַתָּ֞ה דַּבֵּ֨ר אֶל־בְּנֵ֤י יִשְׂרָאֵל֙ לֵאמֹ֔ר אַ֥ךְ אֶת־שַׁבְּתֹתַ֖י תִּשְׁמֹ֑רוּ כִּי֩ א֨וֹת הִ֜וא בֵּינִ֤י וּבֵֽינֵיכֶם֙ לְדֹרֹ֣תֵיכֶ֔ם לָדַ֕עַת כִּ֛י אֲנִ֥י יְהֹוָ֖ה מְקַדִּשְׁכֶֽם׃ \Rashi{\textbf{ואתה דבר אל בני ישראל.}ואתה אע"פ שהפקדתיך לצוותם על מלאכת המשכן, אל יקל בעיניך לדחות את השבת מפני אותה מלאכה: אך את שבתתי תשמרו. אע"פ שתהיו רדופין וזריזין בזריזות המלאכה, שבת אל תדחה מפניה, כל אכין ורקין מעוטין – למעט שבת ממלאכת המשכן: כי אות הוא ביני וביניכם. אות גדלה היא בינינו, שבחרתי בכם בהנחילי לכם את יום מנוחתי למנוחה: לדעת. האמות בה כי אני ה' מקדשכם:}%
% ---------- פרק 31 | פסוק 14 ----------
 \textbf{יד} וּשְׁמַרְתֶּם֙ אֶת־הַשַּׁבָּ֔ת כִּ֛י קֹ֥דֶשׁ הִ֖וא לָכֶ֑ם מְחַֽלְלֶ֙יהָ֙ מ֣וֹת יוּמָ֔ת כִּ֗י כׇּל־הָעֹשֶׂ֥ה בָהּ֙ מְלָאכָ֔ה וְנִכְרְתָ֛ה הַנֶּ֥פֶשׁ הַהִ֖וא מִקֶּ֥רֶב עַמֶּֽיהָ׃ \Rashi{\textbf{מות יומת.}אם יש עדים והתראה: ונכרתה. בלא התראה (מכילתא): מחלליה. הנוהג בה חל בקדשתה:}%
% ---------- פרק 31 | פסוק 15 ----------
 \textbf{טו} שֵׁ֣שֶׁת יָמִים֮ יֵעָשֶׂ֣ה מְלָאכָה֒ וּבַיּ֣וֹם הַשְּׁבִיעִ֗י שַׁבַּ֧ת שַׁבָּת֛וֹן קֹ֖דֶשׁ לַיהֹוָ֑ה כׇּל־הָעֹשֶׂ֧ה מְלָאכָ֛ה בְּי֥וֹם הַשַּׁבָּ֖ת מ֥וֹת יוּמָֽת׃ \Rashi{\textbf{שבת שבתון.}מנוחת מרגוע ולא מנוחת עראי: קדש לה'. שמירת קדשתה לשמי ובמצותי:}%
% ---------- פרק 31 | פסוק 16 ----------
 \textbf{טז} וְשָׁמְר֥וּ בְנֵֽי־יִשְׂרָאֵ֖ל אֶת־הַשַּׁבָּ֑ת לַעֲשׂ֧וֹת אֶת־הַשַּׁבָּ֛ת לְדֹרֹתָ֖ם בְּרִ֥ית עוֹלָֽם׃ %
% ---------- פרק 31 | פסוק 17 ----------
 \textbf{יז} בֵּינִ֗י וּבֵין֙ בְּנֵ֣י יִשְׂרָאֵ֔ל א֥וֹת הִ֖וא לְעֹלָ֑ם כִּי־שֵׁ֣שֶׁת יָמִ֗ים עָשָׂ֤ה יְהֹוָה֙ אֶת־הַשָּׁמַ֣יִם וְאֶת־הָאָ֔רֶץ וּבַיּוֹם֙ הַשְּׁבִיעִ֔י שָׁבַ֖ת וַיִּנָּפַֽשׁ׃ \{ס\} \Rashi{\textbf{וינפש.}כתרגומו ונח, וכל לשון נפש הוא לשון נפש, שמשיב נפשו ונשימתו בהרגיעו מטרח המלאכה; ומי שכתוב בו לא ייעף ולא ייגע (ישעיה מ'), וכל פעלו במאמר, הכתיב מנוחה בעצמו? לשבר האזן מה שהיא יכולה לשמע:}%
% ---------- פרק 31 | פסוק 18 ----------
 \textbf{יח} וַיִּתֵּ֣ן אֶל־מֹשֶׁ֗ה כְּכַלֹּתוֹ֙ לְדַבֵּ֤ר אִתּוֹ֙ בְּהַ֣ר סִינַ֔י שְׁנֵ֖י לֻחֹ֣ת הָעֵדֻ֑ת לֻחֹ֣ת אֶ֔בֶן כְּתֻבִ֖ים בְּאֶצְבַּ֥ע אֱלֹהִֽים׃ \Rashi{\textbf{ויתן אל משה וגו'.}אין מקדם ומאחר בתורה – מעשה העגל קדם לצווי מלאכת המשכן ימים רבים היה – שהרי בי"ז בתמוז נשתברו הלוחות, וביום הכפורים נתרצה הקב"ה לישראל, ולמחרת התחילו בנדבת המשכן והוקם באחד בניסן (תנחומא): ככלתו. ככלתו כתיב – חסר – שנמסרה לו תורה במתנה ככלה לחתן, שלא היה יכול ללמד כלה בזמן מעט כזה. דבר אחר, מה כלה מתקשטת בכ"ד קשוטין – הן האמורים בספר ישעיה – אף תלמיד חכם צריך להיות בקי בכ"ד ספרים: לדבר אתו. החקים והמשפטים שבואלה המשפטים: לדבר אתו. מלמד שהיה משה שומע מפי הגבורה וחוזרין ושונין את ההלכה שניהם יחד: לחת. לחת כתיב שהיו שתיהן שוות (שם):}%
% ---------- פרק 32 | פסוק 1 ----------
 \textbf{א} וַיַּ֣רְא הָעָ֔ם כִּֽי־בֹשֵׁ֥שׁ מֹשֶׁ֖ה לָרֶ֣דֶת מִן־הָהָ֑ר וַיִּקָּהֵ֨ל הָעָ֜ם עַֽל־אַהֲרֹ֗ן וַיֹּאמְר֤וּ אֵלָיו֙ ק֣וּם ׀ עֲשֵׂה־לָ֣נוּ אֱלֹהִ֗ים אֲשֶׁ֤ר יֵֽלְכוּ֙ לְפָנֵ֔ינוּ כִּי־זֶ֣ה ׀ מֹשֶׁ֣ה הָאִ֗ישׁ אֲשֶׁ֤ר הֶֽעֱלָ֙נוּ֙ מֵאֶ֣רֶץ מִצְרַ֔יִם לֹ֥א יָדַ֖עְנוּ מֶה־הָ֥יָה לֽוֹ׃ \Rashi{\textbf{כי בשש משה.}כתרגומו, לשון אחור, וכן בשש רכבו (שופטים ה'), ויחילו עד בוש (שם ג'); כי כשעלה משה להר אמר להם לסוף ארבעים יום אני בא בתוך שש שעות, כסבורים הם שאותו יום שעלה מן המנין הוא, והוא אמר להם שלמים – ארבעים יום ולילו עמו – ויום עליתו אין לילו עמו, שהרי בז' בסיון עלה, נמצא יום ארבעים בשבעה עשר בתמוז. בי"ו בא שטן וערבב את העולם, והראה דמות חשך ואפלה וערבוביה, לומר ודאי מת משה לכך בא ערבוביא לעולם, אמר להם מת משה, שכבר באו שש שעות ולא בא וכו' כדאיתא במסכת שבת (דף פ"ט); ואי אפשר לומר שלא טעו אלא ביום המענן בין קדם חצות בין לאחר חצות, שהרי לא ירד משה עד יום המחרת, שנאמר וישכימו ממחרת ויעלו עלת: אשר ילכו לפנינו. אלוהות הרבה אוו להם (סנהדרין ס"ג): כי זה משה האיש. כמין דמות משה הראה להם השטן, שנושאים אותו באויר רקיע השמים (שבת פ"ט): אשר העלנו מארץ מצרים. והיה מורה לנו דרך אשר נעלה בה, עתה צריכין אנו לאלוהות אשר ילכו לפנינו:}%
% ---------- פרק 32 | פסוק 2 ----------
 \textbf{ב} וַיֹּ֤אמֶר אֲלֵהֶם֙ אַהֲרֹ֔ן פָּֽרְקוּ֙ נִזְמֵ֣י הַזָּהָ֔ב אֲשֶׁר֙ בְּאׇזְנֵ֣י נְשֵׁיכֶ֔ם בְּנֵיכֶ֖ם וּבְנֹתֵיכֶ֑ם וְהָבִ֖יאוּ אֵלָֽי׃ \Rashi{\textbf{באזני נשיכם.}אמר אהרן בלבו: הנשים והילדים חסים על תכשיטיהן, שמא יתעכב הדבר ובתוך כך יבא משה, והם לא המתינו ופרקו מעל עצמן: פרקו. לשון צווי, מגזרת פרק ליחיד, כמו ברכו מגזרת ברך:}%
% ---------- פרק 32 | פסוק 3 ----------
 \textbf{ג} וַיִּתְפָּֽרְקוּ֙ כׇּל־הָעָ֔ם אֶת־נִזְמֵ֥י הַזָּהָ֖ב אֲשֶׁ֣ר בְּאׇזְנֵיהֶ֑ם וַיָּבִ֖יאוּ אֶֽל־אַהֲרֹֽן׃ \Rashi{\textbf{ויתפרקו.}לשון פריקת משא, כשנטלום מאזניהם נמצאו הם מפרקים מנזמיהם, דישקריי"ר בלעז: את נזמי. כמו מנזמי, כמו כצאתי את העיר (שמות ט') – מן העיר:}%
% ---------- פרק 32 | פסוק 4 ----------
 \textbf{ד} וַיִּקַּ֣ח מִיָּדָ֗ם וַיָּ֤צַר אֹתוֹ֙ בַּחֶ֔רֶט וַֽיַּעֲשֵׂ֖הוּ עֵ֣גֶל מַסֵּכָ֑ה וַיֹּ֣אמְר֔וּ אֵ֤לֶּה אֱלֹהֶ֙יךָ֙ יִשְׂרָאֵ֔ל אֲשֶׁ֥ר הֶעֱל֖וּךָ מֵאֶ֥רֶץ מִצְרָֽיִם׃ \Rashi{\textbf{ויצר אתו בחרט.}יש לתרגמו בשני פנים, האחד ויצר – לשון קשירה, בחרט – לשון סודר, כמו והמטפחות והחריטים (ישעיהו ג'), ויצר ככרים כסף בשני חרטים (מלכים ב ה'), והשני ויצר – לשון צורה, בחרט – כלי אמנות הצורפין שחורצין וחורתין בו צורות בזהב, כעט סופר החורת אותיות בלוחות ופנקסין, כמו וכתב עליו בחרט אנוש (ישעיהו ח'), וזהו שתרגם אנקלוס וצר יתיה בזיפא, לשון זיוף, הוא כלי אמנות שחורצין בו בזהב אותיות ושקדים, שקורין בלעז ניי"ל, ומזיפין על ידו חותמות: עגל מסכה. כיון שהשליכו לאור בכור, באו מכשפי ערב רב שעלו עמהם ממצרים ועשאוהו בכשפים; ויש אומרים מיכה היה שם, שיצא מתוך דמוסי בנין שנתמעך בו במצרים, והיה בידו שם וטס שכתב בו משה "עלה שור" "עלה שור" להעלות ארונו של יוסף מתוך נילוס – והשליכו לתוך הכור ויצא העגל (תנחומא): מסכה. לשון מתכת; ד"א קכ"ה קנטרין זהב היו בו כגימטריא של מסכה (שם): אלה אלהיך. ולא נאמר אלה אלהינו, מכאן שערב רב שעלו ממצרים הם שנקהלו על אהרן והם שעשאוהו, ואחר כך הטעו את ישראל אחריו (שם):}%
% ---------- פרק 32 | פסוק 5 ----------
 \textbf{ה} וַיַּ֣רְא אַהֲרֹ֔ן וַיִּ֥בֶן מִזְבֵּ֖חַ לְפָנָ֑יו וַיִּקְרָ֤א אַֽהֲרֹן֙ וַיֹּאמַ֔ר חַ֥ג לַיהֹוָ֖ה מָחָֽר׃ \Rashi{\textbf{וירא אהרן.}שהיה בו רוח חיים, שנאמר בתבנית שור אכל עשב (תהלים ק"ו), וראה שהצליח מעשה שטן ולא היה לו פה לדחותם לגמרי: ויבן מזבח. לדחותם: ויקרא חג לה' מחר. ולא היום, שמא יבא משה קדם שיעבדוהו, זהו פשוטו; ומדרשו, בויקרא רבה, דברים הרבה ראה אהרן – ראה חור בן אחותו שהיה מוכיחם והרגוהו, וזהו ויבן מזבח לפניו – ויבן מזבוח לפניו – ועוד ראה ואמר מוטב שיתלה בי הסרחון ולא בהם, ועוד ראה ואמר אם הם בונים את המזבח, זה מביא צרור וזה מביא אבן ונמצאת מלאכתן עשוי בבת אחת, מתוך שאני בונה אותו ואני מתעצל במלאכתי, בין כך ובין כך משה בא: חג לה'. בלבו היה לשמים, בטוח היה שיבא משה ויעבדו את המקום:}%
% ---------- פרק 32 | פסוק 6 ----------
 \textbf{ו} וַיַּשְׁכִּ֙ימוּ֙ מִֽמׇּחֳרָ֔ת וַיַּעֲל֣וּ עֹלֹ֔ת וַיַּגִּ֖שׁוּ שְׁלָמִ֑ים וַיֵּ֤שֶׁב הָעָם֙ לֶֽאֱכֹ֣ל וְשָׁת֔וֹ וַיָּקֻ֖מוּ לְצַחֵֽק׃ \{פ\} \Rashi{\textbf{וישכימו.}השטן זרזם, כדי שיחטאו: לצחק. יש במשמע הזה גלוי עריות, כמו שנ' לצחק בי (בראשית ל"ט), ושפיכות דמים, כמו שנ' יקומו נא הנערים וישחקו לפנינו (שמואל ב ב'), אף כן נהרג חור (שמות רבה):}%
% ---------- פרק 32 | פסוק 7 ----------
 \textbf{ז} וַיְדַבֵּ֥ר יְהֹוָ֖ה אֶל־מֹשֶׁ֑ה לֶךְ־רֵ֕ד כִּ֚י שִׁחֵ֣ת עַמְּךָ֔ אֲשֶׁ֥ר הֶעֱלֵ֖יתָ מֵאֶ֥רֶץ מִצְרָֽיִם׃ \Rashi{\textbf{וידבר.}לשון קשי הוא, כמו וידבר אתם קשות (בראשית מ"ב): לך רד. מגדלתך – לא נתתי לך גדלה אלא בשבילם (ברכות ל"ב), באותה שעה נתנדה משה מפי בית דין של מעלה (תנחומא): שחת עמך. שחת העם לא נאמר, אלא עמך – ערב רב שקבלת מעצמך וגירתם ולא נמלכת בי, ואמרת טוב שידבקו גרים בשכינה – הם שחתו והשחיתו (שמות רבה):}%
% ---------- פרק 32 | פסוק 8 ----------
 \textbf{ח} סָ֣רוּ מַהֵ֗ר מִן־הַדֶּ֙רֶךְ֙ אֲשֶׁ֣ר צִוִּיתִ֔ם עָשׂ֣וּ לָהֶ֔ם עֵ֖גֶל מַסֵּכָ֑ה וַיִּשְׁתַּֽחֲווּ־לוֹ֙ וַיִּזְבְּחוּ־ל֔וֹ וַיֹּ֣אמְר֔וּ אֵ֤לֶּה אֱלֹהֶ֙יךָ֙ יִשְׂרָאֵ֔ל אֲשֶׁ֥ר הֶֽעֱל֖וּךָ מֵאֶ֥רֶץ מִצְרָֽיִם׃ %
% ---------- פרק 32 | פסוק 9 ----------
 \textbf{ט} וַיֹּ֥אמֶר יְהֹוָ֖ה אֶל־מֹשֶׁ֑ה רָאִ֙יתִי֙ אֶת־הָעָ֣ם הַזֶּ֔ה וְהִנֵּ֥ה עַם־קְשֵׁה־עֹ֖רֶף הֽוּא׃ \Rashi{\textbf{קשה ערף.}מחזרין קשי ערפם לנגד מוכיחיהם וממאנים לשמע:}%
% ---------- פרק 32 | פסוק 10 ----------
 \textbf{י} וְעַתָּה֙ הַנִּ֣יחָה לִּ֔י וְיִֽחַר־אַפִּ֥י בָהֶ֖ם וַאֲכַלֵּ֑ם וְאֶֽעֱשֶׂ֥ה אוֹתְךָ֖ לְג֥וֹי גָּדֽוֹל׃ \Rashi{\textbf{הניחה לי.}עדין לא שמענו שהתפלל משה עליהם והוא אומר הניחה לי? אלא כאן פתח לו פתח והודיעו שהדבר תלוי בו – שאם יתפלל עליהם לא יכלם (ברכות ל"ב):}%
% ---------- פרק 32 | פסוק 11 ----------
 \textbf{יא} וַיְחַ֣ל מֹשֶׁ֔ה אֶת־פְּנֵ֖י יְהֹוָ֣ה אֱלֹהָ֑יו וַיֹּ֗אמֶר לָמָ֤ה יְהֹוָה֙ יֶחֱרֶ֤ה אַפְּךָ֙ בְּעַמֶּ֔ךָ אֲשֶׁ֤ר הוֹצֵ֙אתָ֙ מֵאֶ֣רֶץ מִצְרַ֔יִם בְּכֹ֥חַ גָּד֖וֹל וּבְיָ֥ד חֲזָקָֽה׃ \Rashi{\textbf{למה ה' יחרה אפך.}כלום מתקנא אלא חכם בחכם, גבור בגבור:}%
% ---------- פרק 32 | פסוק 12 ----------
 \textbf{יב} לָ֩מָּה֩ יֹאמְר֨וּ מִצְרַ֜יִם לֵאמֹ֗ר בְּרָעָ֤ה הֽוֹצִיאָם֙ לַהֲרֹ֤ג אֹתָם֙ בֶּֽהָרִ֔ים וּ֨לְכַלֹּתָ֔ם מֵעַ֖ל פְּנֵ֣י הָֽאֲדָמָ֑ה שׁ֚וּב מֵחֲר֣וֹן אַפֶּ֔ךָ וְהִנָּחֵ֥ם עַל־הָרָעָ֖ה לְעַמֶּֽךָ׃ \Rashi{\textbf{והנחם.}התעשת מחשבה אחרת – להיטיב להם על הרעה. אשר חשבת להם:}%
% ---------- פרק 32 | פסוק 13 ----------
 \textbf{יג} זְכֹ֡ר לְאַבְרָהָם֩ לְיִצְחָ֨ק וּלְיִשְׂרָאֵ֜ל עֲבָדֶ֗יךָ אֲשֶׁ֨ר נִשְׁבַּ֣עְתָּ לָהֶם֮ בָּךְ֒ וַתְּדַבֵּ֣ר אֲלֵהֶ֔ם אַרְבֶּה֙ אֶֽת־זַרְעֲכֶ֔ם כְּכוֹכְבֵ֖י הַשָּׁמָ֑יִם וְכׇל־הָאָ֨רֶץ הַזֹּ֜את אֲשֶׁ֣ר אָמַ֗רְתִּי אֶתֵּן֙ לְזַרְעֲכֶ֔ם וְנָחֲל֖וּ לְעֹלָֽם׃ \Rashi{\textbf{זכר לאברהם.}אם עברו על עשרת הדברות, אברהם אביהם נתנסה בעשרה נסיונות ועדין לא קבל שכרו, תנהו לו ויצאו עשרה בעשרה (תנחומא): לאברהם ליצחק ולישראל. אם לשרפה הם, זכר לאברהם שמסר עצמו לשרף עליך באור כשדים, אם להריגה, זכר ליצחק שפשט צוארו לעקדה, אם לגלות, זכר ליעקב שגלה לחרן, ואם אינן נצולין בזכותן, מה אתה אומר לי ואעשה אותך לגוי גדול? ואם כסא של ג' רגלים אינו עומד לפניך בשעת כעסך, קל וחמר לכסא של רגל אחת (ברכות ל"ב): אשר נשבעת להם בך. לא נשבעת להם בדבר שהוא כלה – לא בשמים ולא בארץ ולא בהרים ולא בגבעות – אלא בך, שאתה קים ושבועתך קימת לעולם, שנ' בי נשבעתי נאם ה' (בראשית כ"ב), וליצחק נאמר והקמתי את השבועה אשר נשבעתי לאברהם אביך (שם), וליעקב נאמר אני אל שדי פרה ורבה (שם ל"ה) – נשבע לו באל שדי:}%
% ---------- פרק 32 | פסוק 14 ----------
 \textbf{יד} וַיִּנָּ֖חֶם יְהֹוָ֑ה עַל־הָ֣רָעָ֔ה אֲשֶׁ֥ר דִּבֶּ֖ר לַעֲשׂ֥וֹת לְעַמּֽוֹ׃ \{פ\} %
% ---------- פרק 32 | פסוק 15 ----------
 \textbf{טו} וַיִּ֜פֶן וַיֵּ֤רֶד מֹשֶׁה֙ מִן־הָהָ֔ר וּשְׁנֵ֛י לֻחֹ֥ת הָעֵדֻ֖ת בְּיָד֑וֹ לֻחֹ֗ת כְּתֻבִים֙ מִשְּׁנֵ֣י עֶבְרֵיהֶ֔ם מִזֶּ֥ה וּמִזֶּ֖ה הֵ֥ם כְּתֻבִֽים׃ \Rashi{\textbf{משני עבריהם.}היו האותיות נקראות, ומעשה נסים הוא (שבת ק"ד):}%
% ---------- פרק 32 | פסוק 16 ----------
 \textbf{טז} וְהַ֨לֻּחֹ֔ת מַעֲשֵׂ֥ה אֱלֹהִ֖ים הֵ֑מָּה וְהַמִּכְתָּ֗ב מִכְתַּ֤ב אֱלֹהִים֙ ה֔וּא חָר֖וּת עַל־הַלֻּחֹֽת׃ \Rashi{\textbf{מעשה אלהים המה.}כמשמעו – הוא בכבודו עשאן; ד"א, כאדם האומר לחברו כל עסקיו של פלוני במלאכה פלונית, כך כל שעשועיו של הקב"ה בתורה: חרות. לשון חרת וחרט אחד הוא, שניהם לשון חקוק, אנטלי"יר בלעז:}%
% ---------- פרק 32 | פסוק 17 ----------
 \textbf{יז} וַיִּשְׁמַ֧ע יְהוֹשֻׁ֛עַ אֶת־ק֥וֹל הָעָ֖ם בְּרֵעֹ֑ה וַיֹּ֙אמֶר֙ אֶל־מֹשֶׁ֔ה ק֥וֹל מִלְחָמָ֖ה בַּֽמַּחֲנֶֽה׃ \Rashi{\textbf{ברעה.}בהריעו, שהיו מריעים ושמחים וצוחקים:}%
% ---------- פרק 32 | פסוק 18 ----------
 \textbf{יח} וַיֹּ֗אמֶר אֵ֥ין קוֹל֙ עֲנ֣וֹת גְּבוּרָ֔ה וְאֵ֥ין ק֖וֹל עֲנ֣וֹת חֲלוּשָׁ֑ה ק֣וֹל עַנּ֔וֹת אָנֹכִ֖י שֹׁמֵֽעַ׃ \Rashi{\textbf{אין קול ענות גבורה.}אין קול הזה נראה קול ענית גבורים הצועקים "נצחון" ולא קול חלשים שצועקים "וי אנוסה": קול ענות. קול חרופין וגדופין, המענין את נפש שומען כשנאמרין לו:}%
% ---------- פרק 32 | פסוק 19 ----------
 \textbf{יט} וַֽיְהִ֗י כַּאֲשֶׁ֤ר קָרַב֙ אֶל־הַֽמַּחֲנֶ֔ה וַיַּ֥רְא אֶת־הָעֵ֖גֶל וּמְחֹלֹ֑ת וַיִּֽחַר־אַ֣ף מֹשֶׁ֗ה וַיַּשְׁלֵ֤ךְ מִיָּדָו֙ אֶת־הַלֻּחֹ֔ת וַיְשַׁבֵּ֥ר אֹתָ֖ם תַּ֥חַת הָהָֽר׃ \Rashi{\textbf{וישלך מידו וגו'.}אמר: מה פסח שהוא אחד מן המצוות, אמרה תורה כל בן נכר לא יאכל בו (שמות י"ב), התורה כלה כאן וכל ישראל משמדים ואתננה להם? (שבת פ"ז): תחת ההר. לרגלי ההר:}%
% ---------- פרק 32 | פסוק 20 ----------
 \textbf{כ} וַיִּקַּ֞ח אֶת־הָעֵ֨גֶל אֲשֶׁ֤ר עָשׂוּ֙ וַיִּשְׂרֹ֣ף בָּאֵ֔שׁ וַיִּטְחַ֖ן עַ֣ד אֲשֶׁר־דָּ֑ק וַיִּ֙זֶר֙ עַל־פְּנֵ֣י הַמַּ֔יִם וַיַּ֖שְׁקְ אֶת־בְּנֵ֥י יִשְׂרָאֵֽל׃ \Rashi{\textbf{ויזר.}לשון נפוץ, וכן יזרה על נוהו גפרית (איוב י"ח), וכן כי חנם מזרה הרשת (משלי א'), שזורין בה דגן וקטניות: וישק את בני ישראל. נתכון לבדקן כסוטות; שלש מיתות נדונו שם, אם יש עדים והתראה בסיף – כמשפט אנשי עיר הנדחת שהן מרבין – עדים בלא התראה במגפה, שנאמר ויגף ה' את העם, לא עדים ולא התראה בהדרוקן, שבדקום המים וצבו בטניהם (יומא ס"ו):}%
% ---------- פרק 32 | פסוק 21 ----------
 \textbf{כא} וַיֹּ֤אמֶר מֹשֶׁה֙ אֶֽל־אַהֲרֹ֔ן מֶֽה־עָשָׂ֥ה לְךָ֖ הָעָ֣ם הַזֶּ֑ה כִּֽי־הֵבֵ֥אתָ עָלָ֖יו חֲטָאָ֥ה גְדֹלָֽה׃ \Rashi{\textbf{מה עשה לך העם.}כמה יסורים סבלת שיסרוך עד שלא תביא עליהם חטא זה:}%
% ---------- פרק 32 | פסוק 22 ----------
 \textbf{כב} וַיֹּ֣אמֶר אַהֲרֹ֔ן אַל־יִ֥חַר אַ֖ף אֲדֹנִ֑י אַתָּה֙ יָדַ֣עְתָּ אֶת־הָעָ֔ם כִּ֥י בְרָ֖ע הֽוּא׃ \Rashi{\textbf{כי ברע הוא.}בדרך רע הם הולכין תמיד, ובנסיונות לפני המקום:}%
% ---------- פרק 32 | פסוק 23 ----------
 \textbf{כג} וַיֹּ֣אמְרוּ לִ֔י עֲשֵׂה־לָ֣נוּ אֱלֹהִ֔ים אֲשֶׁ֥ר יֵלְכ֖וּ לְפָנֵ֑ינוּ כִּי־זֶ֣ה ׀ מֹשֶׁ֣ה הָאִ֗ישׁ אֲשֶׁ֤ר הֶֽעֱלָ֙נוּ֙ מֵאֶ֣רֶץ מִצְרַ֔יִם לֹ֥א יָדַ֖עְנוּ מֶה־הָ֥יָה לֽוֹ׃ %
% ---------- פרק 32 | פסוק 24 ----------
 \textbf{כד} וָאֹמַ֤ר לָהֶם֙ לְמִ֣י זָהָ֔ב הִתְפָּרָ֖קוּ וַיִּתְּנוּ־לִ֑י וָאַשְׁלִכֵ֣הוּ בָאֵ֔שׁ וַיֵּצֵ֖א הָעֵ֥גֶל הַזֶּֽה׃ \Rashi{\textbf{ואמר להם.}דבר אחד – "למי זהב" לבד – והם מהרו והתפרקו ויתנו לי: ואשלכהו באש. ולא ידעתי שיצא העגל הזה, ויצא:}%
% ---------- פרק 32 | פסוק 25 ----------
 \textbf{כה} וַיַּ֤רְא מֹשֶׁה֙ אֶת־הָעָ֔ם כִּ֥י פָרֻ֖עַ ה֑וּא כִּֽי־פְרָעֹ֣ה אַהֲרֹ֔ן לְשִׁמְצָ֖ה בְּקָמֵיהֶֽם׃ \Rashi{\textbf{פרע.}מגלה; נתגלה שמצו וקלונו, כמו ופרע את ראש האשה (במדבר ה'): לשמצה בקמיהם. להיות להם הדבר הזה לגנות בפי כל הקמים עליהם:}%
% ---------- פרק 32 | פסוק 26 ----------
 \textbf{כו} וַיַּעֲמֹ֤ד מֹשֶׁה֙ בְּשַׁ֣עַר הַֽמַּחֲנֶ֔ה וַיֹּ֕אמֶר מִ֥י לַיהֹוָ֖ה אֵלָ֑י וַיֵּאָסְפ֥וּ אֵלָ֖יו כׇּל־בְּנֵ֥י לֵוִֽי׃ \Rashi{\textbf{מי לה' אלי.}יבא אלי: כל בני לוי. מכאן שכל השבט כשר (יומא ס"ו):}%
% ---------- פרק 32 | פסוק 27 ----------
 \textbf{כז} וַיֹּ֣אמֶר לָהֶ֗ם כֹּֽה־אָמַ֤ר יְהֹוָה֙ אֱלֹהֵ֣י יִשְׂרָאֵ֔ל שִׂ֥ימוּ אִישׁ־חַרְבּ֖וֹ עַל־יְרֵכ֑וֹ עִבְר֨וּ וָשׁ֜וּבוּ מִשַּׁ֤עַר לָשַׁ֙עַר֙ בַּֽמַּחֲנֶ֔ה וְהִרְג֧וּ אִֽישׁ־אֶת־אָחִ֛יו וְאִ֥ישׁ אֶת־רֵעֵ֖הוּ וְאִ֥ישׁ אֶת־קְרֹבֽוֹ׃ \Rashi{\textbf{כה אמר וגו'.}והיכן אמר? זבח לאלהים יחרם (שמות כ"ב) – כך שנויה במכילתא: אחיו. מאמו והוא ישראל (יומא ס"ו):}%
% ---------- פרק 32 | פסוק 28 ----------
 \textbf{כח} וַיַּֽעֲשׂ֥וּ בְנֵֽי־לֵוִ֖י כִּדְבַ֣ר מֹשֶׁ֑ה וַיִּפֹּ֤ל מִן־הָעָם֙ בַּיּ֣וֹם הַה֔וּא כִּשְׁלֹ֥שֶׁת אַלְפֵ֖י אִֽישׁ׃ %
% ---------- פרק 32 | פסוק 29 ----------
 \textbf{כט} וַיֹּ֣אמֶר מֹשֶׁ֗ה מִלְא֨וּ יֶדְכֶ֤ם הַיּוֹם֙ לַֽיהֹוָ֔ה כִּ֛י אִ֥ישׁ בִּבְנ֖וֹ וּבְאָחִ֑יו וְלָתֵ֧ת עֲלֵיכֶ֛ם הַיּ֖וֹם בְּרָכָֽה׃ \Rashi{\textbf{מלאו ידכם.}אתם ההורגים אותם – בדבר זה תתחנכו להיות כהנים למקום: כי איש. מכם ימלא ידו בבנו ובאחיו:}%
% ---------- פרק 32 | פסוק 30 ----------
 \textbf{ל} וַֽיְהִי֙ מִֽמׇּחֳרָ֔ת וַיֹּ֤אמֶר מֹשֶׁה֙ אֶל־הָעָ֔ם אַתֶּ֥ם חֲטָאתֶ֖ם חֲטָאָ֣ה גְדֹלָ֑ה וְעַתָּה֙ אֶֽעֱלֶ֣ה אֶל־יְהֹוָ֔ה אוּלַ֥י אֲכַפְּרָ֖ה בְּעַ֥ד חַטַּאתְכֶֽם׃ \Rashi{\textbf{אכפרה בעד חטאתכם.}אשים כפר וקנוח וסתימה לנגד חטאתכם, להבדיל ביניכם ובין החטא:}%
% ---------- פרק 32 | פסוק 31 ----------
 \textbf{לא} וַיָּ֧שׇׁב מֹשֶׁ֛ה אֶל־יְהֹוָ֖ה וַיֹּאמַ֑ר אָ֣נָּ֗א חָטָ֞א הָעָ֤ם הַזֶּה֙ חֲטָאָ֣ה גְדֹלָ֔ה וַיַּֽעֲשׂ֥וּ לָהֶ֖ם אֱלֹהֵ֥י זָהָֽב׃ \Rashi{\textbf{אלהי זהב.}אתה הוא שגרמת להם, שהשפעת להם זהב וכל חפצם, מה יעשו שלא יחטאו? משל למלך שהיה מאכיל ומשקה את בנו ומקשטו ותולה לו כיס בצוארו ומעמידו בפתח בית זונות, מה יעשה הבן שלא יחטא:}%
% ---------- פרק 32 | פסוק 32 ----------
 \textbf{לב} וְעַתָּ֖ה אִם־תִּשָּׂ֣א חַטָּאתָ֑ם וְאִם־אַ֕יִן מְחֵ֣נִי נָ֔א מִֽסִּפְרְךָ֖ אֲשֶׁ֥ר כָּתָֽבְתָּ׃ \Rashi{\textbf{ועתה אם תשא חטאתם.}הרי טוב – איני אומר לך מחני, ואם אין מחני; וזה מקרא קצר, וכן הרבה: מספרך. מכל התורה כלה, שלא יאמרו עלי שלא הייתי כדאי לבקש עליהם רחמים:}%
% ---------- פרק 32 | פסוק 33 ----------
 \textbf{לג} וַיֹּ֥אמֶר יְהֹוָ֖ה אֶל־מֹשֶׁ֑ה מִ֚י אֲשֶׁ֣ר חָֽטָא־לִ֔י אֶמְחֶ֖נּוּ מִסִּפְרִֽי׃ %
% ---------- פרק 32 | פסוק 34 ----------
 \textbf{לד} וְעַתָּ֞ה לֵ֣ךְ ׀ נְחֵ֣ה אֶת־הָעָ֗ם אֶ֤ל אֲשֶׁר־דִּבַּ֙רְתִּי֙ לָ֔ךְ הִנֵּ֥ה מַלְאָכִ֖י יֵלֵ֣ךְ לְפָנֶ֑יךָ וּבְי֣וֹם פׇּקְדִ֔י וּפָקַדְתִּ֥י עֲלֵהֶ֖ם חַטָּאתָֽם׃ \Rashi{\textbf{אל אשר דברתי לך.}יש כאן לך אצל דבור במקום אליך, וכן לדבר לו על אדניהו (מלכים א ב׳:י״ט): הנה מלאכי. ולא אני: וביום פקדי וגו'. עתה שמעתי אליך מלכלותם יחד, ותמיד תמיד כשאפקד עליהם עונותיהם ופקדתי עליהם מעט מן העון הזה עם שאר העונות; ואין פרענות באה על ישראל שאין בה קצת מפרעון עון העגל (סנהדרין ק"ב):}%
% ---------- פרק 32 | פסוק 35 ----------
 \textbf{לה} וַיִּגֹּ֥ף יְהֹוָ֖ה אֶת־הָעָ֑ם עַ֚ל אֲשֶׁ֣ר עָשׂ֣וּ אֶת־הָעֵ֔גֶל אֲשֶׁ֥ר עָשָׂ֖ה אַהֲרֹֽן׃ \{ס\} \Rashi{\textbf{ויגף ה' את העם.}מיתה בידי שמים לעדים בלא התראה:}%
% ---------- פרק 33 | פסוק 1 ----------
 \textbf{א} וַיְדַבֵּ֨ר יְהֹוָ֤ה אֶל־מֹשֶׁה֙ לֵ֣ךְ עֲלֵ֣ה מִזֶּ֔ה אַתָּ֣ה וְהָעָ֔ם אֲשֶׁ֥ר הֶֽעֱלִ֖יתָ מֵאֶ֣רֶץ מִצְרָ֑יִם אֶל־הָאָ֗רֶץ אֲשֶׁ֣ר נִ֠שְׁבַּ֠עְתִּי לְאַבְרָהָ֨ם לְיִצְחָ֤ק וּֽלְיַעֲקֹב֙ לֵאמֹ֔ר לְזַרְעֲךָ֖ אֶתְּנֶֽנָּה׃ \Rashi{\textbf{לך עלה מזה.}ארץ ישראל גבוהה מכל הארצות, לכך נאמר עלה; ד"א, כלפי שאמר לו בשעת הכעס "לך רד" אמר לו בשעת רצון "לך עלה" (תנחומא): אתה והעם. כאן לא נאמר ועמך:}%
% ---------- פרק 33 | פסוק 2 ----------
 \textbf{ב} וְשָׁלַחְתִּ֥י לְפָנֶ֖יךָ מַלְאָ֑ךְ וְגֵֽרַשְׁתִּ֗י אֶת־הַֽכְּנַעֲנִי֙ הָֽאֱמֹרִ֔י וְהַֽחִתִּי֙ וְהַפְּרִזִּ֔י הַחִוִּ֖י וְהַיְבוּסִֽי׃ \Rashi{\textbf{וגרשתי את הכנעני וגו'.}ו' אמות הם והגרגשי עמד ופנה מפניהם מאליו (תלמוד ירושלמי שביעית ו'):}%
% ---------- פרק 33 | פסוק 3 ----------
 \textbf{ג} אֶל־אֶ֛רֶץ זָבַ֥ת חָלָ֖ב וּדְבָ֑שׁ כִּי֩ לֹ֨א אֶֽעֱלֶ֜ה בְּקִרְבְּךָ֗ כִּ֤י עַם־קְשֵׁה־עֹ֙רֶף֙ אַ֔תָּה פֶּן־אֲכֶלְךָ֖ בַּדָּֽרֶךְ׃ \Rashi{\textbf{אל ארץ זבת חלב ודבש.}אני אומר לך להעלותם: כי לא אעלה בקרבך. לכך אני אומר לך ושלחתי לפניך מלאך: כי עם קשה ערף אתה. וכששכינתי בקרבכם ואתם ממרים בי, מרבה אני עליכם זעם: אכלך. לשון כליון:}%
% ---------- פרק 33 | פסוק 4 ----------
 \textbf{ד} וַיִּשְׁמַ֣ע הָעָ֗ם אֶת־הַדָּבָ֥ר הָרָ֛ע הַזֶּ֖ה וַיִּתְאַבָּ֑לוּ וְלֹא־שָׁ֛תוּ אִ֥ישׁ עֶדְי֖וֹ עָלָֽיו׃ \Rashi{\textbf{הדבר הרע.}שאין השכינה שורה ומהלכת עמם: איש עדיו. כתרים שנתנו להם בחורב כשאמרו נעשה ונשמע (שבת פ"ח):}%
% ---------- פרק 33 | פסוק 5 ----------
 \textbf{ה} וַיֹּ֨אמֶר יְהֹוָ֜ה אֶל־מֹשֶׁ֗ה אֱמֹ֤ר אֶל־בְּנֵֽי־יִשְׂרָאֵל֙ אַתֶּ֣ם עַם־קְשֵׁה־עֹ֔רֶף רֶ֧גַע אֶחָ֛ד אֶֽעֱלֶ֥ה בְקִרְבְּךָ֖ וְכִלִּיתִ֑יךָ וְעַתָּ֗ה הוֹרֵ֤ד עֶדְיְךָ֙ מֵֽעָלֶ֔יךָ וְאֵדְעָ֖ה מָ֥ה אֶֽעֱשֶׂה־לָּֽךְ׃ \Rashi{\textbf{רגע אחד אעלה בקרבך וכליתיך.}אם אעלה בקרבך ואתם ממרים בי בקשיות ערפכם אזעם עליכם רגע אחד – שהוא שעור זעמו, שנ' חבי כמעט רגע עד יעבר זעם (ישעיהו כ"ו) – ואכלה אתכם, לפיכך טוב לכם שאשלח מלאך: ועתה. פרענות זו תלקו מיד – שתורידו עדיכם מעליכם: ואדעה מה אעשה לך. בפקדת שאר העון אני יודע מה שבלבי לעשות לך:}%
% ---------- פרק 33 | פסוק 6 ----------
 \textbf{ו} וַיִּֽתְנַצְּל֧וּ בְנֵֽי־יִשְׂרָאֵ֛ל אֶת־עֶדְיָ֖ם מֵהַ֥ר חוֹרֵֽב׃ \Rashi{\textbf{את עדים מהר חורב.}את העדי שהיה בידם מהר חורב:}%
% ---------- פרק 33 | פסוק 7 ----------
 \textbf{ז} וּמֹשֶׁה֩ יִקַּ֨ח אֶת־הָאֹ֜הֶל וְנָֽטָה־ל֣וֹ ׀ מִח֣וּץ לַֽמַּחֲנֶ֗ה הַרְחֵק֙ מִן־הַֽמַּחֲנֶ֔ה וְקָ֥רָא ל֖וֹ אֹ֣הֶל מוֹעֵ֑ד וְהָיָה֙ כׇּל־מְבַקֵּ֣שׁ יְהֹוָ֔ה יֵצֵא֙ אֶל־אֹ֣הֶל מוֹעֵ֔ד אֲשֶׁ֖ר מִח֥וּץ לַֽמַּחֲנֶֽה׃ \Rashi{\textbf{ומשה.}מאותו עון והלאה, יקח את האהל. לשון הווה הוא – לוקח אהלו ונוטהו מחוץ למחנה; אמר מנדה לרב מנדה לתלמיד (תנחומא): הרחק. אלפים אמה, כענין שנ' אך רחוק יהיה ביניכם ובניו כאלפים אמה במדה (יהושע ג'): וקרא לו. והיה קורא לו אהל מועד, הוא בית ועד למבקשי תורה: כל מבקש ה'. מכאן למבקש פני זקן כמקבל פני שכינה (תנחומא): יצא אל אהל מועד. כמו יוצא. ד"א והיה כל מבקש ה' – אפלו מלאכי השרת כשהיו שואלים מקום שכינה, חבריהם אומרים להם הרי הוא באהלו של משה:}%
% ---------- פרק 33 | פסוק 8 ----------
 \textbf{ח} וְהָיָ֗ה כְּצֵ֤את מֹשֶׁה֙ אֶל־הָאֹ֔הֶל יָק֙וּמוּ֙ כׇּל־הָעָ֔ם וְנִ֨צְּב֔וּ אִ֖ישׁ פֶּ֣תַח אׇהֳל֑וֹ וְהִבִּ֙יטוּ֙ אַחֲרֵ֣י מֹשֶׁ֔ה עַד־בֹּא֖וֹ הָאֹֽהֱלָה׃ \Rashi{\textbf{והיה.}לשון הווה: כצאת. משה מן המחנה ללכת אל האהל: יקומו כל העם. עומדים מפניו, ואין יושבין עד שנתכסה מהם: והביטו אחרי משה. לשבח – אשרי ילוד אשה שכך מבטח שהשכינה תכנס אחריו לפתח אהלו:}%
% ---------- פרק 33 | פסוק 9 ----------
 \textbf{ט} וְהָיָ֗ה כְּבֹ֤א מֹשֶׁה֙ הָאֹ֔הֱלָה יֵרֵד֙ עַמּ֣וּד הֶֽעָנָ֔ן וְעָמַ֖ד פֶּ֣תַח הָאֹ֑הֶל וְדִבֶּ֖ר עִם־מֹשֶֽׁה׃ \Rashi{\textbf{ודבר עם משה.}כמו ומדבר עם משה; תרגומו ומתמלל עם משה שהוא כבוד שכינה, כמו וישמע את הקול מדבר אליו (במדבר ז'), ואינו קורא מדבר אליו; כשהוא קורא מדבר, פתרונו הקול מדבר בינו לבין עצמו, וההדיוט שומע מאליו, וכשהוא קורא מדבר משמע שהמלך מדבר עם ההדיוט:}%
% ---------- פרק 33 | פסוק 10 ----------
 \textbf{י} וְרָאָ֤ה כׇל־הָעָם֙ אֶת־עַמּ֣וּד הֶֽעָנָ֔ן עֹמֵ֖ד פֶּ֣תַח הָאֹ֑הֶל וְקָ֤ם כׇּל־הָעָם֙ וְהִֽשְׁתַּחֲו֔וּ אִ֖ישׁ פֶּ֥תַח אׇהֳלֽוֹ׃ \Rashi{\textbf{והשתחוו.}לשכינה:}%
% ---------- פרק 33 | פסוק 11 ----------
 \textbf{יא} וְדִבֶּ֨ר יְהֹוָ֤ה אֶל־מֹשֶׁה֙ פָּנִ֣ים אֶל־פָּנִ֔ים כַּאֲשֶׁ֛ר יְדַבֵּ֥ר אִ֖ישׁ אֶל־רֵעֵ֑הוּ וְשָׁב֙ אֶל־הַֽמַּחֲנֶ֔ה וּמְשָׁ֨רְת֜וֹ יְהוֹשֻׁ֤עַ בִּן־נוּן֙ נַ֔עַר לֹ֥א יָמִ֖ישׁ מִתּ֥וֹךְ הָאֹֽהֶל׃ \{פ\} \Rashi{\textbf{ודבר ה' אל משה פנים אל פנים.}ומתמלל עם משה: ושב אל המחנה. לאחר שנדבר עמו היה משה שב אל המחנה ומלמד לזקנים מה שלמד; והדבר הזה נהג משה מיום הכפורים עד שהוקם המשכן, ולא יותר, שהרי בשבעה עשר בתמוז נשתברו הלוחות, ובי"ח שרף את העגל ודן את החוטאים, ובי"ט עלה שנאמר ויהי ממחרת ויאמר משה אל העם וגו' (שמות ל"ב), ועשה שם ארבעים יום ובקש רחמים, שנ' ואתנפל לפני ה' וגו' (דברים ט'), ובראש חדש אלול נאמר לו ועלית בבקר אל הר סיני (שמות ל"ד), לקבל לוחות האחרונות – ועשה שם מ' יום, שנאמר בהם ואנכי עמדתי בהר כימים הראשונים וגו' – מה הראשונים ברצון אף האחרונים ברצון, אמר מעתה אמצעיים היו בכעס – בעשרה בתשרי נתרצה הקב"ה לישראל בשמחה ובלב שלם, ואמר לו למשה סלחתי, ומסר לו לוחות אחרונות, וירד והתחיל לצוותם על מלאכת המשכן, ועשאוהו עד אחד בניסן, ומשהוקם לא נדבר עמו עוד אלא מאהל מועד (תנחומא): ושב אל המחנה. תרגומו ותיב למשריתא, לפי שהוא לשון הווה, וכן כל הענין, וראה כל העם – וחזן, ונצבו – וקימין, והביטו – ומסתכלין, והשתחוו – וסגדין. ומדרשו: ודבר ה' אל משה – שישוב אל המחנה, אמר לו אני בכעס ואתה בכעס, אם כן מי יקריבם? (תנחומא):}%
% ---------- פרק 33 | פסוק 12 ----------
 \textbf{יב} וַיֹּ֨אמֶר מֹשֶׁ֜ה אֶל־יְהֹוָ֗ה רְ֠אֵ֠ה אַתָּ֞ה אֹמֵ֤ר אֵלַי֙ הַ֚עַל אֶת־הָעָ֣ם הַזֶּ֔ה וְאַתָּה֙ לֹ֣א הֽוֹדַעְתַּ֔נִי אֵ֥ת אֲשֶׁר־תִּשְׁלַ֖ח עִמִּ֑י וְאַתָּ֤ה אָמַ֙רְתָּ֙ יְדַעְתִּ֣יךָֽ בְשֵׁ֔ם וְגַם־מָצָ֥אתָ חֵ֖ן בְּעֵינָֽי׃ \Rashi{\textbf{ראה אתה אמר אלי.}ראה – תן עיניך ולבך על דבריך: אתה אמר אלי וגו' ואתה לא הודעתני וגו'. ואשר אמרת לי הנה אנכי שלח מלאך, אין זו הודעה, שאין אני חפץ בה: ואתה אמרת ידעתיך בשם. הכרתיך משאר בני אדם בשם חשיבות, שהרי אמרת לי הנה אנכי בא אליך בעב הענן וגו' וגם בך יאמינו לעולם (שמות י"ט):}%
% ---------- פרק 33 | פסוק 13 ----------
 \textbf{יג} וְעַתָּ֡ה אִם־נָא֩ מָצָ֨אתִי חֵ֜ן בְּעֵינֶ֗יךָ הוֹדִעֵ֤נִי נָא֙ אֶת־דְּרָכֶ֔ךָ וְאֵדָ֣עֲךָ֔ לְמַ֥עַן אֶמְצָא־חֵ֖ן בְּעֵינֶ֑יךָ וּרְאֵ֕ה כִּ֥י עַמְּךָ֖ הַגּ֥וֹי הַזֶּֽה׃ \Rashi{\textbf{ועתה וגו'.}אם אמת שמצאתי חן בעיניך הודעני נא את דרכך. מה שכר אתה נותן למוצאי חן בעיניך: ואדעך למען אמצא חן בעיניך. ואדע בזו מדת תגמוליך – מה היא מציאת חן שמצאתי בעיניך; ופתרון למען אמצא חן – למען אכיר כמה שכר מציאת החן: וראה כי עמך הגוי הזה. שלא תאמר ואעשה אותך לגוי גדול ואת אלה תעזב, ראה כי עמך הם מקדם, ואם בהם תמאס, איני סומך על היוצאים מחלצי שיתקימו, ואת תשלום השכר שלי בעם הזה תודיעני; ורבותינו דרשו במסכת ברכות (דף ז'), ואני לישב המקראות על אפניהם ועל סדרם באתי:}%
% ---------- פרק 33 | פסוק 14 ----------
 \textbf{יד} וַיֹּאמַ֑ר פָּנַ֥י יֵלֵ֖כוּ וַהֲנִחֹ֥תִי לָֽךְ׃ \Rashi{\textbf{ויאמר פני ילכו.}כתרגומו – לא אשלח עוד מלאך, אני בעצמי אלך, כמו ופניך הולכים בקרב (שמואל ב י״ז:י״א):}%
% ---------- פרק 33 | פסוק 15 ----------
 \textbf{טו} וַיֹּ֖אמֶר אֵלָ֑יו אִם־אֵ֤ין פָּנֶ֙יךָ֙ הֹלְכִ֔ים אַֽל־תַּעֲלֵ֖נוּ מִזֶּֽה׃ \Rashi{\textbf{ויאמר אליו.}בזו אני חפץ, כי על ידי מלאך אל תעלנו מזה:}%
% ---------- פרק 33 | פסוק 16 ----------
 \textbf{טז} וּבַמֶּ֣ה ׀ יִוָּדַ֣ע אֵפ֗וֹא כִּֽי־מָצָ֨אתִי חֵ֤ן בְּעֵינֶ֙יךָ֙ אֲנִ֣י וְעַמֶּ֔ךָ הֲל֖וֹא בְּלֶכְתְּךָ֣ עִמָּ֑נוּ וְנִפְלִ֙ינוּ֙ אֲנִ֣י וְעַמְּךָ֔ מִכׇּ֨ל־הָעָ֔ם אֲשֶׁ֖ר עַל־פְּנֵ֥י הָאֲדָמָֽה׃ \{פ\} \Rashi{\textbf{ובמה יודע אפוא.}יודע מציאת החן, הלוא בלכתך עמנו. ועוד דבר אני שואל ממך, שלא תשרה שכינתך עוד על אמות העולם. ונפלינו אני ועמך. ונהיה מבדלים בדבר הזה מכל העם (ברכות ז'), כמו והפלה ה' בין מקנה ישראל וגו' (שמות ט'):}%
% ---------- פרק 33 | פסוק 17 ----------
 \textbf{יז} וַיֹּ֤אמֶר יְהֹוָה֙ אֶל־מֹשֶׁ֔ה גַּ֣ם אֶת־הַדָּבָ֥ר הַזֶּ֛ה אֲשֶׁ֥ר דִּבַּ֖רְתָּ אֶֽעֱשֶׂ֑ה כִּֽי־מָצָ֤אתָ חֵן֙ בְּעֵינַ֔י וָאֵדָעֲךָ֖ בְּשֵֽׁם׃ \Rashi{\textbf{גם את הדבר הזה.}שלא תשרה שכינתי עוד על אמות העולם אעשה, ואין דבריו של בלעם על ידי שרית השכינה, אלא "נופל וגלוי עינים" כגון ואלי דבר יגנב (איוב ד') – שומעין ע"י שליח:}%
% ---------- פרק 33 | פסוק 18 ----------
 \textbf{יח} וַיֹּאמַ֑ר הַרְאֵ֥נִי נָ֖א אֶת־כְּבֹדֶֽךָ׃ \Rashi{\textbf{ויאמר הראני נא את כבדך.}ראה משה שהיה עת רצון ודבריו מקבלים, והוסיף לשאל להראותו מראית כבודו:}%
% ---------- פרק 33 | פסוק 19 ----------
 \textbf{יט} וַיֹּ֗אמֶר אֲנִ֨י אַעֲבִ֤יר כׇּל־טוּבִי֙ עַל־פָּנֶ֔יךָ וְקָרָ֧אתִֽי בְשֵׁ֛ם יְהֹוָ֖ה לְפָנֶ֑יךָ וְחַנֹּתִי֙ אֶת־אֲשֶׁ֣ר אָחֹ֔ן וְרִחַמְתִּ֖י אֶת־אֲשֶׁ֥ר אֲרַחֵֽם׃ \Rashi{\textbf{ויאמר אני אעביר וגו'.}הגיעה שעה שתראה בכבודי מה שארשה אותך לראות, לפי שאני רוצה וצריך ללמדך סדר תפלה; שכשנצרכת לבקש רחמים על ישראל הזכרת לי זכות אבות, כסבור אתה שאם תמה זכות אבות אין עוד תקוה, אני אעביר כל מדת טובי לפניך על הצור ואתה נתון במערה: וקראתי בשם ה' לפניך. ללמדך סדר בקשת רחמים אף אם תכלה זכות אבות, וכסדר זה שאתה רואה אותי – מעטף וקורא י"ג מדות – הוי מלמד את ישראל לעשות כן, וע"י שיזכירו לפני "רחום וחנון" יהיו נענין – כי רחמי לא כלים: וחנתי את אשר אחון. אותן פעמים שארצה לחן: ורחמתי. עת שאחפץ לרחם; עד כאן לא הבטיחו אלא עתים אענה ועתים לא אענה, אבל בשעת מעשה אמר לו הנה אנכי כרת ברית (שמות ל"ד) – הבטיחו שאינן חוזרות ריקם:}%
% ---------- פרק 33 | פסוק 20 ----------
 \textbf{כ} וַיֹּ֕אמֶר לֹ֥א תוּכַ֖ל לִרְאֹ֣ת אֶת־פָּנָ֑י כִּ֛י לֹֽא־יִרְאַ֥נִי הָאָדָ֖ם וָחָֽי׃ \Rashi{\textbf{ויאמר לא תוכל וגו'.}אף כשאעביר כל טובי על פניך איני נותן לך רשות לראות את פני:}%
% ---------- פרק 33 | פסוק 21 ----------
 \textbf{כא} וַיֹּ֣אמֶר יְהֹוָ֔ה הִנֵּ֥ה מָק֖וֹם אִתִּ֑י וְנִצַּבְתָּ֖ עַל־הַצּֽוּר׃ \Rashi{\textbf{הנה מקום אתי.}בהר אשר אני מדבר עמך תמיד, יש מקום מוכן לי לצרכך שאטמינך שם שלא תזק, ומשם תראה מה שתראה, זהו פשוטו; ומדרשו, על מקום שהשכינה שם מדבר, ואומר המקום אתי, ואינו אומר אני במקום, שהקב"ה מקומו של עולם ואין עולמו מקומו (תנחומא):}%
% ---------- פרק 33 | פסוק 22 ----------
 \textbf{כב} וְהָיָה֙ בַּעֲבֹ֣ר כְּבֹדִ֔י וְשַׂמְתִּ֖יךָ בְּנִקְרַ֣ת הַצּ֑וּר וְשַׂכֹּתִ֥י כַפִּ֛י עָלֶ֖יךָ עַד־עׇבְרִֽי׃ \Rashi{\textbf{בעבר כבדי.}כשאעבר לפניך: בנקרת הצור. כמו העיני האנשים ההם תנקר (במדבר ט"ז), יקרוה ערבי נחל (משלי ל'), אני קרתי ושתיתי מים (מלכים ב' יט, כד; ישעיה לז, כה), גזרה אחת להם, נקרת הצור: כרית הצור: ושכתי כפי. מכאן שנתנה רשות למחבלים לחבל, ותרגומו "ואגין במימרי" כנוי הוא לדרך כבוד של מעלה, שאינו צריך לסוכך עליו בכף ממש:}%
% ---------- פרק 33 | פסוק 23 ----------
 \textbf{כג} וַהֲסִרֹתִי֙ אֶת־כַּפִּ֔י וְרָאִ֖יתָ אֶת־אֲחֹרָ֑י וּפָנַ֖י לֹ֥א יֵרָאֽוּ׃ \{פ\}*(ספק פרשה סתומה בכתר ארם צובה) \Rashi{\textbf{והסרתי את כפי.}ואעדי ית דברת יקרי – כשאסלק הנהגת כבודי מנגד פניך, ללכת משם ולהלן: וראית את אחרי. הראהו קשר של תפלין (ברכות ז'):}%
% ---------- פרק 34 | פסוק 1 ----------
 \textbf{א} וַיֹּ֤אמֶר יְהֹוָה֙ אֶל־מֹשֶׁ֔ה פְּסׇל־לְךָ֛ שְׁנֵֽי־לֻחֹ֥ת אֲבָנִ֖ים כָּרִאשֹׁנִ֑ים וְכָתַבְתִּי֙ עַל־הַלֻּחֹ֔ת אֶ֨ת־הַדְּבָרִ֔ים אֲשֶׁ֥ר הָי֛וּ עַל־הַלֻּחֹ֥ת הָרִאשֹׁנִ֖ים אֲשֶׁ֥ר שִׁבַּֽרְתָּ׃ \Rashi{\textbf{פסל לך.}הראהו מחצב סנפירינון מתוך אהלו ואמר לו הפסלת יהיה שלך; משם נתעשר משה הרבה (תנחומא): פסל לך. אתה שברת הראשונות, אתה פסל לך אחרות, משל למלך שהלך למדינת הים והניח ארוסתו עם השפחות, מתוך קלקול השפחות יצא עליה שם רע, עמד שושבינה וקרע כתבתה, אמר אם יאמר המלך להרגה, אמר לו עדין אינה אשתך; בדק המלך ומצא שלא היה הקלקול אלא מן השפחות, נתרצה לה, אמר לו שושבינה כתב לה כתבה אחרת, שנקרעה הראשונה, אמר לו המלך אתה קרעת אותה אתה קנה לה ניר אחר ואני אכתב לה בכתב ידי; כן המלך זה הקב"ה, השפחות אלו ערב רב, והשושבין זה משה, ארוסתו של הקב"ה ישראל, לכך נאמר פסל לך (שם):}%
% ---------- פרק 34 | פסוק 2 ----------
 \textbf{ב} וֶהְיֵ֥ה נָכ֖וֹן לַבֹּ֑קֶר וְעָלִ֤יתָ בַבֹּ֙קֶר֙ אֶל־הַ֣ר סִינַ֔י וְנִצַּבְתָּ֥ לִ֛י שָׁ֖ם עַל־רֹ֥אשׁ הָהָֽר׃ \Rashi{\textbf{נכון.}מזמן:}%
% ---------- פרק 34 | פסוק 3 ----------
 \textbf{ג} וְאִישׁ֙ לֹֽא־יַעֲלֶ֣ה עִמָּ֔ךְ וְגַם־אִ֥ישׁ אַל־יֵרָ֖א בְּכׇל־הָהָ֑ר גַּם־הַצֹּ֤אן וְהַבָּקָר֙ אַל־יִרְע֔וּ אֶל־מ֖וּל הָהָ֥ר הַהֽוּא׃ \Rashi{\textbf{ואיש לא יעלה עמך.}הראשונות ע"י שהיו בתשואות וקולות וקהלות, שלטה בהן עין רעה – אין לך יפה מן הצניעות (שם):}%
% ---------- פרק 34 | פסוק 4 ----------
 \textbf{ד} וַיִּפְסֹ֡ל שְׁנֵֽי־לֻחֹ֨ת אֲבָנִ֜ים כָּרִאשֹׁנִ֗ים וַיַּשְׁכֵּ֨ם מֹשֶׁ֤ה בַבֹּ֙קֶר֙ וַיַּ֙עַל֙ אֶל־הַ֣ר סִינַ֔י כַּאֲשֶׁ֛ר צִוָּ֥ה יְהֹוָ֖ה אֹת֑וֹ וַיִּקַּ֣ח בְּיָד֔וֹ שְׁנֵ֖י לֻחֹ֥ת אֲבָנִֽים׃ %
% ---------- פרק 34 | פסוק 5 ----------
 \textbf{ה} וַיֵּ֤רֶד יְהֹוָה֙ בֶּֽעָנָ֔ן וַיִּתְיַצֵּ֥ב עִמּ֖וֹ שָׁ֑ם וַיִּקְרָ֥א בְשֵׁ֖ם יְהֹוָֽה׃ \Rashi{\textbf{ויקרא בשם ה'.}מתרגמינן וקרא בשמא דה':}%
% ---------- פרק 34 | פסוק 6 ----------
 \textbf{ו} וַיַּעֲבֹ֨ר יְהֹוָ֥ה ׀ עַל־פָּנָיו֮ וַיִּקְרָא֒ יְהֹוָ֣ה ׀ יְהֹוָ֔ה אֵ֥ל רַח֖וּם וְחַנּ֑וּן אֶ֥רֶךְ אַפַּ֖יִם וְרַב־חֶ֥סֶד וֶאֱמֶֽת׃ \Rashi{\textbf{ה' ה'.}מדת רחמים היא, אחת קדם שיחטא, ואחת אחר שיחטא וישוב (ראש השנה י"ז): אל. אף זו מדת רחמים, וכן הוא אומר אלי אלי למה עזבתני (תהילים כ"ב) – ואין לומר למדת הדין למה עזבתני, כך מצאתי במכילתא: ארך אפים. מאריך אפו ואינו ממהר לפרע, שמא יעשה תשובה: ורב חסד. לצריכים חסד – שאין להם זכיות כל כך: ואמת. לשלם שכר טוב לעושי רצונו:}%
% ---------- פרק 34 | פסוק 7 ----------
 \textbf{ז} נֹצֵ֥ר*(בספרי תימן נֹצֵ֥ר בנו״ן רגילה) חֶ֙סֶד֙ לָאֲלָפִ֔ים נֹשֵׂ֥א עָוֺ֛ן וָפֶ֖שַׁע וְחַטָּאָ֑ה וְנַקֵּה֙ לֹ֣א יְנַקֶּ֔ה פֹּקֵ֣ד ׀ עֲוֺ֣ן אָב֗וֹת עַל־בָּנִים֙ וְעַל־בְּנֵ֣י בָנִ֔ים עַל־שִׁלֵּשִׁ֖ים וְעַל־רִבֵּעִֽים׃ \Rashi{\textbf{נצר חסד.}שהאדם עושה לפניו: לאלפים. לשני אלפים דורות: עון ופשע. עונות — אלו הזדונות, פשעים — אלו המרדים שאדם עושה להכעיס: ונקה לא ינקה. לפי פשוטו משמע שאינו מותר על העון לגמרי, אלא נפרע ממנו מעט מעט, ורבותינו דרשו מנקה הוא לשבים ולא ינקה לשאינן שבים (יומא פ"ו): פקד עון אבות על בנים. כשאוחזים מעשה אבותיהם בידיהם, שכבר פרש במקרא אחר לשנאי (שמות כ'): ועל רבעים. דור רביעי; נמצאת מדה טובה מרבה על מדת פרענות אחת לחמש מאות, שבמדה טובה הוא אומר נצר חסד לאלפים (תוספתא סוטה ד'):}%
% ---------- פרק 34 | פסוק 8 ----------
 \textbf{ח} וַיְמַהֵ֖ר מֹשֶׁ֑ה וַיִּקֹּ֥ד אַ֖רְצָה וַיִּשְׁתָּֽחוּ׃ \Rashi{\textbf{וימהר משה.}כשראה משה שכינה עוברת, ושמע קול הקריאה, מיד וישתחו:}%
% ---------- פרק 34 | פסוק 9 ----------
 \textbf{ט} וַיֹּ֡אמֶר אִם־נָא֩ מָצָ֨אתִי חֵ֤ן בְּעֵינֶ֙יךָ֙ אֲדֹנָ֔י יֵֽלֶךְ־נָ֥א אֲדֹנָ֖י בְּקִרְבֵּ֑נוּ כִּ֤י עַם־קְשֵׁה־עֹ֙רֶף֙ ה֔וּא וְסָלַחְתָּ֛ לַעֲוֺנֵ֥נוּ וּלְחַטָּאתֵ֖נוּ וּנְחַלְתָּֽנוּ׃ \Rashi{\textbf{ילך נא אדני בקרבנו.}כמו שהבטחת, מאחר שאתה נושא עון; ואם עם קשה ערף הוא וימרו בך ואמרת על זאת פן אכלך בדרך, אתה תסלח לעונינו וגומר; יש כי במקום אם: ונחלתנו. ותתננו לך לנחלה מיחדת, זו היא בקשת ונפלינו אני ועמך – שלא תשרה שכינתך על האמות:}%
% ---------- פרק 34 | פסוק 10 ----------
 \textbf{י} וַיֹּ֗אמֶר הִנֵּ֣ה אָנֹכִי֮ כֹּרֵ֣ת בְּרִית֒ נֶ֤גֶד כׇּֽל־עַמְּךָ֙ אֶעֱשֶׂ֣ה נִפְלָאֹ֔ת אֲשֶׁ֛ר לֹֽא־נִבְרְא֥וּ בְכׇל־הָאָ֖רֶץ וּבְכׇל־הַגּוֹיִ֑ם וְרָאָ֣ה כׇל־הָ֠עָ֠ם אֲשֶׁר־אַתָּ֨ה בְקִרְבּ֜וֹ אֶת־מַעֲשֵׂ֤ה יְהֹוָה֙ כִּֽי־נוֹרָ֣א ה֔וּא אֲשֶׁ֥ר אֲנִ֖י עֹשֶׂ֥ה עִמָּֽךְ׃ \Rashi{\textbf{כרת ברית.}על זאת: אעשה נפלאת. לשון ונפלינו, שתהיו מבדלים בזו מכל האמות, שלא תשרה שכינתי עליהם:}%
% ---------- פרק 34 | פסוק 11 ----------
 \textbf{יא} שְׁמׇ֨ר־לְךָ֔ אֵ֛ת אֲשֶׁ֥ר אָנֹכִ֖י מְצַוְּךָ֣ הַיּ֑וֹם הִנְנִ֧י גֹרֵ֣שׁ מִפָּנֶ֗יךָ אֶת־הָאֱמֹרִי֙ וְהַֽכְּנַעֲנִ֔י וְהַחִתִּי֙ וְהַפְּרִזִּ֔י וְהַחִוִּ֖י וְהַיְבוּסִֽי׃ \Rashi{\textbf{את האמרי וגו'.}ו' אמות יש כאן, כי הגרגשי עמד ופנה מפניהם:}%
% ---------- פרק 34 | פסוק 12 ----------
 \textbf{יב} הִשָּׁ֣מֶר לְךָ֗ פֶּן־תִּכְרֹ֤ת בְּרִית֙ לְיוֹשֵׁ֣ב הָאָ֔רֶץ אֲשֶׁ֥ר אַתָּ֖ה בָּ֣א עָלֶ֑יהָ פֶּן־יִהְיֶ֥ה לְמוֹקֵ֖שׁ בְּקִרְבֶּֽךָ׃ %
% ---------- פרק 34 | פסוק 13 ----------
 \textbf{יג} כִּ֤י אֶת־מִזְבְּחֹתָם֙ תִּתֹּצ֔וּן וְאֶת־מַצֵּבֹתָ֖ם תְּשַׁבֵּר֑וּן וְאֶת־אֲשֵׁרָ֖יו תִּכְרֹתֽוּן׃ \Rashi{\textbf{אשריו.}הוא אילן שעובדים אותו:}%
% ---------- פרק 34 | פסוק 14 ----------
 \textbf{יד} כִּ֛י לֹ֥א תִֽשְׁתַּחֲוֶ֖ה לְאֵ֣ל אַחֵ֑ר כִּ֤י יְהֹוָה֙ קַנָּ֣א שְׁמ֔וֹ אֵ֥ל קַנָּ֖א הֽוּא׃ \Rashi{\textbf{קנא שמו.}מקנא לפרע ואינו מותר, וזהו כל לשון קנאה – אוחז בנצחונו ופורע מאויביו:}%
% ---------- פרק 34 | פסוק 15 ----------
 \textbf{טו} פֶּן־תִּכְרֹ֥ת בְּרִ֖ית לְיוֹשֵׁ֣ב הָאָ֑רֶץ וְזָנ֣וּ ׀ אַחֲרֵ֣י אֱלֹֽהֵיהֶ֗ם וְזָבְחוּ֙ לֵאלֹ֣הֵיהֶ֔ם וְקָרָ֣א לְךָ֔ וְאָכַלְתָּ֖ מִזִּבְחֽוֹ׃ \Rashi{\textbf{ואכלת מזבחו.}כסבור אתה שאין ענש באכילתו, ואני מעלה עליך כמודה בעבודתם, שמתוך כך אתה בא ולוקח מבנותיו לבניך (עבודה זרה ח'):}%
% ---------- פרק 34 | פסוק 16 ----------
 \textbf{טז} וְלָקַחְתָּ֥ מִבְּנֹתָ֖יו לְבָנֶ֑יךָ וְזָנ֣וּ בְנֹתָ֗יו אַחֲרֵי֙ אֱלֹ֣הֵיהֶ֔ן וְהִזְנוּ֙ אֶת־בָּנֶ֔יךָ אַחֲרֵ֖י אֱלֹהֵיהֶֽן׃ %
% ---------- פרק 34 | פסוק 17 ----------
 \textbf{יז} אֱלֹהֵ֥י מַסֵּכָ֖ה לֹ֥א תַעֲשֶׂה־לָּֽךְ׃ %
% ---------- פרק 34 | פסוק 18 ----------
 \textbf{יח} אֶת־חַ֣ג הַמַּצּוֹת֮ תִּשְׁמֹר֒ שִׁבְעַ֨ת יָמִ֜ים תֹּאכַ֤ל מַצּוֹת֙ אֲשֶׁ֣ר צִוִּיתִ֔ךָ לְמוֹעֵ֖ד חֹ֣דֶשׁ הָאָבִ֑יב כִּ֚י בְּחֹ֣דֶשׁ הָֽאָבִ֔יב יָצָ֖אתָ מִמִּצְרָֽיִם׃ \Rashi{\textbf{חדש האביב.}חדש הבכור, שהתבואה מתבכרת בבשולה:}%
% ---------- פרק 34 | פסוק 19 ----------
 \textbf{יט} כׇּל־פֶּ֥טֶר רֶ֖חֶם לִ֑י וְכׇֽל־מִקְנְךָ֙ תִּזָּכָ֔ר פֶּ֖טֶר שׁ֥וֹר וָשֶֽׂה׃ \Rashi{\textbf{כל פטר רחם לי.}באדם: וכל מקנך וגו'. וכל מקנך אשר תזכר בפטר שור ושה – אשר יפטר זכר את רחמה: פטר. לשון פתיחה, וכן פוטר מים ראשית מדון (משלי י"ז); תי"ו של תזכר לשון נקבה היא, מוסב על היולדת:}%
% ---------- פרק 34 | פסוק 20 ----------
 \textbf{כ} וּפֶ֤טֶר חֲמוֹר֙ תִּפְדֶּ֣ה בְשֶׂ֔ה וְאִם־לֹ֥א תִפְדֶּ֖ה וַעֲרַפְתּ֑וֹ כֹּ֣ל בְּכ֤וֹר בָּנֶ֙יךָ֙ תִּפְדֶּ֔ה וְלֹֽא־יֵרָא֥וּ פָנַ֖י רֵיקָֽם׃ \Rashi{\textbf{ופטר חמור.}ולא שאר בהמה טמאה: תפדה בשה. נותן שה לכהן והוא חלין ביד כהן, ופטר חמור מתר בעבודה לבעלים: וערפתו. עורפו בקופיץ; הוא הפסיד ממון כהן, לפיכך יפסד ממונו: כל בכור בניך תפדה. חמשה סלעים פדיונו קצוב, שנאמר ופדויו מבן חדש תפדה (במדבר י"ח): ולא יראו פני רקם. לפי פשוטו של מקרא דבר בפני עצמו הוא, ואינו מוסב על הבכור – שאין במצות בכור ראית פנים – אלא אזהרה אחרת היא, וכשתעלו לרגל לראות לא יראו פני ריקם – מצוה עליכם להביא עולת ראית פנים. ולפי מדרש בריתא מקרא יתר הוא, ומפנה לגזרה שוה, ללמד על הענקתו של עבד עברי שהוא חמשה סלעים מכל מין ומין, כפדיון בכור; במסכת קדושין (דף י"ז):}%
% ---------- פרק 34 | פסוק 21 ----------
 \textbf{כא} שֵׁ֤שֶׁת יָמִים֙ תַּעֲבֹ֔ד וּבַיּ֥וֹם הַשְּׁבִיעִ֖י תִּשְׁבֹּ֑ת בֶּחָרִ֥ישׁ וּבַקָּצִ֖יר תִּשְׁבֹּֽת׃ \Rashi{\textbf{בחריש ובקציר.}למה נזכר חריש וקציר? יש מרבותינו אומרים על חריש של ערב שביעית הנכנס לשביעית, וקציר של שביעית היוצא למוצאי שביעית, ללמדך שמוסיפין מחל על הקדש, וכך משמעו: ששת ימים תעבד וביום השביעי תשבת, ועבודת ו' הימים שהתרתי לך, יש שנה שהחריש והקציר אסור, ואין צרך לומר חריש וקציר של שביעית, שהרי כבר נאמר שדך לא תזרע וגו' (ויקרא כ"ה), ויש אומרים שאינו מדבר אלא בשבת, וחריש וקציר שהזכר בו לומר לך, מה חריש רשות אף קציר רשות, יצא קציר העמר שהוא מצוה ודוחה את השבת (ראש השנה ט'):}%
% ---------- פרק 34 | פסוק 22 ----------
 \textbf{כב} וְחַ֤ג שָׁבֻעֹת֙ תַּעֲשֶׂ֣ה לְךָ֔ בִּכּוּרֵ֖י קְצִ֣יר חִטִּ֑ים וְחַג֙ הָֽאָסִ֔יף תְּקוּפַ֖ת הַשָּׁנָֽה׃ \Rashi{\textbf{בכורי קציר חטים.}שאתה מביא בו שתי הלחם מן החטים: בכורי. שהיא מנחה ראשונה הבאה מן החדש של חטים למקדש, כי מנחת העמר הבאה בפסח, מן השעורים היא (מנחות פ"ד): וחג האסיף. בזמן שאתה אוסף תבואתך מן השדה לבית; אסיפה זו לשון הכנסה לבית, כמו ואספתו אל תוך ביתך (דברים כ"ב): תקופת השנה. שהיא בחזרת השנה – בתחלת השנה הבאה: תקופת. לשון מסבה והקפה:}%
% ---------- פרק 34 | פסוק 23 ----------
 \textbf{כג} שָׁלֹ֥שׁ פְּעָמִ֖ים בַּשָּׁנָ֑ה יֵרָאֶה֙ כׇּל־זְכ֣וּרְךָ֔ אֶת־פְּנֵ֛י הָֽאָדֹ֥ן ׀ יְהֹוָ֖ה אֱלֹהֵ֥י יִשְׂרָאֵֽל׃ \Rashi{\textbf{כל זכורך.}כל הזכרים שבך; הרבה מצות בתורה נאמרו ונכפלו – ויש מהם שלש פעמים וארבע – לחיב ולענש על מנין לאוין שבהם ועל מנין עשה שבהם:}%
% ---------- פרק 34 | פסוק 24 ----------
 \textbf{כד} כִּֽי־אוֹרִ֤ישׁ גּוֹיִם֙ מִפָּנֶ֔יךָ וְהִרְחַבְתִּ֖י אֶת־גְּבֻלֶ֑ךָ וְלֹא־יַחְמֹ֥ד אִישׁ֙ אֶֽת־אַרְצְךָ֔ בַּעֲלֹֽתְךָ֗ לֵרָאוֹת֙ אֶת־פְּנֵי֙ יְהֹוָ֣ה אֱלֹהֶ֔יךָ שָׁלֹ֥שׁ פְּעָמִ֖ים בַּשָּׁנָֽה׃ \Rashi{\textbf{אוריש.}כתרגומו אתריך, וכן החל רש (דברים ב'), וכן ויורש את האמרי (במדבר כ"א) – לשון גרושין: והרחבתי את גבלך. ואתה רחוק מבית הבחירה, ואינך יכול לראות לפני תמיד, לכך אני קובע לך שלש רגלים הללו:}%
% ---------- פרק 34 | פסוק 25 ----------
 \textbf{כה} לֹֽא־תִשְׁחַ֥ט עַל־חָמֵ֖ץ דַּם־זִבְחִ֑י וְלֹא־יָלִ֣ין לַבֹּ֔קֶר זֶ֖בַח חַ֥ג הַפָּֽסַח׃ \Rashi{\textbf{לא תשחט וגו'.}לא תשחט את הפסח ועדין חמץ קים; אזהרה לשוחט או לזורק או לאחד מבני חבורה (פסח' ס"ג): ולא ילין. כתרגומו; אין לינה מועלת בראש המזבח ואין לינה אלא בעמוד השחר (זבחים פ"ז): זבח חג הפסח. אמוריו, ומכאן אתה למד לכל הקטר חלבים ואברים:}%
% ---------- פרק 34 | פסוק 26 ----------
 \textbf{כו} רֵאשִׁ֗ית בִּכּוּרֵי֙ אַדְמָ֣תְךָ֔ תָּבִ֕יא בֵּ֖ית יְהֹוָ֣ה אֱלֹהֶ֑יךָ לֹא־תְבַשֵּׁ֥ל גְּדִ֖י בַּחֲלֵ֥ב אִמּֽוֹ׃ \{פ\} \Rashi{\textbf{ראשית בכורי אדמתך.}משבעת המינין האמורים בשבח ארצך, ארץ חטה ושערה וגפן וגו' (דברים ח') – ודבש הוא דבש תמרים: לא תבשל גדי. אזהרה לבשר בחלב, ושלש פעמים כתוב בתורה, אחד לאכילה ואחד להנאה ואחד לאסור בשול (חולין קט"ו): גדי. כל ולד רך במשמע, ואף עגל וכבש; ממה שהצרך לפרש בכמה מקומות "גדי עזים" למדת שגדי סתם כל יונקים במשמע (שם קי"ג): בחלב אמו. פרט לעוף שאין לו חלב, שאין אסורו מן התורה אלא מדברי סופרים (שם):}%
% ---------- פרק 34 | פסוק 27 ----------
 \textbf{כז} וַיֹּ֤אמֶר יְהֹוָה֙ אֶל־מֹשֶׁ֔ה כְּתׇב־לְךָ֖ אֶת־הַדְּבָרִ֣ים הָאֵ֑לֶּה כִּ֞י עַל־פִּ֣י ׀ הַדְּבָרִ֣ים הָאֵ֗לֶּה כָּרַ֧תִּי אִתְּךָ֛ בְּרִ֖ית וְאֶת־יִשְׂרָאֵֽל׃ \Rashi{\textbf{את הדברים האלה.}ולא אתה רשאי לכתב תורה שבעל פה (גיטין ס'):}%
% ---------- פרק 34 | פסוק 28 ----------
 \textbf{כח} וַֽיְהִי־שָׁ֣ם עִם־יְהֹוָ֗ה אַרְבָּעִ֥ים יוֹם֙ וְאַרְבָּעִ֣ים לַ֔יְלָה לֶ֚חֶם לֹ֣א אָכַ֔ל וּמַ֖יִם לֹ֣א שָׁתָ֑ה וַיִּכְתֹּ֣ב עַל־הַלֻּחֹ֗ת אֵ֚ת דִּבְרֵ֣י הַבְּרִ֔ית עֲשֶׂ֖רֶת הַדְּבָרִֽים׃ %
% ---------- פרק 34 | פסוק 29 ----------
 \textbf{כט} וַיְהִ֗י בְּרֶ֤דֶת מֹשֶׁה֙ מֵהַ֣ר סִינַ֔י וּשְׁנֵ֨י לֻחֹ֤ת הָֽעֵדֻת֙ בְּיַד־מֹשֶׁ֔ה בְּרִדְתּ֖וֹ מִן־הָהָ֑ר וּמֹשֶׁ֣ה לֹֽא־יָדַ֗ע כִּ֥י קָרַ֛ן ע֥וֹר פָּנָ֖יו בְּדַבְּר֥וֹ אִתּֽוֹ׃ \Rashi{\textbf{ויהי ברדת משה.}כשהביא לוחות אחרונות ביום הכפורים: כי קרן. לשון קרנים, שהאור מבהיק ובולט כמין קרן; ומהיכן זכה משה לקרני ההוד? רבותינו אמרו מן המערה, שנתן הקדוש ברוך הוא ידו על פניו, שנאמר ושכתי כפי (תנחומא):}%
% ---------- פרק 34 | פסוק 30 ----------
 \textbf{ל} וַיַּ֨רְא אַהֲרֹ֜ן וְכׇל־בְּנֵ֤י יִשְׂרָאֵל֙ אֶת־מֹשֶׁ֔ה וְהִנֵּ֥ה קָרַ֖ן ע֣וֹר פָּנָ֑יו וַיִּֽירְא֖וּ מִגֶּ֥שֶׁת אֵלָֽיו׃ \Rashi{\textbf{וייראו מגשת אליו.}בא וראה כמה גדול כחה של עברה שעד שלא פשטו ידיהם בעברה מהו אומר? ומראה כבוד ה' כאש אכלת בראש ההר לעיני בני ישראל (שמות כ"ד) – ולא יראים ולא מזדעזעים, ומשעשו את העגל אף מקרני הודו של משה היו מרתיעים ומזדעזעים (ספרי):}%
% ---------- פרק 34 | פסוק 31 ----------
 \textbf{לא} וַיִּקְרָ֤א אֲלֵהֶם֙ מֹשֶׁ֔ה וַיָּשֻׁ֧בוּ אֵלָ֛יו אַהֲרֹ֥ן וְכׇל־הַנְּשִׂאִ֖ים בָּעֵדָ֑ה וַיְדַבֵּ֥ר מֹשֶׁ֖ה אֲלֵהֶֽם׃ \Rashi{\textbf{הנשאים בעדה.}כמו נשיאי העדה: וידבר משה אלהם. שליחותו של מקום, ולשון הוה הוא כל הענין הזה:}%
% ---------- פרק 34 | פסוק 32 ----------
 \textbf{לב} וְאַחֲרֵי־כֵ֥ן נִגְּשׁ֖וּ כׇּל־בְּנֵ֣י יִשְׂרָאֵ֑ל וַיְצַוֵּ֕ם אֵת֩ כׇּל־אֲשֶׁ֨ר דִּבֶּ֧ר יְהֹוָ֛ה אִתּ֖וֹ בְּהַ֥ר סִינָֽי׃ \Rashi{\textbf{ואחרי כן נגשו.}אחר שלמד לזקנים, חוזר ומלמד הפרשה או ההלכה לישראל; תנו רבנן כיצד סדר המשנה? משה היה למד מפי הגבורה, נכנס אהרן שנה לו משה פרקו, נסתלק אהרן וישב לו לשמאל משה, נכנסו בניו, שנה להם משה פרקם נסתלקו הן, ישב אלעזר לימין משה ואיתמר לשמאל אהרן, נכנסו זקנים, שנה להם משה פרקם, נסתלקו זקנים ישבו לצדדין, נכנסו כל העם, שנה להם משה פרקם; נמצא ביד כל העם אחד, ביד הזקנים שנים, ביד בני אהרן שלשה, ביד אהרן ארבעה וכו'. כדאיתא בערובין (דף נ"ד):}%
% ---------- פרק 34 | פסוק 33 ----------
 \textbf{לג} וַיְכַ֣ל מֹשֶׁ֔ה מִדַּבֵּ֖ר אִתָּ֑ם וַיִּתֵּ֥ן עַל־פָּנָ֖יו מַסְוֶֽה׃ \Rashi{\textbf{ויתן על פניו מסוה.}כתרגומו בית אפי, לשון ארמי הוא, בתלמוד (כתובות ס"ב) סוי לבה, ועוד בכתבות (דף ס'), הוה קא מסוה לאפה, לשון הבטה – היה מסתכל בה, אף כאן מסוה בגד הנתן כנגד הפרצוף ובית העינים; ולכבוד קרני ההוד – שלא יזונו הכל מהם – היה נותן המסוה כנגדן, ונוטלו בשעה שהיה מדבר עם ישראל, ובשעה שהמקום נדבר עמו עד צאתו, ובצאתו יצא בלא מסוה:}%
% ---------- פרק 34 | פסוק 34 ----------
 \textbf{לד} וּבְבֹ֨א מֹשֶׁ֜ה לִפְנֵ֤י יְהֹוָה֙ לְדַבֵּ֣ר אִתּ֔וֹ יָסִ֥יר אֶת־הַמַּסְוֶ֖ה עַד־צֵאת֑וֹ וְיָצָ֗א וְדִבֶּר֙ אֶל־בְּנֵ֣י יִשְׂרָאֵ֔ל אֵ֖ת אֲשֶׁ֥ר יְצֻוֶּֽה׃ \Rashi{\textbf{ודבר אל בני ישראל.}וראו קרני ההוד בפניו; וכשהוא מסתלק מהם …}%
% ---------- פרק 34 | פסוק 35 ----------
 \textbf{לה} וְרָא֤וּ בְנֵֽי־יִשְׂרָאֵל֙ אֶת־פְּנֵ֣י מֹשֶׁ֔ה כִּ֣י קָרַ֔ן ע֖וֹר פְּנֵ֣י מֹשֶׁ֑ה וְהֵשִׁ֨יב מֹשֶׁ֤ה אֶת־הַמַּסְוֶה֙ עַל־פָּנָ֔יו עַד־בֹּא֖וֹ לְדַבֵּ֥ר אִתּֽוֹ׃ \{ס\} \Rashi{\textbf{והשיב משה את המסוה על פניו עד באו לדבר אתו.}וכשבא לדבר אתו נוטלו מעל פניו:}%
\addparasha{חומש שמות | פרשת ויקהל}
% ---------- פרק 35 | פסוק 1 ----------
 \textbf{א} וַיַּקְהֵ֣ל מֹשֶׁ֗ה אֶֽת־כׇּל־עֲדַ֛ת בְּנֵ֥י יִשְׂרָאֵ֖ל וַיֹּ֣אמֶר אֲלֵהֶ֑ם אֵ֚לֶּה הַדְּבָרִ֔ים אֲשֶׁר־צִוָּ֥ה יְהֹוָ֖ה לַעֲשֹׂ֥ת אֹתָֽם׃ \Rashi{\textbf{ויקהל משה.}למחרת יום הכפורים כשירד מן ההר, והוא לשון הפעיל, שאינו אוסף אנשים בידים, אלא הן נאספין על פי דבורו, ותרגומו ואכניש:}%
% ---------- פרק 35 | פסוק 2 ----------
 \textbf{ב} שֵׁ֣שֶׁת יָמִים֮ תֵּעָשֶׂ֣ה מְלָאכָה֒ וּבַיּ֣וֹם הַשְּׁבִיעִ֗י יִהְיֶ֨ה לָכֶ֥ם קֹ֛דֶשׁ שַׁבַּ֥ת שַׁבָּת֖וֹן לַיהֹוָ֑ה כׇּל־הָעֹשֶׂ֥ה ב֛וֹ מְלָאכָ֖ה יוּמָֽת׃ \Rashi{\textbf{ששת ימים.}הקדים להם אזהרת שבת לצווי מלאכת המשכן, לומר שאינו דוחה את השבת (מכילתא):}%
% ---------- פרק 35 | פסוק 3 ----------
 \textbf{ג} לֹא־תְבַעֲר֣וּ אֵ֔שׁ בְּכֹ֖ל מֹשְׁבֹֽתֵיכֶ֑ם בְּי֖וֹם הַשַּׁבָּֽת׃ \{פ\} \Rashi{\textbf{לא תבערו אש.}יש מרבותינו אומרים הבערה ללאו יצאת, ויש אומרים לחלק יצאת (שבת ע'):}%
% ---------- פרק 35 | פסוק 4 ----------
 \textbf{ד} וַיֹּ֣אמֶר מֹשֶׁ֔ה אֶל־כׇּל־עֲדַ֥ת בְּנֵֽי־יִשְׂרָאֵ֖ל לֵאמֹ֑ר זֶ֣ה הַדָּבָ֔ר אֲשֶׁר־צִוָּ֥ה יְהֹוָ֖ה לֵאמֹֽר׃ \Rashi{\textbf{זה הדבר אשר צוה ה'.}לי לאמר לכם:}%
% ---------- פרק 35 | פסוק 5 ----------
 \textbf{ה} קְח֨וּ מֵֽאִתְּכֶ֤ם תְּרוּמָה֙ לַֽיהֹוָ֔ה כֹּ֚ל נְדִ֣יב לִבּ֔וֹ יְבִיאֶ֕הָ אֵ֖ת תְּרוּמַ֣ת יְהֹוָ֑ה זָהָ֥ב וָכֶ֖סֶף וּנְחֹֽשֶׁת׃ \Rashi{\textbf{נדיב לבו.}על שם שלבו נודבו, קרוי נדיב לב. כבר פרשתי נדבת המשכן ומלאכתו במקום צואתן:}%
% ---------- פרק 35 | פסוק 6 ----------
 \textbf{ו} וּתְכֵ֧לֶת וְאַרְגָּמָ֛ן וְתוֹלַ֥עַת שָׁנִ֖י וְשֵׁ֥שׁ וְעִזִּֽים׃ %
% ---------- פרק 35 | פסוק 7 ----------
 \textbf{ז} וְעֹרֹ֨ת אֵילִ֧ם מְאׇדָּמִ֛ים וְעֹרֹ֥ת תְּחָשִׁ֖ים וַעֲצֵ֥י שִׁטִּֽים׃ %
% ---------- פרק 35 | פסוק 8 ----------
 \textbf{ח} וְשֶׁ֖מֶן לַמָּא֑וֹר וּבְשָׂמִים֙ לְשֶׁ֣מֶן הַמִּשְׁחָ֔ה וְלִקְטֹ֖רֶת הַסַּמִּֽים׃ %
% ---------- פרק 35 | פסוק 9 ----------
 \textbf{ט} וְאַ֨בְנֵי־שֹׁ֔הַם וְאַבְנֵ֖י מִלֻּאִ֑ים לָאֵפ֖וֹד וְלַחֹֽשֶׁן׃ %
% ---------- פרק 35 | פסוק 10 ----------
 \textbf{י} וְכׇל־חֲכַם־לֵ֖ב בָּכֶ֑ם יָבֹ֣אוּ וְיַעֲשׂ֔וּ אֵ֛ת כׇּל־אֲשֶׁ֥ר צִוָּ֖ה יְהֹוָֽה׃ %
% ---------- פרק 35 | פסוק 11 ----------
 \textbf{יא} אֶ֨ת־הַמִּשְׁכָּ֔ן אֶֽת־אׇהֳל֖וֹ וְאֶת־מִכְסֵ֑הוּ אֶת־קְרָסָיו֙ וְאֶת־קְרָשָׁ֔יו אֶת־בְּרִיחָ֕ו אֶת־עַמֻּדָ֖יו וְאֶת־אֲדָנָֽיו׃ \Rashi{\textbf{את המשכן.}יריעות התחתונות הנראות בתוכו קרויים משכן: את אהלו. היא אהל יריעות עזים, העשוי לגג: ואת מכסהו. מכסה עורות אילים והתחשים:}%
% ---------- פרק 35 | פסוק 12 ----------
 \textbf{יב} אֶת־הָאָרֹ֥ן וְאֶת־בַּדָּ֖יו אֶת־הַכַּפֹּ֑רֶת וְאֵ֖ת פָּרֹ֥כֶת הַמָּסָֽךְ׃ \Rashi{\textbf{ואת פרכת המסך.}פרכת המחיצה; כל דבר המגן, בין מלמעלה בין מכנגד, קרוי מסך וסכך, וכן סכת בעדו (איוב א'), הנני סך את דרכך (הושע ב'):}%
% ---------- פרק 35 | פסוק 13 ----------
 \textbf{יג} אֶת־הַשֻּׁלְחָ֥ן וְאֶת־בַּדָּ֖יו וְאֶת־כׇּל־כֵּלָ֑יו וְאֵ֖ת לֶ֥חֶם הַפָּנִֽים׃ \Rashi{\textbf{לחם הפנים.}כבר פרשתי, על שם שהיו לו פנים לכאן ולכאן, שהוא עשוי כמין תבה פרוצה:}%
% ---------- פרק 35 | פסוק 14 ----------
 \textbf{יד} וְאֶת־מְנֹרַ֧ת הַמָּא֛וֹר וְאֶת־כֵּלֶ֖יהָ וְאֶת־נֵרֹתֶ֑יהָ וְאֵ֖ת שֶׁ֥מֶן הַמָּאֽוֹר׃ \Rashi{\textbf{ואת כליה.}מלקחים ומחתות: נרתיה. לוצי"יש בלעז, בזיכים שהשמן והפתילות נתונין בהן: ואת שמן המאור. אף הוא צריך חכמי לב, שהוא משנה משאר שמנים, כמו שמפרש במנחות (דף פ"ו), מגרגרו בראש הזית, והוא כתית וזך:}%
% ---------- פרק 35 | פסוק 15 ----------
 \textbf{טו} וְאֶת־מִזְבַּ֤ח הַקְּטֹ֙רֶת֙ וְאֶת־בַּדָּ֔יו וְאֵת֙ שֶׁ֣מֶן הַמִּשְׁחָ֔ה וְאֵ֖ת קְטֹ֣רֶת הַסַּמִּ֑ים וְאֶת־מָסַ֥ךְ הַפֶּ֖תַח לְפֶ֥תַח הַמִּשְׁכָּֽן׃ \Rashi{\textbf{מסך הפתח.}וילון שלפני המזרח, שלא היו שם קרשים ולא יריעות:}%
% ---------- פרק 35 | פסוק 16 ----------
 \textbf{טז} אֵ֣ת ׀ מִזְבַּ֣ח הָעֹלָ֗ה וְאֶת־מִכְבַּ֤ר הַנְּחֹ֙שֶׁת֙ אֲשֶׁר־ל֔וֹ אֶת־בַּדָּ֖יו וְאֶת־כׇּל־כֵּלָ֑יו אֶת־הַכִּיֹּ֖ר וְאֶת־כַּנּֽוֹ׃ %
% ---------- פרק 35 | פסוק 17 ----------
 \textbf{יז} אֵ֚ת קַלְעֵ֣י הֶחָצֵ֔ר אֶת־עַמֻּדָ֖יו וְאֶת־אֲדָנֶ֑יהָ וְאֵ֕ת מָסַ֖ךְ שַׁ֥עַר הֶחָצֵֽר׃ \Rashi{\textbf{את עמדיו ואת אדניה.}הרי חצר קרוי כאן לשון זכר ולשון נקבה, וכן דברים הרבה: ואת מסך שער החצר. וילון פרוש לצד המזרח עשרים אמה אמצעיות של רחב החצר – שהיה חמשים רחב – וסתומין הימנו לצד צפון ט"ו אמה, וכן לדרום, שנאמר, וחמש עשרה אמה קלעים לכתף:}%
% ---------- פרק 35 | פסוק 18 ----------
 \textbf{יח} אֶת־יִתְדֹ֧ת הַמִּשְׁכָּ֛ן וְאֶת־יִתְדֹ֥ת הֶחָצֵ֖ר וְאֶת־מֵיתְרֵיהֶֽם׃ \Rashi{\textbf{יתדת.}לתקע ולקשור בהם סופי היריעות בארץ, שלא ינועו ברוח: מיתריהם. חבלים לקשר:}%
% ---------- פרק 35 | פסוק 19 ----------
 \textbf{יט} אֶת־בִּגְדֵ֥י הַשְּׂרָ֖ד לְשָׁרֵ֣ת בַּקֹּ֑דֶשׁ אֶת־בִּגְדֵ֤י הַקֹּ֙דֶשׁ֙ לְאַהֲרֹ֣ן הַכֹּהֵ֔ן וְאֶת־בִּגְדֵ֥י בָנָ֖יו לְכַהֵֽן׃ \Rashi{\textbf{בגדי השרד.}לכסות הארון והשלחן, המנורה והמזבחות בשעת סלוק מסעות:}%
% ---------- פרק 35 | פסוק 20 ----------
 \textbf{כ} וַיֵּ֥צְא֛וּ כׇּל־עֲדַ֥ת בְּנֵֽי־יִשְׂרָאֵ֖ל מִלִּפְנֵ֥י מֹשֶֽׁה׃ %
% ---------- פרק 35 | פסוק 21 ----------
 \textbf{כא} וַיָּבֹ֕אוּ כׇּל־אִ֖ישׁ אֲשֶׁר־נְשָׂא֣וֹ לִבּ֑וֹ וְכֹ֡ל אֲשֶׁר֩ נָדְבָ֨ה רוּח֜וֹ אֹת֗וֹ הֵ֠בִ֠יאוּ אֶת־תְּרוּמַ֨ת יְהֹוָ֜ה לִמְלֶ֨אכֶת אֹ֤הֶל מוֹעֵד֙ וּלְכׇל־עֲבֹ֣דָת֔וֹ וּלְבִגְדֵ֖י הַקֹּֽדֶשׁ׃ %
% ---------- פרק 35 | פסוק 22 ----------
 \textbf{כב} וַיָּבֹ֥אוּ הָאֲנָשִׁ֖ים עַל־הַנָּשִׁ֑ים כֹּ֣ל ׀ נְדִ֣יב לֵ֗ב הֵ֠בִ֠יאוּ חָ֣ח וָנֶ֜זֶם וְטַבַּ֤עַת וְכוּמָז֙ כׇּל־כְּלִ֣י זָהָ֔ב וְכׇל־אִ֕ישׁ אֲשֶׁ֥ר הֵנִ֛יף תְּנוּפַ֥ת זָהָ֖ב לַיהֹוָֽה׃ \Rashi{\textbf{(על הנשים.}עם הנשים וסמוכין אליהן): חח. הוא תכשיט של זהב עגל, נתון על הזרוע, והוא הצמיד: וכומז. כלי זהב הוא נתון כנגד אותו מקום לאשה, ורבותינו פרשו שם כומז, כאן מקום זמה (שבת ס"ד):}%
% ---------- פרק 35 | פסוק 23 ----------
 \textbf{כג} וְכׇל־אִ֞ישׁ אֲשֶׁר־נִמְצָ֣א אִתּ֗וֹ תְּכֵ֧לֶת וְאַרְגָּמָ֛ן וְתוֹלַ֥עַת שָׁנִ֖י וְשֵׁ֣שׁ וְעִזִּ֑ים וְעֹרֹ֨ת אֵילִ֧ם מְאׇדָּמִ֛ים וְעֹרֹ֥ת תְּחָשִׁ֖ים הֵבִֽיאוּ׃ \Rashi{\textbf{וכל איש אשר נמצא אתו.}תכלת או ארגמן או תולעת שני או עורות אילים או תחשים כלם הביאו:}%
% ---------- פרק 35 | פסוק 24 ----------
 \textbf{כד} כׇּל־מֵרִ֗ים תְּר֤וּמַת כֶּ֙סֶף֙ וּנְחֹ֔שֶׁת הֵבִ֕יאוּ אֵ֖ת תְּרוּמַ֣ת יְהֹוָ֑ה וְכֹ֡ל אֲשֶׁר֩ נִמְצָ֨א אִתּ֜וֹ עֲצֵ֥י שִׁטִּ֛ים לְכׇל־מְלֶ֥אכֶת הָעֲבֹדָ֖ה הֵבִֽיאוּ׃ %
% ---------- פרק 35 | פסוק 25 ----------
 \textbf{כה} וְכׇל־אִשָּׁ֥ה חַכְמַת־לֵ֖ב בְּיָדֶ֣יהָ טָו֑וּ וַיָּבִ֣יאוּ מַטְוֶ֗ה אֶֽת־הַתְּכֵ֙לֶת֙ וְאֶת־הָֽאַרְגָּמָ֔ן אֶת־תּוֹלַ֥עַת הַשָּׁנִ֖י וְאֶת־הַשֵּֽׁשׁ׃ %
% ---------- פרק 35 | פסוק 26 ----------
 \textbf{כו} וְכׇ֨ל־הַנָּשִׁ֔ים אֲשֶׁ֨ר נָשָׂ֥א לִבָּ֛ן אֹתָ֖נָה בְּחׇכְמָ֑ה טָו֖וּ אֶת־הָעִזִּֽים׃ \Rashi{\textbf{טוו את העזים.}היא היתה אמנות יתרה, שמעל גבי העזים טווין אותם (שבת ע"ד):}%
% ---------- פרק 35 | פסוק 27 ----------
 \textbf{כז} וְהַנְּשִׂאִ֣ם הֵבִ֔יאוּ אֵ֚ת אַבְנֵ֣י הַשֹּׁ֔הַם וְאֵ֖ת אַבְנֵ֣י הַמִּלֻּאִ֑ים לָאֵפ֖וֹד וְלַחֹֽשֶׁן׃ \Rashi{\textbf{והנשאם הביאו.}אמר ר' נתן: מה ראו נשיאים להתנדב בחנכת המזבח בתחלה ובמלאכת המשכן לא התנדבו בתחלה? אלא כך אמרו נשיאים, יתנדבו צבור מה שמתנדבין, ומה שמחסרים, אנו משלימין אותו, כיון שהשלימו צבור את הכל – שנ' והמלאכה היתה דים – אמרו נשיאים מה עלינו לעשות? הביאו את אבני השהם וגו', לכך התנדבו בחנכת המזבח תחלה, ולפי שנתעצלו מתחלה נחסרה אות משמם, והנשאם כתיב (ספרי במדבר מ"ה):}%
% ---------- פרק 35 | פסוק 28 ----------
 \textbf{כח} וְאֶת־הַבֹּ֖שֶׂם וְאֶת־הַשָּׁ֑מֶן לְמָא֕וֹר וּלְשֶׁ֙מֶן֙ הַמִּשְׁחָ֔ה וְלִקְטֹ֖רֶת הַסַּמִּֽים׃ %
% ---------- פרק 35 | פסוק 29 ----------
 \textbf{כט} כׇּל־אִ֣ישׁ וְאִשָּׁ֗ה אֲשֶׁ֨ר נָדַ֣ב לִבָּם֮ אֹתָם֒ לְהָבִיא֙ לְכׇל־הַמְּלָאכָ֔ה אֲשֶׁ֨ר צִוָּ֧ה יְהֹוָ֛ה לַעֲשׂ֖וֹת בְּיַד־מֹשֶׁ֑ה הֵבִ֧יאוּ בְנֵי־יִשְׂרָאֵ֛ל נְדָבָ֖ה לַיהֹוָֽה׃ \{פ\} %
% ---------- פרק 35 | פסוק 30 ----------
 \textbf{ל} וַיֹּ֤אמֶר מֹשֶׁה֙ אֶל־בְּנֵ֣י יִשְׂרָאֵ֔ל רְא֛וּ קָרָ֥א יְהֹוָ֖ה בְּשֵׁ֑ם בְּצַלְאֵ֛ל בֶּן־אוּרִ֥י בֶן־ח֖וּר לְמַטֵּ֥ה יְהוּדָֽה׃ \Rashi{\textbf{חור.}בנה של מרים היה:}%
% ---------- פרק 35 | פסוק 31 ----------
 \textbf{לא} וַיְמַלֵּ֥א אֹת֖וֹ ר֣וּחַ אֱלֹהִ֑ים בְּחׇכְמָ֛ה בִּתְבוּנָ֥ה וּבְדַ֖עַת וּבְכׇל־מְלָאכָֽה׃ %
% ---------- פרק 35 | פסוק 32 ----------
 \textbf{לב} וְלַחְשֹׁ֖ב מַֽחֲשָׁבֹ֑ת לַעֲשֹׂ֛ת בַּזָּהָ֥ב וּבַכֶּ֖סֶף וּבַנְּחֹֽשֶׁת׃ %
% ---------- פרק 35 | פסוק 33 ----------
 \textbf{לג} וּבַחֲרֹ֥שֶׁת אֶ֛בֶן לְמַלֹּ֖את וּבַחֲרֹ֣שֶׁת עֵ֑ץ לַעֲשׂ֖וֹת בְּכׇל־מְלֶ֥אכֶת מַחֲשָֽׁבֶת׃ %
% ---------- פרק 35 | פסוק 34 ----------
 \textbf{לד} וּלְהוֹרֹ֖ת נָתַ֣ן בְּלִבּ֑וֹ ה֕וּא וְאׇֽהֳלִיאָ֥ב בֶּן־אֲחִיסָמָ֖ךְ לְמַטֵּה־דָֽן׃ \Rashi{\textbf{ואהליאב.}משבט דן מן הירודין שבשבטים – מבני השפחות – והשוהו המקום לבצלאל למלאכת המשכן והוא מגדולי השבטים, לקים מה שנ' ולא נכר שוע לפני דל (איוב ל"ד):}%
% ---------- פרק 35 | פסוק 35 ----------
 \textbf{לה} מִלֵּ֨א אֹתָ֜ם חׇכְמַת־לֵ֗ב לַעֲשׂוֹת֮ כׇּל־מְלֶ֣אכֶת חָרָ֣שׁ ׀ וְחֹשֵׁב֒ וְרֹקֵ֞ם בַּתְּכֵ֣לֶת וּבָֽאַרְגָּמָ֗ן בְּתוֹלַ֧עַת הַשָּׁנִ֛י וּבַשֵּׁ֖שׁ וְאֹרֵ֑ג עֹשֵׂי֙ כׇּל־מְלָאכָ֔ה וְחֹשְׁבֵ֖י מַחֲשָׁבֹֽת׃ %
% ---------- פרק 36 | פסוק 1 ----------
 \textbf{א} וְעָשָׂה֩ בְצַלְאֵ֨ל וְאׇהֳלִיאָ֜ב וְכֹ֣ל ׀ אִ֣ישׁ חֲכַם־לֵ֗ב אֲשֶׁר֩ נָתַ֨ן יְהֹוָ֜ה חׇכְמָ֤ה וּתְבוּנָה֙ בָּהֵ֔מָּה לָדַ֣עַת לַעֲשֹׂ֔ת אֶֽת־כׇּל־מְלֶ֖אכֶת עֲבֹדַ֣ת הַקֹּ֑דֶשׁ לְכֹ֥ל אֲשֶׁר־צִוָּ֖ה יְהֹוָֽה׃ %
% ---------- פרק 36 | פסוק 2 ----------
 \textbf{ב} וַיִּקְרָ֣א מֹשֶׁ֗ה אֶל־בְּצַלְאֵל֮ וְאֶל־אׇֽהֳלִיאָב֒ וְאֶל֙ כׇּל־אִ֣ישׁ חֲכַם־לֵ֔ב אֲשֶׁ֨ר נָתַ֧ן יְהֹוָ֛ה חׇכְמָ֖ה בְּלִבּ֑וֹ כֹּ֚ל אֲשֶׁ֣ר נְשָׂא֣וֹ לִבּ֔וֹ לְקׇרְבָ֥ה אֶל־הַמְּלָאכָ֖ה לַעֲשֹׂ֥ת אֹתָֽהּ׃ %
% ---------- פרק 36 | פסוק 3 ----------
 \textbf{ג} וַיִּקְח֞וּ מִלִּפְנֵ֣י מֹשֶׁ֗ה אֵ֤ת כׇּל־הַתְּרוּמָה֙ אֲשֶׁ֨ר הֵבִ֜יאוּ בְּנֵ֣י יִשְׂרָאֵ֗ל לִמְלֶ֛אכֶת עֲבֹדַ֥ת הַקֹּ֖דֶשׁ לַעֲשֹׂ֣ת אֹתָ֑הּ וְ֠הֵ֠ם הֵבִ֨יאוּ אֵלָ֥יו ע֛וֹד נְדָבָ֖ה בַּבֹּ֥קֶר בַּבֹּֽקֶר׃ %
% ---------- פרק 36 | פסוק 4 ----------
 \textbf{ד} וַיָּבֹ֙אוּ֙ כׇּל־הַ֣חֲכָמִ֔ים הָעֹשִׂ֕ים אֵ֖ת כׇּל־מְלֶ֣אכֶת הַקֹּ֑דֶשׁ אִֽישׁ־אִ֥ישׁ מִמְּלַאכְתּ֖וֹ אֲשֶׁר־הֵ֥מָּה עֹשִֽׂים׃ %
% ---------- פרק 36 | פסוק 5 ----------
 \textbf{ה} וַיֹּאמְרוּ֙ אֶל־מֹשֶׁ֣ה לֵּאמֹ֔ר מַרְבִּ֥ים הָעָ֖ם לְהָבִ֑יא מִדֵּ֤י הָֽעֲבֹדָה֙ לַמְּלָאכָ֔ה אֲשֶׁר־צִוָּ֥ה יְהֹוָ֖ה לַעֲשֹׂ֥ת אֹתָֽהּ׃ \Rashi{\textbf{מדי העבדה.}יותר מכדי צרך העבודה:}%
% ---------- פרק 36 | פסוק 6 ----------
 \textbf{ו} וַיְצַ֣ו מֹשֶׁ֗ה וַיַּעֲבִ֨ירוּ ק֥וֹל בַּֽמַּחֲנֶה֮ לֵאמֹר֒ אִ֣ישׁ וְאִשָּׁ֗ה אַל־יַעֲשׂוּ־ע֛וֹד מְלָאכָ֖ה לִתְרוּמַ֣ת הַקֹּ֑דֶשׁ וַיִּכָּלֵ֥א הָעָ֖ם מֵהָבִֽיא׃ \Rashi{\textbf{ויכלא.}לשון מניעה:}%
% ---------- פרק 36 | פסוק 7 ----------
 \textbf{ז} וְהַמְּלָאכָ֗ה הָיְתָ֥ה דַיָּ֛ם לְכׇל־הַמְּלָאכָ֖ה לַעֲשׂ֣וֹת אֹתָ֑הּ וְהוֹתֵֽר׃ \{ס\} \Rashi{\textbf{והמלאכה היתה דים לכל המלאכה.}ומלאכת ההבאה היתה דים של עושי המשכן, לכל המלאכה של משכן, לעשות אותה ולהותר: והותר. כמו והכבד את לבו (שמות ח'), והכות את מואב (מלכים ב ג'):}%
% ---------- פרק 36 | פסוק 8 ----------
 \textbf{ח} וַיַּעֲשׂ֨וּ כׇל־חֲכַם־לֵ֜ב בְּעֹשֵׂ֧י הַמְּלָאכָ֛ה אֶת־הַמִּשְׁכָּ֖ן עֶ֣שֶׂר יְרִיעֹ֑ת שֵׁ֣שׁ מׇשְׁזָ֗ר וּתְכֵ֤לֶת וְאַרְגָּמָן֙ וְתוֹלַ֣עַת שָׁנִ֔י כְּרֻבִ֛ים מַעֲשֵׂ֥ה חֹשֵׁ֖ב עָשָׂ֥ה אֹתָֽם׃ %
% ---------- פרק 36 | פסוק 9 ----------
 \textbf{ט} אֹ֜רֶךְ הַיְרִיעָ֣ה הָֽאַחַ֗ת שְׁמֹנֶ֤ה וְעֶשְׂרִים֙ בָּֽאַמָּ֔ה וְרֹ֙חַב֙ אַרְבַּ֣ע בָּֽאַמָּ֔ה הַיְרִיעָ֖ה הָאֶחָ֑ת מִדָּ֥ה אַחַ֖ת לְכׇל־הַיְרִיעֹֽת׃ %
% ---------- פרק 36 | פסוק 10 ----------
 \textbf{י} וַיְחַבֵּר֙ אֶת־חֲמֵ֣שׁ הַיְרִיעֹ֔ת אַחַ֖ת אֶל־אֶחָ֑ת וְחָמֵ֤שׁ יְרִיעֹת֙ חִבַּ֔ר אַחַ֖ת אֶל־אֶחָֽת׃ %
% ---------- פרק 36 | פסוק 11 ----------
 \textbf{יא} וַיַּ֜עַשׂ לֻֽלְאֹ֣ת תְּכֵ֗לֶת עַ֣ל שְׂפַ֤ת הַיְרִיעָה֙ הָֽאֶחָ֔ת מִקָּצָ֖ה בַּמַּחְבָּ֑רֶת כֵּ֤ן עָשָׂה֙ בִּשְׂפַ֣ת הַיְרִיעָ֔ה הַקִּ֣יצוֹנָ֔ה בַּמַּחְבֶּ֖רֶת הַשֵּׁנִֽית׃ %
% ---------- פרק 36 | פסוק 12 ----------
 \textbf{יב} חֲמִשִּׁ֣ים לֻלָאֹ֗ת עָשָׂה֮ בַּיְרִיעָ֣ה הָאֶחָת֒ וַחֲמִשִּׁ֣ים לֻלָאֹ֗ת עָשָׂה֙ בִּקְצֵ֣ה הַיְרִיעָ֔ה אֲשֶׁ֖ר בַּמַּחְבֶּ֣רֶת הַשֵּׁנִ֑ית מַקְבִּילֹת֙ הַלֻּ֣לָאֹ֔ת אַחַ֖ת אֶל־אֶחָֽת׃ %
% ---------- פרק 36 | פסוק 13 ----------
 \textbf{יג} וַיַּ֕עַשׂ חֲמִשִּׁ֖ים קַרְסֵ֣י זָהָ֑ב וַיְחַבֵּ֨ר אֶת־הַיְרִיעֹ֜ת אַחַ֤ת אֶל־אַחַת֙ בַּקְּרָסִ֔ים וַיְהִ֥י הַמִּשְׁכָּ֖ן אֶחָֽד׃ \{פ\} %
% ---------- פרק 36 | פסוק 14 ----------
 \textbf{יד} וַיַּ֙עַשׂ֙ יְרִיעֹ֣ת עִזִּ֔ים לְאֹ֖הֶל עַל־הַמִּשְׁכָּ֑ן עַשְׁתֵּֽי־עֶשְׂרֵ֥ה יְרִיעֹ֖ת עָשָׂ֥ה אֹתָֽם׃ %
% ---------- פרק 36 | פסוק 15 ----------
 \textbf{טו} אֹ֜רֶךְ הַיְרִיעָ֣ה הָאַחַ֗ת שְׁלֹשִׁים֙ בָּֽאַמָּ֔ה וְאַרְבַּ֣ע אַמּ֔וֹת רֹ֖חַב הַיְרִיעָ֣ה הָאֶחָ֑ת מִדָּ֣ה אַחַ֔ת לְעַשְׁתֵּ֥י עֶשְׂרֵ֖ה יְרִיעֹֽת׃ %
% ---------- פרק 36 | פסוק 16 ----------
 \textbf{טז} וַיְחַבֵּ֛ר אֶת־חֲמֵ֥שׁ הַיְרִיעֹ֖ת לְבָ֑ד וְאֶת־שֵׁ֥שׁ הַיְרִיעֹ֖ת לְבָֽד׃ %
% ---------- פרק 36 | פסוק 17 ----------
 \textbf{יז} וַיַּ֜עַשׂ לֻֽלָאֹ֣ת חֲמִשִּׁ֗ים עַ֚ל שְׂפַ֣ת הַיְרִיעָ֔ה הַקִּיצֹנָ֖ה בַּמַּחְבָּ֑רֶת וַחֲמִשִּׁ֣ים לֻלָאֹ֗ת עָשָׂה֙ עַל־שְׂפַ֣ת הַיְרִיעָ֔ה הַחֹבֶ֖רֶת הַשֵּׁנִֽית׃ %
% ---------- פרק 36 | פסוק 18 ----------
 \textbf{יח} וַיַּ֛עַשׂ קַרְסֵ֥י נְחֹ֖שֶׁת חֲמִשִּׁ֑ים לְחַבֵּ֥ר אֶת־הָאֹ֖הֶל לִהְיֹ֥ת אֶחָֽד׃ %
% ---------- פרק 36 | פסוק 19 ----------
 \textbf{יט} וַיַּ֤עַשׂ מִכְסֶה֙ לָאֹ֔הֶל עֹרֹ֥ת אֵילִ֖ם מְאׇדָּמִ֑ים וּמִכְסֵ֛ה עֹרֹ֥ת תְּחָשִׁ֖ים מִלְמָֽעְלָה׃ \{ס\} %
% ---------- פרק 36 | פסוק 20 ----------
 \textbf{כ} וַיַּ֥עַשׂ אֶת־הַקְּרָשִׁ֖ים לַמִּשְׁכָּ֑ן עֲצֵ֥י שִׁטִּ֖ים עֹמְדִֽים׃ %
% ---------- פרק 36 | פסוק 21 ----------
 \textbf{כא} עֶ֥שֶׂר אַמֹּ֖ת אֹ֣רֶךְ הַקָּ֑רֶשׁ וְאַמָּה֙ וַחֲצִ֣י הָֽאַמָּ֔ה רֹ֖חַב הַקֶּ֥רֶשׁ הָאֶחָֽד׃ %
% ---------- פרק 36 | פסוק 22 ----------
 \textbf{כב} שְׁתֵּ֣י יָדֹ֗ת לַקֶּ֙רֶשׁ֙ הָֽאֶחָ֔ד מְשֻׁ֨לָּבֹ֔ת אַחַ֖ת אֶל־אֶחָ֑ת כֵּ֣ן עָשָׂ֔ה לְכֹ֖ל קַרְשֵׁ֥י הַמִּשְׁכָּֽן׃ %
% ---------- פרק 36 | פסוק 23 ----------
 \textbf{כג} וַיַּ֥עַשׂ אֶת־הַקְּרָשִׁ֖ים לַמִּשְׁכָּ֑ן עֶשְׂרִ֣ים קְרָשִׁ֔ים לִפְאַ֖ת נֶ֥גֶב תֵּימָֽנָה׃ %
% ---------- פרק 36 | פסוק 24 ----------
 \textbf{כד} וְאַרְבָּעִים֙ אַדְנֵי־כֶ֔סֶף עָשָׂ֕ה תַּ֖חַת עֶשְׂרִ֣ים הַקְּרָשִׁ֑ים שְׁנֵ֨י אֲדָנִ֜ים תַּֽחַת־הַקֶּ֤רֶשׁ הָאֶחָד֙ לִשְׁתֵּ֣י יְדֹתָ֔יו וּשְׁנֵ֧י אֲדָנִ֛ים תַּֽחַת־הַקֶּ֥רֶשׁ הָאֶחָ֖ד לִשְׁתֵּ֥י יְדֹתָֽיו׃ %
% ---------- פרק 36 | פסוק 25 ----------
 \textbf{כה} וּלְצֶ֧לַע הַמִּשְׁכָּ֛ן הַשֵּׁנִ֖ית לִפְאַ֣ת צָפ֑וֹן עָשָׂ֖ה עֶשְׂרִ֥ים קְרָשִֽׁים׃ %
% ---------- פרק 36 | פסוק 26 ----------
 \textbf{כו} וְאַרְבָּעִ֥ים אַדְנֵיהֶ֖ם כָּ֑סֶף שְׁנֵ֣י אֲדָנִ֗ים תַּ֚חַת הַקֶּ֣רֶשׁ הָאֶחָ֔ד וּשְׁנֵ֣י אֲדָנִ֔ים תַּ֖חַת הַקֶּ֥רֶשׁ הָאֶחָֽד׃ %
% ---------- פרק 36 | פסוק 27 ----------
 \textbf{כז} וּֽלְיַרְכְּתֵ֥י הַמִּשְׁכָּ֖ן יָ֑מָּה עָשָׂ֖ה שִׁשָּׁ֥ה קְרָשִֽׁים׃ %
% ---------- פרק 36 | פסוק 28 ----------
 \textbf{כח} וּשְׁנֵ֤י קְרָשִׁים֙ עָשָׂ֔ה לִמְקֻצְעֹ֖ת הַמִּשְׁכָּ֑ן בַּיַּרְכָתָֽיִם׃ %
% ---------- פרק 36 | פסוק 29 ----------
 \textbf{כט} וְהָי֣וּ תוֹאֲמִם֮ מִלְּמַ֒טָּה֒ וְיַחְדָּ֗ו יִהְי֤וּ תַמִּים֙ אֶל־רֹאשׁ֔וֹ אֶל־הַטַּבַּ֖עַת הָאֶחָ֑ת כֵּ֚ן עָשָׂ֣ה לִשְׁנֵיהֶ֔ם לִשְׁנֵ֖י הַמִּקְצֹעֹֽת׃ %
% ---------- פרק 36 | פסוק 30 ----------
 \textbf{ל} וְהָיוּ֙ שְׁמֹנָ֣ה קְרָשִׁ֔ים וְאַדְנֵיהֶ֣ם כֶּ֔סֶף שִׁשָּׁ֥ה עָשָׂ֖ר אֲדָנִ֑ים שְׁנֵ֤י אֲדָנִים֙ שְׁנֵ֣י אֲדָנִ֔ים תַּ֖חַת הַקֶּ֥רֶשׁ הָאֶחָֽד׃ %
% ---------- פרק 36 | פסוק 31 ----------
 \textbf{לא} וַיַּ֥עַשׂ בְּרִיחֵ֖י עֲצֵ֣י שִׁטִּ֑ים חֲמִשָּׁ֕ה לְקַרְשֵׁ֥י צֶֽלַע־הַמִּשְׁכָּ֖ן הָאֶחָֽת׃ %
% ---------- פרק 36 | פסוק 32 ----------
 \textbf{לב} וַחֲמִשָּׁ֣ה בְרִיחִ֔ם לְקַרְשֵׁ֥י צֶֽלַע־הַמִּשְׁכָּ֖ן הַשֵּׁנִ֑ית וַחֲמִשָּׁ֤ה בְרִיחִם֙ לְקַרְשֵׁ֣י הַמִּשְׁכָּ֔ן לַיַּרְכָתַ֖יִם יָֽמָּה׃ %
% ---------- פרק 36 | פסוק 33 ----------
 \textbf{לג} וַיַּ֖עַשׂ אֶת־הַבְּרִ֣יחַ הַתִּיכֹ֑ן לִבְרֹ֙חַ֙ בְּת֣וֹךְ הַקְּרָשִׁ֔ים מִן־הַקָּצֶ֖ה אֶל־הַקָּצֶֽה׃ %
% ---------- פרק 36 | פסוק 34 ----------
 \textbf{לד} וְֽאֶת־הַקְּרָשִׁ֞ים צִפָּ֣ה זָהָ֗ב וְאֶת־טַבְּעֹתָם֙ עָשָׂ֣ה זָהָ֔ב בָּתִּ֖ים לַבְּרִיחִ֑ם וַיְצַ֥ף אֶת־הַבְּרִיחִ֖ם זָהָֽב׃ %
% ---------- פרק 36 | פסוק 35 ----------
 \textbf{לה} וַיַּ֙עַשׂ֙ אֶת־הַפָּרֹ֔כֶת תְּכֵ֧לֶת וְאַרְגָּמָ֛ן וְתוֹלַ֥עַת שָׁנִ֖י וְשֵׁ֣שׁ מׇשְׁזָ֑ר מַעֲשֵׂ֥ה חֹשֵׁ֛ב עָשָׂ֥ה אֹתָ֖הּ כְּרֻבִֽים׃ %
% ---------- פרק 36 | פסוק 36 ----------
 \textbf{לו} וַיַּ֣עַשׂ לָ֗הּ אַרְבָּעָה֙ עַמּוּדֵ֣י שִׁטִּ֔ים וַיְצַפֵּ֣ם זָהָ֔ב וָוֵיהֶ֖ם זָהָ֑ב וַיִּצֹ֣ק לָהֶ֔ם אַרְבָּעָ֖ה אַדְנֵי־כָֽסֶף׃ %
% ---------- פרק 36 | פסוק 37 ----------
 \textbf{לז} וַיַּ֤עַשׂ מָסָךְ֙ לְפֶ֣תַח הָאֹ֔הֶל תְּכֵ֧לֶת וְאַרְגָּמָ֛ן וְתוֹלַ֥עַת שָׁנִ֖י וְשֵׁ֣שׁ מׇשְׁזָ֑ר מַעֲשֵׂ֖ה רֹקֵֽם׃ %
% ---------- פרק 36 | פסוק 38 ----------
 \textbf{לח} וְאֶת־עַמּוּדָ֤יו חֲמִשָּׁה֙ וְאֶת־וָ֣וֵיהֶ֔ם וְצִפָּ֧ה רָאשֵׁיהֶ֛ם וַחֲשֻׁקֵיהֶ֖ם זָהָ֑ב וְאַדְנֵיהֶ֥ם חֲמִשָּׁ֖ה נְחֹֽשֶׁת׃ \{פ\} %
% ---------- פרק 37 | פסוק 1 ----------
 \textbf{א} וַיַּ֧עַשׂ בְּצַלְאֵ֛ל אֶת־הָאָרֹ֖ן עֲצֵ֣י שִׁטִּ֑ים אַמָּתַ֨יִם וָחֵ֜צִי אׇרְכּ֗וֹ וְאַמָּ֤ה וָחֵ֙צִי֙ רׇחְבּ֔וֹ וְאַמָּ֥ה וָחֵ֖צִי קֹמָתֽוֹ׃ \Rashi{\textbf{ויעש בצלאל.}לפי שנתן נפשו על המלאכה יותר משאר חכמים, נקראת על שמו (שמות רבה):}%
% ---------- פרק 37 | פסוק 2 ----------
 \textbf{ב} וַיְצַפֵּ֛הוּ זָהָ֥ב טָה֖וֹר מִבַּ֣יִת וּמִח֑וּץ וַיַּ֥עַשׂ ל֛וֹ זֵ֥ר זָהָ֖ב סָבִֽיב׃ %
% ---------- פרק 37 | פסוק 3 ----------
 \textbf{ג} וַיִּצֹ֣ק ל֗וֹ אַרְבַּע֙ טַבְּעֹ֣ת זָהָ֔ב עַ֖ל אַרְבַּ֣ע פַּעֲמֹתָ֑יו וּשְׁתֵּ֣י טַבָּעֹ֗ת עַל־צַלְעוֹ֙ הָֽאֶחָ֔ת וּשְׁתֵּי֙ טַבָּעֹ֔ת עַל־צַלְע֖וֹ הַשֵּׁנִֽית׃ %
% ---------- פרק 37 | פסוק 4 ----------
 \textbf{ד} וַיַּ֥עַשׂ בַּדֵּ֖י עֲצֵ֣י שִׁטִּ֑ים וַיְצַ֥ף אֹתָ֖ם זָהָֽב׃ %
% ---------- פרק 37 | פסוק 5 ----------
 \textbf{ה} וַיָּבֵ֤א אֶת־הַבַּדִּים֙ בַּטַּבָּעֹ֔ת עַ֖ל צַלְעֹ֣ת הָאָרֹ֑ן לָשֵׂ֖את אֶת־הָאָרֹֽן׃ %
% ---------- פרק 37 | פסוק 6 ----------
 \textbf{ו} וַיַּ֥עַשׂ כַּפֹּ֖רֶת זָהָ֣ב טָה֑וֹר אַמָּתַ֤יִם וָחֵ֙צִי֙ אׇרְכָּ֔הּ וְאַמָּ֥ה וָחֵ֖צִי רׇחְבָּֽהּ׃ %
% ---------- פרק 37 | פסוק 7 ----------
 \textbf{ז} וַיַּ֛עַשׂ שְׁנֵ֥י כְרֻבִ֖ים זָהָ֑ב מִקְשָׁה֙ עָשָׂ֣ה אֹתָ֔ם מִשְּׁנֵ֖י קְצ֥וֹת הַכַּפֹּֽרֶת׃ %
% ---------- פרק 37 | פסוק 8 ----------
 \textbf{ח} כְּרוּב־אֶחָ֤ד מִקָּצָה֙ מִזֶּ֔ה וּכְרוּב־אֶחָ֥ד מִקָּצָ֖ה מִזֶּ֑ה מִן־הַכַּפֹּ֛רֶת עָשָׂ֥ה אֶת־הַכְּרֻבִ֖ים מִשְּׁנֵ֥י (קצוותו) [קְצוֹתָֽיו]׃ %
% ---------- פרק 37 | פסוק 9 ----------
 \textbf{ט} וַיִּהְי֣וּ הַכְּרֻבִים֩ פֹּרְשֵׂ֨י כְנָפַ֜יִם לְמַ֗עְלָה סֹֽכְכִ֤ים בְּכַנְפֵיהֶם֙ עַל־הַכַּפֹּ֔רֶת וּפְנֵיהֶ֖ם אִ֣ישׁ אֶל־אָחִ֑יו אֶ֨ל־הַכַּפֹּ֔רֶת הָי֖וּ פְּנֵ֥י הַכְּרֻבִֽים׃ \{פ\} %
% ---------- פרק 37 | פסוק 10 ----------
 \textbf{י} וַיַּ֥עַשׂ אֶת־הַשֻּׁלְחָ֖ן עֲצֵ֣י שִׁטִּ֑ים אַמָּתַ֤יִם אׇרְכּוֹ֙ וְאַמָּ֣ה רׇחְבּ֔וֹ וְאַמָּ֥ה וָחֵ֖צִי קֹמָתֽוֹ׃ %
% ---------- פרק 37 | פסוק 11 ----------
 \textbf{יא} וַיְצַ֥ף אֹת֖וֹ זָהָ֣ב טָה֑וֹר וַיַּ֥עַשׂ ל֛וֹ זֵ֥ר זָהָ֖ב סָבִֽיב׃ %
% ---------- פרק 37 | פסוק 12 ----------
 \textbf{יב} וַיַּ֨עַשׂ ל֥וֹ מִסְגֶּ֛רֶת טֹ֖פַח סָבִ֑יב וַיַּ֧עַשׂ זֵר־זָהָ֛ב לְמִסְגַּרְתּ֖וֹ סָבִֽיב׃ %
% ---------- פרק 37 | פסוק 13 ----------
 \textbf{יג} וַיִּצֹ֣ק ל֔וֹ אַרְבַּ֖ע טַבְּעֹ֣ת זָהָ֑ב וַיִּתֵּן֙ אֶת־הַטַּבָּעֹ֔ת עַ֚ל אַרְבַּ֣ע הַפֵּאֹ֔ת אֲשֶׁ֖ר לְאַרְבַּ֥ע רַגְלָֽיו׃ %
% ---------- פרק 37 | פסוק 14 ----------
 \textbf{יד} לְעֻמַּת֙ הַמִּסְגֶּ֔רֶת הָי֖וּ הַטַּבָּעֹ֑ת בָּתִּים֙ לַבַּדִּ֔ים לָשֵׂ֖את אֶת־הַשֻּׁלְחָֽן׃ %
% ---------- פרק 37 | פסוק 15 ----------
 \textbf{טו} וַיַּ֤עַשׂ אֶת־הַבַּדִּים֙ עֲצֵ֣י שִׁטִּ֔ים וַיְצַ֥ף אֹתָ֖ם זָהָ֑ב לָשֵׂ֖את אֶת־הַשֻּׁלְחָֽן׃ %
% ---------- פרק 37 | פסוק 16 ----------
 \textbf{טז} וַיַּ֜עַשׂ אֶֽת־הַכֵּלִ֣ים ׀ אֲשֶׁ֣ר עַל־הַשֻּׁלְחָ֗ן אֶת־קְעָרֹתָ֤יו וְאֶת־כַּפֹּתָיו֙ וְאֵת֙ מְנַקִּיֹּתָ֔יו וְאֶ֨ת־הַקְּשָׂוֺ֔ת אֲשֶׁ֥ר יֻסַּ֖ךְ בָּהֵ֑ן זָהָ֖ב טָהֽוֹר׃ \{פ\} %
% ---------- פרק 37 | פסוק 17 ----------
 \textbf{יז} וַיַּ֥עַשׂ אֶת־הַמְּנֹרָ֖ה זָהָ֣ב טָה֑וֹר מִקְשָׁ֞ה עָשָׂ֤ה אֶת־הַמְּנֹרָה֙ יְרֵכָ֣הּ וְקָנָ֔הּ גְּבִיעֶ֛יהָ כַּפְתֹּרֶ֥יהָ וּפְרָחֶ֖יהָ מִמֶּ֥נָּה הָיֽוּ׃ %
% ---------- פרק 37 | פסוק 18 ----------
 \textbf{יח} וְשִׁשָּׁ֣ה קָנִ֔ים יֹצְאִ֖ים מִצִּדֶּ֑יהָ שְׁלֹשָׁ֣ה ׀ קְנֵ֣י מְנֹרָ֗ה מִצִּדָּהּ֙ הָֽאֶחָ֔ד וּשְׁלֹשָׁה֙ קְנֵ֣י מְנֹרָ֔ה מִצִּדָּ֖הּ הַשֵּׁנִֽי׃ %
% ---------- פרק 37 | פסוק 19 ----------
 \textbf{יט} שְׁלֹשָׁ֣ה גְ֠בִעִ֠ים מְֽשֻׁקָּדִ֞ים בַּקָּנֶ֣ה הָאֶחָד֮ כַּפְתֹּ֣ר וָפֶ֒רַח֒ וּשְׁלֹשָׁ֣ה גְבִעִ֗ים מְשֻׁקָּדִ֛ים בְּקָנֶ֥ה אֶחָ֖ד כַּפְתֹּ֣ר וָפָ֑רַח כֵּ֚ן לְשֵׁ֣שֶׁת הַקָּנִ֔ים הַיֹּצְאִ֖ים מִן־הַמְּנֹרָֽה׃ %
% ---------- פרק 37 | פסוק 20 ----------
 \textbf{כ} וּבַמְּנֹרָ֖ה אַרְבָּעָ֣ה גְבִעִ֑ים מְשֻׁ֨קָּדִ֔ים כַּפְתֹּרֶ֖יהָ וּפְרָחֶֽיהָ׃ %
% ---------- פרק 37 | פסוק 21 ----------
 \textbf{כא} וְכַפְתֹּ֡ר תַּ֩חַת֩ שְׁנֵ֨י הַקָּנִ֜ים מִמֶּ֗נָּה וְכַפְתֹּר֙ תַּ֣חַת שְׁנֵ֤י הַקָּנִים֙ מִמֶּ֔נָּה וְכַפְתֹּ֕ר תַּֽחַת־שְׁנֵ֥י הַקָּנִ֖ים מִמֶּ֑נָּה לְשֵׁ֙שֶׁת֙ הַקָּנִ֔ים הַיֹּצְאִ֖ים מִמֶּֽנָּה׃ %
% ---------- פרק 37 | פסוק 22 ----------
 \textbf{כב} כַּפְתֹּרֵיהֶ֥ם וּקְנֹתָ֖ם מִמֶּ֣נָּה הָי֑וּ כֻּלָּ֛הּ מִקְשָׁ֥ה אַחַ֖ת זָהָ֥ב טָהֽוֹר׃ %
% ---------- פרק 37 | פסוק 23 ----------
 \textbf{כג} וַיַּ֥עַשׂ אֶת־נֵרֹתֶ֖יהָ שִׁבְעָ֑ה וּמַלְקָחֶ֥יהָ וּמַחְתֹּתֶ֖יהָ זָהָ֥ב טָהֽוֹר׃ %
% ---------- פרק 37 | פסוק 24 ----------
 \textbf{כד} כִּכָּ֛ר זָהָ֥ב טָה֖וֹר עָשָׂ֣ה אֹתָ֑הּ וְאֵ֖ת כׇּל־כֵּלֶֽיהָ׃ \{פ\} %
% ---------- פרק 37 | פסוק 25 ----------
 \textbf{כה} וַיַּ֛עַשׂ אֶת־מִזְבַּ֥ח הַקְּטֹ֖רֶת עֲצֵ֣י שִׁטִּ֑ים אַמָּ֣ה אׇרְכּוֹ֩ וְאַמָּ֨ה רׇחְבּ֜וֹ רָב֗וּעַ וְאַמָּתַ֙יִם֙ קֹֽמָת֔וֹ מִמֶּ֖נּוּ הָי֥וּ קַרְנֹתָֽיו׃ %
% ---------- פרק 37 | פסוק 26 ----------
 \textbf{כו} וַיְצַ֨ף אֹת֜וֹ זָהָ֣ב טָה֗וֹר אֶת־גַּגּ֧וֹ וְאֶת־קִירֹתָ֛יו סָבִ֖יב וְאֶת־קַרְנֹתָ֑יו וַיַּ֥עַשׂ ל֛וֹ זֵ֥ר זָהָ֖ב סָבִֽיב׃ %
% ---------- פרק 37 | פסוק 27 ----------
 \textbf{כז} וּשְׁתֵּי֩ טַבְּעֹ֨ת זָהָ֜ב עָֽשָׂה־ל֣וֹ ׀ מִתַּ֣חַת לְזֵר֗וֹ עַ֚ל שְׁתֵּ֣י צַלְעֹתָ֔יו עַ֖ל שְׁנֵ֣י צִדָּ֑יו לְבָתִּ֣ים לְבַדִּ֔ים לָשֵׂ֥את אֹת֖וֹ בָּהֶֽם׃ %
% ---------- פרק 37 | פסוק 28 ----------
 \textbf{כח} וַיַּ֥עַשׂ אֶת־הַבַּדִּ֖ים עֲצֵ֣י שִׁטִּ֑ים וַיְצַ֥ף אֹתָ֖ם זָהָֽב׃ %
% ---------- פרק 37 | פסוק 29 ----------
 \textbf{כט} וַיַּ֜עַשׂ אֶת־שֶׁ֤מֶן הַמִּשְׁחָה֙ קֹ֔דֶשׁ וְאֶת־קְטֹ֥רֶת הַסַּמִּ֖ים טָה֑וֹר מַעֲשֵׂ֖ה רֹקֵֽחַ׃ \{ס\} %
% ---------- פרק 38 | פסוק 1 ----------
 \textbf{א} וַיַּ֛עַשׂ אֶת־מִזְבַּ֥ח הָעֹלָ֖ה עֲצֵ֣י שִׁטִּ֑ים חָמֵשׁ֩ אַמּ֨וֹת אׇרְכּ֜וֹ וְחָֽמֵשׁ־אַמּ֤וֹת רׇחְבּוֹ֙ רָב֔וּעַ וְשָׁלֹ֥שׁ אַמּ֖וֹת קֹמָתֽוֹ׃ %
% ---------- פרק 38 | פסוק 2 ----------
 \textbf{ב} וַיַּ֣עַשׂ קַרְנֹתָ֗יו עַ֚ל אַרְבַּ֣ע פִּנֹּתָ֔יו מִמֶּ֖נּוּ הָי֣וּ קַרְנֹתָ֑יו וַיְצַ֥ף אֹת֖וֹ נְחֹֽשֶׁת׃ %
% ---------- פרק 38 | פסוק 3 ----------
 \textbf{ג} וַיַּ֜עַשׂ אֶֽת־כׇּל־כְּלֵ֣י הַמִּזְבֵּ֗חַ אֶת־הַסִּירֹ֤ת וְאֶת־הַיָּעִים֙ וְאֶת־הַמִּזְרָקֹ֔ת אֶת־הַמִּזְלָגֹ֖ת וְאֶת־הַמַּחְתֹּ֑ת כׇּל־כֵּלָ֖יו עָשָׂ֥ה נְחֹֽשֶׁת׃ %
% ---------- פרק 38 | פסוק 4 ----------
 \textbf{ד} וַיַּ֤עַשׂ לַמִּזְבֵּ֙חַ֙ מִכְבָּ֔ר מַעֲשֵׂ֖ה רֶ֣שֶׁת נְחֹ֑שֶׁת תַּ֧חַת כַּרְכֻּבּ֛וֹ מִלְּמַ֖טָּה עַד־חֶצְיֽוֹ׃ %
% ---------- פרק 38 | פסוק 5 ----------
 \textbf{ה} וַיִּצֹ֞ק אַרְבַּ֧ע טַבָּעֹ֛ת בְּאַרְבַּ֥ע הַקְּצָוֺ֖ת לְמִכְבַּ֣ר הַנְּחֹ֑שֶׁת בָּתִּ֖ים לַבַּדִּֽים׃ %
% ---------- פרק 38 | פסוק 6 ----------
 \textbf{ו} וַיַּ֥עַשׂ אֶת־הַבַּדִּ֖ים עֲצֵ֣י שִׁטִּ֑ים וַיְצַ֥ף אֹתָ֖ם נְחֹֽשֶׁת׃ %
% ---------- פרק 38 | פסוק 7 ----------
 \textbf{ז} וַיָּבֵ֨א אֶת־הַבַּדִּ֜ים בַּטַּבָּעֹ֗ת עַ֚ל צַלְעֹ֣ת הַמִּזְבֵּ֔חַ לָשֵׂ֥את אֹת֖וֹ בָּהֶ֑ם נְב֥וּב לֻחֹ֖ת עָשָׂ֥ה אֹתֽוֹ׃ \{ס\} \Rashi{\textbf{נבוב לחת.}נבוב הוא חלול, וכן ועביו ארבע אצבעות נבוב (ירמיהו נ"ב): נבוב לחת. הלוחות של עצי שטים לכל רוח, והחלל באמצע:}%
% ---------- פרק 38 | פסוק 8 ----------
 \textbf{ח} וַיַּ֗עַשׂ אֵ֚ת הַכִּיּ֣וֹר נְחֹ֔שֶׁת וְאֵ֖ת כַּנּ֣וֹ נְחֹ֑שֶׁת בְּמַרְאֹת֙ הַצֹּ֣בְאֹ֔ת אֲשֶׁ֣ר צָֽבְא֔וּ פֶּ֖תַח אֹ֥הֶל מוֹעֵֽד׃ \{ס\} \Rashi{\textbf{במראת הצבאת.}בנות ישראל היו בידן מראות שרואות בהן כשהן מתקשטות, ואף אותן לא עכבו מלהביא לנדבת המשכן, והיה מואס משה בהן מפני שעשוים ליצר הרע, אמר לו הקב"ה קבל, כי אלו חביבין עלי מן הכל, שעל ידיהם העמידו הנשים צבאות רבות במצרים; כשהיו בעליהם יגעים בעבודת פרך, היו הולכות ומוליכות להם מאכל ומשתה, ומאכילות אותם ונוטלות המראות, וכל אחת רואה עצמה עם בעלה במראה, ומשדלתו בדברים, לומר אני נאה ממך, ומתוך כך מביאות לבעליהן לידי תאוה ונזקקות להם ומתעברות ויולדות שם, שנאמר תחת התפוח עוררתיך (שיר השירים ח'), וזה שנאמר במראת הצובאות. ונעשה הכיור מהם, שהוא לשום שלום בין איש לאשתו – להשקות ממים שבתוכו למי שקנא לה בעלה ונסתרה; ותדע לך שהן מראות ממש, שהרי נאמר ונחשת התנופה שבעים ככר וגו', ויעש בה וגומר, וכיור וכנו לא הזכרו שם, למדת שלא היה נחשת של כיור מנחשת התנופה, כך דרש רבי תנחומא, וכן תרגם אנקלוס במחזית נשיא והוא תרגום של מראות, מירוריי"ש בלעז. וכן מצינו והגליונים (ישעיה ג') – מתרגמינן מחזיתא: אשר צבאו. להביא נדבתן:}%
% ---------- פרק 38 | פסוק 9 ----------
 \textbf{ט} וַיַּ֖עַשׂ אֶת־הֶחָצֵ֑ר לִפְאַ֣ת ׀ נֶ֣גֶב תֵּימָ֗נָה קַלְעֵ֤י הֶֽחָצֵר֙ שֵׁ֣שׁ מׇשְׁזָ֔ר מֵאָ֖ה בָּאַמָּֽה׃ %
% ---------- פרק 38 | פסוק 10 ----------
 \textbf{י} עַמּוּדֵיהֶ֣ם עֶשְׂרִ֔ים וְאַדְנֵיהֶ֥ם עֶשְׂרִ֖ים נְחֹ֑שֶׁת וָוֵ֧י הָעַמּוּדִ֛ים וַחֲשֻׁקֵיהֶ֖ם כָּֽסֶף׃ %
% ---------- פרק 38 | פסוק 11 ----------
 \textbf{יא} וְלִפְאַ֤ת צָפוֹן֙ מֵאָ֣ה בָֽאַמָּ֔ה עַמּוּדֵיהֶ֣ם עֶשְׂרִ֔ים וְאַדְנֵיהֶ֥ם עֶשְׂרִ֖ים נְחֹ֑שֶׁת וָוֵ֧י הָֽעַמּוּדִ֛ים וַחֲשֻׁקֵיהֶ֖ם כָּֽסֶף׃ %
% ---------- פרק 38 | פסוק 12 ----------
 \textbf{יב} וְלִפְאַת־יָ֗ם קְלָעִים֙ חֲמִשִּׁ֣ים בָּֽאַמָּ֔ה עַמּוּדֵיהֶ֣ם עֲשָׂרָ֔ה וְאַדְנֵיהֶ֖ם עֲשָׂרָ֑ה וָוֵ֧י הָעַמֻּדִ֛ים וַחֲשׁוּקֵיהֶ֖ם כָּֽסֶף׃ %
% ---------- פרק 38 | פסוק 13 ----------
 \textbf{יג} וְלִפְאַ֛ת קֵ֥דְמָה מִזְרָ֖חָה חֲמִשִּׁ֥ים אַמָּֽה׃ %
% ---------- פרק 38 | פסוק 14 ----------
 \textbf{יד} קְלָעִ֛ים חֲמֵשׁ־עֶשְׂרֵ֥ה אַמָּ֖ה אֶל־הַכָּתֵ֑ף עַמּוּדֵיהֶ֣ם שְׁלֹשָׁ֔ה וְאַדְנֵיהֶ֖ם שְׁלֹשָֽׁה׃ %
% ---------- פרק 38 | פסוק 15 ----------
 \textbf{טו} וְלַכָּתֵ֣ף הַשֵּׁנִ֗ית מִזֶּ֤ה וּמִזֶּה֙ לְשַׁ֣עַר הֶֽחָצֵ֔ר קְלָעִ֕ים חֲמֵ֥שׁ עֶשְׂרֵ֖ה אַמָּ֑ה עַמֻּדֵיהֶ֣ם שְׁלֹשָׁ֔ה וְאַדְנֵיהֶ֖ם שְׁלֹשָֽׁה׃ %
% ---------- פרק 38 | פסוק 16 ----------
 \textbf{טז} כׇּל־קַלְעֵ֧י הֶחָצֵ֛ר סָבִ֖יב שֵׁ֥שׁ מׇשְׁזָֽר׃ %
% ---------- פרק 38 | פסוק 17 ----------
 \textbf{יז} וְהָאֲדָנִ֣ים לָֽעַמֻּדִים֮ נְחֹ֒שֶׁת֒ וָוֵ֨י הָֽעַמּוּדִ֜ים וַחֲשׁוּקֵיהֶם֙ כֶּ֔סֶף וְצִפּ֥וּי רָאשֵׁיהֶ֖ם כָּ֑סֶף וְהֵם֙ מְחֻשָּׁקִ֣ים כֶּ֔סֶף כֹּ֖ל עַמֻּדֵ֥י הֶחָצֵֽר׃ %
% ---------- פרק 38 | פסוק 18 ----------
 \textbf{יח} וּמָסַ֞ךְ שַׁ֤עַר הֶחָצֵר֙ מַעֲשֵׂ֣ה רֹקֵ֔ם תְּכֵ֧לֶת וְאַרְגָּמָ֛ן וְתוֹלַ֥עַת שָׁנִ֖י וְשֵׁ֣שׁ מׇשְׁזָ֑ר וְעֶשְׂרִ֤ים אַמָּה֙ אֹ֔רֶךְ וְקוֹמָ֤ה בְרֹ֙חַב֙ חָמֵ֣שׁ אַמּ֔וֹת לְעֻמַּ֖ת קַלְעֵ֥י הֶחָצֵֽר׃ \Rashi{\textbf{לעמת קלעי החצר.}כמדת קלעי החצר:}%
% ---------- פרק 38 | פסוק 19 ----------
 \textbf{יט} וְעַמֻּֽדֵיהֶם֙ אַרְבָּעָ֔ה וְאַדְנֵיהֶ֥ם אַרְבָּעָ֖ה נְחֹ֑שֶׁת וָוֵיהֶ֣ם כֶּ֔סֶף וְצִפּ֧וּי רָאשֵׁיהֶ֛ם וַחֲשֻׁקֵיהֶ֖ם כָּֽסֶף׃ %
% ---------- פרק 38 | פסוק 20 ----------
 \textbf{כ} וְֽכׇל־הַיְתֵדֹ֞ת לַמִּשְׁכָּ֧ן וְלֶחָצֵ֛ר סָבִ֖יב נְחֹֽשֶׁת׃ \{ס\} %
\addparasha{חומש שמות | פרשת פקודי}
% ---------- פרק 38 | פסוק 21 ----------
 \textbf{כא} אֵ֣לֶּה פְקוּדֵ֤י הַמִּשְׁכָּן֙ מִשְׁכַּ֣ן הָעֵדֻ֔ת אֲשֶׁ֥ר פֻּקַּ֖ד עַל־פִּ֣י מֹשֶׁ֑ה עֲבֹדַת֙ הַלְוִיִּ֔ם בְּיַד֙ אִֽיתָמָ֔ר בֶּֽן־אַהֲרֹ֖ן הַכֹּהֵֽן׃ \Rashi{\textbf{אלה פקודי.}בפרשה זו נמנו כל משקלי נדבת המשכן, לכסף ולזהב ולנחשת, ונמנו כל כליו לכל עבודתו: המשכן משכן. שני פעמים, רמז למקדש שנתמשכן בשני חרבנין על עונותיהן של ישראל (תנחומא): משכן העדת. עדות לישראל שותר להם הקב"ה על מעשה העגל, שהרי השרה שכינתו ביניהם: עבדת הלוים. פקודי המשכן וכליו היא עבודה המסורה ללוים במדבר – לשאת ולהוריד ולהקים, איש איש למשאו המפקד עליו, כמו שאמור בפ' נשא: ביד איתמר. הוא היה פקיד עליהם למסר לכל בית אב עבודה שעליו:}%
% ---------- פרק 38 | פסוק 22 ----------
 \textbf{כב} וּבְצַלְאֵ֛ל בֶּן־אוּרִ֥י בֶן־ח֖וּר לְמַטֵּ֣ה יְהוּדָ֑ה עָשָׂ֕ה אֵ֛ת כׇּל־אֲשֶׁר־צִוָּ֥ה יְהֹוָ֖ה אֶת־מֹשֶֽׁה׃ \Rashi{\textbf{ובצלאל בן אורי וגו' עשה את כל אשר צוה ה' את משה.}אשר צוה משה אין כתיב כאן, אלא כל אשר צוה ה' את משה, אפלו דברים שלא אמר לו רבו, הסכימה דעתו למה שנאמר למשה בסיני, כי משה צוה לבצלאל לעשות תחלה כלים ואחר כך משכן, אמר לו בצלאל מנהג עולם לעשות תחלה בית ואחר כך משים כלים בתוכו. אמר לו כך שמעתי מפי הקב"ה. אמר לו משה בצל אל היית, כי בודאי כך צוה לי הקב"ה, וכן עשה המשכן תחלה ואחר כך עשה הכלים:}%
% ---------- פרק 38 | פסוק 23 ----------
 \textbf{כג} וְאִתּ֗וֹ אׇהֳלִיאָ֞ב בֶּן־אֲחִיסָמָ֛ךְ לְמַטֵּה־דָ֖ן חָרָ֣שׁ וְחֹשֵׁ֑ב וְרֹקֵ֗ם בַּתְּכֵ֙לֶת֙ וּבָֽאַרְגָּמָ֔ן וּבְתוֹלַ֥עַת הַשָּׁנִ֖י וּבַשֵּֽׁשׁ׃ \{ס\} %
% ---------- פרק 38 | פסוק 24 ----------
 \textbf{כד} כׇּל־הַזָּהָ֗ב הֶֽעָשׂוּי֙ לַמְּלָאכָ֔ה בְּכֹ֖ל מְלֶ֣אכֶת הַקֹּ֑דֶשׁ וַיְהִ֣י ׀ זְהַ֣ב הַתְּנוּפָ֗ה תֵּ֤שַׁע וְעֶשְׂרִים֙ כִּכָּ֔ר וּשְׁבַ֨ע מֵא֧וֹת וּשְׁלֹשִׁ֛ים שֶׁ֖קֶל בְּשֶׁ֥קֶל הַקֹּֽדֶשׁ׃ \Rashi{\textbf{ככר.}ששים מנה, ומנה של קדש כפול היה, הרי הככר ק"כ מנה, והמנה כ"ה סלעים, הרי ככר של קדש שלשת אלפים שקלים, לפיכך מנה בפרוטרוט כל השקלים שפחותין במנינם מג' אלפים, שאין מגיעין לככר:}%
% ---------- פרק 38 | פסוק 25 ----------
 \textbf{כה} וְכֶ֛סֶף פְּקוּדֵ֥י הָעֵדָ֖ה מְאַ֣ת כִּכָּ֑ר וְאֶ֩לֶף֩ וּשְׁבַ֨ע מֵא֜וֹת וַחֲמִשָּׁ֧ה וְשִׁבְעִ֛ים שֶׁ֖קֶל בְּשֶׁ֥קֶל הַקֹּֽדֶשׁ׃ %
% ---------- פרק 38 | פסוק 26 ----------
 \textbf{כו} בֶּ֚קַע לַגֻּלְגֹּ֔לֶת מַחֲצִ֥ית הַשֶּׁ֖קֶל בְּשֶׁ֣קֶל הַקֹּ֑דֶשׁ לְכֹ֨ל הָעֹבֵ֜ר עַל־הַפְּקֻדִ֗ים מִבֶּ֨ן עֶשְׂרִ֤ים שָׁנָה֙ וָמַ֔עְלָה לְשֵׁשׁ־מֵא֥וֹת אֶ֙לֶף֙ וּשְׁלֹ֣שֶׁת אֲלָפִ֔ים וַחֲמֵ֥שׁ מֵא֖וֹת וַחֲמִשִּֽׁים׃ \Rashi{\textbf{בקע.}הוא שם משקל של מחצית השקל: לשש מאות אלף וגו'. כך היו ישראל, וכך עלה מנינם אחר שהוקם המשכן בספר וידבר, ואף עתה בנדבת המשכן כך היו, ומנין חצאי השקלים של שש מאות אלף עולה מאת ככר, כל אחד של שלשת אלפים שקלים, כיצד? שש מאות אלף חצאין הרי הן שלש מאות אלף שלמים, הרי מאת ככר, והשלשת אלפים וה' מאות וחמשים חצאים עולים אלף וז' מאות וע"ה שקלים:}%
% ---------- פרק 38 | פסוק 27 ----------
 \textbf{כז} וַיְהִ֗י מְאַת֙ כִּכַּ֣ר הַכֶּ֔סֶף לָצֶ֗קֶת אֵ֚ת אַדְנֵ֣י הַקֹּ֔דֶשׁ וְאֵ֖ת אַדְנֵ֣י הַפָּרֹ֑כֶת מְאַ֧ת אֲדָנִ֛ים לִמְאַ֥ת הַכִּכָּ֖ר כִּכָּ֥ר לָאָֽדֶן׃ \Rashi{\textbf{לצקת.}כתרגומו לאתכא: את אדני הקדש. של קרשי המשכן, שהם מ"ח קרשים ולהן צ"ו אדנים, ואדני הפרכת ארבעה, הרי מאה, וכל שאר האדנים נחשת כתוב בהם:}%
% ---------- פרק 38 | פסוק 28 ----------
 \textbf{כח} וְאֶת־הָאֶ֜לֶף וּשְׁבַ֤ע הַמֵּאוֹת֙ וַחֲמִשָּׁ֣ה וְשִׁבְעִ֔ים עָשָׂ֥ה וָוִ֖ים לָעַמּוּדִ֑ים וְצִפָּ֥ה רָאשֵׁיהֶ֖ם וְחִשַּׁ֥ק אֹתָֽם׃ \Rashi{\textbf{וצפה ראשיהם.}של עמודים מהן, שבכלן כתיב וצפוי ראשיהם וחשקיהם כסף:}%
% ---------- פרק 38 | פסוק 29 ----------
 \textbf{כט} וּנְחֹ֥שֶׁת הַתְּנוּפָ֖ה שִׁבְעִ֣ים כִּכָּ֑ר וְאַלְפַּ֥יִם וְאַרְבַּע־מֵא֖וֹת שָֽׁקֶל׃ %
% ---------- פרק 38 | פסוק 30 ----------
 \textbf{ל} וַיַּ֣עַשׂ בָּ֗הּ אֶת־אַדְנֵי֙ פֶּ֚תַח אֹ֣הֶל מוֹעֵ֔ד וְאֵת֙ מִזְבַּ֣ח הַנְּחֹ֔שֶׁת וְאֶת־מִכְבַּ֥ר הַנְּחֹ֖שֶׁת אֲשֶׁר־ל֑וֹ וְאֵ֖ת כׇּל־כְּלֵ֥י הַמִּזְבֵּֽחַ׃ %
% ---------- פרק 38 | פסוק 31 ----------
 \textbf{לא} וְאֶת־אַדְנֵ֤י הֶֽחָצֵר֙ סָבִ֔יב וְאֶת־אַדְנֵ֖י שַׁ֣עַר הֶחָצֵ֑ר וְאֵ֨ת כׇּל־יִתְדֹ֧ת הַמִּשְׁכָּ֛ן וְאֶת־כׇּל־יִתְדֹ֥ת הֶחָצֵ֖ר סָבִֽיב׃ %
% ---------- פרק 39 | פסוק 1 ----------
 \textbf{א} וּמִן־הַתְּכֵ֤לֶת וְהָֽאַרְגָּמָן֙ וְתוֹלַ֣עַת הַשָּׁנִ֔י עָשׂ֥וּ בִגְדֵי־שְׂרָ֖ד לְשָׁרֵ֣ת בַּקֹּ֑דֶשׁ וַֽיַּעֲשׂ֞וּ אֶת־בִּגְדֵ֤י הַקֹּ֙דֶשׁ֙ אֲשֶׁ֣ר לְאַהֲרֹ֔ן כַּאֲשֶׁ֛ר צִוָּ֥ה יְהֹוָ֖ה אֶת־מֹשֶֽׁה׃ \{פ\} \Rashi{\textbf{ומן התכלת והארגמן וגו'.}שש לא נאמר כאן, ומכאן אני אומר שאין בגדי שרד הללו בגדי כהנה, שבבגדי כהנה היה שש, אלא הם בגדים שמכסים בהם כלי הקדש בשעת סלוק מסעות, שלא היה בהם שש:}%
% ---------- פרק 39 | פסוק 2 ----------
 \textbf{ב} וַיַּ֖עַשׂ אֶת־הָאֵפֹ֑ד זָהָ֗ב תְּכֵ֧לֶת וְאַרְגָּמָ֛ן וְתוֹלַ֥עַת שָׁנִ֖י וְשֵׁ֥שׁ מׇשְׁזָֽר׃ %
% ---------- פרק 39 | פסוק 3 ----------
 \textbf{ג} וַֽיְרַקְּע֞וּ אֶת־פַּחֵ֣י הַזָּהָב֮ וְקִצֵּ֣ץ פְּתִילִם֒ לַעֲשׂ֗וֹת בְּת֤וֹךְ הַתְּכֵ֙לֶת֙ וּבְת֣וֹךְ הָֽאַרְגָּמָ֔ן וּבְת֛וֹךְ תּוֹלַ֥עַת הַשָּׁנִ֖י וּבְת֣וֹךְ הַשֵּׁ֑שׁ מַעֲשֵׂ֖ה חֹשֵֽׁב׃ \Rashi{\textbf{וירקעו.}כמו לרקע הארץ (תהילים קל"ו), כתרגומו ורדידו, טסין היו מרדדין מן הזהב, אשטנ"דרא בלעז, טסין דקות. כאן הוא מלמדך היאך היו טווין את הזהב עם החוטין, מרדדים טסין דקין, וקוצצין מהן פתילים לארך הטס – לעשות אותן פתילים תערבת עם כל מין ומין בחשן ואפוד שנ' בהן זהב, חוט אחד של זהב עם ו' חוטין של תכלת, וכן עם כל מין ומין, שכל המינים חוטן כפול ו', והזהב חוט שביעי עם כל אחד ואחד:}%
% ---------- פרק 39 | פסוק 4 ----------
 \textbf{ד} כְּתֵפֹ֥ת עָֽשׂוּ־ל֖וֹ חֹבְרֹ֑ת עַל־שְׁנֵ֥י (קצוותו) [קְצוֹתָ֖יו] חֻבָּֽר׃ %
% ---------- פרק 39 | פסוק 5 ----------
 \textbf{ה} וְחֵ֨שֶׁב אֲפֻדָּת֜וֹ אֲשֶׁ֣ר עָלָ֗יו מִמֶּ֣נּוּ הוּא֮ כְּמַעֲשֵׂ֒הוּ֒ זָהָ֗ב תְּכֵ֧לֶת וְאַרְגָּמָ֛ן וְתוֹלַ֥עַת שָׁנִ֖י וְשֵׁ֣שׁ מׇשְׁזָ֑ר כַּאֲשֶׁ֛ר צִוָּ֥ה יְהֹוָ֖ה אֶת־מֹשֶֽׁה׃ \{ס\} %
% ---------- פרק 39 | פסוק 6 ----------
 \textbf{ו} וַֽיַּעֲשׂוּ֙ אֶת־אַבְנֵ֣י הַשֹּׁ֔הַם מֻֽסַבֹּ֖ת מִשְׁבְּצֹ֣ת זָהָ֑ב מְפֻתָּחֹת֙ פִּתּוּחֵ֣י חוֹתָ֔ם עַל־שְׁמ֖וֹת בְּנֵ֥י יִשְׂרָאֵֽל׃ %
% ---------- פרק 39 | פסוק 7 ----------
 \textbf{ז} וַיָּ֣שֶׂם אֹתָ֗ם עַ֚ל כִּתְפֹ֣ת הָאֵפֹ֔ד אַבְנֵ֥י זִכָּר֖וֹן לִבְנֵ֣י יִשְׂרָאֵ֑ל כַּאֲשֶׁ֛ר צִוָּ֥ה יְהֹוָ֖ה אֶת־מֹשֶֽׁה׃ \{פ\} %
% ---------- פרק 39 | פסוק 8 ----------
 \textbf{ח} וַיַּ֧עַשׂ אֶת־הַחֹ֛שֶׁן מַעֲשֵׂ֥ה חֹשֵׁ֖ב כְּמַעֲשֵׂ֣ה אֵפֹ֑ד זָהָ֗ב תְּכֵ֧לֶת וְאַרְגָּמָ֛ן וְתוֹלַ֥עַת שָׁנִ֖י וְשֵׁ֥שׁ מׇשְׁזָֽר׃ %
% ---------- פרק 39 | פסוק 9 ----------
 \textbf{ט} רָב֧וּעַ הָיָ֛ה כָּפ֖וּל עָשׂ֣וּ אֶת־הַחֹ֑שֶׁן זֶ֧רֶת אׇרְכּ֛וֹ וְזֶ֥רֶת רׇחְבּ֖וֹ כָּפֽוּל׃ %
% ---------- פרק 39 | פסוק 10 ----------
 \textbf{י} וַיְמַ֨לְאוּ־ב֔וֹ אַרְבָּעָ֖ה ט֣וּרֵי אָ֑בֶן ט֗וּר אֹ֤דֶם פִּטְדָה֙ וּבָרֶ֔קֶת הַטּ֖וּר הָאֶחָֽד׃ %
% ---------- פרק 39 | פסוק 11 ----------
 \textbf{יא} וְהַטּ֖וּר הַשֵּׁנִ֑י נֹ֥פֶךְ סַפִּ֖יר וְיָהֲלֹֽם׃ %
% ---------- פרק 39 | פסוק 12 ----------
 \textbf{יב} וְהַטּ֖וּר הַשְּׁלִישִׁ֑י לֶ֥שֶׁם שְׁב֖וֹ וְאַחְלָֽמָה׃ %
% ---------- פרק 39 | פסוק 13 ----------
 \textbf{יג} וְהַטּוּר֙ הָֽרְבִיעִ֔י תַּרְשִׁ֥ישׁ שֹׁ֖הַם וְיָשְׁפֵ֑ה מֽוּסַבֹּ֛ת מִשְׁבְּצֹ֥ת זָהָ֖ב בְּמִלֻּאֹתָֽם׃ %
% ---------- פרק 39 | פסוק 14 ----------
 \textbf{יד} וְ֠הָאֲבָנִ֠ים עַל־שְׁמֹ֨ת בְּנֵי־יִשְׂרָאֵ֥ל הֵ֛נָּה שְׁתֵּ֥ים עֶשְׂרֵ֖ה עַל־שְׁמֹתָ֑ם פִּתּוּחֵ֤י חֹתָם֙ אִ֣ישׁ עַל־שְׁמ֔וֹ לִשְׁנֵ֥ים עָשָׂ֖ר שָֽׁבֶט׃ %
% ---------- פרק 39 | פסוק 15 ----------
 \textbf{טו} וַיַּעֲשׂ֧וּ עַל־הַחֹ֛שֶׁן שַׁרְשְׁרֹ֥ת גַּבְלֻ֖ת מַעֲשֵׂ֣ה עֲבֹ֑ת זָהָ֖ב טָהֽוֹר׃ %
% ---------- פרק 39 | פסוק 16 ----------
 \textbf{טז} וַֽיַּעֲשׂ֗וּ שְׁתֵּי֙ מִשְׁבְּצֹ֣ת זָהָ֔ב וּשְׁתֵּ֖י טַבְּעֹ֣ת זָהָ֑ב וַֽיִּתְּנ֗וּ אֶת־שְׁתֵּי֙ הַטַּבָּעֹ֔ת עַל־שְׁנֵ֖י קְצ֥וֹת הַחֹֽשֶׁן׃ %
% ---------- פרק 39 | פסוק 17 ----------
 \textbf{יז} וַֽיִּתְּנ֗וּ שְׁתֵּי֙ הָעֲבֹתֹ֣ת הַזָּהָ֔ב עַל־שְׁתֵּ֖י הַטַּבָּעֹ֑ת עַל־קְצ֖וֹת הַחֹֽשֶׁן׃ %
% ---------- פרק 39 | פסוק 18 ----------
 \textbf{יח} וְאֵ֨ת שְׁתֵּ֤י קְצוֹת֙ שְׁתֵּ֣י הָֽעֲבֹתֹ֔ת נָתְנ֖וּ עַל־שְׁתֵּ֣י הַֽמִּשְׁבְּצֹ֑ת וַֽיִּתְּנֻ֛ם עַל־כִּתְפֹ֥ת הָאֵפֹ֖ד אֶל־מ֥וּל פָּנָֽיו׃ %
% ---------- פרק 39 | פסוק 19 ----------
 \textbf{יט} וַֽיַּעֲשׂ֗וּ שְׁתֵּי֙ טַבְּעֹ֣ת זָהָ֔ב וַיָּשִׂ֕ימוּ עַל־שְׁנֵ֖י קְצ֣וֹת הַחֹ֑שֶׁן עַל־שְׂפָת֕וֹ אֲשֶׁ֛ר אֶל־עֵ֥בֶר הָאֵפֹ֖ד בָּֽיְתָה׃ %
% ---------- פרק 39 | פסוק 20 ----------
 \textbf{כ} וַֽיַּעֲשׂוּ֮ שְׁתֵּ֣י טַבְּעֹ֣ת זָהָב֒ וַֽיִּתְּנֻ֡ם עַל־שְׁתֵּי֩ כִתְפֹ֨ת הָאֵפֹ֤ד מִלְּמַ֙טָּה֙ מִמּ֣וּל פָּנָ֔יו לְעֻמַּ֖ת מַחְבַּרְתּ֑וֹ מִמַּ֕עַל לְחֵ֖שֶׁב הָאֵפֹֽד׃ %
% ---------- פרק 39 | פסוק 21 ----------
 \textbf{כא} וַיִּרְכְּס֣וּ אֶת־הַחֹ֡שֶׁן מִטַּבְּעֹתָיו֩ אֶל־טַבְּעֹ֨ת הָאֵפֹ֜ד בִּפְתִ֣יל תְּכֵ֗לֶת לִֽהְיֹת֙ עַל־חֵ֣שֶׁב הָאֵפֹ֔ד וְלֹֽא־יִזַּ֣ח הַחֹ֔שֶׁן מֵעַ֖ל הָאֵפֹ֑ד כַּאֲשֶׁ֛ר צִוָּ֥ה יְהֹוָ֖ה אֶת־מֹשֶֽׁה׃ \{פ\} %
% ---------- פרק 39 | פסוק 22 ----------
 \textbf{כב} וַיַּ֛עַשׂ אֶת־מְעִ֥יל הָאֵפֹ֖ד מַעֲשֵׂ֣ה אֹרֵ֑ג כְּלִ֖יל תְּכֵֽלֶת׃ %
% ---------- פרק 39 | פסוק 23 ----------
 \textbf{כג} וּפִֽי־הַמְּעִ֥יל בְּתוֹכ֖וֹ כְּפִ֣י תַחְרָ֑א שָׂפָ֥ה לְפִ֛יו סָבִ֖יב לֹ֥א יִקָּרֵֽעַ׃ %
% ---------- פרק 39 | פסוק 24 ----------
 \textbf{כד} וַֽיַּעֲשׂוּ֙ עַל־שׁוּלֵ֣י הַמְּעִ֔יל רִמּוֹנֵ֕י תְּכֵ֥לֶת וְאַרְגָּמָ֖ן וְתוֹלַ֣עַת שָׁנִ֑י מׇשְׁזָֽר׃ %
% ---------- פרק 39 | פסוק 25 ----------
 \textbf{כה} וַיַּעֲשׂ֥וּ פַעֲמֹנֵ֖י זָהָ֣ב טָה֑וֹר וַיִּתְּנ֨וּ אֶת־הַפַּֽעֲמֹנִ֜ים בְּת֣וֹךְ הָרִמֹּנִ֗ים עַל־שׁוּלֵ֤י הַמְּעִיל֙ סָבִ֔יב בְּת֖וֹךְ הָרִמֹּנִֽים׃ %
% ---------- פרק 39 | פסוק 26 ----------
 \textbf{כו} פַּעֲמֹ֤ן וְרִמֹּן֙ פַּעֲמֹ֣ן וְרִמֹּ֔ן עַל־שׁוּלֵ֥י הַמְּעִ֖יל סָבִ֑יב לְשָׁרֵ֕ת כַּאֲשֶׁ֛ר צִוָּ֥ה יְהֹוָ֖ה אֶת־מֹשֶֽׁה׃ \{ס\} %
% ---------- פרק 39 | פסוק 27 ----------
 \textbf{כז} וַֽיַּעֲשׂ֛וּ אֶת־הַכׇּתְנֹ֥ת שֵׁ֖שׁ מַעֲשֵׂ֣ה אֹרֵ֑ג לְאַהֲרֹ֖ן וּלְבָנָֽיו׃ %
% ---------- פרק 39 | פסוק 28 ----------
 \textbf{כח} וְאֵת֙ הַמִּצְנֶ֣פֶת שֵׁ֔שׁ וְאֶת־פַּאֲרֵ֥י הַמִּגְבָּעֹ֖ת שֵׁ֑שׁ וְאֶת־מִכְנְסֵ֥י הַבָּ֖ד שֵׁ֥שׁ מׇשְׁזָֽר׃ \Rashi{\textbf{ואת פארי המגבעת.}תפארת המגבעות – המגבעות המפארות:}%
% ---------- פרק 39 | פסוק 29 ----------
 \textbf{כט} וְֽאֶת־הָאַבְנֵ֞ט שֵׁ֣שׁ מׇשְׁזָ֗ר וּתְכֵ֧לֶת וְאַרְגָּמָ֛ן וְתוֹלַ֥עַת שָׁנִ֖י מַעֲשֵׂ֣ה רֹקֵ֑ם כַּאֲשֶׁ֛ר צִוָּ֥ה יְהֹוָ֖ה אֶת־מֹשֶֽׁה׃ \{ס\} %
% ---------- פרק 39 | פסוק 30 ----------
 \textbf{ל} וַֽיַּעֲשׂ֛וּ אֶת־צִ֥יץ נֵֽזֶר־הַקֹּ֖דֶשׁ זָהָ֣ב טָה֑וֹר וַיִּכְתְּב֣וּ עָלָ֗יו מִכְתַּב֙ פִּתּוּחֵ֣י חוֹתָ֔ם קֹ֖דֶשׁ לַיהֹוָֽה׃ %
% ---------- פרק 39 | פסוק 31 ----------
 \textbf{לא} וַיִּתְּנ֤וּ עָלָיו֙ פְּתִ֣יל תְּכֵ֔לֶת לָתֵ֥ת עַל־הַמִּצְנֶ֖פֶת מִלְמָ֑עְלָה כַּאֲשֶׁ֛ר צִוָּ֥ה יְהֹוָ֖ה אֶת־מֹשֶֽׁה׃ \{ס\} \Rashi{\textbf{לתת על המצנפת מלמעלה.}וע"י הפתילים היה מושיבו על המצנפת כמין כתר; וא"א לומר הציץ על המצנפת, שהרי בשחיטת קדשים שנינו שערו היה נראה בין ציץ למצנפת ששם מניח תפלין, והציץ היה נתון על המצח, הרי המצנפת למעלה והציץ למטה, ומהו על המצנפת מלמעלה? ועוד הקשיתי בה – כאן הוא אומר ויתנו עליו פתיל תכלת, ובענין הצואה הוא אומר ושמת אותו על פתיל תכלת? ואומר אני, פתיל תכלת זה חוטין הן, לקשרו בהן במצנפת – לפי שהציץ אינו אלא מאזן לאזן ובמה יקשרנו במצחו? – והיו קבועין בו חוטי תכלת לשני ראשיו ובאמצעיתו, שבהן קושרו ותולהו במצנפת כשהוא בראשו, ושני חוטין היו בכל קצה וקצה, אחד ממעל ואחד מתחת לצד מצחו, וכן באמצעו, שכך הוא נח לקשר, ואין דרך קשירה בפחות משני חוטין, לכך נאמר על פתיל תכלת, ועליו פתיל תכלת, וקושר ראשיהם השנים כלן יחד מאחוריו למול ערפו, ומושיבו על המצנפת; ואל תתמה שלא נאמר פתילי תכלת הואיל ומרבין הן, שהרי מצינו בחשן ואפוד וירכסו את החשן … בפתיל תכלת, ועל כרחך פחות משנים לא היו שהרי בשני קצות החשן היו ב' טבעות החשן, ובשתי כתפות האפוד היו ב' טבעות האפוד, שכנגדן, ולפי דרך קשירה ד' חוטין היו ומכל מקום פחות משנים א"א:}%
% ---------- פרק 39 | פסוק 32 ----------
 \textbf{לב} וַתֵּ֕כֶל כׇּל־עֲבֹדַ֕ת מִשְׁכַּ֖ן אֹ֣הֶל מוֹעֵ֑ד וַֽיַּעֲשׂוּ֙ בְּנֵ֣י יִשְׂרָאֵ֔ל כְּ֠כֹ֠ל אֲשֶׁ֨ר צִוָּ֧ה יְהֹוָ֛ה אֶת־מֹשֶׁ֖ה כֵּ֥ן עָשֽׂוּ׃ \{פ\} \Rashi{\textbf{ויעשו בני ישראל.}את המלאכה ככל אשר צוה ה' וגו':}%
% ---------- פרק 39 | פסוק 33 ----------
 \textbf{לג} וַיָּבִ֤יאוּ אֶת־הַמִּשְׁכָּן֙ אֶל־מֹשֶׁ֔ה אֶת־הָאֹ֖הֶל וְאֶת־כׇּל־כֵּלָ֑יו קְרָסָ֣יו קְרָשָׁ֔יו בְּרִיחָ֖ו וְעַמֻּדָ֥יו וַאֲדָנָֽיו׃ \Rashi{\textbf{ויביאו את המשכן וגו'.}שלא היו יכולין להקימו; ולפי שלא עשה משה שום מלאכה במשכן, הניח לו הקב"ה הקמתו, שלא היה יכול להקימו שום אדם מחמת כבד הקרשים, שאין כח באדם לזקפן, ומשה העמידו; אמר משה לפני הקב"ה איך אפשר הקמתו ע"י אדם? אמר לו עסק אתה בידך, ונראה כמקימו והוא נזקף וקם מאליו, וזהו שנ' הוקם המשכן (שמות מ') – הוקם מאליו; מדרש ר' תנחומא:}%
% ---------- פרק 39 | פסוק 34 ----------
 \textbf{לד} וְאֶת־מִכְסֵ֞ה עוֹרֹ֤ת הָֽאֵילִם֙ הַמְאׇדָּמִ֔ים וְאֶת־מִכְסֵ֖ה עֹרֹ֣ת הַתְּחָשִׁ֑ים וְאֵ֖ת פָּרֹ֥כֶת הַמָּסָֽךְ׃ %
% ---------- פרק 39 | פסוק 35 ----------
 \textbf{לה} אֶת־אֲר֥וֹן הָעֵדֻ֖ת וְאֶת־בַּדָּ֑יו וְאֵ֖ת הַכַּפֹּֽרֶת׃ %
% ---------- פרק 39 | פסוק 36 ----------
 \textbf{לו} אֶת־הַשֻּׁלְחָן֙ אֶת־כׇּל־כֵּלָ֔יו וְאֵ֖ת לֶ֥חֶם הַפָּנִֽים׃ %
% ---------- פרק 39 | פסוק 37 ----------
 \textbf{לז} אֶת־הַמְּנֹרָ֨ה הַטְּהֹרָ֜ה אֶת־נֵרֹתֶ֗יהָ נֵרֹ֛ת הַמַּֽעֲרָכָ֖ה וְאֶת־כׇּל־כֵּלֶ֑יהָ וְאֵ֖ת שֶׁ֥מֶן הַמָּאֽוֹר׃ %
% ---------- פרק 39 | פסוק 38 ----------
 \textbf{לח} וְאֵת֙ מִזְבַּ֣ח הַזָּהָ֔ב וְאֵת֙ שֶׁ֣מֶן הַמִּשְׁחָ֔ה וְאֵ֖ת קְטֹ֣רֶת הַסַּמִּ֑ים וְאֵ֕ת מָסַ֖ךְ פֶּ֥תַח הָאֹֽהֶל׃ %
% ---------- פרק 39 | פסוק 39 ----------
 \textbf{לט} אֵ֣ת ׀ מִזְבַּ֣ח הַנְּחֹ֗שֶׁת וְאֶת־מִכְבַּ֤ר הַנְּחֹ֙שֶׁת֙ אֲשֶׁר־ל֔וֹ אֶת־בַּדָּ֖יו וְאֶת־כׇּל־כֵּלָ֑יו אֶת־הַכִּיֹּ֖ר וְאֶת־כַּנּֽוֹ׃ %
% ---------- פרק 39 | פסוק 40 ----------
 \textbf{מ} אֵת֩ קַלְעֵ֨י הֶחָצֵ֜ר אֶת־עַמֻּדֶ֣יהָ וְאֶת־אֲדָנֶ֗יהָ וְאֶת־הַמָּסָךְ֙ לְשַׁ֣עַר הֶֽחָצֵ֔ר אֶת־מֵיתָרָ֖יו וִיתֵדֹתֶ֑יהָ וְאֵ֗ת כׇּל־כְּלֵ֛י עֲבֹדַ֥ת הַמִּשְׁכָּ֖ן לְאֹ֥הֶל מוֹעֵֽד׃ %
% ---------- פרק 39 | פסוק 41 ----------
 \textbf{מא} אֶת־בִּגְדֵ֥י הַשְּׂרָ֖ד לְשָׁרֵ֣ת בַּקֹּ֑דֶשׁ אֶת־בִּגְדֵ֤י הַקֹּ֙דֶשׁ֙ לְאַהֲרֹ֣ן הַכֹּהֵ֔ן וְאֶת־בִּגְדֵ֥י בָנָ֖יו לְכַהֵֽן׃ %
% ---------- פרק 39 | פסוק 42 ----------
 \textbf{מב} כְּכֹ֛ל אֲשֶׁר־צִוָּ֥ה יְהֹוָ֖ה אֶת־מֹשֶׁ֑ה כֵּ֤ן עָשׂוּ֙ בְּנֵ֣י יִשְׂרָאֵ֔ל אֵ֖ת כׇּל־הָעֲבֹדָֽה׃ %
% ---------- פרק 39 | פסוק 43 ----------
 \textbf{מג} וַיַּ֨רְא מֹשֶׁ֜ה אֶת־כׇּל־הַמְּלָאכָ֗ה וְהִנֵּה֙ עָשׂ֣וּ אֹתָ֔הּ כַּאֲשֶׁ֛ר צִוָּ֥ה יְהֹוָ֖ה כֵּ֣ן עָשׂ֑וּ וַיְבָ֥רֶךְ אֹתָ֖ם מֹשֶֽׁה׃ \{פ\} \Rashi{\textbf{ויברך אתם משה.}אמר להם יהי רצון שתשרה שכינה במעשה ידיכם, ויהי נעם ה' אלהינו עלינו וגו', והוא אחד מי"א מזמורים שבתפלה למשה (ספרא):}%
% ---------- פרק 40 | פסוק 1 ----------
 \textbf{א} וַיְדַבֵּ֥ר יְהֹוָ֖ה אֶל־מֹשֶׁ֥ה לֵּאמֹֽר׃ %
% ---------- פרק 40 | פסוק 2 ----------
 \textbf{ב} בְּיוֹם־הַחֹ֥דֶשׁ הָרִאשׁ֖וֹן בְּאֶחָ֣ד לַחֹ֑דֶשׁ תָּקִ֕ים אֶת־מִשְׁכַּ֖ן אֹ֥הֶל מוֹעֵֽד׃ %
% ---------- פרק 40 | פסוק 3 ----------
 \textbf{ג} וְשַׂמְתָּ֣ שָׁ֔ם אֵ֖ת אֲר֣וֹן הָעֵד֑וּת וְסַכֹּתָ֥ עַל־הָאָרֹ֖ן אֶת־הַפָּרֹֽכֶת׃ \Rashi{\textbf{וסכת על הארן.}לשון הגנה, שהרי מחצה היתה:}%
% ---------- פרק 40 | פסוק 4 ----------
 \textbf{ד} וְהֵבֵאתָ֙ אֶת־הַשֻּׁלְחָ֔ן וְעָרַכְתָּ֖ אֶת־עֶרְכּ֑וֹ וְהֵבֵאתָ֙ אֶת־הַמְּנֹרָ֔ה וְהַעֲלֵיתָ֖ אֶת־נֵרֹתֶֽיהָ׃ \Rashi{\textbf{וערכת את ערכו.}שתים מערכות של לחם הפנים:}%
% ---------- פרק 40 | פסוק 5 ----------
 \textbf{ה} וְנָתַתָּ֞ה אֶת־מִזְבַּ֤ח הַזָּהָב֙ לִקְטֹ֔רֶת לִפְנֵ֖י אֲר֣וֹן הָעֵדֻ֑ת וְשַׂמְתָּ֛ אֶת־מָסַ֥ךְ הַפֶּ֖תַח לַמִּשְׁכָּֽן׃ %
% ---------- פרק 40 | פסוק 6 ----------
 \textbf{ו} וְנָ֣תַתָּ֔ה אֵ֖ת מִזְבַּ֣ח הָעֹלָ֑ה לִפְנֵ֕י פֶּ֖תַח מִשְׁכַּ֥ן אֹֽהֶל־מוֹעֵֽד׃ %
% ---------- פרק 40 | פסוק 7 ----------
 \textbf{ז} וְנָֽתַתָּ֙ אֶת־הַכִּיֹּ֔ר בֵּֽין־אֹ֥הֶל מוֹעֵ֖ד וּבֵ֣ין הַמִּזְבֵּ֑חַ וְנָתַתָּ֥ שָׁ֖ם מָֽיִם׃ %
% ---------- פרק 40 | פסוק 8 ----------
 \textbf{ח} וְשַׂמְתָּ֥ אֶת־הֶחָצֵ֖ר סָבִ֑יב וְנָ֣תַתָּ֔ אֶת־מָסַ֖ךְ שַׁ֥עַר הֶחָצֵֽר׃ %
% ---------- פרק 40 | פסוק 9 ----------
 \textbf{ט} וְלָקַחְתָּ֙ אֶת־שֶׁ֣מֶן הַמִּשְׁחָ֔ה וּמָשַׁחְתָּ֥ אֶת־הַמִּשְׁכָּ֖ן וְאֶת־כׇּל־אֲשֶׁר־בּ֑וֹ וְקִדַּשְׁתָּ֥ אֹת֛וֹ וְאֶת־כׇּל־כֵּלָ֖יו וְהָ֥יָה קֹֽדֶשׁ׃ %
% ---------- פרק 40 | פסוק 10 ----------
 \textbf{י} וּמָשַׁחְתָּ֛ אֶת־מִזְבַּ֥ח הָעֹלָ֖ה וְאֶת־כׇּל־כֵּלָ֑יו וְקִדַּשְׁתָּ֙ אֶת־הַמִּזְבֵּ֔חַ וְהָיָ֥ה הַמִּזְבֵּ֖חַ קֹ֥דֶשׁ קׇֽדָשִֽׁים׃ %
% ---------- פרק 40 | פסוק 11 ----------
 \textbf{יא} וּמָשַׁחְתָּ֥ אֶת־הַכִּיֹּ֖ר וְאֶת־כַּנּ֑וֹ וְקִדַּשְׁתָּ֖ אֹתֽוֹ׃ %
% ---------- פרק 40 | פסוק 12 ----------
 \textbf{יב} וְהִקְרַבְתָּ֤ אֶֽת־אַהֲרֹן֙ וְאֶת־בָּנָ֔יו אֶל־פֶּ֖תַח אֹ֣הֶל מוֹעֵ֑ד וְרָחַצְתָּ֥ אֹתָ֖ם בַּמָּֽיִם׃ %
% ---------- פרק 40 | פסוק 13 ----------
 \textbf{יג} וְהִלְבַּשְׁתָּ֙ אֶֽת־אַהֲרֹ֔ן אֵ֖ת בִּגְדֵ֣י הַקֹּ֑דֶשׁ וּמָשַׁחְתָּ֥ אֹת֛וֹ וְקִדַּשְׁתָּ֥ אֹת֖וֹ וְכִהֵ֥ן לִֽי׃ %
% ---------- פרק 40 | פסוק 14 ----------
 \textbf{יד} וְאֶת־בָּנָ֖יו תַּקְרִ֑יב וְהִלְבַּשְׁתָּ֥ אֹתָ֖ם כֻּתֳּנֹֽת׃ %
% ---------- פרק 40 | פסוק 15 ----------
 \textbf{טו} וּמָשַׁחְתָּ֣ אֹתָ֗ם כַּאֲשֶׁ֤ר מָשַׁ֙חְתָּ֙ אֶת־אֲבִיהֶ֔ם וְכִהֲנ֖וּ לִ֑י וְ֠הָיְתָ֠ה לִהְיֹ֨ת לָהֶ֧ם מׇשְׁחָתָ֛ם לִכְהֻנַּ֥ת עוֹלָ֖ם לְדֹרֹתָֽם׃ %
% ---------- פרק 40 | פסוק 16 ----------
 \textbf{טז} וַיַּ֖עַשׂ מֹשֶׁ֑ה כְּ֠כֹ֠ל אֲשֶׁ֨ר צִוָּ֧ה יְהֹוָ֛ה אֹת֖וֹ כֵּ֥ן עָשָֽׂה׃ \{ס\} %
% ---------- פרק 40 | פסוק 17 ----------
 \textbf{יז} וַיְהִ֞י בַּחֹ֧דֶשׁ הָרִאשׁ֛וֹן בַּשָּׁנָ֥ה הַשֵּׁנִ֖ית בְּאֶחָ֣ד לַחֹ֑דֶשׁ הוּקַ֖ם הַמִּשְׁכָּֽן׃ %
% ---------- פרק 40 | פסוק 18 ----------
 \textbf{יח} וַיָּ֨קֶם מֹשֶׁ֜ה אֶת־הַמִּשְׁכָּ֗ן וַיִּתֵּן֙ אֶת־אֲדָנָ֔יו וַיָּ֙שֶׂם֙ אֶת־קְרָשָׁ֔יו וַיִּתֵּ֖ן אֶת־בְּרִיחָ֑יו וַיָּ֖קֶם אֶת־עַמּוּדָֽיו׃ %
% ---------- פרק 40 | פסוק 19 ----------
 \textbf{יט} וַיִּפְרֹ֤שׂ אֶת־הָאֹ֙הֶל֙ עַל־הַמִּשְׁכָּ֔ן וַיָּ֜שֶׂם אֶת־מִכְסֵ֥ה הָאֹ֛הֶל עָלָ֖יו מִלְמָ֑עְלָה כַּאֲשֶׁ֛ר צִוָּ֥ה יְהֹוָ֖ה אֶת־מֹשֶֽׁה׃ \{ס\} \Rashi{\textbf{ויפרש את האהל.}הן יריעות העזים:}%
% ---------- פרק 40 | פסוק 20 ----------
 \textbf{כ} וַיִּקַּ֞ח וַיִּתֵּ֤ן אֶת־הָעֵדֻת֙ אֶל־הָ֣אָרֹ֔ן וַיָּ֥שֶׂם אֶת־הַבַּדִּ֖ים עַל־הָאָרֹ֑ן וַיִּתֵּ֧ן אֶת־הַכַּפֹּ֛רֶת עַל־הָאָרֹ֖ן מִלְמָֽעְלָה׃ \Rashi{\textbf{את העדת.}הלוחות:}%
% ---------- פרק 40 | פסוק 21 ----------
 \textbf{כא} וַיָּבֵ֣א אֶת־הָאָרֹן֮ אֶל־הַמִּשְׁכָּן֒ וַיָּ֗שֶׂם אֵ֚ת פָּרֹ֣כֶת הַמָּסָ֔ךְ וַיָּ֕סֶךְ עַ֖ל אֲר֣וֹן הָעֵד֑וּת כַּאֲשֶׁ֛ר צִוָּ֥ה יְהֹוָ֖ה אֶת־מֹשֶֽׁה׃ \{ס\} %
% ---------- פרק 40 | פסוק 22 ----------
 \textbf{כב} וַיִּתֵּ֤ן אֶת־הַשֻּׁלְחָן֙ בְּאֹ֣הֶל מוֹעֵ֔ד עַ֛ל יֶ֥רֶךְ הַמִּשְׁכָּ֖ן צָפֹ֑נָה מִח֖וּץ לַפָּרֹֽכֶת׃ \Rashi{\textbf{על ירך המשכן צפנה.}בחצי הצפוני של רחב הבית: ירך. כתרגומו, צדא, כירך הזה שהוא בצדו של אדם:}%
% ---------- פרק 40 | פסוק 23 ----------
 \textbf{כג} וַיַּעֲרֹ֥ךְ עָלָ֛יו עֵ֥רֶךְ לֶ֖חֶם לִפְנֵ֣י יְהֹוָ֑ה כַּאֲשֶׁ֛ר צִוָּ֥ה יְהֹוָ֖ה אֶת־מֹשֶֽׁה׃ \{ס\} %
% ---------- פרק 40 | פסוק 24 ----------
 \textbf{כד} וַיָּ֤שֶׂם אֶת־הַמְּנֹרָה֙ בְּאֹ֣הֶל מוֹעֵ֔ד נֹ֖כַח הַשֻּׁלְחָ֑ן עַ֛ל יֶ֥רֶךְ הַמִּשְׁכָּ֖ן נֶֽגְבָּה׃ %
% ---------- פרק 40 | פסוק 25 ----------
 \textbf{כה} וַיַּ֥עַל הַנֵּרֹ֖ת לִפְנֵ֣י יְהֹוָ֑ה כַּאֲשֶׁ֛ר צִוָּ֥ה יְהֹוָ֖ה אֶת־מֹשֶֽׁה׃ \{ס\} %
% ---------- פרק 40 | פסוק 26 ----------
 \textbf{כו} וַיָּ֛שֶׂם אֶת־מִזְבַּ֥ח הַזָּהָ֖ב בְּאֹ֣הֶל מוֹעֵ֑ד לִפְנֵ֖י הַפָּרֹֽכֶת׃ %
% ---------- פרק 40 | פסוק 27 ----------
 \textbf{כז} וַיַּקְטֵ֥ר עָלָ֖יו קְטֹ֣רֶת סַמִּ֑ים כַּאֲשֶׁ֛ר צִוָּ֥ה יְהֹוָ֖ה אֶת־מֹשֶֽׁה׃ \{ס\} \Rashi{\textbf{ויקטר עליו קטרת.}שחרית וערבית, כמו שנאמר בבקר בבקר בהיטיבו את הנרות וגומר (שמות ל'):}%
% ---------- פרק 40 | פסוק 28 ----------
 \textbf{כח} וַיָּ֛שֶׂם אֶת־מָסַ֥ךְ הַפֶּ֖תַח לַמִּשְׁכָּֽן׃ %
% ---------- פרק 40 | פסוק 29 ----------
 \textbf{כט} וְאֵת֙ מִזְבַּ֣ח הָעֹלָ֔ה שָׂ֕ם פֶּ֖תַח מִשְׁכַּ֣ן אֹֽהֶל־מוֹעֵ֑ד וַיַּ֣עַל עָלָ֗יו אֶת־הָעֹלָה֙ וְאֶת־הַמִּנְחָ֔ה כַּאֲשֶׁ֛ר צִוָּ֥ה יְהֹוָ֖ה אֶת־מֹשֶֽׁה׃ \{ס\} \Rashi{\textbf{ויעל עליו וגו'.}אף ביום השמיני למלואים, שהוא יום הקמת המשכן, שמש משה והקריב קרבנות צבור, חוץ מאותן שנצטוו בו ביום, שנאמר קרב אל המזבח וגומר (ויקרא ט'): את העלה. עולת התמיד: ואת המנחה. מנחת נסכים של תמיד, כמו שנאמר ועשרן סלת בלול בשמן וגומר (שמות כ"ט):}%
% ---------- פרק 40 | פסוק 30 ----------
 \textbf{ל} וַיָּ֙שֶׂם֙ אֶת־הַכִּיֹּ֔ר בֵּֽין־אֹ֥הֶל מוֹעֵ֖ד וּבֵ֣ין הַמִּזְבֵּ֑חַ וַיִּתֵּ֥ן שָׁ֛מָּה מַ֖יִם לְרׇחְצָֽה׃ %
% ---------- פרק 40 | פסוק 31 ----------
 \textbf{לא} וְרָחֲצ֣וּ מִמֶּ֔נּוּ מֹשֶׁ֖ה וְאַהֲרֹ֣ן וּבָנָ֑יו אֶת־יְדֵיהֶ֖ם וְאֶת־רַגְלֵיהֶֽם׃ \Rashi{\textbf{ורחצו ממנו משה ואהרן.}יום שמיני למלואים השוו כלם לכהנה, ותרגומו ומקדשין מנה, בו ביום קדש משה עמהם:}%
% ---------- פרק 40 | פסוק 32 ----------
 \textbf{לב} בְּבֹאָ֞ם אֶל־אֹ֣הֶל מוֹעֵ֗ד וּבְקׇרְבָתָ֛ם אֶל־הַמִּזְבֵּ֖חַ יִרְחָ֑צוּ כַּאֲשֶׁ֛ר צִוָּ֥ה יְהֹוָ֖ה אֶת־מֹשֶֽׁה׃ \{ס\} \Rashi{\textbf{ובקרבתם.}כמו ובקרבם – כשיקרבו:}%
% ---------- פרק 40 | פסוק 33 ----------
 \textbf{לג} וַיָּ֣קֶם אֶת־הֶחָצֵ֗ר סָבִיב֙ לַמִּשְׁכָּ֣ן וְלַמִּזְבֵּ֔חַ וַיִּתֵּ֕ן אֶת־מָסַ֖ךְ שַׁ֣עַר הֶחָצֵ֑ר וַיְכַ֥ל מֹשֶׁ֖ה אֶת־הַמְּלָאכָֽה׃ \{פ\} %
% ---------- פרק 40 | פסוק 34 ----------
 \textbf{לד} וַיְכַ֥ס הֶעָנָ֖ן אֶת־אֹ֣הֶל מוֹעֵ֑ד וּכְב֣וֹד יְהֹוָ֔ה מָלֵ֖א אֶת־הַמִּשְׁכָּֽן׃ %
% ---------- פרק 40 | פסוק 35 ----------
 \textbf{לה} וְלֹא־יָכֹ֣ל מֹשֶׁ֗ה לָבוֹא֙ אֶל־אֹ֣הֶל מוֹעֵ֔ד כִּֽי־שָׁכַ֥ן עָלָ֖יו הֶעָנָ֑ן וּכְב֣וֹד יְהֹוָ֔ה מָלֵ֖א אֶת־הַמִּשְׁכָּֽן׃ \Rashi{\textbf{ולא יכל משה לבוא אל אהל מועד.}וכתוב אחד אומר ובבא משה אל אהל מועד (במדבר ז')? בא הכתוב השלישי והכריע ביניהם – כי שכן עליו הענן, אמר מעתה כל זמן שהיה עליו הענן לא היה יכול לבא, נסתלק הענן נכנס ומדבר עמו (ספרא):}%
% ---------- פרק 40 | פסוק 36 ----------
 \textbf{לו} וּבְהֵעָל֤וֹת הֶֽעָנָן֙ מֵעַ֣ל הַמִּשְׁכָּ֔ן יִסְע֖וּ בְּנֵ֣י יִשְׂרָאֵ֑ל בְּכֹ֖ל מַסְעֵיהֶֽם׃ %
% ---------- פרק 40 | פסוק 37 ----------
 \textbf{לז} וְאִם־לֹ֥א יֵעָלֶ֖ה הֶעָנָ֑ן וְלֹ֣א יִסְע֔וּ עַד־י֖וֹם הֵעָלֹתֽוֹ׃ %
% ---------- פרק 40 | פסוק 38 ----------
 \textbf{לח} כִּי֩ עֲנַ֨ן יְהֹוָ֤ה עַֽל־הַמִּשְׁכָּן֙ יוֹמָ֔ם וְאֵ֕שׁ תִּהְיֶ֥ה לַ֖יְלָה בּ֑וֹ לְעֵינֵ֥י כׇל־בֵּֽית־יִשְׂרָאֵ֖ל בְּכׇל־מַסְעֵיהֶֽם׃ \Rashi{\textbf{לעיני כל בית ישראל בכל מסעיהם.}בכל מסע שהיו נוסעים, היה הענן שוכן במקום אשר יחנו שם; מקום חניתן אף הוא קרוי מסע, וכן וילך למסעיו (בראשית י"ג), וכן אלה מסעי (במדבר ל"ג), לפי שממקום החניה חזרו ונסעו, לכך נקראו כלן מסעות:}%
